\section{Part D: fundamental groups, descent, and ramification}

\subsection{D.1 and D.2(ii): positive cases and the fixed-boundary obstruction}

\begin{proposition}[Proved under stated hypotheses: fixed-boundary finite generation]
\label{D:fixed-boundary-positive}
Let $X$ be normal, proper and connected over an algebraically closed
field $k$, let $Y\subsetneq X$ be closed, and let $Q(X,Y)$ be
$\pi_1(X\setminus Y)$ modulo the closed normal subgroup generated
by wild inertia at the codimension-one components of $Y$.
The group $Q(X,Y)$ is topologically finitely generated in each of
the following cases: $X$ is projective; $\operatorname{char}k=0$;
or $X$ is regular and $Y$ is a normal-crossings divisor.
Consequently D.2(ii) holds when the proper compactification is
projective, and it holds for every regular proper compactification.
\end{proposition}
\begin{proof}
The zero-dimensional case is immediate. We first treat projective $X$,
of dimension $d\geq2$; the curve case
is the tame-curve theorem \cite[XIII, Corollary 2.12, pp. 392--393]{SGA1}.
Embed $X$ in projective space. Over the algebraic closure $K$ of the
function field of the parameter space of $d-1$ hyperplanes, take the
generic complete-intersection curve $C$ in $X_K$. It is smooth and
geometrically connected, avoids every fixed codimension-two closed
subset of $X$ defined over $k$, and meets the regular locus of each
codimension-one component of $Y$ transversally. Put $C^\circ=C\setminus Y_K$.

Let $W\to X\setminus Y$ be any connected finite \'etale cover
classified by $Q(X,Y)$, and let $\overline W\to X$ be its finite
normalization. Generic Bertini applied to the finite dominant map
$\overline W\to X\hookrightarrow\mathbb P^N$ shows that
$\overline W\times_X C$ is geometrically irreducible. Its nonempty
open part $W\times_XC^\circ$ is therefore connected. The generic
point of the hyperplane parameter space lies outside every proper
closed subset defined over $k$, so the single generic curve has this
property for every cover defined over $k$; no uniform assertion about
a chosen closed point of the parameter space is needed. This is the
generic Bertini argument used in
\cite[X, Lemma 2.10 and Corollary 2.11, p. 274]{SGA1}.

At each boundary component the normalization is tame by hypothesis.
After removing a further codimension-two closed subset, its local
description is a tame Kummer extension followed by an \'etale map;
this follows from the discrete-valuation tame description and
\cite[XIII, Proposition 5.2, pp. 428--429]{SGA1}.
Our generic curve avoids that exceptional set too. The induced cover
of $C^\circ$ is thus tame at its missing points. The resulting
homomorphism $\pi_1^t(C^\circ)\to Q(X,Y)$ is surjective: every
connected finite cover classified by the target stays connected on
restriction. The source is finitely generated by the tame-curve
theorem, proving the first case. In particular, if $Y$ has codimension
at least two there are no boundary inertia groups to kill, and this
proves finite generation of the full $\pi_1(X\setminus Y)$.

In characteristic zero, the characteristic exponent is one, so
\cite[Theorem 4.11, p. 143]{Orgogozo2003} gives finite presentation
of the full fundamental group of every connected finite-type
$k$-scheme. This proves the second case, including the stated
finite-generation assertion.

For regular proper $X$ with normal-crossings boundary, fixed-boundary
tameness equals valuation tameness: at a strict local point,
Abhyankar's lemma dominates a tame cover by adjoining prime-to-$p$
roots of the local equations of the boundary. Every further valuation
extension induced from this cover is tame. The precise comparison is
\cite[Proposition 2.2.3]{AHLS2026}, and
\cite[Theorem 2.6.1]{AHLS2026} gives finite generation for the resulting
tame group without projectivity. Finally, when $X$ is regular and
$Y$ has codimension at least two, purity identifies its open
fundamental group with that of $X$, and
\cite[X, Theorem 2.9, pp. 273--274]{SGA1} gives finite generation.

For the potentially nonnormal scheme $Z$ in D.2(ii), take the disjoint
union of the normalizations of the irreducible components of
$Z_{\mathrm{red}}$. This is a finite proper surjection, and the
inverse image of $Y$ still has codimension at least two in each
component. If $Z$ is projective, the components of that normalization
are projective, so the proved case applies to them. Effective descent
and \cite[IX, Corollary 5.2, pp. 247--248]{SGA1} then give the
assertion for $Z\setminus Y$ itself. This also reduces the general
proper question to its normal integral components, without invoking
the invalid codimension assertion for arbitrary proper maps.
\end{proof}


\begin{corollary}[Known: all surfaces over the algebraic closure of a finite field]
\label{D:finite-field-surfaces}
D.1 holds in dimension at most two over $\overline{\mathbb F}_p$.
So does D.2(ii), including possibly nonnormal proper schemes.
\end{corollary}
\begin{proof}
By \cite[Corollary 2.11; published Russian translation, p.~13]{Artin1962D},
every proper normal surface over $\overline{\mathbb F}_p$ is
projective. The same assertion for proper normal curves is classical.
Apply Proposition~\ref{D:fixed-boundary-positive}. For D.2(ii),
the finite normalizations of the reduced irreducible components
are projective of dimension at most two. Their punctured opens have
finitely generated fundamental groups by the same proposition;
finite effective descent proves the assertion downstairs. Thus a
residual surface case in positive characteristic must be over a
field of positive transcendence degree over $\mathbb F_p$.
\end{proof}

\begin{proposition}[Proved: a quasi-projective neighborhood of the entire boundary]
\label{D:boundary-neighborhood}
Let $X$ be normal, connected and proper over an algebraically closed
field. If the closed subset $Y\subsetneq X$ is contained in one
quasi-projective open subscheme $N\subset X$, then $Q(X,Y)$ is
topologically finitely generated. In particular, D.2(ii) holds when
the omitted set consists of a single closed point, even on a
nonprojective singular proper scheme. More generally, D.2(ii) holds
for a possibly nonnormal proper scheme $Z$ whenever its entire
codimension-two omitted subset has a quasi-projective open
neighborhood in $Z$.
\end{proposition}
\begin{proof}
We use the strengthened form of Chow's lemma in
\cite[Chapter 4, section 4, Proposition 5, p.~30]{Raynaud1972D}.
It supplies a blowup $X_1\to X$ with centre disjoint from the
\emph{prescribed entire open} $N$, with $X_1$ quasi-projective
over the field. Since $X$ is proper, $X_1$ is projective. Its
normalization $X'\to X_1$ is finite, so $X'$ remains projective;
normality of $N$ makes the composite $q:X'\to X$ an isomorphism
over $N$. Put $Y'=q^{-1}(Y)$ and $U'=X'\setminus Y'$.
Every boundary component and its discrete valuation ring are
unchanged, because $q$ is an isomorphism on a neighborhood of
\emph{all} of $Y$.

The map $\pi_1(U')\to\pi_1(X\setminus Y)$ is surjective.
Indeed a connected finite \'etale cover of the normal integral
scheme $X\setminus Y$ is integral, with a field as generic fibre.
After this birational pullback its generic fibre is the same field;
its finite \'etale total space over normal integral $U'$ is therefore
connected. This verifies the Galois-category criterion for surjectivity.
A cover classified by $Q(X,Y)$ pulls back to one classified by
$Q(X',Y')$, by the equality of boundary valuation rings. We thus
obtain a surjection
\[
 Q(X',Y')\longrightarrow Q(X,Y).
\]
The left side is finitely generated by the projective case of
Proposition~\ref{D:fixed-boundary-positive}, proving the assertion.

For a possibly nonnormal $Z$, replace $Z_{\mathrm{red}}$ by the
disjoint union of the normalizations of its irreducible components.
The finite inverse image of the specified quasi-projective
neighborhood is quasi-projective and contains the entire inverse
image of the omitted set. Codimension at least two is preserved
componentwise by this finite map. The normal case just proved applies
to each component, and finite effective descent
\cite[IX, Corollary 5.2, pp.~247--248]{SGA1} proves finite generation
downstairs. A single closed point has an affine open neighborhood;
its whole inverse image under normalization also lies in the finite
inverse image of that affine open, which is affine. This verifies
explicitly the claimed single-point case, including nonnormal $Z$.
\end{proof}


\begin{proposition}[Proved: punctured mutation surfaces]
\label{D:mutation-surfaces}
Let $C$ be a smooth proper connected curve over an algebraically
closed field, and blow up $\mathbb P^1\times C$ at finitely many
points $(b_i,c_i)$ with distinct $b_i$. Let $R_i$ be the strict
transform of the fibre $\{b_i\}\times C$, and put
$U=\operatorname{Bl}_{\{(b_i,c_i)\}}(\mathbb P^1\times C)
\setminus\bigcup_iR_i$. Then projection induces an isomorphism
\[
 \pi_1(U)\simeq\pi_1(C).
\]
Consequently the nonprojective normal mutation surfaces of
\cite[sections 2.2--2.5, author's PDF pp.~3--4]{Schroer1999D},
punctured at all their contracted-curve singularities, satisfy D.2(ii).
These include two-point boundaries having no quasi-projective
neighborhood.
\end{proposition}
\begin{proof}
Put $B=\mathbb P^1\setminus\{b_i\}$. The dense open
$B\times C\subset U$ surjects on fundamental groups because $U$
is normal. The proper product theorem
\cite[X, Corollary 1.7, pp.~264--265]{SGA1} identifies its group
with $\pi_1(B)\times\pi_1(C)$.

For each $i$, take the strict transform $H_i$ of the horizontal
curve $\mathbb P^1\times\{c_i\}$. Its intersection with $U$ is
\[
 H_i\cap U\simeq
 \mathbb P^1\setminus\{b_j:c_j\ne c_i\}.
\]
At every centre lying on the horizontal curve, its strict transform
meets the exceptional curve at a point different from the point
where $R_j$ meets that curve. In particular $b_i$ remains in
$H_i\cap U$. On $B\times C$, this is the constant-$C$ section.
It follows that the inertia subgroup of $\pi_1(B)$ at $b_i$ has
trivial image in $\pi_1(U)$. Choices of paths in the $C$ factor
cause no ambiguity here, since the two factors commute. This holds
for every $i$, so the entire $\pi_1(B)$ factor is killed: its
quotient obtained by filling all $b_i$ is $\pi_1(\mathbb P^1)=1$.
Thus $\pi_1(C)\to\pi_1(U)$ is surjective. The projection $U\to C$
and a complete fibre over any point of $B$ make its composite with
$\pi_1(U)\to\pi_1(C)$ the identity, proving the claimed isomorphism.
Schr\"oer's contraction is an isomorphism off the $R_i$, so its
punctured surface is precisely $U$.
\end{proof}

\begin{proposition}[Proved: tame ramification above positive-genus exceptional curves]
\label{D:lc-elliptic-tame}
Let $(S,s)$ be a normal log canonical surface singularity over an
algebraically closed field $k$.  Let $E$ be an exceptional divisorial
valuation over $s$ whose residue curve has positive geometric genus.
Every finite separable quasi-\'etale cover $S'\to S$ is tame at $E$:
for every extension $F$ of this valuation, its ramification index is
prime to the characteristic exponent of $k$ and
$\kappa(F)/\kappa(E)$ is separable.
In the application below we verify directly the log canonical and
Gorenstein properties of the contraction of one smooth elliptic curve,
so no classification in small characteristic is required.
\end{proposition}
\begin{proof}
We use surface resolution and point-blowup factorization in arbitrary
characteristic, as stated in the Stacks Project,
Theorem~54.14.5 and Lemma~54.17.1, tags 0BGP and 0C5R.
Contraction of a smooth rational exceptional curve of self-intersection
$-1$ is Lemma~54.16.8, tag 0C2M
\cite{StacksDSurfaceResolution}.
The exceptional intersection matrix on a resolution over an affine
neighborhood of the singularity is negative definite by
\cite[section 10.1, Theorem 10.1]{Kollar2013D}; the source explicitly
works with excellent surfaces in arbitrary characteristic in that
section.

We first give the ramification calculation, including the case of an
inseparable residue extension.  Let $F$ be a divisorial valuation in
a finite separable extension of the surface function field, and let
$E$ be its restriction.  Choose models regular at the respective
generic points and a morphism between neighborhoods of those points.
These neighborhoods are smooth over the perfect field $k$.
Choose a uniformizer $t$ at $E$ and a lift $z$ of a separating
transcendence parameter of $\kappa(E)/k$; then $dt\wedge dz$
generates the canonical module at $E$.  Write
$t=u s^e$ at $F$, with $u$ a unit and $s$ a uniformizer.
Choose similarly a separating residue parameter $w$ at $F$.
The order $\delta_F$ of the determinant of differentials is the
coefficient of the relative canonical divisor, and
\[
 d(u s^e)\wedge dz
   =s^{e-1}\bigl(eu\,ds\wedge dz+s\,du\wedge dz\bigr).
\]
Consequently $\delta_F\geq e-1$.  Equality holds precisely when
$e$ is nonzero in $k$ and $d\bar z$ is nonzero in
$\Omega^1_{\kappa(F)/k}$.  For a finite extension of one-variable
function fields over a perfect field, the latter condition is
equivalent to separability of $\kappa(F)/\kappa(E)$.
Thus equality is precisely tame ramification, including the
residue-separability requirement.  The determinant is not identically
zero because the extension of surface function fields is separable.

Since $S'\to S$ is quasi-\'etale, its canonical divisor satisfies
$K_{S'}=\pi^*K_S$ as a rational Cartier divisor: the equality holds
at every codimension-one point, and normality extends the equality
of Weil divisors.  In particular $S'$ is $\mathbb Q$-Gorenstein.
On models extracting a pair of valuations as above, comparison of
canonical divisors gives the exact discrepancy formula
\[
 a(F,S')=e\,a(E,S)+\delta_F.
\]
The restriction of a divisorial valuation through a finite field
extension is again divisorial.  Applying the lower bound to every
such valuation shows that $S'$ is log canonical wherever $S$ is:
$a(F,S')\geq-e+(e-1)=-1$.

We next prove the needed fact about positive-genus exceptional curves
without a classification theorem.  On a minimal resolution of an
arbitrary log canonical surface singularity, write
\[
 K_Z=f^*K_S+\sum_i a_i E_i,\qquad M=(E_i\cdot E_j).
\]
For each exceptional integral curve, adjunction gives
$K_Z\cdot E_i=2p_a(E_i)-2-E_i^2\geq0$.
Indeed $E_i^2<0$; when $p_a(E_i)=0$, the curve is smooth rational
and minimality excludes $E_i^2=-1$.  Thus $Ma=b$ with $b\geq0$.
All off-diagonal entries of $M$ are nonnegative.  If the positive
part $a^+$ were nonzero, negative definiteness and the disjoint
supports of $a^+$ and the nonpositive part $a^-$ would imply
\[
 0\leq (a^+)^{\mathsf T}b
    =(a^+)^{\mathsf T}Ma^++(a^+)^{\mathsf T}Ma^-<0.
\]
Hence every $a_i\leq0$.  If the normalization of $E_i$ has positive
genus, then $p_a(E_i)\geq1$, so
\[
 -E_i^2\leq K_Z\cdot E_i
      =\sum_j a_j(E_j\cdot E_i)
      \leq a_iE_i^2\leq-E_i^2.
\]
The last inequality uses log canonicity $a_i\geq-1$.
Equality forces $p_a(E_i)=1$ and $a_i=-1$.
A positive-genus valuation appearing on any other resolution
survives the contractions to a minimal resolution, because the
curves contracted between smooth surfaces are rational.  Its
discrepancy is therefore also $-1$.

Return to the specified $E$ and a valuation $F$ above it.
The residue curves admit a finite dominant map after normalization.
The curve of $F$ cannot be rational, since a nonconstant map from
$\mathbb P^1$ to a positive-genus smooth proper curve is impossible,
also for inseparable maps.  Equivalently we may apply L\"uroth's
theorem to the corresponding function fields.  Both $E$ and $F$
therefore have discrepancy $-1$ on the respective log canonical
surfaces.  The discrepancy formula now gives
$-1=-e+\delta_F$, hence $\delta_F=e-1$.
The differential calculation proves the asserted tameness.
\end{proof}

\begin{proposition}[Proved under stated hypotheses: a single positive-genus exceptional curve]
\label{D:single-curve-residue-tameness}
Let $k$ be algebraically closed, let $(S,s)$ be a normal surface
singularity of finite type over $k$, and suppose that a projective
resolution $Y\to S$ has reduced exceptional locus a single smooth
connected curve $C$, of genus $g>0$ and self-intersection $C^2=-d<0$.
Let $S'\to S$ be a finite generically separable morphism of normal
integral surfaces, \'etale away from the points over $s$, of degree $n$.
For the extensions $F_i$ of the divisorial valuation $C$, write
$e_i$ for the ramification index, $f_i$ for the residue-field degree,
$f_i^{\mathrm{sep}}$ for its separable degree, and $\delta_i$ for
the coefficient of the relative canonical divisor on smooth models.
Then
\[
 0\leq\sum_i f_i(\delta_i-e_i+1)
 \leq \frac{2g-2}{d}
             \left(n-\sum_i e_i f_i^{\mathrm{sep}}\right).
 \tag{D.PG.1}
\]
In particular the cover is tame at every valuation above $C$ if
every residue extension is separable. If $g=1$, all the residue
extensions are automatically separable and the cover is tame there.
No $\mathbb Q$-Gorenstein assumption on $S$ or $S'$ is required.
\end{proposition}

\begin{proof}
We shrink $S$ to an affine neighborhood of $s$ with no other
singularities. Normalize $Y$ in the function field of $S'$ and
resolve that normal surface by a morphism that is an isomorphism
off the inverse image of $C$, where the normalization is already
smooth. This gives a smooth projective
$S'$-surface $Z$ and a proper generically finite morphism
$f:Z\to Y$ of degree $n$. Resolution by normalized blowups is
\cite[Theorem 54.14.5, Tag 0BGP]{StacksDSurfaceResolution}.
We may make $Z\to S'$ minimal while keeping a morphism to $Y$.
Indeed, any exceptional smooth rational curve of self-intersection
$-1$ maps into $C$, and its image is a point: a nonconstant map
from $\mathbb P^1$ to a smooth curve of positive genus is impossible,
including for an inseparable map, by L\"uroth's theorem.
Contract this curve using
\cite[Lemma 54.16.8, Tag 0C2M]{StacksDSurfaceResolution}.

We spell out descent of $f$ through the contraction $b:Z\to Z_1$.
The scheme-theoretic image $W$ of $(b,f):Z\to Z_1\times_S Y$ is
integral and proper over $Z_1$. Away from the contracted point its
projection is an isomorphism, and over that point its underlying
fibre is a single point because $f$ is constant on the contracted
curve. Thus $W\to Z_1$ is proper and quasi-finite, hence finite.
It is birational, and $Z_1$ is normal, so it is an isomorphism.
This gives the required descended morphism. Each contraction lowers
the number of exceptional components, so the procedure terminates.

On this minimal $Z$, the canonical divisor has nonnegative
intersection with every exceptional integral curve $D$ over $S'$.
The exceptional intersection matrix is negative definite
\cite[Section 10.1, Theorem 10.1]{Kollar2013D}, hence $D^2<0$.
Adjunction gives
\[
                  K_Z\cdot D=2p_a(D)-2-D^2\geq0.
\]
If $p_a(D)=0$, then $D$ is a smooth rational curve; minimality
excludes $D^2=-1$, so the inequality still holds. We use here
the adjunction formula $(K_Z+D)\cdot D=2p_a(D)-2$ for a Cartier
curve on a smooth surface; it follows by applying degree to
$\omega_D=(\omega_Z\otimes\mathcal O_Z(D))|_D$.

Let the $F_i$ now denote the actual curves on $Z$ dominating $C$,
and let $Q_j$ be all the remaining exceptional curves. Their
images under $f$ are points of $C$. We have
\[
             f^*C=\sum_i e_iF_i+\sum_j m_jQ_j,
             \qquad m_j>0,\qquad \sum_i e_if_i=n.
 \tag{D.PG.2}
\]
The last equality is the degree formula for the finite normalization
over the discrete valuation ring at the generic point of $C$;
it includes the inseparable parts of the residue degrees.
By the projection formula and nonnegativity of intersections
between distinct integral curves on a smooth surface,
\[
 -df_i=(f^*C)\cdot F_i
       =e_iF_i^2+\sum_{h\ne i}e_h(F_h\cdot F_i)
                      +\sum_jm_j(Q_j\cdot F_i),
\]
so $e_i(-F_i^2)\geq df_i$.

Let $\widetilde F_i$ be the normalization of $F_i$.
Factoring $\widetilde F_i\to C$ into its purely inseparable
and separable parts, the purely inseparable part preserves the
genus over the perfect field $k$, by
\cite[Proposition 53.13.7 and Lemma 53.13.5, Tags 0CD2 and 0CD0]{StacksDSurfaceResolution}.
Riemann--Hurwitz for the
separable part gives
\[
 2p_a(F_i)-2\geq 2g(\widetilde F_i)-2
                   \geq f_i^{\mathrm{sep}}(2g-2).
\]
For the Riemann--Hurwitz statement we use
\cite[Lemma 53.12.2, Tag 0C1D]{StacksDSurfaceResolution}.
Together with nefness on the $Q_j$, adjunction, and (D.PG.2),
this yields the lower bound
\[
 \begin{split}
 K_Z\cdot f^*C
 &\geq\sum_i e_i\bigl(2p_a(F_i)-2-F_i^2\bigr)\\
 &\geq(2g-2)\sum_i e_i f_i^{\mathrm{sep}}+d\sum_i f_i.
 \end{split}
 \tag{D.PG.3}
\]

The determinant of the differential of $f$ defines an effective
relative canonical divisor $R_f$, with
$K_Z=f^*K_Y+R_f$. The morphism is \'etale outside $f^{-1}C$,
so $R_f$ is supported on the $F_i$ and $Q_j$, with coefficient
$\delta_i$ at $F_i$. As $(f^*C)\cdot Q_j=0$, the projection
formula and adjunction on $C$ give the exact identity
\[
 \begin{split}
 K_Z\cdot f^*C
  &= nK_Y\cdot C-d\sum_i f_i\delta_i\\
  &= n(2g-2+d)-d\sum_i f_i\delta_i.
 \end{split}
 \tag{D.PG.4}
\]
Combining (D.PG.3) and (D.PG.4), and using
$n=\sum_i e_i f_i$, proves the upper inequality in (D.PG.1).

For clarity, the lower inequality and its equality case include
inseparable residue extensions. Choose uniformizers $t,s$ at
$C,F_i$ and separating residue parameters $z,w$ over $k$.
At the corresponding discrete valuation ring write $t=u s^{e_i}$.
Then
\[
 d(us^{e_i})\wedge dz
   =s^{e_i-1}\bigl(e_i u\,ds\wedge dz+s\,du\wedge dz\bigr).
\]
The order of this differential relative to $ds\wedge dw$ is
$\delta_i$, so $\delta_i\geq e_i-1$. Equality holds exactly
when $e_i$ is nonzero in $k$ and the residue extension is
separable, which is precisely tame ramification. The determinant
is nonzero because the extension of the surface function fields
is separable. Thus every summand in (D.PG.1) is nonnegative.
If all $f_i=f_i^{\mathrm{sep}}$, its right-hand side is zero.
If $g=1$, it is zero without any separability assumption.
In either case all summands vanish, giving the conclusion.
\end{proof}

\begin{remark}[Scope of the higher-genus reduction]
For the single-curve singularities obtained by contracting a strict
vertical fibre in the mutation construction, the hypotheses hold
with $d=1$. Consequently a finite cover with wild ramification at
such an exceptional valuation must have an inseparable residue
extension at at least one valuation above it. For a Galois cover,
all residue separability types are conjugate, so every such wild
case has inseparable residue extension. We do not infer finite
generation, a finite image of wild inertia, or a uniform bound on
the orders of finite covers from (D.PG.1). In genus at least two
the remaining inseparable-residue case is still an explicit gate.
\end{remark}


\begin{example}[Refuted: a positive uniform lower bound for fierce excess]
\label{D:fierce-small-excess}
The assertion that Galois extensions with inseparable residue have a
positive uniform lower bound for $(\delta-e+1)/e$ is false.
We fix an algebraically closed field $k$ of characteristic $p>0$.
For every power $q=p^r>1$, there is a finite Galois extension of
two-variable rational function fields over $k$, with a divisorial valuation
having ramification index $e=q(q-1)$, purely inseparable residue degree $q$,
and different exponent
\[
  \delta=q^2+q-3.
\]
Consequently
\[
  \frac{\delta-e+1}{e}=\frac2q\longrightarrow0.
\]
In particular, inseparability of the residue extension does not give a
positive lower bound for this normalized excess different, even at a fixed
characteristic and with a Galois extension between function fields of smooth
surfaces.  This is a local ramification example; we do not assert that it
extends to a quasi-\'etale cover of a prescribed contracted surface.
\end{example}

\begin{proof}
We take algebraically independent elements $s,z$ over $k$, put
$t=s^{q-1}$, and adjoin elements $w,y$ satisfying
\[
 w^q-w=s^{-1},\qquad y^q-y=zs^{-1}.
\]
We write $K=k(t,z)$, $T=k(s,z)$, and $L=T(w,y)$.
The extension $T/K$ has degree $q-1$ and is totally tamely ramified
at $t=0$.  For $\pi=w^{-1}$ the first equation becomes
\[
 \pi^q+s\pi^{q-1}-s=0.
\]
This polynomial is Eisenstein at $s=0$.  Thus $[T(w):T]=q$, the
ramification index is $q$, and its residue field is $k(z)$.
Putting $u=y/w$, we obtain
\begin{equation}
\label{D:fierce-small-excess-coordinates}
 s=\frac{\pi^q}{1-\pi^{q-1}},\qquad
 z=\frac{u^q-\pi^{q-1}u}{1-\pi^{q-1}},\qquad
 L=k(\pi,u).
\end{equation}
The monic equation
\[
 u^q-\pi^{q-1}u-z(1-\pi^{q-1})=0
\]
reduces to $u^q-z$ at $\pi=0$.  This is irreducible over $k(z)$,
since the transcendental element $z$ is not a $p$th power.
The corresponding finite algebra over the discrete valuation ring
$k(z)[\pi]_{(\pi)}$ is local with maximal ideal generated by $\pi$;
it is therefore a discrete valuation ring.  This proves
$[L:T(w)]=q$, ramification index one, and purely inseparable residue
degree $q$.  We have proved $[L:K]=q^2(q-1)$.

For $a,b\in\mathbf F_q$ and $c\in\mathbf F_q^\times$ the formulas
\[
 (w,y,s,z)\longmapsto(cw+a,cy+b,c^{-1}s,z)
\]
give $q^2(q-1)$ distinct automorphisms fixing $t$ and $z$.
They form $\mathbf F_q^2\rtimes\mathbf F_q^\times$, with the
scalar action on the two additive factors.  The extension $L/K$ is
therefore Galois.  Each automorphism preserves the valuation at $\pi=0$:
its action on the displayed coordinates is
\[
 \pi\longmapsto\frac{\pi}{c+a\pi},\qquad
 u\longmapsto\frac{cu+b\pi}{c+a\pi}.
\]
In particular it acts trivially on the residue field $k(u)$.
The ramification index and residue degree over $K$ are consequently
$e=q(q-1)$ and $f=q$, respectively, as also follows by multiplying
the indices in the three steps above.

The different at this divisorial valuation is the order of the
Jacobian determinant for $(\pi,u)\mapsto(t,z)$.
Writing $D=1-\pi^{q-1}$, differentiation in characteristic $p$ gives
\[
 \frac{\partial s}{\partial\pi}
       =-\frac{\pi^{2q-2}}{D^2},\qquad
 \frac{\partial t}{\partial\pi}
       =\frac{\pi^{q^2-2}}{D^q},\qquad
 \frac{\partial z}{\partial u}
       =-\frac{\pi^{q-1}}D.
\]
Since $t$ is independent of $u$, we conclude
\[
 dt\wedge dz
   =-\frac{\pi^{q^2+q-3}}{D^{q+1}}\,d\pi\wedge du.
\]
The factor $D$ is a unit at the valuation, so
$\delta=q^2+q-3$ and the stated normalized excess follows.
\end{proof}

\begin{remark}[What this example leaves open]
In the Galois case, Proposition~\ref{D:single-curve-residue-tameness}
gives the necessary bound
\[
 \frac{\delta-e+1}{e}
   \leq\frac{2g-2}{d}\left(1-\frac1{f_{\mathrm{insep}}}\right).
\]
For every fixed $g>1$ and $d>0$, the numerical ramification data of
Example~\ref{D:fierce-small-excess} satisfy this inequality for all
sufficiently large $q$.  Thus the inequality by itself cannot exclude
fierce residue extensions by imposing a large value of $d$.
Realizability over a contracted surface, and the additional global
conditions in the finite-generation question, remain separate issues.
\end{remark}

\paragraph{Finite certificate.}
The program \nolinkurl{parts/D/code/check_fierce_different.py} and its
output \nolinkurl{parts/D/output/fierce_different_certificate.json}
check $q=2,3,4,5,8,9$.  They verify the finite fields, the faithful affine
semidirect actions and their generation, the two defining equations, all
three derivatives, the full Jacobian identity, and its valuation.
The run passed $1{,}029{,}677$ exact assertions in
$0.5433858470059931$ seconds.  The proof for every prime power is the argument
above, not an extrapolation from these finite checks.


% Written in Part B for independent audit and integration in Part D.
% Configuration lemma and cusp theorem: root/D actual-file audits passed.

\begin{lemma}[Proved: Gorenstein log canonical exceptional configurations]
\label{D:Gorenstein-lc-configurations}
Let $(S,s)$ be a normal Gorenstein log canonical surface singularity
over an algebraically closed field $k$, in arbitrary characteristic.
On its minimal regular resolution, precisely one of the following
possibilities occurs:
\begin{enumerate}
\item the exceptional fibre is empty, or it is a simple normal
crossings tree of smooth rational curves, each of self-intersection
$-2$, and every discrepancy is zero;
\item the exceptional fibre is one smooth elliptic curve, of
discrepancy $-1$;
\item the exceptional fibre is a nodal cycle of rational curves,
and every discrepancy is $-1$.  A single irreducible nodal rational
curve is allowed in this description.
\end{enumerate}
After blowing up the node in the last single-component case, the
cycle has smooth components.  On every further regular resolution,
the elliptic or cyclic core persists; point blowups only subdivide
edges of the cyclic core or attach rational trees.  Every component
of a cyclic core still has discrepancy $-1$.
\end{lemma}
\begin{proof}
We use existence of surface resolutions, contraction of exceptional
smooth rational $(-1)$-curves, and factorization of a proper
birational morphism between regular surfaces into point blowups.
These hold in arbitrary characteristic, including over excellent
strictly henselian local surface rings; see
\cite[Theorem 54.14.5, Lemmas 54.16.8 and 54.17.1;
tags 0BGP, 0C2M and 0C5R]{StacksDSurfaceResolution}.
The intersection matrix of the exceptional curves is negative
definite by \cite[section 10.1, Theorem 10.1]{Kollar2013D}.
The fibre is connected, since a proper birational morphism to a
normal surface has direct image of its structure sheaf equal to
the structure sheaf, and hence has connected fibres.

Write the minimal resolution as $f:Z\to S$ and
\[
 K_Z=f^*K_S+\sum_i a_i E_i,\qquad
 M=(E_i\cdot E_j),\qquad b_i=K_Z\cdot E_i.
\]
Adjunction gives
$b_i=2p_a(E_i)-2-E_i^2\geq0$.
Indeed $E_i^2<0$; if $p_a(E_i)=0$, then $E_i$ is smooth rational
and minimality excludes self-intersection $-1$.
All off-diagonal entries of $M$ are nonnegative.
If $a^+$ denotes the positive part of the vector $a$, and $a^-$
its nonpositive part, then a nonzero $a^+$ would give
\[
 0\leq (a^+)^{\mathsf T}b
 =(a^+)^{\mathsf T}Ma^++(a^+)^{\mathsf T}Ma^-<0,
\]
by negative definiteness and the disjoint supports of the two parts.
Thus every $a_i\leq0$.  Gorensteinness makes the $a_i$ integers,
and log canonicity gives $a_i\geq-1$.  Hence
$a_i\in\{-1,0\}$.

If $a_i=0$ for some $i$, then
\[
 0\leq b_i=\sum_{j\ne i}a_j(E_j\cdot E_i)\leq0.
\]
Every component meeting $E_i$ therefore also has discrepancy zero.
Connectedness shows that all discrepancies are zero.  Adjunction
then gives $p_a(E_i)=0$ and $E_i^2=-2$ for every $i$.
There can be no cycle in the intersection multigraph: the vector
equal to one on a shortest cycle and zero elsewhere would have
self-intersection at least zero.  The same calculation rules out
an intersection number at least two between two components.
A triple intersection would give a three-cycle and is also
impossible.  Thus the fibre is a simple normal crossings rational
tree, as asserted in the first case.

Otherwise all $a_i=-1$, and adjunction becomes
\begin{equation}\label{D:eq-Gorenstein-lc-adjunction}
 2p_a(E_i)-2=-\sum_{j\ne i}E_j\cdot E_i.
\end{equation}
If one component has arithmetic genus at least one, this equation
and connectedness say that it is the only component and has
arithmetic genus one.  Otherwise every component is smooth
rational and has total intersection number two with the remaining
components.

We check the possible singularities of the reduced divisor
$E=\bigcup_i E_i$, rather than assuming a characteristic-zero
classification.  We have $K_Z+E=f^*K_S$, so the pair $(Z,E)$ is
log canonical.  A point of multiplicity $m\geq3$ in $E$ is
impossible: blowing it up creates a divisor of discrepancy
$1-m\leq-2$.  At a singular point of multiplicity two, two
distinct tangent lines give an ordinary node.  If the tangent line
is repeated, the first blowup creates a divisor of discrepancy
$-1$, and the strict transform meets it at a single point with
total intersection multiplicity two.  If the strict transform
has multiplicity at least two there, the reduced total boundary
already has multiplicity at least three.  Otherwise the strict
transform is smooth and tangent to the exceptional curve;
blowing up this tangency gives a point at which three reduced
boundary branches, counted with multiplicity, meet.  In either
case a further blowup creates discrepancy at most $-2$.
This contradicts log canonicity.  This elementary blowup argument
works also in characteristics two and three.
Consequently $E$ is nodal.

An integral nodal curve of arithmetic genus one is either smooth
elliptic or a rational curve with one node.  In the remaining case,
equation~\eqref{D:eq-Gorenstein-lc-adjunction} says that the connected
dual multigraph is two-regular, so it is a cycle.  This includes
two rational curves meeting transversely at two distinct points.
Blowing up the node of an irreducible rational curve gives the
same two-component configuration.  Its new component has
discrepancy $1-2=-1$.

Finally, factorization into point blowups shows the description
on any dominating regular resolution.  A blowup at a node of the
cycle subdivides its edge and has discrepancy $1-1-1=-1$.
A blowup at a smooth point of the reduced fibre attaches a
rational component.  Further blowups of an attached tree still
leave the original core as the unique cyclic, respectively
positive-genus, part.  This proves the final assertion.
\end{proof}

\begin{theorem}[Proved: tame ramification above a Gorenstein cusp cycle]
\label{D:cusp-tame}
Let $(S,s)$ be a normal Gorenstein log canonical surface singularity
over an algebraically closed field $k$.  Suppose that its minimal
resolution has a rational nodal cycle as exceptional fibre.
For every finite separable quasi-\'etale cover $S'\to S$ and every
component $E_i$ of this cycle, every extension $F$ of the associated
divisorial valuation is tame:
\[
 e(F/E_i)\text{ is prime to the characteristic exponent of }k,
 \qquad \kappa(F)/\kappa(E_i)\text{ is separable}.
\]
The same assertion holds for the components inserted by blowing
up nodes of the cycle.  The assertion is valid in every
characteristic; no assumption on the degree of the cover is made.
Here a cover is taken with normal integral source, or is treated
componentwise, and it is \'etale on the punctured germ.  For a
quasi-\'etale cover of an algebraic surface, we obtain this germ
by shrinking away the other finitely many branch points.
\end{theorem}
\begin{proof}
We give the local and group-theoretic comparisons explicitly.
The conclusion descends from a Galois closure, so we may assume
that the cover is Galois.  The Galois closure of its restriction
to the punctured surface is finite \'etale, and normalization
over the excellent surface is finite and quasi-\'etale.

\par\smallskip\noindent\textbf{Canonical divisors and strict localization.}\quad
At every codimension-one point the original cover is \'etale.
Thus compatible rational canonical forms give
$K_{S'}=\pi^*K_S$ as Weil divisors.  Since $K_S$ is Cartier,
the canonical reflexive sheaf on $S'$ is invertible.
A normal surface is Cohen--Macaulay, so $S'$ is Gorenstein.
This proves the property of the covering surface; we do not
assume it merely from the hypothesis on the target.

For a pair of divisorial valuations $F/E$ in a finite separable
surface function-field extension, let $e=e(F/E)$ and let
$\delta_F$ be the order of the determinant of differentials.
The calculation in Proposition~\ref{D:lc-elliptic-tame} gives
\begin{equation}\label{D:eq-cusp-discrepancy}
 a(F,S')=e\,a(E,S)+\delta_F,
 \qquad \delta_F\geq e-1.
\end{equation}
Equality in the second inequality holds exactly when $e$ is
prime to the characteristic exponent and the residue extension is
separable.  For completeness, at the regular generic points
write $t=u z^e$ for uniformizers and choose a separating residue
parameter $x$ on the target.  Then
\[
 d(u z^e)\wedge dx
 =z^{e-1}(eu\,dz\wedge dx+z\,du\wedge dx).
\]
The coefficient at order $e-1$ is nonzero exactly under those
two conditions.  This includes inseparable residue extensions.
The restriction of a divisorial valuation through a finite
function-field extension is again divisorial.  Formula
\eqref{D:eq-cusp-discrepancy} consequently proves that $S'$ is
log canonical.

Replace $\mathcal O_{S,s}$ by its strict henselization $A$.
It is an excellent normal local ring with residue field $k$.
Choose one connected component of the punctured cover, and let
$B$ be the normalization of $A$ in its field $L$.
It is finite over $A$ by excellence.  A finite algebra over a
henselian local ring is a product of local rings; since $B$ is
a domain it is local, and its residue field is $k$.  Thus $B$
is strictly henselian.  The chosen component is Galois with
its decomposition group, which we now denote by $G$.
Normality, the canonical-divisor comparison, Gorensteinness and
log canonicity persist under this ind-\'etale localization.
Equivalently, all the finite-presentation constructions and the
regular models used to compute a fixed discrepancy descend to
an \'etale neighborhood of the original algebraic singularity.
We continue to write $S=\operatorname{Spec}A$ and
$S'=\operatorname{Spec}B$, and put
$U=S\setminus\{s\}$ and $U'=S'\setminus\{s'\}$.
The latter is a connected finite \'etale $G$-cover of $U$.

\par\smallskip\noindent\textbf{The nonzero cycle quotient.}\quad
Take a regular resolution $Y\to S$ whose reduced exceptional
fibre $E$ is a cycle of smooth rational curves.  If the minimal
cycle consists of a single nodal curve, blow up its node once.
This operation leaves each old discrepancy equal to $-1$ and
gives the new component discrepancy $-1$ as well.
For a proper scheme over a henselian local ring, restriction
to the closed fibre is an equivalence on finite \'etale covers;
see \cite[Lemma 58.9.1, tag 0A48]{StacksBCuspLocal}.
Nilpotents do not change this category.  Therefore
\[
 \pi_1(Y)\simeq\pi_1(E)\simeq\widehat{\mathbb Z}.
\]
The last identification follows directly by normalizing the
nodal curve: each normalized component is $\mathbb P^1_k$,
whose finite \'etale covers are trivial, and the gluing at the
nodes is exactly the data of a finite covering of the cycle
graph.  The same gluing description shows that a rational tree
has trivial fundamental group and that attaching rational
trees does not change the fundamental group of a core.

Since $Y$ is normal, the dense-open inclusion $U\hookrightarrow Y$
induces a surjection on fundamental groups: a connected finite
\'etale cover of $Y$ is normal and integral, and remains
connected on its nonempty open preimage of $U$.
The open subgroup $\pi_1(U')\subset\pi_1(U)$ therefore has
finite-index, and hence nonzero, image in
$\widehat{\mathbb Z}$.

\par\smallskip\noindent\textbf{The core on the covering resolution.}\quad
Let $Z_0\to S'$ be the minimal regular resolution, with a node
blown up if necessary to obtain smooth components in its
cyclic case.  Resolve the normalized graph of the rational
map $Z_0\dashrightarrow Y$.  We obtain a regular surface $Z$
with a proper birational morphism $Z\to Z_0$ and a morphism
$Z\to Y$.  It is unchanged over $U'$.  Surface factorization
shows that $Z\to Z_0$ is a sequence of point blowups with
centres in the exceptional fibre.  Let $E'$ be the reduced
exceptional fibre of $Z\to S'$.
Proper henselian comparison is functorial, so the composite
\[
 \pi_1(U')\longrightarrow\pi_1(Z)
 \simeq\pi_1(E')\longrightarrow\pi_1(Y)
 \simeq\pi_1(E)
\]
is the finite-index nonzero homomorphism of the preceding
paragraph.  In particular $\pi_1(E')\to\pi_1(E)$ is nonzero.

Apply Lemma~\ref{D:Gorenstein-lc-configurations} to $S'$.
If its minimal exceptional divisor is empty or a rational tree,
then $E'$ is a point or a rational tree and $\pi_1(E')=1$,
a contradiction.
If its minimal fibre is elliptic, $E'$ consists of that
elliptic core with rational trees attached.  The irreducible
core maps into one component of $E$, or into a point.
Its map on fundamental groups factors through
$\pi_1(\mathbb P^1_k)=1$.  The attached trees contribute
nothing, again a contradiction.
Thus the covering singularity is also a cusp, and $E'$ has
a rational cycle core $C'$ with rational trees attached.
Every component of $C'$ has discrepancy $-1$.

The image of $C'$ must dominate every component of $E$.
Indeed, otherwise its connected reduced image is contained
in a proper connected subcurve of the smooth-component
cycle $E$, hence in a rational chain, or is a point.
Such a subcurve has trivial fundamental group.  Since
$\pi_1(C')\to\pi_1(E')$ is an isomorphism, this would make
$\pi_1(E')\to\pi_1(E)$ zero, contrary to what we have proved.
Consequently, for each target component $E_i$ there is a
component $F_i$ of the core $C'$ mapping dominantly to it.
This is an extension of the corresponding divisorial
valuation, with discrepancy $a(F_i,S')=-1$.

\par\smallskip\noindent\textbf{All extensions of each valuation.}\quad
Normalize $Y$ in $L$.  The extensions of the discrete
valuation at $E_i$ are the prime divisors above it, and the
Galois group $G$ acts transitively on them.  The normal
resolution $Z$ factors through this normalization.
The discrepancy $a(F,S')$ is intrinsic to the valuation
and is invariant under deck transformations of $S'$:
its canonical Cartier class is the pullback of the one on
$S$.  Since one extension $F_i$ has discrepancy $-1$,
all extensions $F/E_i$ have discrepancy $-1$.
The target discrepancy is also $-1$, and
\eqref{D:eq-cusp-discrepancy} now gives
\[
 -1=-e(F/E_i)+\delta_F,
 \qquad \delta_F=e(F/E_i)-1.
\]
The equality criterion proves tameness, with separable
residue extension.  Starting with a target resolution
in which any prescribed cycle node has already been
blown up proves the additional assertion.

Finally, strict henselization has not weakened the
valuation conclusion for the original algebraic cover.
The exceptional curves and their residue function fields
over the closed point remain the same after this base
change, since its residue field is already $k$; the
extension ramification indices and residue extensions
are the ones computed in the corresponding local
components.  More explicitly, the induced maps of the
boundary discrete valuation rings are ind-\'etale with
unchanged residue fields, so these indices and residue
extensions are preserved.  Thus tameness for those components gives
the stated assertion before strict localization.
\end{proof}

\begin{remark}[Exact local scope]
The proof uses strict henselization to make the finite
normal covering ring local and to compare finite \'etale
covers of a proper resolution with its exceptional fibre.
It does not assert an equivalence of punctured algebraic,
strictly henselian and completed neighborhoods, and no
completion comparison is needed.  The Gorenstein
hypothesis is used twice: to make every covering surface
Gorenstein by canonical Cartier pullback, and to make
the minimal-resolution discrepancies integral.  We have
not replaced it by an arbitrary log canonical hypothesis.
\end{remark}



% Additional extension; D actual-file and primary-source audit passed.

\begin{lemma}[Proved: from an SNC core to all exceptional valuations]
\label{D:Gorenstein-lc-valuative-tame}
Let $(S,s)$ be a Gorenstein log canonical surface singularity
over an algebraically closed field, and suppose that it is not
log terminal.  Every finite normal separable cover of $S$ that
is \'etale on the punctured germ is tame at every divisorial
valuation of $k(S)$ centred at $s$, including separability of
the corresponding residue extension.
\end{lemma}
\begin{proof}
By Lemma~\ref{D:Gorenstein-lc-configurations}, a minimal regular
resolution has either one smooth elliptic curve or a rational
nodal cycle as reduced exceptional fibre.  In the latter case
we blow up a possible self-node once.  We obtain a regular
resolution $Z\to S$ whose reduced exceptional divisor $E$ has
simple normal crossings.  Every cover in the statement is
tame along all components of $E$, by
Proposition~\ref{D:lc-elliptic-tame} or
Theorem~\ref{D:cusp-tame}.

We recall precisely the local persistence of this tameness.
Let $x\in E$, let $R=\mathcal O_{Z,x}^{\mathrm{sh}}$, and let
$f_1,\ldots,f_r$, with $r\leq2$, be the regular parameters
defining the boundary components through $x$.
The absolute Abhyankar lemma and its strictly local
corollary, \cite[XIII, Proposition 5.2 and Corollary 5.3,
pp.~428--430]{SGA1}, say that every connected component
of the punctured cover over $R$ is dominated by
\[
 R[T_1,\ldots,T_r]/
       (T_1^{n_1}-f_1,\ldots,T_r^{n_r}-f_r),
 \qquad n_i\text{ prime to the characteristic exponent}.
\]
The associated connected function-field extension is
Galois of prime-to-$p$ degree: all the required roots
of unity are present.  It is therefore tame at every
discrete valuation.  Every subextension is also tame.
This assertion descends from the strict localization:
an \'etale neighborhood of $x$ induces an unramified
extension at a valuation centred at $x$, and tameness
can be checked after such extensions.  One may equally
descend the finite Kummer domination to a finite
\'etale neighborhood and use the multiplicativity of
ramification indices and separability of residue
extensions.

For any specified divisorial valuation centred at $s$,
a further regular resolution extracts it.  A common
regular resolution with $Z$ is obtained from $Z$ by
point blowups, with centres in $E$ and its successive
inverse images.  The preceding local argument proves
tameness at each new exceptional divisor; the reduced
boundary remains simple normal crossings after every
such blowup.  It proves tameness at the specified
valuation, as required.  The tower and subextension
facts for tameness used here are, in particular,
\cite[Lemmas 15.116.5--9, section tag 0EXT]{StacksBIndexTame}.
\end{proof}

\begin{theorem}[Proved under stated hypotheses: non-klt points of prime-to-$p$ index]
\label{D:prime-to-p-index-tame}
Let $(S,s)$ be a normal log canonical surface singularity
over an algebraically closed field $k$, and suppose that it
is not log terminal.  Let $n$ be the local Cartier index of
$K_S$ at $s$.  If $n$ is prime to the characteristic exponent
of $k$, then every finite normal separable cover of $S$
that is \'etale on the punctured germ is tame at every
divisorial valuation centred at $s$.
\end{theorem}
\begin{proof}
We work in the strict henselization $A$ of the local ring
at $s$, as in Theorem~\ref{D:cusp-tame}; passing back to
the algebraic germ preserves the valuation conclusion.
The local index after this base change divides $n$, so
we may still use the given prime-to-$p$ integer $n$.
Write $U=\operatorname{Spec}A\setminus\{s\}$.
The canonical sheaf on $U$ is an invertible sheaf $L$,
and a local trivialization of $\omega_A^{[n]}$ gives
an isomorphism $L^{\otimes n}\simeq\mathcal O_U$.
Here square brackets denote the reflexive power.

\par\smallskip\noindent\textbf{Construction and properties of the cover.}\quad
The algebra
\[
 \mathcal O_U\oplus L^{-1}\oplus\cdots\oplus L^{-(n-1)},
\]
with multiplication specified by that trivialization,
defines a $\mu_n$-torsor over $U$.  Since $n$ is invertible,
it is finite \'etale.  Choose one connected component and
normalize $\operatorname{Spec}A$ in its function field
$K_1$.  Excellence makes the resulting map
$\pi:T\to\operatorname{Spec}A$ finite; it is normal and
quasi-\'etale, and its punctured open is the chosen torsor
component.  Because $k$ contains the $n$th roots of unity,
that component is a connected cyclic Galois cover of
degree $d\mid n$.  Explicitly, its stabilizer in the
constant group $\mu_n(k)$ is a cyclic group acting
simply transitively on each fibre of the component.
The covering normal ring is local, henselian, and has
residue field $k$.

On the torsor the line bundle $L$ has its tautological
trivialization.  Its pullback to every component is
therefore trivial.  Since the map is \'etale on $U$,
the canonical sheaf of the punctured $T$ is trivial.
Normality extends this equality of rank-one reflexive
sheaves across its omitted point, and gives
$\omega_T\simeq\mathcal O_T$.  A normal surface is
Cohen--Macaulay, so $T$ is Gorenstein.  In particular,
this conclusion does not require a hypothesis that
the reflexive cyclic algebra before normalization
was already regular or normal.

The cyclic extension $K_1/k(S)$ has prime-to-$p$ degree.
Its wild inertia is therefore trivial at every discrete
valuation.  Thus all its ramification indices are prime
to $p$, its residue extensions are separable, and the
different determinant has order $e-1$.  For every
divisorial valuation $F/E$ in this extension, the
canonical-divisor comparison yields
\begin{equation}\label{D:eq-index-discrepancy}
 a(F,T)+1=e(F/E)\bigl(a(E,S)+1\bigr).
\end{equation}
It follows that $T$ is log canonical.  Since $S$ is not
log terminal, some exceptional valuation $E$ has
discrepancy $-1$.  Every extension of it then has
discrepancy $-1$ by \eqref{D:eq-index-discrepancy}.
Thus $T$ is not log terminal either.  We have obtained
a Gorenstein non-klt log canonical cover without any
change of residue field.

\par\smallskip\noindent\textbf{Comparison for an arbitrary cover.}\quad
Let $M/k(S)$ be the finite separable extension of any
cover in the statement, and choose a compositum
$M_1=M K_1$.  Normalize $T$ in $M_1$.
It is a finite normal cover of $T$, \'etale on the
punctured germ, because this is true after base change
on that open.  By
Lemma~\ref{D:Gorenstein-lc-valuative-tame}, it is tame
at every exceptional divisorial valuation of $K_1$.
Every extension to $K_1$ of a divisorial valuation of
$k(S)$ centred at $s$ is still divisorial and is centred
at the closed point of $T$.

Consequently, at any such valuation the tower
\[
 k(S)\ \subset\ K_1\ \subset\ M_1
\]
is tame in both steps, and hence in its composite.
The subextension $M/k(S)$ is tame as well.
These implications include residue separability:
ramification indices multiply, and intermediate
fields of a finite separable residue extension are
separable.  See the opened proofs of
\cite[Lemmas 15.116.5 and 15.116.6]{StacksBIndexTame}.
This establishes the assertion for the original
arbitrary cover, not only for the cyclic index cover.
\end{proof}

\begin{corollary}[Proved: mixed omitted points with prime-to-$p$ lc index]
\label{D:prime-to-p-index-boundary}
Let $X$ be a normal proper integral surface over an
algebraically closed field, and let $Y\subset X$ be a
finite set of closed points.  Suppose that each point
of $Y$ is either log terminal, with no restriction on
its Cartier index, or log canonical with Cartier
index prime to the characteristic exponent.
Then $\pi_1(X\setminus Y)$ is topologically finitely
generated.  All singularities retained in the open
surface may be arbitrary.
\end{corollary}
\begin{proof}
Take a fixed regular proper resolution
$q:\widetilde X\to X$; its source is a projective
surface.  A finite \'etale cover of $X\setminus Y$
normalizes to a finite quasi-\'etale cover of $X$.
For every exceptional divisor above a non-klt point
of $Y$, Theorem~\ref{D:prime-to-p-index-tame} makes
the wild inertia image zero.  For each log terminal
point, the image is finite because it factors
through the finite punctured strict-local group
of \cite[Theorem 5.3, p.~739]{CarvajalRojasYasuda2025B}.
The fixed projective modification thus satisfies
Proposition~\ref{D:finite-exceptional-wild}.
As $Y$ has codimension two, that proposition gives
the asserted finite generation of the full group.
\end{proof}

\begin{remark}[Restriction on the index]
We have used the index hypothesis to obtain a separable
cyclic cover tame at every valuation and the exact
discrepancy equality \eqref{D:eq-index-discrepancy}.
When the index is divisible by the characteristic,
the cyclic construction can be inseparable, so the
\'etale canonical-pullback and tame-discrepancy
arguments just used do not apply.  This proof does
not establish preservation of log canonicity in
that case, and the preceding theorem makes no
assertion about it.  The failure of preservation
of \emph{log terminal} singularities by index-one
covers in small characteristic is explicitly
discussed in the introduction and Remark 5.5 of
\cite{CarvajalRojasYasuda2025B}.
\end{remark}


% General log canonical extension, separate from the checked Gorenstein fallback.
% Root, C and D independent actual-file audits passed.

\begin{lemma}[Proved: exceptional configurations for arbitrary lc surfaces]
\label{D:lc-exceptional-configurations}
Let $(S,s)$ be a normal log canonical surface singularity over
an algebraically closed field, in arbitrary characteristic.
On the minimal regular resolution its reduced exceptional
divisor has one of the following forms:
\begin{enumerate}
\item a simple normal crossings tree of smooth rational curves;
\item one smooth elliptic curve, with discrepancy $-1$;
\item a rational nodal cycle, with discrepancy $-1$ on every
component, allowing one irreducible nodal rational curve.
\end{enumerate}
The exceptional divisor is empty if the point is regular.
No Gorenstein hypothesis or restriction on the Cartier
index is imposed.  In the tree case we do not assert that
the singularity is rational or log terminal, or that its
discrepancies are zero.
\end{lemma}
\begin{proof}
Let $f:Z\to S$ be the minimal regular resolution and write
\[
 K_Z=f^*K_S+\sum_i a_i E_i.
\]
The exceptional fibre is connected and its intersection
matrix is negative definite.  The surface facts and exact
sources used here are the ones in
Lemma~\ref{D:Gorenstein-lc-configurations}:
\cite[section 10.1, Theorem 10.1]{Kollar2013D} and
\cite[tags 0BGP, 0C2M and 0C5R]{StacksDSurfaceResolution}.
Adjunction and minimality give
\[
 K_Z\cdot E_i=2p_a(E_i)-2-E_i^2\geq0.
\]
The positive-part argument for the negative definite
intersection matrix therefore gives $a_i\leq0$, exactly
as in the proof of that lemma.  Log canonicity gives
$a_i\geq-1$.  We use only
\begin{equation}\label{D:eq-lc-discrepancy-interval}
                  -1\leq a_i\leq0;
\end{equation}
the coefficients need not be integral.

We first record a useful consequence of a zero discrepancy.
If $a_j=0$, then
\[
 0\leq K_Z\cdot E_j
    =\sum_{i\ne j}a_i(E_i\cdot E_j)\leq0.
\]
Hence every component adjacent to $E_j$ also has
discrepancy zero.  In particular a zero-discrepancy
component cannot meet a component of negative discrepancy.

Suppose that some $E_i$ has arithmetic genus at least one.
From \eqref{D:eq-lc-discrepancy-interval} we obtain
\[
 -E_i^2\leq K_Z\cdot E_i
       =\sum_j a_j(E_j\cdot E_i)
       \leq a_iE_i^2\leq -E_i^2.
\]
All terms are equal.  Thus $p_a(E_i)=1$, $a_i=-1$,
and every component meeting $E_i$ would have discrepancy
zero.  The preceding observation excludes all such
neighbors.  By connectedness, $E_i$ is the whole fibre.
We will determine its possible singularities below.

We may now suppose that every component has arithmetic
genus zero, and hence is a smooth rational curve.
Write
\[
 b_i=-E_i^2\geq2,\qquad c_i=1+a_i\in[0,1].
\]
The lower bound on $b_i$ is minimality.
Form the intersection multigraph, placing
$E_i\cdot E_j$ edges between distinct vertices $i,j$.
This counts intersection multiplicities, before any
normal-crossings property has been established.
If this multigraph contains a cycle, select one with
distinct vertices, allowing a two-vertex cycle consisting
of two selected parallel edges.
Let $I$ be its vertex set and $s_{ij}$ the multiplicity
of its selected edges.  Thus
$\sum_{j\in I}s_{ij}=2$ for every $i\in I$.
For internal extra edges put
$r_{ij}=E_i\cdot E_j-s_{ij}\geq0$.
For an external vertex $j\notin I$ put
$m_{ij}=E_i\cdot E_j$.

Adjunction, or equivalently intersecting
$K_Z+\sum_iE_i=f^*K_S+\sum_i c_i E_i$
with $E_i$, gives, for $i\in I$,
\begin{equation}\label{D:eq-cycle-extra-edges}
 b_i c_i-\sum_{j\in I}s_{ij}c_j
 =-\sum_{j\in I}r_{ij}(1-c_j)
  -\sum_{j\notin I}m_{ij}(1-c_j).
\end{equation}
Summing over the selected cycle yields
\begin{equation}\label{D:eq-cycle-sum}
 \sum_{i\in I}(b_i-2)c_i
 =-\sum_{\{i,j\}\subset I}r_{ij}(2-c_i-c_j)
  -\sum_{\substack{i\in I\\j\notin I}}m_{ij}(1-c_j).
\end{equation}
Every term on the left is nonnegative and every term
on the right is nonpositive.  All are therefore zero.
In particular, each extra internal edge would have
both endpoints $c_i=c_j=1$, and every external
neighbor of the selected cycle would have $c_j=1$.

Negative definiteness guarantees some $b_i>2$ on the
selected cycle.  Otherwise the vector equal to one
on $I$ and zero elsewhere would have square
\[
 -2|I|+2|I|+2\sum_{\{i,j\}\subset I}r_{ij}\geq0,
\]
which is impossible.  Thus $c_i=0$ at some vertex.
All the right-hand terms in
\eqref{D:eq-cycle-extra-edges} have already vanished,
so the equations reduce to
$b_i c_i=\sum_j s_{ij}c_j$.
Nonnegativity propagates this zero around the selected
cycle.  Hence $c_i=0$ for every $i\in I$, or
$a_i=-1$ there.
An extra internal edge is now impossible, since its
endpoints would have to be one.  An external neighbor
would have $a_j=0$, impossible next to $a_i=-1$ by
the zero-discrepancy observation.  Connectedness
therefore shows that the whole exceptional divisor
is this cycle, with all discrepancies $-1$.
This argument covers intersection multiplicities
larger than two; none were discarded as if they
were single edges.

We check nodality in both cases in which all
discrepancies have been shown to be $-1$.
For the reduced exceptional curve $E$ we have
$K_Z+E=f^*K_S$, so $(Z,E)$ is a log canonical pair.
At a point of multiplicity at least three, the
point blowup would have discrepancy at most $-2$.
At a multiplicity-two point with one repeated
tangent, blow up once.  The strict transform
meets the new discrepancy-$-1$ exceptional curve
at one point with intersection multiplicity two.
If it is singular there, their total multiplicity
is at least three.  If it is smooth, it is tangent
to that exceptional curve, and the next blowup
creates a point with total reduced-boundary
multiplicity at least three.  A subsequent blowup
again gives discrepancy at most $-2$.
Thus every singularity of $E$ is an ordinary node,
in every characteristic.
An integral nodal curve of arithmetic genus one
is either smooth elliptic or a rational curve
with one node.  In the rational multigraph cycle
case, nodality gives exactly a rational nodal
cycle, including two components meeting at two
distinct nodes.

Finally, if the rational intersection multigraph
has no cycle, it is a tree.  There can be no
intersection multiplicity greater than one, since
this would give a two-edge cycle, and no triple
point, since this would give a three-cycle.
All components are smooth rational and all
intersections are transverse, so it is a simple
normal crossings rational tree.  This proves all
the alternatives.
\end{proof}

\begin{theorem}[Proved: tame ramification above an arbitrary lc cusp]
\label{D:lc-cusp-tame}
Let $(S,s)$ be a normal log canonical surface singularity
over an algebraically closed field.  Suppose that the
minimal exceptional fibre is a rational nodal cycle.
Every finite normal separable cover of $S$ that is
\'etale on the punctured germ is tame at every component
of that cycle, and at the components obtained by
subdividing its nodes.  No Gorenstein or index
restriction is required.
\end{theorem}
\begin{proof}
We indicate explicitly which argument replaces the
Gorenstein input in Theorem~\ref{D:cusp-tame}, and
repeat the decisive comparison.
The target discrepancies on the cycle are $-1$ by
Lemma~\ref{D:lc-exceptional-configurations}.
Pass to a Galois closure of the punctured cover
and normalize.  The canonical-divisor equality
$K_{S'}=\pi^*K_S$ holds as an equality of rational
Cartier divisors, because the cover is \'etale
in codimension one.  It makes $S'$
$\mathbb Q$-Gorenstein.  For every divisorial
valuation extension $F/E$, the determinant
calculation of Proposition~\ref{D:lc-elliptic-tame}
gives
\[
 a(F,S')=e(F/E)a(E,S)+\delta_F,
 \qquad \delta_F\geq e(F/E)-1.
\]
Consequently $S'$ is log canonical.  No
integrality of its discrepancies is asserted.
We may apply the new configuration lemma to
its minimal resolution.

Work in the strict henselian local germ, choosing
a connected covering component and its decomposition
group.  The finiteness, normality and residue-field
comparisons are exactly those verified in
Theorem~\ref{D:cusp-tame}.  Let $Y\to S$ be a
regular resolution with smooth-component cycle
$E$, after one node blowup if needed.
Proper henselian comparison
\cite[Lemma 58.9.1, tag 0A48]{StacksBCuspLocal}
and normalization of the nodal curve give
\[
 \pi_1(Y)\simeq\pi_1(E)
       \simeq\widehat{\mathbb Z}.
\]
The image of the finite-index subgroup
$\pi_1(S'\setminus\{s'\})$ in this quotient
is again finite-index and nonzero.

Resolve the graph from a minimal regular resolution
of $S'$ to $Y$.  Let $Z\to S'$ be the resulting
regular resolution and let $E'$ be its reduced
exceptional fibre.  Functorial proper henselian
comparison shows that $\pi_1(E')\to\pi_1(E)$
has that nonzero image.  Point-blowup factorization
shows that $E'$ is obtained from its minimal core
by rational tree attachments and cycle-node
subdivisions.
If the covering configuration is a rational tree
or a regular point, then $\pi_1(E')=1$.
If it is elliptic, its irreducible elliptic core
maps into one rational component of $E$, so
its fundamental-group image is again zero.
Both are impossible.  Thus the covering
configuration is a rational cycle, and the
new classification gives discrepancy $-1$
on every component of its cyclic core.

The source cyclic core must dominate every
component of $E$: otherwise its connected image
lies in a proper rational chain of $E$, which
has trivial fundamental group.  This would
contradict the nonzero image just established.
For each component $E_i$, we have therefore
found an extension $F_i$ of its valuation with
$a(F_i,S')=-1$.
The Galois group acts transitively on the
valuation extensions of $E_i$ in the normalized
cover of $Y$.  Discrepancy is invariant under
these deck transformations.  Every extension
$F/E_i$ thus has discrepancy $-1$, and
\[
 -1=-e(F/E_i)+\delta_F.
\]
Equality $\delta_F=e(F/E_i)-1$ says precisely
that the extension is tame, including residue
separability.  The strict henselian comparison
preserves these indices and residue extensions,
so the conclusion holds for the algebraic germ.
The same proof begins with any subdivision of
the target cycle and proves the final assertion.
\end{proof}

\begin{corollary}[Proved: all valuations above an lc elliptic or cyclic core]
\label{D:lc-core-valuative-tame}
If the minimal exceptional fibre of a normal
log canonical surface singularity over an
algebraically closed field is elliptic or cyclic,
then every finite normal separable cover that
is \'etale on the punctured germ is tame at every
divisorial valuation centred at the singular point.
No restriction on the Cartier index is imposed.
\end{corollary}
\begin{proof}
After blowing up a possible self-node, the
minimal core is a simple normal crossings
divisor.  Proposition~\ref{D:lc-elliptic-tame}
and Theorem~\ref{D:lc-cusp-tame} give tameness
on its components.  The Kummer persistence
argument of Lemma~\ref{D:Gorenstein-lc-valuative-tame}
uses only this regular normal-crossings model
and its componentwise tameness.  Applying it
here proves tameness on every further point
blowup, and hence at every specified exceptional
divisorial valuation.  The cited argument is
\cite[XIII, Proposition 5.2 and Corollary 5.3,
pp.~428--430]{SGA1}; it uses no Cartier-index
condition after componentwise tameness is known.
\end{proof}


% Independently audited in full by root, B, and D before integration.
\begin{lemma}[Proved: formal lift group for a punctured surface cover]
\label{D:formal-torus-lift}
Let $k$ be algebraically closed of characteristic $p>0$, let
$R=\bigoplus_{d\geq0}R_d$ be a finitely generated normal graded
$k$-domain of dimension two with $R_0=k$, and let $A$ be its
completion at $R_+$.  Let $B/A$ be the normal extension belonging
to a connected finite Galois cover of the punctured spectrum of
$A$, with group $G$.  The grading action of $\mathbb G_m$ on
$A$ lifts to a continuous action on $B$ after a finite \'etale
isogeny $\mathbb G_m\to\mathbb G_m$.  This lifted torus
centralizes the deck group $G$.
\end{lemma}
\begin{proof}
We first construct the lift group on all test algebras; this
avoids assuming that a pointwise Isom functor is of finite type.
Choose homogeneous parameters $x,y\in R_+$ such that $R$ is
finite over $k[x,y]$, and set $C=k[[x,y]]$.  The rings $A$ and
$B$ are complete normal local rings of dimension two, with
residue field $k$.  They are Cohen--Macaulay and finite free over
$C$, by miracle flatness
\cite[tags 0B3D and 00R4]{C_StacksFormalLift}.
For a $k$-algebra $T$ put
\[
 A_T=A\otimes_C T[[x,y]],\qquad
 B_T=B\otimes_C T[[x,y]].
\]
These are also the completed scalar extensions
$\varprojlim_r(A/\mathfrak m_A^r)\otimes_kT$ and
$\varprojlim_r(B/\mathfrak m_B^r)\otimes_kT$:
the parameter-adic and maximal-ideal-adic topologies are cofinal,
and finite freeness over $C$ identifies the parameter-adic limits.

For each $r$, consider pairs consisting of $c\in\mathbb G_m(T)$
and an augmented $T$-algebra automorphism of
$(B/\mathfrak m_B^r)\otimes_kT$ whose restriction to $A$
is the grading action of $c$, followed by the quotient map.
This is an affine group scheme $E_r$ of finite type over $k$.
Indeed, the automorphisms are matrices preserving a finite
multiplication table, the unit, and the augmentation, with
invertible determinant.  Compatibility with $A$ need only be
imposed on finitely many homogeneous algebra generators of $R$;
these are polynomial equations in the matrix entries and
$c,c^{-1}$.  Augmented automorphisms preserve every power of the
augmentation ideal.  Consequently the truncation maps define
an inverse system, and
\[
 E=\varprojlim_r E_r
\]
is an affine group scheme.  Its $T$-points are exactly the pairs
$(c,\phi)$ with $\phi$ a continuous augmentation-preserving
automorphism of $B_T$ lifting $c$ on $A_T$.  Write
$\rho:E\to\mathbb G_m$ for the projection.

We claim that $\ker\rho$ is the constant finite group scheme
$G$, including on nonreduced test algebras $T$.  Let
\[
 U_T=\operatorname{Spec}(A_T)\setminus V(x,y),\qquad
 V_T=\operatorname{Spec}(B_T)\setminus V(x,y).
\]
The map $V_T\to U_T$ is the base change of the given $G$-torsor.
Because $B_T$ is free over $T[[x,y]]$, the sequence $x,y$ is
regular on $B_T$.  The two-open \v Cech calculation gives
\[
 \Gamma(V_T,\mathcal O)=B_T.
\]
In particular, clopen decompositions of $V_T$ are the idempotent
decompositions of $B_T$.  Completeness lifts idempotents uniquely
from $B_T/(x,y)B_T$; the latter is
$(B/(x,y)B)\otimes_kT$.  Since $B/(x,y)B$ is a local Artin
$k$-algebra with residue field $k$, its nilpotent augmentation
ideal shows that these idempotents are exactly those of $T$
\cite[Section 15.11, tags 09XD, 09XI and 0ALI]{C_StacksFormalLift}.
Every automorphism of the finite \'etale cover $V_T/U_T$
is given by a locally constant deck element on $V_T$; by the
preceding idempotent calculation that element is locally constant
on $\operatorname{Spec}T$.  Thus it belongs to $G(T)$.
The Hartogs equality makes restriction injective, and all such
deck elements extend to $B_T$.  This proves the claim on every
test algebra, not just on geometric points.

The ring $A$ is excellent and strictly henselian.  Gabber's local
finiteness theorem says that $H^1(U,G)$ is finite for every finite
group $G$, where $U$ is its punctured spectrum
\cite[XXI, Theorem 1.4, p. 529]{ILOXXID}.
The divisible abelian group $k^*$ acts on this finite set by
scaling.  Its image in a finite permutation group is both
divisible and finite, hence trivial.  Each scaling therefore
lifts to the cover and to its normal finite algebra $B$.
It follows that $E(k)\to k^*$ is surjective.  In particular,
$k[t,t^{-1}]\to\mathcal O(E)$ is injective.

Faithful flatness of a commutative Hopf algebra over any Hopf
subalgebra, with no finite-type restriction, now makes $\rho$
faithfully flat
\cite[VI, Theorem 11.1, pp. 82--85]{C_MilneAGS2012}.
It is affine, hence fpqc.  The identity
$E\times_{\mathbb G_m}E\simeq E\times G$ and fpqc descent
show that $\rho$ is finite \'etale.  Its identity component
$E^0$ is a connected one-dimensional affine smooth group with
a finite-kernel surjection to $\mathbb G_m$; it is a torus.
Indeed, the only connected smooth affine groups of dimension
one over $k$ are $\mathbb G_a$ and $\mathbb G_m$, and
$\mathbb G_a$ has no nontrivial character
\cite[Corollary 16.16, p. 329]{C_MilneAG2022}.
After choosing its coordinate, $E^0\to\mathbb G_m$ is
$t\mapsto t^n$ with $n>0$ prime to $p$.  Conjugation by this
connected group on the constant finite group $G$ is trivial.
This proves the required lift and centralization.
\end{proof}

\begin{proposition}[Proved: radial tameness for completed graded surface cones]
\label{D:graded-cone-tame}
In the situation of Lemma~\ref{D:formal-torus-lift}, assume the
positive grading is effective, so that its occurring degrees
have greatest common divisor one.  Every extension to
$\operatorname{Frac}(B)$ of the radial divisorial valuation of
$\operatorname{Frac}(A)$ is tame.  Thus a normal surface germ
whose completed resolution along its exceptional curve is the
completion of the zero section in a negative line bundle over
a smooth proper curve has no wild inertia at that exceptional
valuation in any finite cover \'etale on the punctured germ.
\end{proposition}
\begin{proof}
Let $\widetilde T=E^0$ be the lifted torus, acting on $A$ with
weights $n$ times the original weights.  The finite-dimensional
quotients $B/\mathfrak m_A^rB$ are rational torus
representations.  Exactness of weight spaces lets us lift a
finite weight basis of $B/\mathfrak m_AB$ compatibly through
all these quotients.  We obtain eigenvectors $b_1,\ldots,b_s$
which generate the finite $A$-module $B$ by Nakayama's lemma.
If $w_j$ is the weight of $b_j$, projection of an expression
$b=\sum a_jb_j$ to a weight $d$ shows that every eigenvector
of weight $d$ is an $R$-linear combination of these $b_j$.
The coefficient of $b_j$ is the homogeneous component of
$a_j$ with torus weight $d-w_j$, which lies in $R$.

It follows that the direct sum $S$ of the weight spaces in $B$
is a finite graded $R$-algebra and that its weights are bounded
below.  A nonzero element of negative weight would have powers
of arbitrarily negative weight, contradicting this bound and
the fact that $B$ is a domain.  Thus $S$ is nonnegatively graded.
Its degree-zero part is a finite-dimensional domain over the
algebraically closed field $k$, hence equals $k$.
The same eigenvector expansions identify $B$ with the product
of its nonnegative weight spaces: the weight filtration and
the maximal-ideal filtration are cofinal because $S$ is finite
over the positively graded ring $R$.  In particular, $B$ is
the completion of $S$ at its positive ideal.

The group $G$ commutes with $\widetilde T$, hence preserves $S$
and every weight space.  Since $B^G=A$, taking weight spaces
gives $S^G=R$.  Put $L_0=\operatorname{Frac}(S)$ and
$K_0=\operatorname{Frac}(R)$.  Then $L_0/K_0$ is Galois with
group $G$: the action is faithful because $S$ is dense in $B$,
and invariants of the fraction field are fractions of invariants.

Let $g$ be the greatest common divisor of the positive weights
occurring in $S$.  Lowest weight divided by $g$ defines the
radial discrete valuation on $L_0$, and also on
$\operatorname{Frac}(B)$ by the completed weight expansion.
It is preserved by $G$.  Its restriction to the base radial
valuation has ramification index
\[
 e=n/g.
\]
Indeed, the weights of $R$ in $S$ are $n$ times its effective
original weights.  Thus $g\mid n$, and $e$ is prime to $p$.
The residue field of the radial valuation on $L_0$, or on the
completed fraction field, is the field of fractions of
homogeneous expressions of equal degree; equivalently it is
$L_0^{\widetilde T}$.  On the base it is
$K_0^{\mathbb G_m}$.  Since the two actions commute,
\[
 \bigl(L_0^{\widetilde T}\bigr)^G
   =\bigl(L_0^G\bigr)^{\widetilde T}
   =K_0^{\mathbb G_m}.
\]
The residue extension is therefore separable, being an extension
of a fixed field under a finite group.  The exhibited valuation
is $G$-invariant; transitivity of $G$ on prolongations in the
Galois extension shows it is the only prolongation.  Its
prime-to-$p$ ramification index and separable residue extension
prove tameness.

For the final assertion, a negative line bundle $M$ on a smooth
proper curve $C$ has section ring
$R=\bigoplus_{d\geq0}H^0(C,M^{-d})$.  This is a finitely
generated normal ring.  Finite generation follows from ampleness
of $M^{-1}$; normality follows by intersecting the corresponding
divisorial valuation rings inside $k(C)(t)$.  The functions on
the formal zero-section neighborhood are
$\prod_{d\geq0}H^0(C,M^{-d})$, which is the completion of $R$
at its positive ideal.  For a proper normal surface contraction,
degree-zero formal functions identifies this ring with the
completed local ring downstairs: the exceptional-ideal and
maximal-ideal topologies on the resolution are cofinal
\cite[Theorem 30.20.5, tag 02OC]{C_StacksFormalPic}.
Its radial valuation is exactly the valuation of $C$.
The degrees are effective since all sufficiently large degrees
have nonzero sections.  A finite cover of the original
punctured germ pulls back to a disjoint union of connected
finite covers of this completed puncture.  Each is dominated
by a Galois cover, to which the preceding argument applies.
The formal line-bundle identification supplies a common genuine
completed boundary DVR $k(C)[[t]]$: both the original boundary
DVR and the radial DVR on the vertex completion have this
completion, compatibly with the maps from their function fields.
We do not identify vertex completion with DVR completion.
Completion preserves the ramification index and separability
of the residue extension in a finite extension of excellent
DVRs.  Testing the cover over this common complete DVR therefore
descends tameness to the original boundary valuation.
The same argument applies when
we identify the completed resolved germ with this formal
line-bundle neighborhood.
\end{proof}


\subsubsection{A numerical condition forcing a formal cone}

\begin{lemma}[Proved under stated hypotheses: formal linearization of a curve]
\label{D:negative-curve-linearization}
Let $C$ be a smooth proper connected curve in a smooth surface
$Y$ over an algebraically closed field $k$, of any characteristic.
Write $I=\mathcal O_Y(-C)$ and
$N=N_{C/Y}=(I/I^2)^\vee$.  Suppose
\[
 \begin{aligned}
 H^1(C,T_C\otimes N^{-n})&=0 &&(n\geq1),\\
 H^1(C,N^{1-n})&=0 &&(n\geq2).
 \end{aligned}
 \tag{L}
\]
Then the formal completion of $Y$ along $C$ is isomorphic to the
formal completion of the total space of $N$ along its zero
section.  The isomorphism can be chosen to induce the identity
on $C$ and on its conormal bundle.  In particular, if $g$ is
the genus of $C$ and $C^2=-d<0$, the numerical condition
\[
                   d>\max\{0,4g-4\}
 \tag{N}
\]
implies this conclusion.
\end{lemma}

\begin{proof}
Put $C_n=(C,\mathcal O_Y/I^n)$ for $n\geq1$; thus $C_1=C$.
Let $V=\operatorname{Tot}(N)$, let $J$ be its zero-section ideal,
and write $V_n=(C,\mathcal O_V/J^n)$.  We construct compatible
isomorphisms $f_n:C_n\to V_n$ inducing the specified identities.
The distinction between the step $n=1$ and all later steps is
essential for the exact obstruction groups.

First consider the square-zero extension
\[
 0\longrightarrow N^{-1}\longrightarrow\mathcal O_{C_2}
   \longrightarrow\mathcal O_C\longrightarrow0.
\]
Local $k$-algebra sections of its last map exist because $C$ is
smooth.  The difference of two sections is a $k$-derivation
$\mathcal O_C\to N^{-1}$, and adding any such derivation to a
section gives another section.  Thus the sheaf of local sections
is a torsor under $T_C\otimes N^{-1}$.  These statements are the
infinitesimal lifting criterion and square-zero derivation
calculation in
\cite[Lemmas 37.11.7 and 17.28.11, tags 02H6 and 01UP]
     {StacksBNegativeCurve}.
The first vanishing in (L), at $n=1$, makes this torsor trivial.
A global section splits the extension as a square-zero
$\mathcal O_C$-algebra, giving
\[
 \mathcal O_{C_2}\simeq\mathcal O_C\oplus N^{-1}
                       =\mathcal O_{V_2}.
\]
Using the specified identification of $I/I^2$ with $N^{-1}$
gives $f_2$.  In particular, no vanishing of $H^1(C,\mathcal O_C)$
is required at this step.

Now let $n\geq2$ and suppose $f_n$ has been constructed.
Regard it as a morphism $C_n\to V$.  The immersion
$C_n\subset C_{n+1}$ has square-zero ideal
\[
 I^n/I^{n+1}=N^{-n},
\]
annihilated by the ideal of $C$ in $C_n$.  Since $V$ is smooth
over $k$, $f_n$ lifts Zariski locally to $C_{n+1}\to V$.
Differences of local lifts form the sheaf
\[
 \mathcal H_n
  =\mathcal Hom_{\mathcal O_{C_n}}
       (f_n^*\Omega_{V/k},I^n/I^{n+1})
  =(T_V|_C)\otimes N^{-n}.
\]
The sheaf of lifts is a torsor under $\mathcal H_n$;
see \cite[Lemmas 37.9.4--5 and 37.11.7,
tags 04FH, 04FJ and 02H6]{StacksBNegativeCurve}.
The projection $V\to C$ and its zero section give the canonical
splitting
\[
 T_V|_C\simeq T_C\oplus N.
\]
Consequently the exact obstruction group at this step is
\[
 H^1(C,\mathcal H_n)
  =H^1(C,T_C\otimes N^{-n})\oplus H^1(C,N^{1-n}),
                  \qquad n\geq2.
 \tag{O}
\]
It vanishes by (L), so a global lift $f_{n+1}:C_{n+1}\to V$
exists and agrees with the previously chosen $f_n$.

Because this lift restricts to the identity $C\to C$, it sends
the zero-section ideal $J$ into $I/I^{n+1}$.  It therefore
factors through $V_{n+1}$.  Its map on $J/J^2$ agrees with that
of $f_2$ and is the prescribed identity.  Both embeddings are
effective Cartier divisors, so multiplication identifies
\[
 J^r/J^{r+1}\simeq N^{-r}\simeq I^r/I^{r+1}.
\]
The resulting map on every associated graded piece, for
$0\leq r\leq n$, is an isomorphism.  Induction on the finite
ideal filtrations shows that $f_{n+1}:C_{n+1}\to V_{n+1}$ is
an isomorphism.  Passing to the inverse limit of these
compatible isomorphisms proves formal linearization.
All constructions used square-zero extensions and ordinary
derivations; no division by an integer is involved.

We finally verify (N).  Serre duality on the smooth proper
curve, with $\omega_C=\Omega_{C/k}$ and $\deg\omega_C=2g-2$,
gives
\[
 \begin{aligned}
 H^1(C,T_C\otimes N^{-n})^\vee
    &\simeq H^0(C,\omega_C^{\otimes2}\otimes N^n),
       &\deg(\omega_C^{\otimes2}\otimes N^n)&=4g-4-nd,\\
 H^1(C,N^{1-n})^\vee
    &\simeq H^0(C,\omega_C\otimes N^{n-1}),
       &\deg(\omega_C\otimes N^{n-1})&=2g-2-(n-1)d.
 \end{aligned}
\]
The exact primary inputs are
\cite[Lemmas 53.4.1 and 53.5.2, tags 0BS2 and 0BS6]
     {StacksBNegativeCurve}.
Under (N), the first degree is negative for every $n\geq1$,
and the second is negative for every $n\geq2$.
A line bundle of negative degree on a proper integral curve
has no nonzero section, proving (L).
\end{proof}

\begin{corollary}[Proved under stated hypotheses: sufficiently negative exceptional curves]
\label{D:negative-curve-tame}
Let $(S,s)$ be a normal surface germ over an algebraically closed
field of characteristic $p>0$.  Suppose a proper resolution of
this germ has exceptional fibre consisting of one smooth proper
curve $C$, and write $g=g(C)$ and $C^2=-d<0$.  If (L) holds for
$N_{C/Y}$, in particular if $d>\max\{0,4g-4\}$, every finite
normal cover of $S$ which is \'etale on $S\setminus\{s\}$ is
tame at the valuation defined by $C$.  This includes separability
of the residue-field extension and has no restriction on the
degree of the cover or on the Cartier index of $K_S$.
\end{corollary}

\begin{proof}
The resolution is smooth over the algebraically closed field
in a neighborhood of $C$.  Lemma~\ref{D:negative-curve-linearization}
identifies its formal neighborhood with the formal zero section
of the negative line bundle $N_{C/Y}$.  Apply
Proposition~\ref{D:graded-cone-tame}.  Its completion comparison
uses the genuine completed boundary DVR and proves the claimed
tameness for the original valuation.  We do not infer a Cartier
property of the canonical divisor and do not use an index-one
cover.
\end{proof}

\begin{remark}[Scope of the numerical condition]
For $g\geq2$ the numerical threshold is $d>4g-4$.
It is a sufficient condition, not a necessary condition for
linearization or tameness.  The cohomological criterion (L)
also covers borderline degrees when its indicated spaces of
sections vanish.  In particular, a failure of this numerical
inequality is not evidence of a wild cover.  Conversely, the
lemma does not assume in advance that the germ is a cone: it
constructs the required formal identification by compatible
infinitesimal lifts.
\end{remark}


\subsubsection{Mixed surface boundaries with formally linear exceptional curves}

\begin{theorem}[Proved under stated hypotheses: mixed finite boundary with formal cones]
\label{D:formal-cone-mixed-boundary}
Let $k$ be algebraically closed of characteristic $p>0$, let
$S$ be a normal connected proper surface over $k$, and let
$Y\subset S$ be finite.  Suppose that every singular point
$s\in Y$ satisfies at least one of the following conditions:
\begin{enumerate}
\item $S$ is $\mathbb Q$-factorial at $s$;
\item $S$ is log canonical at $s$;
\item a resolution of the germ at $s$ has reduced exceptional
      fibre one smooth curve $C$, and its formal neighborhood
      is isomorphic to the formal zero section of its negative
      normal line bundle $N=N_{C/\widetilde S}$.
\end{enumerate}
Then $\pi_1(S\setminus Y)$ is topologically finitely generated.
The singularities retained in the open are unrestricted.

In condition (3), it is enough that the normal bundle satisfies
\[
 H^1(C,T_C\otimes N^{-n})=0\quad(n\geq1),
 \qquad
 H^1(C,N^{1-n})=0\quad(n\geq2).
\]
In particular, if $C^2=-d<0$ and $g=g(C)$, it is enough that
$d>\max\{0,4g-4\}$.  This includes every single smooth
elliptic exceptional curve.

For a possibly nonnormal proper scheme of dimension at most
two, the same finite-generation conclusion holds if $Y$ has
codimension at least two in every irreducible component and
its finite inverse image on each normal surface component
satisfies the displayed alternatives.
\end{theorem}
\begin{proof}
The cohomological and numerical assertions follow from
Lemma~\ref{D:negative-curve-linearization}.  Thus it suffices
to prove the theorem under alternatives (1)--(3).

Let $T\subset Y$ consist of the singular points which are
$\mathbb Q$-factorial.  We leave these points unchanged.  At
each remaining point of $Y$, choose the local resolution in
(3), or the minimal resolution in (2).  In the latter case,
Lemma~\ref{D:lc-exceptional-configurations} gives either a
rational tree, one smooth elliptic curve, or a rational nodal
cycle.  The rational-tree case would be $\mathbb Q$-factorial
by Proposition~\ref{D:rational-tree-Qfactorial}, so it already
belongs to $T$.  Blow up a node when necessary to give the
cycle smooth components.  At singular points of $S\setminus Y$
choose any regular resolution, and keep regular points
unchanged.

These local choices produce a proper birational map
$q:X\to S$.  Explicitly, choose a neighborhood of each of the
finitely many singular points containing no other singular
point and spread the specified local model to a smaller such
neighborhood.  Each chosen resolution is an isomorphism on
the punctured neighborhood, so it glues with the identity on
the regular locus and with the other resolutions.  Properness
is local on the target.  The resulting $X$ is regular away
from the retained $\mathbb Q$-factorial points and is therefore
$\mathbb Q$-factorial everywhere.  It is proper, hence
projective by the intersection argument in
Lemma~\ref{D:Qfactorial-partial-resolution}.  That argument
derives projectivity from properness and $\mathbb Q$-factoriality;
it does not require the prescribed local resolution to have
been chosen projective in advance.

Put $U=S\setminus Y$ and $U'=q^{-1}(U)$.  If $V\to U$ is a
connected finite \'etale cover, its pullback to the normal
integral $U'$ is connected: every finite \'etale component
dominates $U'$, whereas its generic algebra is the single
field $k(V)$.  Thus $\pi_1(U')\to\pi_1(U)$ is surjective.
The finite normalization of $S$ in $k(V)$ is \'etale away from
$Y$, so at any omitted point its restriction to the punctured
germ is \'etale.

Every divisorial component of $X\setminus U'$ lies over one
of the points of $Y\setminus T$ which we resolved.  At a
formally linear curve, Proposition~\ref{D:graded-cone-tame}
proves tameness at its valuation, including separability of
the residue extension.  At an lc elliptic or cyclic core,
Corollary~\ref{D:lc-core-valuative-tame} proves tameness, also at
the node subdivisions.  Retained points of $T$, and regular
omitted points, have codimension two in $X$ and create no
divisorial boundary condition.  Exceptional curves over
singularities retained in $U$ are contained in $U'$.
Consequently every finite quotient of $\pi_1(U)$ factors
through the fixed-boundary tame quotient of $\pi_1(U')$, and
we have a surjection
\[
                  Q(X,X\setminus U')\longrightarrow\pi_1(U).
\]
Its source is topologically finitely generated by
Proposition~\ref{D:fixed-boundary-positive}, because $X$ is
projective.  This proves the normal case.

For the final assertion, pass to the finite disjoint union
of the normalizations of the reduced irreducible components.
The proved result applies to every normal surface component.
Components of dimension at most one have no omitted points
by the codimension hypothesis, so they are proper and their
fundamental groups are finitely generated.  The finite proper
descent argument in Proposition~\ref{D:fixed-boundary-positive},
using \cite[IX, Corollary 5.2]{SGA1}, gives finite generation
downstairs.
\end{proof}

\begin{corollary}[Proved under stated hypotheses: fixed boundary with isolated formal cones]
\label{D:formal-cone-fixed-boundary}
Let $S$ be a normal connected proper surface over an algebraically
closed field of characteristic $p>0$, and let $Y\subsetneq S$
be any closed subset.  Assume that every singular point
$s\in Y$ which is not $\mathbb Q$-factorial is isolated in
$Y$ and is either log canonical or satisfies condition (3)
of Theorem~\ref{D:formal-cone-mixed-boundary}.  Then
\[
                         Q(S,Y)
\]
is topologically finitely generated.  The cohomological and
numerical sufficient conditions in that theorem apply here
as well.  Singularities outside $Y$ are unrestricted.
\end{corollary}
\begin{proof}
As in the preceding proof, retain every $\mathbb Q$-factorial
singular point on $Y$, resolve the other singular points with
the indicated models, and resolve all singular points outside
$Y$ arbitrarily.  The result is a normal proper
$\mathbb Q$-factorial, hence projective, surface $q:X\to S$.
Put $U=S\setminus Y$ and $U'=q^{-1}(U)$.  The birational
pullback again gives a surjection $\pi_1(U')\to\pi_1(U)$.

Consider a connected finite \'etale cover of $U$ which is
tame at the original codimension-one boundary valuations.
Their strict transforms on $X$ define the same DVRs, so the
pulled-back cover remains tame there.  Every new exceptional
boundary curve lies over a non-$\mathbb Q$-factorial singular
point $s\in Y$.  By hypothesis $s$ is isolated in $Y$; on a
sufficiently small neighborhood the cover is therefore
\'etale on the punctured germ.  Its ramification at the new
exceptional valuations is tame by
Proposition~\ref{D:graded-cone-tame} or
Corollary~\ref{D:lc-core-valuative-tame}, exactly as in the preceding
proof.  Resolutions above points outside $Y$ remain inside
$U'$.  Thus pullback defines a surjection
\[
                         Q(X,X\setminus U')\longrightarrow Q(S,Y).
\]
The projective fixed-boundary theorem makes the source
topologically finitely generated.  This proves the assertion
without dropping the isolation condition on the newly
resolved boundary points.
\end{proof}


% Post-checkpoint 07 extension.  This fragment is separate pending independent audit.
\begin{proposition}[Proved: a Q-factorial extraction for a formally graded surface germ]
\label{D:weighted-graded-extraction}
Let $k$ be algebraically closed of characteristic $p>0$, and let
$(A,\mathfrak m)$ be an excellent normal local $k$-algebra of
dimension two with residue field $k$.  Suppose that its completion
is isomorphic, as a complete local $k$-algebra, to the completion
at $R_+$ of a finitely generated normal positively graded domain
\[
 R=\bigoplus_{d\geq0}R_d,\qquad R_0=k,\qquad \dim R=2.
\]
There is a projective birational morphism
$h:Y\to\operatorname{Spec}A$, an isomorphism on the punctured
spectrum, with the following properties.  The surface $Y$ is
normal and locally Q-factorial, its exceptional locus has exactly
one irreducible curve, and every finite cover of the punctured
spectrum is tame at the valuation of that curve.  The exceptional
valuation, after completion at $\mathfrak m$, is the radial
valuation of the grading.
\end{proposition}
\begin{proof}
We divide all occurring degrees by their greatest common divisor,
so the grading is effective.  Choose $N>0$ such that the Veronese
ring $R^{(N)}$ is generated in degree one, using
\cite[Lemma 10.56.2, tag 0EGH]{C_StacksWeightedCone}.
Put $S=R^{(N)}$ and $C=\operatorname{Proj}S=\operatorname{Proj}R$.
The ring $S$ is normal, $R$ is finite over $S$, and $C$ is a smooth
projective integral curve.  Here normality of the Veronese follows
by intersecting $R$ with its degree-multiple fraction subfield;
normality of $C$ follows from its degree-zero affine charts.

We first construct the model over $R$.  The blow-up
\[
 W=\operatorname{Bl}_{S_+}(\operatorname{Spec}S)
\]
is the total space of $\mathcal O_C(-1)$.  Indeed, for a nonzero
$f\in S_1$ its blow-up chart has ring
\[
 S[S_1/f]=(S[f^{-1}])_0[f].
\]
The element $f$ is algebraically independent over the degree-zero
fraction field, and the changes between these charts are the
linear changes of a line-bundle coordinate.  Write
$F=k(C)$ and $K=\operatorname{Frac}R$.  Effectivity gives a rational
homogeneous element $t$ of degree one, by taking a product of
homogeneous elements with integer exponents.  Thus
\[
 K=F(t),\qquad \operatorname{Frac}S=F(t^N).
\]
Let $Y_R$ be the normalization of $W$ in $K$.  It is finite over
$W$.  The elements of $R$ are integral over $S$, hence regular
on $Y_R$, so this construction gives a projective birational
morphism $Y_R\to\operatorname{Spec}R$, an isomorphism away from
the vertex.  Projectivity follows also by factoring the projective
morphism $Y_R\to\operatorname{Spec}S$ through the separated finite
morphism $\operatorname{Spec}R\to\operatorname{Spec}S$.

At the generic point of the zero section of $W$, a line coordinate
$u$ satisfies $u=f t^N$ for some $f\in F^*$.  The extension
of the $u$-adic valuation to $F(t)$ is totally ramified of degree
$N$, with residue field $F$, and is unique.  This follows directly
from the Eisenstein polynomial $T^N-u/f$ over $F[[u]]$, including
when $p$ divides $N$.  Consequently $Y_R$ has exactly one
exceptional curve.  Its valuation is the degree valuation of $R$.

We verify local Q-factoriality, retaining the characteristic-$p$
possibilities in this construction.  At $c\in C$, choose a
uniformizer $z$ and write $f=z^b\epsilon$ with $\epsilon$ a unit.
Absorb $\epsilon$ into $u$ and multiply $t$ by a suitable power of
$z$.  The local finite normalization is then given by normalizing
\begin{equation}\label{D:weighted-toric-chart}
 T^N=u z^a,\qquad 0\leq a<N.
\end{equation}
We may take an affine neighborhood of $c$ on which $z$ defines an
\'etale morphism to $\mathbb A^1_k$.  Normalization commutes with
this \'etale base change
\cite[Lemma 37.19.2, tag 03GV]{C_StacksWeightedCone}.
The normalization of the displayed affine equation over
$k[u,z]$ has a finite surjection from $\mathbb A^2_k$, given by
\[
 u=r^N,\qquad z=s^N,\qquad T=r s^a.
\]
Indeed, $r$ and $s$ are integral, the polynomial ring $k[r,s]$
is normal, and its field extension over $k(u,z,T)$ has degree
$N$.  To compute the degree, first
$[k(r,s):k(r^N,s^N)]=N^2$, whereas adjoining $r s^a$ to
$k(r^N,s^N)$ has degree $N$, as its $r$-valuation has value one.
The resulting finite map factors through the normalization.
After the \'etale base change in $z$, its source remains regular.

For completeness, a normal local domain admitting a finite
surjection from a regular semilocal scheme is Q-factorial.  Pull
a Weil divisor $D$ back using the ramification indices of the
height-one valuations.  On the regular semilocal source this
divisor is principal.  Taking the product of the field norms on
its components and pushing the divisor down makes $mD$ principal,
where $m$ is the sum of their generic degrees.  The identities
used here are
\[
 f_*f^*D=mD,\qquad
 f_*\operatorname{div}(g)=\operatorname{div}(\operatorname{Nm}(g));
\]
the second is the norm formula for principal divisors
\cite[Lemma 42.18.1, tag 02RS]{C_StacksWeightedCone}, and the first
is the degree formula for a finite extension of DVRs.  Regular
semilocal schemes have trivial Picard groups and are locally
factorial \cite[tags 02M9 and 0AG0]{C_StacksWeightedCone}.
This argument does not require a separable field
extension.  Applied to our local charts it shows that $N$ kills
each local divisor class along the exceptional curve.  The same
argument holds after completing at any of these points: the
finite cover completes to a product of regular complete local
rings.  Thus these completed local rings are Q-factorial as well.

The model just constructed is a normalized blow-up of a
vertex-primary ideal.  More explicitly, put
\[
 J=R_N R=S_+R.
\]
The equality uses that $S$ is generated in degree one, and $J$
is $R_+$-primary because $R$ is integral and graded over $S$.
The blow-up of $J$ is the dominant component of the base change
of $W$ to $R$: the Rees algebra map
\[
 \left(\bigoplus_{n\geq0}S_+^n\right)\otimes_S R
 \longrightarrow \bigoplus_{n\geq0}J^n
\]
is surjective and is an isomorphism on the dense complement of
the vertex.  In particular, this blow-up is finite over $W$.
Its normalization is therefore $Y_R$, by uniqueness of the
normalization of $W$ in $K$.

Identify $\widehat A$ with the completed graded local ring.
The ideal $J\widehat A$ contains a power of the maximal ideal,
so it is the extension of an ideal $I\subset A$: use
$A/\mathfrak m^q=\widehat A/\mathfrak m^q\widehat A$ for a
sufficiently large $q$.  Let $Y$ be the normalization of
$\operatorname{Bl}_I(\operatorname{Spec}A)$.  Its base change to
$\widehat A$ is the normalization of the corresponding completed
blow-up.  This follows directly from
\cite[Lemmas 54.11.3 and 54.11.6, tag 0BG4]{C_StacksWeightedCone}:
away from the special fibre the blow-up is the normal punctured
spectrum.  In particular, the normalization is finite and $Y$
is projective over $A$.

There is exactly one exceptional curve, since the special fibres
before and after completion are identical.  To descend
Q-factoriality, one can work with the completed local rings at
points of this curve.  These rings are unchanged by base change
to $\widehat A$
\cite[Lemma 54.11.1, tag 0BG4]{C_StacksWeightedCone} and, by the
finite regular charts above, are Q-factorial.  If a reflexive
rank-one module $M$ over a normal local ring has an invertible
reflexive power after completion, that power is already
invertible.  Indeed, finite-module duals and double duals commute
with flat base change
\cite[Lemma 17.22.5, tag 01CM]{C_StacksWeightedCone},
and finite locally free modules descend
faithfully flatly.  This proves local Q-factoriality on $Y$;
outside the exceptional fibre its local rings are regular.
The normality needed after the intermediate completion base
changes follows from regularity of completion maps of excellent
rings and normality under regular base change
\cite[tags 0BFK and 0H7S]{C_StacksWeightedCone}.

Finally, pull any finite punctured cover back to $\widehat A$ and
take the connected components and their Galois closures.
Proposition~\ref{D:graded-cone-tame} makes them tame at the
radial valuation.  The original exceptional DVR and the radial
DVR on the completed model have the same actual completion, by
Lemma 54.11.1 applied at the generic point of the special-fibre
curve.  Completion preserves the ramification index and the
separability of the residual extension for finite extensions of
excellent DVRs.  Tameness therefore descends to the original
exceptional valuation, as required.
\end{proof}

\begin{corollary}[Proved: formally graded points in a proper surface boundary]
\label{D:weighted-graded-boundary}
In Theorem~\ref{D:formal-cone-mixed-boundary} and
Corollary~\ref{D:formal-cone-fixed-boundary}, we may add the
following alternative for every omitted non-Q-factorial point:
its completed local ring is a positively graded normal surface
completion as in Proposition~\ref{D:weighted-graded-extraction}.
Thus the punctured fundamental group is topologically finitely
generated in the finite-boundary case, and the fixed-boundary
tame quotient is topologically finitely generated in the general
boundary case with these non-Q-factorial points isolated in the
boundary.  The finite-normalization version of the first theorem
has the same additional alternative on the normalizations.
\end{corollary}
\begin{proof}
At each formally graded omitted point use the preceding
projective partial resolution.  It spreads to an affine
neighborhood by finite presentation.  The ideal may be extended
to a coherent ideal equal to the unit ideal off that point, so
the normalized blow-up is an isomorphism away from the prescribed
point.  Keep the omitted Q-factorial points and resolve the
retained singularities as in the cited mixed-boundary arguments.
The resulting proper normal surface is locally Q-factorial and
hence projective by Lemma~\ref{D:Qfactorial-partial-resolution}.
At a formally graded point the only new boundary divisor is its
radial curve, and all the covers under consideration are tame
there by Proposition~\ref{D:weighted-graded-extraction}.  No
divisors arise from the Q-factorial singularities retained on
the partial resolution.  The remaining boundary divisors are
exactly those already handled in the mixed-boundary theorem.
The projective fixed-boundary tame group consequently surjects
onto the required group.  Its finite generation gives the claim.
The finite-normalization descent is unchanged.
\end{proof}

\begin{corollary}[Proved under stated hypotheses: the three-point surface in every genus]
\label{D:schroer-all-genera}
In Schr\"oer's construction of a normal proper surface $S$ with
$\operatorname{Pic}(S)=0$, let the initial smooth proper curve
$C$ have any positive genus, and let $s_0,s_1,s_2$ be the three
contracted points.  Then
\[
                  \pi_1(S\setminus\{s_0,s_1,s_2\})
\]
is topologically finitely generated.  In particular, the
higher-genus version of this explicit nonprojective surface is
not a counterexample to D.2(ii).
\end{corollary}
\begin{proof}
Characteristic zero is already covered by
Proposition~\ref{D:fixed-boundary-positive}; assume that the
characteristic is positive.  Use the notation and point choices
of \cite[\S\S2.5 and 3.1--3.2, author's PDF pp.~4--5]{Schroer1999D}.
Let $X$ be the blowup of $\mathbb P^1\times C$ at
$(b_0,c_1),(b_1,c_1),(b_2,c_2)$, let $R_i$ be the strict
vertical fibres, and let $A$ be the strict horizontal curve
through $c_1$.  The common open is
\[
 U=X\setminus(R_0\cup R_1\cup A\cup R_2)
   =S\setminus\{s_0,s_1,s_2\}.
\]

We identify the formal neighborhood of $R_i$ for $i=0,1$.
Choose a coordinate $t$ on an affine neighborhood of $b_i$
in $\mathbb P^1$, with $t(b_i)=0$, and a parameter $z$ at
$c_1$ on $C$.  The blowup chart containing the strict transform
of $\{b_i\}\times C$ has coordinates
\[
                t=zu,\qquad R_i=\{u=0\}.
\]
Away from $c_1$ the normal coordinate is $t$.  On the overlap
the transition is $u=t/z$.  These are exactly the linear
transition functions of the line bundle
$\mathcal O_C(-c_1)$.  They identify the completed neighborhood
of $R_i$ with the completion of its zero section.  The other
two blowup centres are disjoint from $R_i$ and do not change
this formal neighborhood.  The line bundle has degree $-1$.
Proposition~\ref{D:graded-cone-tame} therefore shows that any
finite cover of the punctured germ at $s_i$ is tame at the
divisorial valuation $R_i$, including separability of its
residue extension.

Let $V\to U$ be a connected finite \'etale cover.  Its finite
normalization over $S$ is \'etale away from the three omitted
points.  Applying the preceding local calculation at $s_0$
and $s_1$ proves tameness at $R_0$ and $R_1$ for this cover.
There is no restriction on its degree, and no restriction on
the genus of $C$ in this calculation.

Section 3.1 of the cited source constructs a normal projective
surface $T$ by contracting only $A\cup R_2$.  The construction
works for the same arbitrary positive-genus curve; it is the
partial contraction used in Corollary~\ref{D:schroer-elliptic}.
If $q$ is the image of $A\cup R_2$, then
\[
                   U=T\setminus(R_0\cup R_1\cup\{q\}).
\]
The only divisorial boundary components on $T$ are $R_0$ and
$R_1$.  We have proved that every finite \'etale cover of $U$
is tame at both.  The omitted point $q$ has codimension two,
so it adds no condition to this fixed-boundary tame category.
Consequently
\[
                   \pi_1(U)=Q(T,T\setminus U).
\]
The right-hand side is topologically finitely generated by
Proposition~\ref{D:fixed-boundary-positive}, because $T$ is
projective.  We have imposed no log-canonical or formal-cone
condition on the third singularity $q$.
\end{proof}


\subsubsection{The higher-genus three-point example: horizontal wild inertia}

\begin{proposition}[Proved under stated hypotheses: one horizontal tame condition]
\label{D:Schroer-horizontal-tame}
Let $k$ be algebraically closed, let $C/k$ be a smooth proper
connected curve, and choose distinct points $b_0,b_1,b_2\in\mathbb P^1(k)$
and distinct points $c_1,c_2\in C(k)$.  Let
\[
 X=\operatorname{Bl}_{(b_0,c_1),(b_1,c_1),(b_2,c_2)}
       (\mathbb P^1\times C).
\]
Write $R_i$ for the strict transform of $\{b_i\}\times C$, $E_i$
for the corresponding exceptional curve, and $A$ for the strict
transform of $\mathbb P^1\times\{c_1\}$.  Put
\[
 U=X\setminus(A\cup R_0\cup R_1\cup R_2).
\]
Every finite \'etale cover of $U$ which is tame at the divisorial
valuation defined by $A$ is tame at each of the valuations defined
by $R_0,R_1,R_2$.

Consequently, if $P_A$ denotes a wild inertia group at $A$, the
natural quotient map induces an isomorphism
\[
 \pi_1(U)/\overline{\langle\!\langle P_A\rangle\!\rangle}
 \simeq Q(X,X\setminus U).
\]
In particular this quotient is topologically finitely generated.
\end{proposition}

\begin{proof}
We may pass to connected components and then to a finite Galois
closure.  Tameness is preserved under passage to this closure,
composita and subextensions; see
\cite[section 15.116, Lemmas 15.116.5--9]{StacksBIndexTame}.
Thus let $V\to U$ be a connected finite Galois cover with group $G$.
Set
\[
 B=\mathbb P^1\setminus\{b_0,b_1,b_2\},
 \qquad C^\circ=C\setminus\{c_1\}.
\]
The open $U_0=B\times C^\circ$ is dense in $U$.  Since $U$ is
normal and connected, $V$ and its restriction to $U_0$ are integral.

Equip $C$ with the divisorial log structure at $c_1$, and equip
$B$ with the trivial log structure.  The resulting log curve
$C^{\log}$ is proper, log smooth and log geometrically connected
over $k$, while $B$ is connected and log regular.  The primary
logarithmic product theorem gives
\[
 \pi_1(C^{\log}\times B)
   \simeq \pi_1(C^{\log})\times\pi_1(B)
   =\pi_1^t(C^\circ)\times\pi_1(B).
 \tag{*}
\]
Here we use \cite[Proposition 2.4, printed pp.~16--17]{Hoshi2006B}.
The precise identification of covers is
\cite[Proposition A.10, printed pp.~57--58]{Hoshi2006B}: for a log
regular scheme, finite Kummer \'etale covers are equivalent to
finite \'etale covers of its interior whose normalizations are tame
at every codimension-one boundary point.  Thus the sole tameness
condition in $(*)$ is along $B\times\{c_1\}$; ramification in the
$B$ direction is unrestricted.  The hypothesis at $A$ is exactly
this condition on $V|_{U_0}$, since these generic boundary valuations
agree.  Its monodromy therefore has the form
\[
 \rho(\gamma_B,\gamma_C)
     =\rho_B(\gamma_B)\rho_C(\gamma_C),
\]
where the two finite images commute, $\rho_B$ is a representation
of $\pi_1(B)$ and $\rho_C$ a representation of $\pi_1^t(C^\circ)$.
The required choices of geometric base points and paths only
conjugate these representations and do not affect vanishing of
wild inertia.

For $i=0,1$ choose a parameter $t$ at $b_i$ and an \'etale local
parameter $z$ at $c_1$.  In a blowup chart around $E_i$ we have
\[
 z=tu,\qquad E_i=\{t=0\},\qquad A=\{u=0\}.
\]
The point where $R_i$ meets $E_i$ is the point omitted from this
chart.  Choose $\lambda\in k^\times$.  The formal arc
\[
 \operatorname{Spec}k[[s]]\longrightarrow X,
 \qquad t=s,\quad u=\lambda,\quad z=\lambda s
\]
has closed point in $E_i\cap U$ and generic point in $U_0$.
The \'etale parameter on $C$ supplies the indicated lift uniquely
after fixing its closed point.  In particular the entire arc lies
in $U$.  Its pullback of $V$ is finite \'etale over the strictly
henselian ring $k[[s]]$, so its monodromy on $k((s))$ is trivial.

The $C$ component of this monodromy is tame at $s=0$: a cover
classified by $\rho_C$ is tame at $c_1$, and the parameter changes
by $z=\lambda s$.  On the wild inertia of $k((s))$, its contribution
is therefore trivial.  The product formula then forces the
$B$ component to be trivial on that same wild inertia.  Since
$t=s$, this is precisely the assertion that $\rho_B$ is tame at
$b_i$.  We have proved this for both $i=0$ and $i=1$.

At $b_2$ use the strict transform of the horizontal curve
$\mathbb P^1\times\{c_2\}$.  Its intersection with $E_2$ lies in
$U$, because $c_2\ne c_1$ and its tangent direction differs from
that of $R_2$.  An arc on this curve with parameter $t=s$ has
constant $C$ coordinate $c_2$.  Its pullback of $V$ again extends
as a finite \'etale cover of $k[[s]]$.  The $C$ component is now
trivial on the entire local fundamental group.  It follows that
$\rho_B$ is unramified at $b_2$.

Let $V_B\to B$ and $V_C\to C^\circ$ denote finite Galois covers
cut out by the kernels of $\rho_B$ and $\rho_C$, respectively.
The product cover of $U_0$ dominates $V|_{U_0}$.  At the generic
point of $R_i$, the pullback of $V_C$ is \'etale, since that point
maps into $C^\circ$.  The pullback of $V_B$ is tame, by the
conclusions at $b_i$ and tame base-change stability for discrete
valuation rings
\cite[section 15.116, Lemma 15.116.9]{StacksBIndexTame}.
Their compositum, and hence the subextension defining $V$, is
tame at $R_i$.  This proves the assertion on covers.

Finite covers of the left-hand quotient in the statement are
exactly the finite \'etale covers of $U$ tame at $A$; those of the
right-hand quotient are the covers tame at every one of the four
boundary divisors.  We have proved equality of these full
subcategories, so the natural quotient is an isomorphism.
Finally $X$ is smooth and projective, and
Proposition~\ref{D:fixed-boundary-positive} makes the right-hand
group topologically finitely generated.
\end{proof}

\begin{remark}[Scope of the horizontal reduction]
In the higher-genus version of Schr\"oer's three-point construction
\cite[sections 2.5 and 3.1--3.2, author's PDF pp.~4--5]{Schroer1999D},
the same $U$ is the proper surface with trivial Picard group minus
its three contracted points.  The proposition shows that any
finite cover which is wild at $R_0$ or $R_1$ must also be wild at
$A$.  Together with the positive-genus residue criterion, this
isolates the remaining possible covers as ones with simultaneous
horizontal wild ramification and inseparable residue extension
at a vertical exceptional divisor.  By itself this proposition does not
assert that $P_A$ has finite image in $\pi_1(U)$ or prove finite
generation of the full group.  The full group in every positive
genus is now settled by Corollary~\ref{D:schroer-all-genera},
using the separate formal-cone argument.  It uses a logarithmic product theorem with one proper log
factor; no ordinary product formula for two nonproper curves is
being assumed.
\end{remark}


\begin{corollary}[Proved: the elliptic three-point surface with trivial Picard group]
\label{D:schroer-elliptic}
In Schr\"oer's construction of a normal proper surface $S$ with
$\operatorname{Pic}(S)=0$, take the initial curve $C$ to be elliptic
over an algebraically closed field.  Let $s_0,s_1,s_2$ be its three
contracted points.  Then
$\pi_1(S\setminus\{s_0,s_1,s_2\})$ is topologically finitely
generated, in every characteristic.
\end{corollary}
\begin{proof}
We use exactly the construction in
\cite[sections 2.5 and 3.1--3.2, author's PDF pp.~4--5]{Schroer1999D}.
Choose $c_1,c_2\in C$ as in that source, blow up
$\mathbb P^1\times C$ at $(b_0,c_1),(b_1,c_1),(b_2,c_2)$,
and call the resulting smooth projective surface $X$.
Write $R_i$ for the strict vertical fibres, $A$ for the strict
horizontal curve through $c_1$, and $B$ for that through $c_2$.
The disjoint connected configurations $R_0$, $R_1$ and $A\cup R_2$
contract to $s_0,s_1,s_2$.  The common open in question is
\[
 U=X\setminus(R_0\cup R_1\cup A\cup R_2).
\]

We check the local hypothesis at $s_0,s_1$ explicitly.  For
$i\in\{0,1\}$ choose a coordinate $t$ on $\mathbb P^1$ vanishing
at $b_i$, and a nowhere-vanishing regular differential $\eta$ on
the elliptic curve $C$.  The rational form
$(dt/t)\wedge\eta$ has a simple pole on the fibre through $b_i$.
Blowing up $(b_i,c_1)$ adds discrepancy $1$ and pulls back that
pole with multiplicity $1$, so the exceptional rational curve has
coefficient zero.  On a neighborhood of $R_i$ the divisor of the
pulled-back form is exactly $-R_i$.  Its other divisor components
have closed images avoiding $s_i$, so we can shrink around $s_i$.
The descended form has divisor zero on the punctured neighborhood;
normality identifies its reflexive canonical sheaf with
$\mathcal O_S$ there and across $s_i$.  Thus $K_S$ is Cartier near
$s_i$, and the displayed resolution has the single discrepancy
$-1$ on the smooth elliptic curve $R_i$.  It is log canonical:
the boundary consists of that one smooth curve, and subsequent
point blowups preserve discrepancies at least $-1$.

Section 3.1 of the same source constructs a \emph{projective}
normal surface $T$ by contracting only $A\cup R_2$.  Explicitly,
the effective divisor $D=3B+2R_0+2R_1$ satisfies
$D\cdot B=D\cdot R_0=D\cdot R_1=1$, and the associated contraction
is the one asserted there.  Hence
$U=T\setminus(R_0\cup R_1\cup\{q\})$, where $q$ is the image of
$A\cup R_2$.

Every connected finite \'etale cover of $U$ extends by normalization
to a finite cover of $S$ that is \'etale away from its three
omitted points, hence quasi-\'etale.  Proposition~\ref{D:lc-elliptic-tame}
shows that this cover is tame at the two divisorial boundary
components $R_0,R_1$ on $T$.  The remaining boundary point $q$
has codimension two and imposes no condition in $Q(T,T\setminus U)$.
Thus every finite \'etale cover of $U$ already belongs to this
fixed-boundary tame category, and
\[
 \pi_1(U)=Q(T,T\setminus U).
\]
The right-hand side is finitely generated by
Proposition~\ref{D:fixed-boundary-positive}, since $T$ is projective.
We do not require the singularity at $q$ to be log canonical.
Indeed, in the elliptic case the discrepancies on the configuration
$A\cup R_2$ are $-1$ and $-2$, respectively; leaving it as a
codimension-two boundary point is essential to this proof.
\end{proof}


\begin{corollary}[Proved under stated hypotheses: elliptic omitted singularities]
\label{D:elliptic-surface-boundary}
Let $S$ be a normal proper connected surface over an algebraically
closed field, and let $Y\subset S$ be finite.  Assume that every
singular point in $Y$ admits a resolution whose reduced exceptional
locus is one smooth elliptic curve.  Then
$\pi_1(S\setminus Y)$ is topologically finitely generated.
The singularities retained in $S\setminus Y$ are unrestricted.
For a possibly nonnormal proper scheme of dimension at most two,
the same conclusion holds if the finite inverse image of $Y$ on
each normal surface component satisfies this condition and $Y$
has codimension at least two in every component.
\end{corollary}
\begin{proof}
Choose a resolution $q:X\to S$ which agrees with the specified
single-curve resolutions around the singular points of $Y$, is an
isomorphism at its regular points, and resolves the remaining
singularities arbitrarily.  These choices glue: take affine
neighborhoods of each of the finitely many singular points
containing no other singular point; their resolutions are all
isomorphisms on the overlaps with the regular locus.
The resulting $X$ is smooth and proper over $k$, hence projective
by \cite[Lemma 54.16.11, tag 0C5P]{StacksDSurfaceResolution}.
Put $U=S\setminus Y$ and $U'=q^{-1}(U)$.

A connected finite \'etale cover of the normal $U$ stays connected
after the proper birational pullback to $U'$.  Its normalization
over $S$ is finite and \'etale away from $Y$.  At every divisorial
component of $X\setminus U'$ the local situation is the elliptic
case of Proposition~\ref{D:single-curve-residue-tameness}; all these
extensions are therefore tame.  The regular omitted points remain
codimension two and give no divisorial boundary conditions.
Consequently the surjection $\pi_1(U')\to\pi_1(U)$ factors
through a surjection
\[
              Q(X,X\setminus U')\longrightarrow\pi_1(U).
\]
The source is finitely generated by
Proposition~\ref{D:fixed-boundary-positive}, proving the claim.

In the nonnormal case we pass to the finite disjoint union of
normalizations of the reduced irreducible components.  The proved
surface case applies to each corresponding open; any components
of dimension at most one have no omitted point and are proper.
Their fundamental groups are finitely generated.  The finite
proper descent argument of Proposition~\ref{D:fixed-boundary-positive}
then gives finite generation downstairs.
\end{proof}


\begin{proposition}[Proved under stated hypotheses: finite exceptional wild images]
\label{D:finite-exceptional-wild}
Let $X$ be a normal proper integral variety over an algebraically
closed field, let $Y\subsetneq X$ be closed, and put
$U=X\setminus Y$ and $Q=Q(X,Y)$.
Suppose that some normal projective modification $q:X'\to X$
has the following property: for every prime divisor $E$ in
$X'\setminus q^{-1}(U)$, the image $I_E$ of its absolute wild
inertia group in $Q$ is finite.  Then $Q$ is topologically
finitely generated.
It suffices to check the exceptional divisors among these $E$,
because the wild images at the strict transforms of the original
codimension-one boundary are already trivial in $Q$.
\end{proposition}
\begin{proof}
There are finitely many such divisors.  Choose an open normal
subgroup $H\subset Q$ with $H\cap I_E=1$ for all of them.
This is possible by excluding each of the finitely many nonidentity
elements in their finite union.  Normality of $H$ makes the same
condition hold for every conjugate of an $I_E$.
Let $V\to U$ be the connected finite Galois cover classified by
$Q/H$, and let $L=k(V)$.  Normalize $X'$ in $L$, obtaining a
normal projective variety $T$, finite over $X'$.  Over
$U'=q^{-1}(U)$ this normalization is
$V'=V\times_U U'$, which is connected and normal: its generic
fibre is $\operatorname{Spec}L$, and it is finite \'etale over
the normal integral $U'$.
The morphism $V'\to V$ is proper birational.  Consequently
$\pi_1(V')\to\pi_1(V)$ is surjective, since a connected finite
\'etale cover of the normal integral $V$ stays connected after
this birational base change.

We recall the valuation-group comparison needed here.  If $w/v$
is an extension of discrete valuations through a finite separable
field extension $L/K$, choose compatible geometric valuations.
Inside the absolute decomposition group of $v$ we have
\[
 I_w=I_v\cap G_L,\qquad P_w=P_v\cap G_L,
\]
where $I$ and $P$ denote inertia and wild inertia.  The first
identity follows by restricting the action on residue algebraic
closures.  For the second, $P_v\cap G_L$ is a normal pro-$p$
subgroup of $I_w$, and $I_w/(P_v\cap G_L)$ is a subgroup of the
prime-to-$p$ tame inertia quotient $I_v/P_v$; hence it is exactly
the wild inertia subgroup.  This description includes imperfect
residue fields; see the exact sequence in the Stacks Project,
Lemma~58.13.5, tag 0BUA \cite{StacksDTame}.  In characteristic zero all wild groups
are trivial.

Every prime divisor $F$ of $T\setminus V'$ lies above a prime
divisor $E$ of $X'\setminus U'$, because $T\to X'$ is finite.
Under the natural map from $G_L$ to $H$, the preceding identity
puts the image of $P_F$ in $H\cap gI_Eg^{-1}=1$ for a suitable
conjugate.  Thus every finite cover of $V$ classified by a finite
$H$-set, after pullback to $V'$, is tame at every divisorial
boundary component of $T$.

The surjection $\pi_1(V')\to\pi_1(V)\to H$ therefore kills the
closed normal subgroup generated by those wild inertia groups.
It factors as a surjection
\[
 Q(T,T\setminus V')\longrightarrow H.
\]
The source is finitely generated by the projective case of
Proposition~\ref{D:fixed-boundary-positive}.  Hence $H$ is
finitely generated, and so is its finite extension $Q$.

Finally, a proper birational morphism between normal varieties is
an isomorphism at the generic point of every original prime divisor,
so the strict transform has the same discrete valuation.
Its wild inertia image is zero by the definition of $Q(X,Y)$.
This proves the final assertion.
\end{proof}

\begin{remark}
This criterion does not assume that $q$ is an isomorphism over $U$,
nor that the punctured strict-local fundamental groups are finite.
The proof uses only the wild images of the finitely many boundary
valuations on the chosen projective modification.
In any remaining counterexample to finite generation, every such
projective modification must therefore introduce a boundary divisor
with infinite wild image in $Q$.
The criterion by itself is a sufficient-condition refinement, not
an additional separated geometric example.  The elliptic example
of Corollary~\ref{D:schroer-elliptic} supplies the concrete case
where the new images are in fact trivial.
\end{remark}


\begin{corollary}[Proved under stated hypotheses: mixed mild surface boundary]
\label{D:Gorenstein-lc-boundary}
Let $S$ be a normal proper connected surface over an algebraically
closed field, and let $Y\subset S$ be finite.  Suppose that every
point of $Y$ is regular or belongs to at least one of the following
classes:
\begin{enumerate}
\item a log terminal surface singularity;
\item a singularity admitting a resolution whose reduced exceptional
fibre is one smooth elliptic curve, with no $\mathbb Q$-Gorenstein
assumption;
\item a Gorenstein log canonical cusp singularity, whose minimal
exceptional fibre is a rational nodal cycle.
\end{enumerate}
Then $\pi_1(S\setminus Y)$ is topologically finitely generated.
In particular, this holds if every omitted singularity is either
log terminal or Gorenstein log canonical.  The singularities
retained in $S\setminus Y$ are unrestricted.
For a possibly nonnormal proper scheme of dimension at most two,
the same conclusion holds if the finite inverse image of $Y$ on
each normal surface component satisfies these conditions and $Y$
has codimension at least two in every irreducible component.
\end{corollary}
\begin{proof}
Choose a regular resolution $q:X\to S$ which is minimal near the
omitted log terminal and cusp points, is the specified single-curve
resolution near the omitted elliptic points, and is an isomorphism
at the omitted regular points.  At an irreducible nodal cusp core,
blow up its node once.  Resolve the singularities retained in the
open arbitrarily.  As in Corollary~\ref{D:elliptic-surface-boundary},
these finitely many local choices glue along the regular locus.
The regular proper surface $X$ is projective by
\cite[Lemma 54.16.11, tag 0C5P]{StacksDSurfaceResolution}.
Put $U=S\setminus Y$ and $U'=q^{-1}(U)$.

Every connected finite \'etale cover of $U$ stays connected after
pullback to $U'$: its generic fibre is a field and the pullback is
normal, with every connected component dominating the integral
base $U'$.  Its normalization over $S$ is finite and quasi-\'etale.
At every divisorial boundary component over an elliptic point,
Proposition~\ref{D:single-curve-residue-tameness} proves tameness.
At a cusp point, Theorem~\ref{D:cusp-tame} proves tameness, including
at the components introduced by the node blowup.  Thus the image
of wild inertia in $\pi_1(U)$ is zero at all these components.

At an omitted log terminal point, the punctured strict-local
fundamental group is finite by
\cite[Theorem 5.3, p.~739; see also Corollary 6.4, p.~741]
{CarvajalRojasYasuda2025B}.  The theorem holds in every
characteristic, including two and three.  The wild inertia image
at an exceptional divisor over this point factors through the
image of that local group in $\pi_1(U)$, and is therefore finite.
Regular omitted points give only codimension-two boundary on
$X$ and require no divisorial condition.
Consequently all boundary wild images on this fixed projective
modification are finite.  There is no original divisorial
boundary, so $Q(S,Y)=\pi_1(U)$, and
Proposition~\ref{D:finite-exceptional-wild} proves the conclusion.

Lemma~\ref{D:Gorenstein-lc-configurations} proves the asserted
Gorenstein log canonical special case: its all-zero configuration
is canonical, hence log terminal, and its other configurations
are the elliptic and cusp cases just treated.
In the nonnormal case we pass to the finite disjoint union of
normalizations of the reduced irreducible components.  The proved
normal case applies to each surface component, while components
of dimension at most one have no omitted point and are proper.
Finite proper descent, as in
Proposition~\ref{D:fixed-boundary-positive}, gives the conclusion
downstairs.
\end{proof}


\begin{proposition}[Proved under stated hypotheses: rational exceptional trees]
\label{D:rational-tree-Qfactorial}
Let $A$ be the local ring of a normal surface over an algebraically closed
field $k$ of characteristic $p>0$.  Let $f:Y\to\operatorname{Spec}A$
be a projective resolution, isomorphic away from the closed point.
Assume that its reduced exceptional divisor $E=\bigcup_{i=1}^rE_i$
has simple normal crossings, that each $E_i$ is isomorphic to
$\mathbf P^1_k$, and that its dual graph is a tree.
Then $A$ is $\mathbf Q$-factorial.  More precisely, there is an integer
$a\geq0$, depending on this resolution, such that
\[
  |\det(E_i\cdot E_j)|p^a\,\operatorname{Cl}(A)=0.
\]
The proposition does not require that $A$ have a rational singularity.
For a regular local ring the conclusion holds with annihilator one.
\end{proposition}

\begin{proof}
We assume that the exceptional divisor is nonempty, and write
$M=(E_i\cdot E_j)$.  The exceptional intersection matrix is negative
definite, with nonnegative off-diagonal entries
\cite[Section 10.1]{Kollar2013D}.
There is a positive integral exceptional cycle $Z$ such that
$Z\cdot E_i<0$ for every $i$.  Here is a direct linear-algebra argument.
Put $B=-M$ and solve $Bv=\mathbf1$ over $\mathbf Q$.
If $v=v_+-v_-$ is the decomposition into its nonnegative positive and
negative parts and $v_-\ne0$, then
\[
 0<v_-^t\mathbf1=v_-^tBv_+-v_-^tBv_-<0,
\]
because the supports of $v_+$ and $v_-$ are disjoint and the off-diagonal
entries of $B$ are nonpositive.  Thus $v\geq0$.
If $v_i=0$, then $(Bv)_i\leq0$, again a contradiction.
Multiplying $v$ by a positive common denominator gives the coefficients
of the required $Z$.

The invertible sheaf $\mathcal O_Z(-Z)$ has positive degree on every
reduced component of $Z$, and is therefore ample
\cite[Lemma 33.44.15, tag 0B5Y]{C_StacksFormalPic}.
Serre vanishing supplies an integer $N\geq1$ such that
\[
 H^1(Z,\mathcal O_Z(-nZ))=0\qquad(n\geq N)
\]
\cite[Lemma 30.17.1, tag 0B5U]{C_StacksFormalPic}.
The thickening $nZ\subset(n+1)Z$ has square-zero ideal
$\mathcal O_Z(-nZ)$.  Its units--Picard exact sequence shows that
\[
 \operatorname{Pic}((n+1)Z)\longrightarrow\operatorname{Pic}(nZ)
 \quad\hbox{is injective},
\]
and that
\[
 H^0((n+1)Z,\mathcal O_{(n+1)Z}^{\times})
   \longrightarrow H^0(nZ,\mathcal O_{nZ}^{\times})
 \quad\hbox{is surjective}
\]
for every $n\geq N$
\cite[Lemma 37.4.1, tag 0C6R]{C_StacksFormalPic}.
In particular, the first finite thickening $NZ$ detects a line bundle
which is trivial on every higher thickening; the surjectivity on units
also permits compatible choices of its trivializations.

Let $\mathcal I\subset\mathcal O_{NZ}$ be the ideal of $E$.
It is nilpotent.  Choose $a\geq0$ such that
$\mathcal I^{p^a}=0$.  The sheaf of abelian groups $1+\mathcal I$
is killed by $p^a$, since $(1+x)^{p^a}=1+x^{p^a}$ in characteristic $p$.
The units sequence
\[
 1\longrightarrow1+\mathcal I\longrightarrow
 \mathcal O_{NZ}^{\times}\longrightarrow\mathcal O_E^{\times}
 \longrightarrow1
\]
therefore shows that $p^a$ kills
$\ker[\operatorname{Pic}(NZ)\to\operatorname{Pic}(E)]$.

We claim that every invertible sheaf $\mathcal L$ on $Y$ having
degree zero on all $E_i$ satisfies $\mathcal L^{\otimes p^a}\simeq
\mathcal O_Y$.  A degree-zero invertible sheaf on $E$ is trivial:
it is trivial on each $\mathbf P^1$ component, and its gluing constants
at the nodes can be removed successively along the tree.
Thus $\mathcal L^{\otimes p^a}|_{NZ}$ is trivial.  The injectivity
above shows that its restriction to every $nZ$ is trivial, and
surjectivity on units makes these trivializations compatible for
$n\geq N$.

First suppose that $A$ is complete.  The powers of the ideal of $Z$
and the powers of $\mathfrak m_A\mathcal O_Y$ define the same adic
topology, because their radicals agree on the Noetherian scheme $Y$.
The theorem on formal functions in degree zero thus algebraizes the
compatible trivializing sections to a section of
$\mathcal L^{\otimes p^a}$ on $Y$
\cite[Theorem 30.20.5, tag 02OC]{C_StacksFormalPic}.
Its vanishing locus is closed and disjoint from $E$.
Its proper image in the local scheme $\operatorname{Spec}A$ is
therefore empty.  Hence the section is everywhere nonvanishing.

For general $A$ we apply this argument after the faithfully flat
base change $A\to\widehat A$.  The cycles $nZ$ and all calculations
on them are unchanged, since each is annihilated by some power of
$\mathfrak m_A$.  Put $\mathcal M=\mathcal L^{\otimes p^a}$.
Properness and flat base change give
\[
 H^0(Y,\mathcal M)\otimes_A\widehat A
   =H^0(Y_{\widehat A},\mathcal M_{\widehat A})\simeq\widehat A,
\]
where we also use $f_*\mathcal O_Y=\mathcal O_{\operatorname{Spec}A}$
\cite[Lemma 30.5.2, tag 02KH]{C_StacksFormalPic}.
Consequently the finite $A$-module $H^0(Y,\mathcal M)$ is free of
rank one.  Its evaluation map on $Y$ is an isomorphism after the same
faithfully flat base change, and hence is an isomorphism before base
change.  This proves the claim.

Finally let $D$ be a Weil divisor on $\operatorname{Spec}A$, and
let $\widetilde D$ be its strict transform.  Since $Y$ is regular,
$\widetilde D$ is Cartier.  Inverting the matrix $M$ supplies rational
numbers $b_i$ such that
\[
 \left(\widetilde D+\sum_i b_iE_i\right)\cdot E_j=0
 \qquad(1\leq j\leq r).
\]
For $m=|\det M|$ all $mb_i$ are integers.  The Cartier divisor
$m\widetilde D+\sum_i mb_iE_i$ thus defines an invertible sheaf
which is killed by $p^a$, by the claim.  Its $p^a$ multiple is a
principal divisor on $Y$.  Pushing that principal divisor down to
$\operatorname{Spec}A$ shows that $mp^aD$ is principal.  This proves
the asserted bounded exponent and the $\mathbf Q$-factorial conclusion.
\end{proof}

\begin{remark}[Source scope]
For one irreducible exceptional curve, the numerical correction and
descent mechanism also occur in Tanaka's Propositions 3.25 and 3.26
\cite[pp. 23--24]{C_Tanaka2014}.  Those propositions are stated for
one exceptional curve.  The argument above handles the entire rational
tree at once and explicitly accounts for the nilpotent thickenings.
It makes no assertion that a rational exceptional tree forces
$R^1f_*\mathcal O_Y=0$.
\end{remark}


\begin{lemma}[Proved: projective partial resolutions retaining $\mathbf Q$-factorial points]
\label{D:Qfactorial-partial-resolution}
Let $S$ be a normal proper surface over an algebraically closed field.
Let $T$ be a subset of its finite singular locus, and assume that
$\mathcal O_{S,t}$ is $\mathbf Q$-factorial for every $t\in T$.
There is a proper birational morphism $q:X\to S$ which is an
isomorphism near $T$, which resolves every singular point outside
$T$, and for which $X$ is projective and $\mathbf Q$-factorial.
At each resolved point we may prescribe a projective regular
resolution which is an isomorphism on the punctured neighborhood.
\end{lemma}
\begin{proof}
Choose an affine neighborhood of each singular point which contains
no other singular point.  A prescribed projective resolution of
the local germ spreads to such a neighborhood by finite presentation;
after shrinking, properness, regularity and the isomorphism over
the punctured neighborhood hold for this spread-out morphism.
The specified local resolutions, the
identity neighborhoods at points of $T$, and the regular locus
glue along their common regular open subsets.  Properness is local
on the target, so the resulting $q$ is proper.  The surface $X$ is
normal, proper, and regular except at its unchanged points above
$T$.  It is therefore $\mathbf Q$-factorial.

We recall why a normal proper $\mathbf Q$-factorial surface is
projective, including the characteristic-independent argument.
Choose a regular resolution $h:W\to X$.  The proper regular
surface $W$ is projective
\cite[Lemma 54.16.11, tag 0C5P]{StacksDSurfaceResolution}.
Let $H$ be an effective ample Cartier divisor on $W$ containing
no exceptional component, and put $A=h_*H$.
The divisor $A$ is effective and $\mathbf Q$-Cartier.
For every integral curve $C\subset X$, the effective
$\mathbf Q$-Cartier pullback $h^*C$ is nonzero, and the projection
formula gives
\[
 A\cdot C=H\cdot h^*C>0.
\]
Moreover $h^*A=H+F$ for an effective exceptional
$\mathbf Q$-divisor $F$.  Since $h^*A\cdot F=0$, we have
\[
 A^2=(h^*A)^2=h^*A\cdot H=H^2+F\cdot H>0.
\]
The Nakai--Moishezon criterion makes a positive Cartier multiple
of $A$ ample.  This is the argument of
\cite[Lemma 2.2, p. 344]{C_Fujino2012}, also applied explicitly
in positive characteristic in
\cite[proof of Theorem 5.3, p. 50]{C_Tanaka2014}.
It proves the required projectivity of $X$.
\end{proof}

\begin{theorem}[Proved under stated hypotheses: mixed surface boundary]
\label{D:surface-mixed-boundary}
Let $S$ be a normal proper connected surface over an algebraically
closed field $k$ of characteristic $p>0$, and let $Y\subset S$
be a finite closed subset.  Suppose that every point $y\in Y$
satisfies at least one of the following conditions:
\begin{enumerate}
\item $\mathcal O_{S,y}$ is $\mathbf Q$-factorial;
\item the germ at $y$ admits a regular resolution whose reduced
exceptional divisor is one smooth elliptic curve;
\item $S$ is log canonical at $y$.
\end{enumerate}
Then $\pi_1(S\setminus Y)$ is topologically finitely generated.
The singularities retained in $S\setminus Y$ are unrestricted.
Condition (1) includes every point admitting a resolution by a
simple normal crossings rational tree, even when that point is
not log canonical, by Proposition~\ref{D:rational-tree-Qfactorial}.
\end{theorem}
\begin{proof}
Put $U=S\setminus Y$.  We leave unchanged every omitted singular
point which is locally $\mathbf Q$-factorial.  Every other omitted
point has a regular resolution with an elliptic or cyclic
exceptional divisor: for condition (2) this is assumed, and for
condition (3) it follows from
Lemma~\ref{D:lc-exceptional-configurations}, since its rational-tree
alternative is already $\mathbf Q$-factorial.
We use those resolutions at these points, blowing up a self-node
if needed, and resolve all the singularities lying in $U$.
Lemma~\ref{D:Qfactorial-partial-resolution} gives a normal
projective $\mathbf Q$-factorial surface $q:X\to S$ with precisely
these local properties.  Set $U'=q^{-1}(U)$.

Let $V\to U$ be a connected finite \'etale cover.  It is normal
and integral.  Its pullback to $U'$ is connected: it contains a
nonempty dense open identical to a nonempty open of $V$, and a
finite \'etale cover of the integral normal $U'$ is a disjoint
union of integral components, each dominating $U'$.
It follows that $\pi_1(U')\to\pi_1(U)$ is surjective.

The normalization of $S$ in $k(V)$ is finite and \'etale over $U$.
At an omitted point, restrict to a neighborhood containing no
other point of $Y$, and apply the corresponding local result.
For an elliptic exceptional curve, all extensions of its valuation
are tame by Proposition~\ref{D:single-curve-residue-tameness}.
For a log canonical cyclic exceptional divisor, the same assertion
is Theorem~\ref{D:lc-cusp-tame}, including the optional node blowup.
These are all the codimension-one components of $X\setminus U'$.
The unchanged $\mathbf Q$-factorial omitted points have codimension
two and create no further divisorial condition.
Thus every cover pulled back from $U$ is classified by
$Q(X,X\setminus U')$, and we have a surjection
\[
 Q(X,X\setminus U')\longrightarrow\pi_1(U).
\]
The source is topologically finitely generated by the projective
case of Proposition~\ref{D:fixed-boundary-positive}.  The same is
therefore true of its quotient.
\end{proof}

\begin{corollary}[Proved: log canonical omitted points on proper surfaces]
\label{D:lc-surface-boundary}
Let $S$ be a normal proper surface over an algebraically closed
field, and let $Y\subset S$ be finite.  If $S$ is log canonical
at every point of $Y$, then every connected component of
$S\setminus Y$ has topologically finitely generated fundamental
group.  No projectivity, Gorenstein, Cartier-index, or degree
restriction on covers is required.

More generally, let $Z$ be a proper scheme of dimension at most
two and let $Y\subset Z$ have codimension at least two in every
irreducible component.  The same conclusion holds if every point
above $Y$ on the normalization of every surface component
satisfies one of the three conditions in
Theorem~\ref{D:surface-mixed-boundary}.
\end{corollary}
\begin{proof}
In positive characteristic the first assertion is the theorem.
In characteristic zero it follows from
Proposition~\ref{D:fixed-boundary-positive} without any restriction
on the singularities.

For the second assertion we first remove nilpotents and then
take the finite disjoint union of the normalizations of the
irreducible components.  On each normal surface component in
positive characteristic the preceding theorem applies to the finite
inverse image of $Y$; in characteristic zero we use the established
characteristic-zero case of
Proposition~\ref{D:fixed-boundary-positive}.
On each component of dimension at most one the inverse image
of $Y$ is empty, so its fundamental group is that of a proper
scheme and is finitely generated.  These opens form a finite
proper surjective cover of $Z\setminus Y$.  Finite proper descent
for finite \'etale covers, in the form used in
Proposition~\ref{D:fixed-boundary-positive}, then gives finite
generation on each connected component downstairs.
\end{proof}


\begin{corollary}[Proved under stated hypotheses: mixed fixed boundary on surfaces]
\label{D:mixed-fixed-boundary}
Let $S$ be a normal proper connected surface over an algebraically
closed field, and let $Y\subsetneq S$ be closed.
Suppose that every singular point $s\in Y$ at which $S$ is not
locally $\mathbf Q$-factorial is isolated in $Y$ and satisfies
at least one of the following conditions:
\begin{enumerate}
\item it admits a regular resolution with one smooth elliptic
curve as reduced exceptional divisor;
\item $S$ is log canonical at $s$.
\end{enumerate}
Then the fixed-boundary quotient $Q(S,Y)$ of D.1 is topologically
finitely generated.  Singularities outside $Y$ are unrestricted.
In positive characteristic, the omitted locally
$\mathbf Q$-factorial points include all points resolved by
simple normal crossings rational trees, without a log canonical
assumption.
\end{corollary}
\begin{proof}
The characteristic-zero case is already
Proposition~\ref{D:fixed-boundary-positive}, so suppose that
the characteristic is positive.  We retain all the locally
$\mathbf Q$-factorial singular points belonging to $Y$.
We resolve all singularities outside $Y$, and at each remaining
singular point of $Y$ use its single elliptic or cyclic resolution.
The cyclic alternative follows from
Lemma~\ref{D:lc-exceptional-configurations}; its rational-tree
alternative is excluded here by
Proposition~\ref{D:rational-tree-Qfactorial}.
Lemma~\ref{D:Qfactorial-partial-resolution} yields a normal
projective surface $q:X\to S$ with these properties.
Put $U=S\setminus Y$ and $U'=q^{-1}(U)$.

Every connected finite \'etale cover of $U$ remains connected
after pullback to $U'$, by the normal birational argument in
Theorem~\ref{D:surface-mixed-boundary}.
Consider one classified by $Q(S,Y)$.
At a strict transform of a codimension-one component of $Y$,
the associated discrete valuation of the function field is
unchanged, so its pullback is tame there by hypothesis.
Every other divisorial component of $X\setminus U'$ lies above
one of the isolated non-$\mathbf Q$-factorial points of $Y$.
Because that point is isolated in $Y$, the cover is \'etale
on its punctured local neighborhood.  The elliptic theorem
\ref{D:single-curve-residue-tameness}, respectively the lc cusp
theorem~\ref{D:lc-cusp-tame}, therefore makes it tame on all of
these new components as well.
There are no additional boundary divisors: the other exceptional
curves lie over $U$ and belong to $U'$.

Thus the composite homomorphism
$\pi_1(U')\to\pi_1(U)\to Q(S,Y)$ is surjective and kills wild
inertia at every codimension-one component of $X\setminus U'$.
It factors through a surjection
\[
             Q(X,X\setminus U')\longrightarrow Q(S,Y).
\]
The source is finitely generated by the projective case of
Proposition~\ref{D:fixed-boundary-positive}, which proves the result.
\end{proof}


\begin{proposition}[Proved under stated hypotheses: finite regular covers]
\label{D:finite-regular-cover}
Let $X$ be a connected proper scheme of finite type over an
algebraically closed field, and let $Y\subset X$ be closed.
Suppose that there is a finite surjection $V\to X$ with $V$ regular
and with $V\times_XY$ of codimension at least two in every component
of $V$. Then $\pi_1(X\setminus Y)$ is topologically finitely generated.
In particular a finite quotient of a proper regular variety cannot
produce a counterexample to D.2(ii) with a codimension-two omitted set.
\end{proposition}
\begin{proof}
The scheme $V$ is proper. Purity gives an equivalence of finite
\'etale covers of $V$ and of $V^\circ=V\setminus(V\times_XY)$,
so every component of $V^\circ$ has finitely generated fundamental
group by \cite[X, Theorem 2.9]{SGA1}. The finite surjection
$V^\circ\to X\setminus Y$ is effective descent for finite \'etale
covers, and its first two nerve levels have finitely many connected
components. The descent generation theorem
\cite[IX, Corollary 5.2, pp. 247--248]{SGA1} proves the conclusion.
No tameness assumption on the finite map $V\to X$ is required.
For a finite quotient of a regular variety, finiteness of the quotient
map preserves the dimensions of closed subsets, so the inverse image
of a codimension-two set again has codimension at least two.
\end{proof}

\begin{proposition}[Proved under stated hypotheses: local purity along the boundary]
\label{D:local-purity-boundary}
Let $Z$ be a normal connected proper scheme over an algebraically
closed field, let $Y\subsetneq Z$ be closed, and put $U=Z\setminus Y$.
Suppose that for every $y\in Y$ at least one of the following holds:
\begin{enumerate}
\item $\mathcal O_{Z,y}$ is regular of dimension at least two;
\item $\mathcal O_{Z,y}$ is a complete intersection of dimension
at least three.
\end{enumerate}
Here complete intersection means that its completion is a quotient
of a regular complete local ring by a regular sequence.
Restriction is an equivalence
\[
                    \operatorname{FEt}(Z)
                 \xrightarrow{\ \sim\ }\operatorname{FEt}(U).
\]
Consequently $\pi_1(U)\to\pi_1(Z)$ is an isomorphism, and these
groups are topologically finitely generated.
In particular, D.2(ii) holds if $Y$ has codimension at least three
and $Z$ is a local complete intersection at every point of $Y$.
Neither finiteness of $Y$ nor any condition on the singularities
retained in $U$ is required.
\end{proposition}
\begin{proof}
We use local purity in its precise essential-surjectivity form:
every finite \'etale cover of the punctured spectrum of a regular
local ring of dimension at least two, or of a complete-intersection
local ring of dimension at least three, extends to a finite \'etale
cover of its spectrum. These are, respectively,
\cite[Lemma 58.21.3(1), tag 0BMA; Proposition 58.25.3, tag 0BPD]{StacksDLocalPurity}.

The hypotheses imply that $Y$ contains no generic or codimension-one
point of $Z$, so $U$ is dense. Since $Z$ is normal and connected,
it is integral. Let $V\to U$ be a connected finite \'etale cover.
Normalize $Z$ in the finite separable extension $k(V)/k(Z)$, and
write $q:\overline V\to Z$ for the result. The normalization is
finite because $Z$ is of finite type over a field. Its restriction
to $U$ is $V$.

Let $B\subset Z$ be the image of the non-\'etale locus of $q$.
It is closed because $q$ is finite, and it is contained in $Y$.
Suppose $B\ne\varnothing$ and choose a generic point $b$ of an
irreducible component of $B$. Put $A=\mathcal O_{Z,b}$.
Every proper generization of $b$ lies outside $B$. Thus the base
change of $q$ to the punctured spectrum of $A$ is finite \'etale.
By one of the two local-purity statements it extends to a finite
\'etale $A$-scheme $W$.

The ring $A$ is normal, and $W$ is normal. Every component of $W$
dominates $\operatorname{Spec}A$: a finite \'etale morphism is open
and closed on each component, and the local base is connected.
Both $W$ and $\overline V\times_Z\operatorname{Spec}A$ are therefore
the integral closure of $A$ in the same finite product of generic
fields. They are isomorphic. It follows that $q$ is \'etale at
every point above $b$, contrary to $b\in B$. Hence $B$ is empty
and $V$ extends. Disconnected covers extend component by component.

Restriction is fully faithful as well: a morphism between two
finite \'etale covers of $Z$ is determined on their dense open
restrictions, and a morphism on those restrictions extends uniquely
by the integral-closure description in their generic fields.
This proves the asserted equivalence and the isomorphism on
fundamental groups, for any chosen geometric base point of $U$.
Finally, the fundamental group of a proper scheme over an
algebraically closed field is topologically finitely generated by
\cite[X, Theorem 2.9, pp.~273--274]{SGA1}. This proves the final
claim and the stated special case.
\end{proof}


\begin{proposition}[Proved: finite local groups and strongly F-regular compactifications]
\label{D:finite-local-groups}
Let $X$ be normal, connected and proper over an algebraically closed
field, and let $Y\subset X$ have codimension at least two. Put
$U=X\setminus Y$. Suppose that the groups
\[
 L_{\bar x}=\pi_1\bigl(\operatorname{Spec}\mathcal O^{\mathrm{sh}}_{X,\bar x}
                 \setminus Y^{\mathrm{sh}}_{\bar x}\bigr)
 \qquad (\bar x\longrightarrow X)
\]
are finite of uniformly bounded order. Then $\pi_1(U)$ is
topologically finitely generated. In particular, this holds for a
finite omitted set whose punctured strict-local fundamental groups
are finite. It also holds for every strongly F-regular proper $X$
in positive characteristic, with arbitrary codimension-two $Y$.
No quasi-projectivity assumption is required.
\end{proposition}
\begin{proof}
We first prove the bounded-local-group criterion, giving the
comparison that uses properness. Write $G=\pi_1(U)$. By
\cite[Proposition 3.2, arXiv pp.~3--4]{BCGST2019D}, $X$ admits a
finite stratification such that the local groups $L_{\bar x}$
and their maps into $G$ are constant along each stratum, with the
usual conjugacy ambiguity for base points. Thus only finitely many
finite image groups $I_1,\ldots,I_r\subset G$ occur, up to
conjugacy. A finite set of nonidentity elements of a profinite group
can be separated from the identity by one finite quotient.
Consequently there is an open normal subgroup $H\subset G$ with
$H\cap I_a=1$ for every $a$. Normality also gives trivial
intersection with every conjugate of these groups.

Let $V\to U$ be the connected finite \'etale cover corresponding
to $H$, and let $\widetilde X\to X$ be its finite normalization.
The scheme $\widetilde X$ is normal and proper, and
$V=\widetilde X\times_XU$. At a geometric point
$\bar z\to\widetilde X$ above $\bar x$, the punctured strict
local scheme is a connected component of the pullback of
$\operatorname{Spec}\mathcal O^{\mathrm{sh}}_{X,\bar x}
\setminus Y^{\mathrm{sh}}_{\bar x}$ to $V$. Its fundamental
group maps into $H$ with image contained in
$H\cap gI_ag^{-1}=1$ for the appropriate $a$ and $g\in G$.
This is the local-cover comparison in the second proof of
\cite[Corollary 4.2, arXiv pp.~7--8]{BCGST2019D}.

It follows that every finite \'etale cover of $V$ is trivial on
each of these punctured strict local schemes. Normalize
$\widetilde X$ in that cover. On a strict localization its finite
normalization is a disjoint union of copies of the normal strict
local base: the generic splitting is already forced by the
trivial cover on the dense punctured open. It is therefore \'etale
at every point above $Y$, and it was \'etale over $V$. This proves
that restriction gives an equivalence
\[
 \operatorname{FEt}(\widetilde X)\simeq\operatorname{FEt}(V),
 \qquad H\simeq\pi_1(\widetilde X).
\]
Full faithfulness here follows from normality and the common
dense open. Since $\widetilde X$ is proper, its fundamental group
is finitely generated by \cite[X, Theorem 2.9]{SGA1}; a profinite
group containing an open finitely generated subgroup is finitely
generated. This proves the criterion. For a finite $Y$, finiteness
of its finitely many local groups supplies the required uniform
bound; the groups at points outside $Y$ are trivial.

For the strongly F-regular assertion we can apply directly
\cite[Corollary 4.2, arXiv p.~6]{BCGST2019D}. A finite-type scheme
over a perfect field is F-finite, so its hypothesis applies here.
It supplies a finite surjective quasi-\'etale map
$q:X'\to X$ with $X'$ normal such that every finite \'etale cover
of $X'_{\mathrm{reg}}$ extends over $X'$. Equivalently, every
further finite quasi-\'etale cover of $X'$ is \'etale: restrict to
the regular locus, extend there by purity, and use uniqueness of
normalization. The inverse image $q^{-1}Y$ has codimension at
least two. Normalization therefore extends every finite \'etale
cover of $U'=X'\setminus q^{-1}Y$ to a quasi-\'etale cover of
$X'$, which by the stated property is \'etale. Hence
$\pi_1(U')=\pi_1(X')$ is finitely generated. Finally
$U'\to U$ is a finite proper surjection, so effective descent and
\cite[IX, Corollary 5.2]{SGA1} imply finite generation of
$\pi_1(U)$. The cited theorem explicitly states that it requires
no quasi-projectivity assumption.
\end{proof}

\begin{proposition}[Proved under stated hypotheses: arbitrary boundary and finite local images]
\label{D:finite-local-images-arbitrary}
Let $X$ be a normal connected proper scheme over an algebraically
closed field, let $Y\subset X$ be closed of codimension at least
two, and put $U=X\setminus Y$.  For a geometric point $\bar x$ of
$Y$, put
\[
 U_{\bar x}=U\times_X\operatorname{Spec}
                    \mathcal O^{\mathrm{sh}}_{X,\bar x},
 \qquad
 I_{\bar x}=\operatorname{im}\bigl(\pi_1(U_{\bar x})
                                  \longrightarrow\pi_1(U)\bigr).
\]
If every $I_{\bar x}$ is finite, then $\pi_1(U)$ is topologically
finitely generated.  No uniform bound on the orders is required,
and $Y$ need not be finite.  It is enough to check the hypothesis
at geometric closed points of $Y$.
\end{proposition}
\begin{proof}
We use the following precise local-to-global tower theorem:
for a normal integral excellent scheme $X$, a tower of finite
quasi-\'etale Galois covers which is eventually \'etale after
strict localization at each geometric point is eventually
\'etale globally
\cite[Theorem 3.7, pp.~708--709]{CarvajalRojasStabler2023D}.
Its statement imposes no quasi-projectivity hypothesis.  In our
finite-type setting it also follows from regular alterations and
the finite branch-locus stratification in
\cite[Theorem 2 and \S4, proof of (i)$\Rightarrow$(iii),
pp.~6--7]{Stibitz2017D}.  This is a theorem about all finite
quasi-\'etale covers, with no restriction on their degrees.

Write $G=\pi_1(U)$.  For an open normal subgroup $H\subset G$,
let $V_H\to U$ be its connected finite \'etale cover and let
$X_H\to X$ be the finite normalization in $k(V_H)$.  We claim
that some $H$ satisfies
\[
              \operatorname{FEt}(X_H)
                    \simeq\operatorname{FEt}(V_H).
\]
If this failed for every $H$, we could construct successively
open normal subgroups $G=H_0\supset H_1\supset H_2\supset\cdots$
so that every $X_{H_{n+1}}\to X_{H_n}$ were non-\'etale.  Indeed,
choose a finite \'etale cover of $V_{H_n}$ whose finite
normalization over $X_{H_n}$ is not \'etale, and dominate its
generic field by a finite Galois extension of $k(X)$ arising
from the Galois closure of that cover over $U$.  If the
dominating normalization were \'etale over $X_{H_n}$, its
intermediate normal finite cover would also be \'etale, contrary
to the choice.  Every $X_{H_n}\to X$ is quasi-\'etale because
its branch locus is contained in $Y$.

Fix $\bar x$.  The descending subgroups
$I_{\bar x}\cap H_n$ of the finite group $I_{\bar x}$ stabilize.
The connected components of $V_{H_n}\times_U U_{\bar x}$ are
the covers determined by these intersections.  Once two
successive intersections agree, each component upstairs maps
isomorphically to its corresponding component downstairs.
Finite normalization commutes with \'etale localization, and
normality identifies the normalization of a split generic
algebra with a disjoint union of copies of the normal base.
Consequently $X_{H_{n+1}}\to X_{H_n}$ is \'etale over the
strict localization at $\bar x$ for all sufficiently large
$n$.  At points of $U$ it is already \'etale.  The tower theorem
now contradicts our construction.  This proves the claim.

For this $H$, full faithfulness of restriction follows from
normalization in generic fields, so the displayed equivalence
gives $H\simeq\pi_1(X_H)$.  The scheme $X_H$ is proper;
\cite[X, Theorem 2.9]{SGA1} makes its fundamental group
topologically finitely generated.  Adjoining finitely many
coset representatives generates $G$.

Finally, every point $x\in Y$ specializes to a closed point
$y\in Y$, since $X$ is of finite type over a field.  Choose a
geometric lift of $x$ to
$\operatorname{Spec}\mathcal O^{\mathrm{sh}}_{X,\bar y}$.
Strict localization there is a strict localization of $X$ at
$x$, because the map to $X$ is ind-\'etale.  Thus the map
from $\pi_1(U_{\bar x})$ to $G$ factors through
$\pi_1(U_{\bar y})$, up to the usual base-point conjugacy.
It follows that $I_{\bar x}$ is contained in a conjugate of
$I_{\bar y}$.  Finiteness at geometric closed points suffices.
\end{proof}

\begin{corollary}[Proved under stated hypotheses: klt threefold boundary]
\label{D:klt-threefold-boundary}
Let $k$ be algebraically closed of characteristic $p>5$, let
$X$ be a normal connected proper threefold over $k$, and let
$Y\subset X$ be closed of codimension at least two.  Suppose
that $X$ has klt singularities on a neighborhood of every
point of $Y$.  Then $\pi_1(X\setminus Y)$ is topologically
finitely generated.  The retained singularities are unrestricted,
and neither projectivity nor a local-complete-intersection
hypothesis is required.
\end{corollary}
\begin{proof}
For a geometric closed point $\bar y$ of $Y$, take an affine
neighborhood on which the klt hypothesis holds.  Its strict
localization $S$ and the closed subset $Z=Y\times_X S$ satisfy
Setup 1.1 of \cite[p.~694]{CarvajalRojasStabler2023D}:
$Z$ is arbitrary of codimension at least two, not necessarily
just the closed point or the singular locus.  Their Theorem A
\cite[p.~695; Theorem 4.3 and Corollary 4.5,
pp.~713--714]{CarvajalRojasStabler2023D} states that the full
group $\pi_1(S\setminus Z)$ is finite of order prime to $p$.
Its conclusion is not merely finiteness of the maximal
prime-to-$p$ quotient.  The proof of Corollary 4.5 explicitly
excludes a finite quotient of order divisible by $p$ by passing
to a Sylow $p$-subgroup and applying Corollary 2.15 to the
intermediate klt strict-local cover.  In particular the image
$I_{\bar y}$ in the preceding proposition is finite.
That proposition proves the claim.

This argument uses the published local theorem only on affine
neighborhoods and the separately stated scheme-theoretic tower
theorem globally.  It therefore does not depend on whether the
word ``variety'' in a global application includes a
quasi-projectivity convention, or on the additional sentence
in Corollary 4.7 about preserving a prescribed big open.
\end{proof}


\begin{proposition}[Proved: finite global images of local groups suffice]
\label{D:finite-local-images}
Let $X$ be a connected normal proper variety over an algebraically
closed field, let $Y$ be a finite set of closed points of codimension
at least two, and put $U=X\setminus Y$.  For each geometric point
$\bar y$ over $Y$, write
\[
 U_{\bar y}=U\times_X
          \operatorname{Spec}\mathcal O^{\mathrm{sh}}_{X,\bar y}.
\]
Suppose that each image
\[
 I_{\bar y}=\operatorname{im}
       \bigl(\pi_1(U_{\bar y})\longrightarrow\pi_1(U)\bigr)
\]
is finite.  Then $\pi_1(U)$ is topologically finitely generated.
The local groups themselves are allowed to be infinite.
\end{proposition}
\begin{proof}
Write $G=\pi_1(U)$.  Base-point choices change $I_{\bar y}$ only
by conjugacy.  Since $Y$ and all the $I_{\bar y}$ are finite, there
is one open normal subgroup $H\subset G$ meeting every
$I_{\bar y}$ trivially.  Let $V\to U$ be the associated connected
finite etale cover and let $X_H$ be the finite normalization of $X$
in $k(V)$.  It is normal and proper, with $V=X_H\times_XU$.

We show that every finite etale cover of $V$ extends over $X_H$.
For a geometric point $\bar z$ of $X_H$ over $\bar y$, its strict
localization is the appropriate local factor of the normalization
of $\operatorname{Spec}\mathcal O^{\mathrm{sh}}_{X,\bar y}$ in a
connected component of $V\times_UU_{\bar y}$.  This description
follows from finiteness of normalization and its compatibility
with etale localization.  In particular, its punctured open is
that connected component.  Covering-space subgroup comparison
shows that the image of its fundamental group in $H$ is contained
in $H\cap gI_{\bar y}g^{-1}$ for a suitable $g\in G$; this group
is trivial.  Thus any finite etale cover of $V$ becomes a disjoint
union of trivial covers on each punctured strict localization.

Normalize $X_H$ in such a cover.  On each strict localization
above $Y$, the generic algebra is split.  Normality of the strict
local base then identifies its finite normalization with a
disjoint union of copies of that base.  Hence the extension is
etale above $Y$, and it is already etale over $V$.  Uniqueness of
normalization also proves full faithfulness of restriction, so
\[
 \operatorname{FEt}(X_H)\simeq\operatorname{FEt}(V),
 \qquad H\simeq\pi_1(X_H).
\]
This is the same strict-local comparison as the second proof of
\cite[Corollary 4.2, arXiv pp.~7--8]{BCGST2019D}; only finiteness
of the images, rather than of their source groups, is used here.
The proper group $\pi_1(X_H)$ is finitely generated by
\cite[X, Theorem 2.9]{SGA1}.  Adjoining finitely many coset
representatives gives finite generation of $G$.
\end{proof}

\begin{proposition}[Proved: smallness for arbitrary codimension-two boundaries]
\label{D:codim-two-small}
\label{D:finite-boundary-nilpotent}
Let $Z$ be proper over an algebraically closed field $k$, and let
$Y\subset Z$ be closed, of codimension at least two in every
irreducible component. Let $U$ be a connected component of
$Z\setminus Y$. For every connected finite \'etale cover
$V\to U$ and every finite group $F$, the set $H^1(V,F)$ is finite.
Equivalently, every open subgroup of $\pi_1(U)$ is \emph{small}:
it has only finitely many open subgroups of each fixed index.
In positive characteristic, its maximal pronilpotent quotient is
also topologically finitely generated. Normality, projectivity,
and finiteness of the omitted set are not required.
\end{proposition}
\begin{proof}
Put $j:Z\setminus Y\hookrightarrow Z$. Gabber's theorem gives
constructibility of $R^1j_*(F)$ for every finite group $F$,
without any restriction on its order
\cite[Theorem 1.3, p.~529]{ILOXXID}.
The published statement assumes only excellence and the
displayed codimension condition in each component; in
particular it does not require normality.
Apply its proper-composition lemma to
$Z\setminus Y\xrightarrow{j}Z\to\operatorname{Spec}k$.
It says that constructibility of $R^1j_*F$ implies
constructibility of $R^1f_*F$ for the composite $f$, since
$Z\to\operatorname{Spec}k$ is proper
\cite[Lemma 3.2.1(ii), pp.~536--537]{ILOXXID}.
On an algebraically closed point this says exactly that
$H^1(Z\setminus Y,F)$ is finite. Each connected component
therefore has the same finiteness property. The cited proof
uses nonabelian Leray and twisting; it too imposes no
invertibility assumption on $|F|$.

Write $\Pi=\pi_1(U)$. Taking $F=S_n$, degree-$n$ finite
\'etale covers are classified by $H^1(U,S_n)$. Each connected
cover corresponds to at most $n$ conjugate index-$n$ subgroups.
Thus $\Pi$ is small. Every open subgroup $H\subset\Pi$ is
small too: an index-$n$ subgroup of $H$ has index
$n[\Pi:H]$ in $\Pi$. Conversely, for a fixed finite $F$,
the kernel of a continuous $H\to F$ has index at most $|F|$.
There are finitely many possible kernels and finitely many
maps from each resulting finite quotient to $F$. This
proves finiteness of $H^1(V,F)$ for every connected finite
\'etale cover $V\to U$.

Suppose now that $\operatorname{char}k=p>0$. Taking
$F=\mathbb F_p$ shows that the maximal pro-$p$ quotient
of $\pi_1(V)$ has finite generator number $d_p$, by the
pro-$p$ Frattini criterion.

By \cite[Theorem 4.11]{Orgogozo2003}, the maximal prime-to-$p$
quotient of $\pi_1(V)$ has finite generator number $d'$.
Therefore every maximal pro-$\ell$ quotient for $\ell\ne p$
has generator number at most $d'$. The maximal pronilpotent
quotient is the product of all these pro-$\ell$ quotients.
Choosing $\max(d_p,d')$ generators in each factor and combining
them coordinatewise generates that product: in every finite
nilpotent quotient, the Sylow factors have relatively prime
orders, so a subgroup surjecting onto all factors is the whole
quotient. Since $\pi_1(V)$ runs through the open subgroups
of $\pi_1(U)$, the proposition follows.
\end{proof}

\begin{proposition}[Proved: smallness of every fixed-boundary tame quotient]
\label{D:fixed-boundary-small}
Let $X$ be normal, proper and connected over an algebraically
closed field $k$ of characteristic $p>0$, let $Y\subsetneq X$
be closed, and let $Q(X,Y)$ have the fixed-divisor meaning of
Proposition~\ref{D:fixed-boundary-positive}. Then $Q(X,Y)$
is small. Every open subgroup has topologically finitely
generated maximal pronilpotent quotient. No projectivity or
normal-crossings hypothesis is required.
\end{proposition}
\begin{proof}
Put $U=X\setminus Y$. We prove finiteness of the $F$-torsors
on $U$ classified by $Q(X,Y)$ for each fixed finite group $F$.
Let $D_1,\ldots,D_r$ be the codimension-one components of $Y$,
and let $n$ be the prime-to-$p$ part of $|F|$. Weak approximation
for the distinct discrete valuations $v_{D_i}$ gives functions
$f_i\in k(X)^\times$ with
\[
 v_{D_j}(f_i)=\delta_{ij}
 \qquad(1\leq i,j\leq r);
\]
see \cite[Theorem 14.3.1]{SteinWeakD}. Let $q:X'\to X$ be
the finite normalization in the field obtained by adjoining
all $f_i^{1/n}$. When $r=0$ or $n=1$ we use the identity
extension. In every case $X'$ is normal and proper.

Choose a closed subset $T\subset X$ of codimension at least
two containing the singular locus of $X$, the components of
$Y$ of codimension at least two, the singular loci and pairwise
intersections of the $D_i$, and the intersections of the $D_i$
with all the other zero or pole divisors of the $f_j$.
This choice depends only on the functions and the fixed
boundary. On $X_0=X\setminus T$ the boundary consists of
disjoint regular divisors. Near each $D_i\cap X_0$, the function
$f_i$ is a regular parameter times a unit and every $f_j$ for
$j\ne i$ is a unit. Thus $W=q^{-1}(X_0)$ is regular near
its boundary over $Y$: adjoining the root of a parameter
gives a regular ring, while adjoining the roots of these
units is \'etale. Passing to a field factor and normalization
does not change that local description.

Let $P\to U$ be one of the stated $F$-torsors, and put
$U_0=U\setminus T$, $U'_0=q^{-1}(U_0)$. The ramification
indices of $P$ at the $D_i$ divide $n$, since the cover is
tame and its monodromy group is a subgroup of $F$.
The ramification index of $q$ at each $D_i$ is $n$.
Abhyankar's lemma therefore makes the pullback to $U'_0$
unramified at every codimension-one point of $W\setminus U'_0$
\cite[Lemma 15.116.4, tag 0BRM]{StacksDTame}.
Normalize $W$ in that pullback. The normalization is finite
and already \'etale on $U'_0$; near the remaining boundary
the target is regular. Purity makes it \'etale everywhere
\cite[X, Theorem 3.3]{SGA1}. The $F$-action and the torsor
identities extend uniquely by normality. Equivalently, one
may use the regular-divisor form of Abhyankar's lemma
\cite[Lemma 58.31.5, tag 0EYG]{StacksDTame}.

The complement $X'\setminus W=q^{-1}(T)$ has codimension
at least two, since $q$ is finite. Hence
Proposition~\ref{D:codim-two-small} makes $H^1(W,F)$ finite.
It remains to verify that pullback has finite fibres; an
injectivity assertion would be unnecessary and generally false.
For a fixed torsor on $U'_0$, descent data along the finite
proper surjection $U'_0\to U_0$ form a subset of the
$F$-equivariant isomorphisms between its two pullbacks to
$U'_0\times_{U_0}U'_0$. If nonempty, this set is a torsor
under sections of a finite locally constant inner form of
$F$. The double fibre product is of finite type and has
finitely many connected components, so that section set is
finite. Effective descent
\cite[Theorem 5.19(ii)]{Rydh2010} gives finite pullback fibres.
Finally $H^1(U,F)\to H^1(U_0,F)$ is injective, because
$U$ is normal and $U_0$ is dense: isomorphisms between
finite \'etale covers extend uniquely from the generic
point. Thus the original set of fixed-boundary tame
$F$-torsors is finite. Taking symmetric groups proves
smallness of $Q(X,Y)$.

Every open subgroup $H\subset Q(X,Y)$ is small, by the
index argument in Proposition~\ref{D:codim-two-small}.
The inverse image of $H$ in $\pi_1(U)$ corresponds to a
connected finite \'etale cover $V\to U$, and
$\pi_1(V)\twoheadrightarrow H$. Smallness makes
$H$'s maximal pro-$p$ quotient finitely generated, while
\cite[Theorem 4.11]{Orgogozo2003} gives one finite bound
for the generator numbers of all its maximal pro-$\ell$
quotients with $\ell\ne p$. Combining these bounds as
in the preceding proposition proves the pronilpotent assertion.
\end{proof}

\subsubsection{Vertical boundary in a proper curve fibration}

\begin{proposition}[Proved: the homotopy sequence after deleting vertical boundary]
\label{D:vertical-fibration-exact}
Let $k$ be algebraically closed, let $X$ be a smooth connected
proper surface, and let $f:X\to B$ be a morphism to a smooth
connected proper curve.  Let $B^\circ\subset B$ be a nonempty
open over which $f$ is smooth with geometrically connected
one-dimensional fibres.  Let $U\subset X$ be an open containing
$X^\circ=f^{-1}(B^\circ)$.  Suppose that for every
$b\in B\setminus B^\circ$, the open $U$ contains a point of
$X_b$ at which $f$ is smooth.  Then, with compatible geometric
base points, the sequence
\[
 \pi_1(X_{\overline\eta})\longrightarrow\pi_1(U)
       \longrightarrow\pi_1(B)\longrightarrow1
\]
is exact.  If $g$ is the genus of the geometric generic fibre
and $h$ is the genus of $B$, then $\pi_1(U)$ can be generated
topologically by at most $2g+2h$ elements.
\end{proposition}
\begin{proof}
The schemes $X$, $U$ and $X^\circ$ are normal and integral.
Restriction of a connected finite \'etale cover of $U$ to the
dense open $X^\circ$ remains connected.  Thus
\[
 H:=\pi_1(X^\circ)\longrightarrow G:=\pi_1(U)
\]
is surjective.  The proper smooth homotopy sequence gives
\[
 \pi_1(X_{\overline\eta})\longrightarrow H
       \longrightarrow\pi_1(B^\circ)\longrightarrow1
\]
by \cite[X, Corollary 1.4, p. 263]{SGA1}.  Let $N$ be the image
of the first group in $G$.  It is a closed normal subgroup,
and $\pi_1(B^\circ)$ surjects onto $G/N$.

We show that this surjection factors through $\pi_1(B)$.
Let $Q$ be any finite quotient of $G/N$.  It gives a connected
finite \'etale $Q$-torsor $V\to U$.  Its restriction to
$X^\circ$ comes from a finite \'etale $Q$-torsor
$B'^\circ\to B^\circ$, by the proper homotopy sequence.
Normalize $B$ in $k(B'^\circ)$ to obtain the finite map
$B'\to B$.

Fix $b\in B\setminus B^\circ$ and an $f$-smooth closed point
$x\in U\cap X_b$.  Such a closed point exists in the indicated
nonempty smooth locus.  Over the strictly henselian trait
$T=\operatorname{Spec}\mathcal O^{\mathrm{sh}}_{B,b}$, smoothness
gives a section $s:T\to U\times_B T$ specializing to $x$.
To see this directly, choose \'etale coordinates for $f$ at $x$,
lift their residue values to $T$, and use the section of the
resulting \'etale neighborhood of the closed point of $T$;
see \cite[tags 054L and 03QO]{C_StacksFibration}.
The pullback $s^*V$ is finite \'etale over $T$.  On its generic
point it is precisely the base change of $B'^\circ$.
Since $T$ is normal, this finite \'etale cover is the
normalization of $T$ in that generic algebra.  Thus $B'\to B$
is \'etale at $b$.  Doing this at every omitted base point
shows that $B'\to B$ is finite \'etale.

Every finite quotient of $G/N$ therefore kills the kernel of
$\pi_1(B^\circ)\to\pi_1(B)$.  Their intersection is trivial,
so the surjection to $G/N$ factors continuously through
$\pi_1(B)$.  The natural map $G\to\pi_1(B)$ induced by $f$
kills $N$.  These two maps between $G/N$ and $\pi_1(B)$ are
inverse, as is checked after the common surjection from
$\pi_1(B^\circ)$.  This proves exactness.

A smooth proper curve of genus $r$ over an algebraically closed
field has a fundamental group generated by $2r$ elements
\cite[X, Theorem 2.6, pp. 272--273]{SGA1}.
Thus $N$ has at most $2g$ generators and $G/N$ has at most
$2h$ generators.  A generating set of $N$, together with lifts
of a generating set of $G/N$, generates $G$ topologically.
This gives the stated bound.
\end{proof}

\begin{proposition}[Proved: vertical boundary meeting every fibre]
\label{D:vertical-fibration-multiplicity}
Let $k$ be algebraically closed, let $X$ be a normal connected
proper surface, and let $f:X\to B$ be a morphism to a smooth
connected proper curve.  Suppose that over a nonempty open
$B^\circ\subset B$ the morphism is flat with geometrically
reduced, geometrically connected one-dimensional fibres.
Let $U\subset X$ be an open containing $X^\circ=f^{-1}(B^\circ)$
and meeting every fibre of $f$.  Then $\pi_1(U)$ is topologically
finitely generated.  No restriction on the multiplicities of
the exceptional fibres is required.

More precisely, let $N$ be the image of
$\pi_1(X_{\overline\eta})$ in $\pi_1(U)$ and put
$Q=\pi_1(U)/N$.  For each $b\in B\setminus B^\circ$, let $m_b$
be the gcd of the multiplicities of the fibre components whose
generic points lie in $U$.  The image in $Q$ of the full inertia
group at $b$ is finite of order dividing $m_b$.
If all $m_b$ equal one, then $Q=\pi_1(B)$ and the homotopy
sequence in Proposition~\ref{D:vertical-fibration-exact} is exact.
If the geometric generic fibre is smooth of genus $g$, $B$ has
genus $h$, and every $m_b$ is prime to the positive
characteristic, then $\pi_1(U)$ has at most $2g+2h+r$ generators,
where $r=\#(B\setminus B^\circ)$.
\end{proposition}
\begin{proof}
Since $X$ is normal and integral, restriction to $X^\circ$
gives a surjection $H=\pi_1(X^\circ)\to G=\pi_1(U)$.
SGA~1 uses ``separable'' for a flat morphism with geometrically
reduced fibres.  Its proper homotopy sequence therefore applies
to $f|_{X^\circ}$ exactly under the present assumptions
\cite[X, Definition 1.1 and Corollary 1.4, pp. 261--263]{SGA1}.
Thus $N$ is closed and normal, and there is a surjection
\[
                    \pi_1(B^\circ)\longrightarrow Q.
\]
The proper reduced connected geometric generic curve has a
finitely generated fundamental group: normalize its finitely
many irreducible components, use the proper smooth-curve theorem
on each normalization, and apply finite proper descent
\cite[IX, Corollary 5.2; X, Theorem 2.6]{SGA1}.
Consequently $N$ is finitely generated.

Let $Q_0$ be any finite quotient of $Q$.  As in the preceding
proof, the corresponding finite \'etale cover $V\to U$ is over
$X^\circ$ the pullback of a connected finite Galois cover
$B'^\circ\to B^\circ$ with group $Q_0$.  Normalize $B$ in its
function field to obtain $B'\to B$.

Choose a fibre component $D$ over $b$ whose generic point lies
in $U$, and let $m$ be its multiplicity.  Such a component
exists because $U$ meets the fibre and its intersection with
the fibre is open.  Choose a prolongation to $k(V)$ of the
$D$-valuation of $k(X)$.  The first ramification index is one,
since the cover is \'etale on $U$.  The restriction of the
$D$-valuation to $k(B)$ has index $m$.  On the other hand,
$k(B')$ is an intermediate field of $k(V)/k(B)$.  Multiplication
of indices of these discrete value groups yields
\[
                         e(b'/b)\mid m.
\]
This argument applies to every retained component.  Thus
$e(b'/b)$ divides their gcd $m_b$.

Every closed point of $B'$ has residue field $k$.  Since the
base cover is Galois, its full inertia group at $b'$ has order
$e(b'/b)$; here finite extensions of the excellent curve DVRs
satisfy the usual degree formula.  It follows that in every
finite quotient of $Q$, the image of inertia at $b$ has order
dividing $m_b$.  Its image $I_b$ in $Q$ is therefore finite,
with order dividing $m_b$.  Indeed, the orders of all its finite
quotients are bounded; choose a quotient of maximal order and
observe that every refinement has the same kernel on $I_b$.
The intersection of these kernels is trivial.

There are only finitely many omitted base points.  Choose an
open normal subgroup $Q'\subset Q$ meeting every $I_b$ trivially;
this is possible by separating the finitely many nonidentity
elements of these finite groups in finite quotients of $Q$.
Let $B_1^\circ\to B^\circ$ be the connected finite \'etale cover
corresponding to the inverse image of $Q'$ in $\pi_1(B^\circ)$,
and let $B_1$ be its smooth proper compactification.  Every
inertia group of $B_1^\circ$ above a point of
$B\setminus B^\circ$ is an intersection with a conjugate of the
corresponding base inertia.  Its image in $Q'$ is therefore
trivial.  Hence the surjection
\[
                  \pi_1(B_1^\circ)\longrightarrow Q'
\]
factors through $\pi_1(B_1)$: a finite cover of a smooth open
curve extends \'etale to the proper curve exactly when all
these inertia groups act trivially.  The proper smooth-curve
theorem makes $Q'$ finitely generated.  A profinite group with
a finitely generated open subgroup is finitely generated, so
$Q$ is finitely generated.  Since $N$ is finitely generated,
so is $G$.

If $m_b=1$ at every omitted point, all $I_b$ are already
trivial.  The argument in
Proposition~\ref{D:vertical-fibration-exact} then identifies
$Q$ with $\pi_1(B)$ and proves the asserted exactness.
If instead all $m_b$ are prime to the characteristic, every
finite base cover just considered is tame at infinity.
Thus $Q$ is a quotient of $\pi_1^{\mathrm t}(B^\circ)$, which
has at most $2h+r$ generators
\cite[XIII, Corollary 2.12, pp. 392--393]{SGA1}.
When the geometric generic fibre is smooth, $N$ has at most
$2g$ generators, giving the stated numerical bound.
\end{proof}

\begin{theorem}[Proved: finite punctures on a proper surface mapping to a curve]
\label{D:surface-map-to-curve}
Let $S$ be a normal connected proper surface over an
algebraically closed field $k$.  If $S$ admits a nonconstant
morphism to a proper curve, then
\[
                         \pi_1(S\setminus Y)
\]
is topologically finitely generated for every finite closed
subset $Y\subset S$.
The morphism need not be generically smooth: geometrically
reduced singular generic fibres are allowed.
\end{theorem}
\begin{proof}
Replace the target by its reduced image and take Stein
factorization.  We obtain a proper surjective map $f:S\to B$
with $f_*\mathcal O_S=\mathcal O_B$ and geometrically connected
fibres.  The curve $B$ is normal, hence smooth and proper over
$k$, and $F=k(B)$ is algebraically closed in $K=k(S)$.
These assertions follow from Stein factorization and normality
of $S$ \cite[Theorem 37.53.4, tag 03H0]{C_StacksFibration}.
To make the constant-field assertion explicit, the global
functions on the proper normal generic fibre are exactly the
relative algebraic closure of the base field in $K$: an element
algebraic over the base is integral over every local ring of
that fibre and hence is regular there by normality.

We check generic separability in positive characteristic.
The field $F$ has transcendence degree one over the perfect
field $k$, so $[F:F^p]=p$.  Choose $t\in F\setminus F^p$;
then $F^{1/p}=F(t^{1/p})$.  A $p$th root of $t$ in $K$ would
be algebraic over $F$ and therefore lie in $F$, contradicting
the choice of $t$.  It follows that
\[
 K\otimes_F F^{1/p}\simeq K[T]/(T^p-t)
\]
is a field.  The separability criterion
\cite[Lemma 10.44.2, tag 05DT]{C_StacksFibration} implies that
$K/F$ is separable and geometrically reduced.  In characteristic
zero this conclusion is automatic.  Relative algebraic
closedness also gives geometric irreducibility
\cite[Lemma 10.47.8, tag 037P]{C_StacksFibration}.
Thus the generic fibre is geometrically integral.  This does
not assert that it is smooth.

The morphism $f$ is flat: at a point above a closed point of
$B$, its local ring is torsion-free over the DVR of $B$, hence
flat.  Geometric reducedness of the generic fibre spreads to
a nonempty open of $B$, by the openness theorem for flat proper
morphisms
\cite[Lemma 37.26.7, tag 0C0E]{C_StacksFibration}.
Shrink this open further to avoid $f(Y)$ and denote it by
$B^\circ$.  Its fibres are geometrically reduced and connected,
so the proper separable homotopy theorem applies.

Put $U=S\setminus Y$.  Then $f^{-1}(B^\circ)\subset U$.
Every fibre of the dominant map from the integral surface to
the smooth curve has a one-dimensional component; equivalently,
the local fibre parameter cuts a nonzero principal divisor.
Deleting finitely many closed points leaves the generic points
of all these components.  Thus $U$ meets every fibre of $f$.
Proposition~\ref{D:vertical-fibration-multiplicity} gives the
claim.
\end{proof}

\begin{corollary}[Proved: birational and finite-normalization descent of the fibration criterion]
\label{D:vertical-fibration-models}
Let $S$ be a normal connected proper surface over an
algebraically closed field, and let $V\subset S$ be a nonempty
open.  Suppose a proper birational map $q:X\to S$ from a
smooth proper surface has a curve fibration for which
$U=q^{-1}(V)$ satisfies either
Proposition~\ref{D:vertical-fibration-exact} or
Proposition~\ref{D:vertical-fibration-multiplicity}.
Then $\pi_1(V)$ is topologically finitely generated, and inherits any generator bound supplied by the relevant case.  In particular, the same is
true of any fixed-boundary tame quotient of $\pi_1(V)$.

For a possibly nonnormal proper scheme of dimension at most
two, the same conclusion holds if its reduced components have
finite normalizations and their inverse-image opens satisfy
this condition, with lower-dimensional components having
finitely generated fundamental groups.  In particular, the
latter requirement is automatic for a boundary of codimension
at least two in every component.
\end{corollary}
\begin{proof}
A connected finite \'etale cover of the normal integral $V$
has a single function field.  Pulling it back along the
birational map $U\to V$ preserves that generic field.  Every
finite \'etale component over the normal integral $U$
dominates $U$, so the pullback remains connected.  Hence
$\pi_1(U)\to\pi_1(V)$ is surjective.  The first assertion and any stated generator bound follow from the propositions.

For the second assertion use the finite proper surjection
from the disjoint union of the normalizations of the reduced
components, restricted to the indicated opens.  Finite proper
descent of finite \'etale covers, in the form already used in
Proposition~\ref{D:fixed-boundary-positive}, shows that finite
generation descends
\cite[IX, Corollary 5.2]{SGA1}.  Under the final codimension
condition, components of dimension at most one lose no points
and remain proper, so their fundamental groups are finitely
generated.
\end{proof}

\subsubsection{Arbitrary fixed boundary on a surface mapping to a curve}

\begin{proposition}[Proved: the logarithmic fibration criterion]
\label{D:log-fibration-criterion}
Let $k$ be algebraically closed, let $S$ be a normal connected
proper surface, and let $f:S\to B$ be a morphism to a smooth
connected proper curve.  Let $Y\subsetneq S$ be closed and put
$U=S\setminus Y$.  Suppose that there is a nonempty open
$B^\circ\subset B$ such that $S^\circ=f^{-1}(B^\circ)\to B^\circ$
is a smooth proper family of geometrically connected curves,
and $(Y\cap S^\circ)_{\mathrm{red}}$ is a finite \'etale relative
divisor $D^\circ$ in that family, possibly empty.  Then the
fixed-boundary tame group $Q(S,Y)$ is topologically finitely
generated.  The boundary may contain whole fibres, and there
is no restriction on the multiplicities of the bad fibres.

More precisely, let $N$ be the image in $Q(S,Y)$ of the tame
fundamental group of the geometric generic pointed fibre.
Then $N$ is closed, normal, and finitely generated, and
$Q(S,Y)/N$ is a quotient of $\pi_1(B^\circ)$.  In characteristic
$p>0$, the image in this quotient of the wild inertia at any
$b\in B\setminus B^\circ$ is finite, of order dividing the
$p$-part of the gcd of the multiplicities of all components
of $S_b$.
\end{proposition}
\begin{proof}
Give $B^\circ$ the trivial log structure and give $S^\circ$
the log structure of the relative divisor $D^\circ$.
This is a proper log smooth morphism of fs log schemes, its
base is log regular, and its underlying fibres are geometrically
connected.  The logarithmic homotopy theorem gives an exact
sequence
\[
 \pi_1^{\mathrm t}\bigl(S_{\overline\eta}\setminus
                       D_{\overline\eta}\bigr)
 \longrightarrow \pi_1(S^{\circ,\log})
 \longrightarrow \pi_1(B^\circ)\longrightarrow1.
\]
Here the base geometric point has trivial log structure, so
the reduced-covering-point limit in the theorem has the
ordinary geometric point as a terminal object.  We use
\cite[Proposition 2.1 and Theorem 2.3, pp. 13--16]{Hoshi2006B}.
The middle group classifies finite \'etale covers of
$U^\circ=S^\circ\setminus D^\circ$ tame at the codimension-one
points of $D^\circ$, by
\cite[Proposition A.10, pp. 57--58]{Hoshi2006B}.
Thus its notion of tameness is exactly the one imposed by
the fixed-boundary quotient.

Restriction of a connected finite \'etale cover of the normal
integral $U$ to the dense open $U^\circ$ remains connected.
Every cover counted by $Q(S,Y)$ is tame along $D^\circ$.
Consequently the natural homomorphism
\[
             \pi_1(S^{\circ,\log})\longrightarrow G:=Q(S,Y)
\]
is surjective.  The image $N$ of the first group in the exact
sequence is therefore closed and normal in $G$, and there is
a surjection
\[
                    \pi_1(B^\circ)\longrightarrow Q:=G/N.
\]
The tame fundamental group of a smooth geometric curve with
finitely many punctures is finitely generated
\cite[XIII, Corollary 2.12, pp. 392--393]{SGA1}.
Hence $N$ is finitely generated.

We prove that $Q$ has finite wild inertia at every omitted
base point.  Let $Q_0$ be any finite quotient of $Q$.
Its connected torsor $V\to U$ is tame along the original
boundary $Y$, and its restriction to $U^\circ$ comes from a
connected finite Galois cover $B'^\circ\to B^\circ$ with group
$Q_0$.  Normalize $B$ in its function field to obtain the
finite map $B'\to B$.

Fix $b\in B\setminus B^\circ$ and a component $E$ of the
fibre $S_b$ with multiplicity $m_E$.  Let $w$ be any
prolongation of its discrete valuation to $k(V)$, and let
$e_E$ be the ramification index over the $E$-valuation of
$k(S)$.  If the generic point of $E$ belongs to $U$, then
$e_E=1$, since $V\to U$ is \'etale.  Otherwise $E$ is an
original boundary divisor, so $e_E$ is prime to $p$ in positive
characteristic, since this cover is tame there.  Restrict $w$
to the intermediate field $k(B')$.  Multiplication of indices
of the discrete value groups in the two towers gives
\[
                  e(b'/b)\mid m_E e_E.
\]
In characteristic $p>0$, this implies
\[
                v_p\bigl(e(b'/b)\bigr)\leq v_p(m_E)
                \qquad\hbox{for every component }E.
\]
All closed-point residue fields of $B'$ equal $k$.  For this
finite Galois cover of smooth curves the inertia order is
therefore $e(b'/b)$, and the image of wild inertia has order
equal to its $p$-part.  Thus the image of the wild inertia
group $P_b$ in every finite quotient of $Q$ has order dividing
the $p$-part of $\gcd_E m_E$.  Its image $\overline P_b$ in
$Q$ is finite with this same bound: a profinite group whose
finite quotient orders are bounded is finite, as follows by
choosing a finite quotient of maximal order and comparing
all refinements.  In characteristic zero the groups $P_b$
are trivial, and no bound is needed.

There are only finitely many omitted base points.  Choose an
open normal subgroup $Q'\subset Q$ meeting all the finite
groups $\overline P_b$ trivially.  Let
$B_1^\circ\to B^\circ$ be the connected finite \'etale cover
corresponding to the inverse image of $Q'$ in $\pi_1(B^\circ)$,
and let $B_1$ be its smooth proper compactification.
At a point of $B_1\setminus B_1^\circ$, its wild inertia maps
into a conjugate of $P_b$ in $\pi_1(B^\circ)$: its image is
pro-$p$, whereas the quotient of the original inertia by
$P_b$ is pro-prime-to-$p$.  Its further image in $Q'$ is
therefore trivial.  Thus the surjection
\[
                    \pi_1(B_1^\circ)\longrightarrow Q'
\]
factors through $\pi_1^{\mathrm t}(B_1^\circ)$, which is
finitely generated by the same tame-curve theorem.
Consequently $Q'$ and its finite extension $Q$ are finitely
generated.  Since $N$ is finitely generated, so is $G$.
\end{proof}

\begin{theorem}[Proved: every fixed boundary on a proper surface mapping to a curve]
\label{D:surface-map-all-boundaries}
Let $S$ be a normal connected proper surface over an
algebraically closed field.  If $S$ admits a nonconstant
morphism to a proper curve, then $Q(S,Y)$ is topologically
finitely generated for every proper closed subset $Y\subset S$.
In particular, arbitrary horizontal and vertical original
boundary divisors are allowed, with no separability assumption
on the horizontal divisors and no smoothness assumption on
the original generic fibre.
\end{theorem}
\begin{proof}
Take Stein factorization of the nonconstant morphism to its
image.  As in Theorem~\ref{D:surface-map-to-curve}, we obtain
$f:S\to B$ with $B$ a smooth connected proper curve,
$f_*\mathcal O_S=\mathcal O_B$, and $F=k(B)$ algebraically
closed in $K=k(S)$.  The p-rank-one argument in that theorem
shows that $K/F$ is regular.  Hence the normal proper generic
curve $C=S_\eta$ is geometrically integral, although it need
not be smooth.

In positive characteristic, Lemma 33.27.4, tag 0BXT, gives
a finite purely inseparable extension $F_1/F$ for which the
normalization of $C_{F_1}$ is geometrically normal
\cite{C_StacksLogFibration}.  A geometrically normal curve
is smooth: after algebraic closure it is a normal curve,
hence regular and smooth by Lemma 33.43.8, tag 0B8Y, and
smoothness descends under field extensions.  Let $B_1$ be the
normalization of $B$ in $F_1$, and let $S_1$ be the normalization
of $S$ in $KF_1$.  Normalizations here are finite, since all
the schemes are of finite type over a field.  The induced
morphism $S_1\to B_1$ is proper, and its generic curve is
the smooth geometrically integral normalization just chosen.
In characteristic zero we take $F_1=F$, $B_1=B$, and $S_1=S$;
the normal generic curve is already smooth.

Let $E_1,\ldots,E_s$ be the residue fields of the finitely many
closed points cut out by the reduced horizontal boundary
on the generic curve of $S_1/B_1$.  In positive characteristic
choose $r$ sufficiently large that
\[
        E_iF_2/F_2\text{ is separable for every }i,
        \qquad F_2=F_1^{1/p^r}.
\]
This is a finite purely inseparable extension, since $F_1$
has transcendence degree one over the perfect field $k$.
To verify the assertion about the $E_i$, let $E_i^{\mathrm{sep}}$
be the maximal separable subextension of $E_i/F_1$ and choose
$r$ with $E_i^{p^r}\subset E_i^{\mathrm{sep}}$ for all $i$.
Then $E_iF_2$ is a subfield of
$(E_i^{\mathrm{sep}})^{1/p^r}$, which is separable over $F_2$.
In characteristic zero no further extension is needed.

Normalize $B_1$ in $F_2$ and $S_1$ in $KF_2$, obtaining
$B_2$ and $S_2$.  Since the generic curve of $S_1/B_1$ is
already smooth, its base change to $F_2$ is normal; the second
normalization does not change that generic curve.  Its reduced
horizontal boundary therefore has exactly the separable
residue fields just described.  Shrinking a nonempty open
$B_2^\circ\subset B_2$, we may arrange that $S_2\to B_2$ is
a smooth proper geometrically connected curve family there,
that all vertical boundary lies outside it, and that the
reduced horizontal boundary is a finite \'etale relative
divisor.  For the last assertion, finite separability gives
an \'etale dense open on each horizontal boundary curve;
properness makes the horizontal map finite after shrinking,
and we remove its finitely many singularities and pairwise
intersections.  In the smooth curve family the resulting
finite \'etale closed subscheme is a relative Cartier divisor.
Thus Proposition~\ref{D:log-fibration-criterion} applies to
$(S_2,Y_2)$, where $Y_2$ is the inverse image of $Y$.

We finish by checking descent to the original fixed-boundary
quotient.  The composite $S_2\to S$ is finite and radicial.
Indeed, $KF_2/K$ is purely inseparable of some bounded
exponent $q$, and on each normal affine open $\operatorname{Spec}A$
of $S$, every element $a$ of its integral closure in $KF_2$
satisfies $a^q\in K$.  This power is integral over $A$ and
therefore belongs to $A$.  Thus the integral map is universally
injective, and lying over makes it surjective.  It is a
universal homeomorphism.  The same holds for
$U_2=S_2\setminus Y_2\to U=S\setminus Y$.
Topological invariance of the \'etale site gives
\[
                         \pi_1(U_2)\simeq\pi_1(U)
\]
\cite[Theorem 59.45.2, tag 04DY]{C_StacksLogFibration}.
Each codimension-one point of $Y_2$ lies above a
codimension-one point of $Y$, because the finite universal
homeomorphism preserves chains of specializations.  Both
local rings are DVRs.  Tame finite extensions of DVRs remain
tame after arbitrary extension of the base DVR
\cite[Lemma 15.116.9, tag 0EXY; Lemma 58.31.1, tag 0EYD]
{C_StacksLogFibration}.
Consequently every cover counted by $Q(S,Y)$ pulls back to
a cover counted by $Q(S_2,Y_2)$.  Under the displayed
identification this gives a surjection
\[
                         Q(S_2,Y_2)\longrightarrow Q(S,Y).
\]
The source is finitely generated by the proposition, so the
target is finitely generated as well.
\end{proof}

\begin{remark}
The theorem uses a morphism from the given proper surface to
a curve.  Replacing it by a rational pencil on a projective
resolution does not establish the unrestricted surface case:
an exceptional curve over an omitted point may be horizontal
for that pencil.  The original fixed-boundary condition puts
no tameness requirement on such a new horizontal valuation,
so the logarithmic homotopy sequence above does not provide
a finitely generated fibre image for that larger quotient.
\end{remark}

\subsubsection{Generic-fibre generation in arbitrary dimension}

\begin{proposition}[Proved: generation from the geometric generic fibre]
\label{D:generic-fibre-generation}
Let $k$ be algebraically closed, let $X$ be a normal connected
proper $k$-scheme, and let $f:X\to B$ be a dominant morphism
to a smooth connected proper curve.  Suppose the geometric
generic fibre $X_{\overline\eta}$ is integral.  Let
$Y\subsetneq X$ be closed and put $U=X\setminus Y$.
Then the following assertions hold.
\begin{enumerate}
\item If $\pi_1(U_{\overline\eta})$ is topologically finitely
generated and $U$ meets every fibre of $f$, then $\pi_1(U)$
is topologically finitely generated.
\item If $X_{\overline\eta}$ is normal and
$Q(X_{\overline\eta},Y_{\overline\eta})$ is topologically
finitely generated, then $Q(X,Y)$ is topologically finitely
generated.  In this assertion whole fibres may be contained
in $Y$.
\end{enumerate}
No properness assumption is imposed on $U\to B$, and the
generic fibre need not be smooth.
\end{proposition}
\begin{proof}
Put $F=k(B)$.  The arithmetic fundamental-group sequence
over $F$ is exact:
\[
 1\longrightarrow\pi_1(U_{\overline\eta})
  \longrightarrow\pi_1(U_\eta)
  \longrightarrow\operatorname{Gal}(F^{\mathrm{sep}}/F)
  \longrightarrow1;
\]
see \cite[Lemma 58.14.3, tag 0BTX]{C_StacksGenericFibre}.
Moreover, $\pi_1(U_\eta)\to\pi_1(U)$ is surjective.
Indeed, every connected finite \'etale cover of the normal
integral $U$ is integral, and its nonempty localization at
the generic point of the base remains integral.
Consequently the image of $\pi_1(U_{\overline\eta})$ in
\emph{any} quotient $G$ of $\pi_1(U)$ is closed and normal.
Denote this image by $N$, and write $Q=G/N$.  There is a
surjection $\operatorname{Gal}(F^{\mathrm{sep}}/F)\to Q$.

For assertion (1) we take $G=\pi_1(U)$; its subgroup $N$
is finitely generated by hypothesis.  For assertion (2)
we take $G=Q(X,Y)$.  In this case the map
$\pi_1(U_{\overline\eta})\to G$ factors through
$Q(X_{\overline\eta},Y_{\overline\eta})$, so $N$ is again
finitely generated.  To justify this factorization, every
codimension-one boundary point of $X_{\overline\eta}$ lies
above a horizontal codimension-one boundary point of $X$.
This follows from dimension, since the geometric base-field
extension is algebraic.  Both local rings are DVRs, the
first by the assumed geometric normality and the second
by normality of $X$.  Tameness of an original boundary
valuation persists under this extension of DVRs, by
\cite[Lemma 15.116.9, tag 0EXY]{C_StacksLogFibration}.
Thus every finite cover counted by $G$ restricts to a
cover counted by the indicated geometric fixed-boundary
group.

We next show that $Q$ comes from one fixed open curve.
The morphism $U\to B$ has a nonempty smooth locus dominating
$B$.  Indeed its generic fibre is geometrically integral,
hence contains a dense smooth open by
\cite[Lemma 33.25.7, tag 056V]{C_StacksGenericFibre}.
The smooth locus is open and a smooth morphism is open.
Its image therefore contains a nonempty open
$B^\circ\subset B$.  In particular, for every closed
$b\in B^\circ$, there is a retained fibre component of
multiplicity one: take a component through an $f$-smooth
point of $U$ above $b$.

Let $Q_0$ be a finite quotient of $Q$.  The corresponding
connected cover $V\to U$ becomes over the generic fibre
the pullback of a connected finite Galois extension
$F'/F$ with group $Q_0$, by the arithmetic exact sequence.
Thus $k(V)=k(X)F'$.  Normalize $B$ in $F'$ to get
$B'\to B$.  If $E$ is a component of the fibre over $b$
with multiplicity $m_E$, and $e_E$ is the ramification
index of any prolongation of its valuation to $k(V)$,
multiplication of value-group indices gives
\[
                         e(b'/b)\mid m_Ee_E.
\tag{D.GF.1}
\]
At a retained component $e_E=1$, because $V\to U$ is
\'etale.  Therefore $e(b'/b)=1$ for $b\in B^\circ$.
The residue field at every closed point of $B'$ is $k$,
so $B'\to B$ is unramified, hence \'etale, over $B^\circ$.
This holds for every finite quotient of $Q$.  Consequently
the surjection from the absolute Galois group factors as
\[
                         \pi_1(B^\circ)\longrightarrow Q.
\]

In assertion (1), $U$ meets every fibre.  Hence above each
of the finitely many omitted $b$, some fibre component
has its generic point in $U$.  The index (D.GF.1) implies
that the order of the image of the full inertia at $b$
in every finite quotient of $Q$ divides that component's
multiplicity.  Thus the image in $Q$ of the full inertia
is finite.  Choose an open normal subgroup $Q'\subset Q$
meeting these finitely many finite inertia images trivially.
The inverse image of $Q'$ corresponds to a finite \'etale
cover $B_1^\circ\to B^\circ$.  After its smooth proper
compactification $B_1$, the map
$\pi_1(B_1^\circ)\to Q'$ kills all boundary inertia and
factors through $\pi_1(B_1)$.  This group is finitely
generated, so $Q'$, $Q$, and then $G$ are finitely generated.

In assertion (2), a fibre component is either retained,
in which case $e_E=1$, or is an original boundary divisor,
in which case $e_E$ is prime to the positive characteristic
$p$.  Equation (D.GF.1) therefore bounds the $p$-part of
$e(b'/b)$ by the $p$-part of $m_E$ at every component.
The wild inertia image at each omitted $b$ in $Q$ is
finite.  In characteristic zero it is trivial.  Choose
an open normal subgroup $Q'$ avoiding these finite wild
images.  For the corresponding open curve $B_1^\circ$,
every wild inertia group maps into a conjugate of an
original wild inertia group and hence trivially to $Q'$.
Thus $Q'$ is a quotient of $\pi_1^{\mathrm t}(B_1^\circ)$,
which is finitely generated.  This proves finite generation
of $Q$ and $G$ in assertion (2).

The finiteness assertion about a profinite inertia image
uses only its bounded finite quotient orders, as in
Proposition~\ref{D:vertical-fibration-multiplicity}; the
existence of one open normal subgroup avoiding finitely
many finite images follows by separating their nonidentity
elements in finite quotients.  The proper and tame curve
generation theorems used in the last two paragraphs are
\cite[X, Theorem 2.6; XIII, Corollary 2.12]{SGA1}.
\end{proof}

\begin{corollary}[Proved: arbitrary-dimensional vertical boundary]
\label{D:vertical-boundary-all-dimensions}
Let $X$ be a normal connected proper scheme over an
algebraically closed field and let $f:X\to B$ be a
nonconstant morphism to a proper curve.  Let
$Y\subsetneq X$ be a closed subset with $f(Y)$ finite.
Then $Q(X,Y)$ is topologically finitely generated.
If, after Stein factorization of $f$, the open $X\setminus Y$
meets every fibre, its full fundamental group is
topologically finitely generated.  The latter condition
is automatic if $Y$ has codimension at least two in $X$.
\end{corollary}
\begin{proof}
Take Stein factorization, so the new base is a smooth
connected proper curve, $f_*\mathcal O_X=\mathcal O_B$,
and $F=k(B)$ is algebraically closed in $k(X)$.
The p-rank-one argument in
Theorem~\ref{D:surface-map-to-curve} is independent of
$\dim X$: $[F:F^p]=p$, and a non-$p$th-power of $F$
cannot acquire a $p$th root in $k(X)$ without contradicting
relative algebraic closedness.  Thus $k(X)/F$ is regular
and the geometric generic fibre is integral.

Since the boundary is vertical, $U_{\overline\eta}=
X_{\overline\eta}$ is proper and connected.  Its fundamental
group is finitely generated by
\cite[X, Theorem 2.9, pp. 273--274]{SGA1}, with no
normality or smoothness requirement.  For the full group
under the stated fibre-intersection hypothesis, assertion
(1) of Proposition~\ref{D:generic-fibre-generation} applies.
For $Q(X,Y)$ the same proof applies verbatim: its geometric
generic image is a quotient of this proper fundamental
group, and all its finite covers are tame at every deleted
fibre component.  Consequently (D.GF.1) bounds the wild
base inertia and proves finite generation without imposing
geometric normality on the generic fibre.

Finally, every fibre of a dominant morphism from a normal
integral scheme to a smooth curve has a codimension-one
component: a local uniformizer cuts a nonzero principal
divisor.  A subset of codimension at least two cannot
contain such a component.  Thus its complement meets
every fibre.
\end{proof}

\begin{corollary}[Proved: projective normalization of the geometric generic fibre]
\label{D:projective-generic-normalization}
Let $X$ be a normal connected proper scheme over an
algebraically closed field.  Suppose $X$ admits a
nonconstant morphism to a proper curve such that, after
Stein factorization, the normalization of its geometric
generic fibre is projective.  Then $Q(X,Y)$ is
topologically finitely generated for every proper closed
subset $Y\subset X$.  In particular, when $Y$ has
codimension at least two, the full group
$\pi_1(X\setminus Y)$ is topologically finitely generated.
\end{corollary}
\begin{proof}
The same Stein and p-rank-one arguments make the generic
fibre geometrically integral.  By
\cite[Lemma 33.27.4, tag 0BXT]{C_StacksLogFibration}, after
a finite purely inseparable extension of the base function
field its normalization is geometrically normal.  Normalize
the base curve in this extension and normalize $X$ in the
compositum of function fields, obtaining $X'\to B'$.
As in Theorem~\ref{D:surface-map-all-boundaries}, the map
$X'\to X$ is finite and radicial: a fixed $p$-power of
every integral element belongs to the original normal
ring.  In characteristic zero this step is unnecessary.

The geometric generic fibre of $X'$ is the normalization
of the original geometric generic fibre.  Indeed, its
finite birational map to that fibre has normal source,
so it has the defining property of normalization.  It is
therefore projective by hypothesis.  Proposition
\ref{D:fixed-boundary-positive} gives finite generation
of its fixed-boundary tame group for the inverse image
of $Y$.  Assertion (2) of
Proposition~\ref{D:generic-fibre-generation} proves that
$Q(X',Y')$ is finitely generated.

The finite radicial map preserves the \'etale fundamental
group of the open, and each new codimension-one boundary
DVR lies over an original one.  Tame DVR base change
therefore gives a surjection $Q(X',Y')\to Q(X,Y)$,
exactly as in Theorem~\ref{D:surface-map-all-boundaries}.
This proves the assertion.  If the original boundary
has codimension at least two, no wild inertia is killed
in defining $Q(X,Y)$, so it equals the full open
fundamental group.
\end{proof}

\subsubsection{Fibrations over a projective base}

\begin{lemma}[Proved: a uniform bound for residual inseparability]
\label{D:residual-inseparability-bound}
Let $M/F$ be a finitely generated extension of fields of
characteristic $p>0$.  There is a finite purely inseparable
extension $F_0/F$, depending only on $M/F$, with the following
property: for every finite extension $L/F$ in a common
overfield for which $LM/M$ is separable, we have
\[
                 [L:F]_{\mathrm{insep}}\mid [F_0:F].
\]
We may choose any $F_0$ for which $MF_0/F_0$ is separable.
In characteristic zero the corresponding bound is one.
\end{lemma}
\begin{proof}
Choose an integral algebraic $F$-scheme with function field
$M$.  Lemma 33.27.4, tag 0BXT, supplies a finite purely
inseparable $F_0/F$ for which its normalization after base
change is geometrically normal
\cite{C_StacksLogFibration}.  Its function field $MF_0$
is therefore separable over $F_0$.

Put $M_0=MF_0$ and $L_0=LF_0$.  Since $LM/M$ is separable,
its base change $L_0M_0/M_0$ is separable.  A separable
extension $M_0/F_0$ cannot make an inseparable finite
algebraic extension separable.  To verify this assertion,
let $L_0^{\mathrm{sep}}$ be the maximal separable
subextension of $L_0/F_0$.  Then
$L_0M_0/L_0^{\mathrm{sep}}M_0$ is both purely inseparable
and separable, so is trivial.  Thus
$L_0\subset L_0^{\mathrm{sep}}M_0$.  The latter is separable
over $L_0^{\mathrm{sep}}$, so its purely inseparable finite
subextension $L_0/L_0^{\mathrm{sep}}$ is trivial.
Consequently $LF_0/F_0$ is separable.

Now let $L^{\mathrm{sep}}$ be the maximal separable
subextension of $L/F$.  The extension
$LF_0/L^{\mathrm{sep}}F_0$ is purely inseparable and,
by the preceding paragraph, separable, so it too is
trivial.  Hence $L\subset L^{\mathrm{sep}}F_0$ and
\[
 [L:F]_{\mathrm{insep}}
  =[L:L^{\mathrm{sep}}]
  \leq [L^{\mathrm{sep}}F_0:L^{\mathrm{sep}}]
  \leq [F_0:F].
\]
Both sides of the asserted divisibility are powers of
$p$, so this inequality proves it.
\end{proof}

\begin{proposition}[Proved: generic-fibre generation over a projective base]
\label{D:projective-base-generation}
Let $k$ be algebraically closed, let $f:X\to B$ be a
dominant morphism of normal connected proper $k$-schemes,
and suppose $B$ is projective and the geometric generic
fibre $X_{\overline\eta}$ is integral.  Let
$Y\subsetneq X$ be closed and put $U=X\setminus Y$.
Then $Q(X,Y)$ is topologically finitely generated if either
of the following conditions holds:
\begin{enumerate}
\item $\pi_1(U_{\overline\eta})$ is topologically finitely
generated;
\item $X_{\overline\eta}$ is normal and
$Q(X_{\overline\eta},Y_{\overline\eta})$ is topologically
finitely generated.
\end{enumerate}
The dimensions of $X$, $B$, and the fibres are unrestricted.
\end{proposition}
\begin{proof}
Write $F=k(B)$, $K=k(X)$, $G=Q(X,Y)$, and let $N$ be
the image of $\pi_1(U_{\overline\eta})$ in $G$.
The arithmetic fundamental-group exact sequence and
normality of $U$ show that $N$ is closed and normal,
and that $Q=G/N$ is a quotient of
$\operatorname{Gal}(F^{\mathrm{sep}}/F)$.
The proof is exactly the first paragraph of
Proposition~\ref{D:generic-fibre-generation}, using
\cite[Lemma 58.14.3, tag 0BTX]{C_StacksGenericFibre}.
In case (1), $N$ is finitely generated by hypothesis.
In case (2), restrictions of original tame covers to
the geometric generic fibre remain tame at its
codimension-one boundary DVRs, by
\cite[Lemma 15.116.9, tag 0EXY]{C_StacksLogFibration}.
Thus $N$ is a quotient of the finitely generated group
in (2), and is again finitely generated.

We obtain one base open that works for \emph{every}
finite quotient of $Q$.  The smooth locus of $U\to B$
dominates $B$, because the geometrically integral generic
fibre contains a dense smooth open
\cite[Lemma 33.25.7, tag 056V]{C_StacksGenericFibre}.
Let $B^\circ$ be its open image and put
$T=U_{\mathrm{sm}}\times_B B^\circ$.
Then $T\to B^\circ$ is smooth, surjective and
quasi-compact, hence fpqc.

For a finite quotient $Q_0$ of $Q$, its connected cover
$V\to U$ is generically the pullback of a finite Galois
extension $F'/F$, with group $Q_0$.  Let $B'$ be the
normalization of $B$ in $F'$.  Normalization commutes
with smooth base change
\cite[Lemma 37.19.2, tag 03GV]{C_StacksProjectiveBase}.
Consequently $B'\times_B T$ is the normalization of
$T$ in its generic cover, and is therefore $V\times_U T$.
The latter is finite \'etale over $T$.  \'Etaleness is
fpqc local on the target
\cite[Lemma 35.23.31, tag 02VN]{C_StacksProjectiveBase},
so $B'\to B$ is finite \'etale over $B^\circ$.
This proves, including at codimension-two points, that
the surjection from the absolute Galois group factors as
\[
                         \pi_1(B^\circ)\longrightarrow Q.
\]

There are finitely many prime divisors $D$ in
$B\setminus B^\circ$.  For each one choose a prime divisor
$E$ of $X$ dominating $D$.  Such a divisor exists by
properness and dominance: over the DVR at the generic
point of $D$, a nonzero uniformizer cuts a nonempty
principal divisor on the integral source.  Write $m_E$
for its positive valuation at $E$, and set
\[
                        F_D=\kappa(D),\qquad M_E=\kappa(E).
\]
The extension $M_E/F_D$ is finitely generated.  In positive
characteristic fix a finite purely inseparable field
extension $F_{D,0}/F_D$ from
Lemma~\ref{D:residual-inseparability-bound}, and put
$c_D=[F_{D,0}:F_D]$.  These choices are made once, before
any finite quotient is chosen.

Return to a finite quotient $Q_0$, with fields $F'$ and
$k(V)=KF'$.  Choose a prolongation of the $E$-valuation
to $KF'$, and restrict it to a valuation $D'$ of $F'$.
Let $e_E$ be its ramification index over $E$ and $e_D$
the index at $D'/D$.  Multiplication of value-group
indices gives
\[
                         e_D\mid m_Ee_E.
\tag{D.PB.1}
\]
If $E$ is retained in $U$, the cover is \'etale there;
otherwise $E$ is an original boundary divisor, and the
cover is tame there.  In either case $e_E$ is prime to
$p$ and the residual extension above $E$ is separable.
The latter contains the compositum of $M_E$ and
$L=\kappa(D')$.  Hence $LM_E/M_E$ is separable, and
Lemma~\ref{D:residual-inseparability-bound} gives
\[
       [\kappa(D'):F_D]_{\mathrm{insep}}\mid c_D,
       \qquad (e_D)_p\mid (m_E)_p.
\tag{D.PB.2}
\]

For clarity, the second bound alone would not control
fierce inertia.  For a finite Galois extension of these
divisorial DVRs, the full inertia order is
\[
        |I_{D'}|=e_D[\kappa(D'):F_D]_{\mathrm{insep}},
\]
and its wild subgroup has order
\[
        |P_{D'}|=(e_D)_p
                     [\kappa(D'):F_D]_{\mathrm{insep}}.
\]
These formulas include inseparable residue extensions.
The integral closures are finite; the degree formula and
the decomposition action on the normal residue extension
therefore give $|I|=e f_{\mathrm{insep}}$ without a defect
term.  The wild subgroup is a $p$-group and the tame
quotient has order the prime-to-$p$ part of $e$, by
\cite[Lemma 15.114.5, tag 09EE]{C_StacksProjectiveBase}.
Thus (D.PB.2) bounds the image of the absolute wild
inertia at $D$ in every finite quotient of $Q$ by the
fixed integer $(m_E)_p c_D$.  Its image $\overline P_D$
in $Q$ is finite.  This conclusion uses the bounded
finite quotient orders and the surjectivity of wild
inertia onto finite-level wild inertia; the latter is
included in \cite[Lemma 58.13.5, tag 0BUA]
{C_StacksProjectiveBase}.  In characteristic zero all
wild inertia groups are trivial.

Choose an open normal subgroup $Q'\subset Q$ meeting all
the finitely many finite groups $\overline P_D$ trivially.
Its inverse image in $\pi_1(B^\circ)$ defines a connected
finite \'etale cover $B_1^\circ\to B^\circ$.
The normalization $B_1$ of $B$ in its function field is
normal and projective, since it is finite over $B$.
Every codimension-one component of
$B_1\setminus B_1^\circ$ lies over a codimension-one
component of $B\setminus B^\circ$.  Its wild inertia
maps into a conjugate of the corresponding original
wild inertia: it is pro-$p$, maps into original inertia,
and that inertia modulo wild inertia is pro-prime-to-$p$.
Its image in $Q'$ is therefore trivial.  Consequently
\[
       \pi_1(B_1^\circ)\longrightarrow Q'
\]
factors through $Q(B_1,B_1\setminus B_1^\circ)$.
This last group is finitely generated by
Proposition~\ref{D:fixed-boundary-positive}, because
$B_1$ is projective.  Hence $Q'$, its finite extension
$Q$, and finally $G$ are finitely generated.
\end{proof}

\begin{corollary}[Proved: vertical boundary and low-dimensional generic fibres over projective images]
\label{D:projective-image-corollaries}
Let $X$ be a normal connected proper scheme over an
algebraically closed field, and suppose $X$ admits a
nonconstant morphism to a projective scheme.
Replace the target by its image and then by the normal
projective Stein factor $B$.
\begin{enumerate}
\item If $Y\subsetneq X$ is closed and its image is a
proper subset of $B$, then $Q(X,Y)$ is topologically
finitely generated.
\item If the normalization of the reduced geometric
generic fibre is projective, then $Q(X,Y)$ is
topologically finitely generated for every proper
closed subset $Y\subset X$.
\end{enumerate}
In particular (2) applies whenever the generic fibre
has dimension at most one.  In either assertion, if
$Y$ has codimension at least two, the group in question
is the full $\pi_1(X\setminus Y)$.
\end{corollary}
\begin{proof}
The function field $F=k(B)$ is algebraically closed in
$K=k(X)$, by Stein factorization and normality of $X$.
Consequently the geometric generic fibre is irreducible;
it need not be reduced when $\dim B>1$.
By \cite[Lemma 33.27.4, tag 0BXT]{C_StacksLogFibration},
choose a finite purely inseparable extension $F_0/F$
such that the normalization of the generic fibre after
base change is geometrically normal.  Normalize $B$
in $F_0$ and $X$ in $KF_0$ to obtain $X_0\to B_0$.
Then $B_0$ remains normal and projective, and $X_0\to X$
is finite and radicial, by the fixed-power integral
argument in Theorem~\ref{D:surface-map-all-boundaries}.
In characteristic zero we may take the trivial extension.

The geometric generic fibre of $X_0$ is the normalization
of the reduced original geometric generic fibre: the
finite birational map has normal source and hence is
its normalization.  It is geometrically integral.
In case (1) the geometric generic boundary is empty,
and this fibre is proper and connected.  Its full
fundamental group is finitely generated by
\cite[X, Theorem 2.9, pp. 273--274]{SGA1}.
In case (2) its fixed-boundary group is finitely generated
by the projective case of
Proposition~\ref{D:fixed-boundary-positive}.
The respective alternatives of
Proposition~\ref{D:projective-base-generation} therefore
give finite generation of $Q(X_0,Y_0)$.

Topological invariance of the \'etale site and tame
base change at the corresponding boundary DVRs give
a surjection $Q(X_0,Y_0)\to Q(X,Y)$, as proved in
Theorem~\ref{D:surface-map-all-boundaries}.  This proves
both assertions.  A normal proper curve over an
algebraically closed field is projective, as is a
zero-dimensional proper scheme, so the indicated
generic-fibre dimension bound implies (2).
Finally, in codimension at least two there are no
original boundary inertia groups to quotient out.
\end{proof}

\subsubsection{Finite normal covers of the base and iterated fibrations}

For a normal connected proper scheme $B$ over an algebraically
closed field, consider the following property:
\[
\begin{gathered}
\mathcal H(B):\quad\text{for every finite dominant morphism
$B'\to B$ with $B'$ normal and connected,}\\
\text{and every proper closed subset $T\subset B'$, the group
$Q(B',T)$ is topologically finitely generated.}
\end{gathered}
\]
The finite morphisms in this definition need not be separable.
Every normal projective $B$ satisfies $\mathcal H(B)$, because
every such $B'$ is projective.  Moreover $\mathcal H(B)$ implies
$\mathcal H(B_0)$ for every finite dominant normal $B_0\to B$,
by composition of finite morphisms.

\begin{proposition}[Proved: the exact base condition in the fibration argument]
\label{D:hereditary-base-generation}
In Proposition~\ref{D:projective-base-generation}, replace the
assumption that $B$ is projective by $\mathcal H(B)$.
Then the same conclusion holds, under either of its two
geometric generic-fibre hypotheses.
\end{proposition}
\begin{proof}
We specify the only point where the proof needs a property
of a projective base.  All arguments preceding the choice of
$Q'\subset Q$ use normality, properness, the fixed smooth
surjective base open, and the two valuation bounds.  They
therefore apply unchanged.  In particular the full wild
image at each of the finitely many base boundary divisors
is finite, including its inseparable residual contribution.

The normalization $B_1$ in the finite cover corresponding to
$Q'$ is finite and dominant over $B$, and normal and proper.
The same inertia intersection argument gives a surjection
\[
       Q(B_1,B_1\setminus B_1^\circ)\longrightarrow Q'.
\]
Property $\mathcal H(B)$ makes the source finitely generated.
The finite-extension argument for $Q$, and then the extension
by the finitely generated geometric image, proves the assertion.
No projectivity of $B_1$ is needed at this final step.
\end{proof}

\begin{theorem}[Proved: iterated fibrations with projective generic normalizations]
\label{D:iterated-fibration-generation}
Let $k$ be algebraically closed and suppose that
\[
 X_r\longrightarrow X_{r-1}\longrightarrow\cdots
       \longrightarrow X_0
\]
is a sequence of dominant morphisms of normal connected
proper $k$-schemes.  Suppose $X_0$ is projective and, for
each morphism in the sequence, the normalization of the
reduced geometric generic fibre is projective.  This
normalization may have more than one connected component.
Then $\mathcal H(X_i)$ holds for every $i$.

In particular the conclusion holds if every morphism has
generic fibre dimension at most one.  Thus for every finite
normal connected cover $Z\to X_r$ and every proper closed
subset $Y\subset Z$, the group $Q(Z,Y)$ is topologically
finitely generated.  If $Y$ has codimension at least two,
this is the full group $\pi_1(Z\setminus Y)$.
\end{theorem}
\begin{proof}
We prove the assertion by induction.  The projective case
gives $\mathcal H(X_0)$.  Suppose $\mathcal H(X_{i-1})$
has been proved, and fix a finite dominant morphism
$Z\to X_i$ with $Z$ normal and connected.  Stein factorization
of $Z\to X_{i-1}$ gives
\[
                 Z\longrightarrow B'\longrightarrow X_{i-1},
\]
where $B'$ is normal and the second morphism is finite.
The field $k(B')$ is algebraically closed in $k(Z)$.
Consequently the generic fibre is geometrically irreducible;
it need not be geometrically reduced.  This is the
relative-constant-field criterion already used in
Corollary~\ref{D:projective-image-corollaries}.

We verify that finite covers preserve the required generic
projectivity.  Over an algebraic closure of $k(B')$, each
normal component of the reduced generic fibre of $Z$
is finite over the normalization of a component of the
reduced geometric generic fibre of $X_i\to X_{i-1}$.
Indeed the original map is finite.  After reduction and
normalization the dominant map factors through the normal
target component; on affine rings it remains finite, since
an algebra finite over a smaller ring is finite over any
intermediate ring acting on it.  All these normalizations
are finite, as the schemes are of finite type over fields.
The target normal component is projective by hypothesis,
so the source normal component is projective as well.

Choose a finite purely inseparable extension $F_0/k(B')$
after which normalization makes the generic fibre
geometrically normal, using
\cite[Lemma 33.27.4, tag 0BXT]{C_StacksLogFibration}.
Normalize $B'$ in $F_0$ and $Z$ in $k(Z)F_0$, obtaining
$B_0$ and $Z_0$.  The first scheme is finite over $X_{i-1}$,
so $\mathcal H(B_0)$ follows from the induction hypothesis.
The second morphism $Z_0\to Z$ is finite and radicial:
a fixed $p$-power of every integral element of its function
field lies in $k(Z)$ and is integral over the normal local
ring downstairs, hence lies in that ring.  In characteristic
zero we take the trivial field extension.

The geometric generic fibre of $Z_0\to B_0$ is integral,
normal and projective by the preceding two paragraphs.
For $Y_0$ the inverse image of any proper closed $Y\subset Z$,
its fixed-boundary tame group is therefore finitely generated
by the projective case.  Proposition~\ref{D:hereditary-base-generation}
applied to $Z_0\to B_0$ gives finite generation of $Q(Z_0,Y_0)$.
Finite radicial base change preserves the category of finite
\'etale covers of the open and preserves tameness upon pullback
at the original boundary valuations, by
\cite[tags 04DY and 0EXY]{C_StacksLogFibration}.
It consequently gives a surjection
\[
                         Q(Z_0,Y_0)\longrightarrow Q(Z,Y).
\]
This proves $\mathcal H(X_i)$, including inseparable finite
covers, and completes the induction.

For the final sufficient condition, a normal proper scheme
over an algebraically closed field of dimension at most one
is projective, component by component, by
\cite[Lemma 33.43.4, section 0A22]{C_StacksLogFibration}.
Hence the required
generic projectivity is automatic for fibres of dimension
at most one.  If the boundary has codimension at least two,
there are no original divisorial wild inertia groups to kill,
so $Q(Z,Y)=\pi_1(Z\setminus Y)$.
\end{proof}

\begin{remark}
The sequence consists of morphisms defined on the entire
proper schemes.  Replacing a rational pencil by a morphism
on a modification does not give this hypothesis for the
original scheme: the additional horizontal exceptional
valuations can still have uncontrolled wild image.
\end{remark}

\subsubsection{Proper logarithmic and projective-target recognition}

\begin{lemma}[Proved: full Kummer \'etale finite generation on a proper log scheme]
\label{D:proper-log-finite-generation}
Let $k$ be algebraically closed and let $X^{\log}$ be a nonempty
connected fine saturated log scheme whose underlying scheme $X$
is proper over $k$.  Then its full Kummer \'etale fundamental
group is topologically finitely generated.
\end{lemma}
\begin{proof}
We first check that every finite Kummer \'etale cover is tame
relative to $k$ in the precise valuative sense of
\cite[Definition 3.2.4]{AHLS2026}.  Let $(K,V)$ be a strictly
henselian valued field over $k$, and let
$\operatorname{Spec}K\to X$ be a test point.  Properness extends
this map to $\operatorname{Spec}V\to X$.  Equip
$T=\operatorname{Spec}V$ with the pulled-back fine saturated
log structure.  It has a global fine saturated chart: an
\'etale chart near the closed point pulls back along a section
over the strictly henselian local scheme, and an open
neighborhood of its closed point is the entire scheme.

By \cite[Proposition 3.1.8]{AHLS2026}, the full Kummer \'etale
fundamental group of $T$ is a free module of finite rank over
$\widehat{\mathbb Z}'$, with the residue characteristic omitted.
In particular every finite Kummer \'etale cover of $T$ is a
prime-to-residue-characteristic object of its Galois category.
Pullback to $\operatorname{Spec}K$ preserves this property.
This is exactly the condition in the cited definition, for
every test valued field.  Thus
\[
 \operatorname{FEt}^{\mathrm{ket}}(X^{\log})
   =\operatorname{FEt}^{\mathrm{t,ket}}(X^{\log}/k).
\]
Every fine saturated log scheme is of type $V_d$ in the
terminology of that source.  Its Theorem 3.5.1 therefore makes
the group on the right finitely generated, and proves the
assertion for the full Kummer \'etale group.
\end{proof}

\begin{proposition}[Proved: fixed-boundary finite generation for proper log regular pairs]
\label{D:log-regular-boundary}
Let $X^{\log}$ be a connected fine saturated log regular scheme
whose underlying scheme $X$ is proper over an algebraically
closed field.  Let $U$ be its interior, where the log structure
is trivial, and put $Y=X\setminus U$.  Then the group $Q(X,Y)$,
defined using tameness only at the original codimension-one
points of $Y$, is topologically finitely generated.
Neither projectivity nor regularity of the underlying $X$ is
assumed.
\end{proposition}
\begin{proof}
The underlying scheme of a log regular scheme is normal, and
its interior is dense.  Hoshi's exact comparison
\cite[Proposition A.10, pp. 57--58]{Hoshi2006B} identifies the
finite Kummer \'etale covers of $X^{\log}$ with the finite
\'etale covers of $U$ whose normalizations are tame at every
$y\in Y$ with $\dim\mathcal O_{X,y}=1$.  This is precisely the
cover category defining $Q(X,Y)$; there is no additional
exceptional-valuation condition in that comparison.
Lemma~\ref{D:proper-log-finite-generation} now proves the claim.
\end{proof}

\begin{remark}
This is a recognition criterion for the fixed compactification.
For a singular pair without a log regular structure, the
comparison with its original codimension-one tame cover category
is not supplied by the quoted theorem.  Thus this result does
not close the general fixed-boundary gate.  In particular,
placing a logarithmic structure on an arbitrary singular
compactification does not establish the required log regularity.
\end{remark}

\begin{proposition}[Proved: a generically finite morphism to a projective target]
\label{D:generically-finite-projective-target}
Let $X$ be a normal connected proper scheme over an
algebraically closed field.  Suppose $X$ admits a morphism
to a projective scheme that is generically finite onto its
image.  Then $Q(X,Y)$ is topologically finitely generated
for every proper closed subset $Y\subset X$.
In particular, if $Y$ has codimension at least two, the
full group $\pi_1(X\setminus Y)$ is topologically finitely
generated.
\end{proposition}
\begin{proof}
Replace the projective target by the reduced image.
Stein factorization gives a proper birational morphism
\[
                              q:X\longrightarrow Z,
\]
where $Z$ is normal and finite over that image, hence
projective.  More explicitly, the generic fibre of the
original map is the spectrum of the finite extension
$k(X)$ of the image function field.  Its Stein factor
is the normalization of the image in $k(X)$, so the
remaining morphism is birational.  We use
\cite[Theorem 37.53.4, tag 03H0]{C_StacksFibration}.

Let $Z^\circ$ be the largest open over which $q$ is an
isomorphism.  It contains all codimension-zero and
codimension-one points of the normal $Z$, by
\cite[Lemma 33.17.3, tag 0BFP]{C_StacksProjectiveTarget}.
Put $E=Z\setminus Z^\circ$ and
\[
                  W=Z\setminus\bigl(E\cup q(Y)\bigr).
\]
The map $q$ is proper, so $q(Y)$ is closed.  It is a proper
subset of $Z$, since $q$ is birational and $Y$ is proper
in the integral $X$.  Thus $W$ is a nonempty open, and
$q^{-1}(W)\simeq W$ is a dense open of $U=X\setminus Y$.

Every connected finite \'etale cover of the normal
integral $U$ remains connected on this dense open.
Moreover, a codimension-one component of
$Z\setminus W$ cannot be contained in $E$, since $E$
has codimension at least two.  Its generic point is in
$Z^\circ$, and its corresponding divisor on $X$ lies
in $Y$.  The two discrete valuation rings are identical.
Consequently every cover counted by $Q(X,Y)$ restricts
to a cover of $W$ tame at every codimension-one component
of $Z\setminus W$.  We therefore obtain a surjection
\[
                 Q\bigl(Z,E\cup q(Y)\bigr)
                         \longrightarrow Q(X,Y).
\]
The source is finitely generated by the projective case
of Proposition~\ref{D:fixed-boundary-positive}.  This
proves the result.  The codimension-two conclusion follows
because there are then no original boundary inertia groups
to kill in defining $Q(X,Y)$.
\end{proof}

\begin{remark}
The hypothesis is a morphism from $X$ to a projective target.
A projective alteration mapping onto $X$ has the opposite
direction and does not supply this hypothesis.  The proof
uses only unchanged codimension-one DVRs on the target;
it imposes no tameness condition at exceptional divisors
of $X$ whose image has codimension at least two.
\end{remark}

\subsubsection{A wild cover above one low-degree positive-genus curve}

\begin{example}[Refuted: universal tameness above a single smooth exceptional curve]
\label{D:wild-genus2-single-curve}
Let $k$ be algebraically closed of characteristic two.  There is
a normal affine surface $S$, with only one singular point $o$,
whose resolution has exceptional fibre one smooth genus-two
curve $C$ with $C^2=-1$, and a connected finite Galois cover
$T\to S$ with group $\mathbb Z/2\mathbb Z$ which is \'etale
over $S\setminus\{o\}$ and wild at $C$.  At its unique
prolongation the ramification index is one, the residue
extension is purely inseparable of degree two, and the
different exponent is one.

We can take the following explicit hypersurfaces:
\[
 \begin{aligned}
 S&:\quad F(a,b,t)=b^2+(t^5+t^2a^2)b+a^5=0,\\
 T&:\quad H(c,d,t)=d^2+(c^4+t^3+t^4c^2)d
            +tc^5+t^2c^7+t^4c^6+t^6c^5=0,
 \end{aligned}
\]
with map
\[
                 a=c^2+t^2c,\qquad b=t^2d+c^5.
 \tag{M}
\]
In particular the hypothesis of formal linearity in
Proposition~\ref{D:graded-cone-tame} cannot be discarded,
even for a Gorenstein singularity with one smooth exceptional
curve.  This example is a ramification obstruction; it is not
a counterexample to either finite-generation question D.1 or
D.2(ii).
\end{example}

\begin{proof}
All polynomial calculations in this proof take place in
characteristic two.

\par\smallskip\noindent\textbf{The target and a resolution.}\quad
The polynomial $F$ is irreducible.  Indeed, over $k(t)(a)$ put
$h=t^5+t^2a^2$ and divide its quadratic equation by $h^2$.
A root would give an Artin--Schreier expression
$a^5/h^2=w^2+w$.  The left-hand side has a pole of order one
at $a=\infty$.  Every pole of a rational function $w^2+w$
has even order, a contradiction.  The derivatives are
\[
             F_a=a^4,\qquad F_b=t^2(a^2+t^3),
             \qquad F_t=bt^4.
\]
Their common zero on $S$ is only $o=(0,0,0)$.  Thus $S$ is
an integral hypersurface, regular in codimension one and
Cohen--Macaulay, hence normal by Serre's criterion
\cite[Lemma 10.157.4, tag 031S]{StacksBWildCurve}.

We give an explicit proper resolution.  On $k[a,b,t]$ use
weights $(2,5,1)$, let $I_m$ be the ideal spanned by monomials
of weight at least $m$, and form
\[
 \mathcal R=\bigoplus_{m\geq0}I_m s^m.
\]
This is a finite type graded $k[a,b,t]$-algebra: it is generated
by $as,as^2$, by $bs,\ldots,bs^5$, and by $ts$.
The morphism $\operatorname{Proj}\mathcal R\to\mathbb A^3$
is proper, by the relative-Proj theorem
\cite[Lemma 31.31.4, tag 07ZW]{StacksBWildCurve}, and is an
isomorphism away from the origin.  Let $Y$ be the strict
transform of $S$.  Its reduced exceptional fibre is the
weighted projective curve
\[
               b^2+t^5b+a^5=0
                    \quad\hbox{in }\mathbb P(2,5,1).
 \tag{E}
\]
We verify smoothness using charts, so no smoothness assertion
about an ambient weighted projective space is needed.

On the $t$ chart put $a=t^2x$ and $b=t^5y$.  Dividing $F$ by
$t^{10}$ gives
\[
                    y^2+(1+tx^2)y+x^5=0.
 \tag{Y0}
\]
This chart is smooth: its partial derivatives with respect to
$x,y$ are $x^4$ and $1+tx^2$, which never vanish together.
It contains all of the exceptional curve except its point at
infinity.  Near that point both weighted coordinates $a,b$
are nonzero.  The degree-zero ring of the weighted blowup
chart has coordinates
\[
 \rho=\frac b{a^2},\qquad
 \xi=\frac{a^5}{b^2}\in\mathcal O^\times,
 \qquad \tau=\frac{ta^2}{b}.
\]
The inverse monomial formulas are
\[
                  a=\rho^2\xi,
       \qquad b=\rho^5\xi^2,
       \qquad t=\rho\tau.
\]
Their exponent matrix is unimodular.  The weighted inequalities
give exactly the ring $k[\rho,\tau,\xi,\xi^{-1}]$ on this
ambient chart.  Dividing the substituted equation by
$\rho^{10}\xi^2$ gives the strict transform
\[
                  \xi^2(1+\xi+\rho\tau^2)+\tau^5=0.
 \tag{Yinf}
\]
Its derivative with respect to $\xi$ is $\xi^2$, a unit.
The unique exceptional point outside (Y0) is
$(\rho,\tau,\xi)=(0,0,1)$.  The two charts cover the
exceptional fibre, so $Y$ is smooth there.  Away from that
fibre it equals the already smooth $S\setminus\{o\}$.
Thus $Y\to S$ is a proper resolution with only the curve
$C$ of (E) exceptional.

For completeness, $C$ is the smooth proper completion of
$y^2+y=x^5$.  At infinity use $u=1/x$ and $v=y/x^3$; its
equation becomes $v^2+u^3v=u$, smooth at $(u,v)=(0,0)$.
The finite double cover $C\to\mathbb P^1_x$ has
\[
                 \pi_*\mathcal O_C
                    =\mathcal O\oplus\mathcal O(-3),
\]
as follows directly from the bases $1,y$ and $1,v$ and the
transition $y=x^3v$.  Hence $\chi(\mathcal O_C)=1-2=-1$
and $g(C)=2$.  It is integral because $x^5$ cannot be
$w^2+w$ in $k(x)$, by the same odd-pole argument.
On (Yinf), the divisor of the original function $t$ is
locally $\{\rho=0\}+\{\tau=0\}$.  The second curve meets
$C$ transversally once, and there are no other intersections
with the strict transform of $t=0$.  Intersecting the
principal divisor of $t$ with the proper curve $C$ yields
\[
                         0=C^2+1.
\]
Thus $C^2=-1$.

\par\smallskip\noindent\textbf{The connected quadratic extension.}\quad
On the $t$ chart the required extension of function fields is
\[
                            z^2+tz=x.
 \tag{AS}
\]
We verify its nontriviality, including the global gluing which
explains this choice.  Over $k[t]$, the two equations
\[
 y^2+(1+tx^2)y=x^5,
 \qquad v^2+(u^3+tu)v=u
\]
with $u=1/x$, $v=y/x^3$ define a smooth proper genus-two
curve family $\mathcal C\to\mathbb A^1_t$.  Smoothness on
the first chart was checked above; on the second chart the
only new point has $u=v=0$, where the $u$ derivative is one.
It is finite flat of degree two over $\mathbb P^1\times
\mathbb A^1$, so proper.  Again the direct-image module is
$\mathcal O\oplus\mathcal O(-3)$.

The \v Cech class
\[
                             h=\frac y{x^2}
\]
is a nonzero basis vector of $H^1(\mathcal C_t,\mathcal O)$,
including over $k(t)$: the two basis classes are $y/x$ and
$y/x^2$.  The equation of the family gives
\[
                         h^2+th=x+\frac y{x^4}.
\]
Consequently the two covers
\[
 z_0^2+tz_0=x,
       \qquad z_\infty^2+tz_\infty=uv
\]
glue by $z_\infty=z_0+h$.  They form a torsor under the
finite flat group scheme $\ker(F-t:\mathbb G_a\to\mathbb G_a)$.
On $t\ne0$ this is an \'etale group of order two, and the
torsor maps to the nonzero class $h$ under the map to
$H^1(\mathcal O)$ from its defining additive exact sequence.
Thus its generic fibre is nontrivial and connected.  This
proves that (AS) defines a separable quadratic field extension.

Put $c=tz$ and $d=(b+c^5)/t^2$ in that extension.  Direct
substitution gives
\[
          F(c^2+t^2c,t^2d+c^5,t)=t^4H(c,d,t).
 \tag{I}
\]
The algebra defining $T$ is finite over that of $S$: $c$
satisfies the monic equation $c^2+t^2c=a$, and $d$ satisfies
the displayed monic equation $H=0$ over $k[c,t]$.
On $t\ne0$ it is exactly the quadratic cover (AS).  Therefore
it is integral: its localization at $t$ is a domain and the
monic hypersurface equation has no component supported on
$t=0$.

The scheme $T$ is smooth away from the origin.  On $t\ne0$
this follows from the \'etale description over the smooth
target.  On $t=0$ its equation is $d(d+c^4)=0$ and its
$d$ derivative is $c^4$, so only $(c,d,t)=(0,0,0)$ can
be singular.  Being an integral hypersurface, $T$ is normal
by the same $(R_1)+(S_2)$ argument.  Thus it is the finite
normalization of $S$ in the quadratic extension.
The nontrivial deck transformation is explicitly
\[
       c\longmapsto c+t^2,
       \qquad d\longmapsto d+c^4+t^6c+t^8.
\]
It preserves (M) and $H$, and its square is the identity.

\par\smallskip\noindent\textbf{\'Etaleness away from the singular point.}\quad
On $t\ne0$ the derivative of $c^2+t^2c-a$ with respect to
$c$ is $t^2$, so the cover is \'etale.  If $t=0$ and the
image is not $o$, then $a\ne0$, and hence $c\ne0$.
Both surfaces are smooth at such points.  The target has
local \'etale coordinates $(b,t)$, since $F_a=a^4$ is a
unit.  The source has local \'etale coordinates $(c,t)$,
since $H_d=c^4$ at $t=0$.  In these coordinates
\[
                  db=c^4\,dc,
                         \qquad dt=dt
                    \quad\hbox{at }t=0.
\]
The differential is invertible.  Equivalently, the finite map
is flat by miracle flatness and its relative differentials
vanish; see \cite[Lemma 10.128.1, tag 00R4]{StacksBWildCurve}.
It is therefore \'etale at all these points.  The fibre above
$o$ is only the origin, so this proves the claimed
quasi-\'etale property.

\par\smallskip\noindent\textbf{The exceptional valuation.}\quad
At the generic point of $C$ in (Y0), $t$ is a uniformizer
and the residue field is
\[
                 k(C)=k(x,y),\qquad y^2+y=x^5.
\]
The element $x$ is not a square in $k(C)$: the separable map
to $k(x)$ implies $dx\ne0$ (or the displayed curve equation
gives the same differential calculation).  The reduction of
(AS) is therefore the irreducible purely inseparable equation
\[
                              \bar z^2=x.
\]
The algebra obtained by adjoining $z$ to the boundary DVR is
local, finite free of rank two, and has maximal ideal generated
by $t$.  It is consequently a regular DVR and is already the
integral closure.  We obtain
\[
                  e=1,
           \qquad [k(C)(\bar z):k(C)]_{\mathrm{insep}}=2.
\]
The deck action
$z\mapsto z+t$ is trivial on that residue field.  Thus all
of $\mathbb Z/2\mathbb Z$ is inertia and is wild.  The monogenic
different is generated by the derivative $t$ of (AS), so its
exponent is $\delta=1$; see
\cite[Example 49.12.5, tag 0BWB]{StacksBWildCurve}.

The target is Gorenstein because it is a hypersurface.
Adjunction gives $K_Y\cdot C=2g-2-C^2=3$; the discrepancy
of the sole exceptional curve is therefore $-3$.  In
particular the example lies outside the log-canonical case.
It also lies outside the formal-cone case, since that case
has tame exceptional ramification by
Proposition~\ref{D:graded-cone-tame}.  The inequality of
Proposition~\ref{D:single-curve-residue-tameness} is sharp
here: its two sides equal $2$ with $g=2,d=1,n=2,e=1$
and inseparable residue degree two.

The companion program \texttt{check\_wild\_genus2\_curve.py}
performs 26 exact assertions over $\mathbb F_2$, including
(I), the two Jacobians, the involution, both resolution
charts, and the \v Cech identities.  Its recorded runtime
is about $0.15$ seconds.  The proofs of properness, genus,
connectedness, and ramification above are separate from
these finite polynomial certificates.
\end{proof}

% Checked companion to the explicit fierce cover; root, C and D independent audits passed.
\begin{proposition}[Proved: the source resolution of the explicit fierce double cover]
\label{D:wild-genus2-source}
Let $T\to S$ be the cover in Example~\ref{D:wild-genus2-single-curve}.
The surface $T$ has a projective resolution $V\to T$ whose
exceptional fibre is one smooth genus-two curve $C'$, with
$(C')^2=-2$.  There is a finite morphism $V\to Y$ to the
resolution of $S$ in that example.  Its restriction $C'\to C$
is the relative Frobenius, after the displayed identifications
of both curves with $w^2+w=z^5$.  The discrepancy of $C'$
over the Gorenstein surface $T$ is $-2$.
\end{proposition}
\begin{proof}
Use weights $(1,3,1)$ on $(c,d,t)$ and the weighted Rees
construction of Example~\ref{D:wild-genus2-single-curve}.
The algebra is generated by $cs,ds,ds^2,ds^3,ts$ over
$k[c,d,t]$, and its relative Proj is proper and is an
isomorphism away from the origin.  It is projective: a
sufficiently divisible Veronese is generated in degree one,
so the same relative Proj is a closed subscheme of a
projective space over the affine base; see
\cite[Lemma 10.56.2, tag 0EGH]{StacksBWildSource}.
Let $V$ be the strict
transform of $T$.  Its reduced exceptional curve is
\[
                 d^2+t^3d+tc^5=0
                       \quad\hbox{in }\mathbb P(1,3,1).
\]
On the $t$ chart set $c=tz$ and $d=t^3w$.  Dividing the
equation of $T$ by $t^6$ gives
\begin{equation}\label{D:source-t-chart}
 w^2+(1+tz^4+t^3z^2)w+
            z^5+t^3z^7+t^4z^6+t^5z^5=0.
\end{equation}
Its exceptional curve is $w^2+w=z^5$, and the $w$
derivative is one along this curve.  Thus this chart is
smooth on a neighborhood of the exceptional curve.

The only point of the exceptional curve outside this chart
has $t=d=0$ and $c\ne0$ in the weighted projective coordinates.
Since $c$ has weight one, its chart is an ordinary affine
chart.  Put
\[
                   c=r,\qquad t=r\tau,\qquad d=r^3\delta.
\]
The strict-transform equation is
\begin{equation}\label{D:source-c-chart}
 \delta^2+(\tau^3+r+r^3\tau^4)\delta
       +\tau+r^3\tau^2+r^4\tau^4+r^5\tau^6=0.
\end{equation}
At the exceptional point $(r,\tau,\delta)=(0,0,0)$ the
$\tau$ derivative is one.  The two charts therefore prove
that $V$ is smooth near its exceptional fibre.  Off that
fibre the morphism is an isomorphism over the smooth
punctured surface $T$.  They also show that its only
exceptional curve $C'$ is the smooth proper completion of
$w^2+w=z^5$, of genus two by the double-cover computation
in the preceding example.

The original function $c$ equals $tz$ in the first chart.
Its strict transform meets $C'$ at precisely the two points
\[
                    (z,w)=(0,0),\qquad (z,w)=(0,1).
\]
Each intersection is transversal: the $w$ derivative in
\eqref{D:source-t-chart} is a unit there, so $t,z$ are
regular parameters.  In the second chart $c=r$ is only
the exceptional equation, so there is no additional
intersection at infinity.  Intersecting the principal
divisor of $c$ with $C'$ yields
\[
                            0=(C')^2+2.
\]

We give the map to $Y$ on the same charts.  On the $t$
chart of the target it is
\[
                    x=z^2+tz,\qquad y=w+z^5,\qquad t=t.
\]
The equation of $Y$ pulls back to
\eqref{D:source-t-chart}.  On a neighborhood of the source
point at infinity put
\[
             A=1+r\tau^2,\qquad B=1+\tau^2\delta.
\]
Both are units there.  The target's infinity coordinates
are
\[
             \rho=\frac{rB}{A^2},\qquad
             \xi=\frac{A^5}{B^2},\qquad
             \tau_Y=\frac{\tau A^2}{B}.
\]
These formulas follow from $a=r^2A$, $b=r^5B$ and $t=r\tau$,
and agree on overlaps.  Together with the original map off
the exceptional fibre they define a morphism $V\to Y$.
It is proper because $V$ is proper over $S$ and $Y$ is
separated over $S$.  Its fibres are finite away from $C$.
On $C'$ the map is
\[
                      (z,w)\longmapsto(z^2,w^2),
\]
since $w+z^5=w^2$ there.  It is consequently finite on
the exceptional fibre as well.  Properness and
quasi-finiteness make $V\to Y$ finite.  Thus $V$ is also
the normalization of $Y$ in the quadratic extension.

Finally $T$ is Gorenstein, and adjunction gives
$K_V\cdot C'=2g(C')-2-(C')^2=4$.  If $a$ is the
discrepancy of its sole exceptional curve, then
$a(C')^2=4$, whence $a=-2$.
\end{proof}

\begin{corollary}[Proved: the next quadratic step and the local pro-two quotient]
\label{D:wild-genus2-pro2}
Every finite normal cover of $T$ of degree two that is
\'etale off its singular point is tame at $C'$.
Let $U_T$ and $U_S$ denote the punctured spectra of the
strict henselizations at the singular points of $T$ and
$S$, respectively.  Then
\[
                 \pi_1(U_T)^{(2)}=1,\qquad
                 \pi_1(U_S)^{(2)}=\mathbb Z/2\mathbb Z,
\]
where the superscript denotes the maximal pro-two quotient.
This does not assert that the full wild inertia of either
punctured fundamental group is finite.
\end{corollary}
\begin{proof}
Apply Proposition~\ref{D:single-curve-residue-tameness} with
$g=2$, $d=2$ and $n=2$.  If any extension of the exceptional
valuation has inseparable residue, its residue degree is
at least two.  The degree formula then forces a unique
extension, with $e=1$, $f=2$ and $f^{\mathrm{sep}}=1$.
The inequality in that proposition gives
\[
                              2\delta\leq1.
\]
Its equality criterion for tame ramification instead gives
$\delta\geq1$ for an inseparable residue extension, a
contradiction.  Thus all residue extensions are separable,
and the same proposition gives tameness.  This argument
makes no assumption on the number of components in a
resolution of the covering surface.

We explain the strict-local assertion to keep its scope
separate from completion or algebraization claims.  A
finite cover of the punctured strict henselization descends,
by finite presentation, to an \'etale neighborhood
of the singular point.  Normalize that neighborhood and
shrink it if necessary.  The preceding argument applies
there, since the pulled-back resolution has the same
exceptional curve and normal degree.  A degree-two cover
of $U_T$ is therefore tame at $C'$.  In characteristic two
its ramification indices are prime to two and divide two,
so they are all one.  The normalization of the regular
resolution in this cover is unramified at every
codimension-one point, and hence is finite \'etale
by purity of the branch locus
\cite[Lemma 58.21.4, tag 0BMB]{StacksBWildSource}.
Proper henselian
comparison identifies its finite \'etale category
with that of $C'$; see
\cite[Lemma 58.9.1, tag 0A48]{StacksBCuspLocal}.
Here we pass from the scheme-theoretic closed fibre to
its reduction $C'$ using invariance of finite \'etale
covers under nilpotent thickenings.

There is no nontrivial \(\mathbb Z/2\mathbb Z\)-torsor on
$C'$.  Indeed, in the two-chart \v Cech basis
$h_1=w/z$, $h_2=w/z^2$ of $H^1(C',\mathcal O_{C'})$,
the equation $w^2+w=z^5$ gives
\[
                              F(h_1)=h_2,\qquad F(h_2)=0.
\]
The extra terms $z^3,z$ are regular on the finite chart,
and $w/z^4$ is regular on the infinity chart.  Thus $F$
is nilpotent, and the Artin--Schreier exact sequence gives
$H^1_{\text{ét}}(C',\mathbb F_2)=0$.  We conclude
that $\pi_1(U_T)$ has no quotient of order two.  Every
nontrivial finite two-group has such a quotient, so its
maximal pro-two quotient is trivial.

The given double cover remains connected on strict
henselization: its finite normalization is normal, its
closed fibre has only one point, and a finite algebra over
a henselian local ring is a product of its local factors.
This local normalization is strictly henselian and is
ind-\'etale over $\mathcal O_{T,o}$; it is therefore
the strict henselization defining $U_T$.
Thus $\pi_1(U_T)$ is the normal subgroup of index two in
$\pi_1(U_S)$ defined by that cover.  In any finite
two-group quotient of $\pi_1(U_S)$ the image of this
subgroup must be trivial, since $\pi_1(U_T)$ has no
nontrivial two-group quotient.  All such quotients
therefore factor through the given group of order two,
which proves the second assertion.
\end{proof}

\begin{remark}[Exact finite verification]
The companion program
\texttt{check\_wild\_genus2\_source.py} verifies 22 exact
identities over $\mathbb F_2$, including both resolution
charts and their overlap, the finite map on both charts,
the two intersections used for $(C')^2$, and the
Frobenius action on the two cohomology classes.  Its
recorded runtime is about $0.14$ seconds.  The arguments
for properness, connectedness, purity and the profinite
conclusions are the mathematical arguments above.
\end{remark}

\begin{remark}[Remaining scope of the example]
The source resolution and the pro-two calculation rule out
repeating this particular wild quadratic step.  They do
not rule out wild inertia inside a finite group that is
not a two-group, or after passing first to other finite
punctured covers.  A bound for those images does not follow from this calculation
alone.  Corollary~\ref{D:wild-genus2-finite-inertia} below supplies
such a bound by a separate finite formal cone construction;
the general D.1/D.2(ii) statements remain outside that example.
\end{remark}

\subsubsection{Finite formal cone covers and finite exceptional wild inertia}

\begin{lemma}[Proved: a finite formal cone cover]
\label{D:finite-formal-cone-cover}
Let $C$ be a smooth proper connected curve in a smooth surface
$Y$ over an algebraically closed field $k$ of characteristic $p>0$.
Write $N=N_{C/Y}$, with $\deg N<0$, and let
\[
 e_Y\in H^1(C,T_C\otimes N^{-1})
\]
be the class of the square-zero algebra extension
$0\to N^{-1}\to\mathcal O_{C_2}\to\mathcal O_C\to0$.
Let $h:C'\to C$ be a finite morphism of smooth proper connected
curves, of degree $m$, and put $N'=h^*N$.  Suppose that
\begin{align}
 h^*e_Y&=0
     &&\text{in }H^1(C',h^*T_C\otimes N'^{-1}),\label{D:cone-first}\\
 H^1(C',h^*T_C\otimes N'^{-n})&=0
     &&(n\geq2),\label{D:cone-tangent}\\
 H^1(C',N'^{\,1-n})&=0
     &&(n\geq2).\label{D:cone-normal}
\end{align}
Then there is a finite locally free formal morphism of rank $m$
\[
 \widehat{\operatorname{Tot}(N')}_{C'}
                 \longrightarrow \widehat Y_C
\]
whose restriction to the reduced closed curves is $h$, and whose
conormal map $h^*N^{-1}\to N'^{-1}$ is the specified identity.
The ramification index along the closed curves is one.

If $Y\to\operatorname{Spec}A$ is a projective resolution of a
normal excellent surface germ, isomorphic off its closed point,
and its reduced exceptional fibre is precisely $C$, then over
$\widehat A$ this formal morphism algebraizes to a finite locally
free morphism $Z\to Y_{\widehat A}$ of rank $m$, with $Z$ regular
and integral.  Its Stein ring $B=H^0(Z,\mathcal O_Z)$ is a finite
normal local $\widehat A$-algebra of fraction-field degree $m$,
and, as a complete local $k$-algebra,
\[
 B\simeq\widehat{\bigoplus_{j\geq0}H^0(C',N'^{-j})}.
 \tag{FC}
\]
The completion is at the positive-degree ideal.  Moreover,
$Z\to\operatorname{Spec}B$ is a resolution, and its exceptional
curve gives the radial valuation of this completed section cone.
No separability or pure inseparability of
$\operatorname{Frac}(B)/\operatorname{Frac}(\widehat A)$ is asserted.
\end{lemma}

\begin{proof}
Put $V=\operatorname{Tot}(N')$, with zero-section ideal $J$,
and $V_n=(C',\mathcal O_V/J^n)$.  Write $I=\mathcal O_Y(-C)$
and $C_n=(C,\mathcal O_Y/I^n)$.  We construct compatible maps
$f_n:V_n\to C_n$ extending $h$ and the stated conormal map.

We first explain the class in \eqref{D:cone-first}.  On affine
opens of $C$, smoothness supplies $k$-algebra splittings of
$\mathcal O_{C_2}\to\mathcal O_C$.  On overlaps their differences
are derivations $\mathcal O_C\to N^{-1}$, giving the cocycle $e_Y$.
Using one such splitting, pullback by $h$ and the prescribed
map of the square-zero ideals give a local algebra map
$\mathcal O_{C_2}\to h_*\mathcal O_{V_2}$.  The differences of
these maps are the pulled-back cocycle, now in
$\mathcal Hom(h^*\Omega_C,N'^{-1})$ on $C'$.
Changing a local map by a derivation changes this cocycle by
a coboundary.  Thus \eqref{D:cone-first} gives a global $f_2$.
This argument does not assume a global retraction $C_2\to C$,
and it does not identify the prescribed conormal map with the
zero conormal map of absolute Frobenius.

For $n\geq2$, regard $f_n$ as a morphism $V_n\to Y$.
The square-zero ideal of $V_n$ in $V_{n+1}$ is
$J^n/J^{n+1}=N'^{-n}$.  Smoothness of $Y$ gives local lifts,
and their differences form the sheaf
\[
 \mathcal H_n=h^*(T_Y|_C)\otimes N'^{-n}.
 \tag{FL}
\]
The obstruction is its first cohomology.  The tangent-normal
sequence, pulled back along the flat map $h$, is exact and gives
\[
 0\longrightarrow h^*T_C\otimes N'^{-n}
 \longrightarrow\mathcal H_n
 \longrightarrow N'^{\,1-n}\longrightarrow0.
 \tag{FN}
\]
A finite nonconstant map between smooth curves is flat, as
its local modules are torsion free over DVRs.  Conditions
\eqref{D:cone-tangent}--\eqref{D:cone-normal} therefore imply
$H^1(C',\mathcal H_n)=0$ and give a global lift $f_{n+1}$.
It factors through $C_{n+1}$ because it sends $I$ into $J$.
Its conormal map remains that of $f_2$.
The precise square-zero lifting and derivation inputs are
\cite[Lemmas 17.28.11, 37.9.4--5 and 37.11.7,
tags 01UP, 04FH, 04FJ and 02H6]{StacksBNegativeCurve}.

Taking inverse limits gives the claimed formal map.  We check
finiteness and flatness, rather than inferring them solely from
the finite closed-fibre map.  Locally on the target curve choose
a generator $t$ of $I$ and choose $s$ for $J$ on its inverse
image.  The image of $t$ is $s$ times a unit, by the prescribed
conormal isomorphism.  Modulo this parameter the map is $h$,
which is finite locally free of rank $m$.  After restricting
the target open, lift a basis of that finite free module.
Successive reduction modulo $t^r$ and completeness show that
the lifted basis generates the formal algebra.  A relation
between its members reduces to zero modulo $t$, hence all its
coefficients are divisible by $t$.  Both formal rings are
$t$-torsion free; repeating and using separatedness shows that
the relation was zero.  Thus the formal algebra is locally free
of rank $m$.  This also proves the stated ramification index.

For the algebraization assertion, replace $A$ by its completion.
The $\mathfrak m_A$-adic and $I$-adic topologies on the proper
resolution are equivalent: their defining ideals have the same
radical, and a finite affine covering gives cofinal powers.
Our finite formal algebra consequently gives a compatible
system of finite algebras on $Y\times_A A/\mathfrak m_A^r$.
The precise finite-morphism form of Grothendieck existence,
\cite[Lemma 30.28.2, tag 09ZT]{StacksBFiniteCone},
algebraizes it to $Z\to Y$.  In that theorem the closed support
is proper, and its proof algebraizes the multiplication and
unit using full faithfulness; no separate algebraization of
a nonproper formal neighborhood is required.

Flatness and rank $m$ hold along the closed fibre, by faithful
flatness of completed local rings.  The nonflat locus and the
loci of other ranks are closed.  If nonempty they would have
a closed point over the closed point of $\operatorname{Spec}A$,
since $Z$ and $Y$ are proper over this local base.  They are
therefore empty.  At every point of the closed fibre the
completed local ring of $Z$ agrees with that of the smooth
line-bundle model $V$; hence $Z$ is regular there.  Excellence
and the same properness argument make $Z$ regular everywhere.
Its closed fibre is connected, with reduction $C'$; every
connected component of $Z$ meets it.  Thus $Z$ is connected,
and a connected regular scheme is integral in this situation.

Proper coherent cohomology makes $B=H^0(Z,\mathcal O_Z)$ finite
over $A$.  Formal functions
\cite[Lemma 30.20.6, tag 087U]{StacksBFiniteCone}
identifies it with
\[
 H^0(\widehat V_{C'},\mathcal O)
  =\varprojlim_r\bigoplus_{0\leq j<r}H^0(C',N'^{-j})
  =\prod_{j\geq0}H^0(C',N'^{-j}).
 \tag{FS}
\]
The section ring of the ample line bundle $N'^{-1}$ is a
finitely generated normal positive graded domain with degree
zero part $k$.  As in Proposition~\ref{D:graded-cone-tame},
finite generation makes its positive-ideal topology equivalent
to the degree topology.  Thus (FS) is exactly (FC).
In particular $B$ is local and normal.  Normality also follows
directly because a global function on an integral normal
scheme that is integral over its global ring lies in every
local ring and hence is global.

The positive-rank finite locally free morphism $Z\to Y$ is
surjective, so the map $A\to B$ is injective.
On the puncture of $\operatorname{Spec}A$ the original resolution
is an isomorphism, so $Z$ is finite there.  Its Stein morphism
to $\operatorname{Spec}B$ is consequently an isomorphism off
the closed point.  On the generic fibre, finite flatness gives
a field of dimension $m$ over $\operatorname{Frac}A$; flat base
change of $H^0$ identifies it with
$B\otimes_A\operatorname{Frac}A$.  This proves the degree claim
and birationality of the Stein morphism.  At the generic point
of $C'$, a series in (FS) has valuation its least nonzero
homogeneous degree.  This is precisely the radial valuation
of the section cone, and it restricts to the valuation of $C$
with index one.
\end{proof}

\begin{proposition}[Proved: a fixed cone cover bounds all wild images]
\label{D:finite-cone-wild-bound}
In the algebraic situation of Lemma~\ref{D:finite-formal-cone-cover},
the image of the wild inertia of $C$ in the punctured
fundamental group of the surface germ is finite of order at
most $m$.  The same bound holds after strict henselization.
It bounds the wild image in every finite quasi-\'etale cover,
with no solvability or degree restriction on the covering group.
\end{proposition}

\begin{proof}
First work over the complete local ring $A=\widehat A$ and
use the finite local normal cone algebra $B$ constructed above.
The inverse image of the closed point is its unique closed
point.  Thus every finite \'etale cover of the puncture of $A$
pulls back to a finite \'etale cover of the puncture of $B$.
Proposition~\ref{D:graded-cone-tame} makes every such pulled-back
cover tame at the radial valuation $w$ of $B$.

Let $v$ be the valuation of $C$.  The actual finite morphism
$Z\to Y_A$ identifies $w$ as an extension of $v$; the corresponding
complete DVR fraction fields form a finite extension $L/K$ of
degree at most $m$.  We record the wild-inertia comparison
including inseparable extensions.  For a finite separable
extension, compatible choices inside a separable closure give
\[
 I_L=I_K\cap G_L,\qquad P_L=P_K\cap G_L. \tag{WI}
\]
Indeed the first equality follows from the action on the
common algebraic residue closure.  Inertia has a normal
pro-$p$ subgroup with pro-prime-to-$p$ quotient, by
\cite[Lemma 58.13.5, tag 0BUA]{StacksBFiniteCone}.
Intersecting with $G_L$ makes $P_K\cap G_L$ the pro-$p$
Sylow subgroup of $I_L$, which is its wild inertia.
For a finite purely inseparable extension, the equivalence
of separable closures and the uniquely extended valuation
identify the absolute Galois groups, their inertia subgroups,
and their characteristic wild subgroups.  Factoring a general
finite extension into its maximal separable part and a purely
inseparable extension proves the same assertion with
$G_L$ interpreted under this identification.

The preceding tame pullback kills $P_L$ in the monodromy of
every original punctured cover.  Thus, for any such monodromy
representation $\rho$,
\[
 |\rho(P_K)|
 \leq [P_K:P_K\cap G_L]
 \leq [L:K]_{\rm sep}
 \leq m. \tag{WB}
\]
The subgroup $G_L$ need not be normal.  All representations
kill the same intersection in (WI), so the full image of
$P_K$ in the punctured fundamental group is already a
quotient of a set of at most $m$ cosets; equivalently its
finite quotients have uniformly bounded order.  It is a
finite pro-$p$ group of order at most $m$.

Finally suppose the original germ is not complete.  Base
change the proper resolution to the completed local ring.
At the generic point of its closed exceptional curve, the
original DVR and the new DVR have isomorphic genuine DVR
completions, by \cite[Lemma 54.11.1, tag 0BG4]
{StacksBFiniteCone}.  Normalize in any finite cover and make
the same comparison.  The resulting finite extensions of
complete DVRs, and hence their wild inertia images, are the
ones just bounded.  This argument applies to each component
if a cover splits after completion.  Strict henselization
has the same completion here, since the residue field $k$ is
algebraically closed; alternatively each of its finite covers
descends to an \'etale neighborhood with that same completed
resolution.  Thus (WB) applies in all the stated versions.
We have not identified completion at the surface singularity
with completion at its boundary valuation.
\end{proof}

\begin{corollary}[Proved: the threshold $d>2g-2$]
\label{D:negative-curve-finite-wild}
Let a normal surface germ over an algebraically closed field
of characteristic $p>0$ admit a projective resolution with
sole exceptional curve a smooth proper connected curve $C$,
of genus $g$ and self-intersection $-d<0$.
If
\[
                         d>\max\{0,2g-2\},
\]
the exceptional wild inertia has finite image in its full
punctured fundamental group.  More precisely choose a power
$q=p^r$ satisfying
\[
 q(d-2g+2)>2g-2,\qquad qd>2g-2.
 \tag{Q}
\]
Its image has order at most $q$.
\end{corollary}

\begin{proof}
Take the $k$-linear relative Frobenius
$h:C^{(-r)}\to C$, of degree $q$, and put $C'=C^{(-r)}$.
The genus remains $g$.  The line bundle in the first
obstruction has degree $q(d-2g+2)>2g-2$, so its $H^1$
vanishes by Serre duality.  The two line bundles in the
later obstruction have degrees
\[
 q(nd-2g+2),\qquad q(n-1)d,\qquad n\geq2.
 \tag{FD}
\]
Both exceed $2g-2$ by (Q).  Apply the lemma and proposition.
The degree and Serre-duality inputs are
\cite[Lemmas 53.4.1 and 53.5.2, tags 0BS2 and 0BS6]
{StacksBNegativeCurve}.  The conclusion is finiteness of
the full exceptional wild image; it need not be trivial.
\end{proof}

\subsubsection{Torsion equality cases and the explicit fierce cover}

\begin{lemma}[Proved: killing a degree-zero obstruction]
\label{D:torsion-first-obstruction-cover}
Let $C/k$ be a smooth proper connected curve over an
algebraically closed field of characteristic $p>0$, let
$L$ be a torsion line bundle of degree zero, and let
$e\in H^1(C,L)$.  There is a finite map $h:C'\to C$,
obtained by a sequence of finite \'etale covers and relative
Frobenius morphisms, such that $h^*L\simeq\mathcal O_{C'}$
and $h^*e=0$.  If the product of the \'etale degrees is $s$
and the product of the Frobenius degrees is $q$, then
\[
             \deg h=qs,\qquad g(C')-1=s(g(C)-1).
\]
The Frobenius degree can be increased arbitrarily after
these two vanishing properties have been achieved.
\end{lemma}

\begin{proof}
Write the order of $L$ as $ap^b$, with $(a,p)=1$.
Pull back first by relative Frobenius of degree $p^b$;
with the necessary $k$-twists, this raises the line bundle
to its $p^b$-th power.  Its order now divides $a$.
A trivialization of its $a$-th power defines a finite
\'etale $\mu_a$-torsor on which it trivializes, by
\cite[Lemmas 59.28.1 and 59.28.5, section 03PK]
{StacksBFiniteCone}.  Choose a connected component of this
torsor.  We have reduced the assertion to killing one class
in $H^1(D,\mathcal O_D)$ for a smooth proper connected curve
$D$.

Let $F$ denote the semilinear Frobenius on this finite
dimensional space.  Since $k$ is perfect, the images of
$F^r$ are $k$-subspaces.  For sufficiently large $r$ the
images stabilize, and the space is the direct sum of a
nilpotent part and a part on which $F$ is bijective.
Pullback by an appropriate relative Frobenius kills the
nilpotent part of our class.  We explain why the remaining
part is a $k$-linear combination of Frobenius-fixed classes.

On an $r$-dimensional bijective semilinear space, write
$F(x)=A x^{[p]}$ with $A$ invertible.  The polynomial map
$F-1$ is \'etale, with Jacobian $-I$.  Its kernel has
$p^r$ points: after multiplying the equations by $A^{-1}$
they are $x_i^p$ minus linear forms, whose relatively
prime leading monomials give the basis
$\prod x_i^{e_i}$, $0\leq e_i<p$, for the quotient algebra.
This is also the exact calculation in
\cite[Lemma 59.63.2, section 0A3J]{StacksBFiniteCone}.
An $\mathbb F_p$-basis of this kernel is $k$-linearly
independent.  Indeed a shortest normalized $k$-linear
relation, after applying $F$ and subtracting, would give
a shorter one unless all its coefficients lie in
$\mathbb F_p$, contradicting independence over that field.
The fixed vectors therefore form a $k$-basis of the
bijective part.

By the Artin--Schreier exact sequence, each such fixed
class in $H^1(D,\mathcal O_D)$ is the image of a
$\mathbb Z/p\mathbb Z$-torsor on $D$; see
\cite[section 59.63, tag 0A3J]{StacksBFiniteCone}.
Pullback to that finite \'etale torsor kills its class.
A connected component of the fibre product of the
finitely many torsors kills every basis class and hence
the chosen $k$-linear combination.  This proves the
vanishing assertion.  All Frobenius maps here mean
$k$-linear relative Frobenius maps from the inverse
Frobenius twist; the usual semilinear notation records
the induced action under the corresponding scalar twists.
Relative Frobenius preserves the genus, and an \'etale
map of degree $s_i$ multiplies $g-1$ by $s_i$.
Multiplying these identities proves the displayed
formulas.  Further Frobenius pullbacks preserve a zero
class and a trivial line bundle.
\end{proof}

\begin{corollary}[Proved: the torsion equality case]
\label{D:negative-curve-equality-finite-wild}
In the surface setting of
Corollary~\ref{D:negative-curve-finite-wild}, suppose $g\geq2$,
$d=2g-2$, and $\omega_C\otimes N_{C/Y}$ is torsion.
Then the full exceptional wild inertia has finite image.
In particular this holds when the surface germ is
$\mathbb Q$-Gorenstein, without any restriction on its
Cartier index.
\end{corollary}

\begin{proof}
The line bundle
$L=T_C\otimes N^{-1}=(\omega_C\otimes N)^{-1}$
has degree zero and is torsion.  Apply
Lemma~\ref{D:torsion-first-obstruction-cover} to the
first-neighborhood class $e_Y$.  Write the resulting
degree as $qs$ as there, and add one relative Frobenius
if necessary so that $q>1$.
For $N'=h^*N$ and $n\geq2$, both line bundles in
\eqref{D:cone-tangent}--\eqref{D:cone-normal} have degree
\[
                 qs(n-1)d.
\]
This exceeds $sd=2g(C')-2$, since $q>1$.
Their first cohomology vanishes, and
\eqref{D:cone-first} holds by construction.
The finite cone cover and finite-wild propositions apply.

For the last assertion write
$K_Y=f^*K_S+aC$.  Adjunction gives
$a=-1-(2g-2)/d=-2$.
If $rK_S$ is Cartier, it is trivial on the local scheme
$S$.  Consequently $\mathcal O_Y(r(K_Y+2C))$ is
trivial, and restricting to $C$ gives
\[
                    (\omega_C\otimes N)^{\otimes r}
                         \simeq\mathcal O_C.
\]
This proves the torsion hypothesis directly; no cyclic
canonical cover is required.
\end{proof}

\begin{corollary}[Proved: full wild bounds in the explicit genus-two example]
\label{D:wild-genus2-finite-inertia}
For the explicit surfaces $T\to S$ in
Example~\ref{D:wild-genus2-single-curve}, the image of
exceptional wild inertia in the full punctured
fundamental group has order at most four for $T$ and
at most eight for $S$.  These statements also hold for
their punctured strict henselizations.
\end{corollary}

\begin{proof}
The source resolution in
Proposition~\ref{D:wild-genus2-source} has exceptional curve
\[
                   C':\quad w^2+w=z^5,
\]
of genus two and square $-2$.  Let $P_0,P_1$ be its two
points with $z=0$, and let $\infty$ be its point at infinity.
The already checked divisor of the function $c$ gives
\[
                    N_{C'/V}\simeq\mathcal O_{C'}(-P_0-P_1).
\]
On this curve
$\operatorname{div}(z)=P_0+P_1-2\infty$.
Also $\operatorname{div}(dz)=2\infty$.  For completeness,
on the infinity chart $u=z^{-1}$, $v=w/z^3$, the equation
is $v^2+u^3v=u$, and differentiation in characteristic two
gives
\[
                  (1+u^2v)\,du=u^3\,dv.
\]
Here $v$ is a uniformizer, $\operatorname{ord}_\infty u=2$,
and hence $\operatorname{ord}_\infty dz=2$; on the finite
chart $z$ is an \'etale coordinate, so there are no other
zeros or poles.  It follows that
$N_{C'/V}\simeq\omega_{C'}^{-1}$ and
$T_{C'}\otimes N_{C'/V}^{-1}\simeq\mathcal O_{C'}$.

The two-chart calculation in
Corollary~\ref{D:wild-genus2-pro2} gives
$F(w/z)=w/z^2$ and $F(w/z^2)=0$ on
$H^1(C',\mathcal O_{C'})$.  Thus the first-neighborhood
class is killed by the relative Frobenius of degree four.
For that map, the two later line-bundle degrees are both
$8(n-1)>2$ for $n\geq2$.  Lemma~\ref{D:finite-formal-cone-cover}
therefore supplies a fixed cone cover of degree four
over the completed germ $T$.
Proposition~\ref{D:finite-cone-wild-bound} gives the first
bound.

Composing that finite local cone algebra with the checked
quasi-\'etale double cover $T\to S$ gives a fixed finite
local cone cover of degree eight over the completed germ
$S$.  Its actual radial DVR extends the exceptional
DVR of $S$, by the finite morphism of resolutions in
Proposition~\ref{D:wild-genus2-source}.  Every punctured
cover of $S$ again pulls back to a punctured cone cover,
so the identical inertia-index argument gives the
second bound.  The first example already has a nonzero
wild image of order two; the present upper bounds do
not assert that four or eight is attained.
\end{proof}

\begin{remark}[Exact companion certificate]
The program \nolinkurl{parts/D/code/check_finite_cone_genus2.py}
checks sixteen characteristic-two polynomial, divisor-degree,
Frobenius-matrix and obstruction-degree assertions used in
Corollary~\ref{D:wild-genus2-finite-inertia}.
All passed in the independent D rerun in
$0.0381131569884019$ seconds; the output is
\nolinkurl{parts/D/output/finite_cone_genus2.json}.
Root independently read and reran the same program, with all
sixteen assertions passing in $0.04670357100258116$ seconds.
The formal-cover construction and inertia conclusions are
proved separately above.
\end{remark}

\subsubsection{A formal retraction and purely inseparable cone covers}

\begin{lemma}[Proved: linearization after Frobenius along a formal retraction]
\label{D:formal-retraction-cone-cover}
Let $C$ be a smooth proper connected curve of genus $g$ in a
smooth surface $Y$ over an algebraically closed field $k$ of
characteristic $p>0$.  Write $N=N_{C/Y}$ and $\deg N=-d<0$.
Suppose the entire formal completion $\mathfrak Y=\widehat Y_C$
admits a $k$-morphism
\[
                  r:\mathfrak Y\longrightarrow C,
                  \qquad r|_C=\operatorname{id}_C.
\]
Choose a power $q=p^a$ with $qd>2g-2$, and let
$h:C'=C^{(-a)}\to C$ be the $k$-linear relative Frobenius.
Then
\[
 \mathfrak Y\mathbin{\widehat\times}_C C'
       \simeq\widehat{\operatorname{Tot}(h^*N)}_{C'}
\]
as formal schemes over $C'$, inducing the identity on $C'$
and on its conormal bundle.  The resulting morphism from
this formal line bundle to $\mathfrak Y$ is finite locally
free of degree $q$, is radicial, and has normal ramification
index one.
\end{lemma}

\begin{proof}
We first describe the relative formal-coordinate argument.
Put $I=\mathcal O_Y(-C)$.  Over an affine open of $C$ on
which $N$ is trivial, the retraction splits the map to
$\mathcal O_C$.  A lift $t$ of a generator of $I/I^2$
then identifies the restricted formal algebra with
$\mathcal O_C[[t]]$.  Indeed the induced maps on successive
graded pieces are the isomorphisms
$(I/I^2)^{\otimes n}\simeq I^n/I^{n+1}$; completeness and
separatedness give the claimed isomorphism.  The lifts may
be chosen on affine opens because their nilpotent thickenings
are affine, and the transition maps on global sections are
surjective.  All coordinates here respect the fixed map to
$C$.

Consequently the overlap maps have the form
\[
          t_i=a_{ij}t_j+\sum_{n\geq2}b_{ij,n}t_j^n,
\]
where the coefficients are functions on the curve and
$a_{ij}$ is the transition function of the conormal line.
If the maps have been made linear modulo the $n$th power,
their degree-$n$ errors form a cocycle with values in
\[
 \mathcal Hom(N^{-1},N^{-n})=N^{1-n},\qquad n\geq2.
\]
Replacing a coordinate by itself plus a term of degree $n$
changes that cocycle by a coboundary.  Thus the vanishing
of $H^1(C,N^{1-n})$ permits the next linearization step.
There is no tangent-direction error, because the overlap
maps are maps over the fixed curve.  Nor is there a
first-order obstruction: the retraction already identifies
the first square-zero algebra with
$\mathcal O_C\oplus N^{-1}$.

Now form the indicated base change along $h$.  It is finite
locally free and radicial of degree $q$, since $h$ is.
The pulled-back closed curve is $C'$, its normal bundle is
$N'=h^*N$, and the preceding relative coordinates show that
the base-changed formal scheme is regular.  Its genus is
still $g$ and $\deg N'=-qd$.  For $n\geq2$ we have
\[
               \deg(N'^{\,1-n})=qd(n-1)>2g-2.
\]
Serre duality therefore gives all the required vanishings;
see \cite[Lemmas 53.4.1 and 53.5.2, tags 0BS2 and 0BS6]
{StacksBNegativeCurve}.  The compatible coordinate changes
converge, proving the formal isomorphism.  The ideal of
$C$ pulls back to the ideal of $C'$, so the conormal map
is an isomorphism and the normal ramification index is one.
\end{proof}

\begin{proposition}[Proved: tameness from a formal retraction]
\label{D:formal-retraction-tame}
In the preceding lemma, suppose $Y\to S$ is a projective
resolution of a normal excellent surface germ, isomorphic
off its closed point, whose sole reduced exceptional curve
is $C$.  Then every finite normal generically separable
cover of $S$ that is \'etale off the closed point is tame
at the exceptional valuation.  The same assertion holds
over the strict henselization.  There is no lower bound
on $d$ other than $d>0$.

The formal retraction hypothesis holds, in particular,
if a neighborhood of $C$ admits a morphism $f:Y\to B$
to a smooth curve such that $f|_C:C\to B$ is \'etale.
\end{proposition}

\begin{proof}
Pass to the complete local ring of $S$.  The finite formal
cone map in Lemma~\ref{D:formal-retraction-cone-cover}
algebraizes by \cite[Lemma 30.28.2, tag 09ZT]
{StacksBFiniteCone}.  The proof of
Lemma~\ref{D:finite-formal-cone-cover}, beginning with its
algebraization paragraph, applies verbatim to this already
constructed finite formal algebra: its Stein ring is the
completed section cone of $(C',h^*N)$, is finite over the
completed local ring of $S$, and the actual exceptional
valuation on that cone is the radial one.  This use of the
algebraization argument imposes none of that lemma's
tangent-cohomology hypotheses.

We check the stronger ramification conclusion using the
actual complete DVRs.  At the generic point of $C$, the
retraction supplies the coefficient field $k(C)$ and a
normal parameter $t$.  Before and after the formal base
change their fraction fields are, compatibly,
\[
                 K=k(C)((t))\ \subset\ L=k(C')((t)).
\]
The map on coefficient fields is relative Frobenius of
degree $q$, and $t$ is unchanged.  Thus $L/K$ is purely
inseparable.  These are also the genuine completed DVRs
of the algebraized resolutions, by their identical formal
neighborhoods.

Any given punctured cover of $S$ pulls back to a finite
\'etale cover of the punctured section cone, since the
finite local Stein map has only its vertex above the
closed point.  Proposition~\ref{D:graded-cone-tame} makes
that pullback tame at the radial valuation.  A finite
purely inseparable extension identifies the absolute
Galois groups with their compatible valuation, inertia,
and wild-inertia subgroups, as verified in
Proposition~\ref{D:finite-cone-wild-bound}.  The original
wild inertia therefore acts trivially as well.  This
proves tameness, including separability of the original
residue extension, rather than only a bound on its
ramification index.

For the uncompleted germ and its strict henselization,
the same boundary-DVR comparison as in that proposition
applies, using \cite[Lemma 54.11.1, tag 0BG4]
{StacksBFiniteCone}.  In particular no equivalence of
the two different kinds of completion is assumed.

Finally suppose $f|_C:C\to B$ is \'etale.  It is formally
\'etale by \cite[Lemma 37.8.7, tag 04FF]
{StacksFormalRetractionRoot}.  On every infinitesimal
neighborhood $C_n\subset Y$, the map $C_n\to B$ has a
unique lift $C_n\to C$ extending the identity on $C$.
Formal \'etaleness gives these lifts one square-zero
thickening at a time; uniqueness makes them compatible.
Their inverse limit is the required formal retraction.
\end{proof}

\begin{remark}
A splitting of the tangent-normal sequence alone is
weaker than a retraction of the entire formal
neighborhood.  The latter is the exact extra hypothesis
used here.  The argument uses the Frobenius base change
along that existing retraction and therefore does not
construct one for an arbitrary negative exceptional curve.
\end{remark}

\subsubsection{A finite-order retraction suffices for finite wild image}

\begin{proposition}[Proved: finite-jet cone covers]
\label{D:finite-jet-cone-cover}
Let $Y\to S$ be a projective resolution of a normal excellent
surface germ over an algebraically closed field $k$ of
characteristic $p>0$.  Suppose it is an isomorphism off the
closed point and has sole reduced exceptional curve a smooth
proper connected curve $C$.  Put
\[
 I=\mathcal O_Y(-C),\qquad
 C_n=(C,\mathcal O_Y/I^n),\qquad
 N=N_{C/Y},\quad \deg N=-d<0,\quad g=g(C).
\]
Suppose for some integer $r\geq2$ that
\[
 C_r\longrightarrow C
 \quad\hbox{is a $k$-retraction of }C\hookrightarrow C_r,
 \qquad rd>2g-2.                                      \tag{JR}
\]
Choose a power $q=p^a$ with
\[
       qd>2g-2,\qquad q(rd-2g+2)>2g-2.                \tag{JQ}
\]
Then the completed germ admits a fixed finite local cone
cover of fraction-field degree $q$, whose radial valuation
extends that of $C$.  Its full exceptional wild inertia image
has order at most $q$.  The same bound holds for the
uncompleted germ and its strict henselization.
Only the $r$th infinitesimal neighborhood is assumed to
retract; no retraction of the entire formal neighborhood
is required.
\end{proposition}

\begin{proof}
Let $h:C'=C^{(-a)}\to C$ be the $k$-linear relative
Frobenius of degree $q$, and put $N'=h^*N$.
Use the retraction in (JR) to form the ordinary finite
scheme
\[
                      W_r=C_r\times_C C'.
\]
Its projection to $C_r$ is finite locally free and
radicial of degree $q$.  Let $J_r$ be its reduced
closed-curve ideal.  The powers of $J_r$ give
thickenings $W_n$, $1\leq n\leq r$, with
\[
 W_1=C',\qquad J_r^n/J_r^{n+1}=N'^{-n}
                    \quad(0\leq n<r).
\]
Here the equality at $n=0$ means the degree-zero
quotient $\mathcal O_{C'}$.  The retraction of $C_r$
base-changes to one of every $W_n$ onto $C'$.

Write $V=\operatorname{Tot}(N')$, with zero-section
ideal $J$, and $V_n=(C',\mathcal O_V/J^n)$.
We first construct a compatible isomorphism
$W_r\simeq V_r$ over $C'$, inducing the conormal identity.
On affine opens where $N'$ is trivial, the retraction
and a lift of a conormal generator identify $W_r$
with $\mathcal O_{C'}[s]/(s^r)$.
The first-order algebra is therefore split.
Suppose these coordinate identifications have been made
compatible modulo $s^n$, where $2\leq n<r$.
Their errors in degree $n$ are a cocycle in
\[
 \mathcal Hom(N'^{-1},N'^{-n})=N'^{\,1-n}.
 \tag{JC}
\]
Changing one coordinate by a term of degree $n$
adds a coboundary.  Since
\[
             \deg N'^{\,1-n}=qd(n-1)>2g-2,
\]
Serre duality kills this first cohomology group.
Induction for $n=2,\ldots,r-1$ gives $W_r\simeq V_r$;
when $r=2$ there are no such steps.
This is the finite version of the relative-coordinate
argument of Lemma~\ref{D:formal-retraction-cone-cover}.
There are no tangent-direction errors in (JC), since
all maps in this finite stage are maps over $C'$.

Composing the inverse of this isomorphism with
$W_r\to C_r\hookrightarrow Y$ gives
$f_r:V_r\to Y$.  Its closed-curve map is $h$, and its
conormal map is the identity
$h^*N^{-1}\simeq N'^{-1}$.  From this point onward we
use ordinary maps to the smooth surface $Y$, without
requiring them to be maps over a curve.

Suppose $n\geq r$ and $f_n:V_n\to Y$ has been constructed.
For the square-zero extension $V_n\subset V_{n+1}$
with ideal $N'^{-n}$, smoothness gives local lifts.
Their obstruction sheaf is
\[
                  h^*(T_Y|_C)\otimes N'^{-n}.
 \tag{JL}
\]
It has tangent and normal subquotients
$h^*T_C\otimes N'^{-n}$ and $N'^{\,1-n}$, as in
\eqref{D:cone-tangent}--\eqref{D:cone-normal}.
Their respective degrees are
\[
                 q(nd-2g+2),\qquad q(n-1)d.
 \tag{JD}
\]
For every $n\geq r$, the first is at least
$q(rd-2g+2)>2g-2$, and the second is at least
$qd>2g-2$.  Their $H^1$ groups vanish; the exact
tangent-normal sequence therefore kills $H^1$ of
(JL).  The local lifts glue to a global $f_{n+1}$.
It sends $I$ into $J$ and retains the conormal identity
already fixed at the finite stage.
The square-zero deformation and Serre-duality inputs
are the same precise primary statements used in
Lemma~\ref{D:finite-formal-cone-cover}.

The inverse limit of these compatible maps is a formal
map
\[
       \widehat{\operatorname{Tot}(N')}_{C'}
                     \longrightarrow\widehat Y_C.
\]
It is finite locally free of rank $q$: modulo a normal
parameter it is the finite flat map $h$, and the
parameter maps to a generator times a unit.  The
complete basis-lifting argument in
Lemma~\ref{D:finite-formal-cone-cover} applies directly.
Its algebraization proof, starting from this constructed
finite formal algebra, likewise applies without imposing
vanishing for the earlier values $2\leq n<r$.
It supplies an actual finite map over the completed
resolution, a normal local Stein cone algebra of
degree $q$, and its actual radial extension of the
exceptional DVR.

Finally use the inertia-index proof of
Proposition~\ref{D:finite-cone-wild-bound} for this one
fixed finite cone extension.  All original punctured
covers pull back to radially tame cone covers, so the
full wild image has order at most $q$.
That proposition's common-DVR comparison also gives
the assertion before completion and after strict
henselization.  The extension constructed after
order $r$ need not be purely inseparable, so this
argument proves finiteness of wild image, without
asserting that the image is zero.
\end{proof}

\begin{corollary}[Proved: a split first neighborhood]
\label{D:split-first-neighborhood-finite-wild}
Under the same surface hypotheses, suppose that
$C_2=(C,\mathcal O_Y/I^2)$ retracts onto $C$, or
equivalently that its square-zero class
$e_Y\in H^1(C,T_C\otimes N^{-1})$ is zero.
If $d>\max\{0,g-1\}$, the full exceptional wild
inertia image is finite.
One bound for its order is any power $q=p^a$ satisfying
\[
             qd>2g-2,\qquad q(2d-2g+2)>2g-2.
\]
\end{corollary}

\begin{proof}
The square-zero splitting equivalence follows from
the derivation-torsor calculation in
Lemma~\ref{D:finite-formal-cone-cover}.
Apply Proposition~\ref{D:finite-jet-cone-cover} with
$r=2$.  Its numerical condition is $2d>2g-2$,
which is exactly $d>g-1$ in addition to $d>0$.
\end{proof}

\subsubsection{Mixed surface boundaries with finite exceptional wild image}

\begin{theorem}[Proved under stated hypotheses: mixed finite boundary with finite cone covers]
\label{D:finite-cone-mixed-boundary}
Let $k$ be algebraically closed of characteristic $p>0$, let
$S$ be a normal connected proper surface over $k$, and let
$Y\subset S$ be finite.  Suppose every singular point of $Y$
satisfies at least one of the following conditions:
\begin{enumerate}
\item $S$ is $\mathbb Q$-factorial at that point;
\item $S$ is log canonical there;
\item its completed local ring is a completed positively graded
      normal surface cone;
\item a projective resolution of its germ has sole reduced
      exceptional fibre a smooth proper curve $C$, and a finite
      map $h:C'\to C$ satisfies the three conditions
      \eqref{D:cone-first}--\eqref{D:cone-normal}, or some
      $C_r=(C,\mathcal O_{\widetilde S}/I_C^r)$ retracts onto
      $C$ with $r\geq2$ and $r(-C^2)>2g(C)-2$;
\item its completed local ring is isomorphic to one of the
      explicit genus-two completed germs $S$ or $T$ of
      Example~\ref{D:wild-genus2-single-curve};
\item a projective resolution has sole reduced exceptional
      fibre a smooth curve of negative self-intersection, and
      its entire formal neighborhood admits a retraction to
      that curve.
\end{enumerate}
Then $\pi_1(S\setminus Y)$ is topologically finitely generated.
The singularities retained in the open are unrestricted.

In (4), if $C^2=-d<0$ and $g=g(C)$, either of the following
conditions suffices:
\[
 d>\max\{0,2g-2\};\qquad
 g\geq2,\quad d=2g-2,\quad
 \omega_C\otimes N_{C/\widetilde S}\text{ is torsion}.
\]
In the equality case, the torsion condition holds whenever
the germ is $\mathbb Q$-Gorenstein, with arbitrary Cartier index.
If the first neighborhood $C_2$ retracts onto $C$, the weaker
numerical condition $d>\max\{0,g-1\}$ suffices.

The same conclusion holds for a possibly nonnormal proper
scheme of dimension at most two when $Y$ has codimension at
least two in every irreducible component and its inverse image
on every normal surface component satisfies these alternatives.
\end{theorem}
\begin{proof}
Corollaries~\ref{D:negative-curve-finite-wild}
and~\ref{D:negative-curve-equality-finite-wild} prove the stated
sufficient conditions for (4); Proposition~\ref{D:finite-jet-cone-cover}
and Corollary~\ref{D:split-first-neighborhood-finite-wild} supply
its finite-order retraction alternatives.  We prove the general assertion
with alternatives (1)--(6).

We construct a normal proper $\mathbb Q$-factorial modification
$q:X\to S$.  Leave the omitted $\mathbb Q$-factorial singular
points unchanged.  At every other omitted point of type (2),
use its minimal resolution and, for a nodal cyclic core, the
necessary node subdivisions.  A rational-tree exceptional
configuration is already $\mathbb Q$-factorial, by
Proposition~\ref{D:rational-tree-Qfactorial}, so that case was
retained.  At points of type (3), use the partial extraction
of Proposition~\ref{D:weighted-graded-extraction}; it is
$\mathbb Q$-factorial and has just its radial exceptional curve.
At points of type (4), use the specified resolution.  At points of type (6), use the resolution carrying the formal
retraction.  At points
of type (5), the explicit one-curve resolution descends from
completion.  Indeed, use the image weighted Rees algebra
$\bigoplus_m (I_m\mathcal O_S)s^m\subset\mathcal O_S[s]$
of the explicit strict transform, not the whole ambient
base-change blowup.  A sufficiently divisible Veronese is
generated in degree one and equals the Rees algebra of
$J=I_N\mathcal O_S$.  Its Proj is therefore the blowup of
that ideal, which is primary to the closed point; see \cite[Lemma 10.56.2, tag 0EGH]{C_StacksWeightedCone}.
Such an ideal descends from the completion since it contains
a power of the maximal ideal.  Its blowup therefore descends
with the same completed model and the same exceptional curve,
by \cite[Lemmas 54.11.1--3, section 0BG4]{StacksBFiniteCone}.
Regularity descends faithfully flatly.  In particular we do
not replace this one-curve model by a dominating blowup with
additional, uncontrolled exceptional divisors.
Finally resolve all singular points retained in $S\setminus Y$.

There are only finitely many singular points.  The local
modifications are isomorphisms on punctured neighborhoods;
spreading and gluing with the identity therefore give the
claimed proper birational map.  Its source is regular except
at the retained or extracted $\mathbb Q$-factorial points.
It is thus $\mathbb Q$-factorial everywhere and projective by
Lemma~\ref{D:Qfactorial-partial-resolution}.

Put $U=S\setminus Y$ and $U'=q^{-1}(U)$.  Connected finite
\'etale covers of the normal integral $U$ stay connected on
$U'$, since their generic algebra is unchanged.  Consequently
$\pi_1(U')\to\pi_1(U)$ is surjective.  Every exceptional
divisor of $X\setminus U'$ lies over an omitted point of one
of the types just resolved or partially extracted.

For type (2), its wild inertia image in $\pi_1(U)$ is zero by
Corollary~\ref{D:lc-core-valuative-tame}.  For type (3), it is
zero by Proposition~\ref{D:graded-cone-tame} and the weighted
extraction comparison.  For type (4), it is finite by
Proposition~\ref{D:finite-cone-wild-bound} or
Proposition~\ref{D:finite-jet-cone-cover}.  The same is true
for type (5), with bounds four or eight, by
Corollary~\ref{D:wild-genus2-finite-inertia}.
For type (6), the wild image is zero by
Proposition~\ref{D:formal-retraction-tame}.  These are bounds
in the full punctured local group, so they bound the images
in the global group as well.  The completion comparison uses
the same completed divisorial DVR for every original cover.
Retained omitted points have codimension two on $X$ and create
no divisorial boundary.  Exceptional curves above singularities
retained in $U$ lie in $U'$.

We may therefore apply Proposition~\ref{D:finite-exceptional-wild}
to this projective modification.  It chooses one open normal
subgroup avoiding the finitely many finite wild images; after
normalization in its finite cover, the resulting open subgroup
is a quotient of a projective fixed-boundary tame group.
That subgroup, and hence $\pi_1(U)$, is finitely generated.
In particular we have used finiteness of these wild images,
without asserting that the original covers are tame at them.

For a nonnormal proper scheme, pass to the disjoint finite
normalizations of its reduced irreducible components.  The
normal surface cases have just been proved.  Components of
dimension at most one have no omitted points under the
codimension hypothesis, and their proper fundamental groups
are finitely generated.  Finite proper descent, in the form
used in Proposition~\ref{D:fixed-boundary-positive} and
\cite[IX, Corollary 5.2]{SGA1}, proves the assertion downstairs.
\end{proof}

\begin{corollary}[Proved under stated hypotheses: isolated fixed boundary with finite cone covers]
\label{D:finite-cone-fixed-boundary}
Let $S$ be a normal connected proper surface over an
algebraically closed field of positive characteristic, and
let $Y\subsetneq S$ be closed.  Suppose every singular point
of $Y$ which is not $\mathbb Q$-factorial is isolated in $Y$
and satisfies one of alternatives (2)--(6) of
Theorem~\ref{D:finite-cone-mixed-boundary}.  Then $Q(S,Y)$ is
topologically finitely generated.  The numerical and torsion
sufficient conditions stated there apply, and retained
singularities are unrestricted.
\end{corollary}
\begin{proof}
Make the same partial projective modification as in the theorem,
retaining all $\mathbb Q$-factorial boundary singularities.
The strict transforms of original divisorial boundary points
have the identical DVRs, so their wild inertia already has
zero image in $Q(S,Y)$.  Every additional boundary divisor
lies over one of the specified isolated points.  Isolation
ensures that every cover counted by $Q(S,Y)$ restricts to a
finite \'etale cover of that punctured germ.  The preceding
local results consequently give a finite wild image at each
such divisor.  Curves introduced over points outside $Y$
remain in the open.  Proposition~\ref{D:finite-exceptional-wild}
again applies and proves finite generation.  This argument
retains the isolation hypothesis at the newly resolved
boundary points.
\end{proof}

\begin{example}[Proved: smallness and pronilpotent bounds do not imply finite generation]
\label{D:pronilpotent-obstruction}
There is a small profinite group with trivial maximal
prime-to-$2$ quotient such that every open subgroup has
finite maximal pronilpotent quotient, but which is not
topologically finitely generated. Set
\[
 m_n=|A_n|^n+1,\qquad
 G=\prod_{n\geq5}A_n^{m_n}.
\]
The alternating groups $A_n$ for $n\geq5$ are nonabelian
simple groups of even order
\cite[Theorem 1.1, p.~1, with proofs in sections 3--7]{ConradAnD}.
For any fixed finite group $F$, a homomorphism $A_n\to F$
is trivial once $|A_n|>|F|$. A continuous homomorphism
$G\to F$ consequently kills the whole corresponding tail:
finite-support elements are dense there. It factors through
a finite product of finite groups, so there are only finitely
many such homomorphisms. Taking $F=S_r$ proves smallness.

If $d$ elements generate $S^m$ for a finite group $S$,
their coordinate $d$-tuples must all differ. Otherwise the
generated subgroup lies in the proper equality diagonal
between two coordinates. Hence $m\leq |S|^d$.
The finite quotient $A_n^{m_n}$ therefore requires more
than $n$ generators, proving that $G$ is not finitely generated.

An open subgroup $H$ contains the entire product of the
blocks with $n>N$ for some $N$. Every finite nilpotent
quotient kills that tail, because each $A_n$ is perfect
and a finite nilpotent perfect group is trivial. Thus
$H$'s maximal pronilpotent quotient factors through the
finite quotient by that tail. Every odd-order quotient
of $G$ is trivial by simplicity and even order of the
factors. This proves all the asserted restrictions.

The attached \texttt{check\_small\_group\_obstruction.py}
checks the eleven finite coordinate-count certificates
for $5\leq n\leq15$. The separate
\texttt{check\_pronilpotent\_obstruction.py} exhaustively
checks the $A_5$ base facts, including all products,
commutators and nontrivial normal closures. These programs
do not replace the proof for arbitrary $n$ or the infinite
product argument. The construction is an abstract obstruction
to an inference, not a geometric counterexample to D.1 or D.2(ii).
\end{example}

% Post-checkpoint 07.  Independent audit requested before integration.
\begin{proposition}[Proved: a uniform first-cohomology bound does not imply finite generation]
\label{D:uniform-H1-obstruction}
The group $G$ of Example~\ref{D:pronilpotent-obstruction} satisfies
\[
 \dim_{\mathbb F}H^1_{
 \mathrm{cont}}(G,M)\leq\dim_{\mathbb F}M
\]
for every finite field $\mathbb F$ and every finite-dimensional
continuous $\mathbb F[G]$-module $M$.  Nevertheless $G$ is not
topologically finitely generated.  Moreover, for every open subgroup
$H\subset G$ and every finite continuous $\mathbb F[H]$-module $V$,
\[
 \dim_{\mathbb F}H^1_{
 \mathrm{cont}}(H,V)\leq [G:H]\dim_{\mathbb F}V.
\]
Thus even these estimates together with smallness and the previously
proved restrictions on all open subgroups do not settle the
finite-generation question.
\end{proposition}
\begin{proof}
We first give a bound for a finite direct product
$P=\prod_{i=1}^r S_i$ of perfect two-generated finite groups.
For any field $F$ and any finite-dimensional $F[P]$-module $M$,
\begin{equation}\label{D:perfect-product-H1}
 \dim_F H^1(P,M)\leq\dim_F M.
\end{equation}
We may extend scalars to an algebraic closure: kernels and images
in the finite group cochain complex commute with flat scalar
extension.  A composition series and the long exact cohomology
sequence reduce the inequality to irreducible modules.

Over an algebraically closed field every irreducible module for
a finite direct product has the form $\bigotimes_i M_i$, where
each $M_i$ is an irreducible $S_i$-module.  This assertion also
holds in modular characteristic.  Indeed, the ideals obtained
by tensoring the Jacobson radical of any factor group algebra
with the remaining algebras are nilpotent, and hence annihilate
an irreducible module.  After quotienting by them, the assertion
reduces to the tensor product of split matrix algebras.

If all factors act trivially, perfectness gives $H^1(P,M)=0$.
If exactly one factor $S$ acts nontrivially, the other factors
have no homomorphism to the additive group of $M$.  Restriction
therefore identifies $H^1(P,M)$ with $H^1(S,M)$.
A cocycle on $S$ is determined by its values on two generators,
so $\dim Z^1(S,M)\leq2\dim M$.  Since $M^S=0$,
$\dim B^1(S,M)=\dim M$, proving the required inequality.

Suppose that at least two factors act nontrivially.  Write
$P=S\times T$ and $M=V\otimes W$, where $V$ is a nontrivial
irreducible $S$-module and $W$ is a nontrivial irreducible
$T$-module.  Restriction gives an injection
\[
 H^1(P,M)\hookrightarrow H^1(S,M)^T
 =\bigl(H^1(S,V)\otimes W\bigr)^T=0.
\]
For clarity, the injection follows without any semisimplicity
assumption: a cocycle whose restriction to $S$ is a coboundary
can be adjusted to vanish on $S$, after which its values on
$T$ lie in $M^S=0$.  The image is $T$-invariant because the
cocycle identity on commuting elements makes conjugation differ
by a coboundary.  The last equality follows from $W^T=0$.
This proves~\eqref{D:perfect-product-H1}.

The alternating groups $A_n$ are perfect for $n\geq5$ and
two-generated.  Explicit generators are
\[
 \begin{cases}
  (1\,2\,\cdots\,n),\ (1\,2\,3),& n\ \text{odd},\\
  (2\,3\,\cdots\,n),\ (1\,2\,3),& n\ \text{even}.
 \end{cases}
\]
For odd $n$, conjugation gives the consecutive three-cycles,
which generate $A_n$ by induction.  For even $n$, conjugation
gives $(1\,i\,i+1)$; the products of consecutive such cycles
give $(i\,i+1\,i+2)$ on $\{2,\ldots,n\}$.  These generate
$A_{n-1}$, and the remaining cycle joins the letter $1$.
The induction starts with $A_3$, and $A_4$ follows from
$(1\,2\,3),(2\,3\,4)$: their generated group is transitive
and has a point stabilizer of order at least three, hence has
order twelve.

Every finite continuous $G$-module factors through a finite
subproduct $P$, whose complementary product $K$ acts trivially.
Every continuous homomorphism from $K$ to a finite abelian
group vanishes: it vanishes on each perfect factor and then on
their dense finite-support subgroup.  Consequently restriction
and inflation identify $H^1(G,M)$ with $H^1(P,M)$, and
\eqref{D:perfect-product-H1} proves the first assertion.
The failure of finite generation and all the smallness and
pronilpotent assertions were proved in
Example~\ref{D:pronilpotent-obstruction}.  Finally Shapiro's
lemma \cite[II, Proposition 1.11, p.~62, and Proposition 4.2,
p.~87]{MilneCFT2020D}, applied to the finite induced module, gives
\[
 H^1(H,V)=H^1(G,\operatorname{Ind}_H^G V),\qquad
 \dim\operatorname{Ind}_H^G V=[G:H]\dim V,
\]
proving the open-subgroup estimate.

The exact certificate
\nolinkurl{parts/D/code/check_uniform_h1_obstruction.py}
checks the displayed generators for $5\leq n\leq8$, the
four-dimensional deleted permutation module of $A_5$ over
$\mathbb F_2$, and its outer tensor square for $A_5^2$.
It gives $\dim H^1=0$ in both module examples, checks the
first calculation independently with all $14\,400$ bar-cocycle
equations, and verifies $2\,803\,802$ exact assertions.
The mathematical proof above, rather than these finite checks,
establishes the claim for every module and the infinite product.
This is an abstract obstruction to a proof strategy, not a
geometric counterexample to D.1 or D.2(ii).
\end{proof}

\begin{proposition}[Proved: the two uniform bounds needed for profinite generation]
\label{D:cohomology-crown-criterion}
Let $G$ be a profinite group.  For a finite quotient $F$ and a
non-Frattini chief factor $A$ of $F$, let $\delta_F(A)$ be the
number of non-Frattini chief factors in its $F$-equivalence
class.  Equivalently, this is the multiplicity in the crown
quotient $(L_A)_{\delta_F(A)}$; here $L_A$ is the associated
monolithic primitive group and
\[
 (L_A)_m=\{(l_1,\ldots,l_m)\in L_A^m:
                 l_1A=\cdots=l_mA\}.
\]
Then $G$ is topologically finitely generated if and only if
there is an integer $c\geq1$ such that both conditions hold:
\begin{enumerate}
\item for every prime $\ell$ and every finite irreducible
continuous $\mathbb F_\ell[G]$-module $V$,
\[
 \dim_{\mathbb F_\ell}H^1(G,V)
       \leq c\dim_{\mathbb F_\ell}V;
\]
\item for every finite quotient $F$ and every nonabelian
chief factor $A$ of $F$,
\[
 \delta_F(A)\leq |A|^c.
\]
\end{enumerate}
These conditions imply that $G$ is generated by at most $c+2$
elements.  No assertion that either uniform condition is already
known geometrically in the residual D.1/D.2(ii) cases is made.
\end{proposition}
\begin{proof}
We use the finite crown reduction of
\cite[Lemma 2.5 and Proposition 2.6, pp.~10--13]{LucchiniMarionTracey2020D}.
The required inputs are these.  A finite noncyclic group minimal
among its quotients with a given generator number has the form
$L_k$, where $L$ is monolithic primitive with socle $A$.
If $A=S^n$ is nonabelian, then
\[
 d(L)=\max\{2,d(L/A)\},\qquad
 k\leq\psi_L(s)\Longrightarrow d(L_k)\leq s
 \quad(s\geq d(L)),
\]
where
\[
 \psi_L(s)\geq
 \frac{|A|^{s-1}}{2n|\operatorname{Out}(S)|},
 \qquad |\operatorname{Out}(S)|\leq\log_2|S|.
\]
The crown quotient and multiplicity convention are those of
Proposition 2.4, p.~10; the two inequalities are equations
(2.1)--(2.2), p.~12, and the outer-automorphism bound is stated
on p.~13.  These finite-group inputs include the classification
of finite simple groups.  The following deduction is explicit.

If $G$ is $d$-generated, cocycles are determined by their values
on $d$ generators, proving the first bound with $c\geq d$.
Every crown quotient $L_m$ is also $d$-generated.  The $d$-tuple
in each coordinate has the same prescribed image in $(L/A)^d$,
so there are at most $|A|^d$ possible coordinate tuples.
Two identical tuples would force the generated group into a
proper equality diagonal.  Therefore $m\leq|A|^d$, proving
the second bound.

Conversely, fix a finite quotient $F$ and put $d=d(F)$.
There is nothing to prove if $d\leq2$.  Choose a quotient $Q$
of smallest order with $d(Q)=d$.  Every proper quotient of $Q$
has fewer than $d$ generators.  The crown reduction gives
$Q=L_k$, with socle $A^k$ and $H=L/A$; in particular
$d(H)<d$.  The factor $A$ occurs among the chief factors of
$F$ and $k\leq\delta_F(A)$.

First suppose that $A$ is abelian.  Thus $L=A\rtimes H$, and
$A$ is a faithful irreducible $\mathbb F_\ell[H]$-module.
Write $D=\operatorname{End}_H(A)$, $r=\dim_D A$ and
$s=\dim_D H^1(H,A)$.  Since $Q=A^k\rtimes H$ is split,
restriction of cocycles gives
\[
 \dim_D H^1(Q,A)=s+k.
\]
Indeed, every $H$-equivariant homomorphism $A^k\to A$
extends to a cocycle on the semidirect product, and the kernel
is the inflated $H^1(H,A)$.  Inflation from $Q$ to $G$ is
injective.  Condition (1), with dimensions converted from
$\mathbb F_\ell$ to $D$, therefore gives $s+k\leq cr$.

We record directly why this bounds generators, to avoid any
unstated restriction on $d(H)$.  Choose an integer
$b\geq d(H)$ and generators $h_1,\ldots,h_b$ of $H$,
allowing repetitions of the identity.  Evaluating cocycles
on this tuple embeds $Z^1(H,A)$ in $A^b$.  Its $D$-dimension
is $s+\theta r$, where $\theta=0$ if $H=1$ and $\theta=1$
otherwise.  If
\[
 k\leq br-s-\theta r,
\]
choose $k$ independent vectors in $A^b/Z^1(H,A)$ and use
their representatives as the coordinate values of lifts
$(a_i,h_i)\in A^k\rtimes H$.  These lifts generate $Q$.
Otherwise their intersection with $A^k$ is a proper
$H$-submodule, and some nonzero $H$-map $A^k\to A$
annihilates that intersection.  Their image in $A\rtimes H$
would then be a complement, hence the graph of a cocycle,
contradicting the chosen independence.
Taking $b=\max\{d(H),c+1\}$ meets the inequality.
Since $d(H)<d$, it follows that $d\leq c+1$.

If $A=S^n$ is nonabelian, we have
$d(L)=\max\{2,d(H)\}\leq d-1$.
Suppose that $d\geq c+3$ and take $s=d-1$ in the displayed
crown bound.  Since
\[
 2n|\operatorname{Out}(S)|
 \leq2\log_2|A|\leq |A|,
\]
we obtain
\[
 \psi_L(d-1)\geq |A|^{d-3}
      \geq |A|^c\geq k.
\]
This forces $d(Q)\leq d-1$, a contradiction.  Hence
$d(F)\leq c+2$ in every case.

Finally, in the compact space $G^{c+2}$ the sets of tuples
generating each finite quotient are nonempty clopen sets
with the finite-intersection property.  A tuple in their
intersection topologically generates $G$.
\end{proof}

\begin{remark}[The exact remaining cohomological route]
For the groups in D.1 and D.2(ii), our smallness proofs give
finiteness for each fixed finite torsor group and each fixed
finite module.  They provide neither uniform constant in
Proposition~\ref{D:cohomology-crown-criterion}.  In the
codimension-two case, $H^1(\pi_1(U),V)$ is the ordinary
\'etale $H^1(U,\mathcal V)$.  The degree-one Leray sequence
for $j:U\hookrightarrow X$ separates a global term
$H^1(X,j_*\mathcal V)$ from a local term in
$H^0(X,R^1j_*\mathcal V)$.  Constructibility and properness
give pointwise finiteness, without a proved estimate linear
in the rank as $\mathcal V$ varies.  For D.1 one must retain
the fixed-boundary tame cohomology classes: ordinary
$H^1(U,\mathbb F_p)$ can already be infinite on $\mathbb A^1$.

Even a successful uniform linear cohomology estimate would
leave condition (2).  In the preceding alternating-product
example, the finite quotient $A_n^{m_n}$ has
$\delta(A_n)=m_n=|A_n|^n+1$, so that condition fails for
every fixed exponent while condition (1) holds with $c=1$.
Thus a full argument needs a geometric bound on these
nonabelian multiplicities, or a different geometric reduction
that excludes them.  We do not infer such a bound from
Gabber's constructibility theorem alone.
\end{remark}

\begin{example}[Proved: fixed divisors and all divisorial valuations differ]
\label{D:fixed-vs-valuation}
Even for a normal projective surface, tameness at the divisors of one
fixed boundary is strictly weaker than valuation tameness. Let
$k=\overline{\mathbb F}_2$, let $V=\mathbb P^1_k\times\mathbb P^1_k$,
and let $H=C_2$ act on its affine coordinates by
\[
 \sigma(a,b)=(a+1,b+1).
\]
Let $X=V/H$, let $y$ be the image of $(\infty,\infty)$, and put
$U=X\setminus\{y\}$. Then
\[
 Q(X,\{y\})=\pi_1(U)=C_2,
 \qquad \pi_1^{\mathrm{tame,val}}(U/k)=1.
\]
\end{example}
\begin{proof}
The quotient exists as a finite quotient of a projective scheme;
one may use the natural $H$-linearization of $\mathcal O_V(1,1)$
and the invariant section ring. It is normal because invariant
subrings of normal domains are normal. The only geometric fixed
point of the involution is $(\infty,\infty)$, so
$V\setminus\{(\infty,\infty)\}\to U$ is a connected finite
\'etale $H$-torsor, by \cite[V, Corollary 2.3 and Proposition 2.6,
pp. 113--115]{SGA1}. Purity identifies the fundamental group of its
source with $\pi_1(V)$. The projective line has no nontrivial
connected \'etale cover by the genus formula; the proper homotopy
sequence \cite[X, Corollary 1.4, p. 263]{SGA1} applied to the
projection $V\to\mathbb P^1$ therefore gives $\pi_1(V)=1$.
The torsor consequently identifies $\pi_1(U)$ with $H$.
There is no codimension-one component of the omitted set, so
$Q(X,\{y\})=\pi_1(U)$.

Blow up $(\infty,\infty)$ in $V$, and write $x=1/a$, $y'=1/b$
near that point. On the chart $y'=xt$, the action is
\[
 \sigma(x)=\frac{x}{1+x},\qquad
 \sigma(t)=\frac{t(1+x)}{1+xt}.
\]
It preserves the exceptional divisorial valuation $w$ defined by
$x=0$ and acts trivially on its residue field $k(t)$. The inertia
group is therefore all of $C_2$, hence is wild. For a direct check
which also distinguishes ramification index from residue ramification,
the invariant field is
\[
 k(V)^H=k(a+b,a^2+a)=k(v,u),\qquad
 v=\frac{1}{a+b}=\frac{xt}{1+t},\quad
 u=\frac{1}{a^2+a}=\frac{x^2}{1+x}.
\]
The valuation has $w(v)=1$, and $z=u/v^2$ has residue
$(1+t)^2/t^2$. Thus its restriction to $k(v,u)=k(v,z)$ has
residue field $k(t^2)$; the extension of residue fields
$k(t)/k(t^2)$ is purely inseparable of degree two, while the
ramification index is one. This is not a tame extension.
The restricted valuation is divisorial, with centre $y$ on $X$;
the quotient of the blowup by $H$ realizes it as an exceptional
divisor on a normal modification of $X$.
It therefore tests the required additional tameness. Since the only
nontrivial connected cover of $U$
is excluded by this valuation, its valuation-tame group is trivial.
The rational identities, action, fixed locus, and exceptional
residue calculation are checked by the accompanying exact program.
\end{proof}

The remaining fixed-boundary question is in positive characteristic
with a proper nonprojective compactification outside the positive
cases above, and with a boundary having no quasi-projective open
neighborhood. In particular, a single omitted point is settled even
without projectivity. For D.1 this includes regular compactifications whose
boundary is not normal crossings, as well as singular compactifications.
For D.2(ii), purity settles regular compactifications, so the
remaining normal case is singular and lacks the stated finite regular cover.
The example shows concretely why a projective modification cannot
be invoked without controlling its new exceptional boundary divisors.
It does not refute finite generation: its full group is already finite.
The unresolved part records the residual proof obligation.

\subsection{D.2(i): a kernel that is not topologically finitely generated}

\begin{proposition}[Refuted: the kernel question of 29/066]\label{D:kernel-counterexample}
There is a normal connected affine complex surface $X$ and a closed point
$Y\subset X$ of codimension two for which
\[
 K=\ker\bigl(\pi_1^{\mathrm{et}}(X\setminus Y)\longrightarrow
                     \pi_1^{\mathrm{et}}(X)\bigr)
\]
is not topologically finitely generated. In this example $K$ is nevertheless
topologically normally generated by one element of order two.
\end{proposition}
\begin{proof}
Put
\[
 X=\operatorname{Spec}
 \frac{\mathbb C[a,b,c,(a+1)^{-1}]}{(ac-b^2)},\qquad
 Y=\{a=b=c=0\},\qquad U=X\setminus Y.
\]
The ring is the invariant ring of
$\mathbb C[y,z,(y^2+1)^{-1}]$ under $(y,z)\mapsto(-y,-z)$,
with $a=y^2$, $b=yz$, $c=z^2$. Indeed, each invariant monomial has
even total degree, and is a product of these three monomials; their only
relation is $ac=b^2$. The same assertion survives localization at $a+1$.
The invariant subring of a normal domain under a finite group is normal:
an element of its fraction field integral over the invariant ring is
integral over the original ring, hence belongs to that ring, and is
invariant. Thus $X$ is normal and integral. Its dimension is two.
The partial derivatives of $ac-b^2$ vanish simultaneously only at $Y$;
in particular $U$ is smooth and $Y$ has codimension two.

Write $C=\mathbb A^1_{\mathbb C}\setminus\{i,-i\}$ and
$W=(C\times\mathbb A^1)\setminus\{(0,0)\}$. The map $W\to U$
is a degree-two finite \'etale cover: the involution has no fixed point
on $W$. Removing one point from the complex manifold $C\times\mathbb C$
does not change its fundamental group. One way to see this is to use a
small ball about the removed point: its boundary is a simply connected
$3$-sphere, and van Kampen compares the space with and without its
interior. Consequently $\pi_1^{\mathrm{top}}(W)$ is free on two
generators $x_+,x_-$, represented by loops in $C$ about $i,-i$.

Use $(0,1/4)$ as base point upstairs. The path
$[0,1]\to W$, $s\mapsto(0,\tfrac14 e^{\pi i s})$, joins it to its image under
the involution and defines an element $q$ of $\pi_1^{\mathrm{top}}(U)$.
Its square lifts to a full circle in the second coordinate. That circle
contracts in a small punctured four-dimensional ball about the origin,
so $q^2=1$. We choose the two loops about $i,-i$ symmetrically, based at
$0$ in the first coordinate. Conjugation by $q$ exchanges their classes.
The exact sequence of the connected two-sheeted covering therefore gives
\[
 \pi_1^{\mathrm{top}}(U)
 =\langle q,x_+,x_-\mid q^2=1,\ qx_+q^{-1}=x_-\rangle
 \cong C_2*\mathbb Z.
\]
Here no relation in the free kernel has been discarded: this is the
split extension of that kernel by the displayed order-two lift.

The homotopy $(y,z)\mapsto(y,sz)$, $0\leq s\leq1$, is equivariant
and descends on $X(\mathbb C)$ to a deformation retraction onto
$C/\{\pm1\}=\mathbb A^1_a\setminus\{-1\}$. Thus
$\pi_1^{\mathrm{top}}(X)=\mathbb Z$. The element $q$ maps to zero,
while either loop $x_\pm$ maps to a generator: $y^2+1$ has a simple
zero at each of $i,-i$. Hence the map is
$C_2*\mathbb Z\to\mathbb Z$ that kills $C_2$ and is the identity on
the other factor. The complex Riemann existence theorem \cite[XII, Theorem 5.1, p. 333]{SGA1} identifies the
corresponding \'etale groups with the profinite completions and identifies
their map with the completion of this map.

For each positive integer $m$, let
\[
 W_m=(C_2)^m\rtimes C_m
\]
where a generator of $C_m$ cyclically permutes the $m$ factors. Sending
$q$ to the first coordinate involution and $x_+$ to the cyclic shift
defines a surjection $\widehat{C_2*\mathbb Z}\twoheadrightarrow W_m$.
The conjugates $x_+^j qx_+^{-j}$ lie in $K$, and their images are all
$m$ coordinate involutions. The image of $K$ is exactly $(C_2)^m$:
it contains this group, and is contained in it since the quotient map
to $C_m$ factors through $\widehat{\mathbb Z}$. Thus $K$ has elementary
abelian quotients requiring arbitrarily many generators and is not
topologically finitely generated. On the other hand quotienting by the
closed normal subgroup generated by $q$ gives precisely
$\widehat{\mathbb Z}$, proving the last assertion.

The program \texttt{parts/D/code/check\_examples.py} exhaustively checks
the finite wreath-product maps for $1\leq m\leq9$. The arbitrary-$m$
argument above, rather than that finite range, proves non-finite generation.
\end{proof}

\begin{proposition}[Proved under stated hypotheses: the regular replacement]
\label{D:kernel-regular}
If $X$ is a connected regular scheme of finite type over an algebraically
closed field and $Y\subset X$ has codimension at least two, then
$\pi_1(X\setminus Y)\to\pi_1(X)$ is an isomorphism.
\end{proposition}
\begin{proof}
Restriction of finite \'etale covers is fully faithful because $X$ is
normal and the open subset is dense. Normalize $X$ in any connected
finite \'etale cover of $X\setminus Y$; normalization is finite since
$X$ is excellent. The purity of the branch locus for a regular target
\cite[X, Theorem 3.3, pp. 276--277]{SGA1} says that a nonempty branch locus of this generically \'etale finite
map has a codimension-one component. No such component is possible,
since the map is \'etale off $Y$. Thus it is \'etale everywhere.
Restriction is an equivalence and yields the asserted group isomorphism.
\end{proof}


\subsubsection{A rational tail that cannot be filled in}

\begin{example}[Refuted: rational exceptional tails may remain wildly ramified]
\label{D:pure-root-rational-tail-obstruction}
Rationality of the tails introduced by a purely inseparable
modification does not imply that a pulled-back cover extends
\'etale across them.  The following local example proves
this failure even when the normalized pure cover is smooth.
It is not a counterexample to D.1 or D.2(ii).
\end{example}

\begin{proof}
Let $k$ be algebraically closed of characteristic two and put
\[
       Y=\operatorname{Spec}k[x,t],\qquad C=(t=0).
\]
Adjoin $z$ with $z^2=t+x^3$.  This is a purely inseparable
field extension of degree two: $t+x^3$ is not a square in
$k(x,t)$, for example because its $t$-derivative is one.
Its normalization over $Y$ is
\[
       Y'=\operatorname{Spec}k[x,z],\qquad t=z^2+x^3.
\]
Indeed this polynomial ring is finite over $k[x,t]$,
normal, and has the required fraction field.  Thus $Y'$
is smooth and locally factorial.

The inverse image $C'$ of $C$ is the cusp $z^2=x^3$.
Its normalization has coordinates
\[
                      x=w^2,\qquad z=w^3.
\]
At the cusp the element $w=z/x$ is not in
$k[[w^2,w^3]]$, whose nonconstant terms have order at
least two.  Consequently there is no morphism
$C'\to(C')^{\rm nor}$ inverse to normalization.
In particular no retraction of the formal neighborhood
onto this smooth normalized main curve can restrict
to such an inverse on the reduced curve.  Retaining
the locally factorial surface $Y'$ does not resolve
this problem with its main curve.

Blow up the origin of $Y'$ and use the chart
$x=s$, $z=sy$.  On this chart
\[
                  t=s^2(y^2+s).
 \tag{PT}
\]
The exceptional divisor $R=(s=0)$ is rational.  The
strict transform $E$ of $C'$ is $y^2+s=0$ and meets
$R$ only at $y=0$.

On $Y\setminus C$ consider the cover
\[
                         u^2+u=\frac1t.
 \tag{PA}
\]
It is finite \'etale of degree two, since the derivative
with respect to $u$ is one.  It is connected: if
$b^2+b=1/t$ in $k(x,t)$, the $t$-valuation of $b$
would have to be negative, and the left side would
then have an even negative valuation, whereas the
right side has valuation $-1$.  This is also the
usual Artin--Schreier torsor construction
\cite[section 59.63, tag 0A3J]{StacksBFiniteCone}.

Pull this cover back along the pure normalization and
the blowup.  Use (PT) and put
$v=u+1/(sy)$.  In characteristic two the exact
rational-function identity is
\[
 v^2+v
 =\frac{1}{s^2(y^2+s)}
       +\frac1{s^2y^2}+\frac1{sy}
 =\frac{1+y^3+sy}{s\,y^2(y^2+s)}.                    \tag{PW}
\]
At the generic point of $R$, the right side has a
simple pole, with nonzero leading coefficient
$(1+y^3)/y^4$ in $k(y)$.  In any extension of this
DVR valuation, normalized to give $s$ value one,
equation (PW) forces $v$ to have value $-1/2$.
The ramification index is therefore two, since the
field extension has degree two.  In particular it
is wild at this rational exceptional divisor.

The generic point of $R$ is outside $E$.  Thus a
cover that is \'etale on the complement of $E\cup R$
in this chart need not extend \'etale to the complement
of $E$ obtained by filling in $R\setminus E$.
Further embedded resolution does not change the
valuation at the generic point of the strict
transform of $R$.

The curve $C$ in this example is affine and is not
the proper exceptional curve of a contracted
surface germ.  The conclusion concerns precisely
the failed extension inference, not the existence
of a global cover violating any finite-generation
statement.  The exact companion program
\texttt{check\_pure\_root\_tail.py} checks thirteen
finite polynomial and rational-function identities;
the ramification and scope conclusions are the
arguments above.
\end{proof}

\subsection{D.3: the arithmetic result with invertible coefficients}

\begin{proposition}[Known: the invertible-prime arithmetic $K(\pi,1)$ theorem]
\label{D:arithmetic-invertible}
Let $K$ be a number field, let $S$ be a finite set of finite places,
and let $X=\operatorname{Spec}\mathcal O_K\setminus S$. For every prime
$\ell$ invertible on $X$, the pro-$\ell$ completion of the \'etale
homotopy type of $X$ is a $K(\pi,1)$. In particular, all of its higher
homotopy pro-groups, including its second one, vanish. This includes
$\ell=2$ when $K$ has real places, on the ordinary \'etale site.
\end{proposition}
\begin{proof}
Invertibility of $\ell$ means exactly that $S$ contains every finite
place of $K$ over $\ell$. Thus this is \cite[Proposition 2.3]{Schmidt2007};
its meaning as a higher-homotopy vanishing statement is
\cite[Proposition 2.1]{Schmidt2007}. The remark immediately after
Proposition 2.3 distinguishes, and also treats, the modified site at
real places. We can also avoid any convention issue in the real
$\ell=2$ case: adjoining $i$ gives a finite \'etale cover of degree
two, since $2$ is invertible on $X$. Its number field is totally
imaginary, where Proposition 2.3 applies. Corollary 2.2 of the same
source identifies the universal pro-$2$ covers and descends the
$K(\pi,1)$ property back to $X$.
For $K=\mathbb Q$, the allowed primes are exactly the primes in $S$.
When $S$ is empty there are no invertible primes, so a statement
restricted to such coefficients is vacuous.
\end{proof}

\begin{theorem}[Proved: arithmetic tame $K(\pi,1)$ for all permitted coefficients]
\label{D:arithmetic-tame}
Let $X_0=\operatorname{Spec}\mathcal O_K\setminus S$, let $D$ be
any finite reduced divisor on $X_0$, and put $U=X_0\setminus D$.
Let $\mathcal L$ be the set of rational primes invertible on $X_0$;
equivalently, $\ell\in\mathcal L$ if and only if every place of $K$
above $\ell$ lies in $S$. Let $\pi=\pi_1^t(U/X_0)$ classify finite
\'etale covers of $U$ tamely ramified along $D$.
For every finite discrete $\pi$-module $M$ whose order has prime
divisors in $\mathcal L$, and every $q\geq0$, the comparison is an
isomorphism
\[
 H^q(\pi,M)\xrightarrow{\ \sim\ }H^q_{\mathrm{et}}(U,M).
\]
The permitted-prime second tame homotopy pro-group is zero:
$\Pi_2^{t,\mathcal L}(U/X_0)=0$. Thus every pair $(K,S)$ has the
$K(\pi,1)$ property asked for with tame coefficients prime to all
residue characteristics of $X_0$. For $K=\mathbb Q$ and $D$ empty,
this gives $\Pi_2^t(\operatorname{Spec}\mathbb Z-S)=0$ in that
coefficient convention, with allowed primes exactly $S$.
\end{theorem}
\begin{proof}
Let $\widetilde U^t$ denote the inverse system of connected finite
tame covers, and fix $\ell\in\mathcal L$. We prove
\[
 \mathop{\operatorname{colim}}_{Y\to U\ \mathrm{tame}}
           H^q(Y,\mathbb F_\ell)=0\qquad(q>0).
\]
Every connected $Y$ in this system is an open subscheme of the ring
of integers of a finite extension $K'/K$. The prime $\ell$ is
invertible on it. Proposition~\ref{D:arithmetic-invertible} and
\cite[Proposition 2.1(iii)]{Schmidt2007} say that every positive-degree
cohomology class on $Y$ dies after a finite subcover of its universal
pro-$\ell$ covering. Such a subcover has $\ell$-power Galois group.
At every place above $D$ the residue characteristic is different
from $\ell$, so this extension is tamely ramified there: its
ramification index divides its $\ell$-power degree, and the residue
fields are finite and hence perfect. Tameness is preserved by
composition of finite extensions of discretely valued number fields.
The resulting cover is therefore still an object of the tame tower
over $U$. It kills the chosen class and proves the displayed limit
vanishing. No cofinality of the pro-$\ell$ tower inside the whole
tame tower is being asserted or needed.

A finite constant abelian group of permitted order has a finite
filtration with factors $\mathbb F_\ell$, $\ell\in\mathcal L$.
The long exact sequences give the same vanishing for all such
constant coefficients. A finite tame local system becomes constant
at a finite level of the tame tower, so it too has no positive-degree
cohomology on $\widetilde U^t$. The Hochschild--Serre spectral
sequence of this tower now has only its row of degree zero and gives
the asserted comparison in every degree.

The second homotopy statement uses precisely the permitted-coefficient
convention: on 29/088 the manuscript identifies homomorphisms from
$\Pi_2^t$ to a constant finite permitted coefficient group with its
degree-two cohomology on $\widetilde U^t$. The latter groups all
vanish by the argument above. A nonzero profinite abelian pro-group
with permitted primes has a nonzero finite permitted quotient,
so these tests force its second homotopy pro-group to be zero.
The argument simultaneously treats all primes in $\mathcal L$,
and does not infer a simultaneous completion theorem merely from
separate pro-$\ell$ completions. If $\mathcal L$ is empty, only
trivial coefficients and the trivial permitted-prime pro-group occur.
\end{proof}

The divisor $D$ here is the manuscript's additional $D'_0$.
Tameness at the primes removed in $S$ is a different condition: those
primes are not points of $X_0$. Neither this statement nor the
notation $\Pi_2^{t,\mathcal L}$ asserts vanishing for unrestricted
coefficients at residue characteristics that occur on $X_0$.

\subsection{D.4: the exceptional characteristic-two pencil}

\begin{proposition}[Known, with an integral deduction: 29/116]
\label{D:vanishing-conjugacy}
Let $X\subset\mathbb P^N_k$ be smooth projective and connected over an
algebraically closed field of characteristic two, with $\dim X$ odd.
Fix an odd prime $\ell$. For a sufficiently general Lefschetz pencil
of hyperplane sections, its vanishing cycles in the middle
$\mathbb Z_\ell$-cohomology of a smooth member, with the usual Tate
twist, form one orbit under the geometric monodromy group. The same
holds after reduction modulo $\ell^a$, for each $a\geq1$.
Here sufficiently general has its usual precise meaning: the axis
belongs to a suitable nonempty open subset of the Grassmannian.
\end{proposition}
\begin{proof}
Set $n=\dim X-1$, which is even. Deligne's
\emph{La conjecture de Weil. II} \cite[\S4.2, Corollary 4.2.7(ii)]{Deligne1980D},
printed page 221, states that the Picard--Lefschetz reflections are
conjugate for a sufficiently general pencil in precisely this exceptional
case; the passage from generic to general uses \S1.11.
The monodromy group in \S4.2.5 is the image on integral
$\mathbb Z_\ell$-cohomology. The local formula of \S4.2.1 gives,
for a vanishing cycle $\delta$,
\[
 r_\delta(v)=v-\epsilon\langle v,\delta\rangle\delta,
 \qquad \langle\delta,\delta\rangle=2\epsilon,
 \qquad \epsilon\in\{1,-1\}.
\]
Since two is a unit in $\mathbb Z_\ell$, the image of
$r_\delta-1$ is the free rank-one submodule $\mathbb Z_\ell\delta$.
This remains true if the ambient cohomology has torsion. Conjugacy of
two such reflections by an isometry sends their image submodules to
each other, so it sends $\delta$ to $u\delta'$ for a unit $u$.
The self-pairings give $u^2=1$, whence $u=1$ or $-1$.
Finally $r_{\delta'}(\delta')=-\delta'$, so either choice of sign is
itself obtained by monodromy. This proves actual conjugacy of all
signed cycles. Reduction of the same monodromy elements gives the
assertion with finite coefficients. This deduction uses that $n$ is
even and $\ell$ is odd; it is not an unrestricted integral enhancement
of Deligne's Corollary 4.2.8.
\end{proof}

\subsection{D.5: proper descent for infinite constant torsors}

\begin{proposition}[Known: the question of 29/180]\label{D:proper-descent}
Let $S'\to S$ be a proper surjective morphism of noetherian schemes,
and let $G$ be any discrete group. Every descent datum, as a principal
right $G$-bundle, on $G_{S'}=\coprod_{g\in G}S'$ is effective in
schemes. The descended object is a separated \'etale $G_S$-torsor.
No finiteness or residual-finiteness assumption on $G$ is needed.
\end{proposition}
\begin{proof}
Properness over a noetherian base implies finite presentation.
Rydh \cite[Theorem 5.19(ii)]{Rydh2010} gives effective descent for all
\'etale algebraic spaces along this map. It applies to $G_{S'}$ even
when $G$ is infinite. By Proposition 5.4, the descended \'etale space
$P\to S$ is separated; Proposition A.1 makes it a scheme.
Proposition 5.3 descends the equivariant action
$G_{S'}\times_{S'}P_{S'}\to P_{S'}$ uniquely to
$G_S\times_SP\to P$. The group-action identities descend by
faithfulness. The equivariant isomorphism
$P_{S'}\times G_{S'}\to P_{S'}\times P_{S'}$ and its inverse
descend by the same full faithfulness. Surjectivity of $P\to S$
follows after the given surjective base change. Thus $P\to S$ is
an \'etale surjection and satisfies the torsor identity; its pullback
is the original pointed or unpointed descent datum, as appropriate.
\end{proof}

\subsection{D.6: finite presentation of the prime-to-\texorpdfstring{$p$}{p} quotient}

\begin{proposition}[Known: the question of 29/041]\label{D:prime-to-p}
If $k$ is algebraically closed of characteristic $p>0$ and $X$ is a
connected scheme of finite type over $k$, then its maximal prime-to-$p$
\'etale fundamental group is topologically finitely presented.
Resolution of singularities is not assumed.
\end{proposition}
\begin{proof}
This is the assertion of Orgogozo,
\emph{Alt\'erations et groupe fondamental premier \`a $p$},
Theorem 4.11, printed page 143 \cite{Orgogozo2003}.
His hypotheses are a separably closed field of characteristic exponent
$p\geq1$ and an algebraic scheme over it. Our field is separably
closed and our finite-type scheme is algebraic. His proof explicitly
uses alterations in place of resolution. We use the theorem with
these identical hypotheses; no smoothness, properness, or projectivity
is introduced in this application.
\end{proof}

\subsection{D.7: the missing essential-surjectivity statement}

\begin{proposition}[Refuted: the individual functor of 29/203--204]
\label{D:alpha-counterexample}
The axioms R1--R4, full faithfulness and local filtering of $C$, and
the generating property stated on 29/202 do not imply essential
surjectivity of the functor $\alpha$ for each admissible sieve $C_0$.
This failure occurs even with $r(e)=S$ and with all realizations finite
over their bases.
\end{proposition}
\begin{proof}
Take $S=\operatorname{Spec}\mathbb F_3$. Over an \'etale $S$-scheme
$T$, let $\mathscr R_T$ be the category of finite \'etale covers of
$T$, with the evident realization functor. This is a stack of
multigaloisian categories. R1 is the existence and effectiveness of
finite disjoint unions and finite equivalence-relation quotients of
finite \'etale covers. R2 identifies subobjects with open-and-closed
subschemes. R3 is \'etale descent on such an open cover, and R4 is
the locally constant degree decomposition of a finite \'etale map.

Let $C_T$ contain just the terminal cover $T\to T$. These terminal
objects and their pullbacks form a fibrewise terminal category $C$.
Its inclusion $\varphi:C\to\mathscr R$ is cartesian and fully
faithful, and $C$ is locally filtering under either the strict or the
local version of filtering. It is generating in the manuscript's
definition: every finite \'etale cover becomes a constant finite
cover after an \'etale covering of $T$. Every $\varphi(M)\to e_T$
is an isomorphism, so the additional margin hypothesis of 29/204 holds.

Take the whole sieve $C_0=C$, which is admissible. Because this sieve
contains $e_S$, the source $\mathscr R_S^{C_0}$ consists exactly of
constant finite covers of the connected scheme $S$. A cartesian section
of $\mathcal E|C$ is, on the other hand, exactly a finite \'etale
cover of $S$: evaluate at $S$, and obtain all its other values by
pullback. Hence $\alpha$ is the inclusion of constant finite covers
into all finite \'etale covers. The connected degree-two cover
$\operatorname{Spec}\mathbb F_9\to\operatorname{Spec}\mathbb F_3$
is not in its image. We use $\mathbb F_9=\mathbb F_3[T]/(T^2+1)$;
irreducibility, separability and the nine-element field arithmetic are
checked in the accompanying program.

This example distinguishes two different statements. Restriction to
the sieve generated by $\operatorname{Spec}\mathbb F_9$ trivializes
this particular cover; therefore it refutes the asserted equivalence
for each fixed sieve, without by itself refuting an equivalence after
taking the filtered colimit over sieves.
\end{proof}

\begin{proposition}[Refuted: unrestricted reconstruction from R1--R4]
\label{D:alpha-geometric-counterexample}
If only the axioms written on 29/191--192 are imposed, the filtered
colimit reconstruction can also fail. A further hypothesis controlling
finite \'etale covers of the realizations is necessary.
\end{proposition}
\begin{proof}
Let $S=\operatorname{Spec}\mathbb C$, and put
$T=\mathbb G_{m,\mathbb C}$. Take the same stack of finite \'etale
covers on $\operatorname{Et}(S)$, but realize a cover $A\to S'$
as $r(A)=A\times_ST\to S'$. Since $T$ is geometrically connected,
the open-and-closed subsets of $A\times_ST$ correspond exactly to
those of $A$. Thus R2 holds. Products with $T$ preserve the disjoint
unions and the finite quotient diagrams in R1; R3 and R4 are unchanged.
The fibrewise terminal $C$ is still fully faithful, locally filtering,
and generating. The \'etale site of an algebraically closed point has
only disjoint unions of that point. Therefore refining admissible
sieves changes neither the fact that source objects are constant finite
covers nor the fact that the target contains the connected cover
$\mathbb G_m\to\mathbb G_m$, $z\mapsto z^2$.
This cover remains connected after every nonempty \'etale base change
of $S$, and it cannot come from a constant finite cover. No sieve
refinement repairs the obstruction.

The written axioms permit this realization: they do not require
$r(e)=S$. Indeed, the construction on 29/196 starts from an arbitrary
$S$-scheme $Z$ with group action and realizes its final object as
$Z/G$. If one adds $r(e)=S$, the preceding example no longer applies,
but Proposition~\ref{D:alpha-counterexample} still does.
\end{proof}

The two counterexamples settle the literal essential-surjectivity assertion.
We next give a sufficient reconstruction hypothesis, then distinguish
two precise meanings of the manuscript's undefined local filtering
condition. These distinctions also settle the earlier candidate
description of the image by piecewise-constant cartesian sections.

\begin{proposition}[Proved under stated hypotheses: reconstruction with products]
\label{D:alpha-products-image}
In the setting of 29/202--204, assume in addition that $C$ is closed
under nonempty finite products in $\mathscr R$: if $M,N\in C_T$,
there is an object $M\times N\in C_T$ whose image is
$\varphi(M)\times\varphi(N)$. Keep full faithfulness and the margin
hypothesis $\varphi(M)\twoheadrightarrow e_T$.
For every admissible sieve $C_0$, the functor $\alpha$ is an equivalence
onto the full subcategory of cartesian sections $Z$ for which every
$Z(M)\to s(M)$ is constant on its degree strata. This is an exact
necessary and sufficient condition on its essential image.
\end{proposition}
\begin{proof}
We first verify the elementary realization lemma. Over any object
$P\in\mathscr R_T$, realization is an equivalence between objects
constant on degree strata over $P$ and finite covers constant on
degree strata over $r(P)$. R2 lifts all clopen partitions uniquely.
R1 realizes finite constant objects and all maps of finite sets.
A map between finite constant objects is specified by a finite
clopen partition and a map of finite sets on every part, so this
also proves full faithfulness. R4 gives the covering by degree
strata, and R3 glues these objects and morphisms when the ranks are
unbounded. Realization commutes with pullback of such objects:
it preserves finite disjoint sums, and it preserves pullback of a
summand because it preserves that summand and its complement.

The required condition on $Z$ is necessary by the definition of
subordination. Conversely, lift its values and cartesian transition
isomorphisms uniquely to objects $A_M$ over $\varphi(M)$ and
isomorphisms
\[
 A_N\simeq \varphi(u)^*A_M\qquad(u:N\to M).
\]
An admissible sieve contains an object $M$ after an \'etale base
covering; we work over such a base. Since $C_0$ is a sieve, it
contains $M\times M$ and $M\times M\times M$. Its value at
$M\times M$ supplies an isomorphism between the two pullbacks
of $A_M$ to $\varphi(M)\times\varphi(M)$. The cartesian identities
at $M^3$ give the cocycle identity, and the diagonal gives the identity
condition. Epimorphisms in a multigaloisian category are effective
descent morphisms. We therefore obtain an object $X$ of $\mathscr R_T$
with $X\times\varphi(M)\simeq A_M$.

For any other $N\in C_0$, the object $M\times N$ gives compatible
identifications after pullback to $\varphi(M)\times\varphi(N)$.
This is an epimorphic cover of $\varphi(N)$, so descent identifies
$X\times\varphi(N)$ with $A_N$. Thus $X$ belongs to the required
source category and realizes $Z$. A morphism of cartesian sections
is treated by the same descent along $\varphi(M)\to e_T$;
compatibility at $M^2$ is exactly its descent condition. This proves
full faithfulness. The local objects $X$ and their unique transition
isomorphisms glue by the stack property of $\mathscr R$, completing
the proof over the original base.
\end{proof}

\begin{corollary}[Proved under stated hypotheses: finite realizations]
\label{D:alpha-products-finite}
With the product hypothesis of Proposition~\ref{D:alpha-products-image},
assume that every $s(M)\to T$ is finite and retain the generating
and common-refinement hypotheses. Then reconstruction holds after
the filtered colimit over admissible sieves:
\[
 \mathscr R_S\simeq
 \mathop{\mathrm{2\!\operatorname{colim}}}_{C_0}
          \operatorname{Cart}(C_0,\mathcal E).
\]
\end{corollary}
\begin{proof}
A finite \'etale cover of a finite $T$-scheme splits after an
\'etale covering of $T$. Indeed, over a strict henselization of
$T$ the finite scheme is a finite disjoint union of strictly
henselian local schemes. Every finite \'etale cover of each of
these is constant; the finite collection of idempotents and
trivializing isomorphisms descends to an \'etale neighborhood by
finite presentation.

Choose objects $M_i\in C_0$ over an \'etale covering $T_i\to S$,
and refine that covering to split the finitely presented covers
$Z(M_i)\to s(M_i)$. Restrict to the sieve generated by these
base-changed objects. Every value of the restricted cartesian
section is now a pullback of a split value, hence is constant on
degree strata. Proposition~\ref{D:alpha-products-image} supplies
its preimage. It also gives full faithfulness after intersecting
admissible sieves. Finally, the generating condition says exactly
that every object of $\mathscr R_S$ appears in one of the source
subcategories. These facts prove the displayed equivalence.
\end{proof}

\begin{proposition}[Proved under stated hypotheses: genuine local cofiltering]
\label{D:alpha-strict-filtering}
Suppose that local filtering means both common refinements and the
following parallel-arrow condition: for every pair $u,v:N\rightrightarrows M$
in $C_T$, an \'etale covering $T_i\to T$ and arrows $Q_i\to N_{T_i}$
exist which equalize $u$ and $v$. Under full faithfulness, generation,
and $\varphi(M)\twoheadrightarrow e_T$ for every $M$, every
$\varphi(M)$ is isomorphic to $e_T$. In particular the product hypothesis
of Proposition~\ref{D:alpha-products-image} holds, and its exact-image
conclusion and Corollary~\ref{D:alpha-products-finite} apply.
\end{proposition}
\begin{proof}
Fix $P=\varphi(M)$ and decompose $e_T$ into its rank strata $e_n$
for $P\to e_T$. The epimorphism hypothesis gives $e_0=\varnothing$.
Generation supplies, after an \'etale base covering, an object $N$
such that $P\times\varphi(N)\to\varphi(N)$ is constant on every
degree stratum. On its stratum of degree one choose the unique
section; on every stratum of degree $n\geq2$ choose two distinct
constant sections. The rank strata form an effective covering by
R3--R4, so these sections glue to two maps
$a,b:\varphi(N)\rightrightarrows P$ with equalizer
\[
 \operatorname{Eq}(a,b)=\varphi(N)\times e_1.
\]
Full faithfulness lifts $a,b$ to parallel arrows of $C$. The assumed
local equalization gives objects $Q_i$ whose maps to $\varphi(N)$
factor through this equalizer. Hence each map
$\varphi(Q_i)\to e_{T_i}$ factors through $(e_1)_{T_i}$.
It is an epimorphism, so $(e_1)_{T_i}=e_{T_i}$. Descent of this
equality gives $e_1=e_T$. Thus $P\to e_T$ has rank one and is
an isomorphism.

All images of objects of $C_T$ are therefore terminal; full
faithfulness makes every map between two such objects unique.
Whenever this fibre is nonempty, any one of its objects represents
the product of any nonempty finite family of its images. This is
exactly the additional hypothesis in the cited reconstruction
proposition. The assertion is local on the base, as required.
\end{proof}

\begin{proposition}[Refuted under the common-refinement convention]
\label{D:alpha-common-refinement}
If local filtering only supplies common refinements, then even with
finite realizations a cartesian section whose values are all
piecewise constant need not lie in the image of $\alpha$.
The failure can persist after every admissible sieve refinement.
Thus the earlier candidate exact-image theorem and its finite-realization
corollary are false under this convention without an additional hypothesis.
\end{proposition}
\begin{proof}
Let $S=\operatorname{Spec}\mathbb C$, put $H=C_2$, and let the fibre
over $S$ be
\[
 \mathscr R_S=\operatorname{FinSet}\times\operatorname{Fin}(H\text{-sets}),
 \qquad r(A,B)=A\amalg(B/H),
\]
where the finite sets on the right denote disjoint unions of copies
of $S$. Over an arbitrary \'etale $S$-scheme we take these categories
and constructions separately on its components. They form a stack.
Finite unions and quotients are preserved by orbit sets, giving R1.
Subobjects of an $H$-set are unions of orbits, giving R2. Covering
families of such subsets have effective descent, giving R3; and
the fibre cardinality of an equivariant map is constant on each
orbit, giving R4. These are multigaloisian categories.

Let $C$ be the full fibrewise subcategory on
$P=(\{*\},H)$, with $H$ acting regularly on the second component.
Its pullbacks are used over other base objects. Its endomorphisms
over $S$ form $H$. There is just one object, so any two objects
have a common refinement. The object $P$ maps epimorphically to
the final object, and every $(A,B)$ is subordinate to it: the first
component is already a finite constant set, while
$B\times H\to H$ is a trivial equivariant cover. Thus this is
a fully faithful generating system with finite realizations.

The scheme $s(P)$ has two points and every endomorphism of $P$
induces the identity on those points. A cartesian section over $C$
is consequently a pair of finite $H$-sets. The functor $\alpha$ is
\[
 (A,B)\longmapsto(A_{\mathrm{triv}},B).
\]
Indeed, $(B\times H)/H$ identifies with the underlying set of $B$;
right translation on the second factor induces its given $H$-action.
The two-element free $H$-set on the first point, together with a
one-element set on the second, is a cartesian section. Its underlying
covers of the two points are constant, but it is not in this image.
No equivariant bijection identifies a trivial two-element action
with the free action; the accompanying finite computation checks
this directly.

Every nonempty \'etale $S$-scheme has an $S$-section. Hence an
admissible sieve, by restriction along that section, contains $P$
over $S$, and then contains all its pullbacks. No admissible
refinement can remove the counterexample. This category is not
locally cofiltered in the stronger sense: the identity and the
nontrivial automorphism of $P$ cannot be equalized by any arrow
from its only object. The example therefore makes no claim against
Proposition~\ref{D:alpha-strict-filtering}.
\end{proof}

The transcription uses local filtering at 29/204 and 29/207 only
through common refinements; it gives no definition which decides
whether parallel arrows must also be equalized. We have kept the
two readings separate. In particular, a full category of nontrivial
Galois covers with all its automorphisms cannot satisfy the stronger
condition with all objects epimorphic over the final object.

\begin{proposition}[Proved: the root-stack interpretation of Example 2]
\label{D:root-stacks}
Let $D_1,\ldots,D_m$ be effective Cartier divisors on $S$, and let
$n_i$ be positive integers invertible on $S$, as imposed in 29/200--201.
The ramification datum of Example 2 has, over every \'etale $T\to S$,
the category of representable finite \'etale covers of the iterated
root stack
\[
 \mathcal X_T=
 \sqrt[n_1]{(T,D_{1,T})}\times_T\cdots\times_T
 \sqrt[n_m]{(T,D_{m,T})}.
\]
Its realization sends a cover to its coarse space over $T$.
The gerbe of roots of the line bundles $\mathcal O(D_i)$ alone is
not this root stack: the equations and the represented Kummer schemes
in the manuscript retain the divisor sections, which are essential.
\end{proposition}
\begin{proof}
On a base where $D_i=(a_i)$, set
\[
 W_a=\operatorname{Spec}_T
 \mathcal O_T[t_1,\ldots,t_m]/(t_i^{n_i}-a_i),\qquad
 H=\prod_i\mu_{n_i}.
\]
The root-stack definition as roots of a line bundle together with
its section gives $\mathcal X_T=[W_a/H]$; this is the explicit
local chart of Cadman \cite[Definition 2.2.4, Example 2.4.1]{Cadman2007}.
The map $W_a\to[W_a/H]$ is an \'etale surjective $H$-torsor
on the stack. Effective descent of finite \'etale morphisms identifies
finite \'etale covers of $[W_a/H]$ with $H$-equivariant finite
\'etale covers of $W_a$.

A change of equations $b_i=a_i\zeta_i^{n_i}$ acts by
$t_i\mapsto\zeta_i t_i$ on the Kummer charts. These are exactly
the arrows used in the manuscript's groupoid $C$ on 29/201.
Changing trivializations and composing such arrows have exactly the
same cocycle law in the two constructions. Their local equivalences
therefore glue to the stated equivalence of stacks of categories.
The colimit realization on 29/200 is locally the quotient of the
covering scheme by $H$, so is its coarse space. Such quotients exist
as schemes: they are finite over $T$, since the Kummer chart and the
cover of it are finite over $T$.

For the distinction about gerbes, over $S\setminus\bigcup_iD_i$
the root stack is $S$ itself, whereas the gerbe of line-bundle roots
has generic inertia $H$ whenever some $n_i>1$. Thus they cannot be
identified as stacks. Example 2 uses that gerbe together with its
Kummer realization, and the quotient of the realization is the
root stack just described.
\end{proof}


\begin{proposition}[Proved: a precise form of the unwritten Proposition 5]
\label{D:prop5}
Let $B\to S$ be quasi-compact and quasi-separated. Let $I$ be a
filtered index category and let $(Z_i,G_i)$ be an inverse system of
schemes with admissible actions of finite groups. Suppose that
$Z_i\to B$ is finite, $Z_i/G_i=B$, and each transition is equivariant
for a surjection $G_j\twoheadrightarrow G_i$, with
\[
 Z_i\simeq Z_j/\ker(G_j\to G_i).
\]
Over an \'etale $T\to S$, let $\mathscr R_i(T)$ consist of the
$G_i$-equivariant finite \'etale covers of $Z_i\times_ST$
which become piecewise constant after an \'etale covering of $T$.
Then the \'etale stackification of
$T\mapsto\mathop{\mathrm{2\!\operatorname{colim}}}_i\mathscr R_i(T)$,
with realization induced by quotient by $G_i$, is a ramification datum
satisfying R1--R4. Transition functors are fully faithful and exact.
If the original datum is generated \'etale locally by the corresponding
global Galois objects, this construction reconstructs it.
\end{proposition}
\begin{proof}
For a transition, put $H=\ker(G_j\to G_i)$. An object $W_i\to Z_i$
pulls back to $W_j=W_i\times_{Z_i}Z_j$. Since finite-group
quotients commute with flat base change and $W_i\to Z_i$ is
\'etale, $W_j/H=W_i$. An equivariant map between two such pullbacks
descends uniquely to the $H$-quotients. This proves full faithfulness.
Finite limits, finite sums and finite equivalence-relation quotients
of equivariant finite \'etale covers commute with pullback, proving
exactness. Iterated quotient also gives
$W_j/G_j=W_i/G_i$, so realizations agree.

At one level, this is the construction of 29/196, Proposition 4.
Here are direct checks of its axioms. Finite sums and quotient by an
equivalence relation commute with the final group quotient, giving R1.
An invariant clopen part of $W_i$ is exactly the inverse image of a
clopen part of $W_i/G_i$; this gives R2. Covers by such parts have
effective descent for finite \'etale covers, giving R3. Ranks of
finite \'etale covers are locally constant, and their rank strata
cover; hence R4. The condition of being piecewise constant after
an \'etale base covering is stable under the operations just used
and under descent.

Finite diagrams in a filtered colimit occur at a common level, so
R1, R2 and the finite categorical axioms pass to the colimit.
For R3, work over an affine \'etale base $T$. Every realization is
finite over the quasi-compact scheme $B_T$, so any cover of it by
clopen parts has a finite subcover. The finitely many objects, maps
and descent identities then occur at one common index. Effectiveness
at that index proves the claim. The same quasi-compactness bounds
the finitely many ranks needed for R4. Finally \'etale stackification
glues the constructions; realizations glue as finite $B_T$-schemes,
for which effective \'etale descent holds. This verifies all four
axioms for the resulting datum.

For reconstruction, every object of the original datum belongs,
\'etale locally on the base, to the subcategory locally subordinate
to one of the stated Galois objects. The equivariant description of
that subcategory on 29/195 is precisely $\mathscr R_i$. Its proof
is ordinary descent of an object made piecewise constant by the
Galois object, together with R1--R4: clopen partitions lift uniquely by R2,
finite constant objects lift by R1, and R3 glues their degree strata. The fully faithful local equivalences
therefore glue. Without this local generation assumption one only
obtains the generated subdatum, exactly as the warning after
``Prop. 5'' on 29/197 says.
\end{proof}

The manuscript writes only ``Prop. 5 ...'' at this point; there is no
complete assertion whose missing hypotheses can be recovered uniquely.
The natural assertion that the preceding filtered system, under just
the conditions actually written on 29/196--197, always defines a
ramification datum is false. In particular, connectedness of the
common quotient is not a substitute for the hypothesis used in
Proposition~\ref{D:prop5}: connectedness can be lost after restriction
to an \'etale open of the base.

\begin{example}[Refuted: the unqualified filtered construction behind Proposition 5]
\label{D:prop5-counterexample}
There is a quasi-compact noetherian base $S$ and an inverse system
$(Z_i,G_i)$ with every $Z_i$ and $Z_i/G_i$ connected, admissible
finite-group actions, surjective transition homomorphisms, and
$Z_i=Z_j/\ker(G_j\to G_i)$, for which the \'etale stackification
of the filtered categories in Proposition~\ref{D:prop5} fails R3.
The obstructing descent datum uses only rank-two objects.
\end{example}
\begin{proof}
Let $k$ be an algebraically closed field, and put
$S=\operatorname{Spec}k[t]$. We first construct a connected
non-quasi-compact scheme $B\to S$. For each $m\in\mathbb Z$ take
\[
 Q_m=\operatorname{Spec}
 k[u_m,v_m,(u_m-1)^{-1},(v_m-1)^{-1}]/(u_mv_m),
 \qquad t=u_m(u_m-1)+v_m(v_m-1).
\]
Glue $D(v_m)\subset Q_m$ to $D(u_{m+1})\subset Q_{m+1}$ by
$u_{m+1}=1-v_m$. The displayed formulas for $t$ agree. These
are open gluings with no triple intersections, so they define a
scheme. Each $Q_m$ is connected and meets $Q_{m+1}$; thus $B$ is
connected. Geometrically it is an infinite chain of affine lines,
with the point $1$ of one line identified with the point $0$ of the
next. One can also check separatedness directly: the overlap-ring
map from the tensor product of two adjacent chart rings is
surjective, since the two chart rings together provide
$v_m$, $v_m^{-1}$ and $(v_m-1)^{-1}$.

For $T=D(t)\subset S$, all nodes are removed and
\[
 B_T=\coprod_{n\geq1}E_n,
 \qquad E_n\simeq\operatorname{Spec}k[x,x^{-1},(x-1)^{-1}],
 \qquad t=x(x-1),
\]
where we choose a bijection from the integer-indexed components to
the positive integers. Each $E_n\to T$ is finite flat surjective
of degree two: its algebra is given by the monic equation
$x^2-x-t=0$, and $t$ invertible implies both $x$ and $x-1$
invertible.

Set $G_i=(\mathbb Z/2)^i$ and $Z_i=B$, with trivial actions.
The transition maps of schemes are identities and the homomorphisms
$G_j\to G_i$ are coordinate projections. All actions are admissible;
the stated connectedness and quotient conditions are immediate.
In the category over $T$, the clopen components $E_n$ are subobjects
of the terminal object already at level zero, and their realizations
cover $B_T$. On the subobject $E_n$ take the constant two-sheet
cover on which $G_n$ acts by its $n$th-coordinate character.
Regarded as an object over $B_T$, it has rank two on $E_n$ and
rank zero elsewhere, hence is piecewise constant. On all pairwise
intersections the objects are empty, so they carry a descent datum.

There is no glue at a common finite level $i$. Choose $n>i$ and
pass to a level at least $n$ to compare the proposed restriction
with the prescribed object. The $n$th-coordinate involution acts
trivially on every object inflated from level $i$, and nontrivially
on the prescribed two-sheet cover. No equivariant isomorphism can
exist. Nor can \'etale stackification on $T$ repair this: after any
nonempty \'etale base change $T'\to T$, every $E_n\times_TT'$
is still nonempty by finite faithful flatness. Any locally chosen
finite level therefore has the identical contradiction with some
$n$ larger than that level. This disproves effectiveness of the
disjoint clopen descent datum required by R3.

This example satisfies the connectedness and global bounded-rank
conditions discussed on 29/196--197. Those leaves impose
quasi-compactness on $S$, but no such hypothesis on $r(e)=B$ or the
$Z_i$; 29/191 explicitly allows realizations in all schemes over
the base. Thus the additional condition in
Proposition~\ref{D:prop5} is substantive. An alternative construction
would have to complete the filtered categories also for these
clopen descent data, rather than only for \'etale descent on $S$.
\end{proof}


\begin{proposition}[Proved: the general construction with clopen descent]
\label{D:prop5-completion}
For an arbitrary inverse system of admissible finite-group actions
$(Z_i,G_i)$ with common quotient $B$, surjective group transitions,
and $Z_i=Z_j/\ker(G_j\to G_i)$, the following construction always
produces a ramification datum satisfying R1--R4. Start with the
finite-level categories of equivariant finite \'etale covers which
are locally piecewise constant as in Proposition~\ref{D:prop5}.
First complete their filtered union for descent under arbitrary
covers by clopen subsets of $B_T$, and then for \'etale descent
in $T\to S$ (equivalently, use the topology generated by both
kinds of covers). Realization is defined by gluing the finite-level
quotient schemes.

This is universal among ramification data receiving compatible
exact functors from the finite-level categories and preserving
their realizations. The ordinary \'etale stackified filtered union
embeds fully faithfully into this completion. It need not be
essentially surjective, by Example~\ref{D:prop5-counterexample}.
The completion is not asserted to preserve the additional global
bounded-rank condition preceding Proposition 5 in the manuscript.
\end{proposition}
\begin{proof}
We give the descent argument, including the point that replaces
quasi-compactness. At a fixed finite level an object is represented
by a $G_i$-equivariant finite \'etale map $W_i\to Z_i$, and its
realization is $W_i/G_i$. The structural map
\[
 q:W_i/G_i\longrightarrow Z_i/G_i=B
\]
is open and closed and has finite geometric fibres. To see the
first assertion, pull an open or closed subset of $W_i/G_i$ back
to $W_i$, take its image under the finite \'etale map to $Z_i$,
and then take the finite-group quotient. Finite \'etale maps are
open and closed; the quotient topology and saturation by the finite
group preserve both properties. For a geometric point of $B$,
the points above it in $Z_i$ form a single finite orbit. Every
orbit above it in $W_i$ meets the finite fibre over any one of
these points, proving finiteness. We do not assume that $q$ is a
finite morphism for arbitrary rings.

Descent for clopen base covers is ordinary stackification on the
site of clopen subsets of $B_T$. It glues the quotient realizations
as schemes along open-and-closed pieces. Finite diagrams and their
identities can be put at one common index after refining such a
cover. Thus finite limits, sums, equivalence-relation quotients,
and their exactness properties are inherited from the finite-level
categories. This proves the multigaloisian finite categorical axioms
and R1. Clopen parts of a realization lift uniquely at each level,
by the invariant-clopen description. They glue uniquely, giving R2.

For R3, it suffices locally on a clopen base cover to start with an
object having a fixed-level realization $W$. Suppose its clopen
subobjects have realizations $V_\alpha$ covering $W$. For each
geometric point $b$ of $B_T$, finitely many of the $V_\alpha$,
say those indexed by a finite set $J$, cover $q^{-1}(b)$. The set
\[
 B_J=B_T\setminus q\left(W\setminus\bigcup_{\alpha\in J}
 V_\alpha\right)
\]
is a clopen neighborhood of $b$. The $B_J$ cover $B_T$. Over each
one only finitely many pieces are needed, so the objects, maps and
cocycle identities involved in descent occur at a common finite
index after a further clopen refinement. Finite clopen descent
there is effective: one refines a finite clopen cover to a finite
disjoint partition and glues the finite \'etale covers and actions.
The imposed clopen descent then glues these local solutions.
This proves R3 for arbitrary, possibly infinite, clopen coverings
of arbitrary objects, not merely of the final object. Rank strata
and their covering property are checked at each finite level and
are local for these covers, proving R4. For completeness, the
finite-degree axiom is also inherited: a connected target has
connected realization by R2, so exactly one of these clopen rank
strata is nonempty. Its finite rank bounds the number of disjoint
sections after every base change, which is G5. Thus the resulting
fibres are multigaloisian, not merely categories with finite exact
operations. \'Etale stackification in
$T$ preserves these verifications and glues the scheme realizations.

Maps between two fixed-level objects are already sheaves for clopen
covers. Any map that appears only after a later index descends
uniquely by the full faithfulness proved in Proposition~\ref{D:prop5}.
Consequently the original union embeds fully faithfully. Finally,
in any receiving ramification datum, R2 uniquely lifts the clopen
base pieces and R3 glues their objects and morphisms; its \'etale
stack condition supplies the remaining gluing. These operations
force the extension functor and all its compatibilities uniquely.
This proves the universal property.
\end{proof}

\subsection{D.8: G5 and the proposed Hom-finiteness repair}

\begin{proposition}[Refuted: G5 is not redundant, 29/212]
\label{D:G5-Zsets}
The category of countable sets with a $\mathbb Z$-action satisfies
G0--G4 as written on 29/210 and fails both G5 and the finite-fibre
functor condition G5$'$.
\end{proposition}
\begin{proof}
We take an essentially small skeleton. Finite limits, finite sums,
and quotients of equivalence relations are computed on underlying sets
and remain countable. Sums are disjoint and universal. Every map
factors as a surjection followed by the inclusion of its image, an
invariant summand. Surjections are effective descent morphisms:
ordinary set descent produces an equivariant quotient, which is still
countable, and the action and all descent identities descend uniquely.
The one-point terminal object is connected. These verify G0--G4,
including effective descent rather than merely effectiveness of an
epimorphism.

Let $A=\mathbb Z$ with the translation action. For the diagonal action,
$A\times A$ is the disjoint union of the invariant subsets
$\{(a,b):a-b=j\}$, $j\in\mathbb Z$, each isomorphic over the
second factor to $A$. Thus the base change of $A\to1$ by $A\to1$
contains arbitrarily many copies of $A$ as a summand, contradicting G5.

For G5$'$, suppose that an exact functor $F$ to finite sets has
$F(1)=1$. Exactness sends the effective epimorphism $A\to1$ to a
surjection, so $F(A)$ is nonempty. Let $\tau$ be translation by one.
Its action on the finite set $F(A)$ has a positive finite order $m$.
The equalizer of $\tau^m$ and the identity in our category is empty,
whereas their equalizer after applying $F$ is $F(A)$. This contradicts
preservation of finite limits and the initial object.
\end{proof}

\begin{lemma}[Proved: finite-support permutation orbits]\label{D:nominal-orbits}
Let $A$ be a countably infinite set and $P$ its group of permutations
with finite support. A $P$-set is called nominal if each of its
elements is fixed by the pointwise stabilizer $P_B$ of some finite
$B\subset A$. Every transitive nominal $P$-set is isomorphic to
$A^{(n)}/H$ for some $n\geq0$ and subgroup $H\subset\mathfrak S_n$,
where $A^{(n)}$ is the set of ordered $n$-tuples of distinct elements.
If a nominal $P$-set has finitely many orbits, its subset fixed by
$P_B$ is finite for every finite $B$.
\end{lemma}
\begin{proof}
For finite $B,D\subset A$, the subgroups $P_B$ and $P_D$ generate
$P_{B\cap D}$. To verify this it suffices to generate each transposition
outside $B\cap D$. A transposition whose two entries avoid $B$ or
avoid $D$ already lies in one subgroup. In the remaining case its
entries are $b\in B\setminus D$ and $d\in D\setminus B$; choose
$c\notin B\cup D$ and use
$(b\ d)=(b\ c)(c\ d)(b\ c)$. Thus the intersection of two supports
of an element is again a support. A support of minimum cardinality
is consequently its unique least support.

If an element $x$ has least support $B$, its stabilizer contains $P_B$
and preserves $B$ setwise: the least support of $gx=x$ is $gB=B$.
Its stabilizer modulo $P_B$ is therefore a subgroup $H$ of the finite
permutation group of $B$. This gives the asserted orbit description.
Finally an element fixed by $P_D$ has least support contained in $D$.
In an orbit $A^{(n)}/H$, there are at most $|D|!/(|D|-n)!$ such
elements, with zero understood when $n>|D|$. A finite union of orbits
therefore has only finitely many fixed elements.
\end{proof}

\begin{proposition}[Refuted: lead L10]\label{D:G5-Hom-finite}
Even the conjunction of G0--G4, finiteness of all Hom sets, and a finite
decomposition of every object into connected objects does not imply G5.
\end{proposition}
\begin{proof}
Let $\mathcal N$ be the category of nominal $P$-sets with finitely
many orbits and equivariant maps. The preceding lemma shows that it
has an essentially small skeleton. Finite products have finitely many
orbits: for two tuple orbits, the equality pattern among their finitely
many coordinates has only finitely many possibilities. Passing to
finite-group quotients does not increase the number of orbits.
Equalizers and images are invariant subsets, hence unions of some of
the finitely many orbits; quotients remain nominal and have finitely
many orbits. Thus finite limits and finite sums exist, with disjoint
universal sums. An equivariant map factors through its invariant image,
which is a summand.

Every surjection is effective descent inside $\mathcal N$. Indeed,
set-theoretic descent identifies the desired object with a quotient
of the given object upstairs. The equivalence relation is invariant,
so the group action descends. A quotient of an orbit-finite nominal
set is again orbit-finite and nominal. The ordinary equivariant
descent equivalence therefore restricts to this category. The terminal
one-point set is connected; the orbits are precisely the nonempty
connected objects. This proves G0--G4 and the finite connected
decomposition assertion.

An equivariant map from a transitive object is determined by the image
of one element. If that element has finite support $B$, its image lies
in the finite fixed-point set of $P_B$ in the target. There are therefore
only finitely many such maps; finite orbit decompositions prove finiteness
of every Hom set.

Nevertheless, take $X=A$, $Y=1$ and $Y'=A^{(m)}$. These are nonempty
connected objects. For $1\leq i\leq m$ the coordinate map
$A^{(m)}\to A$ defines a section of
$A\times A^{(m)}\to A^{(m)}$. The images of these sections are
pairwise disjoint, and their union is a summand: its complement consists
of pairs $(a,(a_1,\ldots,a_m))$ with $a\notin\{a_1,\ldots,a_m\}$,
an object isomorphic to $A^{(m+1)}$. Thus the splitting multiplicities
are unbounded, and G5 fails.

The category used here is a standard example in the theory of
pre-Galois categories; see Harman and Snowden \cite{HarmanSnowden}.
We have supplied the full elementary verification of the particular
axioms required here. The finite symmetric-group checks in the code
are sanity checks of the coordinate-section formulas only, not a
replacement for the infinite argument.
\end{proof}

\begin{proposition}[Proved: an exact uniform-Hom replacement]
\label{D:G5-uniform-Hom}
Assume G0--G4 and also the following connected-probe condition:
every noninitial object receives a morphism from a noninitial connected
object. Then G5 is equivalent to this condition:
for every $d:X\to Y$ between noninitial connected objects, there is
an integer $N(d)$ such that
\[
 |\operatorname{Hom}_Y(T,X)|\leq N(d)
\]
for every noninitial connected $T$ equipped with a morphism $T\to Y$.
In particular the uniform condition implies G5. Mere finiteness of
each displayed set cannot replace its uniform bound.
\end{proposition}
\begin{proof}
First, G3 implies that every monomorphism is an isomorphism onto a
summand: its effective-descent factor is also a monomorphism and hence
is an isomorphism. Accordingly, the equalizer of two maps from a
connected object is either the entire source or initial. Distinct maps
$T\to X$ over $Y$ have pairwise disjoint graphs in $X\times_YT$.
Finite unions of these graphs are monomorphisms, by disjointness and
universality of finite sums, hence are summands by G3.

If G5 holds with bound $N$, more than $N$ distinct maps would therefore
produce $N+1$ copies of $T$ as a summand after the base change $T\to Y$,
a contradiction. This proves the uniform bound. Conversely, suppose
$X\times_YY'$ contains $N+1$ copies of a noninitial $Y'$ as a summand.
Choose a noninitial connected probe $T\to Y'$. Universal disjointness
preserves those copies under base change, yielding $N+1$ disjoint
sections over $T$ and hence $N+1$ distinct elements of
$\operatorname{Hom}_Y(T,X)$. This contradicts the uniform bound.
The nominal-set example has connected probes (indeed finite orbit
decompositions) and finite individual Hom sets, proving the final claim.
\end{proof}

\subsection{D.9: wild inertia for units and general linear groups}

\begin{proposition}[Refuted: the wild unit comparison of 31/044]
\label{D:wild-units}
For every prime $p$ there is a complete discrete valuation ring $O$
of residue characteristic $p$ and a totally ramified action of
$G=C_p$ such that
\[
 H^1(G,O^\times)\cong\mathbb Z/p\mathbb Z,
 \qquad H^1(G,(O/\mathfrak m)^\times)=0.
\]
Thus the indicated restriction-reduction map is not injective and
is not a bijection.
\end{proposition}
\begin{proof}
Take an algebraically closed field $k$ of characteristic $p$,
$O=k[[t]]$, $L=k((t))$, and
$\sigma(t)=t/(1+t)$, with $\sigma|_k=1$.
Inductively $\sigma^j(t)=t/(1+jt)$, so $\sigma$ has exact order $p$.
The norm
\[
 u=\prod_{j\in\mathbb F_p}\frac{t}{1+jt}
   =\frac{t^p}{1-t^{p-1}}
\]
is invariant and has $t$-adic valuation $p$. The extension
$k((t))/k((u))$ has degree $p$: its ramification index is at least
$p$, while $t$ satisfies $t^p+ut^{p-1}-u=0$. Hence the fixed field
is $K=k((u))$ and its value group inside $\mathbb Z=v_t(L^\times)$
is $p\mathbb Z$. The extension is totally ramified, so $G$ is the
inertia group.

Apply group cohomology to
$1\to O^\times\to L^\times\xrightarrow{v_t}\mathbb Z\to0$.
The valuation group has trivial $G$-action and Hilbert 90 gives
$H^1(G,L^\times)=0$. Exactness identifies $H^1(G,O^\times)$ with
$\mathbb Z/v_t(K^\times)=\mathbb Z/p\mathbb Z$. The class of one
is represented by the cocycle $\sigma^j\mapsto\sigma^j(t)/t$.
Its reduction is trivial. On the residue field the action is trivial,
and $H^1(C_p,k^\times)=\operatorname{Hom}(C_p,k^\times)=0$ because
in characteristic $p$ the only $p$-th root of unity is one.
\end{proof}

\begin{proposition}[Refuted for every rank: the wild $\mathrm{GL}_n$ comparison]
\label{D:wild-GLn}
In the same example, for every integer $n\geq1$ the map of pointed sets
\[
 H^1(G,\mathrm{GL}_n(O))\longrightarrow
 H^1(G,\mathrm{GL}_n(k))
\]
is not injective. It is surjective in this example.
\end{proposition}
\begin{proof}
For $0\leq a<p$ use the cocycles
\[
 c_a(g)=\operatorname{diag}(g(t^a)/t^a,1,\ldots,1).
\]
All reduce to the identity cocycle. Their determinants represent the
distinct classes $a\in H^1(G,O^\times)=\mathbb Z/p\mathbb Z$,
so they cannot be cohomologous as $\mathrm{GL}_n$-cocycles.
For surjectivity, a cocycle on $\mathrm{GL}_n(k)$ is a representation
of $C_p$, since the residue action is trivial. Its matrices, viewed
as constant matrices over the invariant coefficient field $k\subset O$,
are a cocycle upstairs and lift the given class.
\end{proof}

\begin{proposition}[Proved: the tame totally ramified replacement]
\label{D:tame-GLn}
Let $O=k[[t]]$, let $e$ be prime to $\operatorname{char}k$, and assume
that $k$ contains a primitive $e$-th root of unity $\zeta$. Let
$G=C_e$ act by $\sigma(t)=\zeta t$ and trivially on $k$.
For every $n\geq1$, reduction induces a bijection
\[
 H^1(G,\mathrm{GL}_n(O))\xrightarrow{\ \sim\ }
 H^1(G,\mathrm{GL}_n(k)).
\]
Both sides are the set of unordered $n$-tuples of residues in
$\mathbb Z/e\mathbb Z$; for $n=1$ the group is $\mathbb Z/e\mathbb Z$.
\end{proposition}
\begin{proof}
A cocycle defines a semilinear $G$-action on the free module $O^n$.
Its reduction is an ordinary representation over $k$. Since $e$ is
invertible and $k$ contains all characters, choose an eigenbasis of
the reduced module. Lift each vector arbitrarily and apply the
idempotent projector
\[
 \frac1e\sum_{j=0}^{e-1}\zeta^{-aj}\sigma^j
\]
for the prescribed character $\zeta^a$. This is well defined on the
underlying $k[[t^e]]$-module and produces an eigenvector lifting the
chosen reduced one. The lifts form an $O$-basis by Nakayama's lemma.
Thus the cocycle is cohomologous to the corresponding constant diagonal
character matrix. Two such cocycles are cohomologous only if their
reductions are conjugate, and the multiplicities of the distinct
characters classify the latter representations. This proves both
injectivity and surjectivity and the stated classification.
\end{proof}

\subsection{D.10: an exact boundary criterion and its numerical curve form}

We use abelian coefficient sheaves, as required for the higher direct
images in the question. The following criterion applies whenever every
quasi-compact open of the base has finitely many irreducible components;
in particular it applies on every locally noetherian base. It does not
assume that direct image commutes with a change of base; that distinction
is necessary in the wild case. For a completely arbitrary base, Proposition~\ref{D:boundary-arbitrary}
adds a necessary and sufficient constructibility condition on boundary
cohomology. Example~\ref{D:boundary-nonconstructible} shows why that
condition cannot be omitted.

\begin{proposition}[Proved: the boundary vanishing-cycle criterion]
\label{D:boundary-criterion}
Let $f:X\to S$ be proper and smooth. Assume that every quasi-compact
open of $S$ has finitely many irreducible components. Let $Y$ be a
relative normal-crossings divisor, let $j:V=X\setminus Y
\hookrightarrow X$, and put $h=fj$. Let $G$ be a locally constant
constructible sheaf of finite abelian groups annihilated by an integer
invertible on $S$. Put $K=Rj_*G$.

For each geometric point $\bar s$ of $S$ and each geometric point
$\bar t$ of $S_{(\bar s)}$, pull $K$ back to $X_{S_{(\bar s)}}$,
and define
\[
 B_{\bar s\leftarrow\bar t}(G)=
 \operatorname{Cone}\bigl(K|_{X_{\bar s}}\longrightarrow
       i_{\bar s}^{*}Rj_{\bar t*}(K|_{X_{\bar t}})\bigr).
\]
Here $j_{\bar t}:X_{\bar t}\to X_{S_{(\bar s)}}$ and
$i_{\bar s}:X_{\bar s}\to X_{S_{(\bar s)}}$ are the natural maps;
the arrow is the specialization arrow. This boundary complex is
supported on $Y_{\bar s}$. Then all $R^i h_*G$ are locally constant
if and only if
\begin{equation}\label{D:boundary-condition}
 R\Gamma\bigl(Y_{\bar s}, B_{\bar s\leftarrow\bar t}(G)|_{Y_{\bar s}}\bigr)=0
 \quad\hbox{for every }(\bar s,\bar t).
\end{equation}
If the relative dimension is one, this is equivalent to the vanishing
of the boundary complex at each individual point of $Y_{\bar s}$.
Thus in relative dimension one the criterion is precisely local
acyclicity of $Rj_*G$ along the boundary. In higher dimension the
hypercohomological condition, rather than pointwise vanishing, is the
necessary and sufficient condition.
\end{proposition}
\begin{proof}
Proper base change for bounded-below torsion complexes, valid over
arbitrary schemes \cite[Lemma 59.91.12, tag 0DDE]{StacksDGeneral},
and composition of direct images identify the
triangle of global sections of the displayed cone with
\[
 (Rh_*G)_{\bar s}\longrightarrow (Rh_*G)_{\bar t}
 \longrightarrow R\Gamma(X_{\bar s},B_{\bar s\leftarrow\bar t}(G))
 \longrightarrow.
\]
Indeed, proper base change identifies the first term with
$R\Gamma(X_{\bar s},K|_{X_{\bar s}})$. For the second term,
commute proper direct image by $f$ with pullback to $\bar s$
and use $f\circ j_{\bar t}=\bar t\to S_{(\bar s)}$ after
proper direct image over $\bar t$. Global sections over a strictly
local scheme compute its closed geometric stalk. This gives exactly
the usual specialization map, not a proposed identification with
$H^i(V_{\bar s},G)$. See also the specialization formula in
\cite[Lemma 59.75.7]{StacksSpecialization}; over a trait it is the
standard triangle of \cite[XIII, (2.1.8.9)]{SGA7II}.

On $V$, the complex $K$ is $G$ and $h$ is smooth. Smooth local
acyclicity makes the cone zero there. It is consequently supported
on $Y_{\bar s}$, and its global sections are the expression in
\eqref{D:boundary-condition}. Thus that condition is equivalent to
all specialization maps of every $R^i h_*G$ being isomorphisms.
Their stalks are then finite: specialize from a geometric generic
point of an irreducible component through the given point. At such
a generic point, the stalk is the finite cohomology of a finite-type
scheme over a separably closed field, with invertible finite
coefficients. For clarity, no noetherian hypothesis is needed in this
generic-point argument. The strict localization at a geometric generic
point has zero-dimensional underlying space, and its reduction is the
spectrum of a separably closed field. The reduction map is a universal
homeomorphism even when its nilradical is not nilpotent. Topological
invariance of \'etale cohomology
\cite[Proposition 59.45.4, tag 03SI]{StacksDGeneral} therefore reduces
the stalk to the same finite field cohomology.
The finite-stalk specialization criterion
\cite[Lemmas 59.75.5--6]{StacksSpecialization} now proves local
constancy. Conversely, a locally constant sheaf has isomorphic
specialization maps, giving the necessity.

For a relative curve, $Y\to S$ is finite \'etale, locally a
finite disjoint union of sections. A complex supported on the finite
set $Y_{\bar s}$ has hypercohomology equal to the direct sum of
its stalk complexes. It vanishes if and only if every such complex
vanishes. This proves the curve assertion.
\end{proof}


\begin{proposition}[Proved: the exact criterion over an arbitrary base]
\label{D:boundary-arbitrary}
Keep $f,X,Y,j,h,G$ as above, but let $S$ be an arbitrary scheme.
Write $i:Y\hookrightarrow X$ and $g=f i$, and define the boundary
cohomology complex
\[
 C_Y(G)=Rg_*i^*Rj_*G.
\]
Then all $R^q h_*G$ are locally constant if and only if the following
two conditions hold:
\begin{enumerate}
\item Every cohomology sheaf $\mathcal H^q C_Y(G)$ is constructible
(locally on $S$).
\item The boundary hypercohomology in
\eqref{D:boundary-condition} vanishes for every geometric
specialization $\bar s\leftarrow\bar t$.
\end{enumerate}
In relative dimension one the second condition is pointwise vanishing
on the finite geometric boundary. The first condition remains necessary
on an arbitrary base, even for a relative curve.
\end{proposition}
\begin{proof}
Localization for the open--closed decomposition of $X$ gives a triangle
\[
 Rh_!G\longrightarrow Rh_*G\longrightarrow C_Y(G)\longrightarrow.
\]
The functor $Rh_!$ has constructible cohomology on an arbitrary base:
$h$ is compactifiable and of finite presentation, and
\cite[XVII, Theorem 5.3.6, p.~364]{SGA4D} applies to the finite
coefficient ring $\mathbb Z/n$. This theorem explicitly reduces the
arbitrary affine base to a noetherian one by limits; it does not impose
noetherianity on $S$. Its cohomology is locally bounded, by
\cite[XVII, Corollary 5.2.8.1, p.~358]{SGA4D}.
Constructible finite-coefficient sheaves form an abelian category
closed under extensions, locally on every scheme. The long exact
cohomology sequence therefore proves
\[
 \bigl(\mathcal H^q C_Y(G)\text{ constructible for every }q\bigr)
 \ \Longleftrightarrow\
 \bigl(R^q h_*G\text{ constructible for every }q\bigr).
\]
No boundedness assertion about $Rh_*G$ is needed: each term is governed
by a finite segment of this long exact sequence.

The proper-base-change calculation in the preceding proof did not
use its finiteness hypothesis on irreducible components. Thus condition
(2) is equivalent to all specialization maps of every $R^q h_*G$
being isomorphisms. For a constructible sheaf this is equivalent to
local constancy on an arbitrary scheme, by
\cite[IX, Proposition 2.11, p.~20]{SGA4D}.
We spell out the relevant finiteness step. At a geometric point,
finitely many generators and relations spread to an \'etale
neighborhood, giving a map from the constant finite module with the
prescribed stalk to the sheaf, which is an isomorphism at that point.
Its kernel and cokernel are constructible, so their zero loci, and
hence the isomorphism locus, are locally constructible. Because both
sheaves have isomorphic specialization maps, that locus is stable
under generization. A locally constructible generization-stable
subset of a scheme is open. Restricting to this open proves local
constancy. This proves sufficiency.

Conversely, locally constant $R^q h_*G$ have isomorphic specialization
maps. Their stalks are finite by specialization to a geometric generic
point, followed by the reduction and topological-invariance argument
in the preceding proof. They are therefore constructible. The same
localization triangle then proves constructibility of the boundary
cohomology sheaves. This proves necessity.
\end{proof}

\begin{example}[Refuted: specialization alone over an arbitrary base]
\label{D:boundary-nonconstructible}
There exist an affine characteristic-two base $S$, the smooth proper
relative curve $X=\mathbb P^1_S$, the boundary section $Y=\{\infty\}$,
and a rank-one locally constant $\mathbb F_3$-sheaf $G$ on
$\mathbb A^1_S$ such that all the boundary comparison complexes in
\eqref{D:boundary-condition} vanish, while $h_*G$ is not locally
constant. All $R^q h_*G$ for $q>0$ are zero.
\end{example}
\begin{proof}
Let $k=\overline{\mathbb F}_2$, let $K$ be an algebraic closure of
$k((t))$, and let $V$ be the valuation ring of an extension of the
$t$-adic valuation to $K$. It is a strictly henselian rank-one
valuation ring, with residue field $k$, fraction field $K$, and
exactly two prime ideals. Let $A$ be the ring of eventually constant
$k$-valued sequences indexed by the positive integers, equipped with
the eventual-value map $A\to k$. Put
\[
 R=V\times_k A,\qquad S=\operatorname{Spec}R,
 \qquad a=(t,(0,0,\ldots))\in R.
\]
The finite-support sequences form an ideal $I$ of $A$. Each coordinate
idempotent $e_n=(0,\delta_n)$ gives an isolated clopen point $n$ of
$S$ with residue field $k$. Every other prime contains all these
idempotents, hence contains $0\times I$ and corresponds to a prime of
$R/(0\times I)=V$. Thus the remaining points are a generic point
$\eta$ and its specialization $m$. Their local rings are $K$ and
$V$, respectively: localization at $m$ kills every $e_n$ since
$1-e_n$ is a unit there. Every neighborhood of $m$ contains all but
finitely many isolated points $n$, since a sequence whose eventual
value is nonzero is nonzero at all but finitely many coordinates.
Also $D(a)=\operatorname{Spec}K$.

Write $\Lambda=\mathbb F_3$, and let
$G=\mathcal L_\psi(ax)$ be the sign summand of the finite \'etale
Artin--Schreier cover $z^2-z=ax$. For $i_0:V(a)\hookrightarrow S$
and $j_0:D(a)\hookrightarrow S$, we claim
\begin{equation}\label{D:AS-duality}
 Rh_*G \simeq i_{0*}Ri_0^!\Lambda
 \simeq \operatorname{Cone}(\Lambda_S\longrightarrow
 Rj_{0*}\Lambda_K)[-1].
\end{equation}
Here no noetherian approximation of ordinary direct image is asserted.
To prove the formula, first use the universal family
$\mathcal L_\psi(cx)$ on $\mathbb A^1_x\times\mathbb A^1_c$.
Its compactly supported cohomology on every fibre $c\ne0$ is zero.
Indeed the Artin--Schreier cover is itself an affine line, whose
compactly supported cohomology is $\Lambda(-1)[-2]$, and translation
by the deck involution acts trivially on that top cohomology; its
sign summand is zero. At $c=0$ the sheaf is constant, with compact
cohomology $\Lambda(-1)[-2]$. Proper-support base change and localization
therefore identify the universal direct image with
$0_*\Lambda(-1)[-2]$. Arbitrary base change
\cite[XVII, Theorem 5.2.6, p.~357]{SGA4D} now gives
\[
 Rh_!G^\vee\simeq i_{0*}\Lambda(-1)[-2].
\]
Poincar\'e duality for a smooth compactifiable morphism over a
quasi-compact quasi-separated base is valid without noetherianity:
\cite[XVIII, Theorem 3.2.5, pp.~585--586]{SGA4D} identifies the right
adjoint of $Rh_!$ with $h^*(1)[2]$. Internal Hom adjunction, checked
\'etale locally on the base, consequently gives
\[
 R\mathcal Hom_S(Rh_!G^\vee,\Lambda)
 \simeq Rh_*G(1)[2]
 \simeq i_{0*}Ri_0^!\Lambda(1)[2].
\]
Cancelling the twist and shift proves the first isomorphism in
\eqref{D:AS-duality}; localization proves the second.

At an isolated point $n$ the latter complex is $\Lambda$ in degree
zero, since $D(a)$ is empty in that local scheme. At $m$ the strict
localization is $V$, and its intersection with $D(a)$ is
$\operatorname{Spec}K$. Both $V$ and $K$ have constant-sheaf
cohomology $\Lambda$ in degree zero and zero in higher degrees,
and the restriction map is the identity. Thus the complex has zero
stalk at $m$, and also at $\eta$. These are all points, so it is
concentrated in degree zero. Every geometric specialization map is
an isomorphism: the only nontrivial relation between distinct points
is $\eta\rightsquigarrow m$, where both stalks are zero. Nevertheless
$h_*G$ cannot be locally constant near $m$, because its stalk there
is zero and every neighborhood contains isolated points with stalk
$\Lambda$. By Proposition~\ref{D:boundary-arbitrary}, the boundary
cohomology constructibility condition fails. The same computation,
and the finite geometric boundary, show that each boundary comparison
complex itself vanishes. This proves every assertion.
\end{proof}

\begin{proposition}[Known sufficient condition; proved numerical equivalence
on a smooth base]\label{D:curve-Swan}
Assume $f$ has relative dimension one, $S$ is noetherian and excellent,
and $G$ is a locally constant constructible sheaf of vector spaces over
a finite field $\Lambda$ of characteristic $\ell$, invertible on $S$.
Set
\[
 w_G(s)=\sum_{y\in Y_{\bar s}}\operatorname{Sw}_y(G|_{V_{\bar s}}).
\]
If $w_G$ is locally constant, every $R^i h_*G$ is locally constant
and its formation commutes with arbitrary base change. If, moreover,
$S$ is smooth over a perfect field, the converse holds: local constancy
of all $R^i h_*G$ is equivalent to local constancy of $w_G$.
\end{proposition}
\begin{proof}
We work on a connected component, so that the rank of $G$ and the
number of geometric boundary points are constant. Deligne--Laumon
\cite[Theorem 2.1.1, Corollary 2.1.2, Remark 2.1.3,
pp. 185--187]{Laumon1981} says that local constancy of
$\sum_y(\operatorname{Sw}_y(G)+\operatorname{rank}G)$ makes
$j_!G$ universally locally acyclic for $f$, and hence makes
$R^i h_!G$ locally constant. The same applies to the dual local
system $G^\vee$: Swan conductors of dual representations are equal,
because higher inertia is a $p$-group and its representations over
$\Lambda$, of characteristic different from $p$, are semisimple.
Poincar\'e--Verdier duality then makes $R^i h_*G$ locally constant
and compatible with base change. This is the argument of
\cite[Corollary 2.3 and its proof]{Achinger2017}; a finite disjoint
union of sections replaces its single section without altering the
argument. A proper smooth relative curve with a boundary section is
projective locally on the base, and the proper formulation of
Deligne--Laumon already suffices here.

For the converse under the extra assumption, write $d=\dim S$,
locally on a connected component. Absolute duality over the perfect
field gives
\[
 D_S Rh_*G\simeq Rh_!D_VG
             \simeq Rh_!G^\vee(d+1)[2d+2].
\]
The dual of a bounded complex with locally constant finite-dimensional
cohomology on the smooth $d$-dimensional base again has locally
constant cohomology. Thus $Rh_!G^\vee$ has locally constant
cohomology. Proper base change for direct image with proper supports
makes $\chi_c(V_{\bar s},G^\vee)$ locally constant. The
Grothendieck--Ogg--Shafarevich formula, in the normalization of
\cite[\S1.2.1, pp. 178--179]{Laumon1981}, gives
\[
 \chi_c(V_{\bar s},G^\vee)=
 \operatorname{rank}(G)\,\chi_c(V_{\bar s},\Lambda)-w_G(s).
\]
The first Euler characteristic on the right is locally constant for
a smooth proper family of curves with a finite \'etale boundary.
Therefore $w_G$ is locally constant. This establishes both directions
under the stated hypotheses. We do not use the converse without those
hypotheses as a substitute for Proposition~\ref{D:boundary-criterion}.
\end{proof}

\begin{proposition}[Known: tame ramification is sufficient]
\label{D:tame-sufficient}
In every relative dimension, a locally constant constructible $G$
which is tamely ramified along $Y$ relatively to $S$ has locally
constant $R^i h_*G$, and their formation commutes with arbitrary base
change. Here relative tameness means tameness at the generic points
of the boundary in every geometric fibre, as in
\cite[XIII, Definition 2.1.1]{SGA1}; it is not a condition only at
the generic points of the boundary of the total space.
\end{proposition}
\begin{proof}
The relative normal-crossings tame base-change lemma
\cite[XIII, Lemma 2.1.10, p. 105]{SGA7II} says that $Rj_*G$ commutes with
arbitrary base change. Its proof is useful here: \'etale locally,
Abhyankar's lemma gives a finite Kummer cover, still smooth over $S$
with relative normal-crossings boundary, on which $G$ is constant.
The injection of $G$ into the pushforward of its constant pullback,
and repetition on the cokernel, resolves $G$ by such pushforwards.
This is a resolution argument and does not assume that the covering
degree is invertible in the coefficient ring.

For a constant sheaf, the usual alternating resolution on the
intersections of boundary components expresses its extension by zero
in terms of constant sheaves on schemes smooth over $S$. Smooth
local acyclicity proves universal local acyclicity of the extension
by zero; this is \cite[XIII, Lemma 2.1.11, p. 105]{SGA7II} over a trait,
and its proof works for the definition using all geometric
specializations. The Kummer resolution proves the same for tame
local systems and their duals. Proper direct image and relative
Poincar\'e--Verdier duality, or equivalently the same Kummer
calculation for $Rj_*$, therefore show that $Rh_*G$ has locally
constant cohomology. Combining the first paragraph with proper base
change gives arbitrary base-change compatibility.
For a general finite abelian coefficient group, apply the argument
to its finite prime-primary filtration and use the long exact
sequences; local constancy is preserved under kernels, cokernels,
and extensions of finite locally constant sheaves.
\end{proof}

\begin{proposition}[Refuted: pointwise boundary acyclicity is not necessary
in relative dimension two]\label{D:boundary-cancellation}
There is a smooth proper family of surfaces with a smooth relative
boundary divisor and a finite invertible local system $G$ on its
complement such that $Rh_*G=0$, although the boundary vanishing
cycles are nonzero and the Swan conductor jumps.
\end{proposition}
\begin{proof}
Let $k=\overline{\mathbb F}_3$, $S=\mathbb A^1_t$, and take
$\Lambda=\mathbb F_{25}$, which contains the third roots of unity.
Choose a nontrivial character $\psi:\mathbb F_3\to\Lambda^\times$.
On $\mathbb A^1_x\times S$, let $F=\mathcal L_\psi(tx)$ be the
rank-one local system from the Artin--Schreier cover
$z^3-z=tx$. Its Swan conductor at $x=\infty$ is one when $t\ne0$
and zero when $t=0$. Thus Proposition~\ref{D:curve-Swan}, including
its converse over the smooth base $S$, shows that its direct image
is not locally constant. By Proposition~\ref{D:boundary-criterion},
$Rj_*F$ has nonzero boundary vanishing cycles.

Choose an elliptic curve $E'$ over $k$ with a nonzero point of order
two, and let $q:E'\to E=E'/C_2$ be the corresponding finite
\'etale isogeny. The nontrivial character of $C_2$ defines a
rank-one local system $L$ on $E$. We have $R\Gamma(E,L)=0$.
Indeed $q_*\Lambda=\Lambda\oplus L$, since two is invertible
in $\Lambda$. Translations act trivially on $H^i(E',\Lambda)$:
translation by a variable point is a family of automorphisms of
the proper smooth projection $E'\times E'\to E'$ and induces an
endomorphism of its constant cohomology sheaf. Such an endomorphism
is constant on the connected parameter scheme and is the identity
at its origin. Thus the nontrivial $C_2$-isotypic summand of
$R\Gamma(E',\Lambda)$ is zero, as claimed.

Set
\[
 X=\mathbb P^1_x\times E\times S,\quad
 Y=\{\infty\}\times E\times S,\quad
 G=F\boxtimes L\quad\hbox{on }V=\mathbb A^1_x\times E\times S.
\]
Push first along the proper projection $V\to\mathbb A^1_x\times S$.
Proper base change and the projection formula identify its direct
image with $F\otimes_\Lambda R\Gamma(E,L)=0$. Hence $Rh_*G=0$.
On the other hand, smooth base change along $E$ and an \'etale
trivialization of $L$ identify the boundary vanishing cycles locally
with the nonzero vanishing cycles of $Rj_*F$. They are nonzero,
while their global cohomology vanishes by the same elliptic factor.
This proves that the global vanishing condition in
\eqref{D:boundary-condition} cannot in higher dimension be replaced
by pointwise vanishing. It also rules out a necessary condition
requiring each individual boundary conductor to be constant.
\end{proof}


% BEGIN AUDITED POST-CP09 INTEGRATION
\subsection{Further uniform bounds and audited obstructions for D.1 and D.2(ii)}
\subsubsection{Uniform ramification breaks from the different}

\begin{lemma}[Proved: a degree-independent tube bound]
\label{D:different-conductor-bound}
Let $A$ be a complete discrete valuation ring of positive residue
characteristic, with fraction field $K$
and uniformizer $\pi$, and let $B$ be its integral closure in a finite
separable field extension $L/K$. Normalize $v$ by $v(\pi)=1$.
If $\pi^M\Omega_{B/A}=0$ for an integer $M\geq0$, then the
nonlogarithmic Abbes--Saito conductor satisfies
\[
                         c(L/K)\leq2M.                 \tag{D.UB.1}
\]
In particular, writing $e=e(L/K)$ and
$\delta=v_L(\mathfrak D_{B/A})$, we have
\[
                 c(L/K)\leq2\left\lceil\delta/e\right\rceil .
                                                        \tag{D.UB.2}
\]
No separability of the residue extension is assumed.
\end{lemma}

\begin{proof}
We use the tube definition in \cite[sections 3.1 and 6,
Proposition 6.2(i), Definition 6.3 and Proposition 6.4,
preprint pp.5--6 and 15--16]{AbbesSaito2002D}.
Lemma 7.1 of that source, on preprint p.19, gives a presentation
\[
                 B=A[X_1,\ldots,X_r]/(f_1,\ldots,f_r).
\]
Let $J=(\partial f_i/\partial X_j)$. Its transpose presents
$\Omega_{B/A}$. The annihilation assumption therefore gives
a matrix $H$ over $B$ with $JH=\pi^M I_r$.
Indeed, transposing the initial relation gives a scaled left
inverse of $J$, and over the field $L$ this is also a scaled
right inverse. Integral polynomial lifts of $H$ satisfy
\[
                  JH-\pi^M I_r=\sum_i f_i H_i             \tag{D.UB.3}
\]
for polynomial matrices $H_i$ over $A$.

Fix a rational $a>2M$, and work in a complete algebraically
closed valued extension of $K$. On the tube
\[
 X^a=\{x:\ v(x_j)\geq0,\ v(f_i(x))\geq a\text{ for all }i,j\},
\]
(D.UB.3) gives $J(x)H(x)=\pi^M(I_r+E(x))$, with
$v(E(x))\geq a-M>0$. Thus $J(x)$ is invertible and every
entry of its inverse has valuation at least $-M$.
Here valuations of vectors and matrices mean the minimum
of the valuations of their entries.

Apply Newton iteration
$x_{n+1}=x_n-J(x_n)^{-1}f(x_n)$, starting at $x_0=x$.
If $v(f(x_n))\geq a_n$, the correction has valuation
at least $a_n-M$, and integral Taylor expansion gives
$v(f(x_{n+1}))\geq2(a_n-M)$.
With $a_0=a$ and $a_{n+1}=2(a_n-M)$, we have
$a_n=2M+2^n(a-2M)\to+\infty$. Every iterate remains integral
and in $X^a$. Completeness gives a zero $\alpha$ of all the
$f_i$ such that
\[
                          v(x-\alpha)\geq a-M>M.          \tag{D.UB.4}
\]
Since $B[1/\pi]=L$, this zero is one of the finitely many
embeddings into the algebraic closure. We do not assert
convergence inside the noncomplete separable closure.

Distinct zeros $\alpha,\beta$ satisfy
$v(\alpha-\beta)\leq M$. Otherwise put $h=\beta-\alpha$,
$b=v(h)>M$. Taylor expansion gives
$h=-J(\alpha)^{-1}Q(h)$ with $v(Q(h))\geq2b$, which would
force $b\geq2b-M>b$, a contradiction.

Over a finite splitting extension containing elements of
the required rational valuations, intersect $X^a$ with
the finitely many closed polydiscs
\[
                    v(x_j-\alpha_j)\geq a-M\quad(1\leq j\leq r).
\]
These rational subdomains are pairwise disjoint, and
(D.UB.4) says that they cover all geometric points.
They form a finite admissible partition, each member
open and closed in the tube. No component contains two
zeros. Proposition 6.2(i) of the cited source says that
every geometric component contains a zero. Hence the
embeddings map bijectively to the components. Ramification
is bounded by every $a>2M$, which proves (D.UB.1), including
$M=0$.

For a finite flat complete intersection the different is
the zeroth Fitting ideal of $\Omega_{B/A}$ and annihilates it
\cite[section 49.7 and Lemma 49.12.3, Tags 0BVV and 0BWB]{StacksDDifferent}.
If $M=\lceil\delta/e\rceil$, then $v_L(\pi^M)=eM\geq\delta$,
so $\pi^M$ belongs to the different. This proves (D.UB.2).
\end{proof}

\begin{proposition}[Proved: a uniform different bound on a fixed surface resolution]
\label{D:surface-uniform-different}
Let $(S,s)$ be a normal surface germ of finite type over an
algebraically closed field. Let $Y\to S$ be a projective
resolution which is an isomorphism away from $s$, with
exceptional integral curves $E_1,\ldots,E_r$.
Choose positive integers $a_i$ such that
\[
        A=\sum_i a_iE_i,\qquad c_i=-A\cdot E_i>0.           \tag{D.UB.5}
\]
For any connected finite Galois cover of $S-\{s\}$,
let $e_i,\delta_i$ be the ramification index and different
exponent at any prolongation of $E_i$. Then
\[
           \sum_i c_i\,\delta_i/e_i
                         \leq K_Y\cdot A+\sum_i a_i.       \tag{D.UB.6}
\]
The same statement holds for covers of the punctured
strict henselization, with the same intersection numbers.
\end{proposition}

\begin{proof}
The exceptional intersection matrix is negative definite
with nonnegative off-diagonal entries
\cite[section 10.1, Theorem 10.1]{Kollar2013D}.
Put $M=-(E_i\cdot E_j)$ and solve $Ma=(1,\ldots,1)^t$.
To see that $a$ is nonnegative, write $a=a_+-a_-$ with
disjointly supported nonnegative vectors. If $a_-\ne0$, then
$a_-^tMa=a_-^tMa_+-a_-^tMa_-\leq-a_-^tMa_-<0$,
using the off-diagonal signs and positive definiteness.
This contradicts $a_-^tMa=a_-^t(1,\ldots,1)^t>0$.
Clear denominators.
Every coordinate is positive: if one were zero, its row
equation would have nonpositive left side.
This constructs (D.UB.5).

Normalize $S$ in the cover, normalize $Y$ in the same field,
and resolve the latter surface, preserving the smooth open
away from the exceptional locus. We obtain a smooth surface
$Z$, proper over the normal source $S'$, and a proper
generically finite morphism $f:Z\to Y$ of degree $n$
\cite[Theorem 54.14.5, Tag 0BGP]{StacksDSurfaceResolution}.
Contract only smooth rational $(-1)$-curves which map to
points of $Y$, using
\cite[Lemma 54.16.8, Tag 0C2M]{StacksDSurfaceResolution}.
The morphism to $Y$ descends through each contraction:
its graph image over the contracted smooth surface is
proper and quasi-finite, since its fibre over the contracted
point is one point. The projection is finite birational onto
a normal surface, hence is an isomorphism.
Each contraction decreases the number of exceptional
components, so this process terminates.

Let $F_{i\lambda}$ dominate $E_i$, and let $Q_j$ denote the
remaining exceptional curves. All are exceptional over
the finite set above $s$ in $S'$, so each has negative square.
Adjunction gives
\[
        K_Z\cdot Q_j=2p_a(Q_j)-2-Q_j^2\geq0,\qquad
                          K_Z\cdot F_{i\lambda}\geq-1.    \tag{D.UB.7}
\]
For $Q_j$ the only possible negative case would be a
smooth rational $(-1)$-curve, which has been removed.
For $F_{i\lambda}$ use only $p_a\geq0$ and $F_{i\lambda}^2\leq-1$.

The divisor $f^*A$ is effective, has coefficient
$a_i e_{i\lambda}$ along $F_{i\lambda}$, and has nonnegative
coefficients along all the $Q_j$. Consequently
\[
       K_Z\cdot f^*A
          \geq-\sum_i a_i\sum_\lambda e_{i\lambda}
          \geq-n\sum_i a_i,                              \tag{D.UB.8}
\]
because $\sum_\lambda e_{i\lambda}f_{i\lambda}=n$ with
$f_{i\lambda}$ the full residue degree.

The determinant of the differential defines an effective
divisor $R_f$ with $K_Z=f^*K_Y+R_f$. The morphism is
\'etale outside the exceptional locus. At $F_{i\lambda}$
its coefficient is the different exponent
$\delta_{i\lambda}$, also for inseparable residue extensions,
by \cite[Lemma 49.12.6, Tag 0BWJ]{StacksDDifferent}.
Here $f$ is quasi-finite on a neighborhood of the generic
point of each $F_{i\lambda}$, which is where that lemma is
applied; the global effective divisor only uses the
nonzero determinant of the differential.
Every point-mapped curve has intersection zero with $f^*A$.
Projection therefore gives the exact identity
\[
    K_Z\cdot f^*A
       =nK_Y\cdot A-
           \sum_i c_i\sum_\lambda\delta_{i\lambda}f_{i\lambda}.
                                                               \tag{D.UB.9}
\]
Galois conjugacy of the prolongation DVRs makes
$e_{i\lambda},\delta_{i\lambda}$ independent of $\lambda$,
so the inner sum is $n\delta_i/e_i$.
Combining (D.UB.8) and (D.UB.9) proves the assertion.
No equivariant resolution is required: the conjugacy is
intrinsic to the valuations, all of which survive on $Z$.

A finite cover and its finite Galois action over the punctured
strict henselization descend to an \'etale neighborhood.
The exceptional curves and their intersection numbers
are unchanged, since the residue field is algebraically
closed. Applying the same argument there proves the
strict-local assertion.
\end{proof}

\begin{proposition}[Proved: the original fixed-boundary correction]
\label{D:fixed-boundary-uniform-different}
In Proposition~\ref{D:surface-uniform-different}, let $T\subset S$
be any proper closed subset containing $s$, and let $D$ be
the strict transform on $Y$ of its reduced divisorial part.
For every connected finite Galois cover of $S-T$ which is
tame at the original codimension-one points of $T$, we have
\[
         \sum_i c_i\,\delta_i/e_i
                     \leq (K_Y+D)\cdot A+\sum_i a_i.       \tag{D.UB.10}
\]
In particular, $s$ need not be isolated in $T$.
\end{proposition}

\begin{proof}
Repeat the resolution and contraction construction of the
preceding proposition. Resolution of the normal source may
also introduce curves over points outside the exceptional
locus, but these have zero intersection with $f^*A$.
For the point-mapped components in the support of $f^*A$
the same contractions and adjunction estimates apply.
Thus (D.UB.8) remains valid.

Decompose the relative canonical divisor into components
dominating exceptional curves, components mapped to points,
and the horizontal part $R_{\rm hor}$ dominating $D$.
Tameness at an original boundary divisor implies that at
every prolongation there its coefficient is $e-1$ and the
residue extension is separable. Hence the coefficient of
$f_*R_{\rm hor}$ at that component of $D$ is
\[
                    \sum_\lambda(e_\lambda-1)f_\lambda
                          =n-\sum_\lambda f_\lambda\leq n.
\]
There are no other horizontal ramification components.
Since $A$ is effective and exceptional, no component of $D$
is contained in $A$, and intersections are nonnegative.
Projection gives
\[
                 R_{\rm hor}\cdot f^*A
                      =(f_*R_{\rm hor})\cdot A\leq nD\cdot A.
\]
Add this term to the right side of (D.UB.9), then use
(D.UB.8). This gives (D.UB.10).
\end{proof}

\begin{corollary}[Proved: a rank-independent exceptional break bound]
\label{D:surface-uniform-conductor}
Assume that the ground field has positive characteristic.
Put $C_A=(K_Y+D)\cdot A+\sum_i a_i$, where $D=0$ in the
punctured case. Every finite cover in the preceding
propositions, and every finite-monodromy representation
arising from it, has nonlogarithmic and logarithmic
Abbes--Saito breaks at $E_i$ at most
\[
                         2\lceil C_A/c_i\rceil.           \tag{D.UB.11}
\]
For a sole smooth exceptional curve of genus $g\geq2$
and square $-d$ in the punctured case, the stronger
bound is
\[
                         2\left\lceil1+(2g-2)/d\right\rceil.
                                                               \tag{D.UB.12}
\]
No finite-generation conclusion is asserted from these
bounds alone.
\end{corollary}

\begin{proof}
Every $\delta_i$ is nonnegative. Apply the separate
normalized-different bounds and
Lemma~\ref{D:different-conductor-bound} to the completed
boundary DVRs. An arbitrary finite cover is dominated by
its Galois closure, which remains in the same cover class.
For logarithmic breaks use
$G^a_{\log}\subseteq G^a$ from
\cite[Proposition 3.15(1), preprint p.8]{AbbesSaito2002D}.
For the sole positive-genus curve, Galois conjugacy in
(D.PG.1) gives
\[
    \delta/e\leq1-1/e+
       \frac{2g-2}{d}\left(1-\frac1{f^{\rm insep}}\right)
                  \leq1+\frac{2g-2}{d}.
\]
A finite-monodromy representation factors through one
finite Galois cover, and none of these bounds involves
its dimension.
\end{proof}

% Reconstructed from the complete geometry proof audited before the outage.
% Root, B and D: complete proof and Conrad/Abbes--Saito source audits PASS.
\begin{proposition}[Proved: Explicit simple covers with uniformly bounded breaks]
\label{D:psl2-bounded-conductor-covers}
Let $k$ be algebraically closed of characteristic two and let $q=2^r\ge4$.
There are arbitrarily many jointly linearly disjoint
$S_q=\operatorname{PSL}_2(\mathbf F_q)$-covers of $\mathbf G_{m,k}$.
Each compactifies to a map between projective lines with exactly two
branch values. At one value the inertia is tame cyclic of order $q+1$;
at the other it is the Borel group of order $q(q-1)$, with lower groups
\[
 G_0=\mathbf F_q\rtimes\mathbf F_q^\times,
 \qquad G_1=\mathbf F_q,
 \qquad G_2=1.
\]
Thus its positive upper break is $1/(q-1)\le1/3$. The groups $S_q$ are
nonabelian simple and admit three explicit generators.
\end{proposition}
\begin{proof}
In characteristic two every projective linear transformation over
$\mathbf F_q$ has a unique determinant-one representative. Thus
$S_q=\operatorname{SL}_2(\mathbf F_q)=\operatorname{PGL}_2(\mathbf F_q)$,
of order $q(q^2-1)$. Choose a generator $\lambda$ of
$\mathbf F_q^\times$. The matrices
\[
 u(1)=\begin{pmatrix}1&1\\0&1\end{pmatrix},\qquad
 d(\lambda)=\begin{pmatrix}\lambda&0\\0&\lambda^{-1}\end{pmatrix},
 \qquad w=\begin{pmatrix}0&1\\1&0\end{pmatrix}
\]
generate $S_q$. Conjugating $u(1)$ by powers of $d(\lambda)$ gives
$u(\lambda^{2j})$; these parameters span $\mathbf F_q$ additively.
Conjugation by $w$ gives all lower elementary matrices, and elementary
row reduction generates the entire special linear group.

We include the simplicity argument to specify exactly what is used
\citep[Theorem 2.10 and Theorems 3.1--3.5, pp.~2--5]{C_ConradPSL}.
The action on $\mathbf P^1(\mathbf F_q)$ is faithful and doubly
transitive. Hence every nontrivial normal subgroup $N$ is transitive.
Let $U$ be the translation subgroup in the stabilizer $B$ of infinity;
it is abelian and normal in $B$, and its conjugates generate $S_q$.
For each $g$ choose $n\in N$ with $g\infty=n\infty$; then $n^{-1}g\in B$.
In $S_q/N$ every conjugate of the image of $U$ is consequently that
same abelian image. The quotient is abelian. On the other hand
\[
 [d(\lambda),u(b)]=u((\lambda^2-1)b)
\]
and $\lambda^2\ne1$, so the upper and lower elementary groups lie in
the commutator group. Thus $S_q$ is perfect, and $N=S_q$.

Put $A=x^q-x$. The rational function
\[
 J_q(x)=\frac{(x^{q^2}-x)^{q+1}}{(x^q-x)^{q^2+1}}
       =\frac{(A^{q-1}+1)^{q+1}}{A^{q(q-1)}}
\]
is invariant under $S_q$. Indeed for a Mobius transformation
$g(x)=(ax+b)/(cx+d)$ over $\mathbf F_q$,
\[
 g(x)^{q^a}-g(x)
 =\frac{\det(g)(x^{q^a}-x)}{(cx+d)^{q^a+1}}.
\]
The exponents cancel in the displayed quotient. The identity
$x^{q^2}-x=A(A^{q-1}+1)$ proves the second formula. Its numerator and
denominator are coprime, and the reduced degree is
$q(q-1)(q+1)=|S_q|$. Hence $k(x)^{S_q}=k(J_q)$; the quotient is
separable because it is the fixed field of a finite automorphism group.

The poles are exactly $\mathbf P^1(\mathbf F_q)$, and the zeros are
$\mathbf P^1(\mathbf F_{q^2})\setminus\mathbf P^1(\mathbf F_q)$.
At infinity the stabilizer is $B$, of order $q(q-1)$. At a quadratic
point its stabilizer is the nonsplit torus of order $q+1$: identifying
the two-dimensional vector space with $\mathbf F_{q^2}$ realizes it as
$\mathbf F_{q^2}^\times/\mathbf F_q^\times$. Every nonidentity Mobius
transformation has its fixed points in a field extension of degree at
most two. These are therefore the only points with nontrivial
stabilizer, and there are exactly two branch values.

Use $u=1/x$ at infinity. A translation $x\mapsto x+b$, $b\ne0$,
acts by $u\mapsto u/(1+bu)$; its difference from $u$ has order two.
A stabilizer element with nontrivial multiplier has difference of
order one. Hence $G_1$ consists exactly of translations and $G_2=1$.
Herbrand's formula gives upper break $1/(q-1)$. The residue field is
perfect, so this is also the logarithmic break; the corresponding
nonlogarithmic bound is $1+1/(q-1)\le4/3$, by the classical comparison
in Abbes--Saito Proposition 3.15.

For $c\in k^\times$ consider the quotient map $t=cJ_q(x)$, restricted
over $\mathbf G_m$. We show that distinct $c$ give inequivalent
extensions of $k(t)$. The normalizer of $S_q$ in
$\operatorname{PGL}_2(k)$ is exactly $S_q$. A normalizing transformation
preserves the fixed-point set of the elements of order two, which is
$\mathbf P^1(\mathbf F_q)$. Its images of $0,1,\infty$ determine it,
so it is defined over $\mathbf F_q$. It therefore fixes $J_q$.
An isomorphism of two of the extensions would induce such a normalizer
element and would force the two constants to be equal.

Finally distinct Galois extensions with the same nonabelian simple
group are jointly linearly disjoint. We prove this by induction.
A subdirect subgroup of $S_q^m$ which is not the full product imposes,
by Goursat's lemma and simplicity, a diagonal identification between
two factors. Such an identification would identify the corresponding
extensions, contrary to their inequivalence. This proves the assertion
for arbitrarily large finite families, and hence the stated supply of
jointly independent covers.
\end{proof}

\begin{remark}[Finite verification and scope]
The exact certificate accompanying
Proposition~\ref{D:bounded-ramification-generation-obstruction} checks
$q=4,8,16$: finite-field arithmetic, the three generators, both orbits,
stabilizer orders, lower ramification, the invariant identity and the
multiplicity inequalities. The formulas above prove the assertions
for every $q$. These covers supply a counterexample to a proposed
group-theoretic implication from bounded ramification and linear
cohomology; they are not a counterexample to the proper-boundary
questions D.1 or D.2.
\end{remark}

\subsubsection{Bounded ramification and linear cohomology do not force generation}

\begin{proposition}[Refuted: a bounded-ramification generation inference]
\label{D:bounded-ramification-generation-obstruction}
Let $k$ be an algebraically closed field of characteristic $2$.
There is a quotient $G$ of $\pi_1(\mathbf G_{m,k})$ with all
of the following properties.
\begin{enumerate}
\item Every finite quotient cover has nonlogarithmic
Abbes--Saito breaks at $0$ and $\infty$ at most $4/3$.
\item $G$ is small: it has finitely many continuous
homomorphisms to each fixed finite group.
\item The maximal pronilpotent quotient of every open
subgroup of $G$ is finite.
\item For every finite field $\Lambda$ and every
finite-dimensional continuous $\Lambda[G]$-module $M$,
\[
                   \dim_\Lambda H^1(G,M)\le2\dim_\Lambda M.
                                                        \tag{D.BR.1}
\]
For an open subgroup $H$, the corresponding bound is
$2[G:H]\dim_\Lambda M$.
\item $G$ is not topologically finitely generated.
\end{enumerate}
In particular, bounded ramification, smallness, and
all-characteristic rank-linear degree-one cohomology
do not together imply finite generation.
This is an obstruction to that inference; it is not
a counterexample to the original-divisor tame assertions
D.1 or D.2(ii).
\end{proposition}

\begin{proof}
For $i\ge1$, set
\[
 q_i=2^{i+1},\qquad S_i=\operatorname{PSL}_2(\mathbf F_{q_i}),
 \qquad N_i=|S_i|^i+1,\qquad
                   G=\prod_{i\ge1} S_i^{N_i}.             \tag{D.BR.2}
\]
The groups $S_i$ are pairwise nonisomorphic nonabelian
simple groups.  Proposition~\ref{D:psl2-bounded-conductor-covers}
proves, without a classification input, that each is
generated by three elements and supplies arbitrarily many
pairwise inequivalent $S_i$-covers of $\mathbf G_m$.
Their branch points remain $0,\infty$; the wild branch
has classical upper jump $1/(q_i-1)$, and the other branch
is tame.  The same proposition proves linear disjointness
of every finite collection of these covers, including
distinct scalar twists with the same simple group.
Choose $N_i$ twists for each $i$.  The corresponding
map from $\pi_1(\mathbf G_m)$ surjects onto every finite
subproduct in (D.BR.2).  Its image is compact and therefore
closed, so it is the whole product.

At these valuations the residue field is perfect.
The comparison with the classical filtration gives
nonlogarithmic break
\[
                       1+\frac1{q_i-1}\le\frac43
\]
at the wild branch and at most $1$ at the tame branch
\citep[Proposition~3.15(4), preprint p.~8]{AbbesSaito2002D}.
Every finite quotient of $G$ factors through a finite
subproduct.  The ramification subgroup above a bound
acts trivially in a compositum when it acts trivially
in each constituent, and passage to quotients preserves
that property.  Thus the same $4/3$ bound applies to
every finite quotient cover, not only the displayed
simple factors.

Fix a finite group $F$.  A nontrivial homomorphism
$S_i\to F$ is injective.  Therefore it is impossible
when $|S_i|>|F|$.  Only finitely many values of $i$ can
contribute, and each has only the finite multiplicity
$N_i$.  Every continuous homomorphism $G\to F$ hence
factors through one fixed finite subproduct.  This proves
smallness.

If $H$ is open in $G$, it contains a full tail
\[
                     K_m=\prod_{i>m}S_i^{N_i}.
\]
Writing $P_m=\prod_{i\le m}S_i^{N_i}$, we have
$H=H_0\times K_m$ for a subgroup $H_0\subset P_m$:
multiply any chosen lift of an element of $H_0$ by
an element of the contained tail to remove its tail
coordinate.  Every finite nilpotent image of $K_m$
is trivial.  Indeed, a homomorphism to a finite group
uses only finitely many coordinates, and a nonabelian
simple factor has no nontrivial nilpotent quotient.
The maximal pronilpotent quotient of $H$ consequently
factors through the finite group $H_0$.

We give the cohomology argument in all coefficient
characteristics, including characteristics dividing
$|S_i|$.  For a finite product $P=\prod_{\nu=1}^t S_\nu$,
every absolutely simple module over an algebraically
closed coefficient field is an outer tensor product
$V=\bigotimes_\nu V_\nu$ of simple modules for the
factors.  This follows by passing from the finite
group algebras to their semisimple quotients: the
radical of each tensor factor acts trivially on a
simple module, and tensor products of matrix algebras
have the stated simple modules.  Thus no Maschke
assumption is being used.

If no factor acts nontrivially, $V$ is trivial and
$H^1(P,V)=0$, since $P$ is perfect.  If exactly one
factor acts nontrivially, inflation and the perfectness
of the remaining factors identify $H^1(P,V)$ with
$H^1(S_\nu,V_\nu)$.  A cocycle is determined by its
values on three generators, so its space has dimension
at most $3\dim V_\nu$.  Coboundaries have dimension
$\dim V_\nu$, because $V_\nu^{S_\nu}=0$.  Hence
\[
                       \dim H^1(P,V)\le2\dim V.           \tag{D.BR.3}
\]
If at least two factors act nontrivially, write
$P=A\times B$ with one active factor in $A$ and at
least one in $B$, and $V=V_A\otimes V_B$.
The degree-one Hochschild--Serre sequence injects
$H^1(P,V)$ into $H^1(A,V)^B$, since $V^A=0$.
But
\[
 H^1(A,V)^B
       =\bigl(H^1(A,V_A)\otimes V_B\bigr)^B=0.
\]
Thus (D.BR.3) holds in every case.

Extension of the coefficient field commutes with
finite-group cohomology.  A composition series after
algebraic closure and the degree-one long exact
sequence now give (D.BR.3) for every finite-dimensional
module over every finite field.

For a continuous $G$-module $M$, choose $m$ so that
the full tail $K_m$ acts trivially.  Its first
cohomology with these coefficients is
$\operatorname{Hom}_{\rm cont}(K_m,M)=0$, since every
finite tail image is perfect.  Inflation therefore
identifies $H^1(G,M)$ with $H^1(P_m,M)$ and proves
(D.BR.1).  Shapiro applied to
$\operatorname{Ind}_H^G M$ gives the stated open-subgroup
bound.

Finally, if $G$ had $d$ topological generators, choose
$i\ge d$.  Their images would generate $S_i^{N_i}$.
There are at most $|S_i|^d\le|S_i|^i$ possible ordered
$d$-tuples of coordinate entries, but
$N_i=|S_i|^i+1$.  Two coordinates therefore receive
identical tuples.  Every word in the proposed
generators has equal entries at these coordinates,
so the generated subgroup is contained in a proper
diagonal subgroup.  This contradiction proves
non-generation.
\end{proof}

\begin{remark}[Scope and finite certification]
The covers in this example are wild at one original
boundary point.  Accordingly, $G$ is not asserted to
be a quotient of the original-divisor tame group
for the pair $(\mathbf P^1,\{0,\infty\})$.
Its role is to test the proposed abstract deduction
from conductor bounds and cohomology to generation.

The exact finite program
\nolinkurl{parts/D/code/check_psl2_bounded_ramification_v2.py}
checks the fields $q=4,8,16$, the three matrix generators,
the two branch orbits and stabilizers, the lower inertia
filtration, the invariant identities, and the numerical
product inequalities.  Its current output is
\nolinkurl{parts/D/output/psl2_bounded_ramification_v2.json}.
This restored version is separately identified from the
historical program, which had 306,144 assertions and
was independently rerun before the environment replacement.
The general proof is the argument above and the
preceding geometric proposition; the finite computation
does not replace it.
\end{remark}

% Reconstruction includes the finite-free coefficient-ring correction
% accepted by all four proof auditors before restoration.
\begin{proposition}[Proved: Curvewise conductors and rank-linear cohomology]
\label{D:surface-curvewise-conductor}
Let $S$ be a normal proper surface over an algebraically closed field
of characteristic $p>0$, let $Y\subsetneq S$ be closed, and let
$Q=Q(S,Y)$ be the quotient requiring tameness only at the original
codimension-one components of $Y$. There are a smooth projective
surface $\widetilde S$, an SNC divisor $B$, and a fixed effective
integral divisor $R$ supported on $B$, with the following properties.
\begin{enumerate}
\item Pullback to $V=\widetilde S\setminus B$ gives a surjection
$\pi_1(V)\to Q$. Every finite-free local system over a finite local
ring of residue characteristic different from $p$ whose monodromy
factors through $Q$ has logarithmic conductors bounded curvewise by
$R$. In particular, for every morphism $g:C\to\widetilde S$ from a
smooth curve not contained in $B$, its highest logarithmic slope at
$x\in C$ is at most $\operatorname{mult}_x(g^*R)$.
\item There is a constant $c_{S,Y}$, independent of the finite field
$\Lambda$ of characteristic $\ell\ne p$ and of the rank, such that
\[
 \dim_\Lambda H^1(Q,M)\le c_{S,Y}\dim_\Lambda M
\]
for every finite-dimensional continuous $\Lambda[Q]$-module $M$.
For an open subgroup $H\subset Q$ the bound is
$c_{S,Y}[Q:H]\dim M$.
\item These assertions and the available wild Lefschetz theorem do
not by themselves give one finitely generated group surjecting onto
$Q$: the latter theorem's finite-coefficient choices depend on the
rank and on the finite coefficient group.
\end{enumerate}
\end{proposition}
\begin{proof}
Choose a projective resolution and further blowups so that the strict
transform of the original divisorial boundary, all exceptional curves,
and the inverse images of the remaining boundary points have SNC
support $B$. We remove every exceptional curve, including those above
retained singularities. Thus $V$ maps isomorphically to a dense smooth
open of $S\setminus Y$. A connected finite etale cover of the normal
scheme $S\setminus Y$ remains connected on this dense open, which
proves the quotient assertion.

Corollary~\ref{D:surface-uniform-conductor}, applied with its
original-boundary different estimate, bounds the generic logarithmic
break at every component of $B$ by a number depending only on the
fixed pair and resolution. It is a bound for all finite quotients,
not merely for one rank. A representation of such a quotient has no
larger slope. Choose integers dominating these finitely many numbers
and let $R$ be their divisor. For smooth SNC pairs the generic
logarithmic divisor controls pullback to every curve
\citep[Theorem 3.5, Remark 3.9 and Proposition 5.8(3)]{C_HuTeyssierRam2025}.
This gives the first assertion. The finite-local-ring statement here
uses locally constant sheaves with finite free stalks, exactly the
authors' $\operatorname{Loc}_{\rm tf}$ convention; we do not silently
replace it by arbitrary finite modules. The logarithmic conductor is
the highest slope, rather than the rank-weighted Swan conductor.

For finite-field coefficients, the uniform Betti theorem on a smooth
proper compactification bounds $h^1(V,M)$ by a fixed polynomial in
the largest generic logarithmic conductor, multiplied by the rank
\citep[Theorems 7.15 and 7.34, including the proof of 7.34]{C_HuTeyssierBetti2025}.
All its conductor arguments are now fixed, so this is $c_{S,Y}\dim M$.
The continuous group $H^1$ equals etale $H^1$ for a local system,
and inflation along $\pi_1(V)\twoheadrightarrow Q$ is injective.
This proves the second assertion. Shapiro's lemma gives its open
subgroup form by multiplying the induced rank by $[Q:H]$.

Finally we retain the exact coefficient quantifiers in the corrected
wild Lefschetz result: the final version uses finite coefficients,
and its uniform choices are for fixed rank and finite coefficient
group, with dependence through $\operatorname{GL}_r(\Lambda)$
\citep[Theorems 1.4--1.5; version 3]{C_HuTeyssierLefschetz2026}.
No assertion uniform over all these groups follows. The explicit
bounded-break simple-product construction in
Proposition~\ref{D:bounded-ramification-generation-obstruction}
also prevents inferring full finite generation from these conductor
and linear-cohomology bounds alone.
\end{proof}

% Reconstructed from the complete audited author transcriptions after CP09.
% Root, B and D passed the full mathematical transcription and the stated
% primary-source checks. This is a reconstruction, not recovered old bytes.
\begin{theorem}[Proved: A local rank-linear bound in characteristic $p$]
\label{D:local-p-linear-H1}
Let $(R,\mathfrak m,k)$ be a complete excellent normal local ring of
dimension two, of characteristic $p>0$, with algebraically closed residue
field. Put $W=\operatorname{Spec}R\setminus\{\mathfrak m\}$. There is a
constant $c_R$ such that every locally constant sheaf $M$ of rank $r$
over any finite field of characteristic $p$ satisfies
\[
 \dim H^1(W,M)\le c_Rr.
\]
The same conclusion holds for the puncture of an excellent strictly
henselian normal local surface ring with algebraically closed residue
field, with the constant supplied by its completion.
\end{theorem}
\begin{proof}
We first use $\mathbf F_p$ coefficients. The bundle
$\mathcal O_W\otimes_{\mathbf F_p}M$ has a coherent reflexive extension
$E$ to $R$. It is maximal Cohen--Macaulay. Its Frobenius linearization
$F_R^*E\to E$ is an isomorphism on $W$, and its kernel and cokernel
are supported at $\mathfrak m$. Set $H=H^2_{\mathfrak m}(E)$. The top
Cech cokernel for two parameters commutes with every right-exact tensor
functor. Consequently the induced structural map
\[
 R\otimes_{F,R}H\longrightarrow H
\]
is an isomorphism. This assertion does not require flatness of Frobenius:
the finite-length defects of the original linearization have zero top
local cohomology.

Let $H_{\rm nil}$ be the submodule of elements killed by some iterate of
the semilinear Frobenius. The structural isomorphisms show that, for
$z\in H_{\rm nil}$, the element $1\otimes z$ vanishes in every
sufficiently high Frobenius tensor power. Thus $H_{\rm nil}\subset
0_H^*$, with test multiplier $1$. Write $I$ for the injective hull of
$k$. By local duality $H$ embeds in a finite direct sum of copies of
$I$. Persistence of tight closure carries $H_{\rm nil}$ into the
tight closure of zero in this direct sum. The ideal
\[
 J=\operatorname{Ann}_R(0_I^*)
\]
therefore annihilates $H_{\rm nil}$. It is $\mathfrak m$-primary or is
$R$: the non-strongly-$F$-regular locus is contained in the closed point,
since every other local ring is regular. We use precisely the
injective-hull test ideal and its localization/support assertions in
\citet[Theorem 7.1(1), (2), (5)]{C_LyubeznikSmith2001}.

Let $D(-)=\operatorname{Hom}_R(-,I)$ be Matlis duality and set
\[
 P=D(H)=\operatorname{Hom}_R(E,\omega_R),\qquad
 N=D(H/H_{\rm nil}).
\]
Then $P$ is maximal Cohen--Macaulay of rank $r$, and
$JP\subset N\subset P$. Choose once a parameter ideal
$Q\subset\mathfrak mJ$. Since $\mathfrak mJP\subset\mathfrak mN$,
\[
 \dim_k N/\mathfrak mN
 \le \operatorname{length}(P/\mathfrak mJP)
 \le \operatorname{length}(P/QP)
 =r\operatorname{length}(R/Q).
\]
The last equality is the multiplicity formula for a parameter ideal on
a maximal Cohen--Macaulay module over a domain.

Frobenius is injective on $H/H_{\rm nil}$. A Frobenius-fixed element
$z$ of this Artinian module lies in its socle: for $a\in\mathfrak m$,
we have $a^{p^e}z=0$ for large $e$, so injectivity applied to
$F^e(az)$ gives $az=0$. Moreover $\mathbf F_p$-independent fixed vectors
are $k$-independent. Indeed, in a minimal $k$-linear relation we normalize
one coefficient to $1$, apply Frobenius, and subtract; minimality forces
all coefficients to lie in $\mathbf F_p$. It follows that
\[
 \dim_{\mathbf F_p}(H/H_{\rm nil})^{F=1}
 \le \dim_k\operatorname{Soc}(H/H_{\rm nil})
 =\dim_k N/\mathfrak mN\le c_Rr.
\]
A fixed vector in $H_{\rm nil}$ is zero, so the same bound holds for
$H^{F=1}$.

We spell out the Artin--Schreier comparison for a possibly nonfree
coherent module. A semilinear endomorphism of a finite module can be
lifted, after a finite free presentation, to a matrix of additive
polynomials $Cz^{[p]}-z$. Its differential is $-I$. Over an algebraically
closed residue field this additive polynomial map is surjective: its
image is an open subgroup of a connected vector group. Hensel lifting
then solves the equations over a strict local ring. Thus $F-1$ is
etale-locally surjective on $E$, and it is surjective on $E(R)$ here.
On $W$ its kernel is $M$. Since $H^1(W,\widetilde E)=H$, the resulting
long exact sequence embeds $H^1(W,M)$ into $H^{F=1}$. This proves the
assertion over $\mathbf F_p$. Restriction of scalars from
$\mathbf F_{p^f}$ multiplies both the rank and the cohomology dimension
by $f$, which cancels.

Finally let $R_0$ be the indicated strictly henselian ring. A connected
finite etale cover of its puncture normalizes to a finite normal domain
$B$ over $R_0$. Henselianity makes $B$ local. Excellence makes its
completion normal, and that completion is a local domain. Hence the
cover remains connected on the completed puncture. We obtain a
surjection of fundamental groups from the completed puncture to the
original one, and inflation gives the required injection on $H^1$ with
every finite local system. This proves the final assertion.
\end{proof}

\begin{theorem}[Proved: Uniform coherent cohomology for unit reflexive sheaves]
\label{D:unit-reflexive-linear-cohomology}
Let $S$ be a normal proper surface over an algebraically closed field of
characteristic $p>0$. There is a constant $c_S$ such that every reflexive
sheaf $E$ of rank $r$ equipped with an isomorphism
$F_S^{[*]}E\simeq E$ satisfies
\[
 h^0(S,E)\le r,\qquad h^1(S,E)\le c_Sr,\qquad
 h^2(S,E)\le c_Sr.
\]
Here $F^{[*]}$ denotes the reflexive hull of Frobenius pullback. If $E$
is obtained from a finite $\mathbf F_p$ local system on the regular
locus, then its fixed-section sheaf $L$ satisfies
$\dim_{\mathbf F_p}H^1(S,L)\le c_Sr$.
\end{theorem}
\begin{proof}
Frobenius scaling of the numerical first Chern class and second Chern
character gives $c_1(E)=0$ numerically and $\operatorname{ch}_2(E)=0$
\citep[Theorem 4.5(2)]{C_Langer2025}. The rank-linear correction in
Riemann--Roch for proper normal surfaces therefore gives
\[
 |\chi(S,E)-r\chi(S,\mathcal O_S)|\le c_0r
\]
\citep[Theorem 0.1]{C_Langer2026}. Properness, rather than projectivity
of $S$, is the hypothesis used here.

We fix a projective resolution $\pi:Y\to S$ with $Y$ smooth and
projective, and put $F=\pi^{[*]}E$. There is a fixed coherent ideal on
$S$, supported at its singular points, which annihilates the stable
endomorphism groups of all local maximal Cohen--Macaulay modules.
Equivalently, multiplication by each of its elements factors through a
finite free module. We use
\citet[Theorem 1.1 and Remark 2.5]{C_Kimura2025}, or the precise
factorization in \citet[Lemma 1.8]{C_Langer2026}.
If $z\operatorname{id}_E$ factors through a free module, Frobenius
pullback shows that $z^p$ annihilates the cokernel of
$F_S^*E\to F_S^{[*]}E$: take the reflexive hull of the factorization,
and use that the intervening free module already equals its hull.
After principalizing the fixed ideal of these $p$th powers, we obtain
a fixed effective divisor $A$ on $Y$ and injections
\[
 F(-A)\longrightarrow F_Y^*F\longrightarrow F.
\]
No constant in this sandwich depends on $E$ or $r$.

Fix an ample divisor $H$ and a smooth very ample curve $C\subset Y$
not contained in any of the fixed determinant-zero divisors. The maps
remain generically injective, hence injective as bundle maps, on $C$.
We do not require $C$ to be disjoint from $A$. Pulling back a maximal
slope subbundle by Frobenius gives, both on $Y$ and on $C$,
\[
 p\mu_{\max}(F)\le\mu_{\max}(F),\qquad
 p\mu_{\max}(F^\vee)\le\mu_{\max}(F^\vee)+A\cdot H,
\]
with $A\cdot C$ in the curve version. Thus the maximal slope of $F$
is nonpositive and that of $F^\vee$ is bounded by a fixed constant.
We have $h^0(F)\le r$. One direct verification takes the subsheaf
generated by all global sections: its determinant has nonnegative
degree, while the slope bound makes that degree nonpositive. Independent
generating sections therefore have determinant with no divisorial zero;
properness then makes every global section a constant linear combination
of at most $r$ of them.

Let $W=\pi^{[*]}\omega_S$, a fixed line bundle. Duality and restriction
give
\[
 H^2(S,E)^\vee=\operatorname{Hom}_S(E,\omega_S)
 \hookrightarrow H^0(Y,F^\vee\otimes W).
\]
Choose a fixed $m$ so that $F^\vee\otimes W(-mC)$ has negative maximal
slope. Successive restrictions to $C$ reduce the last $H^0$ to $m$
curve bundles whose maximal slopes are bounded by a fixed constant.
For a bundle $G$ on $C$ with maximal slope at most $a$, choose a fixed
effective divisor $D$ of degree greater than $a$; then
$H^0(C,G(-D))=0$ and $h^0(C,G)\le r\deg D$. We obtain
$h^2(S,E)\le c_1r$. Riemann--Roch now gives $h^1(S,E)\le c_Sr$.

The finite-free-presentation and Hensel argument in
Theorem~\ref{D:local-p-linear-H1} gives an exact etale sequence
\[
 0\longrightarrow L\longrightarrow E
 \xrightarrow{F-1}E\longrightarrow0.
\]
On the finite-dimensional space $H^0(S,E)$ the additive polynomial
$F-1$ is likewise surjective. Hence $H^1(S,L)$ injects into
$H^1(S,E)^{F=1}$. Independent fixed vectors over $\mathbf F_p$ are
independent over $k$, by the minimal-relation argument already used.
This proves the last assertion.
\end{proof}

\begin{theorem}[Proved: All-coefficient linear bounds on punctured proper surfaces]
\label{D:punctured-surface-all-characteristics-H1}
Let $S$ be a proper scheme of dimension at most two over an algebraically
closed field of characteristic $p>0$. Let $T$ be a finite closed subset
having codimension at least two in every irreducible component which it
meets, and let $U$ be a connected component of $S\setminus T$. There is
a constant $c_{S,T}$ such that, for every finite field $\Lambda$ and
every finite-dimensional continuous $\Lambda[\pi_1(U)]$-module $M$,
\[
 \dim_\Lambda H^1(\pi_1(U),M)\le c_{S,T}\dim_\Lambda M.
\]
For an open subgroup $H\subset\pi_1(U)$ the bound is
$c_{S,T}[\pi_1(U):H]\dim M$. In particular the maximal prosolvable
quotient of $\pi_1(U)$ is finitely generated. Full finite generation is
still subject to the nonabelian-crown condition of
Proposition~\ref{D:cohomology-crown-criterion}.
\end{theorem}
\begin{proof}
First suppose $S$ is a normal surface and use $\mathbf F_p$ coefficients.
For $j:U\hookrightarrow S$, the canonical reflexive extension of
$\mathcal O_U\otimes M$ is unit. Theorem~\ref{D:unit-reflexive-linear-cohomology}
bounds $H^1(S,j_*M)$ by a constant times the rank. The low-degree Leray
sequence has remaining term $H^0(S,R^1j_*M)$, supported at the finitely
many omitted points. Its stalks are bounded by
Theorem~\ref{D:local-p-linear-H1}; the completion comparison in that
theorem applies to each strict henselian stalk. This proves the
characteristic-$p$ bound, and restriction of scalars proves it over
every finite field of that characteristic. For coefficient characteristic
different from $p$ we apply
Proposition~\ref{D:surface-curvewise-conductor}, including its
rank-independent conductor and rank-linear Betti bounds.

For general $S$, reduction does not change its etale site, and the
normalization $\nu:X\to S_{\rm red}$ is finite. On each normalized
surface component we use the normal case. The components of dimension
at most one are proper and have finitely generated fundamental groups,
so their $H^1$ also has a fixed rank-linear bound. For the finite
surjection $X_U\to U$, effective descent is
\citet[Expose IX, Theorem 4.7]{C_SGA1}. A class in the kernel of pullback
on $H^1$ has a trivialization upstairs; the two pulled-back
trivializations differ by an $M$-translation on
$X_U\times_U X_U$. Its possible descent classes form a subquotient of
$H^0(X_U\times_U X_U,\operatorname{pr}^*M)$. Evaluation at one geometric
point of each connected component bounds this dimension by
$r\#\pi_0(X_U\times_U X_U)$. This proves the required uniform descent
bound, including nonnormal and non-pure-dimensional schemes.

For an open subgroup, Shapiro's lemma replaces $M$ by an induced module
of dimension $[\pi_1(U):H]\dim M$. Finally we use the proved
cohomological crown criterion. A prosolvable group has only abelian
chief factors, so the all-field linear $H^1$ estimate supplies its
finite-generation condition. The argument makes no assertion eliminating
nonabelian crowns of the full group.
\end{proof}

% Complete reconstruction of the local tame-p transcription audited by
% root and D, with the equivariant Matlis calculation retained explicitly.
\begin{lemma}[Proved: A uniform ordinary Frobenius defect in the Kummer rings]
\label{D:Kummer-uniform-Frobenius-defect}
In the local notation of Theorem~\ref{D:local-tame-p-linear-H1}, one
fixed $\mathfrak m$-primary ideal $J\subset R$ has the following
property for every $n$ prime to $p$. For every maximal Cohen--Macaulay
$R_n$-module $E$, multiplication by every element of $JR_n$ factors
through a finite free module. The cokernel of
$F_{R_n}^*E\to F_{R_n}^{[*]}E$ is annihilated by the ideal generated
by the $p$th powers of the elements of $J$.
\end{lemma}
\begin{proof}
Use the fixed ideal $J_0$ of base Jacobian minors displayed in the
proof below. Its extension belongs to the Jacobian ideal of every
$R_n$, independently of $n$. In this equicharacteristic complete
setting the Jacobian ideal is the sum of the Kahler differents for
regular Noether normalizations, and each Kahler different belongs to
the corresponding Noether different
\citep[Section 2, pp.~2--3, and Section 3(3.1), Examples 3.2--3.3,
pp.~6--7]{C_IyengarTakahashiStable}.
We recall the relevant factorization. For a finite regular
normalization $A\subset R_n$, write
$\eta=\sum u_i\otimes v_i$ in the annihilator of the kernel of
$R_n\otimes_A R_n\to R_n$, and let $z=\sum u_iv_i$.
The maps
\[
 E\longrightarrow R_n\otimes_A E\longrightarrow E,
 \qquad e\longmapsto\sum u_i\otimes v_ie,
 \qquad b\otimes e\longmapsto be
\]
are $R_n$-linear and compose to $z\operatorname{id}_E$.
The middle module is free over $R_n$, because $E$ is maximal
Cohen--Macaulay and therefore free over the regular local ring $A$.
Taking finite sums gives the assertion for $J_0R_n$.
Frobenius pullback of the factorization makes $z^p$ times the identity
on $F^*E$ factor through a free module. Taking reflexive hulls of the
maps, the free middle module remains unchanged, and shows that
$z^pF^{[*]}E$ lands in the image of $F^*E$. We may take $J=J_0$,
or any fixed $\mathfrak m$-primary subideal. This proves both claims.
\end{proof}

\begin{theorem}[Proved: A uniform local tame bound in coefficient characteristic $p$]
\label{D:local-tame-p-linear-H1}
Let $(R,\mathfrak m,k)$ be a complete normal local surface ring over an
algebraically closed field of characteristic $p>0$. Let $D$ be a fixed
reduced Weil divisor, and let $Q_D$ be the quotient of the fundamental
group of $\operatorname{Spec}R\setminus(D\cup\{\mathfrak m\})$
which requires tameness at the original height-one components of $D$.
There is a constant $c_{R,D}$ such that every finite-dimensional module
$M$ over a finite field of characteristic $p$ satisfies
\[
 \dim H^1(Q_D,M)\le c_{R,D}\dim M.
\]
The same bound applies to the image in local $H^1$ of classes from any
global fixed-original-divisor tame quotient. The assertion also holds
for excellent strictly henselian surface rings by completion.
\end{theorem}
\begin{proof}
We first use $\mathbf F_p$ coefficients. Choose a fixed reduced principal
divisor $\operatorname{div}(a)$ containing $D$. A general element of
the ideal of $D$ has this property: at each prescribed height-one prime
we choose valuation one, and on the regular puncture the remaining
divisor is generically reduced. We may enlarge $D$ by these auxiliary
components temporarily. Restriction to their complement is injective
on $H^1$, because a connected cover of a normal integral scheme remains
connected on a dense open.

For every $n$ prime to $p$, put
\[
 R_n=R[s]/(s^n-a),\qquad G=\mu_n.
\]
This is a normal Cohen--Macaulay local domain. Indeed it is finite free
over $R$ and therefore Cohen--Macaulay; at a component of $(a)$ it is
an Eisenstein extension of a regular DVR, and elsewhere on the puncture
it is etale. Serre's criterion proves normality, and reduction modulo
the maximal ideal proves locality. These arguments also show that its
puncture is regular.

We require one annihilator independent of $n$. Present
$R=k[[x_1,\ldots,x_v]]/(f_1,\ldots,f_t)$, of codimension $c=v-2$.
Let $I_j$ denote the ideal of $j$-minors. The fixed ideal
\[
 J_0=I_{c+1}\bigl(\operatorname{Jac}(f_1,\ldots,f_t,a)\bigr)
       +aI_c\bigl(\operatorname{Jac}(f_1,\ldots,f_t)\bigr)
       \subset R
\]
is $\mathfrak m$-primary or the unit ideal. Outside the closed point
the ring is regular, and the reduced curve $(a)$ is regular at its
generic points. The Jacobian matrix for $R_n$ has final row
$(-\partial a/\partial x_i,ns^{n-1})$. Its minors contain $J_0R_n$:
the first summand uses no $s$-column, and the second is contained in
$s^{n-1}I_c(\operatorname{Jac}f)$ because $a=s^n$ and $n$ is a unit.
For a finite generically separable regular Noether normalization, the
relative Jacobian elements are completely stable big test elements by
the Lipman--Sathaye theorem in
\citet[printed p.~207]{C_Hochster2007}. The complete-local Jacobian is
the sum of these relative Kahler differents, as in
\citet[Section 3(3.1), Examples 3.2--3.3]{C_IyengarTakahashiStable}.
This proves the assertion for our complete rings; the affine version
is also stated on printed p.~213 of Hochster's notes. Thus the same
fixed $\mathfrak m$-primary ideal $J\subset R$ annihilates the tight
closure of zero in the injective hull over every $R_n$. We may take
$J=J_0$, or any fixed $\mathfrak m$-primary subideal.

Choose $n$ divisible by the orders of the finitely many original tame
inertia images in the representation $M$. The Kummer pullback extends
to a local system $M_n$ on the puncture of $R_n$, by the DVR tame
calculation and purity on that regular puncture. The same $n$ works for
every class in $H^1(Q_D,M)$. Such a class is represented by affine
monodromy in $M\rtimes\operatorname{im}(Q_D\to\operatorname{GL}(M))$.
Its tame inertia is a group of order prime to $p$ and has trivial
intersection with the additive $p$-group $M$. Its projection therefore
has the same order, already dividing $n$. This explains why the cover
used below is independent of the cohomology class.

Write $E_n$ for the unit reflexive extension of
$\mathcal O\otimes M_n$, and set $H_n=H^2_{\mathfrak m_n}(E_n)$.
All constructions are $G$-equivariant. The proof of
Theorem~\ref{D:local-p-linear-H1} gives a Frobenius-injective quotient
$H_n/H_{n,\rm nil}$ with $JH_{n,\rm nil}=0$. We take Matlis duals
over the fixed ring $R$, rather than estimate ranks over $R_n$. Put
\[
 P=\operatorname{Hom}_R(E_n,\omega_R),\qquad
 N=D_R(H_n/H_{n,\rm nil}).
\]
Then $JP\subset N\subset P$. Since $|G|$ is invertible, invariants
are exact. The invariant part
\[
 P^G=\operatorname{Hom}_R(E_n^G,\omega_R)
\]
is maximal Cohen--Macaulay of rank $r$ over $R$: $E_n^G$ has rank $r$
by generic Galois descent. We have
$JP^G\subset N^G\subset P^G$. A fixed parameter ideal
$Q\subset\mathfrak mJ$ therefore gives
\[
 \mu_R(N^G)\le\operatorname{length}(P^G/QP^G)
 =r\operatorname{length}(R/Q).
\]
The Frobenius-fixed invariant vectors inject into the invariant socle
of $H_n/H_{n,\rm nil}$. Exact invariants and Matlis duality identify
its dimension with $\mu_R(N^G)$. Thus the invariant part of
$H^1$ on the upstairs puncture has dimension at most $c_{R,D}r$,
with no factor $n$.

Restriction to the Kummer cover is injective on $H^1$ in characteristic
$p$, by restriction--corestriction, since its degree $n$ is invertible.
Its image is invariant, and the preceding estimate applies to every
class under consideration. This proves the theorem over $\mathbf F_p$.
Restriction of scalars proves the assertion for every finite field of
characteristic $p$.

For a strictly henselian ring we use the completion surjection on
fundamental groups proved in Theorem~\ref{D:local-p-linear-H1}.
The original height-one tame condition remains tame after completion,
so the same injection and estimate apply. The image version for global
classes follows by exactly the same affine-monodromy argument; it does
not assert a bound for every unrestricted local class at a divisor.
\end{proof}

\subsubsection{Linear cohomology bounds for the original tame boundary}

We retain the original boundary throughout this subsection.  Let $S$ be
a normal proper connected surface over an algebraically closed field $k$
of characteristic $p>0$, let $T\subsetneq S$ be closed, and put
$U=S-T$.  Write $D$ for the reduced divisorial part of $T$.  We denote by
$Q(S,T)$ the quotient of $\pi_1(U)$ obtained by killing the wild inertia
at the original generic points of $D$.  Thus no condition is imposed
at the additional exceptional valuations of a resolution.

For a finite-dimensional $\mathbf F_p[Q(S,T)]$-module $M$, of dimension
$r$, choose a finite Galois cover of $U$ which trivializes $M$, normalize
$S$ in that cover, and write $f:S'\to S$ and $G$ for the resulting
finite map and group.  We use the canonical lattice
\[
             E=(f_*\mathcal O_{S'}\otimes_{\mathbf F_p}M)^G.
                                                        \tag{D.TC.1}
\]
This definition is independent of the chosen trivializing cover.
It defines a coherent reflexive sheaf of rank $r$.  Indeed, the
normal finite surface above $S$ is Cohen--Macaulay; kernels between
reflexive sheaves with torsion-free target are reflexive.  In (D.TC.1)
we use invariants only as a kernel.  We do not assume that the order
of $G$ is invertible, or that taking $G$-invariants is exact.

The ordinary Frobenius map induces a generically invertible injection
\[
            \varphi:F_S^{[*]}E\longrightarrow E.          \tag{D.TC.2}
\]
At a tame boundary DVR it is the invariant map on a Kummer lattice.
If a character has weight $j/n$, where $0\le j<n$ and $(n,p)=1$,
its elementary exponent is
\[
                         c_j=\lfloor pj/n\rfloor,
                   \qquad 0\le c_j\le p-1.              \tag{D.TC.3}
\]
Away from $D$, (D.TC.2) is an isomorphism.  Its cokernel
$\mathcal T$ is pure of dimension one: both terms in (D.TC.2)
have depth two at closed points, and the depth lemma excludes
zero-dimensional associated points in the cokernel.

\begin{lemma}[Proved: the summed boundary-lattice index]
\label{D:tame-boundary-summed-index}
Fix a completed normal surface germ $R$ at one of the finitely many
singular or boundary-intersection points.  Choose a reduced element
$a\in R$ whose divisor contains the local original boundary, and a
parameter $x$ such that
\[
                    A_0=k[[x,a]]\subset R
\]
is finite and $x$ is separating on every component of the reduced
curve $R/(a)$.  There is a constant $L$, depending only on these
fixed data, with the following property.

For every tame local system of rank $r$, and every prime-to-$p$
Kummer denominator $n$ which kills its original divisorial inertia,
the pure cokernel $\mathcal T$ of (D.TC.2) admits, as a
$k[[x]]$-lattice, a comparison
\[
         T_{\rm diag}\ \subset\ \mathcal T\
                      \subset\ x^{-L}T_{\rm diag},
       \qquad
         \operatorname{rank}_{k[[x]]}T_{\rm diag}
                      \le (p-1)[R:A_0]r.               \tag{D.TC.4}
\]
The generic graded pieces of $T_{\rm diag}$ are the whole
character eigenspaces, repeated $c_j$ times as in (D.TC.3).
For any one generic filtration of these pieces, the sum of the
positive index corrections after passage to the normalization
of the fixed curve $R/(a)$ is at most $C r$, with $C$ independent
of $n$, of $M$, and of the number of characters.
\end{lemma}

\begin{proof}
Put $A_n=k[[x,s]]$, with $a=s^n$, and
$R_n=R\otimes_{A_0}A_n$.  The reduced-divisor condition and
$(n,p)=1$ make $R_n$ normal in codimension one; it is
Cohen--Macaulay because it is finite free over $R$.  Hence it is
normal.  At the closed point it is local, with residue field $k$.
The pulled-back local system extends over the puncture of $R_n$,
by the regular codimension-one Kummer calculation and normality.
Let $E_n$ be its reflexive unit lattice.  It is free over the
regular ring $A_n$, of rank $d_0r$, where $d_0=[R:A_0]$.

We use the fixed stable-MCM factorization ideal from
Lemma~\ref{D:Kummer-uniform-Frobenius-defect} and
Theorem~\ref{D:local-tame-p-linear-H1}.  Their local
tame-cohomology argument supplies a fixed
power $x^N$ whose Frobenius power annihilates the ordinary
Frobenius defect of every $E_n$, uniformly in $n$ and $r$.
This is the uniform ideal obtained from the base Jacobian
minors and the term $a\operatorname{Jac}(R)$; it is not an
assertion that test ideals commute with Kummer base change.
Indeed, the fixed Jacobian ideal is contained in the sum of
Noether differents.  Each such element factors the identity
on an MCM module through a free module, and its $p$th power
annihilates the ordinary-to-reflexive Frobenius defect
\citep[section~2; section~3, formula~(3.1) and
Examples~3.2--3.3, pp.~2--3 and 6--7]{C_IyengarTakahashiStable}.
Equivalently, after increasing a fixed integer $K$, the cokernel
of $F_{R_n}^*E_n\to E_n$ is killed by $x^K$.

Consider relative Frobenius for the finite free $A_0$-algebra $R$:
\[
             \alpha:F_{A_0}^*R\longrightarrow R.
\]
Its determinant $h(x,a)$ is not zero on the generic point of
any component of $a=0$.  This is exactly where the separating
choice of $x$ is used: the corresponding finite extension of
curve function fields is separable, so relative Frobenius is
an isomorphism there.  Thus $h(x,0)\ne0$.  The adjugate
annihilates the relative cokernel.  The same assertion after
$a=s^n$ gives $h(x,s^n)$ for $R_n/A_n$.

Let $B$ be the matrix, between free $A_n$-modules of rank
$d_0r$, of the composite
\[
       F_{A_n}^*E_n\longrightarrow F_{R_n}^*E_n
                              \longrightarrow E_n.
\]
Tensoring the relative-Frobenius cokernel calculation and
then using the ordinary defect yields
\[
                 x^K h(x,s^n)\operatorname{coker}(B)=0.
                                                        \tag{D.TC.5}
\]
Consequently $x^K h B^{-1}$ has integral entries.
Put $h_0=\operatorname{ord}_x h(x,0)$.  Modulo $s^{np}$,
the inverse of $h(x,s^n)$ is the finite geometric expansion
of its constant term in $s^n$, with at most $p$ terms.
Its $x$-denominator is bounded by $x^{p h_0}$, up to a unit.
Thus one fixed $L$, for example $K+p h_0$ after enlarging
it to absorb the chosen unit convention, satisfies
\[
 B^{-1}\bmod s^{np}\ \in\
     x^{-L}\operatorname{Mat}_{d_0r}
                        (A_n/(s^{np})).                 \tag{D.TC.6}
\]
The bound is independent of the size of the matrix.

Choose a homogeneous $A_n$-basis for $E_n$ under $\mu_n$.
It exists by choosing a homogeneous basis modulo $(x,s)$
and applying Nakayama.  If a basis vector has weight $-j$,
the corresponding invariant generator is $s^jv_j$.
The invariant module
$(F_{A_n}^*E_n)^{\mu_n}$ has generator
$s^{pj\bmod n}Fv_j$, and the map from $F_{A_0}^*E$
to it is diagonal with entries $a^{c_j}$.
The quotient of this diagonal map is the lattice
$T_{\rm diag}$ in the statement.  It is free over $k[[x]]$
and has rank $\sum c_j\le(p-1)d_0r$.

The invariant map induced by $B$ sends the diagonal image
into the ordinary Frobenius image in $E$, hence into the
reflexive image in (D.TC.2).  It therefore gives
$T_{\rm diag}\to\mathcal T$.  This map is an isomorphism
after inverting $x$.  Purity makes $\mathcal T$ torsion-free
over $k[[x]]$, while $T_{\rm diag}$ is free, so it is
injective.  Also $a^{p-1}\mathcal T=0$, since this holds at
each generic boundary point and $\mathcal T$ has no
zero-dimensional submodule.  Equation (D.TC.6), applied
modulo $a^p$, now gives
$x^L\mathcal T\subset T_{\rm diag}$.  This proves (D.TC.4).

For completeness, the final-grade estimate uses a single
filtration, not a separate index bound for each character.
Identify the two lattices over $k((x))$, choose the generic
filtration by boundary component, character and $a$-layer,
and intersect it with both lattices.  Additivity of length
in the resulting finite filtrations gives
\[
 \sum_{\rm grades}\operatorname{length}_{k[[x]]}
        (L_{\rm grade}/L'_{\rm grade})
     =\operatorname{length}_{k[[x]]}(\mathcal T/T_{\rm diag})
     \le L(p-1)d_0r.                                    \tag{D.TC.7}
\]
Let $\overline R_0$ be the normalization of $R/(a)$.
It is finite free over $k[[x]]$, of fixed rank.  For an
inclusion $L'\subset L$ of the lattices occurring in
(D.TC.7), the quotient of their normalized spans is a
quotient of
\[
             \overline R_0\otimes_{k[[x]]}(L/L').
\]
Thus normalization multiplies the total index by at most
$\operatorname{rank}_{k[[x]]}\overline R_0$.  There is no
additional factor equal to the number of grades or to $r$.
This proves the last assertion.
\end{proof}

\begin{lemma}[Proved: a uniform Euler bound for the tame Frobenius quotient]
\label{D:tame-canonical-Euler}
There is a constant $C$, depending only on $(S,T)$, such that
the pure quotient in (D.TC.2) satisfies
\[
                              \chi(S,\mathcal T)\le Cr.  \tag{D.TC.8}
\]
The same bound holds at every stage of a finite cyclic
coefficient-Frobenius orbit of canonical lattices.
\end{lemma}

\begin{proof}
Enlarge a fixed finite set $\Sigma$ so that $S-\Sigma$ is
smooth and its boundary components are smooth and disjoint.
On this open set the assertion is the elementary Kummer
calculation.  A character of weight $j/n$ contributes the
whole weight space in exactly $c_j$ successive normal
layers.  The possible normal twists belong to the fixed
list $0,\ldots,p-2$.

We explain both the global slope control and the correction
at $\Sigma$.  Normalize each fixed boundary curve, obtaining
smooth proper curves $\overline D_i$.  The normal line
bundle on $D_i-\Sigma$ has a fixed extension to
$\overline D_i$.  Choose finitely many fixed rational
normal trivializations.  Kummer root changes between
them differ by roots of these fixed rational functions.
For a weight $j/n\in[0,1)$, the corresponding integral
lattice shift is bounded by the absolute orders of
these functions, independently of $j$ and $n$.

At a point of $\Sigma$, use the local cover $R_n$ of the
preceding lemma.  The module $E_n/sE_n$ is a
maximal-Cohen--Macaulay module on the reduced curve
$R/(a)$.  If $b$ is a non-zero-divisor in its fixed
conductor, then
\[
 b\bigl(\overline R_0\otimes(E_n/sE_n)/{\rm torsion}\bigr)
                  \subset E_n/sE_n
                  \subset
       \overline R_0\otimes(E_n/sE_n)/{\rm torsion}.
                                                        \tag{D.TC.9}
\]
Taking character summands is exact here, since
$(n,p)=1$.  The same fixed conductor and the preceding
normal-trivialization shifts therefore glue the normalized
character lattices to vector bundles $P_{i,j}$ on
$\overline D_i$, with comparison divisors independent
of $n,r$ and $j$.  Equivalently, one may carry out the
gluing on the root gerbe of the chosen normal line bundle
and then twist the character-$j$ part by the universal
character line before descent
\citep[Proposition~3 and Theorem~4, pp.~11--16]{Borne2007D}
\citep[Theorem~6.1, pp.~33--36]{BorneVistoli2012D}.
Formula (D.TC.9) controls
the difference from the chosen normalized local lattice.

The regular Frobenius map on $E_n$ gives, after reducing
modulo $s$, normalizing and accounting for these fixed
comparison divisors, generically invertible maps
\[
       F_{\overline D_i}^*P_{i,j}
                  \longrightarrow
       P_{i,pj\bmod n}(B_i+c_jN_i).                      \tag{D.TC.10}
\]
Here $B_i$ is fixed and $N_i$ denotes the chosen normal
line bundle; changing its sign convention merely changes
the fixed bound.  Regularity of this map uses only the
forward Frobenius map.  We do not assert an inverse
without a pole bound.  The finitely many comparison
divisors in (D.TC.9), and the bound $0\le c_j\le p-1$,
give one constant $C_i$ such that
\[
 p\,\mu_{\max}(P_{i,j})
             \le \mu_{\max}(P_{i,pj\bmod n})+C_i.
\]
Cycling the characters gives
$\mu_{\max}(P_{i,j})\le C_i/(p-1)$.
For coefficient fields $\mathbf F_{p^f}$ the same argument
cycles pairs consisting of a coefficient-Frobenius stage
and a character.  It is unchanged and remains independent
of $f$.

Thus the Euler characteristic of each whole character
bundle, with any of the finitely many normal twists above,
is at most a fixed constant times its rank.  The sum of
these ranks over all normal layers is at most
$(p-1)r$ for each fixed $D_i$.

The actual pure grades of $\mathcal T$ need not equal
these normalized character lattices at the points of
$\Sigma$.  This is precisely the correction estimated
in Lemma~\ref{D:tame-boundary-summed-index}.  On the
generic curve the diagonal source grades are whole
character spaces, since $c_j$ is constant on a character;
the leading term of $B$ maps them integrally into the
target.  Intersect the same generic filtration with
$T_{\rm diag}$ and $\mathcal T$.  Equation (D.TC.7) and
the normalized-span estimate give a total positive
Euler correction $O(r)$, not one $O(r)$ correction for
each grade.  Summing over the fixed finite set $\Sigma$
proves (D.TC.8).
\end{proof}

\begin{proposition}[Proved: linear coherent cohomology for canonical tame lattices]
\label{D:tame-canonical-linear-cohomology}
For each $i=0,1,2$ there is a constant $C_i$, depending
only on $(S,T)$, such that every canonical lattice (D.TC.1)
of rank $r$ satisfies
\[
                              h^i(S,E)\le C_i r.         \tag{D.TC.11}
\]
\end{proposition}

\begin{proof}
The determinant of (D.TC.2) gives
\[
       (1-p)c_1(E)=[\mathcal T]_1
                    =\sum_i m_iD_i,\qquad
                    0\le m_i\le(p-1)r.                 \tag{D.TC.12}
\]
Consequently $K_S\cdot c_1(E)=O(r)$ in the numerical
intersection theory of the normal proper surface.
By the normal-surface Riemann--Roch estimate and Frobenius
compatibility of the Chern character
\citep[Theorem~0.1]{C_Langer2026}
\citep[Theorem~4.5(2)]{C_Langer2025},
\[
 \begin{split}
 \chi(\mathcal T)
   &=\chi(E)-\chi(F_S^{[*]}E)\\
   &=-(p^2-1)\operatorname{ch}_2(E)
                  +\frac{p-1}{2}K_S\cdot c_1(E)+O(r).
 \end{split}                                            \tag{D.TC.13}
\]
The local Riemann--Roch error is bounded by a fixed
constant times the rank.  Combining (D.TC.8),
(D.TC.12), and (D.TC.13) gives
$\operatorname{ch}_2(E)\ge-Cr$, hence
\[
                                  \chi(E)\ge-Cr.         \tag{D.TC.14}
\]
The direction of this inequality is important: an upper
bound for $\chi(\mathcal T)$ supplies the lower bound
for $\chi(E)$ needed to control $h^1(E)$.

We now bound $h^0$ and $h^2$ without assuming that $S$
is projective.  Choose one smooth projective resolution
$\pi:Y\to S$ principalizing the following fixed ideal.
Codimension-one checking in (D.TC.3), followed by
reflexivity, gives
\[
    \mathcal O_S(-(p-1)D)\,E
                       \subset \varphi(F_S^{[*]}E).
\]
Multiply by the fixed stable-MCM Frobenius-defect ideal
\citep[Theorem~1.1 and Remark~2.5(1)]{C_Kimura2025}
\citep[Lemma~1.8]{C_Langer2026}
to land in the image of the ordinary $F_S^*E$.
This works for every reflexive $E$ of every rank.  The
product is one coherent ideal $I$, independent of $E$.
Write $I\mathcal O_Y=\mathcal O_Y(-A)$ and
$F=\pi^{[*]}E$.  Frobenius is flat on the smooth surface
$Y$, and the ordinary map and its fixed-ideal inclusion
give
\[
                     F(-A)\subset F_Y^*F\subset F.       \tag{D.TC.15}
\]
The second map is generically an isomorphism, hence
injective.  The first is checked at codimension-one
DVRs and then extended by reflexivity.

Choose a fixed smooth very ample curve $C$ on $Y$ which
is not contained in the determinant-zero divisors of
these maps.  Intersections with $A$ are allowed.
Both restricted maps are generically injective between
vector bundles on $C$, hence injective.  The usual
pullback inequality for maximal slopes, applied to
(D.TC.15) and its dual, yields
\[
 \mu_{\max,H}(F)\le0,\qquad
 \mu_{\max,H}(F^\vee)\le A\cdot H/(p-1),
                                                        \tag{D.TC.16}
\]
and the identical bounds on $C$, with $H$ replaced by
$C$.  Here $H=C$ may be fixed once and for all.
The maximal-slope subsheaf pulls back to a subsheaf,
so the only Frobenius slope inequality used in
(D.TC.16) is $p\mu_{\max}\le\mu_{\max}(F^*(-))$.

The vanishing $H^0(Y,F(-C))=0$ and evaluation at one
point on $C$ give
\[
                             h^0(S,E)\le h^0(Y,F)\le r.
\]
For $h^2$, let $W=\pi^{[*]}\omega_S$, a fixed line bundle.
Proper duality and reflexive pullback give an injection
\[
 H^2(S,E)^\vee=\operatorname{Hom}_S(E,\omega_S)
                       \hookrightarrow H^0(Y,F^\vee W).
                                                        \tag{D.TC.17}
\]
Put $\beta=A\cdot C/(p-1)+W\cdot C$.
Choose fixed positive integers $m,t$ with
$mC^2>\beta$ and $t>\beta$.
Then $H^0(Y,F^\vee W(-mC))=0$.  On each of the
$m$ restrictions to $C$, evaluation at $t$ distinct
points bounds the space of sections by $tr$,
using the maximal-slope bound $\beta$.
The restriction sequences therefore give
$h^2(S,E)\le mtr$.  Finally
$h^1=h^0+h^2-\chi$ and (D.TC.14) prove (D.TC.11).
\end{proof}

\begin{theorem}[Proved: all-characteristic linear degree-one cohomology for a surface boundary]
\label{D:tame-surface-all-characteristics-H1}
\label{D:surface-original-boundary-p-H1}
\label{D:surface-original-boundary-prosolvable}
For the original fixed-divisor tame quotient $Q=Q(S,T)$
there is a constant $C_{S,T}$ such that, for every finite
field $\Lambda$ and every finite-dimensional continuous
$\Lambda[Q]$-module $M$,
\[
                  \dim_\Lambda H^1(Q,M)
                              \le C_{S,T}\dim_\Lambda M.
                                                        \tag{D.TC.18}
\]
For an open subgroup $H\subset Q$ the bound is
\[
        \dim_\Lambda H^1(H,M)
                    \le C_{S,T}[Q:H]\dim_\Lambda M.       \tag{D.TC.19}
\]
In particular the maximal prosolvable quotient of $Q$
is topologically finitely generated.  We make no
finite-generation assertion here for the full group $Q$.
\end{theorem}

\begin{proof}
We first use $\Lambda=\mathbf F_p$.  On the small
\'etale site of $S$, the canonical semilinear Frobenius
gives an exact sequence
\[
          0\longrightarrow j_*M\longrightarrow E
                   \xrightarrow{\varphi-1}E
                   \longrightarrow0,                  \tag{D.TC.20}
\]
where $j:U\hookrightarrow S$.
For surjectivity, choose local generators of $E$ and
write Frobenius as a matrix on these generators.
The corresponding Lang polynomial has the form
$Cz^{[p]}-z-b$ and Jacobian $-I$.  Over the separably
closed residue field it is surjective: it is an
\'etale additive map of affine space, whose image is
an open subgroup, hence the whole connected group.
Hensel lifting at a strict stalk supplies a solution.
No freeness of $E$ at the singular point is required.
For the kernel, use (D.TC.1).  Frobenius-fixed functions
on the finite normal cover are locally constant
$\mathbf F_p$-valued functions; these extend across
the boundary.  Taking kernels and then $G$-invariants
identifies the result with $j_*M$, without an
exactness assertion about invariants.

The vector space $H^0(S,E)$ is finite dimensional and
$\varphi-1$ is surjective on it by the same Lang argument.
Hence (D.TC.20) injects $H^1(S,j_*M)$ into
$H^1(S,E)^{\varphi=1}$.  Frobenius-fixed vectors which
are independent over $\mathbf F_p$ are independent
over $k$: a shortest normalized $k$-linear relation,
minus its Frobenius transform, forces all coefficients
to lie in $\mathbf F_p$.  Proposition
\ref{D:tame-canonical-linear-cohomology} therefore gives
\[
                      \dim_{\mathbf F_p}H^1(S,j_*M)\le Cr.
                                                        \tag{D.TC.21}
\]

Apply the degree-one Leray sequence for $j$ only to
classes inflated from $H^1(Q,M)$.  These are represented
by finite torsors whose combined monodromy still factors
through the original tame quotient.  At a smooth point
of a smooth boundary divisor their local classes vanish:
the strict-local tame inertia has order prime to $p$,
and its positive cohomology on an $\mathbf F_p$-module
vanishes.  Thus their local images are supported on the
fixed finite set $\Sigma$ of singular, intersection and
isolated omitted points.

At each point of $\Sigma$, the uniform local tame-$p$
image bound in Theorem~\ref{D:local-tame-p-linear-H1}
applies to exactly these classes.  Its
Kummer cover kills the prime-to-$p$ inertia of the
given representation and torsor; the fixed Jacobian
test ideal and equivariant Matlis bound are independent
of its denominator.  The sum of the local dimensions
is therefore $O(r)$.  At a point away from the
divisorial boundary this is the already proved
punctured-surface local bound
(Theorem~\ref{D:local-p-linear-H1}).  We do not apply this
argument to all of $R^1j_*M$, which can also contain
wild classes excluded by $Q$.

Inflation injects $H^1(Q,M)$ into $H^1(U,M)$.
Combining the restricted Leray sequence, (D.TC.21),
and these finitely many local bounds proves (D.TC.18)
in characteristic $p$.  Restriction of scalars gives
the same constant for every finite field of
characteristic $p$.

For $\operatorname{char}\Lambda\ne p$, the uniform
original-boundary conductor theorem, its curvewise
consequence, and the uniform Betti estimate already
proved in Proposition~\ref{D:surface-curvewise-conductor}
give (D.TC.18).  Shapiro applied to
$\operatorname{Ind}_H^Q M$ proves (D.TC.19).
Finally the previously proved cohomology/crown
criterion, restricted to finite solvable quotients,
uses only abelian crowns and gives a uniform bound
on their number of generators.  Compactness then
gives finite generation of the maximal prosolvable
quotient.
\end{proof}

\begin{remark}[The remaining nonabelian gate]
The estimates above control all abelian crowns, in
every coefficient characteristic.  They do not
bound the multiplicities of arbitrary nonabelian
monolithic crowns.  The fixed-rank representation
counts and the bounded-conductor product examples
elsewhere in this paper are kept separate from this
remaining full finite-generation assertion.
\end{remark}

% Full reconstruction of the audited structural slice/finite-descent proof.
\begin{proposition}[Proved: Reduction of a complete local puncture to surfaces]
\label{D:local-surface-reduction}
Let $(A,\mathfrak m,k)$ be a complete excellent normal local domain of
dimension $d\ge2$ with algebraically closed residue field, and let
$U=\operatorname{Spec}A\setminus\{\mathfrak m\}$. There are finitely
many complete normal local surface domains $B_1,\ldots,B_t$, with residue
field $k$, and an integer $e\ge0$ for which there is a continuous
surjection
\[
 \left(\mathop{\ast}_{i=1}^t
 \pi_1(\operatorname{Spec}B_i\setminus\{\mathfrak m_i\})\right)
 \ast\widehat F_e\ \longrightarrow\ \pi_1(U).
\]
The star denotes free profinite product. In equicharacteristic $p>0$,
all the $B_i$ are in the same equicharacteristic setting. Thus any
rank-linear $H^1$ bound, fixed-rank exponential defining-characteristic
representation bound, or uniform quadratic representation bound for
complete normal local surface punctures passes to $U$.
\end{proposition}
\begin{proof}
Choose $d-2$ elements of a system of parameters, let $I$ be their ideal,
and set $C=(A/I)_{\rm red}$. Excellence and catenarity show that every
minimal component of $C$ has dimension two. Write
$W=\operatorname{Spec}C\setminus\{\mathfrak m_C\}$.

A connected finite etale cover of $U$ normalizes to a finite complete
local normal domain $B$ over $A$. Its dimension is $d$ and its
connectedness dimension is $d$, since it is a domain. The ideal $IB$
has cohomological dimension at most $d-2$. The connectedness theorem
\citep[Tag 0ECR, with Tag 0DX9]{C_StacksLocalSlice} therefore shows
that the punctured closed subscheme cut out by $IB$ is connected:
the relevant strict inequality is $d-2<d-1$. This is precisely the
pullback of the given cover to $W$, up to nilpotents. Consequently
$\pi_1(W)\to\pi_1(U)$ is surjective.

Normalize $C$. Its normalization is finite, and it is a disjoint union
of complete normal local surface domains $B_i$; the residue fields are
$k$ because they are finite over the algebraically closed field $k$.
Their punctures give a finite surjection $W'\to W$. We now explain the
group-theoretic content of finite-proper descent. Choose basepoints on
the finitely many connected components of $W'$, and one transport for
each connected component of $W'\times_W W'$. Effective descent of finite
etale covers is \citet[Expose IX, Theorem 4.7]{C_SGA1}. The restrictions
of a cover to the vertex components, together with these finitely many
transports, determine its descent data: an isomorphism between finite
etale covers on a connected scheme is determined by its value at one
geometric point. Choosing a spanning tree in the finite component graph
leaves only finitely many independent transport generators.

More formally the resulting homomorphism from the free profinite product
of the vertex groups and a finite-rank free profinite group has dense
image: an invariant subset in a finite fibre is invariant under all
vertex monodromies and all transports, hence defines compatible
subcovers upstairs and descends to a subcover downstairs. A connected
cover admits no proper such subset. Compactness turns dense image into
surjectivity. This proves the displayed map, and also shows why mere
normal generation would not have been enough.

For completeness, on $H^1$ a cocycle is controlled by its vertex
restrictions and at most $er$ transport coordinates. For representation
classes over $\mathbf F_q$, choose a representative on every vertex;
relative conjugations and transports contribute at most
$|\operatorname{GL}_r(\mathbf F_q)|^{e+t-1}$. Multiplying the finitely
many surface bounds therefore preserves both the fixed-rank
$q^{c_r}$ form and the quadratic $q^{cr^2}$ form. No quantitative
mixed-characteristic assertion is made here.
\end{proof}

% Reconstructed after the environment replacement of 20 November 2026.
% This is a fresh complete transcription, not a byte-for-byte recovery.
\subsection{Uniform quadratic representation growth on local punctures}

Throughout this subsection, \(R\) is a complete excellent normal local
domain of dimension two, of equal characteristic \(p>0\), with an
algebraically closed coefficient field \(k\) and maximal ideal
\(\mathfrak m\). We put \(U=\operatorname{Spec}(R)\setminus\{\mathfrak m\}\).
A relative maximal Cohen--Macaulay module over
\(R_A=R\otimes_k A\), for a finite-dimensional local Artinian \(k\)-algebra \(A\)
with residue field \(k\), means
a finite \(A\)-flat module whose closed fibre is maximal
Cohen--Macaulay. All tensor products with \(A\) here are already complete,
since \(A\) has finite length over \(k\).

\begin{lemma}[Proved: relative Noether-different parameters and prescribed jets]
\label{D:relative-Noether-different-jets}
There is a system of parameters \(x,y\) of \(R\) with the following
properties. Put \(I_b=(x^b,y^b)R_A\).
For every relative maximal Cohen--Macaulay module \(M_A\),
multiplication by \(x\) and by \(y\) factors through finite free
\(R_A\)-modules. Consequently, if \(N_A\) is any finite \(R_A\)-module
on which \(x,y\) is a regular sequence, every homomorphism
\[
 M_A/I_bM_A\longrightarrow N_A/I_bN_A,\qquad b\geq2,
\]
lifts to a homomorphism \(M_A\to N_A\) agreeing with the prescribed
map modulo \(I_{b-1}\).
The target is not required to be flat over \(A\).

For \(H=\operatorname{Hom}_{R_A}(M_A,N_A)\) and
\(K_c=\operatorname{Hom}_{R_A}(M_A,I_cN_A)\), we also have
\[
 I_cH\subset K_c
 \subset (x^c,y^c,x^{c-1}y^{c-1})H
 \subset I_{c-1}H \qquad(c\geq2).
 \tag{L1}
\]
\end{lemma}

\begin{proof}
A normal local surface is Cohen--Macaulay and is regular away from its
closed point. The sum of its Noether differents is therefore
\(\mathfrak m\)-primary
\cite[Theorem 1.1(1), definition on p.~298 and proof on p.~300]
{DaoTakahashi2015B}. Choose parameters \(x,y\) in that sum.
Only finitely many Noether normalizations \(S_i\cong k[[u_i,v_i]]\)
are needed to express them as sums of elements of the corresponding
differents. The ring \(R\) is finite free over each \(S_i\).

We spell out the factorization because its relative form is essential.
For \(S_i\subset R\), let \(\mu:R\otimes_{S_i}R\to R\) be multiplication,
and choose
\(\eta=\sum_\nu u_\nu\otimes v_\nu\) in
\(\operatorname{Ann}_{R\otimes_{S_i}R}(\ker\mu)\).
If \(z=\mu(\eta)\), the maps
\[
 M_A\longrightarrow R_A\otimes_{S_{i,A}}M_A
       \longrightarrow M_A,\qquad
 m\longmapsto\sum_\nu u_\nu\otimes v_\nu m,\quad
 r\otimes m\longmapsto rm
\]
are \(R_A\)-linear and their composite is \(z\,\operatorname{Id}\).
Here \(S_{i,A}=S_i\otimes_k A\).
The middle module is finite free over \(R_A\): indeed the closed fibre
of \(M_A\) is free over \(S_i\), and a lifted basis gives
\(S_{i,A}^{a}\to M_A\). Its cokernel vanishes by nilpotent Nakayama.
Flatness over \(A\) makes its kernel reduce to zero modulo the maximal
ideal of \(A\); nilpotent Nakayama kills that kernel as well.
Taking direct sums of the displayed factorizations proves the assertion
for \(x,y\). It also shows that these elements annihilate
\(\operatorname{Ext}^1_{R_A}(M_A,N)\) for every target \(N\).

Let \(\bar f\) be the prescribed \(b\)-jet. Reduce first modulo \(x^b\).
The factorization of \(y\,\operatorname{Id}_{M_A}\) reduces to a
factorization through a free \(R_A/(x^b)\)-module. The obstruction to
lifting \(y\bar f\) across multiplication by \(y^b\) therefore vanishes.
We obtain
\[
 g:M_A/x^bM_A\longrightarrow N_A/x^bN_A,\qquad
 g\bmod y^b=y\bar f.
\]
The image of \(g\) is contained in \(y(N_A/x^bN_A)\).
Division by the nonzerodivisor \(y\) gives a map \(h\) agreeing with
\(\bar f\) modulo \(y^{b-1}\). Now the \(x\)-annihilation of
\(\operatorname{Ext}^1_{R_A}(M_A,N_A)\) lifts \(xh\) to a map
\(M_A\to N_A\). Its image is in \(xN_A\); divide by \(x\).
The resulting map agrees with \(\bar f\) modulo
\((x^{b-1},y^{b-1})\), as asserted.

For (L1), the Koszul presentation
\[
 0\longrightarrow N_A\longrightarrow N_A^{\,2}
       \longrightarrow I_cN_A\longrightarrow0
\]
and the same annihilation give
\(xK_c,yK_c\subset I_cH\).
The sequence \(x,y\) is regular on \(H\). For example, if \(yh=xh'\),
regularity of \(y\) on \(N_A/xN_A\) implies that every value of \(h\)
lies in \(xN_A\), and division defines a homomorphism \(h/x\).
The usual two-parameter colon calculation, valid on every module
with this regular sequence, gives
\[
 (I_cH:(x,y))=(x^c,y^c,x^{c-1}y^{c-1})H.
\]
This proves all the inclusions.
\end{proof}

\begin{lemma}[Proved: exact realization of relative Artin module data]
\label{D:relative-Artin-module-realization}
Set \(S=k[[x,y]]\), let \(d=\operatorname{rank}_S R\), and fix \(L\geq3\).
Put \(J_L=(x^L,y^L)\subset S\).
An action of
\[
 R_A/I_L^2 \quad\hbox{on}\quad
 (S_A/J_L^2)^{dr}
\]
extending its given \(S_A/J_L^2\)-module structure has a reduction \(T_A\)
modulo \(I_L\) which is the exact reduction of a relative maximal
Cohen--Macaulay \(R_A\)-module \(M_A\).
Its closed fibre has rank \(r\) over \(R\), and it is free of rank \(dr\)
over \(S_A\). Conversely every rank-\(r\) maximal Cohen--Macaulay
\(R\)-module supplies such data.

Over \(k\), the isomorphism class of \(M\) is determined by
\(M/I_LM\). At fixed \(L\), the action data above are described by
\(O(r^2)\) scalar variables and \(O(r^2)\) polynomial equations of
bounded degree, with constants depending on \(R,L\), not on \(r\).
\end{lemma}

\begin{proof}
Write \(C_2=S_A/J_L^2\), \(C_1=S_A/J_L\), and let \(W_A\) be the given
\(C_2\)-free module. A free \(C_2\)-resolution of \(C_1\), tensored
with the \(C_2\)-free algebra \(R_A/I_L^2\), is a free resolution
of \(R_A/I_L\). Thus
\[
 \operatorname{Tor}^{R_A/I_L^2}_i(R_A/I_L,W_A)
   =\operatorname{Tor}^{C_2}_i(C_1,W_A)=0\quad(i>0).
\]
This is a lifting in the full Tor sense, not just a lifting of an
underlying presentation.
Apply \cite[Proposition 3.6, pp.~18--19, including its proof]
{AuslanderDingSolberg1993B} to the regular sequence \(x^L,y^L\).
It says that \(T_A\) is a direct summand of \(N_A/I_LN_A\), where
\[
 N_A=\Omega^2_{R_A}(T_A)\oplus P_A
\]
and \(P_A\) is finite free. Since \(T_A\) is \(A\)-free, the two
syzygies are \(A\)-flat. The closed fibre of \(N_A\) is maximal
Cohen--Macaulay. The basis-lifting argument in the preceding lemma
therefore makes \(N_A\) free over every \(S_{i,A}\).

Let \(\bar e\) be the projector onto \(T_A\) in \(N_A/I_LN_A\).
The prescribed-map lemma lifts it to \(f\in\operatorname{End}_{R_A}(N_A)\)
agreeing modulo \(I_{L-1}\).
Hence \(f^2-f\) lies in the closed two-sided ideal
\[
 K=\{h:h(N_A)\subset I_{L-1}N_A\}.
\]
By (L1), \(K\subset I_{L-2}\operatorname{End}_{R_A}(N_A)\).
It is topologically nilpotent in this complete endomorphism algebra.
The elementary idempotent Newton iteration therefore produces an
idempotent \(e\) with \(e-f\in K\).
Modulo \(I_L\), the projectors \(e'\) and \(\bar e\) differ by the
nilpotent ideal of maps whose images lie in
\(I_{L-1}N_A/I_LN_A\). Idempotents congruent modulo a nilpotent ideal
are conjugate: explicitly
\[
 u=e'\bar e+(1-e')(1-\bar e)
\]
is congruent to \(1\) modulo that ideal and satisfies
\(e'u=u\bar e\). Consequently
\(M_A=eN_A\) has reduction isomorphic to \(T_A\) exactly, not merely
to a smaller truncation. It is a direct summand of a module free over
every \(S_{i,A}\), hence is free over those local rings. Its
\(S_A\)-rank is \(dr\), by reduction, and its closed-fibre \(R\)-rank
is \(r\).

For uniqueness over \(k\), lift a prescribed isomorphism
\(M/I_LM\cong M'/I_LM'\) by the preceding lemma.
The lift is an isomorphism modulo \(\mathfrak m\), hence is surjective.
The two modules are \(S\)-free of the same rank, so it is an isomorphism.
The converse is reduction of the \(S\)-free module \(M\).

Finally choose an \(S\)-basis \(z_1,\ldots,z_d\) of \(R\).
The action is a fixed finite list of \(dr\) by \(dr\) matrices over
the fixed finite-dimensional algebra \(S/J_L^2\).
The multiplication table of \(R\), the unit equation, and the
\(S/J_L^2\)-linearity give \(O(r^2)\) equations of degree at most two
in \(O(r^2)\) scalar coordinates. No determinant inversion of
rank-dependent degree is needed.
\end{proof}


\begin{lemma}[Proved: quadratic witnesses for cyclic Frobenius data]
\label{D:quadratic-cyclic-Frobenius-witnesses}
There are constants \(c,D\), depending only on \(R\), such that for
every \(r,f\geq1\) there is a closed affine \(k\)-scheme \(V_{r,f}\)
in at most \(cfr^2\) scalar coordinates, defined by at most \(cfr^2\)
equations of degrees at most \(D\), with the following properties.
Every \(k\)-point gives a cyclic sequence of rank-\(r\) reflexive modules
and isomorphisms on \(U\)
\[
 F_U^*(M_j|_U)\ \cong\ M_{j+1}|_U
 \qquad(j\in\mathbf Z/f\mathbf Z).
 \tag{L2}
\]
Every such sequence occurs. More strongly, over every local Artinian
\(k\)-algebra \(A\), every point of \(V_{r,f}(A)\) gives actual
relative maximal Cohen--Macaulay modules and such a unit Frobenius
cycle on \(U_A\). The actions at a fixed Artin depth are exactly the
reductions of these modules.
\end{lemma}

\begin{proof}
Fix once and for all
\[
 B=3,\qquad L=pB=3p,\qquad n=p(B-1)=2p,\qquad
 \mathfrak a=(x^p,y^p).
 \tag{L3}
\]
Then \(I_n\subset\mathfrak m\mathfrak a\).
For every reflexive \(M\), the natural map
\(F_R^*M\to F_R^{[*]}M\) has cokernel annihilated by \(\mathfrak a\).
Here is a proof of the uniformity. Apply Frobenius to a factorization
of \(x\operatorname{Id}_M\) through a finite free module.
It factors \(x^p\operatorname{Id}_{F^*M}\) through a free module.
The map from \(F^*M\) to that free module extends to the reflexive hull,
because on \(U\) it is a map between vector bundles and maps into
a free module extend across the closed point. The resulting identity,
first on \(U\) and then by reflexivity, shows
\(x^pF^{[*]}M\subset\operatorname{im}(F^*M)\).
The same argument applies to \(y\).

At each stage take the action matrices of
Lemma~\ref{D:relative-Artin-module-realization} at depth \(I_L^2\).
Their reductions \(T_j\) modulo \(I_L\) are free of rank \(dr\) over
\(S/J_L\).
An arrow
\[
 a_j:T_j/I_BT_j\longrightarrow T_{j+1}
\]
is required to be \(p\)-semilinear. Its matrix \(A_j\) has entries
in \(S/J_L\), and the equations for the fixed \(S\)-basis \(z_i\)
of \(R\) are
\[
 A_j\,\rho_j(z_i)^{[p]}=\rho_{j+1}(z_i)^p A_j.
 \tag{L4}
\]
The notation on the left means entrywise Frobenius; the power on
the right is matrix multiplication and represents multiplication
by \(z_i^p\). The \(S\)-semilinearity is built into the matrix
description. In particular the source is killed by \(I_B\) because
\(I_B^{[p]}=I_L\) kills the target.

For each \(j\) and \(b\in\{x^p,y^p\}\), introduce matrices
\(X_{i,j,b}\) over \(S/J_n\), and impose
\[
 \rho_{j+1}(b)=
   \sum_{i=1}^d \rho_{j+1}(z_i)\,A_jX_{i,j,b}
       \pmod {I_n}.
 \tag{L5}
\]
This says that the image of multiplication by \(\mathfrak a\)
is contained, at depth \(n\), in the \(R\)-span of the semilinear
arrow. The \(z_i\) span \(R\) over \(S\), so these matrices are
precisely inclusion witnesses, with no omitted coefficients.
Equations (L4) have degree at most \(p+1\), and (L5) degree at
most three. All coefficient rings and the number \(d\) are fixed.
Together with the quadratic action equations, there are
\(O(fr^2)\) variables and equations of degree bounded independently
of \(r,f\).

We must prove that every witness, including a nonreduced one, is
effective. Let \(A\) be local Artinian.
Realize its module actions by the preceding lemma, obtaining
\(M_{j,A}\). View the arrow as an \(R_A\)-linear map into
\(F_*M_{j+1,A}/I_BF_*M_{j+1,A}\).
The latter quotient has underlying module \(M_{j+1,A}/I_LM_{j+1,A}\).
The target \(F_*M_{j+1,A}\) is finite and \(x,y\)-regular:
\(R_A\) is \(F\)-finite, and \(x,y\) act on it as \(x^p,y^p\).
It need not be \(A\)-flat. The prescribed-map lemma nevertheless
lifts the arrow to an actual semilinear map, agreeing in the
underlying target modulo \(I_n\).

Let \(C\) be the cokernel of its linearization
\(F^*M_{j,A}\to M_{j+1,A}\).
The lifted equations (L5) imply
\[
 \mathfrak a C\subset I_nC\subset\mathfrak m\mathfrak a C.
\]
Nakayama gives \(\mathfrak a C=0\). On \(U_A\),
\(\mathfrak a\) is the unit ideal, so the map is surjective there.
Both modules are vector bundles of rank \(r\) on \(U_A\), hence
the map is an isomorphism. This proves (L2) over \(A\).
In particular no flatness of the Frobenius-restriction target
was used.

Conversely, a unit cycle gives reflexive modules at all stages.
The uniform defect calculation above supplies (L5), and reduction
of its actual maps supplies (L4). Since each \(M_j\) is \(S\)-free,
it supplies the required action at depth \(I_L^2\).
Thus every desired cycle has a witness.
\end{proof}

\begin{lemma}[Proved: uniqueness of a unit structure on a fixed module]
\label{D:local-complete-Lang}
Let \(q=p^f\), and let \(M\) be a reflexive \(R\)-module with a unit
\(q\)-Frobenius structure on \(U\).
Any two such structures on the same underlying module are isomorphic.
\end{lemma}

\begin{proof}
Put \(\mathcal E=\operatorname{End}_R(M)\), a finite complete
\(R\)-algebra. A chosen unit structure induces a multiplicative
additive \(q\)-semilinear map
\(\tau:\mathcal E\to\mathcal E\), by Frobenius pullback and conjugation.
Maps on \(U\) extend to \(M\), so this is a map on the whole algebra.
It satisfies
\[
 \tau(\mathfrak m^a\mathcal E)\subset
       \mathfrak m^{qa}\mathcal E.
 \tag{L6}
\]
The finite-dimensional \(k\)-algebra
\(\overline{\mathcal E}=\mathcal E/\mathfrak m\mathcal E\)
has a connected smooth unit group: it is the determinant-nonzero
open subset of its underlying affine vector space, using its
left regular representation. The induced \(\bar\tau\) has zero
differential. Every twisted-conjugacy orbit
\[
 g\longmapsto g\,a\,\bar\tau(g)^{-1}
\]
is therefore open, since its orbit map has invertible differential.
Connectedness forces a single orbit. Hence the equation
\(a=g\bar\tau(g)^{-1}\) is solvable modulo \(\mathfrak m\).

The ideal \(\mathfrak m\mathcal E\) lies in the Jacobson radical,
so a solution modulo it lifts to a unit \(g_0\in\mathcal E^\times\).
With the orientation \(g=a\tau(g)\), iterate
\(g_{j+1}=a\tau(g_j)\). By (L6), successive errors tend to zero
\(\mathfrak m\)-adically with exponents multiplied by \(q\).
The complete algebra supplies a unit limit solving the equation.
This conjugates the two unit structures. We have used Lang's
argument only on the finite-dimensional residue algebra, not on
the infinite-dimensional \(k\)-space \(\mathcal E\).
\end{proof}

\begin{theorem}[Proved: uniform quadratic local representation growth]
\label{D:uniform-local-representation-growth}
There exists \(C_R>0\) such that, for every \(r,f\geq1\) and \(q=p^f\),
\[
 \#\bigl(
 \operatorname{Hom}_{\mathrm{cts}}(\pi_1(U),\operatorname{GL}_r(\mathbf F_q))
                  /\operatorname{GL}_r(\mathbf F_q)\bigr)
       \leq q^{C_Rr^2}.
 \tag{L7}
\]
The constant is independent of both the rank and the finite coefficient
field.
\end{theorem}

\begin{proof}
Finite \(\mathbf F_q\)-local systems on \(U\) are equivalent to
vector bundles equipped with a unit \(q\)-Frobenius structure.
For completeness, after trivializing the bundle the fixed-vector
equation is \(v=A v^{[q]}\), with \(A\) invertible. Multiplication
by \(A^{-1}\) makes the leading monomials \(v_i^q\), so the
fixed-vector scheme is finite locally free of degree \(q^r\).
Its Jacobian is invertible, hence it is finite etale.
Over an algebraically closed geometric fibre its fixed vectors
form an \(r\)-dimensional \(\mathbf F_q\)-space spanning the
ambient vector space: apply Frobenius to a shortest normalized
linear relation and subtract to see that the coefficients belong
to \(\mathbf F_q\). The invertible fixed frames therefore form
the required finite etale \(\operatorname{GL}_r(\mathbf F_q)\)-torsor.
Conversely the associated bundle of such a torsor has its evident
unit structure. Its vector bundle extends uniquely to a finite
reflexive module on the normal local surface.

A \(q\)-unit structure can be written as an \(f\)-cycle of
\(p\)-Frobenius structures by taking the successive reflexive
Frobenius pullbacks. Thus every representation class is represented
on \(V=V_{r,f}\) of the preceding lemma.
We prove that each representation class is an open and closed
part of the underlying topological space of \(V\).

The key assertion is stronger than constancy on geometric points.
Fix a closed point \(v\in V(k)\), and consider any point over a
local Artinian \(k\)-algebra \(A\) reducing to \(v\).
Its genuine relative unit cycle consists of \(M_{j,A}\).
Choose a multiple \(e\) of \(f\) so large that the \(p^e\)-power
map kills the maximal ideal of \(A\).
Absolute Frobenius on \(R_A\) then factors through the closed
fibre followed by its \(p^e\)-Frobenius.
Iterating the unit cycle identifies
\[
 M_{j,A}|_{U_A}\cong M_{j,0}|_U\otimes_k A.
 \tag{L8}
\]
Here we use the iterated special-fibre unit structure to identify
the pulled-back special bundle with \(M_{j,0}|_U\).
Both sides satisfy Hartogs extension across \(V(\mathfrak m R_A)\):
the regular sequence \(x,y\) gives vanishing of local cohomology
in degrees zero and one. Hence (L8) extends to an isomorphism
of the complete relative modules.

Only their canonical reductions modulo \(I_L\) are asserted to
be constant. The chosen action modulo \(I_L^2\), the auxiliary
inclusion witnesses, and the chosen arrow lifts can still vary.
This distinction is crucial.

Let \(\mathcal T\) be the affine scheme of tuples of actions of
\(R/I_L\) on the fixed free \(S/J_L\)-module of rank \(dr\).
Reduction of our witness data gives a morphism
\(\pi:V\to\mathcal T\).
The algebraic group
\[
 \mathcal G=
 \operatorname{Res}_{(S/J_L)/k}
          \operatorname{GL}_{dr}^{\,f}
\]
acts by changing the bases at the \(f\) stages.
The preceding Artinian assertion says that every deformation of
\(v\), at every finite Artin order, has its image in the ordinary
\(\mathcal G\)-orbit of \(\pi(v)\).
These ordinary orbits are locally closed.
One can also see their scheme structure directly: the stabilizer
is the unit group of a finite-dimensional endomorphism algebra,
hence is smooth, and the orbit is the corresponding homogeneous
locally closed subscheme.

Choose an ambient open neighbourhood of \(\pi(v)\) in which its
orbit is closed. The ideal of that orbit, pulled back to the
Noetherian local ring \(\mathcal O_{V,v}\), vanishes in every
quotient by a power of its maximal ideal, by (L8).
Krull intersection makes it zero in \(\mathcal O_{V,v}\).
Its finite set of generators therefore vanishes on a neighbourhood
of \(v\). This neighbourhood maps into the orbit.

At closed points, belonging to this orbit means that the canonical
module tuples modulo \(I_L\) are isomorphic. The exact realization
lemma lifts the isomorphisms to the reflexive modules.
Lemma~\ref{D:local-complete-Lang} says that the underlying periodic
module determines its unit \(q\)-structure up to isomorphism.
Thus the neighbourhood just constructed consists of one
representation class.

The class neighbourhoods for distinct classes are disjoint:
a nonempty intersection in a finite-type \(k\)-scheme would have
a closed point. Their union contains all closed points, and hence
all of \(V\), since \(V\) is Jacobson. Quasicompactness gives a
finite clopen partition by classes. In particular the number of
classes is at most the number of irreducible components of \(V\).

Finally, a closed affine scheme cut out by \(m\) equations of
degree at most \(D\) has at most \(D^m\) irreducible components.
Indeed, start with affine space and intersect with the equations
one at a time. A component contained in a hypersurface is retained;
otherwise the hypersurface degree formula bounds the sum of
the degrees of the new components by \(D\) times its old degree.
Induction proves the bound, also when dimensions vary.
Here \(m\leq cfr^2\), so the number is at most \(D^{cfr^2}\).
Increasing \(C_R\) gives (L7).
\end{proof}

\begin{remark}[Scope of the counting theorem]
\label{D:uniform-local-growth-scope}
The theorem bounds isomorphism classes, equivalently conjugacy classes
of representations. It does not assert a bound on the number of raw
homomorphisms without the additional conjugation factor, nor a
uniform bound on generators of the full fundamental group.
The proof uses exact liftability of all Artinian witnesses and
openness of their class fibres. A finite-dimensional parameter
space, or mere determinacy of the special-fibre jet, would not
by itself justify the component count.
\end{remark}

\begin{corollary}[Proved: the bound for higher-dimensional local punctures]
\label{D:uniform-local-growth-all-dimensions}
Let \(A\) be a complete excellent normal local domain of equal
characteristic \(p>0\), of dimension at least two, with algebraically
closed residue field. Its punctured spectrum has a constant \(C_A\)
for which the bound \(q^{C_A r^2}\) holds for every \(q=p^f\) and \(r\).
\end{corollary}

\begin{proof}
The local surface reduction of
\ref{D:local-surface-reduction} gives a surjection
\[
 \left(\mathop{*}_{i=1}^{t}\pi_1(U_i)\right)
                *\widehat F_e\longrightarrow
       \pi_1(\operatorname{Spec}(A)\setminus\{\mathfrak m_A\}),
\]
where the \(U_i\) are finitely many complete normal local surface
punctures with the same algebraically closed residue field.
For each factor choose a representative of each representation
class. Simultaneous conjugacy leaves at most
\(\lvert\operatorname{GL}_r(\mathbf F_q)\rvert^{e+t-1}\)
choices for relative conjugations and for the free generators
(with the harmless empty-factor case treated separately).
Multiplying the finitely many surface bounds and using
\(\lvert\operatorname{GL}_r(\mathbf F_q)\rvert\leq q^{r^2}\)
proves the assertion. Pullback along the displayed surjection
is injective on representation classes.
\end{proof}



% Restored proof avenue: the local atlas is superseded quantitatively by
% D:uniform-local-representation-growth, but its proof is retained.
\subsection{The local minimal-extension atlas}

\begin{lemma}[Proved: a constructible atlas for reflexive modules and Frobenius]
\label{D:local-reflexive-constructible-atlas}
Let \(R,k,U\) be as in
\ref{D:uniform-local-representation-growth}, and fix a rank \(r\).
There is a finite-type parameter scheme \(T_r\) whose geometric
points give all rank-\(r\) reflexive \(R\)-modules, with a
constructible relation that is exactly ordinary isomorphism.
There is also a constructible relation that says exactly
\(M'\cong F_R^{[*]}M\).
No uniform bound on the complexity of the resulting constructible
quotient as \(r\) varies is asserted by this lemma.
\end{lemma}

\begin{proof}
Choose a projective regular resolution
\(\pi:Y\to\operatorname{Spec}(R)\) whose exceptional curves \(E_i\)
form a simple normal crossings divisor, and choose a positive
exceptional integral cycle \(Z\) with \(Z\cdot E_i<0\) for all \(i\).
Put \(I=\mathcal O_Y(-Z)\) and \(Y_n=V(I^n)\).
The \(I\)-adic and \(\mathfrak m\mathcal O_Y\)-adic topologies are
cofinal. Each \(Y_n\) is projective over \(k\).

Every reflexive \(M\) has a minimal vector-bundle extension \(F\)
on \(Y\) with the fixed slope windows
\[
 \alpha_i=g(E_i)-1\leq\mu_{\min}(F|_{E_i})
 \leq\mu_{\max}(F|_{E_i})
 \leq\beta_i=g(E_i)-1-E_i^2.
 \tag{J1}
\]
We use \cite[Definition 1.2 and Lemma 1.3, including the
elementary-transformation proof]{C_Langer2026}.
That proof applies over this complete local base: the relative
Euler characteristics used to terminate the modifications are
finite nonnegative integers, and all modifications take place
on the proper exceptional curves.
For the bundles \(F_Y^*F|_{E_i}\), the bounds of
\cite[Theorem 2.2]{C_Sun2008} and
\cite[Proposition 1.3 and its proof]{C_MehtaPauly2007} give the
enlarged windows
\[
 p\alpha_i-(p-1)\max(g(E_i)-1,0),\qquad
 p\beta_i+(p-1)\max(g(E_i)-1,0).
 \tag{J2}
\]
For genus zero the same assertion follows from splitting on
\(\mathbf P^1\). The semistable inequalities extend to arbitrary
bundles by their Harder--Narasimhan filtrations. All windows are
independent of \(r\).

\par\smallskip\noindent\textbf{Uniform poles and lifting depth.}\quad
Order the multiset of components of \(Z\).
If \(D_j\) is the prefix cycle preceding a copy of \(E_i\), the
resulting filtration of \(\mathcal O_Z\) has that layer
\(\mathcal O_{E_i}(-D_j)\).
For \(F,G\) satisfying any of the finitely many permitted windows,
\[
 H^0\bigl(E_i,\operatorname{Hom}(F_i,G_i)
                ((bZ-D_j)|_{E_i})\bigr)=0
\]
as soon as
\[
 \mu_{\max}(G_i)+(bZ-D_j)\cdot E_i<\mu_{\min}(F_i).
\]
This holds uniformly for all sufficiently large \(b\), since
\(Z\cdot E_i<0\).
The restriction sequence then shows that the global sections of
\(\operatorname{Hom}(F,G)(bZ)\) stabilize as \(b\) increases.
Every morphism on \(U\) has some finite pole order, by localization
along the Cartier divisor and quasicompactness. Thus there is one
\(B\), independent of ranks, such that every such morphism
extends as \(F\to G(BZ)\). This applies to inverse maps as well.

Choose \(N>2B\) sufficiently large that all the following \(H^1\)
groups vanish for \(n\geq N\):
\[
 H^1\bigl(Z,\operatorname{Hom}(F,G)(BZ)\otimes I^n|_Z\bigr),
 \qquad
 H^1\bigl(Z,\operatorname{End}(F)\otimes I^n|_Z\bigr).
 \tag{J3}
\]
The fixed layer filtration reduces this to smooth curves.
Serre duality identifies
\[
 H^1(E_i,\operatorname{Hom}(F_i,G_i)\otimes L)^\vee
     =\operatorname{Hom}(G_i,F_i\otimes L^{-1}\omega_{E_i}).
\]
It vanishes if
\[
 \mu_{\min}(G_i)+\deg L>
          \mu_{\max}(F_i)+2g(E_i)-2.
\]
For \(L=((B-n)Z-D_j)|_{E_i}\), its degree tends to positive
infinity with \(n\). This proves (J3), also for the Frobenius
windows (J2), without a tensor-semistability assertion.

Every vector bundle \(F_N\) on \(Y_N\) with windows (J1) lifts
successively to all \(Y_n\). The obstruction is in an \(H^2\)
group on a one-dimensional proper scheme, hence is zero.
After \(N\), the second vanishing in (J3) makes the lift unique
up to an isomorphism extending any prescribed one.
Projective Grothendieck existence algebraizes the compatible
system to a bundle \(F\) on \(Y\).
It is locally free along the exceptional fibre by faithful
flatness of completion, and everywhere by properness: a
nonempty closed nonfree locus over a local base would meet
the closed fibre. We obtain \(M=(\pi_*F)^{**}\).
This module is independent of the choice of lift up to
isomorphism. The same argument works after an algebraically
closed coefficient-field extension.
The obstruction and algebraization inputs are the
vector-bundle deformation statements and projective existence
theorem in
\cite[Tags 08MF, 08MD, 0885]{StacksLocalAtlasB}.
The first vanishing in (J3) lifts every bounded-pole map at depth
\(N\) to a map between the chosen actual lifts.

\par\smallskip\noindent\textbf{Exact isomorphism tests.}\quad
At depth \(N\), ask for maps
\[
 h_N:F_N\to G_N(BZ),\qquad k_N:G_N\to F_N(BZ)
\]
whose two products equal the canonical inclusions into the
\(2BZ\) twists.
Lift \(h_N,k_N\) by (J3).
As rational endomorphisms their products are
\(\operatorname{Id}+\varepsilon\), with
\(\varepsilon\in I^{N-2B}\operatorname{End}\).
These are invertible on the formal completion, and their
inverses algebraize by full faithfulness of existence.
Equivalently their noninvertibility loci are closed and proper
over the local base and miss its closed fibre.
Thus \(h\) is an isomorphism on \(U\), so \(M\cong M'\).
Conversely every module isomorphism gives maps with poles at most
\(BZ\), hence these equations at depth \(N\).
The finite test is therefore exact.

For Frobenius, replace \(F_N\) by \(F_{Y_N}^*F_N\).
The enlarged windows and the same lifting argument give an
equally exact test for \(M'\cong F_R^{[*]}M\).
Absolute Frobenius commutes with restriction and with base
change of a geometric point; its effect on coefficient
coordinates is their \(p\)-th power.

\par\smallskip\noindent\textbf{Finite-type parameter spaces.}\quad
Embed \(Y_N\) in a fixed projective space, with hyperplane \(H\).
The filtration of \(\mathcal O_{Y_N}\) by the component multiset
of \(NZ\) has finitely many fixed layers.
Choose \(t\) such that
\[
 \alpha_i+((t-1)H-D_j)\cdot E_i>2g(E_i)-2
\]
for each layer. Then \(H^1(Y_N,F_N(t-1))=0\).
All higher cohomology groups vanish because its support has
dimension one. Thus the ambient sheaf is \(t\)-regular and
\(F_N(t)\) is globally generated
\cite[Tag 08A8]{StacksLocalAtlasB}.
Riemann--Roch on the fixed layers gives
\(h^0(F_N(t))\leq c r\), and gives finitely many Hilbert
polynomials at fixed \(r\).

Take the finite union of Quot schemes of quotients of
\(\mathcal O_{Y_N}(-t)^a\), \(a\leq cr\), with those polynomials.
The locally free rank-\(r\) locus is open.
The conditions (J1) are open as well, using the finitely many
bad quotient Hilbert polynomials in this bounded family.
One uses all coherent quotients for their proper Quot images;
torsion only increases the degree and does not invalidate
the slope tests. This gives a finite-type scheme \(T_r\) with
a universal bundle. See \cite[Tag 0DPA]{StacksLocalAtlasB} for
bounded Quot.

On \(T_r\times T_r\), the bounded-pole pairs of maps are
parametrized by affine schemes of finite presentation, by
proper Hom representability
\cite[Tag 08K6]{StacksLocalAtlasB}.
The two inverse-product equations are closed conditions.
Their image is constructible and, by the exact lifting argument,
is precisely module isomorphism. Using the Frobenius pullback
in the source gives the other required constructible relation.
No universal algebraization of the entire parameter family
has been assumed: effectivity and equivalence were proved
on its geometric points, which is what these constructible
relations require.
\end{proof}

\begin{remark}[The independent fixed-rank route]
Applying \ref{D:effective-constructible-quotient} to this atlas
gives a constructible space of module classes and a constructible
reflexive-Frobenius function on it. Its periodic classes are
the finite-field local systems by
\ref{D:local-complete-Lang}.
Gabber's local finiteness theorem, in the precise form
\cite[Theorem 1.4]{DegliseGabber2007LocalB}, makes each finite-target
set finite. Then \ref{D:constructible-periodic-count} gives the
fixed-rank exponential bound. This argument predates and is
quantitatively superseded by the quadratic witness construction.
It is retained because it separates constructibility from the
additional rank-uniform complexity argument.
\end{remark}


% Reconstructed after the environment replacement; this restores the
% constructible-counting route independently of the quadratic witness proof.
\subsection{Constructible quotients and fixed-rank periodic counts}

\begin{lemma}[Proved: effective constructible quotients]
\label{D:effective-constructible-quotient}
Let \(k\) be algebraically closed, let \(T\) be a constructible set
of finite type over \(k\), and let \(\mathcal R\subset T\times T\)
be a constructible equivalence relation. There is a constructible
set \(Q\) and a constructible surjection \(c:T\to Q\) whose fibres
are exactly the equivalence classes. A constructible relation
on \(T\) inducing a partially defined single-valued operation on
classes descends to a partial function on \(Q\) with constructible
graph. The ground field need not be algebraic over a finite field.
\end{lemma}

\begin{proof}
The effective quotient assertion is
\cite[Theorem 2; constructive proof by Proposition 3.1, pp.~10--11]
{Beke2011FixedRankB}. Its scope is the constructible category over
an arbitrary field; the finite-field hypothesis in that paper's
zeta-function applications is not a hypothesis of the quotient
theorem. One concrete construction chooses, uniformly in each
nonempty constructible fibre, a nonempty finite subset of bounded
cardinality, with equal fibres receiving equal subsets. Induction
on the number of affine coordinates gives such a choice. Encoding
an unordered finite tuple by symmetric polynomial invariants
then distinguishes the disjoint equivalence classes.
For the last assertion, the image of the given relation under
\(c\times c\) is constructible by Chevalley's theorem. By the
assumed compatibility and uniqueness it is the graph of the
desired partial function.
\end{proof}

\begin{lemma}[Proved: periodic points of a partial constructible function]
\label{D:constructible-periodic-count}
Let \(Q\) be a constructible set over an algebraically closed field
\(k\), and let \(\Phi\) be a partially defined function on \(Q(k)\)
with constructible graph. Suppose that the set of points with a
defined \(f\)-step orbit returning to the initial point is finite
for each \(f\geq1\). There is a constant \(A\geq1\), depending on
this one graph, such that
\[
 \#\operatorname{Fix}(\Phi^f)\leq A^f\qquad(f\geq1).
 \tag{F1}
\]
No finite field of definition is required.
\end{lemma}

\begin{proof}
A constructible set has a finite disjoint partition by affine
locally closed sets. Adding finitely many coordinate tags embeds
it, in the constructible category, in one affine space
\(\mathbb A^a_k\). Write the graph as a union of \(L\) basic locally
closed pieces
\[
 \Gamma_i=\{(x,y):
 F_{i1}(x,y)=\cdots=F_{is_i}(x,y)=0,\quad g_i(x,y)\ne0\}.
\]
The finite products of the inequalities defining a piece may be
used for \(g_i\). Choose \(s\geq s_i\) and \(D\geq1\) bounding all
degrees, including those of \(u g_i-1\).
For a word \(\boldsymbol i=(i_0,\ldots,i_{f-1})\), impose the
closed affine equations
\[
 F_{i_jb}(x_j,x_{j+1})=0,\qquad
 u_jg_{i_j}(x_j,x_{j+1})=1,\qquad x_f=x_0.
\]
Use only the \(f\) variables \(x_0,\ldots,x_{f-1}\), so the last
equality is substituted rather than introducing more coordinates.
There are at most \((s+1)f\) equations of degree at most \(D\).
A solution is determined by \(x_0\): single-valuedness determines
each successive \(x_j\), and every \(u_j\) is uniquely determined.
The word scheme therefore has finitely many geometric points and
is zero-dimensional.

For clarity, the degree estimate needs no elimination of an
iterated graph. Intersect affine space successively with the
displayed equations. If an irreducible component is contained
in the next hypersurface, retain its degree. Otherwise the
hypersurface degree formula bounds the sum of the degrees of the
new components by \(D\) times its old degree. This formula follows
by projective closure and the exact sequence for multiplication
by a homogeneous equation, or directly from Hilbert polynomials.
Thus the word scheme has at most \(D^{(s+1)f}\) geometric points.
Summing over its \(L^f\) possible words gives (F1) with
\(A=LD^{s+1}\). A point may occur in several words, which only
increases this upper bound. Partiality causes no change:
a cycle solution uses only points where the next arrow is defined.
\end{proof}

\begin{corollary}[Proved: local fixed-rank representation growth]
\label{D:local-fixed-rank-representation-growth}
For a complete excellent normal local surface over an algebraically
closed field of characteristic \(p>0\), and for each fixed \(r\),
there is \(C_{R,r}\) such that the number of rank-\(r\)
\(\mathbf F_{p^f}\)-local systems on its puncture, up to isomorphism,
is at most \(p^{C_{R,r}f}\).
\end{corollary}

\begin{proof}
The uniform theorem \ref{D:uniform-local-representation-growth}
gives this with \(C_{R,r}=C_Rr^2\).
The earlier constructible proof uses the local atlas recorded
below, applies the effective quotient to ordinary reflexive-module
isomorphism, and uses reflexive Frobenius as its constructible
operation. The complete-local Lang lemma
\ref{D:local-complete-Lang} identifies periodic module classes
with unit structures. Finiteness for each fixed finite target
then allows \ref{D:constructible-periodic-count} to count cycles.
This explains the role of the quotient lemma in the independently
retained fixed-rank route; the stronger uniform proof does not
assume a complexity bound for that quotient.
\end{proof}

\begin{corollary}[Proved under stated hypotheses: simple-product quotients]
\label{D:bounded-matrix-simple-product}
Let \(G\) have the fixed-rank representation-class bounds of the
preceding corollary. Fix \(r_0\) and \(\epsilon>0\), and let
\(\mathcal S\) be a collection of pairwise nonisomorphic nonabelian
finite simple groups such that every \(S\in\mathcal S\) admits
a faithful representation
\[
 S\hookrightarrow\operatorname{GL}_{r_0}(\mathbf F_{p^{f_S}}),
 \qquad |S|\geq p^{\epsilon f_S}.
\]
Every quotient of \(G\) which is a product of groups in
\(\mathcal S\) is topologically finitely generated.
This includes any family for which these fixed matrix-size and
field-size inequalities have been verified.
\end{corollary}

\begin{proof}
Write such a quotient as \(\prod_{S\in\mathcal S}S^{m_S}\).
Its coordinate epimorphisms have distinct kernels.
Composing them with the fixed faithful representation of \(S\)
gives distinct conjugacy classes of representations: conjugate
representations have equal kernels. Hence
\[
 m_S\leq p^{C_{r_0}f_S}\leq |S|^{C_{r_0}/\epsilon}.
 \tag{F2}
\]
We recall the precise finite-group input used to turn (F2) into
generation. In the proof of Proposition 2.6(ii), printed p.~12,
of \cite{LucchiniMarionTracey2020D}, used in
\ref{D:cohomology-crown-criterion}, the number of
\(\operatorname{Aut}(S)\)-orbits of generating \(s\)-tuples is at
least
\[
 \frac{|S|^{s-1}}{2|\operatorname{Out}(S)|}\qquad(s\geq2),
\]
and the same passage gives two-generation of \(S\); p.~13 records
\(|\operatorname{Out}(S)|\leq\log |S|\).
Thus one fixed sufficiently large \(s\) supplies at least \(m_S\)
distinct generating-tuple orbits for every \(S\).
Choose those tuples at the coordinate copies.
Goursat's lemma shows that they generate every finite subproduct:
for equal simple factors the distinct automorphism orbits exclude
diagonal relations, and nonisomorphic simple factors have no
common simple quotient. Their \(s\) assembled elements therefore
generate a dense subgroup of the entire product.
This proves the assertion without claiming anything about
unbounded matrix dimension or general monolithic crowns.
\end{proof}


% Reconstruction of the fully audited boundedness and fixed-rank arguments.
\begin{proposition}[Proved: Bounded minimal extensions of periodic reflexive sheaves]
\label{D:global-periodic-reflexive-boundedness}
Let $S$ be a normal proper surface over an algebraically closed field
of characteristic $p>0$. Fix a projective resolution $\pi:Y\to S$,
with smooth projective $Y$, an ample divisor $H$, and an effective
exceptional divisor $Z$ with $Z\cdot E_i<0$ on every exceptional
component. For reflexive sheaves $E$ of rank $r$ satisfying
$F_S^{[p^f]*}E\simeq E$ for some $f\ge1$, their minimal extensions
$F$ to $Y$ have the following bounds, with constants independent of
$r$ and $f$:
\begin{enumerate}
\item their restrictions to the exceptional curves have fixed upper
and lower slope windows;
\item $|\chi(Y,F)|\le Cr$, and $c_1(F)$ is numerically an exceptional
linear combination whose coefficients have absolute values at most $Cr$;
\item $\mu_{\max,H}(F)$ and $\mu_{\max,H}(F^\vee)$ are bounded by a
fixed constant, and the analogous bounds hold after restriction to a
fixed sufficiently ample smooth curve;
\item for each fixed $r$, these extensions form a bounded family.
\end{enumerate}
Ordinary isomorphism of their underlying reflexive sheaves, and the
one-step relation $F_S^{[*]}E\simeq E'$, are constructible relations on
one finite-type parameter space for each fixed rank.
\end{proposition}
\begin{proof}
We use minimal extensions in the sense of
\citet[Definition 1.2 and Lemma 1.3]{C_Langer2026}. On an exceptional
curve $E_i$ of genus $g_i$, their slopes lie in a fixed interval, which
we may take to have endpoints $g_i-1$ and $g_i-1-E_i^2$. The important
point is that the endpoints do not depend on the rank or on the
Frobenius period.

We record the uniform extension calculation. Filter $\mathcal O_Z$ by
a fixed finite sequence whose quotients are line-bundle twists on the
reduced exceptional components. For two bundles with the stated slope
windows, Serre duality and the elementary implication
\[
 \mu_{\min}(A)>\mu_{\max}(B)\quad\Longrightarrow\quad
 \operatorname{Hom}(A,B)=0
\]
give, for all sufficiently large fixed $n$,
\[
 H^0\bigl(Z,\operatorname{Hom}(F,G)\otimes\mathcal O_Z(nZ)\bigr)=0,
 \qquad
 H^1\bigl(Z,\operatorname{Hom}(F,G)\otimes\mathcal O_Z(-nZ)\bigr)=0.
\]
We apply the inequalities to the individual component layers; we do
not use an incorrect assertion that tensor products preserve
semistability in characteristic $p$. The first vanishing says that
every map away from the exceptional set has a fixed pole bound $BZ$.
Indeed such a map first extends with some finite pole order by Cartier
localization, and the exact sequences for successive $Z$-layers lower
that order to $B$. The second vanishing gives a fixed finite stage
which determines formal maps and bundles. Existence of lifts of bundles
uses the vanishing of $H^2$ on the one-dimensional exceptional support;
uniqueness and lifting of prescribed maps use the displayed $H^1$
vanishing. We apply formal GAGA over the completed local ring at
each exceptional image; no algebraization from completion along
an arbitrary divisor of the global surface is required. It algebraizes them
\citep[Tags 0885, 08MD, 08MF]{C_StacksBoundedReflexive}.

The same calculation with one bundle fixed controls the discrepancy
between $E$ and $\pi_*F$. Specifically, the extension identity off the
exceptional set gives
\[
 E=\pi_*F(BZ),\qquad
 E/\pi_*F\hookrightarrow\pi_*\bigl(F(BZ)/F\bigr).
\]
The dimensions of the cohomology of the fixed thickening on the right
are $O(r)$, by its fixed component filtration and the slope windows.
For $R^1\pi_*F$, formal functions and the second vanishing show that
the inverse system stabilizes at one fixed thickening $NZ$; its
$H^1$ is again $O(r)$. Also $R^2\pi_*F=0$, since the fibres have
dimension at most one. Thus
\[
 |\chi(Y,F)-\chi(S,E)|\le Cr.
\]
Frobenius scaling gives numerical $c_1(E)=0$ and
$\operatorname{ch}_2(E)=0$, independently of the period. Proper
Riemann--Roch with the rank-linear correction yields
$|\chi(S,E)|\le Cr$
\citep[Theorem 4.5(2)]{C_Langer2025}
\citep[Theorem 0.1]{C_Langer2026}.

Put $q=p^f$. The isomorphism on $Y\setminus Z$ between $F_Y^{q*}F$
and $F$ gives an actual line-bundle identity
\[
 (\det F)^{\otimes(q-1)}\simeq\mathcal O_Y(D_q)
\]
for an exceptional divisor $D_q$. Hence $c_1(F)$ belongs to the
exceptional numerical span. Its intersection with every $E_i$ is
bounded by a constant times $r$ from the slope windows. The negative
definite exceptional intersection matrix then bounds each coefficient
by $Cr$. This argument uses the Frobenius determinant identity; it
does not assume that numerical triviality downstairs on a nonprojective
surface already identifies the exceptional span upstairs.

The rank-independent one-step slope estimates for Frobenius on a fixed
smooth curve are those of
\citet[Theorem 2.2]{C_Sun2008} and
\citet[Proposition 1.3]{C_MehtaPauly2007}. Iteration bounds the slopes
of the restrictions of $F_Y^{q*}F$ to each exceptional component by
constants times $q$. The preceding pole-lowering argument therefore
extends the periodic isomorphism and its inverse with a pole bound
$B_qZ$, where $B_q\le Cq$. We obtain
\[
 (q-1)\mu_{\max,H}(F)\le B_q Z\cdot H
\]
and the same inequality for the dual. Since $q\ge2$, both slopes are
uniformly bounded. Choose a fixed smooth ample curve not contained in
the exceptional locus or the determinant-zero divisors. Restricting
the same maps gives the same estimates directly on this curve, without
a restriction theorem depending on the rank.

The numerical $c_1$ bounds and the Euler bound give finitely many
Hilbert polynomials for each fixed rank. Langer's boundedness theorem
applies on the fixed smooth projective $Y$
\citep[Theorem 4.2, printed p.~269]{C_LangerBounded2004}.
The full addendum explicitly retains this theorem
\citep{C_LangerAddendum2004}. We therefore choose one finite-type
parameter space with a bounded family of the relevant bundles. We
retain the open slope-window conditions; they can be expressed using
the proper images of finitely many Quot schemes of all coherent
quotients of fixed positive rank and degree below a threshold. Torsion
only raises the degree, and a uniform global-generation bound supplies
a lower bound on the degrees which need be tested. Taking only the
open torsion-free part of a Quot scheme would not justify closedness
of the bad locus.

Finally, bounded-pole Hom schemes on the proper $Y$ are of finite
presentation. Pairs of maps in the two directions, whose products equal
the fixed power of the section defining $Z$, encode ordinary
isomorphisms off $Z$. Their images in the pair of parameter spaces are
constructible. The same construction with $F_Y^*F$ gives the
one-step Frobenius relation: this source has fixed exceptional windows
by the one-step slope bound, so the required pole order is still fixed.
The relation is independent of the chosen minimal extensions because
both sides mean precisely the corresponding isomorphism of reflexive
sheaves downstairs. These assertions concern the represented
$k$-points; no universal algebraized family over a complete local base
is being asserted. This proves the final statement.
\end{proof}

\begin{theorem}[Proved: Fixed-rank growth of defining-characteristic representations]
\label{D:surface-fixed-rank-representation-growth}
Let $S$ be a proper scheme of dimension at most two over an algebraically
closed field of characteristic $p>0$, and let $T$ be finite of codimension
at least two in every component it meets. For a connected component
$U$ of $S\setminus T$ and each fixed $r$, there is $c_{U,r}$ such that
\[
 \#\bigl(\operatorname{Hom}_{\rm cont}
 (\pi_1(U),\operatorname{GL}_r(\mathbf F_{p^f}))/\mathrm{conj}\bigr)
 \le p^{c_{U,r}f}\qquad(f\ge1).
\]
\end{theorem}
\begin{proof}
First let $S$ be normal and work on its regular locus. Use the bounded
atlas in Proposition~\ref{D:global-periodic-reflexive-boundedness}.
Beke's effective quotient theorem for constructible equivalence
relations gives a constructible space $P_r$ whose $k$-points are the
represented reflexive isomorphism classes
\citep[Theorem 2]{Beke2011FixedRankB}. The Frobenius relation descends
to the graph of a partial constructible function $\Phi$ on $P_r$.
It may be undefined when a nonperiodic object leaves the bounded
family. Every periodic orbit is defined throughout and remains in
the family, by the same boundedness proposition.

A fixed periodic object admits at most one unit structure up to
isomorphism. To see this, set $A=H^0(S,\operatorname{End}E)$, a
finite-dimensional $k$-algebra. Its unit group is connected: its
semisimple quotient is a product of general linear groups, and the
kernel from the radical is unipotent. A chosen $q$-unit structure
defines a $q$-semilinear algebra endomorphism $\tau$ of $A$.
The differential of $\tau$ is zero. Every twisted-conjugation orbit
in $A^\times$ is therefore open, since the orbit map has invertible
differential. Connectedness makes it a single orbit. This is the
required Lang surjectivity and uniqueness.

Unit structures correspond to finite etale
$\operatorname{GL}_r(\mathbf F_q)$-frame torsors. The proved smallness
of punctured proper normal surfaces implies that there are finitely
many such classes for each $q$. The exact source is the combination
of ILO Expose XXI, Theorem 1.3, with its proper-composition
Lemma 3.2.1(ii), already used in Proposition~\ref{D:codim-two-small}.
Thus the $f$-cycles of $\Phi$ form a finite set of $k$-points.

Decompose the graph of $\Phi$ into finitely many affine locally closed
pieces, each defined by bounded-degree equations and inequations.
A length-$f$ word of pieces gives a cycle locus in $O_r(f)$ variables
with $O_r(f)$ equations of fixed degree. Replace inequations by inverse
variables. Its set of geometric points is finite, so its reduced
support is zero-dimensional. Successive hypersurface intersection
bounds its number of points exponentially in $f$. There are only
exponentially many words. Hence the total is at most $p^{c_rf}$.
This argument applies to a partial function; alternatively one can add
an absorbing point for the undefined domain.

Purity permits removal of finitely many regular points. Restriction
from a normal scheme to a dense open preserves connectedness of every
connected finite etale cover, so the regular-locus group surjects
onto $\pi_1(U)$. Finally finite normalization handles a nonnormal
surface. For each fixed collection of upstairs representation classes,
descent isomorphisms on the finitely many connected components of the
double fibre product have at most
$|\operatorname{GL}_r(\mathbf F_q)|^c\le q^{cr^2}$ choices; impose the
cocycle equation and apply effective finite-surjective descent
\citep[Expose IX, Theorem 4.7]{C_SGA1}. Proper lower-dimensional
components have finitely generated fundamental groups. The finite
product of these estimates proves the stated scope.
\end{proof}

% Reconstructed from the four exact author transmissions retained by root.
% Root and B: full reconstruction/source audit PASS.
% D: independent targeted Artinian pole and Isom-openness audit PASS.
\begin{lemma}[Proved: Polynomial witnesses for bounded derived diagrams]
\label{D:polynomial-monad-diagrams}
We fix a projective space $P=\mathbf P^N_k$, finitely many vector bundles
on $P$, a bound on complex length, and a finite list of operations
consisting of tensoring with fixed bounded complexes, absolute Frobenius
pullback, composition, and equality of derived morphisms. Suppose the
variable complexes have at most $cr$ copies of each fixed bundle in each
degree. A fixed finite diagram formed from these operations admits a
witness scheme defined by at most $Cr^2$ scalar polynomial equations in
at most $Cr^2$ variables, of degrees at most $D$. The constants depend
on the fixed data and on $p$, but not on $r$. The equations interpret
the diagram correctly over every $k$-algebra; they parameterize
representatives and witnesses for equality, rather than a quotient of
the space of representatives.
\end{lemma}
\begin{proof}
We choose the standard finite affine cover of $P$. For each pair of
fixed bundles, its Cech Hom complex is a complex of $k$-vector spaces
with finite-dimensional cohomology. We split it $k$-linearly into its
cohomology and contractible summands. For two variable bounded complexes,
the total Hom differential has the form $d_0+\delta$, where $d_0$ is
the fixed Cech differential and $\delta$ raises the complex-Hom column.
The column has a fixed finite range. Consequently the elimination of
the contractible summands uses a finite geometric series in $h\delta$,
where $h$ is the chosen contraction. Every resulting differential,
inclusion, projection and homotopy is a polynomial of bounded degree
in the differential coordinates. The finite remaining complex has
dimension at most $Cr^2$: its terms are sums of fixed cohomology spaces
tensored with matrix spaces of dimensions at most $(cr)^2$.

We impose the cycle equation on a degree-zero representative. An equality
of derived morphisms is imposed by adjoining a degree-minus-one variable
whose differential is the difference. Composition is computed by the
Cech cup product and then the finite projection. Tensoring with a fixed
bounded complex enlarges the fixed bundle list and the column range by
fixed amounts. Frobenius replaces coefficients by their $p$th powers
and replaces each fixed bundle by its fixed Frobenius pullback. The
finite list of operations therefore preserves the asserted bounds.
All splittings are over $k$, so their identities remain identities
after tensoring with any $k$-algebra. In particular this calculation
applies to nonreduced Artinian parameter rings. We do not replace the
derived category by its Ext algebra or assume that this algebra has a
strict finite-dimensional differential graded model.
\end{proof}

\begin{theorem}[Proved: A quadratic rank bound for defining-characteristic
representations on a punctured proper surface]
\label{D:global-quadratic-representation-bound}
Let $S$ be a normal proper connected surface over an algebraically closed
field $k$ of characteristic $p>0$, and let $T\subset S$ be a finite closed
subset. We put $U=S\setminus T$ and assume $U\ne\varnothing$. There is
a constant $C_S$ such that, for every $r,f\ge1$, with $q=p^f$,
\[
 \#\bigl(\operatorname{Hom}_{\mathrm{cont}}
   (\pi_1(U),\operatorname{GL}_r(\mathbf F_q))/
   \operatorname{GL}_r(\mathbf F_q)\bigr)
 \le q^{C_Sr^2}.
\]
The same conclusion, with a constant depending on the scheme and boundary,
holds for a proper scheme of dimension at most two whose finite removed
subset has codimension at least two in every irreducible component.
This statement concerns the defining coefficient characteristic. It does
not assert finite generation of the full fundamental group.
\end{theorem}
\begin{proof}
We first work on $S_{\mathrm{reg}}$. We fix a smooth projective resolution
$\pi:Y\to S$, an effective exceptional cycle $Z$ with $Z\cdot E<0$ for
every exceptional component $E$, and an embedding
$i:Y\hookrightarrow P=\mathbf P^N_k$. If the exceptional locus is empty,
we set $Z=0$. We write $V=Y\setminus\operatorname{Supp}Z$.

\emph{Uniform linear sizes.}
A rank-$r$ $\mathbf F_{p^f}$ local system gives reflexive
coefficient-Frobenius stages $E_a$, indexed by
$a\in\mathbf Z/f\mathbf Z$, and isomorphisms
$F_S^{[*]}E_a\simeq E_{a+1}$. We choose their minimal extensions
$F_a$ to $Y$. The proof of
Theorem~\ref{D:global-periodic-reflexive-boundedness} gives the following
bounds uniformly in $r$, $f$ and $a$: fixed upper and lower slope windows
on each exceptional component; $c_1(F_a)$ in the exceptional span with
coefficients $O(r)$; $\chi(F_a)=O(r)$; and fixed upper bounds for the
slopes of $F_a$ and its dual on $Y$ and on one fixed smooth ample curve.
In particular, for every fixed integer $t$,
\[
 h^j(Y,F_a(t))\le C_tr\qquad(0\le j\le2).
\]
Here the bounds for $h^0$ and $h^2$ follow from the fixed-curve restriction
and duality arguments already proved, and Riemann--Roch supplies the bound
for $h^1$. The same conclusion holds after tensoring by any one of
finitely many fixed vector bundles: we take a fixed three-step partial
resolution by sums of fixed line twists, tensor by $F_a$, and use the
vanishing of coherent cohomology in degrees above two. It also holds for
$F_Y^*F_a$ and its finitely many fixed twists. Indeed the one-step pole
comparison bounds its slopes and dual slopes; Frobenius scales the first
and second Chern characters by $p$ and $p^2$, respectively. We use only
this fixed one-step Frobenius, not a complex whose degree grows with $p^f$.

By the uniform pole lemma we may choose one integer $B$ such that the
isomorphisms on $V$ between $F_Y^*F_a$ and $F_{a+1}$, and their inverses,
extend with poles at most $BZ$. We enlarge $B$ once to bound maps in both
directions between any two bundles in the fixed exceptional slope windows.
The same bound will later apply to isomorphisms over Artinian parameter
rings. Put $L=\mathcal O_Y(BZ)$ and let $s$ be its canonical section.

The Beilinson monad of a coherent sheaf $\mathcal F$ on $P$ has terms
\[
 \mathcal B^e(\mathcal F)=
 \bigoplus_j H^j(P,\mathcal F(e-j))\otimes\Omega_P^{j-e}(j-e),
\]
where $0\le j-e\le N$; it is exact except in degree zero, where its
cohomology is $\mathcal F$
\citep[Theorem~6.1, pp.~23--25]{C_EFS2003}. Applied to the pushforward
of any of the bundles just listed, only $j=0,1,2$ occurs. Thus there is a
fixed finite set of bundle types and degrees, and every multiplicity is
at most $Kr$, for one fixed $K$. Per Frobenius stage the number of possible
multiplicity arrays is at most $(Kr+1)^{K_0}$, with $K_0$ fixed.

\emph{The parameter schemes.}
We describe them explicitly in order to keep their degree and size bounds
independent of the rank. Choose fixed finite vector-bundle resolutions
$K_L$ and $K_0$ of $i_*L$ and $i_*\mathcal O_Y$ on $P$. These exist by
the finite graded free resolution over the polynomial coordinate ring.
Fix also a derived representative of the section
$s:\mathcal O_P\to i_*L$.

For each $a$ we use six monads, representing
\[
 F_a,\quad T(F_a),\quad T^2(F_a),\quad
 P_a,\quad T(P_a),\quad T^2(P_a),\qquad
 T(M)=M\otimes L,\quad P_a=F_Y^*F_a.
\]
The $T$-nodes are introduced iteratively. Given a monad $C_M$, we introduce
a monad $C_{T(M)}$ and a derived map
\[
 \beta_M:C_M\otimes K_L\longrightarrow C_{T(M)}.
\]
We require its cone to have fibre cohomology only in degrees at most $-2$.
Once the monads have only degree-zero cohomology, this condition identifies
$H^0(C_{T(M)})$ with $H^0(C_M)\otimes L$. For the $P_a$-node we use in
the same way
\[
 LF_P^*C_{F_a}\otimes K_0\longrightarrow C_{P_a}.
\]
The naive pullback of the vector-bundle complex computes this derived
pullback. Absolute Frobenius need not be flat over a nonreduced parameter
ring; all its negative Tor terms, as well as those from tensoring with
$K_0$ or $K_L$, are retained. Their degree-zero cohomology is the desired
ordinary pullback or tensor product. The cone condition is precisely what
identifies that cohomology with the auxiliary monad.

The fixed section $s$ gives maps $\sigma_M:M\to T(M)$. We introduce four
arrows per stage,
\[
 a_a:P_a\to T(F_{a+1}),\qquad
 b_a:F_{a+1}\to T(P_a),\qquad T(a_a),\qquad T(b_a).
\]
The last two arrows are forced to be the tensor transforms of the first
two by the equations
\[
 T(a_a)\beta_{P_a}
   =\beta_{T(F_{a+1})}(a_a\otimes K_L),\qquad
 T(b_a)\beta_{F_{a+1}}
   =\beta_{T(P_a)}(b_a\otimes K_L).
\]
We impose
\[
 T(b_a)a_a=\sigma_{T(P_a)}\sigma_{P_a},\qquad
 T(a_a)b_a=\sigma_{T(F_{a+1})}\sigma_{F_{a+1}}.
\]
On $V$, division by $s$ makes $a_a$ and $b_a$ inverse isomorphisms.
This iterative six-node diagram makes the tensor identifications
coherent; we do not choose independent, unconnected identifications
for the second twists.

All equations just described are imposed using
Lemma~\ref{D:polynomial-monad-diagrams}. The fixed bundle list consists of
the $\Omega_P^j(j)$, their finitely many twists by the degrees of generators
of $I_Y$, the terms $\Omega_P^j(j)\otimes K_L^b$, and the terms
$F_P^*\Omega_P^j(j)\otimes K_0^b$. The complex lengths are fixed. At each
stage there are six monads, five auxiliary degree-zero identification
maps, four arrow variables, and a fixed number of derived equalities.
Each differential block, cycle space and boundary-witness space has
dimension $O(r^2)$. We consequently obtain at most $Afr^2$ variables
and at most $Afr^2$ scalar equations, all of degrees at most $D$, for
constants $A,D$ depending only on the fixed geometry and $p$.

We impose the scheme-theoretic support condition by choosing fixed
homogeneous generators $g$ of $I_Y$ and requiring multiplication
$g:C(-\deg g)\to C$ to be zero as a derived morphism. A boundary witness
expresses this as polynomial equations. We then take the open locus on
which every monad has cohomology only in degree zero and that cohomology
is a base-flat sheaf $i_*F$, with $F$ locally free of rank $r$ on $Y$.
The relative truncation lemma of Lieblich gives this openness and
compatibility with base change
\citep[Lemma~2.1.4, pp.~4--5]{C_Lieblich2006}. Local freeness follows
from flatness over the parameter base and local freeness on its fibres:
we lift a fibre basis, and apply Nakayama to the cokernel and then, using
flatness, to the kernel. Properness of $Y$ makes this an open condition
on the parameter base. We also impose the open fixed slope-window
conditions on the original $F_a$, and the open cone conditions above.
The slope conditions may be expressed using the proper images of the
finitely many relevant Quot schemes of coherent quotients; torsion can
only increase their degrees, so including all coherent quotients gives
the correct closed bad loci.

We denote the resulting open witness scheme for the chosen arrays by $W$.
It has a universal family of bundles on $Y$ and a universal unit Frobenius
cycle on $V$. Every local system under consideration has a point in one
of these schemes: its six bundles have the stated linear monad bounds,
and the actual tensor maps, inverse arrows and derived equalities provide
all the witnesses. We need no quantitative bound for the complexity of
the open conditions. An open subset of an irreducible scheme remains
irreducible when nonempty.

\emph{Uniform poles over Artinian bases.}
We spell out this point because the bundles chosen on $Y$ can vary even
when their restrictions to $V$ are isomorphic. If $F,G$ lie in the fixed
exceptional slope windows, a fixed filtration of the nonreduced cycle
$Z$ by line-bundle twists on its reduced components gives
\[
 H^0\bigl(Z,\mathcal Hom(F,G)\otimes\mathcal O_Z(nZ)\bigr)=0
 \qquad(n>B),
\]
with $B$ independent of the ranks. Indeed $Z\cdot E<0$, and on each fixed
layer the maximum target slope eventually lies below the minimum source
slope. This uses the Hom vanishing criterion, not semistability of a
tensor product.

Now let $A$ be any local Artinian $k$-algebra and let $F_A,G_A$ be two
families whose closed fibres lie in those windows. The Hom bundle
restricted to each fixed $Z$-layer is flat over $A$. Its maximal-ideal
filtration has graded pieces equal to the closed-fibre sheaf tensored
with vector spaces. The same $H^0$ vanishing therefore holds over $A$.
From
\[
 0\to\mathcal Hom(F_A,G_A)((n-1)Z)
 \to\mathcal Hom(F_A,G_A)(nZ)
 \to\mathcal Hom(F_A,G_A)\otimes\mathcal O_Z(nZ)\to0
\]
we see that the global Hom spaces stabilize at $B$. Every map on $V_A$
first extends with some finite pole order, because localization along
the effective Cartier divisor $Z$ is the filtered colimit of these
twists and global sections commute with this colimit on the
quasi-compact separated scheme $Y_A$. It therefore extends with poles
at most $BZ$. The same holds for its inverse. This bound is independent
of the length of $A$, of $r$, and of $f$. Multiplication by a local
equation of $Z$ is injective on these bundles, even over nonreduced $A$,
so the extension is unique.

\emph{Open unit-isomorphism classes.}
Fix a closed point $w\in W$ and its reference unit cycle. Consider the
functor of isomorphisms from this reference cycle to the variable cycle
on $V$, represented by maps in both directions on $Y$ with poles at most
$BZ$. The proper Hom functor is affine of finite presentation
\citep[Proposition~99.3.10, Tag~08K6]{C_StacksQuantitativeMonads}.
The inverse equations and the one-step Frobenius compatibility equations,
after clearing fixed poles, define an affine finitely presented unit-Isom
scheme over $W$. The compatibility equations involve $pB$ and $B$,
rather than $p^fB$, because we retain all $f$ consecutive $p$-stages.

For a small extension of local Artinian parameter algebras
$A\twoheadrightarrow A/I$, a specified unit-isomorphism over $A/I$
lifts uniquely on $V_A$. Indeed a unit $p^f$-Frobenius vector bundle has
a finite etale $\operatorname{GL}_r(\mathbf F_{p^f})$ frame torsor,
and finite etale covers and their morphisms are invariant under
nilpotent thickenings
\citep[Lemma~58.8.3, Tag~0BQB]{C_StacksQuantitativeMonads}. The frame
torsor can also be checked directly from the equations for
Frobenius-fixed frames: their derivative is invertible, and over every
geometric point they give the usual $\mathbf F_{p^f}$ basis torsor.
The preceding uniform pole argument extends the lifted isomorphism and
its inverse to $Y_A$ with the same bound $B$. Their equations hold
globally because they hold on $V_A$ and localization along $Z$ is
injective. Thus the unit-Isom scheme lifts the specified Artinian point.
The infinitesimal smoothness criterion applies
\citep[Lemma~10.141.2, Tag~02HT]{C_StacksQuantitativeMonads}; its image
contains an open neighbourhood of $w$ consisting of its unit-isomorphism
class.

We repeat this at every closed point of $W$. Distinct unit-isomorphism
classes give disjoint open subsets: an intersection, if nonempty, would
contain a closed point because $W$ is of finite type over $k$. These opens
cover $W$, since their complement is closed and has no closed point.
Quasi-compactness then gives a finite clopen partition. Therefore the
number of classes represented by $W$ is at most its number of connected
components, and hence at most its number of irreducible components.
This argument does not infer algebraicity from formal constancy of
covers on an arbitrary nonproper scheme. The fixed exceptional pole
bound is what makes the unit-Isom functor finitely presented.

An affine scheme cut out by $m$ polynomials of degree at most $D\ge2$
has at most $D^m$ irreducible components. For completeness, we homogenize
the equations and intersect them successively in projective space.
An irreducible component contained in the next hypersurface retains its
degree; for any other component the sum of degrees of its
hypersurface-section components is at most $D$ times its degree, by
the multiplication exact sequence for the hypersurface equation and
its Hilbert polynomial. Induction bounds the sum of degrees, and thus
the number of irreducible components, by $D^m$. Passing to the affine
chart and then to an open subset cannot increase that number. Here
$m\le Afr^2$. The multiplicity-array choices contribute at most
$(Kr+1)^{K_0f}$. Absorbing this factor gives
\[
 \#\{\text{represented unit classes}\}\le\exp(Cfr^2)
 =p^{C'fr^2}.
\]
Unit cycles up to unit-isomorphism are exactly the $\mathbf F_{p^f}$
local systems. This proves the bound on $S_{\mathrm{reg}}$. Removing
the regular points of $T$ does not change its fundamental group by
purity. Restriction from $U$ to $S_{\mathrm{reg}}\setminus T$ is fully
faithful on connected finite covers, since $U$ is normal and the latter
open is dense. Hence its fundamental group surjects onto $\pi_1(U)$,
and the same bound applies to $U$.

Finally, for a possibly nonnormal proper scheme of dimension at most two,
reduction does not change finite etale covers. Its normalization is
finite. On each normalized surface component we use the bound just
proved; proper curve and point components have the analogous bound
from their finitely generated fundamental groups. For fixed pulled-back
local systems, effective finite-surjective descent identifies descent
data with isomorphisms on the double fibre product satisfying the
cocycle equation
\citep[Expose~IX, Theorem~4.7]{C_SGA1}. A section of the resulting Isom
local system is determined by its value at one geometric point of each
connected component. Thus the number of possible descent data is at
most $|\operatorname{GL}_r(\mathbf F_q)|^c\le q^{cr^2}$, for the fixed
finite number $c$ of double-product components. Multiplying these bounds
proves the stated extension. We treat each connected component separately
when necessary.
\end{proof}

\subsection{Fixed-rank representations with tame original boundary}
\begin{proposition}[Proved: Boundedness of the canonical parabolic lattices]
\label{D:tame-parabolic-boundedness}
Let $S$ be a normal proper connected surface over an algebraically closed field $k$ of characteristic $p>0$, and let $D$ be a reduced divisor. Fix a finite set $\Sigma$ containing the singular points of $S$, the singular and intersection points of $D$, and any additional omitted points. Put $S^\circ=S\setminus\Sigma$, $U^\circ=S^\circ\setminus D$, and let $Q^\circ$ be the quotient of $\pi_1(U^\circ)$ classifying covers tame at the generic points of the original components of $D$. Fix a projective resolution $\pi:Y\to S$, an isomorphism over $S^\circ$, and a positive exceptional divisor $Z$ with $Z\cdot E_i<0$ for every exceptional component.
For every fixed rank $r$, one bounded family of vector bundles on $Y$ contains minimal extensions of all canonical reflexive lattices of representations $Q^\circ\to\mathrm{GL}_r(\mathbf F_{p^f})$, all their coefficient-Frobenius conjugates, and all distinct parabolic weight lattices between $E(-D_i)$ and $E$, for every $f\ge1$ and every component $D_i$. The family is independent of the Kummer denominators.
\end{proposition}
\begin{proof}
Write $q=p^f$. For a representation $M$, let $M_a$ be its coefficient-Frobenius conjugates, indexed by $a\in\mathbf Z/f\mathbf Z$, and let $E_a$ be their canonical reflexive lattices on $S$. Thus $E_a|_{U^\circ}=\mathcal O_{U^\circ}\otimes_{\mathbf F_q}M_a$. At a generic boundary DVR we choose a common prime-to-$p$ Kummer denominator $n$. The integral invariant lattice has weight representatives $0\le j<n$, and its Frobenius map has exponent $c_j=\lfloor pj/n\rfloor\le p-1$. We obtain injections
\[\phi_a:F_S^{[*]}E_a\longrightarrow E_{a+1}\]
with pure one-dimensional cokernels $T_a$ supported on $D$.
The canonical-lattice estimates proved above remain uniform for this finite cycle. Indeed, their local ordinary/reflexive Frobenius defect ideal and relative-Frobenius determinant depend only on the fixed surface charts. They apply equally to a map from the $a$th Kummer module to the $(a+1)$st. The normalized eigenpiece maps now run around finite cycles of pairs $(a,j)$. Their one-step upper-slope inequalities have the same fixed error, so taking the maximum around such a cycle gives the same upper-slope bound. The fixed matrix comparison and the additive normalization-index bound for the final cokernel grades are unchanged. Consequently
\[\chi(S,T_a)\le C r\]
for one constant independent of $r,f,n,a$. The same argument gives an upper Euler bound $C\operatorname{rk}G$ for each pure whole-eigenpiece grade $G$, with the sum bounded by $Cr$ along each boundary component.
Taking determinants gives
\[c_1(E_{a+1})-p c_1(E_a)=[T_a]_1,\]
where every coefficient of $[T_a]_1$ lies between $0$ and $(p-1)r$. Iterating the $f$-cycle shows
\[c_1(E_a)=-\sum_i\theta_{a,i}[D_i],\qquad 0\le\theta_{a,i}\le r\]
in the numerical divisor space. In particular $K_S\cdot c_1(E_a)=O(r)$. The uniform normal-surface Riemann--Roch theorem \citep[Theorem~0.1]{C_Langer2026} and Frobenius compatibility of the second Chern character \citep[Theorem~4.5(2)]{C_Langer2025} give
\[\operatorname{ch}_2(E_{a+1})-p^2\operatorname{ch}_2(E_a)\le Cr.\]
After cycling, division by $p^{2f}-1$ gives $\operatorname{ch}_2(E_a)\ge-Cr$, uniformly in $f$.
For clarity, the fixed-resolution sandwich has the adjacent-index form
\[F_{a+1}(-A)\subset F_Y^*F_a\subset F_{a+1},\qquad F_a=\pi^{[*]}E_a,\]
for one fixed effective divisor $A$. It follows from the same fixed boundary ideal and stable-MCM Frobenius-defect ideal used in the canonical-lattice proof. Taking the maximum around the finite cycle, both on $Y$ and on a fixed smooth very ample curve not contained in the determinant-zero divisors, gives the same bounds $h^0(S,E_a)\le r$ and $h^2(S,E_a)\le Cr$. Riemann--Roch therefore yields $|\chi(S,E_a)|\le Cr$.
For a weight lattice $L$ with $E_a(-D_i)\subset L\subset E_a$, the successive pure grades of $E_a/L$ have total upper Euler bound $Cr$. Hence $\chi(S,L)\ge-Cr$. We also have $h^0(S,L)\le r$ and
\[\operatorname{Hom}(L,\omega_S)\hookrightarrow\operatorname{Hom}(E_a(-D_i),\omega_S).\]
Principalize the fixed coherent ideal $\mathcal O_S(-D_i)$ on the resolution. The latter Hom space injects into sections of $F_a^\vee$ tensored by one fixed line bundle, so the preceding dual-slope argument bounds it by $Cr$. Thus $|\chi(S,L)|\le Cr$ as well. Its first Chern class differs from that of $E_a$ by a multiple of $D_i$ between $0$ and $r$.

Choose minimal extensions of these reflexive lattices on $Y$. Their restrictions to the exceptional components lie in the fixed slope windows established earlier \citep[Definition~1.2 and Lemma~1.3]{C_Langer2026}. The uniform relative Euler-characteristic estimate $\chi(\pi,F)\le Cr$ therefore bounds $|\chi(Y,F)|$ by $Cr$. The numerical first Chern classes on $S$ lie in a fixed $O(r)$ divisor span; the exceptional slope windows and the negative-definite intersection matrix then bound $c_1(F)\cdot H$ by $Cr$, for a fixed very ample $H$ on $Y$.
The bounded-pole lemma extends each one-step Frobenius map between minimal extensions to
\[F_Y^*\widetilde E_a\hookrightarrow\widetilde E_{a+1}(BZ)\]
with one fixed $B$. Here the restriction-slope bounds for a Frobenius pullback follow from the already used curve Frobenius estimate; all constants are independent of the rank and of $f$. Cycling gives a uniform upper bound on $\mu_{\max,H}(\widetilde E_a)$. The inclusion $L\hookrightarrow E_a$ likewise extends between their minimal extensions with one fixed exceptional pole, and bounds the maximal slope of the extension of $L$. At fixed rank, the resulting Hilbert polynomials range over a finite set. Langer's boundedness theorem \citep[Theorem~4.2, p.~269]{C_LangerBounded2004}, retained by its addendum \citep{C_LangerAddendum2004}, now gives the asserted bounded family.
\end{proof}

\begin{theorem}[Proved: Fixed-rank growth for the original-boundary tame quotient]
\label{D:tame-parabolic-fixed-rank-growth}
In the preceding situation, for every fixed positive integer $r$ there is a constant $C_r$ such that
\[\#\bigl(\operatorname{Hom}_{\mathrm{cont}}(Q^\circ,\mathrm{GL}_r(\mathbf F_{p^f}))/\mathrm{GL}_r(\mathbf F_{p^f})\bigr)\le p^{C_r f}\qquad(f\ge1).\]
The same bound holds for the original quotient before enlarging the finite omitted set. No uniform quadratic dependence on $r$ is asserted.
\end{theorem}
\begin{proof}
We first construct a finite-type space of unweighted parabolic states. A state consists of a reflexive lattice $E$ and, for each original boundary component $D_i$, an ordered chain
\[E=E_{i,0}\supset E_{i,1}\supset\cdots\supset E_{i,l_i}=E(-D_i),\qquad l_i\le r,\]
whose successive ranks on $D_i$ are positive. On $S^\circ$, this is a flag of subbundles of $E|_{D_i^\circ}$; we retain the full chain, including its extension data, and do not yet assign numerical weights.
The preceding proposition supplies bounded minimal extensions for every lattice occurring in a representation. We take finite unions of their parameter spaces and parameterize the chain maps by bounded-pole Hom schemes on the projective surface $Y$. Conditions on support, endpoint, inclusions, and local freeness of the successive quotients on $S^\circ$ are constructible: they are finite-presentation conditions, after the usual finite flattening stratifications, and their failure loci project to constructible sets. An exact isomorphism of states is an isomorphism of $E$ preserving every member of every chain. The fixed-pole lemma and the inverse-product equations used earlier make this an effective constructible equivalence relation on a finite-type parameter set. Beke's effective quotient theorem \citep[Theorem~2; proof via Proposition~3.1]{Beke2011FixedRankB} therefore gives a constructible quotient $P_r$ whose points are the state-isomorphism classes. We may enlarge the parameter family; only the class set of cycles will matter below.
There are finitely many grade-rank types. A one-step Frobenius type consists, for each $D_i$, of a permutation $\sigma$ of the grades preserving their ranks and digits $c_j\in\{0,\ldots,p-1\}$ nondecreasing in the source order. We also require $\sigma$ to be increasing on every block on which the digit $c_j$ is constant. In local adapted bases \citep[Proposition~3, pp.~11--12]{Borne2007D} $v_j$ and $w_k$ and a local equation $t$ of $D_i$, we impose on the matrix $(a_{kj})$ of a map $F^*E\to E'$ the conditions
\[v_t(a_{kj})\ge c_j+\mathbf1_{k<\sigma(j)},\]
and require the diagonal blocks $k=\sigma(j)$, after division of column $j$ by $t^{c_j}$, to be invertible modulo $t$. Away $D$ the map is required to be an isomorphism.

These local conditions are intrinsic. The normalized Frobenius source lattice is
\[L_c=\sum_j\mathcal O(c_jD_i)\,F^*E_{i,j-1}.\]
Locally it has basis $t^{-c_j}F^*v_j$. For an adapted source change of basis, an entry from a higher-weight column to a lower-weight row is divisible by $t$; after Frobenius it is divisible by $t^p$. Since digit differences are at most $p-1$, the normalized transition matrix is integral. Modulo $t$, entries between unequal digit blocks vanish, while within each digit block the original flag order is preserved. The digit blocks and their subflags therefore define a canonical reordered flag on $L_c|_{D_i}$. The displayed condition is precisely an isomorphism of this flagged bundle with the target flagged bundle.
All these constructions use only finitely many sums, images, restrictions to fixed thickenings, and twists of order at most $p-1$. Thus they are constructible in the bounded families. No root stack is extended across the bad points: maps of reflexive lattices extend uniquely across $\Sigma$, and on $Y$ their poles are controlled by the fixed exceptional-window bound. The finitely many digit twists add only a fixed amount to that bound. We obtain finitely many constructible correspondences
\[R_c\subset P_r\times P_r.\]
Consider a closed word of $f$ such types. Following a grade around the induced permutation cycles, solve
\[\alpha_{\mathrm{next}}=p\alpha-c.\]
On a cycle of length $h$ there is a unique solution, with denominator dividing $p^h-1$. Here $h\le fr$. We discard words for which the solutions fail to lie in $[0,1)$, fail to have the prescribed strict order, or identify two grades prescribed to be distinct. In particular, the all-$(p-1)$ cycle, which gives weight $1$, is excluded. There are at most $A_r^f$ words before this discard.
For a valid word choose a common denominator $n$ prime to $p$. The flagged states become parabolic vector bundles on $(S^\circ,D)$ with these weights, hence vector bundles on the root stack $\mathcal X=\sqrt[n]{(S^\circ,D)}$ by Borne's equivalence \citep[Theorem~4, p.~12, with proof through p.~16]{Borne2007D}. This equivalence is being used only where the boundary components are smooth and disjoint. The edge criterion is exact: if the source and target weights are $\alpha_j$ and $\beta_k$, then
\[\left\lceil p\alpha_j-\beta_k\right\rceil=c_j+\mathbf1_{k<\sigma(j)}.\]
This is exactly the condition for the pulled-back matrix to be regular on the Kummer chart. Its normalized diagonal blocks are invertible precisely when its determinant is a unit there. Hence the map and its inverse are both regular on that chart. Uniqueness on the dense complement glues these maps to genuine Frobenius isomorphisms on $\mathcal X$; the parabolic gluing and tensor compatibility are also explicit in \citet[Theorem~6.1, pp.~33--36]{BorneVistoli2012D}.
Composing the $f$ edges gives a $p^f$-unit vector bundle on $\mathcal X$. Its fixed-frame equations, locally $z^{[p^f]}=C^{-1}z$, define a finite étale $\mathrm{GL}_r(\mathbf F_{p^f})$-torsor. Its restriction to $U^\circ$ has inertia factoring through $\mu_n$ at every original boundary component, so it defines a representation of $Q^\circ$. Conversely, every such representation gives a valid word and a cycle of states.
For a fixed weighted state, its algebra of global parabolic endomorphisms is a finite-dimensional subalgebra of $H^0(S,\mathcal End(E))$. This follows from reflexive extension across the finite omitted set and from preservation of the finitely many lattices. A given $p^f$-unit structure induces a $p^f$-semilinear endomorphism $\tau$ of this algebra. Its unit group is connected. The $\tau$-twisted conjugation action has open orbits, because $d\tau=0$ and the differential in the group variable is invertible. It therefore has one orbit. Thus a weighted state admits at most one unit structure up to isomorphism; choices of intermediate edge isomorphisms do not produce extra representation classes.

Conversely, for a fixed representation and fixed valid word, its finite étale extension to the corresponding root stack is unique: restriction from a normal root stack to its dense open complement is fully faithful, as is seen on its normal Kummer charts. The associated unit bundle and its entire parabolic chain are consequently unique, and all intermediate states are its Frobenius conjugates. Thus the cycle set for a fixed valid word is in bijection with a subset of the representation classes just described. The fixed-original-boundary smallness theorem proved earlier, Theorem~\ref{D:fixed-boundary-small}, makes this set finite. This assertion is made only for valid words.
Choose finite basic affine locally closed descriptions of the fixed constructible correspondences $R_c$. For a word of length $f$, choosing one such piece at each edge gives a cycle locus in $O_r(f)$ variables, defined by $O_r(f)$ equations and inequations of degrees bounded independently of $f$. Introduce an inverse variable for each inequation. Its geometric point set is finite, so the resulting affine scheme is zero-dimensional after reduction. The elementary Bézout bound gives at most $\exp(C_r f)$ points. The number of piece words is also at most exponential in $f$. Summing proves the asserted bound for $Q^\circ$, and hence for its quotient corresponding to the original omitted set.
We emphasize that this argument uses a finite collection of correspondences. It does not require a single-valued Frobenius function on the unweighted state quotient. It also does not provide a rank-independent bound for the constants $C_r$, and therefore does not settle the remaining nonabelian finite-generation problem.
\end{proof}

% Fresh reconstruction after the 09:11 UTC environment replacement.
\subsection{Quadratic representation growth does not force generation}

\begin{example}[Proved: an alternating-product obstruction in every coefficient characteristic]
\label{D:quadratic-growth-alternating-obstruction}
For odd integers \(n\geq5\), put \(m_n=2^{n^2}\), and let
\[
 \Pi=\prod_{\substack{n\geq5\\n\ {\rm odd}}} A_n^{\,m_n}.
\]
This profinite group has the following properties.
\begin{enumerate}
\item It is small: for each finite group \(H\), the set of continuous
homomorphisms \(\Pi\to H\) is finite.
\item Every open subgroup has finite maximal prosolvable quotient.
\item For every finite field \(\mathbf F_q\), of any characteristic,
and every finite-dimensional continuous \(\mathbf F_q[\Pi]\)-module \(V\),
\[
 \dim_{\mathbf F_q}H^1(\Pi,V)\leq2\dim_{\mathbf F_q}V.
 \tag{A1}
\]
\item There is one constant \(C\), independent of the characteristic,
\(q\), and \(r\), such that
\[
 \#\bigl(
 \operatorname{Hom}_{\mathrm{cts}}(\Pi,\operatorname{GL}_r(\mathbf F_q))
             /\operatorname{GL}_r(\mathbf F_q)\bigr)
       \leq q^{Cr^2}.
 \tag{A2}
\]
\end{enumerate}
Nevertheless \(\Pi\) is not topologically finitely generated.
This is an abstract profinite obstruction. No realization as a
geometric fundamental group is asserted.
\end{example}

\begin{proof}
We give the counting argument in all coefficient characteristics.
We use the simplicity of \(A_n\) for \(n\geq5\), with its elementary
proof in \cite[Theorem 4.33 and Lemmas 4.35--4.36, pp.~69--70]
{MilneGT2025B}. For odd \(n\), \(A_n\) is generated by the even
\(n\)-cycle \((1\,2\,\cdots\,n)\) and \((1\,2\,3)\).
Indeed their conjugates include the consecutive three-cycles.
Those generate \(A_n\): inductively they contain \(A_{n-1}\), and
a three-cycle involving the last letter, together with its
\(A_{n-1}\)-conjugates, supplies every three-cycle involving that
letter. The initial cases are immediate. In particular every \(A_n\)
is perfect and two-generated.

\par\smallskip\noindent\textbf{A dimension bound for nontrivial factors.}\quad
Let \(W\) be a nontrivial representation of \(A_n\) over an
algebraically closed field, and put \(a=\dim W\).
Simplicity makes the representation faithful.
In characteristic different from two, use the subgroup
\((C_2)^{\lfloor n/4\rfloor}\) generated by disjoint double
transpositions. In characteristic two, use
\((C_3)^{\lfloor n/3\rfloor}\) generated by disjoint three-cycles.
These groups have order prime to the coefficient characteristic,
so their actions are simultaneously diagonalizable.
For an elementary abelian group of rank \(s\), a faithful diagonal
representation needs at least \(s\) characters: their common
kernel is zero only if they span its character space.
It follows in every characteristic that
\[
 n\leq8a,\qquad a\geq2.
 \tag{A3}
\]
The last inequality also follows from perfectness.

\par\smallskip\noindent\textbf{Absolutely simple modules.}\quad
A continuous finite-dimensional module factors through a finite
product of the coordinate groups. Over an algebraically closed
field, the simple modules for a product are the outer tensor
products of the simple modules of its factors, including in
modular characteristic. To see this without a semisimplicity
assumption on group algebras, take their finite-dimensional
algebras. The ideal generated by their radicals is nilpotent in
the tensor product, while the quotient is a product of matrix
algebras. Its simple modules are exactly the stated tensor
products. We use only these standard semisimple quotient facts
\cite[Theorems 7.25, 7.29, 7.36--7.37, pp.~110--114]{MilneGT2025B}.

Let a simple module have nontrivial dimensions
\(a_1,\ldots,a_t\) at its support coordinates, whose alternating
degrees are \(n_1,\ldots,n_t\). Its dimension is
\(r=\prod_i a_i\). For \(a_i\geq2\),
\[
 t\leq\log_2r,\qquad
 \sum_i a_i^2\leq\prod_i a_i^2=r^2,\qquad
 \sum_i n_i^2\leq64r^2.
 \tag{A4}
\]
The middle inequality follows by induction from
\(a^2+b^2\leq a^2b^2\) for \(a,b\geq2\).
There are at most \((8r)^{1+\lfloor\log_2r\rfloor}\) ordered choices
of the \(n_i\); this overcounts support types and is sufficient.
For each such support, its coordinate-copy choices number at most
\[
 \prod_i m_{n_i}=2^{\sum_i n_i^2}.
\]
For \(A_n\), the number of simple module labels is at most
\(\dim k[A_n]=|A_n|\leq n^n\leq2^{n^2}\).
Thus the label choices have the same bound.
Combining these estimates shows that the number of continuous
absolutely simple \(r\)-dimensional modules is at most
\[
 2^{C_0r^2}
 \tag{A5}
\]
for an absolute constant \(C_0\), uniformly in characteristic.
The rank-one trivial module is included by increasing the constant.

Over \(\mathbf F_q\), simple modules are classified by Frobenius
orbits of the absolutely simple modules, with unique descent.
Here is an algebraic justification valid also in modular
characteristic. A finite-dimensional semisimple algebra over
a finite field is a product of matrix algebras over finite
extensions of that field, since finite division rings are
commutative. A simple module for
\(\operatorname{Mat}_a(\mathbf F_{q^e})\) becomes, after algebraic
closure, the direct sum of the \(e\) distinct Frobenius-conjugate
absolutely simple modules of dimension \(a\). This description
also proves uniqueness. Apply it to the semisimple quotient of
the finite group algebra through which the module factors.
Consequently the number \(s_r(q)\) of simple modules of dimension
\(r\) over \(\mathbf F_q\) satisfies
\[
 s_r(q)\leq\sum_{e\mid r}2^{C_0(r/e)^2}
             \leq q^{C_1r^2},
 \tag{A6}
\]
with \(C_1\) absolute and all \(q\geq2\).

\par\smallskip\noindent\textbf{Uniform extension bounds.}\quad
For two absolutely simple modules \(U,V\), the degree-one
Kunneth formula for a finite product gives
\[
 \operatorname{Ext}^1(U,V)=
 \bigoplus_i
 H^1\!\left(A_{n_i},\operatorname{Hom}(U_i,V_i)\right)
       \otimes\!\!\bigotimes_{j\ne i}
                  \operatorname{Hom}_{A_{n_j}}(U_j,V_j).
 \tag{A7}
\]
One obtains this formula by tensoring the ordinary free bar
resolutions of the factors. Over a field, the degree-zero and
degree-one cohomology of the resulting total complex have exactly
the displayed decomposition. There is no restriction on the
coefficient characteristic.

For a two-generated finite group and a module \(W\), a derivation
is determined by its two values on the generators. Hence
\(\dim H^1(G,W)\leq2\dim W\).
By Schur's lemma, the invariant Hom spaces in (A7) are zero
unless the two simple factors agree, and are one-dimensional
when they do agree. If \(U,V\) differ at two support factors,
(A7) is zero. If they differ at exactly one, it is bounded by
\(2\dim U\dim V\). If \(U=V\), its dimension is at most
\(2\sum_i a_i^2\leq2(\dim U)^2\), by (A4).
Coordinates trivial on both sides contribute
\(H^1(A_n,k)=0\), because \(A_n\) is perfect.

These computations do calculate extensions for the full infinite
product. In an extension of two modules trivial on a coordinate
factor, that factor acts through the additive group
\(\operatorname{Hom}(U,V)\); perfectness kills this action.
Thus no new support coordinate can occur in the extension.
All groups in (A7) may be computed on a common finite product.

Cohomology commutes with extension of a coefficient field on this
finite product. Splitting the finite-field simple modules over
an algebraic closure and summing the preceding bounds therefore
gives
\[
 \dim_{\mathbf F_q}\operatorname{Ext}^1(U,V)
       \leq2\dim U\dim V
 \tag{A8}
\]
for simple finite-field modules. Devissage in both variables,
using the long exact Ext sequences, gives the same inequality
for arbitrary finite-dimensional modules. Taking \(U\) trivial
proves (A1).

To count all \(r\)-dimensional modules, fix an ordered composition
\(r=r_1+\cdots+r_t\). There are at most
\(q^{C_1\sum_i r_i^2}\) choices of the simple factors.
Successively extending them gives at most
\[
 q^{\,2\sum_{i<j}r_ir_j}
\]
choices, by (A8). We count extension classes with fixed
identifications of their ends, which can only overcount
isomorphism classes of middle terms. Since
\[
 C_1\sum_i r_i^2+2\sum_{i<j}r_ir_j
 \leq \max(C_1,1)r^2,
\]
and there are \(2^{r-1}\) ordered compositions of \(r\),
(A2) follows with one larger absolute constant.

\par\smallskip\noindent\textbf{Smallness and failure of generation.}\quad
A homomorphism from \(\Pi\) to a finite group \(H\) is trivial on
every \(A_n\) with \(|A_n|>|H|\), by simplicity. Only finitely many
coordinate factors remain, so there are finitely many such
homomorphisms. This proves smallness.

An open subgroup contains the full tail product after finitely
many coordinates. That tail has no nontrivial finite solvable
quotient, since each of its factors is perfect. Every solvable
quotient of the open subgroup therefore factors through a fixed
finite head quotient. Its maximal prosolvable quotient is finite.

Finally suppose \(s\) elements generated \(\Pi\).
Their projection to \(A_n^{m_n}\) would generate that finite group.
The \(s\)-tuples at any two coordinates must be different: if
they were equal, the generated subgroup would have equal entries
in those two coordinates and could not be the full product.
Thus \(m_n\leq|A_n|^s\).
But \(2^{n^2}>(n!)^s\) for all sufficiently large odd \(n\), since
\(n\log 2-s\log n\to+\infty\).
This contradiction proves that \(\Pi\) is not finitely generated.
\end{proof}

\begin{remark}[What the obstruction excludes]
Even the conjunction of smallness, finite maximal prosolvable quotients
for every open subgroup, all-characteristic rank-linear \(H^1\), and
a uniform all-rank quadratic representation-class bound does not imply
finite generation. A geometric proof must impose a further condition,
for example a sufficiently strong bound on nonabelian simple
multiplicities. The present example is not a counterexample to any
of the geometric statements under investigation.
\end{remark}


\subsubsection{A necessary condition for a finite cone cover}

\begin{proposition}[Proved: A persistent first neighborhood: proved obstruction]
\label{D:persistent-first-neighborhood}
Let $k$ be algebraically closed of characteristic $p>0$, and let
$C$ be a smooth connected projective curve of genus $g\geq2$ in a smooth
surface $Y$.  Put $N=N_{C/Y}$ and $d=-\deg N$, and assume
\[
                         0<d<2g-2.
\]
Let
\[
 e_Y\in H^1(C,T_C\otimes N^{-1})
\]
be the class of the first infinitesimal neighborhood of $C$.
Suppose that every Frobenius pullback of $e_Y$ is nonzero.
Then there is no finite surjective adic morphism
\[
 \phi:\widehat{\operatorname{Tot}(N')}_{C'}
             \longrightarrow\widehat Y_C
\]
where $C'$ is a smooth connected projective curve and $N'$ is a line
bundle on $C'$.  Transverse ramification is allowed in this assertion.

In particular, the hypothesis on Frobenius is satisfied whenever
\[
                   e_Y\ne0,\qquad p(2g-2-d)>2g-2.
\]
This is an obstruction to the finite-cone method.  We make no assertion
here about finite generation of the corresponding punctured fundamental
group.
\end{proposition}

\begin{proof}
We first record the curve-theoretic input.  If $L$ is a negative-degree
line bundle on a smooth projective curve and
$c\in H^1(C,L)$ survives all Frobenius pullbacks, its additive $L$-torsor
is affine.  This is the rank-one case of Brenner's affineness criterion,
valid over an arbitrary algebraically closed field of positive
characteristic: \cite[Proposition 2.1, preprint pp.~5--6]{C_Brenner2006}.
Consequently $h^*c\ne0$ for every finite surjective morphism
$h:C'\to C$ from a smooth projective connected curve.  Otherwise a
section of the pulled-back torsor would give a nonconstant morphism
from the proper curve $C'$ to the affine torsor, which is impossible.
Conversely, if a Frobenius power kills $c$, that power is itself a
finite surjective curve morphism.  Thus for a negative line, finite
annihilation and Frobenius annihilation are equivalent.  We use only
the implication just proved, without a finite-field hypothesis.

Suppose that $\phi$ exists, and let $h:C'\to C$ be its restriction to
the reduced closed fibers.  Finiteness and surjectivity make $h$ finite
surjective; write $m=\deg h$.  Because $\phi$ is adic, the pullback of
the ideal of $C$ has the same radical as the ideal of the zero section
$C'$.  Both formal surfaces are regular.  The divisor of a pulled-back
local equation of $C$ is therefore
\[
                         \phi^*C=eC'
\]
for one positive integer $e$.  In particular its leading coefficient
is a unit, and restriction of the associated invertible ideals gives
\[
                    (N')^{\otimes e}\simeq h^*N,
             \qquad \deg N'=-\frac{md}{e}.             \tag{PN1}
\]
This is the point at which the adic and finite hypotheses exclude an
additional horizontal component in the inverse image of $C$.

For $n\geq1$, write $C'_n$ for the zero section thickened by the
$n$th power of its ideal.  Since a local equation of $C$ maps into
the $e$th power of that ideal, the restriction of $\phi$ to $C'_e$
factors through $C$.  We claim that this factorization is the composite
of the canonical cone retraction $C'_e\to C'$ and $h$.
Both maps agree on $C'$.  At the step from $C'_n$ to $C'_{n+1}$,
the difference of two lifts of a fixed map to the smooth curve $C$
is a global section of
\[
                       h^*T_C\otimes(N')^{-n}.
\]
For $1\leq n<e$, its degree is
\[
               m\left(-(2g-2)+\frac{nd}{e}\right)<0
\]
by (PN1) and $d<2g-2$.  The section space is zero.  Induction proves
the claim, including the vacuous case $e=1$.

Now restrict $\phi$ to $C'_{e+1}$.  The target map factors through
the first neighborhood $C_2$ because $2e\geq e+1$.  Modulo the
$e$th power on the source, its map to $C$ is exactly the retracted
map established above.  The leading normal coefficient in (PN1)
identifies $h^*N^{-1}$ with $(N')^{-e}$.
Choose local splittings of the square-zero extension defining $C_2$.
Their differences represent $e_Y$.  Pulling these splittings back
along $h$, and comparing them with the local expressions of $\phi$
modulo the $(e+1)$st power, shows that the resulting differences are
the coboundary of the order-$e$ tangential corrections.  Hence
\[
 h^*e_Y=0
       \quad\hbox{in }H^1\bigl(C',h^*(T_C\otimes N^{-1})\bigr).
                                                               \tag{PN2}
\]
Equivalently, (PN2) is the obstruction to lifting the already retracted
map $C'_e\to C$ through $C_2$ with its prescribed leading normal map.
The possible lower tangential corrections have been excluded above;
we do not assume their vanishing without the negative-degree argument.
Since $T_C\otimes N^{-1}$ has degree $d-(2g-2)<0$, (PN2) contradicts
the curve-theoretic input.

For the final sufficient condition, set $L=T_C\otimes N^{-1}$ and
$a=-\deg L=2g-2-d>0$.  On a smooth curve over a perfect field the
differential has kernel $\mathcal O_C^p$: this is checked in the
function field using a separating parameter, and regularity follows
by normality.  Consequently there is an injection
\[
           F_*\mathcal O_C/\mathcal O_C\hookrightarrow F_*\Omega_C^1.
\]
Tensoring by $L$ and using the projection formula shows that the
kernel of
\[
                    H^1(C,L)\longrightarrow H^1(C,L^{\otimes p})
\]
is a quotient of a subspace of
$H^0(C,\Omega_C^1\otimes L^{\otimes p})$.
This last space vanishes if $pa>2g-2$.  The same inequality remains
true after replacing $L$ by any $L^{\otimes p^r}$.
All successive Frobenius pullbacks are therefore injective, so a
nonzero $e_Y$ survives every one of them.
\end{proof}

\begin{remark}
We have proved a conditional geometric obstruction and have not
chosen a numerical curve or a contracted algebraic germ.  No finite
example is introduced.  A failure of this particular cone-cover
construction does not imply that its wild inertia is infinite or that
its full fundamental group requires infinitely many generators.
\end{remark}

\begin{corollary}[Proved: residual genus and different bounds]
\label{D:surface-residual-genus-bound}
Use the resolution setup of
\ref{D:surface-uniform-different}, or its fixed-boundary variant
\ref{D:fixed-boundary-uniform-different}. Write the smooth exceptional
curves as $E_i$, choose a positive integral exceptional divisor
$A=\sum_i a_iE_i$ with $A\cdot E_i<0$, and put
\[
 C_A=(K_Y+D)\cdot A+\sum_i a_i,
\]
where $D=0$ in the punctured case and $D$ is the reduced strict
transform of the original divisorial boundary in the tame variant.

Let a normal finite covering in the preceding etale/tame setup
(hence generically separable) have generic degree $n$, and form
$f:Z\to Y$ as in those results, contracting only rational
$(-1)$-curves mapped to points of $Y$. For every component $F$
dominating $E_i$, let $e_F$ be its ramification index and let
$g_F$ be the genus of its normalization. Then
\[
 2\sum_{F/E_i}a_i e_Fg_F\leq n C_A.
 \tag{15}
\]
If $Q$ ranges over the remaining point-mapped exceptional components,
$m_Q$ is its coefficient in $f^*A$, and $g_Q$ is its normalized
genus, the stronger inequality
\[
 2\sum_{F/E_i}a_i e_Fg_F+\sum_Qm_Qg_Q\leq n C_A
 \tag{16}
\]
holds.

Factor the finite residual curve map from the normalization of $F$
to $E_i$ into a purely inseparable map followed by a separable map
$C_F\to E_i$, and let $\Delta_F$ be the different divisor of the
separable map. Then
\[
 \sum_{F/E_i}a_i e_F\deg\Delta_F
 \leq n\left(C_A+2\sum_{g(E_i)=0}a_i\right).
 \tag{17}
\]
All constants depend only on the fixed resolution and boundary.
\end{corollary}

\begin{proof}
Every component of the exceptional fibre has negative integral
self-intersection. Adjunction therefore gives, for a dominating
component,
\[
 K_Z\cdot F=2p_a(F)-2-F^2\geq2g_F-1.
\]
For a point-mapped component $Q$, the contractions stipulated in
the setup ensure $K_Z\cdot Q\geq0$. If $g_Q\geq1$, adjunction
also gives $K_Z\cdot Q\geq2g_Q-1\geq g_Q$; hence this last
inequality holds for every such $Q$.

The coefficient of $f^*A$ along $F/E_i$ is $a_i e_F$.
The Jacobian and projection calculation in the cited results gives
\[
 K_Z\cdot f^*A\leq n(K_Y+D)\cdot A.
\]
Here the vertical Jacobian contribution has nonpositive
intersection with $f^*A$; the tame horizontal contribution is
bounded by $nD\cdot A$. Thus
\[
 2\sum a_i e_Fg_F+\sum_Qm_Qg_Q
 \leq n(K_Y+D)\cdot A+\sum a_i e_F.
\]
At the generic point of $E_i$,
$\sum_{F/E_i}e_F[k(F):k(E_i)]=n$, so
$\sum_{F/E_i}e_F\leq n$. This proves (16), and hence (15).

The ground field is algebraically closed. Purely inseparable maps
between smooth proper curves are iterated relative Frobenius maps,
so $g(C_F)=g_F$; see
\citep[Lemmas~53.13.4--53.13.5 and Proposition~53.13.7]{StacksResidualGenusB}.
If $s_F$ is the separable residual degree,
Riemann--Hurwitz \citep[Lemma~53.12.2]{StacksResidualGenusB} gives
\[
 \deg\Delta_F=2g_F-2-s_F(2g(E_i)-2).
\]
For $g(E_i)\geq1$ the last term is nonpositive. For a rational
$E_i$, its positive contribution is $2s_F$, and
$\sum_{F/E_i}e_Fs_F\leq n$. Multiplying by $a_i$, summing,
and using (15) proves (17).
\end{proof}

\begin{remark}[The remaining counting issue]
These inequalities control the total normalized genus and
separable residual ramification by a constant times the covering
degree. They do not fix the residual branch locus or control all
point-mapped rational configurations: a rational $(-2)$-curve
has zero canonical intersection and contributes nothing to
(16). Consequently no bound on the number of coverings, and no
finite-generation conclusion, is inferred from (15)--(17) alone.
\end{remark}

% Complete actual-file proof audits: root and C, 20 November 2026.
\begin{lemma}[Proved: two-sided slopes for the saturated boundary grades]
\label{D:saturated-boundary-two-sided-slopes}
Use the fixed resolution in Proposition
\ref{D:tame-parabolic-boundedness}, refined once so that the
strict transforms $C_i$ of the original boundary are smooth
and pairwise disjoint, and their union with the exceptional
divisor has simple normal crossings without triple points.
Choose at every coefficient-Frobenius stage one minimal
extension $F_a$ of the canonical reflexive lattice $E_a$.
Extend the full parabolic flag from
$C_i-(C_i\cap\operatorname{Exc}(\pi))$ to $F_a|_{C_i}$
by saturation.  Denote its nonzero graded vector bundles by
$G_{a,i,j}$.

There is a constant $K$, depending only on the fixed
resolution, such that
\[
                -K\le\mu_{\min}(G_{a,i,j})
                    \le\mu_{\max}(G_{a,i,j})\le K
\]
for every rank, coefficient field, Kummer denominator,
stage, boundary component and grade.
\end{lemma}

\begin{proof}
Fix a positive exceptional divisor $Z$ with
$Z\cdot E_\nu<0$ for every exceptional component.
All minimal extensions have restrictions to $E_\nu$
in the same fixed two-sided slope windows.  The curve
Frobenius estimate, obtained from semistability of
Frobenius direct image and adjunction
\citep[Theorem~2.2]{C_Sun2008}
\citep[Proposition~1.3 and its proof]{C_MehtaPauly2007},
puts the restrictions of $F_Y^*F_a$ in fixed windows
as well.  Here $p$ is fixed, so the error in a single
Frobenius step is independent of rank.

On $Y-\operatorname{Exc}(\pi)$ the ordinary map
$\phi_a:F_Y^*F_a\to F_{a+1}$ is regular, and its inverse
has pole order at most $p-1$ along every $C_i$.
The bounded-pole lemma for exceptional slope windows
therefore gives one $B$, independent of all the
varying data, such that both maps extend as
\[
 \begin{split}
 F_Y^*F_a&\longrightarrow F_{a+1}(BZ),\\
 F_{a+1}&\longrightarrow
       F_Y^*F_a\bigl((p-1)\textstyle\sum_i C_i+BZ\bigr).
 \end{split}                                            \tag{D.SF.1}
\]
For the second application, twisting by the fixed
$(p-1)\sum_iC_i$ only shifts the exceptional windows
by fixed intersection numbers.  The bounded-pole
argument consequently applies to this target too.

The saturated flag on a smooth curve is a flag of
subbundles with locally free quotients.  At any point
of $C_i$ choose an adapted basis of its restriction;
that basis lifts to a basis of $F_a$ on a neighborhood
in the smooth surface.  Thus the full flag defines
a locally split parabolic bundle on the pair
$(Y,\sum_iC_i)$.  For a common denominator $n$ prime
to $p$, let $\mathcal F_a$ be the corresponding
bundle on the root stack
\[
                \mathcal X=\sqrt[n]{(Y,\sum_i C_i)}.
\]
Only the strict original boundary is rooted;
the exceptional divisor is not added to the logarithmic
boundary.

We claim that (D.SF.1) normalizes to maps
\[
 F_{\mathcal X}^*\mathcal F_a
        \longrightarrow\mathcal F_{a+1}(BZ),\qquad
 \mathcal F_{a+1}
        \longrightarrow F_{\mathcal X}^*\mathcal F_a(BZ).
                                                        \tag{D.SF.2}
\]
Away from $Z$ these are the genuine Frobenius
isomorphism of the tame local system and its inverse.
It remains to check that normalization does not
increase the exceptional pole order.

At a point of $C_i\cap Z$, take regular parameters
$t,z$, with $C_i=(t)$ and the exceptional component
$(z)$.  Write the coefficient of $BZ$ there as $b$.
On the Kummer chart $t=s^n$, the local ring is regular,
with parameters $s,z$.  A normalized root-matrix entry
is an ordinary matrix entry multiplied by an integral
power of $s$.  From (D.SF.1), after multiplication by
$z^b$, it can have poles only in $s$.  Its regularity
after inverting $z$ follows from the root-stack
Frobenius isomorphism away from the exceptional set.
Consequently it has nonnegative valuation at $(s)$.
Normality, or the intersection of the two localizations
at $s$ and $z$, now says that it is regular.
Exactly the same argument applies to the inverse:
the allowed ordinary pole $t^{-(p-1)}$ contributes
only an $s$-power and no $z$-pole.  This proves
(D.SF.2), with the same $B$, independently of $n$.
At points of $Z$ away from the strict boundary the
claim is simply (D.SF.1).

Let $N_i=\mathcal O_Y(C_i)|_{C_i}$ and let
$P_i=Z|_{C_i}$, an effective fixed divisor on $C_i$.
On the root gerbe over $C_i$, choose the character-line
convention in which its tautological line $\mathcal L$
satisfies $\mathcal L^n=N_i$ and the character-$j$
summand of $\mathcal F_a$ is
$G_{a,i,j}\otimes\mathcal L^j$.
If $pj=b+nc_j$, with $0\le b<n$, then
$0\le c_j\le p-1$.  Since multiplication by $p$ permutes
the characters, (D.SF.2) restricts, character by
character, to generically invertible maps
\[
 \begin{split}
 F_{C_i}^*G_{a,i,j}\otimes N_i^{c_j}
        &\longrightarrow G_{a+1,i,b}(BP_i),\\
 G_{a+1,i,b}
        &\longrightarrow
        F_{C_i}^*G_{a,i,j}\otimes N_i^{c_j}(BP_i).
 \end{split}                                            \tag{D.SF.3}
\]
Both are injective, because they are generically
isomorphisms between vector bundles on a smooth
integral curve.  This step uses the saturated flags
inside the same $F_a|_{C_i}$; it does not choose an
independent extension for each grade.

Put $K_i^0=B\deg P_i+(p-1)|\deg N_i|$.
The first injection gives
\[
 p\,\mu_{\max}(G_{a,i,j})
              \le\mu_{\max}(G_{a+1,i,b})+K_i^0.
\]
Dualizing the second injection gives precisely
the analogous inequality for the dual bundles:
\[
 p\,\mu_{\max}(G_{a,i,j}^{\vee})
              \le\mu_{\max}(G_{a+1,i,b}^{\vee})+K_i^0.
\]
The pairs consisting of the coefficient stage and
the character form finite cycles.  Taking the
maximum in each of the two inequalities around
such a cycle bounds both maxima by $K_i^0/(p-1)$.
Since $\mu_{\min}(G)=-\mu_{\max}(G^\vee)$, taking
the maximum over the finitely many $C_i$ proves
the assertion.
\end{proof}

\begin{remark}[The remaining quantitative presentation step]
The slope windows would give one fixed twist making
every grade, every flag subbundle and every quotient
globally generated with vanishing first cohomology
on each $C_i$.  Their section spaces then have
dimensions bounded by a constant times the rank.
The entire flag can be encoded by one adapted basis
of one such section space, with quadratically many
matrix coordinates.  A full bounded-degree
presentation of the evaluation, quotient and
Frobenius compatibility conditions, and the
Artinian openness argument for its unit classes,
must still be written before a quadratic-rank
representation bound is claimed.
\end{remark}

% Root and B: complete actual-file audit PASS, with the torsion-free cover scope.
\begin{lemma}[Proved: a stable annihilator uniform over finite punctured covers]
\label{D:cover-uniform-MCM-annihilator}
Let $R$ be a complete excellent normal local surface domain. Let
$I\subset R$ be an ideal such that, for every maximal Cohen--Macaulay
$R$-module $M$ and $z\in I$, multiplication by $z$ on $M$ factors
through a finite free $R$-module. Let $A$ be any finite normal,
$R$-torsion-free $R$-algebra which is etale over the punctured spectrum.
Thus every component dominates $\operatorname{Spec}R$, as for
normalization in a finite separable generic algebra. Then
\[
 I^2\subset\mathfrak N(A/R),
\]
where $\mathfrak N$ is the Noether different, defined without a
flatness assumption. Moreover, for every maximal Cohen--Macaulay
$A$-module $M$ and every finite $A$-module $N$,
\[
 I^3\operatorname{Ext}^1_A(M,N)=0.
\]
Equivalently, multiplication by every element of $I^3$ on $M$
factors through a finite free $A$-module. The same $I$ works for
every such cover algebra, irrespective of its degree.
\end{lemma}
\begin{proof}
Set $P=A\otimes_R A$, $T=P^{**}$, and let $K$ be the kernel of
the natural map $j:P\to T$, with duals taken over $R$.
The $R$-module $A$ is maximal Cohen--Macaulay. For $z\in I$, choose
a factorization of $z\operatorname{id}_A$ through $R^m$ and tensor
it with $A$. We obtain a factorization of $z\operatorname{id}_P$
through $A^m$. Since $A^m$ is reflexive, the first map kills $K$;
hence $zK=0$. Its restriction to the puncture extends uniquely to
a map $T\to A^m$. Composing with the second map and then with $j$
equals $z\operatorname{id}_T$ on the puncture, hence everywhere by
reflexivity. Thus
\[
 IK=0,\qquad IT\subset j(P).
\]

On the puncture the finite etale algebra has its separability
idempotent $e$. It is a section of $P$ there, and hence extends to
$T$. Multiplication $P\to A$ also extends to $T\to A$, and sends
$e$ to $1$. For $z,w\in I$, lift $ze$ to $u_0\in P$. For every
$a\in A$ the element $(a\otimes1-1\otimes a)u_0$ belongs to $K$.
Therefore $u=wu_0$ is annihilated by all these differences and its
product is $wz$. The differences generate the kernel of multiplication
$A\otimes_R A\to A$. This is precisely the defining separability
element for membership $wz\in\mathfrak N(A/R)$
\citep[Section 2, pp.~2--3]{C_IyengarTakahashiStable}. Finite sums prove
the first assertion.

Consider an $A$-extension $0\to N\to E\xrightarrow{v}M\to0$.
Its source $M$, viewed over $R$, is maximal Cohen--Macaulay: a system
of parameters of $R$ is also a system of parameters on each finite
component and is regular on $M$. For $c\in I$ there is therefore an
$R$-linear map $s:M\to E$ with $vs=c\operatorname{id}_M$.
Write the preceding separability element as
$u=\sum a_i\otimes b_i$, with product $wz$. The averaged map
\[
 s_A(m)=\sum_i a_i s(b_i m)
\]
is $A$-linear, since $u$ is annihilated by the diagonal differences.
Its composite with $v$ is $cwz\operatorname{id}_M$. Thus every product
of three elements of $I$ kills the extension class. Taking a finite
free presentation of $M$ proves the equivalent factorization assertion.
No trace division by the cover degree is used.
\end{proof}

\begin{remark}[The precise quantitative consequence and remaining gate]
When $I$ is $\mathfrak m$-primary, choose once parameters $x,y\in I^3$.
The prescribed two-parameter lifting proof in
Lemma~\ref{D:relative-Noether-different-jets} then gives the same
Maranda depth for every cover algebra $A$. If $A$ is a domain of
degree $t$ over $R$ and $M$ has $A$-rank $r$, then
\[
 \operatorname{length}_R M/(x^b,y^b)M
 =tr\operatorname{length}_R R/(x^b,y^b).
\]
Thus the canonical finite-jet vector space has dimension linear in
$tr$, with depth independent of the cover. This statement does not
give a presentation with $O(tr^2)$ coordinates or equations. Ordinary
matrices on that vector space still have $O(t^2r^2)$ entries, and
encoding the additional $A$-action can cost more. A linear-index
representation bound therefore remains unproved by this argument.
\end{remark}

\begin{example}[Refuted: bounding rational chains by local degree and different]
\label{D:fixed-different-unbounded-rational-chain}
Let \(k\) be algebraically closed of characteristic two. For \(m\geq1\),
put \(N=2m+1\) and
\[
 X_m=\operatorname{Spec}
 k[x,y,z]/(z^2+xz+y^N),\qquad
 \pi_m:X_m\longrightarrow\mathbb A^2_{x,y}.
\]
These are finite flat, generically separable degree-two normal
coverings, etale away from the fixed divisor \(x=0\).
At its generic point there is one extension of the valuation, with
\[
 e=1,\qquad f_{\mathrm{insep}}=2,\qquad
 f_{\mathrm{sep}}=1,\qquad \delta=1.
 \tag{R1}
\]
Nevertheless the minimal resolution at the point over the origin
has a chain of \(2m\) smooth rational curves, each of square \(-2\);
the absolute determinant of their intersection matrix is \(2m+1\).
The reduced curve over \(x=0\) has local normalization defect
\(\delta_{\mathrm{curve}}=m\).
Thus a fixed smooth base germ, fixed covering degree, and fixed
divisorial different do not bound the length or determinant of
these rational chains. This statement is local along a nonproper
branch curve; it does not supply a covering in the contracted-point
setup of the finite-generation problem.
\end{example}

\begin{proof}
The monic equation in \(z\) makes the covering free of rank two.
Its derivative in \(z\) is \(x\), so the map is etale on \(D(x)\)
and is generically separable.
The change of variables \(u=x+z,\ v=z\) gives
\[
 k[X_m]=k[u,v,y]/(uv+y^N).
 \tag{R2}
\]
This is a domain, for example by viewing the primitive polynomial
\(uv+y^N\) as a linear polynomial in \(u\) over the unique
factorization domain \(k[v,y]\).
Its Jacobian has entries \(v,u,y^{N-1}\), so its only singular
point is the origin. A hypersurface is Cohen--Macaulay, and
Serre's criterion proves normality
\cite[Tag 031S]{StacksRationalChainB}.

The quotient by \(x\) is
\(k[y,z]/(z^2+y^N)\), a domain with normalization
\[
 y=s^2,\qquad z=s^N.
\]
At the generic point its residue field is \(k(s)\), purely
inseparable of degree two over \(k(y)\). The height-one prime
\((x)\) in the normal ring is principal, so \(x\) is a uniformizer
there and \(e=1\). The monogenic different is generated by
the derivative \(x\), by
\cite[Example 49.12.5, tag 0BWB]{StacksRationalChainB}.
Its valuation is one. This proves (R1).
For the curve at the origin, completion gives
\(k[[s^2,s^N]]\subset k[[s]]\). Its missing exponents are exactly
\(1,3,\ldots,2m-1\), so the quotient has \(k\)-dimension \(m\).

We give the resolution explicitly, avoiding any restriction on
the characteristic in a classification theorem.
For \(i=0,\ldots,N-1\), take a smooth chart with coordinates
\(a_i,b_i\) and maps
\[
 u=a_i^{\,i+1}b_i^{\,i},\qquad
 v=a_i^{\,N-i-1}b_i^{\,N-i},\qquad y=a_i b_i.
 \tag{R3}
\]
Neighbouring charts are glued, where \(b_i\ne0\), by
\[
 a_{i+1}=b_i^{-1},\qquad b_{i+1}=a_i b_i^2.
 \tag{R4}
\]
The three formulas (R3) agree on these overlaps.

This glued map is proper. To verify this directly, blow up the
ideal \((u,y)\) on (R2). On the \(u\)-chart put \(y=ua\);
the equation gives \(v=u^{N-1}a^N\), a smooth affine plane.
On the \(y\)-chart put \(u=yb\); the remaining equation is
\(bv=y^{N-1}\). Repeat the same construction on that chart.
The next ideal \((b,y)\) is the unit ideal on its overlap with
the already smooth chart, so it glues to the unit ideal there.
Induction gives exactly (R3)--(R4). Every step is a blowup of
a coherent ideal, hence is projective
\cite[Tag 02NS]{StacksRationalChainB}; it is an isomorphism
away from the original singular point. The resulting smooth
surface is therefore a proper resolution.

For \(j=1,\ldots,N-1\), the exceptional curve \(C_j\) is
\(a_{j-1}=0\) in chart \(j-1\) and \(b_j=0\) in chart \(j\).
The transition \(a_j=b_{j-1}^{-1}\) makes it a projective line.
The normal-coordinate transition in (R4) is
\(b_j=a_{j-1}b_{j-1}^2\); hence its normal bundle is
\(\mathcal O_{\mathbf P^1}(-2)\) and \(C_j^2=-2\).
In the intermediate charts the two adjacent exceptional curves
are the coordinate axes; they meet transversely once.
Nonadjacent curves do not meet. These are all exceptional
components, and none is a rational \((-1)\)-curve, so the
resolution is minimal.

The determinant \(D_t\) of the positive tridiagonal matrix with
diagonal \(2\) and adjacent entries \(-1\) satisfies
\(D_0=1,D_1=2,D_t=2D_{t-1}-D_{t-2}\).
Thus \(D_t=t+1\); taking \(t=N-1=2m\) proves the asserted
determinant. The same explicit formal charts give the resolution
after completing at the origin, with the same exceptional
curves and local invariants.
\end{proof}

\begin{remark}[What global contraction data must still control]
The normalized residual genus in this family is zero and its
separable residual different is zero. Its unbounded contribution
is the arithmetic defect of the nonnormal residual curve.
A global estimate on a proper exceptional fibre may constrain
that defect; the local calculation does not decide this.
It does show that a determinant estimate cannot follow from the
divisorial \(\delta/e\) bound and the degree alone.
\end{remark}


\begin{lemma}[Proved: arithmetic defect over a proper curve when radial indices are one]
\label{D:proper-curve-arithmetic-defect}
Let \(k\) be algebraically closed, let \(Y\) be a smooth integral surface
over \(k\), and let \(E\subset Y\) be a smooth proper connected curve
of genus \(g\). Let \(\pi:Y'\to Y\) be a finite dominant generically
separable morphism of degree \(n\), with \(Y'\) normal and every
irreducible component dominating \(Y\).
Assume that every ramification index over the generic point of \(E\)
is one. Write \(\mathcal B=\pi_*\mathcal O_{Y'}\), let \(D_\pi\)
be its trace-discriminant divisor, and put \(C=Y'\times_Y E\).

Then \(C\) is reduced. If its normalization is
\(\coprod_F F^\nu\), with genera \(g_F\), its total normalization
defect
\(\delta(C)=\dim_k H^0(C,\nu_*\mathcal O_{C^\nu}/\mathcal O_C)\)
satisfies
\[
 \delta(C)=\sum_F(1-g_F)-n(1-g)+\frac{D_\pi\cdot E}{2}.
 \tag{AD1}
\]
Factor each finite map \(F^\nu\to E\) into relative Frobenius
followed by a separable map, of separable degree \(s_F\), and
denote its different divisor by \(\Delta_F\). Then
\[
 2\delta(C)+\sum_F\deg\Delta_F
       =2(g-1)\left(n-\sum_Fs_F\right)+D_\pi\cdot E.
 \tag{AD2}
\]
If the discriminants are supported on a fixed finite set of
divisors \(D_j\) and their coefficients satisfy
\(0\leq\operatorname{ord}_{D_j}(D_\pi)\leq c_j n\), then
\(\delta(C)\leq C_E n\), with \(C_E\) independent of the covering.
\end{lemma}

\begin{proof}
A finite dominant normal surface over a regular surface is flat.
Indeed its local rings are Cohen--Macaulay, its fibres have
dimension zero, and miracle flatness applies
\cite[Tag 00R4]{StacksArithmeticDefectB}.
Thus \(\mathcal B\) is a vector bundle of rank \(n\).
At the generic point of \(E\), the assumed indices say that its
base change to \(E\) is a product of fields, although those fields
may be inseparable over \(k(E)\). Hence \(C\) is generically
reduced. Its structure sheaf is torsion-free over \(\mathcal O_E\)
by flatness, so its nilradical, which is generically zero, is zero.

The trace pairing is a map
\(\mathcal B\to\mathcal B^\vee\), generically invertible by
separability. Its determinant is a nonzero section of
\((\det\mathcal B)^{-2}\); by definition its zero divisor is
\(D_\pi\). Therefore
\[
 (\det\mathcal B)^2=\mathcal O_Y(-D_\pi),\qquad
 \deg(\mathcal B|_E)=-\tfrac12 D_\pi\cdot E.
\]
This is the discriminant construction in
\cite[Tag 0BVH]{StacksArithmeticDefectB}.
Finite pushforward and Riemann--Roch on \(E\) give
\[
 \chi(C,\mathcal O_C)
   =\chi(E,\mathcal B|_E)
   =n(1-g)-\tfrac12 D_\pi\cdot E.
\]
Normalization of a reduced proper curve is finite. Its quotient
sheaf is supported at finitely many points, and its Euler
characteristic is its length. The normalization exact sequence
therefore proves (AD1).

The pure-then-separable factorization preserves the genus of each
normalized source component; this is
\cite[Lemmas 53.13.4--53.13.5 and Proposition 53.13.7]
{StacksResidualGenusB}. Riemann--Hurwitz, with the different, gives
\[
 \deg\Delta_F=2g_F-2-s_F(2g-2)
\]
by \cite[Lemma 53.12.2]{StacksResidualGenusB}.
Adding these equalities to twice (AD1) proves (AD2).
Finally there are at most \(n\) components, all their genera are
nonnegative, and
\[
 |D_\pi\cdot E|\leq n\sum_jc_j|D_j\cdot E|.
\]
These estimates in (AD1) prove the claimed linear bound.
\end{proof}

\begin{remark}[Connection with the fixed different bound]
At a divisorial discrete valuation ring, the discriminant
coefficient is \(\sum_F f_F\delta_F\), where \(f_F\) is the full
residue degree and \(\delta_F\) the different exponent.
One can see the formula directly from the trace-dual lattice:
the valuation of the trace determinant is the length of
\(\mathcal B^\vee/\mathcal B\); decomposing over the integral
closure gives length \(f_F\delta_F\) at each extension.
Thus a fixed bound on every \(\delta_F/e_F\), together with
\(\sum_Fe_Ff_F=n\), gives the coefficient bounds in the lemma.

The condition \(e_F=1\) is what makes the Cartier fibre reduced.
When some index is larger, the fibre has nilpotents at its
generic points and its Euler characteristic no longer directly
equals the normalized-genus contribution minus a nonnegative
normalization defect. No bound for that remaining nonreduced
case, and no finite-generation conclusion, is asserted here.
\end{remark}


% Root: complete actual-file audit PASS; relative refinement of the uniform cover lemma.
\begin{corollary}[Proved: the uniform cover ideal over Artinian parameters]
\label{D:cover-uniform-relative-jets}
Suppose in Lemma~\ref{D:cover-uniform-MCM-annihilator} that $R$ is
equicharacteristic with algebraically closed coefficient field $k$,
and choose $I$ inside the sum of the Noether differents for regular
Noether normalizations of $R$. For every finite local Artinian
$k$-algebra $C$ with residue $k$, set $R_C=R\otimes_k C$ and
$A_C=A\otimes_k C$. If $M_C$ is finite and $C$-flat with
maximal Cohen--Macaulay closed fibre over $A$, multiplication by
every element of $I^3$ on $M_C$ factors through a finite free
$A_C$-module. If $I$ is $\mathfrak m$-primary, one choice of parameters $x,y\in I^3$ gives
the prescribed two-parameter jet-lifting depth uniformly in the
cover $A$, its degree, and the Artinian algebra $C$, for targets
on which $x,y$ is a regular sequence. No flatness over $C$ is
required of that target.
\end{corollary}
\begin{proof}
Only finitely many normalizations $S_i\simeq k[[u_i,v_i]]$ are
needed to express a finite generating set of $I$ as sums of
Noether-different elements. The closed fibre of $M_C$, viewed as
an $R$-module, is maximal Cohen--Macaulay, by the parameter argument
in the preceding lemma. It is consequently free over each $S_i$.
Lift a basis to obtain a map $(S_i\otimes_k C)^a\to M_C$.
Nilpotent Nakayama makes it surjective. Flatness of $M_C$ over $C$
then makes the kernel reduce to zero modulo the maximal ideal of
$C$; nilpotent Nakayama makes the kernel zero as well.

For a Noether-different element of $R/S_i$, its separability-element
factorization therefore base-changes to an $R_C$-linear factorization
of multiplication on $M_C$ through a finite free $R_C$-module. This
is exactly the source-only argument in
Lemma~\ref{D:relative-Noether-different-jets}. In particular $I$
annihilates $\operatorname{Ext}^1_{R_C}(M_C,N_C)$ for every target,
without a flatness condition on that target.

The elements of $I^2$ in $\mathfrak N(A/R)$ have explicit
separability elements $u\in A\otimes_R A$ from the preceding
lemma. Their base changes still annihilate the diagonal ideal of
$A_C\otimes_{R_C}A_C$, and still multiply to the same elements of
$I^2$. Given an $A_C$-extension with source $M_C$, take the
$R_C$-linear $c$-splitting supplied by $c\in I$ and average it by
this base-changed $u$. It is $A_C$-linear and splits the extension
after multiplication by the product in $I^3$. Applying this to a
finite free $A_C$-presentation of $M_C$ gives the asserted
factorization.

When $I$ is $\mathfrak m$-primary, fix parameters $x,y\in I^3$.
They remain a regular sequence on every such relative source.
The two successive divisions in the prescribed-map proof apply
verbatim to any target on which $x,y$ is regular, losing only
one exponent in each coordinate. Neither that proof nor the
averaging uses the degree of $A/R$ or the nilpotence length of $C$.
\end{proof}

% Complete proof; independently audited by root, B, and C, with primary sources checked.
\subsection{A quadratic presentation of the tame boundary flags}

\begin{lemma}[Proved: linear resolutions of an entire bounded-slope flag]
\label{D:filtered-linear-curve-resolution}
Fix a smooth projective curve $C\subset\mathbf P^N_k$ and real numbers
$a\le b$.  There are integers $m,K$ with the following property.
For every flag of vector subbundles
\[
        0=V_0\subset V_1\subset\cdots\subset V_l=V,
            \qquad l\le r=\operatorname{rk}V,
\]
whose nonzero grades $G_j=V_j/V_{j-1}$ have all slopes in $[a,b]$,
the flag is the sheafification of a filtered graded free complex
over $P=k[x_0,\ldots,x_N]$.  This complex has length at most $N+1$,
its term in homological degree $i$ is a sum of at most $Kr$
copies of $P(-m-i)$, and its differential is a matrix of linear
forms preserving coordinate flags in all terms.  The associated
graded complex resolves the corresponding truncated section
modules of the $G_j$.

For fixed ranks and multiplicities, these filtered complexes are
parameterized by a closed affine scheme with $O(r^2)$ variables
and $O(r^2)$ quadratic equations.  Exactness, flatness and the
subbundle conditions are imposed by passing to an open subset
after identifying the total cohomology with a prescribed bundle
on $C$.  The number of multiplicity arrays needed is at most
$\exp(Kr)$.
\end{lemma}

\begin{proof}
Choose $m$ so large that
$a+(m-1)\deg\mathcal O_C(1)>2g(C)-2$.
Every grade, every partial subbundle and every partial quotient
then has vanishing $H^1$ in all twists of degree at least $m-1$.
This follows from the slope bounds and their preservation under
extensions.  A fixed further increase of $m$ ensures global
generation, by applying the same vanishing after subtracting
one point.

For such a bundle $W$ put
\[
          M(W)=\bigoplus_{d\ge m}H^0(C,W(d)).
\]
It is a finitely generated graded $P$-module, with no components
below $m$.  The local-cohomology comparison with its associated
sheaf shows that $H^0_{\mathfrak m}M(W)=0$, that
$H^1_{\mathfrak m}M(W)$ vanishes in degrees at least $m$,
and that $H^2_{\mathfrak m}M(W)$ vanishes in degrees at
least $m-1$.  Higher local cohomology vanishes because the
sheaf is supported on a curve.  Hence $M(W)$ is $m$-regular.
The equivalence of regularity with linear truncation
\citep[Proposition~1.1 and Theorem~1.2, pp.~95--97,
including the proof]{EisenbudGoto1984D}
gives an $m$-linear graded free resolution.
Hilbert's syzygy theorem bounds its length by $N+1$.

Its Betti numbers are linear in the rank bound.  Indeed, for
$d\ge m$ Riemann--Roch gives
\[
       h^0(C,W(d))
        =\operatorname{rk}(W)\deg\mathcal O_C(1)\,d
                      +\deg W+\operatorname{rk}(W)(1-g).
\]
The coefficients of
$(1-t)^{N+1}\sum_{d\ge m}h^0(C,W(d))t^d$
are therefore bounded by $K\operatorname{rk}(W)$.
For an $m$-linear resolution, each homological degree
occurs in its own internal degree $m+i$, so these
coefficients, with alternating signs, are precisely
the Betti numbers.

The $H^1$ vanishing makes
\[
 0\to M(V_{j-1})\to M(V_j)\to M(G_j)\to0
\]
exact.  Apply the graded horseshoe construction successively
to the $m$-linear resolutions of the grades.  In degree $i$
all summands have the same shift $-m-i$; all off-diagonal
differentials consequently still have entries of degree one.
The resulting resolution is filtered by coordinate summands,
and its associated graded is the chosen sum of grade resolutions.
Its total multiplicities are bounded by
$K\sum_j\operatorname{rk}G_j=Kr$, rather than by $Kr^2$.

Choose the coordinate flags once the grade multiplicities
are fixed.  Preservation of them says that specified blocks
of each differential are zero.  The remaining blocks have
$O(r^2)$ scalar coefficients.  The equations $d^2=0$ have
$O(r^2)$ scalar coefficients and degree two, since the
complex length and the number of variables of $P$ are fixed.
The loci where the total and associated-graded complexes
have only base-flat degree-zero cohomology, and where
those cohomology sheaves are vector bundles of the
prescribed ranks on $C$, are open by the relative
perfect-complex truncation criterion already used in
Lemma~\ref{D:polynomial-monad-diagrams}.
The exact filtered complex then gives actual subbundles
and locally free quotients over every Artinian parameter
ring, not only at its reduced points.

For each of the fixed $N+2$ degrees the grade multiplicities
are nonnegative integers with sum at most $Kr$, in at
most $r$ ordered positions.  The number of such arrays
is bounded by a fixed power of
$\binom{Kr+r}{r}$, hence by $\exp(K'r)$.
\end{proof}

\begin{lemma}[Proved: quadratic witnesses for filtered diagrams]
\label{D:filtered-polynomial-diagrams}
The conclusion of Lemma~\ref{D:polynomial-monad-diagrams}
remains valid for split-filtered complexes of the preceding
lemma, provided the total multiplicities in each degree are
$O(r)$ and the complex lengths and bundle types are fixed.
The number of filtration steps may be at most $r$.
Morphisms required to preserve the filtration, and equalities
of these morphisms, have witnesses in $O(r^2)$ variables
and equations of bounded degree.
\end{lemma}

\begin{proof}
We use the same finite affine cover and the same fixed
contractions of its Cech Hom complexes as in
Lemma~\ref{D:polynomial-monad-diagrams}.  In every
matrix space retain only the blocks allowed by the
two coordinate flags.  These form the filtered Hom
complex.  The fixed Cech differential and its
chosen contraction act on the fixed bundle factor
and therefore preserve this block condition.
The variable complex differentials preserve the
flags as well.

The perturbation still raises the complex-Hom
column.  Its nilpotence length is bounded by the
fixed lengths of the complexes, not by the number
of grades.  Thus the finite geometric-series
elimination produces polynomial formulas of the
same bounded degree.  The dimensions of all
remaining spaces are at most those of the
unfiltered Hom spaces, namely $O(r^2)$.
Cycles, boundaries witnessing equality, and
composition are treated exactly as before.

For clarity, this computes the derived morphisms
of the exact category of filtered sheaves with
strict exact sequences.  Locally, a split-filtered
free term is a sum of the elementary filtered
free modules whose filtration starts at one
specified step.  These are projective for the
strict exact structure.  Consequently the
filtered locally free complexes and the Cech
resolution compute the required morphisms.
An associated-graded quasi-isomorphism is the
weak equivalence used here.  We are not replacing
the category by its Ext algebra or making a
growing-length iterated-extension construction.
All contraction identities are over $k$ and
remain valid after tensoring by any parameter
algebra, including a nonreduced Artinian one.
\end{proof}

\begin{theorem}[Proved: quadratic defining-characteristic growth with original tame boundary]
\label{D:tame-global-quadratic-representation-bound}
Let $S$ be a normal proper connected surface over an
algebraically closed field of characteristic $p>0$,
let $T\subsetneq S$ be closed, and let $Q(S,T)$ require
tameness only at the original codimension-one components
of $T$.  There is a constant $C_{S,T}$ such that
\[
 \#\bigl(
  \operatorname{Hom}_{\mathrm{cont}}
       (Q(S,T),\operatorname{GL}_r(\mathbf F_{p^f}))
       /\operatorname{GL}_r(\mathbf F_{p^f})\bigr)
          \le p^{C_{S,T}fr^2}                           \tag{D.TQ.1}
\]
for every $r,f\ge1$.
This is a representation-count bound, not a claim
that the full group is finitely generated.
\end{theorem}

\begin{proof}
Enlarge the finite omitted set as in
Theorem~\ref{D:tame-parabolic-fixed-rank-growth};
the desired group is a quotient of that larger-open
group, so it suffices to count for the latter.
Fix the smooth projective resolution $Y$, the positive
exceptional cycle $Z$, and the smooth disjoint strict
boundary curves $C_i$ of
Lemma~\ref{D:saturated-boundary-two-sided-slopes}.
Put $Y^\circ=Y-\operatorname{Supp}Z$ and
$C=\sum_iC_i$.  All embeddings and ample twists
below are fixed.

\emph{Linear sizes of all the objects.}
For the coefficient-Frobenius stages choose one
minimal extension $F_a$ on $Y$.  The canonical
Euler estimate and the relative Euler estimate
give $\chi(F_a)=O(r)$; their first Chern classes
have intersection $O(r)$ with fixed divisors.
The forward and inverse pole comparisons give
uniform upper slopes for $F_a$ and its dual,
both on $Y$ and on a fixed smooth ample curve.
Consequently
\[
                    h^j(Y,F_a(t))\le C_t r
\]
for every fixed twist and $0\le j\le2$, by the
same restriction, duality and Riemann--Roch proof
as in Theorem~\ref{D:global-quadratic-representation-bound}.
This also applies after one Frobenius pullback
and after tensoring by any of finitely many fixed
bundles.  Only one $p$-Frobenius step is used.

The full saturated flag on each $F_a|_{C_i}$ has
all graded slopes in one fixed two-sided interval
by Lemma~\ref{D:saturated-boundary-two-sided-slopes}.
Lemma~\ref{D:filtered-linear-curve-resolution}
therefore encodes its entire flag by one filtered
linear complex with total term sizes $O(r)$.
It does not introduce $r$ independent rank-$r$
surface lattices.

\emph{A fixed number of digit transforms.}
Fix a one-step digit type from the preceding
fixed-rank theorem.  Its digits are nondecreasing,
belong to $\{0,\ldots,p-1\}$, and its permutation
preserves source order inside each equal-digit block.
At one boundary curve we write the saturated flag in its
decreasing convention, so that the suffix of grades with
digit at least $h$ is a vector subbundle, denoted $V_{\ge h}$.
This reverses the indexing convention of the increasing flag
in Lemma~\ref{D:filtered-linear-curve-resolution}; it changes
no object or slope bound.  There are only $p-1$ relevant thresholds.
Define the elementary transform
\[
 E_h=\ker\bigl(F_a\longrightarrow
             i_*(F_a|_{C_i}/V_{\ge h})\bigr).
                                                        \tag{D.TQ.2}
\]
It is a vector bundle.  In an adapted local basis,
its columns are the original columns for digit
at least $h$, and $t$ times the remaining columns,
where $t$ is an equation of $C_i$.
Put $H=F_Y^*F_a$ and
\[
 L_c=H+\sum_{h=1}^{p-1}F_Y^*E_h(hC_i)
                      \ \subset\ H((p-1)C_i).           \tag{D.TQ.3}
\]
Its local basis is $t^{-c_j}F^*v_j$, so it too is
a vector bundle.  For the fixed finite collection
of disjoint $C_i$, perform these constructions
independently; the total number of auxiliary
bundles per stage is bounded solely by $p$ and
the fixed number of boundary components.

All these auxiliary objects have the same linear
cohomology bounds.  In (D.TQ.2) the quotient is
a curve bundle with fixed slope bounds and rank
at most $r$.  Its degree and Euler characteristic
are $O(r)$; the kernel is bounded between
$F_a(-C_i)$ and $F_a$.  Riemann--Roch and these
inclusions bound its cohomology in all fixed
twists.  For its one-step Frobenius pullback,
the first Chern character scales by $p$ and the
second by $p^2$; the latter is already $O(r)$.
For (D.TQ.3), the quotient of $H((p-1)C_i)$
by $L_c$ has a filtration of fixed length at most
$p-1$ whose curve grades are Frobenius pulls
of whole flag grades with fixed normal twists.
Their slopes, degrees and cohomology are bounded
linearly in their ranks.  The same conclusions
follow for $L_c$.  Moreover, the inclusions
$H\subset L_c\subset H((p-1)C)$ restrict generically injectively
to every exceptional component.  They give fixed two-sided
exceptional slope windows for $L_c$, because those for $H$
follow from one-step Frobenius and the divisor $(p-1)C$ is fixed.
Thus the exceptional pole lemma applies to both normalized
one-step arrows between $L_c$ and $F_{a+1}$ with one fixed
constant $B$, independently of rank, denominator, and stage.
Beilinson monads for these finitely many objects thus have fixed types and
lengths, with total multiplicities $O(r)$.

\emph{The canonical transformed flag.}
Let $W_h$ be the quotient consisting of the
consecutive grades with digit exactly $h$,
including their full internal flag.  The local
basis calculation gives a canonical decomposition
\[
          L_c|_{C_i}\simeq
          \bigoplus_{h=0}^{p-1}
                    F_{C_i}^*W_h\otimes N_i^h,
           \qquad N_i=\mathcal O_Y(C_i)|_{C_i}.           \tag{D.TQ.4}
\]
This is not an independently chosen splitting.
For $h>0$, the natural map
$F_Y^*E_h(hC_i)|_{C_i}\to L_c|_{C_i}$
factors through the natural surjection onto
$F_{C_i}^*W_h\otimes N_i^h$ and identifies it
with its summand.  For $h=0$, use $H|_{C_i}$
and its quotient onto $F_{C_i}^*W_0$.
In the local basis, columns of larger digit
vanish after restriction, columns of smaller
digit acquire a positive $t$-power, and exactly
the columns of digit $h$ remain.  These
factorizations characterize (D.TQ.4) uniquely.
They are preserved by every adapted change of
basis and commute with arbitrary parameter
base change.

Within each of its at most $p$ summands retain
the Frobenius flag.  The prescribed permutation
shuffles the summands and their subflags while
preserving order within each summand.  This
gives the intrinsic transformed full flag on
$L_c|_{C_i}$.  It is the flag used in the exact
ceiling-valuations criterion of
Theorem~\ref{D:tame-parabolic-fixed-rank-growth}.

\emph{Bounded-degree witness schemes.}
We now spell out how the descriptions above
produce polynomial witnesses, including at
nonreduced parameter points.
Use the Beilinson monads of
Theorem~\ref{D:global-quadratic-representation-bound}
for the surface bundles and auxiliary sheaves,
and the filtered linear complexes of
Lemma~\ref{D:filtered-linear-curve-resolution}
for the boundary flags and their Frobenius
digit-block flags.  Every total term size is
$O(r)$ and all complex lengths and fixed bundle
types are bounded.

Identify the total boundary complex with
$F_a|_{C_i}$ using the fixed resolution of
$\mathcal O_{C_i}$ and the same degree-zero
truncation construction used in the preceding
quadratic theorem.  Require the total and
graded complexes to have only base-flat
degree-zero cohomology, locally free on $C_i$
of the prescribed ranks.  Because the total
cohomology is the actual restriction, its
subquotients are annihilated by the fixed ideal
of $C_i$ over every parameter algebra.
These are open conditions on the perfect
complex families after the derived support
and identification equations have been imposed.

The kernels in (D.TQ.2) are encoded by the
derived exact-sequence diagram, with their
auxiliary monads and a surjective quotient map.
To encode (D.TQ.3), use the natural map from
$H\oplus\bigoplus_hF_Y^*E_h(hC_i)$ to
$H((p-1)C_i)$, introduce its bundle image $L_c$
and its coherent quotient, and impose the
factorization, surjectivity onto the image,
and the exact quotient sequence.  Cone
quasi-isomorphisms and open flatness/truncation
conditions express exactly this diagram.
There are only a fixed number of nodes.
The actual local diagonal description above
shows that every genuine representation supplies
such witnesses.

For (D.TQ.4), introduce its direct-sum map and
require the stated factorizations of the
natural maps from the threshold transforms.
The quotient maps are surjective, so this
forces the canonical map, rather than an
arbitrary choice.  Impose that it be an
isomorphism.  The Frobenius internal flags
are identified by filtered derived maps;
their one-step Frobenius, fixed normal twists,
and at most $p$ direct sums are fixed
operations.  The chosen shuffle is only a
fixed block pattern.

Finally introduce the forward and inverse
maps between $L_c$ and $F_{a+1}$ with the
fixed exceptional poles $BZ$, and impose
their inverse equations after clearing
$2BZ$.  Require their restrictions to each
$C_i$ to preserve the transformed and target
flags.  Use the coherent six-node tensor
construction of the preceding quadratic
theorem for the pole twists, so that second
twists and compositions are connected to
the original arrows.

Lemmas~\ref{D:polynomial-monad-diagrams} and
\ref{D:filtered-polynomial-diagrams} express
all these maps and equalities in $O(r^2)$
variables and equations of bounded degree
per stage.  In the filtered case there is
one block-triangular Hom space, not a
separate full-rank map for every flag step.
The finite number of threshold transforms
and digit blocks is bounded by $p$, which
is fixed.  All cone, bundle, slope-window
and exactness conditions pass to an open
subset and need no additional quantitative
degree estimate.

This construction is effective over every
local Artinian parameter algebra.  The
relative truncation conditions supply actual
bundles and strict flags.  Locally the
elementary transforms are diagonal with
entries $t$ or $1$; their Frobenius pullbacks
have entries $t^p$ or $1$.  Since $t$ remains
a non-zero-divisor after any parameter base
change, these transformations have exactly
the required images and quotients even
when the parameter ring is nonreduced.
The cone conditions retain any other
negative Tor terms in the ambient
projective-space computations, as in the
preceding quadratic theorem.

\emph{Valid words and class fibres.}
For a length-$f$ type word, recover the unique
weights from $\alpha_{\rm next}=p\alpha-c$.
Discard the invalid words before making any
unit-class assertion.  For a valid word
the common denominator $n$ is prime to $p$.
The exact local ceiling criterion identifies
every witness cycle on $Y^\circ$ with a
unit $p^f$-Frobenius bundle on
$\sqrt[n]{(Y^\circ,C\cap Y^\circ)}$.
Its fixed-frame torsor is finite etale,
and therefore gives a representation of
the original-boundary tame quotient.
Conversely every representation is supplied
by some word and some witness scheme.

Fix a closed witness point and its reference
unit cycle.  The functor of isomorphisms to
the variable cycle is represented using
proper Hom spaces between the underlying
$F_a$, with maps in both directions having
poles at most $BZ$.  Preservation of the
full flags and compatibility with the
one-step arrows are closed equations.
It is an affine functor of finite
presentation, by the same proper Hom
theorem used in the preceding quadratic
proof.

For a small extension of Artinian parameter
rings, a specified unit-isomorphism lifts
uniquely on the root stack, by invariance
of finite etale covers under nilpotent
thickenings.  This can be checked on its
etale Kummer atlas and the corresponding
descent diagram.  Pushforward to the
ordinary parabolic lattices gives a
specified flag-preserving isomorphism
on $Y^\circ$.  The uniform exceptional
pole lemma extends it and its inverse
to $Y$ with the same $B$, independently
of the Artinian length.  Preservation
of flags persists at the points of
$C_i\cap Z$: their quotients are locally
free on $C_i$ over the parameter ring,
so multiplication by a local parameter
of $Z|_{C_i}$ is injective.  A map
which vanishes off those points is
therefore zero.
All cleared inverse and Frobenius
equations extend for the same reason.
The proper unit-Isom scheme thus lifts
the specified Artinian point.  The
infinitesimal smoothness criterion
makes its image open.

Distinct unit classes give disjoint
open subsets of the witness scheme.
They cover its closed points and hence
the whole finite-type scheme.  Thus
there is a finite clopen partition
by representation classes, and their
number is at most the number of
irreducible components.

\emph{The count.}
For each word and multiplicity array,
the witness scheme is an open subset
of a closed affine scheme defined
by $O(fr^2)$ equations of a fixed
degree $D$.  Its number of irreducible
components is at most $D^{O(fr^2)}$.
The surface monad arrays have only
polynomially many choices per stage.
The boundary grade arrays contribute
at most $\exp(Cr)$ per stage by
Lemma~\ref{D:filtered-linear-curve-resolution}.
The ranks are compositions of $r$,
the digit choices number at most
$p^{Cr}$, and the grade permutations
number at most $(r!)^C$, for a fixed
number of boundary curves.
Over $f$ stages their total is
$\exp(O(fr\log(r+1)))$, which is
absorbed by $\exp(O(fr^2))$.
Summing over all valid choices proves
(D.TQ.1).  The final passage to the
original group is by its quotient
map from the group with enlarged
finite omitted set.
\end{proof}

% END AUDITED POST-CP09 INTEGRATION

% Post-checkpoint-13 fragment; full mathematical and primary-source audits passed.
\subsection{Second cohomology and central lifting obstructions}

\begin{theorem}[Proved: a uniform cross-characteristic second-cohomology bound]
\label{D:surface-tame-cross-characteristic-H2}
Let $S$ be a normal proper connected surface over an algebraically
closed field $k$ of characteristic $p>0$, let $T\subsetneq S$ be
closed, and put $U=S-T$ and $Q=Q(S,T)$.  Thus $Q$ requires tameness
only at the original codimension-one boundary valuations.
There is a constant $c=c(S,T)$ such that, for every finite field
$\Lambda$ of characteristic $\ell\ne p$ and every finite-dimensional
continuous $\Lambda[Q]$-module $M$,
\[
          \dim_\Lambda H^2(Q,M)\le c\dim_\Lambda M.       \tag{D.H2.1}
\]
For an open subgroup $H\subset Q$ and a finite-dimensional continuous
$\Lambda[H]$-module $M$, the corresponding bound is
\[
          \dim_\Lambda H^2(H,M)
                \le c[Q:H]\dim_\Lambda M.               \tag{D.H2.2}
\]
All group cohomology is continuous.
\end{theorem}

\begin{proof}
Write $G=\pi_1(U)$ and $N=\ker(G\to Q)$.  By the definition of
this quotient, $N$ is the closed subgroup generated by the
$G$-conjugates of the images of the wild inertia groups at the
original boundary divisors.  Those groups are pro-$p$.
The $N$-action on $M$ is trivial.  A continuous homomorphism
from $N$ to the additive $\ell$-group $M$ vanishes on every
one of the generating pro-$p$ groups, and hence vanishes.
Consequently
\[
                        H^1(N,M)=0.
\]
The degree-two inflation map therefore gives an injection
\[
                H^2(Q,M)\lhook\joinrel\longrightarrow H^2(G,M).
                                                        \tag{D.H2.3}
\]
This is the inflation criterion in
\citep[II, Proposition~1.34 and Remark~1.35, pp.~71--72;
continuous cochains in II, Proposition~4.2, p.~87]{MilneCFT2020D}.
One can also check the injection directly: a continuous
$M$-extension of $Q$ whose pullback splits over $G$ has a
splitting $G\to E$.  Its restriction to $N$ is a homomorphism
$N\to M$, hence zero, so the splitting descends to $Q$.
This argument also explains why arbitrary prime-to-$p$
finite abelian coefficients will be allowed below.

Let $\mathcal M$ be the locally constant sheaf on $U$
associated with $M$.  There is a second injection
\[
              H^2(G,M)\lhook\joinrel\longrightarrow
                        H^2_{\mathrm{et}}(U,\mathcal M).
                                                        \tag{D.H2.4}
\]
For completeness, use the Cartan--Leray spectral sequence
for the system of all connected finite Galois covers of
$U$, as in \citep[Theorem~14.9, Remark~14.10 and
Proposition~14.13, pp.~96--98]{MilneLEC2013H2D}.
Its geometric degree-zero term is $M$.  Its degree-one
term is zero: after trivializing the coefficient sheaf,
every degree-one class is a finite torsor and dies on a
further finite cover.  In the first-quadrant spectral
sequence, the only possible incoming differential to
$(2,0)$ is therefore zero.  There are no outgoing
differentials from this term.  Its limiting term embeds
as the last filtration piece of total degree two,
which proves (D.H2.4).  We make no assertion that all
higher cohomology is group cohomology.

Choose a fixed smooth projective resolution
$\pi:Y\to S$ such that the reduced inverse image of $T$
and the exceptional curves have simple normal crossings.
Write $V=\pi^{-1}(U)$ and $h:V\to U$.  The map $h$ is
proper and birational, and it is an isomorphism outside
a finite set $\Sigma\subset U$.  The proper birational
map to the normal surface has connected geometric fibres.
Proper base change therefore gives
\[
             h_*h^*\mathcal M=\mathcal M,
\qquad
 (R^1h_*h^*\mathcal M)_{\bar x}
     =M_{\bar x}\otimes_\Lambda H^1(Y_{\bar x},\Lambda).
                                                        \tag{D.H2.5}
\]
The second identity can equivalently be checked after
etale-locally trivializing $\mathcal M$.  We use proper
base change in the precise form
\citep[Theorems~17.7 and~17.10, pp.~116--117]
{MilneLEC2013H2D}.
The sheaf in degree one is supported on $\Sigma$.
Nilpotents on the geometric fibres do not affect their
etale cohomology.  For their reduced simple-normal-crossing
curves, normalization and the node exact sequence show
\[
 \dim_\Lambda H^1(Y_{\bar x},\Lambda)
     =2\sum_jg(E_{x,j})+b_1(\Gamma_x),                    \tag{D.H2.6}
\]
where $E_{x,j}$ are the smooth components and $\Gamma_x$
is their dual graph; a zero-dimensional fibre contributes
zero.  In particular the sum $b$ of these dimensions over
$x\in\Sigma$ is independent of $\Lambda$.

The Leray spectral sequence for $h$
\citep[Theorem~12.7, pp.~82--83]{MilneLEC2013H2D}
now implies
\[
 \dim_\Lambda H^2(U,\mathcal M)
       \le \dim_\Lambda H^2(V,h^*\mathcal M)
                                +b\dim_\Lambda M.       \tag{D.H2.7}
\]
Here the direction is essential.  The term
$E_2^{2,0}=H^2(U,\mathcal M)$ has only the possible
incoming differential from $E_2^{0,1}$; its quotient
by that image is $E_\infty^{2,0}$, a subspace of
$H^2(V,h^*\mathcal M)$.  Thus no degeneration
assertion, and no vanishing of $R^2h_*$, is needed.

By Theorem~\ref{D:surface-uniform-conductor}, the
logarithmic conductors of $h^*\mathcal M$ along the
fixed boundary $Y-V$ have a bound independent of
$\ell$, the coefficient field, and the rank.
The theorem applies to the finite quotient cover
through which the representation factors; a bound
on its ramification filtration bounds every such
module.  At the strict transforms of the original
divisors the ramification is tame, and the remaining
components have the proved uniform exceptional bound.
The passage from these generic logarithmic bounds to the
curvewise bound required by the Betti theorem is exactly
Proposition~\ref{D:surface-curvewise-conductor}, using
\citep[Proposition~5.8(3)]{C_HuTeyssierRam2025}.
Hu--Teyssier's finite-coefficient theorem therefore
provides $b'$ depending only on this fixed pair with
\[
                   \dim_\Lambda H^2(V,h^*\mathcal M)
                            \le b'\dim_\Lambda M;
\]
see \citep[Theorem~7.34 and its proof, p.~41]
{C_HuTeyssierBetti2025}.  This is its $j=2$ case on
a smooth proper compactification; if the boundary
is empty the same proper smooth Betti bound applies.
Combining the two injections and (D.H2.7) proves
(D.H2.1), with $c=b+b'$.  Finally Shapiro's lemma
identifies $H^2(H,M)$ with
$H^2(Q,\operatorname{Ind}_H^Q M)$, whose coefficient
rank is $[Q:H]\dim M$, proving (D.H2.2).
\end{proof}

\begin{corollary}[Proved: polynomially many central obstruction classes]
\label{D:surface-tame-central-obstructions}
In the preceding theorem we may increase $c$ to an integer so
that, for every finite continuous abelian $Q$-module $A$ of
order prime to $p$,
\[
                         |H^2(Q,A)|\le |A|^c.            \tag{D.H2.8}
\]
In particular, for $q=p^f$ and $d=\gcd(r,q-1)$, the central
obstructions to lifting homomorphisms
$Q\to\operatorname{PSL}_r(\mathbf F_q)$ to
$\operatorname{SL}_r(\mathbf F_q)$ take at most $d^c$
values.  The obstructions for
$Q\to\operatorname{PGL}_r(\mathbf F_q)$ and its
$\operatorname{GL}_r(\mathbf F_q)$ lift take at most
$(q-1)^c$ values.

For every fixed integer $B$ there is a single open normal
subgroup $H_B\subset Q$ such that restriction to $H_B$
kills every class in $H^2(Q,\mathbf Z/d\mathbf Z)$
for every $1\le d\le B$ prime to $p$.
\end{corollary}

\begin{proof}
A finite $Q$-module of order prime to $p$ admits a
composition series whose factors are finite-dimensional
$\mathbf F_\ell[Q]$-modules for primes $\ell\ne p$.
For a short exact sequence $0\to A'\to A\to A''\to0$,
the cohomology exact sequence gives
\[
                 |H^2(Q,A)|
                   \le |H^2(Q,A')|\,|H^2(Q,A'')|.
\]
Apply (D.H2.1) to the factors and multiply these
inequalities.  Their cardinalities multiply to $|A|$,
so (D.H2.8) follows.

The two central kernels are respectively the cyclic
scalar groups of order $d$ and $q-1$, with trivial
$Q$-action.  Pulling back the corresponding central
extension defines its continuous degree-two
obstruction class, which vanishes exactly when a
continuous lift exists.  This is the cocycle
classification of extensions
\citep[II, Example~1.18(b), pp.~65--66, and
Proposition~4.2, p.~87]{MilneCFT2020D}.
Thus (D.H2.8) gives the two stated cardinalities.
These count obstruction classes, not representations
within one obstruction class.

For fixed $B$, there are finitely many relevant
coefficient groups and finitely many classes for
each by (D.H2.8).  Every class is represented by a
normalized continuous cocycle factoring through
one finite quotient of $Q$, by the continuous-cochain
statement of the same Proposition~4.2.  Its restriction
to the kernel of that quotient is zero.  Intersect
these finitely many open normal kernels to obtain
$H_B$.
\end{proof}

\begin{remark}[The remaining projective-representation problem]
For unbounded central order, the number of obstruction classes
is now controlled, but a uniform count inside one nontrivial
class is a separate problem.  Applying the linear theorem to
an endomorphism algebra replaces rank $r$ by $r^2$ and yields
an exponent of order $r^4$.  A rank-$r$ twisted unit-bundle
presentation with constants uniform in the central class and
its order has not been supplied by this argument.  No bound
on arbitrary nonabelian crowns or full finite generation is
inferred from (D.H2.1)--(D.H2.8).
\end{remark}


% New research after the restored quadratic count; separate from frozen assembly.
\begin{proposition}[Proved under stated hypotheses: projective representations with bounded central obstruction]
\label{D:bounded-centre-projective-growth}
Let $Q$ be a profinite group and fix a prime $p$. Assume that
$H^2_{\mathrm{cont}}(Q,\mathbf Z/d\mathbf Z)$ is finite for every
positive integer $d$ prime to $p$, with trivial action. Assume also
that every open normal subgroup $H\subset Q$ has a constant $c_H$
such that
\[
 \#\bigl(\operatorname{Hom}_{\mathrm{cont}}
 (H,\operatorname{GL}_r(\mathbf F_{p^f}))/\mathrm{conj}\bigr)
 \le p^{c_Hfr^2}\qquad(r,f\ge1).
\]
For every fixed integer $B\ge1$ there is a constant $C_B$ with the
following property. Put $q=p^f$ and $d=\gcd(r,q-1)$. If $d\le B$,
then
\[
 \#\bigl(\operatorname{Hom}_{\mathrm{cont}}
 (Q,\operatorname{PSL}_r(\mathbf F_q))/
       \operatorname{PGL}_r(\mathbf F_q)\bigr)
 \le q^{C_Br^2}.
\]
Conjugation in the displayed quotient is through the natural normal
subgroup $\operatorname{PSL}_r(\mathbf F_q)\subset
\operatorname{PGL}_r(\mathbf F_q)$. In particular the number of
distinct kernels of epimorphisms to this simple group is bounded by
the same expression whenever the group is nonabelian simple.
\end{proposition}
\begin{proof}
We first construct one subgroup independently of $r$ and $f$. A
continuous normalized two-cocycle from $Q\times Q$ to a finite
discrete group is constant on the pairs of cosets of some open
normal subgroup $N$. Indeed compactness reduces its clopen fibres
to finitely many rectangles, and a common open normal subgroup
refines all their factors. Thus the cocycle factors through
$(Q/N)^2$. Its restriction to $N\times N$ is its value at $(1,1)$,
which is zero. Every continuous degree-two class with finite
coefficients is therefore killed by restriction to an open normal
subgroup; this is also the continuous-cochain calculation in
\citet[II, Proposition 4.2, p.~87]{MilneCFT2020D}.

There are only finitely many prime-to-$p$ integers $d\le B$, and by
hypothesis only finitely many classes for each. Intersect the open
normal subgroups just obtained. We obtain a single open normal
subgroup $H_B\subset Q$ such that
\[
 H^2(Q,\mathbf Z/d\mathbf Z)\longrightarrow
 H^2(H_B,\mathbf Z/d\mathbf Z)
 \quad\hbox{is zero for every such }d.
\]
We require neither an effective bound for its index nor one uniform
in $B$.

The central extension
\[
 1\longrightarrow\mu_d(\mathbf F_q)
 \longrightarrow\operatorname{SL}_r(\mathbf F_q)
 \longrightarrow\operatorname{PSL}_r(\mathbf F_q)
 \longrightarrow1
\]
has cyclic kernel of order $d$, which is prime to $p$. The pullback
of this extension by any continuous representation $\rho$ of $Q$
has a class in $H^2(Q,\mathbf Z/d\mathbf Z)$, after choosing an
identification of its centre with that cyclic group. The class
restricts to zero on $H_B$. A continuous cochain trivializing the
restricted cocycle supplies a continuous homomorphism
\[
 \widetilde\rho_H:H_B\longrightarrow
                   \operatorname{SL}_r(\mathbf F_q)
\]
lifting $\rho|_{H_B}$. This is a lift over the original finite
field; no scalar-field extension is involved.

Every conjugacy class of a restricted projective representation
therefore occurs as the projectivization of a linear representation
of $H_B$. Hence the number of possible restricted classes is at
most $q^{c_{H_B}r^2}$. Put $t=[Q:H_B]$ and choose $t$ coset
representatives, including $1$. After choosing one representative
of a restricted projective conjugacy class, any extension to $Q$
is determined by its values on these coset representatives. There
are at most $|\operatorname{PGL}_r(\mathbf F_q)|^{t-1}$ possible
tuples. Relations can only reduce their number. Conjugation needed
to fix the restriction causes no problem, since every element of
$\operatorname{PGL}_r(\mathbf F_q)$ is represented by a matrix in
$\operatorname{GL}_r(\mathbf F_q)$.

Since $|\operatorname{PGL}_r(\mathbf F_q)|\le q^{r^2}$, the total
is at most $q^{(c_{H_B}+t-1)r^2}$. Finally conjugate homomorphisms
have the same kernel, so the bound for kernels follows.
\end{proof}

\begin{corollary}[Proved: bounded-centre type-A simple quotients]
\label{D:surface-bounded-centre-simple-product}
Let $Q=Q(S,T)$ be the original-boundary tame group of a normal proper
surface over an algebraically closed field of characteristic $p>0$.
Fix $B$. The quotient obtained by taking the product
of all distinct nonabelian simple quotients
$\operatorname{PSL}_r(\mathbf F_{p^f})$ with
$\gcd(r,p^f-1)\le B$ is topologically finitely generated. The ranks
$r$ are not bounded in this assertion.
\end{corollary}
\begin{proof}
Theorem~\ref{D:surface-tame-cross-characteristic-H2} implies the
finiteness assumption of
Proposition~\ref{D:bounded-centre-projective-growth}. An open normal
subgroup $H\subset Q$ corresponds to a finite etale cover $U_H\to U$
of $U=S\setminus T$. Normalize $S$ in this cover to obtain a normal
proper surface $S_H$, and put $T_H=S_H\setminus U_H$. At each
original boundary divisor the cover is tame. The intersection
property for wild inertia shows that every finite representation
of $H$ is tame at the original divisors of $T_H$: its monodromy is
inflated from $H$, and the new wild inertia maps to the old killed
wild inertia. Equivalently there is a surjection
$Q(S_H,T_H)\twoheadrightarrow H$, which is all we need.
Theorem~\ref{D:tame-global-quadratic-representation-bound} therefore
supplies the hereditary quadratic bound used in the proposition.

For these simple groups the order formula
\[
 |\operatorname{PSL}_r(\mathbf F_q)|
 =\frac{q^{r^2-1}}{\gcd(r,q-1)}
       \prod_{i=2}^r(1-q^{-i})
\]
gives $|\operatorname{PSL}_r(\mathbf F_q)|\ge q^{r^2/4}$ for
$r\ge2$. Indeed the product is at least
$1-\sum_{i\ge2}2^{-i}=1/2$, while $2r\le2^r\le q^r$ and
$r^2-r-1\ge r^2/4$ for $r\ge2$. Thus Proposition~\ref{D:bounded-centre-projective-growth}
bounds the multiplicity of every such simple quotient by a fixed
power of its order.

The standard generating-tuple estimate for a nonabelian simple
group, used already in Proposition~\ref{D:bounded-matrix-simple-product},
then gives one number of generators for all the corresponding finite
products. This is the estimate in the proof of Proposition~2.6(ii),
printed pp.~12--13, of \citet{LucchiniMarionTracey2020D}.
Distinct simple quotient kernels give independent factors by
Goursat's lemma, so the resulting profinite quotient is their full
product. Compactness of the finite generating-tuple conditions
gives topological finite generation of that product.
\end{proof}

\begin{remark}[Scope of the remaining central gate]
The fixed subgroup $H_B$ depends on $B$. This argument supplies no
uniform bound as $\gcd(r,q-1)$ grows. In particular it does not
prove an all-rank projective representation bound with one constant
for every type-A simple group. A uniform count for twisted unit
bundles, or another controlled central-lifting argument, is still
needed for that stronger conclusion. Nor does a bound on simple
direct-product quotients by itself handle general nonabelian crowns.
\end{remark}


% New finite-proper descent consequence; awaiting independent actual-file audits.
\begin{corollary}[Proved: bounded-centre projective growth after finite normalization]
\label{D:nonnormal-bounded-centre-projective-growth}
Let $S$ be a proper scheme of dimension at most two over an algebraically
closed field of characteristic $p>0$. Let $T$ be a finite closed subset
of codimension at least two in every irreducible component it meets,
and let $U$ be a connected component of $S\setminus T$. Fix an integer
$B\ge1$. There is a constant $C_{S,T,B}$ such that, for $q=p^f$ and
$r\ge2$ with $\gcd(r,q-1)\le B$,
\[
 \#\bigl(\operatorname{Hom}_{\mathrm{cont}}
 (\pi_1(U),\operatorname{PSL}_r(\mathbf F_q)) /
 \operatorname{PGL}_r(\mathbf F_q)\bigr)
 \le q^{C_{S,T,B}r^2}.
\]
Consequently the quotient formed from all distinct nonabelian simple
quotients of these types is topologically finitely generated. The
ranks are unrestricted, while the natural central order is bounded
by the fixed $B$.
\end{corollary}
\begin{proof}
Reduction does not change finite etale covers. Normalize the reduced
proper scheme. Its normalization is finite, and its inverse image
$U'\to U$ is a finite surjection with finitely many connected
components $U_i$. Each two-dimensional component is a normal proper
surface minus finitely many points. Proposition~\ref{D:bounded-centre-projective-growth}, applied by
the same finite-cover argument as
Corollary~\ref{D:surface-bounded-centre-simple-product}, gives a bound
$q^{c_{i,B}r^2}$ for its projective conjugacy classes. The components
of dimension at most one are proper: the codimension hypothesis
prevents $T$ from meeting those components. Their fundamental groups
are finitely generated, and hence their representations into any
finite group $P$ number at most $|P|^{d_i}$ for a fixed $d_i$.
We apply this with $P=\operatorname{PGL}_r(\mathbf F_q)$ and use
$|P|\le q^{r^2}$.

We give the nonabelian descent count explicitly. A homomorphism to
$\operatorname{PSL}_r(\mathbf F_q)$ determines a torsor for the
constant finite group $P$ by extension of the structure group.
Two such homomorphisms determine isomorphic $P$-torsors precisely
when they are conjugate by $P$. Thus the classes which we wish to
count inject into the set of $P$-torsor classes on $U$. Their
restrictions to $U'$ range over at most
$q^{(\sum_i c_{i,B})r^2}$ tuples, after absorbing the curve and
point bounds. This use of $P$-torsors avoids replacing projective
conjugacy by conjugacy inside $\operatorname{PSL}_r$.

Fix one such tuple of $P$-torsors. The two pullbacks to
$U'\times_U U'$ have a finite etale Isom sheaf whose geometric
fibres, when nonempty, have cardinality $|P|$. A section on a
connected component is determined by its value at one geometric
point: the equalizer of two sections of a finite etale separated
sheaf is open and closed. If $c$ is the fixed number of connected
components of $U'\times_U U'$, there are at most $|P|^c$ possible
descent isomorphisms. The cocycle condition on the triple product
can only reduce this number. Finite-surjective proper descent of
finite etale schemes, together with full faithfulness for the
action morphisms, gives effective descent of these $P$-torsors
\citep[Expose~IX, Theorem~4.7]{C_SGA1}.

We may count all descended $P$-torsors, even if their monodromy
does not lie in $\operatorname{PSL}_r$; this only increases the
bound. Thus the gluing contributes at most $q^{cr^2}$ and proves
the displayed estimate. Distinct epimorphism kernels are separated
by projective conjugacy classes. The order estimate and the
generating-tuple argument in
Corollary~\ref{D:surface-bounded-centre-simple-product} now give
the stated finite generation of the simple-product quotient.
No degree-two cohomology assertion for the nonnormal group is
used in this proof.
\end{proof}


% Root, B and D: full actual-file and exact primary-source audits passed for all three statements.
\begin{theorem}[Proved: constant second cohomology of a complete local surface puncture]
\label{D:local-surface-constant-H2}
Let $(R,\mathfrak m,k)$ be a complete excellent normal local
domain of dimension two, with algebraically closed residue field.
Put $U=\operatorname{Spec}R-\{\mathfrak m\}$ and
$G=\pi_1(U)$. There is an integer $b_R$ such that, for every
positive integer $d$ invertible in $R$,
\[
 |H^2_{\mathrm{cont}}(G,\mathbf Z/d\mathbf Z)|
 \le |H^2_{\mathrm{et}}(U,\mathbf Z/d\mathbf Z)|
 \le d^{b_R}.                                         \tag{D.LH2.1}
\]
The group action on the displayed coefficients is trivial.
In particular these groups are finite. The same bound, with
$d^{b_R}$ replaced by $|A|^{b_R}$, holds for every finite
constant abelian group $A$ whose order is invertible in $R$.
No assertion about arbitrary nonconstant coefficient systems
is needed here.
\end{theorem}
\begin{proof}
The puncture $U$ is regular. By resolution of excellent normal
surfaces followed by embedded resolution, choose a proper
birational morphism
\[
             \pi:Y\longrightarrow\operatorname{Spec}R
\]
which is an isomorphism over $U$, with $Y$ regular and
$E=(\pi^{-1}\{\mathfrak m\})_{\mathrm{red}}$ a strict normal
crossings divisor. If necessary, a blowup of the closed point
first makes the exceptional fibre nonempty. The precise
resolution inputs are Theorem~54.14.5 and Lemma~54.15.6
\citep[Tags 0BGP and 0BIC]{StacksDSurfaceResolution}.
The components $E_1,\ldots,E_v$ are smooth proper curves over
$k$. Let $g_i$ be their genera, and let $e$ be the number of
their intersection points, counting distinct intersection
points separately. There are no triple intersections.

Set $\Lambda=\mathbf Z/d\mathbf Z$ and let $j:U\hookrightarrow Y$.
Gabber's absolute purity theorem, in its regular normal-crossings
form, gives
\[
 \begin{split}
 R^0j_*\Lambda&=\Lambda_Y,\\
 R^1j_*\Lambda&=\bigoplus_i\Lambda_{E_i}(-1),\\
 R^2j_*\Lambda&=\bigoplus_{x\in\operatorname{Nodes}(E)}
                                      \Lambda_x(-2),\\
 R^qj_*\Lambda&=0\quad(q>2).
 \end{split}                                         \tag{D.LH2.2}
\]
Here a sheaf on a closed stratum denotes its direct image to
$Y$. We use exactly Expos\'e~XVI, Theorem~3.1.1 and
Corollary~3.1.3 of \citet{A_ILO2014}; the complete statements
and the proof of the normal-crossings corollary appear on
pp.~242--243 of the preserved arXiv version. That corollary
applies to regular $\mathbf Z[1/d]$-schemes. It does not
require finite type over a field. In particular its hypotheses
apply to this excellent resolution. This is a use of purity
on the regular pair $(Y,E)$, not a constructibility assertion
for an arbitrary open immersion into the singular local scheme.

The base is strictly henselian. Global sections on its small
etale topos are exact, and proper base change consequently gives
\[
                   H^a(Y,\Lambda)=H^a(E,\Lambda).
\]
For the proper base-change input we use
\citep[Theorems~17.7 and 17.10]{MilneLEC2013H2D}.
Nilpotents in the closed fibre do not change its etale topos.
Normalization of the nodal curve $E$ and the node exact
sequence give $|H^2(E,\Lambda)|=d^v$; for each smooth
component $|H^1(E_i,\Lambda)|=d^{2g_i}$.
Tate twists are constant here, since the algebraically closed
residue field contains all $d$th roots of unity and Hensel's
lemma lifts them to $R$.

Apply the Leray spectral sequence
\[
 H^a(Y,R^qj_*\Lambda)\Longrightarrow H^{a+q}(U,\Lambda).
\]
The three terms of total degree two in its second page have
orders $d^v$, $d^{2\sum_i g_i}$ and $d^e$, respectively,
by (D.LH2.2). The corresponding limiting terms are subquotients
of these finite groups. Thus
\[
                |H^2(U,\Lambda)|
                    \le d^{v+2\sum_i g_i+e}.
\]
We may take $b_R=v+2\sum_i g_i+e$. Equivalently, one can obtain
finiteness from the localization triangle and the finite
Mayer--Vietoris complex of supports in the $E_i$, using purity
for the divisors and their intersection points.

Finally the usual degree-two comparison is injective:
\[
                H^2(G,\Lambda)\lhook\joinrel\longrightarrow
                                H^2_{\mathrm{et}}(U,\Lambda).
\]
Indeed the Cartan--Leray spectral sequence for all connected
finite Galois covers has degree-zero geometric term $\Lambda$,
and degree-one geometric term zero: every finite-coefficient
degree-one class is a finite torsor and is killed by a further
finite etale cover. The only possible incoming differential
to $(2,0)$ is therefore zero, and this limiting term is the
last filtration piece in total degree two. This is the same
comparison used in
Theorem~\ref{D:surface-tame-cross-characteristic-H2}; see
\citep[Theorem~14.9, Remark~14.10 and Proposition~14.13]
{MilneLEC2013H2D}. Decomposing a finite constant abelian
group into cyclic groups proves the final assertion.
\end{proof}

\begin{corollary}[Proved: one subgroup kills all bounded local central classes]
\label{D:local-bounded-central-obstructions}
In the preceding setting, for every fixed integer $B$ there
is an open normal subgroup $G_B\subset G$ such that
\[
 H^2(G,\mathbf Z/d\mathbf Z)\longrightarrow
 H^2(G_B,\mathbf Z/d\mathbf Z)
\]
is zero for every $d\le B$ invertible in $R$.
The subgroup is independent of the representation rank and
coefficient extension degree; no bound uniform in $B$ is asserted.
\end{corollary}
\begin{proof}
There are finitely many relevant $d$ and, by the theorem,
finitely many classes for each. Every normalized continuous
two-cocycle with finite coefficients factors through one
finite quotient of $G$: compactness refines its finitely
many clopen fibres by rectangles of cosets of a common open
normal subgroup. Its restriction to that subgroup is the
zero cocycle. Intersect these finitely many kernels.
Equivalently, use the continuous-cochain comparison
\citep[II, Proposition~4.2]{MilneCFT2020D}.
\end{proof}


\begin{corollary}[Proved: prime-to-characteristic cohomology on arbitrary-dimensional strict-local opens]
\label{D:strict-local-open-cohomology-finite}
Let $R$ be a strictly henselian quasi-excellent Noetherian local
ring, let $j:U\hookrightarrow\operatorname{Spec}R$ be an open
immersion, and let $d$ be invertible in $R$. For every constructible
sheaf $\mathcal F$ of $\mathbf Z/d\mathbf Z$-modules on $U$,
all groups $H^q(U,\mathcal F)$ are finite. In particular, if $U$
is nonempty and connected, then $H^2(\pi_1U,\mathbf Z/d\mathbf Z)$ is finite.
For each fixed $B$ one open normal subgroup of $\pi_1U$ kills
all constant central classes with $d\le B$ invertible in $R$.
These assertions impose no dimension-two hypothesis. They do
not assert a uniform exponent in $d$ analogous to (D.LH2.1).
\end{corollary}
\begin{proof}
Since $R$ is Noetherian, its open subscheme $U$ is quasi-compact,
so $j$ is of finite type. Expos\'e~XIII, Theorem~1.1.1(i) of
\citet{A_ILO2014} says exactly that, for a finite-type morphism
into a quasi-excellent Noetherian scheme and constructible
$\mathbf Z/d\mathbf Z$ coefficients with $d$ invertible on the
base, every higher direct image is constructible. Apply this
to $j$. In particular each $R^qj_*\mathcal F$ has a finite
stalk at the closed geometric point.

For a strictly local scheme, taking global sections on the
small etale topos is the closed geometric-stalk functor, hence
is exact. Indeed every etale neighbourhood of that point
has a section through it, so both equality of sections and
existence of local lifts can be tested after pulling back
along such a section. The Leray spectral sequence consequently
identifies
\[
 H^q(U,\mathcal F)
   =\Gamma(\operatorname{Spec}R,R^qj_*\mathcal F)
   =(R^qj_*\mathcal F)_{\overline{\mathfrak m}},
\]
which is finite. The same Cartan--Leray degree-two injection
as above gives the fundamental-group assertion. Finally apply
the finite-cocycle argument of
Corollary~\ref{D:local-bounded-central-obstructions} to the
finitely many groups and classes of bounded order.
\end{proof}


% New post-checkpoint-13 scalar-cover extension. Independent audit requested.
\begin{proposition}[Proved under stated hypotheses: finite matrix groups with bounded scalar kernels]
\label{D:bounded-scalar-projective-growth}
Let $Q$ and $p$ satisfy the two hypotheses of
Proposition~\ref{D:bounded-centre-projective-growth}: constant
prime-to-$p$ second cohomology is finite, and every open normal
subgroup has a quadratic finite-field representation bound.
Fix $B\ge1$. For $q=p^f$, let
\[
 \widetilde S\subset\operatorname{GL}_r(\mathbf F_q),\qquad
 S=\operatorname{im}\bigl(\widetilde S\longrightarrow
                         \operatorname{PGL}_r(\mathbf F_q)\bigr).
\]
Suppose that the scalar kernel of $\widetilde S\to S$ has order at
most $B$. Then, with a constant $C_B$ independent of $r,f$ and
$\widetilde S$,
\[
 \#\bigl\{\operatorname{Hom}_{\rm cont}(Q,S)
             \text{ up to conjugacy in }
                    \operatorname{PGL}_r(\mathbf F_q)\bigr\}
       \le q^{C_Br^2}.                                  \tag{D.SC.1}
\]
We mean the restriction of the ambient conjugacy equivalence
relation to these homomorphisms. We do not require $S$ to be normal
in $\operatorname{PGL}_r(\mathbf F_q)$.

Consequently, fix also $\epsilon>0$ and consider any family of
nonabelian finite simple groups $S$ with such realizations satisfying
$|S|\ge q^{\epsilon r^2}$. The quotient of $Q$ given by all distinct
simple quotient kernels of these types is topologically finitely
generated.
\end{proposition}
\begin{proof}
The scalar kernel is a subgroup of the cyclic group $\mathbf F_q^\times$.
It is therefore central and cyclic of prime-to-$p$ order $d\le B$.
The proof of Proposition~\ref{D:bounded-centre-projective-growth}
constructs one open normal subgroup $H_B\subset Q$ on which every
constant degree-two class with cyclic coefficients of such an order
vanishes. The pullback of
\[
 1\longrightarrow\mathbf Z/d\mathbf Z\longrightarrow
        \widetilde S\longrightarrow S\longrightarrow1
\]
by a homomorphism of $Q$ thus splits continuously over $H_B$.
Every restricted projective representation lifts to a linear
representation of $H_B$ over the original field $\mathbf F_q$.
Its ambient projective conjugacy class is consequently among at
most $q^{c_{H_B}r^2}$ possibilities.

Choose a representative of each such ambient restricted class.
Conjugating a homomorphism of $Q$ so that its restriction equals
this representative may move its full image out of the originally
specified subgroup $S$. This causes no loss in an upper bound:
we count all its possible extensions as homomorphisms into the
ambient projective linear group. If $t=[Q:H_B]$, their values on
$t-1$ fixed nonidentity coset representatives give at most
$|\operatorname{PGL}_r(\mathbf F_q)|^{t-1}$ choices. Thus (D.SC.1)
holds with $C_B=c_{H_B}+t-1$.

Ambient-conjugate epimorphisms have the same kernel, so the number
of distinct kernels for a simple group in the family is at most
$|S|^{C_B/\epsilon}$. The generating-tuple estimate used in
Proposition~\ref{D:bounded-matrix-simple-product} gives a uniform
number of generators for all finite products with these
multiplicities; its precise primary locator is the proof of
Proposition~2.6(ii), pp.~12--13, of
\citet{LucchiniMarionTracey2020D}. Goursat's lemma identifies the
image associated with distinct simple quotient kernels with the
full product on every finite subcollection. Compactness then gives
a generating tuple for the profinite product.
\end{proof}

\begin{corollary}[Proved: unitary quotients with bounded natural centre]
\label{D:bounded-centre-unitary-simple-product}
Let $Q$ satisfy the hypotheses of
Proposition~\ref{D:bounded-scalar-projective-growth}, and fix $B$.
For all $r\ge2$, $q=p^f$, and $\gcd(r,q+1)\le B$, the number of
homomorphisms
\[
           Q\longrightarrow\operatorname{PSU}_r(q)
\]
up to conjugacy in $\operatorname{PGL}_r(\mathbf F_{q^2})$ is
at most $q^{C'_Br^2}$. The quotient obtained by taking all distinct
nonabelian simple quotients of these unitary types is
topologically finitely generated, with no bound on $r$.

In particular these assertions hold for the original-boundary tame
group of a normal proper surface in characteristic $p$, and for
the fundamental group of a proper scheme of dimension at most two
after removing a finite set of codimension at least two in every
irreducible component it meets.
\end{corollary}
\begin{proof}
The natural special unitary group is a subgroup of
$\operatorname{GL}_r(\mathbf F_{q^2})$. Its scalar kernel in the
projective linear group has order $d=\gcd(r,q+1)$: a scalar
$\lambda$ preserves the hermitian form precisely when
$\lambda^{q+1}=1$, and has determinant one precisely when
$\lambda^r=1$. Its projective image is $\operatorname{PSU}_r(q)$.
These statements, and the order formula
\[
 |\operatorname{PSU}_r(q)|
   =\frac{q^{r^2-1}}{d}
          \prod_{i=2}^r\bigl(1-(-1)^i q^{-i}\bigr),      \tag{D.SC.2}
\]
are proved in \citet[Chapter~10, p.~118]{Taylor1992ClassicalB}.
The same source, Theorem~10.20, pp.~127--128, proves simplicity
except for $(r,q)=(2,2),(2,3),(3,2)$; we include only the
nonabelian simple members in the product assertion.

Applying the proposition over the field of size $q^2$ gives
the representation bound $q^{2C_Br^2}$. For completeness its
order comparison is uniform in rank. Odd-index factors in
(D.SC.2) are at least one, and
\[
 \prod_{\substack{2\le i\le r\\i\text{ even}}}(1-q^{-i})
       \ge 1-\sum_{j\ge1}2^{-2j}=\frac23.
\]
Since $d\le r$, $2r\le2^r\le q^r$, and
$r^2-r-1\ge r^2/4$ for $r\ge2$, we obtain
\[
           |\operatorname{PSU}_r(q)|\ge q^{r^2/4}
                         =(q^2)^{r^2/8}.
\]
The simple-product assertion follows with $\epsilon=1/8$.

For a normal proper surface with an original tame boundary,
Theorem~\ref{D:surface-tame-cross-characteristic-H2} supplies
constant second-cohomology finiteness. Each open normal subgroup
is a quotient of the original-boundary tame group of its finite
normal proper covering, exactly as in
Corollary~\ref{D:surface-bounded-centre-simple-product}.
Theorem~\ref{D:tame-global-quadratic-representation-bound}
supplies the required hereditary quadratic counts.

For a possibly nonnormal proper scheme of dimension at most two,
we repeat the finite proper normalization descent of
Corollary~\ref{D:nonnormal-bounded-centre-projective-growth},
now for torsors under $\operatorname{PGL}_r(\mathbf F_{q^2})$.
On normal two-dimensional components the just-proved bound
applies. Lower-dimensional components are proper, so their fixed
finite generator bounds give $q^{O(r^2)}$ classes. For any fixed
tuple of restricted torsors, descent isomorphisms on the fixed
double product have at most
$|\operatorname{PGL}_r(\mathbf F_{q^2})|^c\le q^{2cr^2}$
possibilities, where $c$ is its number of connected components.
Triple-overlap relations only reduce this number. Effective
finite proper descent \citep[IX, Theorem~4.7]{C_SGA1} and the
same order comparison finish the proof. No second-cohomology
assertion for the nonnormal total scheme is used.
\end{proof}

\begin{remark}[Exact central scope]
The constants may depend on the fixed bound $B$. The unitary
statement does not cover an unbounded sequence of values
$\gcd(r,p^f+1)$. The scalar-cover proposition also does not count
arbitrary monolithic nonabelian crowns; its product conclusion
concerns simple quotients only.
\end{remark}


% Direct all-dimensional local application; no higher-dimensional resolution is used.
\begin{corollary}[Proved: bounded scalar obstructions on complete local punctures]
\label{D:local-bounded-scalar-simple-product}
Let $(R,\mathfrak m,k)$ be a complete excellent normal local
domain of dimension at least two, equicharacteristic $p>0$, with
algebraically closed residue field. Put
\[
          U=\operatorname{Spec}R\setminus\{\mathfrak m\},
          \qquad G=\pi_1(U).
\]
Fix $B\ge1$. For every $q=p^f$ and every finite subgroup
$\widetilde S\subset\operatorname{GL}_r(\mathbf F_q)$ whose
scalar kernel has order at most $B$, let $S$ be its image in
$\operatorname{PGL}_r(\mathbf F_q)$. There is a constant
$C_{R,B}$ such that
\[
 \#\{\operatorname{Hom}_{\mathrm{cont}}(G,S)
            \text{ up to ambient projective conjugacy}\}
                  \le q^{C_{R,B}r^2}.                  \tag{D.LSC.1}
\]
If a family of nonabelian simple groups admits such realizations
with $|S|\ge q^{\epsilon r^2}$ for one fixed $\epsilon>0$, the
quotient of $G$ formed by all distinct simple quotient kernels
of these types is topologically finitely generated.

In particular this holds for all the nonabelian simple groups
$\operatorname{PSL}_r(p^f)$ with $\gcd(r,p^f-1)\le B$, and
for all $\operatorname{PSU}_r(p^f)$ with
$\gcd(r,p^f+1)\le B$, with unbounded $r$ in either assertion.
The unitary representation count is understood in
$\operatorname{PGL}_r(\mathbf F_{p^{2f}})$.
\end{corollary}
\begin{proof}
Completeness makes $R$ henselian, and the residue field is
algebraically closed, so $R$ is strictly henselian.
Corollary~\ref{D:strict-local-open-cohomology-finite}, which
uses precisely the constructibility theorem
\citep[XIII, Theorem~1.1.1(i)]{A_ILO2014}, proves that
$H^2(G,\mathbf Z/d\mathbf Z)$ is finite for every $d$ prime
to $p$. We use only finiteness for each $d$, not a bound
uniform in $d$.

We verify the hereditary representation hypothesis carefully.
Let $H\subset G$ be open normal, and let $V\to U$ be the
corresponding connected finite etale cover. The normal connected
scheme $U$ is integral, as is $V$. Normalize $R$ in $k(V)$,
and write the resulting finite normal domain as $R_H$.
Finiteness follows from excellence. Since a finite algebra
over a complete henselian local ring is a product of complete
local rings and $R_H$ is a domain, $R_H$ is local and complete.
Its residue field is a finite extension of $k$, hence is $k$.
It is excellent, normal, equicharacteristic $p$, and has the
same dimension as $R$. Its restriction above $U$ is exactly
$V$, by uniqueness of finite normalization, while its closed
fibre has one underlying point. Consequently
\[
       V=\operatorname{Spec}R_H\setminus\{\mathfrak m_H\},
       \qquad H=\pi_1(V).
\]
Corollary~\ref{D:uniform-local-growth-all-dimensions} therefore
supplies a constant $c_H$ with
\[
 \#\bigl(\operatorname{Hom}_{\mathrm{cont}}
       (H,\operatorname{GL}_r(\mathbf F_{p^f}))/\mathrm{conj}\bigr)
                     \le p^{c_Hfr^2}.
\]
This all-dimensional quadratic theorem uses the finite surface
slice and proper normalization descent of
Proposition~\ref{D:local-surface-reduction}; it requires no
resolution in the dimension of $R$.

Thus both hypotheses of
Proposition~\ref{D:bounded-scalar-projective-growth} hold for
$G$, proving (D.LSC.1) and the simple-product assertion.
The indicated linear and unitary families follow from
Corollaries~\ref{D:surface-bounded-centre-simple-product} and
\ref{D:bounded-centre-unitary-simple-product}'s elementary
matrix-cover and group-order calculations; their geometric
surface hypotheses are not being invoked here.
\end{proof}

\begin{remark}[Local scope]
The prime-to-characteristic cohomology input itself holds on
arbitrary-dimensional strictly local quasi-excellent opens.
The quantitative conclusion here uses the separate quadratic
representation theorem for complete equicharacteristic normal
punctures with algebraically closed residue field. We do not
extend that quantitative theorem to mixed characteristic or
to an arbitrary open subscheme of a local scheme.
\end{remark}


% New post-checkpoint-13 consequence. Full root/B/C/D proof audits passed;
% fresh complete Taylor primary passage audits by root and D.
\begin{corollary}[Proved: symplectic and orthogonal simple quotients]
\label{D:symplectic-orthogonal-simple-product}
Let $Q$ and $p$ satisfy the hypotheses of
Proposition~\ref{D:bounded-scalar-projective-growth}.
For $q=p^f$, consider the following nonabelian simple groups, with
$r$ denoting the dimension of the indicated natural matrix representation:
\[
\begin{array}{c|c|l}
 S&r&\text{range}\\ \hline
 \operatorname{PSp}_{2m}(q)&2m&m\ge1,
       \ (m,q)\notin\{(1,2),(1,3),(2,2)\},\\
 \mathrm{P}\Omega_{2m+1}(q)&2m+1&q\text{ odd},\ m\ge1,
       \ (m,q)\ne(1,3),\\
 \mathrm{P}\Omega^{+}_{2m}(q)&2m&m\ge3,\\
 \mathrm{P}\Omega^{-}_{2m}(q)&2m&m\ge2.
\end{array}                                                   \tag{D.SO.1}
\]
There is a constant $C_Q$, independent of $m,f$ and the row, such that
\[
 \#\{\operatorname{Hom}_{\rm cont}(Q,S)
          \text{ modulo ambient }\operatorname{PGL}_r(\mathbf F_q)
          \text{-conjugacy}\}\le q^{C_Qr^2}.                  \tag{D.SO.2}
\]
The quotient of $Q$ defined by all distinct simple quotient kernels
of these types is topologically finitely generated. There is no
bound on their Lie ranks.

In particular this conclusion applies to the original-boundary tame
group of a normal proper surface over an algebraically closed field
of characteristic $p$. It also applies to the fundamental group of
a proper scheme of dimension at most two after removing a finite
set of codimension at least two in every irreducible component it meets.
It also applies to the puncture of a complete excellent normal local
domain of dimension at least two in equicharacteristic $p$, with
algebraically closed residue field.
\end{corollary}

\begin{proof}
We first verify the precise finite-group inputs. The natural cover
\[
 \operatorname{Sp}_{2m}(q)\longrightarrow
                      \operatorname{PSp}_{2m}(q)
\]
is a quotient by its scalar matrices, which are precisely $I$ and
$-I$. Its kernel therefore has order $\gcd(2,q-1)$.
For an orthogonal geometry $(V,Q_V)$, we use
\[
 \Omega(V)=O(V)'\subset\operatorname{GL}(V),\qquad
 \mathrm{P}\Omega(V)
     =\Omega(V)/(\Omega(V)\cap\mathbf F_q^\times I).
\]
A scalar isometry has square one, since
$Q_V(\lambda v)=\lambda^2Q_V(v)$ and $Q_V$ is nonzero.
Consequently this scalar kernel also has order at most two,
and is trivial in characteristic two. These are the natural
$r$-dimensional matrix lifts. We use the definitions and scalar
calculations in \citet[pp.~69 and 136--137]{Taylor1992ClassicalB}.
Their kernels have prime-to-$p$ order, including when $p=2$.

Theorem~8.8, p.~74, of the same source gives exactly the three
symplectic exceptions in (D.SO.1). For quadratic forms with
nondegenerate polar form, Theorem~11.48, p.~162, gives simplicity
whenever the dimension is at least three and the Witt index is
positive, except for $\mathrm{P}\Omega_3(3)$ and the
four-dimensional plus type. All the orthogonal ranges displayed
in (D.SO.1) therefore consist of nonabelian simple groups.
The omitted low-dimensional groups cause no unmentioned family:
$\operatorname{PSp}_2(q)=\operatorname{PSL}_2(q)$, and
$\mathrm{P}\Omega^{-}_4(q)=\operatorname{PSL}_2(q^2)$.
In even characteristic the odd-dimensional orthogonal geometry
has a degenerate polar form; Theorem~11.9, pp.~143--144, identifies
$O_{2m+1}(q)$ with $\operatorname{Sp}_{2m}(q)$.
Whenever the latter group is perfect its derived group is already
the symplectic group in the first row. The only exceptional even
characteristic parameters are $(m,q)=(1,2),(2,2)$, so they introduce
no unbounded family. We do not apply Theorem~11.48 to those
degenerate polar forms.

For completeness we prove a uniform comparison with group order.
The exact formulas from pp.~70 and 165--166 of
\citet{Taylor1992ClassicalB} are
\[
\begin{split}
 |\operatorname{PSp}_{2m}(q)|
  &=\frac{q^{2m^2+m}}{\gcd(2,q-1)}
                   \prod_{i=1}^{m}(1-q^{-2i}),\\
 |\mathrm{P}\Omega_{2m+1}(q)|
  &=\frac{q^{2m^2+m}}2
                   \prod_{i=1}^{m}(1-q^{-2i})\quad(q\text{ odd}),\\
 |\mathrm{P}\Omega^{\epsilon}_{2m}(q)|
  &=\frac{q^{2m^2-m}}{\gcd(4,q^m-\epsilon)}
       (1-\epsilon q^{-m})
                   \prod_{i=1}^{m-1}(1-q^{-2i}),
       \qquad\epsilon\in\{1,-1\}.
\end{split}                                                   \tag{D.SO.3}
\]
For the orthogonal formulas we also use the spinor-index
calculation and its proof, Theorem~11.51, pp.~164--165.
The exceptional four-dimensional plus group over $\mathbf F_2$
is outside (D.SO.1).

For every $q\ge2$, the elementary inequality
$\prod(1-a_i)\ge1-\sum a_i$, for $0\le a_i\le1$, gives
\[
             \prod_{i=1}^{m}(1-q^{-2i})
                 \ge1-\sum_{i\ge1}2^{-2i}=\frac23.
\]
Thus the first order in (D.SO.3) is at least
$q^{2m^2+m}/3\ge q^{2m^2+m-2}\ge q^{(2m)^2/4}$.
For the second order, $q\ge3$, so it is at least
$q^{2m^2+m-1}\ge q^{(2m+1)^2/8}$ for every $m\ge1$.
For either even-dimensional sign and $m\ge2$,
\[
 \frac{(1-\epsilon q^{-m})}{\gcd(4,q^m-\epsilon)}
       \prod_{i=1}^{m-1}(1-q^{-2i})
       \ge\frac{(3/4)(2/3)}4=\frac18.
\]
Its order is therefore at least
$q^{2m^2-m-3}\ge q^{(2m)^2/8}$, since
$\tfrac32m^2-m-3\ge0$ for $m\ge2$.
In all four rows we have proved
\[
                         |S|\ge q^{r^2/8}.                  \tag{D.SO.4}
\]

Proposition~\ref{D:bounded-scalar-projective-growth}, with
$B=2$ and $\epsilon=1/8$, now proves both assertions for $Q$.
In particular the number of distinct quotient kernels of any
one simple isomorphism type is at most $|S|^{8C_Q}$.
The proposition's uniform generating-tuple argument then applies
to the full product over all these isomorphism types.

For the first geometric application, the second-cohomology
hypothesis follows from
Theorem~\ref{D:surface-tame-cross-characteristic-H2}; the
hereditary quadratic representation bounds follow by finite
normal proper covers from
Theorem~\ref{D:tame-global-quadratic-representation-bound}.
For a possibly nonnormal proper scheme of dimension at most two,
the same normalization descent used in
Corollary~\ref{D:nonnormal-bounded-centre-projective-growth}
applies to the ambient $\operatorname{PGL}_r(\mathbf F_q)$-torsors.
On each fixed connected component of the double fibre product,
a descent isomorphism is determined by one geometric value,
so the total additional number is at most
$|\operatorname{PGL}_r(\mathbf F_q)|^c\le q^{cr^2}$ for a fixed
$c$. Effective finite proper descent
\citep[IX, Theorem~4.7]{C_SGA1} completes this application.
It does not require a second-cohomology theorem for the
nonnormal total scheme.
For the last local application, both abstract hypotheses for its
fundamental group, including the hereditary quadratic bound, have
been verified in
Corollary~\ref{D:local-bounded-scalar-simple-product}.
\end{proof}

\begin{remark}[Scope of the family conclusion]
The scalar kernels of the natural symplectic and orthogonal
lifts are uniformly bounded, independently of their ranks.
The analogous natural unitary lift is treated in
Corollary~\ref{D:bounded-centre-unitary-simple-product}, with
its explicit bound on $\gcd(r,q+1)$. Within these
defining-characteristic classical families, the remaining unbounded
natural centres in linear and unitary groups are not covered
by this argument. No representation-count bound in a coefficient
characteristic different from $p$ is asserted. In particular this
does not settle cross-characteristic Lie families whose orders
are divisible by $p$, or unbounded alternating families.
These conclusions concern products of simple
quotients; they do not supply the required bounds for general
monolithic nonabelian crowns.
\end{remark}


% Auxiliary bounds, not a finite-generation claim or a new requested-row closure.
\begin{proposition}[Proved: arithmetic genus of the full exceptional fibre]
\label{D:exceptional-arithmetic-genus-bound}
Use the one-centre setup of
Propositions~\ref{D:surface-uniform-different} and
\ref{D:fixed-boundary-uniform-different}. Thus $Y\to(S,s)$ is a
fixed smooth projective resolution of a normal surface germ over
an algebraically closed field, $A=\sum_i a_iE_i$ is a fixed
positive exceptional cycle with $A\cdot E_i<0$, and $D$ is the
strict transform of the original reduced divisorial boundary.
Set $D=0$ in the punctured case.

Let a connected normal finite covering in the preceding
etale or tame setup have degree $n$. In particular its extension
of function fields is separable. Resolve its normalization over
$Y$ and contract only the smooth rational $(-1)$-curves mapped
to points of $Y$, as in the cited propositions. Write
\[
 f:Z\longrightarrow Y,\qquad C=f^*A,\qquad
 F=C_{\mathrm{red}},\qquad c=\#\pi_0(F),\qquad
 b=(K_Y+D)\cdot A+\sum_i a_i.
\]
Then $c\le n$ and
\[
             2h^1(F,\mathcal O_F)\le nb+c\le n(b+1).    \tag{D.AG.1}
\]
The possibly nonreduced proper Cartier divisor $C$ satisfies
\[
 2+n\left(A^2-\sum_i a_i\right)
 \le 2p_a(C)
 \le 2+n\left(A^2+(K_Y+D)\cdot A\right).               \tag{D.AG.2}
\]
If we further resolve the reduced exceptional fibre by point
blowups to obtain a strict normal crossings fibre $\widehat F$,
then
\[
 \sum_{G\subset\widehat F}g(G)+b_1(\Gamma_{\widehat F})
                   \le \frac{nb+c}{2},                 \tag{D.AG.3}
\]
where $\Gamma_{\widehat F}$ is its dual multigraph.
\end{proposition}
\begin{proof}
All curves in the support of $C$ are exceptional over the finite
set of source points above $s$. Conversely every such exceptional
curve belongs to this support, since $A$ is positive on all
components over $s$. Write $F_{i\lambda}$ for the curves
dominating $E_i$, and $Q_j$ for those mapped to points of $Y$.
The coefficient of $C$ along $F_{i\lambda}$ is
$a_i e_{i\lambda}$. Its coefficients along the $Q_j$ are
positive integers.

The construction and adjunction estimates in the proof of
Proposition~\ref{D:surface-uniform-different} give
\[
 K_Z\cdot Q_j\ge0,\qquad K_Z\cdot F_{i\lambda}\ge-1,
 \qquad K_Z\cdot C\le n(K_Y+D)\cdot A.                 \tag{D.AG.4}
\]
The last inequality uses the effective determinant-of-differentials
divisor. Its components over the $E_i$ have nonpositive
intersection with $C$, its point-mapped components have zero
intersection, and its horizontal part has intersection at most
$nD\cdot A$ by original-boundary tameness. This argument does
not require the cover to be Galois.

In passing from $C$ to its reduced divisor, lowering a coefficient
on a point-mapped component cannot increase the canonical
intersection. Lowering the coefficient $a_i e_{i\lambda}$ on a
dominating component increases it by at most
$a_i e_{i\lambda}-1$. Therefore
\[
 \begin{split}
 K_Z\cdot F
 &\le K_Z\cdot C+
                \sum_{i,\lambda}(a_i e_{i\lambda}-1)\\
 &\le n(K_Y+D)\cdot A+n\sum_i a_i=nb,
 \end{split}                                          \tag{D.AG.5}
\]
where $\sum_\lambda e_{i\lambda}\le n$ follows from
$\sum_\lambda e_{i\lambda}f_{i\lambda}=n$ with full residue
degrees.

The exceptional intersection matrix is negative definite
\citep[Theorem~10.1]{Kollar2013D}. Distinct connected components
of $F$ are disjoint. Each has a negative integral square,
so $F^2\le-c$. Since $F$ is proper and reduced over the
algebraically closed field, $h^0(\mathcal O_F)=c$.
Adjunction for a proper Cartier divisor on the smooth surface
therefore gives
\[
 2h^1(\mathcal O_F)
   =2c+F^2+K_Z\cdot F\le nb+c.
\]
To check $c\le n$, pass to the strict henselization at $s$.
The finite normalization has at most $n$ local normal factors:
their positive generic degrees add to $n$. The exceptional
fibre of the proper birational morphism over each such normal
local factor is connected, by its pushforward of the structure
sheaf and Stein factorization. This proves the component bound.

For (D.AG.2), projection gives $C^2=nA^2$. The lower canonical
bound from the same construction is
$K_Z\cdot C\ge-n\sum_i a_i$. Combining this and (D.AG.4) with
$2p_a(C)=2+C^2+K_Z\cdot C$ proves both inequalities.

Finally consider one blowup of a point of a reduced divisor of
multiplicity $m\ge1$. Its full reduced total transform is
\[
                 F'=\pi^*F-(m-1)E,
\]
where $E$ is the exceptional $(-1)$-curve. The equalities
$K_{Z'}=\pi^*K_Z+E$, $E^2=-1$ and
$E\cdot\pi^*F=0$ give
\[
             p_a(F')=p_a(F)-\frac{(m-1)(m-2)}2.
\]
The number of connected components is unchanged. Hence $h^1$
of the reduced full fibre does not increase under these blowups.
For a reduced strict normal crossings curve, its normalization
exact sequence gives
$h^1(\mathcal O_{\widehat F})=\sum_G g(G)+b_1(\Gamma_{\widehat F})$.
This proves (D.AG.3). Covers of the punctured strict henselization
are treated by the same finite-presentation descent used in the
cited different bounds.
\end{proof}

\begin{lemma}[Proved: Frobenius saturation of a finite normalization lattice]
\label{D:normalization-Frobenius-saturation}
Let $A$ be a complete discrete valuation ring of characteristic
$p>0$, with uniformizer $\pi$ and fraction field $K$. Let $B$
be its integral closure in a finite separable extension $L/K$.
Write $e$ for the ramification index and $\delta$ for the
different exponent. For every power $q$ of $p$,
\[
  \pi^{\lceil(q-1)\delta/e\rceil}B
                     \subset A\cdot B^q\subset B.      \tag{D.AG.6}
\]
No separability of the residue extension is assumed.

Consequently, let $\rho:Y'\to Y$ be finite, dominant and
generically separable between normal surfaces in characteristic
$p$, with $Y$ smooth. Put $\mathcal B=\rho_*\mathcal O_{Y'}$.
Suppose all ramification is supported on fixed prime divisors
$D_j$, and at every prolongation of $D_j$ we have
$\delta/e\le c_j$. Then $\mathcal B$ is locally free, and the
relative Frobenius map gives actual inclusions
\[
 \mathcal B(-H)\subset
       \operatorname{im}(F_Y^*\mathcal B\longrightarrow\mathcal B)
       \subset\mathcal B,
 \qquad H=\sum_j\lceil(p-1)c_j\rceil D_j.              \tag{D.AG.7}
\]
The divisor $H$ is independent of the degree of the covering.
\end{lemma}
\begin{proof}
Choose an $A$-basis $b_1,\ldots,b_n$ of $B$ and its trace-dual
$K$-basis $b_1^*,\ldots,b_n^*$ of $L$. The inverse different is
the trace-dual lattice, so $v_B(b_i^*)\ge-\delta$; this is the
trace construction of the different
\citep[Section~49.8]{StacksDDifferent}. Since
$\operatorname{Tr}(x^q)=\operatorname{Tr}(x)^q$ in characteristic
$p$, the bases $b_i^q$ and $(b_i^*)^q$ are also trace-dual.
For $b\in B$ we therefore have
\[
        b=\sum_i\operatorname{Tr}_{L/K}
                         \bigl(b(b_i^*)^q\bigr)b_i^q.
\]
Set $M=\lceil(q-1)\delta/e\rceil$. The element
$\pi^M b(b_i^*)^{q-1}$ is integral, because its valuation
is at least $eM-(q-1)\delta\ge0$. Pairing it with $b_i^*$
therefore has integral trace. Thus
\[
 \pi^M\operatorname{Tr}_{L/K}\bigl(b(b_i^*)^q\bigr)\in A
\]
for every $i$, proving (D.AG.6). The proof works componentwise
if the completed normalization is a product of DVRs, with
$M$ the largest corresponding ceiling.

A normal surface is Cohen--Macaulay. Miracle flatness over the
regular surface $Y$ therefore makes $\mathcal B$ locally free
\citep[Tag 00R4]{StacksArithmeticDefectB}. The map
$F_Y^*\mathcal B\to\mathcal B$, $a\otimes b\mapsto ab^p$,
is generically an isomorphism, by separability of the function
field extension. At every codimension-one point, completion
and (D.AG.6) show that its inverse carries
$\mathcal B(-H)$ into $F_Y^*\mathcal B$.
Both are vector bundles on the normal scheme $Y$. A rational
morphism between them which is regular at every codimension-one
point extends everywhere, by the intersection description of
reflexive sheaves \citep[Tag 0EBJ]{C_StacksPicNorm}. This gives the first actual sheaf inclusion
in (D.AG.7); the second is the original Frobenius map. In
particular we have not replaced its image by a reflexive hull
without justification.
\end{proof}

\begin{remark}[Limits of these two bounds]
The arithmetic bound retains dual-graph cycles and normalized
genera but permits arbitrarily long rational trees. Contracting
such trees can create normalization defects on curves in a
finite normal model, so (D.AG.1) does not remove the nilradical
term in the nonreduced-fibre calculation following
Lemma~\ref{D:proper-curve-arithmetic-defect}.
The Frobenius saturation bound concerns the actual canonical
normalization lattice. It does not by itself give uniform
Harder--Narasimhan windows on the exceptional components,
an atlas with size linear in the degree of the cover, or a
factorial bound for the number of finite covers. None of these
statements asserts finite generation of the full group.
\end{remark}


% Auxiliary degree estimate. It does not supply a hereditary object atlas.
\begin{proposition}[Proved: a degree-linear exceptional core]
\label{D:exceptional-degree-linear-core}
Use the fixed surface resolution, strictly anti-nef cycle
$A=\sum_{i=1}^{s}a_iE_i$, and strict original boundary $D$ of
Proposition~\ref{D:fixed-boundary-uniform-different}.  Put
\[
 a=\sum_i a_i,\qquad
 b=(K_Y+D)\cdot A+a,\qquad c=3b+a.
\]
Let a connected finite cover of the indicated open have degree
$n$, and assume it is tame at the original boundary divisors.
The cover need not be Galois.  Normalize $Y$ in its function
field, resolve, and contract all smooth rational $(-1)$-curves
which map to points of $Y$, as in the proof of that proposition.
Write $f:Z\to Y$ for the resulting morphism and
$C=(f^*A)_{\mathrm{red}}$.

Define the \emph{core} $\mathcal K$ to consist of every component
of $C$ dominating an $E_i$, together with every point-mapped
component $Q$ for which $K_Z\cdot Q>0$.  Then
\[
 \#\mathcal K\le n(b+a),\qquad
 \sum_{F\in\mathcal K}(-F^2)\le nc,\qquad
 \sum_{F\in\mathcal K}p_a(F)\le nb.                 \tag{D.GC.1}
\]
Every component outside the core is a smooth rational
$(-2)$-curve.  If $L$ is the sum, with intersection
multiplicity, of all incidences from a core component to a
noncore component, then
\[
                 L<2nc\quad\text{if }\mathcal K\ne\varnothing,
                 \qquad L=0\quad\text{otherwise}.       \tag{D.GC.2}
\]
The sum of intersection multiplicities over unordered pairs of
distinct core components is less than $nc/2$ when the core is nonempty.
There are at most $2nc+n$ connected noncore blocks.
The numbers of vertices within the negative-definite rational
$(-2)$-blocks remain unrestricted by these estimates.
\end{proposition}

\begin{proof}
Write the components dominating $E_i$ as $F_{i\lambda}$ and
the other components of $C$ as $Q_j$.  Their coefficients in
$f^*A$ are $m_{i\lambda}=a_i e_{i\lambda}$ and positive
integers $m_j$, respectively.  The positivity of $m_j$ follows
because $A$ is an effective Cartier divisor whose support
contains the image of $Q_j$.  The degree formula at each base
divisor gives
\[
 \sum_\lambda e_{i\lambda}f_{i\lambda}=n,
 \qquad
 \sum_{i,\lambda}m_{i\lambda}\le na.                    \tag{D.GC.3}
\]
Here the residue degrees are full residue degrees; no
separability of residual extensions is assumed.

The determinant calculation and horizontal tame correction in
Propositions~\ref{D:surface-uniform-different} and
\ref{D:fixed-boundary-uniform-different} give, before any use
of Galois conjugacy,
\[
                 K_Z\cdot f^*A\le n(K_Y+D)\cdot A.
\]
Adjunction and the chosen contractions give
$K_Z\cdot Q_j\ge0$ and $K_Z\cdot F_{i\lambda}\ge-1$.
Consequently the following sum has nonnegative summands:
\[
 \sum_{i,\lambda}m_{i\lambda}
                 (K_Z\cdot F_{i\lambda}+1)
       +\sum_jm_jK_Z\cdot Q_j\le nb.                  \tag{D.GC.4}
\]
In particular $b\ge0$.  There are at most $na$ dominating
components.  Each point-mapped core component contributes at
least one to the second sum in (D.GC.4), proving the first
bound in (D.GC.1).

All these curves are exceptional over points of the normal
cover of $S$, so their intersection matrix is negative
definite and each square is at most $-1$.  For a dominating
component $F$, adjunction gives
\[
 -F^2\le K_Z\cdot F+2,
 \qquad p_a(F)\le (K_Z\cdot F+1)/2.
\]
For a point-mapped core component $Q$, its positive integral
canonical degree gives
\[
 -Q^2\le K_Z\cdot Q+2\le3K_Z\cdot Q,
 \qquad p_a(Q)\le(K_Z\cdot Q+1)/2\le K_Z\cdot Q.
\]
Apply (D.GC.3)--(D.GC.4), using all $m\ge1$, to obtain the
remaining bounds in (D.GC.1).  A noncore point-mapped curve
satisfies
$0=2p_a(Q)-2-Q^2$.  Its square is negative; hence its arithmetic
genus is zero and its square is $-2$.  An integral proper
curve of arithmetic genus zero over an algebraically closed
field is a smooth rational curve.

We give the attachment estimate directly, without needing a
classification of the possible $(-2)$-graphs.  In the negative
intersection matrix, order the core first and the noncore
blocks second, and write
\[
                   M=\begin{pmatrix}M_{\mathcal K}&-T\\
                                      -T^{\mathsf t}&N
                      \end{pmatrix}.
\]
All entries of $T$ are nonnegative integers.  Each diagonal
entry of $N$ is $2$, its off-diagonal entries are nonpositive,
and $N$ is positive definite.  Writing $N=2I-P$ with $P$
symmetric and entrywise nonnegative, positive definiteness
gives $\rho(P)<2$.  Thus
\[
             N^{-1}=\frac12\sum_{j\ge0}(P/2)^j
\]
is entrywise nonnegative and has diagonal entries at least
$1/2$.  The Schur complement
$M_{\mathcal K}-TN^{-1}T^{\mathsf t}$ is positive definite.
For the row $t_F$ corresponding to a core component $F$,
\[
     -F^2>t_FN^{-1}t_F^{\mathsf t}
             \ge\tfrac12\sum_Q(F\cdot Q)^2
             \ge\tfrac12\sum_Q F\cdot Q.
\]
Summing proves (D.GC.2).  Applying positive definiteness of
$M_{\mathcal K}$ to its all-one vector gives the stated bound
on core--core incidences.

Every noncore block in a connected component meeting the
core has at least one attachment, so their number is at most
$L$.  A connected component not meeting the core is one
noncore block.  The whole exceptional fibre has at most $n$
connected components: the finite normal cover has at most
$n$ points over $s$, and proper birational maps to its normal
local germs have connected fibres.  This proves the last
bound.  The same argument applies after strict henselization,
where the connected normalization has only one closed point.
\end{proof}

\begin{remark}[Remaining presentation problem]
These are numerical and intersection-graph bounds.  We do
not assert that contracting the noncore blocks yields an
object atlas with $O(nr^2)$ parameters, or that the constant
in the quadratic representation bound grows linearly with
cover degree.  A uniform treatment of the full module and
Frobenius data along arbitrarily long $(-2)$-blocks remains
necessary for that argument.
\end{remark}


% New post-CP14 research fragment.  Not part of the frozen CP14 assembly.
\subsection{A superperfect-product obstruction to cohomological generation criteria}

\begin{example}[Proved: a superperfect-product obstruction]
\label{D:superperfect-quadratic-obstruction}
For every odd \(n\geq31\), let \(T_n\to A_n\) be the universal
central extension, and put
\[
 m_n=2^{n^2},\qquad
 \Pi=\prod_{\substack{n\geq31\\ n\ {\rm odd}}}T_n^{\,m_n}.
\]
Thus \(T_n\) is the double cover of \(A_n\).
The following assertions hold.
\begin{enumerate}
\item For every finite group \(H\), there are only finitely many continuous
homomorphisms \(\Pi\to H\).  Every open subgroup of \(\Pi\) has finite
maximal prosolvable quotient.
\item For every finite field \(k\), of any characteristic, and every
finite-dimensional continuous \(k[\Pi]\)-module \(M\), we have
\[
 \dim_k H^1(\Pi,M)\leq2\dim_k M,\qquad
 \dim_k H^2(\Pi,M)\leq5\dim_k M.
 \tag{SP1}
\]
For an open subgroup \(J\subset\Pi\), the corresponding constants
may be multiplied by \([\Pi:J]\).
\item There is one absolute constant \(C\), independent of \(q\), its
characteristic, and \(r\), such that
\[
 \#\left(
 \operatorname{Hom}_{\mathrm{cts}}(\Pi,\operatorname{GL}_r(\mathbf F_q))
 /\operatorname{GL}_r(\mathbf F_q)\right)\leq q^{Cr^2}.
 \tag{SP2}
\]
Nevertheless, \(\Pi\) is not topologically finitely generated.
\end{enumerate}
This is an abstract profinite example.  We assert no realization as
a geometric fundamental group.
\end{example}

\begin{proof}
We first specify the literature inputs and then prove the product
and counting assertions.  The Schur multiplier of \(A_n\) is \(C_2\)
for \(n\geq8\), and the universal central extension of a perfect group
has this multiplier as its kernel
\cite[Theorem~6.9.5 and Example~6.9.10, pp.~200--202]{Weibel1994SuperperfectB}.
We use this standard multiplier computation as a cited input.
The group \(T_n\) is perfect and has
\[
 H_1(T_n,\mathbf Z)=H_2(T_n,\mathbf Z)=0.
 \tag{SP3}
\]
For precision, the superperfectness assertion can be deduced from
universality as follows.  Let \(Y\to T_n\) be the universal central
extension of the perfect group \(T_n\); then \(Y\) is perfect.
For an element \(z\) in the kernel of \(Y\to A_n\), its commutators
with \(Y\) lie in the central kernel of \(Y\to T_n\).  Therefore
\(y\mapsto[z,y]\) is a homomorphism from the perfect group \(Y\)
to an abelian group and is zero.  Consequently \(Y\to A_n\) is
central.  The universality of \(T_n\to A_n\) splits \(Y\to T_n\);
the resulting central direct factor is trivial because \(Y\) is
perfect.  Thus \(Y=T_n\), proving (SP3).  This argument uses the
perfectness of \(Y\); a composite of arbitrary central extensions
need not be central.

The simplicity of \(A_n\) implies that every proper normal subgroup
of \(T_n\) lies in its central \(C_2\).  Indeed, if a normal subgroup
surjects onto \(A_n\), its product with the centre is \(T_n\), so
the corresponding quotient is abelian and hence trivial.
For odd \(n\), the elements
\((1\,2\,\cdots\,n)\) and \((1\,2\,3)\) generate \(A_n\):
their conjugates include the consecutive three-cycles, which
generate \(A_n\) by induction.  Any lifts of these two elements
generate \(T_n\).  Their generated subgroup \(L\) satisfies
\(LZ(T_n)=T_n\), hence is normal, and \(T_n/L\) is abelian;
perfectness gives \(L=T_n\).
In particular, for every \(T_n\)-module \(V\),
\[
 \dim H^1(T_n,V)\leq2\dim V,
 \tag{SP4}
\]
since a derivation is determined by its values on two generators.

\par\smallskip\noindent\textbf{The second-cohomology input.}\quad
We use the estimate
\[
 \dim H^2(A_n,V)<3\dim V
 \tag{SP5}
\]
from \cite[Theorem~6.2 and its proof in
Theorems~6.4, 6.5, 6.9, pp.~487--493]{GKKL2007SuperperfectB}.
The strict inequality can be weakened without affecting the argument.
The cited proof uses computational small-degree cases from its
reference~[42]; those computations are literature dependencies here,
not computations reproduced by our certificate.

We claim that \(\dim H^2(T_n,V)\leq5\dim V\) for all finite-dimensional
modules over every field.  In characteristic zero, positive-degree
cohomology of the finite group vanishes by averaging.
By coefficient extension and a composition
series, it suffices to consider an irreducible module over an
algebraically closed field \(k\).  The trivial module has \(H^2=0\)
by (SP3) and the universal-coefficient sequence.  Let \(V\) be
nontrivial.  If \(\operatorname{char}k\ne2\), averaging over the central
\(C_2\) gives either \(H^j(T_n,V)=0\) for every \(j\), when the centre
acts nontrivially, or
\(H^2(T_n,V)=H^2(A_n,V)\), when it acts trivially.
If \(\operatorname{char}k=2\), the normal central \(2\)-subgroup
acts trivially on \(V\).  The degree-two terms of the
Lyndon--Hochschild--Serre spectral sequence give
\[
 \dim H^2(T_n,V)
 \leq \dim H^2(A_n,V)+\dim H^1(A_n,V)+\dim V^{A_n}.
\]
The last term is zero, and \(A_n\) is two-generated, so (SP5)
gives the claimed bound.  This is also the central-subgroup
calculation in \cite[Lemma~4.1 and its proof, p.~481]
{GKKL2007SuperperfectB}.

\par\smallskip\noindent\textbf{Product cohomology and the infinite tail.}\quad
Over an algebraically closed field, every simple module for a finite
product of the \(T_n\)'s is a tensor product of simple modules of its
factors.  This holds in modular characteristic: quotienting the
finite-dimensional group algebras by their radicals reduces the
claim to tensor products of matrix algebras.
Tensoring free bar resolutions and applying the Kunneth formula
over a field gives
\[
 H^a\!\left(\prod_i G_i,\bigotimes_i V_i\right)
 =
 \bigoplus_{\sum_i a_i=a}\bigotimes_i H^{a_i}(G_i,V_i).
 \tag{SP6}
\]
The trivial coefficient has \(H^1(T_n,k)=H^2(T_n,k)=0\).
For a nontrivial simple factor \(V_i\), \(H^0(T_n,V_i)=0\).
It follows that a simple module with one nontrivial support factor
has \(H^1\) of dimension at most twice its dimension and \(H^2\)
of dimension at most five times its dimension.  With two nontrivial
support factors, \(H^1=0\), and the only possible degree-two term is
the tensor product of two \(H^1\)'s, of dimension at most four times
the module dimension.  With at least three support factors,
both \(H^1\) and \(H^2\) vanish.  For the trivial simple module both
groups vanish.

These statements compute the continuous cohomology of the full
infinite product as well.  A finite coefficient module is trivial
on a full tail.  Continuous cochains with finite values on that
tail factor through a finite coordinate product, so its cochain
complex is the filtered union of the corresponding finite-product
complexes.  Formula (SP6) and (SP3) show that the tail has zero
\(H^1\) and \(H^2\) with trivial coefficients.  The
Lyndon--Hochschild--Serre spectral sequence for head times tail
therefore identifies \(H^1\) and \(H^2\) with those of the finite
head acting on the given module.  Coefficient extension and
devissage now prove (SP1).  Shapiro's lemma proves the stated
open-subgroup bounds, since induction from \(J\) multiplies the
coefficient dimension by \([\Pi:J]\).

\par\smallskip\noindent\textbf{A characteristic-independent dimension bound.}\quad
Let \(V\) be a nontrivial simple \(T_n\)-module over an algebraically
closed field of characteristic \(\ell\).  Its kernel is contained
in the centre.  If \(\ell\ne3\), take the elementary abelian subgroup
of \(A_n\) generated by \(\lfloor n/3\rfloor\) disjoint three-cycles.
If \(\ell=3\), take \(\lfloor n/5\rfloor\) disjoint five-cycles.
Each of these odd-order elementary abelian groups lifts to a
subgroup of \(T_n\) with the same order.  To check this directly,
let its exponent be the odd prime \(s\).  A lift of a basis
element can be multiplied by the central involution to have order
\(s\).  Commutators of these lifts are central of order dividing
two and also of order dividing \(s\), so they commute.
The resulting elementary abelian subgroup meets the centre
trivially and acts faithfully on \(V\).

Its order is prime to \(\ell\), hence its action is simultaneously
diagonalizable.  A faithful diagonal representation of
\((C_s)^a\) requires at least \(a\) characters, since they must
span its character space.  Consequently
\[
 \dim V\geq\lfloor n/5\rfloor\geq n/6,\qquad n\geq31.
 \tag{SP7}
\]
In particular every nontrivial simple factor has dimension at least
two.  This argument avoids using a lifted elementary abelian
\(2\)-subgroup, which need not remain abelian.

\par\smallskip\noindent\textbf{Counting simple modules.}\quad
Let an absolutely simple module of dimension \(r\geq2\) have
nontrivial support dimensions \(a_1,\ldots,a_t\), at degrees
\(n_1,\ldots,n_t\).  Then
\[
 r=\prod_i a_i,\qquad t\leq\log_2r,\qquad
 \sum_i a_i^2\leq r^2,\qquad \sum_i n_i^2\leq36r^2.
 \tag{SP8}
\]
The middle inequality follows by induction from
\(a^2+b^2\leq a^2b^2\) for \(a,b\geq2\).
We may order the support, thereby overcounting.  There are at most
\((6r)^{1+\lfloor\log_2r\rfloor}\leq2^{2r^2}\)
choices for its degree sequence.
Here the displayed estimate follows from
\(1+\lfloor\log_2r\rfloor\leq r\) and
\(\log_2(6r)\leq2r\) for \(r\geq2\).
The coordinate-copy choices
number at most \(2^{\sum_i n_i^2}\).
For each coordinate, the number of simple labels is at most
\(\dim k[T_{n_i}]=|T_{n_i}|=n_i!\leq2^{n_i^2}\).
Thus the number of absolutely simple \(r\)-dimensional modules is
at most \(2^{74r^2}\), uniformly in the coefficient characteristic.
For \(r=1\) there is only the trivial module.

For a finite field \(\mathbf F_q\), every simple module becomes
over its algebraic closure a sum of distinct Frobenius-conjugate
absolutely simple modules, with unique descent.  Indeed, the
semisimple quotient of a finite group algebra is a product of
matrix algebras over finite fields, since finite division rings
are commutative.  Splitting each field factor gives precisely
this description.  Hence the number \(s_r(q)\) of simple
\(\mathbf F_q[\Pi]\)-modules of dimension \(r\) satisfies
\[
 s_r(q)\leq\sum_{e\mid r}2^{74(r/e)^2}\leq q^{75r^2}.
 \tag{SP9}
\]
The first inequality may overcount but is enough.

\par\smallskip\noindent\textbf{Extensions and arbitrary representations.}\quad
For two absolutely simple modules \(U,V\), tensoring bar resolutions
gives
\[
 \operatorname{Ext}^1(U,V)=
 \bigoplus_i H^1(G_i,\operatorname{Hom}(U_i,V_i))
 \otimes\!\!\bigotimes_{j\ne i}\operatorname{Hom}_{G_j}(U_j,V_j).
 \tag{SP10}
\]
If they differ at two coordinates, this is zero.  If they differ
at one coordinate, (SP4) bounds its dimension by
\(2\dim U\dim V\).  If \(U=V\), Schur's lemma and (SP8)
bound it by \(2\sum_i a_i^2\leq2(\dim U)^2\).
Trivial coordinates contribute no term because the coordinate
groups are perfect.  In particular, the computation remains valid
for the full product: an extension cannot acquire new support
on a factor trivial on both its submodule and quotient, since
such a factor would act through an additive abelian group.
After coefficient extension and devissage in the second variable,
\[
 \dim_{\mathbf F_q}\operatorname{Ext}^1(U,W)
 \leq2\dim U\dim W
 \tag{SP11}
\]
for every simple \(U\) and every finite-dimensional \(W\).

Fix a composition \(r=r_1+\cdots+r_t\).  By (SP9) there are at
most \(q^{75\sum_i r_i^2}\) choices of simple factors.  Successive
extensions contribute at most \(q^{2\sum_{i<j}r_ir_j}\), by (SP11).
We count extensions with identified ends, which only overcounts
isomorphism classes.  Since
\[
 75\sum_i r_i^2+2\sum_{i<j}r_ir_j\leq75r^2
\]
and there are \(2^{r-1}\) compositions of \(r\), we obtain (SP2),
for example with \(C=76\).  Isomorphism classes of these modules
are exactly conjugacy classes of the representations in (SP2).

\par\smallskip\noindent\textbf{Smallness and failure of finite generation.}\quad
A nontrivial image of \(T_n\) has order at least \(|A_n|\).
Thus a homomorphism to a fixed finite group kills all but finitely
many degrees, and only finitely many coordinate groups remain.
Continuity kills their full tail, proving the asserted finiteness
of homomorphisms.  An open subgroup contains a full tail and is
a product of that tail with a subgroup of a finite head.
Every finite solvable quotient kills the perfect tail, so the
maximal prosolvable quotient is finite.

Finally, \(\Pi\) maps onto \(A_n^{m_n}\) for every odd \(n\geq31\).
If \(d\) elements generated \(\Pi\), the coordinate \(d\)-tuples
in this quotient would have to be pairwise distinct: equal
tuples force a diagonal relation at the corresponding two
coordinates.  Therefore \(m_n\leq|A_n|^d\).
This contradicts \(2^{n^2}>(n!)^d\) for all sufficiently large
odd \(n\), proving non-finite-generation.
\end{proof}

\begin{remark}[Literature and exact scope]
Products of Schur covers of alternating groups already occur in
\cite[Section~1.6, Proposition~1.14, and Section~2, especially
Proposition~2.4 and Theorem~2.5]{CorobCookVannacci2021B}.
We use this literature as context and make no claim of priority
for the construction or its variants.  The present multiplicities
are \(2^{n^2}\), and the assertions proved above are the precise
ones used here.
The example shows that adding an all-characteristic linear
\(H^2\) bound to smallness, finite open prosolvable quotients,
linear \(H^1\), and uniform quadratic representation growth
still does not by itself imply topological finite generation.
It does not refute any geometric statement in this paper.
The finite exact certificate
\nolinkurl{parts/B/post14_research/code/check_superperfect_product.py}
checks the displayed permutation generators, odd-cycle subgroups,
support-square and counting inequalities, and the low-degree tensor
coefficient bounds.  Its output is
\nolinkurl{parts/B/post14_research/output/superperfect_product.json}.
It does not compute Schur multipliers or repeat the cohomology
computations cited in the literature.
\end{remark}


% Post-checkpoint-14 research. Checked auxiliary statements; not in the CP14 PDF.
\begin{proposition}[Proved: contracting the zero-cost exceptional blocks]
\label{D:contract-zero-cost-blocks}
In the algebraic surface-germ setting of
Proposition~\ref{D:exceptional-degree-linear-core}, each connected
noncore block can be contracted to a rational double point.  The
contractions may be performed simultaneously, give a normal scheme
\(\overline Z\), and the morphism \(f:Z\to Y\) descends to
\(\overline Z\to Y\).  Away from the blocks the contraction is an
isomorphism.  The number of resulting rational double points is at
most \(2nc+n\), with the notation of that proposition.
\end{proposition}
\begin{proof}
Every component of a noncore block is a smooth rational curve of
square \(-2\); its intersection matrix is negative definite.
For every nonzero effective integral cycle \(T\) supported on the
block, adjunction gives
\[
 K_Z\cdot T=0,\qquad T^2\in2\mathbf Z_{<0},\qquad
                    p_a(T)=1+T^2/2\le0.                 \tag{D.RC.1}
\]
We apply Artin's contraction criterion
\citep[Theorem~2.3 and its full proof; published Russian
translation, pp.~9--10]{Artin1962D}.  Its hypotheses are precisely
negative definiteness and the inequalities in (D.RC.1); there is no
restriction on the characteristic.  We may first place the smooth
surface \(Z\) in a smooth projective compactification.  Indeed a
proper compactification can be resolved away from its smooth open
\(Z\), and a proper regular surface is projective
\citep[Lemma~54.16.11, tag~0C5P]{StacksDSurfaceResolution}.
Artin's theorem contracts the block on this projective surface to a
normal projective surface, changing nothing outside the block.
Restricting to the image of \(Z\) gives a scheme contraction of the
required open surface.  Distinct blocks are disjoint, so these
contractions can be carried out successively without changing the
remaining blocks or their intersection numbers.
For a cover of the punctured strict henselization of the algebraic
germ, first descend the finite cover, its resolution and the finite
list of indicated curves to an etale neighbourhood of finite type,
as in Proposition~\ref{D:surface-uniform-different}; perform this
construction there and then base change.  We do not use a projective
compactification of an arbitrary complete local nonalgebraic surface.

The resolution of each new singularity has precisely the contracted
smooth rational \((-2)\)-curves and has no exceptional curve of the
first kind.  The characteristic-free criterion for a rational
double point therefore applies
\citep[Proposition~0.4.8 and proof, author's April~2024 text,
p.~69]{C_Enriques2025}.  More explicitly, (D.RC.1) is Artin's
rationality criterion; for the fundamental cycle its arithmetic
genus is zero, hence its square is \(-2\), which gives multiplicity
two.  This is also the argument in the cited proof.

Each contracted block is connected and all its components map to
points of \(Y\).  Their images are therefore one point.  The graph
image of \(Z\to\overline Z\times Y\) is proper over \(\overline Z\)
and has one-point fibres.  It is thus finite and birational over
the normal surface \(\overline Z\), hence isomorphic to it.
This proves descent of \(f\).  The numerical bound is the block
bound of Proposition~\ref{D:exceptional-degree-linear-core}.
\end{proof}

\begin{lemma}[Proved: a uniform generator bound at rational double points]
\label{D:RDP-uniform-local-generators}
Let \(R\) be the complete local ring of a rational double point
over an algebraically closed field of positive characteristic.
Then
\[
          d\bigl(\pi_1(\operatorname{Spec}R-\{\mathfrak m\})\bigr)
                    \le6.                               \tag{D.RC.2}
\]
In particular its homomorphisms to any finite group \(H\) number
at most \(|H|^6\).  The bound is independent of Dynkin length,
characteristic and coindex, and does not assume a tame quotient
presentation of the singularity.
\end{lemma}
\begin{proof}
We use the complete table of local etale groups in
\citet[Convention~1.3(2), p.~81, and Table~7, p.~93]
{C_Matsumoto2022RDP}, recording Artin's calculations.
The unbounded-order families in that table are cyclic, dihedral,
and binary dihedral groups.  They have at most two generators;
the binary dihedral presentation is
\[
 \langle a,b\mid a^{2m}=1,\ b^2=a^m,\ bab^{-1}=a^{-1}\rangle.
\]
Every other listed group has order at most \(120\).  A finite
group of that order has at most six generators: in an irredundant
successive generating sequence, each proper subgroup enlargement
at least doubles the order, while \(2^7>120\).
Images of six fixed generators determine a homomorphism, proving
the last assertion.  The source specifically includes the
characteristic-two non-taut \(D\)-types; no group-scheme quotient
description is used.
\end{proof}

\begin{example}[Proved: a remaining branch invalidates a punctured-local bound]
\label{D:RDP-branch-gluing-obstruction}
Let \(k\) be algebraically closed of characteristic two and put
\[
 R=k[[u,v,w]]/(uv-w^3),\quad
 X=\operatorname{Spec}R,\quad x=\mathfrak m_R,\quad
 D=V(u,w).
\]
The full puncture \(X-\{x\}\) has fundamental group \(C_3\),
but
\[
               H^1_{\mathrm{et}}(X-D,\mathbf F_2)
                        \text{ is infinite}.             \tag{D.RC.3}
\]
Thus a uniform bound for the punctured group of a contracted
rational double point does not bound the covers of its local
complement when other removed curves pass through the point.
\end{example}
\begin{proof}
The ring \(R\) is the invariant ring of \(k[[s,t]]\) under
\(C_3\), acting by \((s,t)\mapsto(\zeta s,\zeta^{-1}t)\).
The invariant monomials are generated by
\(u=s^3,v=t^3,w=st\), with relation \(uv=w^3\).
The action is free away from the origin.  The puncture of
\(\operatorname{Spec}k[[s,t]]\) is simply connected by purity
\citep[Lemma~58.21.3(1), tag~0BMA]{StacksDLocalPurity}, since
every finite etale cover extends over the strictly henselian
regular local scheme.  The connected etale \(C_3\)-cover of
punctures therefore proves the first assertion.  Equivalently it
is the \(A_2\) row of the table used above.

We have \(X-D=D(u)\): modulo \(u\), the equation forces
\(w^3=0\).  For every \(a\in k\), consider the Artin--Schreier
torsor
\[
                   z^2-z=a/u.                           \tag{D.RC.4}
\]
It is finite etale on \(D(u)\).  At the height-one prime
\((u,w)\), the element \(v\) is a unit and \(w\) is a
uniformizer, so the normalized valuation of \(u\) is three.
For \(a\ne0\), the right side of (D.RC.4) has valuation \(-3\).
If \(b\) has negative valuation then
\(v(b^2-b)=2v(b)\) is even; if \(b\) has nonnegative valuation
then so does \(b^2-b\).  Thus \(a/u\) is not an
Artin--Schreier coboundary even in the fraction field.  Applying
this to \(a-a'\) shows that distinct constants give distinct
classes.  In fact (D.RC.4) defines an injection of additive
\(\mathbf F_2\)-spaces \(k\hookrightarrow H^1(D(u),\mathbf F_2)\).
This proves (D.RC.3).

The exact certificate
\nolinkurl{parts/C/code/check_RDP_branch_obstruction.py}
checks the finite monomial, group-action and valuation calculations;
its output is
\nolinkurl{parts/C/output/RDP_branch_obstruction.json}.
The infinite-family argument is the valuation proof just given.
\end{proof}

\begin{remark}[Exact obstruction to compression]
In the core construction the images of the core components usually
pass through the new rational double points.  The relevant cover is
then defined on the local complement of those branches.  It need
not extend to the full puncture, so (D.RC.2) cannot be inserted as a
bound for that local gluing problem.  The example is only a failed
inference for this method: it is not a counterexample to either
original finite-generation question.
\end{remark}


% Post-CP13 research source; deliberately unassembled during the CP14 freeze.
\begin{example}[Refuted: ample Chow sections preserve codimension-two boundaries]
\label{D:Chow-surface-boundary-obstruction}
Let $k$ be any algebraically closed field, let
$S=\mathbf P^3_k$, and let $T=L=V(x_0,x_1)$ be a line.
There is a smooth projective modification $\pi:Y\to S$ and a
very ample invertible sheaf $\mathcal H$ on $Y$ such that the
normal image of a geometric generic hyperplane section has a
divisorial intersection with $L$.  Indeed, this happens for
every integral effective ample divisor on $Y$ other than the
exceptional divisor.  Thus an arbitrary ample-section argument
on a projective Chow modification does not preserve the
codimension-two boundary condition required in D.2(ii).
\end{example}

\begin{proof}
Take $Y=\operatorname{Bl}_L\mathbf P^3$.  The affine blowup
charts identify it with the incidence hypersurface
\[
 Y=V(x_0v-x_1u)\subset\mathbf P^3\times\mathbf P^1,
 \qquad [u:v]\in\mathbf P^1.
\]
For example, on $u\ne0$ the equation is $x_1=x_0(v/u)$,
which is the chart of the blowup in the generator $x_0$;
the other chart uses $x_1$.  These charts also show smoothness
\cite[Tag 0804, Lemma 31.33.2]{StacksDChowSection}.
The exceptional divisor is
$E=\pi^{-1}(L)=L\times\mathbf P^1$.
The restriction
$\mathcal H=\mathcal O_{\mathbf P^3\times\mathbf P^1}(1,1)|_Y$
is very ample by the explicit Segre embedding
\cite[Tag 01WD, Lemma 27.13.6]{StacksDChowSection}.
Writing $H=\pi^*\mathcal O_{\mathbf P^3}(1)$, one also has
$\mathcal H=2H-E$.

Every section of $\mathcal H$ has the form
$u\,a(x)+v\,b(x)$ for linear forms $a,b$, modulo the incidence
equation.  To check completeness, use
\[
 0\longrightarrow\mathcal O
 \longrightarrow\mathcal O(1,1)
 \longrightarrow\mathcal H\longrightarrow0
\]
on $\mathbf P^3\times\mathbf P^1$ and the vanishing of
$H^1(\mathcal O)$, obtained from the standard projective-space
cohomology calculation
\cite[Tag 01XT, Lemma 30.8.1]{StacksDChowSection}.
Off $L$, its image is the quadric
\[
                 x_0a(x)+x_1b(x)=0,
\]
which contains $L$.
The family contains the smooth quadric
\[
 Q=V(x_0x_2+x_1x_3),\qquad
 \widetilde Q=Y\cap V(ux_2+vx_3).
\]
Smoothness holds in every characteristic, since the four
partial derivatives are $x_2,x_3,x_0,x_1$.  The projection
$\widetilde Q\to Q$ is an isomorphism: away from $L$ its
inverse has $[u:v]=[x_0:x_1]$, and near $L$ it has
$[u:v]=[-x_3:x_2]$.  These formulas agree on the overlap by
the equation of $Q$.  In particular,
$\widetilde Q\cap E\to L$ is an isomorphism.
The smooth locus of the projective family of quadrics is
open; the displayed smooth member shows that its geometric
generic member is smooth.  Its line $L$ is Cartier, and its
strict transform is likewise isomorphic to it.  Therefore
normalizing the image does not remove the boundary curve.

There is a numerical explanation independent of this chosen
linear system.  Let $B\subset Y$ be an integral effective
ample divisor different from $E$.  On every exceptional
fibre $F\simeq\mathbf P^1$, the restriction of
$\mathcal O_Y(B)$ has positive degree.  Its defining section
either vanishes identically on $F$ or has a zero there.
Consequently $B$ meets every such fibre, and
$L\subset\pi(B)$.  Since $B\ne E$, its image is a surface.
The finite normalization of that image still has a
one-dimensional inverse image of $L$.  This proves the
assertion for every such ample divisor and, in particular,
for geometric generic hyperplanes of any projective
embedding of $Y$.
\end{proof}

\begin{remark}[Scope and a possible repair]
This example refutes the stated slicing inference, not D.1,
D.2(ii), or the existence of every possible surface reduction.
Indeed, $S$ itself is projective.  A general divisor from the
non-ample system $|H|$ maps to a plane meeting $L$ in one
point, and therefore repairs the boundary dimension in this
example.  The difficulty over a general proper nonprojective
base is obtaining a sufficiently mobile system with the
required zero degree on modification fibres.  Moreover,
connectedness of the restricted original covers alone would
not suffice for D.1: their ramification along divisors newly
created over a codimension-two boundary has no supplied
tameness condition.

The exact symbolic and finite-field certificate is
\nolinkurl{parts/D/code/check_Chow_surface_boundary.py}; its
output is
\nolinkurl{parts/D/output/Chow_surface_boundary.json}.
It verifies the elimination identities, the seven-dimensional
quadrics-through-$L$ linear system, smoothness equations,
and the stated incidence fibres over $\mathbf F_2$,
$\mathbf F_3$, and $\mathbf F_5$.
\end{remark}


% Post-checkpoint-14 theorem. Full author/root/B/C proof audits passed.
\begin{theorem}[Proved: a generic surface detects every original cover]
\label{D:generic-surface-finite-punctures}
Let $k$ be an algebraically closed field of arbitrary characteristic,
let $S$ be a normal proper integral $k$-scheme of dimension $d\ge2$,
and let $T\subset S$ be a finite closed subset. Put $U=S\setminus T$.
There exist an algebraically closed extension $K/k$ of finite
transcendence degree, a normal proper integral surface $X/K$, and a
finite morphism $g:X\to S_K$ onto a closed integral surface, such that
$T_X=g^{-1}(T_K)$ is finite and the composite morphism
\[
       X\setminus T_X\longrightarrow U_K\longrightarrow U
\]
induces a continuous surjection
\[
       \pi_1(X\setminus T_X)\twoheadrightarrow\pi_1(U).
       \tag{D.GS.1}
\]
The extension and surface are fixed independently of the finite cover
of $U$. The target in (D.GS.1) is the fundamental group of the original
$k$-scheme $U$; no surjectivity onto $\pi_1(U_K)$ is asserted.
Consequently, full finite generation for normal proper surfaces with
finite punctures over algebraically closed fields would imply full
finite generation for normal proper schemes of every dimension with
finite punctures.
\end{theorem}

\begin{proof}
For $d=2$ we take $K=k$, $X=S$ and the identity. Suppose $d\ge3$.
Chow's lemma gives a proper modification by a quasi-projective scheme
which is an isomorphism over a dense open. Take the reduced closure of
that open in the modification. Since $S$ is proper, we obtain an
integral projective scheme $Y$ and a proper birational morphism
$q:Y\to S$. Fix a projective embedding $Y\hookrightarrow\mathbf P^N_k$.
There is a nonempty smooth open $W\subset U$ over which $q$ is an
isomorphism; we also regard $W$ as a smooth locally closed subscheme of
$\mathbf P^N_k$. We use only this dense open for the Bertini argument.
In particular, we do not require a resolution of $S$ or normality of a
whole hyperplane section of $Y$
\citep[Tags 0200 and 0FD6]{StacksDGenericSurface}.

We first verify the exact generic-section input, including its application
to a covering. Let $F$ be an algebraically closed field, let
$W_0\hookrightarrow\mathbf P^N_F$ be smooth and integral of dimension
$e\ge3$, and let $V_0\to W_0$ be a connected finite etale cover.
Then $V_0$ is smooth and integral, hence geometrically irreducible.
Pull back the hyperplane linear system to $V_0$. It is basepoint-free.
Choose a closed point $v\in V_0$, with image $w$. Since $V_0\to W_0$
is etale and $W_0\hookrightarrow\mathbf P^N_F$ is an immersion, we can
choose two hyperplanes through $w$ whose local equations have linearly
independent differentials at $v$. Their pulled-back intersection is
smooth of codimension two near $v$. Its irreducible component through
$v$ thus has codimension two in $V_0$. A third hyperplane not passing
through $w$ does not contain that component. These observations verify
conditions (3) and (4) of Stacks, Lemma 37.32.3, as well as its finite-type
and geometric-irreducibility hypotheses. That lemma therefore says that
the pullback of a geometric generic hyperplane is geometrically
irreducible. Meanwhile Lemma 33.47.3 says that the corresponding generic
section of $W_0$ is smooth. Its finite etale pullback is smooth as well,
and hence the geometrically irreducible pullback is geometrically
integral. The section has dimension $e-1$. The same argument applies
to the identity cover of $W_0$
\citep[Tags 0G4F and 0FD6, full proofs]{StacksDGenericSurface}.

Now choose a fixed tower $k=K_0\subset K_1\subset\cdots\subset K_{d-2}$
as follows: $K_i$ is an algebraic closure of
$K_{i-1}(t_{i,0},\ldots,t_{i,N})$, where each displayed set is
algebraically independent over $K_{i-1}$. Let $H_i$ be the hyperplane
with these coefficients. Successively intersect $W$ with these
hyperplanes and extend scalars at each step. The preceding paragraph,
applied to the identity cover, gives a smooth integral surface $W_*$
over $K=K_{d-2}$.

For every connected finite etale cover $V\to U$ defined over $k$, the
restriction $V_W\to W$ is connected: $V$ is normal and connected, hence
integral, and its inverse image of the dense open $W$ is nonempty and
integral. Apply the same paragraph at each successive stage to this
restriction. It follows that
\[
             V\times_U W_*\quad\hbox{is geometrically integral.}
             \tag{D.GS.2}
\]
Each cover may have its own nonempty Bertini open in a coefficient
space. Its generic point belongs to every such open. Thus the already
fixed tower works for all original covers simultaneously, without an
intersection of infinitely many open sets and without a countability
assumption. This is the geometric-generic quantifier also distinguished
in \citet[Remark 1.5 and Corollary 2.3, pp.~2 and 5--6]{Landesman2020D};
we have verified it directly for our successive hyperplanes.

Take the reduced closure $Z$ of $W_*$ in $Y_K$ and the reduced image
$C=q_K(Z)$ in $S_K$. The image is closed because $q_K$ is proper.
Both $Z$ and $C$ are integral surfaces, and $Z\to C$ is birational:
over the open $W_K$ the map $q_K$ is an isomorphism and the section
$W_*$ is closed in $W_K$. In particular,
\[
                        C\cap W_K=W_*.
\]
Let $X=C^\nu$. Normalization is finite for a scheme of finite type over
a field. It is an isomorphism over the normal open $W_*$. Hence $X$ is
a normal proper integral surface containing $W_*$ as a dense open,
and $g:X\to C\hookrightarrow S_K$ is finite
\citep[Tag 035E, Lemmas 29.55.8--11]{StacksDGenericSurface}.
Since $T_K$ is finite, its inverse image $T_X$ under this finite
morphism is finite as well. This is the step for which we use a finite
omitted set. Exceptional curves in $Z$ may be contracted in $C$; they
do not make $T_X$ positive-dimensional.

For a connected finite etale cover $V\to U$, its pullback to
$X\setminus T_X$ is finite etale over a normal integral base. Each of
its connected components is open and closed and maps surjectively to
the base. Every component therefore meets the dense open $W_*$.
Equation (D.GS.2) shows that this restriction is connected, so there is
only one component. The finite-cover criterion for surjectivity of
profinite fundamental groups proves (D.GS.1). The last assertion follows
because a quotient of a finitely generated profinite group is finitely
generated.
\end{proof}

\begin{corollary}[Proved: quantitative surface bounds at finite punctures in every dimension]
\label{D:normal-finite-puncture-quantitative-transfer}
Let $S$ be normal, proper and connected over an algebraically closed
field of characteristic $p>0$, let $\dim S\ge2$, and let $T\subset S$
be finite. Put $G=\pi_1(S\setminus T)$. There are constants $c,C$,
depending on this fixed pair, such that the following hold.

\begin{enumerate}
\item For every finite field $\Lambda$ and finite-dimensional continuous
$\Lambda[G]$-module $M$,
\[
                    \dim_\Lambda H^1(G,M)\le c\dim_\Lambda M.
\]
For each open subgroup $H\subset G$ the bound is
$c[G:H]\dim M$. The maximal prosolvable quotient of $G$ is finitely
generated.
\item For every $r,f\ge1$, with $q=p^f$,
\[
 \#\bigl(\operatorname{Hom}_{\rm cont}(G,\operatorname{GL}_r(\mathbf F_q))
                 /\operatorname{GL}_r(\mathbf F_q)\bigr)
                  \le q^{Cr^2}.
\]
\item Fix $B\ge1$, and put $q=p^f$ for $f\ge1$. Uniformly for finite subgroups
$\widetilde A\subset\operatorname{GL}_r(\mathbf F_q)$ whose scalar
kernel in the projective image $A$ has order at most $B$, the number
of homomorphisms $G\to A$, modulo the restricted ambient
$\operatorname{PGL}_r(\mathbf F_q)$-conjugacy relation, is at most
$q^{C_Br^2}$. In particular the simple-product conclusions proved for
these bounded-scalar families on surfaces hold for $G$ too. This
includes the stated defining-characteristic symplectic and orthogonal
families, and the linear and unitary families with bounded natural
centre.
\end{enumerate}
No degree-two cohomology bound is asserted by this transfer, and these
conclusions do not eliminate arbitrary nonabelian crowns.
\end{corollary}

\begin{proof}
Apply Theorem~\ref{D:generic-surface-finite-punctures} and write
$\Gamma=\pi_1(X\setminus T_X)\twoheadrightarrow G$ for its surjection.
The surface $X/K$ and its finite boundary are fixed, independently of
all coefficients, ranks, and finite covers. Inflation in degree one is
injective, so Theorem~\ref{D:punctured-surface-all-characteristics-H1}
gives (1). Shapiro's lemma gives its open-subgroup form. The maximal
prosolvable quotient is also a quotient of the finitely generated
maximal prosolvable quotient of $\Gamma$.

Precomposition with a surjection is injective on homomorphisms modulo
any fixed target-conjugacy relation: a conjugacy after precomposition
is already a conjugacy on $G$. Thus
Theorem~\ref{D:global-quadratic-representation-bound} gives (2), and
Proposition~\ref{D:bounded-scalar-projective-growth}, applied to the
punctured surface group as in
Corollaries~\ref{D:surface-bounded-centre-simple-product},
\ref{D:bounded-centre-unitary-simple-product} and
\ref{D:symplectic-orthogonal-simple-product}, gives (3).
Equivalently, every collection of distinct simple quotient kernels of
$G$ pulls back to such a collection for $\Gamma$, so its product group
is a quotient of the corresponding surface simple-product group.

There is no analogous unconditional injection on $H^2$ under this
surjection: the five-term sequence can have a nonzero preceding
contribution from $H^1(\ker(\Gamma\to G),M)^G$. The projective counts
were transferred directly, without such an injection.
\end{proof}


% Finite normalization descent from the independently audited generic-surface theorem.
\begin{corollary}[Proved: finite punctures in arbitrary-dimensional proper schemes]
\label{D:nonnormal-finite-puncture-quantitative-transfer}
Let \(S\) be a proper scheme over an algebraically closed field
of characteristic \(p>0\).  Let \(T\) be a finite closed subset
which has codimension at least two in every irreducible component
it meets, and let \(U\) be a connected component of \(S\setminus T\).
Write \(Q=\pi_1(U)\).  Then the following statements hold.
\begin{enumerate}
\item There is a constant \(c_{S,T}\), independent of the finite
coefficient field \(\Lambda\), such that
\[
 \dim_\Lambda H^1(Q,M)\leq c_{S,T}\dim_\Lambda M
\]
for every finite-dimensional continuous \(\Lambda[Q]\)-module \(M\).
The corresponding bound for an open subgroup \(J\) is
\(c_{S,T}[Q:J]\dim_\Lambda M\).
In particular, the maximal prosolvable quotient of \(Q\) is
topologically finitely generated.
\item There is a constant \(C_{S,T}\) such that, for every \(q=p^f\)
and every \(r\geq1\),
\[
 \#\left(\operatorname{Hom}_{\mathrm{cts}}(Q,\operatorname{GL}_r(\mathbf F_q))
          /\operatorname{GL}_r(\mathbf F_q)\right)
 \leq q^{C_{S,T}r^2}.
 \tag{FP1}
\]
\item Fix \(B\geq1\).  For a finite subgroup
\(\widetilde G\subset\operatorname{GL}_r(\mathbf F_q)\), suppose
its scalar subgroup has order at most \(B\), and denote its image
in \(\operatorname{PGL}_r(\mathbf F_q)\) by \(G\).
There is a constant \(C_{S,T,B}\), independent of \(q,r,\widetilde G\),
such that the number of homomorphisms \(Q\to G\), up to conjugacy
in the ambient \(\operatorname{PGL}_r(\mathbf F_q)\), is at most
\[
                         q^{C_{S,T,B}r^2}.                \tag{FP2}
\]
Consequently, for any fixed \(\epsilon>0\), the quotient formed
from all distinct nonabelian simple quotients in such a family
with \(|G|\geq q^{\epsilon r^2}\) is topologically finitely generated.
\end{enumerate}
In particular the bounded-centre linear and unitary families,
and the natural symplectic and orthogonal families established
earlier, satisfy these assertions.  No degree-two cohomology
bound for \(Q\), nor full finite generation of \(Q\), is asserted.
\end{corollary}

\begin{proof}
Reduction does not affect finite etale covers or their local
systems.  Normalize \(S_{\mathrm{red}}\), and let \(V\) be the
inverse image of \(U\).  Excellence of schemes of finite type
over a field makes \(V\to U\) finite and surjective.  Write its
finitely many connected components as \(V_1,\ldots,V_t\).
Each \(V_i\) is a normal proper integral component with finitely
many points removed.  Indeed that complement is the inverse
image of the finite set \(T\); connectedness of a nonempty open
of an integral scheme ensures that a normalized component lies
over a single connected component of \(S\setminus T\).
The codimension hypothesis makes the components of dimension
at most one proper.

For a normal component of dimension at least two,
Theorem~\ref{D:generic-surface-finite-punctures} gives a continuous surjection
\[
 \pi_1(C_i\setminus F_i)\longrightarrow\pi_1(V_i),
 \tag{FP3}
\]
where \(C_i\) is a normal proper surface over an algebraically
closed extension of the ground field and \(F_i\) is finite.
The surface is fixed once \(V_i\) is fixed; it does not depend
on a finite covering group, a representation, or its coefficient
field.  Inflation along (FP3) is injective on \(H^1\).
Precomposition is also injective on representation classes,
whether ordinary linear or ambient projective conjugacy is
used: an equality after a surjective precomposition is an
equality on the target group.

Theorems~\ref{D:punctured-surface-all-characteristics-H1}
and~\ref{D:global-quadratic-representation-bound}, and
Proposition~\ref{D:bounded-scalar-projective-growth}, therefore
give all three bounds on each \(V_i\).  For components of
dimension at most one we use the same already proved
dimension-at-most-two statements with empty omitted set.
Thus there are constants \(c_i,C_i,C_{i,B}\) with the required
uniformity on every normalized component.

Put
\[
                   c=\#\pi_0(V\times_U V).
\]
This is finite, since the double product is noetherian.
We explain both the linear and nonabelian descent estimates;
neither uses cohomology in degree two.

\par\smallskip\noindent\textbf{The first-cohomology bound.}\quad
Regard \(M\) as a finite locally constant sheaf of
\(\Lambda\)-vector spaces on \(U\).  A class in the kernel of
\[
 H^1(U,M)\longrightarrow\bigoplus_i H^1(V_i,M|_{V_i})
\]
is an \(M\)-torsor with a chosen trivialization on \(V\).
The difference of its two pulled-back trivializations is a
section of \(M\) on \(V\times_U V\), satisfying the usual
cocycle condition.  Changing the trivialization adds a
coboundary.  Effective finite-proper descent of finite etale
schemes therefore identifies this kernel with a quotient of
a subspace of \(H^0(V\times_U V,M)\)
\citep[Expose~IX, Theorem~4.7]{C_SGA1}.
A section of a finite local system on a connected scheme is
determined by its value in one geometric stalk.  Consequently
\[
 \dim_\Lambda H^0(V\times_U V,M)\leq c\dim_\Lambda M,
\]
and the desired bound follows with
\(c_{S,T}=c+\sum_i c_i\).
Shapiro's lemma gives the open-subgroup statement.
The prosolvable conclusion is the abelian-crown consequence
of Proposition~\ref{D:cohomology-crown-criterion}.

\par\smallskip\noindent\textbf{Linear representation classes.}\quad
Set \(P=\operatorname{GL}_r(\mathbf F_q)\).
Its representation classes are the isomorphism classes of
torsors for the constant finite group \(P\).
There are at most \(q^{(\sum_i C_i)r^2}\) possible tuples of
their restriction classes to \(V_i\).
Fix representative torsors for one such tuple.
The Isom sheaf between their two pullbacks to \(V\times_U V\)
is finite etale, and each geometric fibre has \(|P|\) elements.
A section on a connected component is determined by one
geometric stalk.  Hence there are at most \(|P|^c\) possible
descent isomorphisms.  The cocycle condition only reduces this
number, and passing to isomorphism classes only reduces it
again.  Finite-proper descent, including full faithfulness for
the group action maps, is the same cited theorem.
Since \(|P|\leq q^{r^2}\), we obtain (FP1).

\par\smallskip\noindent\textbf{Ambient projective conjugacy.}\quad
Now put \(P=\operatorname{PGL}_r(\mathbf F_q)\).
A homomorphism \(Q\to G\) defines a \(P\)-torsor by extension
of structure group.  Two such homomorphisms give the same
torsor class exactly when they are conjugate in \(P\).
This is an equivalence relation on the \(G\)-valued
homomorphisms even when the ambient group does not normalize
the chosen subgroup \(G\).
Restrictions of these torsors to each \(V_i\) still arise
from \(G\)-valued homomorphisms, so their possible tuples
number at most \(q^{(\sum_i C_{i,B})r^2}\).
For each fixed tuple we may count all \(P\)-torsor descent
data, including those whose descended monodromy leaves \(G\);
this is a valid upper bound.  The preceding Isom argument
again contributes at most \(|P|^c\leq q^{cr^2}\), proving (FP2).

Distinct epimorphism kernels are separated by ambient
projective-conjugacy classes.  The order lower bound in the
statement converts (FP2) into a fixed polynomial bound on
the multiplicity of each simple quotient.
The generating-tuple argument already used in
Proposition~\ref{D:bounded-scalar-projective-growth} proves
finite generation of the stated simple-product quotient.
The argument is entirely in degree zero and degree one
descent; it does not infer an \(H^2\) bound from (FP3).
\end{proof}


% Post-checkpoint-14 extension of the finite-puncture reduction.
% Complete author/root/B/C actual-file and scoped primary audits passed.
\begin{theorem}[Proved: surface reduction with controlled exceptional centres]
\label{D:generic-surface-controlled-boundary}
Let $k$ be algebraically closed, let $S$ be a normal proper integral
$k$-scheme of dimension $d\ge2$, and let $T\subsetneq S$ be closed.
Write $U=S\setminus T$. Suppose that there is a proper birational
morphism
\[
                      q:Y\longrightarrow S
\]
with $Y$ normal and projective over $k$, satisfying this condition:
every prime divisor $E\subset Y$ whose image has codimension at least
two in $S$ and is contained in $T$ has zero-dimensional image.

Then there are an algebraically closed extension $K/k$ of finite
transcendence degree, a normal proper integral surface $X/K$, and a
finite morphism $g:X\to S_K$ onto its image, such that, with
$T_X=g^{-1}(T_K)$, the induced map on the original-boundary tame
quotients is surjective:
\[
                 Q(X,T_X)\twoheadrightarrow Q(S,T).
                 \tag{D.CB.1}
\]
Here tameness is required only at the original codimension-one boundary
points in each of the indicated compactifications.

If $T$ has codimension at least two in $S$, then $T_X$ is finite and
there is a surjection of full fundamental groups
\[
                 \pi_1(X\setminus T_X)\twoheadrightarrow\pi_1(U).
                 \tag{D.CB.2}
\]
Thus, in characteristic $p>0$, all the already proved surface bounds
for all-field $H^1$, defining-characteristic quadratic representation
counts, and bounded-scalar projective representation counts hold for
$Q(S,T)$ under the displayed geometric condition. Its maximal
prosolvable quotient and the previously specified bounded-scalar
simple-product quotients are finitely generated. For (D.CB.2) these
conclusions concern the full group $\pi_1(U)$. We assert neither an
$H^2$ transfer nor full finite generation of these groups.
\end{theorem}

\begin{proof}
When $d=2$ the assertion uses $X=S$, $K=k$, so suppose $d\ge3$.
Let $D_1,\ldots,D_s$ be the codimension-one components of $T$ and
let $B_j\subset Y$ be their strict transforms. A proper birational
morphism to a normal scheme is an isomorphism over an open containing
all codimension-one points. Thus every divisorial component of
$q^{-1}(T)$ is either one of the $B_j$ or an exceptional divisor whose
image has codimension at least two. By hypothesis the latter divisors
all have point images. Denote their union by $E_0$ and the finite set
of their images by $P$. The remaining components of $q^{-1}(T)$ have
dimension at most $d-2$
\citep[Tag 0BFP, Lemma 33.17.3 and proof]{StacksDControlledBoundary}.

Choose a dense smooth open $W\subset U$ over which $q$ is an
isomorphism and an embedding $Y\hookrightarrow\mathbf P^N_k$.
Perform the same fixed tower of $d-2$ geometric generic hyperplane
sections used in Theorem~\ref{D:generic-surface-finite-punctures}.
Its final intersection $W_*$ is a smooth integral surface. Every
original connected finite etale cover of $U$ restricts to a connected
cover of $W_*$, by the exact finite-cover Bertini argument in that
theorem.

We also impose the following finite list of geometric generic
conditions. A closed subscheme of $Y$ of dimension at most $d-2$
meets these $d-2$ generic hyperplanes in dimension at most zero.
Indeed at each stage a generic hyperplane contains no
positive-dimensional irreducible component of the preceding cut;
such components drop dimension by one, and a generic hyperplane
avoids all zero-dimensional components. Apply this to the
nondivisorial components of $q^{-1}(T)$.

There is an open $\Omega\subset S$ containing the generic points of
all the $D_j$ such that $q$ is an isomorphism over $\Omega$, the scheme
$\Omega$ is smooth, and the boundary there is a disjoint union of
smooth Cartier divisors $D_j\cap\Omega$; all other components of $T$
are absent. To obtain it, remove the non-isomorphism locus of $q$,
the singular locus of $S$, the singular and pairwise intersection
loci of the $D_j$, and the remaining boundary components, shrinking
only by subsets which avoid each generic point of $D_j$.
For each $j$, the closed subset
\[
                         B_j\setminus q^{-1}(\Omega)
\]
has dimension at most $d-2$. Its intersection with the generic surface
is therefore finite. Apply smooth Bertini successively also to
$q^{-1}(\Omega)$ and to its smooth divisors $B_j\cap q^{-1}(\Omega)$.
The surface section is smooth there and meets each such divisor
transversely at the generic points of its boundary curves. These are
finitely many additional conditions on the same fixed generic tower,
and follow from Stacks, Lemma 33.47.3
\citep[Tag 0FD6]{StacksDGenericSurface}.

Let $Z$ be the reduced closure of $W_*$ in $Y_K$, let $C=q_K(Z)$
with its reduced induced structure, and let $X=C^\nu$. Exactly as in
Theorem~\ref{D:generic-surface-finite-punctures}, $C$ is an integral
proper surface, its normalization is finite, and $X$ contains $W_*$
as a dense open. On $q_K^{-1}(\Omega_K)$ the complete-intersection
section is smooth of pure dimension two; its closed pure
two-dimensional subscheme $Z\cap q_K^{-1}(\Omega_K)$ is a union of
irreducible components, and is therefore smooth too. Consequently
\[
                  \pi_1(X\setminus T_X)\twoheadrightarrow\pi_1(U)
                  \tag{D.CB.3}
\]
is surjective: every original connected cover remains connected on
$W_*$, and each component of its finite etale pullback to the normal
integral scheme $X\setminus T_X$ meets that open.

We describe the codimension-one boundary of $X$. The image of
$Z\cap E_{0,K}$ lies in the finite set $P_K$. Every intersection of
$Z$ with a remaining nondivisorial component of $q^{-1}(T)_K$ is
finite. Hence each curve in $C\cap T_K$ comes from a curve in
$Z\cap B_{j,K}$ for some $j$. Its generic point lies over $\Omega_K$,
where $q_K$ is an isomorphism and $Z$ is smooth and transverse to
$B_{j,K}$. Normalization of $C$ changes nothing near that point.
Thus every codimension-one boundary point of $X$ is of precisely this
transverse original-divisor type. If there are no $D_j$, the entire
set $C\cap T_K$ is finite and so is $T_X$; this proves the extra
boundary assertion in (D.CB.2).

Now let $V\to U$ be a connected original-boundary tame cover. Over
$\Omega$ it is a finite etale cover of the complement of the displayed
smooth divisors, tame at their generic points. The local absolute
Abhyankar lemma gives its precise description: near each such
boundary point, after adjoining an $n$th root of a local boundary
parameter $t$, with $n$ prime to the characteristic exponent, its
normalization is finite etale. Equivalently over a strict localization
the cover is dominated by that Kummer cover
\citep[Expose XIII, Propositions 5.1--5.2 and Corollary 5.3,
pp.~428--430]{C_SGA1}.
After any algebraically closed field extension this description
persists. At a generic boundary point of $X$ the restriction of $t$
is a uniformizer because of transversality. Pulling back the Kummer
cover and its finite etale extension therefore shows that the
pullback of $V$ is tame at this DVR, including separability of the
residual extension. This applies to every codimension-one boundary
point of $X$. No condition is imposed at its isolated boundary
points.

Thus every finite quotient of $Q(S,T)$ pulls back to a tame cover
on $(X,T_X)$. The composite of (D.CB.3) with the quotient to $Q(S,T)$
annihilates the wild inertia groups defining $Q(X,T_X)$, and hence
factors through (D.CB.1). It remains surjective. If the original
boundary has codimension at least two, there are no original
divisorial tameness conditions and (D.CB.3) is exactly (D.CB.2).

The quantitative consequences follow by degree-one inflation and
injectivity of representation classes under surjective precomposition,
using Theorem~\ref{D:tame-surface-all-characteristics-H1},
Theorem~\ref{D:tame-global-quadratic-representation-bound}, and the
surface bounded-scalar results. The argument is the direct transfer
in Corollary~\ref{D:normal-finite-puncture-quantitative-transfer}; it
does not use degree-two inflation.
\end{proof}

\begin{remark}[Exact geometric limitation]
The hypothesis concerns the chosen projective modification and only
its exceptional divisors whose centres are contained in the omitted
set. It holds, for example, if the modification has no exceptional
divisors over that set except ones contracted to finitely many points.
It does not say that an arbitrary proper scheme admits such a
modification for every positive-dimensional boundary. In the blowup
of $\mathbf P^3$ in a line from
Example~\ref{D:Chow-surface-boundary-obstruction}, the exceptional
divisor maps onto the omitted line, so precisely this condition fails.
We have not converted the general positive-dimensional-boundary case
to finite punctures on a surface.
\end{remark}


% Separate extension of finite-puncture descent; no assertion of general finite generation.
\begin{corollary}[Proved: nonnormal schemes with controlled exceptional centres]
\label{D:nonnormal-controlled-boundary-quantitative-transfer}
Let $S$ be a proper scheme over an algebraically closed field $k$
of characteristic $p>0$, and let $T\subset S$ be closed, of
codimension at least two in every irreducible component which it
meets.  Let $U$ be a connected component of $S\setminus T$, and put
$Q=\pi_1(U)$.  Write $S_i$ for the normalizations of the irreducible
components of $S_{\mathrm{red}}$ which meet $U$, and write $T_i$
for the inverse image of $T$ in $S_i$.

For each $S_i$ of dimension at least two, assume that there is a
proper birational morphism $q_i:Y_i\to S_i$, with $Y_i$ normal and
projective, such that every prime divisor on $Y_i$ whose image is
contained in $T_i$ and has codimension at least two in $S_i$ has
zero-dimensional image.  Then all three quantitative conclusions
of Corollary~\ref{D:nonnormal-finite-puncture-quantitative-transfer}
hold for $Q$, with constants depending on this fixed pair and these
modifications.  Explicitly:
\begin{enumerate}
\item There is a constant $c$ such that
\[
             \dim_\Lambda H^1(Q,M)\le c\dim_\Lambda M
\]
for every finite field $\Lambda$ and every finite-dimensional
continuous $\Lambda[Q]$-module $M$.  For an open subgroup $J$ the
constant can be taken to be $c[Q:J]$.  The maximal prosolvable
quotient of $Q$ is topologically finitely generated.
\item There is a constant $C$ such that, for $q=p^f$ and $r\ge1$,
\[
 \#\bigl(\operatorname{Hom}_{\mathrm{cts}}
       (Q,\operatorname{GL}_r(\mathbf F_q))
       /\operatorname{GL}_r(\mathbf F_q)\bigr)
                 \le q^{Cr^2}.
\]
\item Fix $B\ge1$, and put $q=p^f$ with $f\ge1$.  Suppose
$\widetilde G\subset\operatorname{GL}_r(\mathbf F_q)$ has scalar
subgroup of order at most $B$, and let
$G\subset\operatorname{PGL}_r(\mathbf F_q)$ be its projective
image.  The number of homomorphisms $Q\to G$ up to conjugacy in
the ambient projective linear group is at most $q^{C_Br^2}$,
where $C_B$ is independent of $q,r,\widetilde G$.
For every fixed $\epsilon>0$, the quotient formed from all distinct
nonabelian simple quotients in any such family satisfying
$|G|\ge q^{\epsilon r^2}$ is topologically finitely generated.
\end{enumerate}
The omitted set is allowed to have positive dimension.  The
geometric condition on the normalized components is an explicit
hypothesis; it is not asserted for arbitrary projective
modifications.  No degree-two cohomology bound or full finite
generation of $Q$ is asserted.
\end{corollary}

\begin{proof}
Passing to the reduction leaves the finite etale category unchanged.
The normalization of $S_{\mathrm{red}}$ is finite because $S$ is
of finite type over a field.  Its inverse image $V$ of $U$ is a
finite surjective $U$-scheme.  Each nonempty open
$S_i\setminus T_i$ is integral and therefore connected, so it lies
over one connected component of $S\setminus T$.  Consequently the
connected components of $V$ are exactly the nonempty schemes
\[
                         V_i=S_i\setminus T_i
\]
which lie over $U$.

The finite map $S_i\to S$ preserves the dimensions of closed
irreducible subsets.  Thus $T_i$ has codimension at least two in
$S_i$; in particular $T_i$ is empty when $\dim S_i\le1$.
For every $i$ of dimension at least two,
Theorem~\ref{D:generic-surface-controlled-boundary} supplies a fixed
normal proper surface $X_i$ over an algebraically closed extension
of $k$, a finite set $F_i\subset X_i$, and a surjection
\[
                   \pi_1(X_i\setminus F_i)\twoheadrightarrow
                   \pi_1(V_i).
\]
The surface and extension field are independent of the covering
group and the coefficient field.  The established surface bounds
therefore give all three assertions for $\pi_1(V_i)$: use
degree-one inflation for cohomology and injective precomposition
for ordinary or ambient-projective representation classes.
Components of dimension at most one are proper and satisfy the
same bounds by the already established lower-dimensional cases.

For completeness we give the exact estimates which descend these
bounds.  Put $a=\#\pi_0(V\times_UV)$, a finite number.  If a finite
local system $M$ on $U$ has rank $r$, the kernel of
\[
 H^1(U,M)\longrightarrow\bigoplus_iH^1(V_i,M|_{V_i})
\]
is a quotient of a subspace of $H^0(V\times_UV,M)$.  Indeed an
$M$-torsor trivialized on $V$ is specified by its translation
cocycle on the double product, modulo changes of trivialization.
Finite-proper descent is effective and fully faithful for these
finite etale torsors \citep[Expose IX, Theorem 4.7]{C_SGA1}.
A section on a connected component is determined by one geometric
stalk, so this kernel has dimension at most $ar$.  Adding the
finitely many bounds on $V_i$ proves (1).  Shapiro's lemma gives
the stated open-subgroup bound, and
Proposition~\ref{D:cohomology-crown-criterion} gives its
prosolvable consequence.

For (2), put $P=\operatorname{GL}_r(\mathbf F_q)$.  Fix a tuple
of $P$-torsor classes on the components $V_i$ and representatives
for it.  On the double product the Isom scheme between the two
pullbacks is finite etale, with $|P|$ elements in each geometric
fibre.  There are therefore at most $|P|^a$ possible descent
isomorphisms.  The cocycle conditions and passage to isomorphism
classes can only decrease this number.  The componentwise bounds
and $|P|\le q^{r^2}$ yield the required estimate.

For (3), repeat this argument with
$P=\operatorname{PGL}_r(\mathbf F_q)$, restricting the possible
componentwise torsor classes to ones coming from $G$-valued
homomorphisms.  Extension of structure group identifies equality
of the resulting $P$-torsor classes precisely with ambient
$P$-conjugacy.  This is an equivalence relation on the original
$G$-valued homomorphisms even when $P$ does not normalize $G$.
For a fixed tuple it is legitimate to count all $P$-descent data,
including data whose final monodromy does not lie in $G$; this
only enlarges the upper bound.  The same factor $|P|^a$ proves
the estimate.  Distinct epimorphism kernels are separated by
ambient conjugacy, so the order lower bound converts it to a
fixed polynomial bound in $|G|$ on the multiplicities of the
specified simple quotients.  The generating-tuple argument of
Proposition~\ref{D:bounded-scalar-projective-growth} completes the
proof.  All descent used here is in degrees zero and one.
\end{proof}


% Post-checkpoint-14 auxiliary refinement; no representation-count conclusion.
\begin{proposition}[Proved: weighted compression of the exceptional graph]
\label{D:weighted-exceptional-chain-compression}
Retain the notation and hypotheses of
Proposition~\ref{D:exceptional-degree-linear-core}, including the
degree $n$ and constants $a,b,c=3b+a$.  For a component $F$ of the
exceptional support, let $m_F$ be its positive integral coefficient
in $f^*A$.  Mark a noncore vertex if its curve meets the core.
Then
\[
 \sum_{F\in\mathcal K}m_F(-F^2)\le nc,
 \qquad
 \sum_{Q\text{ marked}}m_Q\le nc.                    \tag{D.WC.1}
\]
Every connected noncore block meets the core, so there are at most
$nc$ such blocks.

Each noncore block is a path, or a tree with one trivalent vertex.
In the latter case its three arm lengths are either $(1,1,h)$,
with $h\ge1$, or $(1,2,h)$, with $h\in\{2,3,4\}$.  Let $P$ consist
of the marked vertices together with all trivalent vertices.
Then
\[
                 \#P\le2nc,
                 \qquad \sum_{Q\in P}m_Q\le31nc.      \tag{D.WC.2}
\]
Add a virtual endpoint of weight zero one edge beyond each unmarked
terminal leaf.  Decompose the graph into maximal segments whose
endpoints belong to $P$ or are these virtual endpoints.  The
weights on each segment form an integral arithmetic progression.
The sum of the lengths in edges of all nonconstant segments is
at most $93nc$.  There are at most $3nc$ constant segments, all
with two endpoints in $P$ and the same positive weight at both
ends.  On replacing each constant segment by one edge, the number
of remaining actual noncore vertices is at most $95nc$.

Thus the only unrestricted graph data in this estimate are the
lengths of at most $3nc$ constant-multiplicity rational
$(-2)$-chains.  The proposition does not assert a bound for those
lengths, an object atlas, or a representation-count consequence.
\end{proposition}

\begin{proof}
Use the nonnegative budget from (D.GC.4).  Write its dominating
and positive point-mapped contributions as $B_1$ and $B_2$;
thus $B_1+B_2\le nb$.  The adjunction inequalities used there
give
\[
 \sum_{F\in\mathcal K}m_F(-F^2)
 \le B_1+\sum_{F\text{ dominating}}m_F+3B_2
 \le3nb+na=nc.                                      \tag{D.WC.3}
\]
The pullback cycle is anti-nef on every component of its support:
projection gives $(f^*A)\cdot F<0$ for a dominating component and
$(f^*A)\cdot F=0$ for a point-mapped one.  Therefore
\[
 \Phi:=\sum_{F\in\mathcal K,\ Q\notin\mathcal K}
                         (F\cdot Q)m_Q
 \le\sum_{F\in\mathcal K}m_F(-F^2)\le nc.           \tag{D.WC.4}
\]
Every marked vertex occurs in this sum at least once, proving
the second assertion of (D.WC.1).

Let $N$ be the negative intersection matrix of a connected
noncore block, and let $T$ be its incidence matrix with the core.
Intersection with its point-mapped components gives
\[
                           N m=T^{\mathsf t}m_{\mathcal K}.
\]
If the block had no attachment, positive definiteness of $N$
would give $m=0$, contrary to the positivity of the coefficients
of $f^*A$.  Each block thus contains a marked vertex, and the
number of blocks is at most $\sum_{Q\text{ marked}}m_Q\le nc$.

We record the elementary graph argument, rather than using a
classification of surface singularities.  The matrix is $N=2I-P_0$
with $P_0$ symmetric, nonnegative and integral, and
$\rho(P_0)<2$.  A multiple edge of multiplicity at least two
already violates positive definiteness on its two vertices.
A cycle has adjacency spectral radius two, as does a four-arm
star.  Hence the graph is a tree of valency at most three.
If it contained two trivalent vertices, take two with no
trivalent vertex between them, their joining path, and the
first vertex on each of the four remaining incident edges.
On this principal subgraph, put weight two on the joining path
and weight one on those four leaves.  This is a positive
eigenvector of its adjacency matrix with eigenvalue two,
again impossible.  There is at most one trivalent vertex.

If the arm lengths at that vertex are $u\le v\le w$, elimination
of the three path matrices gives the positive Schur complement
condition
\[
                 \frac1{u+1}+\frac1{v+1}+\frac1{w+1}>1.
                                                               \tag{D.WC.5}
\]
It forces $u=1$.  Either $v=1$, or $v=2$ and $w\in\{2,3,4\}$.
This proves the asserted list by elementary matrix algebra.

There is at most one trivalent vertex per block.  Consequently
$\#P\le\#\{\text{marked vertices}\}+\#\{\text{blocks}\}\le2nc$.
It remains to bound the weight of an unmarked trivalent vertex.
Let its weight be $z$.  On each arm stop at the first marked
vertex, if there is one.  Otherwise stop at the virtual zero
vertex one edge beyond the terminal leaf.  Write $L_i$ for the
distance to the stop and $v_i$ for its weight.  At an unmarked
vertex the pullback intersection equation is
\[
                    2m_Q=\sum_{Q'\sim Q}m_{Q'}.
                                                               \tag{D.WC.6}
\]
Thus weights vary affinely up to the stops, and at the central
vertex (D.WC.6) reads
\[
             z\left(\sum_{i=1}^3\frac1{L_i}-1\right)
                       =\sum_{i=1}^3\frac{v_i}{L_i}.             \tag{D.WC.7}
\]
The coefficient on the left is positive by (D.WC.5), since a
stop can only shorten an arm.  For arms $(1,1,h)$, if neither
short arm contains a mark, then $L_1=L_2=2$, $v_1=v_2=0$, and
(D.WC.7) gives $z=v_3$.  Otherwise one of $L_1,L_2$ is one and
the coefficient is at least $1/2$, giving $z\le2\sum_i v_i$.
For arms $(1,2,h)$, $h\in\{2,3,4\}$, the coefficient is at least
\[
                         \frac12+\frac13+\frac15-1=\frac1{30}.
\]
Hence $z\le30\sum_i v_i$.  These first marked vertices are
distinct within a block and belong to no other block.  Adding
the already marked vertices proves
$\sum_{Q\in P}m_Q\le31\sum_{Q\text{ marked}}m_Q\le31nc$.

After adjoining the specified virtual endpoints, every actual
vertex outside $P$ has valency two.  Equation (D.WC.6) says that
successive weight differences on a maximal segment are equal
integers.  If a segment of length $\ell$ is nonconstant, that
integer is nonzero, so
\[
              \ell\le|m_{\rm right}-m_{\rm left}|
                         \le m_{\rm right}+m_{\rm left}.
\]
Each vertex of $P$ is an endpoint of at most three segments,
and virtual endpoints contribute zero weight.  Summing gives
\[
                \sum_{\text{nonconstant segments}}\ell
                    \le3\sum_{Q\in P}m_Q\le93nc.               \tag{D.WC.8}
\]
A constant segment cannot end at a virtual zero vertex, since
all actual weights are positive.  Its two endpoints lie in
$P$.  The number of constant segments is at most
$3\#P/2\le3nc$.  Once their internal vertices are removed, the
number of remaining actual vertices is at most
$\#P+\sum_{\text{nonconstant}}\ell\le95nc$, as claimed.
\end{proof}


% New post-checkpoint-14 full-generation criterion and toric application.
% Complete author/root/B/C actual-file and current primary source audits passed.
\begin{proposition}[Proved: a projective model with no additional divisors]
\label{D:codimension-one-projective-model}
Let $S$ be a normal proper connected scheme over an algebraically closed
field. Suppose that there is a normal projective integral scheme $P$
and an isomorphism of nonempty opens
\[
                  \phi:S^\circ\xrightarrow{\ \sim\ }P^\circ
\]
such that $P\setminus P^\circ$ has codimension at least two.
No codimension condition is imposed on $S\setminus S^\circ$.
Then, for every
proper closed subset $T\subset S$, the original-boundary tame group
$Q(S,T)$ is topologically finitely generated. In particular, if $T$
has codimension at least two, the full fundamental group
$\pi_1(S\setminus T)$ is topologically finitely generated.

More generally, D.2(ii) holds for a proper scheme whose normalized
irreducible components each admit such a projective model, for every
closed omitted set of codimension at least two in every component it
meets.
\end{proposition}

\begin{proof}
Put $U=S\setminus T$ and
\[
        V=S^\circ\cap U,\qquad V_P=\phi(V),\qquad B=P\setminus V_P.
\]
The schemes $V$ and $V_P$ are isomorphic nonempty opens, and $B$ is
closed. Every codimension-one point $\eta$ of $B$ belongs to
$P^\circ$, since $P\setminus P^\circ$ has codimension at least two.
Its corresponding point in $S^\circ$ has codimension one in $S$ and
belongs to $T$. The local discrete valuation rings at these two
points are identical under $\phi$.

A connected finite etale cover of the normal integral scheme $U$ is
normal and integral. It remains connected on the dense open $V$.
Hence restriction gives a surjection
\[
                     \pi_1(V_P)\twoheadrightarrow\pi_1(U).
\]
If this cover is classified by $Q(S,T)$, then its restriction is
tame at every codimension-one point of $B$, by the equality of
the corresponding boundary DVRs. Consequently the composite map
to $Q(S,T)$ factors through a surjection
\[
                         Q(P,B)\twoheadrightarrow Q(S,T).
                         \tag{D.CP.1}
\]
The source is finitely generated by the projective case of
Proposition~\ref{D:fixed-boundary-positive}. This proves the first
assertion. If $T$ has codimension at least two, no original
codimension-one inertia groups are killed, so $Q(S,T)$ is the full
fundamental group of $U$.

For the final assertion, pass to the reduction and then to the
disjoint union of the finite normalizations of its irreducible
components. The inverse image of the omitted set has codimension
at least two on each such component, because finite morphisms
preserve dimensions of closed irreducible subsets. The proved
normal case gives finite generation on each normalized open.
Finite proper effective descent and its finite-generation
consequence, SGA~1, Expose~IX, Theorem~4.7 and Corollary~5.2,
give finite generation downstairs
\citep[pp.~238--239 and 247--248]{SGA1}.
\end{proof}

\begin{corollary}[Proved: arbitrary boundaries on complete toric varieties]
\label{D:complete-toric-fixed-boundary}
Let $S$ be a complete toric variety over an algebraically closed field
of arbitrary characteristic. For every proper closed subset
$T\subset S$, without a torus-invariance assumption on $T$,
$Q(S,T)$ is topologically finitely generated. Thus D.1 holds for
every complete toric compactification, including nonprojective ones.
If $T$ has codimension at least two, the full group
$\pi_1(S\setminus T)$ is finitely generated.

The corresponding D.2(ii) conclusion also holds for a proper scheme
whose normalized irreducible components are complete toric varieties.
\end{corollary}

\begin{proof}
Fujino and Sato prove that every complete toric variety has a
projective $\mathbf Q$-factorial toric model which is isomorphic to it
in codimension one. Their ground-field convention is explicitly
algebraically closed of arbitrary characteristic
\citep[Theorem~1.1, p.~1; Theorem~1.3, p.~2; full proof in
Section~3, pp.~6--7]{FujinoSato2025D}.

We spell out why the result supplies the common opens required in
Proposition~\ref{D:codimension-one-projective-model}. The toric
birational map identifies the dense tori and their lattices. Being
an isomorphism in codimension one means that the two fans have the
same one-dimensional cones. The common subfan consisting of the
zero cone and these rays defines isomorphic open toric subvarieties
on the two sides. Their complements are unions of orbits of
codimension at least two. Equivalently, this is the common-open
meaning of isomorphism in codimension one used explicitly by the
authors immediately after Theorem~1.1. The model is normal because
it is a toric variety.

Apply Proposition~\ref{D:codimension-one-projective-model} with this
projective model. Its proof applies to every closed $T$; it does not
use invariance of the boundary or introduce extra divisorial
valuations. The nonnormal statement is the final assertion of the
same proposition.
\end{proof}

\begin{remark}[Scope of this full-generation conclusion]
The first proposition requires a dense common open whose complement
on the projective side has codimension at least two. Divisors may be
removed on the $S$ side: removing them does not create additional
boundary DVRs on the projective side. No morphism from $S$ to its
projective model is required. The criterion therefore includes
Proposition~\ref{D:generically-finite-projective-target} after its
Stein factorization, and also the possibly rational codimension-one
toric comparison used above.
The toric source proves this condition in all dimensions. It does
not establish it for arbitrary proper varieties: the source states
the general threefold projectivization question separately as
Problem~1.2. Accordingly the present full-generation class does not
close the unrestricted proper case of D.1 or D.2(ii).
\end{remark}


% Auxiliary geometric-complexity estimate; no cover-atlas or finite-generation conclusion.
\begin{proposition}[Proved: degree-linear size of punctured-cover resolutions]
\label{D:punctured-cover-linear-resolution}
Let $(S,s)$ be a normal algebraic surface germ over an algebraically
closed field of positive characteristic.  Fix a projective resolution
$Y\to S$ which is an isomorphism away from $s$, with nonempty exceptional
curves $E_i$, and a strictly anti-nef positive cycle
\[
 A=\sum_i a_iE_i,\qquad c_i=-A\cdot E_i>0.
\]
Put
\[
 a=\sum_i a_i,\qquad b=K_Y\cdot A+a,\qquad c=3b+a,
 \qquad \lambda=\max_i\frac{b}{a_ic_i}.
 \tag{D.LR.1}
\]
For any connected finite etale cover of $S\setminus\{s\}$ of
degree $n$, normalize $Y$ in its field, resolve, and contract all
smooth rational $(-1)$-curves which map to points of $Y$.
Let $f:Z\to Y$ be the resulting proper generically finite morphism.
The cover need not be Galois.  Then
\[
 \#\operatorname{Irr}\bigl((f^*A)_{\mathrm{red}}\bigr)
 \le n(b+a)+(95+93\lambda)nc.
 \tag{D.LR.2}
\]
In particular the number of exceptional components, the sum of their
arithmetic genera, and the sum of the negatives of their
self-intersections are all bounded linearly in $n$, with constants
depending only on the fixed germ, resolution, and cycle.
The same assertion holds for covers of the punctured strict
henselization, with the same constants.

This is an estimate for resolutions of covers unramified outside
the isolated point.  It does not assert an object atlas, a bound on
the number of covers, or finite generation of the local group.
\end{proposition}

\begin{proof}
Use the notation of Propositions~\ref{D:exceptional-degree-linear-core}
and~\ref{D:weighted-exceptional-chain-compression}, with no original
divisorial boundary.  Write
\[
 C=f^*A=\sum_Hm_HH,\qquad
 R=K_Z-f^*K_Y=\sum_H\delta_HH.
 \tag{D.LR.3}
\]
The second divisor is the divisor of the nonzero determinant of
the differential.  It is effective and supported on the exceptional
fibre, because the cover and the chosen resolution morphism are
etale away from that fibre.  All $m_H$ are positive integers.
At a component $F$ dominating $E_i$, its coefficients satisfy
\[
              m_F=a_i e_F,\qquad \delta_F=d_F,
\]
where $d_F$ is the different exponent of the corresponding
extension of discrete valuation rings, including its full
possibly inseparable residual extension.

\par\smallskip\noindent\textbf{A maximum principle for ramification coefficients.}\quad
Proposition~\ref{D:surface-uniform-different} gives
$d_F/e_F\le b/c_i$ for a Galois cover.  It gives the same separate
bound for every finite cover by taking its finite Galois closure:
in a tower of finite separable DVR extensions $L/M/K$, transitivity
of the different gives
\[
 \frac{d(L/K)}{e(L/K)}
 =\frac{d(M/K)}{e(M/K)}
  +\frac{d(L/M)}{e(L/M)e(M/K)}
 \ge\frac{d(M/K)}{e(M/K)}.
\]
Thus $\delta_F/m_F\le\lambda$ on every dominating component,
independently of the cover degree.

For a point-mapped component $Q$, the chosen contractions give
$K_Z\cdot Q\ge0$, while projection gives $C\cdot Q=0$.
Writing $b_Q=-Q^2>0$, these statements become
\[
 \sum_{H\ne Q}(H\cdot Q)m_H=b_Qm_Q,
 \qquad
 \sum_{H\ne Q}(H\cdot Q)\delta_H\ge b_Q\delta_Q.
\]
Consequently
\[
 \sum_{H\ne Q}(H\cdot Q)m_H
       \left(\frac{\delta_H}{m_H}-\frac{\delta_Q}{m_Q}\right)
                         \ge0.                         \tag{D.LR.4}
\]
Each connected component of the exceptional graph contains a
dominating component.  Otherwise $C$ would have intersection zero
with every component of that connected component, contrary to
negative definiteness of its intersection matrix and the positive
coefficients of $C$.
If the maximum of the finitely many ratios $\delta_H/m_H$ were
larger than $\lambda$, it would occur at a point-mapped vertex.
Every summand in (D.LR.4) would then be nonpositive.  Equality
would force the same maximum at every adjacent vertex, and hence,
by connectedness, at a dominating vertex.  This is impossible.
We have proved the uniform estimate
\[
                         0\le\delta_H\le\lambda m_H
                         \quad\hbox{for every }H.        \tag{D.LR.5}
\]

\par\smallskip\noindent\textbf{The Jacobian cost of a long constant chain.}\quad
Consider one of the constant-multiplicity segments in
Proposition~\ref{D:weighted-exceptional-chain-compression}, of
common weight $m$.  Suppose it has $N-1\ge1$ internal vertices.
These form a chain of smooth rational $(-2)$-curves, and all map
to the same point of $Y$.  Contract just these internal curves.
The contraction argument of
Proposition~\ref{D:contract-zero-cost-blocks} applies to this
negative-definite subchain: every positive cycle $T$ on it has
$K_Z\cdot T=0$ and $T^2$ a negative even integer, so
$p_a(T)=1+T^2/2\le0$.  It gives a normal scheme contraction
$\rho:Z\to\overline Z$, with an $A_{N-1}$ rational double point.
The morphism $f$ descends to $\overline f:\overline Z\to Y$.
Indeed the graph image over $\overline Z$ is proper and
quasi-finite: the contracted chain has one image in $Y$ and all
other fibres are singletons.  The graph projection is therefore
finite birational onto the normal surface $\overline Z$, and
hence an isomorphism.

The completed local ring of an $A_{N-1}$ point in arbitrary
characteristic is
\[
                  B=k[[u,v,w]]/(uv-w^N)
\]
\citep[Theorem 0.4.13, printed p.~72, together with
Proposition 0.4.8]{C_Enriques2025}.
No assertion identifying the two retained end curves with the
coordinate branches is needed.  We use only the underlying
singularity and its minimal resolution.
For clarity, its standard resolution charts, indexed by
$i=0,\ldots,N-1$, are
\[
 u=x^{i+1}y^i,\qquad
 v=x^{N-i-1}y^{N-i},\qquad w=xy.
 \tag{D.LR.6}
\]
The intervening exceptional valuations $v_j$, $1\le j\le N-1$,
therefore satisfy
\[
                     v_j(u)=j,\qquad
                     v_j(v)=N-j,\qquad v_j(w)=1.        \tag{D.LR.7}
\]
These charts can be obtained successively by the normalized
blowup of $(u,w)$: its $u$-chart is smooth, with
$w=uz$ and $v=u^{N-1}z^N$, and its $w$-chart has
$u=wz$ and $zv=w^{N-1}$, so induction applies.
Adjacent charts in (D.LR.6) glue by
$x_{i+1}=y_i^{-1}$ and $y_{i+1}=x_i y_i^2$.
They give a proper resolution whose intervening curves have
normal degree $-2$.  These are the curves of the contracted
chain, since that chain is already its minimal resolution.

A generator of the dualizing module of $B$ is
\[
                         \eta=\frac{du\wedge dw}{u}.
\]
The expression is interpreted in the fraction field and extends
as a dualizing generator.  The hypersurface relation gives
\[
 du\wedge dw=u\eta,\qquad
 dv\wedge dw=-v\eta,\qquad
 du\wedge dv=Nw^{N-1}\eta.
\]
In particular the image of the module of continuous two-forms
in the dualizing module is
\[
                 (u,v,Nw^{N-1})\eta.                    \tag{D.LR.8}
\]
On each chart (D.LR.6), direct differentiation gives
$\rho^*\eta=dx\wedge dy$: the coefficient is
$(i+1)-i=1$, also in positive characteristic.  Thus the
resolution is crepant, with this explicit nonvanishing generator.

Choose regular parameters $z_1,z_2$ at the image point in $Y$.
The completed pullback of $dz_1\wedge dz_2$ is $h\eta$ for
some nonzero
\[
                         h\in(u,v,Nw^{N-1}).
\]
It is nonzero because the original cover is generically separable.
The coefficient $\delta_j$ of $R$ on the $j$th internal curve
is exactly $v_j(h)$, by the displayed crepant formula.  Since
coefficients in $B$ have nonnegative valuations, (D.LR.7)--(D.LR.8)
give, for $j=\lfloor N/2\rfloor$,
\[
 \delta_j\ge\min\{j,N-j,N-1\}=\lfloor N/2\rfloor.
 \tag{D.LR.9}
\]
If the characteristic divides $N$, the last generator of the
ideal is zero and the same inequality still follows from the
first two generators.  Combining (D.LR.9) with (D.LR.5) and
the constant multiplicity $m$ yields
\[
                         N-1\le2\delta_j\le2\lambda m.
                         \tag{D.LR.10}
\]
This argument is independent of the positions of the retained
end curves in the completed $A_{N-1}$ model.

\par\smallskip\noindent\textbf{Summing the chain lengths.}\quad
Let $P$ be the endpoint set in
Proposition~\ref{D:weighted-exceptional-chain-compression}.
There are at most $95nc$ actual noncore vertices after the
internal vertices of constant segments have been suppressed,
and $\sum_{Q\in P}m_Q\le31nc$.
Each endpoint lies on at most three segments.  Both endpoint
weights of a constant segment equal its weight $m$, so
\[
               \sum_{\text{constant segments}}m
                    \le\frac32\sum_{Q\in P}m_Q
                    \le\frac{93}{2}nc.
\]
By (D.LR.10), restoring all their internal vertices adds at
most $93\lambda nc$ vertices.  Adding the core bound
$n(b+a)$ proves (D.LR.2).
All noncore curves have arithmetic genus zero and square $-2$;
the core genus and self-intersection sums are bounded by
Proposition~\ref{D:exceptional-degree-linear-core}.  This proves
the remaining linear estimates.

For the strict-henselian assertion, each finite cover descends
to a finite-type etale neighborhood of the germ.  The chosen
exceptional curves, their intersection numbers, and the
residue field remain unchanged.  The preceding proof applies
there, exactly as in Proposition~\ref{D:surface-uniform-different}.
It produces the same constants.  No assertion about a
degree-linear presentation of the formal neighborhoods or
their covering algebras has been used.
\end{proof}


% Post-review-09 closure lemma; complete D/root/B/C proof and source audits passed.
\begin{proposition}[Proved: finite normal descent of original-boundary generation]
\label{D:finite-normal-tame-descent}
Let $f:Y\to S$ be a finite dominant morphism of normal proper
connected schemes over an algebraically closed field. For a proper
closed subset $T\subset S$, put $T_Y=f^{-1}(T)$ and
$e=[k(Y):k(S)]_{\mathrm{sep}}$. There is a continuous homomorphism
\[
                     Q(Y,T_Y)\longrightarrow Q(S,T)
\]
whose image is open of index at most $e$. Consequently finite
generation of $Q(Y,T_Y)$ implies finite generation of $Q(S,T)$.
If the source can be generated by $d$ elements, the target can be
generated by at most $d+e-1$ elements.

In particular, D.1 holds for every normal proper scheme which
admits a finite dominant morphism from a normal proper scheme
satisfying Proposition~\ref{D:codimension-one-projective-model}.
This includes every normal proper finite image of a complete toric
variety. The finite map need not be separable or tame, and no
compatibility with a torus action is required. If $T$ has
codimension at least two, the assertion concerns the full
fundamental groups.
\end{proposition}

\begin{proof}
Write $U=S\setminus T$ and $U_Y=Y\setminus T_Y$.
Both are nonempty normal integral schemes. Under a finite dominant
map between integral schemes of finite type over a field, a
codimension-one point $y$ of $Y$ maps to a codimension-one point
$s$ of $S$: the two total dimensions agree and the residue-field
extension is finite, so the dimension formula gives equal
codimensions. Thus every original boundary DVR of $(Y,T_Y)$ is
an extension
\[
                     \mathcal O_{S,s}\subset\mathcal O_{Y,y}
\]
of an original boundary DVR of $(S,T)$.

Pullback of a finite etale cover of $U$ is finite etale over $U_Y$.
If its function fields are tame at the original boundary DVRs of
$S$, every field component after base change to $k(Y)$ is tame at
each indicated DVR of $Y$. This is the full tame-base-change
assertion for arbitrary extensions of DVRs, including inseparable
base-field extensions, in
\citep[Lemma~15.116.9(2), tag 0EXY]{StacksDFiniteTameDescent}.
Consequently the homomorphism $\pi_1(U_Y)\to\pi_1(U)$ induces the
displayed homomorphism on the original-boundary quotients. We do
not assume that the finite map itself is tame.

Choose compatible geometric generic points, and put
$K=k(S)$ and $L=k(Y)$. The image of
$\operatorname{Gal}(L^{\mathrm{sep}}/L)$ in
$\operatorname{Gal}(K^{\mathrm{sep}}/K)$ has index
$[L:K]_{\mathrm{sep}}=e$. Indeed, if $L_0$ is the maximal
separable subextension of $L/K$, then $L/L_0$ is purely
inseparable, and restriction identifies its absolute Galois
group with that of $L_0$. The latter is an open subgroup of
index $[L_0:K]=e$ in the absolute Galois group of $K$.

The maps from the two absolute Galois groups to the corresponding
fundamental groups of $U_Y$ and $U$ are surjective, because these
schemes are normal and integral
\citep[Proposition~58.11.3, tag 0BQM; Lemma~58.10.7,
tag 0BQI]{StacksDFiniteTameDescent}. The compatible diagram shows
that the image of $\pi_1(U_Y)$ in $\pi_1(U)$ has index at most
$e$. Passing to $Q(S,T)$ cannot increase this index. The image
of $Q(Y,T_Y)$ is the same image, so it is an open subgroup $H$
of index $h\le e$.

If $Q(Y,T_Y)$ has $d$ topological generators, their images
generate $H$. Adjoining representatives of the $h-1$ nonidentity
left cosets of $H$ generates $Q(S,T)$: every element lies in one
of these cosets and hence in the subgroup generated by $H$ and
these representatives. This proves the stated bound.

Apply Proposition~\ref{D:codimension-one-projective-model}, or
Corollary~\ref{D:complete-toric-fixed-boundary}, to the source
pair $(Y,T_Y)$ for the final assertions. Finiteness of $f$
preserves the codimension of closed irreducible subsets, so
codimension at least two on $S$ gives the same condition on $Y$.
\end{proof}

\begin{remark}[Scope of the descent statement]
The proposition assumes an existing finite morphism of normal
proper schemes; it does not assert existence of abstract group
quotients. A merely proper generically finite map may introduce
boundary divisors above higher-codimension centres, where
original-boundary tameness is not assumed.
\end{remark}


% Complete root/B/D actual-file and exact primary-source audits PASS.
\begin{proposition}[Proved: a fixed canonical-defect annihilator and linear geometric genus]
\label{D:cover-uniform-canonical-defect}
Let $(R,\mathfrak m,k)$ be a complete equicharacteristic normal local
domain of dimension two, with algebraically closed coefficient field $k$.
There is a fixed $\mathfrak m$-primary ideal $J\subset R$ such that the
following holds. Let $A$ be a finite normal local domain over $R$, of
generic degree $n$, with $R\subset A$, and suppose that
\[
 \operatorname{Spec}A\setminus\{\mathfrak m_A\}
 \longrightarrow \operatorname{Spec}R\setminus\{\mathfrak m\}
\]
is etale. For a projective regular resolution
$\pi:Z\to\operatorname{Spec}A$ which is an isomorphism on the puncture,
use the natural trace inclusion $\pi_*\omega_Z\subset\omega_A$. Then
\[
 J\omega_A\subset\pi_*\omega_Z,
 \qquad JH^1(Z,\mathcal O_Z)=0.                         \tag{D.CD.1}
\]
If $Q=(x,y)\subset J$ is one fixed parameter ideal of $R$, then
\[
 p_g(A):=\dim_kH^1(Z,\mathcal O_Z)
 =\operatorname{length}_A(\omega_A/\pi_*\omega_Z)
 \le n\operatorname{length}_R(R/Q).                    \tag{D.CD.2}
\]
The constants and the annihilator are independent of the cover.
For a finite normal $R$-torsion-free algebra, the corresponding assertions
hold componentwise, and the sum of the geometric genera has the same bound
with $n$ the total generic degree. No separability restriction on $n$ is
imposed; etaleness on the puncture supplies generic separability.
\end{proposition}
\begin{proof}
\emph{A fixed different ideal.}
Choose generically separable Noether normalizations
$S_i=k[[x_i,y_i]]\subset R$, finite on $R$, such that the sum of their
Noether differents contains an $\mathfrak m$-primary ideal. Such a finite
choice exists: over the perfect field $k$, the Jacobian ideal is supported
on the isolated singular locus, is the sum of the Kahler differents of
Noether normalizations, and each Kahler different is contained in its
Noether different. A finite selection suffices by Noetherianity
\citep[Section 3, especially (3.1) and Examples 3.2--3.3,
pp.~6--7]{C_IyengarTakahashiStable}.
If $R$ is regular, its identity normalization can be used. Choose once an
$\mathfrak m$-primary ideal
\[
 J\subset\sum_i\mathfrak N(R/S_i).
\]
All rings $R$ and $A$ here are Cohen--Macaulay. In particular $R$ and $A$
are finite free over each $S_i$.

\emph{Trace duality on the punctured cover.}
Put $K=\operatorname{Frac}R$ and $L=\operatorname{Frac}A$.
Using the nondegenerate generic trace pairing, identify
\[
 \operatorname{Hom}_R(A,R)
 =\{c\in L:\operatorname{Tr}_{L/K}(cA)\subset R\}.
\]
This fractional $A$-module is exactly $A$. Indeed, at every height-one
prime of $R$ the finite algebra is etale and its trace dual is itself.
Every height-one prime of $A$ lies over such a prime of $R$, so an element
of the displayed trace dual belongs to every height-one localization of
$A$. Normality then puts it in $A$. Conversely the trace of an integral
element is integral over $R$ and lies in $K$, hence lies in $R$.
We have therefore proved, with its specified generic identification,
\[
 \operatorname{Hom}_R(A,R)=A.                          \tag{D.CD.3}
\]
Finite duality and transitivity identify
$\omega_A=\operatorname{Hom}_R(A,\omega_R)$.

\emph{A Noether different clears the canonical lattice.}
Fix $i$. Choose a generator of the dualizing module of the regular ring
$S_i$ and use finite duality to identify $\omega_R$ with
$\operatorname{Hom}_{S_i}(R,S_i)$ times that generator.
Write $\beta_i$ for the resulting trace section. It is nonzero generically.
The evaluation description of the Noether different is
\[
 \mathfrak N(R/S_i)=
 \operatorname{im}\bigl(\operatorname{Hom}_R(\omega_R,R)
       \longrightarrow R,\quad \varphi\longmapsto\varphi(\beta_i)\bigr),
                                                               \tag{D.CD.4}
\]
with the generator of $\omega_{S_i}$ suppressed in the notation.
This is Stacks Lemma 49.6.7, including its trace normalization
\citep[Tags 0BVT and 0DWM, Lemmas 49.6.7 and 49.15.1--2]
{C_StacksCanonicalDefect}.
The module $\omega_R$ is torsion-free of rank one. Consequently if
$z=\varphi(\beta_i)$, then $z\alpha=\varphi(\alpha)\beta_i$ for every
$\alpha\in\omega_R$, as one checks in its one-dimensional generic fibre.
It follows that
\[
 \mathfrak N(R/S_i)\omega_R\subset R\beta_i.
\]
Apply $\operatorname{Hom}_R(A,-)$ and use (D.CD.3). With the generic
trace identifications and trace transitivity, this gives
\[
 \mathfrak N(R/S_i)\omega_A\subset A\beta_i,            \tag{D.CD.5}
\]
where on the right $\beta_i$ now denotes the trace section for $A/S_i$.
The notation is compatible because
$\operatorname{Tr}_{L/\operatorname{Frac}S_i}
=\operatorname{Tr}_{K/\operatorname{Frac}S_i}\circ
 \operatorname{Tr}_{L/K}$.

\emph{Regularity of this section on the resolution.}
The composite $g_i:Z\to\operatorname{Spec}S_i$ is projective and
generically finite separable. It is an lci morphism of virtual relative
dimension zero. To see this without any finite-type-over-$k$ assumption,
embed $Z$ in a projective space over $S_i$; this is a regular immersion
because both the ambient scheme and $Z$ are regular. The fundamental
class of an lci morphism therefore gives
\[
 \mathcal O_Z\longrightarrow \omega_{Z/S_i}.
\]
After tensoring with the chosen generator of $\omega_{S_i}$, its image of
$1$ is a regular section of $\omega_Z$ whose generic value is precisely
$\beta_i$. Here the normalization is the canonical etale identification:
on the generic finite etale locus, relative duality identifies this image
with the trace functional. Equivalently, a local complete-intersection
presentation gives its Jacobian determinant times the Koszul dualizing
generator; on the etale locus that generator corresponds to the inverse
Jacobian under the trace pairing. These identifications are the ones in
Stacks Lemma 48.29.2 and the finite-flat comparison with duality
\citep[Tags 0E9Z and 0DWM; Lemmas 48.29.2 and 49.15.1--2]
{C_StacksCanonicalDefect}. Thus
$A\beta_i\subset\pi_*\omega_Z$. Summing (D.CD.5) proves the first
inclusion of (D.CD.1). In particular no use of ordinary absolute
differentials on a formal power-series field has been made.

\emph{Duality and length.}
The ring $A$ is complete excellent and has a dualizing complex.
Grauert--Riemenschneider vanishing for normal modifications of local
surfaces gives $R^1\pi_*\omega_Z=0$ in every characteristic
\citep[Tag 0AXD, Proposition 54.7.8, with the proof in Section 54.7]
{C_StacksCanonicalDefect}. We also have $\pi_*\mathcal O_Z=A$ and
$R^j\pi_*\mathcal O_Z=0$ for $j>1$. Put
$M=H^1(Z,\mathcal O_Z)$, a finite-length $A$-module. Duality applied to
$A\to R\pi_*\mathcal O_Z\to M[-1]$ gives
\[
 0\longrightarrow\pi_*\omega_Z\longrightarrow\omega_A
 \longrightarrow D_A(M)\longrightarrow0,               \tag{D.CD.6}
\]
where $D_A$ is Matlis duality. Thus $J$ kills $D_A(M)$ and hence $M$;
it also follows that the quotient in (D.CD.6) has length $p_g(A)$.
As an $R$-module, $\omega_A$ is maximal Cohen--Macaulay of rank $n$.
For any fixed parameter ideal $Q=(x,y)\subset J$,
\[
 \operatorname{length}_R(\omega_A/J\omega_A)
 \le\operatorname{length}_R(\omega_A/Q\omega_A)
 =n\operatorname{length}_R(R/Q).
\]
The last equality is the multiplicity formula for a maximal
Cohen--Macaulay module, or follows by taking ranks over $k[[x,y]]$,
over which both modules are finite free. Every residue field in the
connected finite cover is $k$, so these finite lengths are also the
stated $k$-dimensions. This proves (D.CD.2). A finite normal algebra over
a complete local ring is a finite product of normal local domains;
apply the argument to each dominating component and add their ranks.
\end{proof}

\begin{corollary}[Proved: cohomology of fixed pulled-back exceptional thickenings]
\label{D:cover-thickening-linear-cohomology}
In the algebraic or strict-henselian punctured-cover setting of
Proposition~\ref{D:punctured-cover-linear-resolution}, fix the positive
exceptional cycle $A_0=\sum_i a_iE_i$ on the base resolution $Y$.
For a degree-$n$ cover and the regular model $f:Z\to Y$ used there, put
$C=f^*A_0$. For each fixed integer $t\ge1$ there is a constant $c_t$,
depending only on the base germ, its fixed resolution, $A_0$, and $t$,
such that
\[
 h^0(\mathcal O_{tC})+h^1(\mathcal O_{tC})\le c_t n.
                                                               \tag{D.CD.7}
\]
The cohomology here is that of the full possibly nonreduced Cartier
cycle. The same assertion holds on the completed models.
\end{corollary}
\begin{proof}
Complete the base germ. Excellence makes the normalization and its
completion compatible and normal; finite covers split into complete
normal local factors with total generic degree $n$. Formal functions
and flat base change identify the finite-length $H^1$ of a resolution
with its completed version. Proposition~\ref{D:cover-uniform-canonical-defect}
therefore gives $h^1(Z,\mathcal O_Z)\le c_0 n$ with fixed $c_0$.
For every effective exceptional Cartier cycle $D$ the sequence
\[
 0\longrightarrow\mathcal O_Z(-D)\longrightarrow\mathcal O_Z
 \longrightarrow\mathcal O_D\longrightarrow0
\]
gives a surjection $H^1(Z,\mathcal O_Z)\to H^1(D,\mathcal O_D)$:
$H^2(Z,\mathcal O_Z(-D))=0$ because the proper map to the affine
local surface has fibres of dimension at most one. Hence
$h^1(\mathcal O_{tC})\le c_0n$.

The previously proved exceptional arithmetic calculation gives
\[
 C^2=nA_0^2,\qquad
 -n\sum_i a_i\le K_Z\cdot C\le nK_Y\cdot A_0
\]
after contracting the point-mapped rational $(-1)$-curves, as in
Propositions~\ref{D:exceptional-arithmetic-genus-bound} and
\ref{D:punctured-cover-linear-resolution}. The cycle adjunction formula
now gives
\[
 \chi(\mathcal O_{tC})
 =-\tfrac12\bigl(t^2C^2+tK_Z\cdot C\bigr)=O_t(n).
\]
Since $tC$ is proper of dimension one,
$h^0(\mathcal O_{tC})=\chi(\mathcal O_{tC})+h^1(\mathcal O_{tC})$.
This proves (D.CD.7). All assertions are insensitive to completing the
algebraic or strict-henselian germ, by the same proper base-change and
formal-functions argument.
\end{proof}

\begin{remark}[Scope of the remaining counting problem]
The first proposition applies to arbitrary complete equicharacteristic
normal local surface domains with algebraically closed coefficient field.
The second uses the previously specified algebraic or strict-henselian
resolution setup for its intersection estimates. These are uniform
annihilator and cohomology bounds. They do not provide a presentation of
cover-dependent module objects with $O(nr^2)$ parameters, and no hereditary
quadratic representation-count or finite-generation conclusion is inferred.
\end{remark}


% Complete root/B/D actual-file audits and exact Noether source checks PASS.
\begin{proposition}[Proved: first-power Noether differents on all punctured covers]
\label{D:cover-uniform-Noether-different}
Use the complete equicharacteristic surface ring $R$, the fixed regular
Noether normalizations $S_i\subset R$, and the fixed ideal
$J\subset\sum_i\mathfrak N(R/S_i)$ from
Proposition~\ref{D:cover-uniform-canonical-defect}. For every finite normal
$R$-torsion-free algebra $A$ etale on the puncture, one has
\[
 \mathfrak N(R/S_i)A\subset\mathfrak N(A/S_i)
 \quad\hbox{for every }i,\qquad
 JA\subset\sum_i\mathfrak N(A/S_i).                     \tag{D.ND.1}
\]
Consequently, multiplication by every $z\in J$ on every maximal
Cohen--Macaulay $A$-module factors through a finite free $A$-module.
In particular $J\operatorname{Ext}^1_A(M,N)=0$ for every such source $M$
and every finite target $N$.

The same factorization holds after a finite local Artinian coefficient
extension $T/k$ for any finite $A_T=A\otimes_kT$-module $M_T$ which is
$T$-flat and whose special fibre is maximal Cohen--Macaulay. Thus this
first-power ideal also supplies the relative source-only annihilator
and prescribed-jet lifting input of
Lemma~\ref{D:relative-Noether-different-jets}.
\end{proposition}
\begin{proof}
First suppose $A$ is a domain, and fix $i$. The canonical-defect proof,
including its specified generic trace identifications, established
\[
 z\omega_{A/S_i}\subset A\tau_{A/S_i}
 \quad(z\in\mathfrak N(R/S_i)),                        \tag{D.ND.2}
\]
where $\tau$ is the trace section and is nonzero generically.
Define an $A$-linear map
\[
 \varphi_z:\omega_{A/S_i}\longrightarrow A,
 \qquad \alpha\longmapsto z\alpha/\tau_{A/S_i}.
\]
It takes values in $A$ by (D.ND.2), and
$\varphi_z(\tau_{A/S_i})=z$. Since $A$ is finite free over $S_i$,
the evaluation description of the Noether different applies and gives
$z\in\mathfrak N(A/S_i)$
\citep[Tag 0BVT, Lemma 49.6.7 and its complete proof]
{C_StacksCanonicalDefect}. This proves (D.ND.1); finite normal products
are treated componentwise.

Write an element of $\mathfrak N(A/S_i)$ as the product
$\mu(\eta)$ of a separability tensor
$\eta=\sum_\alpha a_\alpha\otimes b_\alpha\in A\otimes_{S_i}A$
annihilated by all diagonal differences. If $M$ is maximal
Cohen--Macaulay over $A$, then it is finite free over the regular ring
$S_i$. The maps
\[
 M\longrightarrow A\otimes_{S_i}M\longrightarrow M,
 \qquad m\longmapsto\sum_\alpha a_\alpha\otimes b_\alpha m,
 \qquad a\otimes m\longmapsto am
\]
are $A$-linear and have composite $\mu(\eta)\operatorname{id}_M$.
The middle term is finite free over $A$. Sums of these factorizations
prove the assertion for $J$. Lifting the second map from the finite free
middle module through an arbitrary surjection with target $M$ proves the
stated Ext annihilation.

For the relative assertion, put $S_{i,T}=S_i\otimes_kT$. The module
$M_T$ is finite over this complete local ring. It is in fact free over
$S_{i,T}$. Choose lifts of a basis of the free $S_i$-module $M_T/\mathfrak m_T M_T$;
they give a surjection $S_{i,T}^a\to M_T$ by nilpotent Nakayama. In its
kernel sequence, tensoring with $k=T/\mathfrak m_T$ remains exact on
the left because $M_T$ is $T$-flat. The kernel has zero reduction, hence
vanishes by the same Nakayama argument. Base-change the actual tensors
$\eta$ to $T$. The preceding two maps then factor multiplication through
the finite free module $A_T\otimes_{S_{i,T}}M_T$.
No flatness condition on a target module is needed for this source
factorization. Once a chosen parameter pair is regular on the target,
the two successive one-parameter divisions in the already proved
prescribed-jet lemma apply unchanged.
\end{proof}

\begin{remark}
This refines the ideal $I^3$ in
Lemma~\ref{D:cover-uniform-MCM-annihilator} in the present
equicharacteristic, perfect-residue-field setting. It does not replace
an element of a Noether different by an invertible separability
idempotent on a finite Artin quotient. In particular it supplies no
Morita equivalence and no bound with $O(nr^2)$ object coordinates.
\end{remark}

\begin{lemma}[Proved: a weighted genus estimate from geometric genus]
\label{D:exceptional-weighted-genus}
Let $Z\to\operatorname{Spec}A$ be a projective regular resolution of a
normal local surface with algebraically closed residue field. Let
$E_1,\ldots,E_v$ be distinct exceptional integral curves, put
$g_i=p_a(E_i)$ and $b_i=-E_i^2>0$, and suppose
$p_g=\dim_kH^1(Z,\mathcal O_Z)<\infty$. Then
\[
                  \sum_{i=1}^{v}\frac{g_i^2}{b_i}
                  \le 4p_g+2v.                         \tag{D.WG.1}
\]
In the punctured-cover setting of
Proposition~\ref{D:punctured-cover-linear-resolution}, this sum is
$O(n)$ with a constant depending only on the fixed base data.
\end{lemma}
\begin{proof}
For every nonzero effective exceptional cycle $D$, the vanishing of
$H^2(Z,\mathcal O_Z(-D))$ gives
$H^1(Z,\mathcal O_Z)\twoheadrightarrow H^1(D,\mathcal O_D)$.
Since $h^0(\mathcal O_D)\ge1$, it follows that $p_a(D)\le p_g$.
For $D=\sum_i t_iE_i$, cycle adjunction gives
\[
 p_a(D)=1+\sum_i\left[t_i(g_i-1)
               -\frac{b_i t_i(t_i-1)}2\right]
               +\sum_{i<j}t_it_j(E_i\cdot E_j).          \tag{D.WG.2}
\]
All terms in the last sum are nonnegative.

Set $a_i=(g_i-1)_+$. For each $g_i>1$, choose a nearest integer
$t_i$ to $a_i/b_i+1/2$; set $t_i=0$ for $g_i\le1$.
The maximum on the real line of
$q_i(t)=a_it-b_it(t-1)/2$ is
\[
 \frac{a_i^2}{2b_i}+\frac{a_i}{2}+\frac{b_i}{8}.
\]
Rounding changes it by at most $b_i/8$, so
$q_i(t_i)\ge a_i^2/(2b_i)$. If any $a_i>0$, this defines a nonzero
cycle, and (D.WG.2) implies
$\sum_i a_i^2/b_i\le2p_g$. If all $a_i=0$, that inequality is
automatic. Finally, for $g_i\ge1$,
$g_i^2\le2a_i^2+2$, and $b_i\ge1$ by integrality. The terms with
$g_i=0$ vanish, so (D.WG.1) follows. The last assertion combines the
degree-linear bound for $p_g$ in
Proposition~\ref{D:cover-uniform-canonical-defect} with the
degree-linear number of components in
Proposition~\ref{D:punctured-cover-linear-resolution}.
\end{proof}


% Separate post-CP15 research. Not assembled pending independent audits.
\begin{proposition}[Proved: linear-size Eisenstein charts and pointed isomorphism witnesses]
\label{D:local-Eisenstein-linear-atlas}
Let $k$ be an algebraically closed field of arbitrary characteristic.
For integers $e\ge2$ and $d\ge e-1$, put
\[
                  N=\left\lfloor\frac{2d}{e}\right\rfloor+1.
                                                               \tag{D.EA.1}
\]
There is an explicitly defined affine $k$-scheme $\mathcal E_{e,d}$
with the following properties; it is allowed to be empty.
\begin{enumerate}
\item Its coordinates are at most $e(N-1)\le2d$ polynomial
coefficients and at most two auxiliary inverses. Its equations are
linear coefficient vanishings and degree-two inverse equations.
For every $k$-algebra $R$, a point determines a monic Eisenstein polynomial
\[
 f(X)=X^e+\sum_{i=1}^{e}a_i(t)X^{e-i},\qquad
 a_i(t)=\sum_{j=1}^{N-1}a_{ij}t^j.                       \tag{D.EA.2}
\]
The algebra $B_f=R[[t]][X]/(f)$ is finite free of rank $e$ over
$R[[t]]$, is canonically a pointed formal disc $R[[\alpha]]$ over $R$,
and satisfies
\[
 t=\alpha^e u,\qquad f'(\alpha)=\alpha^d v,
 \qquad u,v\in R[[\alpha]]^\times.                       \tag{D.EA.3}
\]
In particular $B_f[1/t]$ is finite etale over $R((t))$, and its
discriminant ideal in $R[[t]]$ is $(t^d)$.
\item Every separable extension of $k((t))$ of degree $e$ and
different exponent $d$, with the valuation on its integer ring
normalized to have value group $\mathbb Z$, is represented by a
$k$-point of this chart. The same assertion holds after every
algebraically closed extension of $k$.
\item Over $\mathcal E_{e,d}\times\mathcal E_{e,d}$ there is an affine
finite-presentation scheme $\mathcal W_{e,d}$ using, in total,
at most $7d+6$ coordinates and equations of degree at most $e^2+1$.
For every $k$-algebra $R$ and every pair $(f,g)$ over $R$, its fibre
has an $R$-point if and only if there exists a continuous
$R[[t]]$-algebra isomorphism $B_g\simeq B_f$ preserving the
specified augmentations to $R$. Thus the witness is effective on
arbitrary, including nonreduced, coefficient rings.
\end{enumerate}
Here a pointed formal disc means $R[[\alpha]]$ with its augmentation
$\alpha\mapsto0$ and its $\alpha$-adic topology. The assertion about
families is deliberately about augmentation-preserving isomorphisms;
it makes no assertion that arbitrary nilpotent families of unpointed
Laurent-series algebras have this presentation. At geometric points
the distinction does not affect the represented local field extensions.
The case $e=1,d=0$ is the single trivial extension and needs no parameters.
\end{proposition}
\begin{proof}
We give the equations and the lifting argument, so that the precision
and the family scope can be checked directly. The usual Eisenstein
coefficient spaces and their fixed-discriminant strata are treated in
\citet[Section 4, particularly Proposition 4.4, Corollary 4.6 and
Lemma 4.7]{Yasuda2025EisensteinB}. That source also distinguishes the
relative etaleness condition from merely having separable field fibres.
Our smaller coefficient cutoff follows from the explicit Newton
argument below; we do not use a local-field counting formula over a
finite residue field.

\emph{The coefficient stratum.}
Always require $a_{e1}$ to be a unit. If the integer $e$ is nonzero
in $k$, the chart is empty unless $d=e-1$. In the latter case no
additional derivative conditions are imposed. Suppose instead that
$e=0$ in $k$. For every $1\le i<e$ for which $e-i$ is nonzero in $k$,
assign the coefficient $a_{ij}$ the derivative weight
\[
                         ej+e-i-1.                    \tag{D.EA.4}
\]
Set it equal to zero whenever this weight is less than $d$.
There must be a pair $(i_*,j_*)$ of these indices with weight exactly
$d$; if there is none the chart is empty. There is at most one such
pair, because the residues $e-i-1$ modulo $e$ are distinct. Require
$a_{i_*j_*}$ to be a unit. These are the asserted equations, with
one inverse for each of the two required units. The required index
$j_*$, when it exists, is less than $N$: its weight is at least
$ej_*$, whereas $eN>2d\ge d$.

The same valuation argument includes the characteristic-zero and
tame cases. In fact the derivative of the leading monomial has
weight $e-1$ when it is nonzero. All other nonzero derivative terms
have weights (D.EA.4), with distinct residues modulo $e$. Thus the
least weight cannot cancel, and its coefficient is a unit under
exactly the stated conditions. In particular these equations encode
the specified different exponent on every field fibre without
introducing resultant equations of growing length.

\emph{Formal discs over arbitrary coefficient rings.}
Write $f(X)=X^e+tU(t,X)$. The constant term $U(0,0)=a_{e1}$ is
invertible. The formal implicit-function recursion, using this unit,
produces a unique series $t=T(\alpha)\in\alpha^eR[[\alpha]]$ solving
$f(\alpha)=0$; its coefficient of $\alpha^e$ is a unit.
Conversely, in the finite free, $t$-adically complete algebra $B_f$,
the image of $U(t,X)$ is a unit: modulo $t$ its constant term is a
unit and $X$ is nilpotent. Hence $t=-\alpha^e/U(t,\alpha)$ in $B_f$.
The $t$- and $\alpha$-adic topologies are cofinal. The two resulting
continuous maps identify $B_f$ with $R[[\alpha]]$ and prove the
first identity in (D.EA.3). This argument uses no reducedness or
Noetherian hypothesis on $R$.

Substitution of $t=\alpha^e u$ in the derivative shows that each
nonzero term has precisely its weight as an $\alpha$-power, times
a coefficient in $R[[\alpha]]$. There is one least term, of weight
$d$, and its coefficient is a unit. This proves the second identity
in (D.EA.3) over $R$ itself. Thus the derivative becomes a unit after
inverting $t$, proving the asserted etaleness. The standard monogenic
trace determinant is the norm of the derivative, up to sign. The norm
of $\alpha$ is $(-1)^e a_e(t)$, which is $t$ times a unit, and the norm
of a unit is a unit. The discriminant is consequently $t^d$ times
a unit, again over arbitrary $R$.

\emph{The elementary Newton estimate.}
We will use a version which works over every coefficient ring.
In $R[[\alpha]]$, let $h$ be a polynomial and let $\beta$ satisfy
\[
 h'(\beta)=\alpha^d v\quad(v\text{ a unit}),\qquad
 h(\beta)\in\alpha^aR[[\alpha]],\qquad a>2d.
                                                               \tag{D.EA.5}
\]
Division of $h(\beta)$ by $h'(\beta)$ is defined, because
$\alpha$ is a nonzerodivisor. Set
$\beta_1=\beta-h(\beta)/h'(\beta)$. Taylor's polynomial identity
gives
\[
 \beta_1-\beta\in\alpha^{a-d}R[[\alpha]],\qquad
 h(\beta_1)\in\alpha^{2(a-d)}R[[\alpha]].               \tag{D.EA.6}
\]
The derivative changes by an element divisible by $\alpha^{a-d}$,
so its quotient by $\alpha^d$ is still a unit. Iteration is therefore
defined; the excess of the error order over $2d$ at least doubles
at every step. Completeness yields a root $\beta_\infty$, with
$\beta_\infty-\beta\in\alpha^{a-d}R[[\alpha]]$.
This proof uses polynomial Taylor identities, not division by
factorials, and works in every characteristic. It is functorial in
the coefficient ring.

\emph{Coverage of local fields.}
Let $L/k((t))$ have the degree and different in the statement.
Its integer ring is a complete discrete valuation ring with residue
field $k$, so the extension is totally ramified and a uniformizer
$\alpha$ has an Eisenstein minimal polynomial $f_\infty$ of degree
$e$. Its derivative has value $d$ for the valuation normalized by
$v_L(\alpha)=1$. The distinct derivative weights just used show
that its coefficients satisfy the same lowest-weight conditions.
Truncate them modulo $t^N$, obtaining the polynomial $g$ of
(D.EA.2). Terms removed have weights at least $eN>d$, so $g$ lies
in $\mathcal E_{e,d}(k)$. In the original integer ring,
\[
 g(\alpha)\in(t^N)=(\alpha^{eN}),\qquad
 g'(\alpha)=\alpha^d\cdot\text{unit},\qquad eN>2d.
\]
The Newton estimate supplies a root of $g$ in $L$. The polynomial
$g$ is Eisenstein of the same degree $e$, so the corresponding
embedding $k((t))[X]/(g)\hookrightarrow L$ is an isomorphism.
There are $e(N-1)=e\lfloor2d/e\rfloor\le2d$ retained coefficients.
The argument works unchanged after an algebraically closed field
extension. The different/derivative identity and the field-valued
Eisenstein description also appear explicitly in
\citet[Section 2 and equation (4.1)]{Yasuda2025EisensteinB}.

\emph{Finite equations for pointed isomorphisms.}
For two polynomials $f,g$ in the chart, work in the free rank-$e$
algebra
\[
 B_f/t^N B_f=(R[t]/t^N)[\alpha]/(f).
\]
Introduce $eN$ coefficients $b_{ij}$ and the candidate element
\[
 \beta=\sum_{i=0}^{e-1}\sum_{j=0}^{N-1}b_{ij}t^j\alpha^i.
                                                               \tag{D.EA.7}
\]
Impose $b_{00}=0$, make $b_{10}$ invertible by one inverse variable,
and impose the $eN$ coefficient equations
\[
                         g(\beta)=0\pmod{(t^N,f)}.     \tag{D.EA.8}
\]
Use the same polynomial representative (D.EA.7) in $B_f$. Its image
in $R[[\alpha]]$ is $\alpha$ times a unit, since $e\ge2$.
The derivative-weight equations for $g$ therefore give
$g'(\beta)=\alpha^d$ times a unit; this assertion remains true
although $\beta$ is not yet a root. Equations (D.EA.8) give
$g(\beta)\in\alpha^{eN}R[[\alpha]]$. Newton's estimate now gives
an exact root $\beta_\infty$. Since $eN-d>d\ge1$, the correction
does not change its invertible linear coefficient. Thus substitution
$X\mapsto\beta_\infty$ gives an augmentation-preserving
$R[[t]]$-algebra isomorphism $B_g\simeq B_f$: a power series with
zero constant term and invertible linear coefficient is an
automorphism of the pointed formal disc over any $R$.

Conversely, such an isomorphism sends the source uniformizer to an
element $\sum_{i<e}b_i(t)\alpha^i$ of $B_f$. Preservation of the
augmentation says $b_0(0)=0$; invertibility on the conormal module
says $b_1(0)$ is a unit. Truncating every $b_i(t)$ modulo $t^N$
produces (D.EA.7) and (D.EA.8). This proves the equivalence on all
$R$-points. We assert existence, not that each approximate witness
is a unique isomorphism or that its Newton lift retains the entire
$t^N$-jet.

Finally, two coefficient charts use at most $4d+4$ coordinates,
and (D.EA.7) with its inverse uses at most $2d+e+1$. Since
$e\le d+1$, the total is at most $7d+6$. To expand (D.EA.8), the
largest initial $\alpha$-power is $e(e-1)$. Replacing
$\alpha^e$ by $-\sum a_i(t)\alpha^{e-i}$ decreases that power by
at least one and adds at most one to coefficient degree at each
step. Starting with degree at most $e+1$ gives the safe bound
$e^2+1$ for every resulting coefficient equation. All other
equations have degree at most two. The number of equations is
$O(d+e)$, from the displayed coefficient and unit conditions.
\end{proof}

\begin{corollary}[Proved: polynomial-complexity charts for bounded local branch data]
\label{D:local-branch-data-linear-complexity}
Fix a finite ordered collection of nontrivial totally ramified branches
over algebraically closed residue fields, with degrees $e_j\ge2$ and
different exponents $d_j$. Put $D=\sum_jd_j$ and suppose $e_j\le n$.
The product of the coefficient charts above uses at most $2D$
coefficient parameters, together with $O(D)$ inverse variables.
The product pointed-isomorphism relation has an existential witness
using $O(D)$ coordinates and equations of degree at most $n^2+1$.
If $D\le cn$, allowing all such ordered degree/different lists and
all their groupings into branch fibres requires at most
$n^{O_c(n)}$ charts, still with $O_c(n)$ coordinates and polynomial
degree in $n$. Unramified branches of degree one add no parameters.
Every counted branch fibre is required to contain a nontrivial branch;
arbitrarily many additional all-unramified marked fibres are not included
in this combinatorial count.
These are charts for local branch data, not a Hurwitz atlas for covers
of a proper curve and not a bound for covers of a surface puncture.
\end{corollary}
\begin{proof}
Add the coordinate bounds in the proposition. Since $d_j\ge e_j-1$,
the number of nontrivial branches is at most $D$ and
$\sum e_j\le2D$. An ordered list of positive integers $d_j$ with
sum at most $D$ has at most $2^{D+1}$ possibilities. For each such
list, choose $2\le e_j\le d_j+1$; a crude upper bound for the
number of choices is $(D+1)^D$. Placing separators to group a list
into branch fibres contributes at most $2^D$. Trivial branches
in each fibre, if total covering degree $n$ is specified, are
then determined by subtracting its nontrivial degrees from $n$.
The same bounds with an additional factor $(n+1)^{D+1}$ allow an
arbitrary bounded total degree instead. These elementary bounds
are $n^{O_c(n)}$ when $D\le cn$; the cases $n=1$ or $D=0$ are
trivial and may be treated separately. Products of the witness
equations retain their degree bound. No global gluing or effective
algebraization statement is needed for these local assertions.
\end{proof}

\begin{remark}[The remaining global step]
To count finite covers of one fixed proper curve, the local data must
still be glued to a cover of its complement. The preceding charts
do not bound the degree of that gluing problem or the number of
global covers having specified formal local extensions. Likewise,
the vanishing of a local deformation obstruction would not by
itself bound the number of connected components of a global
Hurwitz space. Those are additional quantitative gates, retained
explicitly here.
\end{remark}


% Root/B/D full actual-file and complete exact primary-source audits PASS.
\begin{proposition}[Proved: rank-linear total degree on a rational double point]
\label{D:RDP-full-sheaf-degree}
Let $(A,\mathfrak m,k)$ be an algebraic henselian or complete normal
surface local ring over an algebraically closed field, and suppose that
it is a rational double point. Let $\pi:Y\to\operatorname{Spec}A$ be
its minimal resolution, with exceptional curves $E_i$ and fundamental
cycle $Z_0=\sum_i z_iE_i$. For a reflexive $A$-module $M$ of rank $r$,
let
\[
 \mathcal F=\pi^*M/(\text{torsion}).
\]
Then $\mathcal F$ is a globally generated vector bundle, and
\[
 d_i:=\deg_{E_i}(\det\mathcal F)\ge0,\qquad
 \sum_i z_i d_i\le r,\qquad \sum_i d_i\le r.             \tag{D.FS.1}
\]
In particular, for any nonnegative real weights $w_i$,
\[
 \sum_i w_i d_i\le r\max_i w_i.                          \tag{D.FS.2}
\]
All characteristics are allowed, including the bad characteristics for
the individual rational double point types.
\end{proposition}
\begin{proof}
The full-sheaf assertion for algebraic henselian rational surface
singularities, in every characteristic, is
\citet[Section 2.1 and Proposition 2.7 with its full proof,
pp.~4 and 7--8 of the author version]{C_GustavsenIle2018}.
We spell out its duality argument to include the complete case.
The torsion-free sheaf $\mathcal F$ is globally generated because it is
a quotient of $\pi^*M$. Its pushforward equals $M$: both inclusions are
checked on the puncture and then extended by reflexivity of $M$.
A global generating surjection $\mathcal O_Y^a\to\mathcal F$, together
with rationality and the vanishing of $H^2$ for a proper morphism to the
affine local surface with fibres of dimension at most one, gives
$R^1\pi_*\mathcal F=0$. Duality therefore gives
\[
 R\pi_*R\mathcal Hom_Y(\mathcal F,\omega_Y)
       \simeq R\operatorname{Hom}_A(M,\omega_A).
\]
The right side has no positive cohomology since $M$ is maximal
Cohen--Macaulay. The sheaves
$\mathcal Ext_Y^q(\mathcal F,\omega_Y)$, $q>0$, have finite support:
$Y$ is regular and the torsion-free sheaf is locally free in codimension
one and on the puncture. The two-column spectral sequence for the
displayed duality consequently forces all these positive Ext sheaves
to vanish, as well as $R^1\pi_*\mathcal Hom(\mathcal F,\omega_Y)$.
Over a regular surface the resulting depth condition says that
$\mathcal F$ is locally free. This proof uses dualizing modules and
proper duality, so applies unchanged to a complete excellent local
rational double point; it does not use ordinary absolute differentials
over its coefficient field.

For clarity, the cycle facts needed below have an elementary
intersection proof. The curves $E_i$ are smooth rational $(-2)$-curves
with connected negative-definite intersection matrix. Start with one
$E_i$ and repeatedly add an $E_j$ for which the current positive cycle
$D$ satisfies $D\cdot E_j>0$. This procedure stays below any fixed
strictly positive anti-nef exceptional cycle: if the coefficient at $j$
already equalled that of such a bounding cycle, the inequality would
be impossible. It terminates at the minimal positive anti-nef cycle
$Z_0$, whose support is the whole connected exceptional set.
At every step $D^2=-2$. Indeed, if $D\cdot E_j\ge2$, then
$(D+E_j)^2\ge0$, contradicting negative definiteness; hence the positive
intersection is one and the square stays $-2$. Thus
\[
 Z_0^2=-2,\qquad K_Y\cdot Z_0=0,
 \qquad\chi(\mathcal O_{Z_0})=1.                       \tag{D.FS.3}
\]
Rationality and $H^2$-vanishing give
$H^1(\mathcal O_{Z_0})=0$, so $H^0(\mathcal O_{Z_0})=k$.
The sequence for the effective Cartier cycle $Z_0$ now gives
$H^1(\mathcal O_Y(-Z_0))=0$.

Every element of $\mathfrak m$ vanishes on $Z_0$. To check this, for a
nonzero such element write its pullback divisor as $D+H$, where $D$ is
exceptional and $H$ has no exceptional component. The cycle $D$ is
positive on each $E_i$, and
$D\cdot E_i=-H\cdot E_i\le0$. Hence $D\ge Z_0$ by the defining
minimality of $Z_0$. Consequently
$\mathfrak m\mathcal O_Y\subset\mathcal O_Y(-Z_0)$.

Twisting a global generating surjection onto $\mathcal F$ by $-Z_0$
and again using $H^2$-vanishing shows
$H^1(\mathcal F(-Z_0))=0$. Therefore
\[
 M=H^0(Y,\mathcal F)\twoheadrightarrow
 H^0(Z_0,\mathcal F|_{Z_0}),
\]
and the target is killed by $\mathfrak m$. It has dimension at most
$\mu_A(M)$, the minimum number of generators of $M$.
A rational double point has multiplicity two. Over the infinite
residue field choose a two-generated parameter reduction $Q$ of
$\mathfrak m$. The maximal Cohen--Macaulay multiplicity formula gives
\[
 \mu_A(M)\le\operatorname{length}_A(M/QM)=2r.            \tag{D.FS.4}
\]
Equivalently, a generic line through its quadratic tangent cone gives
a parameter ideal with $\operatorname{length}(A/Q)=2$; the same rank
formula then proves (D.FS.4).

The restriction $\mathcal F|_{Z_0}$ has no $H^1$, since
$H^1(Y,\mathcal F)=0$ and the next term is $H^2(\mathcal F(-Z_0))=0$.
The vector-bundle Riemann--Roch formula on the Cartier cycle $Z_0$ and
(D.FS.3) therefore yield
\[
 r+c_1(\mathcal F)\cdot Z_0
 =\chi(\mathcal F|_{Z_0})
 =h^0(\mathcal F|_{Z_0})\le2r.
\]
Each $d_i$ is nonnegative because $\mathcal F$ is globally generated,
and every $z_i\ge1$. This proves (D.FS.1), and (D.FS.2) follows.
\end{proof}

\begin{remark}[What this estimate does and does not control]
On an $A_{N-1}$ chain one can see the estimate directly, including when
the characteristic divides $N$. The character module represented by
$(u,w^a)$ in $uv=w^N$, $1\le a<N$, has invertible full pullback on the
standard crepant charts: the local generators are $u$ before the $a$-th
component and $w^a$ afterwards. Its valuation sequence is
$c_i=\min(i,a)$, with $c_0=0,c_N=a$, so its component degrees are
$2c_i-c_{i-1}-c_{i+1}$, equal to one only at $i=a$.
The proposition avoids a classification of modules and applies also
to $D$ and $E$ configurations.

For a reflexive sheaf on a partial model with rational double points,
this bounds the total degree of its full pullback on each resolved
block by its rank. It does not assert a uniform pole bound for gluing
that full pullback to independently chosen extensions on adjacent
core curves. Such compatibility, together with an effective sparse
object presentation, remains necessary for a degree-linear hereditary
representation count.
\end{remark}


% Separate refinement of the proved pointed witness; no new Laurent-family claim.
\begin{corollary}[Proved: unpointed formal-disc isomorphisms at the same precision]
\label{D:unpointed-Eisenstein-witness}
In Proposition~\ref{D:local-Eisenstein-linear-atlas}, keep the same
coefficient charts and the same integer $N$. In its isomorphism
witness omit the equation $b_{00}=0$, retaining the invertibility
of $b_{10}$ and the equations $g(\beta)=0\pmod{(t^N,f)}$.
For every $k$-algebra $R$, the modified witness has an $R$-point
over $(f,g)$ if and only if there is a continuous
$R[[t]]$-algebra isomorphism $B_g\simeq B_f$, without any
requirement on the specified augmentations to $R$.
The bounds of $7d+6$ coordinates and degree at most $e^2+1$
remain valid. This is a statement about the formal-disc models
already constructed, not a presentation theorem for all families
of finite etale $R((t))$-algebras.
\end{corollary}
\begin{proof}
Use the polynomial representative
\[
       \beta=\sum_{i<e,\,j<N}b_{ij}t^j\alpha^i\in B_f=R[[\alpha]]
\]
from (D.EA.7). The witness equations give $g(\beta)=t^N h$ for
some $h\in B_f$. Evaluation at $\alpha=0$, hence at $t=0$, gives
\[
                              b_{00}^{\,e}=0.           \tag{D.UE.1}
\]
The constant coefficient of $\beta$ as an $\alpha$-series is
$b_{00}$, and its linear coefficient is $b_{10}$, because
$t\in\alpha^eR[[\alpha]]$ and $e\ge2$.

We recall the elementary formal-coordinate fact over arbitrary
coefficient rings. A series with nilpotent constant coefficient
and invertible linear coefficient is a continuous automorphism
of a formal disc. To verify it, let $I$ be the nilpotent principal
ideal generated by its constant coefficient. The simple root
$0$ of its reduction modulo $I$ lifts, by the finite nilpotent
implicit-function recursion, to a root $a\in I$. Translation by
$a$ is a continuous formal-disc automorphism; its substitution
is well defined because $a$ is nilpotent. After that translation
the series has zero constant coefficient and invertible linear
coefficient, and the usual coefficient-by-coefficient inverse
applies. This proves the assertion, including continuity of the
inverse. Applied to (D.UE.1), it identifies $B_f$ with $R[[\beta]]$.

One must still justify that the Newton estimate uses the correct
parameter. Write $g(X)=X^e+tU_g(t,X)$. Its constant
$U_g(0,0)$ is invertible. Modulo $t$ the equation $g(\beta)=t^Nh$
says $\beta^e=0$, so $U_g(t,\beta)$ is a unit in $B_f$.
Since $N\ge2$, we obtain
\[
       \beta^e=t\bigl(-U_g(t,\beta)+t^{N-1}h\bigr)
                    =t\cdot\text{unit}.               \tag{D.UE.2}
\]
The $t$- and $\beta$-adic topologies are consequently cofinal.
The coefficient-weight conditions for $g$, now applied using
(D.UE.2), give
\[
                 g'(\beta)=\beta^d\cdot\text{unit},
       \qquad g(\beta)\in\beta^{eN}R[[\beta]],
       \qquad eN>2d.                                  \tag{D.UE.3}
\]
The all-coefficient-ring Newton argument in the proposition
therefore supplies an exact root $\beta_\infty$ with
\[
           \beta_\infty-\beta\in
                    \beta^{eN-d}R[[\beta]].            \tag{D.UE.4}
\]
Here $eN-d>d\ge1$, so $\beta_\infty$ is a pointed invertible
change of coordinate relative to $\beta$. Composing with the
already established, possibly unpointed, coordinate $\beta$
shows that $X\mapsto\beta_\infty$ defines a continuous
$R[[t]]$-algebra isomorphism $B_g\to B_f$.

Conversely, let such an isomorphism be given and let $\beta$ be
the image of the source coordinate. The source identity
$t=X^e\cdot\text{unit}$ implies that the constant coefficient
of $\beta$ is nilpotent. The inverse isomorphism is also
continuous and given by substitution of its source-coordinate
image. The chain rule for these mutually inverse substitutions
shows that the linear coefficient of $\beta$ is a unit: the
derivative of the inverse evaluated at the nilpotent constant
is defined and multiplies that coefficient to $1$.
Expand $\beta$ in the free $R[[t]]$-basis
$1,\alpha,\ldots,\alpha^{e-1}$ and truncate its coefficient
series modulo $t^N$. This retains its constant and linear
coefficients, gives the invertible $b_{10}$, and preserves
$g(\beta)=0$ modulo $t^N$. It is the required finite witness.
No additional variables or higher-degree equations have been
introduced.
\end{proof}


% Separate post-CP15 candidate. A fibre bound, not an effective Hurwitz atlas.
\begin{proposition}[Proved: a quantitative fibre bound for fixed formal curve covers]
\label{D:fixed-formal-curve-fibre-bound}
Let $k$ be an algebraically closed field, let $E/k$ be a smooth proper
connected curve of genus $h$, and let $B\subset E(k)$ be finite. Fix
an integer $n\ge1$ and, at each $b\in B$, a finite free
$\widehat{\mathcal O}_{E,b}$-algebra $A_b$ of rank $n$ which is a
product of complete discrete valuation rings with residue field $k$
and whose generic algebra is etale. Write $d_b$ for its discriminant
exponent and put $D=\sum_{b\in B}d_b$.
Then the number of isomorphism classes over $E$ of finite generically
etale maps $f:C\to E$ of degree $n$, with $C/k$ smooth proper connected, which
are etale over $E\setminus B$ and satisfy
\[
       \prod_{x\in f^{-1}(b)}\widehat{\mathcal O}_{C,x}
          \simeq A_b
       \quad\text{as }\widehat{\mathcal O}_{E,b}\text{-algebras}
                                                               \tag{D.FC.1}
\]
for every $b$, is at most
\[
                        (n!)^{2h}\,n^{2D}.             \tag{D.FC.2}
\]
For $n=1$ the right side is understood as $1$. In particular, for
fixed $E$ and $D\le cn$, this is $n^{O_{h,c}(n)}$.
No Galois hypothesis, restriction on $p\mid n!$, lower bound on $n$
relative to the source genus, or finiteness hypothesis on $k$ is imposed.
The formal algebras are fixed in full, including their tame and
unramified factors.
\end{proposition}
\begin{proof}
The argument supplies a common characteristic-zero discriminant for
all the covers in one formal-data fibre. The closely related
\citet[Proposition 4.2, pp.\ 4--6]{DeJongKatzCurvesB} treats maps to
$\mathbb P^1$ with the same discriminant and formally isomorphic wild
ramification, under its hypotheses $g(C)\ge2$ and
$\deg(f)\ge2g(C)-1$. Its proof sketches the formal-patching step.
Here we fix every local factor, choose its lift directly, and give
that patching step explicitly. The resulting characteristic-zero
discriminant need not be reduced; this is why the permutation-defect
count below is useful.

\emph{The characteristic-zero count.}
First suppose $k$ has characteristic zero. For a cover in the
statement its discriminant divisor is the fixed divisor
$\Delta=\sum_b d_b[b]$. Let $b_1,\ldots,b_s$ be its support.
The fundamental group of $E\setminus\operatorname{Supp}(\Delta)$
has generators
\[
 x_1,y_1,\ldots,x_h,y_h,\sigma_1,\ldots,\sigma_s,
 \qquad \prod_i[x_i,y_i]\prod_j\sigma_j=1,             \tag{D.FC.3}
\]
where $\sigma_j$ generates inertia at $b_j$; this is the
characteristic-zero case, with its full proof, of
\citet[Expos\'e XIII, Corollary 2.12, pp.\ 392--394]{SGA1}.
That proof also handles arbitrary algebraically closed fields of
characteristic zero by change of algebraically closed field, so no
embedding of our entire field into $\mathbb C$ is assumed.

After labelling a geometric fibre, the cover gives a permutation
representation in $S_n$. If the inertia permutation at $b_j$ has
cycle lengths $e_1,\ldots,e_a$, its discriminant multiplicity is
\[
                    d_{b_j}=\sum_{i=1}^a(e_i-1)=n-a.   \tag{D.FC.4}
\]
Every permutation with index $n-a=d_{b_j}$ is a product of exactly
$d_{b_j}$ transpositions: decompose each of its cycles in this way.
There are at most $\binom n2^{d_{b_j}}\le n^{2d_{b_j}}$
such products. The $2h$ handle generators have at most $(n!)^{2h}$
possible images. Multiplying these estimates, and discarding both
the relation and transitivity when making an upper bound, gives
(D.FC.2). Quotienting labelled representations by conjugacy can
only decrease their number. The same bound therefore counts even
possibly disconnected smooth covers with this discriminant.

\emph{One lift of the target and of each prescribed local algebra.}
Assume now $\operatorname{char}(k)=p>0$, and put $R=W(k)$,
$K=\operatorname{Frac}(R)$. The ring $R$ is a complete discrete
valuation ring with residue field $k$. Choose a smooth proper lift
$\mathcal E/R$ of $E$. Existence follows from the unobstructed
deformation theory of smooth proper curves, followed by
algebraization; an ample line bundle lifts since $H^2(E,\mathcal O_E)=0$.
Equivalently one can use the smoothness statements in
\citet[Tags 0DZQ and 0DZT, Lemmas 93.16.4 and 109.15.1--2]{StacksFormalCurveB}.
This choice is made once, independently of $f$.
Lift the finitely many points $b$ to disjoint sections of
$\mathcal E/R$. Smoothness, henselian lifting and completeness
allow the choices, and their formal neighbourhoods have coordinates
$\widehat{\mathcal E}_{b_R}\simeq\operatorname{Spf}(R[[t_b]])$
lifting the chosen coordinates $\widehat{\mathcal O}_{E,b}=k[[t_b]]$.

For a nontrivial factor of $A_b$, choose a uniformizer and its
monic Eisenstein polynomial over $k[[t_b]]$. Lift every coefficient
to $R[[t_b]]$, still making all nonleading coefficients divisible
by $t_b$ and the constant coefficient divided by $t_b$ a unit.
If desired, Proposition~\ref{D:local-Eisenstein-linear-atlas}
first replaces the polynomial by a finite one representing the
same local extension. The lifted algebra is finite free over
$R[[t_b]]$. It is isomorphic to $R[[s]]$: writing the equation as
$s^e+t_b U(t_b,s)=0$, the constant $U(0,0)$ is a unit, so formal
implicit substitution solves uniquely for $t_b$ as a series in
$s$. The $t_b$- and $s$-adic topologies are cofinal. A degree-one
factor is lifted by the identity algebra. Thus we have fixed
algebras
\[
        \mathcal A_b=\prod_{j}R[[s_{bj}]]
           \quad\text{finite free of rank }n\text{ over }R[[t_b]],
        \qquad \mathcal A_b/p\mathcal A_b\simeq A_b.   \tag{D.FC.5}
\]
These choices need not preserve any local Galois group.

Let $\delta_b(t_b)$ be their trace discriminant. It reduces to
$t_b^{d_b}$ times a unit modulo $p$. In particular
$\mathcal A_b/p^m\mathcal A_b$ becomes finite etale after inverting
$t_b$, for every $m\ge1$: its discriminant is a unit in
$(R/p^m)((t_b))$. Weierstrass preparation also shows that
$R[[t_b]]/(\delta_b)$ is finite free over $R$ of rank $d_b$,
with the zero algebra allowed when $d_b=0$.

\emph{Patching at every nilpotent level.}
Fix any special-fibre cover $f:C\to E$ counted in the statement,
and fix isomorphisms (D.FC.1). Put $R_m=R/p^mR$,
$\mathcal E_m=\mathcal E\times_R R_m$, and let $U_m$ be the
complement of the chosen sections. The given finite etale cover
over $U_1$ has compatible finite etale lifts over all $U_m$.
Moreover every prescribed isomorphism on the reduction lifts
uniquely. This is the equivalence of finite-etale categories for
a nilpotent thickening, not an assertion about a characteristic-zero
punctured disc; see \citet[Tag 0BQB, Lemma 58.8.3]{StacksFormalCurveB}.

On each formal disc use the fixed finite free algebra
$\mathcal A_b/p^m\mathcal A_b$. On its puncture it and the
restriction of the lifted cover of $U_m$ are finite etale.
The chosen special-fibre identification therefore has a unique
lift to an isomorphism of these two algebras on
$R_m((t_b))$. The uniqueness makes these overlap isomorphisms
compatible in $m$.

For completeness, this datum really gives a finite locally free
algebra on $\mathcal E_m$, even when $R_m$ is nonreduced.
Near a branch section choose an affine neighbourhood on which
the section is cut out by one nonzerodivisor. Its complement is
already in $U_m$; its completion is the formal disc just used.
Beauville--Laszlo gluing for this nonzerodivisor gives an
equivalence of finite-projective modules between the affine
neighbourhood and its open/completed gluing data. The exact
statements, with complete proofs, are
\citet[Tag 0BP2 and Tag 0BNI, Theorem 15.92.16 and
Lemmas 15.92.18--19]{StacksFormalCurveB}.
Multiplication and unit descend by full faithfulness, because
tensor products have the corresponding base changes. Their
commutativity, associativity and unit identities can be checked
on the two pieces and consequently hold after gluing. The
resulting module has rank $n$ on both pieces and hence rank $n$
on the neighbourhood. Different branch neighbourhoods meet
inside $U_m$, so the resulting algebras glue there as well.
This constructs compatible finite locally free maps
\[
                     f_m:\mathcal C_m\longrightarrow\mathcal E_m
                                                               \tag{D.FC.6}
\]
with special fibre $f$ and the specified completed local algebras.

Grothendieck existence for finite covers of a proper scheme
algebraizes (D.FC.6) to a finite map
$\mathcal C\to\mathcal E$, proper over $R$; the finite-algebra
form is \citet[Tag 09ZT, Lemma 30.28.2]{StacksFormalCurveB}.
It is finite locally free of rank $n$. Indeed its coherent algebra
is locally free on the formal completion along the closed fibre,
so it is locally free at that fibre by faithful flatness of local
completion. Its nonfree locus is closed and proper over
$\operatorname{Spec}(R)$. Any nonempty such locus meets the closed
fibre, so it is empty. The source is therefore flat over $R$.
Its special fibre is the smooth curve $C$, which makes the
source smooth near that fibre. The nonsmooth locus is again
closed and proper over $R$, and the same argument proves that
$\mathcal C/R$ is smooth everywhere. We have thus constructed
a smooth proper lift of every $f$, using the same local lifts
(D.FC.5), with no bound on $n$ relative to $g(C)$.

\emph{One common global discriminant.}
The trace determinant of $\mathcal C\to\mathcal E$ defines an
effective Cartier divisor $\mathcal\Delta$ on $\mathcal E$.
It has no vertical component, since its reduction is the
nonzero discriminant of the generically etale map $f$.
It is proper and has finite fibres over $R$, hence is finite;
the absence of a vertical component makes it $R$-flat.
Its degree is therefore
\[
                     \deg_R(\mathcal\Delta)=\sum_b d_b=D.         \tag{D.FC.7}
\]
At every nilpotent level its ideal is the unit ideal on $U_m$
and is the fixed ideal $(\delta_b(t_b))$ on the formal disc at
each $b$. The full faithfulness in the preceding gluing step
therefore makes this ideal independent of $f$ at every level.
Full faithfulness of coherent formal algebraization then makes
$\mathcal\Delta$ itself independent of $f$.
If no special-fibre cover exists there is nothing to prove;
otherwise any one of the constructed covers supplies this common
algebraic divisor. In particular we do not assume in advance
that arbitrary formal divisors algebraize on the chosen target.

\emph{Geometric generic fibres distinguish special covers.}
Choose one constructed lift for each special-fibre isomorphism
class. Suppose two such lifts become isomorphic over
$\mathcal E_{\overline K}$, where $\overline K$ is an algebraic
closure of $K$. This isomorphism is defined over a finite
extension $K'/K$. Let $R'$ be the integral closure of $R$ in
$K'$; it is a finite complete discrete valuation ring, and its
residue field is still $k$. Both source models after base change
to $R'$ are smooth, hence regular and normal, and finite locally
free over $\mathcal E_{R'}$. Each is the normalization of the
target in its generic finite algebra: an element integral over
the target is integral over the finite source, and so belongs
to it by normality. Uniqueness of normalization extends the
generic isomorphism to the two entire finite covers over $R'$.
Reducing gives an isomorphism of the original covers over $E$.
Thus different special-fibre classes give different geometric
generic classes. This argument works with normal finite
algebras as well as with fields, so no hidden connectedness
claim after geometric generic base change is needed.

Finally, all geometric generic covers have target
$\mathcal E_{\overline K}$ of genus $h$ and the same discriminant
$\mathcal\Delta_{\overline K}$ of degree $D$. Apply the
characteristic-zero count (D.FC.3)--(D.FC.4). It proves
(D.FC.2), including the cases $B=\varnothing$ or $D=0$.
For fixed $h$ and $D\le cn$, use $n!\le n^n$ to obtain
the last assertion.
\end{proof}

\begin{remark}[Exact quantitative scope]
The proposition bounds the number of global covers in one fibre of
the map to complete local ramification data. In conjunction with
Proposition~\ref{D:local-Eisenstein-linear-atlas}, it separates a
local coefficient space of dimension $O(D)$ from a global fibre
of cardinality $n^{O(n+D)}$. It does not bound the degree of a
finite-type global gluing presentation, its equations, or the
complexity of an additional locus on which a cover extends over
a surface neighbourhood. A fibre-cardinality estimate alone
does not control the number of points in such a constrained
zero-dimensional locus. Those effective gluing and extension
conditions remain unproved here; in particular no new conclusion
about the general finite-generation questions D.1 or D.2 is made.
\end{remark}


% Root/B/D complete actual-file audits PASS; auxiliary, no object-atlas assertion.
\begin{proposition}[Proved: canonical energy on a minimal surface resolution]
\label{D:exceptional-canonical-energy}
Let $A$ be a normal Noetherian local $k$-algebra of dimension two,
with residue field the algebraically closed field $k$, and suppose that
$A$ admits a dualizing complex. Let
$\pi:Z\to\operatorname{Spec}A$ be a projective regular resolution. Suppose
that $p_g=\dim_kH^1(Z,\mathcal O_Z)<\infty$ and that no exceptional
curve is a smooth rational curve of self-intersection $-1$.
Write $E_1,\ldots,E_v$ for the exceptional integral curves, and set
\[
 b_i=-E_i^2,\quad g_i=p_a(E_i),\quad
 N=(-(E_i\cdot E_j))_{i,j},\quad
 \kappa_i=K_Z\cdot E_i=2g_i-2+b_i.
\]
When the exceptional set is empty, all the following sums and energies
are zero. Otherwise $N$ is positive definite with nonpositive
off-diagonal entries, $\kappa_i\ge0$, and
\[
 \mathcal E:=\kappa^{\mathsf t}N^{-1}\kappa
       \le 8p_g+2\sum_i b_i.                           \tag{D.CE.1}
\]
Let $\mathcal B$ be any union of exceptional components with
$\kappa_i=0$, and let $\mathcal C$ be its complement. Every member of
$\mathcal B$ is a smooth rational $(-2)$-curve. For nonempty
$\mathcal C$, form the Schur complement
\[
 S=N_{\mathcal C\mathcal C}
       -N_{\mathcal C\mathcal B}
        N_{\mathcal B\mathcal B}^{-1}N_{\mathcal B\mathcal C},
                                                               \tag{D.CE.2}
\]
with $S=N$ if $\mathcal B$ is empty. It is positive definite, has
nonpositive off-diagonal entries, and
\[
 \sum_{\substack{i\in\mathcal C\\g_i\ge1}}
       g_i^2(S^{-1})_{ii}
 \le \mathcal E,
 \qquad
 \sum_{\substack{i\in\mathcal C\\g_i\ge1}}
       \frac{g_i^2}{S_{ii}}
 \le \mathcal E.                                      \tag{D.CE.3}
\]
These are numerical assertions; they do not assume that the indicated
Schur complement has already been realized by a scheme contraction.
\end{proposition}
\begin{proof}
Negative definiteness is the usual exceptional intersection theorem.
Since an integral proper curve of arithmetic genus zero over the
algebraically closed residue field is a smooth rational curve,
minimality gives $b_i\ge2$ when $g_i=0$. Thus $\kappa_i\ge0$.
For $g_i\ge1$, we also have $\kappa_i\ge g_i$; this uses only
$b_i\ge1$ and does not assert the stronger inequality
$\kappa_i\ge2g_i$.

We recall explicitly the elementary matrix facts used below. For a
positive definite symmetric matrix with positive diagonal and
nonpositive off-diagonal entries, scale the diagonal to write it as
$D^{1/2}(1-P)D^{1/2}$ with $P$ symmetric and entrywise nonnegative.
Positive definiteness gives $\lambda_{\max}(P)<1$, and the largest
absolute eigenvalue of a nonnegative symmetric matrix is its largest
eigenvalue. Hence $(1-P)^{-1}=\sum_{m\ge0}P^m$ is entrywise
nonnegative. In particular $N^{-1}$ is entrywise nonnegative.
Put
\[
                      z_*=\frac{1}{2}N^{-1}\kappa\ge0.
\]
Choose a nearest nonnegative integer $z_i$ to each coordinate of
$z_*$, and write $\epsilon=z-z_*$, so $|\epsilon_i|\le1/2$.
For a nonzero effective exceptional cycle $D_z=\sum_i z_iE_i$,
the fibre-dimension bound gives
$H^2(Z,\mathcal O_Z(-D_z))=0$. The exact sequence for $D_z$
therefore gives a surjection
$H^1(Z,\mathcal O_Z)\to H^1(D_z,\mathcal O_{D_z})$ and hence
$p_a(D_z)\le p_g$. This is the same full, possibly nonreduced,
cycle argument as in Lemma~\ref{D:exceptional-weighted-genus}.
Cycle adjunction gives the exact identity
\[
 p_a(D_z)=1+\frac{\kappa^{\mathsf t}z-z^{\mathsf t}Nz}{2}
       =1+\frac{\mathcal E}{8}
                    -\frac{\epsilon^{\mathsf t}N\epsilon}{2}.
                                                               \tag{D.CE.4}
\]
Write $m_{ij}=E_i\cdot E_j$ for $i<j$. Positive definiteness on
the vector of all ones implies
$2\sum_{i<j}m_{ij}<\sum_i b_i$. Therefore
\[
 \epsilon^{\mathsf t}N\epsilon
 \le \frac{1}{4}\sum_i b_i+\frac{1}{2}\sum_{i<j}m_{ij}
 < \frac{1}{2}\sum_i b_i.                               \tag{D.CE.5}
\]
If $z\ne0$, (D.CE.4)--(D.CE.5) yield
$\mathcal E\le8p_g-8+2\sum_i b_i$, which implies (D.CE.1).

The rounded vector can be zero even when $\kappa\ne0$, so this
case must be handled separately. If $z=0$, then
$0\le(z_*)_i\le1/2$. Since the off-diagonal entries of $N$ are
nonpositive,
\[
 \frac{\mathcal E}{4}=z_*^{\mathsf t}Nz_*
       \le\sum_i b_i(z_*)_i^2
       \le\frac{1}{4}\sum_i b_i.
\]
This proves the stronger bound $\mathcal E\le\sum_i b_i$ in
this case, and finishes (D.CE.1).

If $\kappa_i=0$, adjunction and $b_i>0$ force $g_i=0,b_i=2$.
The Schur complement of a positive definite block is positive
definite. Its off-diagonal entries are nonpositive, because
$N_{\mathcal B\mathcal B}^{-1}$ is entrywise nonnegative and the
two rectangular off-diagonal blocks are entrywise nonpositive.
The block inverse identity gives
\[
 (N^{-1})_{\mathcal C\mathcal C}=S^{-1},\qquad
 \mathcal E=\kappa_{\mathcal C}^{\mathsf t}
                      S^{-1}\kappa_{\mathcal C}.
\]
The inverse $S^{-1}$ is entrywise nonnegative. Dropping its
off-diagonal contributions and using $\kappa_i\ge g_i$ proves
the first inequality of (D.CE.3). Finally
$(S^{-1})_{ii}\ge1/S_{ii}$: this follows from the scalar Schur
complement formula, or from Cauchy--Schwarz applied to
$S^{1/2}e_i$ and $S^{-1/2}e_i$. This proves the second inequality.
\end{proof}

\begin{corollary}[Proved: degree-linear effective genus energy on punctured covers]
\label{D:punctured-cover-linear-energy}
In the algebraic or descended strict-henselian punctured-cover setting
of Proposition~\ref{D:punctured-cover-linear-resolution}, replace the
chosen resolution by a minimal regular resolution of the covering
normal surface. There is a constant depending only on the fixed base
data such that, for every connected cover of degree $n$,
\[
 v+\sum_i b_i+\sum_i g_i+p_g+\mathcal E\le Cn.          \tag{D.CE.6}
\]
Every sum in (D.CE.3), for every choice of zero-canonical components
to eliminate, obeys the same kind of degree-linear bound.
\end{corollary}
\begin{proof}
The starting resolution has $v$, $\sum b_i$ and $\sum g_i$
bounded linearly in $n$ by
Proposition~\ref{D:punctured-cover-linear-resolution}; the canonical
defect theorem, Proposition~\ref{D:cover-uniform-canonical-defect},
gives the same bound for $p_g$, after completion if necessary.
Contract exceptional smooth rational $(-1)$-curves until none remains.
Existence and regularity of these contractions are the surface
contraction facts already used in constructing the preceding
resolution; see \citet[Lemma 54.16.8, Tag 0C2M]
{StacksDSurfaceResolution}. We do not require the resulting minimal
resolution to retain its morphism to the fixed base resolution.

For completeness, these contractions preserve all the needed
numerical bounds even when they destroy a normal-crossings
configuration. If the contracted curve meets the other exceptional
integral curves with intersection multiplicities $a_i$, their new
self-intersection and arithmetic genus are
\[
 (E_i')^2=E_i^2+a_i^2,
 \qquad p_a(E_i')=p_a(E_i)+\frac{a_i(a_i-1)}{2}.          \tag{D.CE.7}
\]
These follow from the total-transform formula
$\rho^*E_i'=E_i+a_iE$ and $K_Z=\rho^*K_{Z'}+E$.
The sum of the positive numbers $b_i$ drops by exactly
$1+\sum_i a_i^2$, while the sum of the arithmetic genera rises
by at most $\frac12\sum_i a_i^2$. Consequently the total genus
increase over the whole contraction sequence is at most half
the original $\sum b_i$. The number of components decreases,
and $p_g$ is unchanged under blowing down a regular $(-1)$-curve.
All four initial bounds therefore persist. Apply
Proposition~\ref{D:exceptional-canonical-energy} to finish.
\end{proof}

\begin{remark}[Scope of the energy estimate]
The diagonal numbers $S_{ii}$ in (D.CE.3) can be small rational
numbers after elimination of long rational blocks. The estimate
controls the corresponding genus contributions without a lower
bound on these individual numbers. It does not construct a
finite-type presentation of bundles or covers, control its
equation degrees, or prove an index-linear representation count.
The hereditary object-atlas problem remains separate.
\end{remark}


% Separate post-CP15 research. Not assembled until independent audits pass.
\begin{example}[Refuted: a constant generator bound at an Azumaya boundary]
\label{D:elliptic-cone-order-generators}
The following constant-generator assertion is false: on one fixed
normal complete local surface, if $\mathcal A=\operatorname{End}(G)$
is Azumaya off the closed point and $H=\operatorname{Hom}(G,F)$ is
locally free of rank one as a right $\mathcal A$-module there, then
$H$ has a bounded number of $\mathcal A$-generators independent of
the common rank of $F$ and $G$. It remains false when both $F$ and
$G$ are reflexive and Frobenius-periodic.

More precisely, put $k=\overline{\mathbb F}_3$ and fix
\[
 R=k[[x,y,z]]/(y^2z-x^3+xz^2).
 \tag{D.EO.1}
\]
For every $r=2^m$, $m\ge1$, there are rank-$r$ reflexive
$R$-modules $G_r,F_r$ and an integer $f_r>0$ such that
\[
 (F_R^{f_r})^{[*]}G_r\simeq G_r,
 \qquad (F_R^{f_r})^{[*]}F_r\simeq F_r.
 \tag{D.EO.2}
\]
Writing $\mathcal A_r=\operatorname{End}_R(G_r)$ and
$H_r=\operatorname{Hom}_R(G_r,F_r)$, we have
\[
 \mu_{\mathcal A_r}(H_r)=3r.
 \tag{D.EO.3}
\]
Here generators are right-module generators, with the action given
by precomposition. On the punctured spectrum, $\mathcal A_r$ is
a split Azumaya algebra of degree $r$ and $H_r$ is locally free
of rank one over it. The completed local ring in (D.EO.1) is the
same for all $r$.
\end{example}
\begin{proof}
Let $C\subset\mathbb P^2_k$ be the cubic
$y^2z-x^3+xz^2=0$, and put $L=\mathcal O_C(1)$. In characteristic
three the three partial derivatives are $z^2,2yz,y^2+2xz$.
Their simultaneous vanishing on the cubic forces $z=y=x=0$,
which is not a projective point. Thus $C$ is smooth. The plane
cubic exact sequence gives $H^0(C,\mathcal O_C)=k$ and genus one,
so $C$ is a smooth geometrically integral genus-one curve; the
point $[0:1:0]$ makes it an elliptic curve. Also $\deg L=3$.
The homogeneous ring
\[
 R^{\mathrm{gr}}=k[x,y,z]/(y^2z-x^3+xz^2)
\]
is a domain, is Cohen--Macaulay, and is regular away from its
vertex $\mathfrak m=(x,y,z)$. It is therefore normal. Its completed
local ring is $R$. Alternatively, normality of $R$ follows directly
from the same Jacobian calculation and Serre's criterion: its
associated graded ring is the domain $R^{\mathrm{gr}}$, so $R$ is
a domain, and its singular locus is only its closed point.

Write $U^{\mathrm{gr}}=\operatorname{Spec}(R^{\mathrm{gr}})
\setminus\{\mathfrak m\}$. The usual degree-zero localization
charts identify
\[
 q:U^{\mathrm{gr}}\longrightarrow C
 \quad\text{with}\quad \operatorname{Tot}(L^{-1})\setminus C,
 \qquad q_*\mathcal O_{U^{\mathrm{gr}}}
       =\bigoplus_{a\in\mathbb Z} L^a.
 \tag{D.EO.4}
\]
These identities can be checked on each of $x\ne0$, $y\ne0$,
$z\ne0$: the nonzero homogeneous coordinate trivializes the
$\mathbb G_m$-torsor. For a degree-zero line bundle $N$ on $C$,
the unique reflexive extension of $q^*N$ gives a finite rank-one
graded $R^{\mathrm{gr}}$-module
\[
 M_N=\Gamma(U^{\mathrm{gr}},q^*N)
    =\bigoplus_{a\in\mathbb Z}H^0(C,N\otimes L^a).
 \tag{D.EO.5}
\]
Coherence, uniqueness, and full faithfulness here are exactly
\citet[Tag 0EBJ, Lemma 31.13.12, including its proof]{StacksDEllipticOrder}:
the missing vertex has depth two. Formula (D.EO.4) gives the
displayed grading without an extra projective-normality assumption.
The same full faithfulness gives
\[
 \operatorname{Hom}_{R^{\mathrm{gr}}}(M_N,M_{N'})
  =\bigoplus_{a\in\mathbb Z}
       H^0(C,N'N^{-1}\otimes L^a),
 \tag{D.EO.6}
\]
where multiplication is composition of the corresponding sections.
This is the full Hom module, not only its degree-zero part.

A nontrivial degree-zero line bundle on a proper integral curve has
no nonzero section: such a section would have an effective divisor
of degree zero and hence would trivialize the bundle. Negative
degree likewise gives no sections. On our genus-one curve,
Riemann--Roch and positive-degree $H^1$-vanishing give
\[
 h^0(C,N\otimes L)=3
 \quad\text{for every }N\in\operatorname{Pic}^0(C).
 \tag{D.EO.7}
\]
For the latter vanishing and degree convention see
\citet[Tags 0E39 and 0AYQ, Lemma 53.22.1 and Definition 33.44.1]{StacksDEllipticOrder}.

Fix $r=2^m$. There is a line bundle $T\in\operatorname{Pic}^0(C)$
of exact order $2r$: multiplication by $2r$ on the elliptic curve
$\operatorname{Pic}^0(C)$ has geometric kernel
$(\mathbb Z/(2r))^2$, since $2r$ is prime to three
\citep[Tag 03RP, Proposition 39.9.11(5)--(7), with proof]{StacksDEllipticOrder}.
Set
\[
 G^{\mathrm{gr}}=\bigoplus_{i=0}^{r-1}M_{T^i},
 \qquad F^{\mathrm{gr}}=\bigoplus_{j=r}^{2r-1}M_{T^j},
 \qquad A^{\mathrm{gr}}=\operatorname{End}(G^{\mathrm{gr}}),
 \quad H^{\mathrm{gr}}=\operatorname{Hom}(G^{\mathrm{gr}},F^{\mathrm{gr}}).
 \tag{D.EO.8}
\]
Let $e_i$ denote the projection onto $M_{T^i}$ in $G^{\mathrm{gr}}$.
By (D.EO.6), all four modules in question have no negative graded
pieces, and
\[
 (A^{\mathrm{gr}})_0=\bigoplus_{i=0}^{r-1}ke_i\simeq k^r,
 \qquad (H^{\mathrm{gr}})_0=0,
 \qquad
 \dim_k (H^{\mathrm{gr}})_1e_i=3r.
 \tag{D.EO.9}
\]
Indeed, the two index sets in (D.EO.8) are disjoint modulo $2r$,
and every one of the $r$ summands contributing to the last space
has dimension three by (D.EO.7).

Complete (D.EO.8) at $\mathfrak m$ and call the completed modules
$G_r,F_r,\mathcal A_r,H_r$. Hom commutes with this flat completion
for finite presented modules, so the completed Hom and End modules
are precisely those in the statement. The completed rank-one
modules remain reflexive: they are maximal Cohen--Macaulay over
the normal two-dimensional hypersurface $R$. On its punctured
spectrum they are invertible. Thus $G_r,F_r$ are locally free of
rank $r$ there, and local choices of bases identify
$\mathcal A_r$ and $H_r$ with the matrix algebra and its regular
right module, respectively.

There is no need to identify the full Jacobson radical of
$\mathcal A_r$. Taking degree zero gives a surjective algebra map
$A^{\mathrm{gr}}\to k^r$ which kills $\mathfrak m A^{\mathrm{gr}}$.
Taking degree one gives a surjective right-module map
\[
 H^{\mathrm{gr}}\longrightarrow (H^{\mathrm{gr}})_1,
 \tag{D.EO.10}
\]
where the target is an $A^{\mathrm{gr}}$-module through $k^r$.
This is a module map because $H^{\mathrm{gr}}$ starts in degree
one. It kills $\mathfrak m H^{\mathrm{gr}}$. Consequently both
maps extend to the completions and retain their surjectivity.
If $\mathcal A_r^s\twoheadrightarrow H_r$, composition with
(D.EO.10) factors as a surjection
\[
 (k^r)^s\twoheadrightarrow (H^{\mathrm{gr}})_1.
\]
Multiplying by any $e_i$ and using (D.EO.9) gives $s\ge3r$.

In fact the lower bound is sharp. Choose a basepoint-free pencil
in $H^0(C,L)$: a nonzero section has only finitely many zeros,
and, since $L$ is globally generated over the infinite field $k$,
one can choose a second section nonzero at all those zeros.
The pencil sequence is
\[
 0\longrightarrow L^{-1}\longrightarrow\mathcal O_C^2
       \longrightarrow L\longrightarrow0.
\]
For every nontrivial degree-zero $N$, tensor by $N L^a$, $a\ge1$.
The group $H^1(C,N L^{a-1})$ vanishes: for $a=1$ this follows
from $H^0(C,N)=0$ and degree-zero Riemann--Roch, and for $a>1$
from positive degree. The resulting multiplication maps show
that $M_N$ is generated by its three-dimensional degree-one part.
All the bundles $T^{j-i}$ occurring in $H^{\mathrm{gr}}$ are
nontrivial. Each column $H^{\mathrm{gr}}e_i$ therefore has $3r$
$R^{\mathrm{gr}}$-generators. For $1\le\alpha\le3r$, combine
the $\alpha$-th chosen generator of every column into a single
matrix $h_\alpha$. Then $h_\alpha e_i$ isolates the chosen
generator of column $i$. These $3r$ matrices generate
$H^{\mathrm{gr}}$ as a right $A^{\mathrm{gr}}$-module, and still
generate after completion. Hence $s=3r$ is sufficient, proving
the exact equality (D.EO.3).

Finally choose $f_r>0$ with $3^{f_r}\equiv1\pmod{2r}$.
Absolute Frobenius pulls an invertible sheaf back to its
$3^{f_r}$-th tensor power, so $q^*T^i$ and $q^*T^j$ are fixed
by this Frobenius pullback. Reflexive extension gives (D.EO.2)
over $R^{\mathrm{gr}}$ and then over $R$. This passage is valid
because finite-module duals commute with flat completion and
Frobenius pullback commutes with the natural completion square.
In particular, the example uses finite-monodromy, prime-to-three
torsion line bundles, not arbitrary varying coherent modules.
\end{proof}

\begin{remark}[Exact scope and finite certificate]
This refutes a proposed constant-generator step for modules over
an order extending an Azumaya algebra across an omitted point.
It does not refute a rank-quadratic presentation of projective
representations, a representation-count bound, or finite generation
of a fundamental group. It also does not exclude choosing a better
anchor for a particular Brauer class: the Azumaya class in this
example is already trivial. The assertion refuted is uniformity
for arbitrary anchors satisfying the stated hypotheses.
The finite certificate
\nolinkurl{parts/D/code/check_elliptic_order_generators.py}, with
output \nolinkurl{parts/D/output/elliptic_order_generators.json},
checks the characteristic-three Jacobian identities, the index-set
and semisimple-module calculations, and Frobenius congruences for
$r=2,4,\ldots,256$. Existence of torsion line bundles and the
geometric module identifications are established by the proof and
the cited primary statements, not claimed as numerical tests.
\end{remark}


% Complete root/B/D actual-file audits PASS; separate from the frozen CP16 PDF.
\begin{lemma}[Proved: an integral anti-nef cycle with controlled energy]
\label{D:integral-antinef-energy-cycle}
Let $N$ be a symmetric positive definite integral matrix with diagonal
$b_i>0$ and nonpositive off-diagonal entries. Let $\kappa$ be an
entrywise nonnegative integral vector. Assume that
\[
 b_i=2\quad\text{if }\kappa_i=0,
 \qquad b_i\le3\kappa_i\quad\text{if }\kappa_i>0,
                                                               \tag{D.IC.1}
\]
and that each connected component of the off-diagonal incidence graph
contains a vertex with $\kappa_i>0$. Put
$B=\sum_i b_i$ and $\mathcal E=\kappa^{\mathsf t}N^{-1}\kappa$.
For every integer $c\ge2$ there is an entrywise positive integral
vector $z$ such that
\[
 \begin{split}
 &Nz\ge0,\qquad
 (Nz)_i\ge(c-3/2)\kappa_i\quad(\kappa_i>0),\\
 &\bigl|(Nz-c\kappa)_i\bigr|\le b_i/2,\qquad
 z^{\mathsf t}Nz\le2c^2\mathcal E+B.
 \end{split}
 \tag{D.IC.2}
\]
In particular
\[
       \sum_i(Nz)_i\le c\sum_i\kappa_i+B/2.            \tag{D.IC.3}
\]
The vector is integral from the outset; no common denominator of
$N^{-1}$ is used.
\end{lemma}
\begin{proof}
Consider the positive definite quadratic function
\[
     F(u)=u^{\mathsf t}Nu-2c\kappa^{\mathsf t}u
     \quad (u\in\mathbb Z^v).
\]
It has a nonempty finite set of global lattice minimizers.
Replacing a vector by its coordinatewise absolute value does not
increase its quadratic term, since the off-diagonal entries of $N$
are nonpositive. It also does not increase its linear term, since
$\kappa\ge0$. Choose a nonnegative global minimizer $z$.
Both $z+e_i$ and $z-e_i$ are competitors in the full lattice,
even when $z_i=0$. The inequalities
$F(z\pm e_i)\ge F(z)$ therefore give
\[
               |(Nz-c\kappa)_i|\le b_i/2.             \tag{D.IC.4}
\]
At a positive-$\kappa$ vertex this implies
$(Nz)_i\ge(c-3/2)\kappa_i>0$.

At a zero-$\kappa$ vertex, (D.IC.1) and integrality imply
$(Nz)_i\ge-1$. If equality holds, replace $z$ by $z+e_i$.
Then
\[
       F(z+e_i)-F(z)=2(Nz)_i+b_i=0.
\]
The new vector is another nonnegative global minimizer and still
satisfies (D.IC.4) at every vertex. Repeat whenever a negative
coordinate of $Nz$ remains. Every step increases one coordinate,
whereas the set of global minimizers is finite; hence the process
terminates. Its output satisfies all the intersection inequalities
in (D.IC.2).

It has full support. If $z_i=0$, then
$(Nz)_i=\sum_{j\ne i}N_{ij}z_j\le0$; since $Nz\ge0$,
every neighbor $j$ of $i$ must also have $z_j=0$.
A zero coordinate consequently propagates through its connected
component. The whole component cannot be zero, because it has a
positive-$\kappa$ vertex where $(Nz)_i>0$.

It remains to bound its quadratic energy. The real minimum of $F$
occurs at $u_*=cN^{-1}\kappa$ and has value $-c^2\mathcal E$.
Round its coordinates to obtain an integer vector $u$ with
$|u_i-(u_*)_i|\le1/2$. As in
Proposition~\ref{D:exceptional-canonical-energy}, positivity of $N$
on the vector of all ones gives
\[
 (u-u_*)^{\mathsf t}N(u-u_*)\le B/2.
\]
Therefore $F(z)\le F(u)\le-c^2\mathcal E+B/2$.
Write $Q=z^{\mathsf t}Nz$. Cauchy--Schwarz for the $N$-inner
product gives $\kappa^{\mathsf t}z\le\sqrt{\mathcal E Q}$, so
\[
       (\sqrt Q-c\sqrt{\mathcal E})^2\le B/2.
\]
It follows that $Q\le2c^2\mathcal E+B$. Finally (D.IC.3) follows
by summing the upper bounds in (D.IC.4).
\end{proof}

\begin{corollary}[Proved: a bounded-cohomology cycle positive on the nonrational core]
\label{D:punctured-cover-antinef-cycle}
In the setting of Corollary~\ref{D:punctured-cover-linear-energy},
suppose each connected exceptional graph contains a component with
$K_Z\cdot E_i>0$. For every fixed integer $c\ge2$ there is an
effective integral exceptional Cartier divisor $D_z$ of full support
such that $\mathcal O_Z(-D_z)$ has nonnegative degree on every
exceptional component and degree at least
$(c-3/2)(K_Z\cdot E_i)$ on each component with positive canonical
intersection. Moreover, with constants independent of the degree-$n$
cover,
\[
 -D_z^2+K_Z\cdot D_z+
       \sum_i\deg(\mathcal O_Z(-D_z)|_{E_i})=O_c(n),     \tag{D.IC.5}
\]
and, for every fixed integer $t\ge1$,
\[
 h^0(tD_z,\mathcal O_{tD_z})+
 h^1(tD_z,\mathcal O_{tD_z})=O_{c,t}(n).                \tag{D.IC.6}
\]
All components on which $\mathcal O_Z(-D_z)$ has degree zero are
smooth rational $(-2)$-curves.
\end{corollary}
\begin{proof}
For the minimal resolution put $\kappa_i=2g_i-2+b_i$ as before.
If $\kappa_i=0$, then $g_i=0,b_i=2$. If $\kappa_i>0$,
then $b_i=\kappa_i-2g_i+2\le\kappa_i+2\le3\kappa_i$.
The lemma therefore applies. On a regular surface every integral
exceptional cycle is Cartier. Its conormal degrees are exactly
$(Nz)_i$, so the asserted positivity and description of the
zero-degree components follow.

The preceding energy corollary bounds
$\mathcal E$, $\sum b_i$, $\sum g_i$ and $p_g$ linearly in $n$.
Thus (D.IC.2)--(D.IC.3) bound $Q=-D_z^2$ and the sum of
conormal degrees linearly as well. Also
$0\le K_Z\cdot D_z=\kappa^{\mathsf t}z
\le\sqrt{\mathcal E Q}=O_c(n)$. This proves (D.IC.5).

For the nonzero effective cycle $tD_z$, the same full-cycle
cohomology argument gives $h^1(\mathcal O_{tD_z})\le p_g$.
Adjunction gives
\[
 \chi(\mathcal O_{tD_z})
       =\frac{t^2Q-t\kappa^{\mathsf t}z}{2}.
\]
Consequently
$h^0+h^1=\chi+2h^1\le t^2Q/2+2p_g=O_{c,t}(n)$,
which includes the nilpotent layers of the cycle.
\end{proof}

\begin{remark}[Zero-canonical components and the remaining gate]
If a whole connected graph has $\kappa=0$, the minimizer is zero
and the full-support assertion of the lemma is inapplicable.
In the geometric situation that component is a rational double
point configuration; its fundamental cycle has energy two, as
proved in Proposition~\ref{D:RDP-full-sheaf-degree}, and can be
handled separately. No claim about an ample line bundle on a
contracted model, formal determination of bundles, or a finite-type
object presentation follows merely from (D.IC.5)--(D.IC.6).
In particular a zero-degree rational block still requires its
own compatibility and gluing argument.
\end{remark}


% Post-checkpoint-16 proof packet. Full root/B/C/D proof and primary-source audits passed.
\subsection{Projective representations with vanishing geometric Brauer class}

Throughout this subsection $k$ is algebraically closed of characteristic
$p>0$.  A constant in an $O(r)$ or $O(fr^2)$ estimate depends on the
fixed geometric pair and the fixed bundles explicitly chosen below,
and not on $r$, $f$, a Kummer denominator, or an Artinian parameter ring.
The Frobenius in a one-step construction is always the single
$p$-Frobenius.

\begin{lemma}[Proved: fixed tensor twists and boundary comparisons]
\label{D:fixed-tensor-boundary-comparison}
Let $R$ be a normal two-dimensional local domain and let $P$ be
a fixed rank-one reflexive $R$-module.  Put
\[
 I_P=\operatorname{im}(P\otimes_R P^\vee\longrightarrow R).
                                                        \tag{D.BZ.1}
\]
For every reflexive $M$, this ideal annihilates both the kernel
and the cokernel of
\[
             M\otimes_RP\longrightarrow (M\otimes_RP)^{**}.
                                                        \tag{D.BZ.2}
\]
The same formula holds over an Artinian coefficient base whenever
the source and target extensions have the depth-two Hartogs property.
It does not require reflexive hulls to commute with arbitrary base
change.

In the fixed completed surface charts and the Kummer covers of
Lemma~\ref{D:tame-boundary-summed-index}, allow a one-step
isomorphism on the punctured cover
\[
       F^{[*]}M_j\otimes^{[*]}P_n \simeq M_{j+1},
       \qquad P_n=(P\otimes_RR_n)^{**}.
                                                        \tag{D.BZ.3}
\]
The uniform inverse-matrix pole bound and the summed normalized
boundary-index estimate in that lemma remain valid.  The additional
constant depends on $P$ and the fixed chart only.  In particular,
there is no factor equal to the Kummer denominator or the number
of boundary characters.
\end{lemma}

\begin{proof}
On the puncture, write $z=\sum_\nu p_\nu\ell_\nu\in I_P$,
with $p_\nu\in P$ and $\ell_\nu\in P^\vee$.  If
$N=(M\otimes P)^{**}$, contraction by $\ell_\nu$ gives a map
$N\to M$: the map on the puncture extends by reflexivity.
Consequently
\[
             zh=\sum_\nu p_\nu(\ell_\nu h)
\]
lies in the image of (D.BZ.2).  If $h$ is in its kernel, every
contraction into the torsion-free module $M$ vanishes, and the
same identity gives $zh=0$.  This proves both assertions.
The ideal is the unit ideal at every height-one point, and hence
is either $R$ or is primary for the maximal ideal.

For the relative assertion, use the stated Hartogs equalities
$j_*M_T=M_T$ and $j_*N_T=N_T$ on the punctured relative spectrum.
They extend the contractions.  The displayed identity extends
because the target also has this Hartogs property.  These are
actual maps of relative modules; a base-change assertion about
double duals is not used.  For relative Cohen--Macaulay extensions,
the Hartogs property follows from the Cech complex of a module-regular
sequence of two parameters, and persists along Artinian coefficient
extensions.

Here is a direct denominator comparison for the second assertion.
Use the notation
$A_0=k[[x,a]]\subset R$, $A_n=k[[x,s]]$, $a=s^n$,
and $R_n=R\otimes_{A_0}A_n$ from
Lemma~\ref{D:tame-boundary-summed-index}.
Choose fixed $v\in P$ and $\ell\in P^\vee$ such that
$b=\ell(v)$ is nonzero at the generic point of every
component of $a=0$.  Such a pair exists by prime avoidance:
$P$ is free of rank one at those finitely many points.
Multiplication by $v$ and contraction by $\ell$ extend to
reflexive modules over $R_n$, and both composites are
multiplication by $b$.  Thus multiplication by $v$ has a
cokernel killed by $b$.

Compose the ordinary Frobenius map of $M_j$ with multiplication
by $v$ and (D.BZ.3).  The ordinary-to-reflexive Frobenius
defect of any $R_n$-MCM module is killed by the fixed ideal
of Lemma~\ref{D:Kummer-uniform-Frobenius-defect}; the proof
there is source-module factorization and does not require
an untwisted unit structure.  Relative Frobenius for $R_n/A_n$
contributes the fixed determinant $h(x,s^n)$ used in
(D.TC.5).  Hence, between the free $A_n$-modules of rank
$[R:A_0]r$, the resulting matrix $B_j$ satisfies
\[
 x^K h(x,s^n)b\,\operatorname{coker}(B_j)=0
                                                        \tag{D.BZ.4}
\]
after increasing a fixed $K$.  Replace $b$ by its
$A_0$-norm.  The adjugate of multiplication by $b$ shows
that this norm belongs to $bR$, and its reduction modulo
$a$ is nonzero.  It therefore supplies a fixed scalar
$b_0(x,a)\in A_0$ with $b_0(x,0)\ne0$ in (D.BZ.4).
Modulo $s^{np}$ the inverse of each of
$h(x,s^n)$ and $b_0(x,s^n)$ uses at most $p$ terms
of its geometric expansion in $s^n$.  Thus
\[
       B_j^{-1}\bmod s^{np}
            \in x^{-L}\operatorname{Mat}
                         (A_n/(s^{np}))
                                                        \tag{D.BZ.5}
\]
for one fixed $L$.

The element $v$ has character zero.  After choosing homogeneous
$A_n$-bases, the invariant Frobenius comparison is consequently
the same diagonal map with entries $a^{\lfloor pj/n\rfloor}$
as before.  The resulting diagonal quotient maps to the pure
coarse cokernel.  It is an isomorphism after inverting $x$:
there $P$ is trivialized by $v$ modulo every required power of
$a$.  Torsion-freeness over $k[[x]]$ makes this an inclusion.
Equation (D.BZ.5) bounds the reverse inclusion, exactly as in
(D.TC.4).  Intersecting one generic filtration with both
lattices bounds the \emph{sum} of all grade indices by their
total index.  Passing to the fixed normalization of $R/(a)$
multiplies that sum by its fixed rank over $k[[x]]$.

Finally, changing $v$ or a local line trivialization changes
the normalized character lattices by the divisors of fixed
rational functions.  These produce fixed comparison divisors
on the normalized boundary curves.  This accounts for their
global gluing.  We have not asserted that $P_n/sP_n$ is a
base-change-compatible reflexive hull.
\end{proof}

\begin{proposition}[Proved: cyclic tensor twists in the quadratic tame atlas]
\label{D:finite-tensor-cyclic-quadratic-atlas}
Fix the normal proper surface pair $(S,T)$ of
Theorem~\ref{D:tame-global-quadratic-representation-bound}.
Enlarge a fixed finite set $\Sigma$ so that
$S^\circ=S\setminus\Sigma$ is smooth and the reduced
original divisor $D$ is smooth with disjoint components there.
Put $U^\circ=S^\circ\setminus D$.
Fix a smooth projective resolution $\pi:Y\to S$, an
isomorphism over $S^\circ$, the exceptional cycle $Z$,
and the strict boundary curves $C_i$ used in that theorem.
Let $\mathscr L$ be a finite collection of line bundles on $Y$.

Allow the coefficient stages to satisfy, on the tame root
stack over $(S^\circ,D)$,
\[
       F^*\mathcal E_j\otimes L_j
                  \simeq\mathcal E_{j+1},
       \qquad L_j\in\mathscr L,\quad j\in\mathbf Z/f\mathbf Z.
                                                        \tag{D.BZ.6}
\]
Here the boundary weights lie in $[0,1)$ with denominators
prime to $p$, and the corresponding ordinary lattices are
extended reflexively to $S$.  There are constants, independent
of the length and weights of the cycle, such that:
\begin{enumerate}
\item minimal extensions $F_j$ on $Y$ have fixed exceptional
and saturated boundary-grade slope windows; $c_1(F_j)$ has
norm $O(r)$ in $N^1(Y)_{\mathbf R}$; and
$|\chi(Y,F_j)|=O(r)$;
\item each fixed twist, one Frobenius pullback, and each of the
fixed auxiliary sheaves in the filtered-diagram construction
has all its cohomology dimensions bounded by $O(r)$;
\item a word of $f$ tensor and boundary-digit types has
polynomial witnesses in $O(fr^2)$ variables and equations of
bounded degree.  Valid unit-cycle isomorphism classes are
open and closed in its witness locus, and their number is at
most $\exp(Cfr^2)$ after summing over the multiplicity arrays.
\end{enumerate}
The constants may depend on $\mathscr L$.  No bound for
an arbitrarily varying tensor line bundle is asserted.
\end{proposition}

\begin{proof}
Write $P_L=(\pi_*L)^{**}$ for the corresponding fixed
rank-one reflexive sheaf on $S$.  There are only finitely
many such sheaves.  On $S^\circ$ the one-step coarse
map in (D.BZ.6) extends reflexively as
\[
 G_j=\bigl(F_S^{[*]}E_j\otimes P_{L_j}\bigr)^{**}
                   \longrightarrow E_{j+1}.
                                                        \tag{D.BZ.7}
\]
It is injective and has a pure one-dimensional cokernel $T_j$.
Its generic exponents on every original divisor lie between
$0$ and $p-1$.  Purity follows from the depth lemma, since
both terms are reflexive on a normal surface.

\emph{Fixed poles, slopes, and first Chern classes.}
Minimal extensions have the fixed exceptional windows of
Langer's extension lemma
\citep[Definition~1.2 and Lemma~1.3]{C_Langer2026}.
One Frobenius pullback has fixed enlarged windows, by the
curve Frobenius bounds already used in
Proposition~\ref{D:global-periodic-reflexive-boundedness}.
Tensoring by a member of $\mathscr L$ shifts these windows
by one of finitely many fixed degrees.  The bounded-pole
argument therefore gives one $B$ for both maps
\[
 \begin{split}
 F_Y^*F_j\otimes L_j&\longrightarrow F_{j+1}(BZ),\\
 F_{j+1}&\longrightarrow
 F_Y^*F_j\otimes L_j\bigl((p-1)\textstyle\sum_iC_i+BZ\bigr).
 \end{split}                                            \tag{D.BZ.8}
\]
These maps are the genuine maps off the exceptional set;
the second is the ordinary inverse with the fixed allowed
boundary pole.

Taking maximal slopes, and then dualizing the second map,
gives
\[
 p\mu_{\max}(F_j)\le\mu_{\max}(F_{j+1})+C,\qquad
 p\mu_{\max}(F_j^\vee)\le\mu_{\max}(F_{j+1}^\vee)+C.
                                                        \tag{D.BZ.9}
\]
The same inequalities hold on a fixed smooth very ample
curve not contained in the determinant-zero divisors.
The maps restricted to that curve are generically injective
and hence injective.  Taking a maximum around the finite
cycle bounds all four kinds of slopes uniformly.

The determinant of (D.BZ.8) gives
\[
 c_1(F_{j+1})-p c_1(F_j)
                 =r c_1(L_j)+D_j,                       \tag{D.BZ.10}
\]
where $D_j$ is supported on the fixed $C_i$ and the
exceptional curves, and each coefficient has absolute
value at most $Cr$.  This follows from the two pole bounds
in (D.BZ.8); at an original boundary component it is also
the sum of the $r$ exponents in $[0,p-1]$.
After one complete cycle, (D.BZ.10) expresses $c_1(F_j)$
as a geometric weighted sum divided by $p^f-1$.
All its summands belong to a fixed finite-dimensional
span and have norm $O(r)$.  Hence $\|c_1(F_j)\|=O(r)$,
uniformly in $f$.  This argument takes place on $Y$ and
does not assume that $P_L$ is Cartier on $S$.

Normalizing (D.BZ.8) on the boundary root charts gives
exactly the two maps of
Lemma~\ref{D:saturated-boundary-two-sided-slopes}, with
the additional line twist $L_j|_{C_i}$.  At $C_i\cap Z$
this line bundle is locally trivial.  The intersection
of the localizations in its two distinct parameters
therefore proves that normalization preserves the same
exceptional pole bound, for the map and its inverse.
On a character grade the two inequalities become
\[
 p\mu_{\max}(G_{j,i,a})
       \le\mu_{\max}(G_{j+1,i,pa})+C,\qquad
 p\mu_{\max}(G_{j,i,a}^{\vee})
       \le\mu_{\max}(G_{j+1,i,pa}^{\vee})+C.
                                                        \tag{D.BZ.11}
\]
The fixed constant absorbs $\deg(L_j|_{C_i})$ and the
normal digit twist.  The pairs consisting of stage and
character form finite cycles.  Thus every saturated
grade of the full flag inside the \emph{same} $F_j|_{C_i}$
has fixed upper and lower slope bounds.

\emph{The Euler estimate, including the fixed tensor defect.}
Lemma~\ref{D:fixed-tensor-boundary-comparison} extends the
local determinant, inverse, and summed-index calculation
to (D.BZ.7).  On normalized complete boundary curves the
forward character maps have only the additional degrees
of the finitely many $P_L$ trivializations.  Taking maxima
around the stage-character cycles gives the same uniform
upper-slope estimates as (D.TC.10).  Each character is
repeated at most $p-1$ times.  The one-filtration index
calculation bounds all positive corrections together by
$Cr$.  The proof of Lemma~\ref{D:tame-canonical-Euler}
therefore gives
\[
                         \chi(S,T_j)\le Cr.              \tag{D.BZ.12}
\]

We also justify quantitatively the tensor change in the
source of (D.BZ.7).  The relative Euler calculation in the
proof of Proposition~\ref{D:global-periodic-reflexive-boundedness}
uses only fixed exceptional slope windows.  For any bundle
$A$ with those windows and
$E=(\pi_*A)^{**}$, it gives
\[
                         |\chi(Y,A)-\chi(S,E)|\le Cr.
                                                        \tag{D.BZ.13}
\]
Indeed $E=\pi_*A(BZ)$ for a fixed $B$; the quotient is
bounded by the cohomology of one fixed exceptional
thickening.  Formal functions and the same layer
vanishing bound $R^1\pi_*A$ by another fixed thickening.
Their component cohomology is $O(r)$ from the upper and
lower slope windows.  This assertion is not restricted
to periodic $E$.

Apply (D.BZ.13) first to $A=F_Y^*F_j$ and then to
$A\otimes L_j$.  Their reflexive pushforwards are
$F_S^{[*]}E_j$ and $G_j$, respectively.  Smooth-surface
Riemann--Roch and (D.BZ.10) give
\[
 \begin{split}
 \chi(S,G_j)-\chi(S,F_S^{[*]}E_j)
 &=\chi(Y,A\otimes L_j)-\chi(Y,A)+O(r)\\
 &=c_1(A)\cdot c_1(L_j)
   +\frac r2\bigl(c_1(L_j)^2-K_Yc_1(L_j)\bigr)+O(r)
   =O(r).
 \end{split}                                            \tag{D.BZ.14}
\]
No formula asserting exact multiplicativity of a
non-Cartier reflexive tensor product on $S$ is used.
Together with (D.BZ.12), normal-surface Riemann--Roch
and Frobenius scaling
\citep[Theorem~0.1]{C_Langer2026}
\citep[Theorem~4.5(2)]{C_Langer2025}
now yield
\[
       \operatorname{ch}_2(E_{j+1})
                   -p^2\operatorname{ch}_2(E_j)\le Cr.
                                                        \tag{D.BZ.15}
\]
The linear first-Chern terms are $O(r)$ by (D.BZ.10).
Cycling the inequality gives
$\operatorname{ch}_2(E_j)\ge-Cr$, with the sign shown.

The slope bounds (D.BZ.9) give $h^0(S,E_j)\le Cr$ by
the fixed-pole comparison and restriction to the fixed
very ample curve.  Proper duality bounds $h^2(S,E_j)$
by $Cr$: $\operatorname{Hom}(E_j,\omega_S)$ injects
into sections of $F_j^\vee$ tensored by a fixed line
bundle and a fixed exceptional pole, and the dual
slope bound applies.  The lower Riemann--Roch estimate
from (D.BZ.15) consequently gives
$h^1(S,E_j)\le Cr$ and $|\chi(S,E_j)|\le Cr$.
Equation (D.BZ.13) proves $|\chi(Y,F_j)|\le Cr$.
The same fixed-curve restrictions, duality, and smooth
Riemann--Roch show
$h^i(Y,F_j(t))\le C_t r$ for all fixed $t$ and $0\le i\le2$.
They also bound one Frobenius pullback and each fixed
tensor auxiliary: their first Chern characters and
second Chern characters are $O(r)$ by (D.BZ.10), the
Euler estimate, and one-step Frobenius scaling.
The threshold-transform auxiliary sheaves have the
same bounds, because there are at most $p$ transforms
and all full-flag grades have the fixed windows
in (D.BZ.11).

\emph{Polynomial witnesses and Artinian effectiveness.}
Use the monads, filtered linear curve resolutions,
and fixed-column derived Hom construction of
Theorem~\ref{D:tame-global-quadratic-representation-bound}.
Insert tensoring by $L_j$ at its one-step Frobenius
node.  This is one of finitely many fixed operations:
tensor the fixed bundle terms of its complexes by
$L_j$, and use the same coherent tensor diagrams
to identify it with each pole-twisted auxiliary.
Their total multiplicities are $O(r)$ by the preceding
cohomology bounds.  Their lengths and bundle types
are fixed.  Therefore the number of scalar variables
and equations per stage remains $O(r^2)$ and their
degrees remain bounded.

The boundary digit construction is unchanged:
digits belong to $[0,p-1]$, and the permutation is
increasing inside each equal-digit block.  Tensoring
by a line has no effect on its canonical block
splitting or on the ceiling-valuation matrix
criterion.  A fixed tensor type word adds at most
$|\mathscr L|^f$ choices.  The boundary words and
multiplicity arrays have the same $\exp(O(fr^2))$
bound as before.

For completeness, the open-class argument is valid
with specified Artinian points.  Composing the cycle
gives
\[
       F^{f*}\mathcal E_0\otimes M\simeq\mathcal E_0,
       \qquad M=\bigotimes_{j=0}^{f-1}
                         L_j^{\,p^{f-1-j}}.
                                                        \tag{D.BZ.16}
\]
On the $(p^f-1)$st root gerbe of $M$, with universal
line $\mathcal L$, this is a unit structure on
$\mathcal E_0\otimes\mathcal L$.
The gerbe is used over the smooth open and its
boundary root stack; it is not extended as an
Azumaya algebra across the singular points of $S$.
Finite etale fixed-frame torsors are invariant
under Artinian nilpotent thickenings.  A specified
unit isomorphism therefore lifts uniquely there,
as is checked on an etale atlas and its descent
diagram.  Untwisting gives the specified maps of
all ordinary stages and their boundary flags.

The fixed exceptional-window pole bound extends
these maps and their inverses to $Y$ with one $B$,
independent of Artinian length.  The layer-vanishing
proof works over an Artinian base by its
maximal-ideal filtration.  Flag preservation
extends at $C_i\cap Z$ because the relative
quotients are vector bundles on $C_i$, so its
local parameter remains a non-zero-divisor.
The actual fixed line twists commute with base
change.  If a reflexive description downstairs
is used, the first part of
Lemma~\ref{D:fixed-tensor-boundary-comparison}
gives the needed tensor comparison through
Hartogs, rather than through an assumed
base-change theorem for reflexive hulls.

Proper bounded-pole Hom schemes with inverse and
cycle equations thus represent a unit-Isom functor
of finite presentation which lifts every specified
small Artinian point.  Its image is open by the
infinitesimal smoothness criterion, as in the
preceding quadratic theorem.  Distinct unit-cycle
classes give disjoint open fibres covering the
witness scheme.  They form a finite clopen
partition, with at most as many members as
irreducible components.  The witness scheme is
an open subset of an affine scheme defined by
$O(fr^2)$ bounded-degree equations.  Bezout and
the array/word estimates prove the last assertion.
\end{proof}

\begin{theorem}[Proved: quadratic growth for Brauer-zero projective representations]
\label{D:Brauer-zero-projective-representation-growth}
Let $S$ be a normal proper connected surface over $k$,
let $T\subsetneq S$ be closed, and put $Q=Q(S,T)$.
For a representation
$\rho:Q\to\operatorname{PGL}_r(\mathbf F_q)$,
$q=p^f$, denote by $\operatorname{br}(\rho)$ the
class in $H^2(U,\mathbf G_m)$ of its induced
algebraic $\operatorname{PGL}_r$-torsor on $U=S-T$.
There is a constant $C_{S,T}$ such that
\[
 \#\left(
   \{\rho:Q\to\operatorname{PGL}_r(\mathbf F_q)
                   \mid \operatorname{br}(\rho)=0\}
                   /\operatorname{PGL}_r(\mathbf F_q)
    \right)
               \le q^{C_{S,T}r^2}                       \tag{D.BZ.17}
\]
for every $r,f\ge1$.  The scalar centre is not required
to have bounded order.
\end{theorem}

\begin{proof}
Enlarge the fixed finite omitted set to $\Sigma$
as in the preceding proposition.  The resulting
original-boundary tame group $Q^\circ$ surjects
onto $Q$.  Conjugacy classes inject under this
surjection, and a zero Brauer class stays zero
on $U^\circ$.  It suffices to count for the
smooth $U^\circ$.

\emph{A fixed finite set of scalar digit twists.}
Restriction $\operatorname{Pic}(Y)\to
\operatorname{Pic}(U^\circ)$ is surjective:
extend a Cartier divisor on the smooth open
by its closure, which is Cartier on the
regular surface $Y$.  The Neron--Severi
group of the smooth projective $Y$ is
finitely generated
\citep[Expose XIII, Theorem~5.1 and proof,
printed pp.~650--651]{SGA6BrauerZeroD}.
The group of $k$-points of $\operatorname{Pic}^0(Y)$
is divisible by every integer prime to $p$:
its reduced connected group is an abelian
variety and multiplication by such an integer
is surjective
\citep[Proposition~39.9.11(3)--(4), Tag 03RP;
Lemma~39.9.8 and proof, Tag 0BFG]{StacksBrauerZeroD}.
Choose line bundles
$L_1,\ldots,L_b$ representing generators
of the Neron--Severi group, including its
finite torsion generators.  For every
$n$ prime to $p$, every class in
$\operatorname{Pic}(U^\circ)/n$ has a
representative
\[
          M_{\mathbf a}=\bigotimes_{i=1}^b L_i^{a_i}
                         \bigm|_{U^\circ},
               \qquad 0\le a_i<n.                       \tag{D.BZ.18}
\]
Only the smooth projective Picard group is used
here; no identical assertion about a complete
local Picard group is being assumed.

Put $n=q-1$.  The finite central extension
\[
 1\longrightarrow\mathbf F_q^\times
  \longrightarrow\operatorname{GL}_r(\mathbf F_q)
  \longrightarrow\operatorname{PGL}_r(\mathbf F_q)
  \longrightarrow1
                                                        \tag{D.BZ.19}
\]
gives the lifting gerbe of $\rho$.
Identify $\mathbf F_q^\times$ with $\mu_n(k)$.
Its class maps to $\operatorname{br}(\rho)$.
The Kummer sequence
\citep[Lemma~59.28.1, Tag 03PK]{StacksBrauerZeroD}
therefore makes the
lifting gerbe, when this class is zero,
bandedly isomorphic to the root gerbe
$\sqrt[n]{M_{\mathbf a}}$ for some
(D.BZ.18).  The tautological lift is a
finite $\mathbf F_q$ local system of
scalar weight one on this gerbe.
The root-gerbe untwisting equivalence
\citep[Sections~2.2.2 and 3.1.1;
Lemmas~3.1.1.7--3.1.1.8 and
3.1.1.11--3.1.1.12]{Lieblich2008BrauerZeroD}
writes its vector bundle as
$E_0\otimes\mathcal L$, where
$\mathcal L^n=M_{\mathbf a}$.
Its $q$-unit structure is precisely
\[
                 F^{f*}E_0\otimes M_{\mathbf a}
                                  \simeq E_0.           \tag{D.BZ.20}
\]

The root gerbe of $M_{\mathbf a}$ extends
over $S^\circ$, including its original
boundary DVRs.  The projective representation
is tame there.  Its wild pro-$p$ inertia
can therefore act on the lift only by
the prime-to-$p$ scalar band, and must act
trivially.  Hence the lifted unit bundle
has a tame parabolic extension on the
boundary root stack.  Conversely a
weight-one unit bundle with this extension
projectivizes to a representation of
$Q^\circ$.  No tameness is imposed at
the additional exceptional valuations.

\emph{Replacing the growing scalar root by fixed digits.}
For each coordinate of $\mathbf a$ define
the cyclic residues and digits by
\[
 a_{i,j+1}=p a_{i,j}-n c_{i,j},\qquad
 0\le a_{i,j}<n,\quad 0\le c_{i,j}\le p-1.
                                                        \tag{D.BZ.21}
\]
Since $p^f\equiv1\pmod n$, the residues
are cyclic.  Let
$L_j=\bigotimes_i L_i^{c_{i,j}}$.
These bundles belong to the one fixed
finite collection
$\mathscr L=\{\bigotimes_iL_i^{c_i}:0\le c_i<p\}$.
Writing
$b_{i,j}=\lfloor p^j a_{i,0}/n\rfloor$
and
$E_j=F^{j*}E_0\otimes\bigotimes_iL_i^{b_{i,j}}$
on the smooth open gives
\[
              F^*E_j\otimes L_j\simeq E_{j+1},
                         \qquad E_f=E_0.               \tag{D.BZ.22}
\]
Their canonical parabolic lattices extend
reflexively to $S$.  Thus every representation
being counted supplies data of the preceding
proposition.

Conversely a word of the fixed scalar digits
determines
$a_{i,0}=\sum_{j=0}^{f-1}p^{f-1-j}c_{i,j}$.
Discard the all-$(p-1)$ word in any coordinate
if it gives $a_{i,0}=n$; no chosen representative
in (D.BZ.18) is lost.  A valid scalar word
and valid boundary word reconstruct the
root gerbe and a weight-one unit bundle
from (D.BZ.22), by (D.BZ.16).
Its projectivization gives the desired
finite-field projective representation.
Choices of the root-gerbe identification
or intermediate edge isomorphisms can only
increase the count of unit-cycle classes.
There are at most $p^{bf}=q^b$ scalar words,
and no equation of degree $q$ or $n$ has
been inserted into the one-step atlas.

For clarity, the finite-class comparison
has no hidden injectivity assertion for
the band of an arbitrary gerbe.  A
fundamental group of a root gerbe may
have a kernel which is only an image of
$\mu_n$.  If a weight-one finite unit
bundle exists, its projective torsor
$P$ identifies that gerbe with the
gerbe of lifts of $P$ through (D.BZ.19).
For the finite monodromy image $\Gamma$, now
let $P$ denote the connected $\Gamma$-torsor
reduction determined by that image, rather
than the full projective-frame torsor.
For its inverse image
$\widetilde\Gamma$ in
$\operatorname{GL}_r(\mathbf F_q)$,
the gerbe is the quotient
$[P/\widetilde\Gamma]$ and its tame
fundamental group is
\[
                  \widetilde\Gamma
                          \times_\Gamma Q^\circ.
                                                        \tag{D.BZ.23}
\]
Indeed its finite covers are exactly
finite covers of $P$ with
$\widetilde\Gamma$-equivariant descent;
the same descent identifies the displayed
pullback group.  Wild inertia lifts
uniquely through the prime-to-$p$
central kernel, so quotienting by its
normal closure gives exactly (D.BZ.23).

A finite central extension of a small
profinite group is small.  For a finite
target, fix its finite image $H$, the
image of the central kernel, and the
restriction to that kernel.  There
are finitely many induced maps from
the quotient to $H$ modulo that central
image.  Two lifts with the same data
differ by a character of the quotient
to that finite central image, again
with only finitely many possibilities.
Theorem~\ref{D:fixed-boundary-small}
therefore supplies the appropriate
finite-class interpretation whenever
the root gerbe has such an object;
empty root classes contribute nothing.
The quantitative open-fibre argument
of the preceding proposition is
stronger and already gives the needed
component count.

For every scalar and boundary word that
argument bounds the unit-cycle classes
by $\exp(Cfr^2)$, including all monad
and flag arrays.  Multiplying by
$q^b$ and the already bounded boundary
type choices only changes $C$.
Projectivization and the initial
surjection to $Q$ prove (D.BZ.17).
\end{proof}

\begin{remark}[Exact scope of the reduction]
This controls projective representations whose geometric
Brauer class vanishes, without a bound on their scalar
centre.  Their obstruction in $H^2(U^\circ,\mu_{q-1})$
may be a nonzero Kummer class; an ordinary finite-field
linear lift is not assumed.  The earlier bounded-centre
argument killed finitely many coefficient orders on one
fixed open subgroup.  Finiteness of $H^2$ separately at
each order does not supply that subgroup uniformly as
$q-1$ grows.  The finite-digit construction supplies this
additional uniformity for the Kummer classes.
It does not control the nonzero Brauer classes,
construct an Azumaya extension across the boundary, or
prove finite generation of the full fundamental group.
Complete local punctures are not included by the
projective Picard argument above.
\end{remark}


% Root/B/D full actual-file audits and independent exact certificate runs PASS.
% Separate post-CP16 packet. No hereditary representation atlas is asserted.
\begin{lemma}[Proved: unframed Hom vanishing for full RDP sheaves]
\label{D:RDP-full-Hom-vanishing}
Let $\pi:Y\to\operatorname{Spec}A$ be the minimal resolution of an
algebraic henselian or complete rational double point over an
algebraically closed field. Let $M,N$ be reflexive $A$-modules of ranks
$r,s$, with full sheaves
$\mathcal F=\pi^*M/(\mathrm{torsion})$ and
$\mathcal G=\pi^*N/(\mathrm{torsion})$. Then
\[
 H^1(Y,\mathcal Hom(\mathcal F,\mathcal G))=0,
 \qquad
 H^0(Y,\mathcal Hom(\mathcal F,\mathcal G))
                  =\operatorname{Hom}_A(M,N).          \tag{D.HF.1}
\]
For any finite closed subscheme $T\subset Y$,
\[
 \dim_k H^1(Y,\mathcal Hom(\mathcal F,\mathcal G)\otimes\mathcal I_T)
                    \le rs\operatorname{length}_k(T). \tag{D.HF.2}
\]
This finite-subscheme assertion does not replace a framing condition
along a whole strict-transform boundary curve.
\end{lemma}
\begin{proof}
Proposition~\ref{D:RDP-full-sheaf-degree}, including its complete-base
duality proof, gives globally generated vector bundles $\mathcal F$
and $\mathcal G$ with $R\pi_*\mathcal F=M$. A rational double point
is Gorenstein and its minimal resolution is crepant, so
$\omega_Y\simeq\mathcal O_Y$ after choosing a generator of
$\omega_A$. Indeed its canonical discrepancy divisor is supported
on the exceptional set, has zero intersection with every $(-2)$-curve,
and is therefore zero by negative definiteness.
The same proper duality calculation gives
\[
 R\pi_*(\mathcal F^\vee)
            \simeq R\operatorname{Hom}_A(M,A)=M^\vee.
\]
The last equality uses the maximal Cohen--Macaulay property of $M$
over the Gorenstein ring $A$. Thus $H^1(Y,\mathcal F^\vee)=0$.
Choose a finite global generating surjection
$\mathcal O_Y^a\twoheadrightarrow\mathcal G$ and tensor its kernel
sequence by $\mathcal F^\vee$. Since all coherent $H^2$ on $Y$
vanish for this proper map with fibres of dimension at most one over
an affine base, its long exact sequence gives the first equality of
(D.HF.1). The equality of global Hom follows either by restriction
and reflexive extension on the puncture, or by pulling an $A$-linear
map through the two torsion-free strict transforms.

Finally tensor the exact ideal sequence of $T$ by the locally free
sheaf $\mathcal Hom(\mathcal F,\mathcal G)$. The vanishing just
proved identifies its framed $H^1$ with a quotient of
$H^0(T,\mathcal Hom(\mathcal F,\mathcal G)|_T)$, which has
dimension $rs\operatorname{length}_k(T)$.
\end{proof}

\begin{proposition}[Proved: an exact two-end framing bound on an A-chain]
\label{D:A-chain-two-end-Hom-bound}
Let $k$ be algebraically closed, let $N\ge2$, and put
\[
       A=k[[u,v,w]]/(uv-w^N).
\]
Let $\pi:Y\to\operatorname{Spec}A$ be the minimal resolution.
Let $B_-$ and $B_+$ be the strict transforms of the two reduced
coordinate branches $(v,w)$ and $(u,w)$, respectively, and put
$B=B_-+B_+$. For any reflexive modules $M,N'$ of ranks $r,s$,
with full sheaves $\mathcal F,\mathcal G$, one has
\[
 \dim_kH^1(Y,\mathcal Hom(\mathcal F,\mathcal G)(-B))
                         \le2rs.                      \tag{D.HF.3}
\]
This bound is independent of the chain length and of the
characteristic, including when it divides $N$.

More precisely, set $I_0=A$ and $I_a=(u,w^a)$ for $1\le a<N$.
There are decompositions
\[
 M\simeq\bigoplus_{a=0}^{N-1}I_a^{r_a},\qquad
 N'\simeq\bigoplus_{a=0}^{N-1}I_a^{s_a},
\]
and the dimension in (D.HF.3) is exactly
\[
                          rs+r s_0-r_0s.               \tag{D.HF.4}
\]
This is a classification and calculation for individual modules;
no decomposition of arbitrary Artinian families is asserted.
\end{proposition}
\begin{proof}
\emph{Character decomposition, including the nonreduced group case.}
Write $T=k[[x,y]]$, with
$u=x^N,v=y^N,w=xy$. The action of $\mu_N$ of weights $1,-1$
has invariant ring $A$. On the complement of the origin it is an
fppf $\mu_N$-torsor: on $D(u)$ one adjoins an $N$th root of the
unit $u$, and on $D(v)$ an $N$th root of $v$.
This description does not require that $N$ be invertible.

Pull $M$ back to this punctured regular cover and extend it
reflexively across the origin. The resulting module $P$ over $T$
is free. Its equivariance extends as well. To see this also when
$\mu_N$ is nonreduced, the matrices of the equivariance and inverse
are defined on the puncture after the finite flat coefficient
extension by $k[\mu_N]$. The parameter pair $x,y$ remains regular
there, so Hartogs extends every entry. The identities extend in
the same way after a further tensor by $k[\mu_N]$.

A $\mu_N$-equivariant module is graded by $\mathbb Z/N\mathbb Z$.
Explicitly, coefficient projections onto the basis $t^j$ of
$k[t]/(t^N-1)$ are orthogonal idempotents by the coaction axioms.
Thus its weight pieces split and taking a fixed weight is exact,
even when the group scheme is nonreduced. This is also the direct
weight-space argument in \citet[Theorem 12.12 and Remark 12.13,
pp.~234--235]{C_MilneAG2022}.
Choose homogeneous lifts of a basis of $P/(x,y)P$.
Nakayama and freeness give a homogeneous free basis of $P$.
Taking invariants therefore decomposes the descended module into
rank-one character modules. These modules are reflexive: they are
direct summands over $k[[u,v]]$ of the free module $T$, and hence
have depth two. They agree with $M$ on the puncture and thus agree
globally by reflexivity.
For $1\le a<N$, the character module generated by
$x^{N-a},y^a$ becomes $I_a$ upon multiplying by $x^a$.
The trivial character gives $I_0$. This proves the decompositions.

\emph{The two end charts.}
Use the standard crepant charts indexed by $0\le i<N$:
\[
 u=X^{i+1}Y^i,\qquad
 v=X^{N-i-1}Y^{N-i},\qquad w=XY.                       \tag{D.HF.5}
\]
They give the minimal $A_{N-1}$ resolution in every characteristic,
as used in Proposition~\ref{D:punctured-cover-linear-resolution}.
The negative-end branch $B_-$ is $Y=0$ in chart $i=0$, with
parameter $u=X$. The positive-end branch $B_+$ is $X=0$ in chart
$i=N-1$, with parameter $v=Y$. They are disjoint, and each maps
isomorphically to its smooth coordinate branch in $\operatorname{Spec}A$.
For $a>0$, the full line bundle $\mathcal F_a$ of $I_a$ is generated
by $u$ on the first chart and by $w^a$ on the last chart; for $a=0$
use the generator $1$ on both charts.

The global Hom modules are the following fractional ideals:
\[
 \operatorname{Hom}_A(I_a,I_b)=
 \begin{cases}
 A,&a=b,\\
 (u,w^{b-a}),&b>a,\\
 (1,v/w^{a-b}),&b<a.
 \end{cases}                                           \tag{D.HF.6}
\]
These equalities can be checked at height one: both sides are
reflexive rank-one modules, and their valuations are the difference
of those of $I_b$ and $I_a$ at the branch $(u,w)$, and zero at the
other height-one primes. Equivalently one checks the inclusions
by multiplying the two displayed generators of $I_a$ and then
uses reflexivity. The first and second rows include $a=0$.

Trivialize the restriction of $\mathcal Hom(\mathcal F_a,\mathcal F_b)$
to $B_-$ and $B_+$ using the ratios of the preceding local generators.
The two-coordinate restriction map from (D.HF.6) then has the
following image inside $k[[u]]\oplus k[[v]]$:
\[
\begin{array}{c|c|c}
 (a,b)&\text{image}&\text{cokernel dimension}\\ \hline
 a=b&\{(f,g):f(0)=g(0)\}&1\\
 a=0<b&k[[u]]\oplus k[[v]]&0\\
 b=0<a&(u)\oplus(v)&2\\
 0<a<b&(u)\oplus k[[v]]&1\\
 0<b<a&k[[u]]\oplus(v)&1.
\end{array}                                             \tag{D.HF.7}
\]
Here is the evaluation explicitly. For $0<a<b$, the generators
$u,w^{b-a}$ restrict to $(u,0)$ and $(0,1)$ after division by the
two local Hom generators. For $0<b<a$, the generators
$1,v/w^{a-b}$ restrict to $(1,0)$ and $(0,v)$.
For $a=0<b$, the generators $u,w^b$ give $(1,0),(0,1)$;
for $b=0<a$, the generators $1,v/w^a$ give $(u,0),(0,v)$.
The coordinate projections $A\to k[[u]],k[[v]]$ are surjective,
so these generate exactly the displayed images.
When $a=b$, the map is the ordinary restriction of $A$ to its
two branches, whose image consists of pairs with the same constant.

The divisor sequence for $B$, tensored by the Hom line bundle,
and Lemma~\ref{D:RDP-full-Hom-vanishing} identify the cokernel in
(D.HF.7) with $H^1$ of the Hom bundle twisted by $-B$.
Summing over all pairs of character summands proves
\[
 \sum_{a,b}r_a s_b\bigl(1+\mathbf1_{a>0,b=0}
                              -\mathbf1_{a=0,b>0}\bigr)
            =rs+r s_0-r_0s,
\]
and this is at most $2rs$.
\end{proof}

\begin{example}[Refuted: a framing bound using only ranks and attachment count]
\label{D:A-chain-internal-framing-obstruction}
There is no bound for the boundary-framed $H^1$ in
(D.HF.3) depending only on the two ranks and the number of transverse
attachment branches if arbitrary internal attachments are allowed.
For every $m\ge1$, put $N=3m$ and use the complete $A_{N-1}$
singularity above. Let $B$ be the strict transform of the smooth
closed formal curve
\[
                    u=w^m,\qquad v=w^{2m}.             \tag{D.HF.8}
\]
For $M=I_m$, $N'=I_{2m}$, both of rank one, its framed group has
\[
 \dim_kH^1\bigl(Y,\mathcal Hom(\mathcal F_m,
                                  \mathcal F_{2m})(-B)\bigr)=m.
                                                               \tag{D.HF.9}
\]
The same phenomenon occurs at fixed characteristic two with both
reflexive modules Frobenius-periodic, by taking any odd $m$.
\end{example}
\begin{proof}
The quotient of $A$ by the two equations in (D.HF.8) is exactly
$k[[w]]$, so the curve is smooth. In chart $i=m-1$ of (D.HF.5)
its strict transform is $Y=1$; it meets $E_m$ transversely at a
point away from the nodes and has local parameter $w=X$ there.
Both full sheaves are generated by $u$ on this chart. Thus their
Hom bundle is trivially generated along $B$.
By (D.HF.6) its global Hom is $I_m=(u,w^m)$. Evaluation on $B$
has image exactly $(w^m)$ in $k[[w]]$. The unframed vanishing
therefore identifies the framed $H^1$ with $k[[w]]/(w^m)$,
proving (D.HF.9).

If $\operatorname{char}(k)=2$ and $m$ is odd, $N=3m$ is prime
to two. Some power $q=2^f$ is one modulo $N$. Under the
character decomposition, reflexive Frobenius pullback multiplies
each character by $q$, so both $I_m$ and $I_{2m}$ are fixed by
that reflexive pullback. No periodicity of their full bundles on
the resolution itself is claimed.

This example concerns the proposed attachment-only estimate, not
finite generation or a hereditary representation-count bound.
The preceding degree-linear total-chain-length estimate might
still absorb such costs if attachment multiplicities and contacts
are controlled. That additional summation and the effective
gluing presentation are not proved here.
\end{proof}

\begin{remark}[Finite exact certificate]
The program
\nolinkurl{parts/C/post16_research/code/check_RDP_Hom_framing.py}
checks the endpoint chart exponents, colon-generator inclusions,
all rank-one cases of (D.HF.7) in its stated finite range, the
internal-branch substitutions, and the fixed-characteristic
Frobenius congruences. The retained author run passed $536\,909$
exact assertions: all endpoint cases for $2\le N\le50$, independent
finite colon-semigroup comparisons for $2\le N\le10$, and the
internal branches for $1\le m\le40$.
The optional \texttt{--output} argument preserves a separate audit run.
It does not replace the duality, descent,
or cohomology proofs and does not claim an Artinian-family
character decomposition.
\end{remark}


% Formal closure of the single post-review10 Brauer-zero main reduction.
% No separate movement event is claimed for these transfers.
\begin{corollary}[Proved: Brauer-zero counts through a detecting surface]
\label{D:Brauer-zero-detecting-surface-transfer}
Let $(S,T)$ satisfy the normal proper hypotheses of
Theorem~\ref{D:generic-surface-controlled-boundary}, in characteristic
$p>0$, and put $U=S\setminus T$ and $Q=Q(S,T)$.
Fix the surface $(X,T_X)$ and the geometric morphism
$h:X\setminus T_X\to U$ supplied there. There is a constant $C$
such that, for $q=p^f$ and $r\ge1$,
\[
 \#\left(\left\{\rho:Q\to\operatorname{PGL}_r(\mathbf F_q):
       \operatorname{br}(h^*\rho)=0\right\}
       /\operatorname{PGL}_r(\mathbf F_q)\right)
                \le q^{Cr^2}.                         \tag{D.BT.1}
\]
In particular this bounds the representations with
$\operatorname{br}(\rho)=0$ on $U$, without any bound on the
scalar centre. The same conclusion holds for every normal proper
scheme with a finite omitted subset and dimension at least two,
using Theorem~\ref{D:generic-surface-finite-punctures}; in that case
$Q=\pi_1(U)$. The chosen detecting surface is independent of $q,r$
and $\rho$.
\end{corollary}
\begin{proof}
The cited geometric theorem gives a continuous surjection
$Q(X,T_X)\twoheadrightarrow Q$. Precomposition is injective on
conjugacy classes: an equality after a surjective precomposition
is the same equality on $Q$. The representations counted on the
left of (D.BT.1) therefore inject into the Brauer-zero classes
counted in Theorem~\ref{D:Brauer-zero-projective-representation-growth}
for the one fixed normal proper surface $X$. That theorem gives
the bound with a constant depending on $(X,T_X)$, hence on the
fixed choice made from $(S,T)$.

The algebraic projective torsor and its lifting gerbe commute with
pullback, so the geometric Brauer obstruction does also. Thus
$\operatorname{br}(\rho)=0$ implies
$\operatorname{br}(h^*\rho)=0$. The converse is not used or asserted.
For a finite omitted subset, the cited finite-puncture theorem
provides precisely the same surjection of full fundamental groups.
No injectivity in degree-two cohomology is inferred from these
group surjections.
\end{proof}

\begin{corollary}[Proved: Brauer-zero counts under finite normalization]
\label{D:Brauer-zero-normalization-transfer}
Let $S$ be proper over an algebraically closed field of
characteristic $p>0$, and let $T\subset S$ have codimension at
least two in every irreducible component that it meets. Let $U$
be a connected component of $S\setminus T$. Assume either that
$T$ is finite or that the normalized components satisfy the
projective-modification hypothesis of
Corollary~\ref{D:nonnormal-controlled-boundary-quantitative-transfer}.

Let $V\to U$ be the restriction of the finite normalization of
$S_{\mathrm{red}}$, with connected components $V_i$. For each
component of dimension at least two, fix one detecting surface
map $h_i:W_i\to V_i$ as in the preceding corollary; here $W_i$
is a normal proper surface with a finite set removed.
There is a constant $C$ such that the number of conjugacy classes
of representations
\[
                \rho:\pi_1(U)\to
                          \operatorname{PGL}_r(\mathbf F_q)
\]
whose geometric Brauer class vanishes after pullback to every
such $W_i$ is at most $q^{Cr^2}$, for all $q=p^f$ and $r\ge1$.
In particular this holds for all representations with zero
geometric Brauer class on $U$. No condition on the scalar-centre
order is imposed.
\end{corollary}
\begin{proof}
The reduction and normalization constructions, the codimension
claim, and the fixed detecting surfaces are exactly those in
Corollaries~\ref{D:nonnormal-finite-puncture-quantitative-transfer}
and~\ref{D:nonnormal-controlled-boundary-quantitative-transfer}.
The preceding corollary bounds the allowed restriction classes
on every $V_i$ of dimension at least two by $q^{C_i r^2}$.
Components of dimension at most one are normal and proper,
since the componentwise codimension hypothesis leaves no
omitted point on them. Their fundamental groups are finitely
generated by the already established projective cases.
If $d_i$ generators suffice, the number of homomorphisms into
$P=\operatorname{PGL}_r(\mathbf F_q)$ is at most
$|P|^{d_i}\le q^{d_i r^2}$. No Brauer restriction is needed on
these components.

Set $a=\#\pi_0(V\times_U V)$, which is finite. Fix one tuple of
allowed $P$-torsor classes on the $V_i$ and choose representative
torsors. The Isom scheme between their two pullbacks to
$V\times_U V$ is finite etale and has $|P|$ elements in each
geometric fibre. A section on a connected component is determined
by its value in one geometric stalk. Hence at most $|P|^a$
descent isomorphisms are possible. The cocycle condition and
passage to isomorphism classes only decrease that number.
Finite proper descent of finite etale schemes, including the
action morphisms, is effective and fully faithful
\citep[Expose IX, Theorem~4.7]{C_SGA1}; this is the same explicit
descent argument used in the two cited normalization corollaries.

Multiplying the component bounds and the descent bound gives
$q^{Cr^2}$ with
$C=a+\sum_{\dim V_i\ge2}C_i+\sum_{\dim V_i\le1}d_i$.
We may count every descent datum on a fixed allowed tuple, even
if its descended Brauer class on $U$ is nonzero. The resulting
overcount is valid for the stated condition, which was imposed
only after the fixed surface pullbacks. Finally zero Brauer class
on $U$ pulls back to zero on each $W_i$. Again no degree-two
cohomology injection or descent theorem is asserted.
\end{proof}

\begin{remark}[Scope and movement accounting]
These are the functorial and finite-normalization closure of the
single Brauer-zero projective-growth theorem. They do not supply
an additional independent main-gate reduction. They also do not
bound projective representations with nonzero Brauer class on
the fixed detecting surfaces, nor arbitrary nonabelian crown
multiplicities, and they do not settle general D.1 or D.2(ii).
\end{remark}


% Four complete exact correction audits PASS. Original fixed-coordinate error and source are preserved.
% Intrinsic branch jets and an explicit adapted-coordinate proof; no hereditary atlas is asserted.
\begin{proposition}[Proved: finite jets for transverse A-chain framings]
\label{D:A-chain-finite-jet-framing}
Let $k$ be algebraically closed, let $N\ge2$, and let
$\pi:Y\to\operatorname{Spec}A$ be the minimal resolution of
$A=k[[u,v,w]]/(uv-w^N)$.  Number its exceptional chain
$E_1,\ldots,E_{N-1}$ so that
$v_{E_j}(u)=j$, $v_{E_j}(v)=N-j$, and $v_{E_j}(w)=1$.
Let $B_1,\ldots,B_q$ be smooth closed formal branches whose strict
transforms on $Y$ meet the exceptional set transversely at
\emph{pairwise distinct smooth points} of the chain, away from its nodes.
We use the same symbols for their strict transforms.
Write $j_i$ for the component met by $B_i$ and put
\[
 d_i=\min(j_i,N-j_i),\qquad
 \ell_i=\sum_{h=1}^q\min(d_i,d_h),\qquad B=\coprod_iB_i.
                                                        \tag{D.JF.1}
\]
There is a finite subscheme $T\subset B$, independent of the
reflexive modules considered below, whose restriction to $B_i$
has the intrinsic ideal $\mathfrak m_{B_i}^{\ell_i}$, where
$\mathfrak m_{B_i}$ is the maximal ideal of the smooth formal branch.
In a standard presentation adapted to these finitely many branches,
chosen in the proof, this ideal is $(w^{\ell_i})$.
For any two reflexive $A$-modules of ranks $r,s$, with full
sheaves $\mathcal F,\mathcal G$, the image of restriction
\[
 H^0(Y,\mathcal Hom(\mathcal F,\mathcal G))
       \longrightarrow H^0(B,\mathcal Hom(\mathcal F,\mathcal G)|_B)
                                                        \tag{D.JF.2}
\]
contains all sections vanishing on $T$.  Consequently a section
on $B$ extends to $Y$ if and only if its restriction to $T$
belongs to the image of the corresponding finite-jet restriction.
Moreover
\[
 h^1\bigl(Y,\mathcal Hom(\mathcal F,\mathcal G)(-B)\bigr)
            \le rs\,\operatorname{length}_k(T),
 \qquad
 \operatorname{length}_k(T)\le 2\,t^{\mathsf t}G_Nt,
                                                        \tag{D.JF.3}
\]
where $t_j=\#\{i:j_i=j\}$ and
\[
 (G_N)_{ab}=\frac{\min(a,b)(N-\max(a,b))}{N}
\]
is the inverse of the negative intersection matrix of the chain.
All characteristics, including those dividing $N$, are allowed.
These assertions concern individual modules on the fixed resolution,
not a simultaneous splitting of arbitrary Artinian families.
\end{proposition}
\begin{proof}
We first choose a standard presentation adapted to the fixed finite
set of branches. This step is necessary: in a preassigned presentation
the strict transform of $(u,w)=0$ meets an end transversely, but $w$
vanishes identically on that branch.

Write $u=x^N$, $v=y^N$, $w=xy$ in
$k[[x,y]]^{\mu_N}$, where $x,y$ have characters $1,-1$.
This invariant-ring description and all changes below also hold when
$p\mid N$: the diagonalizable group $\mu_N$ gives the usual
$\mathbf Z/N$ grading, with no division by $N$.
For $N>2$ choose $a,b\in k$ and set
\[
 x'=x+a y^{N-1},\qquad y'=y+b x^{N-1},\qquad
 u'=(x')^N,\quad v'=(y')^N,\quad w'=x'y'.
                                                        \tag{D.JF.0}
\]
The change on $k[[x,y]]$ has identity tangent map and therefore a
formal inverse. Both summands in $x'$, respectively $y'$, have the
same $\mu_N$ character, so the automorphism and its inverse are
equivariant. It induces a standard presentation
$A=k[[u',v',w']]/(u'v'-(w')^N)$, with
\[
                  w'=w+a v+b u+ab w^{N-1}.
\]

Here is the exact effect on the exceptional chain. At the smooth end
of $E_{N-1}$ the standard chart has coordinates
$(v,t)$ with $t=w/v=x/y^{N-1}$, and
$w=tv$, $u=t^Nv^{N-1}$. The chain node is the point $t=\infty$.
The induced coordinate on the exceptional curve is $t\mapsto t+a$:
indeed
\[
 \frac{x'}{(y')^{N-1}}
       =\frac{t+a}{(1+bv^{N-2}t^{N-1})^{N-1}}.
\]
At $E_1$ the analogous chart $(u,t)$ with $t=w/u=y/x^{N-1}$
gives the independent translation $t\mapsto t+b$.
On an interior component, $1<j<N-1$, the quotients
$x'/x-1$ and $y'/y-1$ have positive exceptional orders
$N-1-j$ and $j-1$, respectively. Thus its leading coordinates
are unchanged. The induced surface automorphism lifts to the
minimal resolution; these calculations also show that it preserves
the component numbering and all three divisorial valuations in
the statement. Choose $a$ and $b$ to avoid the finitely many
negatives of the endpoint attachment coordinates. Then neither
new coordinate-axis strict transform passes through any attachment.

For $N=2$, use instead two independent linear forms $x',y'$ in
$x,y$, each nonzero at all the finitely many attachment directions
in the exceptional $\mathbf P^1$. Such a pair exists because $k$
is infinite. The group $\mu_2$ acts by the same character on $x$
and $y$, even in characteristic two, so this change too is
equivariant and induces a standard presentation. The single
exceptional component is preserved, and the new coordinate axes
avoid every attachment.

In either case, at every attachment the divisor of $w'$ has only
the exceptional component, with multiplicity one. Transversality
therefore makes $w'$ a uniformizer on every branch. This choice
depends only on the branches, not on the reflexive modules.
Rename the adapted coordinates $u',v',w'$ as $u,v,w$ for the
rest of the proof. The intrinsic ideal
$\mathfrak m_{B_i}^{\ell_i}$ is now $(w^{\ell_i})$.
In the standard chart away from the two coordinate axes and the
chain nodes, $u/w^{j_i}$ is a unit. Hence there are fixed units
$\epsilon_i(w)\in k[[w]]^\times$ such that the branch image has
the parametrization
\[
 u=\epsilon_i(w)w^{j_i},\qquad
 v=\epsilon_i(w)^{-1}w^{N-j_i}.                          \tag{D.JF.4}
\]
Its image is a smooth closed disc because the quotient by these
two equations is $k[[w]]$.  Conversely the standard resolution
charts show that (D.JF.4) is the prescribed strict transform.
If $j_i=j_h$ and $i\ne h$, distinct attachment points mean
$\epsilon_i(0)\ne\epsilon_h(0)$.

Set $I_0=A$ and $I_a=(u,w^a)$ for $0<a<N$.
Proposition~\ref{D:A-chain-two-end-Hom-bound} proves the
individual direct-sum decomposition into these modules, including
in the nonreduced $\mu_N$ case.  It therefore suffices to construct
the extension for each ordered pair $I_a,I_b$.
On $B_i$ a generator of the full line bundle of $I_a$ has order
$\min(j_i,a)$, with the convention that this is zero for $a=0$.
Using the fractional Hom ideals in (D.HF.6), the least order
$e_{ab}(j)$ of a nonzero restricted Hom section is
\[
 e_{ab}(j)=
 \begin{cases}
 0,&a=b,\\
 \min(j,b-a)+\min(j,a)-\min(j,b),&b>a,\\
 \min(0,N-j-a+b)+\min(j,a)-\min(j,b),&b<a.
 \end{cases}                                           \tag{D.JF.5}
\]
One of the displayed global generators attains that order, and
its leading coefficient is a unit on the branch.
In every case
\[
                  0\le e_{ab}(j)\le\min(j,N-j).          \tag{D.JF.6}
\]
For $b>a$, the second line is zero if $j\ge b$; otherwise
it is at most both $j$ and $b-j<N-j$.
For $b<a$, put $c=a-b$ and
$\Delta=\min(j,a)-\min(j,b)$, so $0\le\Delta\le\min(j,c)$.
If $c\le N-j$, the third line is $\Delta$.
If $c>N-j$, it is $N-j-c+\Delta$, at most $N-j$ and at
most $j$.  Its nonnegativity also follows directly from the
regularity of these Hom sections, or by the two cases $j<a$
and $j\ge a$.

For each ordered pair $i\ne h$ choose one of the two elements
\[
 u-\epsilon_h(w)w^{j_h},\qquad
 v-\epsilon_h(w)^{-1}w^{N-j_h}                         \tag{D.JF.7}
\]
which has smaller order on $B_i$, and denote it by $f_{ih}$.
Both vanish on $B_h$, while the chosen order on $B_i$ is
exactly
\[
       m_{ih}=\min(j_i,N-j_i,j_h,N-j_h)=\min(d_i,d_h).
                                                        \tag{D.JF.8}
\]
If the component indices differ, their leading exponents differ
and cannot cancel.  If they agree, the distinct constants in
(D.JF.4) prevent cancellation.  The chosen element may depend
on the ordered pair; we are not asserting that the Weil branch
$B_h$ has one defining Cartier equation downstairs.

The product $f_i=\prod_{h\ne i}f_{ih}$ vanishes on every
$B_h$ with $h\ne i$, and its restriction to $B_i$ is a unit
times $w^{c_i}$, where $c_i=\sum_{h\ne i}m_{ih}$.
Multiply $f_i$ by a global Hom generator attaining
$e_{ab}(j_i)$.  This gives a section which vanishes on all
other branches and has order $c_i+e_{ab}(j_i)\le\ell_i$
on $B_i$.  Multiplication by any further series in $w$ is
allowed, since $k[[w]]\subset A$.  In particular, one can
divide by the unit coefficient and supply arbitrary higher
coefficients.  Every section on $B_i$ vanishing to order
$\ell_i$ is therefore the restriction of a global Hom
section which is zero on the other branches.
Summing these sections and then taking the individual character
direct sums proves the asserted image inclusion.

The divisor sequence for $B$ and
Lemma~\ref{D:RDP-full-Hom-vanishing} identify the cokernel of
(D.JF.2) with the displayed $H^1$.  The image inclusion just
proved shows that this cokernel is a quotient of the finite
space of rank-$rs$ sections on $T$, proving its bound and
the extension criterion.

It remains to record the useful total-length estimate.
Let $R_N$ denote the permutation matrix reflecting the chain,
$j\mapsto N-j$.  Direct substitution in the formula for $G_N$
gives
\[
 (G_N+G_NR_N)_{ab}
      =\min\bigl(\min(a,N-a),\min(b,N-b)\bigr).
                                                        \tag{D.JF.9}
\]
Hence $\sum_i\ell_i=t^{\mathsf t}(G_N+G_NR_N)t$.
The positive-definite matrix $G_N$ commutes with $R_N$,
and $R_N^2=1$.  Cauchy--Schwarz for the $G_N$ inner product
gives
$t^{\mathsf t}G_NR_Nt\le t^{\mathsf t}G_Nt$.
This proves the last inequality of (D.JF.3).
For $q=0$, the subscheme $T$ is empty and the assertion is
the unframed vanishing.
\end{proof}

\begin{corollary}[Proved: a degree-linear total jet budget on A-blocks]
\label{D:A-block-total-jet-budget}
Consider a negative-definite exceptional configuration on a
regular surface, partitioned into a complementary set of
curves $\{C_a\}$ and pairwise disjoint $A$-chain blocks.
Assume the intersections between the complementary curves
and each block are transverse, occur at pairwise distinct
smooth points of that block, and are away from its nodes.
After contracting a block, regard these fixed formal branches
on its resolution as the $B_i$ of the preceding proposition.
For comparisons of full sheaves of reflexive modules on the
contracted rational double points, their branch compatibility
is determined by finite subschemes whose total length satisfies
\[
          \sum_{\text{blocks }\alpha}\operatorname{length}(T_\alpha)
                        \le 2\sum_a(-C_a^2).            \tag{D.JF.10}
\]
An entire isolated rational double point component with no
attachments contributes no jet conditions.  If it is not of
type $A$, this last assertion is still the unframed vanishing
of Lemma~\ref{D:RDP-full-Hom-vanishing}; the framed estimate
here is asserted only for the stated $A$-blocks.

In the degree-$n$ punctured-cover setting of
Proposition~\ref{D:punctured-cover-linear-resolution}, whenever
the stated transverse block conditions and full-sheaf comparison
hold, (D.JF.10) is $O(n)$.  The corresponding total number of
linear conditions for maps between ranks $r,s$ is $O(nrs)$.
This is a finite-jet placement and dimension estimate; it
does not construct an effective hereditary object atlas.
\end{corollary}
\begin{proof}
Write the negative intersection matrix in block form, with
principal core matrix $N_C$, block matrices $N_\alpha$, and
off-diagonal incidence matrices $-T_\alpha$.
The entries of $T_\alpha$ are intersection multiplicities;
under the stated hypotheses each intersection point contributes
one branch.  Put
$G_\alpha=N_\alpha^{-1}$ and
$t_\alpha=T_\alpha^{\mathsf t}{\bf1}$.
Positive definiteness of the full matrix gives the positive
Schur complement
\[
        N_C-\sum_\alpha T_\alpha G_\alpha T_\alpha^{\mathsf t}>0.
                                                        \tag{D.JF.11}
\]
If the complementary set is nonempty, evaluating on its
all-ones vector yields
\[
 \sum_\alpha t_\alpha^{\mathsf t}G_\alpha t_\alpha
       <{\bf1}^{\mathsf t}N_C{\bf1}
       \le\sum_a(-C_a^2),                               \tag{D.JF.12}
\]
since the off-diagonal entries of $N_C$ are nonpositive.
Apply (D.JF.3) to each block and sum.
If the complementary set is empty, all blocks have no
attachments and contribute zero, so the same conclusion holds.
The final degree-linear assertion uses the already proved
bound for the total self-intersection budget in the cited
punctured-cover proposition.
\end{proof}

\begin{remark}[Transversality and the remaining family gate]
The requirement of distinct attachment points is essential to
the proof: two different branches at the same point can have
arbitrarily large contact, and the leading-term assertion
(D.JF.8) then fails.  No estimate for such contacts is implicit
in (D.JF.10).  Nor do the character decomposition and the
explicit extension above assert uniform decomposition over
an arbitrary Artinian base.  Those geometric and family
issues, together with an effective global presentation,
remain separate from the finite-jet bound proved here.
\end{remark}


% Root/B/D complete actual-file proof audits PASS.
% Separate post-CP16 auxiliary unit. The resolution is fixed throughout.
\begin{proposition}[Proved: Artinian rigidity and finite-jet framing of full RDP bundles]
\label{D:RDP-Artinian-full-bundle-framing}
Let $Y$ be the fixed minimal resolution of an algebraic henselian
or complete rational double point over an algebraically closed field
$k$.  Let $\mathcal F,\mathcal G$ be full bundles of reflexive modules,
of ranks $r,s$.  Let $\Lambda$ be a local Artinian $k$-algebra with
residue field $k$, and write $Y_\Lambda=Y\times_k\operatorname{Spec}\Lambda$.

Every vector bundle $\mathcal F_\Lambda$ on $Y_\Lambda$ with
identified closed fibre $\mathcal F$ is isomorphic to
$\mathcal F\otimes_k\Lambda$, through an isomorphism inducing that
identification.  More generally, over a small extension
$\Lambda'\twoheadrightarrow\Lambda$ with kernel $J$ annihilated
by the maximal ideal, a specified isomorphism between two such
bundles over $\Lambda$ lifts to an isomorphism over $\Lambda'$.

Suppose in addition that $Y$ resolves an $A$-chain, and fix
the disjoint transverse branches $B$ and the finite subscheme
$T\subset B$ from Proposition~\ref{D:A-chain-finite-jet-framing}.
For any vector-bundle families $\mathcal F_\Lambda,\mathcal G_\Lambda$
with the specified full closed fibres, the image
\[
 H^0(Y_\Lambda,\mathcal Hom(\mathcal F_\Lambda,\mathcal G_\Lambda))
 \longrightarrow
 H^0(B_\Lambda,\mathcal Hom(\mathcal F_\Lambda,\mathcal G_\Lambda)|_{B_\Lambda})
                                                        \tag{D.AF.1}
\]
contains all sections vanishing on $T_\Lambda$.
Its cokernel is a finite free $\Lambda$-module of rank at most
$rs\operatorname{length}_k(T)$; its identification with the
scalar extension of the closed-fibre cokernel may depend on
the chosen trivializations of the two families.

For a small extension as above, a specified global map over
$\Lambda$ and a compatible boundary map over $\Lambda'$
have a simultaneous global lift if and only if their difference
class vanishes in
\[
 \operatorname{coker}\!\left[
 H^0(Y,\mathcal Hom(\mathcal F,\mathcal G))
 \to H^0(B,\mathcal Hom(\mathcal F,\mathcal G)|_B)
 \right]\otimes_k J.                                   \tag{D.AF.2}
\]
This class is determined by the fixed finite jets on $T$.
If the specified map over $\Lambda$ is an isomorphism, any
simultaneous lift is an isomorphism.

These statements require actual vector bundles on the fixed
$Y_\Lambda$.  They do not assert that a relative maximal
Cohen--Macaulay module has a flat, locally free full pullback,
and they do not allow the singularity or its resolution to vary.
\end{proposition}
\begin{proof}
Lemma~\ref{D:RDP-full-Hom-vanishing} gives
\[
 H^1(Y,\mathcal End(\mathcal F))=0,\qquad
 H^1(Y,\mathcal Hom(\mathcal F,\mathcal G))=0.             \tag{D.AF.3}
\]
We recall the elementary deformation calculation explicitly.
Take a finite affine cover of the quasi-compact separated
scheme $Y$ on which the closed-fibre bundles are free.
All finite intersections are affine.  On each member of
the cover, a basis over a nilpotent quotient lifts to a
basis of a vector bundle: lift its vectors and use
Nakayama; the resulting map between equal-rank finite
projective modules is invertible modulo a nilpotent ideal
and therefore invertible.

Consider a small extension $\Lambda'\to\Lambda$ and a
specified isomorphism between the two reductions.
Choose local lifts of its matrices.  They are automatically
invertible.  Their differences on pairwise overlaps lie
in the square-zero kernel and, after using the specified
closed-fibre identifications, form an additive Cech
$1$-cocycle in
$\mathcal Hom(\mathcal F,\mathcal F)\otimes_kJ$
(or in the corresponding Hom between the two identified
full closed fibres).  Equation (D.AF.3) makes this a
coboundary.  Subtracting a $0$-cochain corrects the local
lifts so that they glue to a global isomorphism.
Applying this successively along a composition series
of $\Lambda$ proves the asserted constancy, with any
specified closed-fibre identification retained.
The same argument for arbitrary local matrix lifts,
without the invertibility requirement, lifts global maps
between the two bundle families.
No properness over $k$ or finite-dimensionality of the
global Hom space is needed for this Cech argument.

Choose the resulting isomorphisms of $\mathcal F_\Lambda$
and $\mathcal G_\Lambda$ with their constant families.
Because $\Lambda$ is finite-dimensional and flat over $k$,
the groups of sections on $Y_\Lambda$ and $B_\Lambda$
of their Hom bundles are the corresponding closed-fibre
groups tensored with $\Lambda$.  This follows directly
from the same finite affine Cech complex; on the branches
it is the equality
$k[[w]]\otimes_k\Lambda=\Lambda[[w]]$.
Tensoring the image inclusion in
Proposition~\ref{D:A-chain-finite-jet-framing}
therefore proves (D.AF.1).
Tensoring its finite-dimensional cokernel identifies
the relative cokernel with a finite free $\Lambda$-module
of the asserted rank.  These conclusions are invariant
under the chosen trivializations.

Finally, lift the given global map from $\Lambda$ to
$\Lambda'$ without its boundary constraint, using the
second vanishing in (D.AF.3).  Its boundary restriction
differs from the prescribed boundary lift by a section
of $\mathcal Hom(\mathcal F,\mathcal G)|_B\otimes_kJ$.
A different global lift changes this section by the
image of a global Hom section tensored with $J$.
This proves that the well-defined obstruction is exactly
(D.AF.2), and that its vanishing is necessary and sufficient.
The image inclusion from the finite-jet proposition
makes its class depend only on restriction to $T$.
A global map reducing to an isomorphism modulo a
nilpotent ideal is an isomorphism on every affine chart,
which proves the last assertion.
\end{proof}

\begin{remark}[What the family statement repairs]
On the transverse $A$-blocks in
Corollary~\ref{D:A-block-total-jet-budget}, this puts the
boundary map-lifting obstructions for \emph{actual full-bundle
families} in spaces of total dimension $O(nrs)$.
The statement does not construct the global family of such
bundles, prove that arbitrary module families enter this
scope, or bound the variables and degrees of a global
hereditary object presentation.  Those requirements remain
separate.
\end{remark}


% Separate post-CP16 extension. Root/B/C/D full actual-file proof audits PASS.
\subsection{Brauer-zero projective representations on complete local punctures}

Let $(R,\mathfrak m,k)$ be a complete excellent normal local domain
of dimension two, of equal characteristic $p>0$, with an
algebraically closed coefficient field $k$.  Put
$U=\operatorname{Spec}R\setminus\{\mathfrak m\}$.
All relative coefficient rings below are finite-dimensional local
Artinian $k$-algebras $A$ with residue field $k$; write
$R_A=R\otimes_kA$.  A relative MCM module is finite and $A$-flat
with MCM closed fibre, as in
Lemma~\ref{D:relative-Noether-different-jets}.

\begin{lemma}[Proved: fixed generators for local prime-to-characteristic Picard quotients]
\label{D:local-Picard-prime-to-p-generators}
There are fixed rank-one reflexive $R$-modules $P_1,\ldots,P_b$
such that, for every $n$ prime to $p$, every element of
$\operatorname{Pic}(U)/n$ is represented by
\[
                \bigotimes_{i=1}^b(P_i|_U)^{a_i},
                         \qquad 0\le a_i<n.             \tag{D.LBZ.1}
\]
The modules are independent of $n$.
\end{lemma}

\begin{proof}
Choose a projective regular resolution
$\pi:Y\to\operatorname{Spec}R$, an isomorphism over $U$,
whose reduced exceptional fibre $E$ is a nonempty strict
normal-crossings curve.  A preliminary blowup deals with the
regular case.  This is the same excellent-surface resolution
and embedded-resolution scope as
Theorem~\ref{D:local-surface-constant-H2}.
Write $E_1,\ldots,E_b$ for its smooth proper components.

First, restriction $\operatorname{Pic}(Y)\to\operatorname{Pic}(E)$
is surjective.  For a square-zero thickening $W\subset W'$
with ideal $\mathcal I$, the exact sequence of units gives
\[
 \operatorname{Pic}(W')\longrightarrow\operatorname{Pic}(W)
                         \longrightarrow H^2(W,\mathcal I).
                                                        \tag{D.LBZ.2}
\]
The last group vanishes when the support has dimension one.
Apply this first through the finite nilpotent thickening
from $E$ to the scheme-theoretic closed fibre, and then
through the successive $\mathfrak m$-adic thickenings of $Y$.
Every line bundle on $E$ therefore has a compatible
system of line-bundle lifts, with chosen transition
isomorphisms.  Grothendieck existence algebraizes the system.
It is an invertible sheaf: algebraize the compatible inverse
system too, and use full faithfulness to algebraize their
evaluation isomorphisms.  The two tensor products are then
$\mathcal O_Y$
\citep[Lemma~37.4.1, Tag 0C6Q;
Theorem~30.27.1 and proof, Tag 088E]{StacksLocalBrauerZeroD}.

For $n$ prime to $p$, proper base change over the strictly
henselian ring $R$ gives
\[
                  H^2(Y,\mu_n)\simeq H^2(E,\mu_n).
                                                        \tag{D.LBZ.3}
\]
Indeed taking sections over $\operatorname{Spec}R$
is the closed geometric-stalk functor and is exact;
proper base change applies to every proper morphism of
schemes and an abelian torsion sheaf.  Nilpotents in
the closed fibre have no effect
\citep[Theorem~59.91.11 and complete proof,
Tag 095T]{StacksLocalBrauerZeroD}.
The Kummer injections of $\operatorname{Pic}/n$ into
$H^2(\mu_n)$, together with (D.LBZ.3), show that
$\operatorname{Pic}(Y)/n\to\operatorname{Pic}(E)/n$
is injective.  The surjectivity just proved makes it
an isomorphism.  This argument proves divisibility
of the restriction kernel without an inverse-limit
choice of compatible $n$th roots.

We describe the latter Picard quotient directly.
Normalization of the nodal curve $E$ gives the units
sequence
\[
 1\to\mathcal O_E^\times\to
   \nu_*\mathcal O_{\coprod E_i}^{\times}
       \to\bigoplus_{\text{nodes }x}k_x^\times\to1.
                                                        \tag{D.LBZ.4}
\]
The last map takes the ratio of the two values at a
node.  Its cohomology says that the kernel of the
multidegree map on $\operatorname{Pic}(E)$ is an
extension of $\prod_i\operatorname{Pic}^0(E_i)(k)$
by a quotient of a finite product of $k^\times$.
Both groups are $n$-divisible: for the Jacobians this
is multiplication-surjectivity on an abelian variety,
and for $k^\times$ it is existence of roots in $k$.
An extension of divisible-by-$n$ abelian groups is
divisible by $n$.

The multidegree map is onto $\mathbf Z^b$.  Choose a
smooth $k$-point on each component away from its nodes;
its Cartier divisor on $E$ has the corresponding
unit multidegree.  Consequently
\[
                      \operatorname{Pic}(E)/n
                                  \simeq(\mathbf Z/n)^b.
                                                        \tag{D.LBZ.5}
\]
Lift those $b$ line bundles once to $Y$.
Restriction $\operatorname{Pic}(Y)\to\operatorname{Pic}(U)$
is onto: a rational Cartier divisor of a line bundle
on $U$ extends by closure to a Weil divisor on the
regular Noetherian $Y$, and is Cartier there.
The chosen line bundles therefore generate
$\operatorname{Pic}(U)/n$ for every $n$.
Their restrictions have unique rank-one reflexive
extensions $P_i$ on the normal local surface.
Reducing their integer exponents modulo $n$ gives
(D.LBZ.1).
\end{proof}

\begin{proposition}[Proved: fixed tensor twists in the local quadratic atlas]
\label{D:local-fixed-tensor-cyclic-atlas}
Fix a finite collection $\mathscr P$ of rank-one reflexive
$R$-modules.  For each length-$f$ word $P_j\in\mathscr P$,
consider cycles of rank-$r$ reflexive modules with
isomorphisms on the puncture
\[
       F_U^*(M_j|_U)\otimes P_j|_U
                             \simeq M_{j+1}|_U,
                \qquad j\in\mathbf Z/f\mathbf Z.          \tag{D.LBZ.6}
\]
These cycles have polynomial witnesses in $O(fr^2)$
coordinates and equations of bounded degree, effective
over every local Artinian coefficient ring.
Their isomorphism classes form open and closed parts
of the witness scheme.  In particular, the number
of classes, after summing over the tensor words, is
at most $p^{Cfr^2}$ for a constant depending only on
$R$ and $\mathscr P$.
\end{proposition}

\begin{proof}
\emph{The source need not be MCM in the map-lifting argument.}
We first isolate the precise strengthening of
Lemma~\ref{D:relative-Noether-different-jets}.
Let $K$ be any finite $R_A$-module for which
$x\,\operatorname{Id}_K$ and $y\,\operatorname{Id}_K$
factor through finite free modules.  Let $N$ be a finite
module on which $x,y$ is a regular sequence.  Put
$I_c=(x^c,y^c)R_A$.  Every map
$K/I_cK\to N/I_cN$, $c\ge2$, lifts to a map $K\to N$
agreeing modulo $I_{c-1}N$.

Here is the proof, which does not require regularity
or coefficient flatness of $K$.  Reduce modulo $x^c$.
The factorization of $y\,\operatorname{Id}_K$ remains
a factorization through a free module over $R_A/(x^c)$.
It annihilates the obstruction in $\operatorname{Ext}^1$
to lifting $y$ times the prescribed map through
$N/x^cN\to N/I_cN$.  The lift has image in
$y(N/x^cN)$; divide by the non-zero-divisor $y$.
The resulting map agrees with the prescribed one
modulo $y^{c-1}$.  Next the factorization of
$x\,\operatorname{Id}_K$ lifts $x$ times this map
through $N\to N/x^cN$.  Its image is in $xN$;
division by $x$ gives the asserted lift.
Only the target divisions use regularity.

Choose parameters $u,v$ in the fixed sum of Noether
differents of $R$, as in the cited lemma.  Their
identity factorizations work on every relative MCM
module, including each fixed $P_A=P\otimes_kA$.
For
\[
                          K=F_{R_A}^*M_A\otimes_{R_A}P_A
\]
Frobenius of the $u$-factorization on $M_A$
factors $u^p\,\operatorname{Id}_{F^*M_A}$ through
a free module.  Tensor it with the $u$-factorization
on $P_A$.  The composite factors
$u^{p+1}\operatorname{Id}_K$ through a finite
free module.  The same holds for $v$.
Hence choose once and for all
\[
       x=u^{p+1},\quad y=v^{p+1},\quad
       S_0=k[[x,y]],\quad d=\operatorname{rank}_{S_0}R.
                                                        \tag{D.LBZ.7}
\]
These parameters also have the required factorizations
on all relative MCM modules themselves.  The realization
lemma \ref{D:relative-Artin-module-realization} therefore
applies with this $S_0$ and with depth $L=3$.

\emph{A direct image bound without a relative reflexive hull.}
Suppose $\phi:K|_{U_A}\simeq N_A|_{U_A}$ is an
isomorphism, with $N_A$ relative MCM.  It extends
to $K\to N_A$, since $N_A=j_*N_A|_{U_A}$.
Write $x\operatorname{Id}_K$ as
$K\xrightarrow{\alpha}F\xrightarrow{\beta}K$
with $F$ free.  The map
$\alpha\phi^{-1}:N_A|_{U_A}\to F|_{U_A}$
extends to $N_A\to F$, by the Hartogs property of
the free target.  Composing with
$F\to K\to N_A$ gives $x\operatorname{Id}_{N_A}$:
this equality holds on the puncture and extends
because $N_A$ has the Hartogs property.
Thus
\[
               (x,y)N_A\subset\operatorname{im}(\phi).
                                                        \tag{D.LBZ.8}
\]
This proves the needed uniform image condition
directly.  It makes no assertion that a relative
reflexive hull commutes with base change.

\emph{Equations for the whole tensor map.}
At each stage use the depth-$I_3^2$ action matrices
of Lemma~\ref{D:relative-Artin-module-realization};
their canonical reductions modulo $I_3$ are free
of rank $dr$ over $S_0/(x^3,y^3)$.
For each $P\in\mathscr P$, choose once a finite
presentation with generators $e_1,\ldots,e_a$
and relations
$\sum_\nu c_{\lambda\nu}e_\nu=0$.
All its coefficients and the integers $a$ and
the number of relations are fixed.

A map
$F^*M\otimes P\to N$ modulo $I_3$ is encoded by
$a$ matrices $A_\nu$, specifying the $p$-semilinear
maps $m\mapsto\phi(1\otimes m\otimes e_\nu)$.
For a fixed $S_0$-basis $z_i$ of $R$, impose
\[
 \begin{split}
 A_\nu\rho_M(z_i)^{[p]}&=\rho_N(z_i)^p A_\nu,\\
 \sum_\nu\rho_N(c_{\lambda\nu})A_\nu&=0.
 \end{split}                                            \tag{D.LBZ.9}
\]
The semilinearity over $S_0$ is part of the matrix
description; the left power is entrywise Frobenius,
whereas the right power is ordinary matrix
multiplication.  The second equation imposes
\emph{all} the fixed relations of $P$.
Depth $I_3$ of $M$ suffices: an error in $I_3M$
has Frobenius image in $I_3^{[p]}$, which is
contained in $I_3$ in the target.
Thus (D.LBZ.9) is precisely one map
$K/I_3K\to N/I_3N$, not independently chosen
maps which only approximately respect $P$.

For $b=x,y$, add inclusion-witness matrices over
$S_0/(x^2,y^2)$ and require
\[
       \rho_N(b)=
          \sum_{i,\nu}\rho_N(z_i)A_\nu X_{i,\nu,b}
                                      \pmod {I_2}.       \tag{D.LBZ.10}
\]
The $R$-span of the semilinear generator images
is exactly the image of the linearized tensor
map, so these are the required image witnesses.
The multiplication table of $R$, all truncated
coefficient rings, and the presentations of
$P$ are fixed.  Equations (D.LBZ.9)--(D.LBZ.10)
therefore require $O(r^2)$ variables and equations
per stage, with degrees bounded in terms of
$p$ and the fixed presentations.  No determinant
of growing degree is used.

\emph{Effectiveness over every Artinian base.}
First realize the module actions by relative
MCM modules $M_{j,A}$, using the exact realization
lemma.  Then (D.LBZ.9) gives a genuine map of the
whole tensor source modulo $I_3$.
The source-factorization lifting argument above
gives an actual map
\[
                F^*M_{j,A}\otimes P_{j,A}\longrightarrow
                                                     M_{j+1,A}
\]
agreeing modulo $I_2$.  Its fixed $P$-relations
are exact because we lifted the single map
from $K$, not its matrix columns separately.
If $C$ is its cokernel, (D.LBZ.10) yields
\[
               (x,y)C\subset I_2C
                          \subset\mathfrak m(x,y)C.
                                                        \tag{D.LBZ.11}
\]
Nakayama gives $(x,y)C=0$.
On $U_A$ the ideal $(x,y)$ is the unit ideal,
and the map is a surjection between vector
bundles of the same rank.  It is therefore
an isomorphism.  This proves effectiveness
without requiring $F^*M_{j,A}$ to be $A$-flat.

Conversely, a genuine cycle extends to the
canonical reflexive modules.  The image bound
(D.LBZ.8) gives the inclusion witnesses.
Its actual tensor maps give (D.LBZ.9), and
its modules, free over $S_0$, give the
depth-$I_3^2$ action data.  Thus every cycle
has a point in the witness scheme.

\emph{Open class fibres.}
For a fixed word put $q=p^f$.  Its composite is
a twisted unit structure
\[
           F^{f*}(M_0|_U)\otimes L\simeq M_0|_U,
 \qquad
 L=\bigotimes_{j=0}^{f-1}(P_j|_U)^{p^{f-1-j}}.            \tag{D.LBZ.12}
\]
A sufficiently high iterate, with exponent a
multiple of $f$, kills the maximal ideal of
an Artinian coefficient ring $A$ under
Frobenius.  Every intervening tensor line
is the pullback of a fixed line on $U$,
not a varying line over $A$.  Iterating
(D.LBZ.12), and using its special-fibre
isomorphism, identifies the entire cycle
over $U_A$ with its closed-fibre cycle
tensored by $A$ at the level of the
underlying bundles.  Hartogs extends
these identifications to the relative
MCM modules.  In particular their exact
canonical reductions modulo $I_3$ are
constant up to ordinary change of basis.

There is also uniqueness of the twisted
unit structure on a fixed underlying
module.  In the proof of
Lemma~\ref{D:local-complete-Lang},
replace its conjugation by Frobenius
with conjugation by Frobenius
\emph{tensored by the fixed line $L$}.
It still defines a $q$-semilinear
multiplicative map on
$\operatorname{End}_R(M_0)$: maps on
the puncture extend by reflexivity.
It still sends
$\mathfrak m^a\operatorname{End}(M_0)$
into $\mathfrak m^{qa}\operatorname{End}(M_0)$.
The connected-unit-group argument on
the residue algebra and the contracting
complete-algebra iteration consequently
give exactly the same Lang uniqueness.
For the fixed word the intermediate
modules and edge identifications are
then determined up to cycle isomorphism.

Map the witness scheme to its tuple of
ordinary $R/I_3$ action matrices.
The group
$\operatorname{Res}_{(S_0/(x^3,y^3))/k}
          \operatorname{GL}_{dr}^{\,f}$
has locally closed ordinary module
orbits, with smooth stabilizers given
by unit groups of endomorphism algebras.
Every Artinian deformation of a closed
witness point maps into the orbit of
its special point by the preceding
Hartogs constancy.  Pull back the orbit
ideal in an open neighbourhood where
the orbit is closed.  It vanishes in
every Artinian quotient of the local
ring, hence in that local ring by
Krull intersection.  A neighbourhood
of the point maps into the orbit.
Exact module realization and twisted
Lang uniqueness make it one cycle
class.  This is the same specified-jet
openness argument as in
Theorem~\ref{D:uniform-local-representation-growth}.

The class fibres give a finite clopen
partition.  Bezout bounds the number
of irreducible components of the
witness scheme by $\exp(O(fr^2))$.
There are only $|\mathscr P|^f$
tensor words.  Absorbing this factor
proves the assertion.
\end{proof}

\begin{theorem}[Proved: quadratic Brauer-zero projective growth on a local surface puncture]
\label{D:local-Brauer-zero-projective-growth}
For the complete normal local surface $R$ above, there is
$C_R$ such that for every $q=p^f$ and $r\ge1$,
\[
 \#\left(
  \left\{\rho:\pi_1(U)\to\operatorname{PGL}_r(\mathbf F_q):
              \operatorname{br}(\rho)=0
                    \text{ in }H^2(U,\mathbf G_m)\right\}
       /\operatorname{PGL}_r(\mathbf F_q)\right)
                     \le q^{C_Rr^2}.                   \tag{D.LBZ.13}
\]
The centre is allowed to grow.  There is no requirement
that $R$ be an algebraization of a finite-type singularity.
\end{theorem}

\begin{proof}
Use the fixed generators of
Lemma~\ref{D:local-Picard-prime-to-p-generators}
with $n=q-1$.  The finite central lifting gerbe of
a representation whose geometric Brauer class
vanishes is the root gerbe of
$M_{\mathbf a}=\bigotimes_i(P_i|_U)^{a_i}$
for some $0\le a_i<n$.
The weight-one finite unit bundle untwists to
$E_0$ with
\[
                 F^{f*}E_0\otimes M_{\mathbf a}
                                     \simeq E_0.
                                                        \tag{D.LBZ.14}
\]
This is exactly the Kummer and root-gerbe
construction in
Theorem~\ref{D:Brauer-zero-projective-representation-growth};
it applies to this scheme without a finite-type
over $k$ hypothesis.

Write the scalar exponents in the same cyclic
base-$p$ digits
$a_{i,j+1}=p a_{i,j}-(q-1)c_{i,j}$.
All the one-step twists are among the fixed
finite set of reflexive modules
\[
       \mathscr P=
       \left\{\left(\bigotimes_iP_i^{c_i}\right)^{**}:
                                0\le c_i<p\right\}.
                                                        \tag{D.LBZ.15}
\]
There are at most $p^{bf}$ scalar words.
Formula (D.LBZ.14) gives a cycle (D.LBZ.6)
by the same floor-function recurrence as
(D.BZ.21)--(D.BZ.22).  Conversely a valid
scalar word and such a cycle give a
weight-one $q$-unit bundle on its root
gerbe.  Its fixed-frame equations are
finite etale over every trivializing
chart, so it descends projectively to
the required finite-field projective
representation.  The finite quotient
gerbe description from (D.BZ.23) is
available whenever such an object exists;
the band of an arbitrary root gerbe
is not assumed to inject into its
fundamental group.

The preceding proposition counts all
cycle classes for the fixed finite
alphabet by $p^{Cfr^2}$.  Projectivization
is surjective onto the desired classes,
and any extra choices of gerbe
identification or scalar lifts only
overcount.  Absorbing the scalar-word
factor proves (D.LBZ.13).
\end{proof}

\begin{remark}[Remaining local obstruction]
The local theorem includes nonzero Kummer classes with
zero geometric Brauer image and arbitrary scalar order.
It does not prove vanishing of the Brauer group of the
puncture, count the genuinely twisted projective
representations, or prove finite generation of the
full local fundamental group.
\end{remark}


% Consequence of the single Brauer-zero projective-count reduction.
\begin{corollary}[Proved: simple-product quotients selected by Brauer-zero representations]
\label{D:Brauer-zero-selected-simple-product}
Let $Q$ be a group in one of the global or complete-local cases
of Theorem~\ref{D:Brauer-zero-projective-representation-growth},
Corollaries~\ref{D:Brauer-zero-detecting-surface-transfer} and
\ref{D:Brauer-zero-normalization-transfer}, or
Theorem~\ref{D:local-Brauer-zero-projective-growth}.
Use exactly the Brauer-vanishing condition of the chosen statement.
Fix $\epsilon>0$. For each isomorphism type $H$ in a collection
of nonabelian finite simple groups, choose one faithful embedding
\[
       \iota_H:H\hookrightarrow
                    \operatorname{PGL}_{r_H}(\mathbf F_{q_H}),
       \qquad q_H=p^{f_H},\qquad
                    |H|\ge q_H^{\epsilon r_H^2}.
\]
Let $\mathcal K_H$ consist of the distinct kernels of continuous
epimorphisms $\eta:Q\twoheadrightarrow H$ for which
$\iota_H\eta$ satisfies that Brauer-vanishing condition.
Then
\[
           Q\big/\bigcap_H\bigcap_{K\in\mathcal K_H}K
\]
is topologically finitely generated. Its simple-factor
multiplicities satisfy $|\mathcal K_H|\le |H|^{C/\epsilon}$
for one constant $C$ independent of $H$.
Only the selected kernels are included; no assertion about
the other kernels of the same simple isomorphism types is made.
\end{corollary}
\begin{proof}
The relevant Brauer-zero theorem bounds the ambient projective
conjugacy classes by $q_H^{Cr_H^2}$. Distinct epimorphism kernels
give distinct such classes, because conjugation does not change
a kernel. The same uniform $C$ works for all the chosen fields
and ranks, so
\[
             |\mathcal K_H|\le q_H^{Cr_H^2}
                                  \le |H|^{C/\epsilon}.
\]
For precision, put $a=C/\epsilon$ and
$d=\lceil a\rceil+3$. Equations(2.1)--(2.2) in the proof of
Proposition2.6(ii), printed pp12--13, of
\citet{LucchiniMarionTracey2020D}, with its constant
$\gamma\ge1/2$, give
\[
 \psi_H(d)\ge\frac{|H|^{d-1}}{2|\operatorname{Out}(H)|}
              \ge\frac{|H|^{d-2}}2\ge |H|^a.
\]
Here $\psi_H(d)$ is the largest multiplicity of a direct power
of $H$ generated by $d$ elements. We used the outer-automorphism
bound on p13 of the same source, in particular
$|\operatorname{Out}(H)|\le |H|$, and $d(H)=2$.
Thus every $H^{|\mathcal K_H|}$ is generated by $d$ elements.
Products of different simple isomorphism types have no common
nontrivial simple quotients, so combining these tuples still
generates every finite subproduct, by Goursat's lemma.

Goursat's lemma identifies the image of $Q$ on every finite
subcollection of distinct simple quotient kernels with the full
product. Indeed a proper subdirect product of nonabelian simple
factors identifies two isomorphic simple quotients, and hence
identifies their kernels, contrary to their selection as distinct.
Thus the displayed quotient is the full profinite product of
the selected simple factors. The uniform finite-product generator
bound gives a generating tuple for this product by compactness:
in a fixed power of the product, the closed conditions that a
tuple generate each finite subproduct have the finite intersection
property. This proves the conclusion.
\end{proof}

\begin{remark}[Exact limitation]
This is a direct consequence of the Brauer-zero count, not an
additional movement event. The remaining projective simple
quotient kernels can have nonzero geometric Brauer class on the
relevant surface or local puncture. General monolithic crowns,
alternating multiplicities and the other uncovered families
remain separate.
\end{remark}


% Separate post-CP17 source: not part of the frozen checkpoint-17 manuscript.
\begin{corollary}[Proved: surface-detected Brauer-zero growth on complete local punctures]
\label{D:local-Brauer-zero-surface-transfer}
Let $(R,\mathfrak m,k)$ be a complete excellent normal local domain
of dimension $d\ge2$, of equal characteristic $p>0$, with an
algebraically closed coefficient field $k$. Put
$U=\operatorname{Spec}R\setminus\{\mathfrak m\}$.
There are finitely many fixed complete normal local surface
$k$-algebras $R_1,\ldots,R_t$, with residue field $k$, and
morphisms
\[
 h_i:U_i=\operatorname{Spec}R_i\setminus\{\mathfrak m_i\}
                         \longrightarrow U
\]
with the following property. There is a constant $C_R$ such that,
for every $q=p^f$ and $r\ge1$,
\[
 \#\left(\left\{\rho:\pi_1(U)\longrightarrow
       \operatorname{PGL}_r(\mathbf F_q):
       \operatorname{br}(h_i^*\rho)=0\text{ in }
              H^2(U_i,\mathbf G_m)\text{ for every }i
        \right\}/\operatorname{PGL}_r(\mathbf F_q)\right)
                       \le q^{C_Rr^2}.                 \tag{D.LBT.1}
\]
In particular the bound holds for all representations with
$\operatorname{br}(\rho)=0$ in $H^2(U,\mathbf G_m)$.
The scalar centre is not required to have bounded order.

Fix $\epsilon>0$ and, for each isomorphism type $H$ in a
collection of nonabelian finite simple groups, choose one
embedding $H\hookrightarrow\operatorname{PGL}_{r_H}(\mathbf F_{q_H})$
with $q_H$ a power of $p$ and
$|H|\ge q_H^{\epsilon r_H^2}$. The quotient of $\pi_1(U)$
formed by all distinct simple quotient kernels whose chosen
projective representations satisfy the displayed surface
Brauer-vanishing condition is topologically finitely generated.
\end{corollary}

\begin{proof}
We use the fixed slice in
Proposition~\ref{D:local-surface-reduction}, and recall its
geometric maps in order to keep the Brauer condition precise.
Choose $d-2$ members of a system of parameters, with ideal $I$,
and set
\[
       C=(R/I)_{\mathrm{red}},\qquad
       W=\operatorname{Spec}C\setminus\{\mathfrak m_C\}.
\]
For $d=2$, take $I=0$. Every irreducible component of $C$
has dimension two. Its finite normalization is a disjoint
union of complete normal local surface domains $R_i$.
Their residue fields are finite over $k$ and therefore equal
to $k$; the original coefficient embedding supplies their
coefficient fields. If
\[
            W'=\coprod_{i=1}^t U_i,
\]
then $W'\to W$ is finite and surjective. The maps $h_i$ are
the composites $U_i\to W\to U$. All these choices are fixed
independently of $q$, $r$, and the representation.

The homomorphism $\pi_1(W)\to\pi_1(U)$ is surjective.
Indeed a connected finite etale cover of $U$ normalizes to
a finite normal domain $B$ over $R$. Completeness and
henselianity make $B$ local and complete, with dimension $d$.
The ideal $IB$ has at most $d-2$ generators. The connectedness
and cohomological-dimension lemmas
\citep[Lemmas 52.14.1 and 51.4.3, Tags 0ECR and 0DX9]{C_StacksLocalSlice}
show that $\operatorname{Spec}(B/IB)$ minus its closed point
is connected. Here the connectedness dimension of the domain
$B$ is $d$, and $d-2<d-1$ is the required strict inequality.
Removing nilpotents does not affect connectedness or finite
etale covers. Thus every original connected cover remains
connected after pullback to $W$, which proves the surjection.
In particular, two representations of $\pi_1(U)$ which become
projectively conjugate on $W$ were already projectively
conjugate on $U$.

Put $G_{r,q}=\operatorname{PGL}_r(\mathbf F_q)$ and let
\[
                     a=\#\pi_0(W'\times_W W').
\]
This is a finite integer depending only on the fixed slice
and normalization. For each $i$, Theorem~\ref{D:local-Brauer-zero-projective-growth}
bounds the relevant $G_{r,q}$-torsor classes on $U_i$ by
$q^{C_i r^2}$, where $C_i$ depends only on $R_i$.
Consequently their tuples on $W'$ have at most
$q^{(\sum_i C_i)r^2}$ possibilities.

We give the exact bound for descent fibres. Fix one such
tuple, regarded as a $G_{r,q}$-torsor $P'$ on $W'$. A torsor
on $W$ whose pullback has this isomorphism class can be
identified with $P'$ upstairs and hence supplies a descent
isomorphism between the two pullbacks of $P'$ to
$W'\times_W W'$. On each connected component, a morphism
between finite etale schemes is determined by its value at
one geometric point. If these are equivariant isomorphisms
of $G_{r,q}$-torsors, there are at most $|G_{r,q}|$ possible
values. Thus there are at most $|G_{r,q}|^a$ descent
isomorphisms, before imposing the cocycle equations or
identifying isomorphic descent data. Effective finite-proper
descent is \citet[Expose IX, Theorem 4.7 and its proof]{C_SGA1}.
It follows that each pullback fibre on torsor isomorphism
classes has cardinality at most $|G_{r,q}|^a$.
This bound uses ambient projective torsors and therefore
introduces no comparison between simple-group conjugacy
and ambient projective conjugacy.

Since $|G_{r,q}|\le q^{r^2}$, restriction to $W$ and the
preceding estimates give
\[
        \#\{\text{classes in (D.LBT.1)}\}
              \le q^{(a+\sum_i C_i)r^2}.
\]
Brauer obstruction is functorial under each actual morphism
$h_i$, so original Brauer-zero representations are among
the counted classes. This proves the quantitative assertion.

For the final assertion, distinct simple quotient kernels
give distinct ambient projective conjugacy classes after
the chosen embedding. Their multiplicity for type $H$ is
therefore at most $|H|^{C_R/\epsilon}$. The generating-tuple
estimate and Goursat/compactness argument in the proof of
Corollary~\ref{D:Brauer-zero-selected-simple-product} apply
verbatim, with this fixed constant, and yield a uniform
number of generators for the resulting profinite product.
\end{proof}

\begin{remark}[Scope of the transfer]
Only pullback of Brauer classes is used. No injection from
$H^2(\pi_1(U),-)$ into the cohomology of the detecting
surfaces, and no injectivity of the Brauer pullback, is
asserted. The conclusion concerns complete equicharacteristic
normal local punctures; it gives no mixed-characteristic,
arbitrary-open, or full fundamental-group finite-generation
statement. It is a closure of the established Brauer-zero
reduction, not a separate advance on the remaining genuinely
nonzero-Brauer or general crown problems.
\end{remark}


% Separate post-CP17 research. Proved after full root/B/C/D actual-file and source audits.
\subsection{Two-generator charts with imperfect residue fields}

The coefficient field in this subsection is denoted by $\kappa$.
It need not be perfect.  The charts below fix a finite residue
extension, and make no claim about the poles of their coefficients
when $\kappa$ is a function field.

\begin{lemma}[Proved: a primitive residue element in transcendence degree one]
\label{D:one-variable-residue-primitive}
Let $k$ be perfect, let $\kappa/k$ be a finitely generated field
extension of transcendence degree one, and let $\lambda/\kappa$
be finite.  Then $\lambda=\kappa(\overline u)$ for some
$\overline u\in\lambda$, even when $\lambda/\kappa$ is inseparable.
\end{lemma}

\begin{proof}
Only characteristic $p>0$ requires discussion.  Write $\lambda_s$
for the maximal separable subextension and
$[\lambda:\lambda_s]=q=p^a$.  All the fields in question have
$p$-rank one.  A purely inseparable extension of degree $q$ is
contained in $\lambda_s^{1/q}$, and the latter also has degree
$q$.  Thus $\lambda=\lambda_s^{1/q}$.
Choose a primitive element $v_0$ of the separable extension
$\lambda_s/\kappa$ and $z\in\kappa\setminus\kappa^p$.
Separability gives $dz\ne0$ in $\Omega_{\lambda_s/k}$.
At least one of $v_0$ and $v_0+z$ has nonzero differential.
Call it $v$.  Translation by $z\in\kappa$ preserves primitivity,
and $dv\ne0$ says $v\notin\lambda_s^p$.
If $a>0$, the element $\overline u=v^{1/q}$ generates an extension
of degree $q$ over $\lambda_s$, and $\kappa(\overline u)$ contains
$\lambda_s=\kappa(v)$.  It is therefore $\lambda$.
For $a=0$ use $v$ itself.
\end{proof}

\begin{proposition}[Proved: finite flat sparse charts and exact finite isomorphism witnesses]
\label{D:sparse-DVR-linear-atlas}
Fix a field $\kappa$, integers $e,f\ge1$, a monic irreducible
polynomial $\psi\in\kappa[U]$ of degree $f$, and an integer
$\delta\ge0$.  Put
\[
 n=ef,\qquad N=\left\lfloor\frac{2\delta}{e}\right\rfloor+2.
                                                        \tag{D.SD.1}
\]
There is an affine polynomial chart over $\kappa$ with at most
$(3N+2)n$ coordinates, and a polynomial witness for equivalence
of two chart objects with at most $(8N+4)n+1$ coordinates
including both objects.  All equations may be taken of degree
at most
\[
                         (4N+12)n^3+1.                 \tag{D.SD.2}
\]
The assertions have the following precise meaning.

For every $\kappa$-algebra $R$, a chart point gives a finite free
$R[[t]]$-algebra $C$ of rank $n$, with elements $u,s$ such that
$s^e=t$ times a unit, $s$ is a nonzerodivisor,
$C/(s)=R[U]/(\psi)$, and $C[1/t]$ is finite etale over
$R((t))$.  Its relative K\"ahler different is $(s^\delta)$.
Every $\kappa$-valued object is a complete DVR with uniformizer
$s$, ramification index $e$, residue field $\kappa[U]/(\psi)$,
and different exponent $\delta$.

For two chart points over any $R$, the finite witness has an
$R$-point if and only if their complete algebras are continuously
isomorphic over $R[[t]]$.  This is an existence test, not an
assertion that a prescribed truncated isomorphism lifts without
loss of precision.  No marking of the residue algebra or
augmentation is imposed on the isomorphisms.

Finally, every finite complete DVR $B$ over $\kappa[[t]]$ with
separable fraction field, these invariants, and residue field
generated by a root of $\psi$, is isomorphic to a chart object.
\end{proposition}

\begin{proof}
\emph{The two equations and their topology.}
Let
\[
 \begin{split}
 P&=U^f-\sum_{0\le i<f,\,0\le j<e} a_{ij}(t)U^iS^j,\\
 Q&=S^e-t\sum_{0\le i<f,\,0\le j<e} b_{ij}(t)U^iS^j,
 \end{split}                                            \tag{D.SD.3}
\]
where all $a_{ij},b_{ij}$ have degree less than $N$ in $t$.
Impose $P(0,U,0)=\psi(U)$.  Write $V=\sum b_{ij}U^iS^j$.
Add $f$ coordinates for an inverse of $V(0,U,0)$ in
$R[U]/(\psi)$, and impose the corresponding $f$ product equations.
Define
\[
 C=\varprojlim_m R[t,U,S]/(t^m,P,Q).
                                                        \tag{D.SD.4}
\]
This is the $t$-adic completion, with $U$ an ordinary polynomial
variable.  In particular, $U$ is not assumed topologically
nilpotent.  Over $\kappa$ the chosen irreducible $\psi$ specifies
the unique residue field; over a general $R$ we retain the whole
finite residue algebra $R[U]/(\psi)$.

We prove, including nonreduced coefficient rings, that
\[
                 \{u^is^j:0\le i<f,\ 0\le j<e\}
                                                        \tag{D.SD.5}
\]
is a basis over $R[[t]]$.  First work over the Noetherian universal
coefficient ring $R_0$ for the displayed finitely many coefficients
and inverse equations.  In $R_0[[t]][U,S]$, the sequence
$(t,P,Q)$ is regular.  Indeed $t$ is regular, $\overline P$
is monic in $U$, its quotient is free over $R_0[S]$, and
$\overline Q=S^e$ is regular on that quotient.
At primes containing $(t,P,Q)$, regular sequences permute
\citep[Lemma~10.68.4, complete proof, Tag 00LJ]{StacksSparseDVRB}.
Consequently multiplication by $t$ on the quotient by $(P,Q)$
has kernel supported away from $V(t)$.  Completion is exact on
finite modules over this Noetherian ring, so the completed
quotient $C_0$ is $t$-torsion free.

Modulo $t$, the box (D.SD.5) is a basis over $R_0$:
first reduce by $S^e$, and then by the monic polynomial
$\overline P$ in $U$.  Successively lifting coefficients modulo
$t,t^2,\ldots$ gives a surjection
$R_0[[t]]^n\to C_0$.  The $t$-torsion freeness shows that its
kernel reduces to zero modulo $t$.  The kernel is finite,
and $t$ lies in the Jacobson radical of $R_0[[t]]$; Nakayama
therefore kills it.  This proves the basis assertion over $R_0$.
Modulo each $t^m$ the same quotient commutes with arbitrary
coefficient base change $R_0\to R$.  Passing to the inverse
limit proves (D.SD.5) for every $R$.  More precisely,
$C=C_0\otimes_{R_0[[t]]}R[[t]]$, since the finite free box basis
and its multiplication constants have these reductions.
Thus this is a genuine all-ring flatness proof, independent
of the later field-point DVR argument.

The element $V$ is a unit modulo $(t,S)$ by construction,
hence modulo $t$ because $S$ is nilpotent there, and hence
in $C$ by completeness.  Thus $s^e=tV$.
Multiplication by $t$ is injective on the free module $C$,
so $s$ is a nonzerodivisor, and the $s$-adic and $t$-adic
topologies are cofinal.  The same identity shows
$C/(s)=R[U]/(\psi)$.

\emph{Jacobian stratum and different.}
Put
\[
                 J=\det\frac{\partial(P,Q)}{\partial(U,S)}.
                                                        \tag{D.SD.6}
\]
Add $nN$ coefficients for an element $W$ of $C/t^NC$, and
$f$ coefficients for an inverse of its reduction in $C/(s)$.
Require
\[
                         J=s^\delta W\pmod {t^NC}.
                                                        \tag{D.SD.7}
\]
Every equation here is an equation in the box basis, using
the reduction procedure below.  Because $eN>2\delta$,
the actual Jacobian in $C$ is $s^\delta$ times a unit.
The equations present the relative differential module by
the two-by-two Jacobian matrix.  One may verify this first
over the Noetherian universal ring: finite differential
modules commute with completion, and the usual conormal
sequence applies before completion.  The assertion over
every $R$ then follows by the finite free base change just proved.
Hence
\[
                  \operatorname{Fitt}_0\Omega_{C/R[[t]]}
                                      =(s^\delta).       \tag{D.SD.8}
\]
After inverting $t$, the Jacobian is invertible, so the
finite free algebra is etale.  Over the universal ring,
the fibres at $t=0$ are the zero-dimensional complete
intersections $(\overline P,S^e)$; the other fibres are etale.
Finite flatness and the fibre criterion therefore make this
a finite syntomic morphism
\citep[Lemma~49.11.1(3), Tag 0FKX]{StacksSparseDVRB}.
This property persists under all the base changes above.
Over a DVR the usual different equals this K\"ahler different,
by the finite syntomic Jacobian formula
\citep[Lemmas~49.12.1--49.12.3 and their proofs,
Tag 0BWB]{StacksSparseDVRB}.

For $R=\kappa$, $C/(s)$ is a field and $s$ lies in the Jacobson
radical.  Thus $C$ is local with maximal ideal $(s)$.
It is finite free over the one-dimensional DVR $\kappa[[t]]$,
so has dimension one.  It is a regular local ring, hence a DVR.
The identity $s^e=tV$ and (D.SD.8) give the asserted invariants.
This assertion about field points is separate from flatness
over arbitrary coefficient rings.

\emph{A Newton estimate over arbitrary coefficient rings.}
We will use the following elementary two-variable form.
In a complete ring in which $s$ is regular, suppose a polynomial
pair $F$ has error in $(s^a)$ and its Jacobian determinant at
the chosen point is $s^\delta$ times a unit, where $a>2\delta$.
The adjugate formula gives a Newton correction in
$(s^{a-\delta})$.  Polynomial Taylor expansion, without
division by factorials, puts the new error in
$(s^{2(a-\delta)})$.
The change in the Jacobian determinant is divisible by
$s^{a-\delta}$, so it remains $s^\delta$ times a unit.
Iteration converges, because
$a\mapsto2(a-\delta)$ increases to infinity.
It produces an exact zero differing from the initial point
by $(s^{a-\delta})$.  Regularity of $s$ justifies division by
$s^\delta$; the proof allows arbitrary nilpotents in $R$.

\emph{Every DVR extension occurs.}
Let $B$ be as in the statement.  Choose a uniformizer $s$ and
a lift $u$ of the specified residue generator.  The box
$u^is^j$ is a basis over $\kappa[[t]]$: its reduction modulo
$t$ is a basis by the filtration by powers of $s$, whose
successive residue quotients have dimension $f$ over $\kappa$.
Expand $u^f$ and $s^e/t$ in that basis.  This gives (D.SD.3)
with power-series coefficients, $P(0,U,0)=\psi$, and $V$ a unit.
The preceding basis proof also works with these coefficients
over $\kappa[[t]]$, and the resulting map to $B$ is an
isomorphism modulo $t$, hence an isomorphism of finite free
modules.  Thus these two equations present $B$ topologically.

For completeness, the finite flat algebra $B$ is syntomic:
its closed fibre has the just proved complete-intersection
presentation, and its generic fibre is separable.
The Jacobian of the two equations consequently has order
$\delta$ by the same different formula.
Truncate all coefficient series modulo $t^N$.
At $(u,s)$ the errors have order at least $eN$, and the
Jacobian still has order $\delta$.
Newton gives exact roots of the truncated equations differing
from $(u,s)$ by $(s^{eN-\delta})$.
Here $eN-\delta\ge e$, so the box basis is unchanged modulo
$t$.  The truncated model maps isomorphically to $B$.
Its Jacobian is $s^\delta$ times a unit, so a choice of $W$
and of the two finite inverse witnesses puts it in our chart.

\emph{Finite witnesses for continuous isomorphism.}
For two chart objects $C_1,C_2$, introduce two box elements
$\eta,\beta\in C_1/t^NC_1$, with $nN$ coefficients each.
Require
\[
                     P_2(\eta,\beta)=Q_2(\eta,\beta)=0
                                             \pmod {t^N}.\tag{D.SD.9}
\]
Form the $n$ by $n$ matrix whose columns are the box expansions
of $\eta^i\beta^j$, in the order (D.SD.5).
Add one inverse variable for the constant coefficient in $t$
of its determinant.  Thus (D.SD.9) defines an isomorphism
\[
                         C_2/t^NC_2\simeq C_1/t^NC_1.    \tag{D.SD.10}
\]
Conversely every such truncated isomorphism has these witnesses.

Lift $\eta,\beta$ to $C_1$ by their box polynomials.
The Fitting ideal of the differential module of (D.SD.10)
is invariant under this isomorphism.  Therefore, writing
$s=s_1$, the ideals generated modulo $t^N$ by
$J_2(\eta,\beta)$ and $s^\delta$ agree.
In $C_1$ write
$J_2(\eta,\beta)=s^\delta a+t^N c$ and
$s^\delta=J_2(\eta,\beta)b+t^N d$.
Since $t^N$ is $s^{eN}$ times a unit, cancellation of the
nonzerodivisor $s^\delta$ shows that the relevant coefficient
$a+s^{eN-\delta}c'$ has a multiplicative inverse modulo
$s^{eN-\delta}$.  It is a unit by completeness.
Thus $J_2(\eta,\beta)$ is actually $s^\delta$ times a unit.
The Newton estimate, with $a=eN$, now corrects the two
approximate images to exact roots.  Their corrections lie
in $(s^{eN-\delta})\subset(t)$.
Their box-basis matrix remains invertible modulo $t$,
so the resulting continuous $R[[t]]$-algebra map is an
isomorphism.  Conversely any continuous isomorphism truncates
to (D.SD.9)--(D.SD.10).
This proof imposes neither a pointed condition on $\beta$
nor a prescribed residue-field automorphism.

\emph{Number and degrees of equations.}
There are $2nN$ coefficient variables in (D.SD.3), $f$ for
the inverse of $V$ modulo $s$, $nN$ for $W$, and $f$ for
its inverse modulo $s$.  This gives $(3N+2)n$ as an upper
bound.  Two objects, the two images, and the determinant
inverse give $(8N+4)n+1$.  All identities amount to
$O(Nn)$ box coefficient equations.

Here is an explicit degree bound, so that finite type alone
is not being used quantitatively.  Reduce a monomial modulo
$t^N$ by using $Q$ whenever its $S$-exponent is at least $e$,
then using $P$ when its $U$-exponent is at least $f$.
Every $Q$ step increases the $t$-exponent, so there are at
most $N-1$ such steps in a surviving term.
Each $P$ step lowers the $U$-exponent by at least one;
a $Q$ step raises it by at most $f-1$.
A term of initial $U$-degree $A$ thus uses at most
$A+(N-1)(f-1)$ steps of type $P$ and at most $N-1$
steps of type $Q$.  Each step adds at most one to its
degree in the coefficient parameters.

Before reduction, the substituted equations and box-matrix
entries have $U$-degree at most $n^2$, and degree in
their image parameters at most $n+1$.
Their reduced coefficients therefore have degree at most
$n^2+(N+1)n+1$.
The $n$ by $n$ determinant and its inverse equation have
degree at most $n^3+(N+1)n^2+n+1$.
The Jacobian, unit, and stratum equations have no larger
degree than the more generous common bound (D.SD.2).
This proves all the stated coordinate and degree estimates.
\end{proof}

\begin{corollary}[Proved: linear local coordinates under a normalized different bound]
\label{D:sparse-DVR-normalized-different-complexity}
Fix $\Lambda\ge0$.  In the preceding proposition, if
$\delta/e\le\Lambda$, the object and isomorphism-witness
spaces have $O_\Lambda(n)$ coordinates and equations of
degree $O_\Lambda(n^3)$.  Consequently the number of their
geometric irreducible components is at most
$n^{O_\Lambda(n)}$ (with $n=1$ handled by a fixed constant).
For a one-variable residue field over perfect constants,
Lemma~\ref{D:one-variable-residue-primitive} supplies the
required primitive residue element for every finite DVR
extension.
\end{corollary}

\begin{proof}
The cutoff in (D.SD.1) is at most $2\Lambda+2$.
Apply the explicit preceding estimates and the same elementary
affine component Bezout bound used in
Theorem~\ref{D:uniform-local-representation-growth}.
The residue lemma does not require its extension to be separable.
\end{proof}

\begin{remark}[Unresolved global input]
These are charts over the fixed residue field $\kappa$, for a
fixed residue polynomial $\psi$.  When $\kappa=k(E)$, none of
the estimates bounds the poles on $E$ of $\psi$ or of the
coefficient functions $a_{ij},b_{ij}$.
Nor does a component estimate for these local witness schemes
bound a constrained global extension locus.
Constructing a residue-cover and coefficient presentation
with controlled global poles, and effective compatibility
with the exceptional neighbourhood, remains necessary.
No finite-generation conclusion is drawn from this subsection.
\end{remark}


% Root/B/C/D full actual-file audits passed. Separate post-CP17 closure.
\subsection{Bounded geometric Brauer period}

\begin{theorem}[Proved: quadratic growth at bounded geometric Brauer period]
\label{D:bounded-geometric-Brauer-period}
Fix an integer $B\ge1$.  Let $G$ be one of the following groups,
over an algebraically closed field $k$ of characteristic $p>0$:
\begin{enumerate}
\item $G=Q(S,T)$, where $S$ is a normal proper connected surface,
      $T\subsetneq S$ is closed, and $U=S-T$;
\item $G=\pi_1(U)$, where
      $U=\operatorname{Spec}R-\{\mathfrak m\}$ for a complete
      excellent normal local $k$-algebra $(R,\mathfrak m)$ of
      dimension two with residue field $k$.
\end{enumerate}
For a continuous representation
$\rho:G\to\operatorname{PGL}_r(\mathbf F_q)$, $q=p^f$,
write $\operatorname{br}(\rho)\in H^2(U,\mathbf G_m)$
for the Brauer class of the induced algebraic projective-linear
torsor.  There is a constant $C_{G,B}$, independent of $r$ and $f$,
such that
\[
 \#\left(
  \left\{\rho:G\to\operatorname{PGL}_r(\mathbf F_q):
       \operatorname{ord}(\operatorname{br}(\rho))\le B\right\}
        /\operatorname{PGL}_r(\mathbf F_q)\right)
       \le q^{C_{G,B}r^2}.                         \tag{D.BP.1}
\]
The order of the zero class is one.  The scalar order $q-1$
is unrestricted.  More precisely, one fixed open normal
subgroup $H_B\subset G$ makes every representation counted
in (D.BP.1) geometrically Brauer-zero on its corresponding
cover of $U$.  Its choice can depend on $B$; no uniform
bound on $[G:H_B]$ as $B$ grows is asserted.
\end{theorem}

\begin{proof}
\emph{A finite set of realized geometric classes.}
Put $\operatorname{Br}'(U)=H^2(U,\mathbf G_m)_{\mathrm{tors}}$.
For each integer $m$ prime to $p$, the Kummer sequence makes
$\operatorname{Br}'(U)[m]$ a quotient of $H^2(U,\mu_m)$
\citep[Lemma~59.28.1, Tag 03PK]{StacksBrauerZeroD}.
These cohomology groups are finite.  In the complete-local
case this is the explicit etale-cohomology bound of
Theorem~\ref{D:local-surface-constant-H2}.  In the global
case it follows, with constant coefficients, from the proof
of Theorem~\ref{D:surface-tame-cross-characteristic-H2}:
proper birational Leray bounds $H^2(U,\mathbf F_\ell)$
by the cohomology of the fixed regular open on a smooth
projective resolution and a finite exceptional contribution;
the constant sheaf has zero ramification, so the cited
Betti bound applies.  A composition series for the constant
cyclic coefficient group gives the assertion for every $m$.
Tate twists are constant over $k$.

Let $\mathcal A_B$ be the set of classes $\alpha$ with
order at most $B$ that are \emph{actually realized} as
$\operatorname{br}(\rho)$ by some finite-field projective
representation of $G$, allowing every rank and every
$q=p^f$.  Such a class is killed by $q-1$, since its lifting
class comes from the central extension
\[
 1\longrightarrow\mathbf F_q^\times
 \longrightarrow\operatorname{GL}_r(\mathbf F_q)
 \longrightarrow\operatorname{PGL}_r(\mathbf F_q)
 \longrightarrow1.
\]
Consequently its order is prime to $p$, and
\[
       \mathcal A_B\subset
       \bigcup_{\substack{1\le m\le B\\(m,p)=1}}
                              \operatorname{Br}'(U)[m]
                                                        \tag{D.BP.2}
\]
is finite.  We make no assertion that every element of the
right-hand side is realized by a projective local system.

For each nonzero $\alpha\in\mathcal A_B$, choose one
witnessing representation
$\rho_\alpha:G\to\operatorname{PGL}_{r_\alpha}
(\mathbf F_{q_\alpha})$.  Define
\[
             H_B=\bigcap_{\alpha\in\mathcal A_B,
                            \,\alpha\ne0}\ker(\rho_\alpha).
                                                        \tag{D.BP.3}
\]
An empty intersection means $G$.  This is open and normal.
Let $V\to U$ be the connected finite Galois cover obtained
by composing $\pi_1(U)\to G$ with $G/H_B$.
Each witnessing projective torsor becomes trivial on $V$,
so every class of $\mathcal A_B$ pulls back to zero there.
Therefore every representation in (D.BP.1), not only the
chosen witnesses, has Brauer-zero restriction on $V$.
The fields and ranks of the finitely many witnesses may
differ; after they have been chosen the cover $V$ is fixed.

\emph{The fixed cover lies in the same geometric class.}
In case (1), normalize $S$ in $k(V)$ and call the result
$S'$.  Excellence makes $S'\to S$ finite; $S'$ is a normal
proper connected surface, and, with $T'$ the reduced inverse
image of $T$, its complement is exactly $V$.  Moreover this
cover was chosen through $G=Q(S,T)$, so it is tame at the
original boundary divisors.  Every finite quotient of the
$H_B$-action obtained by restricting a $G$-cover is tame at
the codimension-one points of $T'$: these points lie over
codimension-one points of $T$, and tameness is preserved
under the corresponding extension of discrete valuation
rings.  Thus the natural surjection
$\pi_1(V)\twoheadrightarrow H_B$ factors as
\[
               Q(S',T')\twoheadrightarrow H_B.
                                                        \tag{D.BP.4}
\]
For an arbitrary finite quotient of $H_B$, one may first
replace its open kernel by its core in $G$; this shows that
the assertion about restricted $G$-covers tests every
finite quotient of $H_B$.  The valuation comparison is the
same finite-normalization and tame-base-change argument as
Theorem~\ref{D:finite-normal-tame-descent}, using
\citep[Lemma~15.116.9(2), Tag 0EXY]{StacksDFiniteTameDescent}.
In particular no cover outside the original tame quotient
is used to choose (D.BP.3).

In case (2), the normalization $R'$ of $R$ in $k(V)$ is a
finite complete normal local domain.  Indeed finiteness
follows from excellence, and a finite domain over a
complete henselian local ring has only one local factor.
Its residue field is a finite extension of $k$, hence $k$,
and the given coefficient field remains a coefficient
field.  Its puncture is $V$.  The subgroup $H_B$ is then
$\pi_1(V)$, with the compatible base points understood.

Apply Theorem~\ref{D:Brauer-zero-projective-representation-growth}
to $(S',T')$ in case (1), and
Theorem~\ref{D:local-Brauer-zero-projective-growth} to $R'$
in case (2).  In (1), pullback along (D.BP.4) injects
conjugacy classes of $H_B$-representations into those of
$Q(S',T')$.  Thus there is one constant $C'$ for this fixed
cover such that the possible restricted conjugacy classes
in rank $r$ over $\mathbf F_q$ number at most
$q^{C'r^2}$.

\emph{The restriction fibres.}
Put $e=[G:H_B]$ and $F=\operatorname{PGL}_r(\mathbf F_q)$.
Fix one representative $\sigma:H_B\to F$ of a restricted
conjugacy class.  Every extension class can be conjugated
to have restriction exactly $\sigma$.  After choosing $e$
coset representatives for $H_B$ in $G$, such a homomorphism
is determined by its values on those representatives and
on $H_B$.  Hence there are at most $|F|^e$ extension
classes.  Since $|F|\le q^{r^2}$, the total is at most
$q^{(C'+e)r^2}$, which proves (D.BP.1).
\end{proof}

\begin{corollary}[Proved: simple products selected by bounded geometric period]
\label{D:bounded-Brauer-period-selected-simple-product}
In either case of Theorem~\ref{D:bounded-geometric-Brauer-period},
fix $B$ and $\epsilon>0$.  For each isomorphism type $L$ in
a collection of nonabelian finite simple groups, choose one
faithful embedding
\[
       \iota_L:L\hookrightarrow
                    \operatorname{PGL}_{r_L}(\mathbf F_{q_L}),
       \qquad q_L=p^{f_L},\qquad
                    |L|\ge q_L^{\epsilon r_L^2}.
\]
The quotient of $G$ selected by all distinct kernels of
epimorphisms $\eta:G\twoheadrightarrow L$ for which
$\operatorname{br}(\iota_L\eta)$ has order at most $B$
is topologically finitely generated.
\end{corollary}
\begin{proof}
Distinct kernels give distinct ambient projective conjugacy
classes.  The preceding theorem bounds their multiplicity
by $|L|^{C_{G,B}/\epsilon}$.  Apply exactly the generating-tuple,
Goursat and compactness argument of
Corollary~\ref{D:Brauer-zero-selected-simple-product}, with
$C_{G,B}$ in place of its constant $C$.
\end{proof}

\begin{remark}[Scope of the closure]
This is a finite-cover consequence of the Brauer-zero count,
not a second movement event.  The remaining obstruction in
this projective-representation route has unbounded geometric
Brauer period.  No bound uniform in $B$, no cross-characteristic
representation bound, and no control of general nonabelian
crowns is obtained here.
\end{remark}


% Complete author/root/B/D actual-file audits PASS; auxiliary object-chart input.
\subsection{Uniform generation and fixed-determinant charts on exceptional cycles}

\begin{theorem}[Proved: uniform generation of reflexive pullbacks]
\label{D:reflexive-pullback-uniform-generation}
Let $(A,\mathfrak m,k)$ be an excellent normal local surface
$k$-algebra with algebraically closed residue field $k$ and a
dualizing complex. Let $\pi:Y\to\operatorname{Spec}A$ be a proper
regular resolution, with exceptional curves $E_i$, and let $Z>0$
be an effective exceptional Cartier divisor of full support.
Assume that
\[
               -Z\cdot B\ge 2p_a(B)
       \quad\text{for every nonzero effective exceptional cycle }B.
                                                               \tag{D.UG.1}
\]
If $F$ is a vector bundle on $Y$ admitting a map
$\mathcal O_Y^a\to F$ which is surjective at every exceptional
generic point, then
\[
 H^1(Y,F(-Z))=0,\qquad
 H^1(Y,F(-Z)\otimes I_{P,Y})=0\quad(P\in E),             \tag{D.UG.2}
\]
and $F(-Z)$ is generated by its global sections. Here $P$ ranges
over closed points of the exceptional fibre.

In particular the assertion applies to
$F=(\pi^*M)^{\vee\vee}$ for every finite reflexive $A$-module
$M$, with a single choice of $Z$ independent of $M$ and its rank.
On a minimal resolution with
$\kappa_i=K_Y\cdot E_i\ge0$, it suffices to choose the cycle in
Lemma~\ref{D:integral-antinef-energy-cycle} so that
\[
 -Z\cdot E_i\ge3\kappa_i\quad(\kappa_i>0),
 \qquad -Z\cdot E_i\ge0\quad(\kappa_i=0).              \tag{D.UG.3}
\]
An entirely zero-canonical connected component can instead be
given its rational-double-point fundamental cycle.
\end{theorem}

\begin{proof}
\emph{The point-ideal form of automatic adjunction.}
Let $D$ be an effective exceptional cycle, $L$ a line bundle
on $D$, and $P$ a closed point of $D$. Suppose that
\[
                  \deg_B L\ge2p_a(B)
       \quad\text{for every nonzero subcycle }B\le D.    \tag{D.UG.4}
\]
Then
\[
       \operatorname{Hom}_D(I_{P,D}\otimes L,\omega_D)=0.
                                                               \tag{D.UG.5}
\]
We recall the proof, including the nilpotent layers. A nonzero
map on the left has a maximal scheme-theoretic supporting
subcurve $B$. Automatic adjunction factors it through
$\omega_B=\mathcal Hom_D(\mathcal O_B,\omega_D)$, with a
generically surjective map from its restricted source; see
\citet[Lemma~2.4 and the complete proof of Theorem~1.1,
pp.~192--193]{C_CFHR1999}. The annihilator has no embedded
associated point because the target is Cohen--Macaulay.
Since $Y$ is regular, a Cohen--Macaulay subcurve of its Cartier
divisor $D$ is precisely an effective Cartier subcycle $B\le D$.
The map from $I_{P,D}\otimes\mathcal O_B$ to
$I_{P\cap B,B}$ has a zero-dimensional kernel. That kernel
maps to zero in $\omega_B$. We therefore obtain a generically
surjective morphism
\[
                  I_{P\cap B,B}\otimes L|_B\longrightarrow\omega_B.
\]
Both sides have the same generic lengths; the source is
torsion-free. The morphism is consequently injective and has
zero-dimensional cokernel. Euler characteristics give
\[
 \deg_B L\le 2p_a(B)-2+\ell(P\cap B)\le2p_a(B)-1,
\]
contradicting (D.UG.4). The same argument with the empty cluster
gives $\operatorname{Hom}_D(L,\omega_D)=0$.
This is the full nonreduced-cycle argument, rather than an
inequality tested only on the reduced prime components.

\emph{Insertion of a generically generated vector bundle.}
Restrict the given evaluation map to $D=tZ$, where $t\ge1$.
It remains surjective at every generic Artinian stalk of $D$.
The target $\omega_D$ has no nonzero zero-dimensional subsheaf.
Precomposition therefore gives an injection
\[
 \operatorname{Hom}_D(I_{P,D}\otimes F|_D(-Z),\omega_D)
 \lhook\joinrel\longrightarrow
 \operatorname{Hom}_D(I_{P,D}\otimes\mathcal O_D(-Z),\omega_D)^a.
                                                               \tag{D.UG.6}
\]
Indeed a map in the kernel vanishes at every generic point,
so its image would be a zero-dimensional subsheaf of $\omega_D$.
By (D.UG.1), the right side vanishes by (D.UG.5). Serre duality
on the proper Gorenstein curve $D$ gives
\[
 H^1(D,F(-Z)|_D\otimes I_{P,D})=0,
 \qquad H^1(D,F(-Z)|_D)=0.                              \tag{D.UG.7}
\]

There is one Tor term to retain when passing from $Y$ to $D$:
\[
 \operatorname{Tor}_1^Y(k(P),\mathcal O_D)
       \longrightarrow I_{P,Y}\otimes\mathcal O_D
       \longrightarrow I_{P,D}\longrightarrow0.
                                                               \tag{D.UG.8}
\]
Its image is zero-dimensional. Tensoring with the vector
bundle $F(-Z)|_D$ shows that it does not change $H^1$.
Thus (D.UG.7) also computes the restrictions of
$F(-Z)\otimes I_{P,Y}$.
The ideals of $Z$ and of the closed fibre have the same
radical, so their powers are cofinal. The theorem on formal
functions now makes the completions of both finite $A$-modules
in (D.UG.2) zero. Faithfulness of completion, or Nakayama's
lemma, gives (D.UG.2) itself.

The point-ideal exact sequence implies surjectivity
$H^0(Y,F(-Z))\to F(-Z)\otimes k(P)$ at every $P\in E$.
The evaluation cokernel is coherent. Its proper image in the
local base, if nonempty, would contain the closed point, whereas
its support misses the entire closed fibre. The cokernel is
therefore zero. This proves global generation on $Y$.

For a reflexive $M$, take a finite generating set over $A$.
At a codimension-one regular stalk of $Y$, the natural map
$\pi^*M\to(\pi^*M)^{\vee\vee}$ is onto: the target is the
free quotient of a module over a discrete valuation ring.
Its generating sections therefore have the required generic
surjectivity. No claim that the original reflexive pullback
was already globally generated at every closed point is used.

Finally write $B=\sum m_iE_i$. Adjunction and negative
definiteness give
\[
          2p_a(B)=2+\kappa^{\mathsf t}m+B^2.
\]
If $\kappa^{\mathsf t}m\ge1$, then $B^2\le-1$ and the right
side is at most $\kappa^{\mathsf t}m+1\le2\kappa^{\mathsf t}m$.
Condition (D.UG.3) is stronger than (D.UG.1) in this case.
If $\kappa^{\mathsf t}m=0$, the support consists of smooth
rational $(-2)$-curves. Its integral intersection form is even
and negative definite, so $B^2\le-2$ and $p_a(B)\le0$.
Anti-nefness of $Z$ suffices. A lattice parameter $c\ge5$ in
Lemma~\ref{D:integral-antinef-energy-cycle} gives (D.UG.3).
For a whole zero-canonical component the same parity argument
applies with its fundamental cycle. This proves the last claim.
\end{proof}

\begin{proposition}[Proved: degree-linear fixed-determinant cycle charts]
\label{D:hereditary-fixed-determinant-cycle-chart}
Let $R$ be a fixed complete excellent normal local surface
$k$-algebra of characteristic $p>0$, with algebraically closed
coefficient and residue field $k$. Consider finite normal
quasi-etale cover domains $A/R$ of rank $n$ in the energy
family of Corollary~\ref{D:punctured-cover-linear-energy}.
Here the energy hypothesis means that the stated bounds for
$p_g$, the canonical energy and the total exceptional
self-intersection hold with constants depending only on
the fixed base. In particular one may complete the algebraic
local cover family in that corollary: excellence preserves
regularity of the base-changed resolution, and proper flat
base change preserves the exceptional curves, intersection
numbers and coherent-cohomology lengths. No energy assertion
for arbitrary nonalgebraizable complete covers is added here.
Write $Y$ for the minimal regular resolution of $A$,
and choose $Z$ as in Theorem~\ref{D:reflexive-pullback-uniform-generation}
and Corollary~\ref{D:punctured-cover-antinef-cycle}.
The choice depends on $A$, but not on a module or its rank.
Fix $t\ge1$ and put $D=tZ$.

Let $(M_a)_{a\in\mathbb Z/f\mathbb Z}$ be a cyclic system of
rank-$r$ reflexive $A$-modules with
$F_A^{[*]}M_a\simeq M_{a+1}$. Put
\[
       F_a=(\pi^*M_a)^{\vee\vee},\qquad V_a=F_a(-Z)|_D.
\]
There is a constant $c_t$, independent of $n,r,f$ and this
system, such that
\[
 0\le c_1(F_a)\cdot Z\le c_tnr,
 \qquad h^0(D,\det V_a)\le c_tnr.                    \tag{D.UG.9}
\]
Every $V_a$ is globally generated. For each fixed line bundle
$L=\det V_a$ on $D$, all globally generated rank-$r$ bundles
on $D$ with determinant isomorphic to $L$ occur as the
universal quotients over an open subset of
\[
                 H^0(D,L)^{\oplus(r+1)}.
                                                               \tag{D.UG.10}
\]
In particular the relevant fixed-determinant object chart has
at most $c_tnr(r+1)=O_t(nr^2)$ scalar coordinates. Its universal
quotient is an actual vector bundle, also over nonreduced
parameter schemes.

This is a statement about the objects on one finite exceptional
cycle. It asserts neither that every such bundle is the
restriction of a desired reflexive unit object, nor that a
unit object is determined by this restriction.
\end{proposition}

\begin{proof}
\emph{A fixed ideal and its energy.}
Use the nonzero fixed $\mathfrak m_R$-primary Noether-different
ideal $J$ of Proposition~\ref{D:cover-uniform-Noether-different}.
It factors multiplication on every reflexive $A$-module through
a finite free $A$-module. Let $C$ be the divisorial part of
$J\mathcal O_Y$. The line bundle $\mathcal O_Y(-C)$ is
generated at the exceptional generic points by the fixed
generators of $J$. Hence $C$ is anti-nef.

Principalize $J$ once on a fixed regular modification $Y_0$
of the base, writing $J\mathcal O_{Y_0}=\mathcal O_{Y_0}(-C_0)$.
Resolve the finite normalization above $Y_0$ and take a common
regular resolution $Y'$ dominating $Y$. The divisor of $J$
on $Y'$ is $C'=f^*C_0$, where $f:Y'\to Y_0$ has generic
degree $n$. Projection and the orthogonal decomposition for
the birational map $g:Y'\to Y$ give
\[
 (C')^2=nC_0^2,\qquad C'=g^*C+B_g,\qquad
 (C')^2=C^2+B_g^2\le C^2.
\]
The divisor $B_g$ is exceptional for $g$; no sign of its
coefficients is needed for the last inequality. Consequently
\[
                  -C^2\le n(-C_0^2).
                                                               \tag{D.UG.11}
\]
All these intersections are exceptional intersections over the
local base. The already proved energy estimate for $Z$ and
Cauchy--Schwarz for the negative intersection form imply
\[
       0\le\lambda:=-C\cdot Z
           \le\sqrt{(-C^2)(-Z^2)}=O(n).               \tag{D.UG.12}
\]

\emph{The determinant recurrence.}
The source factorization through a free module shows that
$J^{[p]}$ kills the ordinary-to-reflexive Frobenius cokernel.
Pull back the ordinary Frobenius linearization and extend
across codimension two on regular $Y$. Frobenius is flat on
$Y$, so the resulting inclusions of vector bundles are
\[
        F_{a+1}(-pC)\ \subset\ F_Y^*F_a\ \subset\ F_{a+1}.
                                                               \tag{D.UG.13}
\]
The first inclusion can be checked at every exceptional
discrete valuation ring: the least valuation of a $p$-th
power of a generator of $J$ is $p$ times the coefficient
of $C$. The resulting rational map is regular everywhere
by reflexivity. Thus the effective determinant divisor
$B_a$ of the second inclusion satisfies
\[
       c_1(F_{a+1})=p\,c_1(F_a)+B_a,
                    0\le B_a\le rpC.                 \tag{D.UG.14}
\]
Generic generation implies $\deg(\det F_a|_{E_i})\ge0$.
Put $x_a=c_1(F_a)\cdot Z\ge0$. Since $Z$ is anti-nef,
(D.UG.14) gives
\[
                   p x_a-rp\lambda\le x_{a+1}\le p x_a.
\]
Iterate around the $f$-cycle. Cancellation of $p^f-1$ yields
\[
                   0\le x_a\le\frac{rp\lambda}{p-1}=O(nr),
                                                               \tag{D.UG.15}
\]
independently of the period $f$.

\emph{Determinants on the full thickening.}
Theorem~\ref{D:reflexive-pullback-uniform-generation} gives
global generation of $F_a(-Z)$ and hence of $V_a$ and its
determinant. With $Q=-Z^2$, one has
\[
               \deg_D\det V_a=t(x_a+rQ)=O_t(nr).
\]
The full, possibly nonreduced cycle has
$h^0(\mathcal O_D)=O_t(n)$ by
Corollary~\ref{D:punctured-cover-antinef-cycle} and the
separate fundamental-cycle treatment of zero-canonical
components. Choose a global section of the globally
generated determinant avoiding the finitely many associated
components of the Cohen--Macaulay curve $D$. It is regular,
and its zero scheme is finite. The sequence
\[
     0\longrightarrow\mathcal O_D\longrightarrow\det V_a
                  \longrightarrow T_a\longrightarrow0
\]
has $\ell_k(T_a)=\deg_D\det V_a$. Therefore
\[
       h^0(\det V_a)\le h^0(\mathcal O_D)+\deg_D\det V_a
                         =O_t(nr).
\]
This retains all nilpotent layers; it is not merely a degree
calculation on the reduced fibre.

\emph{The actual object chart.}
Let $V$ be globally generated of rank $r$ on $D$ and fix an
isomorphism $\det V\simeq L$. A general $r+1$-tuple of global
sections generates $V$: the rank-drop condition has codimension
two at a fixed point, while the reduced support is one-dimensional.
The incidence argument avoids every point of that support,
and Nakayama gives the same conclusion on $D$ with its
nilpotents. The kernel of the surjection
$\mathcal O_D^{r+1}\twoheadrightarrow V$ is a line bundle,
necessarily isomorphic to $L^{-1}$. Conversely a tuple of
$r+1$ sections of $L$ generating $L$ defines a subbundle
\[
        L^{-1}\longrightarrow\mathcal O_D^{r+1}
\]
whose quotient is locally free of rank $r$ and has determinant
$L$. The generating condition is open: its complement is the
proper image of the common zero scheme in the product with $D$.
The universal quotient is locally a split quotient of vector
bundles, so it remains a vector bundle after every base change.
The dimension bound follows from (D.UG.9).
\end{proof}

\begin{remark}[The exact remaining realization gate]
The charts in (D.UG.10) are for a fixed determinant line, not
for a moving Picard scheme. They need no scheme of rank-$r$
quotients with a quadratic rank enlargement. Nevertheless,
their open parameter spaces contain bundles which need not
come from a reflexive Frobenius-periodic module. A proof of
the hereditary representation bound still needs an effective
finite-thickening realization and isomorphism criterion,
including all Frobenius arrows and Artinian families, with
the same degree-linear complexity. Reduced-branch framing
does not by itself settle geometric normal thickenings or
the full formal overlap. No representation count or new
finite-generation conclusion is asserted here.
\end{remark}


% Root/B/C/D full actual-file audits passed; auxiliary unbounded-period input.
\begin{proposition}[Proved: Brauer residues and period-degree covers for one exceptional curve]
\label{D:single-curve-Brauer-residue-cover}
Let $(R,\mathfrak m,k)$ be a complete excellent normal local
$k$-algebra of dimension two, with algebraically closed
coefficient and residue field $k$ of characteristic $p>0$.
Put $U=\operatorname{Spec}R-\{\mathfrak m\}$.  Suppose that
there is a projective regular resolution
$\pi:Y\to\operatorname{Spec}R$, an isomorphism over $U$,
whose reduced exceptional fibre is one smooth connected
curve $C$.  Write $g=g(C)$ and $C^2=-d<0$.
For every positive integer $n$ prime to $p$, residue gives
a canonical isomorphism
\[
       \operatorname{Br}'(U)[n]
          \simeq H^1(C,\mathbf Z/n\mathbf Z),
       \qquad |\operatorname{Br}'(U)[n]|=n^{2g},       \tag{D.RC.1}
\]
where $\operatorname{Br}'$ denotes cohomological torsion
Brauer classes.

If $\alpha$ has order $m$ prime to $p$, there is a connected
cyclic finite etale cover $Y'\to Y$ of degree exactly $m$
such that its restriction $V\to U$ kills $\alpha$.
The cover $V$ is the puncture of a complete excellent normal
local ring $R'$ finite over $R$.  The induced resolution
$Y'\to\operatorname{Spec}R'$ has a single smooth connected
exceptional curve $C'$, with
\[
          g(C')-1=m(g-1),\qquad (C')^2=-md.          \tag{D.RC.2}
\]
Thus the splitting cover is etale on the entire resolution,
not only on the puncture.  No assertion about a representation
count uniform in $m$ is part of this proposition.
\end{proposition}

\begin{proof}
The Picard lifting argument in
Lemma~\ref{D:local-Picard-prime-to-p-generators} gives a
surjection $\operatorname{Pic}(Y)\to\operatorname{Pic}(C)$:
the obstructions through the exceptional thickenings lie
in coherent $H^2$ on a curve, and Grothendieck existence
algebraizes a compatible line bundle and its inverse.
Since $R$ is strictly henselian, proper base change gives
\[
           H^a(Y,\mu_n)\simeq H^a(C,\mu_n).
                                                        \tag{D.RC.3}
\]
This is precisely the scope checked in that lemma, including
a possibly nonreduced scheme-theoretic closed fibre.
For a smooth proper curve over $k$, its degree map identifies
$\operatorname{Pic}(C)/n$ with $H^2(C,\mu_n)=\mathbf Z/n$,
and $H^3(C,\mu_n)=0$
\citep[Lemma~59.69.1 and complete proof, Tag 03RQ]
{StacksSingleCurveBrauerD}.
The Kummer square for restriction to $C$ therefore shows
\[
       H^2(Y,\mu_n)=\operatorname{Pic}(Y)/n,
       \qquad \operatorname{Br}'(Y)[n]=0,
       \qquad H^3(Y,\mu_n)=0.                         \tag{D.RC.4}
\]

Also $\operatorname{Pic}(Y)\to\operatorname{Pic}(U)$ is
surjective.  Indeed every line bundle on the regular integral
scheme $U$ is represented by a Cartier divisor; its Weil
divisor closure on the regular scheme $Y$ is Cartier.
Consequently the image of $H^2(Y,\mu_n)$ in $H^2(U,\mu_n)$
is exactly the Kummer image of $\operatorname{Pic}(U)/n$.
The localization sequence and (D.RC.4) give
\[
 \operatorname{Br}'(U)[n]
  =H^2(U,\mu_n)/(\operatorname{Pic}(U)/n)
  \simeq H^3_C(Y,\mu_n)
  \simeq H^1(C,\mathbf Z/n\mathbf Z).                 \tag{D.RC.5}
\]
The last isomorphism is absolute purity for the regular
codimension-one immersion $C\hookrightarrow Y$; its Tate
twist cancels the twist in $\mu_n$.
We use the exact regular-scheme scope of
\citep[Expose~XVI, Theorem~3.1.1]{A_ILO2014}, as in
Theorem~\ref{D:local-surface-constant-H2}; no finite-type
over $k$ hypothesis is imposed on $Y$.
The last group has order $n^{2g}$ by the same curve
cohomology calculation.  The construction is functorial
under finite etale base change of the pair $(Y,C)$.

Now take $n=m$ and let $\chi\in H^1(C,\mathbf Z/m)$
be the residue of $\alpha$.  The isomorphism (D.RC.5)
preserves orders, so the corresponding homomorphism
$\pi_1(C)\to\mathbf Z/m$ is surjective.  Its torsor
$C'\to C$ is therefore connected and has degree $m$.
The proper henselian equivalence for finite etale covers
\citep[Lemma~58.9.1 and complete proof, Tag 0A48]
{StacksSingleCurveBrauerD}
lifts this torsor, including its cyclic action, to a finite
etale $Y'\to Y$.  This cover is connected: any nonempty
open and closed component is proper over $R$ and hence
meets the closed fibre, whereas $C'$ is connected.
It is regular, and thus integral.

Let $L=k(Y')$, and normalize $R$ in $L$ to obtain $R'$.
This is finite by excellence, complete, and local because
it is a finite domain over a henselian local ring.  Its
residue field is $k$.  Equivalently it is the Stein ring
$H^0(Y',\mathcal O_{Y'})$: that ring is finite over $R$,
normal, and has fraction field $L$, so it is exactly the
integral closure.  The map $Y'\to\operatorname{Spec}R'$
is proper and birational, and is an isomorphism away from
the closed point, where it agrees with the finite normal
cover $V\to U$.  Its reduced closed fibre is $C'$.

Pullback of (D.RC.5) along $Y'\to Y$ is the corresponding
residue isomorphism for this new pair.  The pulled-back
character $\chi|_{C'}$ is zero by the defining torsor, so
$\alpha|_V=0$.  Finally $C'\to C$ is etale of degree $m$.
Riemann--Hurwitz gives the first identity in (D.RC.2)
\citep[Lemma~53.12.2, Tag 0C1D]{StacksSingleCurveBrauerD}.
The normal line bundle of $C'$ in $Y'$ is the pullback of
the normal line bundle of $C$ in $Y$.  Its degree is $mC^2$,
which gives the second identity.
\end{proof}

\begin{remark}[Quantitative issue still to be proved]
The splitting degree equals the geometric period, but it
can grow.  Applying a rank-$r$ count on the varying cover
without controlling its constants does not give the desired
uniform rank-quadratic bound.  A residue character of order
$m$ suggests rank-$r/m$ inertia blocks over $C'$, but normal
coordinate powers shift their weights.  A description on
the reduced curve alone supplies no formal-neighbourhood
decomposition or Artinian parameter presentation.
\end{remark}


% Separate post-CP17 research. Full root/B/C/D actual-file audits PASS.
\subsection{Canonical module jets from two fixed parameter thickenings}

\begin{lemma}[Proved: uniform saturation loss for canonical pullbacks]
\label{D:canonical-pullback-uniform-jet-loss}
Let $R$ be a fixed complete excellent normal local surface domain
over an algebraically closed coefficient field $k$ of characteristic
$p>0$. Fix the ideal $J$ of
Propositions~\ref{D:cover-uniform-canonical-defect} and
\ref{D:cover-uniform-Noether-different}. Choose once a generically
separable finite Noether normalization
\[
 S=k[[x,y]]\subset R,\qquad x,y\in J,
 \qquad \mathfrak q=(x,y)\subset S.
                                                        \tag{D.CJ.1}
\]
There is an integer $c\ge0$, depending only on these fixed choices,
with the following property. Let $A/R$ be a finite normal cover domain
of degree $n$, etale on the puncture, let
$\pi:Y\to\operatorname{Spec}A$ be a projective regular resolution
which is an isomorphism on the puncture, and let $M$ be a reflexive
$A$-module. Set
\[
 F=(\pi^*M)^{\vee\vee},\qquad I_b=(x^b,y^b)A,
 \qquad K_b=H^0(Y,I_b\mathcal O_Y\cdot F)\subset M.
\]
Then, for $b\ge2c$,
\[
                 K_b\subset\mathfrak q^{\,b-2c}M.
                                                        \tag{D.CJ.2}
\]
In particular, if $\ell\ge1$ and $b\ge2\ell-1+2c$, then
$K_b\subset I_\ell M$. The constants do not depend on $A$, $M$,
its rank, or the resolution.
\end{lemma}

\begin{proof}
The normalization in (D.CJ.1) can be chosen inside $J$.
Start with a separating finite normalization $k[[u,v]]\subset R$.
For a sufficiently large $p$-power $q$, put
$x=u^q(1+u)$ and $y=v^q(1+v)$. These lie in $J$, generate an
$\mathfrak m_R$-primary ideal, and give a finite normalization.
The monic equations $u^{q+1}+u^q-x=0$ and
$v^{q+1}+v^q-y=0$ have nonzero derivatives $u^q,v^q$;
thus $\operatorname{Frac}k[[u,v]]/\operatorname{Frac}k[[x,y]]$
is separable, as is its composite with the original separating
normalization. One may add this normalization to the fixed finite
list used in the cited Noether-different proposition.

Choose a nonzero $z\in\mathfrak N(R/S)$. This ideal is nonzero
because the generic extension is separable. The first-power
Noether-different result gives
\[
 z\in\mathfrak N(A/S),\qquad
 z\operatorname{Hom}_S(A,S)\subset A\operatorname{Tr}_{A/S}.
                                                        \tag{D.CJ.3}
\]
Here the second inclusion uses the generic trace identification of
the dual module; it also follows directly from the evaluation
description of the Noether different
\citep[Lemma~49.6.7 and complete proof, Tag 0BVT]
{StacksCanonicalJetsB}. Choose a nonzero $c_0\in S\cap zR$;
for example, the constant coefficient of the characteristic
polynomial of multiplication by $z$ on the finite free $S$-module
$R$ has this property by Cayley--Hamilton. Write $c_0=zw$ with
$w\in R$, and put $c=\operatorname{ord}_{\mathfrak q}(c_0)$.

We first prove the following precise scalar estimate:
\[
 a\in H^0(Y,I_b\mathcal O_Y)
                  \quad\Longrightarrow\quad za\in\mathfrak q^bA.
                                                        \tag{D.CJ.4}
\]
Let $v=\operatorname{ord}_{\mathfrak q}$ be the discrete valuation
of $\operatorname{Frac}S$ defined by total order; multiplicativity
holds because $\operatorname{gr}_{\mathfrak q}S=k[X,Y]$ is a domain.
Every extension of $v$ to $\operatorname{Frac}A$ is nonnegative on
$A$, and its valuation ring has a centre on $Y$ by properness
\citep[Lemma~29.43.1, valuative criterion for properness]
{StacksCanonicalJetsB}. At that centre, $a$ is a local linear
combination of $x^b,y^b$. Normalize each extended valuation to
restrict to $v$. Its value on $a$ is therefore at least $b$.

Choose an $S$-basis $e_1,\ldots,e_d$ of $A$, with generic
trace-dual basis $e_1^*,\ldots,e_d^*$. By (D.CJ.3),
$ze_i^*\in A$. Every conjugate of $aze_i^*$ has valuation at
least $b$, so
\[
         \operatorname{Tr}_{A/S}(aze_i^*)\in\mathfrak q^bS.
\]
Indeed its trace belongs to $S$ by integrality, and has total
order at least $b$. Expanding $za$ in the basis $e_i$ proves
(D.CJ.4). The same argument proves
$z\overline{\mathfrak q^bA}\subset\mathfrak q^bA$, but that
integral-closure reformulation is not needed below.

Multiplication by $z$ on $M$ factors through a finite free
$A$-module. Pulling back and taking reflexive hulls gives actual
maps of vector bundles
\[
                 F\xrightarrow{\alpha}\mathcal O_Y^a
                     \xrightarrow{\beta}F,
                 \qquad\beta\alpha=z\operatorname{id}_F.
                                                        \tag{D.CJ.5}
\]
These extensions can be checked at codimension-one discrete
valuation rings and then extended by reflexivity on regular $Y$.
For $m\in K_b$, all coordinates of $\alpha(m)$ belong to
$H^0(Y,I_b\mathcal O_Y)$. Equation (D.CJ.4) consequently gives
\[
 z^2m=\beta(z\alpha(m))\in\mathfrak q^bM,
          \qquad c_0^2m\in\mathfrak q^bM.
\]
The module $M$ is maximal Cohen--Macaulay, hence free over the
regular local ring $S$. Total order in $S$ is additive on
products, so
\[
       (\mathfrak q^bM:c_0^2)=\mathfrak q^{b-2c}M
                   \quad(b\ge2c).
\]
This proves (D.CJ.2). Finally
$\mathfrak q^{2\ell-1}\subset(x^\ell,y^\ell)$ proves the
last assertion. Notice that the valuation argument used the
actual ideal $I_b\mathcal O_Y$; no principalization of that
ideal on $Y$ was assumed.
\end{proof}

\begin{proposition}[Proved: exact two-level recovery of canonical module jets]
\label{D:canonical-module-two-level-recovery}
In the setting of Lemma~\ref{D:canonical-pullback-uniform-jet-loss},
put $r=\operatorname{rank}_A M$ and define the proper parameter-ideal
thickenings and their spaces of sections by
\[
 T_b=V(I_b\mathcal O_Y),\qquad
 V_b(M)=H^0(T_b,F|_{T_b}),\qquad
 W_b(M)=\operatorname{im}\bigl(V_{b+1}(M)\longrightarrow V_b(M)\bigr).
                                                        \tag{D.CJ.6}
\]
The thickenings $T_b$ need not be Cartier or Cohen--Macaulay.
There are natural $A$-module identifications
\[
             W_b(M)=M/K_b,
 \qquad
             W_b(M)/I_\ell W_b(M)\simeq M/I_\ell M
                  \quad(b\ge2\ell-1+2c).
                                                        \tag{D.CJ.7}
\]
They are functorial for homomorphisms of the original reflexive
modules. Thus canonical module jets are recovered from two fixed
parameter-ideal thickenings, with a level independent of the cover
degree and module rank.

Write $d_0=\operatorname{rank}_S R$ and choose $c_R$ such that
$p_g(A)\le c_Rn$, as in
Proposition~\ref{D:cover-uniform-canonical-defect}. Then
\[
 \begin{split}
 h^0(T_b,F|_{T_b})&\le (d_0b^2+2c_R)nr,\\
 h^1(T_b,F|_{T_b})&\le c_Rnr.
 \end{split}                                             \tag{D.CJ.8}
\]
These assertions apply to all actual reflexive modules, without
a Frobenius-periodicity or exceptional-energy hypothesis.
They do not assert effectiveness of an arbitrary family of bundles
on $T_{b+1}$, nor preservation of these images under arbitrary
base change in a parameter scheme.
\end{proposition}

\begin{proof}
First $H^0(Y,F)=M$. One inclusion comes from pullback. Conversely
a section restricts to a section of the bundle associated with
$M$ on the puncture, and Hartogs for the reflexive module $M$
identifies that space of sections with $M$.

We record a useful cohomology estimate with its support argument.
The natural generating sections of $M$ generate $F$ at each
exceptional generic point. Since $k$ is infinite, choose $r$
constant linear combinations whose determinant is a unit at all
those finitely many discrete valuation rings. They give an
injection
\[
                   0\longrightarrow\mathcal O_Y^r
                       \longrightarrow F\longrightarrow Q_F
                       \longrightarrow0.
                                                        \tag{D.CJ.9}
\]
The support of $Q_F$ contains no exceptional component. Its
proper map to $\operatorname{Spec}A$ therefore has finite fibres:
the closed fibre is zero-dimensional, and the resolution is an
isomorphism elsewhere. The support is finite over $A$
\citep[Lemma~37.44.1 and complete proof, Tag 02LS]
{StacksCanonicalJetsB}. Consequently $H^1(Y,Q_F)=0$ and
\[
 H^1(Y,\mathcal O_Y)^r\twoheadrightarrow H^1(Y,F),
 \qquad h^1(Y,F)\le r p_g(A),\qquad JH^1(Y,F)=0.
                                                        \tag{D.CJ.10}
\]
If the exceptional fibre is empty, the same assertion follows
directly from affineness; this also covers the regular case.

Let $D$ be the greatest common divisor of the principal divisors
of $x$ and $y$ on regular $Y$. It is exceptional, since $x,y$
form a parameter ideal. Local unique factorization gives exact
sequences
\[
 0\longrightarrow F(bD)\xrightarrow{\,(-y^b,x^b)\,}
     F^{\oplus2}\xrightarrow{\,(x^b,y^b)\,}
     I_b\mathcal O_Y\cdot F\longrightarrow0.
                                                        \tag{D.CJ.11}
\]
The first arrow is interpreted on the rational sections of
$F(bD)$; its displayed products are regular. Every coherent
sheaf on $Y$ has $H^2=0$, since $\pi$ is proper over an affine
base with fibres of dimension at most one. Hence
\[
          H^1(Y,F)^{\oplus2}\twoheadrightarrow
                      H^1(Y,I_b\mathcal O_Y\cdot F).
                                                        \tag{D.CJ.12}
\]
In the comparison of (D.CJ.11) for $b+1$ and $b$, the map on
the middle terms is $\operatorname{diag}(x,y)$, the map on
the right is inclusion, and the map on the left is
multiplication by $xy$. These maps commute with the displayed
arrows. Since $x,y\in J$, (D.CJ.10)--(D.CJ.12) imply
\[
 H^1(Y,I_{b+1}\mathcal O_Y\cdot F)
                  \longrightarrow H^1(Y,I_b\mathcal O_Y\cdot F)
                        \quad\text{is zero}.
                                                        \tag{D.CJ.13}
\]

Apply cohomology to
$0\to I_b\mathcal O_Y\cdot F\to F\to F|_{T_b}\to0$.
The image of $M$ in $V_b(M)$ is the kernel of its connecting
map to the right-hand group of (D.CJ.13). Naturality and
(D.CJ.13) show that every section from $V_{b+1}(M)$ has zero
connecting image at level $b$. Conversely every section from
$M$ restricts through level $b+1$. It follows that
\[
             W_b(M)=\operatorname{im}(M\to V_b(M))=M/K_b.
\]
The saturation estimate $K_b\subset I_\ell M$ proves the
second identification in (D.CJ.7). All maps in this argument
come from restriction and multiplication by fixed functions,
so the identifications are functorial for actual module maps.

Finally $K_b\supset I_bM$, and $M$ is $S$-free of rank
$d_0nr$. Thus
\[
 \dim_k(M/K_b)\le\dim_k(M/I_bM)=d_0nrb^2.
\]
Equations (D.CJ.10)--(D.CJ.12) bound the possible cokernel of
$M\to V_b(M)$ by $2r p_g(A)$. The same restriction sequence
and $H^2(Y,I_b\mathcal O_Y\cdot F)=0$ give a surjection
$H^1(Y,F)\to H^1(T_b,F|_{T_b})$. These are exactly (D.CJ.8).
\end{proof}

\begin{remark}[The remaining all-family realization problem]
The proposition is an exact pointwise recovery, rather than
an inference from annihilation to stabilization. In particular
it specifies the transition whose image is the correct one.
The schemes $T_b$ are parameter-ideal thickenings and can have
embedded zero-dimensional structure; replacing them silently
by reduced curves or Cartier cycles would change the statement.
To obtain a quantitative hereditary representation atlas one
still has to construct, over arbitrary Artinian parameter
algebras, effective object and arrow diagrams whose section
images realize (D.CJ.7), and control their equation degrees.
Neither a base-change assertion for these images nor a bound
on the resulting number of representation classes is proved here.
\end{remark}


% Post-CP18 source: energy application proved; known correspondence and its scope verified by four independent file audits.
\subsection{Energy-controlled reduction cycles and exact pointwise realization}

\begin{proposition}[Proved: reduction cycles with the established linear energy]
\label{D:energy-controlled-reduction-cycle}
Let $A$ be a complete excellent normal local surface $k$-algebra,
where $k$ is an algebraically closed coefficient and residue field,
and let $\pi:Y\to\operatorname{Spec}A$ be its minimal projective
regular resolution. Suppose its exceptional fibre is nonempty.
Write $\kappa_i=K_Y\cdot E_i\ge0$.
Choose an effective exceptional Cartier divisor $Z>0$ of full
support satisfying
\[
 -Z\cdot E_i\ge3\kappa_i\quad(\kappa_i>0),\qquad
 -Z\cdot E_i\ge0\quad(\kappa_i=0).
                                                        \tag{D.KR.1}
\]
Then $Z$, and every positive integral multiple of $Z$, is a
reduction cycle in the following precise sense:
\[
 \mathcal O_Z(-Z)\text{ is globally generated},\qquad
 H^1(Z,\mathcal O_Z(-Z))=0,\qquad
 \omega_Z^{-1}\text{ is globally generated}.
                                                        \tag{D.KR.2}
\]
In the cover family of Corollary~\ref{D:punctured-cover-linear-energy},
one may choose this $Z$ with the same bounds
\[
 -Z^2+K_Y\cdot Z+\sum_i(-Z\cdot E_i)=O(n),\qquad
 h^0(tZ,\mathcal O_{tZ})+h^1(tZ,\mathcal O_{tZ})=O_t(n)
                                                        \tag{D.KR.3}
\]
for each fixed $t\ge1$. Here $n$ is the cover degree and the
constants depend only on the fixed base and $t$. The assertion
does not enlarge that previously established cover family.
\end{proposition}

\begin{proof}
For a nonzero effective exceptional cycle $B=\sum_i m_iE_i$,
put $a=K_Y\cdot B=\sum_i\kappa_i m_i$. Adjunction gives
\[
                    2p_a(B)=2+a+B^2.                 \tag{D.KR.4}
\]
If $a\ge1$, negative definiteness implies
$2p_a(B)\le a+1\le2a$. If $a=0$, every component of its support
is a smooth rational $(-2)$-curve. The integral intersection
form on that support is even and negative definite, whence
$B^2\le-2$ and $p_a(B)\le0$.

For every subcycle $0<B\le tZ$, (D.KR.1) therefore gives
\[
 \begin{split}
 \deg_B\mathcal O_{tZ}(-tZ)&=-tZ\cdot B\ge2p_a(B),\\
 \deg_B\omega_{tZ}^{-1}&=-(K_Y+tZ)\cdot B
                              \ge2p_a(B).
 \end{split}                                           \tag{D.KR.5}
\]
In the second inequality, when $a>0$ the left side is at least
$(3t-1)a\ge2a$; when $a=0$ it is nonnegative.
The point-ideal automatic-adjunction argument in
Theorem~\ref{D:reflexive-pullback-uniform-generation}, applied
on the full possibly nonreduced curve $tZ$, now proves global
generation of both line bundles and $H^1$-vanishing for the
first. Equivalently, use the complete proof of
\citet[Lemma~2.4 and Theorem~1.1, pp.~192--193]{C_CFHR1999}.
This proves (D.KR.2), also for every $tZ$.

Choose the lattice parameter $c\ge5$ in
Lemma~\ref{D:integral-antinef-energy-cycle}; its lower bound is
stronger than (D.KR.1). Equations (D.KR.3) are then exactly the
bounds already proved in
Corollary~\ref{D:punctured-cover-antinef-cycle}. If a whole
connected exceptional graph has $\kappa=0$, use its rational
double point fundamental cycle, whose energy is two by
Proposition~\ref{D:RDP-full-sheaf-degree}. The same parity
argument proves (D.KR.5) there. Disjoint connected components
can be combined. Thus no extra denominator or enlargement of
the energy budget is needed to obtain a reduction cycle.
\end{proof}

\begin{theorem}[Known correspondence: full-module realization and determination on a reduction cycle]
\label{D:reduction-cycle-full-module-correspondence}
Let $A$ be a complete excellent normal local surface $k$-algebra
with algebraically closed coefficient and residue field $k$,
of arbitrary characteristic. Let $\pi:Y\to\operatorname{Spec}A$
be a projective regular resolution, an isomorphism away from
the closed point, and let $Z>0$ have full exceptional support.
Assume that $\mathcal O_Z(-Z)$ and $\omega_Z^{-1}$ are
generically generated by their global sections and that
$H^1(Z,\mathcal O_Z(-Z))=0$.

For a finite reflexive $A$-module $M$, put
\[
       \mathcal F_M=(\pi^*M)^{\vee\vee},\qquad
       R_Z(M)=\mathcal F_M|_Z.
\]
The functor $R_Z$ is injective on isomorphism classes. A vector
bundle $V$ on $Z$ is in its essential image if and only if:
\begin{enumerate}
\item $V$ is generically generated by global sections;
\item there is a vector bundle $V_2$ on $2Z$, with
$V_2|_Z\simeq V$, for which the boundary map
\[
 \delta_0:H^0(Z,V(Z))\longrightarrow H^1(Z,V)
                                                        \tag{D.KR.6}
\]
is injective. This is the boundary of
$0\to V\to V_2(Z)\to V(Z)\to0$.
\end{enumerate}
In addition, for a full bundle $\mathcal F=\mathcal F_M$
and $V=\mathcal F|_Z$, restriction induces
\[
 H^1(Y,\mathcal End\mathcal F)
       \simeq H^1(Z,\mathcal End V),
 \quad\text{equivalently}\quad
 \operatorname{Ext}^1_Y(\mathcal F,\mathcal F)
       \simeq\operatorname{Ext}^1_Z(V,V).
                                                        \tag{D.KR.7}
\]

The object correspondence is the known Kahn--Gavran theorem.
The proof below records its precise complete-ring scope and the
characteristic-independent marking and tangent-space verification;
no new classification theorem is claimed.
These are pointwise object and tangent-space statements.
They do not assert fullness on morphisms, lifting of a
prescribed isomorphism on $Z$, or an equivalence for arbitrary
Artinian parameter families.
\end{theorem}

\begin{proof}
We reconstruct the relevant argument in the complete algebraic
setting. The original theorem is
\citet[Theorem~1.4, Propositions~1.6 and~1.9, Claim on p.~148,
and Corollary~1.10, pp.~143--149]{C_Kahn1989}.
Its stated conventions include characteristic zero.
\citet[Section~2, Theorem~2 and Propositions~3--4,
pp.~69--71]{C_Gavran2012} extends the correspondence to complete
normal surface rings without that characteristic hypothesis,
but refers back to parts of Kahn's proof. The arguments below
include those parts and use no characteristic-zero vanishing.

\emph{Generic generation and the first cohomology.}
Suppose $\mathcal F$ is a vector bundle on $Y$ whose restriction
$V$ to $Z$ is generically generated. Products of generic
generators of $V$ and $\mathcal O_Z(-Z)$ give a map
\[
      \mathcal O_Z(-Z)^a\longrightarrow V(-jZ)
\]
with zero-dimensional cokernel, for every $j\ge1$.
Its kernel has no $H^2$, and its cokernel has no $H^1$.
Consequently $H^1(Z,V(-jZ))=0$. Filtering the restrictions of
$\mathcal F(-Z)$ to $nZ$ by these sheaves gives
$H^1(nZ,\mathcal F(-Z)|_{nZ})=0$ for every $n$.
Formal functions and completeness give
\[
 H^1(Y,\mathcal F(-Z))=0,
 \qquad H^1(Y,\mathcal F)\simeq H^1(Z,V).              \tag{D.KR.8}
\]
In particular all sections of $V$ lift to $Y$, so
$\mathcal F$ is generated at the exceptional generic points.
Away from the exceptional fibre the evaluation map from the
coherent module $\pi_*\mathcal F$ is onto because $\pi$ is an
isomorphism. This proves generic generation on $Y$.
The same reasoning, applied to $\mathcal O_Y(-jZ)$, supplies
global maps $\mathcal O_Y^a\to\mathcal O_Y(-jZ)$ that are
onto outside finitely many points.

\emph{A finite boundary detects fullness.}
For a vector bundle generically generated on $Y$, fullness is
equivalent to the surjectivity
\[
       H^0(Y,\mathcal F)\longrightarrow
       H^0(Y\setminus Z,\mathcal F).                  \tag{D.KR.9}
\]
Indeed the left module is finite over $A$. If (D.KR.9) holds,
Hartogs makes it reflexive; generic generation identifies its
reflexive pullback with $\mathcal F$ at all codimension-one
points and hence everywhere. Conversely a reflexive module
has exactly this Hartogs property.

Write $F_n=\mathcal F|_{nZ}$. Since $Z$ is Cartier and
$\mathcal F$ is locally free, the local-cohomology computation
for a single nonzerodivisor gives
\[
 H^1_Z(Y,\mathcal F)
        =\varinjlim_n H^0(nZ,F_n(nZ)).                 \tag{D.KR.10}
\]
This follows by taking sections of
$\mathcal F(*Z)/\mathcal F
 =\varinjlim_n F_n(nZ)$; the zeroth local-cohomology sheaf
vanishes. Thus (D.KR.9) is equivalent to injectivity of
$H^1_Z(Y,\mathcal F)\to H^1(Y,\mathcal F)$.
The $n=1$ map, followed by (D.KR.8), is exactly $\delta_0$.

If $\delta_0$ is injective, every map
\[
 \delta_j:H^0(Z,V((j+1)Z))\longrightarrow H^1(Z,V(jZ))
                                                        \tag{D.KR.11}
\]
obtained by twisting the $2Z$ sequence is injective as well.
To see this, dualize a generically surjective global map
$\mathcal O_Y^a\to\mathcal O_Y(-jZ)$. The resulting injection
$\mathcal O_Y(jZ)\to\mathcal O_Y^a$ gives a comparison of
(D.KR.11) with $a$ copies of (D.KR.6); the map on $H^0$
is injective. The comparison exists on $2Z$, so its boundary
maps commute. This proves the assertion.
The successive quotients in (D.KR.10) have boundary maps
whose compositions with restriction are (D.KR.11).
Their injectivity shows that the direct system of $H^0$ in
(D.KR.10) is already constant at $n=1$. We obtain the claimed
injection of local cohomology. The reverse implication follows
by restricting that injection to $n=1$.
Together with (D.KR.8), this proves the full-bundle criterion
using only $V$ and its actual $2Z$ extension.

\emph{Existence of a full extension.}
Start with $V_2$ satisfying the two stated conditions. At every
step, the obstruction to lifting a vector bundle on $nZ$ to
$(n+1)Z$ belongs to
$H^2(Z,\mathcal End V(-nZ))=0$. Choose successive lifts.
The resulting formal vector bundle algebraizes on $Y$ by
Grothendieck existence: the $Z$-adic and closed-fibre adic
topologies are cofinal, and the base is complete. The
algebraized coherent sheaf is locally free because its
nonfree locus is closed and proper over the local base and
misses the closed fibre. The preceding finite-boundary
criterion makes this bundle full. These formal-existence
inputs have the same precise proper complete-base scope as
\citep[Tags 0885, 08MF, 08MD and 02OC]{C_StacksBoundedReflexive}.

\emph{The boundary surjectivity needed for uniqueness.}
Fix a full extension $\mathcal F$ of $V$. For $j\ge1$, the
endomorphism sequence associated to the passage from $jZ$ to
$(j+1)Z$ has a boundary whose restriction is
\[
 \partial'_j:
 H^0(Z,\mathcal End V(-(j-1)Z))
       \longrightarrow H^1(Z,\mathcal End V(-jZ)).     \tag{D.KR.12}
\]
This map is surjective. Here are the details of the relevant
duality argument. The restrictions of
$\mathcal F$, $\mathcal O_Y(-(j-1)Z)$ and
$\omega_Y^{-1}(-Z)$ are generically generated on $Z$.
The last assertion is precisely the hypothesis on
$\omega_Z^{-1}$, by adjunction. The first paragraph lifts
their generic generators to $Y$. Their products therefore
give a generically surjective map
\[
 \mathcal O_Y^a\longrightarrow
                  \mathcal F(-jZ)\otimes\omega_Y^{-1}.
                                                        \tag{D.KR.13}
\]
Serre duality identifies the dual of the injective map
$\delta_0$ with the surjective boundary
\[
 H^0(Z,V^\vee\otimes\omega_Z)
      \longrightarrow H^1(Z,V^\vee(-Z)\otimes\omega_Z).
                                                        \tag{D.KR.14}
\]
Multiplication by (D.KR.13) compares $a$ copies of
(D.KR.14) with (D.KR.12). The map between the target $H^1$
groups is onto: its sheaf map has zero-dimensional cokernel,
and its kernel has no $H^2$. Both comparisons come from maps
on $2Z$, so the boundary square commutes. Surjectivity of
(D.KR.14) proves surjectivity of (D.KR.12). This is the
entire Claim on p.~148 of Kahn in the present algebraic scope.

\emph{Uniqueness of the full bundle.}
For a fixed bundle $F_j$ on $jZ$, isomorphism classes of its
locally free lifts to $(j+1)Z$, with a marking on $jZ$, form
a torsor under $H^1(Z,\mathcal End V(-jZ))$.
Changing the marking gives the actual $\operatorname{Aut}(F_j)$
action on this torsor.
For $j\ge2$, an element of
$H^0(Z,\mathcal End V(-(j-1)Z))$ defines an endomorphism
$h$ of $F_j$ with square zero, so $1+h$ is invertible.
It is the identity on $(j-1)Z$. In local transition matrices,
changing the marking by $1+h$ changes the lifting cocycle by
the boundary $\partial'_j(h)$. There are no higher terms,
because the relevant products of ideal powers vanish.
By (D.KR.12), every lifting class is reached.

At $j=1$, use the actual change-of-marking action, rather
than a pullback-only formula that would forget the induced
identification of the kernel. The group
$\operatorname{Aut}(V)$ is the group of units in the
finite-dimensional algebra $H^0(Z,\mathcal End V)$; it is
a smooth open subscheme of that vector space. The differential
of its orbit map at the identity is the boundary of the
endomorphism sequence, namely $\partial'_1$. This follows
by lifting $1+\epsilon h$ on an affine cover and computing
the discrepancy of the lifted transition matrices.
Thus the orbit map is smooth at the identity and its orbit
is open in the affine torsor of lifts. The full locus is
open there, because it is the injectivity locus of (D.KR.6)
between two fixed finite-dimensional vector spaces.
Every full lift has an open orbit by the same argument.
Two nonempty open subsets of this irreducible affine torsor
intersect, so there is exactly one full orbit.

The corrections for $j\ge2$ are the identity one level
earlier. Therefore the successive isomorphisms stabilize on
each fixed thickening and yield a compatible formal
isomorphism. Formal existence for coherent Hom algebraizes
it and its inverse. This proves uniqueness of the full
bundle, hence injectivity of $R_Z$ on module isomorphism
classes by $M=H^0(Y,\mathcal F_M)$.

\emph{The tangent-space comparison.}
The restriction sequence for $\mathcal End F_{j+1}$ has
kernel $\mathcal End V(-jZ)$. Its $H^0$ boundary contains
the surjective map (D.KR.12), while its kernel has no $H^2$.
It follows that
\[
 H^1((j+1)Z,\mathcal End F_{j+1})
             \simeq H^1(jZ,\mathcal End F_j)
 \quad(j\ge1).
\]
Formal functions gives (D.KR.7). All steps use curve duality,
generic generation, formal existence, or elementary
deformation theory and work in the stated characteristic.
\end{proof}

\begin{remark}[Exact scope of the applied finite-cycle information]
Proposition~\ref{D:energy-controlled-reduction-cycle} puts the
pointwise correspondence on the already constructed cycles
with degree-linear energy. In particular the conditions on a
bundle and its $2Z$ extension are exact pointwise realization
conditions, not merely necessary numerical conditions.
The fixed-determinant charts of
Proposition~\ref{D:hereditary-fixed-determinant-cycle-chart}
can be used on these cycles and on $2Z$ for the periodic
objects covered by that proposition.
Neither (D.KR.7) nor uniqueness on isomorphism classes says
that an arbitrarily prescribed map on $Z$ lifts: the exact
sequence still permits
\[
 H^1(Y,\mathcal End\mathcal F(-Z))
   \simeq\operatorname{coker}\bigl(
 H^0(Y,\mathcal End\mathcal F)\to
 H^0(Z,\mathcal End V)\bigr).
\]
Arbitrary-parameter realization, effective arrow equations,
and the desired degree-linear hereditary representation
count remain separate questions. No finite-generation or
nonabelian-crown conclusion is inferred here.
\end{remark}


% Post-checkpoint18 proof. Root/B/C/D complete actual-file and primary-source audits PASS.
\subsection{Uniform Cartier-cycle recovery in the one-curve case}

\begin{proposition}[Proved: fixed separating parameters with Cartier pullback]
\label{D:single-curve-Cartier-parameter-ideal}
Let $(R,\mathfrak m,k)$ and $\pi:Y\to\operatorname{Spec}R$
satisfy Proposition~\ref{D:single-curve-Brauer-residue-cover}.
Thus $R$ is a complete excellent normal local surface over its
algebraically closed coefficient field $k$ of characteristic $p>0$,
and the reduced exceptional fibre is a single smooth connected
curve $C$ with $C^2=-d<0$.
Let $J\subset R$ be any fixed $\mathfrak m$-primary ideal.
There are an integer $a>0$ and elements $f,g\in J$ such that
\[
 S=k[[f,g]]\subset R\quad\hbox{is finite and generically separable},
 \qquad (f,g)\mathcal O_Y=\mathcal O_Y(-aC).
                                                        \tag{D.CP.1}
\]
For every finite etale map $h:Y'\to Y$, put $C'=h^*C$.
The same parameters and the same integer $a$ satisfy
\[
       (f^b,g^b)\mathcal O_{Y'}=\mathcal O_{Y'}(-abC')
                  \quad\hbox{for every integer }b\ge1.
                                                        \tag{D.CP.2}
\]
In particular this applies to all the period-degree splitting
covers of Proposition~\ref{D:single-curve-Brauer-residue-cover}.
No constant in (D.CP.1)--(D.CP.2) depends on their degrees.
\end{proposition}

\begin{proof}
The restriction of $\mathcal O_Y(-C)$ to $C$ has degree $d>0$.
It is ample on $C$ by the curve degree criterion
\citep[Lemma~33.44.14, Tag 0B5X]{StacksCartierParametersRoot}.
Ampleness is unchanged by the nilpotent thickening from the
reduced curve to the scheme-theoretic closed fibre
\citep[Lemma~32.11.4, Tag 09MW]{StacksCartierParametersRoot}.
The fibre criterion for relative ampleness
\citep[Lemma~30.21.4, Tag 0D2N]{StacksCartierParametersRoot}
makes $\mathcal O_Y(-C)$ relatively ample over a neighbourhood
of the closed point. Every such neighbourhood in
$\operatorname{Spec}R$ is the whole base. Since the base is
affine, this line bundle is ample on $Y$.

For each $t\ge1$ set
\[
                 I_t=H^0(Y,\mathcal O_Y(-tC))\subset R.
                                                        \tag{D.CP.3}
\]
We used $H^0(Y,\mathcal O_Y)=R$, by proper birational normal
descent. These are $\mathfrak m$-primary ideals: their quotient
in $R$ embeds in $H^0(tC,\mathcal O_{tC})$, a finite-length
$R$-module, and $1\notin I_t$.
They decrease, with intersection zero, because a nonzero
element of $R$ has finite valuation at the generic point of $C$.

We spell out the complete-local topological step. For each $N$,
the images of $I_t$ in the finite-dimensional space $R/\mathfrak m^N$
stabilize; call the stable image $V_N$. The maps
$V_{N+1}\to V_N$ are surjective: choose $t$ beyond both
stabilization indices and use the actual quotient map for $I_t$.
If $V_N$ were nonzero for some $N$, choose a nonzero element
there and successive compatible lifts. Completeness gives an
element $x\in R$ with these residues and $x\ne0$.
For every fixed $t$, its residues lie in the images of $I_t$,
so
\[
        x\in\bigcap_N(I_t+\mathfrak m^N)=I_t.
\]
The equality is Krull intersection in the finite module $R/I_t$.
This contradicts $\bigcap_t I_t=0$. Hence for each $N$ some
$I_t$ is contained in $\mathfrak m^N$. Choose $N$ with
$\mathfrak m^N\subset J$. By increasing $a$ we obtain
\[
        I_a\subset J,\qquad
        I_a\mathcal O_Y=\mathcal O_Y(-aC).
                                                        \tag{D.CP.4}
\]
The second equality follows from eventual global generation of
powers of the ample line bundle
\citep[Proposition~28.27.13, condition~(7) and its proof,
Tag 01PR]{StacksCartierParametersRoot}.

We now choose two generators of this pulled-back ideal while
retaining generic separability. Let $L=\mathcal O_Y(-aC)$.
Choose a finite list of elements of $I_a$ whose sections generate
$L$. Fix separating parameters $u,v$ of a finite normalization
$k[[u,v]]\subset R$. Choose $0\ne t\in I_a$ and enlarge the
finite-dimensional $k$-space of sections to include
$t,ut,vt$. Denote the resulting space by $V\subset I_a$.

The pairs in $V\times V$ whose restrictions generate $L|_C$
form a nonempty open subset. Indeed, for any point of $C$, both
sections vanishing imposes two independent linear conditions
on $V\times V$; the incidence locus over the proper curve
therefore has dimension at most $2\dim V-1$, and its image
is closed. A pair outside this image generates at every point
of $C$. The support of the cokernel of
$\mathcal O_Y^2\to L$ is closed and proper over $R$.
If nonempty, its image contains the closed point, so it meets
$C$, a contradiction. Thus such a pair generates $L$ on $Y$.

There is also a nonempty open subset of pairs $(f,g)\in V^2$
for which $df,dg$ are independent in $\Omega_{K/k}$,
where $K=\operatorname{Frac}R$. The differentials $du,dv$
form a $K$-basis, and
\[
       d(ut)-u\,dt=t\,du,\qquad
       d(vt)-v\,dt=t\,dv
\]
show that the differentials of elements of $V$ span this space.
Their two-by-two determinant is a nonzero polynomial in the
two lists of scalar coefficients. Its coefficients lie in a
finite-dimensional $k$-subspace of $K$; expressing them in a
$k$-basis exhibits the nonvanishing condition as a nonempty
Zariski open condition over $k$. Since $k$ is infinite and
algebraically closed, it meets the preceding open subset.
Fix a pair in the intersection.

The zero set of $(f,g)\mathcal O_Y=L$ is exactly $C$.
Since $\pi$ is surjective, $\sqrt{(f,g)R}=\mathfrak m$.
The complete-local finite-parameter argument therefore makes
$R$ finite over $k[[f,g]]$: lift a finite $k$-basis of
$R/(f,g)$ and apply completeness and successive approximation.
The source power-series ring injects, since a nonzero kernel
would give a finite extension of a ring of dimension at most
one, contrary to $\dim R=2$.
Independence of $df,dg$ implies
$\Omega_{K/\operatorname{Frac}S}=0$, so this finite field
extension is separable. Here the separating normalization and
the perfect coefficient field give the usual two-element
$p$-basis description of $\Omega_{K/k}$.
This proves (D.CP.1).

Finally, locally trivialize $L$ with a generator $s$ and write
$f=s\alpha$, $g=s\beta$. The equality $(f,g)\mathcal O_Y=L$
means that $(\alpha,\beta)$ is the unit ideal at every point.
Hence $(\alpha^b,\beta^b)$ is also the unit ideal, and
$(f^b,g^b)\mathcal O_Y=L^{\otimes b}$.
Pulling back along the flat map $h$ gives (D.CP.2), with its
displayed ideal equalities and unchanged constants.
\end{proof}

\begin{corollary}[Proved: fixed Cartier-cycle recovery on etale resolution covers]
\label{D:single-curve-Cartier-module-jet-recovery}
Use the preceding proposition with the fixed canonical-defect
and Noether-different ideal $J$ of
Proposition~\ref{D:cover-uniform-canonical-defect}.
Let $h:Y'\to Y$ be a connected finite etale cover of degree $n$,
and let $A=H^0(Y',\mathcal O_{Y'})$ be its complete normal
local Stein ring. Put $C'=h^*C$.
For every reflexive $A$-module $M$ of rank $r$, set
$F=(\pi'^*M)^{\vee\vee}$ on the resolution
$\pi':Y'\to\operatorname{Spec}A$.
There is a fixed integer $c\ge0$ such that, for
$b\ge2\ell-1+2c$, the natural section image
\[
 W_b=\operatorname{im}\bigl(
 H^0((b+1)aC',F|_{(b+1)aC'})
 \longrightarrow H^0(baC',F|_{baC'})\bigr)
\]
satisfies
\[
           W_b/(f^\ell,g^\ell)W_b
                  \simeq M/(f^\ell,g^\ell)M.
                                                        \tag{D.CP.5}
\]
The identification is functorial for actual module maps.
At each fixed level $b$, the section and first-cohomology
dimensions on $baC'$ are $O(nr)$, with constants depending
only on $R,Y,J,a,b$. In particular the multiple of $C'$
needed to recover a fixed module jet is independent of $n,r$
and of the Brauer period of a splitting cover.
\end{corollary}

\begin{proof}
The proper henselian cover equivalence and Stein argument used
in Proposition~\ref{D:single-curve-Brauer-residue-cover} apply
to every connected finite etale cover of $Y$. They identify
$A$ as the finite normalization of $R$ in $k(Y')$; it is
complete and local, and its puncture is etale over that of
$R$. The smooth curve $C'$ is connected, and $Y'$ is a
projective regular resolution of $A$.
The fixed separating normalization in (D.CP.1) is therefore
available for
Lemma~\ref{D:canonical-pullback-uniform-jet-loss} and
Proposition~\ref{D:canonical-module-two-level-recovery}.
Equation (D.CP.2) identifies their actual parameter-ideal
thickenings $T_b$ with the Cartier cycles $baC'$.
Their fixed trace-clearing integer is the $c$ used here.
Equations (D.CJ.7)--(D.CJ.8) give exactly (D.CP.5) and the
stated cohomology bounds. No new principalization, changing
parameters, or base change of a section image is performed.
\end{proof}

\begin{remark}[Scope of the remaining count]
This identifies the actual finite thickenings in the one-curve
case, including every period-degree cover constructed above.
It does not provide a global twisted Frobenius atlas, a
realization theorem for arbitrary parameter families, a
specified-map lifting theorem, or a representation count
uniform in the geometric Brauer period. Those steps remain
separate obligations.
\end{remark}


% Separate checked post-CP18 source. Root/B/C/D full actual-file audits PASS; root owns assembly.
\subsection{Residual twisted bundles and their exact Frobenius transport}

\begin{proposition}[Proved: the composite-order residual spectral category]
\label{D:residual-twisted-spectral-category}
Let $C$ be a smooth proper connected curve over an algebraically closed
field $k$ of characteristic $p>0$, and let $n$ be prime to $p$.
Let $\mathcal X\to C$ be a banded $\mu_n$-gerbe and let
$\mathcal G\to\mathcal X$ be a second banded $\mu_n$-gerbe.
Call the band of the second gerbe the scalar band. For
$u\in\mathbb Z/n$, let $\mathcal P_u$ be the relative stack of
invertible sheaves on $\mathcal G$ on which the scalar band acts
with character $u$.

There is a finite etale $\mathbb Z/n$-torsor
$h_u:R_u\to C$ such that $\mathcal P_u\to R_u$ is a
$\mathbb G_m$-gerbe. This gerbe is neutral. Choose a corresponding
universal line $\Lambda_u$ on
$\mathcal G_{R_u}=\mathcal G\times_C R_u$, and write
$\pi_u:\mathcal G_{R_u}\to R_u$ and
$j_u:\mathcal G_{R_u}\to\mathcal G$.
Then the functors
\[
 \begin{split}
 \mathcal V&\longmapsto
    (\pi_u)_*\mathcal Hom(\Lambda_u,j_u^*\mathcal V),\\
 W&\longmapsto (j_u)_*(\Lambda_u\otimes\pi_u^*W)
 \end{split}                                                    \tag{D.RF.1}
\]
are inverse equivalences between scalar-weight-$u$ vector bundles on
$\mathcal G$ and vector bundles on $R_u$. Their evaluation maps and
the equivalence commute with arbitrary coefficient-base change,
including nonreduced Artinian bases.

If the class of $R_u$ has order $m$ in
$H^1(C,\mathbb Z/n)$, then $R_u$ has $n/m$ connected components,
each of degree $m$ over $C$. For a bundle $W$ of rank $s_j$ on the
$j$th component, the corresponding bundle in (D.RF.1) has rank
\[
                         m\sum_j s_j.                         \tag{D.RF.2}
\]
This is a statement about the residual gerbe category. It does not
assert a rank drop for a representation of a whole punctured surface
group after restriction to a Brauer-splitting cover.
\end{proposition}

\begin{proof}
Etale-locally on $C$, the composite $\mathcal G\to C$ is the
classifying gerbe of a finite group $D$ fitting into a central
extension
\[
                    1\to\mu_n\to D\to\mu_n\to1.
                                                               \tag{D.RF.3}
\]
Indeed the bands are etale of order invertible on $C$, and one can
choose objects and their identifications etale-locally. A central
extension of a cyclic group by a cyclic group is abelian: a lift of
a quotient generator, together with the central kernel, generates
the whole group. Thus, after a further etale localization, $D$ is
diagonalizable of order $n^2$.

The characters of $D$ restricting to a specified character $u$ of
the scalar kernel form a simply transitive set under the character
group of $D/\mu_n$, which is $\mathbb Z/n$. There are exactly $n$
such characters. Consequently, locally $\mathcal P_u$ is a disjoint
union of $n$ copies of $B\mathbb G_m$. The transition functions act
by translations of this character torsor, and descent gives $R_u$
and its $\mathbb G_m$-gerbe. This is the same character construction
as the prime-order argument in
\citet[Proposition~5.1.6 and its complete proof, pp.~18--19]
{Lieblich2011ResidualD}; the argument just given establishes the
stated composite-order scope without attributing that scope to the
prime-order source.

Each component of $R_u$ is a smooth proper curve over $k$.
The vanishing $H^2(R_u,\mathbb G_m)=0$ therefore neutralizes the
gerbe and supplies $\Lambda_u$
\citep[Theorem~59.68.5 and its complete derivation in
Section~59.68, Tags 03RM and 03RH]{StacksResidualFrobeniusD}.

To check (D.RF.1), use the local presentation $BD$ just described.
The first functor takes the isotypic piece belonging to the selected
character, after tensoring by its inverse. The second functor sums
the resulting pieces over all $n$ characters above $u$. These are
inverse operations. Invariants of a diagonalizable group are the
degree-zero summand of its grading; they are exact and commute with
arbitrary base change. This proves the same assertion for every
coefficient ring, and descent proves the global equivalence and its
base-change compatibility. No decomposition theorem for a wild
finite group is involved.

The image of the monodromy character of a $\mathbb Z/n$-torsor has
order $m$ exactly when its cohomology class has order $m$. Its orbits
on the $n$ labels then have size $m$. This proves the component
statement and the rank formula (D.RF.2).
\end{proof}

\begin{proposition}[Proved: normalized line comparisons and Frobenius transport]
\label{D:residual-twisted-Frobenius-transport}
Use the setting and choices of
Proposition~\ref{D:residual-twisted-spectral-category}.
Fix $f\ge1$ with $q=p^f\equiv1\pmod n$, and put
$u_a=p^a\pmod n$, $0\le a\le f$.
Choose universal lines $\Lambda_a$ on $\mathcal G_{R_{u_a}}$,
with the same choice at $a=f$ as at $a=0$.
Their choices may be normalized so that their rational degrees on
every connected component lie in $[0,1)$.

Tensoring a line by its $p$th power induces an isomorphism
\[
                  \sigma_a:R_{u_a}\longrightarrow R_{u_{a+1}}
                                                               \tag{D.RF.4}
\]
over $C$. There is a line bundle $B_a$ on $R_{u_a}$ with
\[
 \Lambda_a^{\otimes p}
     \simeq \sigma_a^*\Lambda_{a+1}\otimes\pi_{u_a}^*B_a,
       \qquad 0\le\deg(B_a|_Z)\le p-1                    \tag{D.RF.5}
\]
for every connected component $Z$ of $R_{u_a}$; the first pullback
in this display includes the identity on $\mathcal G$.
Under (D.RF.1), absolute Frobenius pullback is exactly the functor
\[
             W\longmapsto
       (\sigma_a)_*\bigl(F_{R_{u_a}}^*W\otimes B_a\bigr).
                                                               \tag{D.RF.6}
\]
This formula, including its structural isomorphisms, commutes with
coefficient-base change when Frobenius also acts on the coefficient
ring.

Although $R_{u_f}=R_{u_0}$, the composite
$\sigma=\sigma_{f-1}\cdots\sigma_0$ need not be the identity.
It is an automorphism of the $\mathbb Z/n$-torsor $R_{u_0}$.
A Frobenius cycle must retain this final automorphism.
\end{proposition}

\begin{proof}
The $n^2$th power of an invertible sheaf on $\mathcal G_{R_u}$
has trivial total inertia and descends to $R_u$ by the local
diagonalizable-group description. Define its rational degree to be
$1/n^2$ times the degree of that descended power. Tensoring by an
ordinary line on a component of $R_u$ adds its integral degree.
Every integral degree occurs because $k$ is algebraically closed.
This permits the normalization into $[0,1)$, component by component;
at the last stage choose the original normalized line again.

On character labels, tensor-$p$ is multiplication by $p$.
It is a bijection because $D$ has order prime to $p$, and it takes
characters above $u_a$ to those above $u_{a+1}$. It thus gives the
isomorphism (D.RF.4); it is equivariant for multiplication by $p$
on the acting group $\mathbb Z/n$.

The two lines in the quotient on the left of (D.RF.5) represent
the same point of the relative Picard space. Their quotient has
trivial total inertia, so it descends to a line $B_a$ on $R_{u_a}$.
On a component, if $\delta_a,\delta_{a+1}\in[0,1)$ are the
normalized rational degrees, then
\[
                   \deg B_a=p\delta_a-\delta_{a+1}.
\]
This is an integer strictly between $-1$ and $p$, hence lies in
$\{0,\ldots,p-1\}$.

Absolute Frobenius pullback of an invertible sheaf is canonically
its $p$th tensor power. Moreover, a finite etale map and absolute
Frobenius form a cartesian square: this is the relative-Frobenius
isomorphism for an etale morphism
\citep[Lemma~41.14.3 and complete proof, Tag 0EBS]
{StacksResidualFrobeniusD}, applied on an atlas and descended.
Therefore Frobenius pullback
commutes with the finite etale pushforward $j_{u_a*}$ in
(D.RF.1). Applying (D.RF.5) to that formula gives (D.RF.6).
One may also check the identity directly in the character
decomposition, where it raises coefficient matrices to $p$th
powers and sends each character to its $p$th power. This check
proves the stated compatibility over every coefficient ring.

After $f$ steps, the action on the translating character group is
multiplication by $q$, which is the identity modulo $n$. Thus the
composite is an automorphism of the torsor, but it can be a
nontrivial translation. For example, take the central extension
$\mu_n\to\mu_{n^2}\to\mu_n$ and $n=q-1$. The weight-one
characters are $1+nj\pmod{n^2}$. Multiplication by $q=1+n$ sends
their label $j$ to $j+1\pmod n$. This proves the warning directly;
the finite arithmetic check described below tests precisely this
identity, not the gerbe geometry.
\end{proof}

\begin{remark}[Scope of this residual comparison]
The propositions give an exact category on the residual curve,
including the Frobenius transport, normalized tensor-line degrees,
and the final character permutation. They do not realize arbitrary
normal-direction thickenings as canonical modules on a contracted
surface. They give neither a uniform finite-thickening atlas nor a
count of projective representations with unbounded geometric
Brauer period. In particular, they do not replace the required
compatibility of normal-layer mixing and punctured Frobenius maps.

The separate program
\nolinkurl{parts/D/post18_research/code/check_residual_Frobenius_labels.py}
checks finite instances of the last character identity and of
the exact graded coefficient/carry formulas used in the research
notes. Its output is
\nolinkurl{parts/D/post18_research/output/residual_Frobenius_labels.json}.
Those finite checks are not used in the general proofs above.
\end{remark}


% Separate post-CP18 research. Full root/B/C/D actual-file and source audits PASS.
\subsection{Artinian realization with a retained full second neighbourhood}

\begin{theorem}[Proved: cohomologically flat full-neighbourhood realization]
\label{D:Artinian-full-neighbourhood-realization}
Let $A$ be a complete excellent normal local surface with
algebraically closed coefficient and residue field $k$. Let
$\pi:Y\to\operatorname{Spec}A$ be a projective regular resolution,
an isomorphism on the puncture, and let $Z$ be a reduction cycle
of full exceptional support in the sense of
Theorem~\ref{D:reduction-cycle-full-module-correspondence}.
The characteristic of $k$ is arbitrary.

Let $T$ be a finite-dimensional local Artinian $k$-algebra with
residue field $k$, put $A_T=A\otimes_kT$ and $Y_T=Y\times_kT$,
and let $V_2$ be a vector bundle on $(2Z)_T$. Write
$V=V_2|_{Z_T}$ and $V_0=V\otimes_Tk$. Assume:
\begin{enumerate}
\item $V_0$ is generically generated by its global sections;
\item for $W=V$ and $W=V(Z)$, both $H^0(Z_T,W)$ and
$H^1(Z_T,W)$ are finite free over $T$, and their formation
commutes with every homomorphism from $T$ to a finite local
Artinian $k$-algebra with residue field $k$;
\item the actual sequence
\[
 0\longrightarrow V\longrightarrow V_2(Z)
       \longrightarrow V(Z)\longrightarrow0
\]
has a split injective boundary
\[
 \delta_0:H^0(Z_T,V(Z))\longrightarrow H^1(Z_T,V).
                                                        \tag{D.AR.1}
\]
\end{enumerate}
There is a vector bundle $\mathcal F_T$ on $Y_T$ restricting
to the given $V_2$. Every such vector-bundle extension has the
following properties. Its module
$M_T=H^0(Y_T,\mathcal F_T)$ is finite over $A_T$, is flat
over $T$, and has maximal Cohen--Macaulay closed fibre.
There are natural identifications
\[
 \begin{split}
 M_T&=H^0(Y_T\setminus Z_T,\mathcal F_T),\\
 (\pi_T^*M_T)^{\vee\vee}&\simeq\mathcal F_T ,
 \end{split}                                           \tag{D.AR.2}
\]
where the double dual is taken over $\mathcal O_{Y_T}$.
Moreover, for every homomorphism $T\to T'$ as in item~(2),
\[
 M_T\otimes_TT'\ \simeq\
 H^0(Y_{T'},\mathcal F_T\otimes_TT').                    \tag{D.AR.3}
\]
In particular these are actual relative maximal
Cohen--Macaulay modules with their full vector-bundle pullbacks,
not only fibrewise realizations.

The hypotheses retain the second neighbourhood and the
cohomology base-change conditions. The theorem does not
assert that arbitrary families of bundles on $Z$ satisfy them,
that a prescribed map on $Z$ lifts, or that a choice of
extension $\mathcal F_T$ is functorial in $V_2$.
\end{theorem}

\begin{proof}
All thickenings and coefficient base changes in this proof
are over the fixed $k$-scheme $Y$. The ring $A_T$ is complete
Noetherian and local: as an $A$-module it is a finite direct
sum of copies of $A$, and its nilpotent coefficient ideal has
quotient $A$. The ideals of $Z_T$ and of the closed fibre of
$Y_T\to\operatorname{Spec}A_T$ have the same radical, so their
adic topologies are cofinal.

\emph{Existence of a vector-bundle extension.}
If $V_j$ is already constructed on $(jZ)_T$, the obstruction
to extending it across $(jZ)_T\subset((j+1)Z)_T$ lies in
\[
 H^2(Z_T,\mathcal End(V)(-jZ))=0.
                                                        \tag{D.AR.4}
\]
Indeed local free lifts exist, and the discrepancy of their
transition matrices is the usual additive degree-two cocycle
in the displayed square-zero ideal. The vanishing holds
because $Z_T$ has dimension one. This is also the
$\operatorname{Ext}^2$ obstruction of
\citep[Lemma~91.10.4 and complete proof, Tag 08MF]
{StacksRelativeKahnB}.
Starting from the prescribed $V_2$, choose compatible lifts
for every $j$. Projective Grothendieck existence algebraizes
the resulting formal coherent sheaf
\citep[Lemma~30.24.3 and complete proof, Tag 0885]
{StacksRelativeKahnB}. At each point of the closed fibre its
completion is free, so the algebraized sheaf is locally free
there. Its nonfree locus is closed and proper over the local
base and misses the closed fibre. It is therefore empty.
This gives $\mathcal F_T$.

\emph{Cohomology away from the first neighbourhood.}
The generic-generation lemma in the proof of
Theorem~\ref{D:reduction-cycle-full-module-correspondence}
gives
\[
                  H^1(Y,\mathcal F_0(-Z))=0,
                                                        \tag{D.AR.5}
\]
where $\mathcal F_0=\mathcal F_T\otimes_Tk$.
For clarity, it only uses generic generation of $V_0$,
generic generation of $\mathcal O_Z(-Z)$, and
$H^1(Z,\mathcal O_Z(-Z))=0$; fullness is not assumed at this
step. An Artinian filtration of any coefficient algebra $T'$
has successive quotients copies of $k$. Since
$\mathcal F_T\otimes_TT'$ is a vector bundle, (D.AR.5) and
the associated exact sequences give
\[
 H^1(Y_{T'},(\mathcal F_T\otimes_TT')(-Z))=0.
                                                        \tag{D.AR.6}
\]
All higher cohomology of coherent sheaves on $Y_{T'}$
vanishes in degrees at least two. One may see this by formal
functions: the closed-fibre thickenings have dimension one,
and the finite higher-cohomology modules over the complete
local base have zero completion. Hence restriction gives
\[
 H^1(Y_{T'},\mathcal F_T\otimes_TT')
             \simeq H^1(Z_{T'},V\otimes_TT').
                                                        \tag{D.AR.7}
\]

\emph{Flatness and base change of the section module.}
A finite affine cover of the separated scheme $Y_T$
computes the cohomology of $\mathcal F_T(-Z)$ by a bounded
Čech complex of flat $T$-modules. Equation (D.AR.6) says that
this complex and its coefficient base changes have no
positive cohomology. Equivalently, its degree-zero
cohomology is flat and commutes with these base changes.
This also follows directly by descending through the bounded
Čech complex: successive kernels of surjections between
flat modules are flat. The restriction sequence and
(D.AR.6) now give
\[
 0\longrightarrow H^0(Y_T,\mathcal F_T(-Z))
 \longrightarrow M_T
 \longrightarrow H^0(Z_T,V)\longrightarrow0.
                                                        \tag{D.AR.8}
\]
The right term is free and commutes with coefficient change
by assumption. Both conclusions therefore hold for $M_T$
as well. This proves flatness and (D.AR.3). Finiteness over
$A_T$ follows from properness.

\emph{The full Hartogs property over $T$.}
The split injection (D.AR.1) remains injective after every
coefficient change under consideration. For each $j\ge1$,
the reduction-cycle generic-generation lemma supplies a
fixed map on $Y$,
\[
 \mathcal O_Y^a\longrightarrow\mathcal O_Y(-jZ),
\]
which is onto at the exceptional generic points. Its dual
restricts to an injection
$\mathcal O_Z(jZ)\to\mathcal O_Z^a$. This injection remains
injective after coefficient extension, since that extension
is flat over $k$. Multiplication gives a commuting comparison
of the boundary
\[
 \delta_j:H^0(Z_T,V((j+1)Z))
                \longrightarrow H^1(Z_T,V(jZ))
                                                        \tag{D.AR.9}
\]
with $a$ copies of (D.AR.1). The comparison is defined on
$(2Z)_T$, and its map on the degree-zero source is injective.
Thus every $\delta_j$ is injective.

Put $F_j=\mathcal F_T|_{(jZ)_T}$. The sequences
\[
 0\longrightarrow F_j(jZ)
 \longrightarrow F_{j+1}((j+1)Z)
 \longrightarrow V((j+1)Z)\longrightarrow0
                                                        \tag{D.AR.10}
\]
show that the transition
$H^0(F_j(jZ))\to H^0(F_{j+1}((j+1)Z))$
is an isomorphism: the next boundary, composed with
restriction, is the injective map (D.AR.9).
Consequently
\[
 H^1_{Z_T}(Y_T,\mathcal F_T)
  =\varinjlim_j H^0((jZ)_T,F_j(jZ))
  =H^0(Z_T,V(Z)).
                                                        \tag{D.AR.11}
\]
The first equality is the Cartier local-cohomology
calculation; $\mathcal F_T$ has no $Z_T$-supported torsion.
Under (D.AR.7), the map from (D.AR.11) to
$H^1(Y_T,\mathcal F_T)$ is exactly $\delta_0$.
It is injective. The localization exact sequence now gives
the first equality in (D.AR.2).

On the closed fibre, the same argument and generic generation
give the ordinary full-bundle criterion. Thus
$M_0=H^0(Y,\mathcal F_0)$ is maximal Cohen--Macaulay and
$\mathcal F_0=(\pi^*M_0)^{\vee\vee}$.
There is no hidden flatness assertion here:
$M_T\otimes_Tk=M_0$ was already proved in (D.AR.3).
For any finite regular Noether normalization
$S=k[[u,v]]\subset A$, the $S$-module $M_0$ is free.
Lift an $S$-basis to $M_T$. Nilpotent Nakayama gives a
surjection $S_T^d\to M_T$; flatness over $T$ makes its kernel
reduce to zero, and nilpotent Nakayama kills the kernel.
Thus $M_T$ is $S_T$-free. In particular $u,v$ form a regular
sequence on it, as required for relative Hartogs.

\emph{Identification of the relative full pullback.}
Assumption~(2) and (D.AR.8) lift a finite set of generic
generators of $V_0$ to global sections of $\mathcal F_T$.
Nilpotent Nakayama implies that
$\pi_T^*M_T\to\mathcal F_T$ is surjective at every exceptional
generic point. Away from the exceptional locus it is an
isomorphism. Let $G$ be its image. The cokernel
$\mathcal F_T/G$ is supported in codimension two, and the
kernel of $\pi_T^*M_T\to G$ is killed locally by a power of
the fixed Cartier equation of $Z$.
The scheme $Y_T$ is Cohen--Macaulay of dimension two:
its local rings are finite free over the corresponding
regular local rings of $Y$. Hence a codimension-two
coherent module has zero $\operatorname{Ext}^0$ and
$\operatorname{Ext}^1$ into $\mathcal O_{Y_T}$.
The exceptional torsion kernel has no homomorphism into
$\mathcal O_{Y_T}$, since the Cartier equation is a
nonzerodivisor. Dualizing the two kernel/image sequences
therefore identifies
$(\pi_T^*M_T)^\vee$ with $\mathcal F_T^\vee$.
Dualizing again proves the second equality in (D.AR.2).
The proof applies to every chosen extension and every stated
coefficient base change. It never invokes a prescribed
$H^0$-isomorphism lift from Kahn's tangent-space comparison.
\end{proof}

\begin{corollary}[Proved: exact two-level recovery for the realized Artinian families]
\label{D:Artinian-canonical-two-level-jets}
Use the fixed $R$, $J$, separating normalization
$S=k[[x,y]]$, scalar $c_0$, and integer $c$ of
Lemma~\ref{D:canonical-pullback-uniform-jet-loss}.
Let $A/R$ be a finite normal cover domain, étale on the
puncture, and let $Y$ be a projective regular resolution.
Let $(M_T,\mathcal F_T)$ be one of the realized families
constructed in
Theorem~\ref{D:Artinian-full-neighbourhood-realization}
on this resolution. In particular its evaluation map is
surjective at every exceptional generic point, its closed
fibre is the canonical full pullback of $M_0$, and
(D.AR.2)--(D.AR.3) hold.

Set
\[
 I_b=(x^b,y^b)A_T,\qquad
 Q_{b,T}=V(I_b\mathcal O_{Y_T}),\qquad
 V_b(T)=H^0(Q_{b,T},\mathcal F_T|_{Q_{b,T}}),
\]
and let $W_b(T)$ be the image of $V_{b+1}(T)\to V_b(T)$.
For $b\ge2\ell-1+2c$ there is a canonical isomorphism
\[
             W_b(T)/I_\ell W_b(T)\ \simeq\
                         M_T/I_\ell M_T.              \tag{D.AR.12}
\]
It is functorial for actual module maps and commutes with
every Artinian coefficient change in
Theorem~\ref{D:Artinian-full-neighbourhood-realization}.
The individual spaces $V_b(T)$ and the raw images $W_b(T)$
are not asserted to commute with base change.
\end{corollary}

\begin{proof}
We give the relative details rather than infer them from the
pointwise theorem. The module $M_T$ is free over $S_T$ by
the lifted-basis argument above, and relative Hartogs gives
$H^0(Y_T,\mathcal F_T)=M_T$.

The actual Noether-different factorizations of
Proposition~\ref{D:cover-uniform-Noether-different}
hold for these $T$-flat sources. In particular the fixed
$z\in\mathfrak N(R/S)$ used to define $c_0$ factors
$z\operatorname{id}_{M_T}$ through a finite free
$A_T$-module. This induces a factorization
\[
       \mathcal F_T\longrightarrow\mathcal O_{Y_T}^a
                          \longrightarrow\mathcal F_T
                                                        \tag{D.AR.13}
\]
with composite $z$. To check the first arrow, pull it back
from $M_T$ where evaluation is onto. The kernel of evaluation
is exceptional torsion, so its map to the free target
vanishes. Thus the first arrow is defined at every
codimension-one point. Hartogs for the locally free source
and target on the Cohen--Macaulay scheme $Y_T$ extends it
across the remaining codimension-two set. The composite is
$z$ since this is true on the puncture.

Put $K_b(T)=H^0(Y_T,I_b\mathcal O_{Y_T}\mathcal F_T)$.
The scalar clearing inclusion in the proof of
Lemma~\ref{D:canonical-pullback-uniform-jet-loss} is an
inclusion of fixed $k$-vector spaces. Flat coefficient
extension over $k$ gives
\[
 zH^0(Y_T,I_b\mathcal O_{Y_T})
                   \subset (x,y)^bA_T.
                                                        \tag{D.AR.14}
\]
Apply (D.AR.13) to a section of $K_b(T)$ and then apply
(D.AR.14) to its coordinates. Exactly as in the pointwise
argument,
$c_0^2K_b(T)\subset(x,y)^bM_T$.
Since $M_T$ is $S_T$-free and the initial form of the fixed
$c_0\in S$ remains a nonzerodivisor in $T[X,Y]$,
\[
 K_b(T)\subset(x,y)^{b-2c}M_T\subset I_\ell M_T
 \quad\text{if }b\ge2\ell-1+2c.                        \tag{D.AR.15}
\]

It remains to verify the adjacent-level cohomology
transition. Choose $r$ closed-fibre sections of $M_0$
that generate $\mathcal F_0$ at each exceptional generic
point and lift them to $M_T$. Their determinant is a
nonzerodivisor on $Y_T$: this follows by the Artinian
filtration from the same assertion on $Y$. Hence there is
an injection $\mathcal O_{Y_T}^r\to\mathcal F_T$.
Its cokernel has support finite over $\operatorname{Spec}A_T$,
because its underlying support misses the exceptional
generic points and is proper with finite fibres. Therefore
\[
 H^1(Y,\mathcal O_Y)^r\otimes_kT
                       \longrightarrow H^1(Y_T,\mathcal F_T)
\]
is surjective. The fixed canonical-defect ideal $J$ kills
the source, so $JH^1(Y_T,\mathcal F_T)=0$.

Let $D=\gcd(\operatorname{div}_Yx,\operatorname{div}_Yy)$.
The two-generator syzygy on $Y$, pulled back over $k$ and
then tensored with the vector bundle $\mathcal F_T$, is
\[
 0\longrightarrow\mathcal F_T(bD)
 \xrightarrow{(-y^b,x^b)}\mathcal F_T^2
 \xrightarrow{(x^b,y^b)}
 I_b\mathcal O_{Y_T}\mathcal F_T\longrightarrow0.
                                                        \tag{D.AR.16}
\]
The comparison from level $b+1$ to $b$ is multiplication
by $xy$ on the first term and $\operatorname{diag}(x,y)$
on the middle term. Since $x,y\in J$ and coherent $H^2$
vanishes on $Y_T$, it follows that
\[
 H^1(Y_T,I_{b+1}\mathcal O_{Y_T}\mathcal F_T)
   \longrightarrow H^1(Y_T,I_b\mathcal O_{Y_T}\mathcal F_T)
 \quad\text{is zero}.                                 \tag{D.AR.17}
\]
The connecting maps for restriction to $Q_{b,T}$ now give,
with no image/base-change interchange,
\[
          W_b(T)=\operatorname{im}(M_T\to V_b(T))
                         =M_T/K_b(T).
                                                        \tag{D.AR.18}
\]
Equations (D.AR.15) and (D.AR.18) prove (D.AR.12).
Functoriality follows from the actual evaluation and
restriction maps. The right side of (D.AR.12) commutes
with coefficient change, proving the stated base-change
assertion for the recovered quotient even when the raw
cohomology spaces and images do not.
\end{proof}

\begin{remark}[What is still required for the quantitative atlas]
The preceding results turn the specified cohomologically
flat $2Z$ data into actual Artinian families and recover
their canonical module jets. They do not provide a
polynomial presentation, with the desired degree-linear
coordinate budget, of the cohomology base-change strata,
their moving determinants, or all compatible Frobenius
arrows. In particular, an arbitrary map on $Z$ is not
declared liftable. No representation-growth or
finite-generation conclusion is inferred here.
\end{remark}


% Post-CP18 source. Four complete actual-file and primary-source audits passed; fixed line bundles only.
\subsection{Finite polynomial Hom complexes for fixed determinant presentations}

\begin{proposition}[Proved: all-coefficient finite Hom and diagram models]
\label{D:fixed-determinant-finite-Hom}
Let $D$ be a projective Cohen--Macaulay scheme of pure dimension
one over an algebraically closed field $k$. Put
$b_D=h^0(D,\mathcal O_D)+h^1(D,\mathcal O_D)$.
Fix globally generated line bundles $L_i$ and positive integers
$r_i$, and put $d_i=\deg_D L_i$. On the open set of generating
tuples in $H^0(D,L_i)^{r_i+1}$ consider the universal sequence
\[
 0\longrightarrow L_i^{-1}
   \xrightarrow{s_i}\mathcal O_D^{r_i+1}
   \longrightarrow V_i\longrightarrow0.
                                                        \tag{D.FH.1}
\]
All constant sheaves in this formula are pulled back to the
coefficient scheme. There are finite graded $k$-vector spaces
$H_{ij}$ and differentials $\Delta_{ij}(s)$ on
$H_{ij}\otimes_k T$, for every coefficient $k$-algebra $T$,
with the following properties.
\begin{enumerate}
\item The construction is given by fixed polynomial formulas
in the coordinates of the tuples $s_i$, commutes with every
homomorphism of coefficient algebras, and computes
\[
 (H_{ij}\otimes_k T,\Delta_{ij})
       \simeq R\operatorname{Hom}_{D_T}(V_i,V_j).
                                                        \tag{D.FH.2}
\]
The degrees of the entries of $\Delta_{ij}$ are at most two.
No flatness or base-change assertion for the cohomology of
this complex is assumed.
\item Its total dimension is at most
\[
 \begin{split}
 B_{ij}={}&b_D(r_i+1)(r_j+1)
 +(b_D+d_j)(r_i+1)\\
 &+(b_D+d_i)(r_j+1)+b_D+d_i+d_j.
 \end{split}                                           \tag{D.FH.3}
\]
It is concentrated in total degrees $-1,0,1,2$.
\item Actual morphisms $V_i\to V_j$ over $T$ are represented
by degree-zero cycles in this finite complex. Equality of
two representatives is expressed by one degree-minus-one
boundary witness. Composition is represented by a polynomial
map of total degree at most five in the object and arrow
coordinates. Inverse witnesses therefore give an exact
polynomial presentation of actual bundle isomorphisms over
every $T$, including nonreduced algebras.
\end{enumerate}

In particular, suppose $b_D\le cn$, $d_i\le cnr_i$, and
$1\le r_i\le r$. For a diagram with $O(f)$ vertices, arrows,
and relations, whose relation lengths are bounded by a fixed
integer, its actual bundles, arrows, and relations admit
these presentations using $O(fnr^2)$ coordinates and scalar
equations of bounded degree. The constants depend on $c$
and the fixed diagram pattern, not on $n,r,f$.
The generating-tuple conditions are open conditions.
The line bundles $L_i$ and the curve $D$ are fixed background
data; this proposition does not parameterize their variation.
\end{proposition}

\begin{proof}
\emph{The line-bundle cohomology budget.}
A general section of a globally generated line bundle on
$D$ avoids its finitely many associated points. It is a
nonzerodivisor, with a zero scheme of length equal to the
degree of the line bundle. The sequences defined by such
sections show that
\[
 \begin{split}
 h^0(L_i)+h^1(L_i)&\le b_D+d_i,\\
 h^0(L_i^{-1})+h^1(L_i^{-1})&\le b_D+d_i,\\
 h^0(L_iL_j^{-1})+h^1(L_iL_j^{-1})
                                      &\le b_D+d_i+d_j.
 \end{split}                                           \tag{D.FH.4}
\]
For the last inequality, tensor the section sequence for
$L_i$ by $L_j^{-1}$, and use the preceding bound for
$L_j^{-1}$. These arguments use the full nonreduced curve;
the Cohen--Macaulay hypothesis rules out embedded associated
points.

\emph{A fixed two-column Cech model.}
Embed $D$ into projective space. Choose two hyperplanes whose
intersection misses $D$ and neither of which contains an
irreducible component. Their complements give an affine
cover $D=U_0\cup U_1$ with affine intersection. Thus the
ordered Cech complex of any coherent sheaf has just the
two terms in Cech degrees zero and one, and computes its
cohomology. The exact primary input is
\citep[Lemmas~30.2.5--30.2.6 and the complete proofs;
Section~30.5, Lemmas~30.5.1--30.5.2]{C_StacksFiniteCechHom}.
The complexes on $D_T$ are obtained by tensoring the fixed
component Cech complexes over $k$ with $T$.

For each of the finitely many line bundles occurring below,
choose, once over $k$, a splitting of its two-term Cech
complex onto its finite-dimensional cohomology. Write its
maps as $\iota,\rho,h$, with
\[
 \rho\iota=1,\qquad
 \iota\rho=1-dh-hd,\qquad
 \rho h=h\iota=h^2=0.                                 \tag{D.FH.5}
\]
Here $d$ is the Cech differential, $h$ lowers Cech degree,
and the target of the splitting has zero differential.
Explicitly, split the kernel in degree zero and split the
image in degree one; invert $d$ between the two complementary
subspaces to define $h$. No finite-dimensionality of the
original spaces of affine sections is required. The maps
are fixed $k$-linear maps, and (D.FH.5) remains valid after
tensoring with every $T$.

Put $P_i=[L_i^{-1}\xrightarrow{s_i}\mathcal O_D^{r_i+1}]$
in internal degrees $-1,0$. Form the total Cech Hom complex
$C_{ij}$ for $\mathcal Hom^\bullet(P_i,P_j)$. Its differential
has the form
\[
                         d+\delta,
\]
with the usual total-complex signs. The component $\delta$
is the internal Hom differential. It raises internal degree,
preserves Cech degree, and is linear in the tuple coordinates.
In particular
\[
 d\delta+\delta d=0,\qquad \delta^2=0,\qquad
                         h\delta h=0.                  \tag{D.FH.6}
\]
The last equality follows because $h$ sends Cech degree one
to zero and is zero on degree zero. Use the direct sums of
the fixed splittings (D.FH.5) in all internal degrees.

\emph{An explicit finite transfer, without variable bases.}
On the finite graded cohomology space define
\[
 \begin{split}
 I&=\iota-h\delta\iota,\qquad
 P=\rho-\rho\delta h,\\
 \Delta&=\rho\delta\iota-\rho\delta h\delta\iota.
 \end{split}                                           \tag{D.FH.7}
\]
The signs correspond to the convention in (D.FH.5).
Substitution using (D.FH.5)--(D.FH.6) gives
\[
 (d+\delta)I=I\Delta,\quad
 P(d+\delta)=\Delta P,\quad PI=1,\quad
 IP=1-(d+\delta)h-h(d+\delta).                          \tag{D.FH.8}
\]
For example the terms beyond (D.FH.7) vanish because
$h\delta h=0$; the cross term
$h\delta\iota\rho\delta h$ also vanishes after replacing
$\iota\rho$ by $1-dh-hd$ and using (D.FH.6).
These identities prove $\Delta^2=0$ and a chain homotopy
equivalence. They hold polynomially over every coefficient
ring, not only at field-valued points. The differential
has degree at most two, while $I$ and $P$ have degree at most
one in the object coordinates.

The finite space here is the direct sum of the cohomology
of the internal Hom terms
\[
 L_j^{-1\,\oplus(r_i+1)},\quad
 \mathcal O_D^{\oplus(r_i+1)(r_j+1)}\oplus L_iL_j^{-1},
 \quad L_i^{\oplus(r_j+1)},                             \tag{D.FH.9}
\]
in internal degrees $-1,0,1$, respectively. In the first
term, the notation means $(L_j^{-1})^{\oplus(r_i+1)}$.
Equations (D.FH.4) give precisely (D.FH.3).
On the open generating locus, $P_i$ is a locally free
resolution of the vector bundle $V_i$. The Cech Hom complex
therefore computes the derived Hom in (D.FH.2), over every
$T$. Exactness of the presentation is preserved by every
coefficient change because its quotient is locally free.

\emph{Composition and equality witnesses.}
The two-column ordered Cech model has its usual associative
composition product, with the total-complex signs. For
cochains in Cech degrees zero and one, the products use
the first or last restriction respectively; a product of
two degree-one Cech cochains is zero. For degree-zero
cycles $a\in H_{ij}\otimes T$ and
$b\in H_{jk}\otimes T$, define
\[
             m_2(b,a)=P_{ik}\bigl(I_{jk}(b)\,I_{ij}(a)\bigr).
                                                        \tag{D.FH.10}
\]
It is a cycle and represents actual composition. Its total
degree is at most five: two arrow coordinates, at most one
object factor from each inclusion, and at most one from the
projection. This product need not be strictly associative
on all finite cochains, and no such assertion is needed;
on degree-zero cohomology it is the associative bundle
composition. A path of any fixed length can instead be
represented by multiplying all its lifted Cech cycles first
and then applying the final projection. Its polynomial
degree is bounded by twice the path length plus one.

Every actual arrow is represented by a cycle, since
$H^0$ of (D.FH.2) is its ordinary Hom module. A proposed
equality between paths is exact precisely when the
difference of their representatives is $\Delta w$ for
some $w$ in degree $-1$. This is a finite list of polynomial
equations with a boundary witness. In particular two arrows
are mutual inverses exactly when their two compositions
differ from the identity cycles by such boundaries.
The identity cycles may be obtained by applying $P$ to the
fixed identity cochains. Thus every assertion about arrows
has its stated all-coefficient meaning; no passage from
fibrewise invertibility to an unproved lifting statement
has been made.

Finally, the object tuple uses
$(r_i+1)h^0(L_i)=O(nr_i^2)$ coordinates. Formula (D.FH.3)
bounds the cycles, boundary witnesses, and their scalar
equations by $O(nr^2)$ for each of the fixed local diagram
patterns. A bounded-length relation needs only its source,
target, and constituent Hom complexes, hence a uniformly
bounded number of them. Summing over $O(f)$ patterns gives
the claimed budget. All splittings, bases, and the finitely
many resulting multiplication coefficients are chosen once
for the fixed line bundles. They are background constants,
not additional matrix-valued moduli coordinates.
\end{proof}

\begin{remark}[What this finite model does and does not provide]
The same construction applied to $P_i$ itself gives a finite
complex of total dimension $O(nr_i)$ computing the cohomology
of $V_i$ over every coefficient algebra; its differential
is linear, since there is only one internal step. Fixed
restriction and tensor diagrams can be treated by the same
finite-width construction. This represents the derived
complexes, but does not assert cohomological flatness or
bound the equations for a chosen cohomology-rank stratum.
Likewise the proposition does not choose moving determinant
lines, realize arbitrary finite-cycle data as full modules,
or lift an arbitrarily prescribed cycle arrow to the
resolution. Those are distinct conditions in an eventual
hereditary unit-object atlas. No representation count or
finite-generation assertion follows from the proposition
alone.
\end{remark}


% Separate checked post-CP18 source. Root/B/C/D full actual-file audits PASS; no representation-count conclusion.
\subsection{Twisted minimal lattices and a depth independent of the residue period}

Throughout this subsection, $R$ is a complete excellent normal local
surface ring containing its algebraically closed residue field $k$ of
characteristic $p>0$. Let $\pi:Y\to\operatorname{Spec}R$ be a
projective regular resolution, an isomorphism off the closed point,
whose reduced exceptional divisor is one smooth connected curve $C$.
Write $g=g(C)$ and $C^2=-d<0$. For an integer $m$ prime to $p$, put
\[
  \mathcal X=\sqrt[m]{(Y,C)},\qquad
  \mathcal D=(\mathcal X\times_Y C)_{\mathrm{red}}.
  \tag{D.TD.1}
\]
Thus $m\mathcal D$ is the pullback of $C$ and
$\mathcal O_{\mathcal D}(m\mathcal D)$ is the pullback of $N_{C/Y}$.
Let $\mathcal G\to\mathcal X$ be a banded $\mu_n$-gerbe, with
$n$ also prime to $p$. Its band is called the scalar band. Hats
below mean completion along $\mathcal D$. A meromorphic map between
formal bundles means a map with some finite pole order along
$\mathcal D$; no larger analytic class of maps is intended.

\begin{lemma}[Proved: a connected mixed-band anchor and minimal slopes]
\label{D:twisted-connected-anchor-minimal-slopes}
Fix a scalar character $u\in\mathbb Z/n$. The relative stack of
scalar-weight-$u$ lines on $\mathcal G_{\mathcal D}$ is a
$\mathbb G_m$-gerbe over a finite etale $\mathbb Z/m$-torsor
$h:R_u\to C$. Suppose this torsor is connected. Then
\[
                 g(R_u)-1=m(g-1).                         \tag{D.TD.2}
\]
Choose a universal line $\Lambda_u$ on
$\mathcal G_{\mathcal D}\times_C R_u$. The resulting spectral
equivalence identifies scalar-weight-$u$ bundles with ordinary
bundles $W$ on $R_u$ and gives $\operatorname{rk}F=m\operatorname{rk}W$.

The universal line extends to the formal gerbe after the unique
finite etale lift $Y_u\to Y$ of $R_u\to C$. Its finite pushforward
is a formal scalar-weight-$u$ bundle $A_u$ of rank $m$, whose
residual spectral bundle is $\mathcal O_{R_u}$.

Let $F_0$ be any scalar-weight-$u$ bundle on $\widehat{\mathcal G}$.
Among the formal bundle lattices in its meromorphic bundle there
is a lattice $F$ whose residual spectral bundle $W$ satisfies
\[
 g(R_u)-1\ \leq\ \mu_{\min}(W)
       \ \leq\ \mu_{\max}(W)
       \ \leq\ g(R_u)-1+d.                              \tag{D.TD.3}
\]
The anchor is used to select this lattice. The statement does not
parameterize $F$ by an ordinary bundle of rank $m\operatorname{rk}F$.
\end{lemma}

\begin{proof}
Here the two band orders need not be equal. Etale-locally on $C$,
the residual gerbe is $BD$ for a central extension
\[
                 1\longrightarrow\mu_n\longrightarrow D
                    \longrightarrow\mu_m\longrightarrow1.
                                                               \tag{D.TD.4}
\]
The group is abelian, because its kernel is central and its
quotient is cyclic. After an etale localization it is
diagonalizable. There are exactly $m$ characters restricting to
$u$ on $\mu_n$, and they form a torsor under $\mathbb Z/m$.
This proves the mixed-band version of the character construction
in Proposition~\ref{D:residual-twisted-spectral-category}.
The Picard gerbe over $R_u$ is neutral since
$H^2(R_u,\mathbb G_m)=0$. Taking the selected isotypic summand
and summing its translates give the same inverse functors as in
that proposition. In particular, they commute with every
coefficient-base change. Connectedness and etale Riemann--Hurwitz
give (D.TD.2).

Proper henselian comparison lifts $R_u\to C$ uniquely to
$Y_u\to Y$; it also lifts its cyclic action
\citep[Lemma 58.9.1, Tag 0A48, complete proof]{StacksSingleCurveBrauerD}.
The obstruction to lifting an invertible sheaf across a
square-zero thickening is in $H^2$ of the corresponding ideal.
In the present situation that cohomology vanishes: the ideal is
supported on a proper one-dimensional stack with diagonalizable
stabilizers, whose exact coarse pushforward reduces the assertion
to a coherent sheaf on $R_u$. Successive lifting therefore gives
a compatible formal line $\widehat\Lambda_u$. This uses the
ordinary square-zero obstruction calculation, checked on an
etale atlas and descended; the exactness of invariants removes
any group-cohomology term. Its finite etale pushforward defines
$A_u$. Restricting to the residual gerbe and using the spectral
equivalence shows that its spectral bundle is exactly
$\mathcal O_{R_u}$.

We explain why a nonnegative integer can be minimized, without
assuming that the twisted lattice itself is a module on
$\operatorname{Spec}R$. Write $q:\widehat{\mathcal G}\to\widehat Y$
for the composite coarse map and put
\[
                        H_F=q_*\mathcal Hom(A_u,F).
                                                               \tag{D.TD.5}
\]
First descend the scalar-weight-zero Hom through the scalar gerbe;
then take root-stack invariants. The result is an ordinary formal
vector bundle on $\widehat Y$ of rank $m\operatorname{rk}F$.
For local freeness, use a root chart with equation $z^m=t$.
A vector bundle has, formally locally along $z=0$, a homogeneous
basis lifted from its character summands. The invariant part of
each shifted free summand is a free module of rank one over the
coarse chart. This also proves exactness on the elementary
transformation sequences used below.

The $C$-adic and closed-fibre-adic topologies on $Y$ are cofinal.
Ordinary projective Grothendieck existence algebraizes $H_F$ to
a vector bundle on $Y$; only this ordinary bundle, not the anchor,
needs to be algebraized here
\citep[Theorem 30.27.1, Tag 088E]{StacksTwistedDepthD}.
All $H_F$ have isomorphic restrictions to $Y\setminus C$:
a meromorphic comparison and its inverse have finite pole order,
and these ordinary Hom maps algebraize after the indicated fixed
Cartier twists. Their common restriction has a coherent reflexive
extension $M$ to $\operatorname{Spec}R$. Indeed, take a coherent
extension and its double dual; normality and the codimension-two
Hartogs property identify it with the punctured direct image.
Consequently
\[
 e(F)=\operatorname{length}_R\bigl(M/\Gamma(Y,H_F)\bigr)
                      +\dim_k H^1(Y,H_F)\ \geq\ 0             \tag{D.TD.6}
\]
is a finite integer. Both terms are supported at the closed point,
by properness and the isomorphism over the puncture. The nonempty
set of these integers has a minimum. Choose $F$ attaining it.

Let $W\twoheadrightarrow Q$ be its minimal-slope
Harder--Narasimhan quotient, and let $\mathcal Q$ be the
corresponding bundle on the residual gerbe. The elementary kernel
\[
           F'=\ker\bigl(F\longrightarrow i_*\mathcal Q\bigr)
                                                               \tag{D.TD.7}
\]
is locally free and meromorphically the same as $F$. This can be
checked on a regular root chart, where the divisor is Cartier.
Applying $q_*\mathcal Hom(A_u,-)$ is exact and its quotient is
\emph{exactly} $i_{C*}h_*Q$. This last assertion is the spectral
identity
\[
 q_{\mathcal D*}\mathcal Hom(A_u|_{\mathcal D},\mathcal Q)
                            =h_*Q.                            \tag{D.TD.8}
\]
There is no additional factor $m$ in its Euler characteristic.
Taking the finite-support Euler difference in (D.TD.6) gives
\[
                   e(F')-e(F)=\chi(R_u,Q)\geq0.
                                                               \tag{D.TD.9}
\]
Thus $\mu_{\min}(W)\geq g(R_u)-1$.

Tensoring by $\mathcal O_{\mathcal D}(\mathcal D)$ acts on the
spectral category by a deck translation $\tau$ of $R_u$, followed
by tensoring by a line $L$ of degree $-d$. To check that degree,
its $m$th normal power is $h^*N_{C/Y}$, of degree $-md$.
The rational degrees of $\Lambda_u$ and $\tau^*\Lambda_u$ agree
because $R_u$ is connected and $\tau$ is an automorphism.
The quotient defining $L$ therefore has degree $-d$.
In particular this operation changes each slope by $-d$.

For the maximal-slope subbundle $S\subset W$, take the inverse
elementary transform $F''$, characterized by
\[
       0\longrightarrow F\longrightarrow F''
          \longrightarrow i_*\bigl(\mathcal S(\mathcal D)\bigr)
          \longrightarrow0.
                                                               \tag{D.TD.10}
\]
Explicitly it is the kernel of $F(\mathcal D)$ mapping to
$(F|_{\mathcal D}/\mathcal S)(\mathcal D)$, so it is again locally
free. Minimality now gives
\[
      0\leq e(F'')-e(F)
             =-\chi\bigl(R_u,\tau_*(S\otimes L)\bigr).
                                                               \tag{D.TD.11}
\]
Hence $\mu_{\max}(W)\leq g(R_u)-1+d$. This proves (D.TD.3).
It is the relative-Euler elementary-transformation argument of
\citet[Definition 1.2 and Lemma 1.3, complete proof]
{C_Langer2026}, with (D.TD.8) supplying the precise twisted
correction rather than an assumed untwisted comparison.
\end{proof}

\begin{proposition}[Proved: uniform normal depth and quadratic Hom cohomology]
\label{D:twisted-period-independent-normal-depth}
Use the preceding setting. Allow scalar characters in a fixed
Frobenius orbit, and suppose the associated spectral torsors are
connected. Normalize their universal lines to rational degrees
in $[0,1)$, as in
Proposition~\ref{D:residual-twisted-Frobenius-transport}.
Consider minimal lattices as in (D.TD.3), their one-step absolute
Frobenius pullbacks, and any fixed finite list of ordinary twists
$bC$, with the list independent of $m,n$ and the ranks.

There are integers $B,N>0$ depending only on $g,d,p$ and that
finite list, and independent of $m,n$ and the ranks, with the
following properties whenever the scalar weights of a Hom agree.

\begin{enumerate}
\item Every meromorphic map between the permitted lattices extends
with poles at most $BC$.
\item Every bundle on the $NC$-thickening of the formal gerbe
whose closed residual spectral bundle has a permitted slope window
extends to a formal bundle. Every specified map on that thickening
between two formal extensions, with any one of the permitted
Cartier twists, lifts to a formal map. Every specified isomorphism
lifts to a formal isomorphism.
These assertions also hold over every finite local Artinian
$k$-algebra, for actual bundle families whose closed residual
bundles have the stated windows.
\item For any fixed integer $a>0$, and permitted bundles of ranks
$r,r'$, there is a constant $K_a$, independent of $m,n,r,r'$, such
that
\[
 \dim_k H^0\bigl(\mathcal G_{aC},\mathcal Hom(F,G)\bigr)
 +\dim_k H^1\bigl(\mathcal G_{aC},\mathcal Hom(F,G)\bigr)
                            \leq K_a rr'.                    \tag{D.TD.12}
\]
The same conclusion permits any fixed Cartier twists from the
list. This is a cohomology budget, not an asserted presentation
by equations of uniformly bounded degree.
\end{enumerate}
\end{proposition}

\begin{proof}
Put $g_m=g(R_u)$ and $\alpha_m=g_m-1=m(g-1)$.
The minimal residual windows are $[\alpha_m,\alpha_m+d]$.
For a one-step Frobenius pullback the exact spectral transport is
\[
                 W\longmapsto\sigma_*\bigl(F_{R_u}^*W\otimes B_u\bigr),
                       \qquad 0\leq\deg B_u\leq p-1.
                                                               \tag{D.TD.13}
\]
The mixed-band proof is unchanged: raise the characters of the
central group (D.TD.4) to their $p$th powers and compare the
normalized universal lines. Their $mn$th powers descend, so the
normalization and degree calculation are valid for unequal band
orders as well. The final translation around a Frobenius cycle
must still be retained.

Write $\gamma_m=\max(g_m-1,0)$. The one-step curve inequalities
give a containing window
\[
 [\,p\alpha_m-(p-1)\gamma_m,
    \ p(\alpha_m+d)+(p-1)\gamma_m+p-1\,].                     \tag{D.TD.14}
\]
Here we use \citet[Theorem 2.2 and proof]{C_Sun2008} and
\citet[Proposition 1.3 and proof]{C_MehtaPauly2007}, with
Harder--Narasimhan filtrations; genus zero follows from splitting.
These are exactly the rank-independent curve bounds already used
in Lemma~\ref{D:local-reflexive-constructible-atlas}.
Ordinary $bC$ shifts the spectral slopes by $-bmd$.
Thus all the finitely many permitted windows lie in
$[L_m,U_m]$ with
\[
                   U_m-L_m\leq a_0m+a_1                    \tag{D.TD.15}
\]
for fixed constants $a_0,a_1$.

Let $T$ be a finite local Artinian coefficient algebra. Filter
successive root thickenings by powers of the root divisor ideal.
A layer of the coarse Hom sheaf, with twist $j\mathcal D$, has
cohomology equal to that of
\[
              \mathcal Hom\bigl(W_F,\tau^j_*(W_G\otimes L_j)\bigr)
                  \quad\hbox{on }R_u,\qquad \deg L_j=-jd.
                                                               \tag{D.TD.16}
\]
This is the same character calculation as (D.TD.8).
Translations preserve slopes and $g_m$. Consequently its $H^0$
vanishes if
\[
                       U_m-jd<L_m.                         \tag{D.TD.17}
\]
For the conormal layer, replace $j$ by $-j$.
Serre duality identifies its $H^1$ dual with an ordinary Hom
whose target slopes are at most $U_m+2g_m-2-jd$.
It therefore vanishes if
\[
                         L_m+jd>U_m+2g_m-2.                 \tag{D.TD.18}
\]
This argument uses the absence of morphisms across a strict
slope gap; it makes no tensor-semistability assertion in positive
characteristic.

Since (D.TD.15) and $2g_m-2=2m(g-1)$ are linear in $m$,
there are fixed $B,N$ for which (D.TD.17) holds for all
$j\geq mB$ and (D.TD.18) for all $j\geq mN$, also after the
finitely many prescribed twists. The pole restriction sequences
then show that global Hom sections stabilize after $BC$.
Every meromorphic map has some finite pole order by definition,
so the stable space contains it. This proves the first assertion.

The obstruction to lifting a bundle to the next root thickening
lies in $H^2$ of its End bundle tensored with that conormal layer.
It is zero because the layer has one-dimensional proper coarse
space and exact invariants. The obstruction to lifting a
specified map is the corresponding $H^1$, which is zero after
$NC$ by (D.TD.18). Lifting successively and taking compatible
formal systems proves existence and the prescribed-map assertion.
A lifted isomorphism is an isomorphism because it is so modulo
the ideal of definition; one can lift its inverse or apply
complete Nakayama locally. Over $T$, filter each locally free
layer by powers of the maximal ideal of $T$. Its subquotients
are copies of the closed-fibre layer. Thus the asserted $H^0$
and $H^1$ vanishings persist over every such $T$, and the same
obstruction arguments lift the specified Artinian data. This
proves the second assertion without assuming a reflexive hull
commutes with arbitrary base change.

Finally write $r=ms$ and $r'=ms'$. The thickening $aC$ has
$ma$ root layers. In each of them, (D.TD.16) is a Hom between
bundles of ranks $s,s'$ on a curve of genus $g_m$, with slope
endpoints bounded in absolute value by a fixed multiple of $m$.
Its $H^0$ has dimension at most $c_0mss'$: evaluation at
$c_0m$ distinct points is injective once subtracting those points
creates the strict gap (D.TD.17). Serre duality and the same
evaluation argument give $h^1\leq c_1mss'$. These estimates
also avoid tensor semistability. Summing the cohomology
inequalities over the $ma$ layers gives
\[
             h^0+h^1\leq a(c_0+c_1)m^2ss'
                            =a(c_0+c_1)rr'.
\]
This proves (D.TD.12).
\end{proof}

\begin{remark}[What the uniform depth does not supply]
The propositions establish formal lattice normalization, finite
normal-depth realization and specified-map lifting, and a
quadratic Hom-cohomology budget under the stated actual-bundle
hypotheses. They do not construct a parameter scheme with
$O(fr^2)$ coordinates and bounded-degree equations for the
graded finite-thickening objects and their Frobenius arrows.
In particular, they do not prove the remaining unbounded-period
projective-representation count or a full-generation statement.
\end{remark}


% Separate checked post-CP19 source. No geometric object atlas is asserted.
\subsection{A weighted component bound for a prospective normal-layer atlas}

\begin{lemma}[Proved: weighted component bound with arbitrary characteristic]
\label{D:weighted-affine-component-bound}
Let $k$ be algebraically closed, let $x_1,\ldots,x_N$ have
positive integral weights $w_1,\ldots,w_N$, and let
$D\geq\max_i w_i$ be an integer. If a closed subscheme
$V\subset\mathbb A^N_k$ is defined by any finite number of
polynomials of weighted degree at most $D$, then
\[
       \#\operatorname{Irr}(V)
          \leq (N+1)\frac{D^N}{\prod_{i=1}^N w_i}.
                                                        \tag{D.WB.1}
\]
The same bound applies to an open subscheme of $V$ and to the
number of nonempty clopen parts in any partition of that open
subscheme. The weights may be divisible by $\operatorname{char}k$.
\end{lemma}

\begin{proof}
We use the ordinary isolated affine intersection bound, including
local scheme lengths, in \citet[Section IV.4.2, equations
(38)--(39), Proposition IV.28, and Corollary VIII.3 following
Theorem VIII.2, pp.~70--71 and 143--144]{Mondal2021WeightedD}.
For $N$ polynomials in $N$ variables of ordinary degrees at
most $D$, the sum of the lengths at their isolated zeroes is
at most $D^N$, even when other components have positive
dimension. The source works over every algebraically closed
field. Its weighted Theorem VIII.2 also directly supplies
the next square-system bound; we record the finite-flat
comparison to make the inseparable-weight scope explicit.

Apply it now to a square weighted system by the finite flat
coordinate-powering morphism
\[
 \mathbb A^N_y\longrightarrow\mathbb A^N_x,
                   \qquad x_i=y_i^{w_i}.                 \tag{D.WB.2}
\]
The monomials $\prod y_i^{a_i}$ with $0\leq a_i<w_i$ form
a basis over $k[x_1,\ldots,x_N]$, so its rank is
$W=\prod_i w_i$, without a separability hypothesis.
A polynomial of weighted degree at most $D$ pulls back to
one of ordinary degree at most $D$. Above an isolated local
zero scheme, the sum of the lengths of the inverse-image local
schemes is exactly $W$ times its length. Thus the number of
isolated geometric points of any square weighted system is
at most $D^N/W$. This includes points with zero coordinates
and inseparable coordinate powers.

We reduce component counting to this square isolated case.
Fix a dimension $d$. Let $I$ be the ideal generated by the
given polynomials. For each $d$-dimensional irreducible
component with generic prime $P$, the regular local ring
$k[x]_P$ has dimension $N-d$, and $I_P$ is primary for its
maximal ideal. Choose $N-d$ constant $k$-linear combinations
of the given generators that form a system of parameters
simultaneously at all these finitely many primes. This
choice is possible by successive prime avoidance: at each
step the bad coefficient choices form finitely many proper
linear conditions over the residue fields, and their
restrictions to the infinite field $k$ do not cover the
coefficient space. Equivalently one can make the choices
over independent generic constants and then specialize in
the resulting nonempty open conditions.

Each selected $d$-dimensional component is consequently an
actual component of the zero scheme of these $N-d$ equations.
Intersect with $d$ general ordinary affine hyperplanes.
They meet each such component in a nonempty finite set in
its good locus, avoiding its intersections with all the
other components. For completeness, this follows by taking
the projective closures: their boundaries and all proper
intersection loci have dimension less than $d$, whereas
$d$ general hyperplanes cut a nonempty zero-dimensional
section of each $d$-dimensional projective component.
The selected points are isolated in the full intersection,
including when larger-dimensional components were present.
Every hyperplane has weighted degree at most $\max_iw_i\leq D$.
We have obtained a square system of $N$ equations of the
required degree, with at least one isolated point for each
original $d$-dimensional component. The square bound gives
\[
       \#\{\text{$d$-dimensional components of }V\}
                         \leq D^N/W.
\]
Summing over $0\leq d\leq N$ proves (D.WB.1).
An open subset of an irreducible scheme is irreducible if
nonempty, so passing to an open cannot increase this number.
Finally every nonempty clopen part contains an irreducible
component, and no irreducible component meets two such parts.
\end{proof}

\begin{corollary}[Proved: the normal-layer product has a quadratic exponent]
\label{D:weighted-normal-layer-budget}
Fix positive constants $A,K$ and an integer $a\geq1$.
Let $m,s,f$ be positive integers, put $r=ms$, and put $q=p^f$
for any fixed prime $p$. Suppose an affine polynomial
presentation has at most $Afms^2$ variables at each weight
$j\in\{1,\ldots,am\}$, and at most $Afms^2$ additional
variables of weight at least one. Suppose the additional
weights are at most $Km$, and every equation has weighted
degree at most $Km$. Enlarging $K$ so that $K\geq a$, its
number of irreducible components is at most
\[
                              q^{C r^2},                 \tag{D.WB.3}
\]
where $C$ depends only on $A,K,a,p$. The same conclusion
holds for any open part and for a partition of it into
nonempty clopen classes.

This statement is conditional on an actual polynomial
presentation with those weights. It does not construct a
presentation of bundles, deformation torsors, or Frobenius
objects from a bound on their tangent spaces.
\end{corollary}

\begin{proof}
Round the fixed constants upward. Pad with unused variables
and equations setting those variables equal to zero until
there are $b=\lceil Afms^2\rceil$ variables at every weight
$1\leq j\leq L=am$ and $b$ additional variables. Padding
does not change the scheme or its components. It does not
violate the degree bound, because every added weight is at
most $Km$. In the denominator of (D.WB.1), conservatively
replace each additional weight by one. If $N=b(L+1)$, the
weighted component bound is at most
\[
        (N+1)\frac{(Km)^{b(L+1)}}{(L!)^b}.              \tag{D.WB.4}
\]
Using $L!\geq(L/e)^L$ yields
\[
 \log\!\left(\frac{(Km)^{b(L+1)}}{(L!)^b}\right)
       \leq bL\log(eK/a)+b\log(Km).
                                                        \tag{D.WB.5}
\]
Both terms are $O(fm^2s^2)$, since $\log m\leq m$ and
$b=O(fms^2)$. Also $\log(N+1)=O(fm^2s^2)$.
Thus (D.WB.4) is at most $\exp(C_0fr^2)$, which is
$q^{(C_0/\log p)r^2}$. This proves (D.WB.3).
In particular, using one common weighted degree $Km$ is
compatible with the bound; it is essential to retain the
product of the variable weights in the denominator.
\end{proof}

\begin{corollary}[Proved: polynomial background coefficients can be separated]
\label{D:weighted-normal-layer-background-budget}
Fix positive constants $A,K,h$ and an integer $a\geq1$.
Let $m,s,f\geq1$, $r=ms$, $q=p^f$, and let $E$ be a
positive integer with $E\leq K(ms)^h$.
Consider a polynomial presentation with the following variables:
\begin{itemize}
\item at most $Afms$ background variables of weight one;
\item at most $Afms^2$ variables at each weight
$Ej$, $1\leq j\leq am$;
\item at most $Afms^2$ further bundle variables of weight $E$.
\end{itemize}
Suppose every equation has weighted degree at most $KEm$,
enlarging $K$ so that $K\geq a$. Then its number of
irreducible components, and the number of nonempty clopen
classes on any open part, is at most $q^{Cr^2}$, with $C$
depending only on $A,K,h,a,p$.

In particular, this numerical bound permits polynomially
growing degrees in a smaller block of background variables,
provided the actual mixed polynomials satisfy the displayed
weighted-degree condition. That condition is not implied
by separate tangent-space dimensions.
\end{corollary}

\begin{proof}
As before, round the constants upward and pad the three
blocks by zero variables. Put
$b=\lceil Afms^2\rceil$, $c=\lceil Afms\rceil$, and $L=am$.
There are $N=b(L+1)+c$ variables. Lemma
\ref{D:weighted-affine-component-bound} bounds the number
of components by
\[
 (N+1)\frac{(KEm)^N}{E^{b(L+1)}(L!)^b}
   =(N+1)\frac{(Km)^{b(L+1)}}{(L!)^b}(KEm)^c.
                                                        \tag{D.WB.6}
\]
The first fraction has logarithm $O(fm^2s^2)$ by the
calculation in (D.WB.5). The extra logarithm is
\[
 c\log(KEm)
       =O\bigl(fms(1+\log(ms))\bigr)
       =O(fm^2s^2),                                    \tag{D.WB.7}
\]
using $E\leq K(ms)^h$ and $1+\log x\leq 2x$ for $x\geq1$.
Also $\log(N+1)=O(fm^2s^2)$. Conversion from
$\exp(O(fr^2))$ to $q^{Cr^2}$ gives the result.
\end{proof}


% Proved post-CP19 auxiliary. Actual fixed-determinant bundles and arrows;
% moving determinant backgrounds and the final unit-object count are separate.
\subsection{Actual normal-layer bundles with fixed residual determinants}

We use the two-affine Cech convention of
Proposition~\ref{D:fixed-determinant-finite-Hom}. On a tame gerbe
over a proper curve, ``affine chart'' below means the inverse image
of an affine coarse open. Such a chart and the intersection of two
such charts are cohomologically affine: the exact invariant functor
followed by the affine section functor is exact. This remains true
after every coefficient-base change. Consequently the ordered
two-chart Cech complex, in degrees zero and one, computes coherent
cohomology and commutes with extension of the ground field to any
$k$-algebra. The scheme version and the base-change calculation are
the complete statements used in
\citep[Tags 0BDX, 01XD and 02KE]{C_StacksFiniteCechHom}; exact
diagonalizable invariants give the stated stack version.

\begin{lemma}[Proved: a polynomial normal-layer presentation of actual bundles]
\label{D:weighted-fixed-residual-bundle-deformations}
Let $\mathcal Z_J$ be a nilpotent thickening of a proper
one-dimensional tame stack $\mathcal Z_0$ over a field $k$.
Suppose its ideal $I$ is locally principal, $I^{J+1}=0$, and
$I^j/I^{j+1}$ is invertible on $\mathcal Z_0$ for $0\leq j\leq J$.
Fix two cohomologically affine opens as above and reference bundles
$K^\circ,A^\circ$ on $\mathcal Z_J$. Put $K_0=K^\circ|_{\mathcal Z_0}$
and $A_0=A^\circ|_{\mathcal Z_0}$.

There is a polynomial family of marked bundles $K$ on
$\mathcal Z_J$ with $K|_{\mathcal Z_0}=K_0$, using
\[
       h^1\bigl(\mathcal Z_0,
          \mathcal End(K_0)\otimes I^j/I^{j+1}\bigr)
                                                        \tag{D.WF.1}
\]
variables of weight $j$ at every $1\leq j\leq J$.
It covers every marked bundle over every $k$-algebra, up to
isomorphism preserving the marking. This assertion is about
actual bundles on the fixed thickening, not only their tangent
spaces.

Over this family there is a finite polynomial complex computing
$R\Gamma\mathcal Hom(K,A^\circ)$ over every coefficient algebra.
Its underlying graded vector spaces are the direct sum, over
$0\leq j\leq J$, of
\[
 H^0\bigl(\mathcal Hom(K_0,A_0)\otimes I^j/I^{j+1}\bigr),\qquad
 H^1\bigl(\mathcal Hom(K_0,A_0)\otimes I^j/I^{j+1}\bigr).
                                                        \tag{D.WF.2}
\]
A coordinate in the first space at level $j$ is assigned weight
$\max(1,j)$. Its cycle equations have weighted degree at most
$j+1$ at output level $j$. The resulting cycles parameterize
\emph{actual global maps} $K\to A^\circ$. Restricting to the open
locus where the residual map is a subbundle inclusion produces
actual locally free cokernels over every coefficient algebra.

For a fixed finite list of such quotient complexes
$P=[K\to A^\circ]$, in internal degrees $-1,0$, their derived
Hom complexes, equality-of-map witnesses, and bounded-length
composition or inverse diagrams also admit finite polynomial
presentations. Their vector-space sizes are bounded by the
corresponding sums of $H^0,H^1$ of the four associated-graded
termwise Hom bundles. If all term and arrow coordinates at layer
$j$ have weight $\max(1,j)$, the weighted degrees are at most
$J+c$, where $c$ depends only on the fixed diagram length.
One-step absolute Frobenius increases this bound by a fixed
factor depending on $p$, without assuming that Frobenius is flat
on the coefficient algebra.
\end{lemma}

\begin{proof}
Write the chosen opens as $U_0,U_1$, with intersection $U_{01}$.
The reference bundles identify their restrictions on the overlap.
For every $j$, choose representatives of a basis of
$H^1(\mathcal End(K_0)\otimes I^j/I^{j+1})$ in the degree-one
Cech space, and lift each representative to
$\Gamma(U_{01},I^j\mathcal End(K^\circ))$.
Lifting is possible by cohomological affineness. Denote the
chosen lifts by $c_{j\nu}$. Glue the two restrictions of
$K^\circ$ by the automorphism
\[
          g=\prod_{j=1}^{J}
                   \left(1+\sum_\nu x_{j\nu}c_{j\nu}\right),
                                                        \tag{D.WF.3}
\]
in increasing order of $j$. Each factor is invertible, since its
nonidentity part is nilpotent. This is a genuine polynomial
gluing family; there is no triple-overlap condition for the
two-open cover. Every coefficient of normal order at most $a$
is a polynomial of weighted degree at most $a$. The same holds
for $g^{-1}$, by the finite geometric series.

Here is coverage over an arbitrary $k$-algebra $T$. A marked
bundle $K_T$ is isomorphic to $K^\circ_T$ on each $U_{i,T}$:
lift the given residual isomorphism one nilpotent layer at a
time. Each obstruction is a first cohomology group on a
cohomologically affine chart and is zero. The lifted maps are
isomorphisms because they are so modulo a nilpotent ideal.
Thus $K_T$ is described by one transition equal to the identity
modulo $I$. At normal order $j$, after the preceding orders
have been normalized, its discrepancy is a Cech one-cochain
in the $j$th associated-graded End bundle. Split that cochain
as a chosen $H^1$ representative plus a Cech coboundary.
The two components of the coboundary are removed by chart
automorphisms equal to the identity modulo $I^j$. They do not
alter the earlier orders. Every splitting used here is a
fixed $k$-linear splitting tensored with $T$, so the argument
works for every $T$, including nonreduced and non-Noetherian
algebras. Induction proves (D.WF.3) covers the stated functor
of marked isomorphism classes.

We next give the finite map construction, including its
weighted-degree assertion. Split the normal filtration on
each Cech section space as $k$-vector spaces. Its associated
graded is the sum of the two-term Cech complexes of the
bundles in (D.WF.2). For each of these complexes choose a
strong contraction onto its cohomology, and sum them. If $d$
is the associated-graded Cech differential, the resulting
maps satisfy
\[
 \rho\iota=1,\quad
 1-\iota\rho=dh+hd,\quad
 h^2=h\iota=\rho h=0.                                  \tag{D.WF.4}
\]
The full differential for $\mathcal Hom(K,A^\circ)$ is
$d+\epsilon$. Its correction $\epsilon$ strictly raises
normal order. This includes both the fixed nonsplit
reference neighbourhood and the variable transition (D.WF.3).
A correction from level $b$ to level $a>b$ has weighted
degree at most $a-b$. The fixed reference corrections have
degree zero. All these statements are identities between
section spaces before taking cohomology.

Set
\[
 \begin{split}
 S&=(1+h\epsilon)^{-1},\qquad
 T_1=(1+\epsilon h)^{-1},\qquad
 \Delta=\rho\epsilon S\iota,\\
 I_1&=S\iota,\qquad P_1=\rho T_1,\qquad H_1=hT_1.
 \end{split}                                           \tag{D.WF.5}
\]
Both inverses are finite geometric sums, because normal order
strictly increases. Expanding the products, using (D.WF.4)
and $(d+\epsilon)^2=0$, gives
\[
 \begin{gathered}
 (d+\epsilon)I_1=I_1\Delta,\qquad
 P_1(d+\epsilon)=\Delta P_1,\qquad P_1I_1=1,\\
 I_1P_1=1-(d+\epsilon)H_1-H_1(d+\epsilon),\qquad
 \Delta^2=0.                                           \tag{D.WF.6}
 \end{gathered}
\]
These are the special deformation-retract identities in
\citet[Sections 2.1--2.4, complete proof, and Section 3.2,
pp.~1--4]{C_Crainic2004Perturbation}, with the sign of the
contracting homotopy reversed to match (D.WF.4).
Here the inverses are finite algebraic sums, so the identities
survive tensoring with every $T$. Equivalently they follow by cancellation of
adjacent terms in the geometric sums; the only uncancelled
terms are the two sides of (D.WF.4). No flatness of a
cohomology module is used.

In a matrix entry from normal level $b$ to $a$, every product
in (D.WF.5) has weighted degree at most $a-b$. Give an input
map coordinate at level $b$ weight $\max(1,b)$. The cycle
equation at output $a$ therefore has degree at most $a+1$.
The same estimate applies to the lifted map $I_1v$.
There is no Cech term in degree $-1$. Thus, on degree-zero
cycles, (D.WF.6) says $I_1P_1=1$, and $P_1I_1=1$ already
holds. The finite cycles and actual global maps are inverse
descriptions, not merely a surjection onto a derived quotient.

The subbundle condition for the residual map is open: its
failure is the proper image of the determinantal closed
locus on $\mathcal Z_0$ (equivalently on its coarse curve).
Locally a nonvanishing maximal minor gives a split inclusion.
If it is a split inclusion modulo the nilpotent ideal $I$,
the same minor is a unit upstairs. Hence its cokernel is
locally free over every $T$, with the required rank.

For quotient complexes, use the same associated-graded Cech
contraction on each termwise Hom bundle. There are now two
types of correction: positive-normal-order terms, which
preserve internal degree, and the quotient-map differential,
which raises internal degree. In the Hom of two complexes
in degrees $-1,0$, the internal degrees are $-1,0,1$.
Consequently every nonzero word in the geometric sums
contains at most two internal differential factors; all
remaining factors strictly raise normal order. The sums in
(D.WF.5) therefore still terminate. Each internal factor at
normal order $a$ has weighted degree at most $a+1$, by the
map construction just proved. The resulting entry from
normal order $b$ to $a$ has weighted degree at most
$a-b+2$. This argument is made over the coordinate ring
where the map cycle equations hold, so the actual total
differential squares to zero there.

The transferred complex computes derived Hom. Its degree-zero
cycles modulo degree-minus-one boundaries are exactly actual
maps between the locally free cokernels. To require equality
of two such maps, introduce a degree-minus-one vector $w$
and the equation that their difference equals $\Delta w$.
For composition, lift cycles with $I_1$, use the ordered
Cech cup product and matrix composition, then apply $P_1$.
For a path of fixed length, multiply all lifted Cech maps
first and project once. The internal width and path length
bound the number of quotient-map factors by a constant;
normal orders add. Thus the asserted $J+c$ bound holds
for actual equality and inverse witnesses. Associativity
is used in the Cech complex; transferred representatives
need only be associative modulo the explicitly retained
boundaries.

All fixed Cartier twists can be included among the reference
termwise Hom bundles. A bounded-pole inverse diagram has
products equal to the canonical maps into the corresponding
Cartier twists, rather than products equal to the identity
before dividing the pole. This introduces no variable
coefficient. Finally absolute Frobenius sends a layer-$j$
coefficient to its $p$th power and normal order $j$ to $pj$.
Its contribution has the same weighted bound, multiplied
by the fixed number $p$. The quotient sequence is locally
split, so its pullback by Frobenius is exact over every
coefficient algebra, even if that Frobenius is not flat.
This proves all assertions.
\end{proof}

\begin{proposition}[Proved: fixed-determinant twisted bundles have a weighted quadratic presentation]
\label{D:weighted-fixed-determinant-normal-presentation}
Use the one-smooth-curve setting and connected spectral torsors
of Proposition~\ref{D:twisted-period-independent-normal-depth}.
Write $r=ms$ and consider any fixed number of types of adjacent
Frobenius, bounded-pole and inverse diagrams. For a cycle of
length $f$, fix the residual spectral determinant lines
$L_0,\ldots,L_{f-1}$, after the fixed ordinary twist specified
below. Each line is fixed geometric data, not a parameter.

For every fixed ordinary depth $a$, the actual bundle and
map data on the $aC$-thickening admit one open part of an affine
polynomial presentation for the fixed reference collection,
with at most $A fms^2$ variables
at each root-normal weight $1\leq j\leq am$ and at most
$A fms^2$ further variables of weight one. All equations
have weighted degree at most $K m$. The constants $A,K$
depend only on the original curve, its normal degree,
$p,a$, and the fixed diagram types; they are independent
of $m,s,f$ and of the fixed lines in the permitted degree
range. Every parameter value produces actual bundles and
maps over every coefficient algebra. Every bundle family
over a finite local Artinian coefficient algebra, with
the prescribed determinant markings and the stated
residual slope windows, occurs up to isomorphism.

This is a presentation with \emph{fixed residual determinant
lines}. It does not count the possible lines, does not
assert a moving-Picard chart, and does not by itself count
unit objects on the punctured surface.
\end{proposition}

\begin{proof}
Let $g_m-1=m(g-1)$ and let $A_u$ be the connected mixed-band
anchor of Lemma~\ref{D:twisted-connected-anchor-minimal-slopes}.
Choose a fixed positive integer $t$ so large that, for every
$m$, shifting the minimal window by $tmd$ gives
\[
 \mu_{\min}\bigl(W\otimes h^*N_{C/Y}^{-t}\bigr)
                          >2g_m-1.                     \tag{D.WF.7}
\]
For example any fixed $t$ with $td>g+2$ suffices, after
enlarging it for the finite list of prescribed windows.
Put $F^+=F(-tC)$ and write $W^+$ for its spectral restriction.
The curve slope criterion and Serre duality give global
generation and $H^1(W^+)=0$. A general $s+1$-dimensional
space of sections generates $W^+$: the rank-drop condition
has codimension two in the section space at a fixed point,
so its incidence over a curve is a proper closed subset.
Thus
\[
       0\longrightarrow L^{-1}
          \longrightarrow\mathcal O_{R_u}^{s+1}
          \longrightarrow W^+\longrightarrow0,
          \qquad L=\det W^+.                           \tag{D.WF.8}
\]
For a family over a finite local Artinian $T$, its closed-fibre
generators lift: the same first-cohomology vanishing persists
by the maximal-ideal filtration, and Nakayama gives generation.
With a determinant marking, the kernel is precisely $L_T^{-1}$.

Every conormal layer of $\mathcal Hom(A_u,F^+)$ has spectral
bundle a translate of $W^+$ tensored by a line of degree $jd$,
$j\geq0$. Its first cohomology is zero by (D.WF.7), also over
$T$ by devissage. Consequently the residual generating map
in (D.WF.8) lifts, with the source $A_u^{s+1}$ kept fixed,
through all root-normal thickenings. A map which is surjective
modulo the ideal of definition remains surjective. Its
kernel is a formal vector bundle $K$ of rank $m$, whose
spectral residual line is $L^{-1}$. Only this kernel needs
to deform; it is not assumed to remain a spectral line
away from the residual gerbe.

Fix a reference lift of $L^{-1}$ on the finite etale resolution
cover, and tensor it with the selected universal scalar line.
Such a formal lift exists because its successive obstructions
lie in $H^2$ of sheaves on a curve. Its finite pushforward
gives a reference rank-$m$ bundle $K^\circ$. Together with
the fixed source $A_u^{s+1}$, it satisfies the hypotheses of
Lemma~\ref{D:weighted-fixed-residual-bundle-deformations}.
The scalar gerbe and the root stack have exact invariants,
so its two-chart argument applies literally. The quotient
presentation covers the actual families just described.

We verify the dimension bounds without a hidden rank-$mr$
parameter space. The residual determinant has
\[
                         |\deg L|\leq c_0ms.            \tag{D.WF.9}
\]
In an End layer of $K^\circ$, the two determinant degrees
cancel: the spectral line is of the form
$L\otimes\tau^{j*}L^{-1}$ times a line of degree $jd$.
Its $h^0+h^1$ is $O(m)$ for $0\leq j\leq am$.
An End layer of $A_u^{s+1}$ is a sum of $(s+1)^2$ lines
of degree $O(m)$ and has total cohomology $O(ms^2)$.
A Hom layer in either direction between $K^\circ$ and
$A_u^{s+1}$ is a sum of $s+1$ lines with degrees $O(ms)$,
so it also has total cohomology $O(ms^2)$.

These estimates use the elementary bound
$h^0(M)+h^1(M)\leq 2g_m+|\deg M|+2$ for a line bundle on
a smooth curve. They therefore require no tensor-semistability
assertion. For adjacent Frobenius terms, the exact spectral
transport is a curve Frobenius pullback, a translation, and
a line of degree between zero and $p-1$. The determinant
degrees are at most multiplied by the fixed $p$.
The same $O(ms^2)$ estimate consequently holds for every
termwise Hom appearing in the fixed list of diagrams,
including fixed ordinary Cartier twists. All constants are
independent of the particular determinant lines.

Apply the preceding lemma at the $ma$ root layers. Kernel
transition variables, quotient-map coordinates, arrow cycles,
and equality or inverse witnesses each use at most a fixed
multiple of $ms^2$ coordinates per layer and per stage.
Layer-zero coordinates join the additional weight-one block.
The finite internal width of the quotient complexes prevents
an extra factor equal to the number of layers: an expansion
can contain at most two internal quotient-map factors between
normal-order steps. Its weighted degree is therefore
$O(ma+1)=O(m)$, including the one-step Frobenius substitutions.
The bounded-pole Frobenius diagrams use the canonical Cartier
inclusions for their inverse products. They do not assert
that a minimal lattice itself is a regular unit lattice at
the exceptional gerbe.

Restrict to the open residual slope windows when needed.
Here is the precise proper-Quot argument. In the finite-type
residual family choose one twist making both $W$ and $W^\vee$
relatively generated. Every positive-rank coherent quotient
of either bundle then has a degree bounded below by a fixed
integer for each rank. A failure of a prescribed lower
slope bound is detected by a quotient of rank $1,\ldots,s$
and degree below the corresponding threshold. There are
only finitely many such Hilbert polynomials. Include all
coherent quotients in these Quot spaces: removing their
torsion can only decrease the degree, so this detects the
same bad locus as torsion-free quotients.

For each polynomial, the Quot space is proper over the
parameter scheme. To see that the exact input applies,
embed the fixed curve projectively and present the twisted
family by a quotient of a fixed finite sum of structure
sheaves. Its Quot functor is a closed subfunctor of that
proper Quot space, given by vanishing of the map from the
presentation kernel to the universal quotient. Properness
of the ambient space is Lemma 108.6.1, Tag 0DPA; the Hom
functor and its zero locus are represented as in Proposition
99.3.10, Tag 08K6. We have read both complete proofs
\citep{C_StacksBoundedReflexive}. Their proper images are
closed. Apply this argument to $W$ and $W^\vee$ to obtain
the two slope inequalities. Thus the window condition is
one open restriction of the single polynomial family.
It introduces no new variables or equations in the closed
polynomial ambient.
The subbundle condition is the open condition already proved
in the lemma. These constructions prove the claimed actual
presentation and all-ring meaning.

For clarity, its bounded-pole finite diagrams can be chosen
deep enough for formal realization by
Proposition~\ref{D:twisted-period-independent-normal-depth}.
First extend the actual finite bundles, then lift their
specified maps. If the retained inverse product equals
the canonical map into $2BC$, the lifted product differs
from it by normal order $aC$. For $a>2B$, division by the
fixed pole gives $1$ plus a positive-order formal endomorphism,
which is invertible. This yields a meromorphic formal
isomorphism. This observation does not identify arbitrary
moving determinant backgrounds, nor replace the separate
finite-type unit-isomorphism and algebraization arguments
needed for a representation count.
\end{proof}


% Separate post-CP19 proved packet. Four actual-file audits passed; background data only.
\subsection{Moving line backgrounds without matrix-valued basis variables}

\begin{lemma}[Proved: compressed minors give a bounded rational chart cover]
\label{D:compressed-minor-background-charts}
Let $X$ be a finite-type scheme over an algebraically closed field
$k$, with $\dim X\leq b$, and let $A$ be an $M\times N$ matrix
of regular functions on $X$. On its scheme-theoretic rank-$\rho$
locus, there are $b+1$ pairs of constant matrices
\[
 U_\alpha\in\operatorname{Mat}_{\rho\times M}(k),\qquad
 V_\alpha\in\operatorname{Mat}_{N\times\rho}(k)
\]
such that the opens $\det(U_\alpha A V_\alpha)\ne0$ cover.
On each open, image, kernel, and cokernel have explicit rational
bases, with no matrix-valued moduli variables. If the entries
of $A$ are rational functions in $b$ ambient coordinates, of
numerator and denominator degree at most $E$, these bases have
degrees bounded by a polynomial in $M,N,E$ after clearing
the existing denominators.

Here the rank-$\rho$ locus means that all $(\rho+1)$-minors
vanish and at least one $\rho$-minor is invertible. These are
scheme conditions, including over nonreduced coefficient rings.
For $\rho=0$ no compression is needed.
\end{lemma}

\begin{proof}
On the reduced rank-$\rho$ locus fix a geometric point $x$.
The polynomial $\det(UA(x)V)$ in the entries of $U,V$ is not
identically zero. Its zero set is therefore a proper hypersurface
in the affine space of compression pairs. For $b+1$ independent
pairs, the incidence set on which all these determinants vanish
has dimension at most
\[
 b+(b+1)(\dim\operatorname{Mat}_{\rho\times M}
          +\dim\operatorname{Mat}_{N\times\rho}-1),
\]
strictly smaller than the dimension of the space of all pairs.
Its constructible image has closure of smaller dimension.
Choose constants outside that closure. The resulting opens
cover all geometric points and hence the entire rank locus,
including its nilpotent structure.

Write $U,V$ for one selected pair and $B=UAV$. The constant
matrices $U,V$ have full rank. Choose a constant complement
$W$ to the columns of $V$, so that $[V\ W]$ is invertible.
Then $AV$ is an image basis and
\[
                       W-VB^{-1}UAW                      \tag{B.MB.1}
\]
is a kernel basis. Indeed, the rank equations and the invertible
minor give $A=AVB^{-1}UA$ over every coefficient ring on this
chart. Likewise the fixed space $\ker U$ is a complement to
$\operatorname{im}A$, giving a cokernel basis and a projection.
The inverse $B^{-1}=\operatorname{adj}(B)/\det B$ proves the
degree assertion. This is a finite polynomial calculation;
no choice of a basis over a residue field is being extended
without its determinant being inverted.
\end{proof}

\begin{proposition}[Proved: rational moving-line cohomology and Cech backgrounds]
\label{D:moving-line-polynomial-background}
Fix a constant $c$ and an integer $h_0$. Let $m,s\geq1$, and
let $X$ be a smooth projective connected curve over an
algebraically closed field $k$, of genus at most $cm$, with a
fixed finite separable map $x:X\to\mathbb P^1$ of degree at
most $cm$. Fix at most $h_0$ moving line bundles with prescribed
degrees of absolute value at most $cms$.

There are ordered-divisor parameter charts with $O(ms)$ ambient
coordinates, equations of degree polynomial in $m,s$, and open
conditions, having the following properties.
\begin{enumerate}
\item Every collection of the prescribed lines over $k$ occurs.
For every finite local Artinian $k$-algebra with residue field $k$,
each collection of line-bundle deformations lifts to such a chart
after a choice of a geometric divisor presentation. The chart
formulas themselves are valid over arbitrary coefficient rings.
\item Consider $O(m)$ line families obtained from this fixed
number of moving lines by tensor products of bounded length,
bounded positive or negative powers, one-step Frobenius in a
fixed characteristic, and fixed deck translations preserving
$x$, together with fixed line twists of degree $O(ms)$.
Their section and cohomology complexes, evaluations, and fixed
finite-pole restriction maps are represented by matrices of
sizes and rational-entry degrees polynomial in $m,s$.
No entries of these matrices are additional moduli coordinates.
\item The base admits a cover by locally closed rational charts
on which the relevant cohomology modules are free and all the
specified matrices have rational splittings. The number of
choices is at most
\[
                      \exp\bigl(O(m\log(ms+1))\bigr).     \tag{B.MB.2}
\]
The equations of these charts have degree polynomial in $m,s$.
They can include all minors required for exact-rank conditions;
no bound on the number of those minors is asserted or needed.
\item On these charts one can choose actual representative,
projection, and homotopy maps for the two-affine line Cech
complexes, compatibly through all higher pole ranges. Through
any prescribed pole range polynomial in $m,s$, all their
coefficients are rational functions of degree polynomial in
$m,s$. This construction introduces no new charts at each
successive pole order.
\end{enumerate}
All implicit polynomials and constants depend only on $c,h_0$,
the bounded tensor-word pattern, and the specified polynomial
pole cutoff. Fixed curve coefficients and fixed line twists are
background constants. For a chain of $f$ phases with a bounded
number of moving lines per phase and only bounded-range
interactions, the ambient background coordinates number
$O(fms)$ and the chart count is
$\exp(O(fm\log(ms+1)))$.

The rank charts cover geometric points. An arbitrary noncoflat
Artinian family is not asserted to factor through one exact-rank
chart. All-ring splitting statements apply to points of the
stated chart schemes. This distinction does not affect the
Artinian surjectivity of the preceding ordered-divisor charts.
\end{proposition}

\begin{proof}
Put $t=\deg x$ and $H=x^{-1}(\infty)$; choose the coordinate
so that $H'=x^{-1}(0)$ is disjoint from $H$. These divisors
need not be reduced. Both have degree $t$.

\emph{A plane chart of controlled degree.}
Choose the least positive $a$ with $at\geq2g(X)+1$. The complete linear series
of $aH$ is very ample. A general $z\in H^0(X,\mathcal O(aH))$
is primitive for the separable extension $k(X)/k(x)$.
Indeed, this linear series generates $k(X)$, and the finitely
many proper intermediate fields of a finite separable extension
cannot contain a general member of its $k$-linear section space.
We may also require the allowed pole order at every point of
$H$. Its minimal polynomial is
\[
 F(x,z)=z^t+c_1(x)z^{t-1}+\cdots+c_t(x),
               \qquad c_i(x)\in k[x],\quad\deg c_i\leq ai.
                                                               \tag{B.MB.3}
\]
Integrality away from infinity gives polynomial coefficients;
their valuations at infinity give the displayed degree bounds.
Since $at\leq2g(X)+1+t$, both degrees of $F$ are $O(m)$.
Its discriminant $\Delta(x)$ is nonzero and has degree $O(m^2)$.
After removing its zeroes, $x=0$, and finitely many other fixed
fibres when required below, the normalization agrees with
\[
             k[x,z,1/(x\Delta)]/(F).                    \tag{B.MB.4}
\]
In fact this algebra is finite etale over the indicated open
of the $x$-line, hence normal, and has the same function field
as $X$. Denote the resulting open curve by $X^\circ$.

\emph{Ordered divisors, including Artinian lifting.}
Choose $E_0=eH$ so that, for every prescribed degree $\ell$,
\[
          d=\ell+\deg E_0\geq2g(X)+1,
                  \qquad d=O(ms).                       \tag{B.MB.5}
\]
Every line $L$ of degree $\ell$ has a presentation
$L=\mathcal O_X(D-E_0)$ with
$D=P_1+\cdots+P_d$, where the $P_i$ belong to $X^\circ$
and their $x$-coordinates are pairwise distinct. To verify
this uniformly, use the very ample line $L(E_0)$. Avoiding
the finite omitted set is open on its section space. A double
zero, or two zeroes at different points with the same
$x$-coordinate, imposes two independent linear conditions
while those points range over a one-dimensional family.
The corresponding bad incidences have smaller dimension.
A general section therefore has the asserted zero divisor.

Each $P_i$ is given by two coordinates $(a_i,b_i)$ satisfying
$F(a_i,b_i)=0$. The conditions
$\Delta(a_i)a_i\ne0$ and $a_i\ne a_j$ are open. This is an
$O(ms)$-coordinate parameter space. For a finite local Artinian
base $T$, vanishing of $H^1(L(E_0))$ on the closed fibre lifts
a chosen section successively through all small extensions.
The lifted zero divisor is finite etale over $T$: at each
closed zero the relative derivative is a unit. Since the
residue field is algebraically closed, it splits uniquely
into sections after the closed zeroes have been ordered.
They retain all the open conditions. The regular section
identifies $L_T(E_0)$ with $\mathcal O(D_T)$. This proves the
stated Artinian lifting assertion directly, without requiring
an unproved surjectivity for cohomology-rank strata.

\emph{Polynomial linear equations for sections.}
Let $h_D(x)=\prod_i(x-a_i)$. Its divisor is
$D+D^{\mathrm c}-dH$, where $D^{\mathrm c}$ is the complement
of the selected points in the corresponding unramified fibres.
Multiplication by $h_D$ identifies sections of $L(qH)$ with
the subspace of
\[
 H^0\bigl(X,\mathcal O((d+q)H-E_0)\bigr)                 \tag{B.MB.6}
\]
that vanishes on $D^{\mathrm c}$. The ambient space is fixed;
for all pole ranges used here its dimension is polynomial
in $m,s$, typically $O(m^2s)$. In the finite etale fibre algebra
$T[z]/F(a_i,z)$, complementary-fibre vanishing is exactly
\[
                  (z-b_i)v(a_i,z)=0.                     \tag{B.MB.7}
\]
Indeed the simple root $b_i$ defines a direct factor $T$,
and the annihilator of $z-b_i$ is precisely that factor.
Thus (B.MB.7) is exact over nonreduced coefficient rings too.
Negative moving points impose evaluation-zero conditions.
For bounded multiplicities, use the corresponding finite
jets: the simple-root factor and its idempotent lift uniquely
modulo $(x-a_i)^c$. Successive coefficient comparison divides
only by the unit $F_z(a_i,b_i)$, and gives rational formulas
of polynomial degree for each bounded $c$.

We give a degree check for the evaluations, rather than assume
that arbitrary regular functions have low-complexity expressions.
The trace-dual basis to $1,z,\ldots,z^{t-1}$ has the form
$Q_j(z)/F_z(z)$, with $\deg_zQ_j<t$ and coefficients drawn
from those of $F$. If a function $v$ has poles only at $H$,
of order at most $qH$, then
\[
 v=\frac{\sum_{j=0}^{t-1}v_j(x)z^j}{F_z(x,z)},
                  \deg v_j\leq q+O(m).                  \tag{B.MB.8}
\]
For clarity, $F_zv$ belongs to $k[x,z]$: trace integrality
puts every element of the normalization in the trace dual
$F_z^{-1}k[x,z]$. The coefficient formula is
$v_j=\operatorname{Tr}(vQ_j)$, and the pole bound gives its
degree. This remains valid with inseparable residue at branch
fibres; only separability of $k(X)/k(x)$ is used.
Fixed finite poles can first be cleared by a fixed polynomial
in $x$. The total degree needed is polynomial in the degree
of their fixed divisor. Evaluating (B.MB.8), or reducing it
modulo $F(a_i,z)$, uses only inverses of $F_z$ or its resultant,
which are units on the chosen open. Adjugate formulas give
polynomial-size matrices of polynomial rational degree.

The same construction applies to any bounded tensor word in
the moving divisors and fixed twists. Representatives of the
fixed twists may be chosen with positive and negative divisor
degrees $O(ms+m)$ by Riemann--Roch. Choose all moving points
away from their finitely many fibres. Multiplication by the
bounded powers of the $h_D$ clears the moving poles; selected
and complementary fibre jets impose all remaining conditions.
One-step Frobenius replaces a moving divisor by $pD$ and raises
coefficient coordinates to their $p$th powers, with $p$ fixed.
A deck translation preserves $x$; its effect on $z$ has the
same pole bound as $z$ and hence a representation (B.MB.8)
of polynomial degree. Thus the listed operations retain the
claimed bounds.

To compute an actual cohomology complex, choose a sufficiently
positive fixed pole divisor so that the ambient line has
vanishing $H^1$, and use its exact section sequence with the
finite divisor of imposed zeroes. This is a two-term complex
of finite free modules before taking cohomology. Alternatively
one may use the finite Cech sequence below. All evaluation
and principal-part maps are actual sheaf maps, so their
matrix models commute with arbitrary coefficient change.
Riemann--Roch and duality bound the total cohomology of each
line of degree $O(ms)$ by $O(ms)$
\citep[Tags 0BS2 and 0BS6]{StacksBNegativeCurve}.

\emph{Rank charts and their number.}
Every raw matrix just constructed has dimensions and entry
degrees bounded by one polynomial in $m,s$. Choose its rank,
impose all next-size minors, and apply
Lemma~\ref{D:compressed-minor-background-charts} to cover
that exact-rank locus by at most $b+1$ rational-splitting
charts, where $b=O(ms)$ is the ambient parameter dimension.
A bounded number of such reductions handles a single
two-term complex and its prescribed evaluation maps.
Each line is processed independently from its raw evaluation
matrices; there is no chain of $O(m)$ successive matrix
inversions feeding into later matrices.
Thus every resulting rational formula still has polynomial
degree. There are polynomially many choices of rank and
$O(ms)$ compression charts for each of $O(m)$ matrices,
which proves (B.MB.2). All defining minors may be retained:
their degrees are polynomial even though their number can
be large. Formula (B.MB.1) shows that the splittings are exact
over every ring satisfying the chart equations.

\emph{Compatible finite-pole Cech homotopies.}
Use $X=(X\setminus H)\cup(X\setminus H')$. For one of the
line families $M$, choose $b_0$ so that
\[
 H^1\bigl(M((b_0-1)H)\bigr)
   =H^1\bigl(M((b_0-1)H')\bigr)=0                       \tag{B.MB.9}
\]
on every fibre. It suffices that both degrees exceed
$2g(X)-2$, so $b_0$ is polynomial in $m,s$.
For $b\geq b_0$, the finite section complex
\[
 C_b^0=H^0(M(bH))\oplus H^0(M(bH')),
       \qquad C_b^1=H^0(M(bH+bH'))                       \tag{B.MB.10}
\]
computes $R\Gamma(M)$. This follows from the exact sheaf
sequence
$0\to M\to M(bH)\oplus M(bH')\to M(bH+bH')\to0$,
since $H,H'$ are disjoint and the two first cohomology groups
vanish. The high-degree section modules and evaluations
are free with base change on the chosen rational charts.

Split $C_{b_0}$ once, by the preceding matrix construction,
onto its free cohomology. Also choose once lifts
$\alpha_1,\ldots,\alpha_t\in H^0(M(b_0H))$ of a basis of
the principal parts at the last $H$-layer, and analogous
$\beta_1,\ldots,\beta_t\in H^0(M(b_0H'))$ at $H'$.
Their existence and coefficient-base change follow from
(B.MB.9). The relevant line on the fixed fibres is independent
of the moving divisor, which is disjoint from those fibres;
one may fix its $k$-basis in advance.

For $b>b_0$, the new principal-part layers are lifted by
\[
          x^{b-b_0}\alpha_i,
                         \qquad x^{-(b-b_0)}\beta_i.     \tag{B.MB.11}
\]
The first family has poles only at $H$, and the second only
at $H'$. In the quotient $C_b/C_{b-1}$ these give matching
bases in degrees zero and one, with the Cech differential
an isomorphism. Hence the complex is the direct sum of
$C_{b_0}$ and compatible contractible pairs, one set for
each new pole order. Define the homotopy on each pair by
that inverse differential, and set its projection to
cohomology equal to zero. The result is a strong deformation
retract of the full two-affine Cech complex, compatible with
every coefficient change on the chart. It satisfies
\[
 \rho\iota=1,\quad \iota\rho=1-dh-hd,\quad
                 h^2=h\iota=\rho h=0.                   \tag{B.MB.12}
\]
It does not increase the pole range above
$\max(b,b_0)$ for an input in $C_b$.

The formulas through a given pole bound are effective rational
formulas in the original point coordinates. Expand at the
two fixed fibres and eliminate the highest principal part,
using (B.MB.11), then repeat. The initial lift matrices are
fixed rational matrices on the chart. Each elimination is
linear in the current coefficients and uses only these fixed
matrices. Thus the degrees grow at most linearly in the
number of pole steps, times their initial degree. The jets
of $h_D^{\pm1}$ have polynomial degree through any polynomial
cutoff: at infinity use
$h_D/x^d=\prod_i(1-a_i/x)$, and at zero use the invertible
constant $h_D(0)$. The same statement holds for bounded powers
and the fixed twists. Consequently products of polynomially
bounded expanded word length, with a bounded number of one-step
Frobenius operations, have polynomial rational coefficient degree
through a polynomial pole cutoff. In the normal-layer application
the word length is at most $am+O(1)$ and there is one fixed-$p$
pullback. This does not assert the same bound for arbitrary
repeated-squaring circuits.
This proves the final assertion, including actual homotopy
data rather than only dimensions or residue-field bases.

The Cech computation and its coefficient description use
the exact affine-intersection and base-change statements in
\citep[Tags 0BDX, 01XD, 02KE]{C_StacksFiniteCechHom}; the
rational matrices, compression argument, and pole-compatible
splittings above are constructed explicitly.
\end{proof}

\begin{remark}[Reference lifts and the limits of the background assertion]
In the one-curve formal surface setting with finite etale
spectral covers, the ordered moving divisor also has an actual
formal lift. Shrink one affine open of the original smooth
curve so that its normal line is trivial. Formal smoothness
and the affine square-zero lifting criterion give an
isomorphism of its formal neighbourhood with the product
with a formal disc. A finite etale spectral cover pulls back
to the corresponding product. Choose all moving divisor
points in this open, as permitted in the proof, and take their
constant transverse Cartier sections. Their divisor lines
glue with the trivial line away from their support; subtract
a fixed formal lift of $E_0$. This produces an actual formal
reference line with the same background point parameters.

On a retained finite normal order, its moving transition
factors are the products of the local equations of these
sections and their inverses. Taylor expansion and inversion
use the same unit denominators as above; through order $am$
their degree in the moving point coordinates is polynomial
in $m,s$. Fixed formal-neighbourhood coefficients remain
background constants. This observation does not assert a
normal-pole bound for arbitrary additional anchor data:
the finite cochain ranges needed for those data must still
be established in the formal object construction.

The proposition provides rational backgrounds and coflat
rank charts. It does not construct the full graded formal
bundle atlas, prove rigidity of its unit objects, or bound
representations. In particular no rank-stratum statement
here repairs the failed open-fullness inference recorded
in the separate genus-two research note.
\end{remark}


% Separate post-CP19 proved source. Coefficient-gerbe transfer and unit-Isom;
% no moving-Picard atlas or unbounded-period representation count is asserted.
\subsection{The actual coefficient gerbe and finite jets of unit isomorphisms}

\begin{proposition}[Proved: the geometric-period root supports the coefficient gerbe]
\label{D:one-curve-coefficient-root-gerbe}
Use the complete local one-smooth-curve setting of
Proposition~\ref{D:single-curve-Brauer-residue-cover}:
$U=\operatorname{Spec}R-\{\mathfrak m\}$,
$Y\to\operatorname{Spec}R$ is a projective regular resolution,
$C$ is its single smooth exceptional curve, $g=g(C)$ and
$C^2=-d<0$. Let $n$ be prime to $p$, and let
$\alpha\in H^2(U,\mu_n)$. Write $\beta$ for its geometric
Brauer class and $m$ for the order of $\beta$, taking $m=1$
when $\beta=0$. Then $m\mid n$, and the actual coefficient
class $\alpha$ extends to a class
\[
 \overline\alpha\in H^2(\mathcal X,\mu_n),\qquad
             \mathcal X=\sqrt[m]{(Y,C)}.                \tag{D.UJ.1}
\]
Let $\mathcal G\to\mathcal X$ be a gerbe for this class and
$\mathcal D$ the reduced root divisor. The scalar-weight-one
spectral torsor over $C$ for $\mathcal G_{\mathcal D}$ is
connected of degree exactly $m$. Its class is the residue of
$\beta$. In particular every scalar-weight-one vector bundle
on $\mathcal G_U$ has rank divisible by $m$.

There is an algebraic scalar-weight-one anchor $A$ of rank $m$
on $\mathcal G$, with all $m$ root characters once at each
geometric point of $C$. Every scalar-weight-one bundle on
$\mathcal G_U$ admits a vector-bundle lattice on $\mathcal G$.
Every formal scalar-weight-one vector bundle along
$\mathcal G_{\mathcal D}$ algebraizes, as do maps with any
fixed divisor of poles. These algebraization assertions also
hold after a finite local Artinian coefficient-base extension.
The anchor is used for existence and descent, not as a rank-$mr$
coordinate presentation.

Finally
\[
 |H^2(U,\mu_n)|=\gcd(n,d)\,n^{2g}.                     \tag{D.UJ.2}
\]
Consequently for $q=p^f$ and $n=q-1$ the number of coefficient
classes of projective $\mathbb F_q$-representations is at most
$dq^{2g}$. This includes their Kummer part, not only their
geometric Brauer classes.
\end{proposition}

\begin{proof}
The root construction extends $\beta$ uniquely to a Brauer
class $\overline\beta$ of order $m$ on $\mathcal X$.
This is exactly the arbitrary prime-to-characteristic period
statement of \citet[Proposition~3.2.1 and complete proof,
pp.~10--11]{Lieblich2011ResidualD}. Its hypotheses are a
regular excellent surface, a normal crossings divisor, and
a primitive $m$th root of unity; all hold here. The later
prime-order convention in that paper's Section~5 is not a
hypothesis of Proposition~3.2.1.

Since $n\overline\beta=0$, Kummer supplies a lift
$\alpha_0\in H^2(\mathcal X,\mu_n)$. The difference
$\alpha-\alpha_0|_U$ is the Kummer class of a line bundle on
$U$. Every such line bundle extends to $Y$: close its Cartier
divisor as a Weil divisor on the regular surface $Y$.
Adding the Kummer class of this extension to $\alpha_0$
proves (D.UJ.1). Thus the construction extends the specified
coefficient gerbe, rather than replacing it by another gerbe
with the same Brauer class.

The edge character of $\overline\beta$ on the residual
root gerbe is its ordinary residue. One can check the
identification at the generic point of $C$: the proof of
Proposition~3.2.1 uses the cyclic algebra $(a,t)_m$, replaces
its generator $y$ by $y/t^{1/m}$, and the resulting residual
character is the torsor $a^{1/m}$. This is precisely the edge
map in \citet[Proposition~4.1.4 and complete proof,
p.~11]{Lieblich2011ResidualD}. The equality at the generic
point determines the finite etale torsor on the normal curve.
The mixed-band character construction in
Lemma~\ref{D:twisted-connected-anchor-minimal-slopes}
identifies that edge character with the scalar-weight-one
spectral torsor. Adding a Kummer line class does not change
this Brauer edge character. By
Proposition~\ref{D:single-curve-Brauer-residue-cover}, its
order is $m$, so the torsor is connected and has degree $m$.

Let $h:Y'\to Y$ be its cyclic finite etale lift. The pullback
of $\beta$ to $U'=h^{-1}U$ is zero by that proposition.
Put $\mathcal X'=\mathcal X\times_Y Y'$. The restriction
$\operatorname{Br}(\mathcal X')\to\operatorname{Br}(U')$
is injective, by the uniqueness assertion and proof of the
same root-construction proposition. Thus
$\overline\beta|_{\mathcal X'}=0$. The scalar gerbe
$\mathcal G'=\mathcal G\times_Y Y'$ consequently has a
weight-one invertible sheaf $\Lambda$. On the residual
curve, its root character is a section of the pulled-back
spectral torsor. Translating by a power of the root line
selects the diagonal section. Tensoring by an ordinary
line bundle on $Y'$ then makes its residual spectral line
$\mathcal O_{C'}$: the required residual line lifts to
$Y'$ by the Picard lifting argument already proved in
Lemma~\ref{D:local-Picard-prime-to-p-generators}.
The finite pushforward $A=h_*\Lambda$ has rank $m$ and
contains all root characters once.

Here are explicit algebraization and extension details.
Write $q:\mathcal G\to Y$ and
$\mathcal E=q_*\mathcal End(A)$. Exact scalar descent and
root invariants make $\mathcal E$ a coherent algebra,
locally free over $\mathcal O_Y$. Locally on a root chart,
$A$ is the sum of the shifted free modules for all root
characters in the scalar-weight-one sector. This sum is a
projective generator of that graded module category.
The two functors
\[
 F\longmapsto q_*\mathcal Hom(A,F),\qquad
 M\longmapsto q^*M\otimes_{q^*\mathcal E}A              \tag{D.UJ.3}
\]
are therefore inverse on the corresponding coherent modules;
$\mathcal Hom(A,F)$ is a right $\mathcal E$-module by
precomposition, and $A$ is a left $q^*\mathcal E$-module.
This assertion is checked on the graded chart and descends.
For a vector bundle $F$, the first sheaf is an ordinary
vector bundle on $Y$: each invariant shifted free root
summand is free on the coarse chart.

Given $F_U$, first extend its ordinary module in (D.UJ.3)
to a coherent module on $Y$, and take its span under the
finite algebra $\mathcal E$ to retain the action. The
second functor gives a coherent extension on $\mathcal G$.
Its reflexive hull is locally free, since every etale chart
of $\mathcal G$ is regular of dimension two. It retains
scalar weight one and restricts to $F_U$.
The connected spectral category applied to that lattice
now gives $\operatorname{rk}F_U=m\operatorname{rk}W$.

For a formal vector bundle, apply the first functor in
(D.UJ.3). Ordinary projective Grothendieck existence
algebraizes the resulting ordinary vector bundle together
with its finite-algebra action. The action identities
algebraize by full faithfulness. Applying the inverse
functor recovers a coherent sheaf whose completion is the
original formal bundle. It is locally free everywhere:
the non-locally-free locus is closed and proper over the
local base, so, if nonempty, it would meet the closed fibre,
where local freeness follows from its completion. Maps
algebraize by the same full faithfulness, after their fixed
Cartier twists. This reasoning applies to every scalar
weight in the Frobenius orbit and to finite Artinian
coefficient bases, since $R\otimes_k T$ is then complete
and the anchor and algebra remain locally free. The
ordinary existence theorem used here is
\citep[Theorem~30.27.1, Tag~088E]{StacksTwistedDepthD}.
No assertion about an Azumaya coarse algebra on the boundary
is made: $\mathcal E$ is used with its actual root-chart
module category.

For the coefficient count, the earlier Picard and proper
base-change calculation gives
$\operatorname{Pic}(Y)/n=H^2(Y,\mu_n)=\mathbb Z/n$,
where the class of a line bundle is the degree of its
restriction to $C$. The kernel of
$\operatorname{Pic}(Y)\to\operatorname{Pic}(U)$ is generated
by $\mathcal O_Y(C)$, whose degree on $C$ is $-d$.
Hence
\[
       |\operatorname{Pic}(U)/n|=\gcd(n,d).
\]
Kummer and the already proved residue isomorphism
$\operatorname{Br}'(U)[n]=H^1(C,\mathbb Z/n)$ give (D.UJ.2).
The injection from continuous group cohomology to etale
cohomology in Theorem~\ref{D:local-surface-constant-H2}
then gives the asserted bound for coefficient classes of
representations.
\end{proof}

\begin{lemma}[Proved: a unit isomorphism is determined by a uniform finite jet]
\label{D:meromorphic-unit-jet-contraction}
Let $\widehat{\mathcal G}$ be the formal root gerbe above,
$I=\mathcal O(-\mathcal D)$, and let the coefficient algebra
have characteristic $p$. Consider two length-$f$ cycles of
formal vector bundles, with actual meromorphic isomorphisms
\[
 \Phi_a:F^*F_a\xrightarrow{\sim}F_{a+1},\qquad
 \Psi_a:F^*G_a\xrightarrow{\sim}G_{a+1},\quad a\bmod f.
                                                               \tag{D.UJ.4}
\]
Here $F^*$ is absolute $p$-Frobenius, with the indicated
scalar weights. Suppose these isomorphisms and their
inverses have root-divisor pole order at most $b$.
For a meromorphic map, ``order at least $t$'' means membership
in $I^t$ times its unshifted Hom bundle; $t$ may be negative.

Fix an integer $M>2b/(p-1)$. Suppose maps $\eta_a:F_a\to G_a$
satisfy
\[
 \Psi_a F^*(\eta_a)\Phi_a^{-1}-\eta_{a+1}
                         \quad\hbox{has order at least }M.
                                                               \tag{D.UJ.5}
\]
There is a unique exact unit-compatible collection
$\eta_a^\infty$ congruent to $\eta_a$ to order $M$.
The conclusion holds over arbitrary coefficient rings.
In particular ordinary maps need not be determined by a
finite jet; this assertion concerns the maps satisfying
all the unit equations.

Suppose also that maps $\zeta_a:G_a\to F_a$ satisfy the
reverse unit equations to order $M$, that all $\eta_a,\zeta_a$
have poles at most $c$, and that both inverse products
are congruent to the identity to order $M-c$. If
\[
                       M-c>2b/(p-1),                  \tag{D.UJ.6}
\]
then the corrected collections are mutually inverse.
Thus simultaneous finite jets of a map and its inverse
recover exactly an actual meromorphic unit isomorphism.
\end{lemma}

\begin{proof}
Set
$T_a(\eta)=\Psi_aF^*(\eta)\Phi_a^{-1}$.
The maps $T_a$ are additive. For any difference of order at
least $t$, their output difference has order at least
\[
                              pt-2b.                    \tag{D.UJ.7}
\]
This follows from $F^*I=I^p$ and the two pole bounds; it
uses no flatness or reducedness of the coefficient algebra.
If $t\geq M$, (D.UJ.7) is at least $t$, and its iterates
tend to infinity. Around the cycle the operator
$T=T_{f-1}\cdots T_0$ satisfies
\[
 \operatorname{ord}(T\eta-T\eta')
       \geq q\operatorname{ord}(\eta-\eta')
                         -2b\frac{q-1}{p-1}.           \tag{D.UJ.8}
\]
The approximate compatibilities imply
$\operatorname{ord}(T\eta_0-\eta_0)\geq M$.
Successive iterates of $T$ are consequently a Cauchy
sequence in the complete formal Hom module, with every
increment of order at least $M$. Their limit is fixed by
$T$, and the other stages are obtained by the successive
$T_a$. They retain the prescribed order-$M$ jets. If two
such fixed collections agree to that order, applying
(D.UJ.8) repeatedly to their difference makes its order
arbitrarily large, so it is zero by separatedness.
All iterates have the same original pole bound: their
differences from the starting maps have positive order.

Apply this construction also to $\zeta$. Composition of
maps with pole bound $c$ loses at most $c$ orders from a
changed factor. Thus the corrected inverse products still
agree with the identity to order $M-c$. They are exact
unit endomorphisms. The identity is another such
endomorphism. By (D.UJ.6) and the uniqueness just proved,
both products equal the identity exactly. This also shows
that the corrected maps retain inverse pole bound $c$.
\end{proof}

\begin{proposition}[Proved: finite-presentation unit-Isom and open unit classes]
\label{D:finite-jet-unit-Isom}
Fix a coefficient gerbe as in
Proposition~\ref{D:one-curve-coefficient-root-gerbe},
$n=q-1$, and the one-curve slope windows of
Proposition~\ref{D:twisted-period-independent-normal-depth}.
Let $S$ be an affine scheme of finite type over $k$
parameterizing actual bundle and bounded-pole unit-cycle
jets on a sufficiently deep fixed ordinary-$C$ thickening.
Require the following precise effectivity conditions:
\begin{enumerate}
\item the finite bundles are locally free over every
coefficient base, and their geometric residual fibres lie
in the stated slope windows;
\item the maps and inverse diagrams are actual maps of
these finite bundles, with inverse products equal to the
canonical Cartier inclusions to the retained precision;
\item the diagram has been retained deeply enough for the
uniform bundle and specified-Hom lifting of
Proposition~\ref{D:twisted-period-independent-normal-depth},
including the fixed twists used in the comparisons below.
\end{enumerate}
These requirements concern actual objects and maps, not
a tangent-space dimension or a derived class without its
boundary witnesses.

Then the functor of meromorphic unit isomorphisms between
two such families is represented by an affine scheme of
finite presentation over their common base. For a fixed
geometric unit object, its Isom scheme over $S$ is etale.
Consequently each nonempty isomorphism-class fibre in $S(k)$
is the $k$-point set of an open and closed subscheme, and
\[
 \#\{\hbox{unit classes represented on }S\}
                            \leq\#\operatorname{Irr}(S).
                                                               \tag{D.UJ.9}
\]
The same conclusion holds after restricting to an open
residual subbundle or slope-window locus. No bound for the
number or degrees of moving-determinant charts is asserted.
\end{proposition}

\begin{proof}
We first spell out the all-base meaning of the lifting
condition. On each conormal layer the relevant Hom sheaf
is a vector bundle on the fixed proper spectral curve,
with the uniform slope estimates already proved. Its first
cohomology vanishes beyond the fixed depth on every
geometric coefficient fibre. These sheaves are flat over
the coefficient base. The perfect cohomology complex, compatible with arbitrary
base change by \citep[Lemma~30.22.1, Tag~07VK]
{StacksUnitJetTransferD}, has a local two-term finite free
representative in degrees zero and one. Its degree-one
differential is surjective locally on the base, by Nakayama. That
surjection splits locally, so its degree-zero kernel is
flat and the vanishing and lifting commute with every
base change. Equivalently the finite Cech complexes in
Proposition~\ref{D:fixed-determinant-finite-Hom} give this
calculation on the charts in use. The uniform pole bound also
holds after arbitrary coefficient base change. Indeed the
negative-degree normal Hom layers have zero $H^0$ on every
geometric fibre. In their two-term free cohomology complexes
the degree-zero differential is fibrewise injective, hence
locally a split injection. Thus $H^0$ remains zero after
every base change. The successive pole-layer exact sequences
then show that no meromorphic map, including one with
nilpotent coefficients, acquires poles beyond the fixed
bound. There is no inference from fibrewise tangent dimensions
to an object atlas.

The second-cohomology obstruction to extending a bundle
through a square-zero normal layer vanishes on the proper
curve over an affine base. We may therefore choose formal
extensions over $S$ of the given actual finite bundles.
The uniform first-cohomology vanishing lifts their specified
maps. An inverse product, after dividing by its fixed pole,
is $1$ plus positive normal order, so its geometric-series
inverse gives actual meromorphic unit arrows. These choices
are used only to define a comparison functor; no quantitative
coordinate cost for their unspecified higher coefficients
is charged to the object presentation.

Let $b=mB$ be a common root pole bound for the one-step
arrows and their inverses, and let $c=mC_0$ be the uniform
root pole bound for meromorphic comparisons between the
permitted lattices. The same bound applies to their inverse
comparisons. Choose $M=mM_0$ with $M_0$ fixed and so large
that specified maps lift at this depth and
\[
                    M-c>2b/(p-1).                     \tag{D.UJ.10}
\]
A comparison variable is a section of the finite sheaf
\[
 \mathcal J_{a}(M,c)=
 \mathcal Hom(F_a,G_a)(c\mathcal D)\big/
                         I^M\mathcal Hom(F_a,G_a),     \tag{D.UJ.11}
\]
and there is a reverse comparison variable in the analogous
sheaf. These are finite bundles on the corresponding
proper finite root thickening; their coarse scalar-zero
pushforwards are coherent and flat over the base. Proper
finite-cycle Hom schemes are affine of finite presentation
\citep[Proposition~99.3.10, Tag~08K6]{StacksUnitJetTransferD}.
Exact root invariants give the same statement on the tame
stack; it can also be seen by retaining the equivariance
equations on a finite root chart. Thus the tuples of the
variables in (D.UJ.11) form an affine scheme of finite
presentation.

Impose the two sets of unit equations (D.UJ.5) modulo true
normal order $M$, and impose the two inverse products
modulo true normal order $M-c$. They are well-defined on
these jets. For the unit equations this is (D.UJ.7), since
$pM-2b\geq M$; for the products it is precisely the loss
of at most $c$ on multiplying by a pole-$c$ map. All the
equations are polynomial maps between finite proper-cycle
Hom schemes. To compute them one only needs finitely many
jets of the chosen unit arrows: for example retain their
coefficients to order greater than $M+pc+3b$. Their true
inverse jets are obtained by the finite truncation of the
geometric series just used. Increasing the original
ordinary depth once supplies all of these jets, independently
of $m,n,r,f$.

Any solution lifts to ordinary formal maps with its specified
jets, by the Hom lifting property. Lemma
\ref{D:meromorphic-unit-jet-contraction} corrects these
maps uniquely, preserving their full order-$M$ jets, to
mutually inverse unit maps. Conversely every meromorphic
unit isomorphism has the pole bound $c$ and produces just
such a solution. The construction is functorial in the
coefficient algebra. We have therefore represented the
\emph{actual} unit-Isom functor by a finite-presentation
closed subscheme, rather than by a space of approximate
maps with uncontrolled lifting choices. Changing the
chosen formal extensions or higher arrow lifts gives the
same functor: compare two choices by a lift of their common
sufficiently deep identity jet and apply the same contraction.

It remains to prove openness, not just infinitesimal
rigidity. For a finite local Artinian coefficient algebra
$T$, Proposition~\ref{D:one-curve-coefficient-root-gerbe}
algebraizes the formal bundles and their pole maps over
$R\otimes_kT$. On the puncture a length-$f$ unit cycle is
equivalent to its first bundle with a $q$-Frobenius unit
structure. Its fixed-frame equations give a finite etale
$\operatorname{GL}_r(\mathbb F_q)$-torsor. This is the
usual Lang-frame construction already used in the local
Brauer-zero count; on an etale atlas the derivative of
the additive fixed-frame equations is invertible, and the
invertible frames form the finite frame torsor. The same
construction descends on the coefficient gerbe, with its
scalar weight retained.

Finite etale covers, including all their morphisms, are
invariant under nilpotent thickenings
\citep[Lemma~58.8.3 and complete proof, Tag~0BQB]
{StacksUnitJetTransferD}; descent on an etale atlas gives
the same assertion for the gerbe. Thus every prescribed
unit isomorphism over a quotient of a finite local Artinian
coefficient algebra lifts uniquely to any nilpotent
extension. Its extension to the bounded-pole jet model is
the unique one constructed above. Since that model is of
finite presentation, the infinitesimal criterion now
applies: the fixed-object Isom morphism is formally etale
and hence etale. One may apply the local Artinian criterion
first at the closed points of the finite-type base and
then use the open etale locus. We use the complete
criteria in \citep[Tags~02HT, 02HM and 00UR]
{StacksUnitJetTransferD}. This argument uses every finite
Artinian extension, not only dual numbers.

The image of an etale morphism is open, so each fixed
unit class has an open image in $S$. These images are
pairwise disjoint: an intersection has a closed geometric
point, and hence supplies an isomorphism of the two fixed
objects. Their union contains every closed point and is
therefore all of the finite-type scheme. The complement
of any one image is the union of the others, so each is
also closed. Every nonempty class image contains an
irreducible component, proving (D.UJ.9). Restricting to
an open locus does not affect the argument.
\end{proof}

\begin{remark}[The quantitative conclusion still requires an object atlas]
The coefficient-class bound, root-gerbe extension and finite
unit-Isom comparison remove distinct transfer ambiguities.
They do not construct charts for the moving spectral
determinant. A uniform unbounded-period representation count
still requires an actual all-base polynomial object
presentation with the weighted size and degree bounds; the
weighted component estimates apply only after that further
construction has been proved.
\end{remark}


% Separate post-CP19 proved source. Controlled reference data, not a count by itself.
\subsection{Controlled references for the actual one-curve coefficient gerbe}

\begin{lemma}[Proved: controlled one-curve gerbe anchors]
\label{D:controlled-one-curve-gerbe-anchors}
Use the setting of Proposition~\ref{D:one-curve-coefficient-root-gerbe}.
Fix an ordinary-$C$ depth $a$, and fix the characteristic $p$.
Let $n=p^f-1$, let $\alpha\in H^2(U,\mu_n)$, and let $m\mid n$
be its geometric Brauer period. Put
$\mathcal X=\sqrt[m]{(Y,C)}$ and let $C'\to C$ be the connected
degree-$m$ residue cover; $m=1$ is allowed.

The actual coefficient gerbe, its scalar-weight anchors at the
$f$ Frobenius stages, and their descent and comparison maps can
be chosen with the following properties.
\begin{enumerate}
\item All reference lines on $C'$ that enter a single root
layer of a termwise Hom complex, or a one-step Frobenius
comparison, have degree bounded in absolute value by $c m$.
Here $c$ depends only on $(Y,C),a,p$ and the fixed finite
diagram pattern. Root characters are reduced modulo $m$.
\item Through root depth $am$, the scalar-zero transition,
descent, and comparison cochains have rational representatives
whose residual pole divisors and plane-coordinate degrees are
bounded by a polynomial in $m$. This assertion concerns chosen
references, not arbitrary formal trivializations.
\item On adding the moving determinant lines of
Proposition~\ref{D:moving-line-polynomial-background}, the
reference data meet that proposition's polynomial background
hypotheses. For diagrams of bounded adjacent-phase range, the
polynomial is independent of $f$ and of $n$. The resulting
identities hold over arbitrary coefficient algebras on the
specified rational charts. No new moduli coordinates are needed
for the fixed reference cochains.
\end{enumerate}
In particular the normal-layer cohomology budget remains
$O(ms^2)$ per layer for spectral rank $s$, even though its
rational coefficient complexity may be a larger polynomial
in $m,s$. The lemma does not assert a representation count
without the effective bundle and unit-Isom presentations.
\end{lemma}

\begin{proof}
We distinguish the degree of a line from the pole degree of a
rational presentation. The former controls cohomology; the
latter controls the coefficients of the polynomial diagrams.
Constants below may increase with the fixed ordinary depth.

\emph{Separate the scalar order from the geometric period.}
Choose a line $H$ on $Y$ with $\deg(H|_C)=1$, using the
Picard lifting in Lemma~\ref{D:local-Picard-prime-to-p-generators}.
Choose $\gamma\in H^2(U,\mu_m)$ mapping to the Brauer image
$\beta$ of $\alpha$. In the Kummer diagram for $m\mid n$,
the map on the middle $\mathbf G_m$ is the identity and the
map on the right $\mathbf G_m$ is the $(n/m)$th power.
Thus pushout along $\mu_m\hookrightarrow\mu_n$ preserves
the Brauer image. The Picard calculation in
Proposition~\ref{D:one-curve-coefficient-root-gerbe} gives
\[
 \alpha=\iota_*\gamma+\delta_n(H^e|_U),\qquad
       0\leq e<\gcd(n,d)\leq d.                         \tag{D.CA.1}
\]
Here and below an empty choice of a nonzero character when
$m=1$ means the identity cover. Formula (D.CA.1) retains
the actual coefficient class. It does not replace it by its
Brauer image. Kummer exactness is the usual one for invertible
orders \citep[Lemma 59.28.1]{StacksControlledAnchorsD}.

Extend $\gamma$ to the $m$th root stack, as in the cited
coefficient-gerbe proposition. Represent the second summand
by the $n$th-root gerbe of $H^e$. If $T$ is its tautological
line, put $u_i$ equal to the residue of $p^i$ in $[0,n-1]$.
Then
\[
 (T^{u_i})^p=T^{u_{i+1}}\otimes H^{e c_i},\qquad
 c_i=\left\lfloor\frac{p u_i}{n}\right\rfloor,
                         0\leq c_i\leq p-1.             \tag{D.CA.2}
\]
There is no ordinary line of degree proportional to $n$ in
this comparison. Scalar-zero Hom cancels the tautological
line; the only ordinary one-step contribution is the bounded
power in (D.CA.2). Its pullback degree on $C'$ is $O(m)$.
All subsequent constructions keep this factor separate.

\emph{Controlled geometry of the residue cover.}
Fix a separating finite map $x:C\to\mathbb P^1$ once.
Its composite on $C'$ is separable of degree $O(m)$, and
$g(C')-1=m(g(C)-1)$. The plane-coordinate construction in
Proposition~\ref{D:moving-line-polynomial-background} therefore
gives a primitive equation of degree $O(m)$, discriminant
degree $O(m^2)$, and rational expressions of polynomial degree
for functions with a polynomially bounded pole divisor.
Deck translations preserve $x$ and satisfy the same bounds.

For completeness the cyclic cover itself requires no
uncontrolled defining functions. Its character is a pair
$(P,P^m\simeq\mathcal O_C)$ with $P\in\operatorname{Pic}^0(C)$.
Riemann--Roch represents $P$ by a difference of two divisors
of degrees bounded on the fixed curve $C$; the chosen
power trivialization then has divisor degree $O(m)$.
All changes to the preceding plane chart are of polynomial
degree. The finitely many fixed poles can be avoided by
the moving divisors, as in the background proposition.

Choose a fixed finite affine cover of $C$ on which its
normal line is trivial. A formal neighbourhood over one
such smooth affine curve is a product with a formal disc:
lift the retraction successively by smoothness, and choose
a generator of the normal ideal; the affine obstruction
groups in positive cohomological degree vanish.
Only the first $a$ ordinary normal coefficients of these
fixed charts are used. Their poles are fixed constants.
Finite etale invariance lifts $C'$ uniquely to these charts;
on a product chart the lift is the corresponding product.
Across charts its transition is the unique etale lift of
the residual transition. In a primitive presentation the
coefficient recursion divides by the invertible residual
derivative. Its inverse has polynomial pole degree by
the discriminant bound. Since there are only $a$ ordinary
normal steps, all these coefficients have polynomial pole
degree in $m$. Equivalently this is the explicit infinitesimal
form of topological invariance of the etale site
\citep[Theorems 59.45.1--2 and Proposition 59.45.4]{StacksControlledAnchorsD}.

Let $Y'\to Y$ be the etale lift, and let $\mathcal X'$ be
the base change of $\mathcal X$; write $D'$ for its root
divisor, so $mD'=C'$ and the degree of $\mathcal O(D')$
on the residual root gerbe is $-d$. If $t=z^m$ on a root
chart, a change of normal coordinate is obtained from its
ordinary change by a unique root of a unit congruent to one.
Through $am$ root layers this still uses only $a$ ordinary
layers. Its degree bound is made explicit below.

\emph{A normalized period-$m$ tautological line.}
The pullback of the geometric Brauer class to $\mathcal X'$
is zero: it is zero on the puncture, and restriction of the
Brauer group of this regular root stack is injective there,
as used in Proposition~\ref{D:one-curve-coefficient-root-gerbe}.
Thus the pulled-back $\mu_m$-gerbe is the root gerbe of
a line $M$ on $\mathcal X'$. The already proved local
Picard calculation, now on $Y'$, and the root-character
description of lines show that its class modulo $m$ has
a representative
\[
       M=(H')^b\otimes\mathcal O(tD'),
                   0\leq b,t<m,\qquad \deg(H'|_{C'})=1.
                                                               \tag{D.CA.3}
\]
Indeed write a line as an ordinary line times a root
character, reduce that character modulo $m$, and use
$\operatorname{Pic}(Y')/m\simeq\mathbb Z/m$ for the
remaining ordinary line. The formal Picard kernel and
the degree-zero residual Picard group are $m$-divisible.
The line $H'$ has a controlled formal representative:
choose one point in a product chart of $C'$ and take its
constant transverse Cartier section, gluing with the
trivial divisor away from that point. Formal GAGA supplies
an algebraic representative if desired. Formula (D.CA.3)
therefore has both residual degree $O(m)$ and controlled
normal transitions. Let $\Lambda$ denote its tautological
weight-one line. This is a choice of a controlled root
gerbe model, not a bound imposed on an arbitrary extension
of a previously selected $\Lambda$.

\emph{Retain and lift the actual descent data.}
Let $\tau$ be a generator of the cyclic cover. On the
residual root gerbe retain the descent data induced by
the original period-$m$ gerbe. In the model (D.CA.3) this
consists of a ratio line $Q$, its power isomorphism
\[
                 Q^m\simeq\tau^*M\otimes M^{-1},       \tag{D.CA.4}
\]
and the actual cyclic coherence, in particular the specified
trivialization of the norm product of the translates of $Q$.
These exist because the original gerbe descends. They are
not replaced by an arbitrarily chosen trivial norm.
The sign of $Q$ in an alternative convention changes none
of the following estimates.

The line $Q$ can have a nonzero root character. Write
$Q=Q_0\otimes\mathcal O(jD')$ with $0\leq j<m$.
Its total rational degree is zero by (D.CA.4), whence
\[
                        \deg Q_0=jd=O(m).               \tag{D.CA.5}
\]
Choose divisor representatives of $Q_0$, and of any
ordinary residual comparison lines, with positive and
negative degrees $O(m)$, using Riemann--Roch on $C'$.
Power isomorphisms and norm coherences are rational
isomorphisms between these specified divisor lines;
their pole degrees are polynomial in $m$ (an $m$th power
and a product of $m$ translates suffice). For powers of
$Q$, first reduce the root character modulo $m$. Their
ordinary line again has degree $O(m)$ because the total
degree is zero, and it can be replaced by a divisor model
of that degree. The comparison with the unreduced product
has polynomial pole degree, at worst adding a polynomial
number of these divisor degrees. This controls all root
characters without a second cover trivializing $D'$.

Lift these divisor lines by constant transverse sections
on the product charts. Such a lift $Q_{\mathrm{ref}}$ need
not yet satisfy (D.CA.4). Its defect is a line trivial on
the residual root gerbe. On the formal Picard kernel,
multiplication by $m$ is an isomorphism: each successive
square-zero kernel is a quotient of the first cohomology
of an additive conormal sheaf, multiplication by $m$ is
invertible there, and compatible roots pass to the inverse
limit. Equivalently use Kummer and invariance of prime-to-$p$
cohomology under nilpotent thickenings. Global residual
units are constants and lift, so no residual-unit boundary
introduces $m$-torsion in this kernel.

Here is the cochain version, including its degree control.
Trivialize the residual defect using the retained power
isomorphism, and extend those trivializations on the affine
charts. Its normal transition functions are $1+u$, with
$u$ in the normal ideal. There is a unique $v\equiv1$
with $v^m=1+u$. If
$v=1+\sum_{j\geq1}v_jz^j$, coefficient comparison at
order $j$ has the form
\[
       m v_j=u_j-P_j(v_1,\ldots,v_{j-1}),               \tag{D.CA.6}
\]
where every monomial of $P_j$ has total normal weight $j$.
The recursion divides only by the unit $m$. No factorial
or binomial denominator is used, even when $j\geq p$.
Thus each $v_j$ is a polynomial in the input coefficients
of weighted degree at most $j$. The same statement holds
for unit inversion. If an initial coefficient has a
polynomial pole bound, the bound at order $j\leq am$ is
still polynomial. Correct $Q_{\mathrm{ref}}$ by this root
of its defect, with the inverse sign when necessary.

Uniqueness preserves all commutative descent identities.
For example, the product of two chosen unit roots is the
chosen root of the product. The norm line has trivial
$m$th power and is residually trivial, hence is trivial
in the formal Picard kernel. Adjust its trivialization
by the unique unit root congruent to one so that its
$m$th power is the prescribed telescoping trivialization.
The remaining cyclic coherence defects have $m$th power
one and reduce to one, and hence equal one. The same
argument on triple overlaps gives exact cocycles.
We have therefore lifted the \emph{actual} residual
cyclic descent, not merely its isomorphism classes of
ratio lines. Nilpotent invariance of $\mu_m$-cohomology
identifies this controlled formal gerbe with the original
one at every retained order; the constructions are
compatible as the order increases. The root corrections
are commutative line/unit corrections, with no assertion
about arbitrary matrix corrections.

All coefficient bounds just used are rational pole bounds
on a fixed two-affine cover of $C'$. The representatives
for the necessary additive conormal classes may be chosen
in the finite-pole Cech complexes of
Proposition~\ref{D:moving-line-polynomial-background}.
Their residual line degrees and genus are $O(m)$, so the
initial pole range is polynomial in $m$. Formula (D.CA.6)
and the compatible pole homotopy increase that range by
at most a polynomial through the retained $am$ layers.
Trace-dual conversion in that proposition turns a
polynomial pole divisor into polynomial plane-coordinate
degree. Consequently evaluating these fixed functions
at moving divisor points has polynomial degree too.
All moving points can avoid the finitely many fixed pole
fibres; the ordered-divisor construction still represents
every required line. This verifies the coefficient bound
for a chosen reference model, including its gluing data.

\emph{Frobenius comparisons and the full scalar band.}
For the period-$m$ factor put $v_i=p^i\bmod m$, with
$0\leq v_i<m$. Use the tautological powers $\Lambda^{v_i}$,
with reduced root characters and the preceding controlled
divisor models. A one-step comparison has ordinary factor
$M^{\lfloor p v_i/m\rfloor}$, where the exponent is at
most $p-1$. Its ordinary and root-character normalized
residual line degrees are $O(m)$. Translated powers of
the actual ratio $Q$ have the controlled representatives
just constructed. The one-step Frobenius raises residual
coordinates to their fixed $p$th powers; it does not
iterate a growing-degree rational formula through all
$f$ phases. Every phase uses its own reduced-character
reference. Hence the same polynomial bound applies to
each bounded-range adjacent-phase block, independently
of $f$.

To recover the full coefficient class, use the contracted
product of the period-$m$ and $n$th-root models in
(D.CA.1). On their product the line
$\Lambda^{v_i}\otimes T^{u_i}$ descends: the contracted
$\mu_m$ acts trivially since $v_i\equiv u_i\pmod m$.
The only extra ordinary Frobenius factor is (D.CA.2).
No degree-$n$ ordinary splitting cover or line is
introduced. Even on local trivializations, a power
$w^{u_i}$ of a unit $w\equiv1$ has, at normal order $j$,
monomials of total normal weight at most $j$ independently
of $u_i$; its numerical binomial coefficients are simply
constants in $k$. When a unit root is required, use
(D.CA.6) with the invertible integer $n$. Thus an arbitrarily
large scalar order cannot enter the pole bound. Scalar-zero
termwise Hom cancels these tautological factors, leaving
exactly the ordinary controlled factors described above.

The residual spectral translations and their Frobenius
comparison maps are those of the retained descent data.
In particular the final comparison after $f$ steps is
retained as prescribed by
Proposition~\ref{D:residual-twisted-Frobenius-transport};
it is not declared to be the identity merely because
$p^f\equiv1\pmod n$.

Finally each normal layer of a termwise Hom has a residual
line of degree $O(m)$ and genus $O(m)$, tensored with the
bounded number of spectral term bundles. Its cohomology
cost is the $O(ms^2)$ cost of the fixed-determinant
construction. The larger polynomial pole range only
affects its rational coefficient matrices. Adding the
moving determinant lines changes the background degree
to a polynomial in $m,s$, exactly as in
Proposition~\ref{D:moving-line-polynomial-background}.
All formulas are identities of line bundles, unit
cocycles, and rational matrices before coefficient
extension. They therefore remain identities over every
coefficient algebra on the stated charts. This proves
all assertions without replacing geometric existence
of an anchor by an unproved quantitative bound.
\end{proof}


% Separate proved count. Controlled anchors and the full synthesis passed
% all four actual-file audits; the exact audited candidate is preserved.
\subsection{An unbounded-period count on a one-curve local surface}

\begin{theorem}[Proved: quadratic projective counts with unrestricted geometric period]
\label{D:one-curve-unbounded-period-projective-count}
Let $R$ be a complete excellent normal local surface ring of
characteristic $p>0$, with an algebraically closed coefficient
and residue field $k$. Suppose a projective regular resolution
$Y\to\operatorname{Spec}R$ has exceptional locus a single
smooth connected curve $C$, with $C^2=-d<0$.
Put $U=\operatorname{Spec}R-\{\mathfrak m\}$ and
$\Gamma=\pi_1(U)$. Using the controlled-anchor construction
of Lemma~\ref{D:controlled-one-curve-gerbe-anchors}, there
is a constant $K_R$ such that, for every $q=p^f$ and
every $r\geq1$,
\[
 \#\bigl(\operatorname{Hom}_{\mathrm{cont}}
       (\Gamma,\operatorname{PGL}_r(\mathbb F_q))
           /\operatorname{PGL}_r(\mathbb F_q)\bigr)
                      \leq q^{K_Rr^2}.                 \tag{D.UP.1}
\]
The bound is uniform in the scalar order $q-1$ and in the
geometric Brauer period. The latter need only divide $r$;
it is not bounded in advance. This theorem is a local
one-smooth-exceptional-curve statement. It does not assert
finite generation of $\Gamma$ or settle general D.1 or D.2(ii).
\end{theorem}

\begin{proof}
The proof uses the controlled-anchor construction of
Lemma~\ref{D:controlled-one-curve-gerbe-anchors}. No existence-only
choice of an anchor is assigned a polynomial degree bound.

\emph{The actual coefficient class.}
Set $n=q-1$. A projective representation has its continuous
central lifting class in $H^2(\Gamma,\mathbb F_q^*)$.
The injection into $H^2(U,\mu_n)$ used in
Theorem~\ref{D:local-surface-constant-H2} retains the whole
class, including its Kummer part. Proposition
\ref{D:one-curve-coefficient-root-gerbe} bounds the number
of possible images by $d q^{2g(C)}$.
Fix one such class $\alpha$. Let $m$ be the period of its
geometric Brauer image. That proposition extends the actual
$\mu_n$-gerbe to the $m$th root along $C$, identifies its
connected degree-$m$ spectral cover $R_\alpha\to C$, and
provides lattices. Every weight-one rank-$r$ bundle on this
gerbe has $r=ms$, with $s\geq1$. If $m\nmid r$, the class
contributes no representations.

It is enough to bound the rank-$r$ unit-object classes on
this fixed coefficient gerbe. Indeed a representative
continuous central cocycle factors through a finite quotient
of $\Gamma$. Its central extension defines the coefficient
gerbe over $U$. Each projective representation of that
lifting class admits a linear representation of this fixed
extension with the central $\mathbb F_q^*$ acting by scalars.
It therefore gives a weight-one $\mathbb F_q$ local system,
or equivalently a $q$-Frobenius unit bundle, on the gerbe.
Projectivization recovers the original representation.
Changing a lift can overcount by scalar characters; it cannot
lose a projective class. Isomorphisms of unit bundles are
isomorphisms of $\mathbb F_q$ local systems, so this comparison
uses conjugacy over $\mathbb F_q$, not merely conjugacy after
extending the coefficient field to $k$.

\emph{Actual finite families covering the unit objects.}
Write a $q$-unit object as a cycle of $f$ one-step
$p$-Frobenius arrows. Choose the minimal lattices of
Lemma~\ref{D:twisted-connected-anchor-minimal-slopes}
at every stage. Lemma~\ref{D:controlled-one-curve-gerbe-anchors}
normalizes their fixed reference data, with residual
comparison-line degree $O(m)$ and polynomial rational-cochain
pole complexity, independently of $n$.
The slope estimates and their normal-layer proofs in
Proposition~\ref{D:twisted-period-independent-normal-depth}
then give a fixed ordinary-$C$ pole bound and a fixed
ordinary depth. If the residual comparison-line bound is
$O(m)$ rather than the narrower normalization used there,
the same proof only enlarges these fixed constants:
$g(R_\alpha)=1+m(g(C)-1)$ and the degree of the
$j$th root conormal layer is $jd$, so the required root
depth remains $O(m)$. No new factor of $m$ enters the number
of ordinary layers.

Fix a sufficiently large ordinary depth $a$, once for all
the bundle, arrow, inverse and comparison constructions.
The fixed ordinary twist used in Proposition
\ref{D:weighted-fixed-determinant-normal-presentation}
gives a residual presentation at each stage
\[
 0\longrightarrow L_i^{-1}\longrightarrow
 \mathcal O_{R_\alpha}^{s+1}\longrightarrow W_i^+
 \longrightarrow0,\qquad |\deg L_i|\leq c ms.          \tag{D.UP.2}
\]
The source anchor stays fixed. Only its rank-$m$ kernel
deforms normally. Consequently every root layer uses
$O(ms^2)$ cohomology coordinates, including all four
termwise Hom bundles, one-step Frobenius, maps, boundary
witnesses, and bounded-length inverse diagrams.
There are at most $am$ root layers. This is the specific
spectral-rank count; replacing it by the generic coarse
rank $mr$ would not give the desired estimate.

For each of the $f$ degrees in (D.UP.2), there are
$O(ms+1)$ choices. Proposition
\ref{D:moving-line-polynomial-background} covers the lines
by ordered divisors, using $O(fms)$ background coordinates.
Its exact-rank rational charts number at most
\[
 \exp\bigl(c_1fm\log(ms+1)\bigr).                     \tag{D.UP.3}
\]
The degrees of all background equations and all rational
cohomology and contraction matrices are polynomial in $m,s$.
The controlled-anchor lemma supplies precisely the fixed
normal transition and comparison coefficients also needed
by these matrices. On each such chart, the proof of
Lemma~\ref{D:weighted-fixed-residual-bundle-deformations}
therefore works over its coordinate ring. Its chosen
representatives and homotopies are actual rational sections,
with denominators inverted on that chart. The identities
are valid after every coefficient-base change.

The exact-rank charts need cover only geometric points for
this counting argument. No assertion that an arbitrary
noncoflat Artinian deformation factors through a chosen
rank chart is used. On each chart itself the splittings
are identities over its ring, and the constructed bundles,
maps and boundary witnesses are actual over every algebra
over that ring. Thus the finite families meet all the
effectivity requirements of Proposition
\ref{D:finite-jet-unit-Isom}. Every original unit object
occurs on at least one of their geometric points.
Every diagram retains its prescribed spectral translations
and comparison lines, including the final residual torsor
automorphism. The congruence $q\equiv1\pmod n$ restores the
scalar character; it does not discard that final geometric
comparison or identify it with the identity.

\emph{The weighted degree bound after moving the background.}
Choose an integer $E\geq1$, polynomially bounded in $m,s$,
which bounds the background degrees of the finitely many
normal transition, cohomology, contraction and comparison
operations through the retained depth. Enlarge it by a
fixed constant to include common denominators and all
background rank equations. Each equation involves only a
bounded number of adjacent phases; denominator clearing
is performed within that block, not simultaneously over
all $f$ phases. Thus $E$ is independent of $f$.
Give each background coordinate
weight one, each bundle or arrow coordinate at root layer
$j\geq1$ weight $Ej$, and each layer-zero or extra bundle
coordinate weight $E$.

In every perturbation word, positive normal orders add
and their sum is at most $am$. The quotient complexes have
internal width two, so there are at most a fixed number of
internal differential factors. Composition has fixed length,
and each arrow uses only one fixed-$p$ Frobenius pullback.
By the expanded-word-length bound in the moving-background
proof, all background coefficients of these expressions,
after clearing denominators, have degree at most $c_2Em$.
Their bundle-coordinate weighted degrees have the same
bound. Enlarging the fixed constant gives equations of
weighted degree at most $c_3Em$.

Denominator invertibility, residual subbundle inclusions
and the permitted slope windows are open restrictions.
They require no extra coordinates in this count. The exact
rank conditions are closed equations of polynomial degree
in background variables alone; their number need not be
bounded. We have obtained open parts of affine polynomial
schemes with the following budgets:
\[
 \begin{array}{ll}
 \text{background variables, weight }1:& A fms,\\
 \text{variables of weight }Ej,\ 1\leq j\leq am:& A fms^2
                  \text{ at each layer},\\
 \text{additional variables of weight }E:& A fms^2,\\
 \text{maximum equation weight}:& KEm,
       \qquad E\leq K(ms)^h.
 \end{array}                                             \tag{D.UP.4}
\]
The constants and the exponent $h$ depend only on the fixed
surface, $p$ and the finite diagram pattern. In particular
they do not depend on $n,m,s,f$ or the coefficient class.

\emph{Components count genuine unit classes.}
Proposition~\ref{D:finite-jet-unit-Isom} gives a finite
presentation model of actual unit isomorphisms on each
family, using both map and inverse jets. Its proof uses
all Artinian coefficient extensions and algebraization,
so the isomorphism-class fibres are open and closed.
The number of classes on a chart is therefore at most its
number of irreducible components. Corollary
\ref{D:weighted-normal-layer-background-budget} applies
to (D.UP.4) and bounds this number by
\[
                            \exp(K_1 f m^2s^2)
                                  =\exp(K_1 f r^2).     \tag{D.UP.5}
\]
This is where the factorial product of the normal weights
is used. An ordinary unweighted degree bound would leave
an unwanted logarithm on the entire quadratic block.

Multiplying by the determinant-degree choices and (D.UP.3)
preserves (D.UP.5), with a larger constant, since
$m\log(ms+1)\leq c(ms)^2$ for $m,s\geq1$.
Since $q=p^f$, this is $q^{K_2r^2}$ for a fixed $K_2$.
Finally multiply by $d q^{2g(C)}$, the total number of
actual coefficient classes. Enlarging $K_R$ absorbs that
factor uniformly for $q\geq2,r\geq1$, proving (D.UP.1)
using the proved controlled-anchor construction.
No assertion about other exceptional configurations, or
about the remaining simple-group and crown gates, is used.
\end{proof}


% Separate post-CP20 auxiliary input. Root/B/C/D actual-file audits passed.
\begin{lemma}[Proved: twisted determinant classes on etale resolution covers]
\label{D:unramified-semilinear-determinant-budget}
Let $(R,\mathfrak m,k)$ be a complete excellent normal local
$k$-algebra of dimension two, with algebraically closed coefficient
and residue field $k$ of characteristic $p>0$. Let
$Y\to\operatorname{Spec}R$ be a projective regular resolution,
an isomorphism over $U=\operatorname{Spec}R-\{\mathfrak m\}$,
whose reduced exceptional fibre is one smooth connected curve $C$
of genus $g$. Let $Y_t\to Y$ be a connected finite etale cover
of degree $t$, and put $C_t=Y_t\times_Y C$ and $U_t=Y_t-C_t$.
Then $C_t$ is smooth and connected, with
$g_t=g(C_t)=1+t(g-1)$.

For an integer $n\geq1$ prime to $p$ and a locally constant
sheaf $\mathcal M$ of free rank-one $\mathbf Z/n\mathbf Z$-modules
on $Y_t$, one has
\[
 \bigl|H^1_{\mathrm{et}}(U_t,\mathcal M|_{U_t})\bigr|
       \leq n^{2g_t+1}.                                 \tag{D.SD.1}
\]
The assertion for $n=1$ means that the group is trivial.
No restriction prime to $p$ is imposed on the order of the
monodromy group of $\mathcal M$.

In particular, let $q=p^f$, and fix a continuous field character
\[
 \chi:\pi_1(U_t)\longrightarrow
             \operatorname{Gal}(\mathbb F_q/\mathbb F_p)
\]
that factors through $\pi_1(Y_t)$. Give $\mathbb F_q^*$ the
action prescribed by $\chi$, and denote this finite module by
$\mathbb F_q^*(\chi)$. Then
\[
 \bigl|H^1_{\mathrm{cont}}
       (\pi_1(U_t),\mathbb F_q^*(\chi))\bigr|
       \leq(q-1)^{2g_t+1}\leq q^{(2g+3)t}.               \tag{D.SD.2}
\]
This bounds the possible semilinear determinant classes for
representations with fixed field character $\chi$, modulo
$\mathbb F_q^*$-conjugacy. The bound is independent of
$d=|\operatorname{im}\chi|$, including when $p\mid d$.
\end{lemma}

\begin{proof}
The proper henselian finite-etale comparison identifies the
category of finite etale covers of $Y$ with that of its closed
fibre; removing nilpotents gives the same category as for $C$
\citep[Lemma 58.9.1, Tag 0A48, complete proof]
{StacksSemilinearDeterminantsB}. Thus the connected cover $Y_t$
restricts to a connected cover $C_t\to C$. This is an etale
degree-$t$ map of smooth proper curves. Its genus is the one in
the statement, by etale Riemann--Hurwitz
\citep[Section 53.12 and Lemma 53.12.2, Tags 0C1B and 0C1D]
{StacksSemilinearDeterminantsB}. If $g=0$, the formula in
particular forces $t=1$; it never assigns a negative genus to
an existing cover.

We may assume $n>1$. Write $j:U_t\hookrightarrow Y_t$ and
$i:C_t\hookrightarrow Y_t$. The regular codimension-one
purity calculation gives
\[
 j_*j^*\mathcal M=\mathcal M,\qquad
 R^1j_*j^*\mathcal M=i_*(\mathcal M|_{C_t}(-1)).          \tag{D.SD.3}
\]
For constant coefficients this is exactly the one-component
case of \citet[Expose XVI, Theorem 3.1.1 and
Corollary 3.1.3]{A_ILO2014}. Its stated regular-scheme scope
applies to this excellent resolution, without a finite-type
over $k$ assumption. For the displayed locally constant
coefficients, trivialize $\mathcal M$ etale locally on $Y_t$
and use the functoriality of the purity isomorphism. The
trivializing cover need not have degree prime to $p$.
Only the coefficient integer $n$ must be invertible.

The degree-one Leray sequence for $j$ now gives the exact sequence
\[
 0\longrightarrow H^1(Y_t,\mathcal M)
 \longrightarrow H^1(U_t,j^*\mathcal M)
 \longrightarrow H^0(C_t,\mathcal M|_{C_t}(-1)).          \tag{D.SD.4}
\]
Every term is etale cohomology. A section of a locally constant
sheaf on the connected curve $C_t$ is determined by its value
at one geometric point. Hence the last term has at most $n$
elements. Equivalently it is the invariant subgroup of a fibre
of order $n$; no trivial-monodromy assertion is used.

The composite $Y_t\to\operatorname{Spec}R$ is proper.
Since $R$ is strictly henselian, its global-section functor
on the etale site is the closed geometric-stalk functor and
is exact. Proper base change, followed by invariance under
the nilpotents in the closed fibre, therefore identifies
\[
          H^1(Y_t,\mathcal M)
             =H^1(C_t,\mathcal M|_{C_t}).                \tag{D.SD.5}
\]
Here proper base change is used for an arbitrary torsion
abelian sheaf, precisely as in
\citep[Theorem 59.91.11, Tag 095T, complete proof]
{StacksSemilinearDeterminantsB}.

By \citet[Theorem X.2.6 and its preceding complete argument,
original pp.~272--273]{SGA1}, the profinite group
$\Pi=\pi_1(C_t)$ has at most $2g_t$ topological generators.
This is the full fundamental group, obtained by surjective
specialization from a characteristic-zero lift. It is not
only its prime-to-$p$ quotient. Let $A$ be a fibre of
$\mathcal M|_{C_t}$, with its continuous $\Pi$-action.
Degree-one torsor classification identifies
$H^1(C_t,\mathcal M|_{C_t})$ with $H^1_{\mathrm{cont}}(\Pi,A)$.
A crossed homomorphism $z:\Pi\to A$ is determined by its
values on the $2g_t$ chosen generators: the identity
$z(ab)=z(a)+a z(b)$ determines it on the abstract subgroup
they generate, and continuity determines it on its closure.
There are at most $n^{2g_t}$ such choices. Passing to
cohomology classes can only decrease this number. Combining
this with (D.SD.4)--(D.SD.5) proves (D.SD.1).

For the last assertion take $n=q-1$ and the locally constant
cyclic module on $Y_t$ defined by $\chi$. Its restriction to
$U_t$ is $\mathbb F_q^*(\chi)$, so degree-one torsor
classification and (D.SD.1) give the first inequality in
(D.SD.2). The second follows from
$2g_t+1=2t(g-1)+3\leq(2g+3)t$. When $q=2$, the
coefficient group is trivial and the same assertion holds.

For clarity, write a semilinear representation with prescribed
field part as $\rho(\gamma)=(A_\gamma,\chi(\gamma))$.
Its determinant $z(\gamma)=\det A_\gamma$ satisfies
\[
 z(\gamma\eta)=z(\gamma)\,
                   \chi(\gamma)(z(\eta)).              \tag{D.SD.6}
\]
Conjugating by a matrix of determinant $a$ changes it to
$a^{-1}z(\gamma)\chi(\gamma)(a)$, a coboundary change.
Conversely every $a\in\mathbb F_q^*$ occurs as a matrix
determinant. Thus the relevant determinant invariant is
exactly a class in the group counted in (D.SD.2), rather
than an arbitrary character on the degree-$d$ cover that
trivializes the field part. No factor $d$ enters this bound.
\end{proof}

\begin{remark}[Counting the field characters themselves]
If the field character is also allowed to vary subject to
extension over $Y_t$, there are at most $f^{2g_t}$ such
characters. Indeed proper henselian comparison gives
$\pi_1(Y_t)\simeq\pi_1(C_t)$, and a homomorphism to the
cyclic group of order $f$ is determined by the same $2g_t$
generators. Since $f\leq p^f=q$, this factor is bounded by
$q^{(2g+2)t}$. This observation imposes no condition that
$f$ or the image order be prime to $p$.
\end{remark}


% Separate post-CP20 proved packet. Restricted hereditary counts and an actual crown family.
\subsection{Hereditary counts for etale resolution covers and permutation crowns}

Fix a complete excellent normal local $k$-algebra $R$ of dimension
 two and characteristic $p>0$, with algebraically closed coefficient
and residue field $k$. Let $Y\to\operatorname{Spec}R$ be a projective
regular resolution, an isomorphism over
$U=\operatorname{Spec}R-\{\mathfrak m\}$, whose exceptional divisor
is one smooth connected curve $C$. Write $g=g(C)$, $C^2=-d<0$,
and $\Gamma=\pi_1(U)$. An \emph{etale resolution cover} means a
connected finite etale morphism $Y'\to Y$, with
$U'=Y'\times_Y U$. Its degree is denoted by $t$.

\begin{theorem}[Proved: a linear-index hereditary count on etale resolution covers]
\label{D:etale-resolution-hereditary-linear-count}
There is a constant $A_R$ such that every etale resolution cover
of degree $t$, every $q=p^f$, and every $r\geq1$ satisfy
\[
 \#\bigl(\operatorname{Hom}_{\mathrm{cont}}
       (\pi_1(U'),\operatorname{GL}_r(\mathbb F_q))
                  /\operatorname{GL}_r(\mathbb F_q)\bigr)
                         \leq q^{A_R t r^2}.             \tag{D.HC.1}
\]
The same $A_R$ works for all such covers, including degrees
divisible by $p$. The assertion is not made for an arbitrary
finite etale cover of $U$ that ramifies above $C$.
\end{theorem}

\begin{proof}
\emph{The covering geometry and its uniform scale.}
Proper henselian comparison identifies finite etale covers of
$Y$ with those of its closed fibre, and nilpotents do not change
that category \citep[Lemma 58.9.1, Tag 0A48, complete proof]
{StacksSingleCurveBrauerD}. Thus
$C'=(Y'\times_Y C)_{\mathrm{red}}$ is a smooth connected
etale degree-$t$ cover of $C$. It is already the scheme-theoretic
pullback of the smooth Cartier divisor $C$. The Stein ring
$R'=H^0(Y',\mathcal O_{Y'})$ is a finite complete local normal
$R$-algebra with residue field $k$. Indeed $Y'$ is integral
and normal, its global functions are integrally closed, and
its closed fibre is connected; a finite algebra over the
henselian local base splits into its local factors, so there
can be only one. The Stein morphism is a regular projective
resolution of $\operatorname{Spec}R'$, isomorphic over $U'$.
Its exceptional divisor is precisely $C'$. In particular,
\[
             g_t-1=t(g-1),\qquad (C')^2=-td,
                         \qquad g_t=g(C').              \tag{D.HC.2}
\]
The genus formula is etale Riemann--Hurwitz, and the
self-intersection formula is the finite-flat projection formula.
When $g=0$, the only connected etale curve cover has $t=1$;
no negative genus is being allowed in (D.HC.2).

Choose a minimal ordinary lattice $F_i$ for each one-step
Frobenius stage of a rank-$r$ local system. The untwisted
case of Lemma~\ref{D:twisted-connected-anchor-minimal-slopes}
gives
\[
 t(g-1)\leq\mu_{\min}(F_i|_{C'})
       \leq\mu_{\max}(F_i|_{C'})
       \leq t(g-1+d).                                  \tag{D.HC.3}
\]
No semistability of a tensor product is used. Pick a fixed
integer $h>0$ with $hd\geq g+2$, and put
$F_i^+=F_i(-hC')$. Its residual minimum slope is larger
than $2g_t-1$, uniformly in $t$. Hence it is globally
generated on $C'$, and every positive conormal twist has
vanishing first cohomology.

We spell out why the pole and realization depth can also be
chosen independently of $t$. Formula (D.TD.14), with
$g_m=g_t$ and the ordinary normal degree $td$, places the
one-step Frobenius pullbacks of the bundles in (D.HC.3)
in intervals whose endpoints have absolute value at most
$a_0t$. The same is true after the fixed twist $-hC'$.
The constant $a_0$ depends only on $g,d,p,h$. A negative
normal Hom layer has zero $H^0$ once
$jtd>2a_0t$. A positive layer has zero $H^1$ once
\[
                jtd>2a_0t+2g_t-2.                      \tag{D.HC.4}
\]
These are precisely the slope-gap and Serre-duality arguments
of Proposition~\ref{D:twisted-period-independent-normal-depth}.
They give fixed integers $B,N$, independent of $t,r,f$, for
poles and specified-map lifting. Enlarging the constants
once includes the pole twists themselves. The paired unit
jet correction of Lemma~\ref{D:meromorphic-unit-jet-contraction}
then uses a fixed ordinary precision $a$, again independent
of $t,r,f$. We choose it to include the inverse-product
precision and all one-step Frobenius arrows. Vanishing
persists over every coefficient base on the actual finite
families, by the two-term proper-curve cohomology argument
in Proposition~\ref{D:finite-jet-unit-Isom}.

\emph{A finite list of residual determinants, with phase compatibility.}
Set $n=q-1$. The determinant of the local system is a
character $\chi:\pi_1(U')\to\mathbb F_q^*$. Its possibilities
are bounded by
\[
 |H^1(U',\mu_n)|
                 \leq n^{2g_t+1}\leq q^{(2g+3)t}.       \tag{D.HC.5}
\]
For $n=1$ the first group is trivial. For general $n$, the
one-divisor localization sequence on the regular pair
$(Y',C')$ injects the quotient of $H^1(U',\mu_n)$ by
$H^1(Y',\mu_n)$ into $H^0(C',\mathbb Z/n)$.
Proper strict-local base change identifies
$H^1(Y',\mu_n)$ with $H^1(C',\mu_n)$, of order $n^{2g_t}$.
This is the degree-one part of the purity calculation used
in Theorem~\ref{D:local-surface-constant-H2}; the exact
excellent-base purity scope is the one already checked in
\citet[Expose XVI, Theorem 3.1.1 and Corollary 3.1.3]
{A_ILO2014}. Thus (D.HC.5) does not presuppose finite
generation of $\pi_1(U')$.

For a fixed $\chi$, the punctured determinant line at the
$i$th $p$-Frobenius stage is the line of $\chi^{p^i}$.
Thus the $f$ stages introduce only \emph{one} character
choice in (D.HC.5). A line on $U'$ extends to the regular
surface $Y'$ by closing a Cartier divisor. Two extensions
differ by $\mathcal O_{Y'}(bC')$, $b\in\mathbb Z$:
a rational trivialization of their ratio has divisor
supported on the sole exceptional prime. Their degrees
on $C'$ consequently differ by the nonzero step $td$.
The determinant degree permitted by (D.HC.3) lies in an
interval of length $rtd$. Hence there are at most $r+1$
possible extension lines at each stage. Fixing an extension
fixes its restriction to $C'$. After the common twist
$-hC'$, the residual determinant $L_i$ still has
\[
                         |\deg L_i|\leq a_1tr,          \tag{D.HC.6}
\]
and there are at most $(r+1)^f$ lists of these lines.
There is no assertion about flatness of arbitrary reflexive
hulls over parameter rings: these are choices for geometric
unit objects, after which the finite family is constructed
from actual bundles.

\emph{One fixed-line presentation per list.}
On $C'$, a general $r+1$-tuple of sections of the globally
generated $F_i^+|_{C'}$ gives
\[
 0\longrightarrow L_i^{-1}\longrightarrow
 \mathcal O_{C'}^{r+1}\longrightarrow F_i^+|_{C'}
 \longrightarrow0.                                    \tag{D.HC.7}
\]
The rank-drop locus for a general tuple has expected
codimension two on a curve; the usual incidence argument
therefore proves surjectivity. Its kernel is the indicated
line by determinants. The positive-conormal $H^1$
vanishing lifts these generators from the fixed formal
source $\mathcal O^{r+1}$ through every normal order.
Only the kernel line deforms. Choose any fixed reference
lift of each $L_i^{-1}$ and apply the actual all-ring
construction of Lemma~\ref{D:weighted-fixed-residual-bundle-deformations}
to $aC'$, together with its finite quotient Hom complexes,
Frobenius and inverse witnesses.

Here the number of normal layers is the fixed integer $a$;
there is no growing residue-period parameter. Every residual
line in one termwise Hom has degree $O(tr)$ and genus
$O(t)$. By curve duality and Riemann--Roch, its total
$H^0,H^1$ dimension is $O(tr)$. The $\mathcal O$ terms
have cohomology $O(t)$; their matrix multiplicities are
at most $(r+1)^2$. Consequently, including every stage,
arrow and boundary witness, this construction has
\[
          \text{at most }A ftr^2\text{ variables},
       \qquad \text{equation degree at most }D_0,       \tag{D.HC.8}
\]
where $A,D_0$ depend only on the fixed original surface
and $p$. Indeed the normal word length is at most $a$,
the internal Hom width is two, the composition diagram
has fixed length, and each arrow uses one $p$th power.
All reference cochains and all cohomology contractions
are fixed $k$-linear maps. Their coefficients are constants,
even if a chosen reference has complicated rational poles.
No moving-divisor evaluation is made, so no polynomial
degree in $t$ is charged to those constants. This is the
reason that the moving-background $t\log t$ loss is absent.

There is one affine polynomial presentation, followed by
open conditions, for each fixed list of $L_i$. The residual
subbundle and slope-window conditions are open by the exact
proper-Quot argument in the cited fixed-determinant
proposition. All its finite bundles and arrows are actual
over every coefficient algebra. The constant-depth
realization above meets the hypotheses of
Proposition~\ref{D:finite-jet-unit-Isom}; thus isomorphism
classes of unit objects are open and closed on this family.
The number of represented classes is bounded by its
number of irreducible components. The ordinary specialization
of Lemma~\ref{D:weighted-affine-component-bound} bounds this
by $\exp(A_2ftr^2)$ from (D.HC.8), independently of the
number of defining equations. Open restrictions do not
increase that number.

Multiply by the $(r+1)^f$ determinant lists and by
(D.HC.5). Since $q=p^f$ and $\log(r+1)=O(r^2)$, the
result is (D.HC.1), after increasing one constant $A_R$.
The comparison is with $\mathbb F_q$-linear local systems,
so it counts conjugacy over $\mathbb F_q$. It does not
replace such conjugacy by arbitrary $k$-linear conjugacy
of monodromy matrices. This proves the claimed uniformity.
\end{proof}

\begin{theorem}[Proved: bounded generators for unramified pure-permutation crowns]
\label{D:unramified-permutation-linear-crowns}
Retain the preceding hypotheses and fix a constant $\kappa>0$.
Let $S$ be a nonabelian finite simple group with a faithful
embedding $S\hookrightarrow\operatorname{GL}_h(\mathbb F_q)$,
where $q$ is a power of $p$, such that
\[
                      h^2\log q\leq\kappa\log|S|.       \tag{D.HC.9}
\]
Let $P\leq\operatorname{Sym}(t)$ be transitive, acting on
$A=S^t$ by pure permutation of the factors, and put
$L=A\rtimes P$. Suppose a continuous surjection
\[
               \Gamma\twoheadrightarrow L_k
                   = A^k\rtimes P                      \tag{D.HC.10}
\]
onto the crown-based power has its induced quotient
$\rho:\Gamma\twoheadrightarrow P$ factoring through
$\pi_1(Y)$. Set $c=\lceil\kappa(A_R+1)\rceil$, increasing
it to one if necessary. Then
\[
             k\leq |S|^{ct}=|A|^c,
          \qquad d(L_k)\leq\max\{2,2g,c+2\}.           \tag{D.HC.11}
\]
Thus these crowns cannot have unbounded generator number
as $S,t,k$ vary. This is a bound on the displayed crown
family, not a finite-generation conclusion for $\Gamma$.
Nontrivial outer actions on $S$, permutation quotients
ramified above $C$, and simple groups not satisfying
(D.HC.9) are not included.
\end{theorem}

\begin{proof}
Fix the quotient $\rho$ while comparing its lifts. Let
$H=\rho^{-1}(P_1)$ be the preimage of a point stabilizer.
It has index $t$ in $\Gamma$. Since $\rho$ extends over
$Y$, its cover is $U'\subset Y'$ for an etale resolution
cover of degree $t$. Apply (D.HC.1) to $H$.

We give the precise nonabelian Shapiro comparison. Write
$X=\Gamma/H$ and identify $A$ with $S^X$, with action
$(\gamma a)_x=a_{\gamma^{-1}x}$. A lift of $\rho$ to
$L$ is a continuous cocycle
$c(\gamma\delta)=c(\gamma)\,\gamma c(\delta)$.
Evaluation at the base coset, restricted to $H$, gives
an ordinary homomorphism $H\to S$. Conversely choose
coset representatives $\gamma_x$ with $\gamma_1=1$.
For $\eta:H\to S$, the formula
\[
 c_\eta(\gamma)_x
       =\eta(\gamma_x^{-1}\gamma\gamma_{\gamma^{-1}x})
                                                               \tag{D.HC.12}
\]
is a continuous cocycle. The product inside $\eta$ lies
in $H$, and multiplying the two such products proves the
cocycle identity. For an arbitrary $c$, set
$a_x=c(\gamma_x)_x$. The cocycle identity gives
\[
 c_\eta(\gamma)_x=a_x^{-1}c(\gamma)_x a_{\gamma^{-1}x}.
\]
Consequently evaluation induces a bijection between lifts
modulo conjugation by $A$ and $\operatorname{Hom}(H,S)/S$.
All choices here are finite coset choices, so continuity
is preserved. If the lift surjects onto $L$, then its
restriction to $\ker\rho$ surjects onto $A$, and the
associated $\eta$ surjects onto $S$.

The $k$ coordinates of (D.HC.10) are surjective lifts of
this same fixed $\rho$. No two are $A$-conjugate: otherwise
the corresponding pair of coordinates would lie in a
proper conjugate diagonal over $P$, contradicting the
surjectivity onto $L_k$. Thus
\[
                 k\leq\#\operatorname{Epi}(H,S)/S.      \tag{D.HC.13}
\]
Under the fixed faithful embedding in (D.HC.9), a
$\operatorname{GL}_h(\mathbb F_q)$-conjugacy class contains
at most
$|N_{\operatorname{GL}_h(\mathbb F_q)}(S)|/|S|$
$S$-conjugacy classes of surjections. Indeed a matrix
conjugating two such images normalizes their common
image $S$. In particular this multiplicity is at most
$q^{h^2}$. Equations (D.HC.1), (D.HC.9) and (D.HC.13)
therefore give
\[
 k\leq q^{(A_Rt+1)h^2}
       \leq |S|^{\kappa(A_R+1)t}\leq |A|^c.            \tag{D.HC.14}
\]

It remains to bound the generator number, rather than
only the number of lifts. The finite group $P$ is a
quotient of $\pi_1(Y)\simeq\pi_1(C)$, so $d(P)\leq2g$.
Here we use the complete smooth-proper-curve generator
argument and Theorem X.2.6 of \citet{SGA1}, original
pp.~272--273: lift to characteristic zero and use the
surjective specialization map. We do not infer this
bound from smallness. The pure-permutation group $L$
is monolithic primitive with socle $A$: transitivity
makes $A$ minimal normal and its centralizer in $L$
is trivial.

We use precisely the crown estimates already stated in
Proposition~\ref{D:cohomology-crown-criterion}. For
$s\geq d(L)=\max\{2,d(P)\}$, the complete proof of
\citet[Proposition 2.6, equations (2.1)--(2.2), pp.~12--13]
{LucchiniMarionTracey2020D} supplies
\[
 k\leq\psi_L(s)\ \Longrightarrow\ d(L_k)\leq s,
 \qquad
 \psi_L(s)\geq\frac{|A|^{s-1}}{2t|\operatorname{Out}(S)|}.
                                                               \tag{D.HC.15}
\]
The same source records
$|\operatorname{Out}(S)|\leq\log_2|S|$. Hence
$2t|\operatorname{Out}(S)|\leq|S|^t=|A|$.
Take $s=\max\{2,2g,c+2\}$ in (D.HC.15). Then
$\psi_L(s)\geq|A|^{s-2}\geq|A|^c\geq k$, proving
(D.HC.11). These finite-group inputs retain their
classification dependence; no elementary replacement is
being claimed. The uniform bound removes exactly the
restricted nonabelian crown family in the statement.
\end{proof}


% Proved after four actual-file proof audits; exact candidate and source-read scopes retained.
\subsection{Field-semilinear counts and unramified field-permutation crowns}

We retain the complete one-smooth-exceptional-curve setting of
Theorem~\ref{D:etale-resolution-hereditary-linear-count}. In this
subsection $C^2=-\delta<0$, $g=g(C)$ and $\Gamma=\pi_1(U)$.
All representations are continuous. For $q=p^f$, write
$\mathcal G_q=\operatorname{Gal}(\mathbb F_q/\mathbb F_p)$.
It acts entrywise on $\operatorname{GL}_r(\mathbb F_q)$.

\begin{lemma}[Proved: twisted Frobenius models of a fixed field character]
\label{D:twisted-deck-Frobenius-semilinear-model}
Let $W$ be a connected locally noetherian scheme over $\mathbb F_q$
with a geometric base point, and fix a
surjective continuous character
$\chi:\pi_1(W)\to\langle\sigma\rangle\subset\mathcal G_q$,
where $\sigma(a)=a^{q_0}$, $q_0=p^e$, $f=ed$ and
$|\langle\sigma\rangle|=d$. Let $v:V\to W$ be its connected
cyclic finite etale cover. Choose the generator $\tau$ of its
deck transformations with the pullback convention for which a
descent map $\tau^*\mathcal L\to\mathcal L$ is
$\sigma$-semilinear. This is the inverse of the corresponding
monodromy permutation convention.

The groupoid of rank-$r$ representations
\[
 \rho:\pi_1(W)\longrightarrow
      \operatorname{GL}_r(\mathbb F_q)\rtimes\langle\sigma\rangle
 \quad\hbox{with field projection }\chi,
                                                        \tag{D.SC.1}
\]
with conjugacies by $\operatorname{GL}_r(\mathbb F_q)$, is
equivalent to the groupoid of vector bundles $E$ on $V$ with
an isomorphism
\[
             \Phi:\tau^*F_{q_0}^*E\xrightarrow{\sim}E.  \tag{D.SC.2}
\]
Here $F_{q_0}$ is absolute Frobenius; a morphism is a bundle
map commuting with (D.SC.2). In the relative version, finite
local systems replace the chosen fundamental-group notation.
That construction works after arbitrary $\mathbb F_q$-coefficient-base
extension. In the surface applications all schemes and coefficient
algebras are over the fixed algebraically closed field $k$.
\end{lemma}

\begin{proof}
Put $T=\tau^*F_{q_0}^*$. The two pullbacks commute, and
$\tau^d=1$, so there is a canonical identification
$T^d=F_q^*$. Define, rather than independently prescribe,
\[
 \Psi=\Phi\circ T\Phi\circ\cdots\circ T^{d-1}\Phi:
                              F_q^*E\xrightarrow{\sim}E.\tag{D.SC.3}
\]
The ordinary unit-bundle/Lang-frame equivalence used in
Proposition~\ref{D:finite-jet-unit-Isom} gives a locally
constant sheaf $\mathcal L$ of rank-$r$ $\mathbb F_q$-vector
spaces on $V$. Concretely its sections are the vectors fixed
by $\Psi$; its frame functor is finite etale and locally
nonempty. This construction and its inverse, extension of
scalars from $\mathbb F_q$ to $\mathcal O_V$, work over any
coefficient algebra over $\mathbb F_q$. No perfectness or
reducedness of that algebra is required: the equations for
Lang frames have invertible linear derivative.

The identity obtained by rotating the $d$ factors in
(D.SC.3) shows that $\Phi$ carries fixed frames to fixed
frames. It induces a $\sigma$-semilinear map
$D:\tau^*\mathcal L\to\mathcal L$. Its $d$-fold deck
composite is the action of $\Psi$ on its fixed sections,
and is therefore the identity. Thus $D$ is genuine cyclic
semilinear descent, with the prescribed field character.
Effective finite etale descent gives (D.SC.1).

Conversely, a representation (D.SC.1) restricts on
$\ker\chi$ to an $\mathbb F_q$-linear local system
$\mathcal L$ on $V$, with this semilinear deck descent.
Set $E=\mathcal O_V\otimes_{\mathbb F_q}\mathcal L$.
The $\sigma$-semilinearity is exactly the scalar twist
needed to extend $D$ to the linear map (D.SC.2) from
$\tau^*F_{q_0}^*E$. Its $d$-fold composite is the ordinary
$F_q$-unit map of $\mathcal L$. These constructions are
inverse on objects and on the stated morphisms, since
frames recover all the descent maps.

This is a categorical deck-cocycle assertion. It does not
assert that an arbitrarily chosen element of $\pi_1(W)$
lifting $\sigma$ has order $d$, or that its semilinear
matrix has $d$th power one. The $d$th power of that element
lies in $\pi_1(V)$; its path transport is already present
in the descended local system. Replacing it by a trivial
pointed group relation would lose representations.
\end{proof}

\begin{theorem}[Proved: an index-linear count with unramified field action]
\label{D:etale-resolution-semilinear-count}
There is a constant $A_R'$ with the following property.
Let $Y_H\to Y$ be a connected finite etale cover of degree
$t$, set $U_H=Y_H\times_Y U$ and $H=\pi_1(U_H)$, and
let $\chi:H\to\mathcal G_q$ be any character factoring
through $\pi_1(Y_H)$. For every $r\geq1$, the number of
representations
\[
 H\longrightarrow\operatorname{GL}_r(\mathbb F_q)
                         \rtimes\mathcal G_q
 \quad\hbox{with field projection }\chi,
                                                        \tag{D.SC.4}
\]
up to conjugation by $\operatorname{GL}_r(\mathbb F_q)$,
is at most $q^{A_R'tr^2}$. The constant is independent of
$q,t,r,\chi$ and the order of its field image. The assumption
that $\chi$ extends over $Y_H$ is part of the statement.
\end{theorem}

\begin{proof}
Write $|\operatorname{im}\chi|=d$, $e=f/d$ and $q_0=p^e$.
The character defines a connected cyclic degree-$d$ finite
etale cover $Y_V\to Y_H$, with deck generator $\tau$,
and $V=Y_V\times_Y U$. Its composite degree over $Y$
is $td$. As in the proof of
Theorem~\ref{D:etale-resolution-hereditary-linear-count},
the Stein ring is a complete normal local surface ring,
and $Y_V$ is a projective regular resolution with sole
smooth exceptional curve $C_V$. We have
\[
       g(C_V)-1=td(g-1),\qquad C_V^2=-td\delta.          \tag{D.SC.5}
\]
The deck map preserves $C_V$ and its normal ideal.

By Lemma~\ref{D:twisted-deck-Frobenius-semilinear-model},
we count (D.SC.2). Factor it into $e$ one-step
$p$-Frobenius arrows, placing $\tau^*$ on the last arrow.
For instance the intermediate punctured bundles may be
$F_p^{i*}E$ for $0\leq i<e$; allowing all isomorphic
one-step presentations gives the same groupoid. This uses
$e$ stages, not $f$ stages on the degree-$td$ cover.

\emph{The determinant is counted before passing to $V$.}
Taking the determinant in (D.SC.4) gives a multiplicative
cocycle in the finite $H$-module
\[
           A_\chi=\mathbb F_q^*,\qquad
                    h\cdot a=\chi(h)(a).
\]
Its isomorphism class lies in $H^1(H,A_\chi)$. The
twisted-coefficient bound of Lemma~\ref{D:unramified-semilinear-determinant-budget}
for this regular one-divisor pair gives
\[
 |H^1(H,A_\chi)|
       \leq(q-1)^{2g(C_H)+1}\leq q^{(2g+3)t}.          \tag{D.SC.6}
\]
For clarity, the ingredients are that $A_\chi$ extends
to a finite local system on $Y_H$, the low-degree
localization sequence has cokernel in
$H^0(C_H,A_\chi(-1))$, and proper henselian base change
reduces $H^1(Y_H,A_\chi)$ to the smooth proper curve
$C_H$. Its fundamental group has at most $2g(C_H)$
generators. This reasoning allows $p\mid d$ because
the coefficient group $A_\chi$ itself has order prime
to $p$; the field-action monodromy is unramified on $Y_H$.
No list of arbitrary ordinary characters on the larger
cover $V$ is used.

Fix one determinant class in (D.SC.6). Its restriction to
$V$ fixes the ordinary punctured determinant line of $E$.
At stage $i$ it fixes its $p^i$th pullback. All extensions
of one such line to $Y_V$ differ by
$\mathcal O_{Y_V}(bC_V)$, $b\in\mathbb Z$. The minimal
residual lattice window has degree width $rtd\delta$,
whereas a change by $C_V$ changes degree by $td\delta$.
There are therefore at most $r+1$ possible determinant
extension lines at each stage, or $(r+1)^e$ lists.

\emph{Uniform actual families with $e$ stages.}
The fixed-line construction in the proof of
Theorem~\ref{D:etale-resolution-hereditary-linear-count}
applies with the cover degree there replaced by $td$.
Its ordinary pole bounds and normal precision $a$ are
independent of that degree: both the genus term and the
positive normal degree in the required slope inequalities
scale by $td$ in (D.SC.5). The same inequalities hold at
the last arrow with $\tau^*$, because it preserves slopes,
degrees, the normal ideal and the conormal twists.
Closing the slope recurrence around this permuted cycle
uses that preservation, not the assertion $\tau=1$.

After the same fixed ordinary twist, each residual bundle
is a quotient of $\mathcal O_{C_V}^{r+1}$ with the fixed
determinant kernel line. Its generators lift through the
chosen normal precision. Applying the actual filtered
Cech construction and the finite quotient-Hom complexes
gives at most $A\,td\,r^2$ coordinates per one-step stage.
Only $a$ ordinary normal layers occur. Reference line
cochains and their linear cohomology splittings are fixed
coefficients for each discrete determinant list.

We can use the pullbacks of two fixed affine complements
on $C$ for these Cech computations; they are affine and
are preserved by $\tau$. Thus the last pullback is an
actual linear operation on the fixed complexes. It adds
neither variables nor factors of variable degree. The
normal perturbation words have length at most $a$,
internal width two, fixed-length compositions and only
one $p$th power per arrow. The complete polynomial
presentation consequently has
\[
     N\leq A\,e\,td\,r^2=A ftr^2,
       \qquad \deg(\text{each equation})\leq D_0,       \tag{D.SC.7}
\]
with $A,D_0$ fixed for $R,p$. Its open conditions are the
same actual subbundle, slope and inverse conditions as
in the fixed-line theorem. This is an all-coefficient
presentation of objects and arrows, not a tangent bound.

The proof of Proposition~\ref{D:finite-jet-unit-Isom}
also applies to this final deck pullback. In the normal
contraction estimate, $\tau^*$ preserves order, so the
one-step inequality remains $j\mapsto pj-2B$. The map
and inverse jets, their specified lifting, and their
algebraization use the same fixed ordinary precision.
After arbitrary Artinian coefficient extension, the
Lang frames of (D.SC.3), together with their finite
semilinear deck descent, are finite etale. Nilpotent
invariance therefore gives exactly the same uniqueness
and lifting of isomorphisms. The represented class
fibres are open and closed. Notice that (D.SC.3) is
used for effectivity, not expanded as an equation of
degree growing with $d$ in (D.SC.7).

By Lemma~\ref{D:weighted-affine-component-bound}, the
number of unit classes in each such open family is at
most $(N+1)D_0^N\leq\exp(A_1ftr^2)$. Multiply by
$(r+1)^e$ and (D.SC.6). Since $e\leq f$, $t\geq1$
and $\log(r+1)=O(r^2)$, their product is
$q^{A_R'tr^2}$ after enlarging a constant. Lemma
\ref{D:twisted-deck-Frobenius-semilinear-model} identifies
the counted isomorphisms with the conjugacies stated in
(D.SC.4). This proves the claim, also when $d=1$.
\end{proof}

\begin{theorem}[Proved: bounded generators for unramified field-permutation crowns]
\label{D:unramified-field-permutation-crowns}
Fix $\kappa>0$. Let $S$ be a nonabelian finite simple
group, embedded faithfully as a subgroup of
$\operatorname{GL}_r(\mathbb F_q)$ stable under entrywise
$\mathcal G_q$, and assume
\[
                    r^2\log q\leq\kappa\log|S|.       \tag{D.SC.8}
\]
Let $P\leq\mathcal G_q^t\rtimes\operatorname{Sym}(t)$
act on $A=S^t$ by its field automorphisms and permutations.
Assume the permutation action is transitive and the
induced map $P\to\operatorname{Out}(A)$ is injective.
Set $L=A\rtimes P$, with socle $A$. If a continuous
surjection
\[
              \Gamma\twoheadrightarrow L_k=A^k\rtimes P
                                                        \tag{D.SC.9}
\]
has its quotient $\beta:\Gamma\twoheadrightarrow P$
factoring through $\pi_1(Y)$, then, for
$c=\max\{1,\lceil\kappa(A_R'+1)\rceil\}$,
\[
             k\leq|A|^c,
        \qquad d(L_k)\leq\max\{2,2g,c+2\}.            \tag{D.SC.10}
\]
The same bound works for all displayed $S,q,t,k,P,\beta$.
This includes field automorphisms of unbounded order.
It excludes quotients $\beta$ ramified over the exceptional
curve and outer actions not of the prescribed field form.
\end{theorem}

\begin{proof}
Let $H\subset\Gamma$ be the preimage of a factor
stabilizer. It has index $t$ and corresponds to an etale
resolution cover $Y_H\to Y$. Its action on the first
factor is a character $\chi:H\to\mathcal G_q$ that
extends over $Y_H$, since $\beta$ extends over $Y$.

We spell out the twisted nonabelian Shapiro map. Put
$X=\Gamma/H$. Choose coset representatives $\gamma_x$,
with $\gamma_1=1$, and identify the copy of $S$ at $x$
with the base copy using the given field transport of
$\beta(\gamma_x)$. In these coordinates the $\Gamma$
action on $A=S^X$ is
\[
 (\gamma a)_x=\chi(h(\gamma,x))(a_{\gamma^{-1}x}),
 \qquad
 h(\gamma,x)=\gamma_x^{-1}\gamma\gamma_{\gamma^{-1}x}\in H.
                                                        \tag{D.SC.11}
\]
The identity
$h(\gamma\delta,x)=h(\gamma,x)h(\delta,\gamma^{-1}x)$
checks the action. Lifts of $\beta$ to $L$, modulo
$A$-conjugation, are $H^1(\Gamma,A)$ for this action.
For a cocycle $z:H\to S$ with
$z(hh')=z(h)\chi(h)(z(h'))$, the formula
\[
                    c_z(\gamma)_x=z(h(\gamma,x))       \tag{D.SC.12}
\]
is a $\Gamma$-cocycle. Conversely evaluate a cocycle at
the base coset and restrict to $H$. If
$a_x=c(\gamma_x)_x$, the cocycle identity gives
\[
 c_z(\gamma)_x=a_x^{-1}c(\gamma)_x
                      \chi(h(\gamma,x))(a_{\gamma^{-1}x}).
                                                        \tag{D.SC.13}
\]
Thus evaluation and (D.SC.12) induce inverse bijections
on the two nonabelian cohomology sets. All choices have
finite index, so continuity is preserved.

The $k$ coordinate lifts in (D.SC.9) are distinct in
$H^1(\Gamma,A)$. Otherwise an $A$-conjugacy between
two would put their image in a proper conjugate
diagonal over $P$, contrary to surjectivity. Under
(D.SC.12), each class is represented by a homomorphism
$h\mapsto(z(h),\chi(h))$ to
$S\rtimes\mathcal G_q$. The fixed embedding into
$\operatorname{GL}_r(\mathbb F_q)\rtimes\mathcal G_q$
is compatible with its field action. An ambient
$\operatorname{GL}_r(\mathbb F_q)$-conjugacy class
contains at most $|\operatorname{GL}_r(\mathbb F_q)|$
classes under $S$: every conjugating matrix determines
at most one homomorphism. No surjectivity to the ambient
linear group is required for this crude bound. Hence
Theorem~\ref{D:etale-resolution-semilinear-count} gives
\[
 k\leq q^{(A_R't+1)r^2}
       \leq |S|^{\kappa(A_R'+1)t}\leq |A|^c.           \tag{D.SC.14}
\]

Transitivity makes $A$ minimal normal in $L$. The
injectivity of the outer action makes its centralizer
trivial, so $A$ is the unique minimal normal subgroup.
It is nonabelian and is not contained in the Frattini
subgroup. A maximal subgroup not containing $A$ is
core-free; thus $L$ is monolithic primitive, as required
by the crown theorem. Moreover $P$ is a quotient of
$\pi_1(Y)\simeq\pi_1(C)$, so $d(P)\leq2g$.

Apply the already checked crown estimate in
Theorem~\ref{D:unramified-permutation-linear-crowns},
namely \citet[Proposition 2.6, equations (2.1)--(2.2),
pp.~12--13]{LucchiniMarionTracey2020D}. For
$s\geq\max\{2,d(P)\}$ it gives
\[
 \psi_L(s)\geq\frac{|A|^{s-1}}{2t|\operatorname{Out}(S)|}
                  \geq |A|^{s-2},
 \qquad k\leq\psi_L(s)\Longrightarrow d(L_k)\leq s.
\]
Take $s=\max\{2,2g,c+2\}$ and use (D.SC.14).
This proves (D.SC.10), retaining the classification
dependence of that finite-group input.

For a concrete unbounded-rank scope, the existing natural
linear/order bound applies to
$S=\operatorname{SL}_{p^a}(\mathbb F_{p^f})$ whenever
this group is nonabelian simple. Its scalar centre is
trivial since $\gcd(p^a,p^f-1)=1$, it is stable under
all field automorphisms, and the previously checked
order estimate gives (D.SC.8) with $\kappa=4$.
The nonsimple small parameters are excluded by the
stated hypothesis. This example is a scope consequence
of the same crown bound, not another movement event.
\end{proof}


% Separate post-CP20 proved input. Root/B/C/D complete actual-file audits passed.
\subsection{Tame radial covers with one exceptional curve}

Let $(R,\mathfrak m,k)$ be a complete excellent normal local
$k$-algebra of dimension two, with algebraically closed coefficient
and residue field $k$ of characteristic $p>0$. Fix a projective
regular resolution $\pi:Y\to\operatorname{Spec}R$, an isomorphism
over $U=\operatorname{Spec}R-\{\mathfrak m\}$, whose reduced
exceptional fibre is one smooth connected curve $C$. Write
$g=g(C)$ and $C^2=-\delta<0$. A connected finite etale cover
$V\to U$ is \emph{radially tame} if its function-field extension
is tame at every extension of the discrete valuation defined
by the generic point of $C$. This condition concerns the cover
$V\to U$, and does not impose tameness on subsequent local
systems on $V$.

\begin{lemma}[Proved: the normalization of a radially tame cover]
\label{D:radial-tame-cover-one-curve}
Let $V\to U$ be a connected radially tame finite etale cover
of degree $t$. Let $h:Y'\to Y$ be the normalization in $k(V)$,
and let $R'$ be the normalization of $R$ in $k(V)$. Then $R'$
is a complete normal local domain with residue field $k$,
$Y'\to\operatorname{Spec}R'$ is a projective regular resolution,
and its puncture is exactly $V$. Its reduced exceptional fibre
is one smooth connected curve $C'$. There are integers $e,u\geq1$
such that $e$ is prime to $p$, the map $C'\to C$ is finite
etale of degree $u$, and
\[
 h^*C=eC',\qquad t=eu,\qquad
 g(C')-1=u(g-1),\qquad (C')^2=-\frac{u\delta}{e}.
                                                        \tag{D.RT.1}
\]
In particular $u\delta/e$ is a positive integer. The map
$h$ is finite locally free. No Galois hypothesis is needed.
\end{lemma}

\begin{proof}
Excellence makes both normalizations finite. We first determine
the normalization at every point of $C$, rather than only at its
generic point. Abhyankar's lemma for a regular divisor gives,
etale locally on $Y$, a disjoint union of the following models
on the complement of $C$:
\[
 \operatorname{Spec}A[1/z][s]/(s^{e_i}-z),\qquad
        e_i\text{ invertible in }A,
                                                        \tag{D.RT.2}
\]
where $z$ is a regular parameter cutting out $C$.
This is the full regular-divisor assertion and proof of
\citep[Lemma 58.31.5, Tag 0EYG]{StacksRadialTameGeometryB};
its proof first works over a strict henselization, kills
ramification by a Kummer root, and uses purity to remove
possible branching at the closed point. It therefore applies
to the present assumption of tameness at the sole boundary
valuation. The same local description is proved in
\citet[Expose XIII, Propositions 5.1--5.2 and
Corollary 5.3, original pp.~428--430]{SGA1}.

The normalization of each model in (D.RT.2) is
$A[s]/(s^{e_i}-z)$. Indeed this ring is regular: its quotient
by $s$ is the regular ring $A/(z)$, and the nonzerodivisor
$s$ replaces the regular parameter $z$. It is consequently
normal, finite free of rank $e_i$ over $A$, and has the
specified generic algebra. Normalization commutes with this
etale localization. It follows that $Y'$ is regular, $h$ is
finite locally free, and the reduced inverse image of $C$
is smooth, with disjoint smooth components. On these charts
its map to $C$ is the identity, so globally it is finite etale
over $C$. The local ramification index is constant on each
connected component of this reduced inverse image.

The ring $R'$ is finite over the complete henselian ring $R$.
A finite algebra over a henselian local ring is a product
of local algebras; since $R'$ is a domain it has one factor.
It is complete, and its finite residue extension of the
algebraically closed field $k$ is $k$ itself. Since $Y'$ is
normal, integral and proper over $\operatorname{Spec}R$, its
Stein ring is $R'$: its global functions form a finite normal
$R$-subalgebra of $k(V)$ with fraction field $k(V)$, hence
the full integral closure. Stein factorization gives
geometrically connected fibres over $\operatorname{Spec}R'$
\citep[Theorem 37.53.4, Tag 03H0, complete proof]
{StacksRadialTameGeometryB}. The map to this finite Stein
base is projective, because the original composite to
$\operatorname{Spec}R$ is projective.

Over $U$, normalization is the given finite etale normal
cover $V$. The unique closed point of $\operatorname{Spec}R'$
is the full inverse image of $\mathfrak m$, so
\[
 V=\operatorname{Spec}R'-\{\mathfrak m'\},\qquad
 Y'-V=(h^{-1}C)_{\mathrm{red}}
 \quad\text{as reduced closed sets}.                   \tag{D.RT.3}
\]
The latter is the support of a connected proper fibre.
Its smooth components were proved disjoint, hence there
is exactly one, denoted $C'$. Thus $C'$ is smooth and
connected, and the ramification index is a single integer
$e$ prime to $p$. The Kummer models now give $h^*C=eC'$.

Let $u=\deg(C'/C)$. The fibre of the finite locally free
map $h$ over a geometric point of $C$ has $u$ local factors,
each of length $e$, so $t=eu$. Etale Riemann--Hurwitz gives
$g(C')-1=u(g-1)$
\citep[Section 53.12 and Lemma 53.12.2, Tags 0C1B and 0C1D]
{StacksSemilinearDeterminantsB}. Finally the intersection
projection formula for the finite degree-$t$ map gives
\[
 e^2(C')^2=(h^*C)^2=tC^2=-eu\delta.
                                                        \tag{D.RT.4}
\]
Both surfaces are regular and $C'$ is a proper integral
Cartier curve, so its self-intersection is an integer.
This proves all the assertions.
\end{proof}

\begin{proposition}[Proved: tame generators and twisted determinant classes]
\label{D:one-curve-radial-tame-generators}
Let $\Gamma=\pi_1(U)$ and let $\Gamma^{\mathrm{rt}}$
be its quotient classifying radially tame finite etale
covers. Then
\[
                d(\Gamma^{\mathrm{rt}})\leq 2g+1.
                                                        \tag{D.RT.5}
\]
Let $A$ be a finite abelian group of order prime to $p$,
endowed with a continuous $\Gamma$-action that factors
through $\Gamma^{\mathrm{rt}}$. Inflation is an isomorphism
and one has
\[
 H^1_{\mathrm{cont}}(\Gamma^{\mathrm{rt}},A)
     \xrightarrow{\ \sim\ }H^1_{\mathrm{cont}}(\Gamma,A),
 \qquad |H^1_{\mathrm{cont}}(\Gamma,A)|
                     \leq |A|^{2g+1}.                 \tag{D.RT.6}
\]
In particular, let $V\to U$ be as in
Lemma~\ref{D:radial-tame-cover-one-curve}, with $t=eu$,
and put $H=\pi_1(V)$ and $g'=1+u(g-1)$. Fix $q=p^f$
and a field character
\[
 \chi:H\longrightarrow\operatorname{Gal}
                          (\mathbb F_q/\mathbb F_p)
\]
that is radially tame on the smooth model $(Y',C')$.
It need not extend etale over $Y'$. Then
\[
 |H^1_{\mathrm{cont}}(H,\mathbb F_q^*(\chi))|
       \leq(q-1)^{2g'+1}\leq q^{(2g+3)u}
                                \leq q^{(2g+3)t}.
                                                        \tag{D.RT.7}
\]
These are the semilinear determinant classes for the fixed
field character. If the character is also allowed to vary
subject to radial tameness, it has at most $f^{2g'+1}$
possibilities. Neither assertion requires the image order
of $\chi$ to be prime to $p$.
\end{proposition}

\begin{proof}
We prove the generator bound on every finite quotient.
Let $G$ be such a quotient, with connected Galois tame cover
$V\to U$. The preceding lemma gives a unique exceptional
component $C'$ in its normalization. Therefore the
decomposition group at its generic point is all of $G$.
The exact decomposition--inertia sequence for this
discrete valuation is
\[
        1\longrightarrow I\longrightarrow G
          \longrightarrow\operatorname{Gal}(k(C')/k(C))
          \longrightarrow1.                           \tag{D.RT.8}
\]
Here $I$ is cyclic of order $e$, prime to $p$, by tame
inertia. The residue extension is separable by the local
Kummer description. The decomposition sequence is
surjective onto its residue automorphism group, and
$|G|/|I|=eu/e=u=[k(C'):k(C)]$, so this residue extension
is indeed Galois. These are exactly the finite-case
decomposition and tame-inertia statements of
\citep[Lemma 58.13.5, Tag 0BUA, complete proof]
{StacksRadialTameGeometryB}. In particular the quotient
in (D.RT.8) is the Galois group of the finite etale map
$C'\to C$.

The full profinite fundamental group of a smooth proper
genus-$g$ curve over an algebraically closed field has at
most $2g$ generators, including its characteristic-$p$
quotients. This follows from the surjective specialization
argument of \citet[Theorem X.2.6 and its complete preceding
argument, original pp.~272--273]{SGA1}. Lift such generators
of the residue quotient in (D.RT.8) and add a generator
of the cyclic inertia group. Hence $d(G)\leq2g+1$.

For completeness, a profinite group whose every finite
quotient is generated by at most $a$ elements is itself
generated by at most $a$ elements. In its compact space
of $a$-tuples, the tuples generating each prescribed
finite quotient form a nonempty clopen set. Finitely
many such conditions are simultaneously satisfied by
generating their common finite quotient. Compactness
therefore supplies a tuple satisfying every condition.
Its closed generated subgroup is the whole group.
This proves (D.RT.5) with $a=2g+1$.

Let $K=\ker(\Gamma\to\Gamma^{\mathrm{rt}})$.
The discrete-valuation tame-inertia description identifies
$K$ with the closed normal subgroup generated by the
images of the wild inertia groups for $C$ and their
conjugates. Each such group is pro-$p$
\citep[Lemma 58.13.5, Tag 0BUA]{StacksRadialTameGeometryB}.
The assumed action makes $K$ act trivially on $A$.
Consequently the restriction to $K$ of a continuous
crossed homomorphism $z:\Gamma\to A$ is an ordinary
homomorphism. It vanishes on every wild inertia group,
since $|A|$ is prime to $p$, hence on their closed normal
generated subgroup. One can also check conjugates
directly from the crossed-homomorphism identity and the
trivial $K$-action. Thus $z$ descends uniquely to
$\Gamma^{\mathrm{rt}}$; coboundaries descend as well.
This proves the inflation assertion without imposing a
prime-to-$p$ condition on the action's image.

A crossed homomorphism to $A$ is determined by its values
on the $2g+1$ chosen topological generators. The cocycle
identity determines its values on the generated abstract
subgroup, and continuity determines the remaining values.
There are at most $|A|^{2g+1}$ choices. Quotienting by
coboundaries cannot increase that number, proving
(D.RT.6).

Apply the same result to the complete local surface $R'$
and its sole exceptional curve $C'$. The module
$\mathbb F_q^*(\chi)$ has order $q-1$, prime to $p$,
and its action is tame by assumption. Formula (D.RT.6)
therefore gives the first inequality in (D.RT.7). The
remaining inequalities follow from
$2g'+1=2u(g-1)+3\leq(2g+3)u$ and $u\leq t$.
When $q=2$ the coefficient group is trivial. The
determinant cocycle and the effect of conjugacy are
exactly those in (D.SD.6): for a prescribed field part
$\chi$, the invariant lies in this twisted $H^1$, not
in an arbitrary character group of the cover killing
$\chi$. Finally a homomorphism from the tame group to
the cyclic group of order $f$ is determined by the same
$2g'+1$ generator values, giving $f^{2g'+1}$. This
argument allows $p$-divisibility in $f$ and in the image
order of $\chi$.
\end{proof}


% Separate post-CP20 proved packet; the audited tame covering geometry is an explicit dependency.
\subsection{A hereditary count across tame radial ramification}

Retain the complete one-smooth-exceptional-curve setting of
Theorem~\ref{D:etale-resolution-hereditary-linear-count}, and write
$C^2=-\delta<0$, $g=g(C)$ and $U=Y-C$. A finite etale cover
$U'\to U$ is called \emph{radially tame} if the normalization
of $Y$ in that cover is tamely ramified at the generic point
of $C$. No tameness condition is imposed on representations
of $\pi_1(U')$ in the following statement.

\begin{theorem}[Proved: index-linear hereditary counts for radially tame covers]
\label{D:tame-resolution-hereditary-linear-count}
There is a constant $A_R^{\mathrm{tame}}$ such that every
connected radially tame finite etale cover $U'\to U$ of
degree $t$, every $q=p^f$, and every $r\geq1$ satisfy
\[
 \#\bigl(\operatorname{Hom}_{\mathrm{cont}}
       (\pi_1(U'),\operatorname{GL}_r(\mathbb F_q))
                         /\operatorname{GL}_r(\mathbb F_q)\bigr)
                \leq q^{A_R^{\mathrm{tame}}t r^2}.
                                                        \tag{D.TH.1}
\]
The constant is independent of the radial ramification index,
the residue degree, and the cover. The normalization is the
smooth one-curve model furnished by
Lemma~\ref{D:radial-tame-cover-one-curve}; arbitrary covers
wildly ramified over $C$ are outside this assertion.
\end{theorem}

\begin{proof}
\emph{The two geometric scales.}
By Lemma~\ref{D:radial-tame-cover-one-curve}, the normalization
$Y'\to Y$ is regular and projective over its complete normal
local Stein ring. Its reduced exceptional fibre is one smooth
connected curve $C'$. There are integers $e,u\geq1$ such that
$e$ is prime to $p$, $C'\to C$ is etale of degree $u$, and
\[
 t=eu,\qquad g'-1=u(g-1),\qquad
 (C')^2=-d',\qquad d'=u\delta/e\in\mathbb Z_{>0}.
                                                        \tag{D.TH.2}
\]
Here $g'=g(C')$. In particular $g=0$ forces $u=1$.
The proof below uses the two scales $u$ and $e$ separately;
replacing every parameter by $t$ before counting would lose
the asserted index dependence.

Choose minimal ordinary bundle lattices $F_i$ on $Y'$ for
the $f$ one-step Frobenius stages of the given rank-$r$
local system. The untwisted case of
Lemma~\ref{D:twisted-connected-anchor-minimal-slopes} gives
\[
       g'-1\leq\mu_{\min}(F_i|_{C'})
       \leq\mu_{\max}(F_i|_{C'})\leq g'-1+d'.
                                                        \tag{D.TH.3}
\]
Their endpoints have absolute value at most a fixed constant
times $u$, since $d'\leq u\delta$. Choose once an integer
$h$ with $h\delta\geq g+2$, and put
$F_i^+=F_i(-heC')$. Its residual minimum slope is at least
$u(g-1+h\delta)>2g'-1$. Thus $F_i^+|_{C'}$ is globally
generated, and its tensor product with any nonnegative
power of the conormal line has vanishing $H^1$.

The one-step Frobenius slope estimate (D.TD.14) places the
pullbacks of (D.TH.3), including this common twist, in
intervals with endpoints of absolute value at most $a_0u$,
where $a_0$ depends only on $g,\delta,p,h$. The positive
and negative Hom-layer arguments in
Proposition~\ref{D:twisted-period-independent-normal-depth}
therefore give uniform thresholds as follows:
\[
 \begin{split}
 j d'>2a_0u&\quad\Longrightarrow\quad
          H^0\text{ of the corresponding negative Hom layer is }0,\\
 j d'>2a_0u+2g'-2&\quad\Longrightarrow\quad
          H^1\text{ of the corresponding positive Hom layer is }0.
 \end{split}                                            \tag{D.TH.4}
\]
These implications are the slope-gap and Serre-duality
calculations for the actual source and target bundles, not
a tensor-semistability assertion. In view of (D.TH.2), fixed
multiples of $e$ suffice for both thresholds. Thus all
meromorphic one-step maps and their inverses have pole bounds
$Be$, and all prescribed-map lifting and stabilization use
ordinary normal depth at most $Ne$, for fixed $B,N$.
After enlarging constants for the finitely many pole twists
and inverse-product diagrams, choose a single integer $a$
such that the whole finite unit presentation is on $aeC'$.
All its cutoffs, including those for one Frobenius pullback,
are fixed multiples of $e$.

For precision, the finite unit-map correction has the same
normal-order inequality as before,
\[
                    j\longmapsto pj-2Be.              \tag{D.TH.5}
\]
Taking the starting precision to be a sufficiently large
fixed multiple of $e$ makes it strictly increasing and
unbounded. The paired inverse jets and their true pole
orders are included in the choice of $a$. Hence
Lemma~\ref{D:meromorphic-unit-jet-contraction} and
Proposition~\ref{D:finite-jet-unit-Isom} apply at this
precision. Their all-base Hom vanishing follows from the
same two-term proper-curve complexes and fibrewise split
surjectivity; increasing $e$ introduces no extra hypothesis
on a coefficient ring.

\emph{One determinant character and a finite line list.}
Set $n=q-1$. The determinant character has at most
\[
 |H^1(U',\mu_n)|\leq n^{2g'+1}
                         \leq q^{(2g+3)u}             \tag{D.TH.6}
\]
possibilities. This is the regular one-divisor localization
and proper strict-local base-change argument in (D.HC.5),
applied to $(Y',C')$; $n=1$ gives the trivial group.
Only one character is chosen: its $p^i$th powers determine
the punctured determinants of all $f$ stages.

For any fixed such punctured line, its extensions to the
regular surface $Y'$ differ by $\mathcal O_{Y'}(bC')$.
Their restriction degrees differ by the step $d'$. The
interval for the determinant in (D.TH.3) has length $rd'$.
Consequently there are at most $r+1$ extension lines per
stage. Fix one list, of at most $(r+1)^f$ possibilities.
After the twist $-heC'$, its residual determinants $L_i$
have degrees of absolute value $O(ur)$, with a constant
independent of $e$. These are discrete choices for the
geometric unit objects. No arbitrary relative reflexive
hull or moving Picard chart is asserted.

\emph{The actual fixed-line presentation and its weights.}
The residual bundle at each stage has a presentation
\[
 0\longrightarrow L_i^{-1}\longrightarrow
       \mathcal O_{C'}^{r+1}\longrightarrow F_i^+|_{C'}
       \longrightarrow0.                              \tag{D.TH.7}
\]
The $r+1$ generating sections lift through the normal
layers because of the positive-conormal $H^1$ vanishing
already established. Thus the source remains the fixed
formal bundle $\mathcal O^{r+1}$; only the kernel line
deforms. Choose fixed reference lifts of the $L_i^{-1}$.
Apply the actual all-coefficient-ring presentation in
Lemma~\ref{D:weighted-fixed-residual-bundle-deformations}
and Proposition~\ref{D:weighted-fixed-determinant-normal-presentation},
in its ordinary, untwisted specialization.

We give the quantitative substitution explicitly, rather
than infer it from deformation-space dimensions. Every
termwise Hom line occurring in the quotient complexes
$[L_i^{-1}\to\mathcal O^{r+1}]$, their one-step Frobenius
pullbacks, their pole twists, and their normal layers of
order $j\leq ae$ has degree $O(ur)$. Indeed a normal
or pole twist of order $O(e)$ changes a line degree by
$O(e d')=O(u)$, and the fixed one-step Frobenius contributes
only the factor $p$. The genus is $O(u)$. Curve duality
and Riemann--Roch bound the total $H^0,H^1$ dimension of
each such line by $O(ur)$; for an $\mathcal O$ term
it is $O(u)$. The four termwise Hom blocks, before the
fixed pole and conormal twists, are
\[
 \mathcal O^{(r+1)^2},\qquad L_i^{\oplus(r+1)},\qquad
 (L_j^{-1})^{\oplus(r+1)},\qquad L_iL_j^{-1}.
\]
For a Frobenius source, replace $L_i$ by $L_i^{\otimes p}$.
The quadratic multiplicity thus occurs only in the
$O(u)$-cohomology block; the $O(ur)$-cohomology blocks
have multiplicity at most $r+1$. The fixed finite diagram therefore
uses at most $Afur^2$ scalar coordinates at each positive
normal order $1\leq j\leq ae$, and at most $Afur^2$
additional residual coordinates. The constant $A$ includes
all maps, inverse maps, cycles and boundary witnesses.

Assign weight $j$ to a variable in positive normal order
$j$, and weight $e$ to every residual variable. The latter
include the residual generating sections and residual
map or boundary coefficients. All reference cochains and
all chosen cohomology contractions are fixed $k$-linear
maps. Their entries are constants, even if their chosen
rational representatives have large poles; they are not
evaluated on moving divisor coordinates.

The actual normal-transition recursion consumes normal
order: a product contributing to order $j$ has sum of its
positive layer orders at most $j$. The same holds after
the finite Cech homological perturbation transfer. The
two-term quotient resolution has internal width two in
its Hom complex, so a transferred word uses at most two
internal differential factors. Thus only a fixed number
of residual factors, each of weight $e$, can occur in
addition to the consumed normal orders. Composition and
inverse-product diagrams have a fixed number of arrows.
Pole clearing shifts the relevant cutoff by $O(e)$;
one absolute $p$-Frobenius multiplies weights by at most
$p$. It follows that every defining polynomial has
weighted degree at most $Ke$, where $K$ is independent
of $e,u,r,f$:
\[
 \begin{gathered}
  \#\{\text{variables of weight }j\}\leq Afur^2
           \quad(1\leq j\leq ae),\\
  \#\{\text{additional variables of weight }e\}
           \leq Afur^2,
  \qquad \deg_w(\text{each equation})\leq Ke.
 \end{gathered}                                         \tag{D.TH.8}
\]
We enlarge $K$ so that $K\geq a$. Variables from the two
blocks may have the same numerical weight; (D.TH.8)
keeps their separate multiplicity counts. Positive-order
transition variables have weight $j$, not $e+j$: adding
$e$ to every positive layer would not give this bound.

The cited construction gives one affine polynomial
presentation followed by open conditions for each fixed
line list. Its bundles, quotient maps and Hom witnesses
are actual over every coefficient algebra; in particular
the weights do not stand for a merely formal moduli
dimension. The open conditions impose residual local
freeness, the slope windows, and the specified inverse
conditions. The uniform finite unit-jet comparison above
identifies its unit objects with the required local
systems. By specified-isomorphism lifting and nilpotent
invariance of finite etale frames, the represented
isomorphism classes are open and closed, exactly as in
Proposition~\ref{D:finite-jet-unit-Isom}.

\emph{The weighted component calculation.}
Pad each positive layer and the additional block by
unused variables set equal to zero, until they each
contain $b=\lceil Afur^2\rceil$ variables. This changes
neither the scheme nor its components and respects the
weight bound. Put $J=ae$ and $M=b(J+1)$. Lemma
\ref{D:weighted-affine-component-bound} bounds the
number of irreducible components by
\[
 (M+1)\frac{(Ke)^{b(J+1)}}{e^b(J!)^b}
     =(M+1)K^b\left(\frac{(Ke)^J}{J!}\right)^b.
                                                        \tag{D.TH.9}
\]
Writing $\mathrm e$ for the base of the natural logarithm,
$J!\geq(J/\mathrm e)^J$ gives
\[
 \log\!\left((M+1)\frac{(Ke)^{b(J+1)}}{e^b(J!)^b}\right)
 \leq\log(M+1)+b\log K+bJ\log(\mathrm e K/a)
                         =O(feur^2)=O(ftr^2).          \tag{D.TH.10}
\]
Thus every fixed-list family represents at most
$\exp(A_1ftr^2)$ unit classes. Open restrictions do not
increase the number of irreducible components. No bound
on the number of defining equations is required by the
weighted component lemma.

Finally multiply by (D.TH.6) and by $(r+1)^f$.
Since $u\leq t$, $\log(r+1)=O(r^2)$ and $q=p^f$,
the result is (D.TH.1), after increasing one constant.
The unit-frame comparison counts actual
$\mathbb F_q$-linear local systems and their
$\mathbb F_q$-linear isomorphisms. The only imposed
ramification condition concerns the covering $U'\to U$;
the represented local systems themselves may be wildly
ramified along $C'$.
\end{proof}


% Proved restricted scope extension; all four actual-file audits passed.
\subsection{Field-permutation crowns with tame radial quotient}

We continue in the complete one-smooth-exceptional-curve
setting, with $C^2=-\delta$, $g=g(C)$ and $\Gamma=\pi_1(U)$.
Let $Q^{\mathrm{rad}}$ denote the quotient of $\Gamma$ by
the closed normal subgroup generated by wild inertia above
the generic point of $C$ on the regular resolution $Y$.
This auxiliary quotient is attached to the exceptional
divisor on $Y$; it is not a redefinition of the original
boundary quotient in D.1.

\begin{theorem}[Proved: index-linear semilinear counts on tame radial covers]
\label{D:tame-resolution-semilinear-count}
There is a constant $A_R^\dagger$ such that the following
holds. Let $U_H\to U$ be a connected radially tame finite
etale cover of degree $t$, let $Y_H$ be the normalization
of $Y$ in it, and let $C_H$ be its unique exceptional
curve. Put $H=\pi_1(U_H)$ and $q=p^f$. Fix a character
$\chi:H\to\mathcal G_q$ whose field cover is tamely
ramified along $C_H$. Then
\[
 \#\left\{\rho:H\to\operatorname{GL}_r(\mathbb F_q)
                         \rtimes\mathcal G_q:
          \operatorname{pr}_{\mathcal G_q}\rho=\chi\right\}
           /\operatorname{GL}_r(\mathbb F_q)
                      \ \leq q^{A_R^\dagger t r^2}.     \tag{D.TS.1}
\]
The notation on the left means the cardinality of the
indicated conjugacy-class set. No bound on the image
order of $\chi$ is imposed. The matrix part of $\rho$
may be wildly ramified along $C_H$.
\end{theorem}

\begin{proof}
Write $a=|\operatorname{im}\chi|$, $\nu=f/a$ and
$q_0=p^\nu$. The degree-$a$ field cover $V\to U_H$
has composite degree $ta$ over $U$. Tameness is stable
under this tower of extensions of the boundary discrete
valuation rings: the ramification indices remain prime
to $p$ and the residue extensions remain separable.
Thus $V\to U$ is radially tame. By
Lemma~\ref{D:radial-tame-cover-one-curve}, its normalization
$Y_V\to Y$ is regular with one smooth exceptional curve
$C_V$. If $E$ and $v$ are respectively its radial
ramification index and residue degree, then
\[
 Ev=ta,\qquad g(C_V)-1=v(g-1),\qquad
                    -C_V^2=v\delta/E.                 \tag{D.TS.2}
\]
The cyclic deck action over $U_H$ extends to $Y_V$ by
normalization and preserves $C_V$ and every power of
its normal ideal.

The determinant module $A_\chi=\mathbb F_q^*$ has
the prescribed field action. Because $\chi$ is tame,
wild inertia acts trivially on $A_\chi$. A degree-one
cocycle restricts to an ordinary homomorphism on the
closed normal wild subgroup; this homomorphism vanishes
on every pro-$p$ inertia group, since $|A_\chi|=q-1$
is prime to $p$, and hence vanishes on their closed
normal span. Inflation therefore identifies
\[
 H^1(H^{\mathrm{rad}},A_\chi)=H^1(H,A_\chi),
\]
where $H^{\mathrm{rad}}$ is the radial tame quotient for
$(Y_H,C_H)$. Proposition
\ref{D:one-curve-radial-tame-generators} gives
$d(H^{\mathrm{rad}})\leq2g(C_H)+1$. Crossed
homomorphisms are determined by their generator values,
so the number of determinant classes is at most
\[
          (q-1)^{2g(C_H)+1}\leq q^{(2g+3)t}.            \tag{D.TS.3}
\]
The genus formula for $C_H\to C$ justifies the last
inequality. For $q=2$ the determinant group is trivial.
This argument includes field-action orders divisible
by $p$; it only requires tame field ramification.

Fix one class in (D.TS.3). Lemma
\ref{D:twisted-deck-Frobenius-semilinear-model} describes
the desired objects on $V$ by
$\Phi:\tau^*F_{q_0}^*E_0\xrightarrow{\sim}E_0$.
Use $\nu$ one-step Frobenius stages, with the deck
pullback on the final arrow. The restricted determinant
class fixes the punctured determinant at every stage.
The same extension-degree argument as in
Theorem~\ref{D:tame-resolution-hereditary-linear-count}
leaves at most $(r+1)^\nu$ lists of residual lines.

We apply the actual construction in that theorem's
proof to the composite tame cover in (D.TS.2). We give
the substitution explicitly; applying its final ordinary
$q$-count to the degree-$ta$ cover would lose a factor
$a$ and would not prove (D.TS.1). The residual slopes
and line degrees are $O(vr)$, ordinary pole and jet
depths are $O(E)$, and each normal-layer Hom block has
dimension $O(vr^2)$. The deck pullback preserves the
slope windows and normal orders. The pullbacks of two
affine complements on $C$ are deck-invariant, so it is
a fixed linear operation on the Cech complexes.

Consequently the fixed-line all-ring presentation has
at most $b=A\nu vr^2$ variables in each positive
normal weight $j$ for $1\leq j\leq J=BE$, and at
most $b$ residual variables of weight $E$. Every
equation has weighted degree at most $KE$, with fixed
$A,B,K$ depending only on the original resolution and
$p$. As in (D.TH.8), the bounded internal width permits
only a fixed number of residual factors per transferred
word, whereas positive normal orders add. The final
deck arrow changes neither fact. All bundles, maps,
boundary witnesses and open conditions are the actual
ones in that fixed-line construction.

After rounding and padding, put $M=b(J+1)$. The
weighted component bound is
\[
 (M+1)\frac{(KE)^{b(J+1)}}{E^b(J!)^b}
             \leq\exp(A_1\nu vE r^2)
             =\exp(A_1ftr^2).                         \tag{D.TS.4}
\]
The first inequality is the full factorial calculation
(D.TH.9)--(D.TH.10); the last equality uses $Ev=ta$
and $\nu=f/a$. Thus the field order is absorbed
without a logarithmic or linear loss in $a$.

For completeness the class-component comparison also
uses the twisted cycle, rather than forgetting descent.
The finite-unit-Isom proof has pole and precision
constants that are fixed multiples of $E$; its
one-step contraction remains $j\mapsto pj-2B_0E$.
Deck pullback preserves that order. All specified-map
lifting and paired inverse jets are therefore valid at
the same precision. The $a$-fold composite of $\Phi$
defines an ordinary $F_q$-unit structure. Its finite
etale frame sheaf has the genuine semilinear deck
cocycle of Lemma
\ref{D:twisted-deck-Frobenius-semilinear-model}, also
over Artinian coefficient bases. Nilpotent invariance
and algebraization give open and closed actual unit
classes, as in the proof of
Theorem~\ref{D:etale-resolution-semilinear-count}.
The $a$-fold composite is used for this comparison;
it is not expanded as an additional high-degree
equation in (D.TS.4).

Finally multiply (D.TS.4) by (D.TS.3) and by
$(r+1)^\nu$. Since $\nu\leq f$, the result is
$q^{A_R^\dagger tr^2}$ after increasing a fixed
constant. The equivalence keeps exactly the
$\mathbb F_q$-linear conjugacies in (D.TS.1).
\end{proof}

\begin{corollary}[Proved: a uniform generator bound for tame field-permutation crowns]
\label{D:tame-field-permutation-crowns}
Use the data $S\subset\operatorname{GL}_r(\mathbb F_q)$,
$\kappa$, $P\leq\mathcal G_q^t\rtimes\operatorname{Sym}(t)$,
$A=S^t$ and $L=A\rtimes P$ of
Theorem~\ref{D:unramified-field-permutation-crowns},
including transitivity and injectivity into
$\operatorname{Out}(A)$. Suppose a surjection
$\Gamma\twoheadrightarrow L_k$ induces a quotient
$\beta:\Gamma\twoheadrightarrow P$ factoring through
$Q^{\mathrm{rad}}$. Then
\[
 k\leq|A|^{c_R},\qquad
 d(L_k)\leq\max\{2,2g+1,c_R+2\},
 \qquad
 c_R=\max\{1,\lceil\kappa(A_R^\dagger+1)\rceil\}.      \tag{D.TS.5}
\]
In particular the quotient $\beta$ may be tamely
ramified along the exceptional curve; it need not
extend etale over $Y$. Quotients with wild radial
ramification and other outer-action types remain
outside the conclusion.
\end{corollary}

\begin{proof}
Let $H$ be the preimage of a point stabilizer for the
transitive permutation action. Its index is $t$,
and its cover is radially tame because it is a
subcover of the tame $P$-cover. The field action
$\chi:H\to\mathcal G_q$ is likewise tame relative
to its normalization $Y_H$: it factors through
the restriction of that same tame cover. Stability
of tameness under base change, subextensions and
normal closure gives the stated boundary condition.

The explicit twisted Shapiro bijection
(D.SC.11)--(D.SC.13) applies without change to this
group action. Distinct crown coordinates give
distinct classes, and the faithful ambient linear
embedding multiplies their count by at most
$q^{r^2}$. Theorem
\ref{D:tame-resolution-semilinear-count} therefore
gives
\[
 k\leq q^{(A_R^\dagger t+1)r^2}
       \leq |S|^{\kappa(A_R^\dagger+1)t}
       \leq |A|^{c_R}.
\]
Proposition~\ref{D:one-curve-radial-tame-generators}
also gives $d(P)\leq2g+1$. The same monolithic
primitive crown estimate (D.HC.15), with
$s=\max\{2,2g+1,c_R+2\}$, proves (D.TS.5).
No independent bound on the number of possible
$P$-quotients is required: the constants work for
each fixed quotient $\beta$ and hence for all
surjections in the stated family.
\end{proof}


% Separate unassembled candidate for the multicomponent count, not a completed count.
\subsection{Brauer classes and residual blocks on an exceptional tree}

\begin{lemma}[Proved: tree residues and the actual coefficient gerbe]
\label{D:SNC-tree-Brauer-geometry}
Let $R$ be a complete excellent normal local surface over an
algebraically closed coefficient and residue field $k$ of
characteristic $p>0$. Let $Y\to\operatorname{Spec}R$ be a
projective regular resolution, an isomorphism over the puncture
$U$, and let $E=\bigcup_{i=1}^v E_i$ be its reduced exceptional
divisor. Suppose $E$ is strict normal crossings and its dual
graph is a tree. Write $g_i=g(E_i)$ and
$N=(-(E_i\cdot E_j))_{ij}$.

For every $n$ prime to $p$, residue gives an isomorphism
\[
 \operatorname{Br}'(U)[n]\simeq
       \bigoplus_i H^1(E_i,\mathbb Z/n),\qquad
 |H^2(U,\mu_n)|=|\operatorname{coker}(N\bmod n)|
                             n^{2\sum_i g_i}.           \tag{D.TR.1}
\]
In particular the latter order is at most
$|\det N|n^{2\sum_i g_i}$. All secondary residues at the
nodes vanish.

A class $\beta$ of prime-to-$p$ order $m$ is killed by a
connected cyclic degree-$m$ cover $Y'\to Y$ which is etale
on all of $Y$. Its restriction to $U$ is the puncture of a
finite complete local normal extension of $R$. The exceptional
divisor upstairs need not have a tree as its dual graph.

Let $\alpha\in H^2(U,\mu_n)$ have Brauer image $\beta$,
and let $m_i$ be the order of the $i$th residue, with
$m_i=1$ for zero residue. Then $m_i\mid m\mid n$, and
the actual class $\alpha$ extends to
\[
       \mathcal X=\sqrt[m_1]{(Y,E_1)}\times_Y\cdots
                         \times_Y\sqrt[m_v]{(Y,E_v)}.  \tag{D.TR.2}
\]
At a node $E_i\cap E_j$, the restriction of its gerbe to
$B(\mu_{m_i}\times\mu_{m_j})$ has zero geometric Brauer
class. Its inertia is therefore a commutative diagonalizable
extension of $\mu_{m_i}\times\mu_{m_j}$ by $\mu_n$.
No nontrivial node commutator is discarded in this statement:
the tree hypothesis forces that commutator to vanish.
\end{lemma}

\begin{proof}
Put $\Lambda=\mathbb Z/n$. Normalization of the proper
nodal curve $E$ gives
\[
 0\longrightarrow\Lambda_E\longrightarrow
   \bigoplus_i\Lambda_{E_i}\longrightarrow
   \bigoplus_{x\in\operatorname{Nodes}(E)}\Lambda_x
   \longrightarrow0.                                  \tag{D.TR.3}
\]
The last map is the difference of the two branches, for
any fixed orientations. Proper strict-local base change,
nilpotent invariance, and the Picard lifting argument of
Lemma~\ref{D:local-Picard-prime-to-p-generators} show
\[
 H^2(Y,\mu_n)\simeq H^2(E,\mu_n)\simeq\Lambda^v,
 \qquad \operatorname{Pic}(Y)/n\simeq\Lambda^v,
 \qquad H^3(Y,\mu_n)=0.                               \tag{D.TR.4}
\]
Here multidegree is surjective on $\operatorname{Pic}(E)$:
choose a smooth point on each component. Its degree-zero
kernel is an extension of the product of the Jacobians by
the graph torus, hence is $n$-divisible. The formal Picard
kernel is uniquely $n$-divisible. In particular
$\operatorname{Br}'(Y)[n]=0$. These arguments impose no
finite-type-over-$k$ hypothesis on the complete local ring.

Apply derived Hom into $\mu_n$ on $Y$ to (D.TR.3),
viewing its closed-stratum sheaves as direct images to $Y$.
Closed-immersion adjunction identifies the resulting
triangle with the Mayer--Vietoris triangle for supports.
Absolute purity gives shifts $[-2]$ for a component and
$[-4]$ for a node. Thus its degree-three exact sequence is
\[
 0\longrightarrow\bigoplus_i H^1(E_i,\Lambda)
 \longrightarrow H^3_E(Y,\mu_n)
 \longrightarrow\Lambda(-1)^{e}
       \xrightarrow{\partial}\Lambda(-1)^v,            \tag{D.TR.5}
\]
where $e$ is the number of nodes. The map $\partial$ is
the signed incidence map: purity is compatible with
composition of the point-to-curve and curve-to-surface
Gysin maps, and the point class has degree one on each
proper component. This uses the exact regular-scheme
scope and compatibility of \citet[Expose XVI, Sections2--3,
Theorem3.1.1 and the proof of Corollary3.1.3]{A_ILO2014}.
The closed-stratum exact sequence used in that proof is
also (D.TR.3). All Tate twists are constant over $k$.

For a tree, $\partial$ is injective over every $\Lambda$.
Indeed a leaf row has just one nonzero entry, equal to
$1$ or $-1$; remove that edge and repeat. Hence
$H^3_E(Y,\mu_n)=\bigoplus_i H^1(E_i,\Lambda)$.
This calculation also explains the secondary-residue
assertion: the node values in (D.TR.5) form a balanced
flow, and a tree has no nonzero such flow.

Every line on the regular puncture extends to $Y$ by
closing a Cartier divisor. Consequently localization
and (D.TR.4) identify
\[
 \operatorname{Br}'(U)[n]
 =H^2(U,\mu_n)/(\operatorname{Pic}(U)/n)
                         =H^3_E(Y,\mu_n).              \tag{D.TR.6}
\]
The kernel of $\operatorname{Pic}(Y)\to\operatorname{Pic}(U)$
is generated by the exceptional divisor lines. Under the
multidegree identification in (D.TR.4), their map is the
intersection matrix $-N$. Therefore
$\operatorname{Pic}(U)/n=\operatorname{coker}(N\bmod n)$.
Smith normal form bounds its order by $|\det N|$;
negative definiteness makes this determinant nonzero.
This proves (D.TR.1). In addition, one fixed finite list
of line bundles on $Y$ represents every
$\operatorname{Pic}(U)/n$: choose representatives of
the finite group $\operatorname{coker}N$ once and use
right exactness of tensoring with $\mathbb Z/n$.

The same normalization sequence gives
$H^1(E,\mathbb Z/m)\simeq\bigoplus_iH^1(E_i,\mathbb Z/m)$
when the graph is a tree. Thus the residue tuple of a
period-$m$ class is a character on $E$ of exact order
$m$. Its torsor is connected and cyclic of degree $m$.
The proper henselian equivalence for finite etale covers
lifts it, including its action, to $Y'\to Y$
\citep[Lemma58.9.1, Tag0A48]{StacksSingleCurveBrauerD}.
Properness and connectedness of the closed fibre show
$Y'$ is connected; regularity then shows it is integral.
Its Stein ring is finite, complete, local and normal,
as in Proposition~\ref{D:single-curve-Brauer-residue-cover}.

The pulled-back class has zero generic residue on every
exceptional component upstairs. It is unramified elsewhere
because it was defined on the puncture. Purity on the
regular surface therefore extends it to $Y'$, whose
prime-to-$p$ Brauer group is zero by the same Picard and
proper-base-change argument as (D.TR.4). This proves
splitting without asserting that the new dual graph is
a tree. Indeed a connected etale cover of a nodal tree
of positive-genus curves can acquire graph cycles.

At the generic point of $E_i$, taking the $m_i$th root
of its parameter kills exactly its cyclic residue.
The proof of \citet[Proposition3.2.1, complete pp10--11]
{Lieblich2011ResidualD} applies with these componentwise
orders: the DVR construction is local on each component,
and regular-stack purity fills the remaining codimension-two
set. It gives a unique Brauer extension on (D.TR.2).
Lift it by Kummer to $H^2(\mathcal X,\mu_n)$ and correct
its restriction by a line from $\operatorname{Pic}(Y)$.
This is possible by the just-proved surjection onto
$\operatorname{Pic}(U)$, and extends the actual $\alpha$.

At a node, restrict to the strictly henselian regular
local surface. Each component residue is etale at the
node and becomes zero on that strictly local branch.
All divisor residues therefore vanish. Purity extends
the local Brauer class over the regular local surface,
where it is zero: an Azumaya algebra splits over a
strictly henselian local ring. Injectivity for the
regular root stack shows its node restriction also has
zero Brauer image. Kummer then says its class on
$B(\mu_{m_i}\times\mu_{m_j})$ is the root class of a
character line. The inertia of that root gerbe is
\[
 \{(h,z): h\in\mu_{m_i}\times\mu_{m_j},\quad
                       z^n=\lambda(h)\},              \tag{D.TR.7}
\]
for that character $\lambda$. This is commutative and
diagonalizable; all its orders are prime to $p$.
It proves the final assertion including the possibly
nonzero scalar Kummer part at the node.
\end{proof}

\begin{lemma}[Proved: the residual tree category and its node gluing]
\label{D:SNC-tree-residual-category}
In the preceding setting let $\mathcal G\to\mathcal X$
be the actual coefficient gerbe. On the normalization
of its reduced exceptional curve, the scalar-weight-one
category on component $i$ is equivalent to vector bundles
on
\[
 \mathcal R_i=\sqrt[m_j]{(R_i,P_{ij})}_{j\sim i},      \tag{D.TR.8}
\]
after choice of a neutralizing line. Here $R_i\to E_i$
is the connected etale cover of degree $m_i$ determined
by the residue, and $P_{ij}$ is the set of its $m_i$
points above the node with $E_j$. For multiple edges the
points are counted separately; the tree assumption has
already excluded parallel edges. The ordinary spectral
rank is $s_i=r/m_i$ for a rank-$r$ scalar-weight-one
bundle, so $m_i\mid r$.

Bundles on the reduced exceptional gerbe are precisely
these component bundles together with equivariant
isomorphisms of their fibres at every node. For fixed
joint character multiplicities $r_{ab}$ at a node,
after choosing homogeneous frames for the character
fibres, the isomorphisms and their inverses have polynomial
witnesses with at most $2\sum_{a,b}r_{ab}^2\leq2r^2$
coordinates and quadratic equations. This frame-atlas
description is valid over arbitrary coefficient algebras;
without frames, the character fibres may be projective
modules. It is an actual
reduced-curve equivalence; it asserts no realization
of normal deformations or final representation count.
\end{lemma}

\begin{proof}
The root-$m_i$ spectral construction on the generic
$E_i$ extends over its whole root-marked curve. Its
character torsor has no monodromy from an adjacent
root inertia: that monodromy is precisely the
commutator in (D.TR.7), which is zero. It is therefore
the pullback of $R_i\to E_i$. Applying the exact
character decomposition of
Proposition~\ref{D:residual-twisted-spectral-category}
on this base yields a Picard gerbe over (D.TR.8).
Its geometric Brauer group is zero. Indeed a root
of a smooth point on a smooth proper curve over $k$
has relative $R^1\pi_*\mathbf G_m=\mathbb Z/m_j$
supported at that point and $R^2\pi_*\mathbf G_m=0$;
the latter is the algebraically closed instance of
\citet[Lemmas4.1.1--4.1.3 and Proposition4.1.4]
{Lieblich2011ResidualD}. The base curve has zero
Brauer group and the point sheaves have zero first
cohomology. The Leray sequence therefore gives
$\operatorname{Br}'(\mathcal R_i)=0$ prime to $p$.
Neutralizing that Picard gerbe gives the claimed
category. Its inverse is pushforward of the universal
weight line tensored with the spectral bundle, exactly
as in the already proved spectral proposition.
The rank formula is $r=m_i s_i$.

At a node all the inertia is the fixed diagonalizable
group (D.TR.7). Scalar-weight-one modules split into
its characters over every $k$-algebra, by the exact
grading projections. Along either branch their
restrictions give the same joint-character list
exactly when the corresponding multiplicities agree.
After passing to the homogeneous fibre-frame atlas,
for a given joint character $(a,b)$ matching fibres
is a matrix of size $r_{ab}$ and an inverse matrix,
with the two product identities. Summing the blocks
gives the asserted coordinate and degree bounds.
The character labels include the fixed central
extension (D.TR.7); they are not replaced by an
unmarked pair of cyclic groups.

Finally the polynomial model of an ordinary node has
the fibre-product description
\[
       A[x,y]/(xy)=A[x]\times_A A[y]                 \tag{D.TR.9}
\]
for every coefficient algebra $A$, and the same identity
holds on compatible localizations. For branch modules
with a specified fibre identification, the gluing module
is the kernel of their difference map to the identified
node fibre. Locally on the coefficient and branch frame
atlases, change one basis by the prescribed invertible
matching matrix; this kernel is then a free module over
the fibre product. This is the elementary Milnor pasting
construction, with inverse given by restriction to the
two branches. The same proof is equivariant for the
diagonalizable finite node group, by choosing homogeneous
frames. These constructions commute with coefficient
changes. Applying them on polynomial node charts and
their etale localizations, followed by etale descent,
then gives vector bundles on the reduced nodal
exceptional gerbe, proving the asserted equivalence.
It does not replace the two-dimensional normal
neighbourhood by its reduced node.
\end{proof}

\begin{remark}[The remaining quantitative step]
On $\mathcal R_i$, normalized root-point weights encode
ordinary rank-$s_i$ bundles and flags at $O(m_i)$ marked
points. Their Hom sheaf is the kernel of the ordinary
Hom sheaf mapping to the forbidden-weight blocks at
those points. Thus its cohomology is represented by a
fixed-width cone of ordinary curve Hom and finite point
evaluations; arbitrary flag-dependent inverse matrices
are not automatically permitted in a quadratic bound.

To finish a tree count one must still construct actual
all-ring normal bundle families on a fixed ordinary
exceptional cycle and prove uniform slope windows and
Hom tails for this residual category. The mixed-root
numerical budget alone is adequate: if $m_i\mid m$,
weight the $i$th component's layer-$j$ variables by
$E(m/m_i)j$ and node matrices by $Em$. With common
equation degree $O(Em)$, each component has the same
factorial cancellation as before, and node variables
cost a constant per coordinate. This observation
does not prove the necessary weighted equations,
normal realization, or a tree representation count.
\end{remark}


% Separate unassembled candidate. No multicomponent representation count is claimed.
\subsection{Marked spectral windows and normal tails on an exceptional tree}

Retain the setting and the actual coefficient gerbe of
Lemmas~\ref{D:SNC-tree-Brauer-geometry} and
\ref{D:SNC-tree-residual-category}. Write
$b_i=-E_i^2>0$, let $\epsilon_i$ be the valence of the
vertex $i$, and put $\eta_i=m_i\epsilon_i$. Thus the
root-marked spectral curve $\mathcal R_i$ has coarse curve
$R_i$ of genus $\gamma_i=1+m_i(g_i-1)$ and has exactly
$\eta_i$ marked points. Let $\mathcal D_i$ denote the
reduced root divisor over $E_i$.

We use normalized root-point weights in $[0,1)$. If $W$
is a vector bundle on $\mathcal R_i$ and $W_0$ its coarse
pushforward, its parabolic degree is
\[
 \deg_{\rm par}W=\deg W_0+
       \sum_{x\text{ marked}}\sum_{0\leq\alpha<1}
                              \alpha\,\dim W_{x,\alpha}.
                                                        \tag{D.TW.1}
\]
This convention makes the root line $\mathcal O(a\mathcal P)$,
$0\leq a<d_x$, have degree $a/d_x$ and coarse pushforward
$\mathcal O$. The character/flag equivalence and tensor
compatibility are the ones of
\citet[Proposition3 and Theorem4, complete proofs]
{Borne2007D} and \citet[Theorem6.1]{BorneVistoli2012D}.

\begin{lemma}[Proved: minimal marked spectral windows]
\label{D:SNC-tree-minimal-marked-windows}
Fix a formal scalar-weight-one bundle $F_0$ on the completed
coefficient gerbe. There is a bundle lattice $F$ in its
meromorphic bundle such that the coarse spectral bundle
$W_{i,0}$ on each $R_i$ satisfies
\[
 \gamma_i-1-\eta_i\ \leq\ \mu_{\min}(W_{i,0})
 \ \leq\ \mu_{\max}(W_{i,0})
 \ \leq\ \gamma_i-1+b_i+2\eta_i.                       \tag{D.TW.2}
\]
The endpoints are $O(m_i)$ with constants depending only
on the fixed resolution. The auxiliary generator in the
proof is used to minimize an integer; it is not an ordinary
parameterization of the twisted bundle by a module of its
coarse Hom rank.
\end{lemma}

\begin{proof}
Put $m=\operatorname{lcm}_i(m_i)$ and $h=m^2$.
We first construct one residual scalar-weight-one generator
$A_E$ of rank $h$. Choose the neutralizing spectral lines
from Lemma~\ref{D:SNC-tree-residual-category}. On
$\mathcal R_i$, put $w_i=h/m_i$. For every marked point
whose root order is $m_j$, the integer $w_i$ is divisible
by $m_j$. Choose a direct sum $V_i$ of $w_i$ root-point
line bundles such that at each such point all $m_j$
weights occur equally often. This can be done independently
at the different points: distribute the weights among
$w_i$ labelled summands and tensor the resulting point
lines in each summand. Every exponent is in
$\{0,\ldots,m_j-1\}$. Consequently
\[
 (V_i)_0=\mathcal O_{R_i}^{w_i},\qquad
                 0\leq\mu_{\rm par}(V_i)\leq\eta_i.
                                                        \tag{D.TW.3}
\]
Let $A_i$ be its inverse spectral image. At a node $i,j$,
each scalar-weight-one joint character occurs in $A_i$
with multiplicity $h/(m_i m_j)$. A neutralizing-line
character merely permutes this uniform list. The same
list occurs in $A_j$, so choose equivariant fibre
isomorphisms and glue by the proved node-pasting lemma.
This gives $A_E$.

The obstruction to extending this bundle across a
square-zero exceptional thickening lies in coherent
$H^2$ of its endomorphism sheaf tensored with the ideal.
Exact tame invariants reduce this to coherent cohomology
on a proper curve, so the obstruction is zero. Successive
extensions give a formal bundle $A$. No uniqueness is
required. For a formal lattice $F$ put
\[
                         H_F=q_*\mathcal Hom(A,F),
                                                        \tag{D.TW.4}
\]
where $q$ is the composite coarse map to $\widehat Y$.
This is an ordinary formal vector bundle. Indeed on a
strict-local crossing chart the root ring has the form
$B[s,t]/(s^{m_i}-x,t^{m_j}-y)$, and a homogeneous basis
for the equivariant Hom has invariant summands generated
over $B$ by the unique monomials with exponents in
$[0,m_i-1]\times[0,m_j-1]$ of the required characters.
The one-root calculation applies away from nodes. The
scalar gerbe first disappears from the weight-zero Hom.
This also shows exactness on the elementary-transform
sequences below.

The exceptional and closed-fibre topologies are cofinal.
Ordinary Grothendieck existence algebraizes $H_F$;
meromorphic comparisons algebraize after fixed Cartier
twists. As in Lemma~\ref{D:twisted-connected-anchor-minimal-slopes},
the ordinary bundles $H_F$ have a common punctured
restriction with coherent reflexive extension $M$ to
$\operatorname{Spec}R$. Hence
\[
 e(F)=\operatorname{length}_R(M/\Gamma(Y,H_F))
                              +\dim_k H^1(Y,H_F)\geq0
                                                        \tag{D.TW.5}
\]
is a finite integer. Choose $F$ minimizing it. This
argument only uses ordinary existence for (D.TW.4),
exactly as in the previously proved one-curve version
\citep[Theorem30.27.1, Tag088E]{StacksTwistedDepthD}.

Let $W_i$ be the spectral bundle of $F|_{\mathcal D_i}$.
For a locally free quotient $W_i\twoheadrightarrow Q$
in the root-curve category, form the elementary kernel
of $F\to\iota_{i*}\mathcal Q$. A quotient locally free
on a Cartier divisor makes that kernel locally free
on the regular root chart. Exact coarse Hom identifies
the quotient of (D.TW.4) with
\[
 (\iota_{E_i})_* (h_i)_*
        (\pi_i)_*\mathcal Hom(V_i,Q),
                                                        \tag{D.TW.6}
\]
where $h_i:R_i\to E_i$ and
$\pi_i:\mathcal R_i\to R_i$. There is no further
factor $m_i$ in the Euler characteristic. Minimality
therefore says
\[
                   \chi(\mathcal R_i,V_i^\vee\otimes Q)
                                  \geq0.              \tag{D.TW.7}
\]

For clarity the numerical estimate needed here follows
directly from ordinary curve Riemann--Roch and (D.TW.1).
Tame pushforward is exact; the normalized weights of a
rank-$a$ bundle have total contribution between $0$
and $a\eta_i$. Parabolic degree is additive under
tensor products in the usual rank-weighted formula.
Using (D.TW.3), for any positive-rank bundle $Q$ we obtain
\[
 \mu_{\rm par}(Q)+1-\gamma_i-2\eta_i
 \ \leq\
 \frac{\chi(\mathcal R_i,V_i^\vee\otimes Q)}
                  {w_i\operatorname{rk}Q}
 \ \leq\ \mu_{\rm par}(Q)+1-\gamma_i.
                                                        \tag{D.TW.8}
\]
Thus every such quotient has parabolic slope at least
$\gamma_i-1$.

Tensoring by the root normal line
$\mathcal O_{\mathcal D_i}(\mathcal D_i)$ transports
the spectral bundle by a deck translation and a line
of parabolic degree $E_i^2=-b_i$. To see the degree
without ignoring the marked root characters, iterate
the transport $m_i$ times. Its line product is
$h_i^*\mathcal O_{E_i}(E_i)$, of degree $m_i E_i^2$.
Deck translation preserves parabolic degree, so each
step has degree $E_i^2$. This also follows by cancelling
the neutralizing line and its translate in rational
degree. For a root-curve subbundle $S\subset W_i$,
the inverse elementary transform has quotient
$\mathcal S(\mathcal D_i)$ and gives
\[
       \chi(\mathcal R_i,V_i^\vee\otimes
                     \operatorname{Nor}_i(S))\leq0.
                                                        \tag{D.TW.9}
\]
The lower estimate in (D.TW.8) consequently gives
$\mu_{\rm par}(S)\leq\gamma_i-1+b_i+2\eta_i$.

Finally take an ordinary subbundle or locally free
quotient of $W_{i,0}$. Give it the induced intersection
or image flags at the marked points. In an adapted
local basis these are strict parabolic subobjects and
quotients, so they define the locally free root-curve
objects used above. Their parabolic slope differs
from their ordinary slope by a number in
$[0,\eta_i]$. Applying the quotient estimate to the
last ordinary Harder--Narasimhan quotient and the
subobject estimate to the first subbundle proves
(D.TW.2). This proves the assertion without assuming
tensor semistability in positive characteristic.
\end{proof}

\begin{proposition}[Proved: fixed ordinary depth for marked tree lattices]
\label{D:SNC-tree-uniform-Hom-depth}
Choose once a positive integral exceptional divisor
$A=\sum_i a_iE_i$ with $c_i=-A\cdot E_i>0$, and
write $\mathcal A=q^*A=\sum_i a_i m_i\mathcal D_i$.
For the minimal lattices of (D.TW.2), and for any fixed
finite list of their ordinary exceptional twists,
there are integers $B,N$ independent of the scalar
order, the residue orders and the ranks such that:
\begin{enumerate}
\item every meromorphic map of matching scalar weights
extends with poles at most $BA$;
\item a bundle on $N\mathcal A$ extends to a formal
bundle; every specified map between formal extensions
on this thickening lifts, and a specified isomorphism
lifts to an isomorphism;
\item the same assertions hold for actual bundle families
over finite local Artinian coefficient algebras;
\item for fixed $a$, the component-layer Hom cohomology
on $a\mathcal A$, for ranks $r,r'$, has total dimension
$O(rr')$, with constant depending on the fixed resolution
and $a$.
\end{enumerate}
The residual bundles in the second and third assertions
are required to satisfy (D.TW.2). This statement supplies
formal lifting and a size bound; it asserts neither a
polynomial object atlas nor a tree representation count.
\end{proposition}

\begin{proof}
Filter a copy of $\mathcal A$ by adding its root
components one at a time, with component $i$ repeated
$a_i m_i$ times. The corresponding factors are line
twists of bundles on the smooth root-marked
$\mathcal R_i$. In the $l$th ordinary-$A$ layer the
positive conormal degree is
\[
                 l m_i c_i+O(m_i).                   \tag{D.TW.10}
\]
The error is uniform: in one copy of $\mathcal A$ a
preceding coefficient of $\mathcal D_j$ lies between
$0$ and $a_jm_j$, whereas its degree on the spectral
curve over $E_i$ is $(m_i/m_j)(E_i\cdot E_j)$ when
$j\ne i$. The self-component error is also $O(m_i)$.
The corresponding negative pole layer has degree
$-l m_i c_i+O(m_i)$.

There is a useful direct comparison at the marked
points. For two root-curve bundles with coarse bundles
$V,W$, their coarse Hom is between
\[
 \mathcal Hom(V,W)(-P_i)\ \subset\
  (\pi_i)_*\mathcal Hom(\mathcal V,\mathcal W)
   \ \subset\ \mathcal Hom(V,W),
                                                        \tag{D.TW.11}
\]
where $P_i$ is the reduced sum of the $\eta_i$ marked
points. Both quotients are zero-dimensional. This is
the adapted homogeneous-basis description: at each
point a matrix entry either has no vanishing requirement
or vanishes to first order, according to the two
normalized weights. Tensoring by a root line can
shift coarse degrees by at most another $\eta_i$.

The windows (D.TW.2), the genus
$\gamma_i=1+m_i(g_i-1)$, and the marked correction
$\eta_i$ are all $O(m_i)$. The slope criterion for a
nonzero map and curve duality applied to (D.TW.11)
therefore imply that every negative Hom layer has
zero $H^0$ once $l\geq B$, and every positive Hom
layer has zero $H^1$ once $l\geq N$, for fixed
$B,N$. For the latter, it suffices to make
$H^1(\mathcal Hom(V,W)(-P_i))$ vanish; duality reduces
this to a Hom from $W$ to
$V\otimes\omega_{R_i}(P_i)$ whose slope gap is
strictly negative. No semistability assertion for
a tensor product is used.

The successive-layer exact sequences now stabilize
sections of every meromorphic Hom after the fixed
pole order $B$. An obstruction to lifting a specified
map beyond $N\mathcal A$ belongs to the $H^1$ just
shown to vanish. Bundle-lifting obstructions belong
to $H^2$ on a proper one-dimensional tame stack and
are zero. The same argument lifts inverse maps;
alternatively a lifted map reducing to an isomorphism
has invertible determinant on the formal neighborhood.

For an Artinian coefficient algebra, filter by its
maximal ideal. Each graded coefficient layer is a
$k$-vector space tensored with one of the preceding
proper-curve Hom layers. The same zero groups give
every prescribed lift, including all coefficient
nilpotents. Thus no fibrewise exclusion of a nilpotent
pole is being assumed. The component filtration also
proves the assertions for the complete exceptional
thickening, including its node overlaps; it is a
filtration of an actual Cartier cycle on the regular
root surface, rather than a replacement by the
normalization of its reduced curve.

Finally put $s_i=r/m_i$ and $s_i'=r'/m_i$ when a
nonzero matching-weight Hom occurs. There are at
most $aa_i m_i$ factors on component $i$. In each
factor the genus, slope windows and line degree
are $O(m_i)$, so curve Riemann--Roch, duality and
(D.TW.11) give
\[
 h^0+h^1=O(m_i s_i s_i')=O(rr'/m_i).
                                                        \tag{D.TW.12}
\]
Their product is $O(rr')$ for that component, and
there are only the fixed number $v$ of components.
This proves the total bound. It does not count
coordinates for arbitrary deformations or prove
polynomial equations for them.
\end{proof}


% Separate post-CP21 proved input. All four complete audits passed; no projective representation count is asserted.
\begin{lemma}[Proved: degree-two classes with tame cyclic coefficients]
\label{D:radial-tame-cyclic-H2}
Retain the complete one-smooth-exceptional-curve setting of
Lemma~\ref{D:radial-tame-cover-one-curve}. Let $\mathcal M$
be a finite locally constant sheaf on $U$ whose fibres are
cyclic groups of order $n$, with $n$ prime to $p$. Suppose
its monodromy is radially tame along $C$. Then
\[
 |H^2_{\mathrm{et}}(U,\mathcal M)|\leq n^{2g+1},\qquad
 |H^2_{\mathrm{cont}}(\pi_1(U),M)|\leq n^{2g+1},
                                                        \tag{D.T2.1}
\]
where $M$ is a geometric fibre with its monodromy action.
For $n=1$ both groups are trivial. No prime-to-$p$
condition is imposed on the order of the entire monodromy
image; only the coefficient order and radial tameness
are assumed.

For a degree-$t=eu$ radially tame cover $U_H\to U$
and a tame field character
$\chi:H=\pi_1(U_H)\to\operatorname{Gal}(\mathbb F_q/\mathbb F_p)$,
where $q=p^f$, this gives
\[
 |H^2_{\mathrm{cont}}(H,\mathbb F_q^*(\chi))|
       \leq(q-1)^{2g(C_H)+1}\leq q^{(2g+3)t}.
                                                        \tag{D.T2.2}
\]
This bounds the scalar obstruction classes for projective
semilinear representations with that fixed field character.
It does not assert inflation in degree two through the
radial tame quotient.
\end{lemma}

\begin{proof}
We assume $n>1$. Write $j:U\hookrightarrow Y$ and
$i:C\hookrightarrow Y$. We first establish the precise
local calculation needed for a tame coefficient sheaf:
\[
 \mathcal M_0=i^*j_*\mathcal M,\qquad
 \mathcal M_1=i^*R^1j_*\mathcal M
 \quad\text{are finite locally constant on }C,
 \qquad R^bj_*\mathcal M=0\quad(b\geq2).
                                                        \tag{D.T2.3}
\]
Their fibres are respectively the invariants and the
coinvariants, with the usual $(-1)$ twist on the latter,
of a tame inertia action on $M$. In particular each has
cardinality at most $n$. This calculation requires more
than a statement about the finite generation of tame
fundamental groups, so we supply the argument.

At a geometric point of $C$, let $A$ be the strict
henselization of the regular local ring of $Y$, and let
$z$ cut out $C$. Then $A/(z)$ is regular and strictly
henselian. For every positive integer $a$ prime to $p$,
put
\[
 X_a=\operatorname{Spec}A[s]/(s^a-z),\qquad
 D_a=V(s),\qquad U_a=X_a-D_a.
                                                        \tag{D.T2.4}
\]
The ring defining $X_a$ is regular and strictly
henselian, and $D_a$ is regular and strictly henselian.
These are the Kummer charts of SGA1 XIII.5.1--5.3 and
Stacks Tag 0EYG, whose exact scope was checked in
Lemma~\ref{D:radial-tame-cover-one-curve}.

For a constant finite abelian group $B$ killed by $n$,
strict-local acyclicity, codimension-one absolute purity,
and localization give
\[
 H^0(U_a,B)=B,\qquad H^1(U_a,B)=B(-1),\qquad
                    H^b(U_a,B)=0\quad(b\geq2).
                                                        \tag{D.T2.5}
\]
Here the purity input is the regular-scheme theorem and
its one-component calculation in
\citet[Expose XVI, Theorem 3.1.1 and
Corollary 3.1.3]{A_ILO2014}. Decomposing $B$ into cyclic
groups reduces (D.T2.5) to those constant coefficients.
Under $U_{ab}\to U_a$ the map on $H^1$ is multiplication
by $b$: the Kummer class of the parameter $s$ pulls back
to the class of $(s')^b$. Equivalently this is the
ramification-index compatibility of the purity residue.

The pullback of $\mathcal M$ becomes constant at a
finite stage of this Kummer tower, by tame Abhyankar.
The filtered limit of its higher cohomology groups on
the tower is zero: the groups in degree at least two
already vanish, and transitions of degree divisible
by $n$ kill degree one. The Cartan--Leray spectral
sequence for the tower therefore identifies its local
cohomology with continuous cohomology of
\[
                   I=\prod_{\ell\ne p}\mathbb Z_\ell(1).
                                                        \tag{D.T2.6}
\]
We use here the tower form of Cartan--Leray, with
filtered cohomology on the inverse limit, as in
\citet[Theorem 14.9 and Remark 14.10, pp.~96--97]
{MilneLEC2013H2D}. This use does not assume that the
kernel of the full fundamental group onto $I$ is
pro-$p$.

For a finite prime-to-$p$ module, continuous cohomology
of the procyclic group $I$ is computed in degrees zero
and one by $[M\xrightarrow{\sigma-1}M]$, after choosing
a compatible tame generator, and is zero in higher
degrees. One way to see the higher vanishing is to
separate each $\ell$-primary summand: the prime-to-$\ell$
factor has exact invariants, and the completed group
algebra of $\mathbb Z_\ell$ has the two-term resolution
with differential $\sigma-1$. The degree-one formula
also follows directly: every value in $M$ gives a
continuous crossed homomorphism after passing to a
cyclic quotient whose order is sufficiently divisible
by the exponent of $M$; coboundaries are $(\sigma-1)M$.
Thus the stalks in (D.T2.3) are $M^I$ and
$M/(\sigma-1)M$, with the residue Tate twist.

The Kummer description descends from the strict
henselization to etale neighbourhoods of points of $C$.
On each such neighbourhood the inertia action is one
fixed finite cyclic action. The above construction is
compatible with etale change of that neighbourhood,
so its invariant and coinvariant modules give locally
constant sheaves on $C$. This proves the lissity, not
just the cardinality of individual stalks, in (D.T2.3).
Outside $C$ the higher direct images vanish because
$j$ is an isomorphism there.

Since $Y$ is proper over the strictly henselian
$\operatorname{Spec}R$, proper base change for torsion
sheaves and exactness of the strict-local global-section
functor identify
\[
 H^a(Y,R^bj_*\mathcal M)
       =H^a(C,i^*R^bj_*\mathcal M).
                                                        \tag{D.T2.7}
\]
Nilpotents in the actual closed fibre do not affect its
etale site. The exact proper-base-change input is
\citep[Theorem 59.91.11, Tag 095T]{StacksSemilinearDeterminantsB}.
The Leray spectral sequence of $j$ and (D.T2.3) now
show that $H^2(U,\mathcal M)$ has a filtration with
subquotients of $H^2(C,\mathcal M_0)$ and
$H^1(C,\mathcal M_1)$.

For clarity we give the required curve bounds without
assuming a prime-to-$p$ monodromy group. For any finite
locally constant abelian sheaf $\mathcal A$ on a smooth
proper connected curve $C$, with prime-to-$p$ fibre
$A_0$, one has
\[
 |H^1(C,\mathcal A)|\leq |A_0|^{2g},\qquad
 |H^2(C,\mathcal A)|\leq |A_0|.
                                                        \tag{D.T2.8}
\]
The first bound follows from the full $2g$-generator
bound for $\pi_1(C)$ and the crossed-homomorphism argument
already used in (D.RT.6). For the second, choose a
connected finite etale cover $v:C_1\to C$ trivializing
$\mathcal A$. The sheaf trace
$v_*v^*\mathcal A\to\mathcal A$ is surjective: etale
locally it is the sum map from a nonempty finite direct
sum. Its kernel $\mathcal K$ is torsion and locally
constant. Since $H^3(C,\mathcal K)=0$, the trace induces
a surjection
\[
 H^2(C_1,(A_0)_{C_1})\twoheadrightarrow
                           H^2(C,\mathcal A).
                                                        \tag{D.T2.9}
\]
The source has cardinality $|A_0|$. The precise inputs
are the trace construction, finite-direct-image
vanishing, constant proper-curve cohomology, and the
cohomological-dimension theorem
\citep[Tags 03SH, 03QP, 03RQ and Theorem 59.83.10
in Tag 03SB]{StacksTameTwistedH2B}. In particular no
division by $\deg v$ is used.

Both fibres in (D.T2.3) have order at most $n$.
Multiplying the two bounds in (D.T2.8) proves the first
inequality of (D.T2.1).

Finally the usual comparison in degree two is injective:
\[
 H^2_{\mathrm{cont}}(\pi_1(U),M)
               \hookrightarrow H^2_{\mathrm{et}}(U,\mathcal M).
                                                        \tag{D.T2.10}
\]
Indeed use Cartan--Leray over the inverse system of
connected finite Galois etale covers of $U$. After a
finite stage the coefficient sheaf is constant, and
every degree-one class on a stage is killed on a
further finite etale cover, namely its torsor, followed
if necessary by a Galois closure. Thus the first
cohomology of the universal pro-cover is zero, and the
five-term sequence gives (D.T2.10). The same tower
form of Cartan--Leray cited above applies to arbitrary
schemes in its stated scope. This proves the second
inequality. Applying it to $(Y_H,C_H)$, whose genus
is $1+u(g-1)$, proves (D.T2.2).
\end{proof}


% Proved post-CP21 result. Fixed partial resolutions replace moving Picard backgrounds.
\subsection{Fixed partial resolutions and hereditary projective counts}

\begin{lemma}[Proved: fixed partial resolutions with exact finite Hom witnesses]
\label{D:fixed-partial-resolution-unit-atlas}
Work in the one-curve coefficient-gerbe setting of
Proposition~\ref{D:weighted-fixed-determinant-normal-presentation},
but do not fix the residual determinant. Allow the underlying
one-curve surface to be the normalization of a radially tame
cover of degree $t=eu$ as in
Lemma~\ref{D:radial-tame-cover-one-curve}. Let $m$ be the
geometric Brauer period, let $r=ms$, and let $R_a$ be the
connected degree-$m$ spectral curve at a Frobenius stage.
Thus $g(R_a)-1=um(g-1)$ and the positive root-normal degree
is $u\delta/e$.

For each fixed coefficient gerbe and its Frobenius stages,
one can choose fixed formal bundles $P_{j,a}$, $0\leq j\leq7$,
with the following properties. Their spectral ranks are
$s+1$ for $j=0$ and two for $j>0$. The normal unit objects
in the fixed minimal slope windows are covered, after one
common positive twist, by complexes
\[
 P_{7,a}\longrightarrow\cdots\longrightarrow
 P_{1,a}\longrightarrow P_{0,a}.                       \tag{D.FR.1}
\]
The equations are $d^2=0$, followed by open residual-rank
and slope conditions. Every point of this presentation over
every coefficient algebra defines an actual locally free
$H^0$ bundle on the chosen finite normal thickening.
Coverage holds for geometric objects and for lifts of
chosen presentations over finite local Artinian coefficient
algebras.

With $J=ae m$ root-normal layers, where $a$ is fixed for
the original resolution, the objects and all fixed-length
unit, inverse and equality diagrams admit a finite
polynomial presentation with at most $Afur^2/m=Afu m s^2$
variables at each positive weight $j\leq J$, at most the
same number of residual variables of weight $em$, and
weighted equation degree at most $Ke m$. Here $f$ is the
number of one-step Frobenius stages, and $A,K,a$ are
independent of $e,u,m,s,f$. No moving Picard variety,
cohomology-basis variables, or inverse of a matrix depending
on a residual determinant is used.
\end{lemma}

\begin{proof}
\emph{Fixed terms and lifting their generators.}
Choose the connected rank-$m$ anchors $A_a$ of
Lemma~\ref{D:twisted-connected-anchor-minimal-slopes},
with residual spectral bundle $\mathcal O_{R_a}$.
They are fixed formal bundles; their coefficients need
not satisfy any additional rational-pole normalization.
Their residual universal lines are normalized as in
Proposition~\ref{D:residual-twisted-Frobenius-transport}.
In particular the absolute degrees of the one-step Frobenius
comparison lines are $O(um)$, with a constant depending only
on the original resolution. This scaled bound is independent
of the scalar band order and of $e,u,m$; the root-normal
degree $u\delta/e$ is at most $u\delta$.

The minimal residual slopes are
$[g(R_a)-1,g(R_a)-1+u\delta/e]$. Twist by a sufficiently
large fixed multiple of the pullback of $-eC'$.
The resulting spectral bundles $W_a$ have slopes in
$[c_1um,c_2um]$, where $c_1$ is chosen so that
$\mu_{\min}(W_a)>2g(R_a)+2$; the displayed constants
may be increased to handle genus zero as well. The
precise assertion needed below is this strict inequality
and a uniform $O(um)$ upper bound, not positivity of
the genus itself. All positive root-normal Hom layers
from $A_a$ to these bundles have zero $H^1$.

Write $L_a=\det W_a$. A general $s+1$-tuple of sections
generates $W_a$, and gives a residual kernel $L_a^{-1}$.
The determinant line $L_a$ is globally generated and
has $H^1=0$: its degree is at least
$s\mu_{\min}(W_a)>2g(R_a)-2$. These properties also
hold after every positive root-normal translation and
twist. Choose a fixed ordinary formal line $B$ whose
degree on each $R_a$ is a positive constant times $um$,
for instance the pullback of a sufficiently large fixed
multiple of $-eC'$. It is invariant under the spectral
deck translations. Choose $b=c s$, with one fixed
integer $c$ large enough that
\[
       L_a^{-1}B^b\quad\hbox{and}\quad L_a
       \quad\hbox{are globally generated and have }H^1=0
                                                        \tag{D.FR.2}
\]
after all the positive normal twists needed here. This
choice is uniform because $|\deg L_a|=O(um s)$ and
$\deg B$ is a fixed positive multiple of $um$.

Put
\[
 P_{0,a}=A_a^{\oplus(s+1)},\qquad
 P_{j,a}=(A_a\otimes B^{-jb})^{\oplus2}\quad(j\geq1).
                                                        \tag{D.FR.3}
\]
Lift the chosen $s+1$ residual generators to the normal
bundle $F_a^+$. This is possible at every order by the
$H^1$ vanishing just stated. Their map remains surjective
by Nakayama, and its kernel $K_{1,a}$ is a bundle of rank
$m$ with spectral line $L_a^{-1}$. Two residual generators
of $L_a^{-1}B^b$ lift to a surjection
$P_{1,a}\to K_{1,a}$. Its kernel has spectral line
$L_aB^{-2b}$. Two generators of $L_a$ then lift to a
surjection $P_{2,a}\to K_{2,a}$. Continuing gives
\[
 \begin{split}
 K_{2h+1,a}|_{R_a}&=L_a^{-1}B^{-2hb},\\
 K_{2h,a}|_{R_a}&=L_aB^{-2hb}\quad(h\geq1).
 \end{split}                                            \tag{D.FR.4}
\]
At odd steps the generating Hom line is
$L_a^{-1}B^b$; at even steps it is $L_a$. Thus (D.FR.2)
proves the lifting at every step, without increasing the
normal precision or the ranks of the fixed terms. The
normal translations may replace a line by a deck translate;
they preserve its degree and the asserted vanishing.

The same argument lifts any specified presentation across
a small extension of local Artinian coefficient algebras.
The obstruction to lifting a chosen map is the indicated
$H^1$ tensored with the square-zero ideal, and is zero.
This assertion concerns actual bundle families with a
chosen closed-fibre presentation. It does not assume
flatness of a relative reflexive hull or global freeness
of a cohomology module over an arbitrary parameter ring.

\emph{Actual bundles from the equations over every ring.}
Place (D.FR.1) in cohomological degrees $[-7,0]$.
At normal order zero require each consecutive map to
have spectral rank at least one everywhere on $R_a$.
This is an open condition on the parameter scheme:
failure is a closed rank locus on the proper residual
curve, and its image is closed. Together with $d^2=0$
it forces every middle rank to be exactly one. Indeed,
at a local point a unit minor splits one rank-one map;
the adjacent product-zero equation and the next unit
minor split its complementary summand. Induction
cancels the contractible pairs. This argument uses unit
minors over the actual coefficient ring, not just ranks
at field-valued points.

Consequently the residual complex is locally the direct
sum of contractible pairs and two vector bundles in its
end degrees. After a nilpotent normal thickening, each
split pair still has a unit minor and can be cancelled.
The remaining end terms are separated by at least two
degrees, so no differential can join them. It follows
that the middle cohomology remains zero and both end
cohomology modules are locally free over every
coefficient algebra. In particular $H^0$ is an actual
bundle of rank $r$, and all these conclusions commute
with base change. The slope-window condition on that
bundle is the open proper-Quot condition already checked
in the fixed-determinant construction. No additional
equations for an arbitrary formal moduli functor are
being inferred from tangent dimensions.

\emph{Exact low-degree Hom, composition and identities.}
Let $P^L_F$ denote the subcomplex in degrees $[-L,0]$
of such a resolution, with $L\geq2$. Its only
cohomology is the bundle $F$ in degree zero and a
bundle $T_F^L$ in degree $-L$. Thus
\[
 T_F^L[L]\longrightarrow P_F^L\longrightarrow F
                   \longrightarrow T_F^L[L+1]         \tag{D.FR.5}
\]
is a canonical triangle. On our proper one-dimensional
tame thickening the derived global Hom of two bundles
has cohomology only in degrees zero and one. This
follows from exact coarse invariants and a two-affine
Cech cover, and remains true over every coefficient
algebra.

Set $M=3$ and $L=7$. Apply derived Hom to (D.FR.5)
for the source and target. The target-end error has
degrees in $[-M,1-M]$. The source-end error against
$P_G^M$ has degrees in $[L-M,L+1]$. Neither can
affect degrees zero, one or two. Therefore
\[
 H^iR\Gamma\mathcal Hom(P_F^L,P_G^M)
       \simeq H^iR\Gamma\mathcal Hom(F,G),
                       \qquad i=0,1,2.                \tag{D.FR.6}
\]
This proof uses locally free end modules, established
above, and is valid over arbitrary coefficient rings.
Equal source and target lengths would not give this
comparison; the separated lengths are essential.

There is a canonical inclusion $\iota_G:P_G^M\to P_G^L$
of complexes, inducing the identity on $G$. Thus arrows
$a:P_F^L\to P_G^M$ and $b:P_G^L\to P_H^M$ compose
by $b\iota_G a$. Equality with another arrow is equality
of their degree-zero cohomology classes in (D.FR.6),
and is imposed by a degree-minus-one boundary witness.
To impose an identity, do not demand a homotopy identity
on an equal-length complex. Instead precompose the
proposed endomorphism with the canonical top-bundle
inclusion $P_{0,F}\to P_F^L$ and compare it with
$P_{0,F}\to P_F^M$. One has
\[
 H^0R\Gamma\mathcal Hom(P_{0,F},P_F^M)
                         =\operatorname{Hom}(P_{0,F},F),
                                                        \tag{D.FR.7}
\]
by the target-end bound with $M\geq2$. Since
$P_{0,F}\to F$ is surjective, this comparison is exactly
the identity condition on $F$. Pole-cleared identities
use the same canonical map multiplied by the fixed
Cartier section. This also gives exact inverse witnesses.

\emph{Finite complexes, weights and dimensions.}
Use the two-affine Cech complexes of the fixed formal
terms (D.FR.3), and split their residual complexes as
in Lemma~\ref{D:weighted-fixed-residual-bundle-deformations}.
Retain the complete finite degree-zero and degree-one
cohomology spaces at each normal layer; do not choose
bases of cohomology depending on the matrices $d$.
The fixed normal perturbation increases normal order.
The variable complex differentials increase internal
degree in a complex of fixed length seven. The same
nilpotent homological perturbation identities therefore
terminate: every word either consumes positive normal
order or increases internal degree, whose total width
is a fixed constant. They give actual finite Hom
complexes, cycles, boundary witnesses and composition
polynomials over every coefficient algebra, by precisely
the identities and coefficient-base-change argument of
the cited fixed-line construction
\citep[Sections2.1--2.4, complete proof, and3.2]
{C_Crainic2004Perturbation}.

For the coordinate count, the block between the two
top terms has spectral matrix multiplicity $(s+1)^2$
and residual line degree $O(um)$ at all layers
$j\leq J=ae m$. Its cohomology contributes $O(um s^2)$.
A block involving exactly one positive-index term has
multiplicity at most $2(s+1)$ and line degree $O(um s)$;
it contributes the same bound. Blocks between two
positive-index terms have multiplicity four and line
degree $O(um s)$. There are only a fixed number of
blocks. These bounds include the fixed pole twists and
one-step Frobenius: $J(u\delta/e)=O(um)$, $b=O(s)$,
and Frobenius has the fixed factor $p$.

Give positive-order variables weight equal to that
order, and every residual variable weight $em$.
Positive normal orders add in a product. The number of
residual complex-differential factors in a transferred
word is bounded by its fixed internal width; all other
factors consume normal order. The finite composition
diagram above has fixed length. Hence every resulting
equation has weighted degree $O(em)$, with no factor
depending on $u$ or $s$. All reference coefficients are
constants. This proves the asserted presentation and
its $O(fum s^2)$ bound per layer.
\end{proof}

\begin{theorem}[Proved: hereditary projective counts without a bounded-period assumption]
\label{D:tame-hereditary-projective-count}
In the setting of Theorem~\ref{D:tame-resolution-hereditary-linear-count}
there is a constant $C_R$ such that, for every connected
radially tame degree-$t$ cover $U'\to U$, every $q=p^f$
and every $r\geq1$,
\[
 \#\operatorname{Hom}_{\mathrm{cont}}
     (\pi_1(U'),\operatorname{PGL}_r(\mathbb F_q))
                      /\operatorname{PGL}_r(\mathbb F_q)
                         \leq q^{C_Rt r^2}.            \tag{D.FR.8}
\]
The cardinality of the displayed conjugacy-class set is
intended. The geometric Brauer period is unrestricted;
the covering is radially tame, but the represented
projective local systems may be wild.
\end{theorem}

\begin{proof}
Write $t=eu$ and retain the smooth model $(Y',C')$ of
Lemma~\ref{D:radial-tame-cover-one-curve}. Put $n=q-1$.
The obstruction to a projective representation is in
$H^2(\pi_1(U'),\mathbb F_q^*)$, which injects into
$H^2(U',\mu_n)$ by the already checked degree-two
fundamental-group comparison. The one-curve purity
calculation gives
\[
 |H^2(U',\mu_n)|=\gcd(n,d')n^{2g'}
       \leq \delta u\,q^{(2g+2)u},\qquad
 d'=u\delta/e,\quad g'-1=u(g-1).                       \tag{D.FR.9}
\]
The harmless factor $\delta u$ is absorbed in
$q^{O(t)}$. For $n=1$ the coefficient class is trivial.

Fix one realized coefficient class $\alpha$. Let $m$
be its geometric Brauer period. A degree-$r$ Azumaya
representative implies $m\mid r$; also $m\mid n$.
Write $r=ms$. The actual coefficient gerbe extends
over the root-$m$ stack along $C'$, as in
Lemma~\ref{D:one-curve-coefficient-root-gerbe}.
All Frobenius scalar characters in its cycle have
connected degree-$m$ spectral curves. Choose fixed
rank-$m$ anchors at these stages. No moving parameter
is introduced by a choice of such reference anchors.

The minimal lattice slopes and one-step Frobenius
bounds are $O(um)$, while one root-normal step has
positive degree $u\delta/e$. Consequently fixed
multiples of $em$ give the poles, the Hom tails and
the paired unit precision, exactly as in the separate
$u,e$ calculation (D.TH.3)--(D.TH.5). Use the actual
fixed-resolution presentation of
Lemma~\ref{D:fixed-partial-resolution-unit-atlas} at
$J=ae m$ root layers for the $f$ stages. It covers
all residual determinant lines, so there is no
additional Picard chart or determinant enumeration.

Pad to $b=\lceil Afum s^2\rceil$ variables at each
weight $1\leq j\leq J$ and $b$ residual variables
of weight $em$. The common weighted degree is $Ke m$.
Lemma~\ref{D:weighted-affine-component-bound} gives
\[
 (b(J+1)+1)\frac{(Ke m)^{b(J+1)}}{(em)^b(J!)^b}
          \leq\exp(A_1f e u m^2s^2)
          =\exp(A_1f t r^2).                           \tag{D.FR.10}
\]
The inequality is the factorial calculation
(D.TH.9)--(D.TH.10), with $e$ replaced by $em$ and
the per-layer factor $u$ replaced by $um$.

All objects and arrows in this presentation are
actual bundles and maps over the normal thickening.
The existing specified-map lifting, paired inverse
correction and formal algebraization therefore apply.
Their constants are the fixed multiples of $em$
already included in the precision. Finite etale
unit frames are unchanged by nilpotent coefficient
thickenings, so the actual unit-isomorphism classes
are open and closed, by
Proposition~\ref{D:finite-jet-unit-Isom}. The parameter
complexes merely present those bundles and maps;
no isomorphism of their extra left-end syzygies is
required. Thus (D.FR.10) bounds the represented
unit classes.

Weight-one $\mathbb F_q$-local systems on the fixed
coefficient gerbe map onto projective representations
with that obstruction class. The map on isomorphism
classes is surjective: choose a cocycle representing
the central extension and adjust the chosen projective
lifts by its scalar coboundary. Additional scalar
characters can give more twisted objects over one
projective object, which only improves this upper
bound. This is the same exact finite-unit comparison
used in Theorem~\ref{D:one-curve-unbounded-period-projective-count}.
Finally multiply (D.FR.10) by (D.FR.9) and use $q=p^f$.
After increasing one constant the result is (D.FR.8).
This is an ordinary projective count; a prescribed
semilinear field character requires its own twisted
coefficient class and deck-cocycle comparison.
\end{proof}


% Separate post-CP21 proved implications. All four audits passed; the projective count remains an explicit hypothesis.
\subsection{Nonsplit outer actions and bounded graph quotients}

\begin{lemma}[Proved: exact transport for a nonsplit automorphism preimage]
\label{D:nonsplit-automorphism-wreath-lifts}
Let $S$ be a finite nonabelian simple group, identify $S$
with $\operatorname{Inn}(S)$, and let $\Omega$ be a finite
set. Put
\[
 W=\operatorname{Aut}(S)^\Omega\rtimes\operatorname{Sym}(\Omega),
 \qquad
 \overline W=\operatorname{Out}(S)^\Omega\rtimes
                                  \operatorname{Sym}(\Omega).
                                                        \tag{D.NW.1}
\]
Let $P\leq\overline W$ act transitively on $\Omega$,
and let $L$ be its full inverse image in $W$.
For a continuous surjection $\beta:\Gamma\to P$ from
a profinite group, choose $o\in\Omega$, put
$H=\operatorname{Stab}_\Gamma(o)$, and define
$\chi:H\to\operatorname{Out}(S)$ by the $o$-coordinate
of $\beta|_H$. Then there is a bijection
\[
 \frac{\{\rho:\Gamma\to L:\overline\rho=\beta\}}
            {S^\Omega\text{-conjugacy}}
 \ \simeq\
 \frac{\{\lambda:H\to\operatorname{Aut}(S):
                         \overline\lambda=\chi\}}
            {S\text{-conjugacy}}.                     \tag{D.NW.2}
\]
All homomorphisms are continuous. No splitting of
$1\to S^\Omega\to L\to P\to1$ is assumed.

Every finite monolithic group with nonabelian socle
$S^\Omega$ has precisely this form under its conjugation
embedding, with $P=L/S^\Omega$. If
$\Gamma\twoheadrightarrow L_k$ is a surjection to its
$k$th crown power inducing the fixed quotient $\beta$,
then $k$ is bounded by the cardinality on the right of
(D.NW.2).
\end{lemma}

\begin{proof}
Use the convention
$(a,\pi)(b,\sigma)_x=a_xb_{\pi^{-1}x}$ for the wreath
product. Choose $g_x\in\Gamma$ with $g_xo=x$, taking
$g_o=1$. Choose $b_x\in\operatorname{Aut}(S)$ lifting
the $x$-coordinate of $\beta(g_x)$, with $b_o=1$.
For $\gamma\in\Gamma$ put $y=\gamma^{-1}x$ and
\[
 h(\gamma,x)=g_x^{-1}\gamma g_y\in H.
                                                        \tag{D.NW.3}
\]
Given a lift $\lambda$ of $\chi$, define
\[
 a_x^\lambda(\gamma)
       =b_x\lambda(h(\gamma,x))b_y^{-1}.
                                                        \tag{D.NW.4}
\]
Together with the prescribed permutation of $\Omega$,
these coordinates define a homomorphism to $W$.
Indeed
$h(\gamma\eta,x)=h(\gamma,x)h(\eta,\gamma^{-1}x)$,
so (D.NW.4) satisfies the wreath cocycle identity.
The same identity applied to $\beta$ shows that its
outer coordinates are exactly those of $\beta$.
Thus the image lies in $L$. Continuity follows because
there are only finitely many cosets and all maps into
the finite groups used here are continuous.

Conversely, for a lift $\rho$ write $a_x^\rho(\gamma)$
for its coordinates and put
\[
 \lambda_\rho(h)=a_o^\rho(h),\qquad
                b_{\rho,x}=a_x^\rho(g_x).
                                                        \tag{D.NW.5}
\]
The first is a homomorphism lifting $\chi$, because
$H$ fixes $o$. The equality $\gamma g_y=g_xh(\gamma,x)$
gives
\[
 a_x^\rho(\gamma)=b_{\rho,x}
       \lambda_\rho(h(\gamma,x))b_{\rho,y}^{-1}.
                                                        \tag{D.NW.6}
\]
The outer images of $b_x$ and $b_{\rho,x}$ agree.
Hence $c_x=b_xb_{\rho,x}^{-1}$ is inner, and
$(c_x)\in S^\Omega$ has $c_o=1$. Conjugating $\rho$
by this tuple normalizes its transports to the fixed
$b_x$, giving exactly (D.NW.4).

Two such normalized homomorphisms are conjugate by
$c=(c_x)\in S^\Omega$ if and only if their evaluations
are conjugate by $c_o\in S$. Necessity follows by
restriction to $H$. For sufficiency use
$c_x=b_xc_ob_x^{-1}\in S$. Conversely any conjugacy
between normalized homomorphisms has this form, by
evaluating it on $g_x$ in coordinate $x$. This proves
the bijection, including its precise conjugacy groups.

For the assertion about monolithic groups, the
centralizer of the socle is normal and intersects the
socle in its trivial centre. If that centralizer were
nontrivial, it would contain the unique minimal normal
subgroup, a contradiction. Thus conjugation is faithful.
An automorphism of $S^\Omega$ permutes its simple
factors and acts by an automorphism on each, giving
the wreath product (D.NW.1). Its inner subgroup is
$S^\Omega$. The image of $L$ contains this full kernel
and has outer image $P$, so it is the entire inverse
image of $P$. Minimality of the socle makes the factor
permutation action transitive.

Finally, the coordinate maps of a surjection
$\Gamma\to L_k\subset L^k$ all induce the same
$\beta$. Two of them cannot be conjugate by the
socle: otherwise the image in those two coordinates
would lie in a conjugate diagonal of $L\times_P L$,
whereas surjectivity gives the whole fibre product.
They therefore yield $k$ distinct classes in (D.NW.2).
\end{proof}

\begin{proposition}[Proved under the stated counting hypothesis: bounded graph reduction for projective crown counts]
\label{D:bounded-graph-projective-crown-reduction}
Work in the complete one-curve geometric setting, and
write $\Gamma=\pi_1(U)$ and $g=g(C)$. Assume the following
\emph{projective semilinear counting input}, with one
constant $A_R$:
for every connected radially tame degree-$t$ cover with
group $H$, every $q=p^f$, every $r\geq1$, and every
tame field character $\chi:H\to\mathcal G_q$, the
number of homomorphisms
\[
 H\longrightarrow\operatorname{PGL}_r(\mathbb F_q)
                                  \rtimes\mathcal G_q
\]
with field part $\chi$, modulo
$\operatorname{PGL}_r(\mathbb F_q)$-conjugacy, is at
most $q^{A_Rtr^2}$.

Fix $b_0\geq1$ and $\kappa>0$. Let $S$ be nonabelian
simple and let $\mathcal B\leq\operatorname{Aut}(S)$
contain $S$. Suppose $\mathcal B$ embeds in a finite
group $J$ with exact sequences
\[
 \begin{split}
 1\to\operatorname{PGL}_r(\mathbb F_q)&\to J
          \to\Theta\to1,\\
 1\to\mathcal G_q&\to\Theta\to\Delta\to1,
             \qquad |\Delta|\leq b_0,
 \end{split}                                           \tag{D.NW.7}
\]
such that the inverse image of $\mathcal G_q$ is the
usual $\operatorname{PGL}_r(\mathbb F_q)\rtimes\mathcal G_q$,
the subgroup $S$ maps into $\operatorname{PGL}_r(\mathbb F_q)$,
and
\[
                    r^2\log q\leq\kappa\log|S|.
                                                        \tag{D.NW.8}
\]
Let $P\leq(\mathcal B/S)^\Omega\rtimes\operatorname{Sym}(\Omega)$
act transitively on $\Omega$, let $t=|\Omega|$, and let
$L$ be its full automorphism preimage, assumed monolithic
primitive as in the standard crown setting. For any
$\Gamma\twoheadrightarrow L_k$ whose induced quotient
$\beta:\Gamma\to P$ is radially tame, one has
\[
 k\leq |S^\Omega|^{c_R},\qquad
 d(L_k)\leq\max\{2,2g+1,c_R+2\},
 \quad c_R=\max\{1,\lceil\kappa b_0(A_R+1)\rceil\}.
                                                        \tag{D.NW.9}
\]
Thus diagonal actions may lie in the projective matrix
part, and a graph quotient of bounded order is removed
at bounded cost. The proposition is a conditional
implication from the stated projective semilinear input;
it does not establish that input.
\end{proposition}

\begin{proof}
The previous lemma reduces $k$ to the number of
$S$-conjugacy classes of lifts
$\lambda:H\to\mathcal B$ of a fixed outer character,
where $H$ has index $t$ in $\Gamma$. Its cover is
radially tame. Indeed it is a subcover of the tame
$P$-cover. The projection of $\lambda$ to $\Theta$
is fixed by the outer character, since $S$ maps into
the projective matrix kernel in (D.NW.7).

Let $H_0$ be the kernel of its further projection to
$\Delta$, and put $b=[H:H_0]\leq b_0$. This subgroup
is independent of the chosen lift. Its cover has
degree $bt$ over $U$ and is again radially tame,
because all these outer characters come from the
given tame quotient. The restriction to $H_0$ has a
fixed tame field part $\chi_0$. The assumed input
therefore bounds its projective conjugacy classes by
$q^{A_Rbtr^2}$.

For each fixed representative of a restriction class,
there are at most
\[
                  |\operatorname{PGL}_r(\mathbb F_q)|^{b-1}
                                                        \tag{D.NW.10}
\]
extensions to $H$ with the prescribed $\Theta$-part.
Choose one representative for each nonidentity coset
of $H_0$ in $H$. Its image lies in one fixed coset of
the projective matrix subgroup, so has at most that
many choices. The subgroup and these representatives
generate $H$; their images determine a homomorphism.
Relations only reduce the number. A projective
conjugation used to normalize the restriction need
not preserve $\mathcal B$: we count all possible
$J$-valued extensions, which is a valid upper bound.

Passing from $S$-conjugacy to projective matrix
conjugacy can identify at most
$|\operatorname{PGL}_r(\mathbb F_q)|$ classes in any
one orbit. To see this, fix a homomorphism in that
orbit and list its conjugates by all projective
matrices; every lift in the fibre appears in this
list before quotienting by $S$. No normalizer
assumption is needed. Since
$|\operatorname{PGL}_r(\mathbb F_q)|\leq q^{r^2}$,
the full number of $S$-classes is at most
\[
 q^{(A_Rbt+b)r^2}
       \leq q^{b_0(A_R+1)tr^2}
       \leq |S^\Omega|^{c_R}.                         \tag{D.NW.11}
\]
Together with the previous lemma this proves the
multiplicity assertion.

The tame generator bound
Proposition~\ref{D:one-curve-radial-tame-generators}
gives $d(P)\leq2g+1$. Apply the same monolithic
primitive crown estimate (D.HC.15) already used in
Corollary~\ref{D:tame-field-permutation-crowns}, namely
\citet[proof of Proposition 2.6(ii), equations
(2.1)--(2.2), pp.~11--13]{LucchiniMarionTracey2020D}.
Taking $s=\max\{2,2g+1,c_R+2\}$ gives the generator
bound in (D.NW.9). This argument nowhere chooses a
splitting of the diagonal outer extension.
\end{proof}


% Separate post-CP21 proved corollary. All four audits passed; only ordinary hereditary projective counting is used.
\subsection{Automorphism crowns with bounded stabilizer field image}

\begin{corollary}[Proved: diagonal and bounded outer-field crowns]
\label{D:bounded-field-automorphism-crowns}
Let $R,Y,C,U,\Gamma=\pi_1(U)$ have the complete one-curve
setting of Theorem~\ref{D:tame-hereditary-projective-count},
and write $g=g(C)$. Fix integers $b_0,b\geq1$ and a real
number $\kappa>0$. There is a constant $c_R$, depending
only on this fixed data, with the following property.

Let $S$ be a finite nonabelian simple group and let
$S\leq\mathcal B\leq\operatorname{Aut}(S)$. Suppose
$\mathcal B$ has an embedding into a finite group $J$
with exact sequences
\[
 \begin{split}
 1\longrightarrow\operatorname{PGL}_r(\mathbb F_q)
   &\longrightarrow J\longrightarrow\Theta\longrightarrow1,\\
 1\longrightarrow\mathcal G_q
   &\longrightarrow\Theta\longrightarrow\Delta\longrightarrow1,
                \qquad |\Delta|\leq b_0,
 \end{split}                                                   \tag{D.BF.1}
\]
where $q=p^f$, $\mathcal G_q=\operatorname{Gal}(\mathbb F_q/\mathbb F_p)$,
$S$ maps into the projective matrix kernel, and
\[
                       r^2\log q\leq\kappa\log|S|.             \tag{D.BF.2}
\]
No splitting of either displayed sequence is required.
Let $\Omega$ be a set of $t$ elements, let
$P\leq(\mathcal B/S)^\Omega\rtimes\operatorname{Sym}(\Omega)$
act transitively, and let $L$ be its full inverse image
in $\mathcal B^\Omega\rtimes\operatorname{Sym}(\Omega)$.
Assume $L$ is monolithic primitive, as in the standard
nonabelian crown setting.

Consider a surjection $\Gamma\twoheadrightarrow L_k$ whose
quotient $\beta:\Gamma\twoheadrightarrow P$ is radially
tame. For a point $o\in\Omega$, put
$H=\operatorname{Stab}_\Gamma(o)$. The fixed outer character
$\chi:H\to\mathcal B/S$ induces a fixed map
$\vartheta:H\to\Theta$. Assume
\[
          |\operatorname{im}(\vartheta)\cap\mathcal G_q|
                                                       \leq b. \tag{D.BF.3}
\]
Equivalently, after killing the bounded graph quotient,
the remaining stabilizer field image has order at most $b$.
Then
\[
 k\leq |S^\Omega|^{c_R},\qquad
 d(L_k)\leq\max\{2,2g+1,c_R+2\}.                            \tag{D.BF.4}
\]
One may take
$c_R=\max\{1,\lceil\kappa b_0b(C_R+1)\rceil\}$, where
$C_R$ is the constant in
Theorem~\ref{D:tame-hereditary-projective-count}.
The diagonal action within the projective matrix group
is unrestricted. Both the permutation degree $t$ and
simple-group rank may grow. The boundedness condition
is (D.BF.3), not bounded defining rank or bounded total
field size.
\end{corollary}

\begin{proof}
The map $\mathcal B\to\Theta$ factors through
$\mathcal B/S$, so $\vartheta$ depends only on the fixed
outer quotient and not on a chosen lift. Put
\[
       H_0=\ker\vartheta,\qquad
       a=[H:H_0]=|\operatorname{im}\vartheta|\leq b_0b.
                                                               \tag{D.BF.5}
\]
The $H_0$-cover of $U$ has degree $at$ and is radially
tame. Indeed $H$ is the inverse image of a point stabilizer
in $P$, and $H_0$ is the inverse image of a subgroup
of that stabilizer under the fixed finite quotient. Thus no new ramification assertion about a
chosen projective representation is needed.

Lemma~\ref{D:nonsplit-automorphism-wreath-lifts} bounds
$k$ by the number of $S$-conjugacy classes of lifts
$\lambda:H\to\mathcal B$ of $\chi$. Restricting such a
lift to $H_0$ gives an ordinary homomorphism into
$\operatorname{PGL}_r(\mathbb F_q)$. Its projective
conjugacy class has at most
\[
                              q^{C_Rat r^2}                    \tag{D.BF.6}
\]
possibilities, by the hereditary projective theorem.
This step uses no projective semilinear counting
statement and no equivariant extension of a gerbe.

Fix a representative of one such restriction class.
Choose one representative of each nonidentity coset
of $H_0$ in $H$. The prescribed $\Theta$-image of each
representative fixes a coset of the projective matrix
kernel in $J$, so there are at most
$|\operatorname{PGL}_r(\mathbb F_q)|^{a-1}$ possible
image tuples. A restriction and these tuples determine
a homomorphism; group relations can only remove tuples.
Conjugating a restriction to the chosen representative
may move its extension out of $\mathcal B$, but counting
all $J$-valued extensions is a valid upper bound.

Within one projective-conjugacy orbit there are at most
$|\operatorname{PGL}_r(\mathbb F_q)|$ distinct
$S$-conjugacy classes that have image in $\mathcal B$.
This elementary upper bound counts the conjugating
matrices themselves and does not require them to
normalize $\mathcal B$ or $S$. Combining it with
(D.BF.6) gives
\[
 \begin{split}
 k&\leq q^{C_Rat r^2}
                 |\operatorname{PGL}_r(\mathbb F_q)|^{a}\\
  &\leq q^{a(C_R+1)t r^2}
   \leq |S^\Omega|^{\kappa b_0b(C_R+1)}.
 \end{split}                                                   \tag{D.BF.7}
\]
We used $t\geq1$, the elementary order bound
$|\operatorname{PGL}_r(\mathbb F_q)|\leq q^{r^2}$,
and (D.BF.2).

The fixed quotient $P$ is radially tame, so
$d(P)\leq2g+1$ by
Proposition~\ref{D:one-curve-radial-tame-generators}.
The standard nonabelian crown estimate, in precisely
the same form used in (D.NW.9), now gives (D.BF.4)
\citep[proof of Proposition2.6(ii), pp.~11--13,
including (2.1)--(2.2)]{LucchiniMarionTracey2020D}.
No splitting of the socle extension, no bounded diagonal
centre, and no assertion about an arbitrary unbounded
field image has been used.
\end{proof}


% Separate candidate: fixed partial resolutions and fixed formal backgrounds.
% This file does not claim a tree representation count.
\subsection{Small marked sources and actual normal component presentations}

Retain the tree and marked spectral notation of
Lemmas~\ref{D:SNC-tree-residual-category} and
\ref{D:SNC-tree-minimal-marked-windows}. All scalar gerbes,
spectral covers, their formal lifts, and neutralizing lines in
this subsection are fixed geometric data. We do not assert a
bound for the coefficients of an arbitrarily chosen such lift.
The root indices and the character multiplicities may vary.

\begin{lemma}[Proved: small sources for marked spectral bundles]
\label{D:SNC-tree-small-marked-source}
Fix a component $i$, write $m=m_i$, $s=r/m$, and let
$\mathcal R$ be its marked spectral curve, with coarse curve
$R$ and reduced marked set $P$. At $x\in P$ write $d_x$
for the adjacent root order and $\mathcal P_x$ for the root
point. After one fixed positive ordinary exceptional twist,
every spectral bundle in the windows of
\ref{D:SNC-tree-minimal-marked-windows} has a presentation
\[
 0\longrightarrow L^{-1}\longrightarrow P_0
       \longrightarrow W\longrightarrow0,
 \qquad (P_0)_0=\mathcal O_R^{s+1},\qquad
 L=\det W_0.                                      \tag{D.TP.1}
\]
The kernel has weight zero at every marked point. The source
$P_0$ has an explicit all-ring presentation using only
homogeneous fibre frames at the marked points.

More precisely, put $h=s+1$ and
\[
 T=\mathcal O_{\mathcal R}
          \left(\sum_{x\in P}(d_x-1)\mathcal P_x\right).
                                                        \tag{D.TP.2}
\]
For each marked point choose the normalized source weights
$0\leq a_{x\ell}<d_x$, $1\leq\ell\leq h$, and put
$b_{x\ell}=d_x-1-a_{x\ell}$. Define the fixed sheaf
\[
 Q_x=\bigoplus_{\ell=1}^{h}
           T/T(-b_{x\ell}\mathcal P_x),
 \qquad Q=\bigoplus_{x\in P}Q_x.                    \tag{D.TP.3}
\]
A zero summand is omitted. A frame $g_x\in\operatorname{GL}_h$
defines the quotient $T^h\to Q_x$ by applying $g_x^{-1}$
and then the displayed component reductions. The kernel
\[
                 P(g)=\ker(T^h\longrightarrow Q)        \tag{D.TP.4}
\]
is a vector bundle over every coefficient algebra, has
coarse bundle $\mathcal O^h$, and realizes the flags defined
by the chosen frames and weights. These frames use
$O(ms^2)$ coordinates and quadratic inverse equations.
The construction extends over the formal component, using
the transverse root divisors in place of the root points.
After applying the fixed inverse spectral construction,
the quotient in (D.TP.1) lifts to that formal component.
Its kernel is a rank-$m$ formal bundle with residual
spectral line $L^{-1}$.
\end{lemma}

\begin{proof}
Choose the positive twist to be $-tA$, where $A$ is the
fixed strictly antinef cycle of
\ref{D:SNC-tree-uniform-Hom-depth}. Its degree on $R$ is
$tm(-A\cdot E_i)$. The lower slope window, genus and
number of marked points are all $O(m)$. A single $t$
therefore makes the coarse bundle generated and makes
the first cohomology of every Hom layer from a bundle
with coarse bundle $\mathcal O^h$ and normalized marked
weights vanish. Indeed that Hom is between $W_0^h(-P)$
and $W_0^h$; duality and the lower slope bound suffice.
In later root-normal layers the conormal degree is
nonnegative, and the marked normalization changes a
coarse degree by at most $|P|$. Enlarge $t$ once for
this further fixed correction. This is independent
of $m,s$ and every adjacent root order.

A general $s+1$-dimensional subspace of sections of the
generated rank-$s$ coarse bundle generates it: at a
point the failure of surjectivity has codimension two,
and the parameter of points has dimension one. Its
kernel is the line $L^{-1}$. At every marked point
split the surjection on fibres, lift a basis adapted
to the target flag, and supplement it by a basis of
the kernel fibre. Give the lifted vectors their
target weights and the kernel vector weight zero.
In particular the positive-weight flag spaces are
the chosen split lifts, not the full preimages of
the target spaces. The resulting strict parabolic
quotient has kernel with weight zero and
gives (D.TP.1) under the exact root/flag equivalence
\citep[Proposition3 and Theorem4]{Borne2007D}
\citep[Theorem6.1]{BorneVistoli2012D}.
Over a finite local Artinian coefficient algebra the
closed-fibre generators lift by the same $H^1$
vanishing and the coefficient-ideal filtration.
Nakayama preserves generation. With the determinant
marked, the kernel is the indicated fixed line.

We verify (D.TP.4) without assuming the coefficient
ring reduced. At $x$ take a root coordinate $z$ with
coarse coordinate $z^{d_x}$. A local generator of $T$
has character $-(d_x-1)$. The kernel of its reduction
modulo $z^{b_{x\ell}}$ is generated by
$z^{b_{x\ell}}$, and therefore has character
$-a_{x\ell}$. Conjugating the diagonal list by $g_x$
gives the asserted homogeneous free basis of $P(g)$.
Its invariant basis vectors are $g_x$ times the
ordinary basis of $\mathcal O^h$. Equivalently $Q_x$
has no invariants: its exponents lie in
$[0,b_{x\ell}-1]\subset[0,d_x-2]$, so none has
character zero after the shift $-(d_x-1)$. Exact
invariants yield $(P(g))_0=\mathcal O^h$.
The flags of this invariant bundle are exactly the
ones given by $g_x$ and the weight list. Frames and
their inverses use $2h^2$ entries at each of the
$O(m)$ points, with equations $g_xg_x^{-1}=1$.

On the lifted formal spectral cover, each marked
point is the intersection with a transverse root
divisor. Distinct such divisors are disjoint in this
formal component. Repeat (D.TP.2)--(D.TP.4) with those
divisors. Frame entries are coefficients from the
parameter algebra and hence define the same maps
on these transverse divisors. Their equations are
Cartier equations, regular even on the nilpotent
normal component. The same homogeneous kernel
calculation proves actual local freeness over every
coefficient algebra. Tensor the fixed neutralizing
line and apply finite pushforward from the fixed
degree-$m$ spectral cover. This is the inverse
spectral reference source on the formal component.

The first paragraph gives $H^1=0$ for all conormal
Hom layers from this source to the chosen positive
lattice. Thus the specified residual surjection
lifts successively through every normal layer,
including over a finite local Artinian coefficient
base. It stays surjective by Nakayama. Its kernel
is a vector bundle of rank $m$, with the stated
residual spectral line. This proves the lemma.
\end{proof}

\begin{proposition}[Proved: fixed-resolution component presentations and exact node gluing]
\label{D:SNC-tree-fixed-normal-presentation}
Fix a cycle of $f$ stages, the discrete marked weight lists, and
the formal background data specified above. Fix an
ordinary depth $a$ and a finite list of adjacent
Frobenius and bounded-pole inverse diagrams. Put
$M=\operatorname{lcm}_i(m_i)$ and $s_i=r/m_i$.
There is a polynomial presentation, with an open
subbundle and slope restriction, for actual bundle
and map data on the Cartier cycle $a\mathcal A$.
Every point over every coefficient algebra produces
actual objects. It covers the corresponding actual
families over finite local Artinian coefficient
algebras after choosing the fibre frames.

For component $i$, at each root-normal order
$1\leq j\leq aa_i m_i$ it uses $O(fm_i s_i^2)$
variables of weight $(M/m_i)j$. Its residual variables
number $O(fm_i s_i^2)$ and may all be assigned weight
$M$. Each original node contributes $O(fr^2)$
further variables of weight $M$. All equations have
weighted degree at most $KM$, where $K$ depends on
the fixed resolution, depth, characteristic and
diagram length, not on the root orders or ranks.

The formal reference data are constants, not moving
parameters. All residual determinant lines are covered.
The proposition does not by itself count coefficient
gerbes or punctured projective representations.
\end{proposition}

\begin{proof}
We first give the component construction, including
the reason that the number of marked points is not
squared. Suppress the component subscript and work
on its fixed spectral cover. As in
Lemma~\ref{D:fixed-partial-resolution-unit-atlas},
choose an ordinary fixed formal line $B$ whose
residual degree is a fixed positive multiple of $m$.
It may be the pullback of a fixed power of $-A$;
in particular it is invariant under spectral deck
translations. If $L=\det W_0$ after the positive
twist of the preceding lemma, then $|\deg L|=O(ms)$.
Choose $b=cs$, with a single fixed $c$ sufficiently
large that $L$ and $L^{-1}B^b$ have the $H^1$
vanishing and two-generator property needed below,
including the $O(m)$ marked corrections. The lower
slope bound for $W_0$ makes the assertion for $L$
uniform as well. Neither $L$ nor a basis of its
cohomology is a parameter in the construction.

Let $A_0$ be a fixed rank-$m$ formal anchor whose
residual spectral line is ordinary $\mathcal O$,
of weight zero at all adjacent root points. Such
a line is obtained by the selected neutralization;
its formal lift exists by $H^2=0$. Put
$P_j=(A_0\otimes B^{-jb})^{\oplus2}$ for $j\geq1$.
Keep the fixed objects $T$ and $Q_x$ of
(D.TP.2)--(D.TP.3), including their normal lifts;
tensor the scalar line and use inverse spectral
pushforward in those expressions when necessary.
Only their maps, including the frame matrices,
will vary. The kernel $K_1$ of the preceding
lemma has spectral line $L^{-1}$ with root weight
zero. Two generators of $L^{-1}B^b$ lift to
$P_1\twoheadrightarrow K_1$; its kernel has
spectral line $LB^{-2b}$. Continuing by two
generators gives the alternating spectral lines
\[
 K_{2v+1}=L^{-1}B^{-2vb},\qquad
 K_{2v}=LB^{-2vb}.                                  \tag{D.TP.5a}
\]
All are ordinary lines at the marked points.
At an odd step the generating Hom line is
$L^{-1}B^b$ and at an even step it is $L$.
Positive root-normal translations add degree
and alter normalized marked degrees by only
$O(m)$. Thus the same $H^1$ vanishing lifts
every specified presentation, including across
every small Artinian coefficient extension.

We obtain a partial complex of fixed terms
\[
 P_7\longrightarrow\cdots\longrightarrow P_1
            \longrightarrow T^h\longrightarrow Q,
 \qquad\text{in degrees }-7,\ldots,0,1.             \tag{D.TP.5}
\]
The last kernel is $P(g)$. Require $d^2=0$
as actual sheaf maps. For the first map
require the factored residual map
$P_1\to P(g)$ to have spectral rank at
least one; impose that condition on the
remaining maps $P_j\to P_{j-1}$ as well.
The condition is not imposed on the
displayed map $P_1\to T^h$: its composite
with the inclusion $P(g)\to T^h$ may
vanish on a marked root fibre. These rank
conditions are open on the proper residual
root curve. Locally the last quotient is
already surjective and $P(g)$ is free.
A unit minor in the preceding map splits
one rank-one summand. The next product-zero
equation and unit minor split its complement
in the rank-two term. Induction cancels the
contractible pairs. Consequently the only
cohomology is a locally free $F$ in degree
zero and a locally free end module in degree
$-7$. The same cancellation holds after any
nilpotent normal thickening and over every
coefficient ring, since each chosen unit
minor remains a unit. This proves actual
effectivity and base change; the generating
argument proves Artinian coverage. No
cohomology group depending on a flag matrix
or an unknown determinant has been inverted.

Here is a precise finite derived-Hom model for
these complexes. Compute the fixed termwise derived
Hom by the two-affine Cech complex and Cartier
resolutions of the individual $Q_x$ summands,
\[
       [\ T(-b\mathcal P_x)\xrightarrow{z_x^b}T\ ].
                                                        \tag{D.TP.6}
\]
The transverse parameter is a non-zero-divisor
even on the normal thickening and after every
coefficient-base change, so this is an actual
resolution there. Distinct supports are disjoint:
$R\mathcal Hom(Q_x,Q_y)=0$ for $x\ne y$.
Contract these zero complexes, and contract each
remaining fixed associated-graded Cech--Cartier
complex to its cohomology, over $k$. These are
$k$-linear strong contractions and remain so after
tensoring with any coefficient algebra.
The contraction is made separately for each pair
of the column objects in (D.TP.5); it preserves
their column degrees. In particular one does not
retain all cross-point terms of a direct sum of
global resolutions as independent variables.

We check the resulting sizes explicitly. At one
root point, each Hom or Ext calculation from
(D.TP.6) has coefficient exponents in an interval
of length at most $d_x$; root invariants require
one prescribed residue class modulo $d_x$.
Thus each pair of summands contributes at most a
fixed number of invariant coefficients. The same
statement holds after one $p$-Frobenius pullback,
when the intervals have length at most $pd_x$,
and after the fixed list of Cartier twists. Hence
all point-supported termwise Hom cohomology has
size $O(ms^2)$ per root-normal order. The top-top
block has multiplicity $h^2$ and line degree
$O(m)$. A block with exactly one positive-index
term has multiplicity at most $2h$ and line
degree $O(ms)$. Blocks with two such terms have
multiplicity four and line degree $O(ms)$.
Here $b=O(s)$, $\deg B=O(m)$,
$\deg_{\rm par}T=O(m)$, and the genus is $O(m)$.
There are only a fixed number of these blocks.
Curve Riemann--Roch gives
the same $O(ms^2)$ total. These are bounds for
the fixed termwise derived-Hom cohomology, not
for the much larger uncontracted Cech spaces.

Now insert the variable transitions and differentials.
The correction to the fixed differential is the
sum of positive-normal-order reference corrections
and the maps between columns of (D.TP.5). The
former raise normal order; the latter raise column
degree. A Hom between these fixed-width complexes
has column degrees between $-8$ and $8$. Since
the fixed contractions preserve column degree,
every nonzero homological-perturbation word
contains only a fixed number of variable column differentials;
all its other factors raise normal order.
Consequently the finite geometric series of
Lemma~\ref{D:weighted-fixed-residual-bundle-deformations}
terminates and supplies exact finite complexes
over every coefficient algebra. The same special
deformation-retract identities apply
\citep[Sections2.1--2.4 and3.2, complete proof]
{C_Crainic2004Perturbation}.

For clarity, the equation that $P_1\to T^h$ has
zero composite in $Q$ is an actual equality of
maps: express its class by a degree-zero cycle
and require it to be a boundary in the finite
derived-Hom complex. Since $P_1$ is locally free,
the degree-zero derived Hom to $Q$ is the
ordinary Hom, so this means exactly zero.
Keep its null homotopy in the total complex:
its lift goes from $P_1$ to the $Q$ column, of internal degree
$-1$ and column increase two. Adding this
column-skipping component makes the total
differential square to zero. This can also
be seen directly in (D.TP.6): the composite
into $T$ is divisible by its regular Cartier
equation, and the division is unique. Its
composition with $P_2$ is zero because the
Cartier equation is a non-zero-divisor and
$d_1d_2=0$. Thus there is no additional
unretained higher compatibility. Equation
(D.TP.5) denotes
this explicit twisted complex in the fixed
Cech--Cartier model when the selected chain
representatives do not compose strictly to
zero. It is quasi-isomorphic to the actual
sheaf complex whose composite is zero.
The added witness still strictly increases
column degree and obeys the same fixed-width
termination and weighted-degree bound.
For exact Hom between the degree-zero bundles,
use source length $L=7$ and target length $M=3$,
as in \ref{D:fixed-partial-resolution-unit-atlas}.
The canonical triangle has the same locally
free left end module as in that proof. Its
target error has cohomological degrees
$[-3,-2]$ and its source error $[4,8]$.
Thus degree-zero Hom of the separated-length
complexes is exactly ordinary Hom of $F,G$,
over every coefficient algebra. Compose
through the canonical inclusion of the
length-three target into its length-seven
complex. To test identity, precompose with
the actual vector bundle $P(g)\twoheadrightarrow F$
and compare with its canonical map. Its
two-column presentation $[T^h\to Q]$ can
be used in the same finite complexes.
This avoids a homotopy-identity condition
on an equal-length partial resolution.
Boundary variables give exact equality
and inverse witnesses. Composition is
lifted to the original Cech--Cartier model
and projected once, retaining its actual
homotopies.

At component normal order $j$, give transition
and map variables weight $\max(1,j)$ initially,
but give residual frame entries and their
inverses weight $m$. Positive-order coefficients
have additive degree in their normal order;
the bounded number of column factors add at most
a fixed multiple of $m$. For each fixed-length
arrow diagram every equation consequently has
degree $O(m)$ through depth $aa_i m$.
After rescaling the positive weights by $M/m$
and giving the residual variables weight $M$,
the bound is $KM$. One-step Frobenius multiplies
normal orders and coordinate degrees by the
fixed number $p$. Exactness of its complex
requires no coefficient-Frobenius flatness:
locally the frame diagonalizes $T^h\to Q$
into reductions modulo $z^b$, whose Frobenius
pullbacks are reductions modulo $z^{pb}$;
these have the same exact Cartier resolutions.
The locally split contractible pairs and the
locally free end modules likewise survive
every pullback. The fixed spectral translations
and root-character carries are retained in
these Frobenius diagrams.

It remains to glue the actual normal components,
not merely their reduced spectral curves.
At an original node $i,j$, write root coordinates
$x,y$, and put $k=aa_i m_i$, $l=aa_jm_j$.
Over every coefficient algebra $B$ there is
the exact fibre-product identity
\[
 B[x,y]/(x^ky^l)=B[x,y]/x^k
       \mathbin{\times}_{B[x,y]/(x^k,y^l)}B[x,y]/y^l.
                                                        \tag{D.TP.7}
\]
It holds by coefficient comparison of monomials,
also when $B$ has nilpotents. Localization and
etale pullback preserve it. Finite projective
module pasting along this fibre product is the
kernel construction of
\citep[Section15.6, Lemmas15.6.4--15.6.9,
especially Tag0D2J]{StacksTreePastingD}.
Its proof and the homogeneous kernel construction
give the equivariant assertion for the fixed
diagonalizable node group. There are no triple
intersections of distinct original components.

Choose homogeneous fibre frames on the rectangle
$B[x,y]/(x^k,y^l)$. The difference of two
scalar-weight-one characters is a character of
$\mu_{m_i}\times\mu_{m_j}$. Each matrix entry
therefore has exponents in one fixed pair of
residue classes modulo $m_i,m_j$. It has at most
\[
                       (aa_i)(aa_j)                  \tag{D.TP.8}
\]
coefficient slots. A gluing isomorphism and its
inverse use $O(r^2)$ variables and quadratic
matrix-product equations. Its products have
only the two bounded root carries. Evaluation
of a component presentation on the rectangle
uses its full perfect complex, including
(D.TP.6), so no false exactness after an
underived restriction of $Q$ is used.
Only the single marked support at this node
remains after restriction; all other $Q_x$
terms are zero. The resulting fixed-width
Hom complexes against a homogeneous free
rectangle bundle also have $O(r^2)$ invariant
slots by (D.TP.8). Their frame and inverse
witnesses therefore impose the actual
restriction identifications at this cost.
Assign these node variables weight $M$.
The fixed-width composition formulas already
proved give weighted degree $O(M)$ for their
compatibility with the component maps and
Frobenius arrows.

Milnor pasting now produces actual vector
bundles on $a\mathcal A$ with the specified
maps, over every coefficient algebra.
Conversely every Artinian family can be
framed on these finite local rectangles,
and its component presentations were covered
above. Proper-Quot openness of the residual
HN windows is the argument already proved
in \ref{D:weighted-fixed-determinant-normal-presentation};
the subbundle conditions are proper images
of determinantal failures. These are open
restrictions, not additional coefficient
variables. Summing over the fixed number
of original components and nodes proves
all the stated bounds and scope.
\end{proof}


% Separate candidate. Uses the complete actual component presentation;
% no finite-generation or hereditary-index assertion is made here.
\subsection{Projective representation counts for exceptional trees}

\begin{lemma}[Proved: finite unit comparison for a marked exceptional tree]
\label{D:SNC-tree-finite-unit-comparison}
Let $R$ be a complete excellent normal local surface
algebra over an algebraically closed coefficient field
$k$ of characteristic $p>0$, with residue field $k$.
Suppose a regular projective resolution $Y\to\operatorname{Spec}R$
is an isomorphism on the puncture $U$ and has a nonempty
SNC exceptional divisor whose dual graph is a tree.
Fix a positive integral strictly antinef cycle
$A=\sum a_iE_i$. Let $n$ be prime to $p$ and let
$\alpha\in H^2(U,\mu_n)$ be a coefficient class realized
by a projective finite local system. Use its actual
mixed-root coefficient gerbe from
\ref{D:SNC-tree-Brauer-geometry}.

The following constants can be chosen depending only
on the fixed resolution and $p$: ordinary-$A$ pole
bounds for one-step Frobenius arrows and meromorphic
comparisons, a normal depth for prescribed maps, and
a paired inverse/unit precision. They apply to all
the minimal marked lattices and all scalar Frobenius
stages, independently of $n$, the ranks, the root
orders and the number of stages.

Any finite-type polynomial family of actual bundle
and paired unit-map jets at that precision, with the
prescribed residual slope windows, has a finite-presentation
functor of its actual unit isomorphisms. The morphism
of isomorphisms from a fixed unit object to the family
is etale. Thus the represented unit-isomorphism
classes are open and closed, and their number is
at most the number of irreducible components of
the parameter scheme. The assertion concerns
actual objects and all nilpotent coefficient
extensions, not only tangent vectors.
\end{lemma}

\begin{proof}
Write $m_i$ for the residue orders, $R_i$ for the
spectral curves, $\gamma_i=g(R_i)$ and
$\eta_i=m_i\operatorname{valence}(E_i)$.
The marked minimal windows are $O(m_i)$ by
\ref{D:SNC-tree-minimal-marked-windows}. We first
check that one-step Frobenius stays in windows
of the same scale. Normalize the rational
parabolic degree of each universal spectral
line into $[0,1)$ by tensoring an ordinary
degree-one line. After the exact spectral
translation, the comparison line between a
$p$th power and the next normalized line
has rational degree between $-1$ and $p$.
At a marked point, a normalized weight $a/d$
is sent to $pa/d$; reducing it modulo one
introduces the integer $\lfloor pa/d\rfloor$
between zero and $p-1$. The resulting coarse
modification is bounded between the unmodified
ordinary Frobenius pullback and its twist by
$(p-1)$ times the reduced marked divisor.
It changes the relevant ordinary slope
endpoints by at most $(p-1)\eta_i$.

For an ordinary bundle $V$ on the smooth
spectral curve the already checked
Sun--Mehta--Pauly argument gives
\[
 \begin{split}
 \mu_{\min}(F^*V)&\geq p\mu_{\min}(V)
                  -(p-1)\max(\gamma_i-1,0),\\
 \mu_{\max}(F^*V)&\leq p\mu_{\max}(V)
                  +(p-1)\max(\gamma_i-1,0).
 \end{split}                                           \tag{D.TC.1}
\]
For positive genus this follows by adjunction
from the semistability of Frobenius pushforward
on the curve; on $\mathbb P^1$ it follows from
splitting. The exact primary statements and
adjunction proof are the ones previously read
in \citet[Theorem2.2]{C_Sun2008} and
\citet[Proposition1.3]{C_MehtaPauly2007}.
There is no assertion that tensor products
of semistable bundles remain semistable.
Since the genus and marked corrections are
$O(m_i)$, all adjacent Frobenius bundles have
enlarged windows $O(m_i)$, uniformly in rank.

The proof of
\ref{D:SNC-tree-uniform-Hom-depth} applies
unchanged to these finitely many enlarged
windows: an ordinary-$A$ layer has positive
degree $l m_i(-A\cdot E_i)+O(m_i)$.
It gives fixed ordinary pole bounds $B,C_0$
for the unit arrows, their inverses, and
meromorphic comparisons, and a fixed
specified-map depth $N$. These bounds
include the actual node intersections
because they use the Cartier-cycle
filtration on the regular root surface.

They also hold after arbitrary coefficient
base change in the finite-type families
under consideration. On a conormal layer
the Hom family is flat on a fixed proper
root-marked curve. Exact tame invariants
give its ordinary proper-curve cohomology.
The perfect cohomology complex commutes
with base change, and locally has a
two-term finite free representative in
degrees zero and one
\citep[Lemma30.22.1, complete proof, Tag07VK]
{StacksUnitJetTransferD}. If its fibre $H^1$
vanishes, its differential is locally a
split surjection; if its fibre $H^0$
vanishes, it is locally a split injection.
The positive-tail lifting and negative-pole
vanishing therefore persist over every
coefficient algebra, including its nilpotents.
Applying the finite component filtration
gives the same assertion on the actual
normal cycle, not just on its normalization.

Put $I=\mathcal O(-\mathcal A)$, where
$\mathcal A=q^*A$ on the mixed root stack.
This ideal is invertible and $F^*I=I^p$.
The unit-contraction argument of
\ref{D:meromorphic-unit-jet-contraction}
therefore applies with ordinary-$A$ order:
if two maps differ to order $t$, one
Frobenius/unit step makes their difference
have order at least
\[
                              pt-2B.                    \tag{D.TC.2}
\]
Choose a fixed precision $N_0$ larger than
the lifting depth and so that
$N_0-C_0>2B/(p-1)$. Around a cycle of $f$
stages the recurrence is
$p^f t-2B(p^f-1)/(p-1)$, whose iterates
tend to infinity at this threshold.
Approximate compatible comparison and
inverse maps correct uniquely, keeping
their full retained jets, to exact mutually
inverse unit maps. This uses only complete
separatedness and the geometric-series
inverse of $1$ plus a positive-order map;
it works over every coefficient algebra.
Increasing one fixed ordinary depth
retains all pole-cleared products and
one-step substitutions needed for this
construction. No root-index multiple
is hidden in (D.TC.2).

We spell out finite-cycle effectivity
before applying algebraization. A bundle
on the retained Cartier cycle lifts
through its successive square-zero
normal extensions: the obstruction is
$H^2(\mathcal End F\otimes I^a/I^{a+1})$,
which is zero on the proper one-dimensional
tame support over an affine coefficient
base. Exact invariants and the two-affine
Cech cover give this vanishing over every
coefficient algebra. The positive-tail
$H^1$ vanishing proved above then lifts
each specified arrow and inverse. After
dividing by the fixed pole, an inverse
product is $1$ plus positive normal
order and is corrected by its geometric
series. The unit-comparison compatibility
is corrected by (D.TC.2), retaining the
prescribed jets. Thus an actual formal
system exists before existence theory
is invoked.

For the algebraization comparison choose
an algebraic generating bundle on the
specified coefficient gerbe as follows.
Put $M=\operatorname{lcm}_i(m_i)$.
The cyclic degree-$M$ etale cover
$Y'\to Y$ in
\ref{D:SNC-tree-Brauer-geometry} kills
its geometric Brauer class on the
puncture. The same regular root-stack
purity and uniqueness proof makes the
pulled-back class zero on the mixed
root stack itself
\citep[Proposition3.2.1 and complete proof]
{Lieblich2011ResidualD}. The pulled-back
scalar gerbe therefore has a weight-one
invertible sheaf $\Lambda$. Push it
forward to obtain an algebraic bundle,
then tensor with the direct sum of
the root line bundles whose exponents
run through all the root characters.
The resulting $A_{\rm alg}$ contains
every scalar-weight-one character at
every root chart. It is an algebraic
projective generator of that graded
module category. Its possibly large
rank has no role in the coordinate
count.

Put $\mathcal E=q_*\mathcal End(A_{\rm alg})$.
It is a coherent locally free algebra
on $Y$. The coarse Hom of a formal
bundle is locally free by the regular
homogeneous-basis calculation. Ordinary
Grothendieck existence algebraizes it
together with its $\mathcal E$ action
and pole maps. Local Morita descent
recovers a bundle on the actual gerbe:
$\mathcal Hom(A_{\rm alg},F)$ is a
right $\mathcal E$-module and is tensored
with the left module $A_{\rm alg}$.
For parameter coefficients $B$ the
formal construction is on the
$\mathfrak m$-adic completion of
$R\otimes_k B$, not on an assumed
complete ordinary tensor product.
The existence theorem is applied to
this complete base when $B$ is of
finite type, and in particular to
the geometric and finite Artinian
bases needed for the unit comparison.
For finite local Artinian $B$ the
ordinary tensor product is already
complete. The preceding formal
base-change and contraction statements
do not require algebraization for
an arbitrary non-Noetherian $B$.

This is the ordinary coarse-Hom proof
of \ref{D:one-curve-coefficient-root-gerbe},
now with an algebraic generating bundle
on all the crossing charts. The same
Morita functor extends a given punctured
bundle coherently by taking the span
under the finite algebra $\mathcal E$.
Its reflexive hull is locally free
on the regular two-dimensional root
charts and retains the scalar weight.
Thus the initial lattices used in the
minimal-window argument also exist.

Now represent comparison and inverse jets
by Hom schemes on the fixed proper finite
cycle. Their representability and finite
presentation are the exact proper-Hom
input of \citep[Proposition99.3.10, Tag08K6]
{StacksUnitJetTransferD}. Impose the finite
unit equations and the two pole-cleared
inverse products at the precision just
chosen. The all-base lifting and unique
contraction identify the resulting
finite-presentation scheme with the
\emph{actual} unit-Isom functor. In
particular it is not merely a space
of approximate maps whose higher
compatibility has been forgotten.

On a nilpotent coefficient extension,
a unit bundle on the punctured gerbe
is its finite etale fixed-frame torsor.
The coefficient band and scalar weight
are retained. Finite etale covers and
all their morphisms are unchanged by
nilpotent thickenings
\citep[Lemma58.8.3, complete proof, Tag0BQB]
{StacksUnitJetTransferD}; the assertion
descends on an etale atlas of the gerbe.
Every specified unit isomorphism over
an Artinian quotient consequently lifts
uniquely. Its unique bounded-pole lift
is the one just represented. The
finite-presentation infinitesimal
criterion gives an etale fixed-object
Isom morphism
\citep[Tags02HT,02HM and00UR]{StacksUnitJetTransferD}.
Its image is open. The complement is
the union of the images for the other
unit classes, so each image is also
closed. A nonempty class consumes at
least one irreducible component. This
proves the stated finite comparison.
\end{proof}

\begin{theorem}[Proved: uniform projective counts for an exceptional tree]
\label{D:SNC-tree-projective-quadratic-count}
Under the hypotheses of the preceding lemma there
is a constant $C_R$ such that, for every $q=p^f$
and every $r\geq1$,
\[
 \#\left(\operatorname{Hom}_{\rm cont}
       (\pi_1(U),\operatorname{PGL}_r(\mathbb F_q))
                    /\operatorname{PGL}_r(\mathbb F_q)\right)
                         \leq q^{C_Rr^2}.              \tag{D.TC.3}
\]
No bound on the geometric Brauer period is imposed.
The curve genera and the number of components belong
to the fixed resolution. This is a projective
representation bound for the fixed local germ;
it asserts neither a uniform hereditary-index
constant nor full finite generation of its
fundamental group.
\end{theorem}

\begin{proof}
Put $n=q-1$, let $G=\sum_i g(E_i)$, and let
$N=-(E_i\cdot E_j)$. A projective representation
has scalar lifting obstruction in
$H^2(\pi_1(U),\mathbb F_q^*)$, which injects
into $H^2(U,\mu_n)$ by the proved degree-two
fundamental-group comparison. The exact
tree computation gives
\[
 |H^2(U,\mu_n)|=|\operatorname{coker}(N\bmod n)|n^{2G}
                      \leq |\det N|\,q^{2G}.           \tag{D.TC.4}
\]
For $n=1$ this coefficient set is trivial.
It suffices to bound representations for
each realized class $\alpha$, uniformly
in that class. We use only realized
classes; choose a witnessing finite
projective local system when constructing
the corresponding lifting gerbe.

Let $m_i$ be the orders of its residue
characters and put $M=\operatorname{lcm}_i(m_i)$.
The Brauer identification for a tree
makes $M$ exactly the geometric Brauer
period. A rank-$r$ twisted bundle has
$m_i\mid r$ on every component by the
connected spectral rank formula. Thus
$M\mid r$, $M\leq r$, and $m_i\mid n$.
Write $s_i=r/m_i$. The coefficient class
extends to the actual mixed-root gerbe;
all scalar Frobenius stages have the
same orders $m_i$. Their covers and
neutralizations may be fixed separately
for this class. They are constants in
the polynomial presentations. The
one-step comparison lines have the
degree bounds proved above; no degree
bound for arbitrary reference cochains
is required, since no determinant or
anchor coefficient is a moving parameter.

Choose the fixed ordinary depth $a$
given by the finite-unit comparison.
For each discrete set of marked and
joint node character multiplicities,
apply \ref{D:SNC-tree-fixed-normal-presentation}
to all $f$ Frobenius stages and to the
paired unit/inverse diagrams. Its
points produce actual bundles and
maps over every coefficient ring,
and all required Artinian families
are covered. Restrict to its open
rank and residual-slope locus.

Here is the complete weighted budget.
For component $i$ put
\[
 J_i=aa_i m_i,\qquad
 b_i=\lceil A_0 f m_i s_i^2\rceil.                     \tag{D.TC.5}
\]
Pad to $b_i$ variables at each
positive weight $(M/m_i)j$, $1\leq j\leq J_i$,
and $b_i$ residual variables of weight
$M$. Add $c\leq A_0fr^2$ node variables
of weight $M$, increasing the constant
for the fixed number of original nodes.
Unused variables do not change the
number of irreducible components.
The total number $v$ of variables
is $O(fr^2)$, and all equations have
weighted degree at most $KM$.
The proved weighted affine component
bound gives
\[
 \begin{split}
 \#\operatorname{Irr}
 &\leq (v+1)
 \frac{(KM)^v}
 {M^{c+\sum_i b_i}
  \prod_i\bigl((M/m_i)^{J_i}J_i!\bigr)^{b_i}}\\
 &=(v+1)K^{c+\sum_i b_i}
       \prod_i\left(\frac{(K m_i)^{J_i}}{J_i!}\right)^{b_i}
 \leq \exp(C_0fr^2).
 \end{split}                                           \tag{D.TC.6}
\]
The last inequality uses $J_i!\geq(J_i/e)^{J_i}$
and the fixed identity $J_i/m_i=aa_i$.
In particular no factor $\log(M/m_i)$
remains: the weight rescaling cancels
that ratio separately on each component.
All the weights are positive integers,
including when the characteristic
divides a weight; the component lemma
already treats that case by finite
flat coordinate powering and lengths.

The number of discrete types does not
spoil (D.TC.6). At a node, scalar-weight-one
characters form a torsor under the
$m_im_j$ characters of its two root
groups. An ordered list of $r$ characters
has at most $(m_im_j)^r$ possibilities.
At the $m_i$ marked spectral points
over a node, the adjacent weight lists
have at most $m_j^{m_is_i}=m_j^r$
possibilities. These overcount compatible
types, which is harmless. There are
only fixed numbers of original nodes
and components and $f$ stages, so the
total is at most
$\exp(C_1fr\log(2r))\leq\exp(C_2fr^2)$.
For $r=1$ the statement also follows
directly from the triviality of
$\operatorname{PGL}_1$.

For each type the previous lemma
identifies actual punctured unit
isomorphism classes with open-and-closed
classes in this finite presentation.
It therefore bounds their number by
(D.TC.6). There is no need to identify
the auxiliary left-end modules of
the partial resolutions: exact
separated-length Hom and the
$P(g)$ identity test represent
isomorphisms of the degree-zero
bundles themselves.

Finally, the scalar-weight-one
$\mathbb F_q$ unit objects on the
fixed lifting gerbe map surjectively
to the projective representation
classes with obstruction $\alpha$.
The witnessing finite quotient
describes that gerbe as its actual
central lifting gerbe; chosen
matrix lifts can be adjusted by
the scalar coboundary of the fixed
cocycle. This is the same exact
coefficient comparison as in
\ref{D:one-curve-unbounded-period-projective-count},
and does not assume that a root
gerbe's band always injects into
its fundamental group. The possibly
larger number of twisted objects
over one projective object only
improves the upper bound.
Multiply (D.TC.6) and the type bound
by (D.TC.4), and use $q=p^f$.
After increasing one constant this
is (D.TC.3), in the exact scope stated.
\end{proof}


% Separate post-CP21 proved packet. All four complete audits passed; coherent full-root extension and compressed sector data.
\subsection{Coherent coefficient gerbes and single-sector normal compression}

\begin{lemma}[Proved: coherent extension of a realized tame-band lifting gerbe]
\label{D:coherent-full-root-coefficient-gerbe}
Let $(Y_H,C_H,U_H)$ be the regular one-curve model of a
connected radially tame cover of the fixed local puncture.
Let $q=p^f$, $n=q-1>1$, and fix a tame field character
$\chi:\pi_1(U_H)\to\operatorname{Gal}(\mathbb F_q/\mathbb F_p)$.
Write $a=|\operatorname{im}\chi|$, $\nu=f/a$, and let
$V\to U_H$ be the corresponding connected cyclic cover.
Its normalization $Y_V\to Y_H$ has unique smooth boundary
$C_V$, with radial inertia order $e_\chi$.
Let $\mathcal M=\mathbb F_q^*(\chi)$, regarded as a finite
etale cyclic band. Suppose a class in
$H^2(\pi_1(U_H),\mathcal M)$ is realized by a projective
semilinear finite representation, and let $\mathscr G_U$
be its actual central lifting gerbe on $U_H$.

The band extends to the root stack
$X_e=\sqrt[e_\chi]{(Y_H,C_H)}$. The specified gerbe extends,
with its band, to
\[
 X_{en}=\sqrt[n]{(X_e,D_e)}
                  =\sqrt[e_\chi n]{(Y_H,C_H)}.          \tag{D.SC.1}
\]
The cover $Y_V\to X_e$ is a representable finite etale
cyclic torsor of degree $a$. Its base change to $X_{en}$
is exactly
\[
                    X_n=\sqrt[n]{(Y_V,C_V)}.            \tag{D.SC.2}
\]
Consequently the pulled-back scalar $\mu_n$-gerbe
$\mathscr G\to X_n$ carries the actual coherent deck
identifications, including their order-$a$ cocycle.
No independent choice of a geometric-period gerbe and
no splitting of this deck coherence is asserted.
\end{lemma}

\begin{proof}
The strict-local tame charts of
Lemma~\ref{D:radial-tame-cover-one-curve} identify the
map with $z=w^{e_\chi}$. They identify the quotient
stack $[Y_V/\operatorname{im}\chi]$ with $X_e$:
etale-locally the inertia is the cyclic group of order
$e_\chi$, and its quotient chart is the indicated root
chart. The canonical root line $\mathcal O(C_V)$ and
its section give the global map. The local identifications
agree by the universal property of the root construction.
Thus $Y_V\to X_e$ is the stated torsor. The character
band is its associated locally constant sheaf, so it
extends to $X_e$. Taking the further $n$th root commutes
with this base change and gives (D.SC.2).

We construct the gerbe extension with all its descent
data, rather than merely extending a Brauer class.
Use regular etale root charts of $X_e$, and shrink
etale-locally to make the band constant. On a strict
localization of such a chart, with its smooth boundary
removed, the constant prime-to-$p$ group has vanishing
$H^2$. This is the exact purity/Kummer calculation in
Lemma~\ref{D:radial-tame-cyclic-H2}. The degree-two
fundamental-group comparison is injective there.
Hence the restricted finite central lifting extension
has zero obstruction and has a continuous lift; its
finite target makes this an actual finite-etale object
of the lifting gerbe. The object descends to an etale
neighborhood of the strict localization. Equivalently,
one uses the finite-presentation and cohomology limit
arguments of \citet[Lemma59.51.2 and Theorem59.51.3,
complete written proofs]{StacksCoherentRootSectorsB}.
We obtain an etale cover $\{W_i\to X_e\}$ with actual
objects of $\mathscr G_U$ on $W_i\cap U_H$.

On an overlap, the isomorphisms between the two objects
form a finite etale $\mathcal M$-torsor $T_{ij}$.
The band is unramified on the overlap. At a strict-local
boundary chart it is constant; wild inertia is pro-$p$
and maps trivially to a group of order $n$. Thus the
torsor is tame and its inertia order divides $n$.
Abhyankar's Kummer normalization and purity extend it,
uniquely as a finite etale torsor, after the $n$th root
of the boundary
\citep[full Lemma58.31.5, Tag0EYG]{StacksRadialTameGeometryB}.
This is an assertion about torsors, not about a normal
subgroup generated by all wild inertia being pro-$p$.

The composition isomorphisms
$T_{ij}\mathbin{\times}^{\mathcal M}T_{jk}\simeq T_{ik}$
extend uniquely. Indeed finite etale morphisms over a
normal root chart are determined by their restriction
to the dense open, and their isomorphisms extend by
normality and finite normalization. The same argument
extends their associativity and unit identities.
Gluing the trivial banded gerbes with these extended
torsors and compositions constructs a gerbe on $X_{en}$.
It is algebraic and etale-locally the classifying stack
of the finite etale band. Its restriction is the specified
lifting gerbe, since the original objects and all their
isomorphism torsors were used. Pullback along the actual
cyclic torsor (D.SC.2) now supplies the deck 2-action
and its full cocycle automatically. None is chosen anew.
\end{proof}

\begin{proposition}[Proved: single-sector Morita lattices and compressed normal maps]
\label{D:single-sector-root-compression}
Let $Y_V,C_V$ be a complete one-curve model and let
$\mathscr G\to X_n=\sqrt[n]{(Y_V,C_V)}$ be a scalar
$\mu_n$-gerbe, with $n$ prime to $p$. Let $m$ be the
period of its geometric Brauer class on the puncture,
put $h=n/m$, and write $X_m=\sqrt[m]{(Y_V,C_V)}$.
Assume its residual Brauer character has order $m$,
as is the case in Lemma~\ref{D:coherent-full-root-coefficient-gerbe}.
Let $D_n,D_m$ be the reduced root divisors, and let
$\rho:X_n\to X_m$ be the canonical map. Then
$\rho^*D_m=hD_n$.

For each invertible scalar weight $u\in(\mathbb Z/n)^*$
and each chosen connected component of its residual
spectral torsor, there is an algebraic rank-$m$ anchor
$A_u$ supported on that component. Its algebra
$\mathcal End(A_u)$ descends to an Azumaya algebra
$\mathcal B_u$ on $X_m$. Every scalar-weight-$u$
bundle on the puncture admits a lattice on $\mathscr G$
whose residual spectral bundle is supported on this
one chosen component. Such lattices, and their actual
families, are Morita equivalent to locally free
$\mathcal B_u$-modules on $X_m$.

For two objects of the same scalar weight in fixed
sectors, their Hom has one fixed $\mu_h$ character.
There is a unique $c\in\{0,\ldots,h-1\}$ such that
\[
 \mathcal Hom(F,G)(-cD_n)=\rho^*\mathcal H_{F,G}
                                                        \tag{D.SC.3}
\]
for a vector bundle on $X_m$; the scalar gerbe has
already canceled in this formula. Locally the coefficient
exponents of maps are $c+h\ell$. Therefore the invertible
ideal
\[
                  I=\mathcal O(-hD_n)=\rho^*\mathcal O(-D_m)
                                                        \tag{D.SC.4}
\]
retains only one compressed normal layer for each value
of $\ell$. Composition has the usual one-step carry
and one $p$-Frobenius has carry at most $p-1$.
These are identities over arbitrary coefficient rings.
They are not decompositions inferred only from geometric
fibres. A final coherent deck operation is retained as
that operation, and need not fix the selected sector.
\end{proposition}

\begin{proof}
The residual character calculation of
Proposition~\ref{D:residual-twisted-spectral-category}
gives a $\mathbb Z/n$-torsor with $h$ connected
components, each of degree $m$. Its monodromy subgroup
is $h\mathbb Z/n$. At a geometric point of the curve,
the weights in one component are a coset of this
subgroup in the character torsor above the scalar
weight $u$.

We can choose the anchor algebraically. The connected
cyclic degree-$m$ residue cover lifts to a finite etale
cover $Y'\to Y_V$, as in
Proposition~\ref{D:single-curve-Brauer-residue-cover}.
It kills the geometric Brauer class. Root-stack purity
and its uniqueness statement show that the pullback
of that Brauer class to $\sqrt[n]{(Y',C')}$ is zero
\citep[Proposition3.2.1 and complete proof]
{Lieblich2011ResidualD}.
The pulled-back scalar gerbe therefore has an invertible
sheaf of scalar weight $u$. Its residual root character
can be translated to the selected component by a power
of the root line. Its residual ordinary line can be
normalized by a line bundle on $Y'$, using the already
proved Picard lifting. Its finite pushforward is an
algebraic rank-$m$ anchor $A_u$, with exactly the
characters of that component once each. The same
construction applies to every invertible scalar weight.

Here is the precise descent calculation. On an etale
chart write the normal root coordinate as $s$, so the
map $X_n\to X_m$ is $t=s^h$. The inverse image of $\mu_h$ in the total inertia
acts on a homogeneous basis of $A_u$ through one
character restricting to $u$ on the scalar band:
differences of any two basis characters are divisible
by $h$. No splitting of that central inertia extension
is required. Consequently the quotient $\mu_h$ acts
trivially on the scalar-weight-zero End algebra. Homogeneous Nakayama
preserves these basis weights throughout the formal
chart. Matrix basis weights of $\mathcal End(A_u)$
are consequently divisible by $h$. Taking invariants
over the ring with coordinate $t=s^h$ gives a free
algebra whose pullback is $\mathcal End(A_u)$.
Multiplication and the identity descend. The local
algebra is Azumaya because this can be tested after
the faithfully flat root chart. These local algebras
and their maps descend uniquely, giving $\mathcal B_u$.
This proof also works over an arbitrary coefficient
ring for actual bundles in the selected sector.

If $E$ is a punctured twisted bundle, set
$M_U=\mathcal Hom(A_u|_U,E)$, a right
$\mathcal B_u|_U$-module. Extend it coherently, take
the span under the finite locally free algebra
$\mathcal B_u$, and take its reflexive hull on the
regular two-dimensional charts of $X_m$. Its module
action and associativity extend uniquely; tensoring
with a locally free algebra preserves reflexivity.
A reflexive module on these regular surfaces is
locally free. Morita evaluation
\[
                  \rho^*M\otimes_{\rho^*\mathcal B_u} A_u
                                                        \tag{D.SC.5}
\]
returns the required lattice of $E$. Its $\mu_h$
character is the selected scalar one. Conversely
Hom from $A_u$ descends for every such lattice, and
local matrix Morita equivalence proves that these
operations are inverse. This is also the exact
all-ring equivalence for actual locally free module
families; no relative reflexive-hull flatness is
asserted. The reflexive-hull argument is used only
to cover geometric punctured objects.

The same homogeneous basis proves (D.SC.3). Define
$c$ by saying that the Hom fibre character on
$\mu_h$ is $-c$. Tensoring by $\mathcal O(-cD_n)$
cancels this character. Its degree-zero invariant
basis descends just as above. A map has coefficient
powers $s^{c+h\ell}$, including negative $\ell$
when poles are allowed. If two offsets are $c_1,c_2$
and the composite offset is $c_3$, then
\[
 c_1+c_2=c_3+h\epsilon,\quad \epsilon\in\{0,1\};
 \qquad pc=c'+h\epsilon_p,\quad0\leq\epsilon_p\leq p-1.
                                                        \tag{D.SC.6}
\]
These follow from division with remainder in the
integer character grading. They remain exact over
nonreduced coefficient algebras. Endomorphisms and
all within-phase object differentials have offset
zero. In particular no $h$ extra object-deformation
layers occur.

A deck automorphism preserves $D_n$, the ideal $I$,
and normal order. Its induced action on the character
torsors can translate the sector labels, especially
when its scalar-band action is nontrivial. Formula
(D.SC.3) is applied with the resulting actual sector
labels. We do not replace this action or its cocycle
by an identity. Finally algebraization of formal
objects and fixed-pole maps follows by using the
algebraic generator formed from $A_u$ and its $h$
root-line translates, then ordinary coarse-Hom
existence and Morita as in
Proposition~\ref{D:one-curve-coefficient-root-gerbe}.
That possibly large generator is used for existence,
not in any coordinate estimate.
\end{proof}

\begin{lemma}[Proved: compressed skew-Frobenius jets with exact deck coherence]
\label{D:compressed-skew-Frobenius-unit-jets}
Retain Lemma~\ref{D:coherent-full-root-coefficient-gerbe}.
Write the total tame degree of $Y_V$ over the original
one-curve resolution as $Ev$, with radial index $E$
and residue degree $v$. Thus
\[
 g(C_V)-1=v(g-1),\qquad -C_V^2=d'=v\delta/E.
                                                        \tag{D.SC.7}
\]
Let $m$ be the geometric Brauer period after this cover,
let $r=ms$, and let $\mathscr T$ denote the actual
coherent deck pullback functor, oriented so that it
carries scalar weight one to scalar weight $p^\nu$.
For the $\nu$ one-step Frobenius stages choose one
sector and one rank-$m$ anchor at each stage, with
the terminal choice $\mathscr T(A_0)$. Then the actual
unit diagrams
\[
 F^*F_j\dashrightarrow F_{j+1}\quad(0\leq j<\nu-1),
 \qquad F^*F_{\nu-1}\dashrightarrow\mathscr T(F_0)
                                                        \tag{D.SC.8}
\]
with paired inverses admit the same all-ring finite
presentation and unit-Isom comparison as the fixed
partial-resolution construction. Uniform compressed
precision is $J=cEm$ for one fixed integer $c$.
At each positive compressed weight $1\leq\ell\leq J$
there are $O(\nu v m s^2)$ variables, and the residual
variables satisfy the same bound and have weight $Em$.
All equations have weighted degree $O(Em)$.
The constants depend only on the original resolution
and $p$, independently of $n,h,a,E,v,m,r,\nu$.

The exact unit objects in (D.SC.8) are the twisted
semilinear local systems on the original coefficient
gerbe with the fixed field character. Their finite
presentation Isom functor is etale at every point
against a fixed object. No order-$a$ matrix identity
is added to the finite presentation: its coherence
comes from the actual deck action and the fixed-frame
construction.
\end{lemma}

\begin{proof}
Choose the normalized universal spectral lines of
Proposition~\ref{D:residual-twisted-Frobenius-transport}.
Their rational degrees on each selected connected
spectral curve lie in $[0,1)$. At the terminal stage
use the deck pullback of the initial normalized line.
It has the same degree. The connected spectral curves
have degree $m$ over $C_V$ and genus
$1+vm(g-1)$. On each of them
\[
                  \deg\mathcal O(D_n)=-d'/h.
                                                        \tag{D.SC.9}
\]
Aligning two sectors by an offset $0\leq c<h$ changes
a comparison degree by at most $d'$. A $p$-power
comparison of normalized universal lines changes
it by at most $p+1$; a deck comparison preserves
degree. Consequently every required one-step line
has absolute degree at most $p+1+d'=O(vm)$.
The constant in this scaled bound is independent of
all cover and coefficient indices. An arbitrary
reference cochain can have complicated coefficients;
those coefficients are constants, not parameters.

Minimal lattices are now chosen in the selected
sector by the ordinary Euler-minimum construction,
using the ideal $I$ of (D.SC.4). Perform elementary
transforms in the descended Azumaya-module category
on $X_m$ and transport them by $A_u$. They remain
in the selected sector, and the quotient Euler
characteristic is that of its single degree-$m$
spectral curve, with no factor $h$. The compressed root
normal line is $\mathcal O(D_m)$ and its spectral
degree is $-d'$. The proof of the connected minimal
window therefore gives the endpoints
$[g(R_j)-1,g(R_j)-1+d']$, just as on $X_m$.
One-step Frobenius and the preceding comparison
lines enlarge them only by $O(vm)$. A fixed positive
twist by a multiple of $-EC_V$ gives the generating
windows. All positive Hom tails and negative pole
layers vanish beyond $O(Em)$ compressed steps,
because each such step has positive degree $d'$
whereas the genus and endpoint errors are $O(vm)$.
The exact layer and Artinian proofs used for
\ref{D:tame-resolution-hereditary-linear-count}
apply after (D.SC.3); an offset adds less than one
compressed pole order.

Apply the fixed partial resolutions of
\ref{D:fixed-partial-resolution-unit-atlas} with
these anchors. Within each phase all terms and
object differentials have one sector, hence offset
zero. Their kernels and quotients remain in that
sector over every coefficient algebra, by the
actual Morita description, homogeneous Nakayama,
and the unit-minor cancellation of that lemma.
The single connected spectral curve still has
rank $s$, so the four Hom-block estimates are
$O(vm s^2)$ per compressed layer. Adjacent phase
Hom and its one-step Frobenius source are aligned
by (D.SC.3); their additional line degrees are the
$O(vm)$ quantities already bounded above.

All maps are kept as actual maps on $\mathscr G$.
The integer reindexing (D.SC.6) does not replace
them by an unrelated graded category. It gives
finite normal cochain models with $J$ compressed
layers. Their multiplication and Frobenius use
the displayed carries, fixed reference maps and
the actual final deck functor. The finite
homological-perturbation construction, with its
fixed source/target internal lengths, works
unchanged over every coefficient algebra.
Within-phase variable differentials have offset
zero. The fixed-length arrow/equality diagrams
contain only a bounded number of cross-phase
factors; their carries add a bounded number of
compressed orders, and one Frobenius has bounded
factor $p$. Thus positive orders have additive
weights, residual variables can be assigned
weight $Em$, and all equations have weighted
degree $O(Em)$. No integer $h=n/m$ is charged as
a number of layers, a polynomial degree, or a
number of sector choices per object. One sector
per phase was fixed before making the atlas.

For exact unit comparison use the invertible
ideal $I$ itself, not its uncompressed generator.
It satisfies $F^*I=I^p$. Pole-bounded conjugation
therefore improves order $t$ to at least $pt-2b$
with $b=O(Em)$. The final deck step preserves $I$.
The earlier paired-jet contraction, with a fixed
larger compressed precision, consequently gives
actual specified unit isomorphisms. It also gives
the proper finite-cycle Hom presentation of the
actual Isom functor. Normal and coefficient
nilpotents are treated by the same exact graded
and proper-cohomology arguments as before.

Finally, iterate (D.SC.8) exactly $a$ times using
the coherent identification $\mathscr T^a\simeq1$.
This produces a $q$-Frobenius unit object. Its
finite etale $\operatorname{GL}_r(\mathbb F_q)$
fixed-frame torsor inherits the semilinear deck
action supplied by the original $\nu$-step arrow.
The $a$th-power condition holds on fixed frames
because their $q$-Frobenius action is the identity,
and because the original gerbe deck cocycle has
been retained. Conversely the semilinear finite
frame torsor of a local system gives (D.SC.8).
This is the exact skew-unit/frame argument of
\ref{D:twisted-deck-Frobenius-semilinear-model},
applied on an etale atlas of the coefficient gerbe;
all its maps descend by the given band and deck
coherence. Iteration is used only for this
categorical verification. The finite equations
still have only $\nu$ one-step stages, not an
expanded $a$-fold matrix word.

Finite etale torsors and every specified equivariant
isomorphism are unchanged by nilpotent thickenings.
Their unique lifts respect the deck cocycle, since
both sides restrict to the same specified map.
The finite-presentation proper-Hom model and the
paired contraction therefore give an etale
fixed-object Isom morphism exactly as in
\ref{D:finite-jet-unit-Isom}. This proves the
asserted comparison and the precise presentation
scope. It does not omit the actual coefficient
class in favor of its geometric Brauer image.
\end{proof}


% Separate candidate: full projective semilinear count and its crown consequence.
\subsection{Unbounded field actions in projective hereditary counts}

\begin{theorem}[Proved: radially tame projective semilinear hereditary count]
\label{D:tame-projective-semilinear-count}
Let $R$ be a complete excellent normal local surface algebra over
an algebraically closed field $k$ of characteristic $p>0$, with
residue field $k$. Suppose a regular projective resolution
$Y\to\operatorname{Spec}R$ is an isomorphism on the puncture
$U$ and has one smooth connected exceptional curve $C$.
Put $g=g(C)$, $C^2=-\delta<0$, and $\Gamma=\pi_1(U)$.
There is a constant $A_R$ with the following property.

Let $H\leq\Gamma$ be an open subgroup of index $t$ whose
connected cover $U_H\to U$ is radially tame along $C$.
Let $q=p^f$, $r\geq1$, and fix a continuous character
$\chi:H\to\mathcal G_q=\operatorname{Gal}(\mathbb F_q/\mathbb F_p)$
whose cover is tame along the exceptional curve of $Y_H$.
Let $N_\chi(H,r,q)$ denote the number of continuous homomorphisms
$H\to\operatorname{PGL}_r(\mathbb F_q)\rtimes\mathcal G_q$
with field part $\chi$, modulo conjugacy by
$\operatorname{PGL}_r(\mathbb F_q)$. Then
\[
                   N_\chi(H,r,q)\leq q^{A_Rtr^2}.       \tag{D.PS.1}
\]
No bound on $|\operatorname{im}\chi|$, the
scalar order $q-1$, or the geometric Brauer period is assumed.
Only the cover of $H$ and the specified field character are
required to be radially tame; the projective representation
itself may have wild geometric inertia.
\end{theorem}

\begin{proof}
For $r=1$ there is exactly one homomorphism with the prescribed
field part. If $q=2$, the field group and the scalar group are
trivial, and the hereditary linear count gives the assertion.
We may therefore take $r\geq2$ and $n=q-1>1$.

Write $a=|\operatorname{im}\chi|$ and $\nu=f/a$.
The central lifting sequence with its specified field action is
\[
 1\longrightarrow\mathbb F_q^*
 \longrightarrow\operatorname{GL}_r(\mathbb F_q)\rtimes\mathcal G_q
 \longrightarrow\operatorname{PGL}_r(\mathbb F_q)\rtimes\mathcal G_q
 \longrightarrow1.                                    \tag{D.PS.2}
\]
After pullback the kernel has coefficient $H$-module
$\mathcal M=\mathbb F_q^*(\chi)$ after pullback along a
representation with field part $\chi$; it is central over
the linear subgroup but need not be central in the full
semilinear group. Its obstruction is a class
$\alpha\in H^2(H,\mathcal M)$.
Lemma~\ref{D:radial-tame-cyclic-H2} gives
\[
                  |H^2(H,\mathcal M)|\leq q^{c_Rt}      \tag{D.PS.3}
\]
for a constant depending only on the original curve.
This is the full group-cohomology obstruction and includes
its Kummer part. It does not use inflation of degree-two
cohomology through the normal subgroup generated by wild
inertia. We count only classes realized by a representation.

Fix one such $\alpha$ and its actual lifting gerbe.
Let $V\to U_H$ be the field cover. The total tame cover
$Y_V\to Y$ has radial index $E$ and residue degree $v$,
so
\[
 Ev=ta,\qquad \nu Ev=tf,\qquad
 g(C_V)-1=v(g-1),\qquad -C_V^2=v\delta/E.               \tag{D.PS.4}
\]
These are the exact tame cover identities, with $E,v$
measured over the original resolution. In particular
their constants are uniform as $t$ and $a$ grow.

Lemma~\ref{D:coherent-full-root-coefficient-gerbe} extends
the specified twisted-band gerbe first over the full root
supporting its finite isomorphism torsors. Its pullback to
$\sqrt[n]{(Y_V,C_V)}$ retains the actual deck identifications
and their cocycle. Let $m$ be the geometric Brauer period
there. Since an actual rank-$r$ twisted bundle exists,
the connected residual character calculation gives $m\mid r$;
also $m\mid n$. Put $s=r/m$ and $h=n/m$.

Proposition~\ref{D:single-sector-root-compression} permits
one fixed connected spectral sector at each scalar phase,
without excluding any punctured object. The construction
extends its Morita module and keeps the coefficient gerbe;
it does not replace that gerbe by another class with the
same Brauer image. Its fixed rank-$m$ anchors, their
reference coefficients and the coherent deck operation
are geometric constants. No number of their coefficients
is charged as a variable count. The normal ideal
$I=\mathcal O(-hD_n)$ has one layer for each compressed
root order. Cross-phase maps retain their actual sector
offset, with only bounded carries in products and
one-step Frobenius.

We spell out the counting estimate so that the cancellation
of the field-cover order is explicit. By
Lemma~\ref{D:compressed-skew-Frobenius-unit-jets} the
actual $\nu$-stage bundle, arrow and paired inverse data
have a polynomial presentation over $k$ with
\[
 J=cEm,\qquad b=\lceil c_1\nu v m s^2\rceil             \tag{D.PS.5}
\]
variables at each positive weight $1,\ldots,J$, and at
most $b$ residual variables of weight $Em$. Enlarge
the fixed constants to include the bounded number of
Hom, equality and inverse witnesses. The equation
weighted degree is at most $KEm$. We may pad each
block with unused variables, since a product with an
affine space does not change the number of irreducible
components. The total number $w$ of variables then
satisfies $w\leq c_2\nu Ev r^2$.

The weighted component estimate gives
\[
 \begin{split}
 \#\operatorname{Irr}
 &\leq (w+1)\frac{(KEm)^{b(J+1)}}{(J!)^b(Em)^b}\\
 &=(w+1)K^b
           \left(\frac{(KEm)^J}{J!}\right)^b
 \leq\exp(c_3\nu Ev r^2).
 \end{split}                                          \tag{D.PS.6}
\]
The last inequality follows from $J!\geq(J/e)^J$ and
$J/(Em)=c$, which is fixed. Since $\nu v m s^2\geq1$,
the ceiling in $b$ is absorbed by the same constant.
There is no residual factor involving $h$, $\log m$,
or the field-image order $a$. Restricting to the open
subbundle and slope locus cannot increase the number
of irreducible components. Fixed partial resolutions
cover every determinant without a moving Picard factor.
All sectors were fixed once; the single-sector Morita
coverage requires no sum over $h$ choices for each object.

The same lemma supplies an actual finite-presentation
unit-Isom functor with open-and-closed class fibres.
Thus (D.PS.6) bounds the number of represented twisted
unit classes. This comparison holds on all finite
Artinian coefficient extensions, after actual formal
bundle and map lifting. It is not inferred from a
tangent dimension or from the component estimate alone.

For clarity, only $\nu$ one-step Frobenius arrows occur
in the equations. The last arrow has target the actual
deck pullback of the initial bundle. This deck operation
is $k$-linear on coefficient variables and preserves
the compressed normal ideal; its scalar-band action is
part of the fixed geometric data. Iterating the skew
cycle $a$ times, with its supplied gerbe coherence,
produces the $q$-unit structure used to identify finite
frames. On those frames the required deck descent holds
by their fixed-frame equation. This iteration proves
the categorical equivalence; it is not an additional
degree-$a$ matrix equation in the parameter scheme.
Nor is it an equation asserting that a based coset lift
has order $a$ in the fundamental group.

The twisted semilinear finite frames on the fixed lifting
gerbe map onto the projective semilinear representations
with obstruction $\alpha$. Indeed the defining central
extension supplies the gerbe, and two representative
cocycles for $\alpha$ differ by a scalar coboundary;
adjusting the chosen matrix lifts by that coboundary
gives the prescribed coefficient gerbe. Conversely its
scalar-weight-one frames project to the required
projective representation. Distinct twisted objects
may project to the same projective class, which only
improves this upper bound. All band and deck coherence
is retained in this construction.

Finally, (D.PS.4) turns (D.PS.6) into
$\exp(c_3tf r^2)\leq q^{c_4tr^2}$.
Multiplication by (D.PS.3) and enlargement of one
constant proves (D.PS.1), uniformly in every listed
index. This proves the theorem in exactly its stated
radial-tame-cover scope.
\end{proof}

\begin{corollary}[Proved: projective automorphism crowns with unbounded field image]
\label{D:tame-projective-automorphism-crowns}
Retain the fixed local surface and $\Gamma$ of the preceding
theorem. Fix $b_0\geq1$ and $\kappa>0$. Let $S$ be a finite
nonabelian simple group and $S\leq\mathcal B\leq\operatorname{Aut}(S)$.
Assume $\mathcal B$ embeds into a finite group $J$ with
\[
 \begin{split}
 1\longrightarrow\operatorname{PGL}_r(\mathbb F_q)
    &\longrightarrow J\longrightarrow\Theta\longrightarrow1,\\
 1\longrightarrow\mathcal G_q
    &\longrightarrow\Theta\longrightarrow\Delta\longrightarrow1,
                   \qquad |\Delta|\leq b_0,
 \end{split}                                           \tag{D.PS.7}
\]
where $q=p^f$, the inverse image of $\mathcal G_q$ is the
usual $\operatorname{PGL}_r(\mathbb F_q)\rtimes\mathcal G_q$,
$S$ maps into the projective matrix kernel, and
$r^2\log q\leq\kappa\log|S|$.
Let $\Omega$ have $t$ elements, let
$P\leq(\mathcal B/S)^\Omega\rtimes\operatorname{Sym}(\Omega)$
act transitively, and let $L$ be its full inverse image in
$\mathcal B^\Omega\rtimes\operatorname{Sym}(\Omega)$.
Suppose $L$ is monolithic primitive.

For every continuous surjection $\Gamma\twoheadrightarrow L_k$
whose quotient to $P$ is radially tame, one has
\[
 k\leq |S^\Omega|^{c_R},\qquad
 d(L_k)\leq\max\{2,2g+1,c_R+2\},\qquad
 c_R=\max\{1,\lceil\kappa b_0(A_R+1)\rceil\}.
                                                        \tag{D.PS.8}
\]
The diagonal action and the stabilizer field-image order
are unrestricted. The hypotheses on the faithful projective
embedding, its order comparison, the bounded graph quotient
and radial tame ramification remain part of the statement.
\end{corollary}

\begin{proof}
The preceding theorem proves exactly the uniform projective
semilinear counting hypothesis of
Proposition~\ref{D:bounded-graph-projective-crown-reduction}.
Every subgroup used there is the inverse image of a subgroup
of the fixed tame outer quotient, so its cover and its
specified field character satisfy the theorem's hypotheses.
That proposition therefore applies without any remaining
counting assumption. Its exact nonsplit wreath lifting
bijection bounds the number of distinct crown coordinates
by the corresponding stabilizer lift classes. Killing the
graph part costs an index at most $b_0$; the projective
extension and conjugacy costs are at most $q^{b_0r^2}$.
Together these give $q^{b_0(A_R+1)tr^2}$, which is at most
$|S^\Omega|^{c_R}$ by the assumed order comparison.
The radial quotient has at most $2g+1$ generators, and the
proved monolithic crown inequality then gives (D.PS.8)
\citep[proof of Proposition2.6(ii), equations(2.1)--(2.2),
pp.~11--13]{LucchiniMarionTracey2020D}.
No bound on the field-image order enters these constants.
\end{proof}


% Proved post-CP21 corollary. Closure of the same projective crown reduction.
\subsection{Bounded outer quotients without an ambient graph action}

\begin{corollary}[Proved: projective crowns from a bounded-index normal subgroup]
\label{D:normal-subgroup-projective-crowns}
Retain the complete one-smooth-exceptional-curve setting of
Theorem~\ref{D:tame-projective-semilinear-count}. Write
$\Gamma=\pi_1(U)$, $g=g(C)$, and let $A_R$ be the constant
in that theorem. Fix an integer $b_0\geq1$ and a real
number $\kappa>0$.

Let $S$ be a finite nonabelian simple group and suppose
\[
 S\leq\mathcal B_0\mathrel{\trianglelefteq}\mathcal B
       \leq\operatorname{Aut}(S),\qquad
              [\mathcal B:\mathcal B_0]\leq b_0.
                                                        \tag{D.NP.1}
\]
Here $S$ is identified with its inner automorphism group.
Assume there is an injective homomorphism
\[
 \iota:\mathcal B_0\hookrightarrow
       \operatorname{PGL}_r(\mathbb F_q)
                  \rtimes\mathcal G_q,
 \qquad \iota(S)\subset\operatorname{PGL}_r(\mathbb F_q),
 \qquad \mathcal G_q=\operatorname{Gal}(\mathbb F_q/\mathbb F_p),
                                                        \tag{D.NP.2}
\]
where $q=p^f$, such that
\[
                       r^2\log q\leq\kappa\log|S|.
                                                        \tag{D.NP.3}
\]
No embedding of $\mathcal B$ in an ambient projective
semilinear group or extension of its outer action to that
ambient group is assumed.

Let $\Omega$ be a set of $t$ elements. Let
$P\leq(\mathcal B/S)^\Omega\rtimes\operatorname{Sym}(\Omega)$
act transitively, and let $L$ be its full inverse image in
$\mathcal B^\Omega\rtimes\operatorname{Sym}(\Omega)$.
Assume $L$ is monolithic primitive, with socle $S^\Omega$.
For every continuous surjection $\Gamma\twoheadrightarrow L_k$
whose induced quotient $\beta:\Gamma\twoheadrightarrow P$
is radially tame, one has
\[
 \begin{split}
 k&\leq |S^\Omega|^{c_R},\\
 d(L_k)&\leq\max\{2,2g+1,c_R+2\},\\
 c_R&=\max\{1,\lceil\kappa(A_Rb_0+1)+b_0-1\rceil\}.
 \end{split}                                           \tag{D.NP.4}
\]
The field-image order inside $\mathcal B_0$ is unrestricted.
There is also no independent order bound on $\mathcal B$.
The normal-subgroup hypothesis (D.NP.1), the projective
size comparison (D.NP.3), and tameness of the outer quotient
remain explicit hypotheses.
\end{corollary}

\begin{proof}
Choose $o\in\Omega$, let $H$ be the inverse image in
$\Gamma$ of its point stabilizer in $P$, and let
$\chi:H\to\mathcal B/S$ be the fixed $o$-coordinate
of $\beta|_H$. Then $[\Gamma:H]=t$.
The transport formulas of
Lemma~\ref{D:nonsplit-automorphism-wreath-lifts} take
values in $\mathcal B$ when the chosen transport lifts
do. They therefore give the exact bijection
\[
 \frac{\{\Gamma\to L\text{ lifting }\beta\}}{S^\Omega}
       \simeq
 \frac{\{\lambda:H\to\mathcal B\text{ lifting }\chi\}}{S},
                                                        \tag{D.NP.5}
\]
where both quotients are conjugacy quotients. In particular
surjectivity to $L_k$ forces its $k$ coordinates to have
distinct classes on the right. It suffices to bound that
set; no splitting of the socle extension is used.

Put
\[
 H_0=\chi^{-1}(\mathcal B_0/S),\qquad
          a=[H:H_0]\leq b_0.
                                                        \tag{D.NP.6}
\]
This is a normal open subgroup of $H$ determined by the
fixed outer quotient, independently of the lift
$\lambda$. Its cover of $U$ has degree $at$ and is
radially tame. The field projection of (D.NP.2) kills
$S$, so it factors through a fixed homomorphism
$\varphi:\mathcal B_0/S\to\mathcal G_q$. Thus every
restriction $\lambda|_{H_0}$ has the same field part
\[
                    \chi_0=\varphi\circ\chi|_{H_0}.
                                                        \tag{D.NP.7}
\]
The cover defined by this character is radially tame
as well. Both its subgroup and $H_0$ are inverse images
of subgroups of the fixed finite tame quotient $P$;
the tame tower and normalization statements of
\ref{D:radial-tame-cover-one-curve} apply. This step
places no tameness requirement on the geometric part
of a chosen lift.

Theorem~\ref{D:tame-projective-semilinear-count} bounds
ambient $\operatorname{PGL}_r(\mathbb F_q)$-conjugacy
classes of the restricted homomorphisms by
$q^{A_Rat r^2}$. Within one ambient orbit, there are
at most $|\operatorname{PGL}_r(\mathbb F_q)|$ distinct
$S$-conjugacy classes of actual $\mathcal B_0$-valued
restrictions. Indeed, fix any one homomorphism in the
ambient orbit and list its conjugates by all projective
matrices. Every actual restriction in the orbit occurs
in this list before quotienting by $S$. This upper
bound does not require the conjugating matrices to
normalize $\mathcal B_0$ or $S$. Consequently the number
of restriction classes modulo $S$ is at most
\[
 q^{A_Rat r^2}|\operatorname{PGL}_r(\mathbb F_q)|
                       \leq q^{(A_Rat+1)r^2}.          \tag{D.NP.8}
\]

For each such class choose an \emph{actual}
$\mathcal B_0$-valued representative. Normalize a
restriction to it using an element of $S$. Since
$S\subset\mathcal B$, this conjugation preserves the
condition that the full lift is $\mathcal B$-valued
and leaves its outer character $\chi$ unchanged.
Choose one representative of each nonidentity coset
of $H_0$ in $H$. Its image under a lift must lie in
the fixed coset of $S$ prescribed by $\chi$, so has
at most $|S|$ possibilities. A restriction and these
$a-1$ images determine the homomorphism; relations
only remove choices. Thus the extension cost is
\[
                              |S|^{a-1}.               \tag{D.NP.9}
\]
In particular it is unnecessary to bound $|\mathcal B|$
or to realize the conjugation by these elements as an
automorphism of the ambient projective matrix group.

Combining (D.NP.5), (D.NP.8) and (D.NP.9), then using
(D.NP.3), gives
\[
 \begin{split}
 k&\leq q^{(A_Rat+1)r^2}|S|^{a-1}\\
  &\leq |S|^{\kappa(A_Rat+1)+a-1}\\
  &\leq |S^\Omega|^{\kappa(A_Rb_0+1)+b_0-1}
   \leq |S^\Omega|^{c_R}.
 \end{split}                                           \tag{D.NP.10}
\]
Here $a\leq b_0$ and $t\geq1$ are the only estimates
used to pass to the third line. The radial tame
quotient $P$ has at most $2g+1$ generators by
\ref{D:one-curve-radial-tame-generators}. The exact
nonabelian crown inequality already used in
\ref{D:tame-projective-automorphism-crowns} therefore
gives the generator bound in (D.NP.4)
\citep[proof of Proposition2.6(ii), equations(2.1)--(2.2),
pp.~11--13]{LucchiniMarionTracey2020D}.
This is a closure of the same restricted crown reduction;
it does not assert full finite generation of $\Gamma$.
\end{proof}


% Separate post-CP22 candidate. No general-graph count is inferred from this file alone.
\subsection{Secondary residues and noncommuting node inertia}

\begin{lemma}[Proved: arbitrary SNC residues and mixed-root coefficient gerbes]
\label{D:SNC-general-Brauer-geometry}
Let \(R\) be a complete excellent normal local surface algebra
over an algebraically closed coefficient and residue field \(k\)
of characteristic \(p>0\). Let \(Y\to\operatorname{Spec}R\)
be a projective regular resolution, an isomorphism on the
puncture \(U\), whose nonempty reduced exceptional divisor
\(E=\bigcup_{i=1}^{v}E_i\) is strict normal crossings.
Its connected dual graph \(G\) may have cycles or parallel
edges. Write \(g_i=g(E_i)\), \(b_1=b_1(G)\), and
\(N=(-(E_i\cdot E_j))_{ij}\).
For every \(n\) prime to \(p\), there is an exact sequence
\[
 0\longrightarrow\bigoplus_iH^1(E_i,\mathbb Z/n)
 \longrightarrow\operatorname{Br}'(U)[n]
 \longrightarrow H_1(G,\mathbb Z/n)(-1)
 \longrightarrow0,                                    \tag{D.GG.1}
\]
and
\[
 |H^2(U,\mu_n)|
 =|\operatorname{coker}(N\bmod n)|\,
       n^{\,2\sum_i g_i+b_1}
 \leq|\det N|\,n^{\,2\sum_i g_i+b_1}.                  \tag{D.GG.2}
\]
The last map in (D.GG.1) records the secondary residues.

Fix \(\alpha\in H^2(U,\mu_n)\) whose Brauer image
\(\beta\) is represented by an Azumaya algebra on \(U\);
in particular this holds for a projective finite local
system. Let \(m_i\) be the order of the generic residue
of \(\beta\) on \(E_i\), with \(m_i=1\) for zero residue.
Then \(M=\operatorname{lcm}_i m_i\) is the order of
\(\beta\), and the actual coefficient class \(\alpha\)
extends to the mixed-root stack
\[
 \mathcal X=\prod_Y\sqrt[m_i]{(Y,E_i)}.                \tag{D.GG.3}
\]
At an edge \(x\in E_i\cap E_j\), let \(h_x\) be the
order of its secondary residue. Then
\(h_x\mid\gcd(m_i,m_j)\). The inertia of the extended
\(\mu_n\)-gerbe at \(x\) is a central extension
\[
 1\longrightarrow\mu_n\longrightarrow D_x
 \longrightarrow\mu_{m_i}\times\mu_{m_j}
 \longrightarrow1,                                   \tag{D.GG.4}
\]
whose scalar-weight-one commutator has order \(h_x\).
It need not be commutative. There is an algebraic
scalar-weight-one generating bundle on the coefficient
gerbe. Its rank is not used in a coordinate bound.
\end{lemma}

\begin{proof}
The Picard, proper-base-change and support calculations
in Lemma~\ref{D:SNC-tree-Brauer-geometry}, up to its
use of the leaf argument, did not assume a tree.
For completeness, their exact output here is
\[
 \operatorname{Pic}(Y)/n\simeq H^2(Y,\mu_n)
       \simeq(\mathbb Z/n)^v,\qquad H^3(Y,\mu_n)=0,
 \qquad\operatorname{Br}'(Y)[n]=0.                    \tag{D.GG.5}
\]
Multidegrees on the nodal curve are surjective, its
degree-zero Picard group is an extension of Jacobians
by a torus, and the formal Picard kernel is uniquely
\(n\)-divisible. These facts and ordinary existence
give (D.GG.5) in the complete excellent scope, as in
the cited lemma; the graph torus causes no failure
of divisibility.

The same absolute-purity support triangle gives
\[
 0\to\bigoplus_iH^1(E_i,\mathbb Z/n)
 \to H^3_E(Y,\mu_n)\to(\mathbb Z/n)(-1)^{\#\mathrm{edges}}
 \xrightarrow{\partial}(\mathbb Z/n)(-1)^v.           \tag{D.GG.6}
\]
Here \(\partial\) is the signed incidence map, by the
composition compatibility of the point and divisor
Gysin maps. The exact regular-scheme purity scope is
\citet[Expose XVI, Sections2--3, Theorem3.1.1 and the
complete proof of Corollary3.1.3]{A_ILO2014}.
Choose a spanning tree. Its incidence columns identify
the subgroup of vertex tuples with sum zero, over
every \(\mathbb Z/n\); each remaining edge gives one
free cycle parameter. Thus
\(\ker\partial=H_1(G,\mathbb Z/n)\simeq(\mathbb Z/n)^{b_1}\).
Localization, (D.GG.5), and
\(\operatorname{Pic}(Y)\twoheadrightarrow\operatorname{Pic}(U)\)
identify \(H^3_E(Y,\mu_n)\) with
\(\operatorname{Br}'(U)[n]\). The image of the exceptional
divisor lines is the intersection matrix \(-N\).
Consequently the kernel of \(H^2(U,\mu_n)\to H^3_E(Y,\mu_n)\)
has order \(|\operatorname{coker}(N\bmod n)|\).
This proves (D.GG.1)--(D.GG.2); Smith normal form gives
the determinant bound. No splitting of (D.GG.1) is needed.

Purity identifies the generic residues with cyclic
characters of \(k(E_i)\), unramified outside the nodes.
At a node their two secondary residues are opposite.
Their common order divides each generic character order,
giving \(h_x\mid m_i,m_j\). If all generic residues of
a multiple of \(\beta\) vanish, purity extends it to
\(Y\), where (D.GG.5) makes it zero. This proves the
assertion about \(M\).

We use the actual Azumaya extension proof of
\citet[Proposition3.2.1, complete pp10--11]
{Lieblich2011ResidualD}, not merely a cohomological
extension assertion. At the generic point of \(E_i\),
the \(m_i\)-th root kills its residue. The DVR order
construction in that proof is local on the component.
It therefore applies with these componentwise orders.
On the regular mixed-root surface, reflexive extension
of the resulting Azumaya order from codimension one is
locally free. Its multiplication map to its endomorphism
algebra is an isomorphism in codimension one, hence
everywhere by the determinant criterion. This is the
same codimension-two purity step in that proof.
The unique Brauer extension has the same period as
its generic restriction. Kummer therefore lifts it to
\(H^2(\mathcal X,\mu_n)\); correcting by a line from
\(\operatorname{Pic}(Y)\) makes its restriction exactly
\(\alpha\). In particular the scalar Kummer data have
not been discarded.

We record why the commutator in (D.GG.4) is exactly
the secondary residue. On the strictly henselian
regular two-dimensional local ring at \(x\), with
parameters \(z_i,z_j\), the support calculation gives
\[
 H^2(\operatorname{Spec}A[1/(z_i z_j)],\mu_n)
       \simeq\mathbb Z/n(-1).
\]
The cup product of the Kummer classes of \(z_i,z_j\)
is a generator: its successive divisor and point
residue is \(1\). Thus the local Brauer class is a
power of this symbol, of order \(h_x\). The cyclic
algebra calculation after taking the roots in (D.GG.3)
identifies that symbol with the commutator of the two
inertia lifts. This is the calculation of
\citet[Lemma4.3.1, Proposition4.3.2 and the explicit
matrices in Proposition4.3.7, pp13--15]
{Lieblich2011ResidualD}. Those statements are written
there for equal prime orders; the mixed-order scalar
algebra needed here is computed explicitly below.
Changing the coefficient lift by a Kummer line changes
the powers of the two lifts, not this commutator.

Finally choose an algebraic twisted bundle \(V\)
representing the extended Azumaya algebra on the
actual \(\mu_n\)-gerbe. Such a bundle is obtained by
pulling the tautological bundle from the associated
\(\mathbf G_m\)-gerbe. Tensor \(V\) with the direct
sum of all root-character lines
\(\mathcal O(\sum_i a_i\mathcal D_i)\), \(0\leq a_i<m_i\).
At a node the ordinary character group of
\(\mu_{m_i}\times\mu_{m_j}\) acts transitively on the
scalar-weight-one irreducibles: in the explicit
matrix-algebra computation below it acts transitively
on the eigenvalue labels of \(X^{h_x},Y^{h_x}\).
The analogous assertion at a generic divisor is the
cyclic character calculation. Hence this bundle,
denoted \(G_0\), contains every relevant inertia
irreducible at every point. On a finite root chart
the equivariant locally free module is projective,
because the group order is invertible. Fibrewise
generation and Nakayama make it a generator.
Thus coarse Hom from \(G_0\) and evaluation give the
ordinary finite-algebra Morita description used for
existence and algebraization. No étale cover of all
of \(Y\) splitting the ramified residues is asserted.
\end{proof}

\begin{proposition}[Proved: ramified spectral curves and exact noncommuting node pasting]
\label{D:SNC-general-residual-category}
Retain (D.GG.3)--(D.GG.4) and fix the coefficient gerbe.
The scalar-weight-one category on normalized component
\(i\) is, after a fixed neutralizing line, the category
of vector bundles on
\[
 \mathcal R_i=
   \sqrt[m_j/h_x]{(R_i,P_{ix})}_{x\in E_i\cap E_j},
                                                        \tag{D.GG.7}
\]
where \(R_i\to E_i\) is the connected cyclic cover of
degree \(m_i\) determined by its generic residue.
It is tame of ramification index \(h_x\) at \(x\);
\(|P_{ix}|=m_i/h_x\). The ordinary spectral rank is
\(s_i=r/m_i\), and
\[
 \begin{split}
 2g(R_i)-2
   &=m_i(2g_i-2)+
                   \sum_{x\in E_i\cap E_j}m_i(1-1/h_x),\\
 \#\operatorname{Marks}(\mathcal R_i)
   &=\sum_{x\in E_i\cap E_j}m_i/h_x
     \leq m_i\operatorname{valence}(i).
 \end{split}                                           \tag{D.GG.8}
\]
Both the genus and the marked-point count are \(O(m_i)\)
with constants depending only on \(Y,E\).

At a node \(x\in E_i\cap E_j\), every scalar-weight-one
irreducible of \(D_x\) has dimension \(h_x\), and there
are \(m_i m_j/h_x^2\) such irreducibles.
Reduced-curve bundles are exactly the component bundles
with \(D_x\)-equivariant identifications at each node.
On the multiplicity-frame atlas, for a fixed node type
\(V=\bigoplus_\lambda V_\lambda^{\oplus a_\lambda}\),
the identifications and inverse identifications use
\[
       2\sum_\lambda a_\lambda^2\leq 2r^2             \tag{D.GG.9}
\]
coordinates and quadratic equations, over every
coefficient algebra.

On the root-node rectangle
\[
 B[s,t]/(s^{a a_i m_i},t^{a a_j m_j}),                 \tag{D.GG.10}
\]
the equivariant Hom between rank-\(r\) and rank-\(r'\)
objects has at most
\((a a_i)(a a_j)rr'\) coefficient slots after choosing
equivariant frames. This statement retains the actual
noncommuting \(D_x\)-action. It concerns actual
families; it does not infer their existence from
geometric fibres alone.
\end{proposition}

\begin{proof}
First compute the finite node algebra for arbitrary
orders, including all scalar Kummer powers. Put
\(m_i=b,m_j=c,h=h_x\). Choosing lifts of the two
quotient generators, their action on the fixed scalar
weight is governed by
\[
 X^b=\lambda,\quad Y^c=\mu,\quad
 XY=\zeta YX,\qquad \lambda,\mu\in k^*,\quad
             \operatorname{ord}(\zeta)=h.             \tag{D.GG.11}
\]
The monomials \(X^aY^d\), \(0\leq a<b,0\leq d<c\),
form a basis of this scalar block of the group algebra.
Its centre is generated by \(X^h,Y^h\): commuting with
both generators forces both exponents to be divisible
by \(h\). It is the split étale algebra of dimension
\((b/h)(c/h)\) given by
\((X^h)^{b/h}=\lambda\), \((Y^h)^{c/h}=\mu\).
For each pair of its eigenvalues, the corresponding
block has dimension \(h^2\). The usual diagonal
matrix with eigenvalues \(1,\zeta,\ldots,\zeta^{h-1}\)
and the cyclic shift matrix, multiplied by suitable
fixed scalars, realize it as \(M_h(k)\).
They satisfy the prescribed powers, not powers chosen
by deleting the Kummer part. This proves
\[
        k[D_x]_{\text{scalar }1}
               \simeq\prod_{bc/h^2}M_h(k).             \tag{D.GG.12}
\]
It also proves the transitivity used above: multiplying
\(X,Y\) by ordinary \(b\)-th and \(c\)-th roots of unity
acts transitively on the eigenvalues of \(X^h,Y^h\).
All these identities remain valid after tensoring
with any \(k\)-algebra. Multiplicity modules can be
projective; the stated frame atlas trivializes them.

Now apply the relative character construction of
Proposition~\ref{D:residual-twisted-spectral-category}
to the generic root inertia over \(E_i\).
The adjacent root inertia at \(x\) translates these
characters by their commutator, so its permutation
orbits have length \(h_x\). The character torsor is
therefore finite étale over the root-marked \(E_i\).
Its coarse cover is the cyclic residue cover \(R_i\);
the stabilizer at a point above \(x\) is the kernel
of this permutation, of order \(m_j/h_x\).
The strict-local chart is exactly the tame chart
\(z=u^{h_x}\) with the residual root of order
\(m_j/h_x\). This proves (D.GG.7), including the number
of points and their orders, without claiming that
\(R_i\to E_i\) is étale at a node.

The residual category on this spectral root curve
is a Picard gerbe. Its prime-to-\(p\) Brauer class
vanishes: the coarse smooth proper curve has zero
Brauer group, a cyclic point root has
\(R^1\pi_*\mathbf G_m\) a point-supported
\(\mathbb Z/d\) and \(R^2\pi_*\mathbf G_m=0\), and
the point-supported sheaves have zero \(H^1\).
This is the same exact Leray argument, using
\citet[Lemmas4.1.1--4.1.3 and Proposition4.1.4]
{Lieblich2011ResidualD}, as in the tree category.
A neutralizing line gives the stated equivalence,
whose inverse is finite pushforward of that line
tensored with the spectral bundle. The rank is
\(r=m_i s_i\). Tame Riemann--Hurwitz on \(R_i\to E_i\)
gives (D.GG.8); each of the \(m_i/h_x\) points has
different exponent \(h_x-1\).

The spectral torsor and its cyclic action lift uniquely
through nilpotent normal thickenings, by the exact
étale-site invariance already used for the tree
construction. Here the base is the formal component
with the adjacent root divisors retained; it is not
the coarse surface with those roots removed.
Consequently ramification of the coarse \(R_i\) does
not obstruct this finite étale formal lift. The
neutralizing Picard line also lifts: its obstruction
is coherent \(H^2\) on a tame proper curve, which is
zero, and the prime-to-\(p\) formal Picard kernel is
divisible. Fix such lifts as geometric references.
No polynomial bound for their fixed coefficients
is asserted or needed here.

At a node, (D.GG.12) describes precisely how the two
branch spectral fibres induce a \(D_x\)-representation.
They glue if and only if these representations are
isomorphic. For fixed isotypic multiplicities, Schur
matrix blocks give (D.GG.9). Over an arbitrary
coefficient algebra this is ordinary Morita for the
split algebras in (D.GG.12), after passing to the
frame atlas, not semisimplicity asserted fibrewise.
The fibre-product ring
\(B[s,t]/(st)=B[s]\times_B B[t]\), and the explicit
kernel construction for projective modules over it,
give actual node pasting. Exact invariants preserve
this construction because \(|D_x|\) is invertible.
The full module statements are
\citet[Lemmas15.6.4--15.6.9, complete proofs]
{StacksTreePastingD}. Descent from the root charts
therefore gives the asserted reduced-curve category.

Finally the scalar centre acts trivially on
\(\operatorname{Hom}(V,W)\), even though it acts
nontrivially on \(V,W\). Hence this Hom representation
factors through the abelian group
\(\mu_{m_i}\times\mu_{m_j}\) and splits into its exact
character summands. The variables \(s,t\) carry the
two corresponding root characters. For each Hom
character, equivariance fixes one residue class for
each exponent modulo \(m_i,m_j\). The rectangle in
(D.GG.10) contains exactly \(a a_i\) and \(a a_j\)
choices in those two classes. Summing the dimensions
of the Hom character spaces proves the slot bound.
An actual equivariant bundle on this nilpotent
rectangle can be framed from its origin representation:
lift a frame equivariantly using exact invariants;
Nakayama makes the lift an isomorphism. This works
over every coefficient algebra locally on the
projective-multiplicity frame atlas. Thus the slot
calculation describes actual all-ring maps and their
compositions, while retaining the noncommuting node
inertia on the bundles themselves.
\end{proof}


% Separate post-CP22 proved corollary. All four complete audits passed; recognition within the same projective crown reduction.
\subsection{Classical defining-characteristic automorphism crowns}

\begin{corollary}[Proved: classical families with unrestricted diagonal, field and graph actions]
\label{D:classical-defining-characteristic-crowns}
Retain the complete one-smooth-exceptional-curve setting of
Theorem~\ref{D:tame-projective-semilinear-count}, and let
$A_R$ be its hereditary counting constant. There is one
constant $c_R$, depending only on that germ, such that the
following holds for every nonabelian simple classical group
$S$ in defining characteristic $p$ listed below.
For every subgroup
$S\leq\mathcal B\leq\operatorname{Aut}(S)$, every transitive
\[
 P\leq(\mathcal B/S)^\Omega\rtimes\operatorname{Sym}(\Omega),
 \qquad t=|\Omega|,
                                                        \tag{D.CA.1}
\]
and its full preimage
$L\leq\mathcal B^\Omega\rtimes\operatorname{Sym}(\Omega)$,
assumed monolithic primitive with socle $S^\Omega$, a
surjection $\Gamma\twoheadrightarrow L_k$ with radially
tame outer quotient $\Gamma\twoheadrightarrow P$ satisfies
\[
 k\leq |S^\Omega|^{c_R},\qquad
 d(L_k)\leq\max\{2,2g(C)+1,c_R+2\}.                    \tag{D.CA.2}
\]
For the standard simple ranges in the table one may take
\[
                     c_R=\max\{1,\lceil96A_R+21\rceil\}.
                                                        \tag{D.CA.3}
\]
There is no bound on classical rank, field degree, scalar
centre, diagonal action, field-image order, or graph action.
The outer quotient in (D.CA.1) is still required to be
radially tame. This conclusion does not include groups
whose defining characteristic differs from $p$.

Here $q=p^f$; $d$ is the natural projective dimension and
$Q$ is the coefficient-field order used in the counting
theorem. The last column is the index of the normal
inner-diagonal-field subgroup $\mathcal A_0$ in
$\operatorname{Aut}(S)$, in Steinberg's standard root
coordinates.
\[
\begin{array}{c|c|c|c}
 S\text{ and range}&d&Q&[\operatorname{Aut}(S):\mathcal A_0]\\ \hline
 \operatorname{PSL}_d(q),\ d\geq2,
       (d,q)\notin\{(2,2),(2,3)\}
       &d&q&1\ (d=2),\ 2\ (d\geq3)\\
 \operatorname{PSU}_d(q),\ d\geq3,
       (d,q)\ne(3,2)
       &d&q^2&1\\
 \operatorname{PSp}_{2\ell}(q),\ \ell\geq2,
       (\ell,q)\ne(2,2)
       &2\ell&q&2\ (\ell=2,\ p=2),\ 1\ \text{otherwise}\\
 \operatorname{P}\Omega_{2\ell+1}(q),\ \ell\geq3,
       q\text{ odd}
       &2\ell+1&q&1\\
 \operatorname{P}\Omega^+_{2\ell}(q),\ \ell\geq4
       &2\ell&q&6\ (\ell=4),\ 2\ (\ell\geq5)\\
 \operatorname{P}\Omega^-_{2\ell}(q),\ \ell\geq4
       &2\ell&q^2&1
\end{array}                                             \tag{D.CA.4}
\]
Low-dimensional repetitions are read through their
linear, unitary or symplectic identification. Odd-dimensional
orthogonal groups in characteristic two are read through
the symplectic row. The twisted triality groups ${}^3D_4(q)$
are exceptional groups, not the minus orthogonal row of
this table. They are not included by (D.CA.4).
\end{corollary}

\begin{proof}
We use the known automorphism theorem with its exact
normal-subgroup conclusion. For the groups in his
Section~2, Steinberg defines $\widehat S$ to be the group
of inner and diagonal automorphisms and
$\mathcal A_0=\langle\widehat S,F\rangle$, where $F$ is
the field-automorphism subgroup. His statements~3.2--3.6
assert
\[
 S\mathrel{\trianglelefteq}\widehat S
       \mathrel{\trianglelefteq}\mathcal A_0
       \mathrel{\trianglelefteq}\operatorname{Aut}(S),
                                                        \tag{D.CA.5}
\]
with graph automorphisms giving the final quotient.
The complete written proofs are in Sections~4--7
\citep[pp.~607--614, especially statements~3.2--3.6 on
p.~608 and their proofs]{Steinberg1960ClassicalAutomorphismsB}.
We apply that theorem rather than assuming that an
arbitrary collection of diagonal and field representatives
is normal.

For clarity we explain its natural projective realization.
In the split classical matrix models, the adjoint group
acts faithfully on the natural projective space. Type $A$
uses projectivities, type $C$ symplectic similitudes modulo
scalars, type $B$ the odd orthogonal model when $p\ne2$,
and type $D$ proper orthogonal similitudes modulo scalars.
These models and their natural root matrices are defined
over the prime field. The projective action remains
faithful in characteristic two in the even-dimensional
orthogonal and symplectic cases; for type $B$ in
characteristic two we use the corresponding type $C$
finite group.

One can check explicitly that the entire diagonal group
in Steinberg's definition is present in this projective
model. Write $h_i$ for its prescribed nonzero values on
the simple roots. In type $A$ choose a diagonal matrix
with consecutive ratios $h_i$. For type $C$, in a
hyperbolic basis choose
\[
 D=\operatorname{diag}(a_1,\ldots,a_\ell,
                         c/a_1,\ldots,c/a_\ell),
 \quad a_\ell=1,\quad c=h_\ell^{-1},\quad
 a_i/a_{i+1}=h_i.
                                                        \tag{D.CA.6}
\]
The last-root value is $a_\ell^2/c=h_\ell$. For type $D$
use the same display with
$c=h_{\ell-1}/h_\ell$ and the consecutive-ratio
conditions; the last simple root is
$e_{\ell-1}+e_\ell$, whose value is
$a_{\ell-1}a_\ell/c=h_\ell$.
For odd type $B$ use the diagonal matrix with an additional
middle entry $1$, take $a_\ell=h_\ell$, and replace
$c/a_i$ by $a_i^{-1}$. The last-root value is then
$a_\ell=h_\ell$. These matrices have their entries in
the stated field and act on each root group by exactly
Steinberg's formula~3.1. No extraction of a square root
in the coefficient field is needed.

The twisted classical groups are taken in the standard
split matrix model over $K=\mathbb F_{q^2}$, with the
commuting diagram/field condition used to define
${}^2A$ or ${}^2D$. The diagonal characters are precisely
the self-conjugate characters of Steinberg's Section~3.
The preceding projective torus construction is equivariant
for that diagram action. Its self-conjugate projectivities
therefore normalize the twisted group. Entrywise
$p$th-power automorphisms of $K$ also normalize it.
This is the reason for retaining $Q=q^2$ in the minus
orthogonal row: its field subgroup in these coordinates
has order $2f$, not $f$. We do not extend a semilinear
representation from $\mathbb F_q$ to $\mathbb F_{q^2}$
by assuming a section of their Galois groups.

In both split and twisted cases, the scalar quotient of
these natural matrices realizes $\widehat S$ faithfully.
Indeed an element acting trivially by conjugation on
$S$ centralizes its root matrices projectively. The
scalar factor on a root unipotent is $1$, since its
only eigenvalue is $1$. The natural root matrices have
scalar simultaneous commutant: in type $A$ this follows
at once from the elementary matrix units; in the form
models, the root matrices joining hyperbolic basis
vectors force the same diagonal scalar on all paired
basis vectors, and the short-root matrices give the
same scalar on the extra odd-dimensional vector.
The twisted versions give the same calculation after
using the paired root groups and extending scalars.
Equivalently this is the absolute irreducibility of the
natural classical module. Thus the projective conjugation
kernel is trivial. These standard matrix realizations
are compatible with the classical models and definitions
in \citet[Chapters~7--11]{Taylor1992ClassicalB}.

The entrywise field action normalizes this projective
matrix group. The resulting subgroup of
$\mathrm{P}\Gamma\mathrm{L}_d(\mathbb F_Q)$ is isomorphic
to $\mathcal A_0$. To check injectivity, a kernel element
would express a field automorphism as an inner-diagonal
one. The uniqueness of the field part in Steinberg~3.2
forces that field automorphism to be trivial; the
remaining projective element is trivial by the preceding
commutant calculation. We have therefore obtained an
actual embedding
\[
 \mathcal A_0\hookrightarrow
       \operatorname{PGL}_d(\mathbb F_Q)
                    \rtimes\operatorname{Gal}(\mathbb F_Q/\mathbb F_p),
 \qquad S\subset\operatorname{PGL}_d(\mathbb F_Q).
                                                        \tag{D.CA.7}
\]

Steinberg~3.6 gives the exact last column of (D.CA.4).
In split type $D_4$ the quotient is $S_3$, including
triality. In type $C_2$ in characteristic two it has
order two; the exceptional graph representative can
have square equal to a field automorphism, so no split
semidirect product with a graph involution is used.
For the twisted classical types the diagram operation
is already absorbed by the field subgroup over
$\mathbb F_{q^2}$, and the final graph quotient is
trivial. In every row its order is at most six.

The order comparisons already proved from the exact
classical order formulas give
\[
 \begin{array}{ll}
 |\operatorname{PSL}_d(q)|\geq q^{d^2/4},
       &|\operatorname{PSU}_d(q)|\geq(q^2)^{d^2/8},\\
 |S|\geq q^{d^2/8}
       &\text{for the symplectic and orthogonal rows}.
 \end{array}                                           \tag{D.CA.8}
\]
See the proofs of
\ref{D:surface-bounded-centre-simple-product},
\ref{D:bounded-centre-unitary-simple-product}, and
\ref{D:symplectic-orthogonal-simple-product}; their
numerical order inequalities do not require any bound
on the scalar centre. In particular, for every row,
\[
                         d^2\log Q\leq16\log|S|.
                                                        \tag{D.CA.9}
\]
The factor two for $Q=q^2$ is included in this inequality.

For an arbitrary $\mathcal B$ in the statement put
$\mathcal B_0=\mathcal B\cap\mathcal A_0$. By (D.CA.5),
this is normal in $\mathcal B$, contains $S$, and has
index at most six. Equations (D.CA.7) and (D.CA.9)
verify every embedding and size hypothesis of
Corollary~\ref{D:normal-subgroup-projective-crowns} with
$b_0=6$ and $\kappa=16$. That corollary gives (D.CA.2)
and (D.CA.3). It does not require the graph automorphisms
themselves to act on the natural projective space.
This is essential for triality and removes any hidden
ambient-extension hypothesis.
\end{proof}

\begin{remark}
If one includes in one's terminology a fixed finite
collection of simple derived exceptions omitted from
(D.CA.4), they can be added after enlarging $c_R$.
For each such $S$, the faithful permutation action of
$\operatorname{Aut}(S)$ on the set $S$ gives an embedding
in $\operatorname{GL}_{|S|}(\mathbb F_p)$. Its projective
kernel is trivial, since a scalar permutation matrix
is the identity. Apply
\ref{D:normal-subgroup-projective-crowns} with
$\mathcal B_0=\mathcal B$, coefficient field
$\mathbb F_p$, and the maximum of the finitely many
ratios $|S|^2\log p/\log|S|$. This finite-exception
observation makes no additional classification claim.
\end{remark}


% Separate post-CP22 candidate. The constants are for one fixed local surface.
\subsection{Quadratic projective counts with arbitrary SNC dual graph}

\begin{proposition}[Proved: uniform marked lattices and exact general-SNC presentations]
\label{D:SNC-general-weighted-presentation}
Retain the fixed resolution and coefficient gerbe of
\ref{D:SNC-general-Brauer-geometry}. Put
\[
 \gamma_i=g(R_i),\quad
 \eta_i=\sum_{x\in E_i\cap E_j}m_i/h_x,\quad
 b_i=-E_i^2,\quad M=\operatorname{lcm}_i m_i.
\]
Choose a fixed positive integral exceptional cycle
\(A=\sum_i a_iE_i\) with \(c_i=-A\cdot E_i>0\).
Every scalar-weight-one punctured vector bundle has a
formal lattice whose coarse spectral bundles satisfy
\[
 \gamma_i-1-\eta_i\leq\mu_{\min}(W_{i,0})
 \leq\mu_{\max}(W_{i,0})
 \leq\gamma_i-1+b_i+2\eta_i.                            \tag{D.GP.1}
\]
All endpoints are \(O(m_i)\). The same conclusion with
enlarged \(O(m_i)\) endpoints holds for the adjacent
one-step Frobenius and fixed ordinary exceptional twists
used in a unit diagram.

There are fixed integers \(B,N\), independent of the
coefficient order, residue orders and ranks, giving
the original-\(A\) meromorphic Hom pole bound and the
positive Hom-tail vanishing, including prescribed lifts
over finite local Artinian coefficient algebras, exactly
as in \ref{D:SNC-tree-uniform-Hom-depth}.
For a fixed ordinary precision \(a\), a cycle of \(f\)
stages and fixed compatible node types, the actual
bundle, arrow and paired inverse data on \(aA\) have a
polynomial presentation with an open subbundle/slope
condition. It covers the corresponding Artinian
families, and every all-ring point produces actual
objects. For \(r=m_i s_i\), its variables have the bounds
\[
 \begin{array}{c|c|c}
 \text{block}&\text{number}&\text{weight}\\ \hline
 \text{component }i,\ 1\leq j\leq aa_i m_i
       &O(fm_i s_i^2)&(M/m_i)j\\
 \text{component }i,\ \text{residual}
       &O(fm_i s_i^2)&M\\
 \text{each original node}&O(fr^2)&M .
 \end{array}                                           \tag{D.GP.2}
\]
All equations have weighted degree at most \(KM\),
where \(K\) depends only on the fixed resolution,
precision, characteristic and diagram lengths.
All determinant lines are covered; there is no moving
Picard parameter. Fixed formal reference coefficients
are constants, not variables in this estimate.
\end{proposition}

\begin{proof}
We give the changes needed in the full marked-lattice
and presentation proofs, so that neither a tree
condition nor commuting node inertia is used silently.
Let \(H=M^2\) and \(w_i=H/m_i\). At a marked point
above \(x\in E_i\cap E_j\), its root order is
\(d_x=m_j/h_x\), and
\[
           \frac{w_i}{d_x}
                   =\frac{Hh_x}{m_i m_j}\in\mathbb Z.   \tag{D.GP.3}
\]
Choose a rank-\(w_i\) spectral bundle \(V_i\) which is
a sum of normalized root-point lines, with every one
of the \(d_x\) weights occurring equally often at
each mark. The choices at different marks are
independent. Its coarse bundle is \(\mathcal O^{w_i}\),
and its parabolic slopes lie in \([0,\eta_i]\).
Under inverse spectral pushforward, its node
representation contains every scalar-weight-one
irreducible of dimension \(h_x\) with multiplicity
\(Hh_x/(m_i m_j)\). Indeed the \(m_i/h_x\) points
correspond to the orbits of first-branch characters,
and their \(m_j/h_x\) residual weights distinguish
all irreducibles in those orbits. The same uniform
list occurs on the other branch. The exact node
pasting of \ref{D:SNC-general-residual-category}
therefore glues them to a rank-\(H\) residual bundle.
Coherent \(H^2\) on a tame proper curve vanishes, so
this bundle lifts to a formal bundle \(A_0\).

The algebraic generator \(G_0\) of
\ref{D:SNC-general-Brauer-geometry} algebraizes these
formal bundles by ordinary coarse-Hom existence and
Morita. Its possible large rank is only an existence
tool. For any lattice \(F\),
\(q_*\mathcal Hom(A_0,F)\) is an ordinary formal
vector bundle. At a node its scalar centre cancels;
its Hom is a representation of the abelian root
quotient, and invariant monomials in the two
prescribed residue classes give a free coarse basis.
The same calculation proves exactness on the
elementary-transform sequences. Thus the finite
nonnegative integer
\[
 e(F)=\operatorname{length}
       \bigl(M_F/\Gamma(Y,q_*\mathcal Hom(A_0,F))\bigr)
              +\dim_kH^1(Y,q_*\mathcal Hom(A_0,F))       \tag{D.GP.4}
\]
is defined as in \ref{D:SNC-tree-minimal-marked-windows};
\(M_F\) is the common reflexive extension of the
punctured coarse-Hom bundle. Choose a minimizer.

For a locally free spectral quotient \(Q\) on
\(\mathcal R_i\), exact coarse Hom identifies the
elementary-transform quotient with
\((h_i)_*(\pi_i)_*\mathcal Hom(V_i,Q)\).
The finite map \(h_i:R_i\to E_i\) is now ramified,
but its Euler characteristic is still exactly that
on \(R_i\), with no extra \(m_i\) or \(h_x\) factor.
Ordinary curve Riemann--Roch and normalized root
weights give
\[
 \mu_{\rm par}(Q)+1-\gamma_i-2\eta_i
 \leq
 \frac{\chi(\mathcal R_i,V_i^\vee\otimes Q)}
             {w_i\operatorname{rk}Q}
 \leq\mu_{\rm par}(Q)+1-\gamma_i.                      \tag{D.GP.5}
\]
Minimality makes this Euler characteristic nonnegative.
The inverse transform for a subbundle has the opposite
inequality after tensoring with the root normal line.
This line transports the spectrum by its cyclic deck
translation and has parabolic degree \(E_i^2=-b_i\).
Indeed its \(m_i\)-fold product is
\(h_i^*\mathcal O_{E_i}(E_i)\), of degree \(m_iE_i^2\);
deck translation preserves parabolic degree also for
a ramified cyclic cover. The normalized universal
line cancels from this calculation. Giving ordinary
HN subbundles and quotients their induced root flags,
whose degree corrections lie in \([0,\eta_i]\), proves
(D.GP.1). This is the full Euler argument of the
tree proof, with the changed genus and point orders.
The root/flag equivalence and its tensor compatibility
are \citet[Proposition3 and Theorem4]{Borne2007D}
and \citet[Theorem6.1]{BorneVistoli2012D}.

The degree estimates needed for normal layers likewise
use the degree of \(h_i\), not its being étale.
The restriction of \(-A\) to \(R_i\) has degree
\(m_i c_i\). At a node its adjacent root point has
degree \(h_x/m_j\); there are \(m_i/h_x\) such points.
Thus the sum of these fractional corrections is
\(m_i/m_j\), exactly as required by pullback of the
original intersection point. A full original-\(A\)
layer has degree \(m_i c_i\), while all incomplete
root-layer and marked corrections have size \(O(m_i)\).
The sandwich between coarse parabolic Hom and
ordinary Hom twisted by the marked set is unchanged.
Duality reduces the relevant \(H^1\) vanishing to
the slope inequality
\[
 \mu_{\min}(W'_{i,0})>
 \mu_{\max}(W_{i,0})+2\gamma_i-2-\ell m_i c_i
                       +O(m_i).                       \tag{D.GP.6}
\]
By (D.GG.8), (D.GP.1) and fixed positivity of \(c_i\),
one fixed \(\ell\) suffices. The negative-layer
\(H^0\) calculation is the dual pole bound. This
proves fixed original-\(A\) depth and poles.

Normalize universal spectral lines to rational
parabolic degree in \([0,1)\). A one-step Frobenius
comparison changes that degree by at most \(p+1\).
At a marked root point it multiplies characters by
\(p\), with carries at most \(p-1\); the total
ordinary correction is \(O(\eta_i)\).
The already proved rank-independent Frobenius HN
estimate, using the full Sun/Mehta--Pauly inputs in
\ref{D:twisted-period-independent-normal-depth}, enlarges the
ordinary slopes by at most
\((p-1)\max(\gamma_i-1,0)\) besides the factor \(p\).
For genus zero the splitting calculation applies.
Thus all adjacent sources still have \(O(m_i)\)
windows and satisfy the same fixed-depth bound.
No tensor-semistability assertion is being used.

These vanishings are valid on arbitrary coefficient
bases on the indicated open slope locus. The fixed
proper-curve cohomology complexes are finite free
locally on the parameter base. Fibrewise \(H^0\)
vanishing or \(H^1\) vanishing gives respectively
a locally split injective or surjective differential,
and hence the same vanishing after every base change.
This is the cohomology-and-base-change argument of
\ref{D:SNC-tree-finite-unit-comparison}, using
Tag07VK. In particular nilpotent poles are excluded;
one is not testing meromorphic maps only on reduced
fibres. The layer exact sequences give the specified
Artinian lifts.

We next verify the exact object presentation. After
a fixed positive ordinary twist, the coarse spectral
bundle has \(s_i+1\) generating sections. At each
marked point lift a basis adapted to its target
flag, and assign the kernel vector weight zero.
With \(d_x=m_j/h_x\), put
\[
 T_i=\mathcal O\!\left(\sum_x(d_x-1)\mathcal P_x\right),
 \quad
 Q_x=\bigoplus_{\ell=1}^{s_i+1}
     T_i/T_i(-(d_x-1-a_{x\ell})\mathcal P_x).
                                                        \tag{D.GP.7}
\]
A frame \(g_x\) defines \(T_i^{s_i+1}\to Q_x\);
its kernel \(P(g)\) is a homogeneous free bundle
with coarse bundle \(\mathcal O^{s_i+1}\) over
every coefficient algebra. The exact local
calculation and the zero-weight kernel are those
of \ref{D:SNC-tree-small-marked-source}; they depend
only on a root point of order \(d_x\), not on
unramifiedness of the coarse cover.
There are \(O(m_i)\) marks. Their transverse root
divisors are disjoint on the formal spectral
component. That formal component retains adjacent
roots, as specified in the preceding proposition,
so the finite étale spectral lift and all these
Cartier equations actually exist.

The full fixed partial-resolution construction of
\ref{D:SNC-tree-fixed-normal-presentation} now applies
with these exact inputs. To state its effectivity
check explicitly, take \(L=\det W_{i,0}\) and a
fixed positive ordinary line \(B_i\) of degree
a fixed multiple of \(m_i\). For \(b=cs_i\) the
alternating kernel lines are
\(L^{-1}B_i^{-2vb}\) and \(LB_i^{-2vb}\).
The generating Hom lines are \(L^{-1}B_i^b\) and
\(L\); their \(H^1\) vanishes, including the marked
corrections. Hence the fixed rank-two terms through
length seven cover every determinant and every
specified Artinian lift. Impose actual \(d^2=0\)
and the open rank conditions on the factored
map \(P_1\to P(g)\), not on \(P_1\to T_i^{s_i+1}\).
Local unit-minor cancellation gives an actual
locally free cokernel and locally free end term
over every ring.

For maps, use the fixed termwise two-affine
Cech--Cartier complexes. Contract the Hom between
distinct point supports before retaining variables.
Each point-supported invariant Hom/Ext block has
one prescribed residue class in an exponent
interval of length at most \(d_x\), or at most
\(pd_x\) after one Frobenius. Thus its size is
bounded independently of \(d_x\). The top-top,
top-rank-two and rank-two/rank-two blocks have
multiplicity/degree respectively
\[
 O(s_i^2),O(m_i);\qquad
 O(s_i),O(m_i s_i);\qquad
 O(1),O(m_i s_i).
 \]
Together with genus \(O(m_i)\) and \(O(m_i)\)
marks, their total cohomology is \(O(m_i s_i^2)\)
per normal layer. The strong contractions are
over \(k\) and remain exact after every coefficient
extension. The finite homological-perturbation
identities used in the cited presentation preserve
column degree. Positive normal terms raise order;
the variable internal differential crosses only
a fixed number of columns. Their words therefore
have degree at most normal order plus a fixed
constant. The length-seven/length-three comparison
retains exact actual degree-zero Hom; identities
are tested on the generating \(P(g)\), rather
than by a spurious identity between equal-length
partial complexes. These are the explicit
all-ring finite-Hom constructions already proved
there, using \citet[Sections2.1--2.4 and3.2]
{C_Crainic2004Perturbation}.

At a node use the actual noncommuting type
(D.GG.12) and its multiplicity frames. The
Cartier-cycle ring identity
\[
 B[s,t]/(s^{aa_i m_i}t^{aa_j m_j})
 =
 B[s,t]/(s^{aa_i m_i})
 \mathbin{\times}_{B[s,t]/(s^{aa_i m_i},t^{aa_j m_j})}
 B[s,t]/(t^{aa_j m_j})                                \tag{D.GP.8}
\]
and Milnor kernel pasting give actual bundles from
the component bundles and overlap isomorphisms.
Restriction of the presentation complexes is
derived restriction of their full Cartier
resolutions, not underived restriction of \(Q_x\).
Their amplitude is still fixed. Scalar-zero Hom
on the rectangle has the exact slot bound of
(D.GG.10); each original edge therefore costs
only \(O(fr^2)\) coefficients. All equivariance,
inverse and composition equations are actual
matrix equations on these finite modules.
There is no need to commute the two actions
on the original bundle. Multiple edges and
graph cycles merely add their distinct overlap
data; no triple intersections occur in an SNC
surface divisor.

Finally assign the weights in (D.GP.2).
Within component \(i\), normal multiplication
adds order, so rescaling by \(M/m_i\) makes every
retained equation have degree \(O(M)\).
The fixed internal column length and one-step
Frobenius change the constant only. At a node
the two normal exponents have weights \(M/m_i\)
and \(M/m_j\); their retained maxima give
\(O(M)\) as well. The fixed-amplitude derived
restrictions and overlap equations obey the
same bound. Open local-freeness and HN conditions
do not add polynomial degrees. This proves
the full table, including actual all-ring
effectivity and Artinian coverage. The number
of components and original edges is fixed.
\end{proof}

\begin{theorem}[Proved: quadratic projective counts for arbitrary SNC local surfaces]
\label{D:SNC-general-projective-quadratic-count}
Let \(R,k,U,Y,E\) satisfy
\ref{D:SNC-general-Brauer-geometry}, with no tree
assumption on the SNC exceptional divisor.
There is a constant \(C_R\) such that, for every
\(q=p^f\) and \(r\geq1\),
\[
 \#\operatorname{Hom}_{\rm cont}
       \bigl(\pi_1(U),\operatorname{PGL}_r(\mathbb F_q)\bigr)
       /\operatorname{PGL}_r(\mathbb F_q)
                         \ \leq q^{C_R r^2}.           \tag{D.GP.9}
\]
All geometric Brauer periods and all secondary
node residues are allowed. The constant is for
this fixed local surface. No uniform hereditary
index bound or full finite generation is asserted.
\end{theorem}

\begin{proof}
For \(r=1\) the assertion is immediate. For \(q=2\)
the scalar kernel is trivial and the already proved
local linear representation count applies.
Assume \(r\geq2\) and \(n=q-1>1\).
The obstruction of a projective representation
lies in \(H^2(\pi_1(U),\mathbb F_q^*)\), which
injects into \(H^2(U,\mu_n)\) by the already proved
degree-two fundamental-group comparison. In
particular (D.GG.2) bounds the number of possible
full coefficient classes by \(q^{c_R}\).
Only realized classes are used. Their Brauer
images are Azumaya, because a projective finite
frame torsor gives its matrix-algebra bundle.
Thus the mixed-root coefficient gerbe of
\ref{D:SNC-general-Brauer-geometry} is available.
It restricts to the actual lifting gerbe,
including its Kummer class.

Fix one such class and one extension gerbe.
The scalar-weight-one finite local system on
its puncture has a vector bundle lattice:
extend it coherently using the algebraic
generator \(G_0\), and take the reflexive hull
on the regular two-dimensional root charts.
It is locally free there, and all actions
extend uniquely. This step is only coverage
of geometric objects; no relative hull/base
change assertion is made. The spectral rank
formula gives \(m_i\mid r\) for every \(i\),
so \(M=\operatorname{lcm}m_i\mid r\).

Choose the minimizing lattices of (D.GP.1).
Use the \(f\) one-step Frobenius phases, whose
scalar weights are \(1,p,\ldots,p^{f-1}\).
Multiplication by \(p\) preserves every residue
order and every commutator order. It induces
the actual relabeling of the spectral covers
and of the node representations; no node
commutator is erased. The proposition gives
uniform ordinary-\(A\) poles and precision for
all phases and their paired inverse maps.

We spell out the finite-to-formal comparison.
A bundle on the retained Cartier cycle lifts
through each further square-zero layer because
the obstruction lies in coherent \(H^2\) on
the tame one-dimensional proper support, hence
is zero. This remains true over an affine
coefficient algebra: a fixed two-affine cover
of the coarse curve computes the relevant
cohomology, and finite tame invariants are
exact. The proved positive Hom-tail vanishing
lifts prescribed arrows. If the two inverse
products equal the required Cartier power
to retained precision, their errors lie in
a strictly positive normal ideal; the
geometric series restores actual inverse
maps. Write \(I=\mathcal O(-A)\) on the root
stack. Absolute Frobenius satisfies
\(F^*I=I^p\). Pole-bounded conjugation therefore
improves an error of original-\(A\) order \(t\)
to at least \(pt-2B\), and one full cycle to
\[
           p^f t-2B\frac{p^f-1}{p-1}.                 \tag{D.GP.10}
\]
Taking \(t>2B/(p-1)\) gives the same unique
paired-jet contraction as in
\ref{D:SNC-tree-finite-unit-comparison}.
It preserves every prescribed retained jet.
All these calculations are made in the
\(\mathfrak m_R\)-adic completion of
\(R\otimes_k B\) for a coefficient algebra \(B\);
for finite local Artinian \(B\) this is the
usual complete tensor base. They do not
invoke an existence theorem over an
uncompleted arbitrary parameter tensor ring.

For this algebraization the coefficient algebra
\(B\) is finite type over \(k\), or finite local
Artinian; the universal polynomial parameter
family is of finite type. Its completed tensor
base is Noetherian. Over arbitrary coefficient
algebras the assertions made above are the
finite-cycle effectivity and base-change
identities, not an application of Noetherian
Grothendieck existence.
Algebraization uses the algebraic \(G_0\)
constructed from the mixed-root Azumaya
extension. Ordinary coarse-Hom Grothendieck
existence, with its finite-algebra action
and relations, followed by evaluation,
algebraizes the compatible formal bundles
and maps on the actual coefficient gerbe.
Here Hom from \(G_0\) is a right module over
its coarse End algebra, and \(G_0\) is the
left module used in evaluation. This avoids
both an assumed global étale Brauer splitter
and circular algebraization of the
Euler-minimizing formal anchor. The rank of
\(G_0\) has entered no coordinate estimate.

The actual specified unit-Isom functor is
of finite presentation. Proper finite-cycle
Hom schemes, their inverse/equality equations,
and the unique contraction just proved
represent it, exactly as in
\ref{D:finite-jet-unit-Isom}. The uniform
pole bound holds after every coefficient
base change by the \(H^0\) argument in the
proposition. Finite étale fixed-frame torsors
and their specified isomorphisms lift uniquely
across nilpotent coefficient extensions.
Hence the fixed-object Isom morphism is
formally étale and of finite presentation,
so étale; its image is open. The class
fibres are consequently open and closed,
and their number is bounded by the number
of irreducible components of the parameter
open. These are the exact finite-Isom
criteria already checked in Tags08K6,
0BQB,02HT and02HM in the cited unit lemma.
This is an all-Artinian comparison, not a
tangent-dimension argument.

Apply the weighted component bound to
(D.GP.2). Pad the component blocks to
\(b_i=\lceil K_1 f m_i s_i^2\rceil\)
variables at every positive layer and at
the residual layer, and put \(J_i=aa_i m_i\).
The \(O(fr^2)\) node variables have weight
\(M\). For a constant \(K_2\), the relevant
weighted product, apart from its harmless
linear factor in the total variable count,
is bounded by
\[
 K_2^{\,O(fr^2)+\sum_i b_i}
       \prod_i
       \left(\frac{(K_2 m_i)^{J_i}}{J_i!}\right)^{b_i}
                     \leq \exp(K_3fr^2).              \tag{D.GP.11}
\]
Indeed the rescaling \(M/m_i\) cancels
componentwise, \(J_i/m_i=aa_i\) is fixed,
and \(b_iJ_i=O(fm_i^2s_i^2)=O(fr^2)\).
There are only the fixed number of
components and nodes. No
\(\log(M/m_i)\) loss is introduced.

The discrete types cost the same scale.
At a node the number of scalar irreducible
labels is \(m_i m_j/h_x^2\leq r^2\), and
their common dimension is \(h_x\).
Choosing multiplicities with total rank
\(r\) is bounded by the number of
multisets of at most \(r\) labels from
at most \(r^2\) choices, namely
\(\exp(O(r\log(r+1)))\). On a branch
there are \(m_i/h_x\) marked points,
each with \(s_i\) weights of order
at most \(m_j/h_x\leq r\); their total
weight choices obey the same bound.
Imposing compatibility only decreases
it. Over all \(f\) phases this is
\(\exp(O(fr\log(r+1)))\leq\exp(O(fr^2))\).
Graph cycles introduce no new discrete
choices beyond these fixed edges and
their actual overlap matrices.

Thus, for each realized coefficient class,
the number of punctured scalar-weight-one
unit classes is at most \(\exp(K_4fr^2)\).
Their finite \(\mathbb F_q\)-frames project
onto the projective representation classes
with that obstruction: identifying the
fixed coefficient gerbe changes a chosen
lift only by its scalar coboundary. Distinct
twisted objects may project to the same
projective object, which improves the
upper bound. Multiplying by the
\(q^{c_R}\) coefficient-class bound and
using \(q=p^f\) proves (D.GP.9).

The construction is for the fixed resolution.
It proves no linear dependence of these
constants on the degree of a varying
punctured cover, and therefore gives no
unrestricted nonabelian crown or full
finite-generation conclusion.
\end{proof}


% Closure of the same arbitrary-SNC local projective reduction.
\begin{corollary}[Proved: projective growth on complete local punctures in all dimensions]
\label{D:complete-local-projective-growth-all-dimensions}
Let \((R,\mathfrak m,k)\) be a complete excellent local
\(k\)-algebra of characteristic \(p>0\), with algebraically
closed coefficient and residue field \(k\).
Assume every irreducible component of \(\operatorname{Spec}R\)
has dimension at least two. Put
\(U=\operatorname{Spec}R\setminus\{\mathfrak m\}\).
For every connected component \(V\) of \(U\), one has
\[
 \#\operatorname{Hom}_{\rm cont}
       \bigl(\pi_1(V),\operatorname{PGL}_r(\mathbb F_{p^f})\bigr)
          /\operatorname{PGL}_r(\mathbb F_{p^f})
                    \leq p^{C_R f r^2}                \tag{D.AP.1}
\]
for one constant \(C_R\), every \(f\geq1\), and every
\(r\geq1\). Normality and reducedness of \(R\) are not
required. No Brauer-vanishing or bounded-period condition
is imposed.

Fix \(\epsilon>0\). For each isomorphism type \(S\) in
a collection of nonabelian finite simple groups, choose
one embedding
\[
 S\hookrightarrow\operatorname{PGL}_{r_S}(\mathbb F_{q_S}),
 \qquad q_S=p^{f_S},\qquad |S|\geq q_S^{\epsilon r_S^2}.
\]
The quotient of \(\pi_1(V)\) formed by all distinct
simple quotient kernels of those types is topologically
finitely generated. Its multiplicity for type \(S\) is
at most \(|S|^{C_R/\epsilon}\). This is a simple-product
conclusion, not full finite generation of \(\pi_1(V)\).
\end{corollary}

\begin{proof}
First let \(R\) be a normal domain of dimension two.
It has a projective regular resolution, an isomorphism
off its closed point. Indeed the normal complete ring
satisfies condition (4) of
\citet[Theorem54.14.5, Tag0BGP]{StacksDSurfaceResolution},
whose proof gives normalized blowups at the singular
points. Excellence makes their normalizations finite,
so these maps are projective. Apply
\citet[Lemma54.15.6, Tag0BIC]{StacksDSurfaceResolution}
to the exceptional closed fibre. Its curve components
have finite normalization by excellence; the lemma
gives a further sequence of point blowups with SNC
support. These modifications remain over the closed
point. If the exceptional divisor is empty, blow up
the closed point once. The resulting nonempty
exceptional components are smooth proper curves
over the perfect field \(k\). Thus
\ref{D:SNC-general-projective-quadratic-count} applies.

Now let \(R\) be a normal domain of dimension \(d\geq2\).
Use the fixed slice of
\ref{D:local-surface-reduction}: choose \(d-2\) members
of a parameter system, with ideal \(I\), and put
\[
 C=(R/I)_{\mathrm{red}},\qquad
 W=\operatorname{Spec}C\setminus\{\mathfrak m_C\}.
 \]
Every component of \(C\) has dimension two.
For every connected finite étale cover of \(U\),
normalization gives a finite complete local normal
domain \(B\) over \(R\). The ideal \(IB\) has at most
\(d-2\) generators. The exact connectedness and
cohomological-dimension lemmas
\citep[Tags0ECR and0DX9]{C_StacksLocalSlice}
therefore make its punctured slice connected.
Consequently \(\pi_1(W)\twoheadrightarrow\pi_1(U)\).
This single fixed slice detects all original covers.

Normalize \(C\) and write its puncture as
\(W'=\coprod_i U_i\), where each \(U_i\) is the
puncture of a complete normal local surface
\(k\)-algebra. Their residue fields are \(k\).
Put \(a=\#\pi_0(W'\times_W W')\).
For \(G=\operatorname{PGL}_r(\mathbb F_q)\), the
surface bounds give at most \(q^{(\sum_i C_i)r^2}\)
tuples of \(G\)-torsor isomorphism classes on \(W'\).
Each such tuple has at most \(|G|^a\) descents to
\(W\). To see the exact fibre bound, fix one tuple
upstairs. An equivariant isomorphism between its
two pullbacks to a connected component of the
double product is determined by its value at one
geometric point, with at most \(|G|\) possible
values. Cocycle conditions and quotienting by
automorphisms only decrease the number.
Finite-proper effectiveness is
\citet[Expose IX, Theorem4.7 and its proof]{C_SGA1}.
Since \(|G|\leq q^{r^2}\), this gives
\(q^{(a+\sum_iC_i)r^2}\) classes on \(W\).
Surjectivity to \(\pi_1(U)\) makes restriction
injective on ambient projective conjugacy
classes. This proves (D.AP.1) for every normal
complete local domain of dimension at least two.

For the stated nonnormal ring first remove its
nilpotents, which does not change finite étale
torsors. Its finite normalization is a finite
disjoint union of complete normal local domains
\(A_j\), all of dimension at least two by the
hypothesis. Their residue fields remain \(k\).
Their punctures give a fixed finite surjection
\(U'\to U\). The normal-domain bounds just proved,
and the identical finite-proper torsor descent
argument on \(U'\times_U U'\), give
\[
                     |H^1(U,G)|\leq q^{C_Rr^2}.
 \]
There are finitely many connected components of
\(U\). A torsor on any one of them extends by
the trivial torsor on the others; thus its
isomorphism classes inject into \(H^1(U,G)\).
This proves (D.AP.1) exactly as stated.
No group-\(H^2\) or Brauer injectivity is used
in either descent step.

Distinct simple quotient kernels give distinct
ambient projective conjugacy classes under the
chosen embedding. Hence their multiplicity is
at most \(q_S^{C_Rr_S^2}\leq |S|^{C_R/\epsilon}\).
The exact generating-tuple and Goursat/compactness
argument in the proof of
\ref{D:Brauer-zero-selected-simple-product}
now applies without its former Brauer restriction.
For example \(d=\lceil C_R/\epsilon\rceil+3\)
suffices, by the same
\citet[Proposition2.6(ii), equations(2.1)--(2.2),
pp12--13]{LucchiniMarionTracey2020D}.
It gives the stated simple-product quotient.

The lower bound two on every component dimension
is retained deliberately: the proof does not
supply the corresponding representation count
for a one-dimensional punctured component.
There is no mixed-characteristic assertion,
no uniform index-linear hereditary constant,
and no unrestricted crown conclusion.
\end{proof}


% Separate post-CP22 proved corollary. All four complete mathematical audits passed; named exceptional recognition within the same crown reduction.
\subsection{Exceptional defining-characteristic automorphism crowns}

\begin{corollary}[Proved: all exceptional defining-characteristic families]
\label{D:exceptional-defining-characteristic-crowns}
Retain the complete one-smooth-exceptional-curve setting and
hereditary constant $A_R$ of
Theorem~\ref{D:tame-projective-semilinear-count}.
The crown conclusion of
Corollary~\ref{D:classical-defining-characteristic-crowns}
also holds for every nonabelian simple exceptional group in
(D.EC.1) below, with arbitrary diagonal, field, graph and
permutation actions and the same requirement that the outer
quotient be radially tame. Here $q$ has characteristic $p$,
$D$ is a faithful projective dimension, and $Q$ is the actual
coefficient-field order:
\[
\begin{array}{c|c|c|c}
 S\text{ and simple range}&D&Q&
       [\operatorname{Aut}(S):\mathcal A_0]\\ \hline
 G_2(q),\ q\geq3&14&q&2\ (p=3),\ 1\ (p\ne3)\\
 F_4(q),\ q\geq2&52&q&2\ (p=2),\ 1\ (p\ne2)\\
 E_6(q),\ q\geq2&78&q&2\\
 E_7(q),\ q\geq2&133&q&1\\
 E_8(q),\ q\geq2&248&q&1\\
 {}^2E_6(q),\ q\geq2&78&q^2&1\\
 {}^3D_4(q),\ q\geq2&28&q^3&1\\
 {}^2B_2(q),\ q=2^{2a+1},\ a\geq1&10&q&1\\
 {}^2G_2(q),\ q=3^{2a+1},\ a\geq1&14&q&1\\
 {}^2F_4(q),\ q=2^{2a+1},\ a\geq1&52&q&1
\end{array}                                                    \tag{D.EC.1}
\]
The groups in this table are the simple adjoint groups;
thus any centre of the corresponding universal group has
already been divided out. One common size constant is
\[
 \kappa_E=3\cdot248^2=184512,
 \qquad
 c_R=\max\{1,\lceil\kappa_E(6A_R+1)\rceil+5\}.            \tag{D.EC.2}
\]
Explicitly, for $S\leq\mathcal B\leq\operatorname{Aut}(S)$,
a transitive subgroup
\[
 P\leq(\mathcal B/S)^\Omega\rtimes\operatorname{Sym}(\Omega),
 \qquad t=|\Omega|,
\]
its full preimage
$L\leq\mathcal B^\Omega\rtimes\operatorname{Sym}(\Omega)$,
assumed monolithic primitive with socle $S^\Omega$, and a
surjection $\Gamma\twoheadrightarrow L_k$ whose quotient
$\Gamma\twoheadrightarrow P$ is radially tame, one has
\[
 k\leq|S^\Omega|^{c_R},\qquad
 d(L_k)\leq\max\{2,2g(C)+1,c_R+2\}.                       \tag{D.EC.3}
\]
The three additional simple derived exceptions
$G_2(2)'$, ${}^2G_2(3)'$ and ${}^2F_4(2)'$ can be included
after enlarging $c_R$ by a fixed constant. The excluded
Suzuki group ${}^2B_2(2)$ is solvable. Together with the
classical corollary and its finite-exception convention,
this recognizes the defining-characteristic finite simple
Lie families in this radially tame outer-quotient crown
class. It makes no assertion about cross-characteristic
families or wild outer quotients.
\end{corollary}

\begin{proof}
We first treat the first seven rows. Steinberg's 1960
construction realizes each group as a group of actual
linear automorphisms of the corresponding Chevalley Lie
algebra $\mathfrak g_K$. His field $K$ is $\mathbb F_q$
for the untwisted groups, $\mathbb F_{q^2}$ for ${}^2E_6(q)$,
and $\mathbb F_{q^3}$ for ${}^3D_4(q)$. The self-conjugate
root-lattice characters realize the diagonal automorphisms
on this same Lie algebra, and field automorphisms act on
its coefficients. Statements~3.2--3.6 give the normal
sequence
\[
 S\mathrel{\trianglelefteq}\widehat S
   \mathrel{\trianglelefteq}\mathcal A_0
   \mathrel{\trianglelefteq}\operatorname{Aut}(S),
 \qquad \mathcal A_0=\widehat S\rtimes
                  \operatorname{Gal}(K/\mathbb F_p),           \tag{D.EC.4}
\]
where $\widehat S$ consists of the inner-diagonal
automorphisms. In particular the field part is unique.
The final quotient has the indices shown in (D.EC.1)
\citep[Section~2, statements~3.2--3.6, and their complete
written proofs in Sections~4--7]{Steinberg1960ClassicalAutomorphismsB}.

Here is the precise faithfulness check, which does not
require irreducibility of the adjoint Lie module in bad
characteristic. Let $\widetilde S$ be the subgroup of
$\operatorname{Aut}_{K,\mathrm{Lie}}(\mathfrak g_K)$
generated by $S$ and these diagonal matrices. Conjugation
gives a surjection $\widetilde S\to\widehat S$.
Write an element in its kernel as $s h$, with $s\in S$
and $h$ diagonal. The diagonal matrix normalizes both
the positive and negative root groups $U,U^-$.
Since $s h$ centralizes $S$, the element $s$ normalizes
both $U$ and $U^-$. The Bruhat normalizer calculation
therefore gives $s\in H$, the intersection of their
normalizers. This is exactly the calculation in
Steinberg~4.2--4.4, with its twisted analogue following
4.12. Thus $s h$ is itself a diagonal root-character
matrix. Centralizing every root element forces all
its root multipliers to be $1$, by the injective root
parameterization and formula~3.1. In a twisted orbit,
the uniqueness of the ordered product of ordinary root
factors forces the multiplier of each member to be $1$;
one varies the common field parameter, which occurs in
those factors with its prescribed conjugates. The
root-character matrix is then the identity on every
root vector and on the Cartan generators. Hence the
surjection has trivial kernel.

There is also no scalar kernel on projectivizing.
A scalar $c\operatorname{Id}$ preserving the Lie bracket
would satisfy
$c[x,y]=c^2[x,y]$ for every $x,y$. The bracket is nonzero
in every characteristic in the table: each of these
Cartan matrices has an off-diagonal entry $-1$, so an
appropriate relation $[H_\alpha,X_\beta]=-X_\beta$
survives reduction in every characteristic. Thus $c=1$.
Combining this with the uniqueness of the field part in
Steinberg~3.2 gives an actual embedding
\[
 \mathcal A_0\hookrightarrow
       \mathrm{P}\Gamma\mathrm{L}_{D}(K),
 \qquad S\subset\operatorname{PGL}_{D}(K).                   \tag{D.EC.5}
\]
No statement about semisimplicity of $\mathfrak g_K$ or
its adjoint representation has entered this argument.
The exceptional long/short-root automorphisms of $G_2$
in characteristic three and $F_4$ in characteristic two
lie in the bounded graph quotient; their squares need
not be treated as the identity.

The displayed dimensions are the dimensions of the
underlying Chevalley Lie algebras, namely rank plus the
number of roots. In particular $D\leq248$. One may see
the bound directly from the standard root descriptions:
$E_8$ has $112$ roots $\pm e_i\pm e_j$ and $128$
half-integral roots with an even number of minus signs;
$E_6,E_7$ are its root subsystems. The $F_4$ roots consist
of $24$ roots $\pm e_i\pm e_j$, $8$ roots $\pm e_i$,
and $16$ half-integral roots; $G_2$ has $12$ roots.
The $D_4$ and $C_2$ systems have $24$ and $8$ roots,
respectively. This also verifies the last three Lie
module dimensions without using their smaller minimal
representations.

We now treat the Suzuki and Ree rows, for which the
1960 article alone is not the automorphism input.
Use Steinberg's original Yale notes, Section~11.
There $G$ is the ambient Chevalley group of type $C_2$
(in characteristic two), $G_2$ (in characteristic three),
or $F_4$ (in characteristic two). The twisted group is
$S=\langle U_\sigma,U^-_\sigma\rangle$, where the
long/short-root operation is combined with a field
operation. On a finite field $K$ its defining condition
is $p\theta^2=1$ in Steinberg's convention. Thus
$|K|=p^{2a+1}$ and the groups are the three rows in
(D.EC.1), with $K=\mathbb F_q$, not a field of order
$\sqrt q$. Lemma~63 gives their injective root-coordinate
formulas, and Lemma~64 identifies the fixed torus with
the torus in the generated twisted group. Theorem~36
states that every automorphism of $S$ is inner,
diagonal, and field, with no additional graph factor
\citep[Section~11, pp.~177--186 and Theorem~36 with its
complete written proof sketch on pp.~195--199]
{Steinberg1968YaleExceptionalB}.

In these three cases the diagonal automorphisms are
already inner. For $G_2$ and $F_4$, the root and weight
lattices agree. For $C_2$, their quotient has order two,
and $K^*$ has odd order. Consequently every root-lattice
character has a unique extension to the weight lattice;
the extension commutes with the defining twist by
uniqueness. A self-conjugate diagonal element therefore
belongs to the fixed torus. Lemma~64 puts this torus
inside $S$. Thus Theorem~36 gives
$\operatorname{Aut}(S)=S\,\operatorname{Gal}(K/\mathbb F_p)$.
For completeness the product is semidirect. If a field
operation were inner, then, since it preserves $U_\sigma$
and $U^-_\sigma$, its inner representative would lie in
their common torus by the twisted Bruhat calculation.
On a simple relative root group the first coordinate
$t\in K$ would then be transformed both as $t\mapsto
c t$ and as $t\mapsto t^{p^j}$. The coordinate formulas
in Lemma~63 and the torus formulas at the start of the
proof of Theorem~36 give exactly these transformations.
Taking $t=1$ gives $c=1$, and varying $t$ makes the
field operation the identity. For ${}^2F_4$ one uses
one of its simple relative root groups; its first
coordinate again ranges over all $K$.

The ambient adjoint Lie realization restricts faithfully
to these simple groups. In type $C_2$, the universal
and adjoint groups have the same finite points over
$K$, since their possible central kernel has order two
and $|K^*|$ is odd. The other two ambient types have
trivial centre. Entrywise field operations normalize
the twisted groups. The nonzero-bracket scalar argument
above and the just-proved field-part injectivity yield
\[
 \operatorname{Aut}(S)\hookrightarrow
              \mathrm{P}\Gamma\mathrm{L}_{D}(\mathbb F_q),
 \qquad S\subset\operatorname{PGL}_{D}(\mathbb F_q),
                                                        \tag{D.EC.6}
\]
with $D=10,14,52$, respectively. These three rows therefore
have graph index one in the sense needed here.

The simple ranges and omitted groups are exactly the
known simplicity statements in the same source:
Theorem~5 and its proof, pp.~47--53, for untwisted
adjoint groups, and Theorem~34 with its written proof
sketch and following remark, pp.~187--188, for twisted
groups. For the first seven rows only $G_2(2)$ is
omitted. For the last three rows the omitted parameters
are $2,3,2$, respectively. The source's notation
${}^2A_2(4)$ uses the ambient field order and therefore
refers to the familiar nonsimple unitary group at
$q=2$; it creates no exception in the exceptional
families of (D.EC.1). We use these stated literature
theorems as known inputs. The calculations omitted
from their proof sketches, and the Chevalley/Bruhat
facts cited in them, are not being claimed as new
independent reproofs here.

Every row has $|S|\geq q$. In an untwisted group an
ordinary root group has $q$ elements. In ${}^2E_6$ and
${}^3D_4$ a fixed simple root gives a group with
$\mathbb F_q$ parameter. In the Suzuki and Ree cases,
Lemma~63 supplies at least $q$ distinct unipotent
elements by varying the first coordinate and fixing
the others. These unipotent elements survive the
adjoint quotient; alternatively the construction above
already takes place in the faithful adjoint group.
As $D\leq248$ and $Q\leq q^3$, we obtain
\[
             D^2\log Q\leq3\cdot248^2\log q
                         \leq\kappa_E\log|S|.                \tag{D.EC.7}
\]
For an arbitrary intermediate group $\mathcal B$ take
$\mathcal B_0=\mathcal B\cap\mathcal A_0$ in the first
seven rows, and $\mathcal B_0=\mathcal B$ in the last
three. It is normal in $\mathcal B$, contains $S$,
and has index at most six. Equations (D.EC.5)--(D.EC.7)
verify all hypotheses of
Corollary~\ref{D:normal-subgroup-projective-crowns} with
$b_0=6$ and $\kappa=\kappa_E$. This proves (D.EC.2)--(D.EC.3).

Finally the three simple derived exceptions are fixed
finite groups. The faithful permutation action of
$\operatorname{Aut}(S)$ on the finite set $S$, over the
fixed ground field $\mathbb F_p$, has trivial projective
kernel. The finite-exception argument after
\ref{D:classical-defining-characteristic-crowns}
therefore includes all of them by enlarging $c_R$.
It requires neither a low-dimensional exceptional
representation nor a new automorphism classification
for those finitely many groups.
\end{proof}


% Proved post-CP22 example. This is an obstruction to automatic tameness only.
\subsection{A wildly ramified cyclic cover with one smooth exceptional curve}

\begin{example}[Proved: failure of automatic radial tameness]
\label{D:one-curve-wild-swap-cover}
Let $k$ be any algebraically closed field of characteristic $2$, and
let $C/k$ be a smooth projective connected curve of genus $g\geq2$.
There is a complete excellent normal local $k$-domain $R$ of dimension
two, with residue field $k$, and a projective regular resolution
\[
       \pi_R:Y_R\longrightarrow\operatorname{Spec}R
                                                        \tag{D.WS.1}
\]
whose reduced exceptional divisor is a single smooth curve $D$ of
genus $g$ and square $D^2=1-g$. Its puncture
$U=\operatorname{Spec}R\setminus\{\mathfrak m_R\}$ has a connected
finite \'etale $\mathbb Z/2$-torsor which is not radially tame along
$D$. In the corresponding extension of divisorial discrete valuation
rings, the ramification index is $1$ and the residue extension is
purely inseparable of degree $2$. The cover does not extend to a
finite \'etale cover of $Y_R$.

Consequently the hypothesis of a radially tame outer or permutation
quotient in the one-curve crown results does not follow merely from
the existence of one smooth exceptional curve. This example makes
no claim against the representation-count bounds or finite generation.
\end{example}

\begin{proof}
Put $W=C\times C$, let $\sigma$ interchange the two factors, and let
\[
       f:W\longrightarrow Y=\operatorname{Sym}^2 C
                                                        \tag{D.WS.2}
\]
be the degree-two quotient. Write $\Delta\subset W$ for the diagonal
and $D\subset Y$ for its reduced image. We include the local
calculation because ramification index alone does not detect the
wild ramification here.

At a point $(c,c)$, choose the same uniformizer on both factors.
The completed local quotient is
\[
 \begin{aligned}
 k[[x,y]]^{\langle\sigma\rangle}&=k[[u,v]],
       &u=x+y,\quad v=xy,\\
 k[[x,y]]&=k[[u,v]][x]/(x^2+ux+v).
 \end{aligned}                                           \tag{D.WS.3}
\]
The symmetric-power-series identity follows from the elementary
symmetric-polynomial identity by completion; the two adic topologies
are cofinal by the displayed monic relation. Finite-group invariants
commute with this flat completion, since they are the kernel of the
finite collection of maps $\sigma-1$. The right-hand presentation is
finite free with basis $1,x$ and identifies the complete local rings.
Away from $\Delta$ the constant group $\mathbb Z/2$ acts freely, so
the quotient is \'etale there. These descriptions show that $Y$ is a
smooth projective surface and that $f$ is finite flat of degree two.

The equation of $D$ is $u=0$, whereas that of $\Delta$ is $x+y=0$.
In particular
\[
       f^*D=\Delta,\qquad
       f_\Delta:\Delta\longrightarrow D
       \text{ is relative Frobenius of degree }2.
                                                        \tag{D.WS.4}
\]
Indeed modulo $u$ the map is $k[[v]]\to k[[x]]$, $v\mapsto x^2$.
The same description of the generic residue fields identifies $D$
with the Frobenius twist of $C$; thus $D$ is smooth of genus $g$.

We construct the contracting line bundle on $Y$ itself. Choose an
effective canonical divisor $K_C=\sum_c n_c c$, and let
$X_c=\{c+z:z\in C\}\subset\operatorname{Sym}^2C$. Define
\[
                  H=\sum_c n_c X_c,\qquad L=H+2D.
                                                        \tag{D.WS.5}
\]
Each $X_c$ is Cartier on the smooth surface, and divisor pullback gives
\[
            f^*H=\operatorname{pr}_1^*K_C+
                         \operatorname{pr}_2^*K_C.
                                                        \tag{D.WS.6}
\]
This also holds at $2c$: there $X_c$ has equation $v=0$, whose
pullback is $xy=0$. Thus this construction uses an actual divisor;
it does not assume descent of a line bundle through a wild quotient
from triviality of stabilizer actions. The pullback of $H$ is ample,
because $\deg K_C=2g-2>0$, and finite surjective descent of ampleness
makes $H$ ample.

The normal bundle of $\Delta$ is $T_C$. Intersection theory and
(D.WS.4)--(D.WS.6) therefore give
\[
 \begin{aligned}
 \Delta^2&=2-2g,& D^2&=1-g,\\
 H\cdot D&=2g-2,& H^2&=4(g-1)^2,\\
 L\cdot D&=0,& L^2&=4g(g-1)>0.
 \end{aligned}                                           \tag{D.WS.7}
\]
For every irreducible curve $B\ne D$, the two terms
$H\cdot B>0$ and $D\cdot B\geq0$ imply $L\cdot B>0$.
Hence $L$ is nef and big, with null locus exactly $D$.
Moreover
\[
 f_\Delta^*(L|_D)
       \simeq\omega_C^{\otimes2}\otimes T_C^{\otimes2}
       \simeq\mathcal O_C.
                                                        \tag{D.WS.8}
\]
The finite map $f_\Delta$ is flat of degree two between smooth
curves. Applying its norm to (D.WS.8) gives
$(L|_D)^{\otimes2}\simeq\mathcal O_D$. Thus $L|_D$ is semiample.

\citet[Theorem0.2, or directly Proposition1.6 with the effective
Cartier divisor $2D$]{C_Keel1999WildSwap} now implies that $L$ is
semiample. The direct application uses $L=H+2D$ with $H$ ample;
Proposition1.6 and its descent lemma are proved on pp.~261--264.
The related contraction of the diagonal of $C\times C$ by
$\omega_{C\times C}(2\Delta)$ is already discussed in
\citet[Section3, Theorem3.0 and Lemma3.2]{C_Keel1999WildSwap}.
We use that construction here only to produce the stated wild-cover
obstruction, and make no novelty claim for the diagonal contraction.

Take a basepoint-free positive power of $L$, its morphism to the
projective image, and its normal Stein factor. This gives a
projective birational morphism
\[
                  \pi:Y\longrightarrow X
                                                        \tag{D.WS.9}
\]
to a normal projective surface, with $\pi_*\mathcal O_Y=\mathcal O_X$.
Its contracted curves are exactly the curves of $L$-degree zero,
so $D$ is the only one. Birationality follows equally from bigness
and the normal Stein factor. Proper birational quasi-finite maps
to a normal scheme are isomorphisms; thus $\pi$ is an isomorphism
off $D$. Its connected fibres show that the set-theoretic fibre
over $x=\pi(D)$ is exactly $D$.

Apply Stein factorization once more, now to $\pi f$:
\[
             W\xrightarrow{\rho} Z\xrightarrow{h}X.
                                                        \tag{D.WS.10}
\]
Here $Z$ is normal, $h$ is finite of generic degree two, and
$\rho$ has geometrically connected fibres; see
\citet[Tag03H0, complete Noetherian Stein-factorization theorem]
{C_StacksWildSwap}. Over $X\setminus\{x\}$ this is the original
finite map $f$, so $\rho$ is an isomorphism there and $h$ is a
$\mathbb Z/2$-torsor there. The set-theoretic fibre of $\pi f$
over $x$ is $\Delta$. It is connected and surjects onto the finite
set $h^{-1}(x)$, so that set consists of a single point $z$.

Set
\[
       R=\widehat{\mathcal O}_{X,x},\qquad
       A=\widehat{\mathcal O}_{Z,z}.
                                                        \tag{D.WS.11}
\]
The local rings before completion are excellent normal local
domains essentially of finite type over $k$. Their completions
are normal local domains; in particular they do not split into
several branches. These are the normal-completion consequences
of excellence in \citet[Tags07QS,0C23 and0C22]
{C_StacksWildSwap}. Since $z$ is the unique point above $x$,
completion of the finite algebra gives exactly the local
$R$-algebra $A$, rather than a product of several completions.
The degree remains two. The involution extends through Stein
factorization and completion, and its invariants are $R$:
flat completion commutes with the kernel defining invariants.
The finite map
\[
 \operatorname{Spec}A\setminus\{\mathfrak m_A\}
                 \longrightarrow U
                                                        \tag{D.WS.12}
\]
is therefore \'etale of degree two with the given involution.
Its source is a nonempty open of the integral scheme
$\operatorname{Spec}A$, hence is connected. Thus (D.WS.12) is
a connected $\mathbb Z/2$-torsor, not a cover which splits after
completion.

Base change of (D.WS.9) to $\operatorname{Spec}R$ gives (D.WS.1).
The completion map is regular by excellence. Its finite-type
base changes are regular, and a regular map over a regular
scheme has regular source; see \citet[Tags07C1 and07QI]
{C_StacksWildSwap}. Hence $Y_R$ is regular. Its exceptional
fibre is still the single smooth curve $D$, with the same
normal bundle and self-intersection. Similarly $W_R$ is regular.
The finite map $W_R\to Y_R$ is the normalization in the function
field of (D.WS.12). Its divisor above $D$ is still $\Delta$,
and (D.WS.4) is unchanged by this base change. At their generic
points it follows that
\[
      e(\Delta/D)=1,\qquad
      [k(\Delta):k(D)]_{\rm insep}=2.
                                                        \tag{D.WS.13}
\]
The full involution acts trivially on the residue field, so
the inertia group has order two. This is fierce wild
ramification: tameness also requires a separable residue
extension, not just a prime-to-characteristic ramification
index. It excludes extension of the punctured torsor to a
finite \'etale cover of $Y_R$, since the finite normal
extension is unique.

The Artin--Schreier description can be seen without an
inertia convention. In (D.WS.3), after inverting $u$, put
$w=x/u$. Then
\[
                        w^2+w=\frac{v}{u^2}.
                                                        \tag{D.WS.14}
\]
In the completed divisorial field the leading residue $v$
is not a square in $k((v))$. An Artin--Schreier correction
of pole order one could cancel the pole of order two only
if its leading coefficient squared to $v$; larger pole
orders produce a strictly larger uncancellable leading
pole. Thus this local equation also displays the failure
of an unramified extension.

Finally, this is genuinely an obstruction already among
$p$-cyclic classes. Normality and proper birationality give
$\Gamma(U,\mathcal O_U)=\Gamma(Y_R,\mathcal O_{Y_R})=R$.
The map $a\mapsto a^2-a$ is surjective on the strictly
henselian ring $R$. The Artin--Schreier exact sequence
therefore identifies $H^1(-,\mathbb F_2)$ on both spaces
with the Frobenius-fixed part of $H^1(-,\mathcal O)$.
The class of (D.WS.12) is outside the image of
$H^1(Y_R,\mathbb F_2)\to H^1(U,\mathbb F_2)$, by
(D.WS.13). Thus finiteness of the latter cohomology group
does not force its classes to extend across $D$.
\end{proof}


% Proved post-CP22 corollary. Auxiliary repair for the swap obstruction.
\begin{corollary}[Proved: exact wild image in the symmetric-square example]
\label{D:wild-swap-exact-image}
In Example~\ref{D:one-curve-wild-swap-cover}, let
$\Gamma_R=\pi_1(U)$ and fix a compatible geometric
divisorial basepoint at $D$. The image of its absolute
wild inertia in $\Gamma_R$ has order exactly two.
Moreover
\[
                         d(\Gamma_R)\leq2g+2.
                                                        \tag{D.WE.1}
\]
Thus the example disproves automatic radial tameness
while its punctured fundamental group is finitely
generated. The assertion about order two concerns
the image of one divisorial wild-inertia subgroup;
it does not identify its normal closure.
\end{corollary}

\begin{proof}
Use the notation $W=C\times C$, $Z$, $A$ and $R$
of that example, and put
$U_A=\operatorname{Spec}A\setminus\{\mathfrak m_A\}$.
The connected degree-two torsor gives a normal open
subgroup
\[
              \Gamma_A=\pi_1(U_A)\mathrel{\trianglelefteq}
                 \Gamma_R,\qquad[\Gamma_R:\Gamma_A]=2.
                                                        \tag{D.WE.2}
\]
The regular resolution $W_R\to\operatorname{Spec}A$
has sole reduced exceptional curve $\Delta\simeq C$.
The first projection $W_R\to C$ restricts to the
identity on $\Delta$, so gives a retraction of its
entire formal neighborhood. Its normal degree is
$\Delta^2=2-2g<0$. Proposition~\ref{D:formal-retraction-tame}
therefore says that every finite \'etale cover of
$U_A$ is radially tame along $\Delta$.

Let $K$ be the completed divisorial field at $D$
and $L$ the one at $\Delta$. Then $L/K$ is the
separable degree-two extension calculated in the
example. Choose their absolute Galois groups
compatibly, so $G_L\subset G_K$, and let $P_K,P_L$
be their wild inertia groups. The valuation comparison
(WI) in Proposition~\ref{D:finite-cone-wild-bound}
gives
\[
                   P_L=P_K\cap G_L.
                                                        \tag{D.WE.3}
\]
The maps to the punctured fundamental groups commute
with $G_L\subset G_K$ and $\Gamma_A\subset\Gamma_R$.
Every finite quotient of $\Gamma_A$ kills $P_L$ by
the preceding formal-retraction conclusion. Thus the
full image of $P_K$ in $\Gamma_R$ has order at most
\[
             [P_K:P_K\cap G_L]\leq[G_K:G_L]=2.
                                                        \tag{D.WE.4}
\]
This is a bound on the full profinite image, not just
on each quotient separately: the same fixed subgroup
$P_K\cap G_L$ is killed in all the quotients.

The residue extension in $L/K$ is purely inseparable,
so the entire Galois group $\mathbb Z/2$ is inertia.
As it is a $2$-group, it is its own wild inertia.
Consequently $P_K$ maps onto the quotient
$\Gamma_R/\Gamma_A=\mathbb Z/2$, proving that the
upper bound in (D.WE.4) is attained.

Finally $\Gamma_A$ equals its radially tame quotient.
The curve $\Delta$ has genus $g$, so
Proposition~\ref{D:one-curve-radial-tame-generators}
gives $d(\Gamma_A)\leq2g+1$. Adjoin one lift of
the generator of $\Gamma_R/\Gamma_A$ to a topological
generating tuple of $\Gamma_A$. Its closed generated
subgroup contains $\Gamma_A$ and surjects onto the
quotient, hence equals $\Gamma_R$. This proves
(D.WE.1). All inputs beyond the explicit example
are the already proved formal-retraction and tame
generator statements; no general wild-cover atlas
or unrestricted crown assertion is inferred.
\end{proof}


% Separate post-CP23 proved input. All four complete mathematical audits passed.
\subsection{The different controls the invariant Frobenius lattice}

\begin{lemma}[Proved: a rank-independent invariant-lattice bound]
\label{D:different-Frobenius-invariant-lattice}
Let $k$ be algebraically closed of characteristic $p>0$, let
$A=k[[t]]$, and let $B=k[[s]]$ be the integral closure of $A$
in a finite Galois extension $L/k((t))$, with group $G$.
Write $e=[L:k((t))]=v_s(t)$ and
$\mathfrak D_{B/A}=s^\delta B$ for its different.
Fix $q=p^f$, a representation $\rho:G\to\mathrm{GL}_r(\mathbb F_q)$,
and its coefficient twists $\rho_a=\rho^{[p^a]}$, indexed
cyclically by $a\in\mathbb Z/f\mathbb Z$.
Give $W_a=B^r$ the semilinear action
$g(w)=\rho_a(g)g(w)$, where the second occurrence of $g$
acts on the coefficients, and put $E_a=W_a^G$.
Then $E_a$ is a free $A$-module of rank $r$.
Coordinatewise $p$-th power induces an injective linearization
\[
 \phi_a:F_A^*E_a\longrightarrow E_{a+1},\qquad
 t^cE_{a+1}\subset\operatorname{im}\phi_a\subset E_{a+1},
 \quad c=\left\lceil\frac{p\delta}{e}\right\rceil .
 \tag{D.DF.1}
\]
The same $c$ works at every phase and is independent of $r$ and $f$.
More generally the $p^h$-power map from phase $a$ to phase $a+h$
satisfies (D.DF.1) with $c_h=\lceil p^h\delta/e\rceil$.

If $G^v=1$ for all upper indices $v>b$, where $b\ge0$, then
\[
                 \frac\delta e\le b+1,
       \qquad c\le\lceil p(b+1)\rceil.                 \tag{D.DF.2}
\]
For the trivial extension $\delta=0$ and every displayed map
is an isomorphism. These statements concern the actual invariant
lattices for a fixed local Galois representation. They impose no
effectiveness assertion for arbitrary families of local types.
\end{lemma}

\begin{proof}
All actions in the proof fix $k$, so they preserve its specified
copy of $\mathbb F_q$. The extension has no residual part because
$k$ is algebraically closed. In particular it is totally ramified,
$B$ is finite free of rank $e$ over $A$, and its generic Galois
group has order $e$.

\emph{A trace calculation with no division by $|G|$.}
Choose an integral $A$-basis $b_1,\ldots,b_e$ of $B$ and its
generic trace-dual basis $b_1^*,\ldots,b_e^*$ in $L$.
By the definition of the complementary module and the different,
\[
                 s^\delta b_i^*\in B.                 \tag{D.DF.3}
\]
These trace identifications, including the equality with
$\operatorname{Hom}_A(B,A)$, are
\citep[Section~49.8, especially Lemmas~49.8.2--3 and their
complete proofs, Tag~0BW0]{StacksDifferentLatticeB}.
The trace-dual identity for a Galois extension is
\[
                 \sum_i b_i g(b_i^*)=
                   \begin{cases}1&g=1,\\0&g\ne1.\end{cases}
                                                               \tag{D.DF.4}
\]
For example it follows by multiplying the embedding matrix
$(g(b_i))_{g,i}$ by the transpose of its trace-dual embedding
matrix; the product is the identity.
For $w\in W_a$ and $d=s^\delta$, (D.DF.3)--(D.DF.4) give
\[
 d w=\sum_i b_i\left(\sum_{g\in G}g(d b_i^*w)\right).
                                                               \tag{D.DF.5}
\]
Each inner sum is integral and invariant, hence belongs to $E_a$.
The factor $d$ is not assumed invariant; its conjugate occurs
\emph{inside} each inner sum. Thus, writing
$L_a=\operatorname{im}(B\otimes_AE_a\to W_a)$, we have
\[
                    s^\delta W_a\subset L_a\subset W_a.
                                                               \tag{D.DF.6}
\]

The invariant module $E_a$ is finite and torsion-free over the
DVR $A$, so it is free. For completeness, independent invariant
vectors over $K=k((t))$ remain independent over $L$: a shortest
$L$-linear dependence, normalized to have one coefficient $1$,
can be subtracted from each of its $G$-translates. Minimality
forces every coefficient into $L^G=K$, a contradiction.
It follows that $\operatorname{rank}_AE_a\le r$ and that
$B\otimes_AE_a\to W_a$ is injective. Equation (D.DF.6) makes
its generic image all of $L^r$, giving equality of the ranks.
In particular this is integral descent by a lattice calculation,
not an application of averaging with $1/|G|$.

\emph{Frobenius and faithfully flat descent.}
The map $w\mapsto w^{[p]}$ carries invariants for $\rho_a$ to
invariants for $\rho_{a+1}$, since
$(\rho_a(g)g(w))^{[p]}=\rho_{a+1}(g)g(w^{[p]})$.
Let $L_{a,p}$ be the $B$-span of this image inside $W_{a+1}$.
Apply coordinatewise Frobenius to the generators in (D.DF.6).
The constant coordinate vectors of $B^r$ show that
\[
 s^{p\delta}W_{a+1}\subset L_{a,p}\subset L_{a+1}.
                                                               \tag{D.DF.7}
\]
If $ec\ge p\delta$, multiplication by $t^c$ takes $L_{a+1}$
into the leftmost module. Under the injective identifications
$B\otimes_AE_{a+1}=L_{a+1}$, this is precisely the scalar
extension of $t^cE_{a+1}\subset\operatorname{im}\phi_a$.
The finite free extension $B/A$ is faithfully flat, so the
inclusion descends to $A$.
The linearization is generically an isomorphism by (D.DF.7),
and its free source has no torsion, so it is injective.
Repeating the identical argument with $p^h$ proves the more
general assertion. No replacement of $\mathbb F_q^r$ by its
$fr$-dimensional restriction of scalars is made.

\emph{The upper-break comparison.}
We include the calculation to specify the endpoint convention.
Write $G_i$ for the lower ramification groups, and put
$G_u=G_{\lceil u\rceil}$ for real $u\ge0$. The Herbrand
function and upper indexing are
\[
 \varphi(v)=\int_0^v\frac{|G_u|}{e}\,du,
                 \qquad G^{\varphi(v)}=G_v.
                                                               \tag{D.DF.8}
\]
This is the convention of
\citet[Section~10, printed pp.~28--29]{Kedlaya2010DifferentLatticeB}.
The monogenic description of $B/A$ is Lemma~10.1 there.
If $P$ is the minimal polynomial of a uniformizer $s$, then
the monogenic different formula gives
\[
 \delta=v_s(P'(s))
     =\sum_{g\ne1}v_s(s-g(s))
     =\sum_{i\ge0}(|G_i|-1).                           \tag{D.DF.9}
\]
Here the derivative formula is
\citep[Lemmas~49.9.3 and~49.12.1--3, Example~49.12.5,
Tags~0BW6 and~0BWB, with their complete proofs]
{StacksDifferentLatticeB}; the last equality simply counts
the integers $i$ with $i+1\le v_s(s-g(s))$ for each $g\ne1$.
Thus the use of the different is compatible with (D.DF.3).

If the extension is nontrivial, let $a\ge0$ be the largest
integer with $G_a\ne1$. Its highest upper break is
$b_0=\varphi(a)\le b$; in the tame case $a=b_0=0$.
Equations (D.DF.8)--(D.DF.9) yield the exact equality
\[
 \frac\delta e
 =1-\frac1e+\frac1e\sum_{i=1}^a(|G_i|-1)
 =1+b_0-\frac{a+1}{e}\le 1+b.
                                                               \tag{D.DF.10}
\]
The trivial extension is immediate. This proves (D.DF.2).
\end{proof}


% Proved after four complete actual-file audits; assembled by root.
\subsection{A fixed detecting curve and its completed local sources}

\begin{lemma}[Proved: finite punctures give fixed local types on one curve]
\label{D:finite-puncture-detecting-curve-local-types}
Let $S$ be a normal proper integral surface over an algebraically
closed field $k$ of characteristic $p>0$, let $T$ be a finite closed
subset, and put $U=S\setminus T$ and $\Gamma=\pi_1(U)$.
There are a finite set $F\supset T$ of closed points of $S$, an
algebraically closed extension $\Omega/k$ of finite transcendence
degree, a smooth projective connected curve $C/\Omega$, a finite
nonempty marked set $P\subset C$, and a morphism
$C^\circ=C\setminus P\to U$ with the following properties.
\begin{enumerate}
\item The induced homomorphism $\pi_1(C^\circ)\to\Gamma$ is
surjective. Its target is the original group over $k$; no
surjectivity onto $\pi_1(U_\Omega)$ is asserted.
\item Put $R_s=\widehat{\mathcal O}_{S,s}$ and
$\Lambda_s=\pi_1(\operatorname{Spec}R_s\setminus\{\mathfrak m_s\})$
for $s\in F$. These are fixed complete excellent normal local
surface rings over the original field $k$. For every $x\in P$
there is an $s=s(x)\in F$ and a fixed morphism
\[
 \operatorname{Spec}\Omega((t_x))\longrightarrow
 \operatorname{Spec}R_s\setminus\{\mathfrak m_s\}
 \longrightarrow U .                                      \tag{D.DC.1}
\]
The corresponding homomorphisms of fundamental groups, with
compatible paths, give the local restriction of
$\pi_1(C^\circ)\to\Gamma$. Thus each full local representation
type at $x$ is a restriction of one representation of the fixed
group $\Lambda_s$, for every finite target group. If $s\notin T$,
the composite $\Lambda_s\to\Gamma$ is trivial.
\item There are fixed real numbers $b_x\ge0$, independent of the
finite quotient of $\Gamma$, which bound the highest classical
upper break of each induced local quotient at $x$. In particular,
if $L/\Omega((t_x))$ is its nontrivial Galois splitting extension,
with ramification index $e$ and different exponent $\delta$, then
\[
                    \delta/e\le b_x+1.                    \tag{D.DC.2}
\]
\end{enumerate}
All choices in this statement are independent of ranks, coefficient
fields and individual finite covers. Item~(2) concerns actual full
local types, not only their ramification breaks.
\end{lemma}

\begin{proof}
Choose the smooth projective surface $Y$ and SNC boundary $B$
used in Proposition~\ref{D:surface-curvewise-conductor} for the
pair $(S,T)$. Thus $Y\to S$ is proper birational, and after
enlarging a finite set $F\supset T$ its complement
$Y\setminus B$ maps isomorphically onto $S\setminus F$.
This uses normality in dimension two: a proper birational
morphism to $S$ is an isomorphism at the generic point of every
divisor, so its exceptional image is a finite set of closed points.
Chow's lemma and surface resolution produce this projective $Y$
even when $S$ itself is not projective. These are the same
choices and facts used in the cited proposition and in
Theorem~\ref{D:generic-surface-finite-punctures}.
If necessary add one smooth closed point to $F$ and blow it up;
this makes $B$ nonempty. The added exceptional curve corresponds
to a retained smooth point of $U$ and introduces no ramification
of a cover of $U$.

Fix a projective embedding $Y\hookrightarrow\mathbf P^N_k$.
Let $\Omega$ be an algebraic closure of
$k(z_0,\ldots,z_N)$, with the $z_i$ algebraically independent,
and let $C$ be the hyperplane section with these coefficients.
The geometric generic hyperplane is smooth and geometrically
irreducible. It avoids the finitely many nodes of $B$ and meets
every component of $B$ transversely. Indeed apply smooth Bertini
also to each smooth component of $B$; their zero-dimensional
sections are then reduced. Set $P=C\cap B_\Omega$. It is finite
and nonempty, since a hyperplane has positive intersection with
every projective curve. Smooth and irreducible Bertini here are
\citep[Lemmas~33.47.3 and~37.32.1--3, Tags~0FD6 and~0G4C,
with the written proofs]{StacksDGenericSurface}.

We check the simultaneous-cover quantifier explicitly. Let
$V\to U$ be any connected finite etale cover defined over $k$.
Its restriction $V_0$ to $S\setminus F$ is smooth and integral:
the total space $V$ is connected and normal, hence integral,
and a nonempty open of an integral scheme is integral.
Pull the hyperplane system back from $Y$ to $V_0$. It is
basepoint-free. At any closed point $v\in V_0$, etaleness and
the immersion in $\mathbf P^N$ allow two hyperplanes through
its image whose pullbacks have independent differentials at $v$.
Their intersection has a codimension-two component through $v$.
A third hyperplane avoiding the image of $v$ does not contain
that component. These are precisely conditions~(3) and~(4) of
Stacks, Lemma~37.32.3. Hence the geometric generic pullback
$V_0\times_{Y}C$ is geometrically irreducible, and is smooth
because it is etale over $C^\circ$. It is therefore connected.
The generic coefficients lie in the Bertini open for each
original cover separately. The one chosen generic point works
for all of them; no intersection of infinitely many open sets
at a closed point, or countability assumption, is used.
The connected-cover criterion proves item~(1).

For $x\in P$ let $s$ be its image in $F$. The map on local rings
from $\mathcal O_{S,s}$ to $\mathcal O_{C,x}$, followed by completion,
takes $\mathfrak m_s$ into $(t_x)$ in $\Omega[[t_x]]$.
It is therefore continuous and extends uniquely to
\[
                 R_s\longrightarrow\Omega[[t_x]].          \tag{D.DC.3}
\]
Its generic point does not map to the closed point of
$\operatorname{Spec}R_s$: the generic point of $C$ lies in
$Y\setminus B$, so some element of $\mathfrak m_s$ has nonzero
image in $\Omega[[t_x]]$, and hence invertible image in
$\Omega((t_x))$. This proves the first arrow in (D.DC.1).
The punctured spectrum of $R_s$ maps to $S\setminus T$.
Indeed the only point of $\operatorname{Spec}R_s$ above $s$
is its maximal ideal, and $\operatorname{Spec}\mathcal O_{S,s}$
has no other closed point of $S$ in its image. The composition
agrees with the punctured formal branch of $C^\circ\to U$,
by the uniqueness in (D.DC.3). This proves item~(2), including
compatibility for arbitrary finite torsors and their morphisms.
There is no appeal to invariance of punctured fundamental groups
under extending $k$, and no completed tensor product is needed.
The rings $R_s$ are normal domains by excellence and normality
of the finite-type local rings before completion. They have
coefficient and residue field $k$.
If $s\notin T$, a finite etale cover of $U$ restricts to one of
$\operatorname{Spec}R_s$. This strictly henselian local scheme
has only split finite etale covers, proving the last assertion.

For completeness we distinguish the coefficient characteristic
in the ramification argument. Apply the proof of
Proposition~\ref{D:surface-curvewise-conductor} over the fixed
algebraically closed field $\Omega$, to $(S_\Omega,T_\Omega)$
and the fixed SNC resolution $Y_\Omega$. It supplies a fixed
effective divisor $R$ on $Y_\Omega$ controlling all finite-free
systems of coefficient characteristic $\ell\ne p$ whose
monodromy comes from $\pi_1(U_\Omega)$. It is unnecessary
that its numerical coefficients equal those chosen over $k$.
Every original finite cover pulls back to such a cover over
$\Omega$. For its finite group $H$, take the faithful regular
representation $\mathbb F_\ell[H]$. The curvewise assertion
bounds its highest logarithmic slope at $x$ by
$b_x=\operatorname{mult}_x(C^*R)$, independently of $H$ and
its order. Over the perfect residue field $\Omega$, logarithmic
upper numbering is classical upper numbering. Faithfulness of
the regular representation means that every upper group above
$b_x$ in the induced quotient is trivial. This gives item~(3).
The curvewise comparison is exactly the SNC assertion of
\citet[Theorem~3.5, Remark~3.9 and Proposition~5.8(3)]
{C_HuTeyssierRam2025}; it uses the highest slope, not the
rank-weighted Swan conductor. Finally the elementary
different/Herbrand calculation in
Lemma~\ref{D:different-Frobenius-invariant-lattice} gives
(D.DC.2). In particular the bound uses the ramification index
of the curve's totally ramified local extension. No replacement
of a surface ramification index by a degree with a possibly
inseparable residue factor is made.
\end{proof}


% Proved after four actual-file audits. Separate post-CP23 research; ready for integration.
\subsection{Quadratic counts on a curve with prescribed full local types}

\begin{lemma}[Proved: finite integral jets determine cyclic unit identifications]
\label{D:curve-integral-unit-jet-comparison}
Let $A$ be any algebra over a field of characteristic $p>0$.
Write $F$ for coefficientwise $p$-th power, including $t\mapsto t^p$,
on $A[[t]]$. Fix $f,r\geq1$ and cyclic families of matrices
$U_a,V_a\in\operatorname{Mat}_r(A[[t]])$, indexed by
$a\in\mathbb Z/f\mathbb Z$. Assume that they are invertible over
$A((t))$ and that $t^cU_a^{-1}$ and $t^cV_a^{-1}$ are integral,
for one fixed integer $c\geq0$.
Choose integers $N\geq1$ and $L=N+c$ with $(p-1)N>c$.
Suppose that invertible matrix jets $h_a$ modulo $t^L$ satisfy
\[
 h_{a+1}U_a-V_aF(h_a)=0\pmod {t^L}.
 \tag{D.CL.1}
\]
There is a unique family of integral invertible matrices $H_a$
congruent to $h_a$ modulo $t^N$ and satisfying
\[
                  H_{a+1}U_a=V_aF(H_a).
 \tag{D.CL.2}
\]
The construction commutes with every homomorphism of coefficient
algebras. Only the retained $N$-jet, not necessarily the whole
$L$-jet, is preserved. Conversely every integral unit identification
supplies jets satisfying (D.CL.1).
\end{lemma}

\begin{proof}
Lift the jets arbitrarily to integral invertible matrices $H^{(0)}_a$.
Such lifts exist since invertibility of a matrix over $A[[t]]$ is
equivalent to invertibility of its constant term. Define
\[
 (\mathcal T H)_{a+1}=V_aF(H_a)U_a^{-1}.
 \tag{D.CL.3}
\]
The first difference $\mathcal T H^{(0)}-H^{(0)}$ is divisible by
$t^{L-c}=t^N$. If two families differ by $t^j$, their images under
$\mathcal T$ differ by $t^{pj-c}$. Thus $\mathcal T$ preserves
the integral $t^N$-ball about $H^{(0)}$ and is contracting there,
since $pN-c>N$. Every iterate remains invertible by its unchanged
constant term. The $t$-adic limit exists in every coefficient ring
and satisfies (D.CL.2). The same estimate proves uniqueness in
that ball. Different arbitrary lifts of the prescribed jets have
the same limit, and uniqueness proves compatibility with base
change. No division by $p$, no flatness of Frobenius on $A$, and
no bound on the nilpotence of $A$ are used.
\end{proof}

\begin{theorem}[Proved: fixed full local linear types have quadratic growth]
\label{D:curve-fixed-local-linear-growth}
Let $X$ be a smooth proper connected curve over an algebraically
closed field $k$ of characteristic $p>0$, let $D$ be a nonempty
finite reduced divisor, and put $X^\circ=X\setminus D$.
Fix a real number $b\geq0$. There is a constant $C$, depending
only on $(X,D,b,p)$, with the following property.
For $q=p^f$, $r\geq1$, and each $x\in D$, fix a continuous
local representation
\[
 \tau_x:\operatorname{Gal}(\overline{k((t_x))}/k((t_x)))
                  \longrightarrow\operatorname{GL}_r(\mathbb F_q)
\]
up to conjugacy. Suppose that every quotient defined by $\tau_x$
has trivial upper ramification groups above $b$.
Then the number of conjugacy classes of continuous representations
$\pi_1(X^\circ)\to\operatorname{GL}_r(\mathbb F_q)$ whose full
local types are the prescribed $\tau_x$ is at most
\[
                              q^{Cr^2}.
 \tag{D.CL.4}
\]
The constant is independent of the local splitting extensions,
their degrees, the actual matrices of the local types, $r$ and $f$.
\end{theorem}

\begin{proof}
We give both the finite presentation and the integral local
comparison needed for the component count. Bounds for ramification
breaks alone are not used as a replacement for full local types.

\emph{Canonical lattices and bounded slopes.}
For a global representation in question, normalize $X$ in a finite
Galois cover trivializing it. Write its finite group as $G$ and its
coefficient twists as $V_a=V^{[p^a]}$, for $a\in\mathbb Z/f\mathbb Z$.
The coherent sheaves
\[
                 E_a=(\pi_*\mathcal O_{X'}\otimes_k V_a)^G
 \tag{D.CL.5}
\]
are torsion-free, hence vector bundles of rank $r$ on the regular
curve $X$. Here $V_a$ denotes the corresponding $k$-vector space
with its specified finite-field matrices. Invariants are kernels;
no exactness of $G$-invariants in characteristic $p$ is asserted.
At $x\in D$, their completions are exactly the invariant lattices
of Lemma~\ref{D:different-Frobenius-invariant-lattice} for $\tau_x$
and its coefficient twists. This follows by the decomposition of
the completed normalization into its local factors and taking
invariants in the induced module. In particular these local
lattices depend functorially only on the local representation.
Coordinatewise Frobenius gives maps
\[
 F_X^*E_a\xrightarrow{\phi_a}E_{a+1},\qquad
 E_{a+1}(-cD)\xrightarrow{\psi_a}F_X^*E_a,
 \quad c=\lceil p(b+1)\rceil,
 \tag{D.CL.6}
\]
whose two composites are multiplication by the canonical Cartier
section of $cD$. The different estimate in the cited lemma proves
these assertions at every marked completion; they are isomorphisms
on $X^\circ$, so they hold globally. The same conclusion includes
trivial local types.

For a vector bundle $E$ on a curve one has
$p\mu_{\max}(E)\leq\mu_{\max}(F_X^*E)$, by pulling back its
maximal-slope subbundle. An injection of bundles bounds the maximal
slope of the source by that of the target. Applying this first to
$\phi_a$ and then to the dual of $\psi_a$, and cycling the $f$
inequalities, gives
\[
 -\frac{c\deg D}{p-1}\leq\mu_{\min}(E_a)
                  \leq\mu_{\max}(E_a)\leq0.
 \tag{D.CL.7}
\]
This argument requires no semistability assertion for tensor
products in positive characteristic. Every fixed twist of these
bundles has $h^0+h^1=O(r)$, with constants depending only on the
fixed curve and the twist.

\emph{An actual polynomial atlas with fixed terms.}
We spell out the ordinary-curve specialization of the fixed
partial-resolution construction in
Lemma~\ref{D:fixed-partial-resolution-unit-atlas}. There is no
root stack, normal layer or moving determinant parameter here.
Choose a fixed line bundle $H$ of sufficiently large positive degree
so that $E_a^+=E_a\otimes H$ has minimum slope greater than
$2g(X)+2$ whenever (D.CL.7) holds. A general $(r+1)$-tuple of
sections generates $E_a^+$: the failure locus of a general such
tuple has expected codimension two on the curve. Its kernel is
a line $L_a^{-1}$, where $L_a=\det E_a^+$ is globally generated
and has $H^1=0$. Choose a fixed positive line bundle $B$ and
$b_r=\kappa r$ with one fixed integer $\kappa$ large enough that
$L_a^{-1}B^{b_r}$ is also globally generated with $H^1=0$ for every
allowed $E_a$. These choices are possible because $\deg L_a=O(r)$
and the lower bound on its degree is strictly above $2g(X)-2$.

Use the fixed terms
\[
 P_0=\mathcal O_X^{r+1},\qquad
 P_j=(B^{-j b_r})^{\oplus2}\quad(1\leq j\leq7).
 \tag{D.CL.8}
\]
Two generators at each later step give exact middle terms with
alternating line syzygies
$L_a^{-1}B^{-2v b_r}$ and $L_aB^{-2v b_r}$.
The generating line at an odd step is $L_a^{-1}B^{b_r}$ and
at an even step is $L_a$. Their $H^1$ vanishing also lifts any
specified presentation across a small extension of local Artinian
coefficient algebras. This proves coverage, including coverage of
such prescribed lifts, without choosing a basis of a moving Picard
variety.

Parameterize the actual global differentials between (D.CL.8),
impose $d^2=0$, and impose the open rank-one conditions on every
consecutive map. Over any coefficient algebra the unit-minor
cancellation proves that the middle cohomology is zero and the two
end cohomology modules are bundles. Thus the degree-zero cohomology
is an actual rank-$r$ bundle $E_a^+$ over the proper curve. The
rank conditions are open on the coefficient scheme because their
failure locus on the proper curve has closed image. The slope
condition (D.CL.7) is the proper-Quot open condition already proved
in the cited atlas. It introduces no equations in the ambient
coordinate count. Untwist by the fixed $H$ to obtain $E_a$.
All these assertions hold over arbitrary coefficient algebras;
rank is imposed by unit minors on the actual local rings.

Here are the exact Hom witnesses used for (D.CL.6). Truncate the
source presentation at length seven and the target at length three.
Their end-error triangles give error intervals $[-3,-2]$ and
$[4,8]$, disjoint from degrees zero, one and two. Consequently
finite derived Hom computes exactly the ordinary Hom of the
bundles in degree zero. Compose through the inclusion of the
length-three target in its length-seven presentation. To test an
identity, precompose with its surjective top term $P_0$; do not
impose an identity on an equal-length partial complex. A degree
minus-one boundary witness tests equality exactly. This is the
separated-length argument (D.FR.5)--(D.FR.7), applied on $X$.

For clarity, the finite complexes are obtained by splitting the
two-affine Cech complexes of the fixed line-bundle terms over $k$.
The only perturbations are the variable internal differentials.
Every perturbation word increases internal degree, whose width is
bounded by seven and by the fixed composition diagram. Hence the
homological perturbation formulas terminate after a fixed number
of factors. They compute cycles, boundaries and composition over
every coefficient algebra, by the explicit identities previously
proved in that construction and
\citep[Sections~2.1--2.4 and~3.2]{C_Crainic2004Perturbation}.
This is a finite algebraic construction, not a dimension estimate
for a deformation space.

Every top-to-top Hom block has multiplicity $O(r^2)$ and fixed line
degree. A block with one positive-index term has multiplicity
$O(r)$ and line degree $O(r)$; a block with two such terms has
bounded multiplicity and line degree $O(r)$. Their total $h^0+h^1$
is $O(r^2)$. Fixed twists, the poles $cD$, and one-step Frobenius
multiply these bounds only by a constant. Frobenius pulls back
locally split complexes exactly over every coefficient algebra,
without a flatness assumption on that algebra, and replaces
variable coefficients by their $p$-th powers. We therefore obtain
all $f$ objects, maps $\phi_a,\psi_a$, and their exact inverse
identities with $O(fr^2)$ variables and equations of degree at
most a constant $d_0$ depending only on the fixed data and $p$.

\emph{Full local types from fixed-length jets.}
Choose the same $N,L=N+c$ as in
Lemma~\ref{D:curve-integral-unit-jet-comparison}. For each $x$ and
phase fix a basis of the canonical local lattice of $\tau_x^{[p^a]}$,
and write its Frobenius matrix as $V_{x,a}(t_x)$. The actual
coefficients of these matrices are constants, however complicated
the splitting cover. Only their first $L$ coefficients occur in
the presentation. Their inverse-pole bound is the common $c$.

On $X\times\operatorname{Spec}A$ impose frames of $E_a$ on
$D_L=\coprod_x\operatorname{Spec}k[t_x]/(t_x^L)$, with inverse
witnesses, and impose (D.CL.1) for the frame matrices and
$V_{x,a}$. These add $O(fr^2)$ variables of degree bounded in
terms of $p$ and $L$. To make the frame construction explicit,
for a local presentation $P_1\xrightarrow d P_0\to E_a\to0$
on $D_L$, use matrices $u:P_0\to A[t]/(t^L)^r$ and
$v:A[t]/(t^L)^r\to P_0$ with
\[
 u d=0,\qquad uv=1,\qquad vu-1=d w.
 \tag{D.CL.9}
\]
The last equality is witnessed by a matrix $w:P_0\to P_1$.
The locally split middle exactness makes these precisely a frame
of the quotient with a chosen lift. All sizes are $O(r)$ and $L$
is fixed. Restriction of arrows uses the separated partial
complexes before taking their cohomology; their locally free end
terms remain so after restriction. The same finite boundary
witnesses therefore express the induced jet maps exactly. There
are no omitted negative Tor terms. Frames of Frobenius pullbacks
are obtained by coefficientwise $p$-th power; the induced maps
are polynomials of bounded degree.

Let $Z$ be the resulting open locus of the affine polynomial
witness scheme. Its underlying affine equations have at most
$A_0fr^2$ variables and degree at most $d_1$, independent of the
local types. Every geometric point gives actual bundles and maps
(there is no formal algebraization step on this proper curve).
The jet comparison lemma gives actual integral cyclic unit
identifications of their completed local lattices with the fixed
local models. On $X^\circ$, the $f$ cyclic arrows give a unit
$q$-Frobenius bundle. The fixed-frame equations give a finite etale
$\operatorname{GL}_r(\mathbb F_q)$-torsor, exactly as in the
previous unit-bundle comparisons. Conversely every representation
being counted has its canonical family (D.CL.5), lies in (D.CL.7),
and is represented in $Z$. The local contractions commute with
coefficient base change.

\emph{Why unit-isomorphism classes are open.}
Fix a represented unit object $E^0$. The functor of isomorphisms
from $E^0$ to the universal family on the proper curve, with all
inverse and cyclic unit equations, is of finite presentation.
Indeed proper Hom of flat coherent bundles is affine of finite
presentation, and the displayed finite equations cut out this
functor
\citep[Proposition~99.3.10, Tag~08K6, and the stated inverse equations]
{StacksUnitJetTransferD}.
We check the required lifting at a specified Artinian point, rather
than only its tangent space. Let $A'\to A$ be a nilpotent small
extension of finite local Artinian $k$-algebras and let a unit
isomorphism be specified over $A$. Nilpotent invariance of finite
etale covers lifts it uniquely on $X^\circ_{A'}$
\citep[Lemma~58.8.3, Tag~0BQB, complete proof]
{StacksUnitJetTransferD}.
At each puncture the retained jets identify both objects, integrally
and compatibly over $A'$ and $A$, with their fixed canonical local
unit lattices. After these identifications, the local reduction
of the specified isomorphism is an automorphism of the constant
local $\mathbb F_q$-system. Its unique constant lift preserves the
canonical invariant lattice: the invariant-lattice construction
(D.CL.5) is functorial for every intertwiner. Full faithfulness of
nilpotent etale invariance says that the restriction of the global
lift is exactly this integral local lift. It therefore has no
pole at any point of $D$; its inverse has none either. Faithful
flatness of completion, applied to the coherent Hom bundles,
extends the isomorphism uniquely over all of $X_{A'}$.

This proves lifting of every specified Artinian point for this
finite-presentation Isom functor. It is formally etale, hence
etale near these points, by the Artinian local criterion and the
finite-presentation etaleness criterion
\citep[Tags~02HT, 02HM and~00UR, complete written proofs]
{StacksUnitJetTransferD}. In particular its image contains an
open neighborhood of every point in the same unit class.
Any unit isomorphism on $X^\circ$ between represented objects
preserves their canonical completed lattices by the same local
functoriality, so it is the proper-curve isomorphism just used.
Thus each represented class has an open image in $Z$.
Distinct classes give disjoint open images: an isomorphism after
an algebraically closed coefficient extension would give a point
of a finite-type Isom scheme, hence already a $k$-point.
Two disjoint nonempty opens cannot meet the same irreducible
component. The number of classes is therefore at most the number
of irreducible components of $Z$.

The ordinary-weight specialization of
Lemma~\ref{D:weighted-affine-component-bound} bounds this number
by $(A_0fr^2+1)d_1^{A_0fr^2}$, since passing to an open locus
cannot increase the number of irreducible components. Enlarging
the constant gives (D.CL.4). The local identifications are essential
in this argument: bounded poles by themselves would not justify
a nilpotence-independent extension of a unit isomorphism.
\end{proof}

\begin{corollary}[Proved: fixed full local projective types have quadratic growth]
\label{D:curve-fixed-local-projective-growth}
Under the hypotheses on $(X,D,b)$ above, fix full local conjugacy
classes $\sigma_x$ with values in $\operatorname{PGL}_r(\mathbb F_q)$,
all with highest upper break at most $b$. The number of global
projective conjugacy classes with these types is at most $q^{C'r^2}$,
with $C'$ depending only on $(X,D,b,p)$. More generally, allowing
at most $q^{a_xr^2}$ types at the point $x$ gives at most
\[
                            q^{(C'+\sum_x a_x)r^2}
 \tag{D.CL.10}
\]
global projective classes.
\end{corollary}

\begin{proof}
Put $n=q-1$. Since $D$ is nonempty, $X^\circ$ is a smooth affine
curve. Its etale $H^2(X^\circ,\mu_n)$ vanishes
\citep[Lemma~59.69.3, Tag~03RR, complete proof]{StacksCurveFixedLocalD}.
The already established degree-two injection from fundamental-group
cohomology to etale cohomology therefore gives
$H^2(\pi_1(X^\circ),\mathbb F_q^*)=0$; see also the precise
Cartan--Leray comparison previously used with
\citep[Theorem~14.9 and Remark~14.10]{MilneLEC2013H2D}.
The obstruction to lifting a continuous projective representation
through $\operatorname{GL}_r(\mathbb F_q)$ is consequently zero.
Thus every such representation has a continuous linear lift over
the same finite field. For $q=2$ the scalar kernel is trivial.

If the fibre being counted is empty there is nothing to prove.
Otherwise one global projective representation and its linear lift
just constructed supply a local linear lift of each $\sigma_x$.
After using a lifted projective conjugating matrix, any other
linear lift differs from it by a continuous scalar character into
$\mathbb F_q^*$. Kummer theory gives exactly $n$ such characters:
$k((t))^*/k((t))^{*n}\simeq\mathbb Z/n$, since the residue field
is algebraically closed and the units are $n$-divisible. All are
tame. Further, the prime-to-$p$ central scalar kernel adds no wild
upper group: for positive upper indices the kernel lies in a
$p$-group, and upper numbering is compatible with quotients.
Every linear lift therefore still has highest break at most $b$.
There are at most $n^{\deg D}$ possible tuples of linear local
types. Apply the theorem to each tuple and map linear conjugacy
classes to projective conjugacy classes. The factor $n^{\deg D}$
is absorbed in $q^{C'r^2}$. Summing over the stated allowed local
types proves (D.CL.10). No injectivity on global Brauer groups,
fixed neutralizing cover, or bound on the scalar order is assumed.
\end{proof}


% Separate post-CP23 proved reduction. All four complete mathematical audits passed.
\subsection{Factorial cover counts from the degrees of local permutation orbits}

\begin{lemma}[Proved: permutation types with bounded local orbit square-sum]
\label{D:local-permutation-orbit-energy}
Let $\Lambda$ be a profinite group and suppose that, for some prime $p$
and real $a\ge0$, every $d\ge1$ satisfies
\[
 \#\bigl(\operatorname{Hom}_{\mathrm{cont}}
       (\Lambda,\operatorname{GL}_d(\mathbb F_p))
           /\operatorname{GL}_d(\mathbb F_p)\bigr)
                       \le p^{a d^2}.                    \tag{D.OE.1}
\]
For a finite continuous $\Lambda$-set $X$, write $d_1,\ldots,d_v$
for the cardinalities of its orbits and set
$\mathcal E_\Lambda(X)=\sum_jd_j^2$.
The number of isomorphism classes of $\Lambda$-sets of cardinality $n$
with $\mathcal E_\Lambda(X)\le E$ is at most
\[
                         2^{n-1}p^{(a+1)E}.             \tag{D.OE.2}
\]
In particular this number is $n^{O_{p,a,K}(n)}$ when
$E\le K n\log(n+1)$, where $\log$ denotes the natural logarithm.
\end{lemma}
\begin{proof}
The permutation matrices give an injective homomorphism
$S_d\hookrightarrow\operatorname{GL}_d(\mathbb F_p)$.
For one ambient linear-conjugacy class, its intersection with
$\operatorname{Hom}(\Lambda,S_d)$ has at most
$|\operatorname{GL}_d(\mathbb F_p)|\le p^{d^2}$ elements:
all the homomorphisms in that intersection are conjugates of any
one chosen representative by elements of the finite general linear
group. Passing to $S_d$-conjugacy classes cannot increase this
cardinality. Consequently
\[
 \#\bigl(\operatorname{Hom}_{\mathrm{cont}}(\Lambda,S_d)/S_d\bigr)
                           \le p^{(a+1)d^2}.             \tag{D.OE.3}
\]
We do not assert that a modular permutation module determines its
underlying $\Lambda$-set. The factor $p^{d^2}$ explicitly allows
for this possible failure.

Choose an ordering of the orbits of an $n$-element set. Their positive
cardinalities form a composition of $n$, and there are $2^{n-1}$
such compositions. For one composition satisfying $\sum_jd_j^2\le E$,
(D.OE.3), even without imposing transitivity on each selected factor,
bounds the number of choices by
$\prod_j p^{(a+1)d_j^2}\le p^{(a+1)E}$.
Every set under consideration occurs among these choices; their
ordering only overcounts. This proves (D.OE.2).
\end{proof}

\begin{theorem}[Proved: a factorial bound and an alternating simple-product quotient]
\label{D:finite-puncture-local-orbit-factorial-count}
Let $S$ be a normal connected proper surface over an algebraically
closed field $k$ of characteristic $p>0$, let $T\subset S$ be finite,
and put $U=S\setminus T$ and $G=\pi_1(U)$.
Choose once the detecting-curve data of
Lemma~\ref{D:finite-puncture-detecting-curve-local-types}.
Thus there are an algebraically closed extension $\Omega/k$, a smooth
proper connected curve $C/\Omega$, a finite subset $B\subset C$,
a morphism $C\setminus B\to U$, and a surjection
\[
                       \pi_1(C\setminus B)\twoheadrightarrow G.
                                                               \tag{D.OE.4}
\]
For each $b\in B$ fix the punctured completed normal surface group
$\Lambda_b$ and the factorization of its local curve group supplied
by that lemma:
\[
 \pi_1(\operatorname{Spec}\Omega((t_b)))
                      \longrightarrow\Lambda_b\longrightarrow G.
                                                               \tag{D.OE.5}
\]
All these choices are fixed independently of $n$ and of the cover.
For a finite continuous $G$-set $X$ define
\[
          \mathcal E(X)=\sum_{b\in B}
                 \sum_{O\in\Lambda_b\backslash X}|O|^2.
                                                               \tag{D.OE.6}
\]
Here the orbits are for the \emph{whole completed-local group}
$\Lambda_b$, before restricting to the curve inertia group in
(D.OE.5).

For every real $K>0$ there is a constant $A=A(S,T,K)$, also allowed
to depend on the fixed detecting data, such that the number of
isomorphism classes of connected degree-$n$ finite etale covers of
$U$ whose associated $G$-set satisfies
\[
                       \mathcal E(X)\le K n\log(n+1)           \tag{D.OE.7}
\]
is at most $n^{A n}$ for every $n\ge2$.
In particular, the same conclusion holds if every $\Lambda_b$-orbit
has cardinality at most $K_0\log(n+1)$ for one fixed $K_0$.

For $n\ge5$, let $\mathcal K_n$ consist of the distinct kernels of
continuous epimorphisms $G\twoheadrightarrow A_n$ for which the
natural $n$-letter $G$-set satisfies (D.OE.7). Then there is one
constant $c$ such that
\[
                   |\mathcal K_n|\le |A_n|^c\quad(n\ge5),
       \qquad
          G\Big/\bigcap_{n\ge5}\bigcap_{N\in\mathcal K_n}N
             \text{ is topologically finitely generated}.    \tag{D.OE.8}
\]
Only these selected simple quotient kernels are asserted to occur
in the displayed product. No bound on unrestricted alternating
quotients of a completed local surface group is asserted.
\end{theorem}
\begin{proof}
Put $w=|B|$ and $h=g(C)$.
By Theorem~\ref{D:uniform-local-representation-growth}, applied to
each of the finitely many fixed completed normal surface rings,
there are constants $a_b\ge0$ for which (D.OE.1) holds for
$\Lambda_b$ in every dimension $d$. Put
$a_*=\max(\{a_b:b\in B\}\cup\{0\})$.
For a tuple of compositions of $n$, one at each $b$, the argument of
Lemma~\ref{D:local-permutation-orbit-energy} shows that the number
of tuples of local $\Lambda_b$-sets with total energy at most
$K n\log(n+1)$ is bounded by
\[
             2^{w(n-1)}
                   p^{(a_*+1)K n\log(n+1)}.             \tag{D.OE.9}
\]
We do not count the smaller curve-inertia orbits as if they were
independent $\Lambda_b$-representations.

Fix one of the local tuples counted in (D.OE.9). Restriction along
(D.OE.5) specifies an isomorphism class of degree-$n$ finite etale
$\Omega((t_b))$-algebra at each puncture. Its integral closure over
$\Omega[[t_b]]$ is therefore specified as a full formal local
algebra, including all its factors. An isomorphism of the original
$\Lambda_b$-sets restricts to an isomorphism of these local algebras;
no labelling of their sheets is needed.

The detecting-curve lemma and
Proposition~\ref{D:surface-curvewise-conductor} give numbers
$\beta_b\ge0$, independent of the original cover, such that the
highest upper ramification break of its local Galois closure is
at most $\beta_b$. Hence, putting
$\lambda=\sum_{b\in B}(\beta_b+1)$, the sum $D$ of the discriminant
exponents of the degree-$n$ local algebras satisfies
\[
                               D\le\lambda n.             \tag{D.OE.10}
\]
Here is the precise passage from a Galois closure to a permutation
cover. For a totally ramified local Galois extension $L/K_b$ of
degree $e$, Lemma~\ref{D:different-Frobenius-invariant-lattice}
gives $\delta_{L/K_b}/e\le\beta_b+1$.
For an intermediate field $F/K_b$ of degree $d$, transitivity of
the different gives
$\delta_{L/K_b}=\delta_{L/F}+e(L/F)\delta_{F/K_b}$, so
$\delta_{F/K_b}\le d(\beta_b+1)$.
The residue field is algebraically closed, and consequently the
discriminant exponent of $F/K_b$ is $\delta_{F/K_b}$.
Summing over all factors proves (D.OE.10). Equivalently this is
the conductor bound for the permutation representation.

Every connected original cover stays connected on $C\setminus B$
by (D.OE.4). Its normalization over $C$ is a smooth proper connected
curve, since a normal curve over the perfect field $\Omega$ is
smooth. Moreover the restriction on isomorphism classes is
injective: the functor on finite continuous sets associated to a
surjective homomorphism of profinite groups is fully faithful.
We may therefore apply
Proposition~\ref{D:fixed-formal-curve-fibre-bound} to the tuple of
formal algebras just fixed. It bounds the number of possible
original covers by
\[
                             (n!)^{2h}n^{2D}
                             \le n^{(2h+2\lambda)n}.       \tag{D.OE.11}
\]
This uses the proved fixed-formal-extension fibre theorem; it does
not assume a general bound on affine-curve covers from their genus
or branch set alone.

Multiplying (D.OE.9) and (D.OE.11), and using
$\log(n+1)\le2\log n$ for $n\ge2$, proves the first assertion, for
example with
\[
        A=w+2(a_*+1)K\log p+2h+2\lambda.                 \tag{D.OE.12}
\]
If $w=0$ the local tuple is empty and the same argument is the
ordinary proper-curve bound. If every local orbit has cardinality
at most $K_0\log(n+1)$, then for each $b$ its square-sum is at most
$K_0 n\log(n+1)$; thus (D.OE.7) holds with $K=wK_0$.

Finally a surjection to $A_n$ gives a transitive $n$-letter $G$-set
whose kernel is precisely the kernel of that surjection. Distinct
selected kernels therefore give distinct isomorphism classes of
these $G$-sets, even if different representations with one kernel
can give different sets. It follows that
$|\mathcal K_n|\le n^{An}$.
For $n\ge5$ the elementary estimate
\[
                           |A_n|=n!/2\ge n^{n/8}           \tag{D.OE.13}
\]
follows from $n!\ge(n/2)^{n/2}$ and
$n\log n\ge8\log2$. Thus one can take $c=\lceil8A\rceil$.

The already used simple-product argument applies with this
polynomial multiplicity bound. Explicitly,
\citet[Proposition~2.6(ii), equations (2.1)--(2.2), pp.~12--13]
{LucchiniMarionTracey2020D} gives a common generator bound
$d=\lceil c\rceil+3$ for all the powers $A_n^{|\mathcal K_n|}$.
Distinct simple quotient kernels have full subdirect product,
by Goursat's lemma; distinct alternating degrees give distinct
simple isomorphism types. Hence the quotient in (D.OE.8) is the
full profinite product of these selected simple factors.
The uniform finite-product generator bound gives a generating
$d$-tuple for the full product by compactness, exactly as in
Corollary~\ref{D:Brauer-zero-selected-simple-product}.
This proves (D.OE.8).
\end{proof}

\begin{remark}[Exact scope of the reduction]
The local condition concerns completed surface punctures, not merely
DVR inertia and not a hypothetical decomposition of a modular
permutation module. It permits wild local actions; it makes no
claim that such actions are automatically tame or that all their
orbit degrees are small. The unrestricted alternating gate remains
precisely where local permutation orbit square-sums can be of order
$n^2$. Neither the quadratic ambient representation theorem nor
etale-algebra rigidity alone removes that remaining loss.
\end{remark}


% Proved: four complete actual-file audits; one explicit known-input citation added.
\subsection{Unrestricted projective counts at global finite punctures}

\begin{theorem}[Proved: global finite-puncture projective growth]
\label{D:global-finite-puncture-projective-growth}
Let $S$ be a normal proper connected surface over an algebraically
closed field $k$ of characteristic $p>0$, let $T\subset S$ be finite,
and put $\Gamma=\pi_1(S\setminus T)$. There is a constant $C_{S,T}$
such that, for every $q=p^f$ and every $r\geq1$,
\[
 \#\bigl(\operatorname{Hom}_{\rm cont}
       (\Gamma,\operatorname{PGL}_r(\mathbb F_q))
                 /\operatorname{PGL}_r(\mathbb F_q)\bigr)
                          \leq q^{C_{S,T}r^2}.           \tag{D.GFP.1}
\]
Neither the scalar-centre order nor the geometric Brauer period is
bounded. The surface need not be projective. The constant is for the
fixed pair; no bound linear in the degree of an arbitrary finite
cover is asserted.
\end{theorem}

\begin{proof}
Choose once the curve $C^\circ=C\setminus P$ over $\Omega$,
the surjection $\pi_1(C^\circ)\twoheadrightarrow\Gamma$, and the
fixed completed local groups $\Lambda_s$ of
Lemma~\ref{D:finite-puncture-detecting-curve-local-types}.
The marked set is nonempty. Write $s(x)$ for the local source of
the formal branch at $x\in P$ and let $b=\max_x b_x$ be the
fixed upper-break bound in that lemma.

For each fixed complete normal surface ring $R_s$, choose the
constant $a_s$ in
Theorem~\ref{D:SNC-general-projective-quadratic-count}, equivalently
Corollary~\ref{D:complete-local-projective-growth-all-dimensions}.
It bounds its projective representation classes in every rank by
$q^{a_s r^2}$. A homomorphism
$\Gamma\to\operatorname{PGL}_r(\mathbb F_q)$ restricts first to
$\Lambda_{s(x)}$ and then to the formal curve group at $x$.
Conjugate representations remain conjugate after restriction.
Consequently at most $q^{a_{s(x)}r^2}$ full local projective types
can occur at that point. We may count independently at different
points even when their local source is the same: this merely
overcounts possible tuples. Only types actually realized by an
original representation need be retained; all of them have
highest upper break at most $b$.

Corollary~\ref{D:curve-fixed-local-projective-growth} therefore
bounds the possible projective classes on $C^\circ$ by
\[
                    q^{(C_C+\sum_{x\in P}a_{s(x)})r^2}.
                                                               \tag{D.GFP.2}
\]
Precomposition with the fixed surjection
$\pi_1(C^\circ)\twoheadrightarrow\Gamma$ is injective on
conjugacy classes: a conjugacy after precomposition is the same
equality on the whole target group $\Gamma$. This proves (D.GFP.1).
The curve, local rings, upper-break constants and $C_C$ were all
chosen independently of $q,r$ and the representation.
We have counted the actual projective representations through
restriction and same-field linear lifting on an affine curve.
No vanishing or injectivity assertion for the global Brauer group
of $S\setminus T$ has been used.
\end{proof}

\begin{corollary}[Proved: global dimensional and finite-normalization closure]
\label{D:global-projective-finite-puncture-closure}
Let $S$ be proper over an algebraically closed field of characteristic
$p>0$, and let $T\subset S$ be closed with codimension at least two
in every irreducible component that it meets. Assume either
\begin{enumerate}
\item $T$ is finite; or
\item every normalized component of dimension at least two admits
the projective modification with controlled exceptional centres
in Theorem~\ref{D:generic-surface-controlled-boundary}, for the
inverse image of $T$.
\end{enumerate}
For every connected component $U$ of $S\setminus T$ there is a
constant $C_U$ such that
\[
 \#\bigl(\operatorname{Hom}_{\rm cont}
      (\pi_1(U),\operatorname{PGL}_r(\mathbb F_q))
                 /\operatorname{PGL}_r(\mathbb F_q)\bigr)
                            \leq q^{C_Ur^2}             \tag{D.GFP.3}
\]
for all $q=p^f$ and $r\geq1$. Normality and reducedness of $S$
are not required. No Brauer-period or bounded-scalar condition is
imposed in either case.
\end{corollary}

\begin{proof}
First take a normal proper integral component of dimension at
least two. In case~(1),
Theorem~\ref{D:generic-surface-finite-punctures} gives one fixed
normal proper surface over an algebraically closed extension,
with a finite omitted set, whose punctured fundamental group
surjects onto the original group. In case~(2), use the
codimension-at-least-two clause (D.CB.2) of
Theorem~\ref{D:generic-surface-controlled-boundary}; its detecting
surface also has only finitely many punctures. Apply
Theorem~\ref{D:global-finite-puncture-projective-growth} to that
one surface. Surjective precomposition proves (D.GFP.3) for the
normal component. This step does not require a uniform constant
over all field extensions or all surfaces: the detecting surface
and its ground field are fixed once for the original component.

For a general $S$, reduction does not change finite etale covers.
Take the finite normalization of the reduced irreducible
components and restrict it over $U$:
\[
                         V=\coprod_i V_i\longrightarrow U.
                                                               \tag{D.GFP.4}
\]
This is a finite proper surjection. The codimension assumption
persists on the normalized components, and in case~(1) the
inverse image of $T$ is finite. In case~(2) the stipulated
modifications give exactly the normal-component conclusion
of the preceding paragraph. Components of dimension at most
one meet no omitted point, by the codimension assumption.
They are proper normal curves or points; their fundamental groups
are finitely generated
\citep[Expose~X, Theorem~2.6, pp.~272--273]{C_SGA1}.
If $d_i$ generators generate those groups, their representation counts
are bounded by
$|\operatorname{PGL}_r(\mathbb F_q)|^{d_i}\leq q^{d_i r^2}$.

Set $a=\#\pi_0(V\times_U V)$, a fixed finite integer, and fix
one tuple of torsor classes on the $V_i$. Choose representatives.
Their two pullbacks to $V\times_U V$ have a finite etale
equivariant Isom scheme of degree
$|P|$, where $P=\operatorname{PGL}_r(\mathbb F_q)$.
A section on a connected component is determined by its value
at one geometric point, so there are at most $|P|^a$ possible
descent isomorphisms. The cocycle condition and passage to
isomorphism classes only decrease this number. Finite proper
descent of finite etale covers, with their action morphisms,
is effective and fully faithful
\citep[Expose~IX, Theorem~4.7]{C_SGA1}; this is the same descent
calculation used in
Corollary~\ref{D:Brauer-zero-normalization-transfer}.
Multiplying the fixed component bounds by
$|P|^a\leq q^{ar^2}$ proves (D.GFP.3). No cohomological
degree-two inflation or Brauer-descent injection is asserted.
\end{proof}

\begin{corollary}[Proved: all defining-characteristic Lie simple products in this scope]
\label{D:global-unrestricted-projective-simple-products}
Let $\Gamma$ be any group covered by (D.GFP.1) or (D.GFP.3).
Suppose a family $\mathscr S$ of pairwise nonisomorphic nonabelian
finite simple groups admits embeddings
\[
 S\hookrightarrow\operatorname{PGL}_{r_S}(\mathbb F_{q_S}),
 \qquad q_S=p^{f_S},\qquad
 r_S^2\log q_S\leq\kappa\log|S|,                       \tag{D.GFP.5}
\]
with one constant $\kappa$. Then the quotient of $\Gamma$ by
the intersection of all kernels of epimorphisms onto groups in
$\mathscr S$ is topologically finitely generated.
This applies to all defining-characteristic finite simple Lie
families, including unbounded natural scalar centres and all
classical, exceptional, Suzuki and Ree families. The finitely
many simple derived exceptions can be included as well.
\end{corollary}

\begin{proof}
Let $C$ be the constant in the applicable projective count.
Distinct quotient kernels with target $S$ give distinct ambient
projective conjugacy classes after the chosen embedding: an
ambient conjugacy preserves the kernel. Thus their multiplicity
$m_S$ satisfies
\[
                         m_S\leq q_S^{Cr_S^2}
                                      \leq |S|^{C\kappa}.
                                                               \tag{D.GFP.6}
\]
The exact multiplicity-to-generation estimate already used in
Corollary~\ref{D:Brauer-zero-selected-simple-product}, namely
\citet[Proposition~2.6(ii), equations~(2.1)--(2.2), pp.~12--13]
{LucchiniMarionTracey2020D}, gives a common generator bound for
the powers $S^{m_S}$. Goursat's lemma gives the full subdirect
product for distinct simple quotient kernels, and the different
simple isomorphism types are independent. Compactness therefore
gives a generating tuple for the full profinite product with
the same uniform bound.

The faithful projective realizations and size comparisons for
the named Lie families were checked in
Corollaries~\ref{D:classical-defining-characteristic-crowns} and
\ref{D:exceptional-defining-characteristic-crowns}.
We use only their embeddings of the simple group and their
order comparisons, with constants $16$ and $3\cdot248^2$.
Their one-curve and radially tame outer-quotient hypotheses
were needed for the crown conclusions there, and are not
hypotheses of these purely finite-group embeddings.
The finite derived exceptions have faithful projective
permutation realizations over $\mathbb F_p$ and are absorbed
by increasing $\kappa$ over a fixed finite list.
The present statement is a simple-product conclusion.
It does not bound arbitrary nonabelian crowns, alternating
multiplicities, or cross-characteristic families failing
(D.GFP.5), and does not prove full finite generation of $\Gamma$.
\end{proof}

\begin{remark}[Movement and geometric scope]
These three statements close one global projective-count gate:
finite-puncture proper surfaces and their stated dimensional and
finite-normalization consequences now allow unrestricted geometric
Brauer period. The curve presentation, local restriction lemma,
transfers and named simple-family consequence belong to that
single reduction. Original divisorial boundaries, arbitrary
positive-dimensional omitted loci outside the stated modification
hypothesis, and constants linear in arbitrary cover degree remain
separate proof problems. General D.1, D.2(ii) and residual L8
are not settled by these statements.
\end{remark}


% Proved enlargement of the same alternating gate reduction; four complete audits.
\subsection{Removing controlled local orbits from the square-sum condition}

\begin{corollary}[Proved: fixed local quotients and arbitrary nilpotent-image orbits]
\label{D:fixed-local-quotient-alternating-count}
Keep the surface, detecting curve, and fixed completed-local groups
$\Lambda_b$ of Theorem~\ref{D:finite-puncture-local-orbit-factorial-count}.
For each marked point $b$ choose finitely many fixed continuous
quotients
\[
             \Lambda_b\twoheadrightarrow Q_{b,\nu}
                          \quad(1\le\nu\le J_b)
                                                               \tag{D.LQ.1}
\]
that are topologically finitely generated. For a finite $G$-set $X$,
call a $\Lambda_b$-orbit controlled if its permutation action factors
through at least one of (D.LQ.1). Define
\[
 \mathcal E_{\rm rem}(X)=
       \sum_b\sum_{\substack{O\in\Lambda_b\backslash X\\
                              O\text{ not controlled}}}|O|^2.
                                                               \tag{D.LQ.2}
\]
For fixed $K>0$, connected degree-$n$ covers with
$\mathcal E_{\rm rem}(X)\le K n\log(n+1)$ satisfy a bound
$n^{A'n}$, where $A'$ depends only on the fixed geometric data,
$K$, and the chosen quotients. The distinct natural $A_n$ quotient
kernels satisfying this condition have multiplicity at most
$|A_n|^{c'}$ for one constant $c'$, and their selected
simple-product quotient is topologically finitely generated.

One may always take $Q_{b,1}$ to be the maximal pronilpotent
quotient of $\Lambda_b$. Consequently, the conclusions hold when
(D.LQ.2) sums only over orbits whose whole completed-local
permutation image is \emph{not nilpotent}. There is no restriction
on the sizes of the nilpotent-image orbits. In particular these
can have arbitrarily large $p$-group image. This assertion does
not replace ``nilpotent'' by ``solvable''.
\end{corollary}

\begin{proof}
The pronilpotent finite-generation assertion is known:
\citet[Theorem~1.2]{Suzuki2023} proves it for connected punctured
spectra of complete noetherian local rings with algebraically
closed residue field and all irreducible components of dimension
at least two. We record the short derivation in the present
surface case to identify a fixed generator bound from the exact
cohomology inputs already proved. Let
$\Lambda=\pi_1(\operatorname{Spec}R\setminus\{\mathfrak m\})$
for one of the fixed completed normal surface rings.
Theorem~\ref{D:local-p-linear-H1}, with the constant rank-one
sheaf, bounds $\dim_{\mathbb F_p}H^1(\Lambda,\mathbb F_p)$.
For primes $\ell\ne p$, use the fixed SNC resolution $Y$ with
exceptional components $E_1,\ldots,E_v$ and $e$ nodes from the
proof of Theorem~\ref{D:local-surface-constant-H2}.
Writing $j$ for the puncture inclusion, the same purity calculation
in (D.LH2.2) gives
\[
 R^1j_*\mathbb F_\ell
                   =\bigoplus_{i=1}^v\mathbb F_{\ell,E_i}(-1).
                                                               \tag{D.LQ.3}
\]
This is exactly the regular normal-crossings corollary in
\citet[Expos\'e XVI, Theorem~3.1.1 and Corollary~3.1.3,
pp.~242--243]{A_ILO2014}; it applies to the excellent resolution,
without a finite-type-over-a-field hypothesis. Its proof was used
in the cited local cohomology theorem.
Proper strict-local base change identifies
$H^1(Y,\mathbb F_\ell)$ with $H^1(E_{\rm red},\mathbb F_\ell)$.
The normalization/node exact sequence bounds its dimension by
$2\sum_i g(E_i)+e$; allowing the extra $e$ is harmless even when
the dual graph has fewer independent cycles. The degree-one
edge sequence for $j$, together with (D.LQ.3), therefore gives
\[
       \dim_{\mathbb F_\ell}H^1(\Lambda,\mathbb F_\ell)
                  \le 2\sum_i g(E_i)+e+v\quad(\ell\ne p).
                                                               \tag{D.LQ.4}
\]
The degree-one etale/group comparison is the classification of
finite constant torsors. Tate twists in (D.LQ.3) are constant over
the algebraically closed residue field. Thus there is an integer
$d_R$ bounding these degree-one dimensions for \emph{all} primes.
No bound for arbitrary nonconstant prime-to-$p$ systems is asserted
or needed.

For clarity, this implies that each maximal pro-$\ell$ quotient
$\Lambda^{(\ell)}$ is generated by at most $d_R$ elements.
Its elementary abelian quotient dual to
$H^1(\Lambda,\mathbb F_\ell)$ is finite; choose lifts of a basis.
If their generated closed subgroup were proper, a finite
$\ell$-group quotient would still have proper generated image.
A maximal subgroup of that finite $\ell$-group has index $\ell$,
and hence supplies a nonzero homomorphism to $\mathbb F_\ell$
vanishing on the chosen lifts, contrary to the basis choice.
This is the usual finite-group Frattini argument, applied to every
finite quotient; it assumes no prior finite generation.

The maximal pronilpotent quotient is
\[
                        \Lambda^{\rm nil}
                            =\prod_{\ell}\Lambda^{(\ell)}.
                                                               \tag{D.LQ.5}
\]
Indeed finite nilpotent groups are direct products of their Sylow
subgroups, so every homomorphism into such a group factors through
the displayed product. The natural image is the whole product:
its image on every finite collection of finite Sylow quotients is
surjective by the coprime-order subdirect-product argument, and
its image is compact. Choose a generating $d_R$-tuple in each
factor, padding it with identities. Combining these tuples
coordinatewise generates every finite subproduct, again by
coprime orders, and hence generates (D.LQ.5) topologically.
This proves the claimed fixed finitely generated quotient.

We now prove the counting assertion for any choices (D.LQ.1).
Let $d_{b,\nu}$ be a number of topological generators of
$Q_{b,\nu}$ and set
\[
             J=\sum_bJ_b,\qquad
             d_0=\sum_{b,\nu}d_{b,\nu},\qquad w=|B|.
                                                               \tag{D.LQ.6}
\]
At each $b$, assign every controlled orbit to one quotient through
which it factors. Join all the orbits assigned to the same quotient
into one set; its action still factors through that quotient.
For an $m$-element set, the number of such actions up to isomorphism
is at most $(m!)^{d_{b,\nu}}$, since the images of a fixed generating
tuple determine a continuous homomorphism to $S_m$.
No nilpotency restriction need be imposed when obtaining this
upper bound. The sizes of these joined sets have at most
$(n+1)^{J_b}$ possibilities. The unassigned part has an ordered
orbit composition, with at most $2^n$ possibilities.
By Lemma~\ref{D:local-permutation-orbit-energy}, its actual
permutation types contribute at most
$p^{(a_*+1)\mathcal E_b}$, where $a_*$ is the fixed maximum of
the local linear-count constants and $\mathcal E_b$ is its
remaining square-sum. Summing over all choices therefore bounds
the number of tuples of full local types by
\[
          (n+1)^J2^{wn}(n!)^{d_0}
                 p^{(a_*+1)K n\log(n+1)}.                \tag{D.LQ.7}
\]
Assignment choices overcount and require no canonical selection.

The detecting curve and the full local algebras are the same as
in Theorem~\ref{D:finite-puncture-local-orbit-factorial-count}.
The different bound is still $D\le\lambda n$, independent of
whether an orbit is controlled. Each fixed tuple consequently
has at most $(n!)^{2h}n^{2\lambda n}$ global covers by the
fixed-formal fibre theorem. Multiplying by (D.LQ.7) gives
$n^{A'n}$ for $n\ge2$, for instance with
\[
 A'=J+w+d_0+2(a_*+1)K\log p+2h+2\lambda.                \tag{D.LQ.8}
\]
Here $(n+1)^J\le n^{Jn}$, $2^{wn}\le n^{wn}$ and
$\log(n+1)\le2\log n$. The same selected-kernel argument,
$|A_n|\ge n^{n/8}$, and simple-product generation proof used in
the cited theorem give the remaining conclusions, with
$c'=\lceil8A'\rceil$.

Finally, a finite permutation action factors through
$\Lambda^{\rm nil}$ if and only if its image is nilpotent:
one direction follows from the definition of pronilpotent, and
the other from (D.LQ.5). This proves the concrete enlargement.
It uses no finite-generation assertion for the maximal
prosolvable quotient of a complete local puncture.
\end{proof}


% Separate post-CP23 proved local-pair input. All four complete audits passed.
\subsection{Central cohomology for a complete local pair}

\begin{lemma}[Proved: constant second cohomology for an original divisorial tame quotient]
\label{D:local-divisor-tame-central-H2}
Let $(R,\mathfrak m,k)$ be a complete excellent normal local
$k$-algebra of dimension two, where $k$ is algebraically closed of
characteristic $p>0$. Let $D$ be a nonempty reduced divisor on
$\operatorname{Spec}R$, put $U=\operatorname{Spec}R\setminus D$,
and write $G=\pi_1(U)$. Let $Q$ be the quotient of $G$ classifying
covers tame at the generic points of the original divisor $D$.
Choose a proper log resolution $\pi:Y\to\operatorname{Spec}R$
which is an isomorphism over $U$, with nonempty exceptional fibre
$E=(\pi^{-1}\mathfrak m)_{\rm red}$, and with
$E+D'$ a strict normal-crossings divisor; $D'$ is the strict
transform of $D$. Write $v$ for the number of components $E_i$,
$g_i$ for their genera, $e$ for the number of exceptional nodes,
and $h$ for the number of points of $E\cap D'$.
For every $n\ge1$ prime to $p$, with trivial action on the
coefficients, there are injections and a bound
\[
 H^2_{\rm cont}(Q,\mu_n(k))\lhook\joinrel\longrightarrow
 H^2_{\rm cont}(G,\mu_n(k))\lhook\joinrel\longrightarrow
 H^2_{\rm et}(U,\mu_n),\qquad
 |H^2_{\rm et}(U,\mu_n)|\le n^{v+2\sum_i g_i+e+h}.
 \tag{D.LDH.1}
\]
In particular the number of constant central lifting classes is
bounded by a fixed power of the scalar order. No assertion that
the kernel $G\to Q$ is itself pro-$p$ is made.
\end{lemma}

\begin{proof}
Blowing up further if needed makes $E$ nonempty and the total
boundary strict normal crossings. These operations are over the
closed point and preserve $U$. Since $R$ is normal, its puncture
is regular, and the resolution and embedded-resolution inputs
are the same as those of
Theorem~\ref{D:local-surface-constant-H2}.
Each $E_i$ is a smooth proper connected curve over $k$.
The horizontal components $D'_a$ are regular and proper over
$\operatorname{Spec}R$; their closed fibres are finite sets.
There are no horizontal-horizontal intersections: any such point
would lie over $\mathfrak m$, and hence on $E$, contradicting the
absence of triple intersections. All boundary nodes are therefore
exactly the $e+h$ points listed in the statement.

Let $j:U\hookrightarrow Y$ and put $\Lambda=\mathbf Z/n$.
The normal-crossings form of absolute purity gives
\[
 \begin{aligned}
 R^0j_*\Lambda&=\Lambda_Y,\\
 R^1j_*\Lambda&=
   \bigoplus_i\Lambda_{E_i}(-1)
       \oplus\bigoplus_a\Lambda_{D'_a}(-1),\\
 R^2j_*\Lambda&=\bigoplus_{x\in\operatorname{Nodes}(E+D')}
                                      \Lambda_x(-2),
 \qquad R^qj_*\Lambda=0\quad(q>2).
 \end{aligned}
 \tag{D.LDH.2}
\]
The exact input is \citet[Expose~XVI, Theorem~3.1.1 and
Corollary~3.1.3, pp.~242--243]{A_ILO2014}; the corollary's full
proof identifies the derived direct image with the exterior
algebra on its first direct image. Its scope is regular
$\mathbf Z[1/n]$-schemes, so no finite-type-over-$k$ assumption
on the complete local base is needed.

The base is strictly henselian. Sections of an etale sheaf on
that base equal its closed geometric stalk and are exact.
Proper base change consequently identifies the cohomology of
$Y$, and of each proper horizontal component, with the
cohomology of its closed fibre
\citep[Theorem~59.91.11, Tag~095T, complete written proof]
{StacksSemilinearDeterminantsB}.
Nilpotents do not affect etale cohomology. Thus
\[
 |H^2(Y,\Lambda)|=n^v,
 \quad |H^1(E_i,\Lambda)|=n^{2g_i},
 \quad H^1(D'_a,\Lambda)=0.
 \tag{D.LDH.3}
\]
For the first equality use the normalization/node exact sequence
of the proper nodal curve $E$; its node terms have no positive
cohomology. The smooth-component calculation is
\citep[Lemma~59.69.1, Tag~03RQ]{StacksLocalPairCohomologyB}.
The Tate twists are constant because all prime-to-$p$ roots of
unity lie in $k$ and lift by Hensel's lemma.

The three terms of total degree two in
$H^a(Y,R^qj_*\Lambda)\Rightarrow H^{a+q}(U,\Lambda)$ have
orders $n^v$, $n^{2\sum g_i}$ and $n^{e+h}$. Their limiting
terms are subquotients, which proves the cardinality bound.
An identification $\Lambda\simeq\mu_n$ gives the stated form.

For the second injection in (D.LDH.1), use the Cartan--Leray
spectral sequence for the directed tower of connected finite
Galois covers of $U$. Every finite constant degree-one torsor
on one cover is killed by a further finite etale cover.
The degree-one term in that tower is therefore zero, so the
edge map in degree two is injective. This is precisely the
comparison explained in the proof of
Theorem~\ref{D:local-surface-constant-H2}; see also
\citep[Theorem~14.9 and Remark~14.10, complete written passage]
{MilneLEC2013H2D}.

Finally let $N=\ker(G\to Q)$. It is topologically normally
generated by the images of the pro-$p$ wild inertia groups
at the original boundary valuations. Every continuous
homomorphism $N\to\mu_n(k)$ is zero on each conjugate of
these groups, hence is zero. The five-term sequence gives
injectivity of inflation. Equivalently, if a central extension
of $Q$ by $\mu_n(k)$ splits after pullback to $G$, the splitting
homomorphism is trivial on $N$ by the preceding argument and
therefore descends to $Q$. This proves the first injection
without replacing the normally generated kernel by a pro-$p$
group.
\end{proof}

\begin{lemma}[Proved: the horizontal residue order is its single node order]
\label{D:horizontal-residue-order-at-node}
In the preceding setting, each irreducible horizontal component
$D'_a$ is the spectrum of a complete discrete valuation ring
isomorphic to $k[[z_a]]$. It meets $E$ at one transverse point
$x_a$. For $n$ prime to $p$, the secondary-residue map at that
point is, after fixing the usual Tate-twist identification,
an isomorphism
\[
 H^1(k(D'_a),\mathbf Z/n)\xrightarrow{\ \partial_{x_a}\ }
                                \mathbf Z/n.
 \tag{D.LDH.4}
\]
In particular every horizontal residue character has the same
order as its secondary residue at $x_a$.
\end{lemma}

\begin{proof}
The map from $D'_a$ to the original irreducible branch is
proper and birational. Its fibres are finite: away from the
closed point this is automatic, and over the closed point
$D'_a\cap E$ is a finite set. Hence it is finite
\citep[Lemma~37.44.1, Tag~02LS]{StacksLocalPairCohomologyB}.
Its coordinate ring is a finite integral domain over a complete
local ring. Henselian decomposition makes it local, and
finiteness makes it complete. Regularity and dimension one
make it a complete discrete valuation ring. Its residue field
is $k$, and the given coefficient field yields $k[[z_a]]$.
There is consequently exactly one point above the closed point;
normal crossings makes the intersection transverse.

Kummer theory identifies $H^1(k((z_a)),\mu_n)$ with
$k((z_a))^*/k((z_a))^{*n}$. Every unit of $k[[z_a]]$ has an
$n$th root by Hensel's lemma, since $n$ is invertible and $k$
is algebraically closed. Valuation therefore identifies this
quotient with $\mathbf Z/n$. The Kummer boundary at the
closed point is that valuation modulo $n$, which proves
(D.LDH.4). A different choice of compatible roots of unity or
orientation changes it by an automorphism and does not change
orders.
\end{proof}

\begin{remark}
These are cohomological and residue inputs for a local pair.
They do not provide an object atlas, a representation count,
a hereditary index estimate, or finite generation of the
whole original-boundary tame group.
\end{remark}


% Proved: closure of the same restricted alternating gate; four complete audits.
\subsection{Dimensional and finite-normalization closure of the alternating bound}

\begin{corollary}[Proved: alternating counts detected on fixed local surface groups]
\label{D:alternating-dimensional-normalization-closure}
Let $S$ be proper over an algebraically closed field of characteristic
$p>0$, and let $T\subset S$ be closed of codimension at least two
in every irreducible component that it meets. Assume either that
$T$ is finite, or that every normalized component of dimension at
least two has a projective modification satisfying the controlled
exceptional-centre hypothesis of
Theorem~\ref{D:generic-surface-controlled-boundary} for the inverse
image of $T$. Let $U$ be a connected component of $S\setminus T$,
and put $G=\pi_1(U)$.

There is a fixed finite collection of homomorphisms
\[
 \Lambda_{i,b}\longrightarrow G,\qquad
 \Lambda_{i,b}=\pi_1\bigl(\operatorname{Spec}R_{i,b}
                                      \setminus\{\mathfrak m_{i,b}\}\bigr),
                                                               \tag{D.ANC.1}
\]
where the $R_{i,b}$ are complete normal local surface rings with
algebraically closed coefficient fields. Their precise construction
is given below; they are fixed before $n$ or a cover is chosen.
For each group in (D.ANC.1), choose a finite collection of fixed
topologically finitely generated quotients. Call one of its orbits
on a finite $G$-set controlled if the orbit action factors through
one of these quotients. Fix $K>0$.

The isomorphism classes of degree-$n$ finite etale covers $X\to U$
such that
\[
 \sum_{i,b}\ 
 \sum_{\substack{O\in\Lambda_{i,b}\backslash X\\
                         O\text{ not controlled}}}|O|^2
                              \le K n\log(n+1)            \tag{D.ANC.2}
\]
are at most $n^{C n}$ for a constant $C$ depending only on the
fixed data. The covers in this assertion may be disconnected.
Consequently the distinct natural $A_n$ quotient kernels satisfying
(D.ANC.2) have multiplicity at most $|A_n|^c$ for one constant $c$,
and the corresponding selected alternating simple-product quotient
of $G$ is topologically finitely generated.

One may jointly choose the maximal pronilpotent and maximal
pro-prime-to-$p$ quotients. Thus (D.ANC.2) may charge only orbits
whose \emph{whole fixed local surface-group image} is neither
nilpotent nor of order prime to $p$, without a size restriction
on the other orbits. No condition is imposed merely on the smaller
curve-inertia orbits. Normality and reducedness of $S$ are not
required.
\end{corollary}

\begin{proof}
Reduce $S$ and normalize its irreducible components. Over $U$ this
gives a fixed finite proper surjection
\[
                         V=\coprod_i V_i\longrightarrow U
                                                               \tag{D.ANC.3}
\]
with connected normal $V_i$. For a component of dimension at least
two, choose once a normal proper surface $Z_i$ over an algebraically
closed extension of $k$, with a finite set $F_i$, and a surjection
\[
        \pi_1(Z_i\setminus F_i)\twoheadrightarrow\pi_1(V_i).
                                                               \tag{D.ANC.4}
\]
These are supplied by
Theorem~\ref{D:generic-surface-finite-punctures} in the finite case,
and by the codimension-at-least-two clause of
Theorem~\ref{D:generic-surface-controlled-boundary} in the other
case. In dimension two the original normal component itself may
be used. Apply
Lemma~\ref{D:finite-puncture-detecting-curve-local-types} to each
fixed pair $(Z_i,F_i)$. Its fixed completed-local groups are the
$\Lambda_{i,b}$ in (D.ANC.1); compose their maps with (D.ANC.4)
and the map from $V_i$ to $U$. Choose geometric base points and
connecting paths once. Changing those paths conjugates the maps
and changes neither orbit sizes nor factorization through the
chosen quotients. All the completed local rings here are attached
to these fixed \emph{detecting surfaces}; no assertion that they
are completions of the original higher-dimensional $S$ is made.

We first extend the surface count to disconnected covers. Keep one
fixed detecting curve $C^\circ=C\setminus B$, with $w=|B|>0$ and
$h=g(C)$, for a surface $(Z_i,F_i)$. The local-type enumeration
in the proof of
Corollary~\ref{D:fixed-local-quotient-alternating-count} counted
arbitrary finite local sets, so it already bounds the tuples of
full rank-$n$ local algebras under (D.ANC.2) by $n^{A_0 n}$ for
a fixed $A_0$. The discriminant sum is still at most $\lambda n$:
the upper-break bound applies to each actual cover and the
different estimate is additive over its local factors.

For a fixed tuple of full local algebras, consider a cover of
$C^\circ$ with $k\le n$ connected components. Temporarily order
those components. At a marked point its completed algebra has
at most $n$ connected factors. Assigning those factors to the
$k$ labelled global components has at most $n^n$ possibilities.
Doing this at every marked point, and allowing the $n$ choices of
$k$, costs at most $n^{w n+1}$. Assignments with inconsistent
degrees or empty global components are discarded. Since $B$ is
nonempty, every actual connected component has positive degree
at each marked point and is included in this enumeration.

For a consistent assignment, let $d_j$ and $D_j$ be the degrees
and discriminant sums of its connected components. Then
$\sum_jd_j=n$ and $\sum_jD_j\le\lambda n$. The connected
fixed-formal-fibre bound in
Theorem~\ref{D:fixed-formal-curve-fibre-bound} gives at most
\[
 \prod_j(d_j!)^{2h}d_j^{2D_j}
                  \le (n!)^{2h}n^{2\lambda n}             \tag{D.ANC.5}
\]
such ordered tuples of global covers. Forgetting the ordering
can only lower the count. This proves a bound $n^{A_i n}$ for
all degree-$n$ covers of $Z_i\setminus F_i$ satisfying the local
condition, with $A_i$ fixed. The original detecting-curve
surjection is injective on finite-set isomorphism classes even
for nontransitive actions: an equivariant bijection after
precomposition is equivariant for the entire quotient group.
The same applies to (D.ANC.4), giving the bound for the relevant
degree-$n$ covers of $V_i$.

Components of dimension at most one meet no omitted point, by
the codimension hypothesis. They are proper normal curves or
points. Their fundamental groups have fixed finite generating
sets by \citet[Expose~X, Theorem~2.6, pp.~272--273]{C_SGA1}.
If $d_i$ elements suffice, all their degree-$n$ covers have at
most $(n!)^{d_i}\le n^{d_i n}$ isomorphism classes. They need
no groups (D.ANC.1) and no energy charge.

Fix one tuple of rank-$n$ covers on the $V_i$. On
$V\times_U V$, the finite etale Isom scheme between the two
pullbacks has degree $n!$. A section on a connected component
is determined by its value at one geometric point. Thus there
are at most $(n!)^a$ descent isomorphisms, where the fixed integer
$a=\#\pi_0(V\times_U V)$ is independent of $n$. The cocycle
condition and passage to isomorphism classes only decrease the
count. Finite proper descent of finite etale covers is effective
and fully faithful
\citep[Expose~IX, Theorem~4.7]{C_SGA1}. Multiplying the fixed
component bounds by $(n!)^a\le n^{a n}$ proves the asserted
$n^{Cn}$ bound. Rank one has only the trivial cover of connected
$U$, so the stated form also holds for $n=1$.

The selected alternating-kernel and simple-product conclusions
now follow exactly as in
Theorem~\ref{D:finite-puncture-local-orbit-factorial-count}:
distinct kernels give distinct natural permutation classes,
$|A_n|\ge n^{n/8}$ for $n\ge5$, and the uniform polynomial
multiplicity bound gives one generating bound by the cited
LMT estimate, Goursat's lemma and compactness.

For the concrete final choice, each ring in (D.ANC.1) has a
finitely generated maximal pronilpotent quotient by
Corollary~\ref{D:fixed-local-quotient-alternating-count}. The other
quotient is also a \emph{known} finitely generated group:
\citet[\S9.1 and Theorem~9.2, pp.~122--123]
{C_GrothendieckMurre1971PrimeToP} prove finite presentation of the
maximal prime-to-$p$ quotient for a connected open in the spectrum
of a complete noetherian local ring of dimension two with
algebraically closed residue field. The finite presentation is in
the pro-prime-to-$p$ category; its stated definition includes a
finite set of topological generators, which is all that is used
here. Every ring in (D.ANC.1) satisfies those hypotheses. A finite
action factors through this quotient if and only if its image
has order prime to $p$. Choosing these two fixed quotients
therefore controls precisely the union of the nilpotent-image
and prime-to-$p$-image orbits, as asserted.
\end{proof}

\begin{remark}
This is a closure of the same restricted alternating count.
The local homomorphisms and quotients are part of the fixed
detecting data. The conclusion does not discard unrestricted
nonsolvable local images, and does not prove an unrestricted
factorial cover bound or finite generation of the whole $G$.
\end{remark}


% Proved after four complete actual-file audits. Separate post-CP23 local-pair research.
\subsection{Projective counts with an original local divisorial tame boundary}

\begin{lemma}[Proved: fixed coefficient gerbes and finite horizontal data]
\label{D:local-pair-coefficient-root-comparison}
Let $(R,\mathfrak m,k)$ be a complete excellent normal local
$k$-algebra of dimension two, with $k$ algebraically closed of
characteristic $p>0$. Let $D\ne0$ be a fixed reduced divisor,
$U=\operatorname{Spec}R\setminus D$, and let $Q(R,D)$ classify
finite etale covers of $U$ tame at the generic points of $D$.
Fix a projective regular log resolution with total boundary
$E+D'$, where $E$ is the nonempty proper SNC exceptional fibre.
Write $D'_a$ for the horizontal strict transforms.
For $q=p^f$, put $n=q-1$. A rank-$r$ projective representation
of $Q(R,D)$ determines a class
\[
        \alpha\in H^2(Q(R,D),\mu_n(k)),
        \qquad \beta\in\operatorname{Br}'(U)[n].
 \tag{D.LPT.1}
\]
There are at most $n^{c_0}$ choices for $\alpha$, for a constant
$c_0$ depending only on the fixed resolution. Let $m_i$ be the
order of the residue of $\beta$ at $E_i$, and let $h_a$ be its
order at $D'_a$. Then
\[
 h_a\mid m_{i(a)},\qquad
 M:=\operatorname{lcm}_i m_i=\operatorname{per}(\beta),
                     \qquad M\mid r,\quad M\mid n.
 \tag{D.LPT.2}
\]
The convention is $M=1$ for the zero Brauer class.

For each $a$ choose a generator of prime-to-$p$ geometric tame
inertia and write $g_a\in\operatorname{PGL}_r(\mathbb F_q)$ for
its image. Its order $e_a$ is prime to $p$.
After retaining one matrix lift of each $g_a$, there are at most
$q^{c_1r^2}$ possible horizontal data. For each such tuple and
$\alpha$, choose
\[
             N_a=n\operatorname{lcm}(h_a,e_a).
 \tag{D.LPT.3}
\]
One fixed $\mu_n$-gerbe over the mixed root stack
\[
 \mathcal X=Y[\sqrt[m_i]{E_i},\sqrt[N_a]{D'_a}]
 \tag{D.LPT.4}
\]
then carries the scalar-weight-one unit bundle of every represented
object with those data, after a banded identification over $U$.
The horizontal orders $N_a$ need not be bounded in terms of $r$.
They will not be used as numbers of coordinates.
\end{lemma}

\begin{proof}
The class count and the injection into etale cohomology are
Lemma~\ref{D:local-divisor-tame-central-H2}.
The associated degree-$r$ Azumaya algebra on $U$ has Brauer class
$\beta$, so its period divides $r$. The local Picard and proper
henselian comparison used in the complete-local Brauer proofs give
$\operatorname{Br}'(Y)[n]=0$: $\operatorname{Pic}(Y)\to
\operatorname{Pic}(E)$ is onto, proper base change identifies
$H^2(Y,\mu_n)$ with $H^2(E,\mu_n)$, and the latter is the Kummer
quotient of $\operatorname{Pic}(E)$.
Purity therefore detects $\beta$ by its residues along $E+D'$.
At a horizontal component, Lemma~\ref{D:horizontal-residue-order-at-node}
identifies its residue character with its secondary residue at
its unique node. The opposite secondary residue comes from the
adjacent $E_{i(a)}$, so $h_a\mid m_{i(a)}$.
Thus the least common multiple of the exceptional residue orders
already equals the period, proving (D.LPT.2).

First form the smaller coefficient root stack with orders $m_i$
on $E_i$ and $h_a$ on $D'_a$. The same regular-root purity and
mixed-order node calculation as in
Lemma~\ref{D:SNC-general-Brauer-geometry} extend $\beta$ over it.
The local input is
\citep[Proposition~3.2.1, Lemma~4.3.1 and Propositions~4.3.2,
4.3.7]{Lieblich2011ResidualD}; the arbitrary composite-order node
algebra is the internal calculation in that lemma, not a
prime-order statement attributed to the source. The Kummer
sequence \citep[Lemma~59.28.1, Tag~03PK]{StacksBrauerZeroD},
and extension of a line bundle across a Cartier boundary
on a regular root stack, extend the actual $\mu_n$ coefficient
class $\alpha$, not just its Brauer image. Fix such an extension
$\mathcal G_0$ once. This is the coefficient-gerbe extension
argument already proved for the general-SNC local count.

We specify why this fixed extension suffices for all horizontal
representations. On $Y^0=Y\setminus E$, a projective torsor whose
horizontal inertia has order $e_a$ extends to the root-$e_a$
stack along $D'_a\cap Y^0$. On a strict-local chart this is
Abhyankar's Kummer lemma, followed by purity; the exact input is
\citep[Expose~XIII, Propositions~5.1--5.2 and Corollary~5.3, pp.~428--430]{C_SGA1}.
The lifting gerbe of this extended projective torsor is a
$\mu_n$-gerbe. On the common roots
$l_a=\operatorname{lcm}(h_a,e_a)$ its restriction to $U$ has the
same coefficient class as $\mathcal G_0$.
The difference of the two classes is supported on the disjoint
horizontal root divisors in $Y^0$. Localization and codimension-one
purity identify the degree-two support contribution with
$\bigoplus_a H^0(\mathcal D_a,\mathbf Z/n)$; its image is the
Kummer first Chern class of the root-divisor line.
This can equally be checked by the Kummer sequence: after extending
an $n$th root of its restriction to $U$, a difference line is a
Cartier combination of the root divisors. The purity calculation
applies on etale root charts and descends
\citep[Expose~XVI, Theorem~3.1.1 and the normal-crossings
Corollary~3.1.3]{A_ILO2014}.
On replacing $l_a$ by $N_a=nl_a$, each such root-divisor line pulls
back to the $n$th power of the new root-divisor line. Its Kummer
class is zero. The two actual coefficient gerbes are consequently
bandedly isomorphic on the enlarged root stack over $Y^0$.
Pull $\mathcal G_0$ back to (D.LPT.4) and use this identification.
Absolute Frobenius is natural for this identification, so the
scalar-weight-one unit bundle transfers with its unit structure.
No independently chosen incoherent gerbe equivalences are needed
at the $f$ stages.

A lift of a prime-to-$p$ element of $\operatorname{PGL}_r(\mathbb F_q)$
to $\operatorname{GL}_r(\mathbb F_q)$ has prime-to-$p$ order as
well: its $p$-part would lie in the scalar group of order $n$.
Its order divides $ne_a$, so it acts on the root-$N_a$ chart.
There are at most $|\operatorname{GL}_r(\mathbb F_q)|$ choices
per original horizontal branch. Their coefficientwise $p$-th
powers determine all $f$ phases, including the final character
permutation. This proves the asserted $q^{c_1r^2}$ bound, with
$c_1$ the fixed number of branches. It also avoids counting
arbitrary root characters independently at every stage.
Incompatible coefficient classes or matrices are simply discarded.
\end{proof}

\begin{proposition}[Proved: the marked normal presentation is uniform in horizontal root orders]
\label{D:local-pair-unbounded-horizontal-root-atlas}
Fix the coefficient class, its root gerbe, and horizontal matrix
data of the preceding lemma. Put $s_i=r/m_i$ and choose a fixed
positive integral exceptional cycle $A=\sum a_iE_i$ with
$c_i=-A\cdot E_i>0$. There are fixed original-$A$ pole and
normal-depth constants, independent of $q,r,m_i,N_a$.
At a fixed sufficient original depth $\ell$, the actual cyclic
bundle, Frobenius, inverse and equality data have a polynomial
presentation, with open subbundle and slope conditions, whose
variables may be assigned the following weights:
\[
\begin{array}{c|c|c}
\text{block}&\text{number per indicated layer}&\text{weight}\\ \hline
 i,\ 1\le j\le\ell a_i m_i&O(fm_i s_i^2)&(M/m_i)j\\
 i,\ \text{residual}&O(fm_i s_i^2)&M\\
 \text{each exceptional--exceptional node}&O(fr^2)&M.
\end{array}
 \tag{D.LPT.5}
\]
Every equation has weighted degree at most $KM$, with $K$ fixed.
Horizontal intersections are included as marked points in the
component rows; they do not add a component to the proper normal
cycle. In particular no factor $N_a$ occurs in (D.LPT.5).
The finite-cycle construction is effective over every coefficient
algebra and covers specified finite local Artinian families.
\end{proposition}

\begin{proof}
We verify precisely the extra uniformity beyond
Proposition~\ref{D:SNC-general-weighted-presentation}.
All fixed-reference coefficients below are constants, not moving
background variables.

\emph{Residual categories and a normalized Euler anchor.}
Let $R_i$ be the connected degree-$m_i$ cyclic spectral cover of
$E_i$ defined by its residue. At an exceptional node with secondary
order $h_x$ the residual root order on $R_i$ is $m_j/h_x$,
as before. At a horizontal node it is now $N_a/h_a$.
There are $m_i/h_x$ points over the node. Tame Riemann--Hurwitz gives
\[
 2g(R_i)-2=m_i(2g(E_i)-2)
               +\sum_x m_i(1-1/h_x),\qquad
 \eta_i:=\sum_xm_i/h_x=O(m_i).
 \tag{D.LPT.6}
\]
The sums include both sorts of nodes; their number on the original
resolution is fixed. Thus the genus and the number of spectral
marks remain $O(m_i)$, independently of every $N_a$.
The mixed-node central extension has irreducibles of dimension
$h_x$, with radical characters giving the residual root-point
weights. This is the same finite algebra calculation for the two
root orders, whether or not one of them is large. Scalar-zero
Hom carries an ordinary character of the abelian root quotient.
All these decompositions and invariant Hom calculations hold over
arbitrary coefficient bases on homogeneous frame atlases, because
the root orders are invertible in characteristic $p$.

For the qualitative Euler minimization alone choose a rank $H$
divisible by every required product of exceptional and horizontal
root orders; for instance $H=(M\prod_aN_a)^2$ suffices.
On $R_i$ choose a spectral bundle of rank $w_i=H/m_i$ whose
coarse bundle is $\mathcal O^{w_i}$ and whose normalized weights
at every mark occur uniformly. At an exceptional node the same
list of irreducibles occurs on the two proper branches, with
multiplicity $Hh_x/(m_im_j)$. At a horizontal mark every residual
weight occurs $w_i/(N_a/h_a)$ times. These integers exist by the
choice of $H$. Exact node pasting gives an anchor on the proper
residual support, and coherent $H^2=0$ lifts it formally.
An algebraic generating anchor is obtained from an Azumaya
representative on the regular root stack and all its root-character
translates, exactly as in
Lemma~\ref{D:SNC-general-Brauer-geometry}; its possibly enormous
rank is only an existence device for coarse-Hom Morita and
algebraization.

For a lattice $F$, minimize the nonnegative relative Euler defect
of $q_*\mathcal Hom(A_0,F)$ as in (D.GP.4). Its restriction off
$E$ is fixed on the root stack, including the horizontal divisors.
Its coarse Hom is locally free there, so its reflexive extension
over $R$ is fixed and both terms of that defect are finite.
For an elementary-transform quotient $Q_i$ on the spectral root
curve, the exact Euler correction is
$\chi(R_i,\pi_*\mathcal Hom(V_i,Q_i))$, with no extra cover
factor. Dividing by $w_i\operatorname{rk}Q_i$ gives the same bounds
\[
 \mu_{\rm par}(Q_i)+1-g(R_i)-2\eta_i
 \ \le\ \frac{\chi(\mathcal Hom(V_i,Q_i))}
                    {w_i\operatorname{rk}Q_i}
 \ \le\ \mu_{\rm par}(Q_i)+1-g(R_i).
 \tag{D.LPT.7}
\]
Indeed every normalized weight lies in $[0,1)$, however large its
denominator. The anchor's parabolic slopes lie in $[0,\eta_i]$.
The factor $w_i$ cancels exactly; no estimate proportional to $H$
is used in the coordinate count. The inverse elementary transform
uses root-normal parabolic degree $E_i^2$. Giving HN subbundles
their induced flags therefore gives coarse windows
\[
 g(R_i)-1-\eta_i\le\mu_{\min}(W_{i,0})
       \le\mu_{\max}(W_{i,0})
       \le g(R_i)-1-E_i^2+2\eta_i.
 \tag{D.LPT.8}
\]
All endpoints are $O(m_i)$ with constants independent of $N_a$.
This is the normalized marked Euler argument, not an assertion
that the Euler anchor has rank bounded by a polynomial in $r$.

A full original-$A$ conormal layer has degree $m_ic_i$ on $R_i$.
Every incomplete root layer through a fixed original depth has
parabolic degree $O(m_i)$; passing to normalized coarse weights
changes it by at most $\eta_i$. The one-step Frobenius carries
at every marked point are at most $p-1$, even for order $N_a/h_a$.
The universal spectral lines can be normalized to rational
parabolic degree in $[0,1)$, so their Frobenius comparison has
bounded parabolic degree before these $O(\eta_i)$ corrections.
The established rank-independent curve Frobenius HN estimate in
Lemma~\ref{D:twisted-period-independent-normal-depth} consequently
gives enlarged $O(m_i)$ windows for all adjacent Frobenius terms.
Duality and the ordinary/parabolic Hom sandwich then give fixed
original-$A$ negative-$H^0$ pole bounds and positive-$H^1$ tail
vanishing. The fixed proper-curve cohomology complexes give these
vanishings after every coefficient base change, by the same
split-differential argument using Tag~07VK in the general-SNC
proof. They therefore give prescribed Artinian lifts and exclude
nilpotent poles. This uses the positivity of the fixed $c_i$,
not the magnitude of a horizontal root index.

\emph{Actual small marked sources, not a rank-$H$ object atlas.}
For each component take the connected rank-$m_i$ reference anchor
and the spectral source $P(g)$ of
Lemma~\ref{D:SNC-tree-small-marked-source}. Its coarse rank is
$s_i+1$. At a marked root point of any order $d$ put
$T=\mathcal O((d-1)\mathcal P)$ and use the quotients
\[
                 T/T(-(d-1-a_\nu)\mathcal P),
                  \qquad 0\le a_\nu<d.
 \tag{D.LPT.9}
\]
A homogeneous frame gives their quotient map from $T^{s_i+1}$.
Its kernel is an actual free homogeneous source over every
coefficient ring, with coarse bundle $\mathcal O^{s_i+1}$.
Give the kernel of its generating quotient weight zero, and use
the alternating rank-two syzygies to form the fixed partial
complex of length seven. The rank condition is imposed on
$P_1\to P(g)$, not on its possibly vanishing composite into
$T^{s_i+1}$. Unit-minor cancellation proves all-ring local
freeness and exact middle terms. The usual $H^1$ generator
vanishing proves coverage of every specified Artinian lift.
Neither the moving determinant nor the large Euler anchor is
used as an object parameter.

We isolate the quantitative issue at an unbounded root point.
For normalized characters $a,b\in[0,d)$, a coarse parabolic Hom
entry is either ordinary or required to vanish on the one coarse
marked Cartier divisor. This is independent of $|a-b|$ and $d$.
Equivalently, the Cartier resolutions of (D.LPT.9) have intervals
of exponents of length at most $d$, and one residue class modulo
$d$ is prescribed by invariance. Each pair contributes at most
one coefficient per normal layer. After a single $p$-Frobenius
pullback the interval has length at most $pd$, so at most $p$
coefficients occur. Transverse Cartier equations remain regular
on every normal thickening and after arbitrary coefficient base
change. Division by them gives the unique support null homotopy;
its next composite is zero by regularity. Distinct marked supports
have zero cross-Hom and are contracted separately. Thus no
square of the number of marks, and no factor $d$, is retained.

The fixed termwise Hom cohomology is therefore $O(m_i s_i^2)$
per exceptional root layer: point blocks have the preceding
bound at $O(m_i)$ marks, while curve blocks have the genus and
line-degree bounds $O(m_i)$ or $O(m_is_i)$ exactly as in
Proposition~\ref{D:SNC-tree-fixed-normal-presentation}.
The two-affine Cech--Cartier contractions preserve column degree.
The only variable perturbations either raise root-normal order
or raise column degree in a fixed-width complex. The finite
homological-perturbation formulas consequently have weighted
degree $O(m_i)$ through $\ell a_i m_i$ layers, including the
fixed-number Cartier-division witnesses. Source length seven,
target length three and the top-source identity test give the
exact ordinary Hom of the quotient bundles. This is the full
all-ring construction in that proposition; the interval count
above verifies that it remains uniform when $d=N_a/h_a$ is
arbitrarily large.
The known inputs are the exact tensor/root equivalence
\citep[Proposition~3 and Theorem~4]{Borne2007D},
\citep[Theorem~6.1]{BorneVistoli2012D}, and the finite perturbation
identities
\citep[Sections~2.1--2.4 and~3.2]{C_Crainic2004Perturbation}.
They are used in their invertible-denominator scopes.

Only exceptional--exceptional nodes must be pasted to assemble
the proper Cartier cycle $\ell A$. Their finite rectangles have
root lengths $\ell a_i m_i$ and $\ell a_jm_j$. Scalar-zero Hom
has exactly the bounded residue-class slots of the general-SNC
proof. Milnor pasting, with the full perfect restriction of the
point terms and all inverse witnesses, costs $O(fr^2)$ at each
original node
\citep[Lemmas~15.6.4--15.6.9, Tag~0D2J]{StacksTreePastingD}.
A horizontal intersection is already a marked point on one
component thickening; no horizontal component or extra infinite
rectangle is pasted. Rescale component weights by $M/m_i$ and
give residual and exceptional-node variables weight $M$.
This proves (D.LPT.5) and the common bound $KM$.

Finally the finite-cycle data have the same formal effectivity
as in Proposition~\ref{D:SNC-general-weighted-presentation}.
The proper tame one-dimensional support has coherent $H^2=0$,
so actual bundles extend successively. The established $H^1$
tails lift prescribed arrows, and the paired unit contraction
raises ordinary-$A$ order by $t\mapsto pt-2B$ after a fixed
precision. It restores the exact unit identities while retaining
the specified jets. For a finite-type coefficient algebra, or a
finite local Artinian one, perform this over the completion of
$R\otimes_k A_{\rm coeff}$ and apply ordinary coarse-Hom
Grothendieck existence with the algebraic generator, followed by
Morita. All arbitrary-ring assertions before this step concern
finite cycles. This is not an application of Noetherian existence
to an arbitrary coefficient ring.
The same bounded-pole proper Hom functor represents the actual
unit-Isom functor in finite presentation. Specified nilpotent
unit isomorphisms lift uniquely off $E$, including on the horizontal
root divisors, by etale invariance; the all-base $H^0$ pole bound
and the contraction give their finite-precision representatives.
Thus its class images are open, by the already checked
finite-unit argument and
\citep[Tags~08K6, 0BQB, 02HT, 02HM and~00UR]{StacksUnitJetTransferD}.
No property in this comparison depends on $N_a$ except that it is
invertible in $k$ and is fixed within the chosen presentation.
\end{proof}

\begin{theorem}[Proved: quadratic projective growth for an original local tame pair]
\label{D:local-divisorial-tame-projective-growth}
For every fixed pair $(R,D)$ in the preceding lemma there is a
constant $C_{R,D}$ such that, for all $q=p^f$ and $r\geq1$,
\[
 \#\bigl(\operatorname{Hom}_{\rm cont}
          (Q(R,D),\operatorname{PGL}_r(\mathbb F_q))
                    /\operatorname{PGL}_r(\mathbb F_q)\bigr)
                         \le q^{C_{R,D}r^2}.
 \tag{D.LPT.10}
\]
Tameness is imposed only at the original height-one boundary.
Represented monodromy may be wild at exceptional divisors.
For empty divisorial boundary, the same assertion for the
punctured spectrum is the already proved general-SNC theorem.
No constant uniform in the degree of arbitrary finite covers,
nonabelian crown bound or full finite-generation assertion is
included.
\end{theorem}

\begin{proof}
There are at most $n^{c_0}$ realized coefficient classes by
Lemma~\ref{D:local-divisor-tame-central-H2}, and at most
$q^{c_1r^2}$ horizontal matrix tuples by
Lemma~\ref{D:local-pair-coefficient-root-comparison}.
Fix both. The coefficient gerbe and all its scalar Frobenius
stages are now fixed, and the latter lemma places every represented
object on its fixed root stack. Its residue orders divide $r$;
write $r=m_is_i$.

The remaining character and rank allocations cost
$\exp(O(fr^2))$ possibilities. At an exceptional node the
mixed-order argument and weak-composition bound of the general-SNC
theorem apply. At a horizontal node the fixed matrix has at most
$r$ eigencharacters. After the secondary order $h_a$ is fixed,
its $h_a$-orbits are assigned to the at most $m_i/h_a\le r$
spectral points, with total multiplicity at most $r/h_a$.
Labeling these slots overcounts all assignments by $r^{O(r)}$;
subflag orders and their phase permutations obey the same bound.
The actual horizontal matrices determine their $p$-twists at all
$f$ stages. There is no independent factor $N_a^r$, either once
or at every phase. Keep only compatible allocations.

For each choice use Proposition~\ref{D:local-pair-unbounded-horizontal-root-atlas}.
Its component-wise weighted factorial estimate is exactly the
one in the general-SNC count: with
$b_i=\lceil C_0fm_i s_i^2\rceil$ and $J_i=\ell a_i m_i$,
the positive-layer denominator contains
\[
 \prod_i\bigl((M/m_i)^{J_i}J_i!\bigr)^{b_i},
\]
while residual and exceptional-node variables of weight $M$
cancel the same factor in the common degree $KM$.
Using $J_i!\ge(J_i/e)^{J_i}$, the logarithm of the resulting
component bound is
\[
            O\!\left(\sum_i b_iJ_i+fr^2\right)=O(fr^2).
 \tag{D.LPT.11}
\]
The component lemma used here is
Lemma~\ref{D:weighted-affine-component-bound}; open slope and
subbundle restrictions do not increase the number of components.
All constants depend only on the fixed exceptional resolution
and original boundary.

Actual unit-Isom class images in this presentation are open by
the last paragraph of the proposition. Distinct classes have
disjoint images, so their number is bounded by its irreducible
components. Every projective representation being counted has
at least one scalar-weight-one unit realization in one of these
presentations. Conversely its fixed-frame torsor descends
projectively over the horizontal root stack off $E$, and hence
restricts to a representation tame along the original divisor.
Only this surjection from unit classes to represented projective
classes is needed; uniqueness of a linear lift is not asserted.
Multiplying (D.LPT.11) by the coefficient-class, horizontal-matrix
and discrete-type factors yields $\exp(O(fr^2))=q^{O(r^2)}$.
This proves (D.LPT.10).
\end{proof}


% Proved: four complete actual-file audits; local-pair input separately promoted.
\subsection{Global projective growth with the original divisorial boundary}

\begin{lemma}[Proved: one curve detects the original tame quotient with fixed local sources]
\label{D:tame-detecting-curve-local-pairs}
Let $S$ be a normal proper integral surface over an algebraically
closed field $k$ of characteristic $p>0$, let $T\subsetneq S$ be
closed, and let $Q=Q(S,T)$ require tameness only at the original
codimension-one components of $T$. There are a smooth projective
connected curve $C$ over a fixed algebraically closed extension
$\Omega/k$, a finite nonempty set $P\subset C$, and a morphism
$C^\circ=C\setminus P\to S\setminus T$ such that
\[
                      \pi_1(C^\circ)\twoheadrightarrow Q.
                                                        \tag{D.GTP.1}
\]
Every marked local type of a representation of $Q$ comes from
one of the following fixed sources:
\begin{enumerate}
\item a complete normal local surface puncture group
$\Lambda_s=\pi_1(\operatorname{Spec}R_s\setminus\{\mathfrak m_s\})$;
\item the original divisorial tame quotient $Q(R_s,D_s)$ of a
fixed complete normal local surface pair with $D_s\ne0$;
\item the pro-prime-to-$p$ cyclic group, or the trivial group.
\end{enumerate}
Here $R_s=\widehat{\mathcal O}_{S,s}$ is over the original field
$k$. The source maps are independent of the representation, its
rank and its finite field. There are fixed numbers $b_x\ge0$
bounding the highest classical upper break of every realized
local quotient at $x\in P$.
\end{lemma}

\begin{proof}
Write the reduced boundary as the union of its reduced divisorial
part $D$ and a finite set outside $D$. Choose a smooth projective
log resolution $\pi:Y\to S$. Enlarge a finite set $F\subset S$
so that it contains the non-isomorphism image, all isolated
boundary points, and all singular or intersection points of the
original boundary curves. Blow up further over $F$ as needed.
Then $\pi$ is an isomorphism away from $F$, its exceptional
fibre $E$ has support $\pi^{-1}(F)$, and
$B=E+D'$ is SNC, where $D'$ is the strict transform of $D$.
The complement $Y\setminus B$ identifies with
$(S\setminus T)\setminus F$. If $B$ is empty, add a retained
smooth point to $F$ and blow it up. These are the choices in
Proposition~\ref{D:surface-curvewise-conductor}.

Take one geometric generic hyperplane of a fixed projective
embedding of $Y$, over
$\Omega=\overline{k(z_0,\ldots,z_N)}$, and put $P=C\cap B_\Omega$.
The curve is smooth and geometrically irreducible, meets the
boundary transversely and avoids its nodes. The set $P$ is finite
and nonempty. The exact basepoint-free Bertini argument in
Lemma~\ref{D:finite-puncture-detecting-curve-local-types} applies
to every connected original cover of $S\setminus T$: its
restriction to the dense smooth open is integral, the two
independent hyperplane differentials verify the codimension-two
condition, and the third hyperplane avoids that component.
Thus the pullback to the one fixed geometric-generic curve is
connected. This proves surjectivity onto the original full
fundamental group and hence (D.GTP.1). The primary Bertini
inputs are \citep[Lemmas~33.47.3 and~37.32.1--3, Tags~0FD6 and
0G4C]{StacksDGenericSurface}. No surjectivity onto every cover
after extending $k$ is asserted.

First take $x\in C\cap E_\Omega$, lying over $s\in F$.
The completed formal branch gives the continuous map
\[
                 R_s\longrightarrow\Omega[[t_x]]
                                                        \tag{D.GTP.2}
\]
by exactly the completion argument in that lemma. If $s$ is
an isolated boundary point outside $D$, its generic branch
factors through the punctured spectrum of $R_s$, giving source~(1).
If $s\in D$, let $D_s$ be the reduced inverse image of $D$ in
$\operatorname{Spec}R_s$. The generic point of the branch is
outside $D_s$, since the generic point of $C$ is outside $B$.
The complement of $D_s$ maps to $S\setminus T$; the closed point
of the complete local scheme already lies on $D_s$.

Restriction of every original tame cover to this complement is
tame at each height-one branch of $D_s$. To check this precise
transfer, let $\mathfrak p$ be such a height-one prime of $R_s$.
Its contraction to $\mathcal O_{S,s}$ is a height-one prime of
$D$. It contains the ideal of $D$, so is not zero, and cannot
be the maximal ideal: the extended maximal ideal is the maximal
ideal of $R_s$, of height two. Thus there is an inclusion of DVRs
\[
 (\mathcal O_{S,s})_{\mathfrak p\cap\mathcal O_{S,s}}
                               \longrightarrow (R_s)_{\mathfrak p}.
                                                        \tag{D.GTP.3}
\]
Faithful flatness of completion gives injectivity. Tame
extensions stay tame after any extension of DVRs
\citep[Lemma~15.116.9(2), Tag~0EXY]{StacksDFiniteTameDescent}.
Consequently the map to $Q$ kills each defining wild inertia
group and factors through the fixed $Q(R_s,D_s)$, giving
source~(2). If $s\notin T$, all original covers extend over
the strictly henselian $\operatorname{Spec}R_s$ and split
there; the corresponding local map to $Q$ is trivial.

Now take a marked point on a strict transform $D'_a$ away from
$E$. Its image in the original $k$-curve $D_a$ is the generic
point. Indeed a geometric generic hyperplane avoids every
fixed closed $k$-point of $D'_a$: its defining linear relation
cannot vanish on algebraically independent coefficients.
The formal branch therefore gives an extension from the
original boundary DVR to $\Omega[[t_x]]$; transversality gives
ramification index one. The same tame-base-change lemma shows
that every pulled-back original cover is tame there.
Since $\Omega$ is algebraically closed, the full tame local
group is pro-prime-to-$p$ cyclic. This is also the local
Kummer description in \citep[Expose~XIII, Propositions~5.1--5.2
and Corollary~5.3]{C_SGA1}. It proves source~(3).

Finally apply Proposition~\ref{D:surface-curvewise-conductor}
over the fixed field $\Omega$ to the fixed pair and resolution.
An original tame cover remains tame on the original divisors
after this base change, by the same DVR argument. The proposition
therefore bounds its curvewise highest logarithmic slopes by
one fixed divisor. For a finite quotient $H$, use the faithful
regular representation $\mathbb F_\ell[H]$, with $\ell\ne p$.
The bound is independent of its rank and hence of $H$.
Logarithmic upper numbering on $\Omega((t_x))$ agrees with
classical upper numbering. Faithfulness gives the asserted
$b_x$ for the entire finite quotient. This is precisely the
highest-slope argument already checked in the finite-puncture
detecting-curve lemma, using
\citep[Theorem~3.5, Remark~3.9 and Proposition~5.8(3)]
{C_HuTeyssierRam2025}. Only realized local types need this bound.
\end{proof}

\begin{theorem}[Proved: quadratic projective growth for the original tame surface quotient]
\label{D:global-original-boundary-projective-growth}
For every pair $(S,T)$ in the preceding lemma there is a
constant $C_{S,T}$ such that, for every $q=p^f$ and $r\ge1$,
\[
 \#\bigl(\operatorname{Hom}_{\rm cont}
       (Q(S,T),\operatorname{PGL}_r(\mathbb F_q))
                    /\operatorname{PGL}_r(\mathbb F_q)\bigr)
                           \le q^{C_{S,T}r^2}.
                                                        \tag{D.GTP.4}
\]
No geometric Brauer-period, scalar-order, projectivity, or
exceptional-tameness assumption is imposed. Tameness is required
only at the original codimension-one boundary. The constant is
for this fixed pair, with no linear dependence on arbitrary
finite-cover degree asserted.
\end{theorem}

\begin{proof}
Use the one curve and fixed local sources just constructed.
At a point of source type~(1), the arbitrary-SNC local theorem
\ref{D:SNC-general-projective-quadratic-count} bounds the possible
projective types by $q^{a_xr^2}$ for a fixed $a_x$.
For type~(2), use
Theorem~\ref{D:local-divisorial-tame-projective-growth}.
For type~(3), a homomorphism from the pro-prime-to-$p$ cyclic
group is determined by one prime-to-$p$ element of the target,
so there are at most
$|\operatorname{PGL}_r(\mathbb F_q)|\le q^{r^2}$ choices;
the trivial source gives one choice.

After restriction to the completed branch these bounds remain
valid. Multiplying them independently at the finitely many
marks only overcounts the possible tuples. Every tuple actually
realized has the fixed highest-break bound of the lemma.
Corollary~\ref{D:curve-fixed-local-projective-growth} therefore
gives at most $q^{(C_C+\sum_xa_x)r^2}$ projective classes on
$C^\circ$. Precomposition with (D.GTP.1) is injective on
conjugacy classes, proving (D.GTP.4). No global degree-two
inflation, Brauer vanishing, or global linear lifting has been
used; the same-field lift is only the one in the affine curve
theorem with its fixed full local types.
\end{proof}

\begin{corollary}[Proved: controlled-boundary dimensions and all defining-characteristic simple products]
\label{D:global-tame-projective-controlled-closure}
Let $S$ be normal, proper and integral of dimension at least two
over an algebraically closed field of characteristic $p>0$,
and let $T\subsetneq S$ be closed. Suppose either that $T$ is
finite, or that the pair admits the projective modification
with controlled exceptional centres in
Theorem~\ref{D:generic-surface-controlled-boundary}.
Then (D.GTP.4) holds for $Q(S,T)$ with a constant for the fixed
pair. The quotient formed by all defining-characteristic finite
simple Lie quotient kernels is topologically finitely generated,
including all classical, exceptional, Suzuki and Ree families
and the fixed simple derived exceptions.
\end{corollary}

\begin{proof}
The finite-puncture case is already
Corollary~\ref{D:global-projective-finite-puncture-closure}.
In the other case the stated controlled-boundary theorem gives
one fixed surjection $Q(X,T_X)\twoheadrightarrow Q(S,T)$ from
a normal proper surface pair over a fixed algebraically closed
extension. Apply the preceding theorem to that fixed pair.
Injectivity of representation classes under precomposition
transfers the count with the same constant.

The purely finite-group embeddings and size estimates in
Corollary~\ref{D:global-unrestricted-projective-simple-products}
now apply unchanged. Distinct simple quotient kernels give
distinct ambient projective classes, hence polynomial
multiplicities in the simple-group order. Its exact LMT
generating-tuple estimate, Goursat argument and compactness
give the asserted simple-product quotient. This implication
does not transfer a crown theorem or prove full finite
generation of $Q(S,T)$.
\end{proof}

\begin{remark}[Scope of the combined projective reduction]
The local-pair construction and the present detecting-curve
transfer extend the same global projective reduction from
finite punctures to arbitrary original surface boundaries.
They retain every exceptional wild image. Hereditary constants
linear in arbitrary cover degree, unrestricted alternating and
cross-characteristic simple multiplicities, general nonabelian
crowns, and higher-dimensional boundaries outside the stated
reduction remain separate gates. General D.1, D.2(ii) and
dependent residual L8 remain open.
\end{remark}


% Proved enlargement of the alternating gate; four complete actual-file audits.
\subsection{Local permutation orbits with small matrix models}

\begin{lemma}[Proved: counting fixed permutation models through actual matrix homomorphisms]
\label{D:permutation-matrix-model-count}
Let $\Lambda$ be the fundamental group of the punctured spectrum of
a fixed complete normal local surface ring over an algebraically
closed field of characteristic $p>0$. For every $N\ge2$, choose
a finite collection of faithful permutation models
\[
 H_{N,j}\hookrightarrow S_N\qquad(j\in\mathcal J_N),
 \qquad
 H_{N,j}\hookrightarrow G_{N,j},
 \quad
 G_{N,j}\in\{\operatorname{GL}_{r_{N,j}}(\mathbb F_{q_{N,j}}),
              \operatorname{PGL}_{r_{N,j}}(\mathbb F_{q_{N,j}})\},
                                                        \tag{D.MO.1}
\]
where every $q_{N,j}$ is a power of $p$. The two displayed
embeddings are part of the fixed model. Assume that one constant
$B$ satisfies
\[
 \log\max\{1,|\mathcal J_N|\}\le B N\log(N+1),\qquad
 \max_{j\in\mathcal J_N}r_{N,j}^{2}\log q_{N,j}
                              \le B N\log(N+1).          \tag{D.MO.2}
\]
An empty collection is allowed, with the empty maximum in
(D.MO.2) interpreted as zero. Then the number of isomorphism
classes of transitive $N$-element $\Lambda$-sets whose actions are
conjugate into at least one model $H_{N,j}\le S_N$ is at most
\[
                         \exp\bigl(C N\log(N+1)\bigr)   \tag{D.MO.3}
\]
for a constant $C$ depending only on $\Lambda$ and $B$.
This bound counts permutation sets, not merely their linearizations.

The collections in (D.MO.1)--(D.MO.2) may include, simultaneously,
all the following standard actions in characteristic $p$:
\begin{enumerate}
\item $\operatorname{AGL}_d(\mathbb F_q)$ on the $N=q^d$ vectors,
      for every $d\ge1$ and every power $q$ of $p$;
\item $\operatorname{PGL}_d(\mathbb F_q)$ on all flags of any one
      fixed nonempty type of proper nonzero subspaces of
      $\mathbb F_q^d$, for every $d\ge2$ and every power $q$ of $p$.
\end{enumerate}
In particular the second family includes projective points and
every Grassmannian. For the union of these standard families,
the logarithm of the number of local $N$-element types is bounded
by $C_p\bigl(1+\log^2(N+1)\bigr)$, with a constant also depending
on $\Lambda$.
\end{lemma}

\begin{proof}
The local linear and projective representation bounds are
Theorem~\ref{D:uniform-local-representation-growth} and
Theorem~\ref{D:SNC-general-projective-quadratic-count}.
After increasing a fixed constant $a_\Lambda$, the number of
conjugacy classes of homomorphisms to either matrix group in
(D.MO.1) is at most $q^{a_\Lambda r^2}$. Each conjugacy class has
at most the order of the ambient group, and that order is at
most $q^{r^2}$. Consequently
\[
       |\operatorname{Hom}_{\rm cont}(\Lambda,G_{N,j})|
                     \le q_{N,j}^{(a_\Lambda+1)r_{N,j}^{2}}.
                                                        \tag{D.MO.4}
\]
Injection of $H_{N,j}$ into $G_{N,j}$ makes its actual
homomorphisms a subset of those counted in (D.MO.4).

Every permutation set covered by the assertion admits a
bijection with the standard $N$-element set under which its
action lands in one $H_{N,j}$. Choose one such bijection. It
produces an actual homomorphism $\Lambda\to H_{N,j}$.
Conversely every such homomorphism supplies one permutation
set; retaining only transitive homomorphisms can only lower
the count. Thus the disjoint union of the actual Hom sets
surjects onto the required isomorphism classes. There is no
additional cost for choosing affine structures, vector-space
structures, flags, or bijections: all choices have already
been allowed by this surjection. In particular no injection
from permutation classes into matrix-conjugacy classes is
asserted. Equations (D.MO.2) and (D.MO.4) give (D.MO.3), for
example with $C=B(a_\Lambda+2)$.

For the affine action, the faithful matrix realization is
\[
 (v,A)\longmapsto
 \begin{pmatrix}A&v\\0&1\end{pmatrix}
                  \in\operatorname{GL}_{d+1}(\mathbb F_q).
                                                        \tag{D.MO.5}
\]
Writing $q=p^f$ and $N=q^d$, we have
\[
     (d+1)^2\log q\le4d^2f\log p
                    \le\frac{4}{\log p}\log^2N.          \tag{D.MO.6}
\]
There are at most $\log_pN$ pairs $(d,f)$ with $p^{fd}=N$.

For a flag type, choose one of its dimensions $k$, with
$1\le k\le d-1$. Forgetting the other subspaces is a
surjection onto the Grassmannian. Counting ordered bases
and cancelling the stabilizer gives its usual Gaussian
binomial cardinality, and directly
\[
 N\ \ge\ \genfrac{[}{]}{0pt}{}{d}{k}_q
   =\prod_{i=0}^{k-1}\frac{q^{d-i}-1}{q^{k-i}-1}
   \ge q^{k(d-k)}\ge q^{d-1}.                            \tag{D.MO.7}
\]
The factorwise middle inequality follows by dividing numerator
and denominator by their leading powers of $q$.
The action of $\operatorname{PGL}_d$ on this flag set is
faithful: a linear map fixing all such flags fixes every
$k$-subspace, hence every line by intersecting the
$k$-subspaces through that line, and is therefore scalar.
It follows from (D.MO.7) that
\[
       d^2\log q\le\frac{4}{\log p}\log^2N.              \tag{D.MO.8}
\]
For fixed $N$, there are at most
$(1+\log_pN)^2$ possible pairs $(d,f)$. For each pair, the
number of nonempty types is at most $2^{d-1}$, which is
at most $N^{\log_p2}\le N$. This bounds the number of
possible models by a polynomial in $N$, independently
of which types actually have cardinality exactly $N$.
The matrix exponents (D.MO.6) and (D.MO.8), together with
this model count, prove the final logarithmic-square
bound and, in particular, (D.MO.2). A singleton has only
the trivial permutation action and can be included
separately without changing any estimate.
\end{proof}

\begin{theorem}[Proved: factorial counts with matrix-modelled local orbits]
\label{D:alternating-compressed-local-orbits}
Keep any of the fixed geometric and detecting-local data of
Corollary~\ref{D:alternating-dimensional-normalization-closure}.
For each fixed whole local surface group $\Lambda_{i,b}$,
retain the finite collection of finitely generated quotients
allowed there, and choose permutation models satisfying
(D.MO.1)--(D.MO.2), with fixed constants. Call a local orbit
admissible if either its action factors through one of those
fixed quotients, or its permutation action is conjugate into
one of the models of its degree. Fix $K>0$.

Then the isomorphism classes of degree-$n$ finite etale covers
of $U$, including disconnected covers, for which
\[
    \sum_{i,b}\
    \sum_{\substack{O\in\Lambda_{i,b}\backslash X\\
                     O\text{ not admissible}}}|O|^2
                                \le K n\log(n+1)         \tag{D.MO.9}
\]
are at most $n^{A n}$ for one constant $A$ depending only
on the fixed data. Distinct natural $A_n$ quotient kernels
satisfying (D.MO.9) have multiplicity at most $|A_n|^c$
for one constant $c$, and their selected alternating
simple-product quotient is topologically finitely generated.

In particular, one may leave uncharged all nilpotent-image
orbits, all prime-to-$p$-image orbits, and all orbits
affinely or projectively modelled by the standard families
in Lemma~\ref{D:permutation-matrix-model-count}, regardless
of their sizes. The linear complements of affine images
need not be solvable. All conditions concern the whole
fixed local surface-group action, not merely curve inertia.
\end{theorem}

\begin{proof}
It is enough to give the changed enumeration of full local
types. Every later step of
Corollary~\ref{D:alternating-dimensional-normalization-closure}
uses only a factorial bound for that enumeration, the
unchanged degree-linear discriminant bound on the fixed
detecting curves, and effective finite proper descent.
We nevertheless specify how each of those counts is retained.

There are finitely many local groups; write their number
as $w$. Write $J$ for the total number of chosen fixed
quotients and $d_0$ for the sum of fixed generator bounds
for those quotients. Assign an orbit admitting a quotient
factorization to one such quotient, and join all orbits
assigned to it. For each local group, the joined sizes
and their permutation actions contribute, in total over
all groups, at most
\[
                       (n+1)^J(n!)^{d_0}.                \tag{D.MO.10}
\]
This is the actual-generator count from
Corollary~\ref{D:fixed-local-quotient-alternating-count}.
Ambiguous assignments only overcount the required types.

Order the remaining orbits at each local group. There are
at most $2^n$ ordered compositions of their total size,
and at most $2^n$ ways to declare each part modelled or
uncontrolled. If a modelled part has size $m$, its
isomorphism types are bounded by
$\exp(C_*m\log(m+1))$, where $C_*$ is the maximum of
the finitely many constants in (D.MO.3). Since the total
size is at most $n$ at each local group,
\[
       \sum_{\text{modelled parts}}m\log(m+1)
                            \le w n\log(n+1).           \tag{D.MO.11}
\]
For the uncontrolled parts, the exact permutation-class
bound of Lemma~\ref{D:local-permutation-orbit-energy}
is $p^{(a_*+1)\sum m^2}$, including its ambient
linear-conjugacy fibre allowance. Here $a_*$ is a
fixed maximum of the local linear-count constants.
By (D.MO.9)--(D.MO.11), all tuples of full local types
therefore number at most
\[
 (n+1)^J(n!)^{d_0}4^{wn}
 \exp\!\left(\bigl[C_*w+(a_*+1)K\log p\bigr]
                                      n\log(n+1)\right).
                                                        \tag{D.MO.12}
\]
This is $n^{A_0n}$ for $n\ge2$, with $A_0$ fixed.
Every actual orbit is represented by at least one of
the assignments used; no assertion about canonical
structures or distinctness of assignments is needed.

The detecting curves and their ramification bounds do
not change when an orbit is counted by a different
method. For each fixed full local tuple, the connected
fixed-formal-fibre estimate is
$(n!)^{2h}n^{2D}$, with $D\le\lambda n$, by
Theorem~\ref{D:fixed-formal-curve-fibre-bound}.
For disconnected covers, assign local connected
factors to at most $n$ labelled global components.
The same factor assignment as in (D.ANC.5) costs at
most $n^{b_C n+1}$, where $b_C$ is the number of marked
points on this particular detecting curve, and bounds the
product of connected
fibre counts by $(n!)^{2h}n^{2\lambda n}$.
All exponents here are fixed by the detecting curve.
The detecting-surjection argument remains injective
on arbitrary finite-set isomorphism classes.
The finitely many proper curve components contribute
their fixed $(n!)^{d_i}$ bounds. Finite normalization
descent contributes at most $(n!)^a$ for the fixed
number of connected components of the double fibre
product, exactly as in the cited closure corollary.
Multiplying these bounds with (D.MO.12) gives $n^{An}$.
The degree-one case is the unique trivial cover.

Finally, distinct selected natural $A_n$ quotient
kernels give distinct degree-$n$ permutation classes.
The inequality $|A_n|\ge n^{n/8}$ for $n\ge5$
and the already proved LMT/Goursat/compactness
argument in
Theorem~\ref{D:finite-puncture-local-orbit-factorial-count}
give the polynomial multiplicity and one uniform
generator bound for the selected simple product.
This proves the whole assertion.
\end{proof}

\begin{remark}
This enlarges the same restricted alternating count:
a standard affine orbit of size $p^d$, for example,
can be uncharged even though its square $p^{2d}$
would violate a degree-$p^d$ square-sum bound.
The statement neither claims that every local
permutation action has one of these models nor
asserts a geometric realization of every standard
model as a local quotient. It makes no
finite-generation assertion for the full local
prosolvable group or for the whole global group.
\end{remark}


% Proved auxiliary: four full actual-file audits; separate post-CP24 research.
\subsection{Fixed parameter ideals as Cartier cycles on all covering models}

\begin{proposition}[Proved: Cartier two-level recovery on arbitrary complete covers]
\label{D:hereditary-Cartier-parameter-jets}
Let $R$ be a fixed complete excellent normal local surface over an
algebraically closed coefficient and residue field $k$ of characteristic
$p>0$. Fix the ideal $J$ and a separating parameter normalization
$S=k[[x,y]]\subset R$, with $x,y\in J$, as in
Lemma~\ref{D:canonical-pullback-uniform-jet-loss}. Put
$Q=(x,y)R$. There is one fixed regular projective modification
$Y_0\to\operatorname{Spec}R$ and an effective exceptional Cartier
cycle $C_0$ of full support such that
\[
                    Q\mathcal O_{Y_0}=\mathcal O_{Y_0}(-C_0).
                                                        \tag{D.HC.1}
\]
For every finite normal cover domain $A/R$ of degree $t$, etale on
the puncture, choose a regular projective resolution $\pi:Y\to
\operatorname{Spec}A$ dominating the normalization above $Y_0$.
Write $f:Y\to Y_0$ and $C=f^*C_0$. Then, for every $b\ge1$,
\[
 (x^b,y^b)\mathcal O_Y=\mathcal O_Y(-bC),\qquad
                    C^2=tC_0^2.                       \tag{D.HC.2}
\]
The cycle $C$ is anti-nef and has full exceptional support.

For a reflexive $A$-module $M$ of rank $r$, put
$F=(\pi^*M)^{\vee\vee}$ and
\[
 V_b=H^0(bC,F|_{bC}),\quad
 W_b=\operatorname{im}(V_{b+1}\longrightarrow V_b),\quad
 K_b=H^0(Y,F(-bC)).
\]
With the same constant $c$ as in the cited saturation lemma,
\[
 W_b=M/K_b,\qquad
 W_b/(x^\ell,y^\ell)W_b\simeq M/(x^\ell,y^\ell)M
                 \quad(b\ge2\ell-1+2c).                \tag{D.HC.3}
\]
For each fixed $b$, both $h^0(bC,F|_{bC})$ and
$h^1(bC,F|_{bC})$ are $O_b(tr)$, uniformly over all such complete
cover rings. In particular
\[
 -C^2+|K_Y\cdot C|+h^0(\mathcal O_{bC})+
                         h^1(\mathcal O_{bC})=O_b(t).  \tag{D.HC.4}
\]
The jet identifications are functorial for actual module maps.
They do not assert effectiveness or cohomology-and-base-change for
arbitrary parameter families.
\end{proposition}

\begin{proof}
Blow up $Q$ and resolve the resulting excellent surface, keeping the
puncture unchanged. On the blowup the extended ideal is invertible;
its pullback remains an invertible ideal on the regular resolution.
These are the usual blowup properties
\citep[Lemmas~31.33.4--31.33.5, Tag~01OF]{StacksHereditaryCartierD}.
Existence of regular projective surface resolutions is the established
surface-resolution input. This proves (D.HC.1).
Every finite normalization above $Y_0$ is proper over
$\operatorname{Spec}A$, so it admits a regular projective resolution.
One may moreover dominate any chosen regular resolution of $A$:
eliminate the rational map to that proper normalization by finitely
many blowups of closed points
\citep[Lemma~54.4.3, Tag~0C5H]{StacksHereditaryCartierD}.

Locally write $x=su$ and $y=sv$, where $s$ generates the invertible
ideal $Q\mathcal O_Y$. Then $(u,v)=\mathcal O_Y$.
The ideal $(u^b,v^b)$ also equals $\mathcal O_Y$: it is contained
in no maximal ideal. Thus $(x^b,y^b)=(s^b)$, proving the first
identity in (D.HC.2) without a flatness assumption on $f$.
Projection for Cartier intersections under a proper generically
finite map proves the second. The line bundle $\mathcal O_Y(-C)$
is generated by $x,y$, so it has nonnegative degree on every
exceptional curve. Since $Q$ is primary to the closed point,
its divisor contains every exceptional component.

Now apply Proposition~\ref{D:canonical-module-two-level-recovery}
on this particular regular model. Its parameter-ideal schemes
$T_b$ are exactly the Cartier schemes $bC$, by (D.HC.2).
Its exact transition-image identity and saturation bound give
(D.HC.3), with no change of the fixed integer $c$.
If $d_0=\operatorname{rank}_S R$ and $p_g(A)\le c_Rt$, its
explicit bounds are
\[
 h^0(bC,F|_{bC})\le(d_0b^2+2c_R)tr,\qquad
 h^1(bC,F|_{bC})\le c_Rtr.                              \tag{D.HC.5}
\]
These use only the fixed canonical-defect annihilator and apply
to arbitrary complete covers; they use no linear resolution-size
or algebraization-of-the-cover assumption.
For $M=A$ they give the asserted scalar cohomology bounds.
Adjunction for the proper Cartier cycle $C$ says
$2\chi(\mathcal O_C)=-C^2-K_Y\cdot C$. Hence
\[
 |K_Y\cdot C|\le -C^2+
                    2\bigl(h^0(\mathcal O_C)+h^1(\mathcal O_C)\bigr),
\]
which proves (D.HC.4). It is unnecessary to bound the number of
blowups used to make (D.HC.1) hold or to dominate the chosen
covering resolution.
\end{proof}

\begin{proposition}[Proved: generation and determinant budgets on the Cartier jet cycles]
\label{D:hereditary-Cartier-generation-budget}
Use the preceding proposition. Suppose in addition that the minimal
regular resolution $Y_{\min}$ of $A$ has a full-support anti-nef
cycle $Z_0$ satisfying
\[
 -Z_0\cdot B\ge2p_a(B)\quad(0<B\text{ exceptional}),
 \qquad -Z_0^2+|K_{Y_{\min}}\cdot Z_0|=O(t),             \tag{D.HC.6}
\]
with constants for the fixed base family. The established
energy/generation family of
Corollary~\ref{D:punctured-cover-antinef-cycle} and
Theorem~\ref{D:reflexive-pullback-uniform-generation} supplies
this hypothesis, including its separate treatment of wholly
zero-canonical components. That established family has the
algebraic or descended strict-henselian scope stated there;
no extension of its energy assertion to every nonalgebraizable
complete cover is made here.

Choose $Y$ above to dominate $Y_{\min}$, write $\rho:Y\to Y_{\min}$,
and put
\[
                          W=\rho^*Z_0+C.
                                                        \tag{D.HC.7}
\]
Then $W\ge C$, is anti-nef of full support, and
\[
 -W^2+|K_Y\cdot W|=O(t),\qquad
                 -W\cdot B\ge2p_a(B)\quad(0<B\text{ exceptional}).
                                                        \tag{D.HC.8}
\]
Consequently every canonical full pullback $F$ satisfies that
$F(-W)$ is globally generated.

For a cyclic Frobenius system of reflexive modules
$(M_a)_{a\in\mathbf Z/f\mathbf Z}$ of rank $r$, put
$F_a=(\pi^*M_a)^{\vee\vee}$. On each fixed Cartier jet cycle
$bC$, the bundles $V_a=F_a(-W)|_{bC}$ obey
\[
 h^0(bC,\det V_a)=O_b(tr),
                                                        \tag{D.HC.9}
\]
and each fixed-determinant class is covered by an actual open
quotient chart with $O_b(tr^2)$ coordinates, as in
Proposition~\ref{D:hereditary-fixed-determinant-cycle-chart}.
This assertion concerns actual objects on the finite Cartier
cycle. It supplies neither a period-uniform projective anchor
nor an arbitrary-family unit-object realization or count.
\end{proposition}

\begin{proof}
The regular birational map $\rho$ factors into blowups of closed
points \citep[Lemma~54.17.1, Tag~0C5R]{StacksHereditaryCartierD}.
The all-subcycle inequality is preserved by pulling back through
each blowup. Indeed, for a nonzero effective cycle $B$ upstairs,
write $B_0=\rho_*B$. If $B_0\ne0$, then
$B=\rho^*B_0-kE$ for an integer $k$, and adjunction gives
\[
                  p_a(B)=p_a(B_0)-k(k-1)/2\le p_a(B_0).
                                                        \tag{D.HC.10}
\]
Projection identifies the intersection with the pulled-back
cycle. If $B_0=0$, then $B=aE$ with $a\ge1$ and
$p_a(B)=1-a(a+1)/2\le0$, while the pullback has degree zero
on $B$. This proves the pullback assertion, by induction.
Adding the anti-nef cycle $C$ can only increase the left side
of the all-subcycle inequality. It proves the second assertion
of (D.HC.8), and the cited uniform-generation theorem applies
to every canonical $F$.

For the energy, Cauchy--Schwarz in the negative exceptional
intersection form gives
\[
 -W^2\le\left(\sqrt{-Z_0^2}+\sqrt{-C^2}\right)^2=O(t).
\]
The canonical intersection with a pulled-back Cartier divisor
is unchanged under point blowup. Thus
$K_Y\cdot\rho^*Z_0=K_{Y_{\min}}\cdot Z_0$.
Together with (D.HC.4) this proves the first assertion of
(D.HC.8), even if $Y$ has arbitrarily many additional rational
$(-1)$-curves. No positivity of their individual canonical
intersections is assumed.

Since $Q\subset J$, the ordinary-to-reflexive Frobenius defect
is killed by $(x^p,y^p)$. On $Y$ this is the invertible ideal
$\mathcal O_Y(-pC)$. The same codimension-one calculation as in
(D.UG.13)--(D.UG.14) therefore gives
\[
 F_{a+1}(-pC)\subset F_Y^*F_a\subset F_{a+1},\qquad
 c_1(F_{a+1})=p c_1(F_a)+B_a,
                         \quad0\le B_a\le rpC.         \tag{D.HC.11}
\]
Generic generation makes $x_a=c_1(F_a)\cdot C\ge0$.
Anti-nefness of $C$ yields
$p x_a-rp(-C^2)\le x_{a+1}\le p x_a$.
Cycling and cancelling $p^f-1$ proves
\[
                  0\le x_a\le\frac{rp}{p-1}(-C^2)=O(tr).
                                                        \tag{D.HC.12}
\]
Also $0\le-W\cdot C\le\sqrt{(-W^2)(-C^2)}=O(t)$.
Hence the globally generated determinant of $V_a$ has degree
\[
              \deg_{bC}\det V_a=b(x_a-rW\cdot C)=O_b(tr).
\]
A regular section of that determinant gives
$h^0(\det V_a)\le h^0(\mathcal O_{bC})+
\deg_{bC}\det V_a$, proving (D.HC.9).
The fixed-determinant chart is the established quotient
$L^{-1}\to\mathcal O_{bC}^{r+1}$ on its generating open.
Its universal quotient is locally split and remains a vector
bundle over arbitrary coefficient rings. The asserted size is
$(r+1)h^0(L)=O_b(tr^2)$.
\end{proof}


% Proved: original-boundary closure; four complete actual-file audits.
\subsection{Alternating counts with the original divisorial tame boundary}

\begin{theorem}[Proved: matrix-model alternating bounds for original tame boundaries]
\label{D:original-boundary-matrix-alternating-count}
Let $S$ be normal, proper and integral over an algebraically
closed field of characteristic $p>0$, with $\dim S\ge2$, and let
$T\subsetneq S$ be closed. In dimension greater than two, assume
that either $T$ is finite or $(S,T)$ satisfies the controlled
exceptional-centre hypothesis of
Theorem~\ref{D:generic-surface-controlled-boundary}. Put
$Q=Q(S,T)$, imposing tameness only along the original
codimension-one components of $T$.

There is one fixed finite collection of homomorphisms
\[
                         \Lambda_x\longrightarrow Q     \tag{D.OBM.1}
\]
in which each $\Lambda_x$ is a complete normal local surface
puncture group, the original divisorial tame quotient of a
complete normal local surface pair, a pro-prime-to-$p$ cyclic
group, or the trivial group. These are precisely the fixed
sources supplied by Lemma~\ref{D:tame-detecting-curve-local-pairs}
after the stated surface reduction, as specified in the proof.

For each source, choose finitely many fixed topologically
finitely generated quotients, and choose finite permutation
models satisfying (D.MO.1)--(D.MO.2), with fixed constants.
Call a whole $\Lambda_x$-orbit on a finite $Q$-set admissible
if its action factors through one of those quotients or is
conjugate into one of its prescribed permutation models.
For every fixed $K>0$, the number of isomorphism classes of
degree-$n$ finite continuous $Q$-sets satisfying
\[
       \sum_x\
       \sum_{\substack{O\in\Lambda_x\backslash X\\
                       O\text{ not admissible}}}|O|^2
                              \le K n\log(n+1)           \tag{D.OBM.2}
\]
is at most $n^{A n}$, with $A$ depending only on the fixed
data. Equivalently this counts all rank-$n$ finite etale
covers of $S\setminus T$ belonging to the original tame
quotient, including disconnected covers.

For the first two kinds of source one may jointly choose
the maximal pronilpotent and maximal pro-prime-to-$p$
quotients. For the last two, one may choose the whole source.
One may simultaneously include all standard affine point
and projective partial-flag models from
Lemma~\ref{D:permutation-matrix-model-count}. Thus arbitrary
sizes of all those orbit kinds are uncharged in (D.OBM.2).
The local linear complements need not be solvable.

The distinct natural $A_n$ quotient kernels satisfying
(D.OBM.2) have multiplicity at most $|A_n|^c$ for one
constant $c$, and the corresponding selected alternating
simple-product quotient of $Q$ is topologically finitely
generated. No exceptional-divisor tameness is imposed.
\end{theorem}

\begin{proof}
In dimension two use the fixed detecting curve and sources
of Lemma~\ref{D:tame-detecting-curve-local-pairs}. Its
surjection is onto the original $Q$, and its local source
maps and highest-upper-break constants are fixed before
any finite set or quotient is chosen. In higher dimension,
Theorem~\ref{D:generic-surface-finite-punctures} for finite
$T$, or Theorem~\ref{D:generic-surface-controlled-boundary}
in the other case, supplies one normal proper surface
pair $(S_0,T_0)$ over a fixed algebraically closed extension
and a surjection
\[
                         Q(S_0,T_0)\twoheadrightarrow Q.
                                                        \tag{D.OBM.3}
\]
Apply the same detecting-curve lemma to this fixed surface
pair and compose its source maps with (D.OBM.3). In this
case the complete local rings are attached to the chosen
detecting surface, not necessarily to the original
higher-dimensional $S$. The finite collection (D.OBM.1)
is thereby fixed in every case.

We first verify the precise local counting inputs, including
the newly allowed local pairs. For a complete puncture,
Theorem~\ref{D:SNC-general-projective-quadratic-count}
gives the projective representation bound. For a local
pair, it is
Theorem~\ref{D:local-divisorial-tame-projective-growth}.
For a pro-prime-to-$p$ cyclic source, a homomorphism to
a finite group is determined by the image of one
topological generator, and the trivial source requires
one choice. Consequently every fixed source has a
constant $a_x$ for which
\[
 |\operatorname{Hom}_{\rm cont}
       (\Lambda_x,\operatorname{PGL}_r(\mathbb F_q))|
                          \le q^{a_x r^2}
                    \qquad(q=p^f).                     \tag{D.OBM.4}
\]
Here the passage from projective-conjugacy classes to
actual homomorphisms costs at most
$|\operatorname{PGL}_r(\mathbb F_q)|\le q^{r^2}$,
which has been absorbed into $a_x$.

There is a faithful homomorphism
\[
 \operatorname{GL}_r(\mathbb F_q)\longrightarrow
 \operatorname{PGL}_{r+1}(\mathbb F_q),\qquad
                 A\longmapsto[\operatorname{diag}(A,1)].
                                                        \tag{D.OBM.5}
\]
Indeed a scalar diagonal matrix with final entry one is
the identity only when its scalar is one. Thus (D.OBM.4)
also bounds actual linear Hom sets by $q^{4a_xr^2}$.
This uses an embedding of actual finite groups, not a
claim that all projective representations lift linearly.
Applying the same embedding to the natural permutation
matrices also bounds the number of arbitrary
$m$-element local permutation types by
\[
                              p^{c_xm^2}                \tag{D.OBM.6}
\]
for a fixed $c_x$. As in the matrix-model lemma, counting
actual homomorphisms automatically pays any further
permutation-versus-matrix-conjugacy fibre.
Equations (D.OBM.4)--(D.OBM.6) give all the local
enumeration inputs of
Theorem~\ref{D:alternating-compressed-local-orbits}.
They are uniform over all models in (D.MO.1)--(D.MO.2).

We next justify the concrete finitely generated local
quotients. For a puncture these were proved and identified
in Corollary~\ref{D:fixed-local-quotient-alternating-count}
and Corollary~\ref{D:alternating-dimensional-normalization-closure}.
For a local pair write
\[
  \Lambda=Q(R,D),\quad
  U_R=\operatorname{Spec}R\setminus D,\quad
  G_R=\pi_1(U_R),\quad G_R\twoheadrightarrow\Lambda.
                                                        \tag{D.OBM.7}
\]
The degree-one inflation map for a quotient is injective
with trivial coefficients. Theorem~\ref{D:local-tame-p-linear-H1}
therefore supplies the required finite bound in
coefficient characteristic $p$ directly for $\Lambda$.
For every prime $\ell\ne p$, use the fixed SNC resolution
and notation $v,g_i,e,h$ from
Lemma~\ref{D:local-divisor-tame-central-H2}.
Its exact formula (D.LDH.2) and the degree-one Leray
sequence give
\[
 \begin{split}
 \dim_{\mathbb F_\ell}H^1(\Lambda,\mathbb F_\ell)
 &\le\dim_{\mathbb F_\ell}H^1(U_R,\mathbb F_\ell)\\
 &\le 2\sum_i g_i+e+v+h.                                
 \end{split}                                             \tag{D.OBM.8}
\]
More explicitly, proper strict-local base change identifies
$H^1(Y,\mathbb F_\ell)$ with the cohomology of the reduced
exceptional fibre; the normalization/node sequence bounds
its dimension by $2\sum_i g_i+e$. The $H^0$ of the first
direct image of the puncture inclusion contributes at
most one dimension for each exceptional or horizontal
component, hence at most $v+h$. Each horizontal component
has one exceptional intersection, by
Lemma~\ref{D:horizontal-residue-order-at-node}.
The purity input is exactly
\citet[Expose~XVI, Theorem~3.1.1 and Corollary~3.1.3,
pp.~242--243]{A_ILO2014}, already used in (D.LDH.2), with
no finite-type-over-the-coefficient-field restriction.

The $p$-bound and (D.OBM.8) give one fixed bound on
$H^1(\Lambda,\mathbb F_\ell)$ for all primes. The finite
Frattini argument and coprime-product argument written
in the proof of
Corollary~\ref{D:fixed-local-quotient-alternating-count}
then show that the maximal pronilpotent quotient of
$\Lambda$ has a fixed finite generating set. This
argument assumes no prior finite generation and makes
no assertion about its maximal prosolvable quotient.

The maximal pro-prime-to-$p$ quotient is finitely
generated as well. The complete local open $U_R$ is
connected because $R$ is a domain. It satisfies exactly
\S9.1 of \citet[Theorem~9.2, pp.~122--123]
{C_GrothendieckMurre1971PrimeToP}: a connected complement
of a closed set in the spectrum of a complete
noetherian local ring of dimension two with algebraically
closed residue field. Their known theorem gives a
finite presentation, in the pro-prime-to-$p$ category,
of the maximal prime-to-$p$ quotient of $G_R$; its
definition includes a finite generating set.
The maximal prime-to-$p$ quotient of $\Lambda$ is a
quotient of that group, so is finitely generated.
Finite actions factor through the two chosen
quotients precisely when their images are nilpotent
or of order prime to $p$, respectively.

It remains to perform the global count. Enumerate the
tuples of full local permutation types by the proof
of Theorem~\ref{D:alternating-compressed-local-orbits},
using (D.OBM.4)--(D.OBM.6) for the present sources.
The fixed quotient generator counts, finite model
counts, ordered orbit compositions and the residual
square sum (D.OBM.2) give at most $n^{A_0n}$ tuples.
All counts apply to the whole $\Lambda_x$-orbits;
no decomposition of a smaller inertia subgroup is
substituted for one of these orbits.

Let $C^\circ=C\setminus P$ be the fixed detecting
curve. For every actual tuple, the highest-break
constants of
Lemma~\ref{D:tame-detecting-curve-local-pairs}
give total discriminant at most $\lambda n$, with
$\lambda$ fixed. Different transitivity and additivity
over the local factors justify this also for
nongalois and disconnected covers, exactly as in
Theorem~\ref{D:finite-puncture-local-orbit-factorial-count}.
For each tuple, the fixed-formal-fibre theorem and
the disconnected factor assignment bound give at
most
\[
              n^{|P|n+1}(n!)^{2g(C)}n^{2\lambda n}
                                                        \tag{D.OBM.9}
\]
curve covers. This is a fixed factorial-order bound.
The surjection $\pi_1(C^\circ)\twoheadrightarrow Q$
is injective on isomorphism classes of finite
$Q$-sets, even disconnected ones. Multiplying
(D.OBM.9) by the full-type bound proves $n^{An}$
for $n\ge2$; degree one is the trivial cover.
Finally the same natural-permutation kernel
injection, $|A_n|\ge n^{n/8}$, and existing
LMT/Goursat/compactness argument prove the
selected alternating conclusions.
\end{proof}

\begin{remark}
This is the original-divisorial and controlled-dimensional
closure of the same restricted alternating reduction.
Its extra geometric scope concerns normal proper pairs.
It does not assert a finite-normalization descent theorem
for arbitrary nonnormal divisorial tame pairs, an
unrestricted factorial bound, a hereditary constant
linear in finite-cover degree, or finite generation of
the whole $Q(S,T)$.
\end{remark}


% Proved: intrinsic recognition; four complete actual-file audits.
\subsection{Intrinsic small linear images of local permutation orbits}

\begin{lemma}[Proved: subgroup and point-stabilizer enumeration for small linear images]
\label{D:intrinsic-small-linear-permutation-images}
Fix a prime $p$ and a profinite group $\Lambda$ for which
\[
 |\operatorname{Hom}_{\rm cont}
       (\Lambda,\operatorname{GL}_D(\mathbb F_p))|
                      \le p^{aD^2}\qquad(D\ge1)          \tag{D.IM.1}
\]
with one constant $a$. For integers $N\ge2$ and $D\ge1$,
let $\mathcal T_{N,D}(\Lambda)$ be the isomorphism classes
of transitive $N$-element continuous $\Lambda$-sets whose
finite permutation image admits a faithful representation
over $\mathbb F_p$ of dimension at most $D$. Then
\[
 \log\max\{1,|\mathcal T_{N,D}(\Lambda)|\}
 \le \frac{2(\log p)^2}{\log2}D^4
                       +(a+2)(\log p)D^2.               \tag{D.IM.2}
\]
No linear representation or identification of the
permutation image is required as part of the given set.

More explicitly, there is a fixed collection of faithful
degree-$N$ permutation models, each embedded into
$\operatorname{GL}_D(\mathbb F_p)$, that contains every
action counted here. The logarithm of the maximum of
one and its number is at most
\[
            \frac{2(\log p)^2}{\log2}D^4
                              +2(\log p)D^2.             \tag{D.IM.3}
\]
Consequently, if positive integers $D_N$ satisfy
$D_N^4\le B N\log(N+1)$ for a fixed $B$, these
collections satisfy the hypotheses (D.MO.1)--(D.MO.2)
of Lemma~\ref{D:permutation-matrix-model-count}, after
enlarging its constant by an amount depending only
on $p$ and $B$.
\end{lemma}

\begin{proof}
Put $G=\operatorname{GL}_D(\mathbb F_p)$, $M=|G|$, and
$l=\lceil\log_2\max\{2,M\}\rceil$. Every subgroup
of a group of order at most $M$ has a generating tuple
of length at most $l$. To see this without assuming
any bound on its minimal generators, start with the
trivial subgroup and repeatedly adjoin an element
outside the subgroup already generated. Each proper
enlargement multiplies its order by at least two.
Padding the resulting tuple by identities gives
length exactly $l$.

Thus $G$ has at most $M^l$ subgroups. For each
$H\le G$, the same argument bounds the number of
subgroups $K\le H$ by $M^l$. Retain all ordered
pairs
\[
        (H,K),\qquad H\le G,\quad K\le H,\quad
        [H:K]=N,\quad \bigcap_{h\in H}hKh^{-1}=1.
                                                        \tag{D.IM.4}
\]
There are at most $M^{2l}$ such pairs. For each,
choose once a numbering of its left cosets $H/K$.
Left multiplication gives a faithful model
$H\hookrightarrow S_N$, and its other embedding
is the inclusion $H\hookrightarrow G$. These are
the promised fixed models.

Indeed, given a transitive set in the assertion,
write $H_0\le S_N$ for its actual permutation
image and $K_0$ for the stabilizer of a point.
The core of $K_0$ in $H_0$ is trivial because
$H_0$ acts faithfully. Choose any faithful
$\mathbb F_p$ representation of $H_0$ of
dimension at most $D$, and pad by a trivial
summand if necessary. Its images $H,K$ of
$H_0,K_0$ satisfy (D.IM.4). The original set
is isomorphic to the coset action of $H$ on
$H/K$, with the homomorphism $\Lambda\to H$
induced by that chosen representation.

In particular, the disjoint union over (D.IM.4)
of the actual sets
$\operatorname{Hom}_{\rm cont}(\Lambda,H)$
surjects onto $\mathcal T_{N,D}(\Lambda)$
after retaining its transitive actions. Its
cardinality is at most
\[
                  M^{2l}p^{aD^2},                       \tag{D.IM.5}
\]
by (D.IM.1) and the inclusion $H\le G$.
Choices of a faithful representation, point
stabilizer, coset numbering, or isomorphism
of the image are already covered by this
surjection; none needs an additional factor.
No injection of permutation isomorphism
classes into linear conjugacy classes is used.

If $M=1$, the collection (D.IM.4) is empty
for $N\ge2$, and the assertions follow.
If $M\ge2$, then $l\le\log M/\log2+1$
and $\log M\le D^2\log p$. Hence
\[
       2l\log M
       \le \frac{2}{\log2}(\log M)^2+2\log M
       \le \frac{2(\log p)^2}{\log2}D^4
                                      +2(\log p)D^2.    \tag{D.IM.6}
\]
This proves (D.IM.3), and (D.IM.5) proves
(D.IM.2). Finally $D_N^2\le D_N^4$ for
positive integers, so (D.IM.3) and
$D_N^2\log p\le B(\log p)N\log(N+1)$
give both matrix-model hypotheses.
The argument is elementary and uses no
classification of finite or primitive groups.
\end{proof}

\begin{corollary}[Proved: uncharged orbits with intrinsically small defining-characteristic images]
\label{D:intrinsic-linear-orbit-alternating-count}
In either Theorem~\ref{D:alternating-compressed-local-orbits}
or Theorem~\ref{D:original-boundary-matrix-alternating-count},
fix a sequence of positive integers $D_N$ with
$D_N^4\le B N\log(N+1)$. Enlarge the admissible
whole local orbits by allowing every degree-$N$
orbit whose finite image admits a faithful
$\mathbb F_p$ representation of dimension at
most $D_N$. Then all conclusions of the
respective theorem remain valid, with the
square-sum condition imposed only on the
remaining orbits.

In particular, for any fixed constants $c>0$
and $b\ge0$, all local orbit images of
faithful $\mathbb F_p$ dimension at most
$c(1+\log(N+1))^b$ may be left uncharged,
regardless of the particular transitive
action of that image. This is an intrinsic
condition on the whole local permutation
image, rather than a requirement that its
action be one prescribed affine or flag
action.
\end{corollary}

\begin{proof}
For complete puncture sources, (D.IM.1)
follows by multiplying the proved local
linear conjugacy-class count by the order
of the ambient linear group. For local
original-divisorial tame sources, it follows
from (D.OBM.4)--(D.OBM.5), which first
bound actual projective Hom sets and then
use $A\mapsto[\operatorname{diag}(A,1)]$.
For pro-prime-to-$p$ cyclic or trivial
sources the actual Hom bound follows
directly from one or zero generators.
Thus every one of the finitely many fixed
sources has an exponent $a$ as in
(D.IM.1).

Apply Lemma~\ref{D:intrinsic-small-linear-permutation-images}
at each degree, and adjoin its models to
the already chosen ones. The logarithm
of the number of models and the matrix
cost still satisfy (D.MO.2), with a new
fixed constant. The two cited theorems
therefore give the same factorial cover
bound, polynomial selected $A_n$ kernel
multiplicity, and finite generation of
the selected alternating simple product.
The whole local-group orbit remains the
unit of this condition.

For the final assertion take
$D_N=\max\{1,\lceil c(1+\log(N+1))^b\rceil\}$.
The ratio
\[
                  \frac{D_N^4}{N\log(N+1)}              \tag{D.IM.7}
\]
is bounded for $N\ge2$: its numerator
is bounded by a fixed power of
$1+\log(N+1)$, while its denominator
grows exponentially in $\log N$.
Hence a constant $B$ exists, depending
only on $c,b$, as required.
\end{proof}

\begin{remark}
This is a recognition enlargement within
the same matrix-model alternating criterion.
It does not assert that arbitrary local
permutation images have a small faithful
representation. No primitive-group
classification, hereditary cover estimate,
or unrestricted alternating finite-generation
claim is needed or inferred.
\end{remark}


% Four-audit scope transfer: known algebraization input and proved complete-cover budgets.
\subsection{One henselian algebraization for all complete punctured covers}

\begin{lemma}[Known: fixed-base algebraization of the covering category]
\label{D:fixed-henselian-algebraization}
Let $(R,\mathfrak m,k)$ be a complete normal local surface ring
with algebraically closed coefficient and residue field $k$.
There is one normal strictly henselian local $k$-algebra $R^h$,
the henselization of a local ring essentially of finite type over
$k$, together with a fixed $k$-isomorphism
\[
                         \widehat {R^h}\simeq R.        \tag{D.AC.1}
\]
Completion induces an equivalence between the categories of finite
$R^h$-algebras etale away from the closed point and finite
$R$-algebras etale away from the closed point. It restricts to an
equivalence on normal local domains and preserves generic degree.
In particular the henselian base in this assertion is fixed before
any finite punctured cover is chosen.
\end{lemma}

\begin{proof}
The cited algebraization inputs are
\citet[Theorems~3.8--3.9 and~3.11, with Lemmas~3.12 and~3.17,
pp.~32--38]{Artin1969ApproximationD}. We explain the hypotheses
and the normal-domain restriction, rather than reprove those
known theorems.

Choose a power-series presentation $R=k[[X_1,\ldots,X_e]]/I$.
Over the perfect field $k$, the continuous Jacobian ideal has
support in the isolated singular locus. This is the same
formal Jacobian input used in the fixed-different construction
of Proposition~\ref{D:cover-uniform-canonical-defect}; its
continuous-differential interpretation is
\citet[Example~3.3, p.~7]{C_IyengarTakahashiStable}.
Thus the formal smoothness hypothesis outside the closed point
in Artin's Theorem~3.8 holds. That theorem gives an algebraic
local $k$-ring whose completion is $R$. Its residue field is
$k$, and its henselization $R^h$ has the same completion.
Faithfully flat descent of normality makes $R^h$ normal
\citep[Lemma~10.164.3, Tag~033G]{StacksCompleteBudgetD}.
Since $k$ is algebraically closed, $R^h$ is strictly henselian.
If $R$ is regular, the usual polynomial local ring and its
henselization provide the same conclusion directly.

Apply Artin's Theorem~3.11 to this fixed henselian ring and
its maximal ideal. It gives the asserted equivalence for
finite algebras, including all algebra homomorphisms. Let a
normal local cover domain $A$ over $R$ descend to a finite
algebra $A^h$ over $R^h$. The map
\[
              A^h\longrightarrow A^h\otimes_{R^h}R\simeq A
                                                               \tag{D.AC.2}
\]
is faithfully flat. Consequently $A^h$ is normal and is a
domain. A finite algebra over a henselian local ring is a
product of local rings; being a domain, $A^h$ is one local
ring. Its residue field is the finite extension $k$ of $k$.
The $\mathfrak m_{R^h}$-adic and maximal-ideal-adic topologies
on $A^h$ are cofinal, so (D.AC.2) identifies $A$ with its
local completion. Conversely, an excellent normal local
$A^h$ has normal local completion, hence a domain, and
completion preserves etaleness on the puncture. Generic
degree is preserved by flat base change: on the puncture
the finite algebra is locally free of its constant generic
rank, which is unchanged in (D.AC.2). This proves the
normal-domain and degree assertions.
\end{proof}

\begin{proposition}[Proved: degree-linear geometric budgets for all complete covers]
\label{D:all-complete-cover-geometric-budgets}
Let $R$ be as in Lemma~\ref{D:fixed-henselian-algebraization},
in characteristic $p>0$. For every finite normal local
domain $A/R$ of degree $t$ that is etale on the puncture,
the resolution-size, canonical-energy and integral-cycle
bounds of
Proposition~\ref{D:punctured-cover-linear-resolution},
Corollary~\ref{D:punctured-cover-linear-energy}, and
Corollary~\ref{D:punctured-cover-antinef-cycle}
hold with constants depending only on $R$ and fixed base
choices. No algebraizability hypothesis on the individual
complete cover needs to be added.

More explicitly, there is a regular projective resolution
with the degree-linear size bounds of the first cited
proposition. On the minimal regular resolution, with the
notation of the energy corollary,
\[
 v+\sum_i(-E_i^2)+\sum_i p_a(E_i)+p_g(A)
                +\kappa^{\mathsf t}N^{-1}\kappa=O(t).
                                                               \tag{D.AC.3}
\]
The empty exceptional set is interpreted as contributing
zero to the exceptional sums. There is also a full-support
anti-nef generation cycle $Z_0$ on that minimal resolution
such that
\[
 -Z_0\cdot B\ge2p_a(B)\quad(0<B\text{ exceptional}),
 \qquad -Z_0^2+|K\cdot Z_0|=O(t),                    \tag{D.AC.4}
\]
with the already specified fundamental-cycle treatment
of wholly zero-canonical configurations. When the minimal
exceptional set is empty take $Z_0=0$.

Consequently the additional energy hypothesis in
Proposition~\ref{D:hereditary-Cartier-generation-budget}
is satisfied for every such complete cover. On its
Cartier model one obtains the actual cycles
$W=\rho^*Z_0+C\ge C$ and all the energy, generation and
fixed-determinant finite-cycle chart bounds stated there,
uniformly in the degree $t$.
\end{proposition}

\begin{proof}
Fix the single $R^h$ of Lemma~\ref{D:fixed-henselian-algebraization}.
Every $A$ descends to one normal local $A^h/R^h$ of the
same degree. The strict-henselian clause of
Proposition~\ref{D:punctured-cover-linear-resolution}
already has constants determined by one fixed base
resolution and exceptional cycle. In particular its
constants do not change with the etale neighbourhood to
which a particular cover descends. Apply it to all the
$A^h$ over this fixed $R^h$.

We check precisely what completion preserves. Let
$Y^h\to\operatorname{Spec}A^h$ be any regular projective
resolution that is an isomorphism on the puncture and put
\[
                   Y=Y^h\times_{A^h}\operatorname{Spec}A.
                                                               \tag{D.AC.5}
\]
The complete local rings at corresponding points of the
closed fibres agree. In particular, a proper finite-type
scheme over a Noetherian local ring is regular if and
only if its base change to the completion is regular:
\citet[Lemmas~54.11.1--54.11.2, Tag~0BG4]
{StacksCompleteBudgetD}. Hence (D.AC.5) is a regular
projective resolution. It remains an isomorphism off the
closed fibre. The closed fibres, with their scheme
structures over $k$, are literally the same.

Each exceptional integral curve and each effective
exceptional cycle therefore corresponds to the same
proper $k$-scheme on both sides. To see the assertion
also for a thickened cycle, its structure sheaf is killed
by some power of $\mathfrak m_{A^h}$; the quotients of
$A^h$ and $A$ by that power agree. The normal line
bundles of the corresponding Cartier curves agree as
well. Thus all intersection numbers, arithmetic genera,
and cohomology of exceptional cycles are unchanged.
Canonical intersections are unchanged by adjunction,
$K\cdot E_i=2p_a(E_i)-2-E_i^2$; no assertion about
ordinary absolute differential forms after completion
is required.

Flat coherent base change gives
\[
 H^1(Y^h,\mathcal O_{Y^h})\otimes_{A^h}A
                    \simeq H^1(Y,\mathcal O_Y)
                                                               \tag{D.AC.6}
\]
by \citet[Lemma~30.5.2, Tag~02KH]{StacksCompleteBudgetD}.
The module on the left before tensoring has finite
length, so its $k$-dimension is unchanged. This proves
preservation of $p_g$. The same closed-fibre identification
shows that minimality of a regular resolution is
preserved: the exceptional smooth rational $(-1)$-curves
are exactly the same curves before and after completion.
One may equivalently complete a minimal resolution of
$A^h$. Therefore every entry of the intersection matrix,
the canonical vector and its energy is unchanged.

Apply the established strict-henselian size and energy
theorems and complete their resolutions. This proves
(D.AC.3) and the stated size assertion. The integral
anti-nef cycle construction and uniform-generation
criterion depend only on the same intersections,
arithmetic genera, and complete-cycle inequalities.
Their fixed sufficiently large choice of the core
positivity parameter yields (D.AC.4), with wholly
zero-canonical components treated by the already
proved rational-double-point fundamental cycle.
All constants remain those for the one fixed base
$R^h$. Finally apply
Proposition~\ref{D:hereditary-Cartier-generation-budget}.
This proves its conclusions for all the present
complete covers without altering its proof or silently
enlarging any earlier theorem's stated hypotheses.
\end{proof}

\begin{remark}[Scope of this transfer]
Artin's algebraization and finite-algebra comparison are
known inputs. The conclusion removes the former scope
qualification on the geometric budgets by using one
fixed henselian base for all covers. It does not produce
an arbitrary-family module or twisted-bundle atlas, a
uniform unit-Isom presentation, or a hereditary
$q^{Ctr^2}$ representation bound.
\end{remark}


% Separate post-CP24 source. Four complete audits passed. Ordinary unit modules only; no hereditary representation count.
\subsection{Ordinary-unit determinant lists and whole-cycle Hom budgets}

\begin{lemma}[Proved: degree-linear determinant lists for ordinary unit cycles]
\label{D:hereditary-ordinary-unit-determinants}
Fix a complete normal local surface $(R,\mathfrak m,k)$ over an
algebraically closed coefficient field $k$ of characteristic $p>0$.
For every finite normal local cover $A/R$ of degree $t$, etale on
the puncture, one can choose a projective regular resolution
$\pi:Y\to\operatorname{Spec}A$, and full-support effective
anti-nef Cartier cycles $C\le W$, with the following properties.
The choices depend on $A$ and fixed base data, but not on a
representation. The number $v$ of exceptional prime curves is
$O(t)$. The actual Cartier powers and all bounds of
Proposition~\ref{D:hereditary-Cartier-generation-budget} hold on
this same model.

Write $U_A=\operatorname{Spec}A\setminus\{\mathfrak m_A\}$,
let $q=p^f$, and let $\mathcal E$ be the vector bundle with
ordinary $q$-unit structure associated to any continuous
$\mathbb F_q$ representation of $\pi_1(U_A)$ of rank $r\ge1$.
For $0\le a<f$, let
\[
 M_a=j_*(F_{U_A}^{a*}\mathcal E),\qquad
 F_a=(\pi^*M_a)^{\vee\vee},\qquad
 V_{a,b}=F_a(-W)|_{bC},\quad L_{a,b}=\det V_{a,b},
                                                        \tag{D.UD.1}
\]
where $b\ge1$ is fixed and $j:U_A\hookrightarrow\operatorname{Spec}A$.
Here $F_{U_A}$ denotes the absolute $p$-Frobenius, and
$F_{U_A}^{a*}$ is its $a$-fold pullback.
These are the canonical reflexive extensions and their full
pullback bundles. The bundles $V_{a,b}$ and lines $L_{a,b}$ are
globally generated, and
\[
 h^0(\mathcal O_{bC})+h^1(\mathcal O_{bC})=O_b(t),\qquad
 \deg L_{a,b}+h^0(L_{a,b})=O_b(tr).                    \tag{D.UD.2}
\]
There are constants $c_0,c_1$, depending only on the fixed base
data, such that the number of ordered tuples of isomorphism
classes of the lines $(L_{a,b})_{0\le a<f}$ arising in this way
is at most
\[
 q^{c_0t}\exp\!\bigl(c_1ft\log(r+1)\bigr)
                       \le p^{c_2ftr^2}.               \tag{D.UD.3}
\]
The same bound covers any fixed finite list of thicknesses $b$:
once the determinant lines on $Y$ are chosen, their restrictions
at all those levels are determined. This statement enumerates
the determinants of geometric ordinary unit objects. It does
not assert that an arbitrary infinitesimal deformation of a
line or bundle has a constant determinant.
\end{lemma}

\begin{proof}
\emph{One model with both Cartier powers and linearly many components.}
Use Lemma~\ref{D:fixed-henselian-algebraization} to fix a single
henselian algebraization of $R$. The fixed $\mathfrak m$-primary
Noether-different ideal $J$ used in the Cartier-jet construction
descends to it: any ideal containing a power of the maximal
ideal is determined in the corresponding finite quotient.
Choose a separating parameter pair $x,y$ in that descended
ideal. Such a choice may be made in the henselian model.
Indeed, an $\mathfrak m$-primary ideal contains sufficiently high
powers of every local coordinate. Enlarge a finite list of
ideal generators by sufficiently high consecutive powers of
coordinates in a fixed algebraic neighbourhood; their ratios
recover those coordinates in the fraction field. Their generic
differentials span the two-dimensional differential space,
also after its separable henselian extension. Over the infinite
perfect field, two general linear combinations of this enlarged
finite list are both a system of parameters and separating.
Parameter avoidance and independence
of two generic differentials are nonempty open conditions.
After completion this is a separating finite normalization
$k[[x,y]]\subset R$, with $(x,y)\subset J$.

Resolve the one fixed henselian base and principalize $(x,y)$
there. Denote the resulting fixed regular model by $Y_0$ and
write $(x,y)\mathcal O_{Y_0}=\mathcal O_{Y_0}(-C_0)$.
If necessary include the blowup of the closed point, so its
exceptional set is nonempty. Fix also a strictly anti-nef positive
cycle on that exceptional set. Apply
Proposition~\ref{D:punctured-cover-linear-resolution} on this
fixed base model, and then the completion comparison in
Proposition~\ref{D:all-complete-cover-geometric-budgets}.
For every $A$ this yields a regular model $Y$ dominating the
finite normalization over $Y_0$, with $v=O(t)$. It dominates
the minimal resolution of $A$.

Set $C=f^*C_0$ on $Y$ and let
$W=\rho^*Z_0+C$ be the cycle in
Proposition~\ref{D:hereditary-Cartier-generation-budget},
where $\rho$ is the morphism to the minimal resolution and
$Z_0$ is its generation cycle. On this actual model
\[
 (x^b,y^b)\mathcal O_Y=\mathcal O_Y(-bC),\qquad
 -C^2=O(t),\qquad -W^2+|K_Y\cdot W|=O(t).             \tag{D.UD.4}
\]
The model is not subsequently replaced by arbitrary blowups.
Thus the component bound and Cartier statements are simultaneous.
The cited Cartier-generation proposition gives (D.UD.2).

\emph{The number of punctured scalar characters.}
For every integer $n$ prime to $p$, we first prove
\[
       |H^1(U_A,\boldsymbol\mu_n)|
                         \le n^{2p_g(A)+v}.             \tag{D.UD.5}
\]
For $n=1$ the assertion is immediate. Put $E=Y_{\rm exc,red}$.
Remove from $Y$ the finite set of singular points of $E$ and
intersections of distinct components, obtaining $Y^\circ$.
Its boundary is a disjoint union of smooth connected open
curves $E_i^\circ$. Purity in codimension two identifies
$H^1(Y^\circ,\boldsymbol\mu_n)$ with
$H^1(Y,\boldsymbol\mu_n)$. Codimension-one purity and the
degree-one localization sequence give
\[
 0\longrightarrow H^1(Y,\boldsymbol\mu_n)
 \longrightarrow H^1(U_A,\boldsymbol\mu_n)
 \longrightarrow\bigoplus_{i=1}^{v}
                       H^0(E_i^\circ,\mathbb Z/n).
                                                        \tag{D.UD.6}
\]
Only the last term's cardinality, at most $n^v$, is needed.
The exact regular-scheme purity input is
\citet[Expose XVI, Theorem 3.1.1 and Corollary 3.1.3]
{A_ILO2014}; its scope includes these excellent resolutions.

Proper strict-local base change, and invariance under nilpotents,
identify $H^1(Y,\boldsymbol\mu_n)$ with
$H^1(E,\boldsymbol\mu_n)$
\citep[Theorem 59.91.11, Tag 095T]{StacksSemilinearDeterminantsB}.
For completeness, normalize $E$. If the normalized components
have genera $g_i$, and $\Gamma$ is the finite incidence graph
having a vertex for each component and each singular point and
an edge for each branch above a singular point, the normalization
sequence for constant etale sheaves gives
\[
 |H^1(E,\mathbb Z/n)|\le
                       n^{2\sum_i g_i+b_1(\Gamma)}.     \tag{D.UD.7}
\]
The graph contribution follows directly from the difference map
on branches. For each smooth proper normalized component the
bound $n^{2g_i}$ follows from the full $2g_i$ generator theorem
\citep[Expose X, Theorem 2.6 and its preceding argument]
{C_SGA1}. No assertion about the order of a coefficient monodromy
group is involved; the coefficients here are constant.

If a singular point has $s$ branches, its normalization defect
is at least $s-1$, since the quotient already distinguishes the
$s$ residual constants modulo their common diagonal. The
coherent normalization sequence consequently gives
\[
       2\sum_i g_i+b_1(\Gamma)
                              \le2h^1(E,\mathcal O_E).
\]
The map $H^1(Y,\mathcal O_Y)\to H^1(E,\mathcal O_E)$ is
surjective: the next coherent cohomology group of
$\mathcal O_Y(-E)$ is zero, since the proper morphism to the
affine surface has fibres of dimension at most one. Hence
$h^1(E,\mathcal O_E)\le p_g(A)$. Equations (D.UD.6)--(D.UD.7)
prove (D.UD.5). By the established all-complete budgets,
$2p_g(A)+v\le c_0t$.

The determinant of an ordinary $\mathbb F_q$ local system is
one character with values in $\mathbb F_q^*$. Since $k$ contains
the $n=q-1$ roots of unity, these characters are counted by
(D.UD.5). Thus at most $q^{c_0t}$ initial determinant
characters occur. The punctured determinant at phase $a$ is
the $p^a$-Frobenius pullback of that initial determinant;
there is no independent scalar-character choice at every phase.

\emph{Integral exceptional multidegrees.}
Fix this initial scalar character and hence every punctured
phase determinant. Let
\[
 K_a=\bigl(\pi^*\det(M_a)\bigr)^{\vee\vee}
\]
be its fixed full rank-one pullback on the regular surface $Y$.
Here $\det(M_a)=(\bigwedge^r M_a)^{\vee\vee}$ is the reflexive
determinant.
Both $K_a$ and $\det F_a$ restrict to the same line on $U_A$.
Choosing a trivialization of their ratio there expresses it
as an integral exceptional divisor:
\[
                     \det F_a\simeq K_a(D_a),
             \qquad D_a=\sum_i z_{a,i}E_i,\quad z_{a,i}\in\mathbb Z.
                                                        \tag{D.UD.8}
\]
This uses only that a line bundle on a regular integral scheme
with a trivialization on a dense open is given by its Cartier
divisor, here supported on the exceptional fibre.

Put $d_{a,i}=\deg_{E_i}(\det F_a|_{E_i})$.
The evaluation $\pi^*M_a\to F_a$ is surjective at every
exceptional generic point: over its discrete valuation ring,
the torsion-free pullback equals its double dual. Its exterior
power therefore gives a nonzero section of $\det F_a|_{E_i}$,
so $d_{a,i}\ge0$. Write $C=\sum_i m_iE_i$, with every
$m_i\ge1$. The actual Frobenius determinant recurrence of the
Cartier-generation proposition gives
\[
 0\le c_1(F_a)\cdot C
       =\sum_i m_i d_{a,i}
       \le \frac{rp}{p-1}(-C^2)\le c_3tr.              \tag{D.UD.9}
\]
In particular $\sum_i d_{a,i}\le c_3tr$.

For fixed $K_a$, a multidegree vector determines $\det F_a$
uniquely. Indeed the difference of two possible divisors in
(D.UD.8) has zero intersection with every $E_i$, so is zero
by negative definiteness. There are therefore at most
\[
             \binom{\lfloor c_3tr\rfloor+v}{v}
                    \le \exp\bigl(c_4t\log(r+1)\bigr) \tag{D.UD.10}
\]
possibilities per phase. For example, increase a constant
$c_5$ so $v\le\lfloor c_5t\rfloor$, add zero coordinates,
and use $\binom{a+b}{b}\le(e(a+b)/b)^b$ with
$a=\lfloor c_3tr\rfloor$, $b=\lfloor c_5t\rfloor$.
The regular case and empty exceptional set cause no problem;
the chosen principalized model may be used throughout.

Finally $L_{a,b}=\det F_a|_{bC}\otimes\mathcal O_{bC}(-rW)$,
with the second factor fixed. Multiplying (D.UD.10) over
$f$ phases and then the scalar-character count proves the
first bound in (D.UD.3). The second follows from
$\log q=f\log p$, $r\ge1$, and $\log(r+1)=O(r^2)$.
All choices at further fixed thicknesses are restrictions of
the same lines on $Y$, so introduce no new determinant factor.
\end{proof}

\begin{proposition}[Proved: whole-Cartier-cycle object and Hom budgets]
\label{D:hereditary-ordinary-unit-finite-cycle-Hom}
Use the model and cycles of
Lemma~\ref{D:hereditary-ordinary-unit-determinants}, and fix $b\ge1$.
Put $D=bC$. For any fixed determinant lines $L_i$ occurring in
that lemma, with respective ranks $r_i\le r$, the following is
an actual finite-cycle presentation over every coefficient
$k$-algebra $T$.

There is an open subset of the affine space
\[
                         H^0(D,L_i)^{r_i+1}             \tag{D.UD.11}
\]
whose points specify exact locally split sequences
\[
 0\longrightarrow (L_i^{-1})_T
   \xrightarrow{s_i}\mathcal O_{D_T}^{r_i+1}
   \longrightarrow V_i\longrightarrow0,              \tag{D.UD.12}
\]
with $V_i$ locally free of rank $r_i$ and determinant $(L_i)_T$.
Every globally generated geometric bundle of rank $r_i$ and
determinant $L_i$ occurs. The object coordinates have number
$O_b(tr_i^2)$.

For each pair of these universal quotient bundles, their actual
derived Hom is represented by a fixed graded finite-dimensional
$k$-space tensored with $T$, with a differential polynomial of
degree at most two in the object coordinates. Its total dimension
is $O_b(tr_ir_j)$ and its degrees lie in $[-1,2]$.
It is functorial in all coefficient rings, without any assumption
that its cohomology is flat. In particular
\[
 h^0(D,\mathcal Hom(V_i,V_j))+
 h^1(D,\mathcal Hom(V_i,V_j))=O_b(tr_ir_j)              \tag{D.UD.13}
\]
at every geometric point. Actual arrows, equality of arrows,
and inverse relations have finite polynomial witnesses. For a
fixed pattern with $O(f)$ objects, arrows and bounded-length
relations, the combined number of coordinates and scalar
equations is $O_b(ftr^2)$, and their degrees are bounded
independently of $t,r,f$. The chosen determinant lines and all
linear splittings used below are fixed background data.

For the geometric canonical bundles in (D.UD.1), twisting both
arguments by $-W$ cancels in Hom. Thus (D.UD.13) is also the
actual $H^0+H^1$ bound for
$\mathcal Hom(F_a,F_{a'})|_{bC}$, not a bound inferred from
a sandwich of lattices on the surface.
\end{proposition}

\begin{proof}
The cycle $D$ is projective and Cohen--Macaulay of pure dimension
one, including its full nonreduced structure. Write
$b_D=h^0(\mathcal O_D)+h^1(\mathcal O_D)$ and
$d_i=\deg L_i$. Equation (D.UD.2) gives
$b_D=O_b(t)$, $d_i=O_b(tr_i)$.
A globally generated line has a section avoiding all associated
generic points. Its multiplication map is a nonzerodivisor,
and its zero scheme is finite of length its degree. The exact
section sequences imply
\[
 \begin{split}
 h^0(L_i)+h^1(L_i)&\le b_D+d_i,\\
 h^0(L_j^{-1})+h^1(L_j^{-1})&\le b_D+d_j,\\
 h^0(L_iL_j^{-1})+h^1(L_iL_j^{-1})&\le b_D+d_i+d_j.
 \end{split}                                           \tag{D.UD.14}
\]
For the last inequality, multiplication by a regular section of
$L_j$ embeds $L_iL_j^{-1}$ into $L_i$ with finite cokernel of
length $d_j$; its long exact cohomology sequence proves the
displayed bound. These arguments do not replace $D$ by its
reduction.

The generating-tuple locus in (D.UD.11) is open, since failure
to generate is a closed subset of the proper product with $D$.
At every local ring on this locus one tuple entry is a unit
after trivializing $L_i$. The injection in (D.UD.12) therefore
splits locally and its quotient is free. This proves exactness
and its preservation by arbitrary coefficient change, including
nonflat changes and nonreduced coefficient rings.
Conversely a globally generated geometric rank-$r_i$ bundle
on a curve is generated by $r_i+1$ general global sections:
the rank-drop incidence for $r_i+1$ sections has codimension
two at each point, while the curve has dimension one.
The kernel of the resulting surjection is a line with
determinant $L_i^{-1}$. Dualizing its inclusion gives a
generating tuple in (D.UD.11). This proves the covering
assertion on geometric bundles, with
$(r_i+1)h^0(L_i)=O_b(tr_i^2)$ coordinates.

Here is the actual derived calculation and its base-change
scope. Regard the two terms of (D.UD.12) as a complex
$P_i=[L_i^{-1}\to\mathcal O_D^{r_i+1}]$ in degrees $-1,0$.
The terms of $\mathcal Hom^\bullet(P_i,P_j)$ are
\[
 (L_j^{-1})^{\oplus(r_i+1)},\quad
 \mathcal O_D^{\oplus(r_i+1)(r_j+1)}\oplus L_iL_j^{-1},
 \quad L_i^{\oplus(r_j+1)}.                            \tag{D.UD.15}
\]
Their total cohomology dimension is at most
\[
 \begin{split}
 B_{ij}={}&b_D(r_i+1)(r_j+1)+(b_D+d_j)(r_i+1)\\
          &+(b_D+d_i)(r_j+1)+b_D+d_i+d_j
                  =O_b(tr_ir_j).
 \end{split}                                           \tag{D.UD.16}
\]

To obtain finite complexes, not just these numerical estimates,
use Proposition~\ref{D:fixed-determinant-finite-Hom}. We recall
the construction to specify exactly what commutes with base
change. Choose a two-affine cover of $D$ with affine intersection.
For each fixed line in (D.UD.15), split its two-term Cech complex
over $k$ onto its finite cohomology. Denote the inclusion,
projection and homotopy by $\iota,\rho,h$, respectively, so
$\iota\rho=1-dh-hd$, $\rho\iota=1$ and
$h^2=h\iota=\rho h=0$. These fixed $k$-linear identities
remain true after tensoring with every $T$.

In the total Cech Hom complex, the internal differential
$\delta$ is linear in $s_i,s_j$, preserves Cech degree,
and anticommutes with $d$. Since $h$ lowers Cech degree
from one to zero, $h\delta h=0$. The finite transferred
maps are exactly
\[
 I=\iota-h\delta\iota,\quad
 P=\rho-\rho\delta h,\quad
 \Delta=\rho\delta\iota-\rho\delta h\delta\iota.
                                                        \tag{D.UD.17}
\]
The identities proved in the cited proposition give a chain
homotopy equivalence with differential $d+\delta$ over every
$T$. Thus the finite complex computes
$R\operatorname{Hom}_{D_T}(V_i,V_j)$, has total size
(D.UD.16), and has differential of degree at most two.
The two-affine Cech and flat-over-$k$ tensor comparisons are
the exact inputs in
\citep[Lemmas 30.2.5--30.2.6 and 30.5.1--30.5.2]
{C_StacksFiniteCechHom}. No assertion of cohomological
flatness is needed or deduced.

For clarity, a degree-zero cycle represents an actual bundle
map. Equality of two such maps is equivalent to their
difference being $\Delta$ of a degree-minus-one witness.
Composition is obtained by applying $I$ to the two cycles,
multiplying in the Cech Hom complex, and applying $P$.
Its total polynomial degree is at most five. A path of any
fixed length is treated by first multiplying all its lifted
cycles and then projecting; equality is again one boundary
witness. This is an exact presentation of arrows and inverse
relations over arbitrary coefficient rings, not a test merely
on geometric fibres. Summing (D.UD.16) over the fixed diagram
pattern proves the asserted coordinate, equation and degree
bounds. At a geometric point, cohomology of this finite complex
has dimension at most its total size, giving (D.UD.13).
Finally, tensoring both arguments by the same invertible line
does not change their Hom bundle; apply this to $\mathcal O_D(-W)$.
\end{proof}

\begin{remark}[The remaining realization condition]
The determinant list covers geometric ordinary unit objects.
For each fixed line choice, (D.UD.12)--(D.UD.17) describe actual
finite-cycle quotient bundles and their actual arrows over all
coefficient rings. These are distinct assertions. They do not
show that every Artinian deformation of a canonical bundle
admits one of the chosen quotient presentations, that arbitrary
finite-cycle data are full pullbacks of flat maximal
Cohen--Macaulay modules, or that an arbitrarily prescribed
finite-cycle isomorphism lifts to the resolution with uniform
pole loss. In particular the coflat/split-boundary hypotheses
in the existing relative realization theorem have not been
removed. No twisted-PGL determinant assertion, hereditary
representation count, or finite-generation consequence is
claimed here.
\end{remark}


% Proved recognition; the primitive order theorem is a Known literature input.
\subsection{Primitive local images with defining-characteristic socles}

\begin{corollary}[Proved: primitive defining-characteristic socles have small intrinsic matrix models]
\label{D:primitive-defining-characteristic-orbit-models}
Fix the characteristic $p$ and a finite collection $\mathcal F$
of nonabelian finite simple groups. There is a constant
$c_{p,\mathcal F}$ with the following property. Let $G\le S_N$
be a faithful primitive permutation group, $N\ge2$, whose socle
is a direct product of nonabelian simple groups, each either
\begin{enumerate}
\item a defining-characteristic-$p$ simple Lie group in the
      families and finite-exception conventions of
      Corollaries~\ref{D:classical-defining-characteristic-crowns}
      and~\ref{D:exceptional-defining-characteristic-crowns},
      not isomorphic to an alternating group; or
\item a member of the fixed collection $\mathcal F$.
\end{enumerate}
Then $G$ has a faithful $\mathbb F_p$ representation of dimension
\[
                  D\le c_{p,\mathcal F}(1+\log N)^2.       \tag{D.PM.1}
\]
Consequently all such whole local permutation orbits may be
left uncharged in
Corollary~\ref{D:intrinsic-linear-orbit-alternating-count}.
For any one of its fixed local source groups, the number of
these degree-$N$ local types is at most
\[
                 \exp\bigl(C(1+\log(N+1))^8\bigr),         \tag{D.PM.2}
\]
and the corresponding global factorial and selected
alternating-product conclusions remain valid.

This recognition applies to every primitive action satisfying
the stated socle condition, not only its natural projective
action. The previously allowed affine actions in characteristic
$p$ can of course be retained simultaneously. The collection
$\mathcal F$ may contain any fixed finite set of alternating,
sporadic, or other simple groups.
\end{corollary}

\begin{proof}
We first extract a finite-group representation consequence of
the already proved automorphism recognitions. For every simple
Lie group $S$ in the stated families there is a subgroup
$\mathcal A_0\trianglelefteq\operatorname{Aut}(S)$, of index at
most six, and a faithful embedding
\[
       \mathcal A_0\hookrightarrow
          \operatorname{P}\Gamma\operatorname L_d
                        (\mathbb F_{p^f}),
       \qquad d^2 f\log p\le \kappa\log|S|,              \tag{D.PM.3}
\]
where $\kappa$ is one fixed constant. These are the actual
finite-group embeddings and order estimates (D.CA.7)--(D.CA.9)
and (D.EC.5)--(D.EC.7). They do not depend on the geometric
tameness hypothesis in the crown application. The fixed simple
derived exceptions are handled by the faithful permutation
representation of their automorphism group, as in those proofs.

There is a faithful $\mathbb F_p$-linear representation
\[
 \begin{split}
 \operatorname{P}\Gamma\operatorname L_d(\mathbb F_{p^f})
       &\longrightarrow
          \operatorname{GL}_{d^2f}(\mathbb F_p),\\
 ([g],\sigma)&:\ X\longmapsto g\,\sigma(X)\,g^{-1}
                    \quad\bigl(X\in M_d(\mathbb F_{p^f})\bigr).
 \end{split}                                             \tag{D.PM.4}
\]
Scalar representatives give the same map. If the map is the
identity, its restriction to the centre
$\mathbb F_{p^f}I_d$ makes $\sigma$ the identity. Then $g$
commutes with every matrix, so $g$ is scalar. This proves
faithfulness, including when $p$ divides the field-automorphism
order.

Induce the restriction of (D.PM.4) from $\mathcal A_0$ to
$\operatorname{Aut}(S)$. This gives a faithful representation
of dimension at most $6d^2f$ in every characteristic.
Indeed an element acting trivially on the induced module
fixes each coset summand, belongs in particular to
$\mathcal A_0$, and acts trivially on its identity-coset
summand; faithfulness there makes the element the identity.
Neither semisimplicity nor division by the index is needed.
For each of the finitely many groups in $\mathcal F$,
use its automorphism group's faithful permutation action on
the set $S$, realized by permutation matrices over
$\mathbb F_p$. After increasing a constant $c_0$, we have
therefore proved
\[
 \operatorname{Aut}(S)\hookrightarrow
             \operatorname{GL}_{d_S}(\mathbb F_p),
       \qquad d_S\le c_0\log|S|                           \tag{D.PM.5}
\]
for every permitted simple factor $S$. The constant depends
only on $p$ and $\mathcal F$, including the fixed exceptions.

Put $L=\operatorname{Soc}(G)=\prod_i S_i^{k_i}$, with the
$S_i$ pairwise nonisomorphic. It has trivial centre.
The centralizer $C_G(L)$ is normal in $G$ and intersects
$L$ in $Z(L)=1$. If nontrivial, it would contain a minimal
nontrivial normal subgroup of the finite group $G$, which
by the definition of the socle lies in $L$. This is a
contradiction. Hence conjugation embeds $G$ into
\[
             \operatorname{Aut}(L)
          =\prod_i\bigl(\operatorname{Aut}(S_i)
                                  \wr S_{k_i}\bigr).     \tag{D.PM.6}
\]
Here automorphisms permute the isomorphic simple direct
factors and preserve their isomorphism types. Acting by
(D.PM.5) on each direct summand, and by permutation of the
summands, realizes the displayed group faithfully on a
space of dimension
\[
             \sum_i k_i d_{S_i}
                \le c_0\log|L|\le c_0\log|G|.             \tag{D.PM.7}
\]

To relate the last order to the permutation degree, use the
Known primitive-order theorem of
\citet[Theorem~1.1; author version p.~2, complete written
proof in Section~2, pp.~4--8]{C_Maroti2002PrimitiveOrders}.
It leaves three cases: the bound
$|G|<N^{1+\lfloor\log_2N\rfloor}$; four fixed Mathieu
actions; or an alternating-product action
\[
      A_m^r\le G\le S_m\wr S_r,\qquad
                         N=\binom{m}{k}^{\,r}.           \tag{D.PM.8}
\]
This theorem, including its classification inputs, is used
as a known result, rather than reproved here.

The first case gives $\log|G|=O((1+\log N)^2)$ directly,
and the four fixed actions do also. In (D.PM.8), suppose
first $m\ge5$. The subgroup $A_m^r$ is normal in $G$.
The products of its factors in each $G$-orbit form minimal
normal subgroups of $G$: a nontrivial normal subgroup of
one such product is a product of simple factors, and
$G$-invariance forces a union of whole factor orbits.
Thus $A_m$ occurs as a simple factor of $L$. By our
hypothesis it must belong to $\mathcal F$. It follows
that $m$ is bounded by a constant $m_0$ depending only
on $\mathcal F$. Enlarge $m_0$ to at least four, also
covering any smaller $m$ in (D.PM.8).

Since $N\ge2$, the subset size in (D.PM.8) is proper
and nonempty. Consequently $N\ge2^r$ and
\[
        \log|G|\le r\log(m_0!)+\log(r!)
             \le r\log(m_0!)+r\log\max\{1,r\}
             =O_{\mathcal F}((1+\log N)^2).               \tag{D.PM.9}
\]
Together with (D.PM.7), this proves (D.PM.1).
We do not need to classify any additional coincidences
between finite simple Lie and alternating groups: the
unbounded Lie family in the hypothesis explicitly
excludes alternating isomorphism types, and any fixed
exceptions may be placed in $\mathcal F$.

Finally insert the integer ceiling of (D.PM.1) into
Lemma~\ref{D:intrinsic-small-linear-permutation-images}.
Its $D^4$ cost gives (D.PM.2). This is a permissible
fixed polylogarithmic dimension in
Corollary~\ref{D:intrinsic-linear-orbit-alternating-count},
which gives the global conclusions with the previously
stated geometric and residual-orbit conditions.
\end{proof}

\begin{remark}[Remaining local image types]
The statement does not cover arbitrary primitive images
with unbounded alternating socle, general
cross-characteristic simple socle, or arbitrary
imprimitive extensions. Nor does it infer a bound on
all local permutation actions from quadratic matrix
counts. It is a named recognition within the same
restricted alternating reduction, with no additional
movement assigned to its individual families.
\end{remark}


% Separate post-CP24 auxiliary extension. Four audits passed; no family-realization or representation-count claim.
\begin{corollary}[Proved: ordinary-unit budgets on the generation cycle itself]
\label{D:hereditary-generation-cycle-Hom-budget}
In the setting and on the same chosen model of
Lemma~\ref{D:hereditary-ordinary-unit-determinants}, fix $b\ge1$
and replace the finite cycle $bC$ by $bW$. Then
\[
 h^0(\mathcal O_{bW})+h^1(\mathcal O_{bW})=O_b(t),\qquad
 \deg\det(F_a(-W)|_{bW})+
 h^0\bigl(\det(F_a(-W)|_{bW})\bigr)=O_b(tr).
                                                        \tag{D.UW.1}
\]
The ordinary-unit determinant tuples on $bW$ satisfy the same
bound (D.UD.3). For each fixed tuple of determinant lines, the
actual generating-quotient and finite-Hom presentations of
Proposition~\ref{D:hereditary-ordinary-unit-finite-cycle-Hom}
therefore hold on $bW$, with $O_b(ftr^2)$ coordinates and
scalar equations of bounded degree over every coefficient
$k$-algebra.

These statements apply separately to any fixed finite list
of thickenings of $C$ and $W$, with one determinant enumeration
on $Y$. They do not assert realization of finite-cycle objects
or lifting of their prescribed isomorphisms.
\end{corollary}

\begin{proof}
Put $Q_C=-C^2$ and $Q_W=-W^2$. Both are $O(t)$, and
$|K_Y\cdot W|=O(t)$, by
Proposition~\ref{D:hereditary-Cartier-generation-budget}
and its all-complete application. The negative intersection
form on the exceptional lattice gives
\[
 0\le -C\cdot W\le\sqrt{Q_CQ_W}=O(t).                   \tag{D.UW.2}
\]
The first inequality uses anti-nefness of $C$ and effectiveness
of $W$; the second is the ordinary Cauchy--Schwarz inequality
for the positive definite negative intersection form.

Write $x_a=c_1(F_a)\cdot W$. The nonnegative exceptional
multidegrees established in the determinant lemma give
$x_a\ge0$. The actual Frobenius determinant divisors satisfy
\[
 c_1(F_{a+1})=p\,c_1(F_a)+B_a,\qquad 0\le B_a\le rpC,
                                                        \tag{D.UW.3}
\]
cyclically in the $f$ phases. Since $W$ is anti-nef, (D.UW.3)
implies
\[
 p x_a-rp(-C\cdot W)\le x_{a+1}\le p x_a.
                                                        \tag{D.UW.4}
\]
Iterating the left inequality around the cycle and dividing
by $1-p^{-f}$ gives
\[
 x_a\le \frac{rp}{p-1}(-C\cdot W)=O(tr).                 \tag{D.UW.5}
\]
Consequently the globally generated bundle
$F_a(-W)|_{bW}$ has globally generated determinant of degree
\[
 b\bigl(c_1(F_a)\cdot W-rW^2\bigr)=O_b(tr).              \tag{D.UW.6}
\]

The exact sequence for the effective Cartier divisor $bW$,
together with the vanishing of degree-two coherent cohomology
on this proper surface over the affine local base with fibres
of dimension at most one, gives
\[
 h^1(bW,\mathcal O_{bW})\le
 h^1(Y,\mathcal O_Y)=p_g(A)=O(t).
                                                        \tag{D.UW.7}
\]
Adjunction for a possibly nonreduced exceptional Cartier cycle
gives
\[
 \chi(\mathcal O_{bW})
           =-\tfrac12\bigl(b^2W^2+bK_Y\cdot W\bigr)=O_b(t).
                                                        \tag{D.UW.8}
\]
Thus $h^0+h^1=\chi+2h^1=O_b(t)$. A regular section of
the generated determinant line in (D.UW.6) bounds its $h^0$
by $h^0(\mathcal O_{bW})$ plus its degree. This proves
(D.UW.1) on the full nonreduced cycle.

The proof of (D.UD.3) enumerates the determinant lines
$\det F_a$ on $Y$ before restricting them. Their restrictions
to $bW$, followed by the fixed twist $\mathcal O_{bW}(-rW)$,
therefore introduce no further choice. Finally $bW$ is again
a projective pure one-dimensional Cohen--Macaulay scheme.
The proof of the quotient/Hom proposition uses only this
property, global generation, and the scalar and line-degree
bounds (D.UW.1). Applying its same locally split quotient
sequences and fixed Cech contractions proves all the asserted
all-coefficient finite presentations. No cohomological
flatness, full-module realization, or prescribed surface-map
lifting is inferred from these finite complexes.
\end{proof}


% Separate post-CP24 proved auxiliary; full root/B/C/D audits. Retained coflat/split-boundary hypotheses; no counting claim.
\subsection{Relative realization on weak cycles and the uniform Cartier model}

\begin{proposition}[Proved: weak-cycle full-neighbourhood realization]
\label{D:weak-cycle-Artinian-full-realization}
In Theorem~\ref{D:Artinian-full-neighbourhood-realization}, replace
the hypothesis that $Z$ is a reduction cycle by the following
weaker conditions: $Z$ is an effective Cartier divisor of full
exceptional support, $\mathcal O_Z(-Z)$ is generically generated,
and
\[
                     H^1(Z,\mathcal O_Z(-Z))=0.
                                                        \tag{D.WR.1}
\]
Retain exactly its three hypotheses on the actual bundle $V_2$
on $(2Z)_T$: generic generation of the closed first restriction,
cohomological freeness and all Artinian base changes for $V$
and $V(Z)$, and the split injectivity of its actual boundary
\[
 H^0(Z_T,V(Z))\longrightarrow H^1(Z_T,V).
                                                        \tag{D.WR.2}
\]
Then all conclusions of that theorem hold. In particular every
vector-bundle extension $\mathcal F_T$ is the full pullback of
the $T$-flat finite module $M_T=H^0(Y_T,\mathcal F_T)$, its
closed fibre is maximal Cohen--Macaulay, and its section module
has relative Hartogs and the asserted Artinian base changes.
The extension may also be required to retain any already given
bundle on $(dZ)_T$, for a fixed $d\ge2$, whose restriction to
$2Z$ satisfies these conditions. No generation assumption on
$\omega_Z^{-1}$ is needed.

When the fixed Noether-different data of
Corollary~\ref{D:Artinian-canonical-two-level-jets} are available,
that corollary applies to these realized families as well,
with the same integer loss and the same qualification about
the raw cohomology spaces and images.
\end{proposition}

\begin{proof}
We check every use of the cycle hypothesis in the existing
realization argument. All other hypotheses, including the
complete excellent local base and the finite local Artinian
coefficient algebra $T$, remain unchanged.

The curve $Z$ is Cohen--Macaulay of pure dimension one. Generic
generation of $\mathcal O_Z(-Z)$ gives a section regular at
every associated generic point. Its multiplication is injective
with zero-dimensional cokernel. Tensoring by powers of
$\mathcal O_Z(-Z)$, and using (D.WR.1), proves
\[
 H^1(Z,\mathcal O_Z(-jZ))=0\quad(j\ge1).
                                                        \tag{D.WR.3}
\]
The powers are generically generated too. The successive layers
of $\mathcal O_Y(-jZ)$ on the formal neighbourhood are these
lines with indices at least $j$. Formal functions therefore
give $H^1(Y,\mathcal O_Y(-jZ))=0$. Consequently the restriction
of global sections to $Z$ is surjective, and for every $j\ge1$
there is an actual map
\[
               \mathcal O_Y^{a_j}\longrightarrow
                                  \mathcal O_Y(-jZ)
                                                        \tag{D.WR.4}
\]
surjective at each exceptional generic point. Here the integer
$a_j$ is not used as a quantitative parameter.

Write $V_0=V\otimes_Tk$. Its generic generation similarly
gives an injection $\mathcal O_Z^r\to V_0$ with
zero-dimensional cokernel: choose $r$ general global sections,
which generate at all the finitely many associated generic
points. Tensoring with $\mathcal O_Z(-jZ)$ and applying
(D.WR.3) gives $H^1(Z,V_0(-jZ))=0$ for $j\ge1$.
For any vector-bundle extension $\mathcal F_0$ to $Y$, formal
functions now give
\[
                    H^1(Y,\mathcal F_0(-Z))=0.
                                                        \tag{D.WR.5}
\]
Thus the initial tail-vanishing step of the earlier theorem
uses only the present weak conditions. Artinian filtrations
give the same vanishing after every specified coefficient
change. The flat Cech argument and the retained cohomology
base-change hypotheses then give flatness and base change of
$M_T$, exactly as in (D.AR.6)--(D.AR.8).

For clarity, the Hartogs step does not assert that local
cohomology with support is zero. Dualizing (D.WR.4) on $Z$
gives an injection $\mathcal O_Z(jZ)\to\mathcal O_Z^{a_j}$.
It stays injective after coefficient extension over $k$.
Multiplication on the actual second neighbourhood compares
the boundary for $V((j+1)Z)$ with copies of (D.WR.2).
The map on its degree-zero source is injective, so every
boundary $\delta_j$ in (D.AR.9) is injective. The Cartier
local-cohomology calculation consequently gives
\[
 H^1_{Z_T}(Y_T,\mathcal F_T)=H^0(Z_T,V(Z)),\qquad
 H^1(Y_T,\mathcal F_T)=H^1(Z_T,V).
                                                        \tag{D.WR.6}
\]
The natural map between these groups is the injective
boundary (D.WR.2). The localization exact sequence therefore
gives $H^0(Y_T,\mathcal F_T)=H^0(Y_T\setminus Z_T,\mathcal F_T)$.
This is precisely the required Hartogs conclusion.

There is no further hidden full-bundle criterion requiring
$\omega_Z^{-1}$. On the closed fibre, finiteness and the
Hartogs equality show directly that $M_0$ is a reflexive,
hence maximal Cohen--Macaulay, module over the normal surface.
The evaluation is an isomorphism off the exceptional locus
and is generically onto it, by the lifted generic sections.
Its codimension-two cokernel and exceptional torsion kernel
give $\mathcal F_0=(\pi^*M_0)^{\vee\vee}$. Over $T$ the
same argument applies: $Y_T$ is Cohen--Macaulay, its fixed
Cartier equations are nonzerodivisors, and codimension-two
modules have vanishing $\operatorname{Ext}^0$ and
$\operatorname{Ext}^1$ into its structure sheaf. This is the
dualization argument already written in the earlier theorem,
and proves the relative full-pullback identification.

Finally, existence of an extension of the prescribed $dZ$
data uses only the vanishing of the degree-two vector-bundle
obstruction on the one-dimensional support, followed by
Grothendieck existence over the complete Noetherian Artinian
coefficient base. These are the exact previously cited inputs
\citep[Tags 08MF and 0885]{StacksRelativeKahnB}. Starting at
$d\ge2$ changes neither argument. The subsequent proof of
the two-level jet corollary uses only the resulting relative
maximal Cohen--Macaulay module, its Hartogs equality, full
evaluation and base changes. All have just been proved, so
the same source factorizations, scalar clearing and adjacent
cohomology transition apply without alteration.
\end{proof}

\begin{corollary}[Proved: conditional realization on the uniform generation cycle]
\label{D:hereditary-weak-cycle-relative-realization}
Fix the complete normal local surface $R$ and the data of
Proposition~\ref{D:hereditary-Cartier-generation-budget}.
For every finite normal punctured cover $A/R$ of degree $t$,
choose its model and cycles $C\le W$ as in that proposition
and Proposition~\ref{D:all-complete-cover-geometric-budgets}.
They may be chosen independently of every representation.
Then $W$ satisfies (D.WR.1) and the weak generic-generation
condition. Consequently any finite local Artinian bundle on
$(dW)_T$, $d\ge2$, satisfying the three retained hypotheses
on $(2W)_T$, realizes an actual relative maximal
Cohen--Macaulay module with all the conclusions of
Proposition~\ref{D:weak-cycle-Artinian-full-realization}.

For each fixed module-jet level $\ell$, choose the fixed
integer $b\ge2\ell-1+2c$ from the established scalar clearing
bound, and take $d\ge\max\{2,b+1\}$. The retained $(dW)_T$
bundle then determines its module jet as follows. Put
$Q_j(T)=H^0(jC_T,\mathcal F_T|_{jC_T})$ and
$B_b(T)=\operatorname{im}[Q_{b+1}(T)\to Q_b(T)]$.
The exact formula is
\[
 B_b(T)/(x^\ell,y^\ell)B_b(T)
       \simeq M_T/(x^\ell,y^\ell)M_T.                 \tag{D.WR.7}
\]
Here $b,c,d$ depend only on the fixed base and $\ell$, not on
$t$, the module rank, or the Artinian coefficient algebra.
The recovered quotient commutes with the stated coefficient
changes; the raw cohomology spaces and images need not do so.
\end{corollary}

\begin{proof}
The all-complete budget theorem supplies the cycle family
required by the Cartier-generation proposition. Apply
Theorem~\ref{D:reflexive-pullback-uniform-generation} to
$\mathcal O_Y$ and $W$. It gives global generation of
$\mathcal O_Y(-W)$ and $H^1(Y,\mathcal O_Y(-W))=0$.
Its restriction to $W$ is globally generated. From
\[
 0\longrightarrow\mathcal O_Y(-2W)
  \longrightarrow\mathcal O_Y(-W)
  \longrightarrow\mathcal O_W(-W)\longrightarrow0
\]
and $H^2(Y,\mathcal O_Y(-2W))=0$ we obtain
$H^1(W,\mathcal O_W(-W))=0$. This proves the weak-cycle
conditions on the actual possibly nonminimal Cartier model.
The proposition gives realization, and its jet conclusion
applies because $(x^j,y^j)\mathcal O_Y=\mathcal O_Y(-jC)$
for every $j$. Since $C\le W$ and $d\ge b+1$, both retained
Cartier cycles in (D.WR.7) are closed subschemes of $dW$.
Thus the exact relative two-level formula uses only the
stated finite neighbourhood and proves the assertion.
\end{proof}

\begin{remark}[The quantitative family condition remains]
This extension removes a strong-cycle assumption from an
already conditional realization theorem. It does not remove
cohomological freeness, its base changes, or the split-boundary
condition. Nor does it provide degree-linear coordinates for
those conditions, a lift of a prescribed finite-cycle map, or
uniformly bounded lifts of specified unit-isomorphisms over
nilpotent coefficient extensions. The finite-cycle object and
Hom budgets therefore remain separate from a hereditary
representation count.
\end{remark}


% Separate post-CP24 proof. Complete author, root, C and D audits PASS; actual arrows between already realized families.
\subsection{Intrinsic finite-cycle lifting of canonical Hom maps}

\begin{proposition}[Proved: fixed-loss lifting of intrinsic Hom jets]
\label{D:intrinsic-Hom-fixed-loss-jets}
Let $\pi:Y\to\operatorname{Spec}A$ be a projective regular
resolution of a complete normal local surface over an
algebraically closed coefficient field $k$. Let $C$ be a
nonzero effective exceptional Cartier divisor whose support is
the entire closed fibre, and suppose
two elements $x,y$ generate its invertible ideal:
\[
                 (x,y)\mathcal O_Y=\mathcal O_Y(-C).
                                                        \tag{D.HJ.1}
\]
Let $J$ be an ideal of the fixed base ring acting on $A$, with
$x,y\in J$ and $JH^1(Y,\mathcal O_Y)=0$.
Let $T$ be a finite-dimensional local Artinian $k$-algebra with
residue field $k$, and let $F_T,G_T$ be vector bundles on $Y_T$.
Assume that, for fixed positive integers $a,b$, each of
$x^a\operatorname{Id}_{F_T}$ and $y^a\operatorname{Id}_{F_T}$
factors through a finite free $\mathcal O_{Y_T}$-module,
and each of $x^b\operatorname{Id}_{G_T}$ and
$y^b\operatorname{Id}_{G_T}$ does also.

Put $\ell=a+b+1$ and
$H_T=\mathcal Hom(F_T,G_T)$. For every integer $d$ and every
$k\ge\max\{0,d\}$, the inclusion induces the zero map
\[
 H^1\bigl(Y_T,H_T((d-k-\ell)C)\bigr)
       \longrightarrow H^1\bigl(Y_T,H_T((d-k)C)\bigr).
                                                        \tag{D.HJ.2}
\]
Consequently, for $k\ge\max\{1,d\}$,
\[
 \begin{split}
 &\operatorname{im}\bigl[
 H^0((k+\ell)C_T,H_T(dC)|_{(k+\ell)C_T})
       \longrightarrow H^0(kC_T,H_T(dC)|_{kC_T})\bigr]\\
 &\hspace{12mm}=
 \operatorname{im}\bigl[
 H^0(Y_T,H_T(dC))
       \longrightarrow H^0(kC_T,H_T(dC)|_{kC_T})\bigr].
 \end{split}                                           \tag{D.HJ.3}
\]
Thus a map on the higher cycle has a global map with precisely
its prescribed lower-cycle restriction. The assertion holds
after every Artinian coefficient change, using the pulled-back
factorizations. It does not assert that either raw image
commutes with coefficient change.

When $a=b=1$ and $d=0$, the loss is three. If a given higher-cycle
map is an isomorphism, its global lift matching the lower cycle
is an isomorphism on $Y_T$.
\end{proposition}

\begin{proof}
At a local trivialization of $\mathcal O_Y(-C)$, write
$x=su$, $y=sv$ with $(u,v)=1$. Then $(u^j,v^j)=1$ for every
$j\ge1$: otherwise a prime would contain both $u$ and $v$.
Thus $x^j,y^j$ generate $\mathcal O_Y(-jC)$ and give the
locally split pencil sequence
\[
 0\longrightarrow\mathcal O_Y(jC)
 \longrightarrow\mathcal O_Y^{\,2}
 \xrightarrow{(x^j,y^j)}\mathcal O_Y(-jC)
 \longrightarrow0.                                    \tag{D.HJ.4}
\]
All coherent cohomology on $Y_T$ in degrees at least two
vanishes. Indeed properness gives finite cohomology modules
over the complete local base $A\otimes_kT$; formal functions
compute their completions on thickenings of the one-dimensional
closed fibre, where those degrees vanish. This is the same
exact fibre-dimension argument used in the relative realization
and Cartier-jet theorems.

The scalar base-change equality over the field $k$ shows
$JH^1(Y_T,\mathcal O_{Y_T})=0$. Apply (D.HJ.4) after base
change and use $H^2(\mathcal O_{Y_T}(jC))=0$. It gives a
surjection from two copies of $H^1(\mathcal O_{Y_T})$ onto
$H^1(\mathcal O_{Y_T}(-jC))$, so
\[
       JH^1(Y_T,\mathcal O_{Y_T}(-jC))=0
                              \quad(j\ge0).             \tag{D.HJ.5}
\]

Dualizing the assumed factorization on $F_T$ gives a
factorization of $x^a\operatorname{Id}_{F_T^\vee}$ through
a free bundle. Tensoring it with the factorization on $G_T$
therefore factors $x^{a+b}\operatorname{Id}_{H_T}$ through
a free bundle. The same holds for $y^{a+b}$. After tensoring
by $\mathcal O_{Y_T}((d-k)C)$, the induced cohomology maps
factor through sums of
$H^1(\mathcal O_{Y_T}((d-k)C))$. Since $k\ge d$, (D.HJ.5)
annihilates these groups by both $x$ and $y$. Hence
\[
 x^\ell H^1(Y_T,H_T((d-k)C))
 =y^\ell H^1(Y_T,H_T((d-k)C))=0.                         \tag{D.HJ.6}
\]

Tensor the $j=\ell$ sequence (D.HJ.4) by
$H_T((d-k)C)$. Its degree-one cohomology gives a surjection
\[
 H^1(Y_T,H_T((d-k)C))^{\,2}
       \longrightarrow H^1(Y_T,H_T((d-k-\ell)C)).
                                                        \tag{D.HJ.7}
\]
Composing (D.HJ.7) with the inclusion to the shallower
cohomology is multiplication by $x^\ell,y^\ell$ on its two
summands. Both vanish by (D.HJ.6), proving (D.HJ.2).

Now compare the restriction exact sequences for $kC_T$ and
$(k+\ell)C_T$ applied to the vector bundle $H_T(dC)$.
The connecting homomorphism of a section restricted from
the larger cycle is the image of its larger-cycle connecting
class under (D.HJ.2), hence is zero. The smaller-cycle exact
sequence then gives a global lift with that precise restriction.
Conversely every global section restricts through the larger
cycle. This proves (D.HJ.3).

All arguments use actual vector bundles, actual factorizations,
and exact sequences over $T$. Pulling them back to another
allowed Artinian algebra reproduces the proof; no cohomology
base-change inference about the displayed images is made.
Finally, if $d=0$ and the lower-cycle map is an isomorphism,
the determinant of its global lift is invertible at every
point of the closed fibre, since $C$ has full exceptional
support. Its noninvertible locus is closed and proper over
the local base and misses that fibre, so is empty. The
global map is therefore an isomorphism.
\end{proof}

\begin{corollary}[Proved: canonical and Frobenius arrow windows on arbitrary covers]
\label{D:hereditary-canonical-Hom-arrow-window}
Use the fixed Noether-different and Cartier data for the
finite punctured covers $A/R$ of the preceding hereditary
budget theorems. Let $F_T,G_T$ be the canonical full pullbacks
of already realized relative maximal Cohen--Macaulay modules,
for example the families satisfying the retained hypotheses
of Proposition~\ref{D:hereditary-weak-cycle-relative-realization}.
Then (D.HJ.3) holds with loss three for their ordinary Hom
maps. If their closed fibres are the ordinary-unit bundles
in the preceding determinant theorem and the families carry
its fixed-determinant quotient presentations at the retained
levels, then at each fixed pair $k,k+3$ the actual finite-cycle
Hom data have the $O_k(tr_Fr_G)$ size of the quotient-Hom
theorem. No such size assertion is made for arbitrary
maximal Cohen--Macaulay closed fibres.

For the one-step Frobenius Hom
$\mathcal Hom(F_{Y_T}^*F_T,G_T)$, the loss is $p+2$.
For the opposite bounded-pole Hom
$\mathcal Hom(G_T,F_{Y_T}^*F_T)(pC)$, the same loss applies
for $k\ge p$. Here $F_{Y_T}$ is the absolute $p$-Frobenius,
including its action on the Artinian coefficient ring.
These are actual prescribed lower-jet lifts of the indicated
arrows. No unit-compatibility equation or arbitrary object
realization is asserted by this corollary.
\end{corollary}

\begin{proof}
Proposition~\ref{D:cover-uniform-Noether-different} and its
relative source factorization give factorizations of
$x\operatorname{Id}$ and $y\operatorname{Id}$ on both
realized modules through finite free $A_T$-modules. Pull
back and apply double dual. Functoriality gives the same
factorizations on their canonical full bundles through
finite free $\mathcal O_{Y_T}$-modules. The fixed ideal
annihilates scalar $H^1(\mathcal O_Y)$ by
Proposition~\ref{D:cover-uniform-canonical-defect}.
The proposition therefore applies with $a=b=1$.
Its finite-cycle data are the actual Hom spaces computed
by Proposition~\ref{D:hereditary-ordinary-unit-finite-cycle-Hom};
the source modules of the factorization are not extra
matrix-valued coordinates in that presentation.

Absolute Frobenius pulls the first factorization back to
one for $x^p\operatorname{Id}_{F_{Y_T}^*F_T}$, and similarly
for $y^p$. Pullback preserves the free middle term and
the equality of the composite, regardless of flatness of
Frobenius on $T$. Thus the first Frobenius Hom has exponents
$a=p,b=1$, and the opposite one has $a=1,b=p$.
The twist $pC$ in the latter is exactly $d=p$ in
(D.HJ.2)--(D.HJ.3), explaining the condition $k\ge p$.
This proves the claimed uniform arrow windows without
assuming that their Frobenius or composition relations
lift automatically.
\end{proof}

\begin{remark}
The factorizations are available for already realized
canonical families. The proposition does not impose them
on arbitrary finite-cycle bundles, and does not remove
the coflat or split-boundary realization conditions.
Likewise a lift of an individual arrow is not asserted
to preserve independently lifted relations. These
distinctions are necessary in any later unit-cycle
representation count.
\end{remark}


% Separate post-CP25 proof; complete author/root/C/D file and source audits PASS. No hereditary representation count asserted.
\subsection{Reduction cycles comparable with one fixed parameter cycle}

\begin{lemma}[Proved: a strictly anti-nef separating parameter pencil]
\label{D:strict-Cartier-separating-parameters}
Let $(R,\mathfrak m,k)$ be a complete excellent normal local
surface with algebraically closed coefficient and residue field
$k$ of characteristic $p>0$. Let $Y_0\to\operatorname{Spec}R$
be a fixed projective regular resolution, an isomorphism off the
closed point, with nonempty exceptional fibre $E=\bigcup_iE_i$.
For every fixed $\mathfrak m$-primary ideal $J\subset R$ there
are a positive integral exceptional cycle $C_0$ and $x,y\in J$
such that
\[
 \begin{gathered}
 -C_0\cdot E_i>0\quad\hbox{for every }i,\qquad
 (x,y)\mathcal O_{Y_0}=\mathcal O_{Y_0}(-C_0),\\
 k[[x,y]]\subset R\quad\hbox{is finite and generically separable}.
 \end{gathered}                                             \tag{D.CR.1}
\]
For any morphism $f:Y\to Y_0$ for which the pullback of $C_0$
is Cartier, put $C=f^*C_0$. Then
\[
 (x^j,y^j)\mathcal O_Y=\mathcal O_Y(-jC)
                       \quad(j\ge1).                     \tag{D.CR.2}
\]
All choices in (D.CR.1) are made on the fixed base resolution.
\end{lemma}

\begin{proof}
Put $N=(-(E_i\cdot E_j))_{ij}$. This matrix is positive definite
with nonpositive off-diagonal entries. The unique real minimizer
of
\[
             \frac12v^{\mathsf t}Nv-\sum_i v_i
\]
is $N^{-1}{\bf1}$. Replacing $v$ by its coordinatewise absolute
value cannot increase this expression, so uniqueness makes
the minimizer nonnegative. No coordinate is zero: at such a
coordinate the equation $Nv={\bf1}$ would have nonpositive
left side. Since $N$ has integral entries, multiplication by
one positive integer gives a positive integral vector $d$
with $Nd>0$. Thus $D=\sum_i d_iE_i$ is positive and strictly
anti-nef.

The line bundle $\mathcal O_{Y_0}(-D)$ has positive degree on
every irreducible component of the closed fibre. The
componentwise curve criterion and the fibre criterion for
relative ampleness make it relatively ample, and hence ample
because the base is affine
\citep[Lemma~33.44.15, Tag 0B5Y, and Lemma~30.21.4,
Tag 0D2N]{StacksComparableCyclesB}. These statements include
the nonreduced closed fibre; alternatively first use invariance
of ampleness under nilpotent thickenings, Tag 09MW.

Write
\[
 I_a=H^0(Y_0,\mathcal O_{Y_0}(-aD))\subset
      H^0(Y_0,\mathcal O_{Y_0})=R.                         \tag{D.CR.3}
\]
Every $I_a$ is $\mathfrak m$-primary: $R/I_a$ embeds into the
finite-length module $H^0(aD,\mathcal O_{aD})$, and the
positive exceptional valuation excludes $1$. The ideals
decrease and their intersection is zero, since a nonzero
element has finite valuation at any fixed exceptional
component.

For completeness, this sequence is cofinal with the
$\mathfrak m$-adic topology. The images of $I_a$ in each finite
length quotient $R/\mathfrak m^n$ stabilize. Call the stable
image $V_n$. The maps $V_{n+1}\to V_n$ are onto: compute
both stable images using the same sufficiently large $a$.
If $V_n\ne0$ for some $n$, compatible lifts give, by completeness,
a nonzero element of $R$ lying in every
$I_a+\mathfrak m^n$ for every $n$. Krull intersection in
$R/I_a$ places it in every $I_a$, a contradiction. Thus some
$I_a$ lies in any prescribed $\mathfrak m^n$. For a sufficiently
large fixed $a$, therefore,
\[
        I_a\subset J,\qquad
        I_a\mathcal O_{Y_0}=\mathcal O_{Y_0}(-aD).
                                                               \tag{D.CR.4}
\]
The second equality follows from eventual generation of
powers of an ample line bundle
\citep[Proposition~28.27.13, condition~(7) and its full proof,
Tag 01PR]{StacksComparableCyclesB}.

Set $C_0=aD$ and $L=\mathcal O_{Y_0}(-C_0)$.
Choose finitely many elements of $I_a$ generating $L$.
Choose separating parameters $u,v$ for a finite regular
normalization of $R$, and one nonzero $t\in I_a$. Enlarge the
finite-dimensional $k$-space of the chosen sections to include
$t,ut,vt$; call the resulting space $V\subset I_a$.
For every closed point of the proper one-dimensional reduced
fibre, both sections of a pair in $V^2$ vanishing in the line
fibre imposes two independent linear conditions. The incidence
locus has dimension at most $2\dim V-1$, and its projection is
closed by properness. Hence a nonempty open set of pairs
generates $L$ at every point of $E$. The cokernel of
$\mathcal O_{Y_0}^2\to L$ is coherent and proper over the
local base. If its support were nonempty, its image would
contain the closed point, contradicting generation along $E$.
Thus these pairs generate $L$ on all of $Y_0$.

The differentials of elements of $V$ span the two-dimensional
space $\Omega_{\operatorname{Frac}R/k}$: the identities
$d(ut)-u\,dt=t\,du$ and $d(vt)-v\,dt=t\,dv$ prove this.
Here the separating regular normalization and the perfect
coefficient field supply the two-element $p$-basis calculation,
as in Proposition~\ref{D:single-curve-Cartier-parameter-ideal}.
Independence of the differentials of a pair is a nonempty
Zariski open condition on $V^2$; expressing the finitely many
coefficients of their determinant in a $k$-basis makes this
an open condition over $k$. It intersects the preceding
nonempty open set. Choose $(x,y)$ there.

Their common zero set on $Y_0$ is exactly the exceptional
fibre, so surjectivity of $Y_0\to\operatorname{Spec}R$ gives
$\sqrt{(x,y)R}=\mathfrak m$. The complete-local parameter
argument makes $R$ finite over $k[[x,y]]$; injectivity follows
from dimension. The independent differentials give generic
separability. This proves (D.CR.1). Finally, locally write
$x=s\alpha$, $y=s\beta$ for a generator $s$ of
$\mathcal O_{Y_0}(-C_0)$. Since $(\alpha,\beta)=1$, also
$(\alpha^j,\beta^j)=1$. This proves (D.CR.2) on $Y_0$ and
after every stated pullback.
\end{proof}

\begin{proposition}[Proved: uniformly comparable strong reduction cycles on all punctured covers]
\label{D:uniform-comparable-reduction-cycle}
Fix $R$ and its first-power Noether-different ideal $J$ as in
Proposition~\ref{D:cover-uniform-Noether-different}. Choose
$Y_0$ as the completion of a fixed henselian algebraic resolution,
using Lemma~\ref{D:fixed-henselian-algebraization}, and choose
$x,y,C_0$ by the preceding lemma. If necessary blow up once
before making these choices so that the exceptional fibre is
nonempty.

There are fixed integers $\beta\ge3$, $\alpha\ge1$, depending
only on these base choices, with the following property.
For each connected finite etale cover of the puncture of degree
$t$, let $A$ be the finite normal complete local covering ring.
Normalize $Y_0$ in its field, resolve, and contract all smooth
rational $(-1)$-curves mapping to points of $Y_0$. Denote the
resulting regular projective morphism by $f:Y\to Y_0$.
Put
\[
 C=f^*C_0,\qquad R_f=K_Y-f^*K_{Y_0},\qquad
                   Z=\alpha C-\beta R_f.                 \tag{D.CR.5}
\]
The canonical divisors in this formula are chosen compatibly;
$R_f$ is the actual effective Jacobian divisor. Then
\[
 \begin{gathered}
 C\le Z\le\alpha C,\qquad Z\text{ is anti-nef},\\
 -Z\cdot B\ge2p_a(B),\qquad
 -(K_Y+Z)\cdot B\ge2p_a(B),\\
 \text{for every nonzero effective exceptional cycle }B.
 \end{gathered}                                           \tag{D.CR.6}
\]
Consequently $Z$ and all its positive multiples are reduction
cycles: $\mathcal O_Z(-Z)$ and $\omega_Z^{-1}$ are globally
generated and $H^1(Z,\mathcal O_Z(-Z))=0$. For every finite
reflexive $A$-module $M$, its canonical full bundle
$F=(\pi^*M)^{\vee\vee}$ satisfies
\[
             F(-Z)\text{ is globally generated},\qquad
                    H^1(Y,F(-Z))=0.                      \tag{D.CR.7}
\]
The same cycle has the bounds
\[
 \begin{gathered}
 -Z^2+|K_Y\cdot Z|+
       \sum_H(-Z\cdot H)=O_R(t),\\
 h^0(bZ,\mathcal O_{bZ})+
       h^1(bZ,\mathcal O_{bZ})=O_{R,b}(t)
                         \quad(b\ge1\text{ fixed}).
 \end{gathered}                                           \tag{D.CR.8}
\]
Here $H$ ranges over the exceptional prime curves on this
specified relative-minimal model. No assertion is made for
an arbitrary sequence of additional blowups.

In particular the pointwise reduction-cycle correspondence,
the qualified all-Artinian realization theorem, and the
parameter-cycle jet and Hom-window theorems can all be applied
on this same $Y$, with
$bC\le bZ\le\alpha bC$ and fixed $\alpha$. This does not
remove the coflat or split-boundary hypotheses on a parameter
family, does not lift independent relations automatically,
and does not assert a hereditary representation count.
\end{proposition}

\begin{proof}
The fixed henselian algebraization descends every finite
punctured cover, its degree, and the chosen base resolution.
Completion preserves the exceptional fibre, intersections,
normal degrees, and geometric genus, as proved in
Proposition~\ref{D:all-complete-cover-geometric-budgets}.
Thus the relative resolution construction and the exact
maximum principle (D.LR.5) in
Proposition~\ref{D:punctured-cover-linear-resolution} apply
on this fixed base. They give a constant $\lambda\ge0$,
independent of $A$ and $t$, such that
\[
                  0\le R_f\le\lambda C.                  \tag{D.CR.9}
\]
This is the coefficientwise estimate for the actual
Jacobian divisor, including inseparable residue extensions,
not an estimate for a replacement divisor. The contractions
in the statement are essential: a curve $Q$ mapping to a
point then satisfies $K_Y\cdot Q\ge0$.

Write $c_i=-C_0\cdot E_i>0$ and $k_i=K_{Y_0}\cdot E_i$.
Fix $\beta\ge3$, and choose one integer $\alpha$ so large that
\[
        \alpha\ge\beta\lambda+1,\qquad
        \alpha c_i-\beta k_i\ge\beta\quad\text{for all }i.
                                                               \tag{D.CR.10}
\]
The choices are independent of the cover. By (D.CR.9),
$C\le Z\le\alpha C$; in particular $Z$ is an effective
integral Cartier cycle with full exceptional support.

Let $H$ be an exceptional prime on $Y$. If it maps to a
point, projection and (D.CR.5) give
\[
                   -Z\cdot H=\beta K_Y\cdot H\ge0.
\]
If it dominates $E_i$ with full residue-field degree $u_H$,
projection gives
\[
 -Z\cdot H
    =\beta K_Y\cdot H+
             u_H(\alpha c_i-\beta k_i).                 \tag{D.CR.11}
\]
Every exceptional prime on the regular surface has negative
integral square. Adjunction gives
$K_Y\cdot H=2p_a(H)-2-H^2\ge-1$.
If this number is negative it is exactly $-1$, and $H$ is
a smooth rational $(-1)$-curve. Thus (D.CR.10)--(D.CR.11),
and $u_H\ge1$, prove the stronger componentwise inequality
\[
          -Z\cdot H\ge
                  \beta\max\{K_Y\cdot H,0\}.             \tag{D.CR.12}
\]
It proves anti-nefness and handles the remaining radial
$(-1)$-curves without discarding them.

We verify both required inequalities on every nonreduced
cycle, rather than only on prime components. For
$B=\sum_H m_HH>0$, set
\[
 A_B=\sum_Hm_H\max\{K_Y\cdot H,0\},\qquad
 n_B=\sum_{K_Y\cdot H=-1}m_H,\qquad Q_B=-B^2>0 .
\]
Then
\[
 \begin{split}
 2p_a(B)&=2+A_B-n_B-Q_B,\\
 -Z\cdot B&\ge\beta A_B,\\
 -(K_Y+Z)\cdot B&\ge(\beta-1)A_B+n_B.
 \end{split}                                             \tag{D.CR.13}
\]
If $A_B\ge1$, use $Q_B\ge1$ and $\beta\ge3$ to see that
both lower bounds dominate $2+A_B-n_B-Q_B$.
If $A_B=0$ and $n_B\ge1$, then $n_B+Q_B\ge2$, so the
desired right side is nonpositive. If $A_B=n_B=0$, then
$K_Y\cdot B=0$, and adjunction makes $B^2$ even. Negative
definiteness therefore gives $Q_B\ge2$, with the same
conclusion. This proves (D.CR.6).

For every positive integer $b$, the first inequality remains
true for $bZ$ because $Z$ is anti-nef. The second remains
true because
$-(K_Y+bZ)\cdot B=-(K_Y+Z)\cdot B-(b-1)Z\cdot B$.
Apply the full nonreduced automatic-adjunction and point-ideal
argument of Theorem~\ref{D:reflexive-pullback-uniform-generation}
on every subcycle of $bZ$. It proves generation and
$H^1$-vanishing for $\mathcal O_{bZ}(-bZ)$ and generation
of $\omega_{bZ}^{-1}$, exactly as in
Proposition~\ref{D:energy-controlled-reduction-cycle}.
This uses the previously checked
\citet[Lemma~2.4 and Theorem~1.1, complete written proof,
pp.~192--193]{C_CFHR1999}.
The same uniform-generation theorem applies to every
canonical full $F$, whose global evaluation is surjective
at all exceptional generic points, and gives (D.CR.7).
No algebraization of $M$ itself is needed for that theorem.

It remains to check that comparison with $C$ also controls
the numerical budgets. Projection gives
$C^2=tC_0^2$. Since $C,Z$ are anti-nef and $Z\le\alpha C$,
\[
       0<-Z^2\le-\alpha Z\cdot C
                         \le\alpha^2(-C^2)=O_R(t).
                                                               \tag{D.CR.14}
\]
Every coefficient of $C$ is a positive integer, so
$\sum_H(-Z\cdot H)\le-Z\cdot C=O_R(t)$.
Also (D.CR.9) and anti-nefness give
\[
                |R_f\cdot Z|
                  =-R_f\cdot Z\le\lambda(-C\cdot Z)
                  =O_R(t).                              \tag{D.CR.15}
\]
If $C_0=\sum_i a_iE_i$, then positivity of pushforward and
$f_*C=tC_0$ imply $0\le f_*Z\le\alpha tC_0$. Hence
\[
 |f^*K_{Y_0}\cdot Z|
    =|K_{Y_0}\cdot f_*Z|
    \le\alpha t\sum_i a_i|k_i|=O_R(t).
                                                               \tag{D.CR.16}
\]
Together (D.CR.15)--(D.CR.16) prove $|K_Y\cdot Z|=O_R(t)$.
No coefficientwise comparison of squared divisors has been
used in these estimates.

The already proved first-power canonical-defect bound gives
$p_g(A)=h^1(Y,\mathcal O_Y)=O_R(t)$.
Coherent $H^2$ vanishes on $Y$ over the complete local base;
therefore restriction surjects
$H^1(Y,\mathcal O_Y)\to H^1(bZ,\mathcal O_{bZ})$.
Cycle Riemann--Roch gives
\[
       \chi(\mathcal O_{bZ})
           =-\frac{b^2Z^2+bK_Y\cdot Z}{2}=O_{R,b}(t).
                                                               \tag{D.CR.17}
\]
These two bounds prove the remaining part of (D.CR.8).

Finally (D.CR.2) identifies every parameter ideal
$(x^j,y^j)\mathcal O_Y$ with $\mathcal O_Y(-jC)$ on this
same resolution. The inclusions $bC\le bZ\le\alpha bC$
are actual Cartier-subscheme inclusions with constants fixed
before the cover. Thus the stated existing realization and
jet results have compatible finite levels. Their hypotheses
on families and relations are retained, not inferred from
these cycle comparisons.
\end{proof}


% Proved after four complete actual-file/source audits and exact precision refresh. Actual unit families only; no coflat object-atlas construction.
\subsection{Intrinsic integral unit jets and canonical-family openness}

\begin{proposition}[Proved: intrinsic finite jets of integral unit maps]
\label{D:intrinsic-integral-unit-Hom-jets}
Let $A$ be a complete Noetherian local $k$-algebra, with $k$
algebraically closed of characteristic $p>0$. Let $Y$ be
projective over $A$, and let $C$ be an effective Cartier
divisor whose support is the entire closed fibre. Assume
that every $jC$ is a proper $k$-scheme. Put $I=\mathcal O_Y(-C)$.
For an affine finite-type $k$-scheme $S=\operatorname{Spec}B$,
use the completion of $A\otimes_k B$ along $\mathfrak m_A$,
and the corresponding proper base change of $Y$.

Suppose that two actual length-$f$ cycles of vector bundles
$F_a,G_a$, with indices modulo $f$, are given on this
completed base change. Their arrows are meromorphic
isomorphisms
\[
 \phi_a:F_Y^*F_a\dashrightarrow F_{a+1},\qquad
 \psi_a:F_Y^*G_a\dashrightarrow G_{a+1},
                                                        \tag{D.IU.1}
\]
regular away from $C$, where $F_Y$ is absolute $p$-Frobenius,
including coefficient Frobenius. Assume a fixed bound $bC$
for every arrow and its inverse. Set
$H_a=\mathcal Hom(F_a,G_a)$ and choose $c\ge0$ such that
\[
 T_a(\eta)=\psi_a\,F_Y^*(\eta)\,\phi_a^{-1}
       \quad\hbox{maps }H_a\hbox{ into }H_{a+1}(cC).
                                                        \tag{D.IU.2}
\]
Take $c$ to work also for the reverse comparisons and the
two endomorphism cycles; $c=2b$ always does. Choose an integer $N\ge1$ with
\[
                         pN-c>N.                       \tag{D.IU.3}
\]
Then the functor of integral formal unit morphisms between
these cycles is represented by an affine scheme of finite
presentation over $S$, using only their restrictions to a
fixed finite thickening. More precisely, an integral unit
morphism is uniquely determined by the tuple
$\eta_a\in H^0(NC,H_a|_{NC})$. The exact finite equations are
\[
 \iota_N(\eta_{a+1})=T_{a,N}(\eta_a)
 \quad\hbox{in }H^0((N+c)C,H_{a+1}(cC)|_{(N+c)C}).
                                                        \tag{D.IU.4}
\]
Here $\iota_N$ is induced by $H_{a+1}\hookrightarrow
H_{a+1}(cC)$, and $T_{a,N}$ is obtained by Frobenius
pullback to $pNC$, composition with the actual pole maps,
and restriction to $(N+c)C$.

The functor of integral formal unit isomorphisms is also
affine of finite presentation: retain a reverse tuple with
the reverse unit equations and require the two products
to be the identity on $NC$. These descriptions are functorial
over every coefficient $B$-algebra, in the category of formal
bundles. Over finite-type coefficient algebras and finite
local Artinian coefficient algebras, the formal maps
algebraize on the proper Noetherian completed bases by
Grothendieck existence. No arbitrary-ring algebraization
or global lifting of an arbitrarily prescribed ordinary
Hom jet is asserted.
\end{proposition}

\begin{proof}
Every coefficient base change in the first part is interpreted
on the formal completion along $C$. Thus it is defined by the
compatible restrictions to the proper finite-level thickenings $jC$,
also when the coefficient algebra is not Noetherian. The
Cartier equation remains a nonzerodivisor: the finite
thickenings and their layer sequences are obtained by flat
coefficient extension over the field $k$.

We first type (D.IU.4) without dividing on a nonreduced
scheme. For a Cartier ideal,
\[
                 I^{N+c}H_{a+1}(cC)=I^N H_{a+1}.
                                                        \tag{D.IU.5}
\]
Consequently the inclusion of lattices induces a canonical
injection
\[
 H_{a+1}/I^N H_{a+1}\longrightarrow
 H_{a+1}(cC)/I^{N+c}H_{a+1}(cC).
                                                        \tag{D.IU.6}
\]
This is $\iota_N$. Frobenius sends $I^N$ into $I^{pN}$,
so an $N$-jet of a map determines its Frobenius pullback
on $pNC$. The pole maps in (D.IU.2), interpreted as honest
maps into the indicated twists, give a section of
$H_{a+1}(cC)$ on that thickening. Condition (D.IU.3)
implies $pN\ge N+c$, so restriction defines $T_{a,N}$.
There is therefore no arbitrary choice of a residual
fraction in (D.IU.4).

\emph{Local correction and gluing.}
Work on a common formal affine trivialization of the finitely
many bundles, and choose any integral formal lifts
$\eta_a^{(0)}$ of a solution of (D.IU.4). Such local lifts
exist by lifting matrix entries. Regard $T=(T_a)_a$ as the
cyclic operator on the tuple, so its output at phase $a+1$
uses the input at phase $a$. By (D.IU.5), the finite equation
says exactly
\[
                T\eta^{(0)}-\eta^{(0)}\in I^N\prod_a H_a.
                                                        \tag{D.IU.7}
\]
In particular the first iterate is integral. The operator
is additive and sends a difference of order at least $j$
to one of order at least $pj-c$. Thus for
$\eta^{(v+1)}=T\eta^{(v)}$ the successive differences
have order at least
\[
       N_v=p^vN-c\frac{p^v-1}{p-1}\longrightarrow\infty.
                                                        \tag{D.IU.8}
\]
Every iterate is integral and preserves the prescribed
$N$-jet. Completeness gives an integral limit $\eta$ with
$T\eta=\eta$. If two integral fixed tuples have the same
$N$-jet, their difference lies in every $I^{N_v}H_a$ and
is zero by separatedness. On an overlap two choices of
local lifts therefore give the same fixed tuple. The local
limits glue to the unique global formal unit morphism.
The construction and this uniqueness commute with every
coefficient change. It uses local formal lifts, not a
surjectivity assertion for global ordinary Hom sections.

For reverse tuples the same argument applies. The product
of two exact unit morphisms is an exact unit endomorphism.
If its $N$-jet is the identity, the just-proved uniqueness
makes it the identity formally. Thus the two finite inverse
equations are equivalent to an actual unit isomorphism.

\emph{Finite presentation.}
For each phase, the $N$-jet Hom functor is affine of finite
presentation over $S$: the finite cycle is proper and flat
over $S$, and both restricted bundles are finitely presented
and flat over $S$. The same statement applies to the target
Hom functor on $(N+c)C$, with its fixed line twist. These
are precisely the hypotheses of
\citet[Proposition 99.3.10, Tag 08K6]{StacksIntrinsicUnitHomD}.
Absolute Frobenius, the fixed pole maps and restriction
define natural polynomial transformations of these finite
Hom functors. Hence (D.IU.4) is an equality of two morphisms
between affine finite-presentation schemes. Its zero scheme
is again affine of finite presentation. Composition and the
two inverse equations are morphisms of the same finite
proper-cycle Hom functors, giving the asserted Isom scheme,
as in the paired-inverse construction of
\citet[Proposition 99.4.3, Tag 08K9]{StacksIntrinsicUnitHomD}.
This represents the actual integral formal functor by the
local contraction, rather than the image of a scheme of
uncontrolled higher-jet witnesses.

Finally, if the coefficient algebra is finite type over $k$
or finite local Artinian, the completed base ring is
Noetherian. The $C$-adic and closed-fibre topologies are
cofinal. The full faithfulness in Grothendieck existence
therefore identifies the formal maps with algebraic maps
on the proper base change of $Y$
\citep[Theorem 30.27.1, Tag 088E]{StacksIntrinsicUnitHomD}.
Their unit and inverse identities algebraize as well.
This is the only algebraization assertion used here.
\end{proof}

\begin{proposition}[Proved: specified Artinian unit lifts for canonical families]
\label{D:canonical-family-unit-Isom-openness}
In Proposition~\ref{D:intrinsic-integral-unit-Hom-jets},
suppose now that $A$ is a complete normal local surface and
$Y$ is a projective regular resolution, an isomorphism on
its puncture $U$. Assume that the actual parameter families
have the following property at every finite local Artinian
coefficient point, including every subsequent Artinian
coefficient change. For each phase there is a finite,
coefficient-flat module $M_{a,T}$ with maximal
Cohen--Macaulay closed fibre such that
\[
 M_{a,T}=H^0(U_T,F_{a,T}|_{U_T}),\qquad
 F_{a,T}=(\pi_T^*M_{a,T})^{\vee\vee},
                                                        \tag{D.IU.9}
\]
and these module identifications commute with the specified
coefficient changes. Require the same condition for every
other family being compared. The coflat, split-boundary
families in Proposition~\ref{D:weak-cycle-Artinian-full-realization}
supply (D.IU.9) whenever actual compatible unit arrows have
also been constructed; those arrows are not supplied by
that realization proposition alone.

Every prescribed unit isomorphism over a quotient of a
finite local Artinian coefficient algebra then lifts uniquely,
and its lift is integral on $Y_T$. For a fixed geometric
unit object, its integral unit-Isom scheme over $S$ is etale.
Consequently its isomorphism-class image is open and closed,
and
\[
 \#\{\hbox{unit classes represented by }S(k)\}
                         \le \#\operatorname{Irr}(S).
                                                        \tag{D.IU.10}
\]

For ordinary canonical cycles on the Cartier models of the
hereditary budget theorems, the forward arrows are regular
and the inverse arrows have pole at most $pC$. In comparing
two such cycles one may therefore take $c=p$ and $N=3$ in
(D.IU.2)--(D.IU.4), independently of the cover degree,
rank, number of phases and nilpotence order. This is an
isomorphism criterion for already actual families. It does
not construct a degree-linear parameter scheme of all such
families or remove the retained coflat/split-boundary
conditions.
\end{proposition}

\begin{proof}
Let $T\to T_0$ be a nilpotent quotient of finite local
Artinian coefficient algebras, and let a unit isomorphism
over $T_0$ be specified. On $U_T$, a length-$f$ ordinary
unit cycle is equivalent to its first bundle with a
$q=p^f$ Frobenius unit structure. The fixed-frame
construction gives its finite etale
$\operatorname{GL}_r(\mathbb F_q)$-torsor, as in the
earlier local unit and unit-Isom proofs. In local frames
this is the usual Lang equation, with invertible derivative;
the finite etale frame construction retains the maps and
their compositions. Finite etale covers and all their
morphisms are invariant under nilpotent thickenings
\citep[Lemma 58.8.3, Tag 0BQB]{StacksIntrinsicUnitHomD}.
Thus the specified unit isomorphism on $U_{T_0}$ lifts
uniquely to $U_T$.

The first identity in (D.IU.9) extends this punctured map
to an $A_T$-linear isomorphism $M_{a,T}\to M'_{a,T}$ at
every phase. Apply pullback and double dual to obtain an
integral isomorphism $F_{a,T}\to G_{a,T}$ on $Y_T$.
The base-change clause in (D.IU.9) makes its reduction
the prescribed isomorphism. It commutes with the unit
arrows: the two meromorphic maps agree on $U_T$, and the
fixed Cartier equation is a nonzerodivisor on their
locally free targets after a pole twist. Hence they agree
everywhere. The same observation proves uniqueness of
the integral extension. This is where canonical relative
fullness enters. Bounded poles of the unit arrows alone
would not provide this integral nilpotent lift.

The first proposition represents these integral unit
isomorphisms by a finite-presentation scheme. We have just
proved the required unique lifts for every finite local
Artinian extension at each closed point, not merely for
dual numbers. The local Artinian smoothness criterion,
together with uniqueness and vanishing of relative
tangent spaces, makes the morphism etale at those points;
the open etale locus is the whole scheme since it is
finite type over the algebraically closed field. These
are the exact criteria in
\citep[Tags 02HT, 02HM and 00UR]{StacksIntrinsicUnitHomD}.

Fixing a $k$-valued unit object, its Isom scheme has open
image in $S$. Such images for distinct classes are
disjoint: a nonempty intersection is open in a finite-type
$k$-scheme and has a closed $k$-point, which would identify
the two fixed objects. The union of these images contains
every closed point, hence all of $S$. Each complement is
the union of the other open images, so each class image
is also closed. Every nonempty class image contains an
irreducible component, giving (D.IU.10).

For the last assertion use the actual canonical Frobenius
sandwich of the Cartier-generation theorem. Its relative
pole statement can also be checked directly. The forward
punctured map extends to the canonical target module by
relative Hartogs, and then to the full bundles by generic
evaluation and codimension-two Hartogs. At each Artinian
coefficient base this uses the stated coefficient-flat maximal
Cohen--Macaulay target modules and their relative depth-two
Hartogs property; no regularity of a Frobenius source module
is presumed. The source bundle
$F_{a,T}$ has factorizations of $x\operatorname{Id}$ and
$y\operatorname{Id}$ through free bundles. Frobenius gives
such factorizations for $x^p$ and $y^p$ on $F_Y^*F_{a,T}$.
Compose the inverse punctured arrow with the first free-target
map. Relative module Hartogs and full pullback extend that
composite from the canonical target bundle. The second map
of the factorization then shows that both $x^p\phi_a^{-1}$
and $y^p\phi_a^{-1}$ are regular. Since $x^p,y^p$ generate
$\mathcal O_Y(-pC)$, the inverse has pole at most $pC$.
The same argument works at every Artinian coefficient base.
For a finite-type parameter base, these bounds hold on the
actual family as well: relative to its initially fixed finite
pole bound, each obstruction to the sharper bound is a
coherent section on a fixed proper finite-level thickening.
It vanishes after every Artinian coefficient specialization;
Krull intersection at the closed coefficient points therefore
makes it zero. Thus the target
forward arrow in (D.IU.2) has no pole and the source inverse
has pole at most $pC$, so $c=p$. For every prime $p\ge2$,
$3p-p>3$. Thus $N=3$ satisfies (D.IU.3), and the finite
equations are on $(p+3)C$. The integral Artinian comparison
just proved prevents any increase of this depth with the
nilpotence order.
\end{proof}

\begin{remark}[The unresolved quantitative step]
The finite Hom schemes here provide qualitative finite
presentation and exact specified-isomorphism lifting for
actual families. No coordinate or degree estimate for
cohomology base-change strata or split-boundary conditions
is inferred. The known $O(ftr^2)$ finite-cycle quotient/Hom
budgets and the present openness theorem therefore do not,
by themselves, give an $O(ftr^2)$ effective object atlas or
a hereditary representation bound.
\end{remark}


% Proved after four complete actual-file audits. Generated fixed-determinant families; no assumed relative fullness.
\subsection{Generated fixed-determinant unit families without a coflatness hypothesis}

\begin{lemma}[Proved: generation and determinants from one weak cycle]
\label{D:weak-cycle-generated-determinant-lifting}
Let $Y$ be projective over a complete Noetherian local
$k$-algebra, and let $Z$ be an effective Cartier divisor
with the entire closed fibre as support. Assume that $Z$
is a proper one-dimensional $k$-scheme and that
\[
             H^1(Z,\mathcal O_Z(-jZ))=0\qquad(j\ge1).
                                                        \tag{D.GD.1}
\]
For a $k$-algebra $B$, work on the formal coefficient base
change of $Y$ along $Z$. Let $\mathcal F$ be a formal
vector bundle of rank $r$. Suppose that
$V=\mathcal F(-Z)|_{Z_B}$ is supplied with a surjection
$\mathcal O_{Z_B}^{h}\to V$. Then $\mathcal F(-Z)$ is
generated by $h$ formal global sections extending those
specified generators.

Restriction on isomorphism classes of line bundles from
the formal coefficient base change to $Z_B$ is injective.
In particular, if $L$ is a fixed line bundle on $Y$ and
an isomorphism $\det V\simeq L|_{Z_B}$ is supplied, it
extends to an isomorphism
\[
             \det(\mathcal F(-Z))\simeq L_B
\quad\hbox{on the formal coefficient base change}.       \tag{D.GD.2}
\]
The extending generators and determinant isomorphism need
not be unique. After choosing them over $B$, their pullbacks
have the same properties over every coefficient $B$-algebra.
For finite-type or finite local Artinian coefficient algebras,
these assertions algebraize on the proper Noetherian
completed base.
\end{lemma}

\begin{proof}
Put $I=\mathcal O_Y(-Z)$. The layer $I^j/I^{j+1}$ is
$\mathcal O_Z(-jZ)$. Flat coefficient extension over the
field and (D.GD.1) give
$H^1(Z_B,\mathcal O_{Z_B}(-jZ))=0$ for every $B$.
The proper curve has a two-affine cover with affine
intersection, which remains such a cover after coefficient
extension. Consequently quasi-coherent cohomology in degree
two vanishes over every coefficient ring. The kernel of the
specified generating quotient is a vector bundle. Its
sequence, tensored by $\mathcal O_{Z_B}(-jZ)$, therefore gives
\[
                    H^1(Z_B,V(-jZ))=0\qquad(j\ge1).
                                                        \tag{D.GD.3}
\]
The transition from $jZ_B$ to $(j+1)Z_B$ for
$\mathcal F(-Z)$ has kernel $V(-jZ)$. Thus (D.GD.3)
lifts each of the $h$ chosen global sections successively
through all the thickenings. Choose compatible lifts.
They generate on every thickening by Nakayama, since they
generate on $Z_B$. Their limit is the asserted formal
surjection. This uses the specified quotient itself, not
cohomology base change for an unspecified family.

For $j\ge1$, the kernel of
$\mathcal O_{(j+1)Z_B}^{*}\to\mathcal O_{jZ_B}^{*}$
is, under $1+u\leftrightarrow u$, the additive sheaf
$\mathcal O_{Z_B}(-jZ)$. Its $H^1$ is zero. Thus an
isomorphism between two line bundles on $jZ_B$ lifts to
$(j+1)Z_B$ whenever the two line bundles already exist
there: the obstruction to gluing local lifts lies in this
$H^1$. Starting from the prescribed isomorphism on $Z_B$
and choosing compatible lifts proves injectivity and
(D.GD.2). It does not prove uniqueness of the isomorphism;
units congruent to one may change it.

All constructions take place on the finite thickenings
and commute with pullback once the lifts have been chosen.
For the stated Noetherian coefficient bases, full faithfulness
of Grothendieck existence algebraizes these maps and line
isomorphisms. The formal surjection is an actual surjection:
its coherent cokernel has zero formal completion, and
properness over the adically complete base forces that
cokernel to vanish. This is the same proper-complete scope
as Proposition~\ref{D:intrinsic-integral-unit-Hom-jets},
using \citet[Theorem 30.27.1, Tag 088E]{StacksIntrinsicUnitHomD}.
\end{proof}

\begin{proposition}[Proved: meromorphic unit-Isom openness for generated fixed-determinant families]
\label{D:generated-fixed-determinant-unit-openness}
Let $A$ be a complete excellent normal local surface with
algebraically closed coefficient and residue field $k$ of
characteristic $p>0$. Let $\pi:Y\to\operatorname{Spec}A$
be a projective regular resolution, an isomorphism over
the puncture $U$. Let $C,Z$ be effective Cartier cycles
with full exceptional support and $Z\le\alpha C$ for a
fixed integer $\alpha$. Assume that $Z$ satisfies (D.GD.1).

Let $S=\operatorname{Spec}B$ be of finite type over $k$.
On the proper base change over the completion of
$A\otimes_k B$, suppose there is an actual length-$f$
cycle of rank-$r$ bundles $F_a$ with meromorphic Frobenius
unit arrows, regular and invertible on $U_S$, whose arrows
and inverses have a fixed finite pole bound along $C$.
For every phase assume that
\[
 F_a(-Z)|_{Z_S}\text{ has a specified generating quotient},
 \qquad
 \det(F_a(-Z))|_{Z_S}\simeq L_a|_{Z_S},                 \tag{D.GD.4}
\]
where each $L_a$ is one fixed line bundle on $Y$.
No relative maximal-Cohen--Macaulay pushforward, coflatness,
or canonical full-pullback property of $F_a$ is assumed.

For each fixed punctured unit class represented by $S(k)$,
there is a finite integer $d$, allowed to depend on that
fixed comparison object, such that its punctured unit-Isom
functor over all finite local Artinian coefficient points
is exactly the functor of unit isomorphisms with both
meromorphic directions bounded by $dC$. The latter bounded
functor is represented over $S$ by an affine scheme of
finite presentation and is etale. Its image is precisely
the class in question and is open and closed. Consequently
\[
 \#\{\text{punctured unit classes represented by }S(k)\}
                    \le\#\operatorname{Irr}(S).          \tag{D.GD.5}
\]
The pole bound here need not be independent of rank, and
no coordinate bound for the comparison scheme is asserted.
Neither bound is needed for (D.GD.5).
\end{proposition}

\begin{proof}
By the preceding lemma, every $F_a(-Z)$ is globally generated
on the completed proper family, and
\[
             \det F_a\simeq D_a:=L_a(rZ)_S.              \tag{D.GD.6}
\]
Choose such generators and determinant isomorphisms once.
Their pullbacks work at every coefficient point. The
assertion also applies if the initial family was supplied
formally: Grothendieck existence on the completed finite-type
base algebraizes the bundles, these maps, and the bounded
meromorphic unit arrows. Arbitrary coefficient rings below
are used for the finite-level formal functors, not for an
additional algebraization assertion.

\emph{A fixed canonical comparison target.}
Fix a $k$-point with bundle cycle $F_a^0$. Its restriction
to $U$ determines a finite reflexive $A$-module $M_a^0$:
extend the coherent bundle across the closed point and take
its reflexive hull. Put
\[
                    P_a=(\pi^*M_a^0)^{\vee\vee}.
                                                        \tag{D.GD.7}
\]
This is a vector bundle on the regular surface and agrees
with $F_a^0$ on $U$. All $P_a$ in this proof are fixed
objects over $k$; their constant coefficient extensions
are used as targets only. In particular,
\[
 H^0(U_T,P_a|_{U_T})=M_a^0\otimes_kT,
 \qquad H^0(U_T,\mathcal O_{U_T})=A\otimes_kT           \tag{D.GD.8}
\]
for every finite local Artinian $k$-algebra $T$.
These identities are the usual normal-surface Hartogs
identities followed by flat coefficient extension over
$k$. They do not assert normality of the nonreduced ring
$A\otimes_kT$.

Consider any punctured unit isomorphism from $F_{a,T}$
to the fixed cycle $F_a^0|_{U_T}$, and identify the latter
with $P_a|_{U_T}$. Compose it with the global generators
of $F_{a,T}(-Z)$. Each resulting section on $U_T$ extends
by (D.GD.8) to a section of $P_{a,T}$: an element of
$M_a^0\otimes_kT$ has its canonical evaluation section in
that full pullback. The relations between the generators
vanish on $U_T$, hence everywhere in the locally free
target. Indeed, after a Cartier pole twist the defining
equation of $C$ is a nonzerodivisor, also after coefficient
extension. The resulting map therefore factors through
the generating quotient and gives an actual map
\[
        F_{a,T}(-Z)\longrightarrow P_{a,T},\qquad
        h_{a,T}:F_{a,T}\longrightarrow P_{a,T}(Z).
                                                        \tag{D.GD.9}
\]
This argument requires global generation of the source
after the stated twist and Hartogs for the fixed canonical
target. It requires no relative fullness of the source.

Apply the same construction to the identity of the fixed
punctured object, obtaining
$h_a^0:F_a^0\to P_a(Z)$. Its determinant is a nonzero
meromorphic line isomorphism
\[
                 \tau_a:D_a\dashrightarrow\det P_a
                                                        \tag{D.GD.10}
\]
whose divisor is supported on the exceptional fibre.
Using (D.GD.6), the ratio
$\det h_{a,T}/\tau_a$ is an invertible function on $U_T$.
Both that function and its inverse extend by (D.GD.8),
so the ratio is a unit of $A\otimes_kT$. In particular
no nilpotent coefficient can change the exceptional
orders of this determinant comparison.

The entries of $h_{a,T}$ have poles bounded by $Z$.
Their adjugate entries have poles bounded by $(r-1)Z$.
Dividing by the determinant adds only the fixed positive
part of $\operatorname{div}(\tau_a)$. Choose $e_a$ with
$\operatorname{div}(\tau_a)^+\le e_aC$. Thus, in the
meromorphic sense,
\[
 \begin{split}
 \operatorname{pole}(h_{a,T})&\le Z,\\
 \operatorname{pole}(h_{a,T}^{-1})&\le(r-1)Z+e_aC.
 \end{split}                                             \tag{D.GD.11}
\]
The same bound holds for $h_a^0$ and its inverse. Composing
with $h_a^0$ compares $F_{a,T}$ with $F_a^0$, with each
direction bounded by $rZ+e_aC$. Consequently one fixed
integer
\[
                       d=r\alpha+\max_a e_a             \tag{D.GD.12}
\]
works for every such punctured isomorphism and every finite
Artinian coefficient base. It is independent of the
nilpotence order. Its dependence on the chosen fixed
object and its determinant comparison is harmless.

\emph{The finite-presentation bounded comparison functor.}
We spell out the meromorphic version of the already proved
integral jet argument. If $H_a=\mathcal Hom(F_a,F_a^0)$,
use the lattice $H_a(dC)$ in the meromorphic Hom sheaf.
If the unshifted comparison operator $T$ has loss $cC$,
then it has loss
\[
                       c_d=c+(p-1)d                    \tag{D.GD.13}
\]
relative to this shifted lattice. This follows because
Frobenius multiplies the pole $dC$ by $p$.
For the product comparison use the endomorphism lattice
$\mathcal End(F_a)(2dC)$ and the analogous loss
$c+2(p-1)d$. Choose a common jet depth $N$ with
$pN-\max\{c_d,c+2(p-1)d\}>N$, increasing $c$ once to
work for both directions and all endomorphism cycles.
The intrinsic equations of
Proposition~\ref{D:intrinsic-integral-unit-Hom-jets}, now
on these lattices, represent exactly the meromorphic
unit maps with pole at most $dC$. Require reverse tuples
and the two product equations with the identity section
in the $2dC$ lattice. Products are determined on the
retained finite cycle; uniqueness of the fixed formal
endomorphism makes the two products exactly the identity.
No inversion on a nonreduced Cartier cycle is used.

Thus proper finite-cycle Hom schemes and finite equations
give an affine finite-presentation scheme of these bounded
unit isomorphisms. This is precisely the finite-level
representability and local contraction already proved in
Proposition~\ref{D:intrinsic-integral-unit-Hom-jets}; its
primary inputs are
\citet[Tags 08JY, 08K6, 08K9 and 088E]{StacksIntrinsicUnitHomD}.
The scheme represents the formal functor over arbitrary
coefficient algebras, and the algebraic bounded functor
on the stated proper Noetherian completed bases.

\emph{Specified Artinian lifts and class images.}
Given a nilpotent quotient $T\to T_0$ and a specified
bounded unit isomorphism over $T_0$, finite-etale fixed
frames on $U_T$ lift its punctured isomorphism uniquely.
This is the nilpotent invariance of finite etale covers
and all their morphisms, as used in
Proposition~\ref{D:canonical-family-unit-Isom-openness}
\citep[Tag 0BQB]{StacksIntrinsicUnitHomD}.
The argument (D.GD.8)--(D.GD.12) puts that unique lift
inside the same fixed pole bound $dC$. Its reduction is
the prescribed bounded meromorphic map, since maps into
a Cartier pole lattice agreeing on $U_{T_0}$ agree
everywhere. Hence the bounded comparison scheme has
unique lifts over every finite local Artinian extension.

The exact finite-presentation Artinian criteria in
\citet[Tags 02HT, 02HM and 00UR]{StacksIntrinsicUnitHomD}
make it etale at every closed point, and therefore
etale everywhere on this finite-type $k$-scheme. Its
image is open. Over each closed coefficient point, the
same determinant argument shows that every punctured
unit isomorphism is inside the chosen bound, so this
image is exactly the punctured class of $F^0$.
Images for different classes are disjoint; any nonempty
open intersection would have a closed $k$-point. They
cover all closed points and thus all of $S$. Their
complements are unions of the other open images, making
each image closed as well. Counting nonempty class
images by irreducible components proves (D.GD.5).
\end{proof}

\begin{remark}[Role in the hereditary count]
The preceding propositions replace relative fullness by
actual generation after a fixed negative cycle twist and
a fixed determinant. They do not construct compatible
Frobenius objects from finite jets. A separate effective
object theorem and a polynomial parameter budget are still
required before applying the component bound to a hereditary
representation count. In particular no bound on the size
of the meromorphic comparison scheme has been used.
\end{remark}


% Separate post-CP25 proof. Four complete actual-file audits PASS. Formal unit lattices; canonical fullness is not inferred.
\subsection{Formal effectivity of finite bounded-pole Frobenius data}

\begin{proposition}[Proved: finite Frobenius data and formal object correction]
\label{D:finite-Frobenius-object-effectivity}
Let $A$ be a complete Noetherian local $k$-algebra, where $k$
is an algebraically closed field of characteristic $p>0$.
Let $Y\to\operatorname{Spec}A$ be projective, and let $C$ be
an effective Cartier divisor supported on the entire closed
fibre. Assume that every $jC$ is proper over $k$, and put
$I=\mathcal O_Y(-C)$. For an affine finite-type $k$-scheme
$S=\operatorname{Spec}B$, use the completed base change and
its formal completion along $C$ as in
Proposition~\ref{D:intrinsic-integral-unit-Hom-jets}.
The formal assertions below also apply after every change
to an arbitrary coefficient $B$-algebra.

Fix integers $b\ge0$, $N\ge1$, and $f\ge1$ with
\[
                 pN-b>N,\qquad L=N+b.                  \tag{D.FO.1}
\]
On $LC\times_k S$, suppose that rank-$r$ vector bundles
$E_a$, indexed modulo $f$, and maps
\[
 \alpha_a:F^*E_a\longrightarrow E_{a+1},\qquad
 \beta_a:E_{a+1}\longrightarrow (F^*E_a)(bC)             \tag{D.FO.2}
\]
are given. Here $F$ is absolute $p$-Frobenius of the finite
thickening, including coefficient Frobenius. Require the
two exact finite equalities
\[
 \alpha_a(bC)\,\beta_a=\iota_{E_{a+1}},\qquad
 \beta_a\,\alpha_a=\iota_{F^*E_a},                      \tag{D.FO.3}
\]
where $\iota_V:V\to V(bC)$ is the canonical Cartier map.
It is not asserted to be injective on this nonreduced
finite thickening.

These finite data produce vector bundles $\mathcal F_a$
on the formal completion and actual arrows
\[
 \phi_a:F^*\mathcal F_a\longrightarrow\mathcal F_{a+1}
                                                                    \tag{D.FO.4}
\]
which are regular, are isomorphisms off $C$, and have
inverse pole at most $bC$. There are specified
identifications $\mathcal F_a|_{NC}\simeq E_a|_{NC}$
under which the forward maps and the inverse maps into
the $bC$-twists agree with (D.FO.2) on $NC$.

The construction is independent, up to a unique integral
unit isomorphism inducing these specified $N$-jet
identifications, of its local choices. A morphism of the
finite data intertwining (D.FO.2) gives a unique integral
unit morphism inducing its prescribed $N$-jet. These maps
retain identities and compositions. The construction and
these comparisons commute with all coefficient changes
in this precise formal, marked sense.

For finite-type or finite local Artinian coefficient
algebras, the constructed formal bundles and arrows
algebraize over the corresponding proper Noetherian
completed base. The assertion retains the prescribed
data on $NC$; it does not claim to retain the entire
input bundle on $LC$. It does not assert that the
constructed lattices are canonical full pullbacks or
that their punctured pushforwards are coefficient flat.
\end{proposition}

\begin{proof}
All arbitrary coefficient changes are defined by the
compatible finite-level thickenings $jC$. Their Cartier
layer sequences remain exact after coefficient extension
over the field $k$. We may work on formal affine opens
where $C$ has equation $t$, and where every $E_a$ has a
frame on $LC$. These opens exist by local freeness on
the finite thickening; nilpotent thickenings of affine
opens are affine
\citep[Lemma 37.2.3, Tag 06AD]{C_StacksFiniteObjectAffineness}.
Their inverse systems have surjective
transition maps on functions. Their formal rings are
complete and separated for the Cartier filtration.
Absolute Frobenius preserves their underlying open
sets and sends $I^j$ into $I^{pj}$.

\emph{Lifting the local Frobenius matrices.}
Write the matrices of $\alpha_a$ and $\beta_a$ as
$\overline A_{a,i}$ and $t^{-b}\overline B_{a,i}$ on
such a formal open $i$. The notation $t^{-b}$ specifies
the frame of the line bundle $\mathcal O(bC)$; no
element is inverted in the finite ring. The finite
equalities (D.FO.3) are exactly
\[
 \overline A_{a,i}\overline B_{a,i}=t^b\operatorname{Id},
 \qquad
 \overline B_{a,i}\overline A_{a,i}=t^b\operatorname{Id}
                  \pmod {t^L}.                       \tag{D.FO.5}
\]
Choose arbitrary integral formal lifts $A_{a,i},B_{a,i}$
of their entries. Since $L-b=N\ge1$, (D.FO.5) gives
\[
 A_{a,i}B_{a,i}=t^b(\operatorname{Id}+t^N R_{a,i}),
 \quad
 B_{a,i}A_{a,i}=t^b(\operatorname{Id}+t^N R'_{a,i})      \tag{D.FO.6}
\]
for integral matrices $R_{a,i},R'_{a,i}$. Both parenthesized
matrices are invertible by their convergent geometric
series. Thus $A_{a,i}$ is invertible after inverting $t$,
and its inverse has pole at most $b$. Moreover
\[
          t^b A_{a,i}^{-1}\equiv B_{a,i}\pmod {t^N}.    \tag{D.FO.7}
\]
These statements hold over arbitrary coefficient rings;
they do not use field points to infer invertibility.

\emph{Correcting the overlap matrices.}
Use the convention that the finite transition matrix
$\overline g_{a,ij}$ sends coordinates on open $j$ to
coordinates on open $i$. Choose integral formal lifts
$g^{(0)}_{a,ij}$. They are invertible, since their
reductions modulo $I$ are invertible. Compatibility of
the finite map $\alpha_a$ gives
\[
 A_{a,i}F(g^{(0)}_{a,ij})
       -g^{(0)}_{a+1,ij}A_{a,j}\in I^L
                                                        \tag{D.FO.8}
\]
on the overlap, in its matrix coordinates. Define the
cyclic correction operator
\[
 \bigl(T_{ij}(g)\bigr)_{a+1}
                 =A_{a,i}F(g_a)A_{a,j}^{-1}.            \tag{D.FO.9}
\]
The inverse pole bound and (D.FO.8) show that
$T_{ij}(g^{(0)})-g^{(0)}\in I^{L-b}=I^N$.
In particular the first iterate is integral. A difference
of order $j$ is sent to one of order at least $pj-b$.
Consequently the successive differences in the iterates
have orders at least
\[
        N_v=p^vN-b\frac{p^v-1}{p-1}\longrightarrow\infty.
                                                        \tag{D.FO.10}
\]
All iterates remain integral and preserve the original
$N$-jet. Their unique limit $g_{a,ij}$ is integral and
invertible, agrees with $\overline g_{a,ij}$ on $NC$,
and is fixed by (D.FO.9). Two fixed tuples with the same
$N$-jet coincide, by the same order calculation and
separatedness.

On a triple overlap, the tuple
$g_{a,ij}g_{a,jk}$ is a fixed tuple for $T_{ik}$: the
middle factors $A_{a,j}^{-1}A_{a,j}$ cancel exactly.
Its $N$-jet equals that of $g_{a,ik}$, because the
original finite transition matrices satisfy their
cocycle identity. Uniqueness therefore gives
\[
       g_{a,ij}g_{a,jk}=g_{a,ik},\qquad
       g_{a,ii}=\operatorname{Id},\qquad
       g_{a,ji}=g_{a,ij}^{-1}.                          \tag{D.FO.11}
\]
The corrected matrices glue actual formal vector bundles.
Their fixed-point equations say exactly that the local
$A_{a,i}$ glue to the arrows (D.FO.4). Regularity and
the bound $bC$ for their inverses are local and have
already been proved. Equations (D.FO.7) and (D.FO.11)
give the prescribed restrictions of both arrows on $NC$.

\emph{Comparisons and coefficient changes.}
For two choices of local matrix lifts of the same finite
data, the local identity matrix has unit-equation defect
in $I^{L-b}=I^N$. Correct it by
\[
                 h_{a+1}\longmapsto
             A'_{a,i}F(h_a)A_{a,i}^{-1}.                \tag{D.FO.12}
\]
The same contraction produces a unique integral unit
comparison with identity $N$-jet. Comparisons on overlaps
agree by uniqueness, as do their composites. They are
isomorphisms because their $N$-jets are isomorphisms.
More generally, lift a given morphism of finite data in
local frames. Its finite compatibility again gives
defect in $I^N$, so (D.FO.12) produces the asserted
formal unit morphism with its specified $N$-jet.
Products and identities are retained because they are
fixed tuples with the prescribed products and identity
jets. This also proves independence of refinements of
the chosen cover. Base change preserves the finite
equalities, and every calculation at a fixed order uses
only finitely many additions, products and inverses of
matrices congruent to the identity. The limits therefore
commute with coefficient change in the category of
compatible formal restrictions. The unique comparison
just constructed makes this assertion independent of
choices made after the change of coefficients.

Finally, for the stated Noetherian completed bases,
the Cartier and closed-fibre topologies are cofinal.
Grothendieck existence algebraizes the formal bundles,
the regular arrows and their inverse maps into the
fixed line twists
\citep[Theorem 30.27.1, Tag 088E]{StacksIntrinsicUnitHomD}.
Full faithfulness preserves the two inverse identities.
No arbitrary-ring algebraization is used.
\end{proof}

\begin{remark}[Scope of this object construction]
The hypotheses are actual equations on one fixed finite
cycle, including both finite inverse equations. Thus this
construction does not infer a simultaneous Frobenius
lift from the existence of unrelated bundle extensions.
Conversely, it does not imply canonical fullness. For
example, for $f=1$ the bundle $\mathcal O_Y(-C)$ has
the natural regular map
$F^*\mathcal O_Y(-C)=\mathcal O_Y(-pC)\to
\mathcal O_Y(-C)$, with inverse pole $(p-1)C$.
Its restriction to the puncture is the trivial unit line.
For a nonzero exceptional $C$ on a normal surface
resolution its canonical full extension is
$\mathcal O_Y$, not $\mathcal O_Y(-C)$.
In particular, quantitative finite-cycle coordinates and
formal object effectivity do not themselves give the
canonical-family hypotheses of
Proposition~\ref{D:canonical-family-unit-Isom-openness}.
No representation count or new finite-generation
conclusion is asserted here.
\end{remark}


% Separate post-CP25 proof. Four complete actual-file audits PASS. Ordinary hereditary count on arbitrary punctured covers.
\subsection{Degree-linear ordinary representation growth on arbitrary punctured covers}

\begin{theorem}[Proved: hereditary ordinary representation growth]
\label{D:arbitrary-cover-ordinary-hereditary-growth}
Let $(R,\mathfrak m,k)$ be a complete excellent normal local
surface with algebraically closed coefficient and residue field
$k$ of characteristic $p>0$. Put
$U_R=\operatorname{Spec}R\setminus\{\mathfrak m\}$.
There is a constant $A_R$ with the following property.
For every connected finite etale cover $U_A\to U_R$ of
degree $t$, every $q=p^f$, and every $r\ge1$,
\[
 \left|\operatorname{Hom}_{\rm cont}
       (\pi_1(U_A),\operatorname{GL}_r(\mathbb F_q))
          /\operatorname{GL}_r(\mathbb F_q)\right|
                         \le q^{A_Rtr^2}.                \tag{D.HU.1}
\]
There is no tameness condition on the finite cover.
The constant is independent of its degree, its ramification
along a resolution, $r$, and $f$.

More precisely, after choosing one of at most
\[
              q^{c_0t}\exp(c_1ft\log(r+1))               \tag{D.HU.2}
\]
fixed ordered determinant tuples, all the relevant ordinary
unit objects are covered by a single open subset of an affine
polynomial zero scheme with at most $c_2ftr^2$ coordinates.
Its equations have degree at most a constant $D_R$ independent
of $t,r,f$. Every point of this open parameter scheme produces
an actual formal, and on finite-type coefficient bases actual
algebraic, Frobenius unit object on $U_A$.
Its punctured unit classes have open and closed images.
The parameter families need not be relatively canonical,
coflat, or flat families of maximal Cohen--Macaulay modules.
\end{theorem}

\begin{proof}
The finite normalization $A$ is a complete normal local domain
with residue field $k$: it is finite over the complete henselian
base, and connectedness of the punctured cover selects one
normal local factor. Use the one fixed henselian algebraization
and all-cover geometric budget results to choose the specified
regular projective model $\pi:Y\to\operatorname{Spec}A$.
On the fixed base model choose the separating parameter pencil
inside the cover-uniform Noether-different ideal. Apply
Proposition~\ref{D:uniform-comparable-reduction-cycle}.
It provides actual cycles
\[
 C=f_A^*C_0,\qquad Z=\alpha C-\beta R_{f_A},\qquad
                C\le Z\le\alpha C,                     \tag{D.HU.3}
\]
where $\alpha,\beta$ and the base data depend only on $R$.
The number of exceptional components and $p_g(A)$ are $O_R(t)$
on this same model. In addition,
\[
 \begin{gathered}
 (x^j,y^j)\mathcal O_Y=\mathcal O_Y(-jC),\qquad
 -C^2,\ -Z\cdot C,\ -Z^2,\ |K_Y\cdot Z|=O_R(t),\\
 H^1(Z,\mathcal O_Z(-jZ))=0\quad(j\ge1),\qquad
 \mathcal O_Y(-C),\ \mathcal O_Y(-Z)
          \text{ are globally generated}.
 \end{gathered}                                        \tag{D.HU.4}
\]
Here generation and the conormal vanishings follow from the
all-subcycle inequalities in that proposition and the checked
uniform-generation theorem. They apply to every fixed multiple
of $Z$. Generation of $\mathcal O_Y(-C)$ also follows directly
from the parameter pencil.

For clarity, the scalar budgets on any fixed $\ell C$ are
simultaneous with (D.HU.4). Indeed
$K_Y=f_A^*K_{Y_0}+R_{f_A}$, with
$0\le R_{f_A}\le\lambda C$ and $C$ anti-nef. Projection
therefore gives $|K_Y\cdot C|=O_R(t)$, as well as
$-C^2=O_R(t)$. Coherent $H^2$ vanishes over the affine
surface, so $H^1(Y,\mathcal O_Y)$ surjects onto
$H^1(\ell C,\mathcal O_{\ell C})$. Adjunction now gives
\[
 h^0(\mathcal O_{\ell C})+h^1(\mathcal O_{\ell C})
                              =O_{R,\ell}(t).            \tag{D.HU.5}
\]
No additional blowup of this chosen model is made.

\emph{Geometric determinant lists and canonical input objects.}
A continuous $\mathbb F_q$ representation is equivalently an
ordinary $q$-unit bundle on $U_A$, or a length-$f$ cycle of
ordinary $p$-unit bundles. At phase $a$ take its reflexive
extension $M_a$ and its canonical full pullback
$F_a=(\pi^*M_a)^{\vee\vee}$. The actual canonical arrows
satisfy
\[
 \phi_a:F_Y^*F_a\longrightarrow F_{a+1},\qquad
 \phi_a^{-1}:F_{a+1}\longrightarrow(F_Y^*F_a)(pC).
                                                               \tag{D.HU.6}
\]
These are regular maps to the indicated bundles. They are
the canonical Frobenius sandwich supplied by the first-power
Noether factorization and the parameter pencil, as in
Proposition~\ref{D:canonical-family-unit-Isom-openness}.
For these geometric input objects no relative fullness
assertion is needed.

The proof of Lemma~\ref{D:hereditary-ordinary-unit-determinants}
applies on this model. We recall why its constants remain
unchanged when the generating cycle is now $Z$.
There are at most $q^{c_0t}$ initial punctured determinant
characters, because its bound
$|H^1(U_A,\boldsymbol\mu_{q-1})|
 \le(q-1)^{2p_g(A)+v}$ uses only the geometric genus and the
number $v=O_R(t)$ of exceptional components. All later
punctured determinants are the prescribed $p$-power phases
of that one character. Generic evaluation of a full bundle
gives nonnegative integral exceptional determinant degrees.
Pairing the cyclic determinant recurrence of (D.HU.6) with
the anti-nef cycle $C$ gives
\[
 0\le c_1(F_a)\cdot C
              \le\frac{rp}{p-1}(-C^2)=O_R(tr).          \tag{D.HU.7}
\]
Every coefficient of $C$ is a positive integer. Thus the
unweighted sum of these degrees is also $O_R(tr)$.
For a fixed punctured determinant, the negative definite
intersection matrix makes its exceptional multidegree
determine its extension line on $Y$ uniquely. The same
binomial enumeration as in the cited lemma gives at most
$\exp(c_1t\log(r+1))$ determinant lines per phase.
This proves (D.HU.2) for the ordered global lines
\[
                    L_a=\det(F_a(-Z)).                 \tag{D.HU.8}
\]
These are choices of actual lines on $Y$, not moving Picard
parameters. The cycle $Z$ is fixed before the representation.
By (D.HU.3) and the comparable-cycle generation theorem,
every $F_a(-Z)$, and hence every $L_a$, is globally generated.

\emph{One fixed finite input level.}
Choose once
\[
       b=p,\qquad N\ge\max\{3,\alpha\},\qquad L=N+p,
       \qquad D=LC.                                    \tag{D.HU.9}
\]
Then $pN-p>N$ and $Z\subset NC\subset D$ as actual
Cartier subschemes. Fix one ordered line tuple (D.HU.8).
On $D$ the nonnegative degrees of its restrictions satisfy
\[
 \deg_D(L_a)=L\bigl(c_1(F_a)\cdot C-rZ\cdot C\bigr)
                                                =O_R(tr).
                                                               \tag{D.HU.10}
\]
In conjunction with (D.HU.5), the regular-section exact
sequence for a generated line gives
$h^0(D,L_a)+h^1(D,L_a)=O_R(tr)$.

Use the actual quotient presentation of
Proposition~\ref{D:hereditary-ordinary-unit-finite-cycle-Hom}:
on the generating-tuple open locus in
$H^0(D,L_a)^{r+1}$ there is a locally split exact sequence
\[
 0\longrightarrow L_a^{-1}\xrightarrow{s_a}
       \mathcal O_D^{r+1}\longrightarrow V_a
       \longrightarrow0,
       \qquad E_a=V_a(Z).                               \tag{D.HU.11}
\]
The universal quotients are actual vector bundles after
every coefficient change, and their determinants are the
specified $L_a$. Every geometric canonical bundle
$F_a(-Z)|_D$ occurs, because a generated rank-$r$ bundle
on this proper Cohen--Macaulay curve has $r+1$ general
generators. The total number of object coordinates is
$O_R(ftr^2)$.

\emph{Exact arrow coordinates and their degree bound.}
We must encode the finite data of
Proposition~\ref{D:finite-Frobenius-object-effectivity},
not merely independent extensions of objects. Introduce
actual maps on $D$
\[
 \alpha_a:F^*E_a\to E_{a+1},\qquad
 \beta_a:E_{a+1}\to(F^*E_a)(pC),                         \tag{D.HU.12}
\]
and impose both exact equalities
\[
 \alpha_a(pC)\beta_a=\iota_{E_{a+1}},\qquad
 \beta_a\alpha_a=\iota_{F^*E_a}.                         \tag{D.HU.13}
\]
These are the two Cartier inclusion maps on $D$; no
inversion of the Cartier equation on $D$ occurs.

Here is a full check of the coordinate bound for these
particular arrows. Frobenius pullback of (D.HU.11) is
exact over every ring because the sequence is locally
split. Its line term is $L_a^{-p}$, and its differential
is the $p$-power of the tuple $s_a$. The two residual
line twists in (D.HU.12) are respectively
\[
                    (1-p)Z,\qquad(p-1)Z+pC.              \tag{D.HU.14}
\]
The source, target and product Hom complexes consequently
involve only bounded powers and products of
$L_a$, $L_{a+1}$, $\mathcal O_D(-Z)$,
$\mathcal O_D(-C)$ and their inverses. These four kinds
of positive lines are generated, with degrees $O_R(tr)$,
$O_R(tr)$, $O_R(t)$ and $O_R(t)$ respectively.
Every needed line can therefore be written $P\otimes Q^{-1}$,
where $P,Q$ are generated and
$\deg P+\deg Q=O_R(tr)$.
Regular sections on the full nonreduced Cohen--Macaulay
curve give
\[
 h^0(D,PQ^{-1})+h^1(D,PQ^{-1})
 \le h^0(\mathcal O_D)+h^1(\mathcal O_D)+\deg P+\deg Q
                                  =O_R(tr).             \tag{D.HU.15}
\]

Apply the fixed two-affine Cech construction of
Proposition~\ref{D:fixed-determinant-finite-Hom} to these
fixed line terms. The Hom between the two-term quotient
complexes has four types of terms: a fixed-cycle line term with
$(r+1)^2$ copies and cohomology $O_R(t)$, two line terms with $r+1$ copies, and
one line term with one copy. By (D.HU.5) and (D.HU.15),
their total finite cohomology dimension is $O_R(tr^2)$.
All Cech splittings are chosen once over $k$, with no
variable cohomology bases. The transferred differential
and the inclusion/projection maps are the exact
all-ring polynomial maps of that proposition.
Substitution of the Frobenius tuples raises object
coefficients only to their $p$-th powers; the internal
complex width is still two. Thus the degrees of these
maps and of composition are bounded by a constant
depending on $p$ alone.

A degree-zero cycle in each finite Hom complex represents
an actual map (D.HU.12). Impose its cycle equation.
For (D.HU.13), use the exact transferred composition and
one degree-minus-one boundary witness for each equality.
This is equivalent over every coefficient ring to equality
of the actual maps, even when cohomology is not flat.
There are a fixed number of such complexes, maps and
length-two relations per phase. The combined coordinates
and scalar equations are therefore $O_R(ftr^2)$ and
have a uniform finite degree bound $D_R$.
The chosen lines, Cech splittings, Cartier sections and
all multiplication coefficients are fixed background
constants; they may depend on the cover and line tuple,
but they introduce no variable coordinates or degree cost.

Let $X_{\boldsymbol L}$ be the affine zero scheme of
these polynomial equations in the stated coordinates.
Let $S_{\boldsymbol L}$ be its open generating-tuple
locus. This is one open subset of one polynomial zero
scheme, with no uncontrolled finite union of charts.
The properness defining the open quotient locus is used
only to obtain an open subset; its equations do not enter
the component estimate.

\emph{Actual families and the class comparison.}
Every point of $S_{\boldsymbol L}$ is actual finite
bounded-pole Frobenius data (D.HU.11)--(D.HU.13).
Proposition~\ref{D:finite-Frobenius-object-effectivity}
produces its formal unit family and retains the specified
marking on $NC$. On affine parameter opens the family
algebraizes over the proper Noetherian completed base.
On overlaps the unique integral unit comparison inducing
the specified identity $N$-jet satisfies its cocycle
identity. Thus these local constructions glue, and no
parameter-cover cardinality is charged in the estimate.

Because $Z\subset NC$, the constructed family retains
the actual generating quotient
$\mathcal O_{Z}^{r+1}\to F_a(-Z)|_Z$ from (D.HU.11),
and its specified determinant line $L_a|_Z$.
Lemma~\ref{D:weak-cycle-generated-determinant-lifting}
therefore lifts these generators formally and makes
$\det(F_a(-Z))$ the fixed global line $L_a$, locally on
the parameter scheme. The resulting objects need not
be canonical full pullbacks. That property is unnecessary:
Proposition~\ref{D:generated-fixed-determinant-unit-openness}
applies to the actual family and proves that its punctured
unit classes have disjoint open and closed images.
The comparison pole bounds may depend on the fixed
reference class and rank. Only qualitative finite
presentation is used there; no parameter of its Isom
scheme is added to $X_{\boldsymbol L}$.
Consequently
\[
 \#\{\text{punctured unit classes represented by }
                       S_{\boldsymbol L}\}
          \le\#\operatorname{Irr}(S_{\boldsymbol L})
          \le\#\operatorname{Irr}(X_{\boldsymbol L}).    \tag{D.HU.16}
\]

Every original representation occurs in these families.
Indeed its canonical objects and actual arrows restrict
to (D.HU.11)--(D.HU.13) for one chosen determinant tuple.
The true formal object is one permissible choice of local
matrix lifts in the object-effectivity construction.
Its choice-independence supplies an actual integral unit
isomorphism to the constructed object, retaining the
prescribed $N$-jet. Thus coverage is coverage of the
original punctured unit class, not just of its finite
truncation.

\emph{Counting components.}
The affine component bound of
Proposition~\ref{D:weighted-affine-component-bound}, with
all variable weights equal to one, gives, with
$M=\dim(\text{ambient affine space})\le c_2ftr^2$,
\[
 \#\operatorname{Irr}(X_{\boldsymbol L})
            \le(M+1)\max\{2,D_R\}^{M}
            \le p^{c_3ftr^2}.                          \tag{D.HU.17}
\]
It applies to arbitrary rectangular systems and counts
positive-dimensional components, so no gauge quotient or
zero-dimensional assumption has been inserted.
Multiplying this bound by (D.HU.2), and using
$\log(r+1)=O(r^2)$ and $\log q=f\log p$, gives
\[
 q^{c_0t}\exp(c_1ft\log(r+1))\,p^{c_3ftr^2}
                                      \le q^{A_Rtr^2}.
\]
Ordinary $q$-unit classes on $U_A$ are exactly the
conjugacy classes of continuous $\mathbb F_q$
representations. This proves (D.HU.1).
\end{proof}

\begin{remark}[Scope of the hereditary conclusion]
The theorem counts ordinary defining-characteristic
linear representations on arbitrary finite punctured
covers. No projective determinant enumeration, twisted
gerbe atlas, factorial bound for unrestricted alternating
quotients, or full profinite finite-generation conclusion
is inferred. The proof avoids coflat rank strata by
constructing all bounded-pole unit objects represented
by its finite equations and using generated
fixed-determinant comparison families.
\end{remark}


% Proved after four complete actual-file audits. Exact crown consequence of the arbitrary-cover ordinary hereditary count.
\subsection{Faithful-linear crowns over arbitrary permutation quotients}

\begin{theorem}[Proved: arbitrary-quotient linear crowns and bounded residual outer actions]
\label{D:arbitrary-quotient-faithful-linear-crowns}
Let $R$ be a complete excellent normal local surface with
algebraically closed coefficient and residue field $k$ of
characteristic $p>0$, and put
$\Gamma=\pi_1(\operatorname{Spec}R\setminus\{\mathfrak m\})$.
Let $A_R$ be the constant in
Theorem~\ref{D:arbitrary-cover-ordinary-hereditary-growth}.
Fix an integer $b_0\ge1$ and a real number $\kappa>0$.

Let $S$ be a finite nonabelian simple group, identified
with its inner automorphism group, and suppose
\[
 S\le\mathcal B_0\mathrel{\trianglelefteq}\mathcal B
          \le\operatorname{Aut}(S),\qquad
 [\mathcal B:\mathcal B_0]\le b_0.                       \tag{D.AL.1}
\]
Assume there is a faithful linear embedding
\[
 \mathcal B_0\hookrightarrow\operatorname{GL}_r(\mathbb F_q),
 \qquad q=p^f,\qquad
                   r^2\log q\le\kappa\log|S|.           \tag{D.AL.2}
\]
No embedding of $\mathcal B$ into an ambient matrix group
is required.

Let $\Omega$ have $t$ elements and let
\[
 P\le(\mathcal B/S)^\Omega\rtimes\operatorname{Sym}(\Omega)
                                                        \tag{D.AL.3}
\]
have transitive permutation action. Let $L$ be its full
inverse image in
$\mathcal B^\Omega\rtimes\operatorname{Sym}(\Omega)$,
and assume $L$ is monolithic primitive with socle
$N=S^\Omega$. Write $L_k$ for its $k$th crown-based power,
that is the fibre product of $k$ copies of $L$ over $P$.
For every continuous surjection $\Gamma\twoheadrightarrow L_k$
one has
\[
 \begin{split}
 k&\le |N|^{c_R},\\
 d(L_k)&\le\max\{2,d(P),c_R+2\},\\
 c_R&=\max\{1,\lceil\kappa(A_Rb_0+1)+b_0-1\rceil\}.
 \end{split}                                            \tag{D.AL.4}
\]
There is no ramification or resolution-extension condition
on the induced quotient $\Gamma\twoheadrightarrow P$.
Its generator number $d(P)$ is retained explicitly.

In particular, take $\mathcal B_0=\mathcal B=S$.
For every faithful $S\hookrightarrow\operatorname{GL}_r(\mathbb F_q)$
satisfying (D.AL.2), every transitive
$P\le\operatorname{Sym}(t)$, and every quotient
\[
       \Gamma\twoheadrightarrow (S^t)^k\rtimes P
                                                        \tag{D.AL.5}
\]
with pure permutation action, the same conclusions hold
with $c_R=\max\{1,\lceil\kappa(A_R+1)\rceil\}$.
If a chosen subfamily has $d(P)\le d_0$, then its crowns
have at most $\max\{2,d_0,c_R+2\}$ generators. No uniform
bound on $d(P)$, and no full finite-generation conclusion
for $\Gamma$, is asserted without that hypothesis.
\end{theorem}

\begin{proof}
Fix the induced quotient $\beta:\Gamma\twoheadrightarrow P$.
Choose $o\in\Omega$. Let $H$ be the inverse image in
$\Gamma$ of the point stabilizer in $P$, and let
$\chi:H\to\mathcal B/S$ be the $o$-coordinate of
$\beta|_H$. Then $[\Gamma:H]=t$.

\emph{The exact restriction and extension count.}
The transport calculation of
Lemma~\ref{D:nonsplit-automorphism-wreath-lifts}, with
transports chosen in $\mathcal B$, gives
\[
 \frac{\{\Gamma\to L\text{ lifting }\beta\}}
            {N\text{-conjugacy}}
 \simeq
 \frac{\{\lambda:H\to\mathcal B\text{ lifting }\chi\}}
            {S\text{-conjugacy}}.                       \tag{D.AL.6}
\]
For precision, choose coset transports $g_x$ with $g_xo=x$
and lifts $b_x\in\mathcal B$ of the corresponding outer
coordinates. A stabilizer lift $\lambda$ gives the
$x$-coordinate
\[
       b_x\lambda(g_x^{-1}\gamma g_{\gamma^{-1}x})
                                           b_{\gamma^{-1}x}^{-1}.
                                                        \tag{D.AL.7}
\]
The two transport products multiply to the product for
$\gamma\eta$, giving the wreath homomorphism; evaluation
at $o$ is its inverse modulo the stated conjugacies.
The checked calculation does not require a splitting of
$1\to N\to L\to P\to1$, or any condition on the cover
associated with $\beta$.

The $k$ coordinates of a surjection to $L_k$ have distinct
classes in (D.AL.6). Otherwise two would be conjugate by
$N$, so their image would be contained in a proper
conjugate diagonal of $L\times_P L$, contradicting
surjectivity. It therefore suffices to bound the set
on the right of (D.AL.6).

Put
\[
 H_0=\chi^{-1}(\mathcal B_0/S),\qquad a=[H:H_0]\le b_0.
                                                        \tag{D.AL.8}
\]
The subgroup is normal in $H$ and depends only on the
fixed outer character, not on a chosen lift. Its
connected punctured cover has degree $at$ over $R$.
Theorem~\ref{D:arbitrary-cover-ordinary-hereditary-growth}
therefore applies with no condition on that cover and gives
\[
 \#\bigl(\operatorname{Hom}(H_0,\operatorname{GL}_r(\mathbb F_q))
              /\operatorname{GL}_r(\mathbb F_q)\bigr)
                                      \le q^{A_Rat r^2}.
                                                        \tag{D.AL.9}
\]
Within one ambient orbit there are at most
$|\operatorname{GL}_r(\mathbb F_q)|\le q^{r^2}$
$S$-conjugacy classes of actual $\mathcal B_0$-valued
restrictions. Indeed fix one homomorphism in the ambient
orbit and list all its conjugates by the finite matrix
group. Every actual restriction in that orbit appears
in the list before quotienting by $S$. The matrices in
this list are not required to normalize $\mathcal B_0$
or to extend to automorphisms of $\mathcal B$.
Thus there are at most
\[
                         q^{(A_Rat+1)r^2}              \tag{D.AL.10}
\]
restriction classes modulo $S$.

Choose an actual $\mathcal B_0$-valued representative
for each such restriction class. An $S$-conjugation
normalizes a restriction to that representative and
preserves both the $\mathcal B$-valued condition and
the prescribed outer character. Choose representatives
for the $a-1$ nonidentity cosets of $H_0$ in $H$.
Each representative must map into the one coset of $S$
specified by $\chi$, so it has at most $|S|$ possible
images. The fixed restriction and these images determine
the homomorphism. Relations can only remove choices.
The extension cost is therefore at most $|S|^{a-1}$,
without a bound on $|\mathcal B|$.

Combining these estimates and (D.AL.2) yields
\[
 \begin{split}
 k&\le q^{(A_Rat+1)r^2}|S|^{a-1}\\
  &\le |S|^{\kappa(A_Rat+1)+a-1}\\
  &\le |S^\Omega|^{\kappa(A_Rb_0+1)+b_0-1}
   \le |N|^{c_R}.
 \end{split}                                            \tag{D.AL.11}
\]
Only $a\le b_0$ and $t\ge1$ were used in the third
line. In particular no generator estimate for $P$ has
entered the multiplicity bound.

\emph{Generators with the quotient term retained.}
Use the exact nonabelian crown estimate already checked
in Theorem~\ref{D:unramified-permutation-linear-crowns}
and Corollary~\ref{D:normal-subgroup-projective-crowns}.
For $s\ge d(L)=\max\{2,d(P)\}$, it gives
\[
 \psi_L(s)\ge\frac{|N|^{s-1}}
                         {2t|\operatorname{Out}(S)|},
 \qquad k\le\psi_L(s)\ \Longrightarrow\ d(L_k)\le s.
                                                        \tag{D.AL.12}
\]
This is the complete-proof input from
\citet[Proposition 2.6(ii), equations (2.1)--(2.2),
pp.~11--13]{LucchiniMarionTracey2020D}, with its stated
classification dependence. The same source records
$|\operatorname{Out}(S)|\le\log_2|S|$. Since $S$ is
nonabelian simple, $|S|\ge60$; hence
\[
                  2t\log_2|S|\le |S|^t=|N|.             \tag{D.AL.13}
\]
For example the $t=1$ inequality holds for every integer
$|S|\ge60$, and passage from $t$ to $t+1$ multiplies
the right side by at least $60$ while multiplying the
left side by at most two. Thus
$\psi_L(s)\ge|N|^{s-2}$.
Taking $s=\max\{2,d(P),c_R+2\}$ in (D.AL.12) and using
(D.AL.11) proves the second conclusion of (D.AL.4).

For the pure permutation specialization, $N=S^t$ is
minimal normal by transitivity, and its centralizer in
$N\rtimes P$ is trivial. Any additional minimal normal
subgroup would centralize $N$, so $N$ is the unique
minimal normal subgroup. A nonabelian minimal normal
subgroup is not contained in the Frattini subgroup;
a maximal subgroup not containing it is therefore
core-free. Thus the group is monolithic primitive,
so it meets the stated crown hypotheses. Substituting
$b_0=1$ in (D.AL.11)--(D.AL.12) gives (D.AL.5) and the
asserted conditional uniform generator bound.
\end{proof}

\begin{remark}[Concrete scope and the residual gate]
The previously checked natural linear and order estimates
supply, for instance, $S=\operatorname{SL}_{p^u}(\mathbb F_{p^f})$
whenever this group is nonabelian simple, with the natural
rank $r=p^u$ and $\kappa=4$. Its scalar centre is trivial
because $\gcd(p^u,p^f-1)=1$; the nonsimple small parameters
are excluded. This is an unbounded-rank instance of the
pure permutation case, using the same audited recognition
as Theorem~\ref{D:unramified-field-permutation-crowns}.

The new scope is every complete normal surface germ and
every induced permutation or displayed bounded-outer
quotient, including those ramified above a resolution.
The earlier radial-tame and one-curve restrictions are
absent. The fixed faithful linear embedding and its
quadratic order comparison remain essential. Arbitrary
projective representations, field-semilinear actions
not represented by (D.AL.2), and unrestricted alternating
socles are not included. The explicit term $d(P)$ also
prevents reading this result as a proof of full finite
generation of $\Gamma$.
\end{remark}


% Post-CP25 proved after four complete actual-file audits. Hereditary linear counts applied to actual imprimitive permutation towers.
\subsection{Iterated imprimitive towers in the alternating count criterion}

\begin{proposition}[Proved: hereditary counting through imprimitive wreath towers]
\label{D:hereditary-wreath-permutation-count}
Let $\Lambda$ be the fundamental group of the punctured
spectrum of a fixed complete excellent normal local surface
over an algebraically closed field of characteristic $p>0$.
For every integer $m\ge2$, choose a finite collection
\[
 P_{m,j}\le S_m,\qquad
 P_{m,j}\hookrightarrow
       \operatorname{GL}_{r_{m,j}}(\mathbb F_{q_{m,j}}),
                 \quad j\in\mathcal J_m,                \tag{D.WO.1}
\]
where $q_{m,j}$ is a power of $p$ and the embeddings are
part of the fixed models. Suppose that one constant $B$
satisfies
\[
 \log\max\{1,|\mathcal J_m|\}\le Bm\log(m+1),\qquad
 r_{m,j}^2\log q_{m,j}\le Bm\log(m+1).                  \tag{D.WO.2}
\]
Empty model collections are allowed.

A tower certificate of degree $N\ge2$ consists of an ordered
factorization $N=m_1\cdots m_h$, $m_i\ge2$, together with,
at each level, either a chosen model $P_i=P_{m_i,j_i}$
or the unrestricted group $P_i=S_{m_i}$. Call these levels
modelled and unrestricted respectively. Put
\[
 T_i=\prod_{j>i}m_j,\qquad
 W_h=P_h,\quad W_i=P_i\wr W_{i+1}\quad(i<h),\qquad
 \mathcal E=\sum_{i\text{ unrestricted}}T_i m_i^2.
                                                        \tag{D.WO.3}
\]
The wreath products use their standard imprimitive actions,
so $W_1\le S_N$. A transitive $\Lambda$-set admits this
certificate if its action is permutation-conjugate into
$W_1$. Define $e_{\mathcal J}(X)$ to be the minimum of
$\mathcal E$ over its certificates. A singleton has cost
zero. Every transitive $N$-set has a certificate and
\[
                     0\le e_{\mathcal J}(X)\le N^2,     \tag{D.WO.4}
\]
by using the one-level unrestricted model $S_N$.

There is a constant $C$, depending only on $\Lambda$ and
$B$, such that, for all $N\ge2$ and integers $M\ge0$,
\[
 \#\{\text{transitive $N$-element $\Lambda$-sets up to
 isomorphism}:e_{\mathcal J}(X)\le M\}
       \le\exp\bigl(C(N\log(N+1)+M)\bigr).             \tag{D.WO.5}
\]
In particular, arbitrary-height iterated modelled wreath
actions have a factorial-order count without an orbit-size
restriction. The estimate concerns actual permutation
sets, not only their associated linear representations.
\end{proposition}

\begin{proof}
Let $A$ be the constant of
Theorem~\ref{D:arbitrary-cover-ordinary-hereditary-growth}
for this fixed local source. For every open subgroup
$H\le\Lambda$ of index $T$ and every fixed injection
$P\hookrightarrow\operatorname{GL}_r(\mathbb F_q)$,
that theorem and the size of an ambient conjugacy orbit give
\[
 |\operatorname{Hom}_{\rm cont}(H,P)|
       \le q^{(AT+1)r^2}\le q^{(A+1)Tr^2}.              \tag{D.WO.6}
\]
The first inequality counts actual homomorphisms to $P$:
its injection makes them a subset of the ambient actual
Hom set. It does not identify permutation-conjugacy classes
with linear-conjugacy classes. For an unrestricted level
use the faithful permutation-matrix inclusion
$S_m\hookrightarrow\operatorname{GL}_m(\mathbb F_p)$.
Its cost in the logarithm of (D.WO.6) is at most
$(A+1)T m^2\log p$.

We give the exact induction through the tower. Fix a
transitive homomorphism
$\gamma:\Lambda\to W_{i+1}\le S_{T_i}$ and one point
$o$ of its $T_i$-element set. Put
$H=\gamma^{-1}((W_{i+1})_o)$; then $[\Lambda:H]=T_i$.
Nonabelian Shapiro for the pure permutation base gives
\[
 \frac{\{\Lambda\to P_i^{T_i}\rtimes W_{i+1}
                       \text{ lifting }\gamma\}}
           {P_i^{T_i}\text{-conjugacy}}
       \simeq
       \frac{\operatorname{Hom}_{\rm cont}(H,P_i)}{P_i}.
                                                        \tag{D.WO.7}
\]
This formula does not require simplicity of $P_i$.
To verify it directly, choose $g_x\in\Lambda$ with
$\gamma(g_x)o=x$ and $g_o=1$. A homomorphism
$\lambda:H\to P_i$ supplies the base coordinate
\[
                 \lambda(g_x^{-1}g g_{\gamma(g)^{-1}x})
                                                        \tag{D.WO.8}
\]
at $x$ for $g\in\Lambda$. The inner word stabilizes $o$.
Multiplication for $g g'$ follows by cancellation of the
middle transport. Evaluation on $H$ at $o$ is inverse
after normalizing the transports by base conjugation.
The remaining choice at $o$ is exactly $P_i$-conjugation.
This proves (D.WO.7), including continuity, since all
targets are finite and the displayed coset partition is
finite and open.

For each $W_{i+1}$-conjugacy class of transitive top
homomorphisms choose one actual representative $\gamma$.
Every lift whose top is in that class can be conjugated
to a lift of this representative, using the canonical
$W_{i+1}$ subgroup of the wreath product. Thus the number
of $W_i$-conjugacy classes above it is at most the right
side of (D.WO.7). Lifts whose total permutation action is
not transitive may be discarded. Induction gives, for
one fixed tower, a logarithmic bound
\[
 (A+1)\sum_{i\text{ modelled}}
                  T_i r_i^2\log q_i
       +(A+1)\log p\,\mathcal E.                       \tag{D.WO.9}
\]
This bounds wreath-conjugacy classes, hence also the
permutation-set classes obtained by enlarging conjugacy
to $S_N$. Every subgroup used in (D.WO.6) is an actual
degree-$T_i$ punctured cover of the original fixed source.
There is no radial-tameness hypothesis in this application.

Since $T_i m_i\le N$, (D.WO.2) implies
\[
 \sum_{i\text{ modelled}}T_i r_i^2\log q_i
 \le BN\sum_i\log(m_i+1)
 \le cB N\sum_i\log m_i
       =cB N\log N,                                    \tag{D.WO.10}
\]
where one may take $c=\log 3/\log 2$.
The model choices for a fixed ordered factorization have
logarithm at most
\[
       B\sum_i m_i\log(m_i+1)+h\log2
                                      =O_B(N\log N).    \tag{D.WO.11}
\]
The extra term permits either a modelled or unrestricted
level. The number of ordered factorizations is at most
$(1+\log_2N)N^{\lfloor\log_2N\rfloor}$: its length is
at most $\log_2N$, and this estimate overcounts by allowing
every integer between two and $N$ in every slot.
Its logarithm is $O(N\log(N+1))$.

Retain only certificates with $\mathcal E\le M$.
Equations (D.WO.9)--(D.WO.11), followed by this finite
union bound, prove (D.WO.5). There is no independent
factor for labelling a block system at each stage.
A tower with given degrees has one standard imprimitive
set; any identification with it merely produces another
actual homomorphism into the same fixed wreath model,
which is already covered by (D.WO.7).
\end{proof}

\begin{theorem}[Proved: alternating counts with local tower energy]
\label{D:alternating-local-wreath-energy}
Use the fixed complete-local surface detecting groups and
global finite-puncture geometric data of
Corollary~\ref{D:alternating-dimensional-normalization-closure}.
For each of these finitely many whole local groups choose
catalogues satisfying (D.WO.1)--(D.WO.2). Retain all
previously admissible local orbits of
Theorem~\ref{D:alternating-compressed-local-orbits},
including its fixed finitely generated quotient and
single matrix-model choices. Give these cost zero.
For every other local orbit use its tower cost
$e_{\mathcal J}$ from (D.WO.3)--(D.WO.4).

For any fixed $K>0$, degree-$n$ finite etale covers whose
whole detecting-local actions satisfy
\[
       \sum_{s}\ \sum_{O\in\Lambda_s\backslash X}
                         e_s(O)\le K n\log(n+1)        \tag{D.WO.12}
\]
have at most $n^{An}$ isomorphism classes, including
disconnected covers. Here $A$ depends only on the fixed
geometric data, catalogues, quotient bounds, and $K$.
Their distinct natural $A_n$ quotient kernels have
multiplicity at most $|A_n|^c$ for one fixed $c$, and
their selected alternating simple-product quotient is
topologically finitely generated.

The catalogues can be chosen so that an orbit has cost
zero whenever it admits a complete block decomposition
whose primitive constituents are all either
\begin{enumerate}
\item affine groups in their natural actions in
      characteristic $p$, contained in
      $\operatorname{AGL}_d(\mathbb F_{p^a})$; or
\item the faithful primitive groups with nonabelian
      socles permitted by
      Corollary~\ref{D:primitive-defining-characteristic-orbit-models},
      for one fixed finite collection of simple exceptions.
\end{enumerate}
There is no restriction on the height of that block
decomposition. More generally only the weighted squared
degrees of its unrestricted levels need be charged.
The one-level certificate in (D.WO.4) shows that this
condition includes the earlier whole-orbit square-sum
criterion.

All actions and block constituents in this statement
belong to the whole fixed complete-local puncture groups
or their actual finite-index stabilizers. No inference
is made from the smaller curve inertia actions alone.
This theorem uses the finite omitted-set and its already
proved dimensional/finite-normalization scope. It does
not assert a hereditary theorem for original-divisorial
local-pair quotients.
\end{theorem}

\begin{proof}
We first check the named zero-cost collections without
adding a classification assumption to the counting
proposition. The affine groups have the faithful
$(d+1)$-dimensional matrix realization over
$\mathbb F_{p^a}$ from
Lemma~\ref{D:permutation-matrix-model-count}; their
degree-$m$ catalogue and matrix cost satisfy (D.WO.2).
Subgroups in that fixed action are automatically covered
by the ambient affine model.

For the second collection, the cited primitive recognition
gives a faithful $\mathbb F_p$ representation of dimension
$D_m\le c(1+\log m)^2$. Enumerate the actual subgroup
and core-free point-stabilizer pairs in
$\operatorname{GL}_{\lceil D_m\rceil}(\mathbb F_p)$
as in Lemma~\ref{D:intrinsic-small-linear-permutation-images}.
This is a finite collection of degree-$m$ permutation
models containing every such primitive action, with
logarithmic model count $O((1+\log(m+1))^8)$ and
matrix cost $O((1+\log(m+1))^4)$.
Both are bounded by $B m\log(m+1)$ for one fixed $B$,
since each fixed power of a logarithm divided by
$m\log(m+1)$ is bounded on the integers $m\ge2$.
The known primitive-order/classification scope in the
cited recognition is unchanged.

For completeness, a transitive permutation action with
such primitive block constituents embeds in the stated
iterated imprimitive wreath product. Choose a minimal
nontrivial block. Its setwise stabilizer acts transitively
and primitively on it: a smaller stabilizer block would
have disjoint translates outside it and would itself be
a smaller block of the whole group. Choose transports
from one block to every other block. Restriction of a
permutation to these transported copies, together with
its action on the block set, gives the faithful wreath
embedding by the same calculation as (D.WO.8).
Repeat with the transitive action on the block set.
At a primitive stage stop. This terminates, and uses
one permitted model at each stage, with no unrestricted
level. Conversely the counting assertion only needs the
existence of a certificate, so no uniqueness of this
block decomposition is required.

We now supply the changed local enumeration. The number
of fixed local groups is finite, say $w$. The previous
zero-cost orbits are counted exactly as in
Theorem~\ref{D:alternating-compressed-local-orbits}:
orbits assigned to the same fixed finitely generated
quotient may be joined and counted by its fixed generators,
while the single matrix-model orbit types have their
factorial-order bound. Assigning an ambiguous orbit to
one type only overcounts. Ordered orbit-size compositions
and these finite tags cost $\exp(O(n\log(n+1)))$ over
all $w$ groups.

Put $M=\lfloor K n\log(n+1)\rfloor$.
For the remaining orbits of degrees $d_i$, assign integer
budgets $M_i\in\{0,\ldots,M\}$ with
$\sum_i M_i\le M$. There are at most $(M+1)^{wn}$ such
assignments, whose logarithm is $O_K(n\log(n+1))$.
Proposition~\ref{D:hereditary-wreath-permutation-count}
then bounds their types by
\[
 \exp\!\left(C\sum_i d_i\log(d_i+1)+C\sum_i M_i\right)
                           \le\exp(O(n\log(n+1))).     \tag{D.WO.13}
\]
Singletons contribute their unique trivial type. Thus
the complete whole-source local permutation data still
have a factorial-order bound.

Every subsequent geometric step of the existing
finite-puncture proof is unchanged. The one fixed
detecting curve sees only restrictions of these whole
local types. The uniform curvewise different estimate
still gives total discriminant $O(n)$. For each fixed
tuple of full local extensions, the proved formal-fibre
bound is $(n!)^{2g}n^{2D}$ with $D=O(n)$.
If the detecting curve has $b_C$ marked points, the
disconnected factor assignments cost at most
$n^{b_Cn+1}$; finite normalization descent contributes
at most $(n!)^a$ for a fixed integer $a$. The fixed
detecting-surface maps in the dimensional extension
retain the same injection of original-cover classes.
Multiplying these factors by (D.WO.13) proves $n^{An}$.
These are precisely the audited comparisons of
Corollary~\ref{D:alternating-dimensional-normalization-closure}
and Theorem~\ref{D:alternating-compressed-local-orbits}.

Finally the distinct natural alternating quotient
kernels inject into the counted permutation classes.
The same elementary factorial comparison and the
nonabelian simple-power generator bound used in
Theorem~\ref{D:finite-puncture-local-orbit-factorial-count}
give multiplicity $|A_n|^c$ uniformly for $n\ge5$.
Goursat's lemma and compactness identify the quotient
selected by these kernels with its product of simple
powers, and give its stated finite generation.
No assertion about unselected alternating kernels or
the entire fundamental group is made.
\end{proof}


% Separate post-CP25 proved input. Full B/root/C/D actual-file audits PASS; twisted scalar H1 only.
\subsection{Scalar determinants for arbitrary field characters}

\begin{lemma}[Proved: twisted scalar classes on arbitrary punctured covers]
\label{D:arbitrary-field-character-scalar-H1}
Let $R$ be a complete excellent normal local surface with
algebraically closed coefficient and residue field $k$ of
characteristic $p>0$, and put
$\Gamma=\pi_1(\operatorname{Spec}R\setminus\{\mathfrak m\})$.
Let $H\subset\Gamma$ be open of index $t$, let $q=p^f$,
and fix a continuous field character
\[
 \chi:H\longrightarrow\operatorname{Gal}(\mathbb F_q/\mathbb F_p).
                                                               \tag{D.WD.1}
\]
Write $d=|\operatorname{im}\chi|$, $q_0=p^{f/d}$, and
$n=q-1$. Let $A_V$ be the normal complete local ring of
the connected punctured cover corresponding to $\ker\chi$;
its degree over $R$ is $td$. Let $v_V$ be the number of
exceptional components on its minimal regular resolution.
Use $v_V=0$ if $A_V$ is regular. Then
\[
 \bigl|H^1(H,\mathbb F_q^{\times}(\chi))\bigr|
       \le(q_0+1)^{2p_g(A_V)+v_V}
       \le q^{c_Rt},                                    \tag{D.WD.2}
\]
where $c_R$ is independent of $H,q,\chi$ and $d$.
The action on the coefficient group is the prescribed field
action. There is no tameness assumption, and $d$ is allowed
to be divisible by $p$.

In particular (D.WD.2) bounds the possible twisted scalar
determinant classes of rank-$r$ representations into
$\operatorname{GL}_r(\mathbb F_q)\rtimes\operatorname{im}\chi$
with fixed field character, independently of $r$.
It does not assert a representation count for that group.
\end{lemma}

\begin{proof}
Put $K=\ker\chi$ and $M=\mathbb F_q^\times$.
Choose the generator of $H/K$ whose action on $M$ is
$z\mapsto z^{q_0}$. Finite cyclic Hilbert 90 here is
an elementary calculation: $M$ is cyclic of order
$n=q_0^d-1$, the norm is the power
$n/(q_0-1)$, and both its kernel and the image of the
power $q_0-1$ have order $n/(q_0-1)$. Thus
$H^1(H/K,M)=0$. Inflation--restriction gives an injection
\[
 H^1(H,M(\chi))\lhook\joinrel\longrightarrow
     \ker\bigl(q_0-\tau^*:H^1(U_V,\boldsymbol\mu_n)
                              \to H^1(U_V,\boldsymbol\mu_n)\bigr).
                                                               \tag{D.WD.3}
\]
The identification $M=\mu_n(k)$ is fixed once. We choose
the deck-pullback orientation of $\tau$ to make (D.WD.3)
the displayed equation; replacing the deck generator by
its inverse changes none of the arguments below.
The deck action is $k$-linear, so it acts trivially on
the constant group $\mu_n(k)$. This distinguishes its
geometric action from the prescribed field action $q_0$.

\emph{An equivariant resolution with the existing degree budget.}
Take the minimal regular resolution $Y\to\operatorname{Spec}A_V$.
Every automorphism of $A_V$ lifts uniquely to $Y$.
We recall why the uniqueness statement does not require
the group order to be invertible. Resolve the graph between
two resolutions with no exceptional $(-1)$-curves, and
factor both projections into point blowups. On a common
resolution minimal with respect to their total number,
a last exceptional $(-1)$-curve for one projection cannot
be contracted by the other, since then its contraction
would factor both projections. If it is not contracted,
its image in the other resolution has self-intersection
at least $-1$. Negative definiteness makes that square
exactly $-1$; the blowdown formula then says that no centre
met the curve, so its image is a smooth rational
$(-1)$-curve, a contradiction. Thus both projections are
isomorphisms. Existence is obtained by contracting the
exceptional $(-1)$-curves of a resolution. The point-blowup
factorization used here is
\citet[Lemma 54.17.1, Tag 0C5R]{StacksHereditaryCartierD};
the standard contraction and negative-definiteness inputs
are the same as in the already checked minimal-resolution
and all-complete budget proofs. This also proves coherence
of the lifted finite cyclic action by uniqueness.

Lemma~\ref{D:fixed-henselian-algebraization} and
Proposition~\ref{D:all-complete-cover-geometric-budgets}
give
\[
                         2p_g(A_V)+v_V=O_R(td).          \tag{D.WD.4}
\]
If $A_V$ is regular, codimension-two purity makes
$H^1(U_V,\mu_n)=0$, and there is nothing further to prove.
Otherwise let $E=Y_{\mathrm{closed,red}}$. It is a connected
proper reduced curve. Proper base change over the strictly
henselian local base gives, equivariantly,
\[
                    H^1(Y,\mu_n)\simeq H^1(E,\mu_n).
                                                               \tag{D.WD.5}
\]
Here and below cohomology of finite coefficients is etale.
The exact proper-base-change input is
\citet[Theorem 59.91.11, Tag 095T]{StacksSemilinearDeterminantsB}.

\emph{The localization and normalization filtration.}
Remove from $Y$ the finite, $\tau$-stable set consisting
of singular and intersection points of $E$.
Codimension-two purity leaves $H^1(Y,\mu_n)$ unchanged.
The remaining exceptional components are disjoint smooth
curves. The degree-one localization sequence and
codimension-one purity therefore give an equivariant
exact sequence
\[
 0\longrightarrow H^1(E,\mu_n)
  \longrightarrow H^1(U_V,\mu_n)
  \longrightarrow J_n\longrightarrow0,
 \qquad J_n\subset(\mathbb Z/n)^{v_V}.                 \tag{D.WD.6}
\]
The action on the right-hand ambient module is the
permutation of the exceptional components; the Gysin
orientation is canonical and the deck action is $k$-linear.
These are the same precise purity and localization inputs
used in Lemma~\ref{D:hereditary-ordinary-unit-determinants};
no claim about the wild kernel being pro-$p$ is involved.

Let $\nu:E^\nu\to E$ be normalization, let $g_i$ be the
genera of its smooth proper components, and let $G_E$
be the bipartite incidence graph with vertices the
components and the singular points, and one edge for
each branch. The exact sheaf sequence
$0\to\mu_n\to\nu_*\mu_n\to\mathcal Q_n\to0$
has finite support in its last term. At a singular point
with $b$ branches that term is $(\mathbb Z/n)^b$ modulo
its diagonal copy. Its cohomology gives the equivariant
exact sequence
\[
 0\longrightarrow H^1(G_E,\mathbb Z)\otimes\mathbb Z/n
 \longrightarrow H^1(E,\mu_n)
 \longrightarrow\prod_i\operatorname{Jac}(E_i^\nu)[n](k)
 \longrightarrow0.                                    \tag{D.WD.7}
\]
The graph incidence cokernel is a free integral lattice,
so this calculation is valid for every $n$, not just
for a field of coefficients. Kummer theory identifies
the smooth-component groups with the displayed Jacobian
torsion, canonically
\citep[Lemma 59.69.1, Tag 03RQ]{StacksLocalPairCohomologyB}.
The cyclic action permutes the components and branches
and induces an automorphism of finite order on the
product Jacobian, preserving its identity.

Writing $b_1=b_1(G_E)$, the structure-sheaf normalization
sequence gives
\[
 2\sum_i g_i+b_1\le2h^1(E,\mathcal O_E)
                         \le2p_g(A_V).                 \tag{D.WD.8}
\]
Indeed the normalization defect at a point is at least
its number of branches minus one; this proves the first
inequality directly from the two normalization sequences.
The second follows from
$H^1(Y,\mathcal O_Y)\twoheadrightarrow H^1(E,\mathcal O_E)$,
since coherent $H^2$ over the affine surface vanishes.

\emph{The finite-order kernel estimate.}
Let $T$ be a finite-order automorphism of an integral
lattice $\Lambda$ of rank $b$. Smith normal form gives
\[
 |\ker(q_0-T:\Lambda/n\Lambda\to\Lambda/n\Lambda)|
 \le|\det(q_0-T)|\le(q_0+1)^b.                         \tag{D.WD.9}
\]
The last inequality holds because the complex roots of
the characteristic polynomial of $T$ are roots of unity.
It applies to the graph lattice and to the exceptional
permutation lattice. For the subgroup $J_n$ the same
bound follows by inclusion in that permutation lattice.

If $J$ is an abelian variety of dimension $g$ and $T$
is an automorphism with $T^d=1$, then $[q_0]-T$ is an
isogeny: its product with
$[q_0]^{d-1}+[q_0]^{d-2}T+\cdots+T^{d-1}$ is $[n]$.
The characteristic-polynomial/degree theorem and its
Tate-module identification show
\[
 |\ker([q_0]-T:J[n](k)\to J[n](k))|
 \le\deg([q_0]-T)
 =\left|\det(q_0-T\mid V_\ell J)\right|
 \le(q_0+1)^{2g}.                                      \tag{D.WD.10}
\]
We use precisely
\citet[Theorem 10.9 and Proposition 10.20, with the written
proofs through Lemmas 10.21--10.22, pp.~46--52]{Milne2008AVScalarB}.
No semisimplicity in characteristic dividing $d$ is used:
the Tate module has characteristic zero, and kernel sizes
on finite invariant extensions satisfy
$|\ker u_M|\le|\ker u_{M_1}|\,|\ker u_{M_2}|$ for
$0\to M_1\to M\to M_2\to0$.

Apply this last observation to (D.WD.6)--(D.WD.7), using
(D.WD.8)--(D.WD.10). It bounds the right side of (D.WD.3)
by $(q_0+1)^{2p_g(A_V)+v_V}$, proving the first inequality
of (D.WD.2). Finally $q_0+1\le q_0^2$, (D.WD.4), and
$q_0^d=q$ give the second inequality with one constant
$c_R$. For a semilinear matrix representation, determinant
is a crossed homomorphism into $M(\chi)$; conjugation
changes it by a coboundary. Thus the stated determinant
consequence is exactly the same cohomology group.
\end{proof}


% Proved after four complete actual-file and exact-source audits. Arbitrary covers and arbitrary field characters.
\subsection{Hereditary semilinear growth with arbitrary field monodromy}

\begin{lemma}[Proved: equivariant comparable cycles on the field cover]
\label{D:equivariant-cover-comparable-cycles}
Retain the fixed complete normal surface, base resolution $Y_0$,
parameter cycle $C_0$, and constants of
Proposition~\ref{D:uniform-comparable-reduction-cycle}.
Let $U_V\to U_R$ be a connected finite etale cover of degree $s$,
and let a finite group $G$ act on this cover over $U_R$.
The regular model used in that proposition can be chosen with
the lifted $G$-action. On this same model the actual cycles
\[
 C=g^*C_0,\qquad Z=\alpha C-\beta R_g,\qquad
                 C\le Z\le\alpha C                         \tag{D.SH.1}
\]
are $G$-invariant. All scalar, generation, and canonical
Frobenius bounds of that proposition and
Theorem~\ref{D:arbitrary-cover-ordinary-hereditary-growth}
hold with their original constants and cover degree $s$.
No invertibility assumption on $|G|$ is needed.
\end{lemma}

\begin{proof}
Let $X$ be the finite normalization of $Y_0$ in the cover's
function field. The action extends to $X$ by the integral
closure property. The model in the comparable-cycle theorem
is the minimal regular resolution of $X$: a curve maps to a
point of $Y_0$ exactly when it is exceptional over $X$, since
$X\to Y_0$ is finite. Thus contracting precisely the
point-mapped exceptional curves of the first kind gives
the relative-minimal model stipulated in that theorem.

We spell out the equivariance deduction from the standard
surface facts, rather than require an equivariant resolution
algorithm. Resolve a rational map between two minimal regular
resolutions of $X$ by a shortest sequence of point blowups.
The other projection is also a sequence of point blowups.
These are the known statements of
\citet[Lemmas 54.4.3 and 54.17.1, Tags 0C5H and 0C5R]
{StacksSkewMinimalModelD}. Let $E$ be the last exceptional
$(-1)$-curve of the first projection. If the other projection
contracts $E$, its factorization through that contraction
contradicts shortest length, by
\citet[Lemma 54.16.1, Tag 0C5J]{StacksSkewMinimalModelD}.
Otherwise its image is a curve exceptional over $X$.
Successive blowdowns increase its square from $-1$ by sums
of nonnegative squares. Negative definiteness of the
exceptional intersection form makes the final square
negative, so every increase is zero. Hence no blowdown
centre meets the curve; its final image is still a smooth
rational $(-1)$-curve. This contradicts minimality.
There were therefore no blowups, and interchanging the
resolutions shows that the map is an isomorphism.
An isomorphism over the normal base is unique because it is
the identity on the common dense regular locus.

Every automorphism of $X$ consequently lifts uniquely; the
group identities follow from uniqueness. This argument is
independent of the characteristic and of the group order.
The map $g:Y\to Y_0$ is invariant. Its actual effective
Jacobian divisor $R_g$ is invariant by functoriality of the
map of canonical sheaves. Thus both cycles in (D.SH.1) are
invariant. This is exactly the same relative-minimal model
on which the bound $R_g\le\lambda C$ and the all-complete
degree-linear geometric budgets were proved. No additional
equivariant blowups, and no enlargement of their constants,
have been introduced.
\end{proof}

\begin{theorem}[Proved: arbitrary-cover hereditary semilinear growth]
\label{D:arbitrary-cover-semilinear-hereditary-growth}
Let $(R,\mathfrak m,k)$ be a complete excellent normal local
surface with algebraically closed coefficient and residue
field $k$ of characteristic $p>0$. Put
$\Gamma=\pi_1(\operatorname{Spec}R\setminus\{\mathfrak m\})$.
There is a constant $A_R^{\rm sl}$ such that, for every open
subgroup $H\subset\Gamma$ of index $t$, every $q=p^f$,
every $r\ge1$, and every fixed continuous character
\[
 \chi:H\longrightarrow
                 \operatorname{Gal}(\mathbb F_q/\mathbb F_p),
                                                               \tag{D.SH.2}
\]
the representations
\[
 H\longrightarrow \operatorname{GL}_r(\mathbb F_q)
          \rtimes\operatorname{Gal}(\mathbb F_q/\mathbb F_p)
 \quad\text{with field projection }\chi
                                                               \tag{D.SH.3}
\]
have at most $q^{A_R^{\rm sl}tr^2}$ conjugacy classes under
$\operatorname{GL}_r(\mathbb F_q)$.
There is no tameness condition on either cover, no bound
on the field-image order, and no assumption that the
field character extends over a resolution.
\end{theorem}

\begin{proof}
Write $d=|\operatorname{im}\chi|$, $\nu=f/d$, and
$q_0=p^\nu$. The subgroup $K=\ker\chi$ gives a connected
cyclic degree-$d$ cover $U_V\to U_H$. Its total degree over
$U_R$ is $td$. The finite normalization $A_V$ is complete,
normal, and local with residue field $k$. Choose the deck
generator $\tau$ with the pullback convention of
Lemma~\ref{D:twisted-deck-Frobenius-semilinear-model}.
By that proved category equivalence, (D.SH.3) is exactly
the category of bundles $E$ on $U_V$ with
\[
              \Phi:\tau^*F_{q_0}^*E\xrightarrow{\sim}E.
                                                               \tag{D.SH.4}
\]
The equivalence includes morphisms and arbitrary coefficient
base extension. The $d$-fold composite defines the ordinary
$F_q$-unit structure; it is not an independently imposed
matrix equation or a false order-$d$ relation on a chosen
element of $H$.

Use Lemma~\ref{D:equivariant-cover-comparable-cycles} for
the cyclic deck group on the degree-$td$ cover. We therefore
have a deck-equivariant model $Y$ and invariant $C,Z$ with
all the ordinary theorem's budgets $O_R(td)$.
In particular, $C\le Z\le\alpha C$,
$\mathcal O_Y(-C)$ and $\mathcal O_Y(-Z)$ are generated,
and $H^1(Z,\mathcal O_Z(-jZ))=0$ for every $j\ge1$.

\emph{Short Frobenius cycles and determinant choices.}
Factor (D.SH.4) into $\nu$ one-step arrows: for
$0\le a<\nu-1$ they are $F_p^*E_a\to E_{a+1}$,
and the last arrow is $\tau^*F_p^*E_{\nu-1}\to E_0$.
The intermediate punctured bundles can be $F_p^{a*}E$.
For geometric input representations take their reflexive
extensions and canonical full pullbacks $F_a$ on $Y$.
The forward arrows are regular and the inverse arrows
have poles at most $pC$, exactly as in the ordinary proof.
This includes the last arrow: canonical pullback commutes
with the automorphism $\tau$, and $\tau^*C=C$.

The determinant of (D.SH.3) is a crossed character with
coefficients $\mathbb F_q^\times(\chi)$. By
Lemma~\ref{D:arbitrary-field-character-scalar-H1} there
are at most $q^{c_Rt}$ such classes. This enumeration is
made on $H$, before passing to the field cover. Fix one
class; its restriction to $U_V$ fixes the ordinary
punctured determinant of $E$, and its $p^a$-powers fix
the determinants at the other $\nu$ stages.

Put $x_a=c_1(F_a)\cdot C$. These numbers are nonnegative.
The determinant divisors of the forward arrows lie
between zero and $rpC$. At the last stage
$c_1(\tau^*F_{\nu-1})\cdot C=x_{\nu-1}$.
The same cyclic recurrence as in the ordinary theorem
therefore gives
\[
     0\le x_a\le\frac{rp}{p-1}(-C^2)=O_R(td r).
                                                               \tag{D.SH.5}
\]
Every coefficient of $C$ is a positive integer. Hence
the unweighted exceptional determinant multidegrees have
nonnegative entries with sum $O_R(td r)$.
There are $O_R(td)$ components, and the negative definite
intersection matrix makes the multidegrees determine
the extension of a fixed punctured line uniquely.
The binomial enumeration in the ordinary theorem bounds
the ordered lists of global lines
$L_a=\det(F_a(-Z))$ by
\[
 q^{c_Rt}\exp\bigl(c_1\nu td\log(r+1)\bigr)
       =q^{c_Rt}\exp\bigl(c_1ft\log(r+1)\bigr).         \tag{D.SH.6}
\]
No list of all ordinary characters on the degree-$td$
cover is substituted for the twisted scalar count.

\emph{The actual finite polynomial parameter scheme.}
Choose $N\ge\max\{3,\alpha\}$, $L=N+p$, and $D=LC$,
as in the ordinary proof. All these Cartier subschemes
are deck-invariant. For a fixed ordered line tuple use
the generated quotient charts
\[
 0\longrightarrow L_a^{-1}\longrightarrow
 \mathcal O_D^{r+1}\longrightarrow V_a\longrightarrow0,
             \qquad E_a=V_a(Z).                        \tag{D.SH.7}
\]
They are actual locally split sequences over every
coefficient algebra. The scalar cohomology and line-degree
budgets on $D$ are $O_R(td)$ and $O_R(td r)$.
Their object coordinates are $O_R(td r^2)$ per stage.

For the first $\nu-1$ arrows impose the two inverse
relations of the ordinary theorem. For the last arrow
put $P=\tau^*F_p^*E_{\nu-1}$ and introduce maps
\[
 a:P\longrightarrow E_0,\qquad
 b:E_0\longrightarrow P(pC),\qquad
 a(pC)b=\iota_{E_0},\quad ba=\iota_P.                 \tag{D.SH.8}
\]
Both equalities are equalities of maps to the indicated
Cartier twists on $D$; the Cartier equation is not inverted
on a nonreduced finite cycle.

Here is the precise effect of the deck pullback on the
ordinary coordinate proof. It replaces the fixed line
$L_{\nu-1}$ by $\tau^*L_{\nu-1}$ in the last source.
This line is generated and has the same degrees as the
original line. It fixes the geometric lines
$\mathcal O_D(-C)$ and $\mathcal O_D(-Z)$.
The pulled-back quotient tuple has the same parameter
coefficients; its only nonlinear operation is the one
$p$-th power from Frobenius. All reference sections and
their deck pullbacks are fixed background data.
One may compute all required Hom complexes on any fixed
two-affine cover of $D$, treating the pulled-back fixed
lines as additional reference lines. No deck-equivariant
choice of cohomology bases is required.

The four-term quotient-Hom calculation of
Proposition~\ref{D:fixed-determinant-finite-Hom}, with the
bounded powers of generated lines used in the ordinary
proof, still has cohomology dimension $O_R(td r^2)$.
Its exact all-ring cycle, composition, and boundary-witness
equations have the same bounded polynomial degrees.
There are a fixed number of complexes per stage, including
(D.SH.8). Thus the whole affine zero scheme has
\[
 \#\text{coordinates},\ \#\text{scalar equations}
       \le c_2\nu td r^2=c_2ft r^2,
 \qquad \deg(\text{each equation})\le D_R,             \tag{D.SH.9}
\]
with constants independent of $d,t,r,f$.
Its generating-tuple conditions form one open subset.
The fixed reference cochains may depend on the cover and
the line tuple, but their coefficients are constants and
do not change variable degrees or coordinate counts.

\emph{Effectivity of the skew cycle over coefficient bases.}
The local matrix construction of
Proposition~\ref{D:finite-Frobenius-object-effectivity}
extends verbatim at the level of its specified formulas.
For precision, take a finite affine trivializing cover
closed under $\tau$; the number of this qualitative local
cover is not a parameter count. At the last step the
correction operator on transition matrices is
$g\mapsto A_i\,\tau^*F_p^*(g)\,A_j^{-1}$, with the
appropriate permuted affine indices. The inverse still
has pole at most $pC$. Since $\tau$ preserves the ideal
of $C$, an error of order $j$ becomes an error of order
at least $pj-p$. Thus $pN-p>N$ gives the same convergent
correction, unique with its prescribed $NC$ marking.
Products and inverse transitions satisfy the same
correction and marking, so uniqueness proves the actual
cocycle relations. This works over every coefficient
algebra, with no averaging and no invertibility of $d$.
The resulting formal bundles have the actual skew unit
arrows. On finite-type coefficient bases they algebraize
over the proper Noetherian completed tensor base by the
same Grothendieck-existence argument. Arbitrary-ring
finite-cycle and formal assertions remain separate from
that Noetherian algebraization assertion.

The retained $NC$ marking includes $Z$. It therefore
retains the generating quotients and the determinant
lines in (D.SH.7). Lemma
\ref{D:weak-cycle-generated-determinant-lifting} lifts
these generators and identifies the formal determinant
with the fixed global line. This argument does not depend
on the Frobenius cycle length or on its final deck pullback.
It makes no relative canonicality or coflatness assertion.

\emph{The actual punctured class comparison.}
The proof of
Proposition~\ref{D:generated-fixed-determinant-unit-openness}
applies to these generated, fixed-determinant families.
Indeed target Hartogs and adjugates bound both directions
of an Artinian punctured isomorphism by a fixed multiple
of $C$, independently of its nilpotence order.
The multiple may depend on the reference class and rank;
it is used only for qualitative finite presentation.
The intrinsic finite unit-Hom equation has the same
order improvement at the final arrow because $\tau^*$
preserves $C$. The shifted Hom and product lattices,
paired inverse equations, and uniqueness of their formal
corrections are consequently unchanged.

For the infinitesimal lifting assertion, iterate the skew
cycle $d$ times only to identify its finite etale frame
sheaf with the ordinary $F_q$-unit frame sheaf. The
semilinear deck descent is the one in the proved category
equivalence (D.SH.4). A prescribed unit isomorphism over
a nilpotent quotient lifts uniquely by finite etale
invariance; its deck compatibility follows by the same
uniqueness. The preceding generated-family argument puts
this lift inside the fixed meromorphic bound. Hence the
represented comparison Isom functor is locally of finite
presentation and formally etale, so its class images are
open and closed, exactly as in the ordinary theorem.
The $d$-fold composite is not expanded in the polynomial
equations (D.SH.9).

Every original representation is covered. Its canonical
input objects give (D.SH.7)--(D.SH.8) for one determinant
tuple, and its genuine formal transition matrices are
one admissible choice in the correction construction.
Choice-independence identifies the constructed object
with that original punctured skew unit class.

Finally the rectangular affine component bound, applied
to (D.SH.9), gives at most
$(M+1)\max\{2,D_R\}^M\le p^{c_3ft r^2}$ classes for
each line tuple, where $M\le c_2ft r^2$.
Multiplying by (D.SH.6), and using
$\log(r+1)=O(r^2)$ and $\log q=f\log p$, proves
$q^{A_R^{\rm sl}tr^2}$.
The category equivalence identifies these classes with
exactly the conjugacies in (D.SH.3). This also treats
$d=1$ and allows $p\mid d$.
\end{proof}

\begin{remark}[Scope of the field extension]
The theorem removes the resolution-extension and radial-tame
conditions on the field character in the earlier semilinear
counts. It uses the newly proved twisted scalar bound rather
than an all-characteristics local $H^1$ theorem. Its target
is the ordinary general semilinear group with a prescribed
field projection. No projective lifting or twisted-gerbe
representation count is inferred.
\end{remark}


% Proved after four complete actual-file audits. Closure of the arbitrary-cover hereditary reduction.
\subsection{Faithful semilinear crowns over arbitrary outer quotients}

\begin{corollary}[Proved: arbitrary-quotient faithful semilinear crowns]
\label{D:arbitrary-quotient-faithful-semilinear-crowns}
Retain the complete normal surface group $\Gamma$ of
Theorem~\ref{D:arbitrary-cover-semilinear-hereditary-growth}
and its constant $A_R^{\rm sl}$. Fix $b_0\ge1$ and
$\kappa>0$. Let $S$ be a nonabelian finite simple group,
and suppose
\[
 S\le\mathcal B_0\mathrel{\trianglelefteq}\mathcal B
       \le\operatorname{Aut}(S),\qquad
 [\mathcal B:\mathcal B_0]\le b_0,
 \qquad
 \mathcal B_0\hookrightarrow
      \Gamma\operatorname{L}_r(\mathbb F_q),\qquad q=p^f,
                                                            \tag{D.AS.1}
\]
where the embedding is faithful and
$\Gamma\operatorname{L}_r(\mathbb F_q)
=\operatorname{GL}_r(\mathbb F_q)\rtimes
\operatorname{Gal}(\mathbb F_q/\mathbb F_p)$.
Assume
\[
                       r^2\log q\le\kappa\log|S|.        \tag{D.AS.2}
\]
The image of $S$ lies in the linear subgroup, as follows
also from its perfectness. No ambient representation of
$\mathcal B$ is assumed.

Let $\Omega$ have $t$ elements, and let
$P\le(\mathcal B/S)^\Omega\rtimes\operatorname{Sym}(\Omega)$
have transitive permutation action. Let $L$ be its full
inverse image in
$\mathcal B^\Omega\rtimes\operatorname{Sym}(\Omega)$,
and assume $L$ is monolithic primitive with socle
$N=S^\Omega$. For every continuous quotient
$\Gamma\twoheadrightarrow L_k$, where $L_k$ is the
$k$th crown-based power over $P$, one has
\[
 \begin{split}
 k&\le |N|^{c_R^{\rm sl}},\\
 d(L_k)&\le\max\{2,d(P),c_R^{\rm sl}+2\},\\
 c_R^{\rm sl}
   &=\max\{1,\lceil
           \kappa(A_R^{\rm sl}b_0+1)+b_0-1\rceil\}.
 \end{split}                                               \tag{D.AS.3}
\]
The quotient covers and their field characters are arbitrary.
In particular, the conclusion imposes no radial-tameness,
resolution-extension, or bounded-field-order condition.
The term $d(P)$ remains part of the generator bound.
\end{corollary}

\begin{proof}
The exact nonsplit transport argument of
Theorem~\ref{D:arbitrary-quotient-faithful-linear-crowns}
is unchanged. Fix $\beta:\Gamma\twoheadrightarrow P$,
let $H$ be the inverse image of a point stabilizer, and
write $\eta:H\to\mathcal B/S$ for its outer coordinate.
Then $[\Gamma:H]=t$, and lifts $\Gamma\to L$ over
$\beta$, modulo $N$-conjugacy, correspond to lifts
$H\to\mathcal B$ of $\eta$, modulo $S$-conjugacy.
The coordinates of a surjection to $L_k$ have distinct
classes in that set.

Set $H_0=\eta^{-1}(\mathcal B_0/S)$ and
$a=[H:H_0]\le b_0$. This subgroup is fixed before a
lift is chosen, and its punctured cover has degree $at$.
Compose the embedding in (D.AS.1) with the field projection.
It is trivial on $S$, so factors through
$\mathcal B_0/S$. Consequently every restriction
$H_0\to\mathcal B_0$ under consideration has one and
the same prescribed field character
$\chi_0:H_0\to\operatorname{Gal}(\mathbb F_q/\mathbb F_p)$.
Its order is not bounded by $b_0$ and need not be prime
to $p$; neither restriction is needed.

The arbitrary-cover semilinear count now bounds the
ambient classes of these restrictions, modulo
$\operatorname{GL}_r(\mathbb F_q)$, by
$q^{A_R^{\rm sl}at r^2}$. One ambient orbit contains
at most $|\operatorname{GL}_r(\mathbb F_q)|\le q^{r^2}$
$S$-conjugacy classes of actual $\mathcal B_0$-valued
restrictions: list every ambient conjugate before taking
the smaller conjugacy quotient. This does not require
the conjugating matrix to normalize $\mathcal B_0$.
The field character remains fixed under linear conjugacy.
Thus there are at most
\[
                     q^{(A_R^{\rm sl}at+1)r^2}           \tag{D.AS.4}
\]
restriction classes modulo $S$.

Normalize each restriction to an actual representative
using an element of $S$ itself. This preserves both
$\mathcal B$-valuedness and the prescribed outer map.
Each of the $a-1$ nonidentity coset representatives for
$H/H_0$ has its image in one fixed coset of $S$ in
$\mathcal B$. There are at most $|S|$ choices for that
image, independently of the order of $\mathcal B$.
The extension cost is therefore $|S|^{a-1}$, precisely
as in the linear theorem. It follows that
\[
 \begin{split}
 k&\le q^{(A_R^{\rm sl}at+1)r^2}|S|^{a-1}\\
  &\le |S|^{\kappa(A_R^{\rm sl}at+1)+a-1}
   \le |S^\Omega|^{c_R^{\rm sl}}.
 \end{split}                                               \tag{D.AS.5}
\]
Only $a\le b_0$ and $t\ge1$ enter the last step.

The generator conclusion uses exactly the already checked
nonabelian crown estimate, not an extra representation
argument. For $s\ge d(L)=\max\{2,d(P)\}$,
\[
 \psi_L(s)\ge\frac{|N|^{s-1}}
                        {2t|\operatorname{Out}(S)|}
             \ge |N|^{s-2},\qquad
 k\le\psi_L(s)\ \Longrightarrow\ d(L_k)\le s.
                                                            \tag{D.AS.6}
\]
The precise known input is
\citet[Proposition 2.6(ii), equations (2.1)--(2.2),
pp.~11--13]{LucchiniMarionTracey2020D}; the second
inequality is the same $|S|\ge60$ and
$|\operatorname{Out}(S)|\le\log_2|S|$ calculation
in the preceding linear theorem. Taking
$s=\max\{2,d(P),c_R^{\rm sl}+2\}$ proves (D.AS.3).
\end{proof}

\begin{remark}[Concrete scope and remaining restrictions]
For a simple group $S\hookrightarrow\operatorname{GL}_r(\mathbb F_q)$
stable under entrywise field automorphisms, whenever
$S\rtimes\operatorname{Gal}(\mathbb F_q/\mathbb F_p)$
acts faithfully on $S$, it supplies $\mathcal B_0$ in
(D.AS.1). A bounded residual graph extension is permitted
by the displayed index condition. The previously checked
natural action for
$S=\operatorname{SL}_{p^u}(\mathbb F_{p^f})$, with the
nonsimple small parameters omitted, is a concrete
unbounded-rank and unbounded-field-order instance:
its scalar centre is trivial, and the existing natural
order comparison allows $\kappa=4$.
This uses the same checked field/graph recognition as
the earlier one-curve crown results.

Thus the field action and permutation quotient may both
ramify arbitrarily over the resolution of any fixed
complete normal surface. The faithful semilinear
embedding and quadratic simple-order comparison remain
explicit hypotheses. This does not replace them by a
projective representation, include unrestricted
alternating socles, or bound $d(P)$ without an additional
assumption. It is a closure of the one arbitrary-cover
hereditary reduction, not a second geometric reduction.
\end{remark}


% Separate post-CP26 proved recognition. Recognition within the proved matrix/wreath criterion, not an independent movement.
\subsection{Non-natural alternating primitive actions in the matrix catalogue}

\begin{corollary}[Proved: alternating socles outside the natural primitive action]
\label{D:non-natural-alternating-primitive-models}
Fix the characteristic $p$ and a finite collection $\mathcal F$
of nonabelian finite simple groups. For each $N\ge2$, there
is a finite collection of faithful permutation models
\[
 P_{N,j}\le S_N,\qquad
 P_{N,j}\hookrightarrow\operatorname{GL}_{r_{N,j}}(\mathbb F_p)
                 \quad(j\in\mathcal J_N)                 \tag{D.NA.1}
\]
and one constant $B=B_{p,\mathcal F}$ such that
\[
 \log\max\{1,|\mathcal J_N|\}\le BN\log(N+1),\qquad
 r_{N,j}^2\log p\le BN\log(N+1).                         \tag{D.NA.2}
\]
These models contain, up to permutation conjugacy, every
faithful primitive $G\le S_N$ with nonabelian socle whose
simple factors are each either
\begin{enumerate}
\item an alternating simple group $A_m$, $m\ge5$;
\item a defining-characteristic-$p$ simple Lie group in
      the exact families and finite-exception conventions
      of Corollaries~\ref{D:classical-defining-characteristic-crowns}
      and~\ref{D:exceptional-defining-characteristic-crowns}; or
\item a member of $\mathcal F$,
\end{enumerate}
provided that the permutation action is not the natural
$A_N$ or $S_N$ action with $N\ge5$.
Any fixed finite set of these excluded degrees can be
restored by adding the corresponding fixed symmetric models.

More explicitly, the models are the union of a catalogue
with faithful matrix dimension $O((1+\log N)^2)$ and
the standard alternating-product ambient groups
\[
 P_{m,k,\ell}=S_m\wr S_\ell,
 \qquad N=\binom{m}{k}^{\ell},\qquad
 m\ge5,\quad 1\le k\le m/2,\quad (k,\ell)\ne(1,1).
                                                            \tag{D.NA.3}
\]
The action in (D.NA.3) is the product action on ordered
$\ell$-tuples of $k$-element subsets. This ambient group
has the faithful block-permutation matrix model of rank
$m\ell$ over $\mathbb F_p$, with the sharper bound
\[
                         (m\ell)^2\le4N.                \tag{D.NA.4}
\]
There are at most $N^3$ parameter triples in (D.NA.3)
for a fixed $N$. Every subgroup in one of these fixed
actions is covered by the same model.

Consequently these primitive actions can be used as
zero-cost constituents at arbitrary height in
Theorem~\ref{D:alternating-local-wreath-energy}.
The same factorial-order cover count and selected
alternating simple-product conclusion follow under its
unchanged residual tower-energy and geometric conditions.
Within the stated nonabelian-socle families, this leaves
only the unbounded natural alternating and symmetric
primitive actions outside this recognition. It makes no
new assertion about general cross-characteristic simple
socles or affine actions in characteristic different
from $p$.
\end{corollary}

\begin{proof}
We first extend the small intrinsic representation argument
to alternating factors. The Known automorphism description is
\[
                 \operatorname{Aut}(A_m)=S_m
                       \qquad(m\ge5,\ m\ne6),           \tag{D.NA.5}
\]
where the isomorphism is induced by conjugation. We use
\citet[Theorems 11.2 and 11.7, with Lemmas 11.3--11.5,
pp.~42--44]{C_Lyons1999AlternatingAutomorphisms}.
The source's printed overlap table is not needed as a
literal formula: the verification note accompanying this
proof gives the centralizer calculation and the corrected
orientation-independent subgroup-order test.
For $m\ne6$, the permutation matrices of $S_m$ give
a faithful $\mathbb F_p$ representation of dimension $m$.
For $A_6$, use the faithful action of its automorphism
group on the $360$ elements of $A_6$; no particular
description of its extra outer automorphisms is needed.
It follows that one absolute constant $c_0$ satisfies
\[
 \operatorname{Aut}(A_m)\hookrightarrow
       \operatorname{GL}_{d_m}(\mathbb F_p),\qquad
                         d_m\le c_0\log|A_m|             \tag{D.NA.6}
\]
for every $m\ge5$. For example $c_0=64$ works:
$m\le2\log(m!/2)$ for $m\ge5$, and
$360<64\log360$.

Combine (D.NA.6) with the already proved automorphism
representations (D.PM.3)--(D.PM.5) for the other permitted
simple groups. Thus every permitted $S$ has
\[
 \operatorname{Aut}(S)\hookrightarrow
       \operatorname{GL}_{d_S}(\mathbb F_p),\qquad
                         d_S\le c\log|S|,               \tag{D.NA.7}
\]
with one constant depending only on $p,\mathcal F$ and
the fixed exception conventions. The nonabelian socle
$L=\prod_i S_i^{a_i}$ has trivial centralizer in $G$:
a nontrivial normal centralizer would contain a minimal
normal subgroup of $G$, hence one lying in $L$, contrary
to $Z(L)=1$. Therefore the actual conjugation embedding
\[
 G\hookrightarrow\operatorname{Aut}(L)
   =\prod_i\bigl(\operatorname{Aut}(S_i)\wr S_{a_i}\bigr)
                                                            \tag{D.NA.8}
\]
and block permutations of the representations (D.NA.7)
give a faithful representation of dimension at most
$c\sum_i a_i\log|S_i|\le c\log|G|$.
This is the same full-socle argument as in
Corollary~\ref{D:primitive-defining-characteristic-orbit-models};
it works in every characteristic, without semisimplicity.

Now apply the Known primitive-order theorem
\citep[Theorem 1.1 and complete written Section 2,
pp.~2 and 4--8 of the author version]{C_Maroti2002PrimitiveOrders}.
In its small-order case,
$|G|<N^{1+\lfloor\log_2N\rfloor}$, so (D.NA.8)
has dimension $O((1+\log N)^2)$.
The four fixed Mathieu actions are absorbed by enlarging
the constant and using their fixed permutation matrices.
No estimate about unbounded sporadic or other unspecified
families is needed.

In the remaining case, the theorem places $G$ in
$S_m\wr S_\ell$ with its product action of degree
$N=\binom{m}{k}^{\ell}$, containing $A_m^\ell$.
Replace $k$ by $m-k$ if necessary; taking complements
is an equivariant identification of the subset actions.
Thus $1\le k\le m/2$. If $m\le4$, the faithful block
model has dimension $m\ell\le4\log_2N$, since
$\binom{m}{k}\ge2$. These cases can be included in
the preceding intrinsic catalogue.
Assume henceforth that $m\ge5$.
If $(k,\ell)=(1,1)$, then $N=m$ and
$A_N\le G\le S_N$. This is exactly one of the two
natural actions excluded in the statement.
Every other case has the model (D.NA.3).

The permutation action of $S_m$ on proper nonempty
$k$-subsets is faithful: the intersection of all such
subsets containing a point is that point, so a permutation
fixing every subset fixes every point. Its product wreath
action is faithful as well, by varying one coordinate
of a tuple at a time. As an abstract group
$S_m\wr S_\ell$ also acts faithfully on $\ell$ labelled
blocks of $m$ basis vectors. Its permutation matrices
give the required rank-$m\ell$ injection in every
characteristic. We are using two different faithful
actions of the same fixed group; no assertion identifying
their conjugacy relations is involved.

Here are the uniform numerical bounds. Put
$Q=\binom{m}{k}$. If $k\ge2$, then
\[
 Q\ge\frac{m(m-1)}2\ge\frac25m^2,\qquad Q\ge10,
 \qquad \ell^2\le Q^{\ell-1}.                           \tag{D.NA.9}
\]
The last inequality follows by induction from $\ell=1$,
since $((\ell+1)/\ell)^2\le4\le Q$. Consequently
$(m\ell)^2\le(5/2)Q^\ell\le4N$.
If $k=1$, our exclusion gives $\ell\ge2$.
Induction from $\ell=2$ gives
\[
 \ell^2\le4m^{\ell-2},\qquad
 \left(\frac{\ell+1}{\ell}\right)^2
                        \le\frac94\le m.               \tag{D.NA.10}
\]
This again yields $(m\ell)^2\le4m^\ell=4N$.
Both inductions hold for all the stated integers;
no finite numerical test is used in place of them.
For fixed $N$, every eligible $m,k,\ell$ lies between
one and $N$, so the number of triples is at most $N^3$.
The logarithm of that model count is at most $3\log N$,
and each model's matrix cost is at most $4N\log p$.

It remains to make the small-order collection an actual
fixed catalogue. Choose the integer ceiling of
$c'(1+\log N)^2$ as its matrix dimension and apply
Lemma~\ref{D:intrinsic-small-linear-permutation-images}:
enumerate the subgroups of that fixed general linear
group and their core-free point-stabilizer subgroups
of index $N$. This supplies the actual degree-$N$
permutation models, not merely abstract isomorphism
types. Their logarithmic number is
$O((1+\log(N+1))^8)$ and their individual matrix
cost is $O((1+\log(N+1))^4)$.
These bounds, the at most $N^3$ product models, and
any fixed small-degree additions satisfy (D.NA.2).
For clarity, every fixed power of a logarithm is bounded
by a constant times $N\log(N+1)$ on $N\ge2$:
with $x=\log(N+1)$ this follows from the boundedness
of $x^a e^{-x}$ for each fixed $a$. Thus the catalogue
constant is independent of the degree $N$.

Insert this single union of models into
Proposition~\ref{D:hereditary-wreath-permutation-count}.
All the identified primitive constituents are modelled
levels, so contribute zero to its unrestricted-level
energy. Theorem~\ref{D:alternating-local-wreath-energy}
then supplies exactly the asserted arbitrary-height
and geometric consequences. Its count uses actual Hom
sets followed by permissible permutation conjugacy,
so it does not assume that a modular permutation module
determines its underlying permutation set.
The original-divisorial hereditary scope is not enlarged.
\end{proof}

\begin{remark}[Scope of the recognition]
The primitive classification and the automorphism theorem
are Known inputs. The new use here is the explicit
matrix-model cost and its insertion into the already
proved tower criterion. The general natural alternating
local action is not made inexpensive by this argument:
its standard permutation matrix model still has rank $N$.
Those orbits retain the previous allowed quotient/model
conditions or their unrestricted tower-energy charge.
This is part of the same restricted alternating reduction,
not an additional independent movement or a proof of
unrestricted alternating multiplicity bounds.
\end{remark}


% Post-CP26 proved auxiliary refinement. Four full actual-file/source audits; no independent main-gate movement.
\subsection{Controlled normal-crossings refinement of the inherited cycles}

\begin{lemma}[Proved: a linear bound for embedded curve refinement]
\label{D:linear-SNC-refinement-centres}
Let $Y$ be a regular excellent surface containing a reduced proper
curve $B$ over an algebraically closed field $k$. Assume the closed
points of $B$ have residue field $k$. Put
\[
 G=h^1(B,\mathcal O_B),\qquad
 v=\#\operatorname{Irr}(B),\qquad c=h^0(B,\mathcal O_B),
 \qquad
 \delta=\dim_k H^0(B,\nu_*\mathcal O_{\widetilde B}/\mathcal O_B),
                                                        \tag{D.SR.1}
\]
where $\nu$ is normalization. There is a sequence of point
blowups $\rho:Y'\to Y$, with every centre a non-SNC point of
the current reduced total transform, such that
$B'=(\rho^{-1}B)_{\rm red}$ is an SNC divisor and
\[
 n_{\rm blow}\le2\delta+7G
             \le9G+2v-2c.                              \tag{D.SR.2}
\]
The final divisor has $v+n_{\rm blow}$ components,
the same normalization genera as $B$ together with rational
exceptional components, and at most
$G+v+n_{\rm blow}-c$ nodes. The assertions hold in every
characteristic.
\end{lemma}

\begin{proof}
We use only point blowups of regular two-dimensional local rings.
Their exceptional curve is $\mathbf P^1_k$, with normal bundle
$\mathcal O(-1)$, and the surface remains regular. The exact
cohomology calculation is
\[
 \pi_*\mathcal O(aF)=\mathcal O,
 \qquad
 \operatorname{length}R^1\pi_*\mathcal O(aF)
       =\frac{a(a-1)}2\quad(a\ge0),                    \tag{D.SR.3}
\]
with no higher direct images. These are the full written
quadratic-transformation results in
\citet[Lemmas 54.3.1, 54.3.2 and 54.3.4,
Tags 0AGQ, 0AGR and 0AGS]{StacksSNCRefinementRoot}:
$\mathcal O(F)=\mathcal O(-1)$ in their convention.
In particular $R\pi_*\mathcal O=\mathcal O$.

Suppose a reduced current boundary $D$ has multiplicity $m$
at a centre. Its reduced total transform is
$D'=\pi^*D-(m-1)F$. The ideal sequence and (D.SR.3),
after tensoring with the locally invertible $\mathcal O(-D)$,
give
\[
 0\longrightarrow\mathcal O_D
  \longrightarrow\pi_*\mathcal O_{D'}
  \longrightarrow Q\longrightarrow0,
 \quad \dim_k Q=\frac{(m-1)(m-2)}2,
 \quad R^1\pi_*\mathcal O_{D'}=0.                      \tag{D.SR.4}
\]
This is a calculation with the reduced total transform, not
with the strict transform alone. The connected components are
unchanged: the new exceptional curve joins exactly the strict
branches through the old centre. Hence $h^0$ is unchanged,
and taking Euler characteristics yields
\[
 h^1(D',\mathcal O_{D'})
   =h^1(D,\mathcal O_D)-\frac{(m-1)(m-2)}2.            \tag{D.SR.5}
\]

Call a centre high if $m\ge3$, and low if $m=2$.
A non-SNC point of a reduced curve has $m\ge2$.
Set $a_m=(m-1)(m-2)/2$. Along every finite sequence
of the stipulated blowups, (D.SR.5) gives
\[
 \sum_{\rm high}a_m\le G,
 \qquad \#\{\text{high centres}\}\le G,
 \qquad \sum_{\rm high}m\le3G,                        \tag{D.SR.6}
\]
because $a_m\ge1$ and $m\le3a_m$ for $m\ge3$.

Here is the needed local control of low centres. On the
blowup at a multiplicity-two point, the strict transform
meets $F$ in a scheme of total length two, given by the
degree-two tangent cone. If there are two distinct tangent
directions, each resulting point has a single smooth strict
branch transverse to $F$ and is SNC. Otherwise there is
one possible unresolved child. If its strict transform has
multiplicity at least two, the reduced total boundary there
has multiplicity at least three, because it also contains $F$.
The only unresolved child still of multiplicity two thus
consists of $F$ and one smooth strict branch with intersection
multiplicity two. They are tangent. After blowing up this
tangency, their two strict transforms and the new exceptional
curve pass through the same point. Its total multiplicity is
at least three. There are no other unresolved children.
Thus a path of consecutive low centres has length at most two.
This uses only the length-two tangent cone over $k$, not a
derivative, separability assumption, or division by two.

Organize the centres as a rooted forest. Its roots lie over
the initial non-SNC points, of which there are at most
$\delta$: each such point is singular on the reduced curve
and contributes at least one to its normalization defect.
After a multiplicity-$m$ blowup, all unresolved children
lie at intersections of the strict transform with $F$.
There are at most $m$ such points, since the intersection
scheme has length $m$. Away from these points the new
exceptional curve is smooth, and the old boundary is unchanged.

Delete the high vertices from any finite part of this forest.
Each remaining component is a path of at most two low vertices.
There are at most $\delta+\sum_{\rm high}m$ such paths:
each begins at an initial root or immediately below a high
vertex. Therefore
\[
 n_{\rm low}\le2\delta+2\sum_{\rm high}m
                   \le2\delta+6G.                     \tag{D.SR.7}
\]
Together with (D.SR.6) this proves the first bound in
(D.SR.2) for every finite prefix. Repeatedly blowing up
a remaining non-SNC point must terminate, since an infinite
sequence would violate that uniform bound. Its endpoint
is the asserted SNC divisor.

The normalization sequence for $B$ gives
\[
 \delta=G-\sum_i g(\widetilde B_i)+v-c\le G+v-c.       \tag{D.SR.8}
\]
This proves the second bound. Every blowup adds exactly one
rational component and leaves the normalization of each old
component unchanged. At the SNC endpoint the normalization
defect is exactly the number of nodes. Apply (D.SR.8) there,
using $h^1(B',\mathcal O_{B'})\le G$ and the unchanged $c$,
to obtain the stated node bound.
\end{proof}

\begin{corollary}[Proved: SNC refinement with inherited degree-linear cycle budgets]
\label{D:SNC-refinement-inherited-cycles}
Use the degree-$t$ cover model $g:Y\to Y_0$ and actual
Cartier cycles $C,Z$ of
Proposition~\ref{D:uniform-comparable-reduction-cycle}.
There is a projective point-blowup refinement
$\rho:Y'\to Y$ with SNC exceptional support and $O_R(t)$
centres, components, nodes and total normalization genus.
Put $C'=\rho^*C$ and $Z'=\rho^*Z$. Then
\[
 C'\le Z'\le\alpha C',\qquad
 (Z')^2=Z^2,\qquad K_{Y'}\cdot Z'=K_Y\cdot Z,          \tag{D.SR.9}
\]
and the corresponding equalities hold for $C$.
Generation of $\mathcal O_Y(-C)$ and $\mathcal O_Y(-Z)$
pulls back. For every positive integer $b$ and every
line bundle $L$ on $Y$,
\[
 \begin{aligned}
 \rho_*\mathcal O_{bZ'}&=\mathcal O_{bZ},
       &R^1\rho_*\mathcal O_{bZ'}&=0,\\
 H^i(bZ',\rho^*L|_{bZ'})&=H^i(bZ,L|_{bZ})
       &&(i=0,1).
 \end{aligned}                                           \tag{D.SR.10}
\]
The same statements hold with $C$ in place of $Z$.
In particular all the specified fixed-level scalar budgets
and the weak-cycle vanishings
$H^1(Z',\mathcal O_{Z'}(-jZ'))=0$ for $j\ge1$ persist.
Generated actual bundles and their fixed determinant data
pull back, and an actual Frobenius arrow with inverse pole
at most $pC$ pulls back with inverse pole at most $pC'$.
No canonical-fullness assertion for these pulled-back
families is required or made.

If the original Jacobian divisor satisfies $R_g\le\lambda C$,
the specified refinement also satisfies
\[
 K_{Y'/Y}\le\rho^*E_{\rm red}-E'_{\rm red},\qquad
 R_{g\rho}\le(\lambda+1)C'-E'_{\rm red},              \tag{D.SR.11}
\]
where $E_{\rm red}$ and $E'_{\rm red}$ are the respective
reduced exceptional supports. This does not redefine
$Z'$ as a linear combination involving the new Jacobian.
\end{corollary}

\begin{proof}
The reduced exceptional curve $E=E_{\rm red}$ has
$O_R(t)$ components by the already proved cover budget.
Also $h^1(E,\mathcal O_E)\le h^1(Y,\mathcal O_Y)=O_R(t)$:
the ideal sequence surjects onto $H^1(E,\mathcal O_E)$,
because $H^2(Y,\mathcal O_Y(-E))=0$ for a proper surface
over this affine complete local base with fibres of
dimension at most one. This is the same coherent-dimension
input used in the earlier exceptional-curve scalar bound.
Apply Lemma~\ref{D:linear-SNC-refinement-centres}.
The sum of normalization genera is at most $h^1(E,\mathcal O_E)$
by its normalization sequence. The required combinatorial
bounds therefore follow with constants independent of $t$.

Pullback of effective Cartier divisors preserves their
inequalities. The projection formula gives equality of
self-intersections. At a point blowup the canonical
divisor changes by the new exceptional curve, whose
intersection with every pulled-back divisor is zero.
Iterating proves the canonical-intersection equality.
Pullback preserves a specified generating surjection,
its determinant isomorphism and the stated inverse-arrow
identities. Frobenius commutes with pullback.

For one point blowup (D.SR.3) at $a=0$ gives
$R\rho_*\mathcal O_{Y'}=\mathcal O_Y$; composition
preserves this statement. Tensor the ideal sequence
of the actual Cartier divisor $bZ$ with this identity.
The projection formula, or the same sequence term by
term, gives $R\rho_*\mathcal O_{bZ'}=\mathcal O_{bZ}$.
Tensoring by $L$ and taking cohomology proves (D.SR.10).
This retains the full nonreduced cycles. Substituting
$b=1$ and $L=\mathcal O_Y(-jZ)$ proves the weak-cycle
vanishing. The fixed scalar equalities also hold after
coefficient extension over $k$, by flat base change;
no arbitrary nonflat blowup comparison is asserted.

Finally, let $B_i$ be the current reduced total boundary
and let a stipulated centre have multiplicity $m\ge2$.
For its new exceptional curve $F$,
\[
 \pi^*B_i-B_{i+1}=(m-1)F\ge F=K_{Y_{i+1}/Y_i}.
                                                        \tag{D.SR.12}
\]
Pull back the previous inequalities and add this one.
Induction gives the first inequality of (D.SR.11).
The Jacobian chain rule gives
$R_{g\rho}=\rho^*R_g+K_{Y'/Y}$.
Since every coefficient of $C$ on $E_{\rm red}$ is a
positive integer, $E_{\rm red}\le C$. Thus
\[
 R_{g\rho}\le\lambda C'+\rho^*E_{\rm red}-E'_{\rm red}
             \le(\lambda+1)C'-E'_{\rm red},
\]
as claimed. Repeated additional blowups at smooth
boundary points are not included in this argument.
\end{proof}

\begin{remark}[Remaining projective obligation]
These results give an SNC model with inherited scalar
and ordinary-family budgets. They do not bound sums of
products of the multiplicities of adjacent components,
construct a uniformly controlled twisted anchor, or
give degree-linear root-invariant Hom presentations.
Those are separate requirements of the hereditary
projective count. This refinement is an auxiliary
geometric clarification, not an independent main-gate
movement or a projective representation estimate.
\end{remark}


% Separate post-CP26 proved auxiliary; four complete audits and exact layout refresh. No hereditary projective count.
\subsection{Full-bundle degrees and canonical orders with a fixed Morita conductor}

\begin{lemma}[Proved: a degree bound for every full bundle on a weak cycle]
\label{D:full-bundle-cycle-degree}
Let $A$ be a complete excellent normal local surface with
algebraically closed coefficient and residue field $k$.
Let $\pi:Y\to\operatorname{Spec}A$ be a projective regular
resolution, an isomorphism on the puncture. Let $Z>0$ be a
Cartier exceptional cycle of full support such that
$\mathcal O_Z(-Z)$ is generically generated and
$H^1(Z,\mathcal O_Z(-Z))=0$.
For a finite reflexive $A$-module $M$ of rank $s$, set
$F=(\pi^*M)^{\vee\vee}$ and $V=F|_Z$. Then
\[
 0\le c_1(F)\cdot Z
 \le s\bigl(2h^1(\mathcal O_Z)-h^0(\mathcal O_Z)-Z^2\bigr).
                                                        \tag{D.OC.1}
\]
In the degree-$t$ cover family with its comparable strong
cycle, the right side is $O_R(st)$, without assuming any
Frobenius structure on $M$.
\end{lemma}

\begin{proof}
Use the already verified weak-cycle part of the
Kahn--Gavran correspondence, as isolated in
Proposition~\ref{D:weak-cycle-Artinian-full-realization} and
Theorem~\ref{D:reduction-cycle-full-module-correspondence}.
Here are the exact consequences used. The restriction $V$
is generically generated, and
\[
 H^1(Y,F(-Z))=0,\qquad H^1(Y,F)=H^1(Z,V).
                                                        \tag{D.OC.2}
\]
Indeed generic generators of $V$, together with those of
$\mathcal O_Z(-Z)$, give the conormal $H^1$ vanishings;
filter the higher cycles and apply formal functions.
The full bundle has $H^0(Y,F)=H^0(Y\setminus Z,F)$.
Consequently the actual $2Z$ boundary
\[
 H^0(Z,V(Z))\longrightarrow H^1(Y,F)=H^1(Z,V)
                                                        \tag{D.OC.3}
\]
is injective: in the sequence
$0\to F\to F(Z)\to V(Z)\to0$, every section of $F(Z)$
already belongs to $H^0(Y,F)$ by this Hartogs equality.
No vanishing of $H^1_Z(Y,F)$ is being asserted.

The full, possibly nonreduced curve $Z$ is Cohen--Macaulay.
Its associated points are its finitely many generic points.
Generic generation and Nakayama, followed by simultaneous
avoidance of finitely many proper determinant conditions
over the infinite field $k$, choose $s$ global sections of
$V$ forming a basis at all these points. Their determinant
is a nonzerodivisor on $Z$. Thus
\[
 0\longrightarrow\mathcal O_Z^s\longrightarrow V
   \longrightarrow T\longrightarrow0
                                                        \tag{D.OC.4}
\]
has finite-length cokernel, with
$\operatorname{length}T=\deg V=c_1(F)\cdot Z\ge0$.
It gives $h^1(V)\le s h^1(\mathcal O_Z)$.
Riemann--Roch on the Cartier curve gives
\[
 \begin{split}
 \deg V+sZ^2+s\chi(\mathcal O_Z)
   &=\chi(V(Z))\le h^0(V(Z))\\
   &\le h^1(V)\le s h^1(\mathcal O_Z).
 \end{split}                                           \tag{D.OC.5}
\]
Rearrangement proves (D.OC.1). The established scalar and
energy budgets on the comparable cycle prove the final
assertion. This argument uses the actual full pullback;
it does not apply to an arbitrary intermediate lattice
having only the same generic fibre.
\end{proof}

\begin{proposition}[Proved: canonical full orders and bounded Morita products]
\label{D:canonical-order-Morita-conductor}
Use the fixed base $R$, degree-$t$ normal punctured cover
$A/R$, fixed Noether ideal $J$, and original comparable-cycle
model $Y$ of Proposition~\ref{D:uniform-comparable-reduction-cycle}.
Write $U=\operatorname{Spec}A\setminus\{\mathfrak m_A\}$.
The separating elements $x,y\in J$ satisfy
$(x,y)\mathcal O_Y=\mathcal O_Y(-C)$.
Let $\mathcal G\to U$ be a banded finite prime-to-$p$ gerbe.
Let $\mathcal A,\mathcal B$ be vector bundles of the same
scalar weight and the same rank $h$, with $h$ invertible in $k$.
Their endomorphism and mutual Hom bundles descend to ordinary
bundles on $U$. Take their unique reflexive extensions to $A$:
\[
 \begin{aligned}
 P_{A,0}&=j_*\mathcal End(\mathcal A),&
 P_{B,0}&=j_*\mathcal End(\mathcal B),\\
 Q_0&=j_*\mathcal Hom(\mathcal A,\mathcal B),&
 R_0&=j_*\mathcal Hom(\mathcal B,\mathcal A),
 \end{aligned}                                           \tag{D.OC.6}
\]
where $j:U\hookrightarrow\operatorname{Spec}A$ and descent
from $\mathcal G$ is understood. Put $P_A,P_B,Q,R_1$ for
their canonical full pullbacks to $Y$.
Then $P_A,P_B$ are locally free unital associative algebras,
$Q$ is a left $P_B$, right $P_A$ bimodule, and $R_1$ has the
opposite actions. The actual composition maps obey
\[
 \begin{split}
 \mathcal O_Y(-C)P_B
       &\subset\operatorname{im}(Q\otimes R_1\to P_B),\\
 \mathcal O_Y(-C)P_A
       &\subset\operatorname{im}(R_1\otimes Q\to P_A).
 \end{split}                                           \tag{D.OC.7}
\]
These are inclusions of actual coherent sheaf images on $Y$.
Under the generic trace identifications one also has
\[
 P_A\subset P_A^\vee\subset P_A(C),\qquad
 R_1\subset Q^\vee\subset R_1(C),\qquad
 Q\subset R_1^\vee\subset Q(C).                       \tag{D.OC.8}
\]
Every punctured twisted bundle $\mathcal F$ of this scalar
weight gives a genuine right $P_A$-module
$\operatorname{full}(j_*\mathcal Hom(\mathcal A,\mathcal F))$.

All four full bundles in (D.OC.6) have rank $h^2$,
nonnegative degree against $Z$ bounded by $O_R(h^2t)$,
and their twists by $-Z$ are globally generated on the
original comparable-cycle model. The same algebra, module,
conductor, generation and degree statements hold for their
actual pullbacks, together with $C,Z$, along any regular
point-blowup refinement. This last statement does not
replace those pullbacks by newly canonicalized full bundles.
\end{proposition}

\begin{proof}
Reflexive extensions from the puncture are finite and unique
on this normal surface. The algebra and composition maps
extend by target reflexivity. At every height-one point of
$Y$, the full pullback is precisely the torsion-free image
of the ordinary pullback in its generic vector space.
Products of pulled-back elements remain in that image.
Thus multiplication, units, all bimodule actions and
composition extend at every such point, and hence extend
over $Y$ by reflexivity. A reflexive sheaf on the regular
surface is locally free. Associativity, the unit equations
and the bimodule equations hold on the dense puncture and
therefore everywhere. This proves an order statement, not
an Azumaya assertion on the exceptional divisor.

Trace composition identifies $R_0=Q_0^\vee$ and
$P_{A,0}=P_{A,0}^\vee$. These identifications hold on $U$
by the matrix trace pairing and extend reflexively.
By Proposition~\ref{D:cover-uniform-Noether-different},
for each $z\in J$ multiplication by $z$ on $Q_0$ factors
through a finite free $A$-module. Equivalently, there are
$q_i\in Q_0$ and $r_i\in R_0$ such that
\[
 z\operatorname{id}_{Q_0}
       =\sum_i q_i\,\operatorname{tr}(r_i\,\cdot\,).
                                                        \tag{D.OC.9}
\]
On an fppf trivialization of the two rank-$h$ bundles,
partial trace of this equality gives
\[
                  \sum_i q_i r_i=hz\,1_{P_{B,0}}.
                                                        \tag{D.OC.10}
\]
For completeness, the tensor for the identity on the
$h\times h$ matrix space is
$\sum_{a,b}E_{ab}\otimes E_{ba}$; its composition is
$\sum_{a,b}E_{ab}E_{ba}=hI_h$.
This identity is invariant under changes of both frames.
It descends and extends to $A$. Since $h$ is invertible,
$z1_{P_{B,0}}$ lies in the composition image. Multiplication
by $P_{B,0}$ gives the entire $zP_{B,0}$.
Use the opposite factorization for $R_0$ to obtain the
corresponding assertion for $P_{A,0}$.

The sections $q_i,r_i$ pull back to actual sections of
$Q,R_1$, and the equations (D.OC.9)--(D.OC.10) hold on $Y$.
The image of composition is stable under the appropriate
algebra action. Thus it contains $xP_B$ and $yP_B$, and
similarly for $P_A$. The exact Cartier pencil now proves
(D.OC.7), including at closed points; no saturation of
these images has been substituted.

Here is the dual comparison used in (D.OC.8).
For any reflexive source $L_0$, write
$z\operatorname{id}_{L_0}=vu$ through a finite free module.
Let $L=\operatorname{full}(L_0)$ and
$L'=\operatorname{full}(L_0^\vee)$.
The entries of $u$ are actual global sections of $L'$.
For any local $\lambda\in L^\vee$,
\[
 z\lambda=\lambda vu
   =\sum_i\lambda(v_i)u_i\in L'.                      \tag{D.OC.11}
\]
The natural evaluation pairing gives $L'\subset L^\vee$.
Taking $z=x,y$ proves
$L'\subset L^\vee\subset L'(C)$.
Apply this to $Q_0,R_0,P_{A,0}$, with their trace dualities,
to obtain (D.OC.8).

The right module action on the stated full Hom bundle is
proved by the same height-one argument as multiplication.
Its handedness is precomposition by $P_A$.
The degree bounds are Lemma~\ref{D:full-bundle-cycle-degree};
generation by $-Z$ is the already proved uniform canonical
full-pullback generation on this original model.
For a point-blowup refinement, pull back the actual algebras,
actions, pairings and factorizations. The displayed tensor
expressions still prove the conductor inclusions, and
pullback preserves the given generating surjections.
Intersection with a pulled-back cycle preserves the degree.
The identity $R\rho_*\mathcal O=\mathcal O$ also preserves
the global sections and their punctured Hartogs equality.
No new fullness assertion on the refined model is needed.
\end{proof}

\begin{remark}[Exact limitation of the order reduction]
For bounded $h$ the orders and comparison modules above have
bounded rank, linear degree budgets and the fixed conductor
$C$. They are Azumaya and Morita equivalent on $U$.
On the exceptional divisor they may be non-Azumaya orders.
The assertions do not lift arbitrary truncated order actions,
associativity equations and Frobenius comparisons simultaneously
to formal objects. Nor does a divisor sandwich alone imply
cohomology bounds or generation for an arbitrary intermediate
root-stack lattice. Those effectivity and comparison issues
remain before a hereditary projective representation count.
\end{remark}


% Separate post-CP26 checked source. All four full audits passed; local order actions only, no global Morita-unit effectivity or projective count.
\subsection{A uniform separability conductor and local Newton lifting for full orders}

\begin{lemma}[Proved: scaled separability of canonical full endomorphism orders]
\label{D:full-order-scaled-separability}
Use the complete degree-$t$ cover ring $A$, fixed Noether ideal $J$,
regular comparable-cycle model $Y$, and separating pencil
$(x,y)\mathcal O_Y=\mathcal O_Y(-C)$ of
Proposition~\ref{D:canonical-order-Morita-conductor}.
Let $\mathcal V$ be a rank-$h$ twisted vector bundle on the puncture,
with $h$ invertible in the coefficient field. Put
\[
 P_0=j_*\mathcal End(\mathcal V),\qquad
 P=(\pi^*P_0)^{\vee\vee}.
                                                        \tag{D.SN.1}
\]
The endomorphism bundle is understood to have descended to the
punctured ordinary scheme. Then, for each $u\in J$, there is an
actual tensor $e_u\in P_0\otimes_A P_0$ such that
\[
 (a\otimes1)e_u=e_u(1\otimes a)\quad(a\in P_0),
 \qquad \mu(e_u)=u^2\,1.
                                                        \tag{D.SN.2}
\]
Here the first identity means
$\sum a e_u'\otimes e_u''=\sum e_u'\otimes e_u''a$;
the second tensor factor has this right action.
On any open subset of $Y$ where $C=(z)$, after shrinking to one
of the two pencil charts there is an actual tensor
$e\in P\otimes P$ with
\[
             (a\otimes1)e=e(1\otimes a),\qquad
                         \mu(e)=z^2\,1.                \tag{D.SN.3}
\]
All these identities persist after arbitrary coefficient extension
over the original field and completion along $C$. Frobenius
pullback has the same property with $z^2$ replaced by $z^{2p}$.
The assertion concerns a scaled separability element; $P$ need
not be Azumaya on $C$.
\end{lemma}

\begin{proof}
The matrix trace pairing identifies $P_0$ with $P_0^\vee$ on the
puncture and hence on $\operatorname{Spec}A$ by reflexivity.
Proposition~\ref{D:cover-uniform-Noether-different} gives a
factorization of $u\operatorname{id}_{P_0}$ through a finite free
$A$-module. Thus an actual tensor $c_u\in P_0\otimes P_0$
maps to $u\operatorname{id}_{P_0}$ under
\[
 a\otimes b\longmapsto
             \bigl(v\longmapsto a\operatorname{tr}(bv)\bigr).
                                                        \tag{D.SN.4}
\]
On the puncture this tensor is $u$ times the trace Casimir. In a
matrix frame the Casimir is
$\sum_{i,j}E_{ij}\otimes E_{ji}$; it centralizes as in
(D.SN.2), and its product is $hI_h$. Consequently
\[
 \mu(c_u)=hu\,1,
 \quad (a\otimes1)c_u-c_u(1\otimes a)
                   \in H^0_{\mathfrak m}(P_0\otimes P_0).
                                                        \tag{D.SN.5}
\]
The product identity extends from the puncture because $P_0$ is
torsion free. The displayed supported submodule is annihilated
by $J$. Indeed, for $v\in J$, tensor a factorization
$v\operatorname{id}_{P_0}:P_0\to A^N\to P_0$ with $P_0$.
The intermediate module $P_0^N$ has depth two, so a section
supported at the closed point maps to zero there. The composite
is multiplication by $v$, proving the assertion about $J$.
It follows that
\[
                          e_u=h^{-1}u c_u              \tag{D.SN.6}
\]
is an actual tensor satisfying both identities (D.SN.2).
No centralization identity in a reflexive tensor hull is
substituted for an identity in the ordinary tensor product.

Pull back these tensors to $P\otimes P$ on $Y$. At every
height-one point the full pullback is the torsion-free image of
the ordinary pullback, so it is generated there by pulled-back
elements of $P_0$. The centralization identities therefore hold
for all elements of $P$ at every height-one point. Both sides
are morphisms between vector bundles, so the identities hold on
$Y$. Their multiplication identities also pull back literally.
Locally write $x=za$ and $y=zb$. The functions $a,b$ generate the
unit ideal. On the chart where $a$ is invertible use $a^{-2}e_x$;
on the other chart use $b^{-2}e_y$. This proves (D.SN.3).

Every identity is an equality of fixed tensor morphisms. Scalar
extension and adic completion preserve it. Frobenius pullback
of locally free modules preserves the displayed tensor identities
and replaces $z^2$ by $z^{2p}$, also over nonreduced coefficient
rings. None of these steps requires division by a factorial.
\end{proof}

\begin{proposition}[Proved: explicit Newton lifting and averaging for scaled-separable actions]
\label{D:scaled-separable-action-Newton}
Let $S$ be a commutative ring, complete and separated for $(z)$,
where $z$ is a nonzerodivisor. Let $P$ be a finite projective
unital associative $S$-algebra of positive rank at every point,
and suppose that an actual tensor
\[
 e=\sum_i a_i\otimes b_i\in P\otimes_S P,
 \qquad (a\otimes1)e=e(1\otimes a),\qquad
                       \sum_i a_i b_i=z^s\,1            \tag{D.SN.7}
\]
is fixed, for an integer $s\ge0$. Let $V$ be a finite projective
$S$-module. Suppose a unital $S$-linear map
$\rho:P\to\operatorname{End}_S(V)$ has multiplicativity defect
in $z^N\operatorname{End}_S(V)$, with $N>2s$.
Then the explicit iteration below produces a unital algebra
map $\rho_\infty$ with
\[
                \rho_\infty\equiv\rho\pmod {z^{N-s}}.
                                                        \tag{D.SN.8}
\]
In particular a sufficiently precise truncated action on a
locally free module admits, locally after lifting its underlying
unital linear map, an exact formal action with the stated loss
of precision. It does not assert that the whole original
$N$-jet is retained.

If $\rho$ on $V$ and $\sigma$ on $W$ are already exact actions,
and $f:V\to W$ is $S$-linear with intertwining defect divisible
by $z^N$, $N\ge s$, then
\[
 \operatorname{Av}(f)=z^{-s}\sum_i\sigma(a_i)f\rho(b_i)
                                                        \tag{D.SN.9}
\]
is an integral exact intertwiner, congruent to $f$ modulo
$z^{N-s}$. If $N>s$ and $f$ is invertible modulo $z$, then
$\operatorname{Av}(f)$ is an isomorphism. Consequently two
exact actions on the same module which are congruent modulo
$z^N$, $N>s$, are conjugate by an automorphism congruent to the
identity modulo $z^{N-s}$.

The constructions work with nilpotent coefficient rings and
are compatible with coefficient changes for which $z$ remains
regular and the prescribed adic completions are taken. For
the fixed full orders of the preceding lemma the local action
loss has $s=2$, and for their Frobenius pullbacks it has $s=2p$.
This is local algebra-action and intertwiner effectivity, not
simultaneous global object, Morita-unit or cocycle effectivity.
\end{proposition}

\begin{proof}
Write
\[
 d(a,b)=\rho(a)\rho(b)-\rho(ab),\qquad
 H(d)(a)=\sum_i\rho(a_i)d(b_i,a).
                                                        \tag{D.SN.10}
\]
All these are $S$-multilinear maps. Since $d$ is divisible by
$z^N$ and $z$ is regular on the finite projective target,
$c=z^{-s}H(d)$ is a well-defined integral linear map divisible
by $z^{N-s}$. It vanishes at $1$. The following exact identity
requires neither multiplicativity of $\rho$ nor a field:
\[
 \begin{split}
 (\delta_\rho H(d))(a,b)
   &:=\rho(a)H(d)(b)-H(d)(ab)+H(d)(a)\rho(b)\\
   &=z^s d(a,b)
      +\sum_i d(a_i,b_i)d(a,b)
      +\sum_i d(a,a_i)d(b_i,b).
 \end{split}                                           \tag{D.SN.11}
\]
To verify it, replace $\rho(a)\rho(a_i)$ by
$\rho(aa_i)+d(a,a_i)$ and use centralization of $e$ to
replace $\sum\rho(aa_i)d(b_i,b)$ by
$\sum\rho(a_i)d(b_i a,b)$. The purely associative defect
identity
\[
 \rho(u)d(a,b)-d(ua,b)+d(u,ab)-d(u,a)\rho(b)=0
                                                        \tag{D.SN.12}
\]
then gives $\sum\rho(a_i)\rho(b_i)d(a,b)$, besides the
last sum in (D.SN.11). Finally use
$\sum\rho(a_i)\rho(b_i)=z^s I+\sum d(a_i,b_i)$.
This proves the formula without an omitted Hochschild
cohomology obstruction argument.

Set $\rho'=\rho-c$. Its defect is
\[
 d'=d-\delta_\rho c+c\,c
       \in z^{\,2N-2s}\operatorname{Hom}_S
             (P\otimes P,\operatorname{End}_S(V)).      \tag{D.SN.13}
\]
Indeed the two quadratic sums in (D.SN.11), after division
by $z^s$, have order at least $2N-s$, and $c\,c$ has order
at least $2N-2s$. The new map is still unital. Iterate with
\[
 N_j=2^j(N-2s)+2s,
 \qquad \rho_{j+1}=\rho_j-z^{-s}H(d_j).               \tag{D.SN.14}
\]
The corrections have orders at least $N_j-s$, tending to
infinity. Completeness gives a unital limit, separatedness
makes its defect zero, and the first correction bound proves
(D.SN.8). The iteration is an explicit construction, not a
claim of uniqueness of the lifted action.

For exact actions put
$D(a)=\sigma(a)f-f\rho(a)$. Then (D.SN.7) gives the
integral expression
\[
 \operatorname{Av}(f)
       =f-z^{-s}\sum_i\sigma(a_i)D(b_i),              \tag{D.SN.15}
\]
which proves its claimed congruence. Centralization of $e$
gives $\sigma(a)\operatorname{Av}(f)
=\operatorname{Av}(f)\rho(a)$ exactly. If the correction
has positive order and $f$ is invertible modulo $z$, the
usual convergent inverse, or the local determinant criterion,
makes $\operatorname{Av}(f)$ invertible over $S$.
Apply this with $f=I$ for the conjugacy assertion.

For clarity, an underlying unital linear lift of a truncated
action is available locally: finite projectivity lifts the
linear map, and $1\in P$ is locally unimodular, so the value
at $1$ can be fixed. If a global linear lift is supplied, the
same formulas give a global action on that affine base.
All divisions above are divisions of elements already known
to be in the displayed power of $(z)$. Therefore they are
unique and commute with the indicated coefficient maps.
In the geometric application $S$ is an affine Cartier-adic
completion after extension of a coefficient field; $z$ remains
regular even when the coefficient algebra is nonreduced.

One may average successively for two commuting algebra
actions. The two averaging operators preserve each other's
exact intertwining property, giving total loss at most the
sum of their exponents. This is useful for balanced bimodule
maps. It still does not preserve arbitrary independently
imposed product relations between several maps. Such Morita
inverse identities and overlap cocycles require a separate
argument before these local statements can be used in a
hereditary projective count.
\end{proof}


% Separate post-CP26 proved geometric comparison. No finite-family effectivity or hereditary projective count.
\subsection{Bounded-period root anchors and canonical full orders}

\begin{proposition}[Proved: a stable twisted anchor and its induced lattices]
\label{D:stable-root-anchor-canonical-comparison}
Use the fixed complete surface base, a degree-$t$ normal punctured
cover $A/R$, and its original comparable-cycle model
$Y_0\to\operatorname{Spec}A$ from
Proposition~\ref{D:uniform-comparable-reduction-cycle}.
Let $\rho:Y\to Y_0$ be a regular point-blowup refinement, and
write $C=\rho^*C_0$, where
$(x,y)\mathcal O_{Y_0}=\mathcal O_{Y_0}(-C_0)$ with $x,y\in J$.
All full bundles below mean the actual pullbacks from $Y_0$;
they are not newly canonicalized on $Y$.

Suppose the reduced exceptional support on $Y$ is SNC and
$\mathcal X\to Y$ is a tame root stack on that support. Let
$\mathcal G\to\mathcal X$ be a banded $\mu_n$-gerbe,
$(n,p)=1$, extending a realized coefficient class on the
puncture $U$. Fix a scalar weight and a rank-$h$ twisted
bundle $\mathcal A_U$ which admits a vector-bundle extension
to $\mathcal G$. Put
\[
 P_0=j_*\mathcal End(\mathcal A_U),\qquad
 P=\rho^*\operatorname{full}_{Y_0}(P_0),
 \qquad j:U\hookrightarrow\operatorname{Spec}A.
                                                        \tag{D.BA.1}
\]
Descent of scalar-weight-zero bundles to $U$ is understood.
There exists a vector-bundle extension $\mathcal A_1$ of
$\mathcal A_U$, of the same rank $h$, with an actual left
$P$-action. If $q:\mathcal G\to Y$ is the coarse map, then
$\mathcal B_1=q_*\mathcal End(\mathcal A_1)$ is locally free
over $Y$ and, inside its fixed generic endomorphism algebra,
\[
                  P\subset\mathcal B_1\subset P^\vee
                             \subset P(C).             \tag{D.BA.2}
\]
The middle order need not be Azumaya on the exceptional divisor.

For every punctured twisted bundle $\mathcal F_U$ of the
same scalar weight, set
\[
 \begin{aligned}
 Q_0&=j_*\mathcal Hom(\mathcal A_U,\mathcal F_U),&
 R_0&=j_*\mathcal Hom(\mathcal F_U,\mathcal A_U),\\
 Q&=\rho^*\operatorname{full}_{Y_0}(Q_0),&
 R_1&=\rho^*\operatorname{full}_{Y_0}(R_0).
 \end{aligned}                                           \tag{D.BA.3}
\]
There is a vector-bundle extension $\mathcal F$ on
$\mathcal G$, of the original rank, for which the ordinary
locally free coarse Hom bundle satisfies
\[
 Q\subset q_*\mathcal Hom(\mathcal A_1,\mathcal F)
       \subset R_1^\vee\subset Q(C).                  \tag{D.BA.4}
\]
Every identification and inclusion restricts to the specified
one on $U$. These are assertions about individual geometric
bundles and actual sheaf lattices, not an assertion about
arbitrary relative reflexive hulls.
\end{proposition}

\begin{proof}
The canonical order and module actions are those of
Proposition~\ref{D:canonical-order-Morita-conductor}.
We recall the dual comparison to fix its precise scope.
For any finite reflexive module $L_0$, put
$L=\operatorname{full}_{Y_0}(L_0)$ and
$L'=\operatorname{full}_{Y_0}(L_0^\vee)$.
The natural evaluation gives $L'\subset L^\vee$.
For $z\in J$, the first-power Noether factorization writes
$z\operatorname{id}_{L_0}=vu$ through a finite free module.
The entries $u_i$ are sections of $L'$; hence for a local
$\lambda\in L^\vee$,
\[
                 z\lambda=\sum_i\lambda(v_i)u_i\in L'.
                                                        \tag{D.BA.5}
\]
Using $z=x,y$ proves $L^\vee(-C_0)\subset L'\subset L^\vee$.
Pull back these actual maps to $Y$.
The trace pairings on the puncture extend reflexively over
the normal surface, so $P_0=P_0^\vee$ and $R_0=Q_0^\vee$.
Thus, on the refined model as well,
\[
 P\subset P^\vee\subset P(C),\qquad
 R_1\subset Q^\vee\subset R_1(C),\qquad
 Q\subset R_1^\vee\subset Q(C).                       \tag{D.BA.6}
\]
Neither a boundary Azumaya property nor a division by $h$
is used in these dual inclusions.

Start with any vector-bundle extension $\mathcal A^{(0)}$
of $\mathcal A_U$. Form the coherent image
\[
 \mathcal L=\operatorname{im}
   \bigl(q^*P\otimes\mathcal A^{(0)}
                   \longrightarrow\mathcal A^{(0)}(*E)\bigr),
             \qquad \mathcal A_1=\mathcal L^{\vee\vee}.
                                                        \tag{D.BA.7}
\]
The map is the prescribed generic action. Coherence of the
image follows locally by clearing finitely many denominators;
on a finite cover they can be cleared by some exceptional
Cartier multiple. No uniform bound for that preliminary
denominator is needed or claimed. The unit makes
$\mathcal A^{(0)}\subset\mathcal L$. Associativity of $P$
makes $\mathcal L$ stable under $q^*P$.

The gerbe and the root stack have regular two-dimensional
etale scheme charts. On these charts a coherent reflexive
sheaf is locally free, and maps to a reflexive target extend
uniquely once they extend at height one
\citep[Tags 0AVC and 0B3N; Lemmas 31.13.12--31.13.15]
{C_StacksRootAnchorReflexive}.
Reflexive hull and these maps commute with the etale
transition maps, so the construction descends to
$\mathcal G$. It gives a rank-$h$ vector bundle and an
actual $P$-action on $\mathcal A_1$. The algebra identities
hold everywhere because they hold on the dense puncture.

Scalar-weight-zero Hom descends from the gerbe to the root
stack. Its coarse pushforward is locally free: on a root
chart the coordinate algebra is finite locally free over
the regular coarse chart, a bundle is projective over that
algebra, and taking invariants under the finite tame root
group is a direct summand. Thus its invariants are projective
over the coarse chart. This also proves coherence and the
local-freeness assertions for the later coarse Hom bundles.

The action embeds $P$ into $\mathcal B_1$. If $b$ is a
section of $\mathcal B_1$ and $a$ a section of $P$, the
trace of their product on $\mathcal A_1$ is regular on
$\mathcal G$ and descends to $\mathcal O_Y$.
It agrees on $U$ with the ordinary matrix trace pairing
used in (D.BA.6). Consequently $b\in P^\vee$.
This proves (D.BA.2), including at the closed points.

For (D.BA.4), use the right $P$-module $Q$ to define
\[
 \mathcal F=
   \bigl((q^*Q\otimes_{q^*P}\mathcal A_1)_{
                  \mathrm{torsion\,free}}\bigr)^{\vee\vee}.
                                                        \tag{D.BA.8}
\]
Here torsion-free quotient means the image in the specified
generic twisted vector space. On $U$, ordinary matrix
Morita evaluation identifies that space with $\mathcal F_U$.
The regular-chart argument just given makes (D.BA.8) a
vector bundle of its original rank. Tensor evaluation
gives the actual map
$Q\to q_*\mathcal Hom(\mathcal A_1,\mathcal F)$.
It is injective because it is an isomorphism generically.

The canonical composition $R_1\otimes Q\to P$ also gives
\[
 q^*R_1\otimes
   (q^*Q\otimes_{q^*P}\mathcal A_1)
                 \longrightarrow\mathcal A_1.
\]
It kills torsion. At height one the torsion-free quotient
is already reflexive; hence this map extends to (D.BA.8),
by the same height-one extension criterion. It provides
$R_1\to q_*\mathcal Hom(\mathcal F,\mathcal A_1)$.
Pair this map with $q_*\mathcal Hom(\mathcal A_1,\mathcal F)$
by composition and trace on $\mathcal A_1$. The pairing is
regular and is the perfect generic pairing $Q_0\times R_0$.
Therefore the coarse Hom lattice lies in $R_1^\vee$.
The last inclusion in (D.BA.4) is (D.BA.6).
\end{proof}

\begin{corollary}[Proved: bounded-rank generating anchors after a bipartite refinement]
\label{D:bounded-period-generating-root-anchor}
In the preceding cover setup, suppose the geometric Brauer
period is $M\le B$, with $(M,p)=1$. After an $O_R(t)$
point-blowup SNC refinement and one further blowup of each
SNC node, use the common $M$-root stack of the entire
exceptional divisor. The actual coefficient gerbe admits
a scalar-weight-one generating bundle $\mathcal A_{\rm gen}$
of rank at most $B^4$.

More precisely, choose a rank-$h$ punctured anchor with
$h=M^2$ by the qualitative residual-anchor construction
in Proposition~\ref{D:SNC-general-weighted-presentation},
and choose $\mathcal A_1$ as above. There are $d=M^2$
line twists, each canonically trivial on $U$, such that
their sum tensored with $\mathcal A_1$ is a generating
bundle. Put $P^{[d]}=\operatorname{Mat}_d(P)$ and
$Q^{[d]}=Q^{\oplus d}$. Under their canonical punctured
identifications, its actual ordinary coarse order satisfies
\[
 P^{[d]}(-C)\subset
 q_*\mathcal End(\mathcal A_{\rm gen})
                  \subset P^{[d]}(2C),                \tag{D.BA.9}
\]
and the induced extension $\mathcal F$ of (D.BA.8) satisfies
\[
 Q^{[d]}(-C)\subset
 q_*\mathcal Hom(\mathcal A_{\rm gen},\mathcal F)
                  \subset Q^{[d]}(C).                 \tag{D.BA.10}
\]
All pole constants are independent of the covering degree,
coefficient order, rank of $\mathcal F_U$, and number of
exceptional components. They remain valid for any scalar
weight prime to the coefficient order.
\end{corollary}

\begin{proof}
Lemma~\ref{D:linear-SNC-refinement-centres} and
Corollary~\ref{D:SNC-refinement-inherited-cycles} give
the first SNC refinement with $O_R(t)$ nodes and the
inherited Cartier cycles. Blow up each of its nodes once.
All original components are now pairwise disjoint, and all
new exceptional components are pairwise disjoint. Each node
joins one component of each kind. This is an actual bipartite
divisor, also when the old graph had cycles or parallel edges.
These are $O_R(t)$ extra multiplicity-two blowups; all inherited
cycle and cohomology identities still hold. Write the two
sets of components as $I_0,I_1$.

The actual coefficient class extends to its mixed residue-root
stack by Lemma~\ref{D:SNC-general-Brauer-geometry}. Its residue
orders divide $M$, so pull it to the common $M$-root stack.
The geometric-period-$M$ residual anchor of rank $M^2$,
followed by formal lifting and algebraization, is precisely
the qualitative existence construction (D.GP.3)--(D.GP.5).
Use its restriction as $\mathcal A_U$. No finite etale
splitter of degree bounded in $M$ is inferred.

Let $\mathcal D_i$ be the $M$-root divisor over $E_i$, and put
\[
 \mathcal T_{a,b}=\mathcal O_{\mathcal X}
       \left(a\sum_{i\in I_0}\mathcal D_i
                    +b\sum_{i\in I_1}\mathcal D_i\right),
 \quad 0\le a,b<M,\qquad
 \mathcal A_{\rm gen}
       =\bigoplus_{a,b}\mathcal A_1\otimes\mathcal T_{a,b}.
                                                        \tag{D.BA.11}
\]
Their effective rational divisor coefficients on $Y$ are
$a/M$ or $b/M$, all less than one. Every coefficient of
$C$ along the exceptional support is a positive integer.
The canonical root sections therefore give regular maps
\[
 \mathcal O\longrightarrow\mathcal T_{a,b}
                      \longrightarrow\mathcal O(C),
 \qquad
 \mathcal O(-C)\longrightarrow
       \mathcal T_{a',b'}\mathcal T_{a,b}^{-1}
                      \longrightarrow\mathcal O(C),   \tag{D.BA.12}
\]
all restricting to the chosen identity on $U$.

At a node, the twists in (D.BA.11) give every character of
$\mu_M\times\mu_M$ once; at a one-component point they give
every character of $\mu_M$. This still generates in the
presence of a noncommuting central gerbe inertia. Indeed let
$D$ be that finite tame inertia group, $Z$ its scalar centre,
and $H=D/Z$ its abelian root quotient. For any nonzero
representation $V$ of the specified scalar weight,
\[
 V\otimes k[H]\simeq
       \operatorname{Ind}_Z^D\operatorname{Res}_Z^D V.
                                                        \tag{D.BA.13}
\]
For every irreducible $W$ of that scalar weight, Frobenius
reciprocity gives
$\operatorname{Hom}_D(W,V\otimes k[H])
 =\operatorname{Hom}_Z(W,V)\ne0$.
Semisimplicity holds because all inertia orders are prime to
$p$. Hence every scalar-weight-one inertia simple occurs.
Away from the exceptional divisor the scalar centre alone
remains, and a nonzero bundle suffices. On the finite root
charts, this fibrewise statement and Nakayama give a projective
generator, exactly as in the generating-bundle proof of
Lemma~\ref{D:SNC-general-Brauer-geometry}.
Thus (D.BA.11) is an actual generating bundle of rank
$hM^2=M^4\le B^4$.

For each matrix block in its endomorphism algebra,
(D.BA.12), exact tame pushforward, and (D.BA.2) give
\[
 P(-C)\subset\mathcal B_1(-C)
 \subset q_*\mathcal Hom(
       \mathcal A_1\mathcal T_{a,b},
       \mathcal A_1\mathcal T_{a',b'})
 \subset\mathcal B_1(C)\subset P(2C).
                                                        \tag{D.BA.14}
\]
Summing the blocks proves (D.BA.9). In particular the claim
is not $P^{[d]}\subset q_*\mathcal End(\mathcal A_{\rm gen})$:
matrix units between unequal inertia weights can require
one Cartier factor.

Put $E_1=q_*\mathcal Hom(\mathcal A_1,\mathcal F)$.
The positivity in the first part of (D.BA.12) gives
\[
 E_1(-C)\subset
 q_*\mathcal Hom(\mathcal A_1\mathcal T_{a,b},\mathcal F)
                  \subset E_1.                        \tag{D.BA.15}
\]
Combine this with (D.BA.4) and sum all $M^2$ blocks to
obtain (D.BA.10). All constructions and trace identities
work for any fixed scalar weight coprime to $n$, giving
the last assertion.
\end{proof}

\begin{remark}[What the comparison does not establish]
The canonical bundles $P,Q$ are generated after the inherited
strong-model twist, and have the previously proved degree
budgets. The intermediate coarse lattices in (D.BA.9)--(D.BA.10)
are not thereby proved generated or cohomologically bounded:
a divisor sandwich alone does not bound their instability.
Nor has this construction made the arbitrary choices of
reflexive hulls flat in Artinian families. The exact fixed
generating-bundle Morita equivalence is available once the
anchor has been chosen, but arbitrary finite nilpotent order
actions are not asserted to produce locally free root-stack
lattices. Simultaneous Frobenius comparison, finite-data
effectivity, and a degree-linear whole-cycle presentation
remain separate requirements. These two results are geometric
comparison inputs, not a hereditary projective count or an
independent review movement.
\end{remark}


% Separate post-CP27 proved packet; four full actual-file audits and exact local-frame precision refresh. Morita only on the puncture.
\subsection{Finite Morita--Frobenius data and actual punctured unit objects}

\begin{theorem}[Proved: finite Morita--Frobenius object effectivity]
\label{D:finite-Morita-Frobenius-object-effectivity}
Use the original comparable-cycle model $Y\to\operatorname{Spec}A$
and $C\le Z\le\alpha C$ of
Proposition~\ref{D:uniform-comparable-reduction-cycle}.
The same statement holds for their actual point-blowup pullbacks.
Fix $q=p^f$, $n=q-1$, and a banded $\mu_n$-gerbe
$\mathcal G\to U$ on the puncture. Suppose fixed rank-$h$
anchors $\mathcal A_a$ of scalar weight $p^a$ are given,
$0\le a<f$, with $p\nmid h$; put $\mathcal A_f=\mathcal A_0$.
All indices below are cyclic. On $Y$ take the canonical full
orders and comparison bundles
\[
 \begin{aligned}
 P_a&=\operatorname{full}\bigl(j_*\mathcal End\mathcal A_a\bigr),\\
 K_a&=\operatorname{full}\bigl(j_*\mathcal Hom(
                   \mathcal A_{a+1},F^*\mathcal A_a)\bigr),\\
 R_a&=\operatorname{full}\bigl(j_*\mathcal Hom(
                   F^*\mathcal A_a,\mathcal A_{a+1})\bigr)(-pC).
 \end{aligned}                                         \tag{D.MF.1}
\]
Here $F$ is absolute $p$-Frobenius. The scalar-weight-zero
Hom bundles are descended to $U$ before taking their full
pullbacks. On a refinement use their actual pullbacks.

Set $b=p+1$. Choose an integer
\[
 N\ge\max\{\lceil\alpha\rceil,4p+10\},\qquad
                       L=N+3p+5.                      \tag{D.MF.2}
\]
For every coefficient $k$-algebra $T$, consider the following
actual data on $LC_T$: vector bundles $E_a$ of rank $hr$,
unital right $P_a$-actions, and balanced intertwining maps
\[
 \mu_a:F^*E_a\otimes K_a\longrightarrow E_{a+1},\qquad
 \nu_a:E_{a+1}\otimes R_a\longrightarrow F^*E_a.
                                                        \tag{D.MF.3}
\]
Restriction of Frobenius pullback from $pLC_T$ is understood.
The actions and the balance and intertwining relations hold on
$LC_T$. Require also the two finite right-inverse equations
(D.MF.6) below, for the fixed pencil elements $x,y$.

These data have an actual formal realization on the
$C$-adic completion, retaining their specified $NC_T$ marking.
The realization consists of vector bundles, exact order
actions, and maps (D.MF.3) which are inverse Morita unit maps
on the formal puncture. Over a finite-type Noetherian
coefficient base, or a finite local Artinian coefficient
base, ordinary Grothendieck existence on the completion of
$A\otimes_kT$ algebraizes them. Azumaya Morita equivalence
on $U$ then gives an actual scalar-weight-one $f$-stage
unit bundle on $\mathcal G$ and its finite-field frames.
If two actual realizations have identical marked data on
$NC_T$, they are isomorphic on the puncture. No assertion
that a reconstructed root-stack lattice is flat is used.
\end{theorem}

\begin{proof}
\emph{The actual fixed Morita context.}
The right module convention is important:
$K_a$ is left $F^*P_a$, right $P_{a+1}$, and $R_a$ has the
opposite actions. Put
$P_a^+=\operatorname{full}(j_*\mathcal End(F^*\mathcal A_a))$.
There are actual inclusions
\[
                 F^*P_a\subset P_a^+\subset F^*P_a(pC).
                                                        \tag{D.MF.4}
\]
To verify this precision comparison, let $M$ be any reflexive
$A$-module. Frobenius applied to a factorization of
$x\operatorname{id}_M$ through a free module factors
$x^p\operatorname{id}_{F^*M}$ through a free module.
Extend the first map to $(F^*M)^{**}$ by target Hartogs.
Its composite still equals $x^p$; the second map lands in
the ordinary Frobenius image. Pull back these actual maps
to $Y$. At height one $F^*$ is flat on the regular surface,
so the ordinary image is $F^*\operatorname{full}(M)$.
The same argument for $y$ and the identity
$(x^p,y^p)\mathcal O_Y=\mathcal O_Y(-pC)$ prove (D.MF.4),
including at closed points. The natural other inclusion
is obtained from the actual Frobenius pullback map at
height one and reflexive extension.

Proposition~\ref{D:canonical-order-Morita-conductor}, followed
by the $-pC$ twist in (D.MF.1), therefore gives regular
pairings
$K_a\otimes R_a\to F^*P_a$ and
$R_a\otimes K_a\to P_{a+1}$.
They satisfy the Morita associativity identities. Their
images contain the corresponding orders times
$\mathcal O_Y(-bC)$, and they are Morita equivalences on
$U$. In particular for $u=x,y$ there are fixed actual
finite tensors
\[
 c_{a,u}=\sum_i r_{a,u,i}\otimes k_{a,u,i}\in R_a\otimes K_a,
 \qquad \operatorname{ev}(c_{a,u})=u^b1_{P_{a+1}}.
                                                        \tag{D.MF.5}
\]
One obtains them by multiplying the actual Noether-factorization
sections in (D.OC.10) by $u^p$ in the opposite factor.
They need not be centralizing in the unbalanced tensor
product. Their image in the balanced tensor on $U$ is
$u^b$ times the inverse Morita unitor.
The input right-inverse equations are exactly
\[
 \sum_i\mu_a\bigl(\nu_a(v\otimes r_{a,u,i})
                          \otimes k_{a,u,i}\bigr)=u^b v
       \quad(u=x,y)\quad\text{on }LC_T.               \tag{D.MF.6}
\]
These are equations of actual ordinary bundle maps, with
no localization of a nonreduced cycle.

\emph{Exact local actions and maps.}
Work on affine formal charts where one of $x,y$ generates
$\mathcal O(-C)$; call that generator $z$, absorbing a
unit into the fixed tensors. Refine these charts so that
all input bundles $E_a$ have frames on $LC_T$, and use
those frames as free formal lifts. Such formal affine
charts exist by local freeness on the finite thickening;
its higher nilpotent thickenings are affine, with surjective
transition maps on functions
\citep[Lemma 37.2.3, Tag 06AD]{C_StacksFiniteObjectAffineness}.
This is exactly the chart construction in
Theorem~\ref{D:finite-Frobenius-object-effectivity}.
All the following divisions are by powers of a
nonzerodivisor on these finite free local modules.
After arbitrary coefficient extension $z$ stays regular
also in the separated $z$-adic completion: this follows
by taking the inverse limits of the exact multiplication
sequences for $S/z^j$.

Lift each unital linear action and apply
Proposition~\ref{D:scaled-separable-action-Newton}.
The full orders have separability exponent $2$, so the
exact actions retain precision $L-2$. Their Frobenius
pullbacks have exponent $2p$.
Lift $\mu_a,\nu_a$ as ordinary linear maps. Average on
the two commuting actions to make them balanced and
intertwining. For example, balancing a map $f$ uses
$z^{-s}\sum f(ma_i,b_i k)$ for a scaled separability
tensor $\sum a_i\otimes b_i$; right-module averaging
uses $z^{-s}\sum f(ma_i)b_i$. Centralization gives the
exact relations, and subtracting the original map is
a sum of its relation defects divided by $z^s$.
The two averages commute because the two actions on
the fixed bimodule commute. The total loss is at most
$2+2p$. Thus the exact maps retain precision
\[
                         M=L-2p-4=N+b.                \tag{D.MF.7}
\]
This step imposes exact module relations, not merely
independently chosen ordinary maps.

\emph{Making the Morita inverse exact.}
Use the averaged maps and the tensor (D.MF.5) to form
\[
 T_a=z^{-b}\sum_i\mu_a\bigl(\nu_a(-\otimes r_i)
                                  \otimes k_i\bigr).
                                                        \tag{D.MF.8}
\]
Equation (D.MF.6) and (D.MF.7) show that this is an
integral endomorphism equal to $I$ modulo $z^N$.
On the punctured chart it is $P_{a+1}$-linear because
both maps are exact module maps and the balanced tensor
is the Morita unitor. The equality expressing linearity
then holds everywhere by $z$-torsion-freeness.
The convergent inverse of $T_a$ is integral. Replace
$\nu_a$ by its composite with $T_a^{-1}$ on its source;
this changes it only modulo $z^N$.
The induced map from the balanced Morita source to
$E_{a+1}$ now has an actual right inverse on the puncture.
Both are projective modules of the same ordinary rank
$hr$: this is matrix Morita equivalence for the Azumaya
orders on $U$. A right inverse between equal-rank finite
projectives is an inverse, as can be checked after local
frames by determinants. Thus $\mu_a$ is a genuine Morita
isomorphism there. Its inverse has an unbalanced
representative with pole at most $bC$, given by the
corrected $\nu_a$ and $z^{-b}c_a$.
Any ordinary tensor identities for these inverse maps
which have regular targets hold on the whole formal
chart by torsion-freeness. No centralizing identity for
an unbalanced coevaluation tensor was assumed.

\emph{Correcting the overlaps.}
Lift the marked bundle transition maps between charts.
Averaging makes them exact $P_a$-intertwiners and retains
precision at least $L-4$. They remain invertible.
On the puncture the exact Morita maps define the
cyclic transport operator
\[
 \mathcal T(g)_{a+1}
       =\mu_{a,i}\,(F^*g_a\otimes1)\,\mu_{a,j}^{-1}.
                                                        \tag{D.MF.9}
\]
It is additive, functorial for compositions and inverses,
and sends an error of order $j$ to one of order at least
$pj-b$. The original finite compatibility of the maps,
retained to precision $M$, shows
$\mathcal T(g)-g\in\mathcal I_C^N$.
Since $pN-b>N$, iteration converges to integral
intertwiners with the same $NC_T$ marking.
A fixed point in that congruence class is unique.
Products of two overlap limits and the direct overlap
limit are both fixed points with the same marked jet;
therefore they agree. The inverse and identity equations
follow in exactly the same way. This gives actual bundle
cocycles, exact global order actions and globally compatible
Morita maps. Compatibility of the regular unbalanced maps
extends from the puncture by torsion-freeness.
This is a simultaneous object and cocycle construction.

For comparison of two realizations with identical $NC_T$
data, average lifts of their marked identity maps.
The averaged maps retain precision $N-2$; their unit
transport defect has order at least $N-b-2$.
The choice (D.MF.2) implies
$p(N-b-2)-b>N-b-2>0$.
The same iteration therefore yields inverse integral
unit intertwiners, proving the claimed punctured
isomorphism. Exact retention of the entire comparison
$N$-jet is not asserted.

The construction is formal over every coefficient ring.
For a finite-type Noetherian $T$ use the completion of
$A\otimes_kT$ at the original maximal ideal; for finite
local Artinian $T$ use the same complete base directly.
The Cartier and closed-fibre topologies are cofinal.
Ordinary Grothendieck existence algebraizes the vector
bundles and all fixed tensor morphisms. Local freeness
follows from its finite-cycle marking and properness over
the complete base. These are precisely the Noetherian
scope and morphism-effectivity inputs of the earlier
finite-object theorem; no Noetherian existence theorem
is being applied to an arbitrary coefficient ring.
Finally use Azumaya Morita only on $U$. The right
$P_a$-module gives the weight-$p^a$ twisted bundle
$E_a\otimes_{P_a}\mathcal A_a$ of rank $r$.
The maps (D.MF.3) become its Frobenius isomorphisms.
The last weight is $q\equiv1\pmod n$, so this is the
required cyclic unit object. The already established
finite-field frame correspondence applies, also with
nilpotent coefficients.
\end{proof}

\begin{proposition}[Proved: generated Morita-unit classes are open]
\label{D:generated-Morita-unit-Isom-openness}
Keep the fixed orders and Morita context above. Let an
actual family of such ordinary order-module unit bundles
be given over a finite-type coefficient base, on the
complete proper surface family. Suppose every $E_a(-Z)$
is generated and every determinant line $\det E_a$ is
fixed formally. It suffices to retain on $Z$ a specified
generating quotient and a fixed determinant line: the
weak-cycle lifting of
Lemma~\ref{D:weak-cycle-generated-determinant-lifting}
gives exactly these conditions.
Then the locus of every punctured scalar-weight-one
unit isomorphism class is open and closed. For a fixed
comparison class, the unit-Isom functor is represented
by an affine scheme of finite presentation and is etale
over the parameter scheme on its image. Its pole bound
may depend on the class, rank and determinant chart;
its coordinates are not charged to the quantitative
object atlas.
\end{proposition}

\begin{proof}
We first give a finite presentation for integral unit
intertwiners between two already actual systems $E,E'$.
Let $H_a=\mathcal Hom(E_a,E'_a)$ be the ordinary Hom bundle.
For $u=x,y$ let $e_{a,u}\in P_a\otimes P_a$ be the actual
scaled separability tensor with product $u^2$.
Write $A_{a,u}(g)$ for its unnormalized right-module
average, so $u^{-2}A_{a,u}(g)$ is an exact $P_a$-intertwiner
on the open set where $u$ is invertible.
Follow this average by the Morita transport from $E$ to
$E'$. Using the inverse representative (D.MF.5), one obtains
a regular, additive Frobenius-semilinear operator $S_{a,u}$
with
\[
 S_{a,u}=u^d\mathcal T_{a,u},\qquad d=2p+b,
                                                        \tag{D.MF.10}
\]
where $\mathcal T_{a,u}$ first averages and then transports.
There are no denominators in $S_{a,u}$: the $u^{2p}$
factor clears the Frobenius average and $u^b$ clears the
Morita inverse. All ingredients are actual global tensor
morphisms, so $S_{a,u}$ is global, even where $u$ is not
a local generator of $C$.

Choose $n_0$ with $pn_0>n_0+d$. A jet
$\eta_a\in H^0(n_0C_T,H_a)$ defines both sides of
\[
 S_{a,u}(\eta_a)=u^d\eta_{a+1}
       \quad\text{on }(n_0+d)C_T,
                \qquad u=x,y.                         \tag{D.MF.11}
\]
The Frobenius side uses the jet on $pn_0C_T$;
the other side is independent of a lift because
$u^d\mathcal I_C^{n_0}\subset\mathcal I_C^{n_0+d}$.
Thus these are intrinsic finite-cycle equations.
On the chart where $u$ generates $C$, a local formal
lift satisfies $\mathcal T_u(\eta)-\eta\in\mathcal I_C^{n_0}$.
This operator sends error order $j$ to at least $pj-d$.
Iteration converges to an integral fixed point. Its image
is an exact $P$-intertwiner and an actual unit map on the
puncture, since the normalized average is already an
exact intertwiner there. Different charts have the same
marked jet; their limits, now exact unit intertwiners,
coincide by the uniqueness estimate for the true Morita
transport. Conversely every unit intertwiner satisfies
(D.MF.11). Consequently the finite-cycle solution functor
represents the actual integral unit-Hom functor, over all
coefficient rings in the formal setting.

The ordinary Hom functors on these proper finite cycles
are affine of finite presentation, by the fully checked
Hom theorem used in
Proposition~\ref{D:intrinsic-integral-unit-Hom-jets}.
Equations (D.MF.11) cut out a finite-presentation equalizer.
Paired maps and finite inverse equations give unit-Isom;
uniqueness of unit maps with a sufficiently deep marked
jet makes those finite inverse equations exact.
There is no appeal to representability of a raw higher-jet
image functor.

For a bounded meromorphic map, apply the same construction
to $H_a(\ell C)$. The loss becomes
$2p+b+(p-1)\ell$; clearing that power gives exactly the
same global finite-cycle equations. Product equations are
written in the $2\ell C$ twist, as in the previously proved
shifted unit-Isom construction. This provides finite
presentation for prescribed bounded poles without
localizing a finite nilpotent cycle.

It remains to prove a pole bound for every specified
Artinian unit isomorphism to a fixed punctured class.
Choose the canonical full ordinary comparison bundles
$E_a^{\rm can}$ for a fixed representative of that class.
The argument of
Proposition~\ref{D:generated-fixed-determinant-unit-openness}
is an argument about ordinary bundles and maps, and
applies verbatim to these coarse modules: generators of
$E_a(-Z)$ and target Hartogs extend a punctured map to
$E_a\to E_a^{\rm can}(Z)$.
The determinant line is fixed formally by its restriction
to $Z$. Relative Hartogs for the target determinant and
its inverse shows that the determinant ratio to a fixed
comparison is a unit. Its exceptional divisor is therefore
fixed through every Artinian lift. Adjugates bound the
inverse poles by $(hr-1)Z$ plus that fixed exceptional
divisor. These bounds are independent of the nilpotence
order. The same argument at each phase gives one finite
pole bound for the comparison class on this determinant
chart. It does not require the parameter fibres to be
canonical full bundles or relative reflexive hulls.

On the puncture, the corresponding twisted unit bundles
have finite etale $\mathbf F_q$-frames on the fixed gerbe.
A specified isomorphism lifts uniquely through every
nilpotent coefficient extension by finite-etale invariance.
The preceding pole bound puts this unique lift in the
finite-presentation bounded-pole Isom functor. Thus that
functor is formally etale, and hence etale. Its image is
open. Distinct unit classes have disjoint such images,
and every geometric point has an actual unit object;
the complement of any class image is the union of the
other open images. The class images are therefore clopen.
All the invoked Hom, etale-invariance, formal-etaleness
and Krull/Hartogs inputs are exactly the fully audited
ones in the intrinsic-unit and generated-family propositions.
No arbitrary root-stack lattice or flattening assertion
has entered this proof.
\end{proof}

\begin{remark}[Scope]
The effectivity theorem uses actual ordinary vector bundles
with exact full-order actions; it does not claim that their
saturation by a matrix algebra on a root chart is flat.
Morita reconstruction occurs only on the puncture, where
the fixed orders are Azumaya. Quantitative coordinate and
class bounds require the separate determinant lists and
whole-cycle Hom presentations. These two statements alone
are not a hereditary projective counting theorem.
\end{remark}


% Separate post-CP27 proved synthesis. Both Morita-effectivity dependencies and all four final exact-file audits pass.
\subsection{A hereditary projective count with bounded geometric period}

\begin{theorem}[Proved: bounded-period hereditary projective growth]
\label{D:bounded-period-hereditary-projective-growth}
Let $R$ be a fixed complete excellent normal local surface
algebra over an algebraically closed coefficient and residue
field $k$ of characteristic $p>0$, and put
$\Gamma=\pi_1(\operatorname{Spec}R\setminus\{\mathfrak m_R\})$.
Fix an integer $B\ge1$.
Then there is a constant $A_{R,B}$ such that, for every open
subgroup $H\subset\Gamma$ of index $t$, every $q=p^f$ and
every $r\ge1$, the number of ambient projective conjugacy
classes of continuous representations
\[
                   H\longrightarrow\operatorname{PGL}_r(\mathbb F_q)
                                                        \tag{D.BH.1}
\]
whose geometric Brauer obstruction has order at most $B$ is
at most $q^{A_{R,B}tr^2}$.
The full scalar lifting class is retained; its coefficient
group $\mathbb F_q^*$ is not assumed to have bounded order.
The subgroup $H$ is arbitrary: its cover need not be tame
along an exceptional divisor. No semilinear projective,
unbounded-geometric-period, or full finite-generation
assertion is included.
\end{theorem}

\begin{proof}
We spell out the geometric choices, the exact whole-cycle
equations, and the comparison needed for a component count.
We apply Theorem~\ref{D:finite-Morita-Frobenius-object-effectivity}
and Proposition~\ref{D:generated-Morita-unit-Isom-openness}
in their proved all-coefficient, retained-jet and Cartier-conductor
scopes. The local action-Newton lemma alone does not
replace either result.

\emph{The fixed cover and the coefficient classes.}
Let $A/R$ be the normal finite cover ring corresponding to
$H$, and write $U_A$ for its puncture. The cover degree is $t$.
The cases $r=1$ or $q=2$ follow from the ordinary hereditary
count, so assume $r\ge2$ and $n=q-1>1$.
The lifting obstruction injects from
$H^2(H,\mathbb F_q^*)$ into $H^2(U_A,\mu_n)$, with the
same exact fundamental-group comparison used in
Theorem~\ref{D:SNC-general-projective-quadratic-count}.
Use the controlled SNC refinement only for this cohomology
calculation and for qualitative anchor existence. Its numbers
of vertices and nodes and the sum of normalization genera
are $O_R(t)$. In (D.GG.2) use the elementary bound
$|\operatorname{coker}(N\bmod n)|\le n^v$, rather than a
bound involving $|\det N|$. Thus
\[
                |H^2(U_A,\mu_n)|\le n^{c_Rt}\le q^{c_Rt}.
                                                        \tag{D.BH.2}
\]
Only classes realized by (D.BH.1) are retained.

Fix one such coefficient class $\alpha$, with geometric
Brauer period $M\le B$. Its period is prime to $p$.
The actual lifting gerbe $\mathcal G_U$ carries the fixed
scalar band $\mu_n$. The residual-anchor construction in
\ref{D:SNC-general-weighted-presentation}, or its explicit
geometric realization in
\ref{D:bounded-period-generating-root-anchor}, supplies a
punctured scalar-weight-one anchor of rank
\[
                            h=M^2\le B^2.               \tag{D.BH.3}
\]
Here only the rank-$M^2$ initial anchor is used, not a
boundary generating bundle. There is no bounded-degree
etale splitting-cover assertion. For the $f$ Frobenius
phases choose rank-$h$ anchors $\mathcal A_a$ of scalar
weight $p^a$, $0\le a<f$. The actual coefficient gerbe
and its Frobenius identifications are retained, including
the last phase. No chosen anchor is assumed itself to be
a finite local system or Frobenius-periodic.

Return now to the original regular comparable-cycle model
$Y\to\operatorname{Spec}A$. It need not have SNC support.
Use its cycles $C\le Z\le\alpha_R C$, fixed separating
pencil in $J$, and the proved degree-linear scalar,
intersection and full-bundle generation budgets.
The anchor need only exist on $U_A$: all the sheaves used
for coordinates below are ordinary full pullbacks on $Y$.
In particular the presentation introduces no root-node
rectangles and no term $\sum a_i a_jr^2$.

\emph{Actual full orders, modules and correspondences.}
Let $P_a$ be the canonical full pullback of
$j_*\mathcal End(\mathcal A_a)$.
For a geometric weight-one unit object with phases
$\mathcal F_a$, let $E_a$ be the full pullback of
$j_*\mathcal Hom(\mathcal A_a,\mathcal F_a)$.
Use the right-module convention. The ordinary ranks are
$h^2$ and $hr$, and the right $P_a$-action on $E_a$ is
actual, not merely an action after reflexive saturation.
The full comparison bimodules and their opposites are
formed from
$\mathcal Hom(\mathcal A_{a+1},F^*\mathcal A_a)$ and
the opposite Hom, with the last phase interpreted on the
actual same gerbe. Their ranks are $h^2$.

Proposition~\ref{D:canonical-order-Morita-conductor} gives
their actual compositions and conductor inclusions.
For a Frobenius source, use the ordinary order $F^*P_a$;
the canonical full extension of its punctured algebra
lies between $F^*P_a$ and $F^*P_a(pC)$ by the same
source-only factorization that bounds canonical one-step
Frobenius inverse poles. Tensor the opposite comparison
module by $\mathcal O(-pC)$ if necessary. Its pairings
then land in $F^*P_a$ and $P_{a+1}$, and their images
contain $\mathcal O(-(p+1)C)$ times the respective orders.
Equivalently one may take exactly the scaled context
specified in
\ref{D:finite-Morita-Frobenius-object-effectivity}.
All scaling exponents here depend only on $p$, not on
$t,r,f$ or the particular coefficient class.

The canonical full orders and unscaled comparison modules
have twists by $-Z$ which are globally generated. By
Lemma~\ref{D:full-bundle-cycle-degree}, their degrees on
$Z$ are $O_{R,B}(t)$. The same facts for $E_a$ are
\[
 E_a(-Z)\text{ globally generated},\qquad
 0\le c_1(E_a)\cdot Z=O_{R,B}(tr).
                                                        \tag{D.BH.4}
\]
Tensoring by any of the fixed multiples of $C,Z$ or
taking one Frobenius pullback changes the required
budgets by fixed factors only. The line bundles
$\mathcal O(-C)$ and $\mathcal O(-Z)$ are globally
generated. These assertions use full ordinary bundles;
no divisor sandwich is used to claim generation of an
intermediate root-stack lattice.

\emph{The finite list of determinant lines.}
Suppose the fixed class $\alpha$ is represented in rank
$r$, and choose one reference twisted unit object
$\mathcal F^*$. For the common anchors,
\[
 \left.\frac{\det E_a}{\det E_a^*}\right|_{U_A}
       =\left(\frac{\det\mathcal F_a}
                        {\det\mathcal F_a^*}\right)^h.
                                                        \tag{D.BH.5}
\]
The scalar weights cancel in the ratio. The ratio in
parentheses is an ordinary rank-one finite
$\mathbb F_q$-local system. Its phase-$a$ value is the
absolute $p^a$-Frobenius pullback of its initial value.
The established constant prime-to-$p$ scalar bound on
the cover therefore gives at most $q^{c_Rt}$ choices
for all punctured determinant ratios together.
This does not assert that an individual twisted
determinant is an ordinary character.

For a fixed punctured determinant, put
$d_{ai}=c_1(E_a)\cdot E_i$ on the original model.
Generic generation of full pullbacks gives $d_{ai}\ge0$.
Since $C_i\ge1$, $C\le Z$, and (D.BH.4) holds,
\[
             \sum_i d_{ai}\le c_1(E_a)\cdot C
                            \le c_1(E_a)\cdot Z
                            \le c_{R,B}tr.             \tag{D.BH.6}
\]
There are $O_R(t)$ exceptional components. Thus the
number of such degree vectors for all phases is at most
\[
       \exp\bigl(c_{R,B}ft\log(r+1)+c_{R,B}ft\bigr).
                                                        \tag{D.BH.7}
\]
Two extensions of the same punctured line differ by
an integral exceptional divisor. Its intersection
vector determines it uniquely by negative definiteness.
Consequently (D.BH.5)--(D.BH.7) give actual finite lists
of determinant lines on $Y$, not merely numerical
classes in a moving Picard parameter space. Once these
lists are fixed, so are the lines
$L_a=\det(E_a(-Z))$ and all their restrictions below.
Retain only lists occurring for at least one canonical
geometric input in the class under consideration. Thus
every retained $L_a$ is actually globally generated;
no generation assertion is inferred from a numerical
degree budget alone. Taking this subset adds no count.

\emph{Whole-cycle objects and their exact equations.}
Choose the fixed precision $D=LC$, with $L$ sufficiently
large for the two effectivity results and for
$Z\subset NC\subset LC$ at the retained depth. This
$L$ is independent of $t,r,f,\alpha$, for fixed $R,B,p$.
On the proper Cohen--Macaulay curve $D$, a globally
generated rank-$R_a$ bundle, $R_a=hr$, has a generating
quotient by $R_a+1$ sections. Thus each geometric
$V_a=E_a(-Z)|_D$ is covered by
\[
       0\longrightarrow L_a^{-1}|_D
       \xrightarrow{s_a}\mathcal O_D^{R_a+1}
       \longrightarrow V_a\longrightarrow0.          \tag{D.BH.8}
\]
Use the open subbundle condition on $s_a$.
Over every coefficient algebra it is locally split
and produces an actual vector bundle; it is not a
presentation only of its geometric fibres.
The scalar and determinant budgets give
$h^0(L_a|_D)+h^1(L_a|_D)=O_{R,B}(tr)$, so (D.BH.8)
uses $O_{R,B}(tr^2)$ coordinates per phase.

Each fixed order, comparison module, and fixed tensor
word used in an action or inverse identity has bounded
rank in $B$. It has a fixed complex of bounded length
whose terms are line bundles, obtained by the same
rank-plus-one quotient construction after the controlled
twists just described. Its line degrees are $O_{R,B}(t)$.
These reference sections and matrices are constants,
not parameter variables.
Every line occurring after tensoring this fixed complex
with either determinant end of (D.BH.8) is of the form
$P\otimes Q^{-1}$, where $P,Q$ are globally generated
and have degrees $O_{R,B}(tr)$.
A regular section of $Q$ embeds this line into $P$;
the corresponding regular-section sequences and
Riemann--Roch give
\[
 h^0(PQ^{-1})+h^1(PQ^{-1})=O_{R,B}(tr).
                                                        \tag{D.BH.9}
\]
When neither determinant end occurs, the sharper bound
is $O_{R,B}(t)$. These statements hold on the full
nonreduced curve; regular means a nonzerodivisor at
every generic point, as in the proved full-bundle
degree and ordinary Hom lemmas.

There is one global action map
$E_a\otimes P_a\to E_a$ per phase. Impose its unitality
and associativity, and impose the actual balanced
correspondence maps, their linearity, and the scaled
inverse identities of the fixed Morita context.
There is no separate arbitrary endomorphism matrix
for each of a possibly large set of algebra generators.
For every required Hom complex the top-top, mixed,
and bottom-bottom blocks have respectively
\[
 \begin{array}{c|c|c}
 &\text{multiplicity}&h^0+h^1\text{ of each line factor}\\ \hline
 \text{top--top}&O_B(r^2)&O_{R,B}(t)\\
 \text{mixed}&O_B(r)&O_{R,B}(tr)\\
 \text{bottom--bottom}&O_B(1)&O_{R,B}(tr).
 \end{array}                                           \tag{D.BH.10}
\]
Associativity has a fixed tensor word in the bounded-rank
order, and still only the two variable object ends.
Composition through an intermediate object is computed
as a morphism, not by retaining a third independent
object matrix index. Hence every such complex has
total size $O_{R,B}(tr^2)$.

The fixed two-affine Cech contractions and finite Hom
construction of
\ref{D:fixed-determinant-finite-Hom}
now apply to these bounded-length locally split complexes.
The internal differential is linear in the object sections,
or of degree $p$ after one Frobenius. Perturbation is finite
because the internal widths are bounded, and all its word
lengths are bounded in $B,p$. Fixed Cech splittings are
tensored with the coefficient algebra; no varying bases
of cohomology or determinantal flattening strata are
introduced. An arrow is a degree-zero cycle, and equality
of actual maps is imposed by a degree-minus-one boundary
witness. Products and fixed pairings are transferred by
the same finite chain maps. Thus, for all $f$ phases, the
actual finite data have a polynomial presentation/open with
\[
  N_{\rm var}\le c_{R,B}ftr^2,\qquad
  \deg(\text{each equation})\le D_{R,B}.               \tag{D.BH.11}
\]
The constants are independent of $f$ as well: each equation
uses only a bounded number of adjacent phases. The final
phase is one $p$-Frobenius arrow with the fixed anchor
comparison, not a polynomial of degree $p^f$.
The scaled separability tensors and coevaluations are
fixed morphisms; the number of terms in one expression
for them adds no variable and increases no polynomial
degree.

\emph{Where simultaneous effectivity is used.}
Apply Theorem~\ref{D:finite-Morita-Frobenius-object-effectivity}
to these exact finite data. It supplies actual formal ordinary vector bundles
$E_a$ with actual $P_a$-actions and actual Morita-unit
correspondences, retaining the prescribed $NC$ marking.
For every original canonical geometric input, that input
and the constructed realization have identical marked
$NC$ data. The final comparison assertion of the same
effectivity theorem makes their punctured unit objects
isomorphic. Thus the construction covers the original
classes, not merely some class attached to each finite jet.
For Noetherian finite-type parameter bases, ordinary
Grothendieck existence over the completed tensor base
algebraizes the bundles and the finitely many structure
maps and equations. Over arbitrary coefficient rings,
only the finite-level and formal assertions are used.
On $U_A$ the orders are Azumaya. Their ordinary Morita
equivalence recovers actual twisted vector bundles and
the specified Frobenius unit structure. No root-stack
local freeness is asserted for an arbitrary boundary
order module. In particular no relative reflexive hull
or coefficient-flat boundary Morita equivalence is
required here.

The retained data contain the actual generating quotient
of $E_a(-Z)|_Z$ and its fixed determinant. The weak-cycle
vanishings lift these generators and fix the determinant
on the entire formal family, exactly as in
\ref{D:weak-cycle-generated-determinant-lifting}.
Thus the next comparison applies to every point of the
parameter open, including its nonreduced coefficient
points, not merely to the initially canonical objects.

\emph{Generated-source, fixed-determinant comparison.}
Fix one represented punctured twisted unit class. Extend
its ordinary coarse Hom modules reflexively over $A$ and
take their canonical full pullbacks $E_a^*$. These are
fixed comparison targets only. A specified punctured
unit isomorphism over a finite local Artinian coefficient
ring extends from generators to
$E_a\to E_a^*(Z)$ by ordinary module Hartogs, exactly
as in \ref{D:generated-fixed-determinant-unit-openness}.
Its $P_a$-linearity persists: the two maps obtained by
action before or after the isomorphism agree on the
puncture, and their difference has a locally free target
on which the Cartier equation is regular. They therefore
agree everywhere, also over the nonreduced base.

The global determinant line of the source is fixed.
Divide the determinant of this meromorphic map by the
fixed determinant comparison at the chosen class.
The quotient and its inverse are functions on the
puncture and extend by ordinary scalar Hartogs. They
are units. The adjugate formula therefore bounds both
directions of every specified Artinian comparison by
one fixed multiple $dC$, independent of nilpotence.
The integer $d$ may depend on the fixed comparison
object and on $r$; it is not a coordinate-budget input.

Proposition~\ref{D:generated-Morita-unit-Isom-openness}
represents this bounded
$P$-linear Morita-unit Isom functor by actual finite
equations and proves its all-Artinian comparison.
Finite etale twisted frames on the fixed coefficient
gerbe, and their specified isomorphisms, are invariant
under nilpotent thickenings. Hence the fixed-class Isom
morphism is etale and of finite presentation, and its
image is precisely that class and is open and closed.
Each irreducible component of the parameter open can
therefore represent at most one punctured unit class.
This use of the comparison theorem cannot be replaced
by a tangent-space calculation or by separately lifted
maps whose product identities were never restored.

\emph{Counting components and frames.}
The affine component bound, with all weights one, gives
for each fixed determinant list and coefficient class
\[
 (N_{\rm var}+1)\max\{2,D_{R,B}\}^{N_{\rm var}}
                         \le p^{c_{R,B}ftr^2}.          \tag{D.BH.12}
\]
Open restrictions do not increase the number of
irreducible components. Multiply this by (D.BH.2),
the scalar choices of (D.BH.5), and the degree-vector
choices (D.BH.7). Since $r\ge2$ and $q=p^f$, the
result is $q^{A_{R,B}tr^2}$ after enlarging the constant.
Finite $\mathbb F_q$ frames of the scalar-weight-one
twisted unit objects project onto the projective
representation classes with the specified obstruction,
as in the exact coefficient-gerbe comparison of
\ref{D:SNC-general-projective-quadratic-count}.
Several twisted lifts may give the same projective
object, which only improves the upper bound. No
normalizer-conjugacy identification is assumed.
This proves the assertion.
\end{proof}

\begin{remark}[The bounded-centre ordinary subfamily]
For a projective representation into $\operatorname{PGL}_r(\mathbb F_q)$,
choose lifts of its matrices. If $c(g,h)$ is the scalar multiplier,
taking determinants shows that $c(g,h)^r$ is a scalar coboundary.
The full coefficient obstruction is therefore killed by $r$, and
also by $q-1$ because that is its coefficient-group exponent.
Its geometric Brauer period divides $\gcd(r,q-1)$.
Consequently the stated theorem covers every ordinary projective
representation in the range
$\gcd(r,q-1)\le B$, with rank unbounded. This observation does
not supply a projective semilinear theorem or remove an outer
quotient generator term in a subsequent crown argument.
\end{remark}



% Separate post-CP27 proved corollary. Four complete audits passed; closure of the same bounded-period hereditary projective enlargement.
\subsection{Bounded-period projective crowns over arbitrary permutation quotients}

\begin{corollary}[Proved: arbitrary-quotient projective crowns with bounded multiplier period]
\label{D:bounded-period-arbitrary-quotient-projective-crowns}
Let $R$ and $\Gamma$ be as in
Theorem~\ref{D:bounded-period-hereditary-projective-growth}.
Fix constants
$B_*,b_0\ge1$ and $\kappa>0$.
Let $L$ be a finite monolithic primitive group with
nonabelian socle $N=S^t$, where $S$ is a finite nonabelian
simple group, and put $P=L/N$.
Fix a continuous epimorphism $\Gamma\to P$.
Choose a socle factor, and let $\mathcal B\le\operatorname{Aut}(S)$
be the image of its stabilizer in $L$. Suppose
\[
 S\le\mathcal B_0\le\mathcal B,
 \quad [\mathcal B:\mathcal B_0]\le b_0,
 \quad \iota:\mathcal B_0\hookrightarrow
                     \operatorname{PGL}_r(\mathbb F_q),
 \quad r^2\log q\le\kappa\log|S|,                    \tag{D.BC.1}
\]
where $q$ is a power of the defining characteristic $p$.
Identify $S$ with its inner automorphism group.

Let $H\subset\Gamma$ be the preimage of the point
stabilizer of that factor. The fixed outer map determines
$\eta:H\to\mathcal B/S$, and put
$H_0=\eta^{-1}(\mathcal B_0/S)$.
Require that, for every stabilizer lift $\theta:H\to\mathcal B$
of $\eta$ which occurs in the crown comparison, the
projective representation
\[
                  \iota\theta|_{H_0}:H_0\to
                                  \operatorname{PGL}_r(\mathbb F_q)
                                                        \tag{D.BC.2}
\]
has geometric Brauer period at most $B_*$ on the puncture
corresponding to $H_0$. This condition holds for all such
lifts if $\gcd(r,q-1)\le B_*$.

If the crown-based power $L_k$ is an epimorphic image
of $\Gamma$ over the chosen map to $P$, then
\[
 k\le |N|^{c_R},\qquad
 d(L_k)\le\max\{2,d(P),c_R+2\},                       \tag{D.BC.3}
\]
where one may take
\[
 c_R=\max\left\{1,
   \left\lceil\kappa(A_{R,B_*}b_0+1)+b_0-1\right\rceil\right\}.
                                                        \tag{D.BC.4}
\]
There is no tameness or etale-extension assumption on
the permutation quotient cover. The term $d(P)$ is
retained. The embedding in (D.BC.1) is an ordinary
projective embedding; a semilinear embedding does not
satisfy this hypothesis by itself.
\end{corollary}

\begin{proof}
The exact nonsplit transport proof of
\ref{D:nonsplit-automorphism-wreath-lifts}, and the
restriction argument of
\ref{D:normal-subgroup-projective-crowns}, identify the
relevant $N$-conjugacy classes of lifts over $P$ with
$S$-conjugacy classes of stabilizer lifts
$\theta:H\to\mathcal B$ having the prescribed outer map
$\eta$. We recall the points needed for this use.
The centralizer of the unique nonabelian socle is trivial,
so conjugation gives the faithful actual automorphism
wreath embedding. Choose actual coset transports between
the socle factors. Restriction to the chosen factor gives
the stabilizer lift, and changing the transports changes
it by inner $S$-conjugation. The reverse construction
uses the same actual transports and does not split the
outer extension. Two crown coordinates in the same
$N$-conjugacy class would force their pair to lie in a
diagonal graph subgroup, contradicting surjectivity onto
the crown-based power. Therefore $k$ is at most the
number of these stabilizer lift classes.

Put $a=[H:H_0]\le b_0$. The open subgroup $H$ has
index $t$ in $\Gamma$, and $H_0$ has index $at$.
Every lift maps $H_0$ into $\mathcal B_0$, by its fixed
outer map. The assumed multiplier-period condition and
the hereditary projective theorem give at most
\[
                         q^{A_{R,B_*}at r^2}            \tag{D.BC.5}
\]
ambient $\operatorname{PGL}_r(\mathbb F_q)$ conjugacy
classes of the restricted maps. Faithfulness of $\iota$
identifies an actual restricted matrix homomorphism with
at most one homomorphism to $\mathcal B_0$.
Each ambient orbit contains at most
$|\operatorname{PGL}_r(\mathbb F_q)|\le q^{r^2}$
actual homomorphisms, hence no more than that many
$S$-conjugacy classes. This uses the ambient orbit
cardinality, without any normalizer-conjugacy assertion.
Thus the restricted $S$-class count is at most
\[
       q^{(A_{R,B_*}at+1)r^2}
                   \le |S|^{\kappa(A_{R,B_*}at+1)}.     \tag{D.BC.6}
\]

For each restricted $S$-class choose an actual
representative. Conjugate a full lift by $S$ so that
its restriction is this representative. Choose $a-1$
nonidentity coset representatives for $H_0$ in $H$.
Together with $H_0$ they determine every homomorphism
from $H$. The fixed outer map places the image of
each such representative in one fixed coset of $S$
in $\mathcal B$, giving at most $|S|$ choices each.
Consequently extension costs at most $|S|^{a-1}$.
This argument needs neither normality of $H_0$ nor
normality of $\mathcal B_0$; $S$ itself is normal
in every group displayed in (D.BC.1).
Combining with (D.BC.6) gives
\[
 \#\{\theta\}/S
   \le |S|^{\kappa(A_{R,B_*}at+1)+a-1}
   \le |S|^{t\{\kappa(A_{R,B_*}b_0+1)+b_0-1\}}
   \le |N|^{c_R}.                                    \tag{D.BC.7}
\]
The exact crown generation estimate already used in
the cited transport and normal-subgroup corollaries,
from the fully audited finite-group input, now gives
the second inequality of (D.BC.3). Its hypotheses are
the nonabelian simple socle, the bound $k\le|N|^{c_R}$,
and the actual quotient $P$; no bound for $d(P)$ is
inserted.

Finally, the determinant of chosen linear lifts of
a projective representation makes its full scalar
obstruction $r$-torsion, while its coefficient group
has exponent $q-1$. Its geometric image therefore
has period dividing $\gcd(r,q-1)$, proving the stated
sufficient condition.
\end{proof}

\begin{remark}[A precise central-lift alternative]
Another sufficient condition in (D.BC.2) is the following
explicit one. Suppose the pullback to $\mathcal B_0$ of
the scalar extension
$\operatorname{GL}_r(\mathbb F_q)\to
\operatorname{PGL}_r(\mathbb F_q)$ is the pushout of
a fixed central extension
$1\to C\to\widetilde{\mathcal B}_0\to\mathcal B_0\to1$
with $\exp C\le B_*$.
Every multiplier in (D.BC.2) is then the pushforward
of a class in $H^2(H_0,C)$ and is killed by $\exp C$.
Its geometric period is at most $B_*$.
Merely naming a small abstract centre without this
actual compatible central lift would not prove the
required coefficient-class assertion.
\end{remark}

\begin{remark}[A verified ordinary type-A instance]
The natural projective model for the simple groups
$S=\operatorname{PSL}_r(\mathbb F_q)$ is among the
fully checked defining-characteristic models of
\ref{D:classical-defining-characteristic-crowns}.
Apart from the stated nonsimple small cases there,
it has the uniform order comparison with $\kappa=16$.
Thus one may take any
$S\le\mathcal B_0\le\operatorname{PGL}_r(\mathbb F_q)$
and any stabilizer image $\mathcal B\le\operatorname{Aut}(S)$
containing it with $[\mathcal B:\mathcal B_0]\le b_0$,
provided $\gcd(r,q-1)\le B_*$.
The natural inclusion is the actual faithful projective
embedding; the cover is allowed to be arbitrary.
No recognition of an unbounded field-automorphism image
inside an ordinary projective group is being asserted.
This is one concrete scope of the same bounded-period
enlargement, not a second counting reduction.
\end{remark}


% Post-CP27 proved input. Four complete actual-file and exact-source audits; genuinely twisted coefficients retained.
\subsection{Second cohomology for arbitrary field characters}

\begin{lemma}[Proved: arbitrary field-character lifting classes]
\label{D:arbitrary-field-character-H2}
Let \(R\) be a fixed complete excellent normal local surface
algebra over an algebraically closed coefficient and residue
field \(k\) of characteristic \(p>0\), and put
\(\Gamma=\pi_1(U_R)\), where
\(U_R=\operatorname{Spec}R-\{\mathfrak m_R\}\).
Let \(H\subset\Gamma\) be open of index \(t\), let
\(q=p^f\), and fix a continuous field character
\[
 \chi:H\longrightarrow\operatorname{Gal}(\mathbb F_q/\mathbb F_p).
                                                        \tag{D.WH.1}
\]
Write \(a=|\operatorname{im}\chi|\), \(\nu=f/a\),
\(q_0=p^\nu\), and \(n=q-1\).
Let \(U_H\) be the puncture corresponding to \(H\), and
let \(V=U_V\to U_H\) be the connected cyclic cover
corresponding to \(\ker\chi\). Its complete normal local
cover ring \(A_V\) has degree \(ta\) over \(R\).
Let \(\mathcal M_\chi\) be the locally constant sheaf
with fibre \(\mathbb F_q^\times\) and its prescribed
field action.

There is a deck-equivariant SNC resolution of \(A_V\)
with \(v_V=O_R(ta)\) exceptional components such that
\[
 \begin{aligned}
 |H^2_{\mathrm{cont}}(H,\mathbb F_q^\times(\chi))|
 &\le |H^2_{\mathrm{et}}(U_H,\mathcal M_\chi)|\\
 &\le(q_0+1)^{\,4p_g(A_V)+3v_V}
 \le q^{c_Rt}.
 \end{aligned}                                         \tag{D.WH.2}
\]
The constant \(c_R\) is independent of \(H,q,\chi,a\).
Neither the cover nor the field character is assumed
tame along an exceptional divisor; in particular
\(p\mid a\) is allowed.
\end{lemma}

\begin{proof}
The case \(n=1\) is immediate. Hence assume \(n>1\).
Fix \(\mathbb F_q^\times=\mu_n(k)\) once, after passing
to \(V\). Choose the geometric deck-pullback convention
\(\tau\) so that the generator of the cyclic descent
group acts on \(H^j(V,\mu_n)\) as
\[
                         \sigma=q_0(\tau^*)^{-1}.
                                                        \tag{D.WH.3}
\]
Thus its invariants are the kernel of \(q_0-\tau^*\).
The geometric automorphism \(\tau\) is \(k\)-linear
and fixes \(\mu_n(k)\). The scalar \(q_0\) in
(D.WH.3) is the coefficient action, not a geometric
action on roots of unity.

\emph{A controlled equivariant SNC model.}
Begin with the minimal regular resolution of \(A_V\).
Every deck automorphism lifts uniquely, by the
minimal-resolution argument in
Lemma~\ref{D:arbitrary-field-character-scalar-H1}.
Consequently these lifts satisfy the group identities,
even when \(p\mid a\).
The all-complete geometric budgets give
\(p_g(A_V)=O_R(ta)\) and \(v_0=O_R(ta)\) for the
number of components on this model.
If \(A_V\) is regular, first blow up its closed point;
this is a single invariant centre and gives a
nonempty proper exceptional curve.

Let \(E_0\) denote the resulting reduced exceptional
curve, and put \(G_0=h^1(E_0,\mathcal O_{E_0})\).
The ideal sequence and coherent \(H^2=0\) give
\(G_0\le p_g(A_V)\). Apply the blowup procedure of
Lemma~\ref{D:linear-SNC-refinement-centres}, at every
step blowing up the entire finite set of non-SNC
points. This set is deck-stable. Its distinct centres
are disjoint, so simultaneous blowup is the same
sequence of individual point blowups in any order.
The proof of that lemma bounds every finite prefix
whose centres are non-SNC, by
\(2\delta_0+7G_0\). It therefore bounds this equivariant
procedure as well, without a factor \(a\).
The resulting SNC model \(Y\) is deck-equivariant,
has \(v=v_V=O_R(ta)\) components and \(O_R(ta)\)
nodes, and has the same geometric genus \(p_g(A_V)\).
No arbitrary pencil or additional orbit of regular
boundary points is used.

Write \(E=\bigcup_{i=1}^v E_i\), \(g_i=g(E_i)\),
and let \(G\) be its connected dual graph, with
\(b_1=b_1(G)\). Put
\(\Lambda=\mathbb Z^{\{E_i\}}\) and
\(N_E=(-(E_i\cdot E_j))_{ij}\).
The deck action on \(\Lambda\) is a permutation
and commutes with \(N_E\).
The exact support and Picard calculation in
Lemma~\ref{D:SNC-general-Brauer-geometry} is natural
for isomorphisms of the regular pair. It gives
the following two equivariant exact sequences:
\[
 \begin{aligned}
 0&\longrightarrow Q_n\longrightarrow H^2(V,\mu_n)
             \longrightarrow B_n\longrightarrow0,\\
 Q_n&=\operatorname{coker}(N_E:\Lambda/n\Lambda
                                      \to\Lambda/n\Lambda),\\
 0&\longrightarrow\prod_i\operatorname{Jac}(E_i)[n](k)
       \longrightarrow B_n
       \longrightarrow H_1(G,\mathbb Z)\otimes\mathbb Z/n
       \longrightarrow0.
 \end{aligned}                                         \tag{D.WH.4}
\]
Here \(B_n=\operatorname{Br}'(V)[n]\).
The constant Tate twist on the last term has been
suppressed using the fixed roots of unity.
This calculation uses proper strict-local base change,
the Kummer sequence, and the divisor/point support
complex of the regular SNC pair. The exact purity
input is \citet[Expose XVI, Theorem 3.1.1 and the
complete proof of Corollary 3.1.3]{A_ILO2014}.
Its Gysin maps and signed incidence map are natural;
no deck-invariant spanning tree is required.
The Jacobian identifications are those of
\citet[Lemma 59.69.1, Tag 03RQ]{StacksLocalPairCohomologyB}.
In particular (D.WH.4) is an equivariant filtration,
not a choice of a splitting of the Brauer sequence.

\emph{Kernel estimates on the equivariant pieces.}
Write \(u=q_0-\tau^*\). For a lattice \(L\) of rank
\(d\) with finite-order automorphism \(T\), Smith
normal form gives
\[
 |\ker(q_0-T:L/nL\to L/nL)|
       \le|\det(q_0-T)|\le(q_0+1)^d.                   \tag{D.WH.5}
\]
Indeed the Smith invariant factors bound the kernel
modulo \(n\), and every complex eigenvalue of \(T\)
is a root of unity.
There is one quotient in (D.WH.4) for which
left exactness alone would not suffice. Set
\(P_n=\Lambda/n\Lambda\) and
\(I_n=\operatorname{im}(N_E\bmod n)\).
Both are \(u\)-stable. The kernel--cokernel sequence
of \(0\to I_n\to P_n\to Q_n\to0\), together with
\(|\operatorname{coker}(u|I_n)|=|\ker(u|I_n)|\), gives
\[
 |\ker(u|Q_n)|
 \le|\ker(u|P_n)|\,|\operatorname{coker}(u|I_n)|
 \le|\ker(u|P_n)|^2
 \le(q_0+1)^{2v}.                                     \tag{D.WH.6}
\]
Thus no estimate for the determinant of the
intersection matrix is introduced.

Let \(J=\prod_i\operatorname{Jac}(E_i)\).
Its deck automorphism has finite order dividing \(a\).
The endomorphism \([q_0]-\tau^*\) is an isogeny:
its product with
\([q_0]^{a-1}+[q_0]^{a-2}\tau^*+\cdots+(\tau^*)^{a-1}\)
is \([q_0^a-1]=[n]\).
The characteristic-polynomial and degree formula give
\[
 |\ker(u|J[n](k))|
 \le\deg([q_0]-\tau^*|J)
 \le(q_0+1)^{2\sum_i g_i}.                            \tag{D.WH.7}
\]
We use \citet[Theorem 10.9 and its written proof,
and Proposition 10.20 with Lemmas 10.21--10.22,
printed pp.46--52]{Milne2008AVScalarB}.
The eigenvalue estimate is on the characteristic-zero
Tate space. No semisimplicity in characteristic
dividing \(a\), or division by \(a\), is used.

For finite \(u\)-modules in an exact sequence
\(0\to A_1\to A\to A_2\to0\), one always has
\(|\ker(u|A)|\le|\ker(u|A_1)|\,|\ker(u|A_2)|\).
Applying this to (D.WH.4), and applying (D.WH.5)
to the integral graph lattice, yields
\[
 |\ker(u:H^2(V,\mu_n)\to H^2(V,\mu_n))|
       \le(q_0+1)^{2v+2\sum_i g_i+b_1}.               \tag{D.WH.8}
\]

We also need a bound on a degree-one cokernel.
The equivariant localization and normalization
filtration used in the proof of
Lemma~\ref{D:arbitrary-field-character-scalar-H1} is
\[
 \begin{gathered}
 0\longrightarrow H^1(E,\mu_n)
 \longrightarrow H^1(V,\mu_n)\longrightarrow J_n'
 \longrightarrow0,\qquad J_n'\subset\Lambda/n\Lambda,\\
 0\longrightarrow H^1(G,\mathbb Z)\otimes\mathbb Z/n
 \longrightarrow H^1(E,\mu_n)\longrightarrow J[n](k)
 \longrightarrow0.
 \end{gathered}                                        \tag{D.WH.9}
\]
For the first line one removes the finite invariant
node set, uses codimension-two purity for \(H^1\),
and then divisor purity for the disjoint smooth
component complements. The last term is a subgroup,
so its \(u\)-kernel is bounded directly by that on
\(\Lambda/n\Lambda\).
Since all these groups are finite, the same kernel
and cokernel orders give
\[
 |\operatorname{coker}(u:H^1(V,\mu_n)\to H^1(V,\mu_n))|
 \le(q_0+1)^{v+2\sum_i g_i+b_1}
 \le(q_0+1)^{v+2p_g(A_V)}.                            \tag{D.WH.10}
\]
The last inequality follows from the normalization
sequence:
\(h^1(E,\mathcal O_E)=\sum_i g_i+b_1\le p_g(A_V)\).
It also shows \(2\sum_i g_i+b_1\le2p_g(A_V)\).

\emph{The genuinely twisted descent calculation.}
Apply Cartan--Leray to the cyclic finite etale cover
\(V\to U_H\), with the sheaf \(\mathcal M_\chi\).
Its exact scheme scope and written proof are
\citet[Theorem 14.9 and Remark 14.10, pp.96--97]
{MilneLEC2013H2D}. The spectral sequence is
\[
 H^i(C_a,H^j(V,\mu_n))\Longrightarrow
                          H^{i+j}(U_H,\mathcal M_\chi),
                                                        \tag{D.WH.11}
\]
where the action is (D.WH.3), not just geometric
pullback. In degree zero the coefficient module is
\(\mathbb F_q^\times\) with action \(q_0\).
For the cyclic resolution, its norm is the power
\(n/(q_0-1)\). The norm image has order \(q_0-1\)
and equals the fixed subgroup. Therefore
\[
                       H^2(C_a,\mathbb F_q^\times)=0.
                                                        \tag{D.WH.12}
\]
This computation remains valid when \(p\mid a\).

For any finite \(C_a\)-module \(A\), the cyclic
resolution identifies
\(H^1(C_a,A)=\ker(\mathrm{Norm})/(\sigma-1)A\),
a subgroup of \(A/(\sigma-1)A\).
Since \(\sigma-1=(\tau^*)^{-1}u\), (D.WH.10)
bounds the total-degree-two term with \((i,j)=(1,1)\).
The \((0,2)\) term is the kernel bounded in
(D.WH.8), and (D.WH.12) makes the \((2,0)\)
term zero. Every limiting term is a subquotient
of its second-page term, so no vanishing of the
remaining differentials is assumed. Multiplying
the two bounds gives
\[
 |H^2(U_H,\mathcal M_\chi)|
 \le(q_0+1)^{\,2v+2\sum_i g_i+b_1+v+2p_g(A_V)}
 \le(q_0+1)^{\,4p_g(A_V)+3v}.                         \tag{D.WH.13}
\]
Since \(p_g(A_V)+v=O_R(ta)\), \(q_0+1\le q_0^2\),
and \(q_0^a=q\), this is at most \(q^{c_Rt}\).
In particular the degree \(a\) of the field cover
has been absorbed by the smaller Frobenius exponent.

Finally the degree-two comparison
\[
 H^2_{\mathrm{cont}}(H,\mathbb F_q^\times(\chi))
       \lhook\joinrel\longrightarrow
                  H^2_{\mathrm{et}}(U_H,\mathcal M_\chi)
                                                        \tag{D.WH.14}
\]
is the same Cartan--Leray injection as in the proof
of Lemma~\ref{D:radial-tame-cyclic-H2}; that argument
for the comparison itself does not use tameness.
Pass first to a finite cover making the coefficient
sheaf constant. Every degree-one class on a finite
stage is a finite torsor and dies after a further
finite etale cover, followed by a Galois closure.
Thus the universal pro-cover has no degree-one
cohomology with these coefficients. The five-term
sequence, or its degree-two filtration, gives
(D.WH.14). This proves all of (D.WH.2).
\end{proof}

\begin{remark}[Period and phase scope]
For a prescribed field character, the scalar obstruction
belongs to the twisted group
\(H^2(H,\mathbb F_q^\times(\chi))\).
Its pullback to \(V\) has constant band \(\mu_n\)
and hence an ordinary geometric Brauer image.
A bounded-geometric-period hypothesis for the
semilinear construction refers to this latter image.
No ordinary \(\mathbf G_m\)-Brauer image of the
nonconstant band on \(U_H\) is assumed.
The lemma counts the full twisted coefficient classes.
It is a cohomological input; the actual short-phase
Morita atlas and its specified-Isom comparison are
separate requirements for a semilinear projective count.
\end{remark}


% Proved after four complete actual-file audits and exact H2 dependency-discharge refreshes. Coherent band and deck data retained.
\subsection{Bounded-period projective growth with arbitrary field monodromy}

\begin{lemma}[Proved: coherent skew Morita--Frobenius effectivity]
\label{D:skew-Morita-Frobenius-effectivity}
Let $U_V\to U_H$ be the cyclic field cover attached to a
continuous character
$\chi:\pi_1(U_H)\to\operatorname{Gal}(\mathbb F_q/\mathbb F_p)$,
where $q=p^f$, $a=|\operatorname{im}\chi|$, $\nu=f/a$ and
$q_0=p^\nu$. Let $\mathcal G_H$ be an actual lifting gerbe
with band $\mathbb F_q^\times(\chi)$ on $U_H$, and let
$\mathcal G_V$ be its pullback, with the resulting constant
band $\mu_{q-1}$. Write $\mathscr T$ for its coherent deck
pullback, oriented to carry scalar weight one to scalar
weight $q_0$. It covers a deck automorphism $\tau$ of $U_V$,
and its given coherence includes $\mathscr T^a\simeq1$.

Use the deck-equivariant comparable-cycle model $Y$ of
Lemma~\ref{D:equivariant-cover-comparable-cycles}. Its cycles
$C,Z$ and its actual separating pencil $x,y$ are invariant.
Choose rank-$h$ anchors $\mathcal A_j$ of scalar weight $p^j$
for $0\le j<\nu$, with $p\nmid h$, and put
$\mathcal A_\nu=\mathscr T\mathcal A_0$. Define full orders
and comparison bundles as in (D.MF.1), using this terminal
anchor. Thus
\[
 P_\nu=\tau^*P_0,\qquad
 K_j=\operatorname{full}j_*\mathcal Hom
              (\mathcal A_{j+1},F^*\mathcal A_j),\qquad
 R_j=\operatorname{full}j_*\mathcal Hom
              (F^*\mathcal A_j,\mathcal A_{j+1})(-pC).
                                                        \tag{D.SM.1}
\]
Here all scalar-weight-zero sheaves descend before their
full pullback, and $F$ is absolute $p$-Frobenius.

Take $b,N,L$ exactly as in (D.MF.2). On $LC_T$ over any
coefficient algebra $T$, take the same finite data and
actual action, balance, intertwining and two right-inverse
equations (D.MF.3), (D.MF.6), for $0\le j<\nu$, except that
\[
             E_\nu=\tau^*E_0,\qquad
 \mu_{\nu-1}:F^*E_{\nu-1}\otimes K_{\nu-1}
                                     \longrightarrow\tau^*E_0 .
                                                        \tag{D.SM.2}
\]
These data have actual formal realizations with their
specified $NC_T$ marking and with genuine skew Morita
unit maps on the puncture. Two realizations with the same
marked $NC_T$ data give isomorphic punctured unit objects.
For finite-type Noetherian or finite local Artinian
coefficient bases, the ordinary bundles and their structure
maps algebraize over the same completed tensor base as in
Theorem~\ref{D:finite-Morita-Frobenius-object-effectivity}.

For actual families with every $E_j(-Z)$ generated and
every determinant fixed formally, the analogue of
Proposition~\ref{D:generated-Morita-unit-Isom-openness}
holds: the fixed-object meromorphic unit-Isom functor has
finite presentation and is etale, and its class image is
open and closed. The punctured objects are exactly the
twisted projective-semilinear frame objects on the specified
coefficient gerbe. No gerbe extension to $Y$, root-lattice
flatness, averaging by the order $a$, or order-$a$ matrix
equation is assumed.
\end{lemma}

\begin{proof}
We check the changes to the already proved simultaneous
construction, including its categorical comparison.
The actual banded gerbe is pulled back from $U_H$, so
its deck functors and all their coherence maps are given
by descent; they are not selected from a geometric Brauer
class alone. Pullback commutes with absolute Frobenius.
The terminal scalar weight is $p^\nu=q_0$, so every Hom
in (D.SM.1) has weight zero, as required.
In particular
\[
 \mathcal Hom(\mathscr T\mathcal A_0,\mathscr T\mathcal F_0)
       =\tau^*\mathcal Hom(\mathcal A_0,\mathcal F_0)
                                                        \tag{D.SM.3}
\]
after descent to $U_V$. Full pullback commutes with the
automorphism of the same regular model. This proves the
typing of the terminal order and module; no extra
Frobenius operation occurs in $\tau^*$.

The source-only Frobenius comparison (D.MF.4) and the
actual Morita conductor hold for every consecutive pair
of the anchors in (D.SM.1), including the terminal pair.
The opposite $-pC$ twist gives the same conductor
$bC=(p+1)C$. The global tensors for $x^b,y^b$ are actual
Noether-factorization tensors, exactly as in (D.MF.5).
The deck group fixes $x,y$, since these functions come
from the fixed base surface. There is thus no change to
the two unlocalized inverse equations on $LC_T$.

Choose a finite affine formal trivializing cover closed
under $\tau$; its size is not a parameter count. The
finite-cycle frames, affine nilpotent-thickening lifts,
and arbitrary-coefficient Cartier regularity used in the
proof of Theorem~\ref{D:finite-Morita-Frobenius-object-effectivity}
apply on this refinement. Use the corrected actions at
the initial stage and their deck pullbacks at the terminal
stage, rather than independently correcting a duplicate
terminal action. The two scaled-separability exponents
are still $2$ and $2p$. Action Newton retains $L-2$,
the two commuting averages retain
$M=L-2p-4=N+b$, and the exact inverse normalization
changes the opposite map only modulo $NC_T$. These are
the same actual algebra identities, with the last target
replaced by its deck pullback. The induced balanced
source and target have equal ordinary rank $hr$ on
$U_V$, so the right inverse is again an actual Morita
isomorphism.

For $j<\nu-1$ the overlap operator is (D.MF.9).
At the last stage it is, with the corresponding affine
indices understood,
\[
 \mathcal T(g)_0=(\tau^{-1})^*
 \left(\mu_{\nu-1,i}
       (F^*g_{\nu-1}\otimes1)\mu_{\nu-1,j}^{-1}\right).
                                                        \tag{D.SM.4}
\]
It preserves the exact order-intertwining relation and
is functorial for products and inverses. Since $\tau$
preserves the ideal of $C$, every error of order $u$ is
sent to order at least $pu-b$, including at (D.SM.4).
The initial overlap defect is still $M-b=N$; this
uses the unchanged forward map, even though the opposite
map was corrected only to precision $N$.
The same strict inequalities in (D.MF.2) give convergence,
unique overlap limits, actual cocycles, and comparison
of two realizations with identical $NC_T$ data. The
proof uses no averaging over the deck group, so it
also covers $p\mid a$.
All regular tensor equations glue by Cartier
torsion-freeness. Formal effectivity over arbitrary
coefficient algebras and Noetherian Grothendieck existence
therefore have exactly their previous distinct scopes.
Azumaya Morita is used only on $U_V$.

For clarity, the resulting twisted diagrams are
\[
 F^*\mathcal F_j\xrightarrow{\sim}\mathcal F_{j+1}
       \ (j<\nu-1),\qquad
 F^*\mathcal F_{\nu-1}\xrightarrow{\sim}
                                      \mathscr T\mathcal F_0 .
                                                        \tag{D.SM.5}
\]
Their composite is a map
$F_{q_0}^*\mathcal F_0\to\mathscr T\mathcal F_0$.
For weight-zero sheaves this is
$F_{q_0}^*L\to\tau^*L$, equivalently
$(\tau^{-1})^*F_{q_0}^*L\to L$. Thus the generator in
the ordinary skew-frame convention of
Lemma~\ref{D:twisted-deck-Frobenius-semilinear-model}
is $\tau^{-1}$ with our terminal-target convention.
Iterate it $a$ times, using the \emph{given} identification
$\mathscr T^a\simeq1$. This defines an $F_q$-unit map.
Its finite etale $\mathbb F_q$-frame torsor inherits the
semilinear deck transport, and the $a$-fold deck identity
holds on its fixed frames. Conversely finite semilinear
frames on the actual lifting gerbe give (D.SM.5).
This is precisely the etale-atlas categorical argument
at the end of Lemma~\ref{D:compressed-skew-Frobenius-unit-jets}
and Lemma~\ref{D:twisted-deck-Frobenius-semilinear-model};
that argument uses the punctured coefficient gerbe and
its coherent descent, not its extension over a root
stack or any tameness assumption. Morphisms and
coefficient base changes are respected.
The $a$-fold iteration is used for this identification
only, never as an expanded polynomial equation.
An arbitrary pointed lift of the deck generator need
not have order $a$.

Finally consider comparison of actual families.
The global averaging operators and the finite $S_x,S_y$
equations in
Proposition~\ref{D:generated-Morita-unit-Isom-openness}
still exist. At the terminal stage their target is
$\tau^*$ of the initial Hom bundle, or equivalently
one applies $(\tau^{-1})^*$ afterward. Both pencil
elements and $C$ are invariant, so their loss remains
$d=2p+b$; a meromorphic Hom lattice shifted by $\ell C$
has loss $d+(p-1)\ell$ as before. Consequently their
intrinsic finite fixed-point equations, paired inverse
equations and proper Hom finite-presentation argument
are unchanged.

The target-Hartogs and determinant-adjugate comparison
uses only the generated source, fixed determinant,
and canonical fixed target at each of the $\nu$ stages.
It therefore supplies the same reference-dependent
pole bound for all specified Artinian isomorphisms,
independent of nilpotence. The bound is qualitative,
not an input to the parameter dimension. On the
coefficient gerbe finite etale frames and specified
isomorphisms lift uniquely across nilpotent thickenings.
Their deck compatibility follows from uniqueness.
Thus the represented comparison is formally etale
and of finite presentation. Its image is open, and
its complement is the union of the other open class
images. This proves the stated clopen comparison.
\end{proof}

\begin{theorem}[Proved: bounded-period hereditary projective semilinear growth]
\label{D:bounded-period-semilinear-hereditary-projective-growth}
Let $(R,\mathfrak m,k)$ be a fixed complete excellent
normal local surface over an algebraically closed
coefficient and residue field of characteristic $p>0$,
and put $\Gamma=\pi_1(\operatorname{Spec}R\setminus\{\mathfrak m\})$.
Fix $B\ge1$.
Then there is $A_{R,B}^{\rm ps}$ such that, for every
open subgroup $H\subset\Gamma$ of index $t$, every
$q=p^f$, every $r\ge1$, and every fixed field character
$\chi:H\to\operatorname{Gal}(\mathbb F_q/\mathbb F_p)$,
the homomorphisms
\[
 H\longrightarrow\operatorname{PGL}_r(\mathbb F_q)
                 \rtimes\operatorname{Gal}(\mathbb F_q/\mathbb F_p)
       \quad\text{with field projection }\chi
                                                        \tag{D.SM.6}
\]
have at most $q^{A_{R,B}^{\rm ps}tr^2}$ projective
conjugacy classes in the following specified range:
the geometric Brauer obstruction after restriction
to $K=\ker\chi$ has order at most $B$.
There is no tameness assumption on either cover and
no bound on the field-image order. The range includes
every representation (D.SM.6) if $\gcd(r,q-1)\le B$.
Unbounded geometric period is not asserted.
\end{theorem}

\begin{proof}
The case $r=1$ is immediate, and $q=2$ is the ordinary
hereditary count. Otherwise put $n=q-1$,
$a=|\operatorname{im}\chi|$, $\nu=f/a$, and $q_0=p^\nu$.
The field cover $U_V\to U_H$ has degree $a$ and total
degree $ta$ over the original puncture.

The exact sequence
$1\to\mathbb F_q^\times\to\operatorname{GL}_r(\mathbb F_q)
\rtimes\operatorname{Gal}(\mathbb F_q/\mathbb F_p)
\to\operatorname{PGL}_r(\mathbb F_q)
\rtimes\operatorname{Gal}(\mathbb F_q/\mathbb F_p)\to1$
associates to (D.SM.6) a class with its actual coefficient
action in $H^2(H,\mathbb F_q^\times(\chi))$.
The group-to-etale comparison is injective.
Lemma~\ref{D:arbitrary-field-character-H2} therefore
bounds the number of realized actual coefficient classes by
\[
                 |H^2(U_H,\mathbb F_q^\times(\chi))|
                                  \le q^{c_Rt}.        \tag{D.SM.7}
\]
Fix such a class and its actual lifting gerbe.
It is not replaced by a gerbe for its geometric image.

On the field cover its band becomes $\mu_n$.
If its geometric Brauer period is $M\le B$, the rank
$h=M^2$ anchor construction used in
Theorem~\ref{D:bounded-period-hereditary-projective-growth}
gives rank-$h$ anchors for every invertible scalar
weight. Choose weights $p^j$, $0\le j<\nu$, and take
the terminal anchor to be $\mathscr T\mathcal A_0$.
Only the existence on the puncture is used.
In particular $p\nmid h$. The anchors need not have
a finite-field unit structure or a separately chosen
deck linearization.

Use the equivariant comparable-cycle model on the
degree-$ta$ cover. On the same original model the
canonical full orders $P_j$, comparison bundles
$K_j,R_j$, and modules $E_j$ of geometric inputs have
the actual actions and conductors of (D.SM.1)--(D.SM.2).
Their ranks are $h^2,h^2,hr$, respectively.
The full-bundle degree and generation theorem gives
\[
 E_j(-Z)\ \text{generated},\qquad
 0\le c_1(E_j)\cdot Z=O_{R,B}(ta r),\qquad
 h^0(\mathcal O_{\ell C})+h^1(\mathcal O_{\ell C})
                                      =O_{R,\ell}(ta).
                                                        \tag{D.SM.8}
\]
Fixed order and comparison terms have degree
$O_{R,B}(ta)$. These bounds concern the ordinary
full bundles on this model; there is no root-node
rectangle or intermediate-lattice generation argument.

\emph{Determinants before passage to the field cover.}
For a nonempty fixed coefficient/rank class choose
one twisted skew-unit reference $\mathcal F^*$.
With the same anchors the punctured ratios satisfy
\[
 \frac{\det E_j}{\det E_j^*}
   =\left(\frac{\det\mathcal F_j}
                    {\det\mathcal F_j^*}\right)^h .
                                                        \tag{D.SM.9}
\]
The ratios in parentheses have scalar weight zero.
Their $\nu$ Frobenius arrows and terminal coherent
deck arrow give one rank-one semilinear local system
on $U_H$. Its class lies in
$H^1(H,\mathbb F_q^\times(\chi))$.
Lemma~\ref{D:arbitrary-field-character-scalar-H1}
bounds all initial choices by $q^{c_Rt}$.
Each of the remaining phases is determined by one
absolute Frobenius pullback. No enumeration of
ordinary characters on the larger degree-$ta$ cover
is substituted for this crossed-character bound.

There are $O_R(ta)$ exceptional components. The
nonnegative multidegrees of each canonical full $E_j$
have sum $O_{R,B}(ta r)$ by (D.SM.8) and $C\le Z$.
For a fixed punctured line its extension is uniquely
determined by this vector, by negative definiteness.
The ordinary determinant enumeration thus gives at most
\[
 q^{c_Rt}
    \exp\!\left(c_{R,B}\nu ta\{\log(r+1)+1\}\right)
       =q^{c_Rt}
    \exp\!\left(c_{R,B}ft\{\log(r+1)+1\}\right)          \tag{D.SM.10}
\]
lists of actual global determinant lines.
Retain only lists occurring for a canonical geometric
input, so the corresponding determinant lines
$L_j=\det(E_j(-Z))$ are actually generated.
The terminal list is $\tau^*L_0$ and is not chosen
independently. The fixed anchor comparison contributes
no further determinant-local-system choice.

\emph{The finite polynomial model.}
Choose $N,L$ as in the preceding lemma, independent
of $a,t,r,f$ for the fixed original geometry and $B$,
and put $D=LC$. For each retained determinant list use
the actual quotient presentations
\[
 0\longrightarrow L_j^{-1}|_D\longrightarrow
       \mathcal O_D^{hr+1}\longrightarrow V_j\longrightarrow0,
       \qquad E_j=V_j(Z).
                                                        \tag{D.SM.11}
\]
The open subbundle condition gives locally split
sequences over every coefficient algebra.
Introduce one actual order action per stage and the
balanced forward/opposite Morita maps, and impose
their action, associativity, balance, and two scaled
right-inverse equations. All maps have the precise
typing in the preceding lemma.

The full fixed-reference quotient/Hom proof of
Theorem~\ref{D:bounded-period-hereditary-projective-growth}
applies with cover degree $ta$. Its top--top, mixed,
and bottom--bottom blocks have multiplicities
$O_B(r^2),O_B(r),O_B(1)$ and corresponding line
cohomology budgets $O_{R,B}(ta),O_{R,B}(ta r),
O_{R,B}(ta r)$. Fixed reference tensor words have
bounded rank and bounded internal width in $B$.
Only two variable object ends occur, including in
associativity and composition identities. Consequently
each stage has $O_{R,B}(ta r^2)$ object and equation
coordinates and bounded polynomial degrees.

At the terminal stage the target quotient, determinant,
order and reference maps are their actual $\tau$
pullbacks. The parameter coefficients are the same
as at stage zero; the only nonlinear coefficient
operation is the adjacent $p$-Frobenius.
The pulled-back fixed line is generated, with the
same cohomology bounds because $C,Z$ are invariant.
One computes its Hom complexes on any fixed two-affine
cover, treating the pulled-back lines and maps as fixed
reference data. A deck-stable trivializing refinement
is used only for qualitative effectivity and is not
multiplied into this coordinate count.
The all-ring finite Hom and boundary-witness construction
therefore yields one polynomial zero scheme and open
subbundle locus with
\[
 N_{\rm var}\le c_{R,B}\nu ta r^2=c_{R,B}ft r^2,
 \qquad
 \deg(\text{each equation})\le D_{R,B}.
                                                        \tag{D.SM.12}
\]
Each equation uses a bounded number of adjacent phases.
Neither a degree-$q$ polynomial nor an expanded
$a$-fold deck word appears.

\emph{Effectivity, class comparison, and coverage.}
Lemma~\ref{D:skew-Morita-Frobenius-effectivity} turns
every point of this finite model into an actual formal
skew unit object with the prescribed $NC$ marking.
Noetherian completed-base existence algebraizes it
where required. Its retained quotient on $Z\subset NC$
lifts globally, and its determinant remains the fixed
line by the weak-cycle generator/Picard lemma.
This holds for the whole parameter family, not only
for its initially canonical geometric points.

Every original projective-semilinear representation
occurs. Its canonical full input supplies one retained
determinant list and the finite data (D.SM.11).
That genuine input and the constructed realization
have the same $NC$ marking. The comparison in the
preceding lemma identifies their punctured skew
unit objects, so finite reconstruction does not
change the represented class.

For a fixed represented class, use its canonical
full coarse modules as fixed comparison targets.
The generated-source Hartogs extension, fixed
determinant ratio and adjugate bound give a uniform
pole bound for every specified Artinian comparison,
independent of nilpotence. The bound may depend on
this class and rank, which is harmless for qualitative
finite presentation. The preceding lemma's global
finite $S_x,S_y$ equations and paired inverses
represent these comparisons; unique nilpotent
lifting of finite etale frames preserves the actual
coefficient-band and deck coherence.
The fixed-object Isom map is therefore etale of
finite presentation, with clopen class image.
Each irreducible component of the parameter open
contains at most one such class.

The component bound applied to (D.SM.12) is
$(N_{\rm var}+1)\max\{2,D_{R,B}\}^{N_{\rm var}}
\le p^{c_{R,B}ft r^2}$.
Multiplying by the coefficient classes (D.SM.7)
and determinant lists (D.SM.10) gives
$q^{A_{R,B}^{\rm ps}tr^2}$.
Fixed frames on the actual lifting gerbe project
onto the projective-semilinear representation classes
with that coefficient obstruction. Extra choices
of a twisted lift only increase the set being counted,
so they do not invalidate this upper bound.

Finally restriction to $K$ is an ordinary projective
representation into $\operatorname{PGL}_r(\mathbb F_q)$.
Its scalar multiplier is killed by both $r$ (take
determinants of lifts) and $q-1$ (the coefficient
exponent). Its geometric period therefore divides
$\gcd(r,q-1)$, proving the sufficient condition.
The construction includes $a=1$ and $p\mid a$.
\end{proof}

\begin{remark}[Exact remaining condition]
The period bound is imposed on the ordinary geometric
Brauer class on the field cover. It is not a claimed
ordinary Brauer class for the nonconstant band on $U_H$,
nor a bound on the full coefficient obstruction.
The twisted $H^2$ input is supplied by
Lemma~\ref{D:arbitrary-field-character-H2} in its proved
arbitrary-cover scope. Neither result here asserts
unbounded-period projective growth.
\end{remark}


% Proved after four complete actual-file audits. The original detecting curve is fixed once; all cover-dependent costs are displayed.
\subsection{Hereditary projective growth from one fixed detecting curve}

\begin{lemma}[Proved: a numerical bound for fixed local projective types on a curve]
\label{D:curve-fixed-local-growth-numerical}
There is a constant \(c_p\), depending only on the positive
characteristic \(p\), with the following property.
Let \(X\) be a smooth proper connected curve of genus \(g\)
over an algebraically closed field of characteristic \(p\),
let \(P\ne\varnothing\) be a finite reduced divisor, and put
\(X^\circ=X-P\).
For \(q=p^f\), fix full local projective representation
classes \(\sigma_x\) in \(\operatorname{PGL}_r(\mathbb F_q)\)
at the points \(x\in P\).
Suppose that the finite Galois splitting extension of
\(\sigma_x\) has normalized different at most \(d_x\ge0\):
its different exponent divided by its ramification index
is at most \(d_x\). Use zero for the trivial extension.
Then the number of global projective conjugacy classes
with these full local types is at most
\[
 q^{\,c_p\left(1+g+|P|+\sum_{x\in P}d_x\right)r^2}.
                                                        \tag{D.GH.1}
\]
The constant is independent of the curve, the local extensions,
their degrees, and the field containing their coefficients.
\end{lemma}

\begin{proof}
We make quantitative the actual atlas and local comparison
in Theorem~\ref{D:curve-fixed-local-linear-growth}.
The affine curve \(X^\circ\) has
\(H^2(X^\circ,\mu_{q-1})=0\), so every projective
representation lifts to \(\operatorname{GL}_r(\mathbb F_q)\)
over the same field, exactly as in
Corollary~\ref{D:curve-fixed-local-projective-growth}.
If the counted fibre is nonempty, one such global lift
supplies one lift of each local type. At each point all
other linear lifts differ by a character of
\(k((t_x))^\times/k((t_x))^{\times(q-1)}\); there are
exactly \(q-1\) such characters. For \(q=2\) there is no
scalar choice.

We need a normalized-different assertion, rather than an
unjustified inference about a highest upper break.
Let \(L/K\) and \(\widetilde L/K\) be the splitting
extensions of a projective local type and any linear lift.
Then \(L\subset\widetilde L\), and the group of
\(\widetilde L/L\) is a subgroup of
\(\mathbb F_q^\times\). It is therefore tame.
All these complete fields have the same algebraically
closed residue field. If \(e=e(L/K)\) and
\(\widetilde e=e(\widetilde L/L)\), transitivity of the
different gives
\[
 \frac{\delta(\widetilde L/K)}{e\widetilde e}
  =\frac{\delta(L/K)}e+
       \frac{\widetilde e-1}{e\widetilde e}
  \le d_x+1.                                           \tag{D.GH.2}
\]
In this equal-characteristic scope the transitivity formula
also follows immediately from the chain rule: for
\(t=F(s)\), \(s=G(w)\), the valuation of
\((F\circ G)'(w)=F'(G(w))G'(w)\) is the sum of the
two different valuations with the indicated ramification
factor. The derivative formula for the different is
\citet[Lemmas 49.12.1--3 and Example 49.12.5,
Tag 0BWB]{StacksDifferentLatticeB}.

Fix one tuple of linear local lifts. Put
\[
 c_x=\left\lceil p(d_x+1)\right\rceil,\qquad
 B=\sum_{x\in P}c_xx,\qquad
 \Delta=\deg B,\qquad
 A=1+g+\Delta+|P| .
                                                        \tag{D.GH.3}
\]
Lemma~\ref{D:different-Frobenius-invariant-lattice}
gives the canonical invariant lattices for all coefficient
phases, with forward Frobenius maps holomorphic and their
inverse maps of pole at most \(B\).
Thus the canonical bundles \(E_j\), \(j\in\mathbb Z/f\),
on \(X\) have actual arrows
\[
 F_X^*E_j\longrightarrow E_{j+1},
 \qquad E_{j+1}(-B)\longrightarrow F_X^*E_j,
                                                        \tag{D.GH.4}
\]
with both inverse products the Cartier inclusion.
The slope argument of (D.CL.7), with \(B\) in place of
\(cD\), gives
\[
       -\frac{\Delta}{p-1}\le\mu_{\min}(E_j)
                  \le\mu_{\max}(E_j)\le0.               \tag{D.GH.5}
\]
This argument uses the two actual injections and their
duals; it assumes no tensor-product semistability.

Here are choices of all presentation degrees.
Choose a line bundle \(H\) of degree
\[
 h=2g+4+\left\lceil\frac{\Delta}{p-1}\right\rceil
       =O_p(A).
                                                        \tag{D.GH.6}
\]
Such a line exists, for example as a divisor supported
at one rational point. The bundle \(E_j^+=E_j\otimes H\)
has minimum slope at least \(2g+4\), and hence has
vanishing \(H^1\) and is generated. A general
\((r+1)\)-tuple of sections generates it, with line
kernel \(L_j^{-1}\), where \(L_j=\det E_j^+\).
The degree bounds are
\[
       r(2g+4)\le\deg L_j\le rh.
                                                        \tag{D.GH.7}
\]
Consequently both \(L_j\) and \(L_j^{-1}H^{3r}\)
are generated and have vanishing \(H^1\).
The fixed partial-resolution construction used in
(D.CL.8) therefore applies with
\[
 P_0=\mathcal O_X^{r+1},\qquad
 P_i=(H^{-3ir})^{\oplus2}\quad(1\le i\le7).             \tag{D.GH.8}
\]
The alternating line syzygies are exactly those in that
construction. All differential coefficients are
parameters in spaces of global sections of these fixed
terms; no moving determinant or variable cohomology
basis is introduced.

Impose \(d^2=0\) and the actual open rank conditions.
Local unit-minor cancellation makes the degree-zero
cohomology a bundle over every coefficient algebra.
Use length seven for source complexes, length three for
targets, and the same top-generator identity test as in
Theorem~\ref{D:curve-fixed-local-linear-growth}.
The two end-error intervals are disjoint from the Hom
degrees used for maps and identities. Thus the finite
Hom cycles and degree-minus-one boundary witnesses
express precisely (D.GH.4) and their inverse products.

We verify the dependence on the changing curve.
For any line of degree \(d\) on \(X\),
\(h^0+h^1\le2g+2+|d|\); this follows from a regular
section when one exists and Riemann--Roch, and is
immediate when there is no section.
Every top--top Hom block has multiplicity \(O(r^2)\)
and line degree \(O_p(h+\Delta)\). A block with one
positive-index term has multiplicity \(O(r)\) and
line degree \(O_p(rh+\Delta)\); a block with two
positive-index terms has bounded multiplicity and
the same degree bound. The twists \(H\), \(B\), and
one-step Frobenius in (D.GH.4) only change the
absolute constants depending on \(p\). Hence each
phase has
\[
                   O_p(A r^2)                         \tag{D.GH.9}
\]
object and Hom coordinates.

As in the proved fixed-type atlas, split the two-affine
Cech complexes of the fixed line terms over the ground
field. The internal widths are bounded by seven and by
the fixed map-composition diagram. The finite homological
perturbation identities therefore use a bounded number
of factors, independent of \(g,\Delta,r,f\).
They commute with every coefficient-base change and
give polynomial equations of degree bounded by a constant
depending only on \(p\). The only coefficient power in
one phase is the \(p\)-th power. No expanded \(q\)-Frobenius
identity occurs.

For each \(x\), choose
\[
 N_x=\left\lfloor\frac{c_x}{p-1}\right\rfloor+1,
 \qquad L_x=N_x+c_x .
                                                        \tag{D.GH.10}
\]
Then \((p-1)N_x>c_x\) and
\(\sum_xL_x=O_p(\Delta+|P|)\).
Frame the quotients of (D.GH.8) on
\(\coprod_x\operatorname{Spec}k[t_x]/(t_x^{L_x})\)
using the explicit matrices \(u,v,w\) of (D.CL.9),
and impose the finite cyclic unit equations with the
fixed canonical local models. These add
\[
                  O_p\!\left(fr^2\sum_xL_x\right)
                         =O_p(f A r^2)                \tag{D.GH.11}
\]
variables and equations of bounded degree.
Increasing the truncation length increases the number
of coefficient variables, not the polynomial degree.
Lemma~\ref{D:curve-integral-unit-jet-comparison}
gives the exact integral local identifications over
all coefficient algebras.

Every geometric point of the resulting polynomial
zero scheme and open rank locus is an actual global
unit bundle with these fixed canonical local types,
and every representation counted is covered.
The entire specified-Artinian Isom proof in
Theorem~\ref{D:curve-fixed-local-linear-growth}
applies unchanged: a unit isomorphism lifts uniquely
on \(X^\circ\), its local lift preserves the fixed
canonical invariant lattices, and faithful flatness
of completion extends it over the proper curve.
The finite-presentation Isom is etale and its distinct
class images are disjoint open subsets. This uses
the fixed full local types, not just the different
bounds.

There are at most \(C_p f A r^2\) affine variables
and the equation degrees are at most \(D_p\), with
constants depending only on \(p\).
The component bound, including its prefactor, is
\[
 (C_p f A r^2+1)\max\{2,D_p\}^{C_p f A r^2}
                         \le p^{C'_p f A r^2}.         \tag{D.GH.12}
\]
Passing to an open locus does not add irreducible
components. Finally multiply by the at most
\((q-1)^{|P|}\) tuples of linear local lifts.
Since \(\Delta\le p\sum_xd_x+(p+1)|P|\),
this proves (D.GH.1). The case of an empty counted
fibre, and the trivial projective rank-one case,
are included without any additional choices.
\end{proof}

\begin{theorem}[Proved: global hereditary growth with bounded local periods]
\label{D:global-finite-puncture-hereditary-bounded-local-period}
Let \(S\) be a fixed normal proper connected surface
over an algebraically closed field \(k\) of characteristic
\(p>0\), let \(T\subset S\) be finite, and put
\(\Gamma=\pi_1(U)\), \(U=S-T\).
Choose the fixed finite set \(F\supset T\), complete
local sources \(R_s\), and detecting curve of
Lemma~\ref{D:finite-puncture-detecting-curve-local-types}.
Fix \(B_*\ge1\).

For an open \(H\subset\Gamma\) of index \(t\), let
\(U_H\to U\) be its connected finite etale cover.
For each \(s\in F\), write the pullback to the fixed
puncture \(U_{R_s}\) as its connected components:
\[
 U_H\times_U U_{R_s}=\coprod_i U_{s,i},
 \qquad n_{s,i}=\deg(U_{s,i}/U_{R_s}),\qquad
                       \sum_i n_{s,i}=t.               \tag{D.GH.13}
\]
Each \(U_{s,i}\) is the puncture of its complete
normal local normalization of \(R_s\).
There is a constant \(A_{S,T,B_*}\) such that the
homomorphisms
\(\rho:H\to\operatorname{PGL}_r(\mathbb F_q)\),
\(q=p^f\), for which every restriction to
\(\pi_1(U_{s,i})\) has geometric Brauer period at most
\(B_*\), have at most
\[
                       q^{A_{S,T,B_*}t r^2}             \tag{D.GH.14}
\]
projective conjugacy classes.

It suffices to require the period condition at
\(s\in T\); the other restrictions are trivial.
In particular the conclusion covers every such
representation if \(\gcd(r,q-1)\le B_*\).
A bound \(B_*\) on the global geometric Brauer
period on \(U_H\) is another sufficient condition.
The detecting curve is chosen once for \((S,T)\).
No constant belonging to a new unspecified curve
or to a new surface resolution is used.
\end{theorem}

\begin{proof}
Write \(C^\circ=C-P\) for that one fixed detecting
curve over its fixed algebraically closed extension
\(\Omega/k\), and let \(b_x\) be its fixed upper-break
bounds at \(x\in P\).
The pullback \(C_H^\circ=C^\circ\times_U U_H\)
is connected, since
\(\pi_1(C^\circ)\twoheadrightarrow\Gamma\).
Let \(C_H\) be its smooth proper completion and
\(P_H=C_H-C_H^\circ\). Then \(C_H\to C\) is a
finite separable morphism of degree \(t\), and
\(\pi_1(C_H^\circ)\twoheadrightarrow H\).
All these are pullbacks of the original detector;
no further Bertini choice is made.

For \(y\in P_H\) lying over \(x\in P\), put
\[
 K_x=\Omega((t_x)),\quad K_y=\Omega((s_y)),\quad
 e_y=[K_y:K_x],\quad \delta_y=\delta(K_y/K_x).
                                                        \tag{D.GH.15}
\]
The residue fields are \(\Omega\), so the displayed
degree is the ramification index and
\(\sum_{y\mid x}e_y=t\).
Choose a finite Galois quotient of \(\Gamma\) whose
kernel is contained in \(H\). Its local splitting
field at \(x\) contains every \(K_y\) after compatible
embeddings. The original detecting-curve lemma bounds
its normalized different by \(b_x+1\).
The tower formula (D.GH.2), without its tame
specialization, implies
\[
             0\le\delta_y\le e_y(b_x+1).
                                                        \tag{D.GH.16}
\]
The Riemann--Hurwitz formula with the actual different
therefore gives
\[
 \begin{aligned}
 2g(C_H)-2
   &=t(2g(C)-2)+\sum_{y\in P_H}\delta_y,\\
 |P_H|&\le t|P|,\qquad g(C_H)+|P_H|+1=O_{S,T}(t).
 \end{aligned}                                         \tag{D.GH.17}
\]
The exact wild different formula and its derivation
are \citet[Section 53.12 and Lemma 53.12.2,
Tag 0C1B]{StacksResidualGenusB}. We have not used the
tame replacement \(\delta_y=e_y-1\).

The same argument also controls all representations
of \(H\), rather than only the degree-\(t\) cover.
Given any finite representation of \(H\), its kernel
has an open normal core in \(\Gamma\).
The corresponding finite Galois quotient of
\(\Gamma\) simultaneously trivializes that
representation and the coset cover.
Let \(L/K_x\) be its local splitting field.
It is Galois, contains \(K_y\), and satisfies
\(\delta(L/K_x)/e(L/K_x)\le b_x+1\).
Transitivity gives the exact equality and bound
\[
 \frac{\delta(L/K_y)}{e(L/K_y)}
 =e_y\,\frac{\delta(L/K_x)}{e(L/K_x)}-\delta_y
 \le d_y,\qquad
                  d_y=e_y(b_x+1)-\delta_y\ge0.          \tag{D.GH.18}
\]
Passing to a Galois quotient over \(K_y\) can only
decrease its normalized different, again by
transitivity. Thus \(d_y\) bounds the projective
local splitting type of every representation
under consideration.
In particular
\[
                 \sum_{y\in P_H}d_y
                      \le t\sum_{x\in P}(b_x+1)
                      =O_{S,T}(t).                    \tag{D.GH.19}
\]
For any same-field linear lift on \(C_H^\circ\),
the scalar tame extension adds at most one to
each normalized different by (D.GH.2).
Consequently the local inverse-Frobenius depths
\(\lceil p(d_y+1)\rceil\), and the retained and
testing jet lengths (D.GH.10), have total \(O_{S,T,p}(t)\).
Using the maximum local depth at all marked points
would lose this estimate and is unnecessary.

It remains to count the possible tuples of full local
projective types. Each source \(R_s\) is fixed before
\(H\) is chosen. Its completed finite cover factors
in (D.GH.13) have total degree \(t\), even when
the maps of local fundamental groups to \(\Gamma\)
are not injective.
Theorem~\ref{D:bounded-period-hereditary-projective-growth}
therefore bounds all representation classes on
these factors by
\[
 \prod_{s\in F}\prod_i q^{A_{R_s,B_*}n_{s,i}r^2}
                  =q^{t r^2\sum_{s\in F}A_{R_s,B_*}}.
                                                        \tag{D.GH.20}
\]
The factorization of the original formal branches
through \(U_{R_s}\) pulls back to factorizations
through the appropriate \(U_{s,i}\).
The finite etale algebra over \(K_x\) is the product
of the fields \(K_y\), so these are precisely all
the full local types on \(C_H\).
Choose compatible paths once for this fixed cover.
One class on a factor \(U_{s,i}\) determines, by
restriction, the classes at every curve branch
using that factor. Hence (D.GH.20) counts each
surface factor once. There is no extra factor
equal to the number of branches or the local
cover degree. Independent choices of different
surface factors can overcount, which is harmless.

For each retained tuple apply
Lemma~\ref{D:curve-fixed-local-growth-numerical}
to \(C_H,P_H\), with the \(d_y\) of (D.GH.18).
Equations (D.GH.17)--(D.GH.19) bound its exponent
by a constant times \(t r^2\).
Multiplying by (D.GH.20) proves the same bound
for the corresponding projective classes on
\(C_H^\circ\). Precomposition with the surjection
\(\pi_1(C_H^\circ)\twoheadrightarrow H\) is injective
on conjugacy classes, proving (D.GH.14).

If \(s\notin T\), every original cover of \(U\)
is split over \(\operatorname{Spec}R_s\), as in
the detecting-curve lemma; the induced local
representations are trivial. A global geometric
period bound restricts to every factor in
(D.GH.13). Finally the ordinary projective
multiplier is killed by both \(r\), by taking
determinants of linear lifts, and \(q-1\), by
its coefficient exponent. Its geometric image
therefore has period dividing \(\gcd(r,q-1)\).
These prove the stated sufficient conditions.
\end{proof}

\begin{remark}[Scope of the hereditary transfer]
The new uniformity concerns arbitrary finite covers of
one fixed finite-puncture proper surface and fixed
bounded local geometric periods. It is stronger than
reapplying a fixed-surface theorem separately upstairs:
the genus, total different, local jet length, and
surface-source indices have all been bounded linearly
in the original index. It does not claim unrestricted
local period, a semilinear projective transfer, or
finite generation of the whole global group.
This is a transfer of the same bounded-period
hereditary projective enlargement, not a separate
independent reduction.
\end{remark}


% Proved after four complete actual-file audits. Closure of the bounded-period hereditary projective-semilinear count.
\subsection{Bounded-period projective semilinear crowns}

\begin{corollary}[Proved: projective semilinear crowns over arbitrary permutation quotients]
\label{D:bounded-period-arbitrary-quotient-projective-semilinear-crowns}
Let $R,\Gamma$ be as in
Theorem~\ref{D:bounded-period-semilinear-hereditary-projective-growth}.
Fix integers $B_*,b_0\ge1$ and $\kappa>0$.
Let $L$ be a finite monolithic primitive group with
nonabelian socle $N=S^t$, where $S$ is nonabelian simple,
and put $P=L/N$. Fix a continuous epimorphism
$\beta:\Gamma\to P$. Choose one simple socle factor,
and let $\mathcal B\le\operatorname{Aut}(S)$ be the
image of its stabilizer in $L$.
Suppose
\[
 \begin{gathered}
 S\le\mathcal B_0\le\mathcal B,\qquad
 [\mathcal B:\mathcal B_0]\le b_0,\\
 \iota:\mathcal B_0\lhook\joinrel\longrightarrow
     \operatorname{PGL}_r(\mathbb F_q)
        \rtimes\operatorname{Gal}(\mathbb F_q/\mathbb F_p),
 \qquad r^2\log q\le\kappa\log|S|,
 \end{gathered}                                         \tag{D.BS.1}
\]
where $q$ is a power of $p$. Here $S$ is identified
with its group of inner automorphisms.
Normality of $\mathcal B_0$ in $\mathcal B$ is not
required.

Let $H\subset\Gamma$ be the inverse image of the point
stabilizer of the chosen factor. Its prescribed outer
map is $\eta:H\to\mathcal B/S$; put
$H_0=\eta^{-1}(\mathcal B_0/S)$.
Perfectness of $S$ makes the field projection of
$\iota$ factor through $\mathcal B_0/S$. Consequently
the map $\eta|_{H_0}$ specifies one field character
$\chi_0:H_0\to\operatorname{Gal}(\mathbb F_q/\mathbb F_p)$.
Write $K_0=\ker\chi_0$.

For a crown quotient $\Gamma\twoheadrightarrow L_k$
over $\beta$, require that each of its stabilizer
lifts $\theta:H\to\mathcal B$, after restriction to
$K_0$, have geometric Brauer period at most $B_*$
under the ordinary projective map $\iota\theta|_{K_0}$.
Then
\[
 \begin{aligned}
 k&\le |N|^{c_R},\\
 d(L_k)&\le\max\{2,d(P),c_R+2\},\\
 c_R&=\max\left\{1,
  \left\lceil\kappa(A_{R,B_*}^{\rm ps}b_0+1)+b_0-1
                                              \right\rceil\right\}.
 \end{aligned}                                         \tag{D.BS.2}
\]
There is no tameness or resolution-extension assumption
on the permutation cover or on $\chi_0$, and its image
order is unrestricted. The quadratic order comparison,
bounded residual index, and the displayed $d(P)$ term
remain hypotheses or conclusions exactly as written.

The period condition is automatic if
$\gcd(r,q-1)\le B_*$.
It is also automatic for $B_*\ge2$ under the following
explicit algebraic condition. The full inverse image
$\widetilde{\mathcal B}_0$ of $\iota(\mathcal B_0)$ in
the general semilinear group consists of semisimilitudes
of one fixed nonzero bilinear or quadratic form over
$\mathbb F_q$. In the bilinear case this means that
every $(A,\sigma)\in\widetilde{\mathcal B}_0$ satisfies
\[
 B(A\sigma v,A\sigma w)
       =\lambda(A,\sigma)\,\sigma(B(v,w)),
       \qquad \lambda(A,\sigma)\in\mathbb F_q^\times;
                                                        \tag{D.BS.3}
\]
in the quadratic case replace the two sides by
$Q(A\sigma v)$ and
$\lambda(A,\sigma)\sigma(Q(v))$.
This is a condition on the actual semilinear embedding,
not a recognition assertion for an unspecified
automorphism group.
\end{corollary}

\begin{proof}
The faithful conjugation embedding of $L$ into the
automorphism wreath product and the exact nonsplit
transport of
Lemma~\ref{D:nonsplit-automorphism-wreath-lifts}
identify its lift classes over $\beta$, modulo $N$,
with stabilizer lifts $\theta:H\to\mathcal B$ of
$\eta$, modulo $S$. The actual outer image determines
the full inverse image containing the socle, so no
outer extension is split in this identification.
Surjectivity onto $L_k$ forces its $k$ coordinate
lifts to give distinct such $S$-classes: two inner
conjugate coordinates would have image in a diagonal
graph subgroup of $L\times_P L$.
It is enough to bound the classes satisfying the
stated period condition.

Put $a=[H:H_0]\le b_0$. Then $[\Gamma:H]=t$ and
$[\Gamma:H_0]=at$. All these subgroups are determined
by $\beta$ and the fixed subgroup $\mathcal B_0$.
The field group is abelian, so its projection kills
the perfect group $S$. The projective-semilinear
restrictions of every lift to $H_0$ therefore have
the same $\chi_0$, including its entire field image.
The period hypothesis is precisely the one in
Theorem~\ref{D:bounded-period-semilinear-hereditary-projective-growth}
for this field character. That theorem gives at most
$q^{A_{R,B_*}^{\rm ps}at r^2}$ ambient
$\operatorname{PGL}_r(\mathbb F_q)$ conjugacy classes
of the restricted maps.

Faithfulness of $\iota$ identifies an actual matrix
homomorphism with at most one $\mathcal B_0$-valued
restriction. An ambient projective conjugacy orbit
contains at most $|\operatorname{PGL}_r(\mathbb F_q)|
\le q^{r^2}$ actual homomorphisms and hence no more
than that many $S$-classes. This bound does not
require the ambient conjugating matrices to normalize
$\iota(\mathcal B_0)$ or $\iota(S)$. Thus the number
of restricted $S$-classes is at most
\[
 q^{(A_{R,B_*}^{\rm ps}at+1)r^2}
             \le |S|^{\kappa(A_{R,B_*}^{\rm ps}at+1)}.
                                                        \tag{D.BS.4}
\]

For each class select an actual restriction and
normalize a full lift to it by conjugating with $S$.
This preserves its values in $\mathcal B$ and its
outer map. Choose one representative of every
nonidentity coset of $H_0$ in $H$. Each representative
has its image in one fixed $S$-coset prescribed by
$\eta$, so has at most $|S|$ possibilities. Its
restriction and these $a-1$ images determine the
whole homomorphism. The resulting cost is at most
$|S|^{a-1}$. Normality of $H_0$ is unnecessary:
coset representatives together with $H_0$ generate
$H$, and every relation can only remove choices.
No ambient action of all of $\mathcal B$ has been
assumed.

Combining these bounds gives
\[
 \begin{aligned}
 k&\le |S|^{\kappa(A_{R,B_*}^{\rm ps}at+1)+a-1}\\
  &\le |S|^{t\{\kappa(A_{R,B_*}^{\rm ps}b_0+1)+b_0-1\}}
   \le |N|^{c_R}.
 \end{aligned}                                         \tag{D.BS.5}
\]
Apply the exact nonabelian crown estimate already
used in
Corollary~\ref{D:arbitrary-quotient-faithful-semilinear-crowns},
from \citet[proof of Proposition 2.6(ii),
equations (2.1)--(2.2), pp.11--13]{LucchiniMarionTracey2020D}.
With $s=\max\{2,d(P),c_R+2\}$ it gives $d(L_k)\le s$.
In particular no bound for $d(P)$ is inserted.

The first sufficient period condition follows by
taking determinants of linear lifts after restriction
to $K_0$: their scalar multiplier is killed by $r$
and by $q-1$.
We verify the second condition at the level of the
full twisted coefficient class, rather than only its
geometric image. Nonzeroness of the form makes the
factor $\lambda$ in (D.BS.3) unique. Substitution
into the form identity twice gives
\[
 \lambda(gh)=\lambda(g)\,\sigma_g(\lambda(h)),
 \qquad \lambda(cI)=c^2 .
                                                        \tag{D.BS.6}
\]
The same calculation holds for a quadratic form,
including in characteristic two. Thus $\lambda$,
together with the field projection, is an actual
homomorphism
$\widetilde{\mathcal B}_0\to
\mathbb F_q^\times\rtimes\operatorname{Gal}(\mathbb F_q/\mathbb F_p)$
whose restriction to the scalar kernel is the square
map.

Choose lifts $\widetilde\theta(g)$ for the restriction
of a stabilizer lift to $H_0$, and write its scalar
multiplier as
$\widetilde\theta(g)\widetilde\theta(h)
=c(g,h)\widetilde\theta(gh)$.
Applying (D.BS.6) yields
\[
 c(g,h)^2=
 \frac{\lambda(\widetilde\theta(g))\,
       \chi_0(g)(\lambda(\widetilde\theta(h)))}
      {\lambda(\widetilde\theta(gh))}.
                                                        \tag{D.BS.7}
\]
Hence the actual class in
$H^2(H_0,\mathbb F_q^\times(\chi_0))$ is killed by
$2$. Its restriction to $K_0$ and its geometric
Brauer image are killed by $2$ as well. This proves
the sufficient condition without an additional
bound on the scalar kernel, the field-image order,
or the permutation quotient.
\end{proof}

\begin{remark}[Recognition is a separate input]
The semisimilitude condition may be checked in a
specified natural matrix model before applying this
corollary. The statement does not assert that the
whole automorphism group of any named simple family
has such a model, nor that an unbounded graph or
residual outer quotient has been removed.
It is a closure of the same bounded-period hereditary
projective-semilinear reduction, not a separate
counting movement or a finite-generation theorem.
\end{remark}


% Proved after four complete actual-file/source audits. Known cohomology with the actual scalar-band actions retained.
\subsection{Twisted scalar coefficients on affine curves and formal branches}

\begin{lemma}[Proved: cyclic scalar bands on a smooth affine curve]
\label{D:affine-curve-field-twisted-H2}
Let \(X\) be a smooth connected affine curve over an
algebraically closed field of characteristic \(p>0\).
Let \(\mathcal M\) be a finite locally constant sheaf
of cyclic groups of order \(n\), with \(p\nmid n\),
and let \(M\) be its geometric fibre with its actual
\(\pi_1(X)\)-action. Then
\[
 H^2_{\mathrm{et}}(X,\mathcal M)=0,\qquad
 H^2_{\mathrm{cont}}(\pi_1(X),M)=0.                 \tag{D.CS.1}
\]
In particular fix \(q=p^f\) and a field character
\(\chi:\pi_1(X)\to\operatorname{Gal}(\mathbb F_q/\mathbb F_p)\).
Every continuous homomorphism
\[
 \rho:\pi_1(X)\longrightarrow
 \operatorname{PGL}_r(\mathbb F_q)\rtimes
                 \operatorname{Gal}(\mathbb F_q/\mathbb F_p)
                                                        \tag{D.CS.2}
\]
with field projection \(\chi\) lifts over the same
finite field to \(\operatorname{GL}_r(\mathbb F_q)
\rtimes\operatorname{Gal}(\mathbb F_q/\mathbb F_p)\),
with the same field projection. Neither the order of
the monodromy of \(\mathcal M\) nor the order of the
image of \(\chi\) is required to be prime to \(p\).
\end{lemma}

\begin{proof}
The etale assertion is a special case of the known
affine-curve vanishing theorem
\citep[Proposition 14.12, with the preceding
prime-to-characteristic coefficient convention]
{MilneLEC2013H2D}. We give a direct verification for
cyclic coefficients that retains their action.
It uses only constant-coefficient affine-curve
vanishing and prime-to-coefficient averaging.

The primary decomposition of the cyclic fibre is
characteristic, hence decomposes \(\mathcal M\) into
locally constant cyclic \(\ell\)-primary summands.
It suffices to consider one of order \(\ell^e\),
where \(\ell\ne p\). Its filtration by powers of
\(\ell\) has rank-one \(\mathbb F_\ell\)-local systems
as its successive quotients. For any such quotient
\(\mathcal L\), its monodromy character has image
\[
       G\subset\mathbb F_\ell^\times,\qquad
                       |G|\mid\ell-1.                 \tag{D.CS.3}
\]
The connected finite etale Galois cover \(v:X'\to X\)
defined by the kernel trivializes \(\mathcal L\).
The curve \(X'\) is again smooth and affine.
The Cartan--Leray spectral sequence
\citep[Theorem 14.9 and its complete written proof]
{MilneLEC2013H2D} has terms
\[
 H^i\bigl(G,H^j_{\mathrm{et}}(X',v^*\mathcal L)\bigr)
                 \ \Longrightarrow
 H^{i+j}_{\mathrm{et}}(X,\mathcal L).                 \tag{D.CS.4}
\]
Every coefficient module here is an
\(\mathbb F_\ell\)-vector space. Since \(|G|\) is
invertible in \(\mathbb F_\ell\), averaging by
\(|G|^{-1}\sum_{g\in G}g\) makes the invariant
functor exact, so the terms with \(i>0\) vanish.
The remaining degree-two term vanishes by constant
affine-curve cohomology
\citep[Lemma 59.69.3, Tag 03RR, complete written proof]
{StacksCurveFixedLocalD}. Thus \(H^2(X,\mathcal L)=0\).
Induction along the \(\ell\)-power filtration, using
the long exact sequence in cohomology, proves
\(H^2(X,\mathcal M)=0\). This induction uses only
the degree-two groups of the subobject and quotient;
no vanishing in degree three is being assumed.
The averaging degree in (D.CS.3) need not be prime
to the geometric characteristic \(p\).

For the continuous group statement use the tower
Cartan--Leray comparison of
\citet[Theorem 14.9 and Remark 14.10]{MilneLEC2013H2D}.
First pass to a finite Galois cover trivializing the
coefficient sheaf. Every class in degree one at any
later finite stage is killed on a further finite
etale cover: take its finite torsor, one connected
component, and then a Galois closure. Consequently
the degree-one cohomology in the universal-cover
direct limit is zero. The five-term sequence gives
the injection
\[
 H^2_{\mathrm{cont}}(\pi_1(X),M)
       \lhook\joinrel\longrightarrow
 H^2_{\mathrm{et}}(X,\mathcal M).                      \tag{D.CS.5}
\]
The action on \(M\) is not replaced by a trivial
action in this comparison. This proves (D.CS.1).

Finally the kernel of the natural general-to-projective
semilinear map is \(\mathbb F_q^\times\).
Conjugation through \(\rho\) acts on that kernel by
the specified \(\chi\), so the lifting obstruction
belongs to the actual group
\(H^2(\pi_1(X),\mathbb F_q^\times(\chi))\).
Choose set-theoretic lifts on the finite image of
\(\rho\). Their multiplier is a continuous cocycle
for this module. By (D.CS.1) it is a continuous
coboundary; adjusting the lifts by that cochain
produces the required homomorphism. Its target and
scalar kernel are finite, so continuity is retained.
The case \(q=2\), or \(n=1\), is included with the
trivial scalar module.
\end{proof}

\begin{lemma}[Proved: the exact local field-character scalar group]
\label{D:local-field-character-crossed-H1}
Let \(K=\Omega((t))\), where \(\Omega\) is algebraically
closed of characteristic \(p>0\), and put
\(G_K=\operatorname{Gal}(K^{\mathrm{sep}}/K)\).
Let \(q=p^f\) and fix a continuous character
\[
 \chi:G_K\longrightarrow
             \operatorname{Gal}(\mathbb F_q/\mathbb F_p)
\]
whose image has order \(a\mid f\). Write
\(q_x=p^{f/a}\) and
\(M_\chi=\mathbb F_q^\times(\chi)\).
Then
\[
 \left|H^1_{\mathrm{cont}}(G_K,M_\chi)\right|
                           =q_x-1\le q-1.             \tag{D.CS.6}
\]
No tameness assumption on the field-kernel extension
is required, including when \(p\mid a\).
For a fixed local projective-semilinear homomorphism
with field character \(\chi\), its same-field linear
lifts, if any, have at most \(q_x-1\) conjugacy classes
under \(\operatorname{GL}_r(\mathbb F_q)\).
The same bound holds if the projective map is
specified only up to projective conjugacy.
\end{lemma}

\begin{proof}
Let \(K'/K\) be the actual cyclic extension cut out
by \(\chi\), let \(C_a=\operatorname{Gal}(K'/K)\),
and let \(n=q-1=q_x^a-1\).
The field \(K'\) is complete with residue field
\(\Omega\); after choosing a uniformizer \(s\),
successive expansion using the existing coefficient
field \(\Omega\) identifies it with \(\Omega((s))\).
The coefficient action on \(G_{K'}\) is trivial.
A generator of \(C_a\) can be chosen to act on
\(M=\mathbb F_q^\times\) by \(z\mapsto z^{q_x}\).

We record both cyclic cohomology groups explicitly.
In additive notation \(M=\mathbb Z/n\), the two
maps in the cyclic resolution are
\[
 q_x-1,\qquad
 T=1+q_x+\cdots+q_x^{a-1}
                         =\frac{n}{q_x-1}.            \tag{D.CS.7}
\]
The standard periodic resolution has alternating
maps \(\gamma-1\) and \(1+\gamma+\cdots+\gamma^{a-1}\);
its exactness follows by comparing coefficients in
the free group ring. It gives
\[
 H^1(C_a,M)=\ker T/(q_x-1)M,\qquad
 H^2(C_a,M)=M^{C_a}/TM.                               \tag{D.CS.8}
\]
Here \(TM\) and \(M^{C_a}\) both have order \(q_x-1\)
and the former is contained in the latter.
Similarly \(\ker T\) and \((q_x-1)M\) both have order
\(T\) and the latter is contained in the former.
Thus both groups in (D.CS.8) are zero. This is the
finite-field norm and Hilbert-90 computation without
division by \(a\); it remains valid when \(p\mid a\).

The five-term sequence for the actual extension
\(1\to G_{K'}\to G_K\to C_a\to1\) consequently gives
\[
 H^1(G_K,M_\chi)\ \simeq
                  H^1(G_{K'},M)^{C_a}.               \tag{D.CS.9}
\]
This is the same Cartan--Leray sequence for the
Galois cover \(\operatorname{Spec}K'\to
\operatorname{Spec}K\) cited above.
Choose an identification \(M=\mu_n(\Omega)\).
Kummer theory
\citep[Lemma 59.28.1 and its following long exact
sequence, Tag 03PK]{StacksControlledAnchorsD}
identifies
\[
 H^1(G_{K'},M)=K'^\times/K'^{\times n}
                   \xrightarrow[\ v_s\ ]{\ \simeq\ }
                              \mathbb Z/n.            \tag{D.CS.10}
\]
Indeed \(\Omega^\times\) is \(n\)-divisible, and
every principal unit has a unique \(n\)-th root
with prescribed constant root, by successive
coefficients since \(p\nmid n\). The valuation
therefore gives the displayed isomorphism.

The geometric deck automorphism of \(K'\) preserves
its valuation and acts trivially on (D.CS.10).
The coefficient action in (D.CS.9) still multiplies
that cyclic group by \(q_x\). The invariant subgroup
is consequently
\(\ker(q_x-1:\mathbb Z/n\to\mathbb Z/n)\),
of order \(q_x-1\). This proves (D.CS.6).
There has been no passage to constant coefficients
on \(G_K\), and no assumption on a global wild kernel.

To verify the assertion about lifts, fix one linear
lift. A second lift of the same projective map differs
from it by a continuous crossed character with values
in \(M_\chi\). Scalar conjugation changes that character
by a coboundary. Hence the number of linear conjugacy
classes is bounded by (D.CS.6); further nonscalar
conjugations can only identify classes. A projective
conjugating matrix has a lift in
\(\operatorname{GL}_r(\mathbb F_q)\), reducing a
prescribed projective conjugacy class to this same
calculation. We have not asserted the existence of
a local lift in this last argument; in the curve
application it is supplied by the preceding lemma.
\end{proof}

\begin{remark}[Known cohomology and the retained coefficient action]
These are cohomological inputs to the semilinear
curve presentation. Affine-curve vanishing is Known;
the explicit rank-one filtration and the local cyclic
calculation record its precise scalar-band use here.
They assert neither a new representation-counting
theorem nor an independent reduction of the main
problem. In particular the field-kernel degree is
not a constant-coefficient counting factor.
\end{remark}


% Proved after four complete actual-file audits. The exact curve scalar-cohomology inputs are proved and their dependency audits retained.
\subsection{Global hereditary semilinear growth with bounded local periods}

\begin{lemma}[Proved: a numerical curve bound with a prescribed field character]
\label{D:curve-fixed-local-semilinear-growth-numerical}
There is a constant \(c_p^{\rm sl}\), depending only on
the positive characteristic \(p\), with the following
property. Let \(X\) be a smooth proper connected curve
of genus \(g\) over an algebraically closed field of
characteristic \(p\), let \(P\ne\varnothing\) be a finite
reduced divisor, and write \(X^\circ=X-P\).
For \(q=p^f\), prescribe a continuous field character
\[
 \chi:\pi_1(X^\circ)\longrightarrow
                   \operatorname{Gal}(\mathbb F_q/\mathbb F_p)
                                                        \tag{D.GS.1}
\]
and full local projective-semilinear representation
classes \(\sigma_x\) with this field projection.
Suppose that the finite Galois splitting extension of
\(\sigma_x\) has normalized different at most \(d_x\ge0\).
Then the homomorphisms to
\(\operatorname{PGL}_r(\mathbb F_q)\rtimes
\operatorname{Gal}(\mathbb F_q/\mathbb F_p)\)
with the prescribed field character and full local types
have at most
\[
 q^{\,c_p^{\rm sl}(1+g+|P|+\sum_{x\in P}d_x)r^2}
                                                        \tag{D.GS.2}
\]
\(\operatorname{PGL}_r(\mathbb F_q)\)-conjugacy classes.
The constant is independent of the field-image order,
which may be divisible by \(p\), and of all local
splitting degrees and coefficients.
\end{lemma}

\begin{proof}
Put \(n=q-1\) and
\(M_\chi=\mathbb F_q^\times(\chi)\), with the actual
field action. Lemma~\ref{D:affine-curve-field-twisted-H2}
gives
\[
 H^2(X^\circ,M_\chi)=0,
 \qquad H^2(\pi_1(X^\circ),M_\chi)=0.
                                                        \tag{D.GS.3}
\]
The second assertion uses the degree-two comparison
injection, not a replacement of \(M_\chi\) by a constant
sheaf. Consequently every representation under
consideration lifts over the same field to
\(\operatorname{GL}_r(\mathbb F_q)\rtimes
\operatorname{Gal}(\mathbb F_q/\mathbb F_p)\), with
the same field character.

We count the possible full local linear lifts on
\(X\), before making the field cover.
Let \(a_x\) be the image order of the local field
character, and set \(q_x=p^{f/a_x}\).
If a local lift exists, the possible lifts of a fixed
projective-semilinear homomorphism form a torsor under
crossed homomorphisms to \(M_{\chi_x}\).
Scalar linear conjugation identifies coboundaries.
Lemma~\ref{D:local-field-character-crossed-H1} gives
\[
 \bigl|H^1(G_{k((t_x))},M_{\chi_x})\bigr|=q_x-1
                                                     \le q-1.
                                                        \tag{D.GS.4}
\]
This is the genuine twisted calculation. On the cyclic
field-kernel extension the Kummer group is \(\mathbb Z/n\)
by valuation, geometric deck pullback acts trivially,
and the fixed subgroup for the coefficient action has
size \(q_x-1\); the finite-field Hilbert--90 term is zero.
It is valid also when \(p\mid a_x\).
For classes rather than fixed homomorphisms, lift any
projective conjugating matrix to a linear matrix first;
the same crossed-cocycle bound then applies.
If the global fibre is nonempty, a global linear lift
supplies one local lift at every point. Thus at most
\((q-1)^{|P|}\) tuples need be considered. For \(q=2\)
this factor is one.

Fix one tuple of full linear-semilinear local lifts.
Write
\[
 a=|\operatorname{im}\chi|,\qquad
 \nu=f/a,\qquad q_0=p^\nu,\qquad
 A=1+g+|P|+\sum_{x\in P}d_x.
                                                        \tag{D.GS.5}
\]
Let \(V^\circ\to X^\circ\) be the connected degree-\(a\)
cover defined by \(\ker\chi\), let \(V\) be its smooth
proper completion, and let \(P_V=V-V^\circ\).
Choose the deck generator \(\tau\) so that a skew
map is written \(F_{q_0}^*E\to\tau^*E\).
Thus \(\tau^{-1}\) is the deck generator used for
the coefficient automorphism \(z\mapsto z^{q_0}\)
in Lemma~\ref{D:twisted-deck-Frobenius-semilinear-model}:
applying \(\tau^{-1*}\) gives its map
\(\tau^{-1*}F_{q_0}^*E\to E\).
This is the actual cover determined by (D.GS.1).
No split lift of its generator to the fundamental
group is chosen or required.

At \(x\), the local field-kernel extension
\(K'_x/K_x\) is cyclic of degree \(a_x\), with
different exponent \(\delta_x'\). There are
\(a/a_x\) points of \(P_V\) over \(x\).
Since this field extension is a quotient of the full
projective-semilinear splitting extension,
\[
       0\le\delta_x'\le a_x d_x.
                                                        \tag{D.GS.6}
\]
Riemann--Hurwitz with the actual wild differents gives
\[
 \begin{aligned}
 2g(V)-2&=a(2g-2)+\sum_{x\in P}(a/a_x)\delta_x',\\
 |P_V|&=\sum_{x\in P}a/a_x\le a|P|,\\
 g(V)+|P_V|+1&\le C a A
 \end{aligned}                                         \tag{D.GS.7}
\]
for an absolute constant \(C\).

We also need the different of the actual linear
system after this field cover. Let \(L_x/K_x\)
split the fixed projective-semilinear local type.
Then \(K'_x\subset L_x\), and exact transitivity gives
\[
 \frac{\delta(L_x/K'_x)}{e(L_x/K'_x)}
 =a_x\frac{\delta(L_x/K_x)}{e(L_x/K_x)}-\delta_x'
 \le a_xd_x-\delta_x'.
                                                        \tag{D.GS.8}
\]
The linear splitting field over \(L_x\) has scalar
group of order dividing \(q-1\), so is tame. Applying
the tame part of (D.GH.2) over \(K'_x\), rather than
multiplying its former bound unnecessarily by \(a_x\),
shows that the restricted linear type has normalized
different at most
\[
             d'_y=a_xd_x-\delta_x'+1
                  \quad (y\mid x).
                                                        \tag{D.GS.9}
\]
The derivative chain-rule proof of different
transitivity in Lemma~\ref{D:curve-fixed-local-growth-numerical}
applies to all these equal-characteristic complete
DVRs. In particular
\[
 \sum_{y\in P_V}d'_y
       \le a\sum_{x\in P}(d_x+1).
                                                        \tag{D.GS.10}
\]
Set
\[
 c_y=\lceil p d'_y\rceil,\qquad
 B_V=\sum_{y\in P_V}c_y y,\qquad
 \Delta_V=\deg B_V=O_p(aA).
                                                        \tag{D.GS.11}
\]
These integers and this divisor are deck-invariant.
There is no use of a largest local depth multiplied
by the total number of branches.

The canonical invariant lattices of the restricted
linear systems, together with their actual semilinear
deck comparison, give \(\nu\) coefficient phases
\(E_0,\ldots,E_{\nu-1}\) on \(V\), with
\(E_\nu=\tau^*E_0\), and maps
\[
 F_V^*E_j\longrightarrow E_{j+1},\qquad
 E_{j+1}(-B_V)\longrightarrow F_V^*E_j,
              \qquad 0\le j<\nu,
                                                        \tag{D.GS.12}
\]
whose two products are the Cartier inclusion.
Functoriality of the invariant lattice under field
and coefficient automorphisms gives the last
comparison on the entire proper curve.
The cyclic slope argument of (D.GH.5) works without
change because \(\tau\) preserves slopes. Thus
\[
             -\frac{\Delta_V}{p-1}
                \le\mu_{\min}(E_j)
                \le\mu_{\max}(E_j)\le0.
                                                        \tag{D.GS.13}
\]

Choose a line \(H_0\) on \(X\) of integer degree
\(h_0=\lceil C_p A\rceil\), with \(C_p\) fixed
sufficiently large depending only on \(p\), and put
\(H=\pi^*H_0\) on \(V\).
Equations (D.GS.7) and (D.GS.11) permit this choice
with
\[
 h=\deg H=ah_0=O_p(aA),\qquad
 h-\Delta_V/(p-1)\ge2g(V)+4.
                                                        \tag{D.GS.14}
\]
The line \(H\) has its actual deck linearization.
Each \(E_j^+=E_j\otimes H\) is generated with vanishing
\(H^1\), and the determinant \(L_j=\det E_j^+\) satisfies
\(r(2g(V)+4)\le\deg L_j\le rh\).
Both \(L_j\) and \(L_j^{-1}H^{3r}\) are generated
with vanishing \(H^1\).

Use exactly the fixed partial-resolution terms
\[
 P_0=\mathcal O_V^{r+1},\qquad
 P_i=(H^{-3ir})^{\oplus2}\quad(1\le i\le7)
                                                        \tag{D.GS.15}
\]
of the preceding numerical curve lemma.
These are deck-invariant fixed bundles.
There are only \(\nu\) independent complexes.
The terminal complex is the actual \(\tau\)-pullback
of the initial complex, not a new independent object.
Impose the differential equations, the open unit-minor
rank conditions, and the forward and inverse equations
of (D.GS.12). Length-seven source and length-three
target complexes, with the top-generator identity
test, express actual bundle maps and their products
over every coefficient algebra. No determinant
line is enumerated or made a background parameter.

The estimate \(h^0+h^1\le2g(V)+2+|\deg|\) for a
line gives the following three Hom costs, including
the twists in (D.GS.12) and one-step Frobenius:
top--top blocks have multiplicity \(O(r^2)\) and
line degree \(O_p(aA)\); mixed blocks have
multiplicity \(O(r)\) and degree \(O_p(aAr)\);
bottom--bottom blocks have bounded multiplicity
and degree \(O_p(aAr)\).
Consequently each phase, including its exact
map and boundary witnesses, has \(O_p(aAr^2)\) coordinates.
The fixed-width two-affine Cech contractions and
finite homological perturbation formulas of
Lemma~\ref{D:curve-fixed-local-growth-numerical}
commute with arbitrary coefficient-base change.
Their polynomial degrees are bounded by a constant
depending only on \(p\). Pullback by the fixed
automorphism \(\tau\) is linear on these fixed
coefficient spaces and does not introduce a new
polynomial iteration. One may enlarge the fixed
cochain comparison diagram by this one pullback;
its internal width is still bounded independently
of \(a,f,g,r\).

Here is the exact local condition in this skew
presentation. For \(y\in P_V\), put
\[
 N_y=\left\lfloor\frac{c_y}{p-1}\right\rfloor+1,
 \qquad L_y=N_y+c_y,
 \qquad \sum_y L_y=O_p(aA).
                                                        \tag{D.GS.16}
\]
These lengths are deck-invariant.
The fixed full local semilinear lift at \(x\)
supplies canonical lattices and actual comparison
maps on all the branches \(y\mid x\), including
their deck permutation. Choose bases for these
fixed reference models and their \(\nu\) phases.
Use the quotient-frame matrices of (D.CL.9) at
order \(L_y\), and impose the finite unit
comparison equations with those reference models.
At the last arrow the equation uses the actual
\(\tau\)-pullback reference frame. Substitution
by the fixed power series for \(\tau\) is a fixed
linear operation on the truncated coefficient
spaces; even wild deck action does not increase
the polynomial degree in the atlas variables.

For clarity, the exact contraction argument is
the one in Lemma~\ref{D:curve-integral-unit-jet-comparison}
on the product of these local discs and phase
indices. Transport by a forward unit arrow and
the inverse reference arrow sends an error of
order \(j\) to order at least \(pj-c_y\).
After the last phase apply \(\tau^{-1*}\) and
the corresponding branch permutation.
Deck pullback preserves valuations and the
integers \(N_y,c_y\), so it preserves this
improvement. Since \(pN_y-c_y>N_y\), the
successive corrections converge over arbitrary
coefficient rings, uniquely retaining the
specified \(N_y\)-jet. Products and inverses
have the same uniqueness property. These finite
equations therefore give exact integral local
skew-unit comparisons, not just the linear
restriction type after the field cover.
They use \(O_p(\nu r^2\sum_y L_y)\)
coordinates and equations of bounded degree.

Every geometric point is consequently an actual
skew unit bundle with the fixed full local types.
Conversely every counted linear-semilinear
representation yields these canonical bundles
and is covered by the presentation.
Composing the \(\nu\) arrows gives
\(F_{q_0}^*E_0\simeq\tau^*E_0\) on \(V^\circ\).
The exact frame equivalence in
Lemma~\ref{D:twisted-deck-Frobenius-semilinear-model}
recovers the semilinear local system on \(X^\circ\).
Iterating with the actual deck coherence gives
the \(q\)-unit and its finite etale frames.
No additional order-\(a\) matrix identity is
imposed, and no \(q_0\)-power or \(q\)-power
equation in the atlas variables is expanded.

The specified-Artinian Isom argument also retains
the actual skew data. A unit isomorphism lifts
uniquely on \(X^\circ\) by nilpotent invariance
of finite etale objects. Pull it back to
\(V^\circ\). The fixed local canonical skew
models and the preceding integral comparison
extend it over each completed disc, and
faithful flatness of completion extends it
over \(V\). Deck compatibility holds by its
validity on the dense open and uniqueness.
The finite-presentation proper-curve Isom is
therefore etale over the parameter space,
and distinct representation classes have
disjoint open class images. This uses the
full local types, not merely their different
bounds.

The total number of affine variables is at most
\[
       C_p\nu a A r^2=C_p f A r^2,
                                                        \tag{D.GS.17}
\]
and the polynomial degrees are at most a constant
\(D_p\). This is the exact cancellation of the
field-cover degree against the number of phases.
The component bound, with its prefactor, gives
\[
 (C_p f A r^2+1)\max\{2,D_p\}^{C_p f A r^2}
                              \le p^{C'_p f A r^2}.
                                                        \tag{D.GS.18}
\]
The open rank locus adds no irreducible components.
Multiplying by the at most \((q-1)^{|P|}\) local
linear-lift tuples proves (D.GS.2).
Counting linear-semilinear classes bounds the
projective-semilinear classes because the lift
map is surjective, with conjugation respected.
\end{proof}

\begin{theorem}[Proved: global hereditary semilinear growth with bounded local periods]
\label{D:global-finite-puncture-hereditary-semilinear-local-period}
Let \(S,T,U,\Gamma,F,R_s\) be the fixed data of
Theorem~\ref{D:global-finite-puncture-hereditary-bounded-local-period},
and fix \(B_*\ge1\).
Let \(H\subset\Gamma\) have index \(t\), let
\(q=p^f\), and fix a continuous field character
\(\chi:H\to\operatorname{Gal}(\mathbb F_q/\mathbb F_p)\).
For each connected factor \(U_{s,i}\) of the
degree-\(t\) pullback of the fixed puncture
\(U_{R_s}\), let \(\chi_{s,i}\) be the induced
local field character, and put
\[
 K_{s,i}=\ker\bigl(\chi_{s,i}:
                \pi_1(U_{s,i})\to
                    \operatorname{Gal}(\mathbb F_q/\mathbb F_p)\bigr).
                                                        \tag{D.GS.19}
\]
Count the projective-semilinear homomorphisms
\(\rho:H\to\operatorname{PGL}_r(\mathbb F_q)\rtimes
\operatorname{Gal}(\mathbb F_q/\mathbb F_p)\)
with field projection \(\chi\) such that, for
every \(s,i\), the ordinary projective restriction
to \(K_{s,i}\) has geometric Brauer period at most
\(B_*\).
There is a constant \(A^{\rm sl}_{S,T,B_*}\),
independent of \(t,r,f,\chi\), for which their
\(\operatorname{PGL}_r(\mathbb F_q)\)-conjugacy
classes number at most
\[
                    q^{A^{\rm sl}_{S,T,B_*}t r^2}.
                                                        \tag{D.GS.20}
\]
It suffices to impose the local condition over
\(s\in T\). It is automatic if
\(\gcd(r,q-1)\le B_*\), or if the ordinary
geometric Brauer period on the global field-kernel
cover of \(U_H\) is at most \(B_*\).
There is no tameness or bounded-order assumption
on the field character or on the original cover.
No ordinary geometric Brauer class is assigned
on \(U_H\) to the nonconstant scalar band.
\end{theorem}

\begin{proof}
Use the one fixed detecting curve
\(C^\circ=C-P\) over \(\Omega\) and its connected
degree-\(t\) pullback \(C_H^\circ\), as in the
ordinary global theorem. Its fundamental group
surjects onto \(H\), so its field character has
the same image as \(\chi\). The field cover used
by the preceding lemma is precisely the
pullback of the actual global field-kernel
cover; it is connected. Thus all cover and
deck choices in that lemma have their required
geometric origin.

Every finite projective-semilinear representation
of \(H\), including its field character, has an
open normal core in \(\Gamma\). The core quotient
simultaneously trivializes that representation
and the degree-\(t\) coset cover. The original
detector therefore bounds its local normalized
different over \(K_x=\Omega((t_x))\) by \(b_x+1\).
For each \(y\in P_H\) above \(x\), exact
transitivity gives the same nonnegative budget
as (D.GH.18):
\[
 d_y=e_y(b_x+1)-\delta(K_y/K_x),\qquad
 \sum_{y\in P_H}d_y=O_{S,T}(t).
                                                        \tag{D.GS.21}
\]
This bounds the splitting field of the full
local projective-semilinear representation,
not only its linear restriction.
The original-cover Riemann--Hurwitz calculation
also gives
\[
           1+g(C_H)+|P_H|+\sum_y d_y=O_{S,T}(t).
                                                        \tag{D.GS.22}
\]
The subsequent field cover and its \(\nu\) phases
are accounted for inside the numerical lemma;
no extra field-image factor is multiplied into
(D.GS.22).

For each fixed complete surface source \(R_s\),
write \(n_{s,i}\) for the degrees of its connected
pullback factors. Then \(\sum_i n_{s,i}=t\).
The prescribed \(\chi\) fixes every
\(\chi_{s,i}\), including its actual image
order, which can differ from the global one.
The period hypothesis is exactly the one in
Theorem~\ref{D:bounded-period-semilinear-hereditary-projective-growth}
for the factor \(U_{s,i}\) and that local
field character. It supplies at most
\[
 \prod_{s\in F}\prod_i
      q^{A^{\rm ps}_{R_s,B_*}n_{s,i}r^2}
       =q^{t r^2\sum_{s\in F}A^{\rm ps}_{R_s,B_*}}
                                                        \tag{D.GS.23}
\]
tuples of full local projective-semilinear classes.

The formal-branch maps in the fixed detector
pull back to the appropriate \(U_{s,i}\).
One factor representation determines all
the full branch types that use it.
Choose compatible paths once for this fixed
cover and prescribed field character; the
induced field actions and local splitting
fields are then exactly the ones just counted.
As in the ordinary transfer, every surface
factor is counted once, regardless of how
many curve branches restrict through it.
No injectivity of the local group map to
\(\Gamma\) is needed.

For each tuple apply
Lemma~\ref{D:curve-fixed-local-semilinear-growth-numerical}
on \(C_H,P_H\), with the pulled-back character
and the bounds (D.GS.21).
Equation (D.GS.22) bounds its exponent by
\(O_{S,T,p}(t r^2)\).
Multiplication by (D.GS.23), followed by the
conjugacy-class injection induced by
\(\pi_1(C_H^\circ)\twoheadrightarrow H\),
proves (D.GS.20).

At added sources \(s\notin T\) the map to
\(\Gamma\) is trivial, so both the field
character and representation restrictions
are trivial. A global period bound on the
field-kernel cover restricts to each local
field-kernel component: that pullback is a
disjoint union of copies of the connected
cover determined by \(K_{s,i}\).
Finally, choose set-theoretic semilinear lifts
of the group elements in the finite image of \(\rho\).
The determinants of their matrix parts show
that the full twisted scalar obstruction is
killed by \(r\); its coefficient exponent
also kills it by \(q-1\).
After restriction to a field kernel its
ordinary geometric image has period dividing
\(\gcd(r,q-1)\).
These prove both sufficient conditions while
retaining the nonconstant-band distinction.
\end{proof}

\begin{remark}[Scope of the global semilinear closure]
The local periods in the theorem are measured
only after the actual local field-kernel covers.
Neither arbitrary geometric period nor a
hereditary theorem for a divisorial boundary
is asserted here. The affine-curve twisted
cohomology input, the downstairs local lift
count, the \(\nu a=f\) phase cancellation,
and the once-per-surface-factor comparison
are all needed. This is a global closure of
the same bounded-period hereditary enlargement;
it does not by itself prove finite generation
of the full global group.
\end{remark}


% Proved after four complete actual-file audits. Exact form-model recognition within the same hereditary crown reduction.
\subsection{Symplectic and orthogonal crowns over arbitrary quotients}

\begin{corollary}[Proved: defining-characteristic symplectic and orthogonal crowns]
\label{D:arbitrary-quotient-symplectic-orthogonal-crowns}
Let $R$ be a complete excellent normal local surface ring
with algebraically closed coefficient and residue field
of characteristic $p>0$, and put
$\Gamma=\pi_1(\operatorname{Spec}R\setminus\{\mathfrak m\})$.
Write $A^{\rm ps}_{R,2}$ for the constant of
Theorem~\ref{D:bounded-period-semilinear-hereditary-projective-growth}
with period bound two. Let $S$ be one of the nonabelian
simple groups in the following table, where $q=p^f$.
The symbols $d$ and $Q$ specify the actual faithful
projective-semilinear model used in the proof.
\[
\begin{array}{c|c|c|c}
 S\text{ and range}&d&Q&[\operatorname{Aut}(S):\mathcal A_0]\\ \hline
 \operatorname{PSp}_{2\ell}(q),\ \ell\ge2,
       (\ell,q)\ne(2,2)
       &2\ell&q&2\ (\ell=2,\ p=2),\ 1\ \text{otherwise}\\
 \operatorname{P}\Omega_{2\ell+1}(q),\ \ell\ge3,
       q\text{ odd}
       &2\ell+1&q&1\\
 \operatorname{P}\Omega^+_{2\ell}(q),\ \ell\ge4
       &2\ell&q&6\ (\ell=4),\ 2\ (\ell\ge5)\\
 \operatorname{P}\Omega^-_{2\ell}(q),\ \ell\ge4
       &2\ell&q^2&1
\end{array}                                                    \tag{D.SO.1}
\]
Here $\mathcal A_0$ is the normal inner-diagonal-field
subgroup in the root-coordinate model of
Corollary~\ref{D:classical-defining-characteristic-crowns}.

Let $L$ be any finite monolithic primitive group with
socle $N=S^t$, put $P=L/N$, and fix any continuous
epimorphism $\beta:\Gamma\to P$.
If $\Gamma\twoheadrightarrow L_k$ is a crown quotient
over $\beta$, then
\[
 \begin{aligned}
 k&\le |N|^{c_R},\\
 d(L_k)&\le\max\{2,d(P),c_R+2\},\\
 c_R&=\max\{1,\lceil96A^{\rm ps}_{R,2}+21\rceil\}.
 \end{aligned}                                                \tag{D.SO.2}
\]
The constant is independent of $q$, classical rank,
$t$, the field-image orders and the prescribed
permutation quotient. No radial-tameness or
resolution-extension condition on that quotient
is imposed. The term $d(P)$ is retained.

Odd-dimensional orthogonal groups in characteristic
two, when denoted separately, are included through
their symplectic identification in the same range.
The table makes no assertion about omitted low-rank
repetitions or nonsimple exceptions.
\end{corollary}

\begin{proof}
We first identify the exact prior recognition input.
The proof of
Corollary~\ref{D:classical-defining-characteristic-crowns},
especially (D.CA.5)--(D.CA.9), establishes for every
row of (D.SO.1) an actual faithful embedding
\[
 \mathcal A_0\lhook\joinrel\longrightarrow
 \operatorname{PGL}_d(\mathbb F_Q)\rtimes
                   \operatorname{Gal}(\mathbb F_Q/\mathbb F_p),
 \qquad S\subset\operatorname{PGL}_d(\mathbb F_Q),
 \qquad d^2\log Q\le16\log|S|.                         \tag{D.SO.3}
\]
Its normal-subgroup and exact graph-index assertions
are the known theorem of
\citet[Section~2, statements~3.2--3.6 and their
written proofs in Sections~4--7, pp.~607--614]
{Steinberg1960ClassicalAutomorphismsB}.
The earlier proof constructs the natural projective
diagonal matrices, proves the scalar-commutant
faithfulness, treats the twisted coefficient field,
and checks the quadratic order inequalities.
Only these recognition assertions are used here;
the radial hypothesis of the earlier counting
corollary is not used as an input to a new count.

We verify that these very embeddings satisfy the
form condition of
Corollary~\ref{D:bounded-period-arbitrary-quotient-projective-semilinear-crowns}.
This verification concerns the full inverse image
of $\mathcal A_0$ in $\Gamma\mathrm L_d(\mathbb F_Q)$.
The forms below have coefficients in $\mathbb F_p$.
The matrix definitions and quadratic-form convention
are those of \citet[pp.~54--55, 68--70 and 136--139]
{Taylor1992ClassicalB}; in characteristic two we use
the quadratic form itself, not just its polar form.

For the symplectic row, order a hyperbolic basis as
$(e_1,\ldots,e_\ell,f_1,\ldots,f_\ell)$ and use
\[
 J=\begin{pmatrix}0&I\cr-I&0\end{pmatrix},\qquad
 B(v,w)=v^{\mathsf t}Jw.
                                                               \tag{D.SO.4}
\]
The matrices representing $S$ are isometries of $B$.
The diagonal representatives constructed in (D.CA.6)
have the form
\[
 D=\operatorname{diag}(a_1,\ldots,a_\ell,
                         c/a_1,\ldots,c/a_\ell),
             \qquad D^{\mathsf t}JD=cJ.                \tag{D.SO.5}
\]
They realize the entire diagonal automorphism group
by the already checked simple-root values, with no
square-root extraction. Thus the full inner-diagonal
matrix group consists of symplectic similitudes.
Entrywise field automorphisms preserve $J$, since
its entries lie in the prime field. Its field
extension is consequently a group of actual
semisimilitudes of this same $B$.

For split even orthogonal type, use the prime-field
quadratic form
\[
                  Q_+(x,y)=\sum_{i=1}^{\ell}x_i y_i.
                                                               \tag{D.SO.6}
\]
The natural root matrices representing the inner
group preserve $Q_+$. The same diagonal matrix
(D.SO.5), with the type-$D$ simple-root values
specified in (D.CA.6), satisfies
$Q_+(Dv)=cQ_+(v)$. These identities are polynomial
identities in every characteristic, including two.
Entrywise field automorphisms also preserve this
prime-field quadratic form semilinearly.

For odd orthogonal type in odd characteristic,
use
\[
                  Q_0(x,z,y)=z^2+\sum_{i=1}^{\ell}x_i y_i.
                                                               \tag{D.SO.7}
\]
The type-$B$ diagonal representatives in the prior
recognition proof are
\[
 \operatorname{diag}(a_1,\ldots,a_\ell,1,
                         a_1^{-1},\ldots,a_\ell^{-1}).
\]
They preserve $Q_0$, as do the natural root matrices.
Again all entrywise field automorphisms preserve
the displayed prime-field form semilinearly.

For the minus orthogonal row we retain exactly the
earlier $K=\mathbb F_{q^2}$ root-coordinate model.
In Steinberg's Section~2, the twisted group of
type ${}^2D_\ell$ is a subgroup of the split
type-$D_\ell$ matrix group over $K$, defined by
the commuting diagram/field condition. Its
natural matrices therefore preserve the split
quadratic form (D.SO.6) over $K$.
The self-conjugate diagonal characters are realized
by the same type-$D$ projective torus construction
from the earlier proof; their representatives
satisfy $Q_+(Dv)=cQ_+(v)$ over $K$.
Entrywise $p$-power automorphisms preserve $Q_+$
and normalize this twisted group. We keep their
actual subgroup of order $2f$ over $K$, so there
is no assertion of a section between two Galois
groups or of a model over the smaller field.

In each case the full inverse image of the
projective-semilinear subgroup is generated by
these actual representatives and scalar matrices.
Every scalar $uI$ has form multiplier $u^2$.
The full inverse image thus consists of
semisimilitudes of one fixed nonzero bilinear
or quadratic form, as required. The crossed
multiplier calculation (D.BS.6)--(D.BS.7) applies:
the full twisted scalar obstruction of every
stabilizer lift is killed by two. In particular,
after its actual field-kernel cover the ordinary
geometric Brauer period is at most two.
This conclusion is uniform in the field-image
order and does not depend on the size of a chosen
central covering group.

Let $\mathcal B\le\operatorname{Aut}(S)$ be the
image of the stabilizer of a chosen simple factor
in $L$, and put
$\mathcal B_0=\mathcal B\cap\mathcal A_0$.
It contains $S$, is normal in $\mathcal B$, and
has index at most six by (D.SO.1).
For split $D_4$ the quotient includes the full
triality group of order six. For $C_2$ in
characteristic two it has order two; we do not
assume that a graph representative is an
involution rather than having square a field
automorphism. No graph automorphism is required
to act on the natural projective space.
The field and diagonal subgroups used in
(D.SO.3) are the already recognized normal
subgroup, rather than arbitrary selected
representatives of outer classes.

Apply
Corollary~\ref{D:bounded-period-arbitrary-quotient-projective-semilinear-crowns}
with $B_*=2$, $b_0=6$ and $\kappa=16$.
Its period condition is automatic by the form
calculation, and its constant is
\[
 \max\{1,\lceil16(6A^{\rm ps}_{R,2}+1)+5\rceil\}
        =\max\{1,\lceil96A^{\rm ps}_{R,2}+21\rceil\}.
\]
This gives both conclusions in (D.SO.2) for the
arbitrary prescribed quotient $\beta$.
\end{proof}

\begin{remark}[Exact family scope]
The conclusion concerns the stated symplectic and
orthogonal natural models. In a unitary Hermitian
model a scalar $u$ has multiplier $u^{q+1}$, so
the square-multiplier argument does not apply.
No unbounded-centre type-$A$ recognition is
asserted here. The subgroup-index and order
comparisons are supplied by the earlier exact
classical recognition, and the new application
is a closure of the same bounded-period
hereditary crown reduction. The displayed
$d(P)$ term prevents an unproved assertion of
finite generation of the entire local group.
\end{remark}


% Proved after four complete actual-file audits. Global closure of the same bounded-period hereditary crown reduction.
\subsection{Global symplectic and orthogonal crowns over arbitrary quotients}

\begin{corollary}[Proved: global symplectic and orthogonal crown bounds]
\label{D:global-arbitrary-quotient-symplectic-orthogonal-crowns}
Let $X$ be a fixed normal proper connected surface over an
algebraically closed field of characteristic $p>0$, let
$T\subset X$ be finite, and put $\Gamma=\pi_1(X-T)$.
Let $S$ be any of the nonabelian simple groups, in precisely
the ranges and actual coefficient-field models of
Corollary~\ref{D:arbitrary-quotient-symplectic-orthogonal-crowns}.
Write $d,Q$ for its model dimension and field cardinality.
Let $L$ be a finite monolithic primitive group with socle
$N=S^t$, put $P=L/N$, and prescribe a continuous epimorphism
$\beta:\Gamma\to P$.

Let $A^{\rm sl}_{X,T,2}$ be the constant of
Theorem~\ref{D:global-finite-puncture-hereditary-semilinear-local-period}
with period bound two. Every crown quotient
$\Gamma\twoheadrightarrow L_k$ over $\beta$ satisfies
\[
 \begin{aligned}
 k&\le |N|^{c_{X,T}},\\
 d(L_k)&\le\max\{2,d(P),c_{X,T}+2\},\\
 c_{X,T}&=\max\{1,\lceil96A^{\rm sl}_{X,T,2}+21\rceil\}.
 \end{aligned}                                               \tag{D.GO.1}
\]
There is no tameness, resolution-extension or bounded
field-image assumption. The surface and finite omitted
set remain fixed; the constant is independent of the
classical rank, defining field, $t$ and prescribed
permutation quotient. The term $d(P)$ is retained.
\end{corollary}

\begin{proof}
Use the actual recognition proved in
Corollary~\ref{D:arbitrary-quotient-symplectic-orthogonal-crowns}.
Its normal inner-diagonal-field subgroup
$\mathcal A_0\subset\operatorname{Aut}(S)$ has index
at most six and an actual faithful embedding
\[
 \begin{gathered}
 \mathcal A_0\lhook\joinrel\longrightarrow
 \operatorname{PGL}_d(\mathbb F_Q)\rtimes
                  \operatorname{Gal}(\mathbb F_Q/\mathbb F_p),\\
 S\subset\operatorname{PGL}_d(\mathbb F_Q),
 \qquad d^2\log Q\le16\log|S|.
 \end{gathered}                                             \tag{D.GO.2}
\]
Its full inverse image consists of semisimilitudes
of one fixed nonzero prime-field bilinear or quadratic
form, with scalar multiplier $u^2$.
In particular the minus orthogonal row keeps
$Q=q^2$ and its actual field subgroup, exactly as in
that recognition proof. No additional family or
smaller-field model is introduced here.

Choose a simple socle factor. Let $H\subset\Gamma$
be the inverse image under $\beta$ of its point
stabilizer, so $[\Gamma:H]=t$. Let
$\mathcal B\subset\operatorname{Aut}(S)$ be its
stabilizer image in $L$, and let
$\eta:H\to\mathcal B/S$ be the fixed outer map.
Set
\[
 \mathcal B_0=\mathcal B\cap\mathcal A_0,
 \qquad H_0=\eta^{-1}(\mathcal B_0/S),
 \qquad a=[H:H_0]\le6 .                                  \tag{D.GO.3}
\]
Since $S$ is perfect, the field projection on
$\mathcal B_0$ factors through $\mathcal B_0/S$.
Thus $\eta$ fixes a single field character $\chi_0$
on $H_0$ for every stabilizer lift under consideration.
The connected cover corresponding to $H_0$ has degree
$at$ over $X-T$.

We verify the precise hypothesis of the global
hereditary theorem on this cover. For a stabilizer
lift $\theta:H\to\mathcal B$, choose set-theoretic
semilinear matrix lifts $\widetilde\theta(g)$ of
its restriction to $H_0$, and write
\[
 \widetilde\theta(g)\widetilde\theta(h)
       =c(g,h)\widetilde\theta(gh).
\]
The form multiplier $\lambda$ has the exact crossed
multiplicativity and scalar-square properties from
(D.BS.6). Consequently
\[
 c(g,h)^2=
 \frac{\lambda(\widetilde\theta(g))\,
       \chi_0(g)(\lambda(\widetilde\theta(h)))}
      {\lambda(\widetilde\theta(gh))} .                   \tag{D.GO.4}
\]
The full obstruction class in
$H^2(H_0,\mathbb F_Q^\times(\chi_0))$ is therefore
killed by two. This conclusion is obtained before
passing to any geometric Brauer group.

Use the fixed complete local sources of the global
hereditary theorem. For each connected factor of
the pullback of such a source to the degree-$at$
cover, restrict this class to its fundamental group,
and then to the kernel of the actual restricted
field character. Functorial restriction preserves
(D.GO.4). On that kernel the scalar band is constant,
so its ordinary coefficient class and its geometric
Brauer image are killed by two.
This is exactly the local period condition of
Theorem~\ref{D:global-finite-puncture-hereditary-semilinear-local-period}.
It does not require the restricted field image to
have the same order as the global image, any
injectivity of a local fundamental-group map,
or any comparison of Brauer groups in the opposite
direction. No ordinary Brauer class is assigned to
the nonconstant band on the degree-$at$ cover itself.

That theorem gives at most
$Q^{A^{\rm sl}_{X,T,2}atd^2}$ ambient projective
conjugacy classes for the restricted maps on $H_0$.
The exact nonsplit transport of
Lemma~\ref{D:nonsplit-automorphism-wreath-lifts}
and the orbit and extension argument of
Corollary~\ref{D:bounded-period-arbitrary-quotient-projective-semilinear-crowns}
are purely profinite and apply to the present
$\Gamma$. We spell out the resulting numerical
bound. An ambient conjugacy orbit contains at most
$|\operatorname{PGL}_d(\mathbb F_Q)|\le Q^{d^2}$
actual homomorphisms, hence at most that many
$S$-classes of $\mathcal B_0$-valued restrictions.
Using (D.GO.2), their number is at most
\[
 Q^{(A^{\rm sl}_{X,T,2}at+1)d^2}
       \le |S|^{16(A^{\rm sl}_{X,T,2}at+1)}.             \tag{D.GO.5}
\]
Fix an actual representative of each restricted
$S$-class. Conjugating by $S$ normalizes every full
lift in that class to this restriction and preserves
its prescribed outer map. The images of $a-1$
nonidentity coset representatives of $H_0$ in $H$
then have at most $|S|^{a-1}$ choices, because each
lies in one prescribed $S$-coset. The restriction
and these images determine the whole homomorphism.
This argument uses actual $S$-normalization, not
an action of the whole stabilizer on the ambient
projective space.

Surjectivity onto the crown power forces its
coordinate lifts to represent distinct stabilizer
$S$-classes, by the same diagonal-graph argument
in the exact transport lemma. Hence
\[
 \begin{aligned}
 k&\le |S|^{16(A^{\rm sl}_{X,T,2}at+1)+a-1}\\
  &\le |S|^{t(96A^{\rm sl}_{X,T,2}+21)}
   \le |N|^{c_{X,T}}.
 \end{aligned}                                              \tag{D.GO.6}
\]
The exact nonabelian crown inequality used in the
preceding crown corollary, from the proof of
\citet[Proposition~2.6(ii), equations~(2.1)--(2.2),
pp.~11--13]{LucchiniMarionTracey2020D}, then gives
$d(L_k)\le\max\{2,d(P),c_{X,T}+2\}$.
This establishes (D.GO.1) without bounding $d(P)$.
\end{proof}

\begin{remark}[Scope of the global recognition closure]
This applies only to a fixed normal proper surface
with a finite omitted set. No hereditary theorem
for a divisorial boundary is used or asserted.
The named simple families and their low-rank
exclusions are exactly those of the local
recognition corollary. In particular the unitary
Hermitian scalar exponent $q+1$ and unbounded-centre
type $A$ are not included. This is a global closure
of the same bounded-period hereditary crown
reduction, not a separate movement or a finite-
generation theorem for $\Gamma$.
\end{remark}


% Proved after four complete actual-file/source audits. Explicit failure of a uniform algebraic period shortcut; no geometric realization or main counting movement.
\subsection{An unbounded unitary scalar obstruction}

\begin{example}[Proved: no uniform scalar exponent from a unitary form]
\label{D:unitary-unbounded-scalar-obstruction}
Fix a prime $p$. Let $q$ be a power of $p$, and let
$m\ge3$ be odd with $m\mid q+1$.
Put $K=\mathbb F_{q^2}$ and let $\overline a=a^q$.
There is an injective homomorphism
\[
 \rho_m:(\mathbb Z/m)^2
       \lhook\joinrel\longrightarrow\operatorname{PSU}_m(q)
       \lhook\joinrel\longrightarrow\operatorname{PGL}_m(K)
                                                               \tag{D.UO.1}
\]
whose pullback of the scalar extension
\[
 1\longrightarrow K^\times\longrightarrow
 \operatorname{GL}_m(K)\longrightarrow
 \operatorname{PGL}_m(K)\longrightarrow1
                                                               \tag{D.UO.2}
\]
has class of exact order $m$ in
$H^2((\mathbb Z/m)^2,K^\times)$, with trivial
coefficient action. The same class is obtained by
restriction of the full projective-semilinear scalar
extension, with trivial field character.

For this fixed $p$, such odd $m$ are unbounded.
Consequently there is no fixed positive exponent
that kills the scalar obstruction of every map
into the natural projective unitary models as
rank and field vary. This remains so within the
nonabelian simple unitary ranges, by taking $m$
sufficiently large.
\end{example}

\begin{proof}
We use the standard unitary definitions and scalar
quotient of \citet[pp.~114--116, including Lemma~10.1
and its complete proof]{Taylor1992ClassicalB}.
In particular the norm-one subgroup of $K^\times$
is cyclic of order $q+1$. Choose $\zeta\in K^\times$
of exact order $m$; then $\zeta^{q+1}=1$.
Give $K^m$ the standard nondegenerate Hermitian form
\[
             H(v,w)=\sum_{j\in\mathbb Z/m}v_j\overline{w_j}.
                                                               \tag{D.UO.3}
\]
On its coordinate basis define the shift and diagonal
matrices
\[
 Xe_j=e_{j+1},\qquad Ye_j=\zeta^j e_j
                         \quad(j\in\mathbb Z/m).       \tag{D.UO.4}
\]
The definition of $Y$ is independent of the integer
representative of $j$. Clearly $X^m=Y^m=I$.
The matrix $X$ permutes an orthonormal basis, while
\[
 H(Yv,Yw)=\sum_j\zeta^j\overline{\zeta}^{\,j}
                              v_j\overline{w_j}
                         =H(v,w).
\]
Thus both matrices preserve the displayed Hermitian
form. Since $m$ is odd, their determinants are
\[
 \det X=(-1)^{m-1}=1,
 \qquad \det Y=\zeta^{m(m-1)/2}=1.                     \tag{D.UO.5}
\]
The last equality uses the integer $(m-1)/2$.
These formulas hold also in characteristic two.
Consequently $X,Y\in\operatorname{SU}_m(q)$.

A calculation on the coordinate basis gives
\[
       XY=\zeta^{-1}YX,
       \qquad XYX^{-1}Y^{-1}=\zeta^{-1}I.              \tag{D.UO.6}
\]
The scalar on the right belongs to the scalar
kernel of $\operatorname{SU}_m(q)\to
\operatorname{PSU}_m(q)$: it has norm one and its
$m$th power is one.
The projective images $\overline X,\overline Y$
therefore commute. If $X^aY^b$ is scalar for
$0\le a,b<m$, its permutation of the coordinate
lines forces $a=0$. The remaining diagonal matrix
has entries $1,\zeta^b,\ldots,\zeta^{(m-1)b}$;
it is scalar only if $\zeta^b=1$, hence $b=0$.
This proves that
$\rho_m(a,b)=\overline X^{\,a}\overline Y^{\,b}$
is the injection (D.UO.1), including its exact
projective kernel assertion.

Choose the set-theoretic lift
$s(a,b)=X^aY^b$ for $0\le a,b<m$.
There are no scalar factors from reducing exponents
modulo $m$, since $X^m=Y^m=I$.
Equation (D.UO.6) gives the scalar cocycle
\[
 \begin{aligned}
 s(a,b)s(a',b')
   &=c\bigl((a,b),(a',b')\bigr)s(a+a',b+b'),\\
 c\bigl((a,b),(a',b')\bigr)&=\zeta^{ba'}.
 \end{aligned}                                              \tag{D.UO.7}
\]
where the arguments of $s$ on the right are taken
modulo $m$. In particular $c^m=1$ pointwise, so
its class $\alpha_m$ is killed by $m$.

Conversely, every multiplicative coboundary for
the trivial action on the abelian group
$(\mathbb Z/m)^2$ is symmetric in its two arguments.
The alternating ratio of (D.UO.7), evaluated on
$g=(1,0)$ and $h=(0,1)$, is
\[
                   \frac{c(g,h)}{c(h,g)}=\zeta^{-1}.
                                                               \tag{D.UO.8}
\]
If $e\alpha_m=0$ for a positive integer $e$,
then $c^e$ is a coboundary and its alternating
ratio is one. Equation (D.UO.8) implies
$\zeta^e=1$, so $m\mid e$.
Together with the upper bound this proves that
$\alpha_m$ has exact order $m$ in the group with
coefficients $K^\times$ itself. No injectivity
from a smaller coefficient group is being assumed.
For the full semilinear extension the field
projection of $\rho_m$ is trivial, so restriction
recovers exactly (D.UO.2) and the same class.

For explicit unbounded choices at fixed $p$, if
$p=2$ take
\[
                       q=2^f,\qquad m=2^f+1
                       \quad(f\ge2).
                                                               \tag{D.UO.9}
\]
If $p$ is odd, take any positive even $f$ and put
\[
                       q=p^f,\qquad m=(p^f+1)/2.
                                                               \tag{D.UO.10}
\]
Here $p^f\equiv1\pmod8$, so $m$ is odd.
In both cases $m\mid q+1$, $p\nmid m$, and
$m$ tends to infinity. We may keep only values
$m\ge5$, so that the unitary groups lie in the
nonabelian simple ranges already specified in
Corollary~\ref{D:classical-defining-characteristic-crowns}.
The resulting exact orders $m$ rule out a uniform
scalar exponent.
\end{proof}

\begin{remark}[What this obstruction does and does not establish]
For the standard Hermitian form a scalar $uI$ has
multiplier $u^{q+1}$. That multiplier identity kills
the full scalar obstruction by $q+1$, and the
example shows that its consequence cannot be
replaced by an absolute exponent solely from the
existence of this natural unitary form.
The construction is an explicit algebraic test of
that shortcut, with no novelty claim.

No homomorphism from one fixed complete-local or
global geometric fundamental group is constructed
here. In particular there is no assertion that
$\alpha_m$ has a geometric Brauer image of order
$m$ on one fixed punctured surface, nor that an
unrestricted hereditary representation bound is
false. The result does not remove the unbounded-
period counting gate and is not a main reduction.
\end{remark}


% Proved after four complete actual-file/source audits. Qualitative anchor refinement only; no hereditary quantitative atlas or main movement.
\subsection{A generating anchor of rank equal to the geometric period}

\begin{proposition}[Proved: a rank-period generator after subdividing root nodes]
\label{D:rank-period-generating-root-anchor}
Let \(R\) be a complete excellent normal local surface algebra
over an algebraically closed coefficient and residue field \(k\)
of characteristic \(p>0\), and put
\(U=\operatorname{Spec}R-\{\mathfrak m\}\).
Let \(n\) be prime to \(p\), and let
\(\alpha\in H^2(U,\mu_n)\) be realized by a positive-rank
scalar-weight-one twisted bundle. Write \(\beta\) for its
geometric Brauer image and \(M\) for the order of \(\beta\).
Fix a regular resolution with SNC exceptional divisor
\(E=\bigcup_i E_i\). Let \(m_i\) be the order of the
generic residue of \(\beta\) on \(E_i\); order one is
allowed. At a node \(x\in E_i\cap E_j\), let \(h_x\)
be the order of its secondary residue.

Blow up every node once, writing \(Y'\to Y\) for the
result. On \(Y'\), give the strict transform of \(E_i\)
root order \(m_i\), and give the new exceptional curve
over \(x\) root order \(h_x\). Omit roots of order one.
Let \(\mathcal X'\) be this regular mixed-root stack.
There is an actual \(\mu_n\)-gerbe
\(\mathcal G'\to\mathcal X'\) restricting to the prescribed
coefficient gerbe of \(\alpha\) on \(U\), and a
scalar-weight-one locally free sheaf \(\mathcal A\) of
rank exactly \(M\) on \(\mathcal G'\) with the following
properties.

\begin{enumerate}
\item \(\mathcal A\) contains every scalar-weight-one
geometric inertia simple at every point. Thus it is a
generating bundle for the weight-one category.
\item If \(\mathcal F\) is a scalar-weight-one locally
free sheaf of rank \(r\), then \(M\mid r\).
Put \(s=r/M\), let \(q:\mathcal G'\to Y'\) be the
coarse map, and set
\[
 \mathcal B=q_*\mathcal End(\mathcal A),\qquad
 \mathcal E=q_*\mathcal Hom(\mathcal A,\mathcal F).
                                                        \tag{D.RA.1}
\]
The ordinary algebra \(\mathcal B\) is locally free of
rank \(M^2\), and \(\mathcal E\) is locally free as a
right \(\mathcal B\)-module of rank \(s\), locally for
the etale topology on \(Y'\).
For this fixed \(\mathcal A\), the corresponding Morita
functors are exact and commute with extension of
coefficient rings. Their restriction identifies
weight-one vector-bundle families of rank \(r\) with
locally free right \(\mathcal B\)-module families of
rank \(s\).
\item If the original SNC model has \(v\) components
and \(e\) nodes, the construction adds \(e\) curves
and replaces its nodes by \(2e\) nodes. It therefore
uses \(O_R(t)\) centres on the controlled SNC models
of degree-\(t\) covers. Pullbacks of the previously
chosen cycles retain their already proved cohomology
and intersection budgets.
\end{enumerate}
The construction of \(\mathcal A\) is qualitative.
It asserts no uniform bound for the cohomology,
generation, or formal pole constants of the
particular order \(\mathcal B\).
\end{proposition}

\begin{proof}
\emph{The known period-index input and the actual coefficient
gerbe.}
Theorem 2.1 of the author preprint of
\citet{C_CTOP2000LocalPeriodIndex} states that a central
division algebra of prime-to-\(p\) period \(M\) over the
fraction field of an excellent henselian two-dimensional
local domain with separably closed residue field is
cyclic of index \(M\). We use that precise local
theorem, not a function-field theorem for algebraic
surfaces. Its complete written proof is on pp.~11--12
of the retained preprint; its unramified vanishing
input is Theorem 1.3 and Corollary 1.9 there.
This period-index result is Known.

The regular puncture has injective Brauer restriction
to its function field \(K\). Hence the generic image
of \(\beta\) has the same period \(M\), and its index
is \(M\). The actual generic \(\mu_n\)-gerbe has a
weight-one twisted vector space of rank \(M\):
over a field the Kummer map has zero Picard term,
and the index is the least such positive rank.
In particular, this statement does not replace
\(\alpha\) by a chosen class with the same Brauer image.
Every weight-one vector space of rank \(r\) over this
generic gerbe has \(M\mid r\).

The residue calculation and the vanishing of the
prime-to-\(p\) Brauer group of the regular proper-local
model give
\[
           M=\operatorname{lcm}_i m_i,\qquad
           h_x\mid m_i,\ m_j.                         \tag{D.RA.2}
\]
These are precisely the assertions already proved in
\ref{D:SNC-general-Brauer-geometry}.
We next compute the new residues before using the
mixed-root extension of that result.

\emph{What one node blowup does.}
Work at a node \(x\) with regular parameters \(u,v\).
In the strictly henselian completed local ring, the
prime-to-\(p\) ramified class has the form
\[
                         c(u,v)_n,\qquad c\in\mathbb Z/n.
                                                        \tag{D.RA.3}
\]
Indeed the two-dimensional residue complex first
writes it as an unramified class plus two
unit-parameter symbols and \(c(u,v)_n\).
The units are \(n\)-th powers in this complete
strictly local ring, and its unramified Brauer class
vanishes. This is the local computation in the
complete proof of the cited period-index theorem.
The order of \(c\) is \(h_x\).

In the blowup chart \(v=uw\), the symbol becomes
\(c(u,uw)_n=c(u,w)_n+c(u,u)_n\).
The second summand is \(c(u,-1)_n=0\), since
\(-1\) is an \(n\)-th power in \(k\).
The exceptional residue is therefore the class of
\(w^{c}\) or its inverse in
\(k(w)^\times/k(w)^{\times n}\); the sign depends only
on the residue convention. Its order is exactly
\(h_x\), since the divisor of \(w\) has coefficients
\(1,-1\). The two new secondary residues again have
order \(h_x\). The old generic residues are unchanged.
Thus the only new root nodes have orders
\[
                    (m_i,h_x)\quad\hbox{or}\quad
                    (h_x,m_j),
 \qquad\hbox{commutator order }h_x.                   \tag{D.RA.4}
\]
If \(h_x=1\), the exceptional curve carries no root;
there is only the single old root inertia at each
intersection point. This case must not be treated
as a node with two old root orders.

The actual mixed-root Azumaya extension and the
Kummer correction by an extending line bundle, proved
in \ref{D:SNC-general-Brauer-geometry}, now give an
actual \(\mu_n\)-gerbe \(\mathcal G'\) restricting to
\(\alpha\). The local Azumaya construction is the
one in \citet[Proposition 3.2.1, complete pp.~10--11]
{Lieblich2011ResidualD}; only its already verified
componentwise-root application is used.
In particular the scalar coefficient band remains
\(\mu_n\), even when \(M<n\).

\emph{Extending the generic rank-\(M\) object.}
The existing algebraic generating bundle on
\(\mathcal G'\), supplied by
\ref{D:SNC-general-Brauer-geometry}, is used here only
as a qualitative coherence tool. Apply its coarse
Hom functor to the generic rank-\(M\) object.
Extend the resulting finite generic module to a
coherent module over the corresponding coherent
coarse algebra: first choose a coherent ordinary
lattice, then take the algebra span.
The exact inverse Morita functor produces a coherent
weight-one extension of the generic object.
Remove its torsion and take its double dual.

The root stack and gerbe have regular two-dimensional
etale scheme charts. A coherent reflexive module on
such a chart is locally free. Etale flat pullback
commutes with the duals, so these hulls descend,
remain scalar-weight one, and retain rank \(M\).
This gives \(\mathcal A\).
The reflexive extension and regular-surface facts
are exactly the full proofs recorded under
\citet[Tag 0AVC, Lemma 15.24.19, and Tag 0B3N,
Lemma 31.13.15]{C_StacksRootAnchorReflexive}.
This construction is made over \(k\); it is not
a claim that relative reflexive hulls preserve
flatness over arbitrary coefficient rings.

\emph{Generic multiplicities are forced.}
On the generic residual gerbe of a component of
root order \(m\), the weight-one inertia characters
form the cyclic torsor attached to its residue
character. That torsor is connected of degree \(m\),
because the residue has exact order \(m\).
Consequently a rank-\(d\) weight-one bundle has
every one of the \(m\) geometric characters with
the same multiplicity \(d/m\).
This also follows directly by Galois permutation
of the character eigenspaces. The exact residual
category used here is
\ref{D:SNC-general-residual-category}.
Multiplicities are locally constant along the
whole component, including its smooth point on
each new node branch.

\emph{At a subdivided node the marginal determines
every simple.}
We give the mixed-order algebra, to avoid inferring
a generating property just from total rank.
At a geometric node of type \((m,h)\) in
(D.RA.4), the central extension of
\(\mu_m\times\mu_h\) by the scalar band has
scalar-weight-one commutator of exact order \(h\).
After scalar normalizations its twisted group
algebra is generated by \(X,Y\) with
\[
        X^m=1,\qquad Y^h=1,\qquad
        XY=\zeta YX,\quad \operatorname{ord}(\zeta)=h.
                                                        \tag{D.RA.5}
\]
The scalar normalizations are made over the
algebraically closed geometric residue field.
They do not assert a splitting of the central
inertia extension itself.

The \(m\) eigencharacters of \(X\) break into
disjoint orbits of size \(h\) under \(Y\).
For each orbit there is exactly one irreducible
module, of dimension \(h\): choose a vector in
one eigenspace, apply \(Y\), and normalize its
cyclic product using \(Y^h=1\).
Different orbits give nonisomorphic modules,
and all modules are semisimple since the inertia
order is prime to \(p\).
Thus there are \(m/h\) simples. Each restricts
to exactly one orbit, once at each of its
\(h\) branch characters. This is the explicit
mixed-order special case of the node algebra
already proved in
\ref{D:SNC-general-residual-category}; compare the
central-extension and commutator computations in
\citet[Lemma 4.3.1 and Proposition 4.3.2, complete
pp.~13--14]{Lieblich2011ResidualD}.

For \(\mathcal A\), every branch character has
multiplicity \(M/m\). Equation (D.RA.5) therefore
forces every node simple to occur with multiplicity
\(M/m>0\). For \(\mathcal F\), the same multiplicity
is \(r/m=s(M/m)\). The other new node type is
identical after interchanging the two branches.
At a point with only one root inertia the same
conclusion is the generic character calculation;
at a point without root inertia it is immediate.
This proves the generating assertion and, more
precisely, that every geometric inertia fibre of
\(\mathcal F\) is isomorphic to that of
\(\mathcal A^{\oplus s}\).

\emph{Fixed-anchor Morita and coefficient bases.}
The preceding inertia assertion can be checked on
strictly local root charts. There the finite tame
inertia representation splits into its character
or simple factors, and taking invariants is exact.
An isomorphism of the two residual representations
lifts to an equivariant map
\(\mathcal A^{\oplus s}\to\mathcal F\);
Nakayama and equal rank make it an isomorphism
near that geometric point. The same argument
applies to coefficient algebras: the fixed finite
inertia has invertible order, so invariants and
the invariant Hom commute with coefficient
extension, and the geometric fibre isomorphism
lifts on the corresponding local frame opens.
No saturation of a varying module is involved.

The scalar centre cancels in
\(\mathcal End(\mathcal A)\); its remaining
equivariance factors through the abelian root
quotient. A homogeneous basis, adjusted by the
unique root monomial in each residue class,
shows that its coarse invariant module is free
of rank \(M^2\) on each local root chart.
This is the same all-ring monomial calculation
as in the existing mixed-root presentation.
The local isomorphism
\(\mathcal F\simeq\mathcal A^{\oplus s}\)
therefore identifies
\(\mathcal E\simeq\mathcal B^{\oplus s}\).

Equivalently, on these charts \(\mathcal A\) is
a projective generator for the weight-one
equivariant category: its Hom is exact and it
contains all residual simples. The usual
evaluation and coevaluation maps give the
Morita equivalence. The maps and their inverse
identities commute with restriction and hence
descend. All the operations are fixed finite
module maps, invariant direct summands, and
tensor products, so the equivalence commutes
with coefficient extension. In the locally
free subcategories its inverse is visibly
locally \(\mathcal A^{\oplus s}\).
This establishes the precise family assertion
in item~2 for the already fixed anchor.

Finally, the original SNC nodes are pairwise
distinct and have no triple points. One blowup
adds one smooth rational component and replaces
that node by two transverse nodes. Old strict
transforms are mutually disjoint, as are new
exceptional curves. Thus the stated centre
and graph counts are exact. The already proved
pullback comparison for the inherited cycles
applies to these ordinary point blowups.
\end{proof}

\begin{remark}[The quantitative question that remains]
The index theorem supplies a generic cyclic splitting
extension \(K(f^{1/M})/K\), but its proof chooses
\(\operatorname{div}(f)\) equal to the ramification
divisor plus an auxiliary divisor avoiding finitely
many points. It does not control that auxiliary
divisor's degrees, and that extension need not be
etale on \(U\). No bounded-geometry splitting cover
is inferred here.

More importantly, \(\mathcal B=q_*\mathcal End(\mathcal A)\)
in (D.RA.1) need not be the canonical full pullback
of its punctured algebra. The same distinction
applies to \(\mathcal E\).
The full-bundle degree and generation theorems
for the canonical ordinary pullbacks do not
automatically transfer to these particular
coarse lattices. Even the proved bounded
divisorial sandwiches do not imply generation
of an intermediate lattice or flatness of a
relative saturation.

A possible whole-cycle presentation would require,
in addition, controlled \(\mathcal B\)-generating
quotients of \(\mathcal E\), short exact syzygies
in that locally free module category, and fixed
Hom complexes with total size \(O_R(tr^2)\).
The local rank identity
\(M^2(r/M)^2=r^2\) is consistent with that target,
but is not a proof of any of those global
requirements. This proposition is an auxiliary
qualitative refinement, not removal of the
unbounded-period hereditary counting gate.
\end{remark}


% Separate post-CP28 Proved auxiliary after four full actual-file/source audits. Qualitative class openness for already effective finite-type charts; no new quantitative object atlas.
\subsection{Exceptional Picard rigidity and qualitative unit-class openness}

\begin{lemma}[Proved: the infinitesimal Picard kernel of a surface resolution]
\label{D:exceptional-infinitesimal-Picard-kernel}
Let $A$ be a complete excellent normal local surface algebra
over an algebraically closed coefficient and residue field $k$.
Let $\pi:Y\to\operatorname{Spec}A$ be a regular resolution,
an isomorphism on $U=\operatorname{Spec}A-\{\mathfrak m\}$,
and let $E$ be its closed-fibre support. Then
\[
 H^0(D,\mathcal O_D(D))=0\quad(D>0\text{ supported on }E),
 \qquad H^1_E(Y,\mathcal O_Y)=0.                       \tag{D.IP.1}
\]
In particular $H^1(Y,\mathcal O_Y)\to H^1(U,\mathcal O_U)$
is injective.

For every finite local Artinian $k$-algebra $T$ with residue
field $k$, the map
\[
 \ker\bigl(\operatorname{Pic}(Y_T)\to\operatorname{Pic}(Y)\bigr)
 \longrightarrow
 \ker\bigl(\operatorname{Pic}(U_T)\to\operatorname{Pic}(U)\bigr)
                                                               \tag{D.IP.2}
\]
is injective. Here the subscript $T$ means scalar extension
over $k$.

More precisely, let $L_T,L'_T$ be line bundles on $Y_T$ and
let $\varphi_U:L_T|_{U_T}\simeq L'_T|_{U_T}$ be specified.
Let $D_0$, an integral exceptional divisor of either sign,
be the divisor of its closed-fibre meromorphic extension.
Then $\varphi_U$ has a unique meromorphic extension on $Y_T$
whose divisor is exactly the constant Cartier divisor
$(D_0)_T$. In particular its pole divisor, and the pole
divisor of its inverse, acquire no extra contribution
depending on the nilpotence order of $T$.
\end{lemma}

\begin{proof}
If $\pi$ is the identity resolution, then $Y=\operatorname{Spec}A$
is regular, there is no positive exceptional divisor, and
$H^1_E(Y,\mathcal O_Y)=H^1_{\mathfrak m}(A)=0$ by depth two.
Also $H^1(Y,\mathcal O_Y)=0$; the Artinian arguments below
apply directly, with $D_0=0$. Otherwise the closed-fibre
support is a divisor, and we prove (D.IP.1) as follows.

We use the negative definiteness of the exceptional
intersection form and the exact Riemann--Roch formula for
line bundles on proper nonreduced exceptional cycles,
as in Proposition~\ref{D:exceptional-canonical-energy}.
Here is a direct proof that retains the nilpotent lengths.

Suppose that $0\ne s\in H^0(D,\mathcal O_D(D))$.
Locally on the regular surface write $D=(e)$ and
$s=f/e\pmod{\mathcal O_Y}$. The regular local ring is a
unique factorization domain. If $g$ is a greatest common
divisor of $e$ and $f$, write $e=ge'$ and $f=gf'$.
Then
\[
              \operatorname{ann}(s)=(e):f=(e'),
              \qquad \gcd(e',f')=1.                  \tag{D.IP.3}
\]
These annihilator ideals glue, so they define an effective
Cartier subcycle $0<B\le D$. The identity $f/e=f'/e'$
shows that $s$ factors through the natural subsheaf
$\mathcal O_B(B)\subset\mathcal O_D(D)$.
The resulting morphism
\[
                  \mathcal O_B\longrightarrow\mathcal O_B(B)
                                                               \tag{D.IP.4}
\]
is injective: locally multiplication by $f'$ is injective
modulo $e'$, by (D.IP.3). At the generic point of every
component of $B$, $f'$ is a unit; thus the map is an
isomorphism on the entire Artinian generic local ring,
not just on its reduced residue field. Its cokernel $Q$
is zero-dimensional. Consequently
\[
  B^2=\chi(\mathcal O_B(B))-\chi(\mathcal O_B)
      =\operatorname{length}_k Q\ge0,                 \tag{D.IP.5}
\]
contrary to negative definiteness. This proves the first
assertion of (D.IP.1), without an inference from Hartogs
alone and without a characteristic-zero vanishing theorem.

Choose any effective Cartier cycle $C$ with full support
$E$. The usual local calculation for a Cartier divisor is
\[
 \mathcal H^0_E(\mathcal O_Y)=0,\qquad
 \mathcal H^1_E(\mathcal O_Y)
       =\varinjlim_{a\ge1}\mathcal O_{aC}(aC).
                                                               \tag{D.IP.6}
\]
It is simply $\mathcal O_Y[1/z]/\mathcal O_Y$ on an affine
chart with equation $z$ for $C$. The local-to-global
sequence for cohomology with support identifies
$H^1_E(Y,\mathcal O_Y)$ with the global sections of the
second sheaf in (D.IP.6). Global sections commute with
this filtered colimit on the quasi-compact separated
scheme $Y$. Each term has zero global sections by the
first assertion of (D.IP.1), proving the second one.
The localization sequence gives the stated $H^1$
injection.

For the Artinian statement, induct over small extensions
$0\to I\to T\to T_0\to0$ with $\mathfrak m_T I=0$.
On both $Y$ and $U$ the kernel of the reduction map on
unit sheaves is the additive sheaf $I\otimes_k\mathcal O$.
Moreover
\[
 H^0(Y_T,\mathcal O_{Y_T})
   =H^0(U_T,\mathcal O_{U_T})=A\otimes_k T.            \tag{D.IP.7}
\]
Normal-surface Hartogs gives the equality over $k$;
finite-dimensional flat scalar extension gives (D.IP.7).
Units lift across the nilpotent quotient
$A\otimes T\to A\otimes T_0$. The unit-sheaf exact
sequences therefore identify the two infinitesimal
Picard kernels with
$I\otimes H^1(Y,\mathcal O_Y)$ and
$I\otimes H^1(U,\mathcal O_U)$, respectively.
The just-proved injection and induction give (D.IP.2).
No normality of $A\otimes_kT$ is being asserted.

Finally put
$N_T=L'_T\otimes L_T^{-1}\otimes\mathcal O_{Y_T}(-D_0)$.
Its closed fibre is trivial, and the prescribed
$\varphi_U$ trivializes its restriction to $U_T$.
By (D.IP.2), $N_T$ is trivial. Compare any global
trivialization with the prescribed one on $U_T$.
Their ratio is a unit of $A\otimes_k T$ by (D.IP.7),
and hence is a unit on $Y_T$. Adjusting by this unit
produces precisely the prescribed meromorphic section
of $L'_T\otimes L_T^{-1}$ with divisor $(D_0)_T$.
It is unique because a Cartier equation of $E$ is a
nonzerodivisor on $Y_T$ and on every invertible sheaf
there. This proves the final assertion.
\end{proof}

\begin{proposition}[Proved: unit-class openness without fixed determinant charts]
\label{D:effective-formal-unit-family-class-openness}
Keep the preceding surface and suppose $\pi$ is projective.
Fix an effective exceptional divisor $H_0$ such that
$\mathcal O_Y(-H_0)$ is relatively ample, and fix an
effective Cartier cycle $C$ with full exceptional support.
Let $S$ be a finite-type $k$-scheme carrying an already
effective formal family of vector bundles $F_a$ on the
completion of $Y\times_k S$ along $E\times_k S$, with a
fixed finite number of phases.

Assume that its unit data belong to either of the following
already established finite-comparison settings:
ordinary Frobenius unit cycles of
Proposition~\ref{D:intrinsic-integral-unit-Hom-jets}, or
the fixed-order Morita unit cycles of
Proposition~\ref{D:generated-Morita-unit-Isom-openness}.
In the latter case the actual orders, balanced
correspondences and linearity relations are retained.
In either setting require the stated fixed Cartier pole
losses and the actual finite-etale fixed-frame comparison
on the puncture. This is a hypothesis about effective
objects and their exact comparison operators, not merely
about a formal tangent complex.

Then the loci in $S$ of isomorphism classes of the
punctured finite unit objects are open and closed.
In particular
\[
 \#\{\text{punctured unit classes occurring on }S\}
                         \le\#\operatorname{Irr}(S). \tag{D.IP.8}
\]
No fixed determinant line is required in the object
chart. No quantitative bound for a generating twist
or for the comparison-Isom scheme is required either.
This proposition does not construct the finite-type
object chart or bound its equations.
\end{proposition}

\begin{proof}
We give the source-generation, determinant and specified
Artinian-lifting arguments, since these are the points
at which a merely pointwise rigidity assertion would
be insufficient.

\emph{Qualitative generation after actual algebraization.}
Work on an affine open $S=\operatorname{Spec}B$.
Write
$A_B=\varprojlim_j(A/\mathfrak m^j)\otimes_kB$.
This is a Noetherian complete ring: a fixed Cohen
presentation of $A$ identifies it with a quotient of
$B[[x_1,\ldots,x_e]]$. Put $Y_B=Y\times_A\operatorname{Spec}A_B$.
The $C$-adic and $\mathfrak m$-adic topologies on this
proper model are cofinal, since their supports agree.
Projective Grothendieck existence therefore algebraizes
the actual formal bundles $F_a$ and their bounded-pole
maps on $Y_B$; use the precise Noetherian complete scope
of \citet[Tag 0885, Lemma 30.24.3, complete proof]
{StacksRelativeKahnB}.

The algebraized sheaves are vector bundles. Indeed their
nonfree loci are closed and miss the closed formal fibre;
properness makes their images closed in $\operatorname{Spec}A_B$.
Every maximal ideal of $A_B$ contains
$\mathfrak m A_B$, so a nonempty such image could not
miss that fibre. Thus there is no nonfree locus.

Ampleness and eventual generation, in the exact form of
\citet[Tag 01PR, Proposition 28.27.13(7)--(8), complete proof]
{StacksCartierParametersRoot}, give an integer $b$ and
finite generating quotients
\[
           \mathcal O_{Y_B}^{N_a}\twoheadrightarrow
                         F_a(-bH_0)                  \tag{D.IP.9}
\]
for all phases. These are actual sheaf quotients.
They remain surjective after every finite Artinian
specialization $B\to T$. Neither flat base change for
$H^0$ nor the ability to lift an arbitrarily prescribed
section is needed. The integer $b$ and the $N_a$ may
depend on this whole finite-type chart. They will not
be inserted as variables in that chart or charged to
its component estimate.

\emph{A bound for every specified Artinian comparison.}
Fix a closed coefficient point $s\in S(k)$ and its
punctured unit object. For each phase let $P_a$ be the
canonical full pullback of the reflexive extension of
the corresponding ordinary punctured bundle. In the
Morita setting this is the ordinary module obtained
from the fixed anchor. It is used only as a fixed
constant target. As in (D.GD.8),
\[
 H^0(U_T,P_a|_{U_T})=M_a\otimes_kT,
 \qquad H^0(U_T,\mathcal O_{U_T})=A\otimes_kT.        \tag{D.IP.10}
\]
Consider a finite local Artinian specialization centred
at $s$ and any specified punctured unit isomorphism to
the constant object at $s$. Compose it with (D.IP.9).
Each resulting punctured section extends to $P_{a,T}$
by (D.IP.10). Relations in the generating quotient
vanish on $U_T$ and hence everywhere in this locally
free target, since an exceptional Cartier equation
remains a nonzerodivisor after scalar extension over
$k$. Thus the comparison extends as
\[
              h_{a,T}:F_{a,T}\longrightarrow P_{a,T}(bH_0).
                                                               \tag{D.IP.11}
\]
This is the same target-only Hartogs argument as in
Proposition~\ref{D:generated-fixed-determinant-unit-openness}.
If order-linearity is required, its two sides agree
on $U_T$ and therefore agree meromorphically everywhere.

Take the determinant of the specified punctured
comparison. Lemma~\ref{D:exceptional-infinitesimal-Picard-kernel}
shows that its exceptional divisor is exactly its
closed-fibre divisor $D_a$, for every nilpotence order.
Different closed-fibre automorphisms of the same
punctured unit object do not change $D_a$: the ratio
of their determinants and its inverse extend by
(D.IP.10), so that ratio is a unit on $Y$.
The adjugate identity now gives
\[
 \operatorname{pole}(h_{a,T})\le bH_0,
 \qquad
 \operatorname{pole}(h_{a,T}^{-1})
                  \le(\operatorname{rk}F_a-1)bH_0+D_a^+.
                                                               \tag{D.IP.12}
\]
Apply the same bound to the constant comparison at $s$.
Since $C$ has full support, some integer $d_s$ bounds
both directions of the resulting comparison between
$F_{a,T}$ and the actual constant lattice $F_{a,s}$
by $d_s C$, for all phases and all such Artinian bases.
This $d_s$ depends on the chart and the closed object,
but not on the nilpotence order or the chosen unit
automorphism at that object. In particular no fixed
determinant line in the parameter chart was used.

\emph{The finite comparison scheme and openness.}
The existing intrinsic comparison construction applies
to the pole lattices
$\mathcal Hom(F_a,F_{a,s})(d_sC)$ and their reverse
Hom lattices. A pole shift by $d_sC$ changes a
one-step Frobenius loss $cC$ to
$(c+(p-1)d_s)C$; for products use the $2d_sC$ lattice.
Choose a finite depth exceeding these fixed losses.
The exact finite-jet equations, the two inverse-product
equations and local contraction then represent the
bounded unit-Isom functor by a finite-presentation
scheme over $S$. In the Morita case retain the
regular $S_x,S_y$ equations and the actual balanced
relations from the cited Morita-Isom proposition.
No raw image of global section modules is substituted
for this functor. This use of the previous finite-Hom
theorems requires no upper bound on the number of
coordinates of the comparison scheme.

Finite-etale fixed frames and their morphisms are
unchanged by nilpotent thickenings
\citep[Tag 0BQB, Lemma 58.8.3, complete proof]
{StacksIntrinsicUnitHomD}. Consequently a specified
unit isomorphism over an Artinian quotient has a
unique punctured lift. Equations (D.IP.11)--(D.IP.12)
place that lift inside the same bound $d_s C$ whenever
the closed point is $s$. Restriction to the quotient
is the prescribed bounded meromorphic map: two such
maps agreeing on the puncture agree in the Cartier
pole lattice. Thus the bounded comparison scheme is
formally etale at every point over $s$ representing
such an isomorphism. The finite-presentation Artinian
criteria already used in
\ref{D:generated-fixed-determinant-unit-openness}
make it etale there. Its etale locus has open image
in $S$, containing $s$, and every point of that
image has the same punctured unit class.

Repeat at every closed point. The union of these
open neighbourhoods for one class is its class
locus on closed points. Loci for different classes
are disjoint: a nonempty intersection of two opens
in a finite-type $k$-scheme would contain a closed
$k$-point. They cover $S$, since a nonempty closed
complement would contain such a point. Each class
locus is therefore both open and closed. Every
irreducible component lies in one locus, proving
(D.IP.8). The construction is Zariski local on $S$,
so the affine-open reduction loses no classes.
\end{proof}

\begin{remark}[The quantitative obligation remains]
The argument begins with an actual effective formal
unit-family on a finite-type parameter scheme.
It cannot be applied to a collection of geometric
objects, a deformation dimension, a generic Morita
rank, or arbitrary independently lifted order maps.
The large Serre twist in (D.IP.9) is used only in
the comparison proof. It does not change the
Frobenius poles or jet depth of the object chart.

For the proposed unbounded-period order-linear route,
one must still construct a finite chart of actual
locally free order-modules and balanced unit maps
with $O_R(ftr^2)$ coordinates and controlled equation
degrees. In particular the rank-period generating
anchor and this openness result do not imply the
needed whole-cycle order-module generation, short
syzygy presentations, or cohomology estimates.
No unbounded-period representation count is asserted.
\end{remark}


% Proved after four complete actual-file/source audits. Explicit geometric-period diagnostic on one fixed complete elliptic cone; no representation-count or simple-quotient conclusion.
\subsection{Unbounded unitary geometric periods on one elliptic cone}

\begin{lemma}[Proved: explicit tame Heisenberg covers of a fixed cone]
\label{D:elliptic-cone-Heisenberg-covers}
Let $k$ be algebraically closed of characteristic $p>0$, let
$(E,O)$ be an elliptic curve over $k$, and put
$N=\mathcal O_E(-3O)$. Let $R$ be the completion at the vertex
of the affine cone for the plane embedding given by
$N^{-1}=\mathcal O_E(3O)$, and put
$U=\operatorname{Spec}R\setminus\{\mathfrak m_R\}$.
This is one fixed complete normal surface germ. Its resolution
has the single smooth exceptional curve $E$, with $E^2=-3$.

For every odd $m$ with $(m,3p)=1$ there is a connected finite
\'etale Galois cover $U_m\to U$ whose group has the presentation
\[
 G_m=\langle x,y,z\mid x^m=y^m=z^m=1,\quad
 z\text{ central},\quad [x,y]=z^3\rangle,
 \qquad |G_m|=m^3.                                      \tag{D.EC.1}
\]
On the corresponding regular resolution cover, the exceptional
curve is again $E$; its map to the original exceptional curve is
$[m]:E\to E$. The generic inertia group is $\langle z\rangle$
of order $m$, and the residue extension has group $E[m]$ of
order $m^2$. After choosing a primitive $m$th root $\zeta\in k$,
the generator $z$ may be chosen to act on a normal parameter
by $z(s)=\zeta s$.
\end{lemma}

\begin{proof}
\emph{Known inputs.}
We use the theta-group construction and its nondegenerate
commutator pairing for an ample line bundle of separable type
from \citet[\S1, pp.~288--295, especially Theorem~1 and its proof]
{D_Mumford1966ThetaCone}.
The multiplication and symmetric-pullback formulas are
\citet[Lemmas~39.9.7--39.9.9 and Proposition~39.9.11]
{D_StacksEllipticThetaCone}.
For an elliptic curve, the identification of $E$ with
$\operatorname{Pic}^0(E)$ through $P\mapsto\mathcal O_E(P-O)$
also follows from the genus-one case of Lemma~44.6.7 of that
reference: the degree-one Abel map has one-point projective-space
fibres. These are cited results; no prime-order assumption is
introduced below.

Write $Y_0=\operatorname{Tot}(N)$ and let $E\subset Y_0$ be its
zero section. The usual blowup charts of a homogeneous plane
cone identify $Y_0$ with the blowup of its vertex. To make the
geometry explicit, if $F(X_0,X_1,X_2)$ is the smooth cubic
equation of $E$, the chart with direction $X_i\ne0$ has equation
$F(1,u,v)=0$ and the free fibre coordinate $t=X_i$.
The other charts glue by the transition functions of
$\mathcal O_E(-1)=N$.
The hypersurface cone is a domain and Cohen--Macaulay; its
puncture is smooth by the same charts. It is therefore normal
by the codimension-one normality criterion. Its excellent
local ring has normal, hence integral, completion
\citep[Lemma~15.53.6, Tag~0C23]{D_StacksEllipticThetaCone}.
Completion preserves the regular proper resolution and its
closed fibre, by Lemmas~54.11.1--54.11.2 of the same reference.
Thus the asserted fixed complete germ and self-intersection
are as stated.

Put $M=N^{\otimes m}$. Since $N$ is symmetric,
\[
        [m]^*N\simeq N^{\otimes m^2}=M^{\otimes m}.
                                                               \tag{D.EC.2}
\]
Fix such an isomorphism $\psi$. Fibrewise $m$th power followed
by $\psi$ defines
\[
 f_m:\operatorname{Tot}(M)\longrightarrow\operatorname{Tot}(N),
 \qquad (P,w)\longmapsto([m]P,\psi_P(w^{\otimes m})).
                                                               \tag{D.EC.3}
\]
It factors as the finite locally free $m$th-power map to
$\operatorname{Tot}(M^{\otimes m})$, then the base change of
the finite \'etale isogeny $[m]$, of degree $m^2$.
Consequently $f_m$ is finite locally free of degree $m^3$,
and is \'etale off the zero section. Its source is smooth
and integral, and its inverse image of $E$ has reduced
support the single curve $E$.

For $u\in E[m]$, translation by $u$ preserves $M$:
under $E\simeq\operatorname{Pic}^0(E)$ the difference
$\tau_u^*M\otimes M^{-1}$ is the point $3mu=0$, up to the
irrelevant sign convention for translation.
A linear lift of this translation to $M$ has an $m$th tensor
power on $M^{\otimes m}$. Compare it, via $\psi$, with the
canonical deck lift on $[m]^*N$. The discrepancy is an
automorphism of a line bundle on the proper connected curve
$E$, and is therefore a scalar in $k^\times$.
Rescaling the lift removes this discrepancy. There are exactly
$m$ such rescalings. These compatible lifts form a group
\[
       1\longrightarrow\mu_m\longrightarrow G_m
         \longrightarrow E[m]\longrightarrow1.        \tag{D.EC.4}
\]
The kernel acts by fibre scalars and is central.
Off the zero section the action is free: a nonzero translation
has no fixed point on the base, while a nontrivial fibre scalar
has no fixed nonzero vector. Its order equals $\deg(f_m)$.
Thus (D.EC.3), off its zero section, is the corresponding
connected $G_m$-torsor.

Here is the exact commutator and exponent calculation, including
characteristic two. Apply the known theta-pairing theorem to
$M_0=\mathcal O_E(mO)$, whose degree $m$ is prime to $p$.
The translation kernel of $M_0$ is exactly $E[m]$.
Its commutator pairing is therefore a perfect alternating
pairing on $(\mathbb Z/m)^2$.
Since $M=M_0^{\otimes-3}$, tensoring lifts shows directly
that the commutator pairing for $M$ on $E[m]$ is the $(-3)$rd
power of this pairing. It is still perfect because $(3,m)=1$.
We have not applied a separable-type assertion to a line of
degree divisible by $p$ when $p=3$.

Give $N$ its symmetric inversion lift normalized to be the
identity on the fibre over $O$. Give $M=N^m$ the induced lift.
Both inversion structures in (D.EC.2) are normalized at $O$,
so $\psi$ respects them: their possible discrepancy is a
constant equal to one at $O$.
Conjugation by this inversion therefore defines an automorphism
$\iota$ of $G_m$ which induces $u\mapsto-u$ on $E[m]$ and
fixes $\mu_m$ pointwise. For $g\in G_m$ one has $g^m\in\mu_m$,
and $\iota(g)=g^{-1}c$ for some $c\in\mu_m$.
It follows that
\[
               g^m=\iota(g^m)=\iota(g)^m=g^{-m}.
                                                               \tag{D.EC.5}
\]
As $m$ is odd, $g^m=1$. This argument does not divide by two
in $k$ and does not invoke the characteristic restrictions
of the later symmetric-theta sections of Mumford's paper.

Choose the central generator $z$ to have the indicated
normal-parameter character $\zeta$.
Perfection of the commutator pairing permits a basis of
$E[m]$ whose lifts $x,y$ satisfy $[x,y]=z^3$.
The exponent calculation gives $x^m=y^m=1$.
The elements $x^ay^bz^c$, $0\le a,b,c<m$, are distinct,
since their images in $E[m]$ first determine $a,b$ and
then their central coordinate determines $c$.
They exhaust $G_m$. This proves (D.EC.1).

Finally we check that the torsor stays connected on the
fixed completed puncture. Compose $f_m$ with
$Y_0\to\operatorname{Spec}A$, where $A$ is the cone ring,
and take its Stein factorization. Its finite algebra
$A_m=\Gamma(\operatorname{Tot}(M),\mathcal O)$ is a normal
domain: global sections of an integral normal scheme are
integrally closed in their fraction field. Its field has
degree $m^3$ over $\operatorname{Frac}A$, since the Stein map
is an isomorphism generically. There is only one point of
$\operatorname{Spec}A_m$ over the vertex. Indeed its preimage
in $\operatorname{Tot}(M)$ is the connected zero section,
and the proper Stein map is surjective.
The puncture of this finite normalization is the torsor just
constructed, since $f_m$ is finite there. These uses of Stein
factorization are in its Noetherian form
\citep[Theorem~37.53.4, Tag~03H0]{D_StacksEllipticThetaCone}.

The completed algebra $A_m\otimes_A R$ is thus the completion
of one excellent normal local domain. By Tag~0C23 it is a
normal local domain, not a product. Its generic degree remains
$m^3$ by flat base change. Its puncture is the connected finite
\'etale $G_m$-torsor $U_m\to U$.
The resolution obtained by base change from
$\operatorname{Tot}(M)$ is regular by Tags~0BG4, and its
single exceptional curve and its map to $E$ are unchanged.
At the generic exceptional point, (D.EC.3) has ramification
index $m$ and residue degree $m^2$. Its inertia is exactly
the fibre-scalar subgroup $\langle z\rangle$.
\end{proof}

\begin{proposition}[Proved: unbounded actual unitary Brauer periods on one fixed puncture]
\label{D:fixed-elliptic-cone-unitary-geometric-periods}
For the one fixed complete surface $R$ of
Lemma~\ref{D:elliptic-cone-Heisenberg-covers}, there are
continuous finite-image representations
\[
 \rho_j:\pi_1(U)\longrightarrow\operatorname{PSU}_{m_j}(q_j)
             \longrightarrow\operatorname{PGL}_{m_j}(\mathbb F_{q_j^2})
                                                               \tag{D.EC.6}
\]
whose associated geometric Brauer classes on $U$ have
exact order $m_j$, where $m_j\to\infty$ at fixed $p$.
One can take, for $j\ge1$,
\[
 q_j=p^{2^j},\qquad
 m_j=\begin{cases}q_j+1,&p=2,\\ (q_j+1)/2,&p\ne2.\end{cases}
                                                               \tag{D.EC.7}
\]
Their images are the tame Heisenberg groups of order $m_j^3$
in (D.EC.1). In particular this is an assertion about actual
projective local systems and their geometric Brauer periods;
it is not just an assertion about an abstract scalar extension.
\end{proposition}

\begin{proof}
Fix $m,q$ from (D.EC.7). Then $m$ is odd, $m\ge5$,
$(m,6p)=1$, and $m\mid q+1$. For the assertion about three,
if $p=3$ then $q+1\equiv1\pmod3$; otherwise $q\equiv1\pmod3$.
For odd $p$, the even exponent also gives $q\equiv1\pmod8$,
so $(q+1)/2$ is odd. These observations prove all the stated
coprimality assertions. The sequence is plainly unbounded.

Put $K=\mathbb F_{q^2}\subset k$, and choose $\zeta\in K$
of exact order $m$. Its Hermitian norm is one.
On the standard orthonormal basis indexed by $j\in\mathbb Z/m$
define
\[
 Xe_j=e_{j+1},\qquad Ze_j=\zeta^j e_j,\qquad
 Ye_j=\zeta^{-3j(j-1)/2}e_j.                            \tag{D.EC.8}
\]
The quadratic exponent is well defined modulo $m$ since
$m$ is odd. All three matrices preserve the standard
nondegenerate Hermitian form. Their determinants are
\[
 \det X=1,\qquad \det Z=\zeta^{m(m-1)/2}=1,
 \qquad \det Y=\zeta^{-3m(m-1)(m-2)/6}=1.               \tag{D.EC.9}
\]
For the last equality, $(m,6)=1$ implies that
$(m-1)(m-2)/6$ is an integer.
Thus $X,Y,Z\in\operatorname{SU}_m(q)$, with the exact
unitary and projective conventions of
\citet[pp.~114--116, Lemma~10.1]{Taylor1992ClassicalB}.
Direct calculation gives
\[
 \begin{aligned}
 &X^m=Y^m=Z^m=I,\\
 &[X,Y]=\zeta^{-3}Z^3,\qquad [X,Z]=\zeta^{-1}I,
          \qquad [Y,Z]=I .
 \end{aligned}                                                \tag{D.EC.10}
\]
Hence $x\mapsto\overline X$, $y\mapsto\overline Y$,
$z\mapsto\overline Z$ defines a homomorphism
$G_m\to\operatorname{PSU}_m(q)$.
If $X^aY^bZ^c$ is scalar, its coordinate-line permutation
first gives $a=0$. The second finite difference, modulo $m$,
of its diagonal exponent $-3b j(j-1)/2+cj$ is $-3b$.
Thus $b=0$, since $(3,m)=1$, and then the first difference
gives $c=0$. The homomorphism is injective.
Composing it with the quotient supplied by the lemma gives
(D.EC.6). The stated $m,q$ lie in the usual nonabelian simple
unitary ranges; simplicity is not used in the argument.

Let $\alpha\in\operatorname{Br}(U)$ be the class of the
Azumaya algebra obtained from the $G_m$-torsor by descent
of $\operatorname{Mat}_m(\mathcal O_{U_m})$ under these
projective matrices. The section
$x^ay^bz^c\mapsto X^aY^bZ^c$, with exponents between zero
and $m-1$, has scalar multiplication defects in $\mu_m$,
by (D.EC.10). Consequently its obstruction on $U$ is the
image of a class in $H^2(U,\mu_m)$, and
\[
                              m\alpha=0.               \tag{D.EC.11}
\]

We prove the opposite inequality in the \emph{geometric}
Brauer group. Let $H=\langle x,z\rangle\simeq(\mathbb Z/m)^2$.
Restrict the torsor to its intermediate quotient by $H$,
and complete the exceptional discrete valuation.
The resulting Galois field extension $\mathcal L/\mathcal K$
has group $H$, inertia $\langle z\rangle$, and residue
extension cyclic of degree $m$ with generator $x$.
Here the unique exceptional curve upstairs ensures that
the decomposition group is all of $H$.
Thus $\mathcal L^{\langle z\rangle}/\mathcal K$ is an
unramified extension of degree $m$, and
$\mathcal L^{\langle x\rangle}/\mathcal K$ is a totally
ramified cyclic extension of degree $m$.

Choose a Kummer generator $\xi$ of the unramified extension
with $x(\xi)=\zeta\xi$. Multiplying it by a power of a
uniformizer of $\mathcal K$ makes it a unit. Its residue
still generates the degree-$m$ residue extension: reduction
preserves the primitive eigencharacter, and $m$ is invertible.
Write $a=\xi^m\in\mathcal K^\times$.
Choose a uniformizer $s_0$ of the totally ramified field.
The action of $z$ on its normal line is $\zeta$, by our
choice in the lemma. The eigenspace projection
\[
       s=\frac1m\sum_{i=0}^{m-1}\zeta^{-i}z^i(s_0)
                                                               \tag{D.EC.12}
\]
is a uniformizer, satisfies $z(s)=\zeta s$, and is fixed
by $x$. Therefore $b=s^m\in\mathcal K^\times$ has
$v_{\mathcal K}(b)=1$.

The actual descended algebra can now be identified without
any inference from the order of an abstract group-cohomology
class. On $\mathcal L\otimes_k\operatorname{Mat}_m(k)$ the
semilinear descent actions of $x,z$ are their actions on
$\mathcal L$ followed by conjugation by $X,Z$, respectively.
The two matrices
\[
             A=\xi Z,\qquad B=s^{-1}X
                                                               \tag{D.EC.13}
\]
are invariant under both actions, since
$XZX^{-1}=\zeta^{-1}Z$ and $ZXZ^{-1}=\zeta X$.
They satisfy
\[
             A^m=aI,\qquad B^m=b^{-1}I,
                       \qquad AB=\zeta BA.             \tag{D.EC.14}
\]
The $m^2$ matrices $Z^iX^j$, $0\le i,j<m$, are a basis
of the matrix algebra over $\mathcal L$: the shifts give
the different diagonals and the primitive-root Vandermonde
matrix separates their diagonal coefficients.
Consequently $A^iB^j$ are linearly independent and span
the descended algebra over $\mathcal K$.
It is precisely the symbol algebra $(a,b^{-1})_\zeta$.

We use the known cyclic-algebra norm criterion, with no
primality restriction on $m$:
\citet[Proposition~4.7.3 and Corollaries~4.7.4--4.7.5,
pp.~109--111]{D_GilleSzamuely2006Cone}.
It identifies the class of this algebra, up to the harmless
inverse choice of cyclic generator, with the class of
$b^{-1}$ modulo norms from the unramified degree-$m$
extension $\mathcal K(\xi)/\mathcal K$.
Its $e$th multiple is therefore trivial only if $b^{-e}$
is such a norm. Every norm from this unramified extension
has valuation divisible by $m$: for $u\ne0$, all $m$
conjugates have the same integral valuation, and the
valuation of their product is $m v(u)$.
Since $v(b)=1$, this forces $m\mid e$.
Together with (D.EC.11), it proves that the restriction of
$\alpha$ to this actual valued-field extension, and hence
$\alpha$ itself on $U$, has exact order $m$.

In particular Proposition~\ref{D:single-curve-Brauer-residue-cover}
identifies its residue with a class of exact order $m$ in
$H^1(E,\mathbb Z/m)$. This last observation is a consequence
of the established geometric order, not an appeal to the
prime-order-only eigensheaf assertions in the literature.
The fixed germ $R$ and fixed exceptional curve $E$ have
not varied with $m$ or $q$. This completes the proof.
\end{proof}

\paragraph{Scope of the diagnostic.}
The image of (D.EC.6) is the explicitly displayed nilpotent
group $G_m$, not the whole simple unitary group. Thus this
does not construct epimorphisms onto the ambient simple
groups, does not contradict any quadratic representation
bound, and does not close an unbounded-period hereditary
count. It proves that a uniform bound on geometric Brauer
period cannot be inferred merely from a unitary form,
even for actual tame projective local systems on one fixed
complete simple elliptic surface puncture.

